\section{Common conventions}
\label{sec:II:conventions}
All spaces and coherent sheaves are complex analytic unless a polynomial
ring is specified. We use $B=\PP^1_t$, $(p_1,p_2,p_0)=(0,1,\infty)$,
the ordered lattice
$\Lambda=\Z\langle\widehat\gamma,\widehat u,\widehat w,\widehat\delta\rangle$,
and the marked periods
\[
\Pi_c=\begin{pmatrix}6\mu&\tau&1&0\\\beta^{\rm part}+c&\mu&0&1\end{pmatrix},
\qquad w=e^{-\pi i c},\qquad (\beta^{\rm part}+\tau)(p_0)=0.
\]
Base meridians are counterclockwise. The analytic root decks are clockwise,
$s_j(g_jz)=e^{-2\pi i/m_j}s_j(z)$, with $g_1g_2g_0=1$;
the resulting homology monodromies are $A_1,A_2,M_0$.
The primitive invariant vectors are
$\varepsilon=(1,2,-4,0)^t$ and
$\varepsilon'=(1,3,-3,0)^t$.
For a finite affine datum we write
$v_1=r_1\varepsilon+k_1\widehat\delta$ and
$v_2=r_2\varepsilon'+k_2\widehat\delta$; the coefficients
$r_1,r_2$ are also denoted $l_1,l_2$ in the peripheral presentation.
Define
\[
 D_{\rm part}=\operatorname{Im}\beta^{\rm part}
       -\frac{6(\operatorname{Im}\mu)^2}{\operatorname{Im}\tau},
 \qquad M=\max_{B\setminus\{p_0\}}D_{\rm part}
       =-\frac3\pi\log6.
\]
The admissible parameter inequality is $\operatorname{Im}c<-M$;
the maximum of $D_{\rm part}$ is attained only at $t=1$, as proved in \cref{fam:src:C.2.3}.
The standard multiplicities are $m_1=3,m_2=4$,
$f^*p_j=m_jS_j$, and $W=f^{-1}(p_0)$ is reduced, with normalization
$\eta:\mathrm dP_6\to W$.
For general finite affine twists we use
$p=12l_0-4l_1-3l_2$, retaining $\pi_1(X)=\Z/|p|$.
The complete marked transport and construction are those of Paper I.
An unqualified fibre is its full Cartier scheme; its reduction is always
specified explicitly. Homology and cohomology retain their displayed
coefficient groups, and complex manifolds carry their complex orientations.


\section{Index, Picard components and endpoint cohomology}\label{sec:II:01}


\begin{proposition}\label{cdp:st:C.1.1}
\StatementRecord{cdp:st:C.1.1}{proved}{Historical events: Claude Fable 5.2 (Anthropic), effort high, 16 November 2026: carried as background without fault found. GPT-6 Astra (OpenAI), ChatGPT Pro mode, effort Pro, S6\_09\_PAPERS, 17 November 2026: targeted proof check. Also checked in this writing. Pre-continuation writing events: GPT-6 Astra (OpenAI), ChatGPT Pro mode, effort Pro, S6\_10\_PAPERS, 18 November 2026: targeted proof check. Continuation metadata events: GPT-6 Astra (OpenAI), ChatGPT Pro mode, effort Pro, S6\_10\_PAPERS continuation, 19 November 2026: used and rechecked in this run; status-vocabulary reconciliation; metadata reconciliation only; no new mathematical check. \emph{Pre-continuation recorded proof scope.} The exponential sequence with finite H2, rational Riemann-Roch degree-six term, Poincare duality, and parity of the middle pairing were re-derived}{the hypotheses, notation and definitions stated here; the explicit constructions and calculations in the proof}

Let $Y$ be a connected compact complex threefold with
$b_1(Y)=b_2(Y)=0$. Every holomorphic line bundle $L$ on $Y$ satisfies
\[
 \chi(Y,T_Y\otimes L)=\frac12\int_Yc_3(T_Y)
 =1-\frac{b_3(Y)}2.
\]
The orientation in the integral is the complex orientation. The group
$\Pic(Y)$ has finitely many connected components, each a torsor for the
complex vector group $H^1(Y,\OO_Y)$, and
$\pi_0\Pic(Y)\simeq H^2(Y;\Z)$ through the first Chern class.
\end{proposition}


\begin{proof}
\noindent\textbf{Step 1.} The universal coefficient theorem gives
$H^1(Y;\Z)=\Hom(H_1(Y;\Z),\Z)=0$ from $b_1(Y)=0$.
The group $H^2(Y;\Z)$ is finite because $Y$ has a finite CW model and
$b_2(Y)=0$. In the exponential sequence, normalized by
$\exp(2\pi i\,\cdot)$, the relevant maps are
\[
0=H^1(Y;\Z)\longrightarrow H^1(Y,\OO_Y)
 \xrightarrow{\exp} \Pic(Y)
 \xrightarrow{c_1} H^2(Y;\Z).
\]
They are group homomorphisms; the middle map is holomorphic on the
Picard group. Consequently its image $\Pic^0(Y)$ is isomorphic to the
complex vector space $H^1(Y,\OO_Y)$. The next map in the exponential
sequence is $H^2(Y;\Z)\to H^2(Y,\OO_Y)$; it is zero because its
source is finite and its target is a complex vector space. Thus
$c_1$ is surjective, and the component group is exactly $H^2(Y;\Z)$.

\noindent\textbf{Step 2.} Poincar\'e duality gives $H^4(Y;\Q)=0$. Thus
$c_1(T_Y)$, $c_2(T_Y)$ and $c_1(L)$ vanish in rational cohomology.
The degree-six component of
$\operatorname{ch}(T_Y)\exp(c_1(L))\operatorname{td}(T_Y)$ is then
$c_3(T_Y)/2$. The Hirzebruch--Riemann--Roch formula for the Dolbeault
index gives the first equality. Chern--Gauss--Bonnet identifies
$\int_Yc_3(T_Y)$ with $e(Y)$. Since $Y$ is a connected closed oriented
six-manifold, its Betti numbers satisfy
$b_0=b_6=1$, $b_1=b_5=0$, and $b_2=b_4=0$. Hence
$e(Y)=2-b_3(Y)$. The pairing on $H^3(Y;\Q)$ is nondegenerate and
alternating, so $b_3(Y)$ is even, consistently with the integral index.
The exact symbolic expansion is recorded in
\texttt{code/D\_exact\_checks.py} and its output.
\end{proof}






\begin{proposition}\label{cdp:st:C.1.3}
\StatementRecord{cdp:st:C.1.3}{proved}{Historical events: Claude Fable 5.2 (Anthropic), effort high, 16 November 2026: carried as background without fault found. GPT-6 Astra (OpenAI), ChatGPT Pro mode, effort Pro, S6\_09\_PAPERS, 17 November 2026: targeted proof check. Also checked in this writing. Pre-continuation writing events: GPT-6 Astra (OpenAI), ChatGPT Pro mode, effort Pro, S6\_10\_PAPERS, 18 November 2026: targeted proof check. Continuation metadata events: GPT-6 Astra (OpenAI), ChatGPT Pro mode, effort Pro, S6\_10\_PAPERS continuation, 19 November 2026: used and rechecked in this run; status-vocabulary reconciliation; metadata reconciliation only; no new mathematical check. \emph{Pre-continuation recorded proof scope.} The generic endpoint identity and analytic semicontinuity propagation on each Picard component were re-derived; the endpoint-vanishing result remains an explicitly cited input}{the hypotheses, notation and definitions stated here; \cref{cdp:st:C.1.2}; \cref{cdp:st:C.1.1}}

For every $Y$ and $D$ in Statement \ref{cdp:st:C.1.2} and every $L\in\Pic(Y)$,
\[
 h^2(Y,T_Y\otimes L)\geq
 \max\left\{0,1-\frac{b_3(Y)}2\right\}.
\]
In particular, $b_3(Y)=0$ gives $H^2(Y,T_Y\otimes L)\ne0$ for
every $L$. A threefold in this class possessing a twist with vanishing
second tangent cohomology has $b_3(Y)\geq2$.
\end{proposition}


\begin{proof}
\noindent\textbf{Step 1.} On the dense open subset of each Picard component supplied by
Statement \ref{cdp:st:C.1.2}, Statement \ref{cdp:st:C.1.1} becomes
\[
 h^2(Y,T_Y\otimes L)-h^1(Y,T_Y\otimes L)
 =1-b_3(Y)/2.
\]
The asserted lower bound holds there. The jump locus for the lower
bound on $h^2$ is closed analytic in the explicit family above and
contains a dense subset of the component. It therefore contains the
entire component. This proves the bound for all line bundles,
including torsion bundles and the trivial bundle. The final assertion
is its direct numerical consequence.
\end{proof}



Statements \ref{cdp:st:C.1.1}--\ref{cdp:st:C.1.3} prove the stated index obstruction with rational Betti-number
vanishing. No integral trivialization of either Chern class is needed
for the index. The existence of the divisor enters in the endpoint
vanishing, and a positive multiple of its actual Cartier divisor
enters the Picard family. These statements supply an index
obstruction; they make no assertion that the failed degree-two
vanishing holds.




\label{cdp:sec:E}

\begin{proposition}[Vertical divisors and the field of meromorphic functions]\label{cdp:st:E1}
\StatementRecord{cdp:st:E1}{proved}{Historical events: Claude Fable 5.2 (Anthropic), effort high, 16 November 2026: carried as background without fault found. GPT-6 Astra (OpenAI), ChatGPT Pro mode, effort Pro, S6\_09\_PAPERS, 17 November 2026: targeted proof check. Also checked in this writing. Pre-continuation writing events: GPT-6 Astra (OpenAI), ChatGPT Pro mode, effort Pro, S6\_10\_PAPERS, 18 November 2026: targeted proof check. Continuation metadata events: GPT-6 Astra (OpenAI), ChatGPT Pro mode, effort Pro, S6\_10\_PAPERS continuation, 19 November 2026: used and rechecked in this run; status-vocabulary reconciliation; metadata reconciliation only; no new mathematical check. \emph{Pre-continuation recorded proof scope.} The actual restriction of the divisor class to a Kahler fibre, positive divisor integral, and proper graph descent of the meromorphic function were re-derived}{the hypotheses, notation and definitions stated here; the explicit constructions and calculations in the proof}

Let $Y$ be a connected compact complex manifold of dimension $n\geq2$
with $H^2(Y,\mathbb R)=0$. Let $g:Y\to C$ be a surjective holomorphic
map to a connected compact Riemann surface, with connected fibres.
Suppose that its smooth fibres over a nonempty open subset of $C$ are
K\"ahler. Every effective divisor of $Y$ is supported on finitely many
fibres, and pullback induces an isomorphism of fields
\[
 g^*: \mathcal M(C)\xrightarrow{\ \sim\ }\mathcal M(Y).
\]
In particular $a(Y)=1$. Reducedness of the other fibres is unnecessary.
\end{proposition}


\begin{proof}
\noindent\textbf{Step 1.} For a prime divisor $D$, the image $g(D)$ is analytic by properness. It
is a point or all of $C$. In the latter case choose a regular value $c$
in the stated open subset for which $F=g^{-1}(c)$ is not contained in
$D$. The canonical section of $\mathcal O_Y(D)$ restricts to a nonzero
section with nonempty effective zero divisor $D_F$ on $F$. For a
K\"ahler form $\omega$ on $F$, the divisor class formula gives
\[
 0=\int_F c_1(\mathcal O_Y(D)|_F)\wedge\omega^{n-2}
   =\int_{D_F}\omega^{n-2}>0.
\]
The first equality follows from the restriction homomorphism
$H^2(Y,\mathbb R)\to H^2(F,\mathbb R)$; the last follows from positivity
on every component with its positive multiplicity. Thus $g(D)$ is a
point. A divisor has finitely many prime components, proving the first
assertion.

\noindent\textbf{Step 2.} For $h\in\mathcal M(Y)$ remove from $C$ the finitely many values
supporting its polar divisor. On the inverse image of the complement,
$h$ is holomorphic: a meromorphic function on a manifold having no
codimension-one pole extends across its codimension-two indeterminacy
set by the local normality of the structure sheaf. Its restriction to
each compact connected fibre is constant. The identity
$g_*\mathcal O_Y=\mathcal O_C$, supplied by Stein factorization and
the normality of $C$, therefore gives a holomorphic descended function
there. To check meromorphic extension at the removed values, take the
closure of the graph of $h$ in $Y\times\mathbb P^1$ and apply the proper
map $g\times1$. Its image is a one-dimensional analytic subset of
$C\times\mathbb P^1$, generically the graph of the descended function.
Its component dominating $C$ has generic degree one, and hence is the
graph of a meromorphic function on $C$. Pullback equals $h$ on a dense
open subset and therefore in the meromorphic function field. Pullback
is injective because $g$ is surjective. This proves the field
isomorphism and its asserted transcendence degree.
\end{proof}






\section{Integral specialization, coherent restrictions and component counts}\label{sec:II:02}


\begin{proposition}\label{cdp:st:D.1.2}
\StatementRecord{cdp:st:D.1.2}{proved}{Historical events: Claude Fable 5.2 (Anthropic), effort high, 16 November 2026: carried as background without fault found. Claude Fable 5.2 (Anthropic), effort high, 17 November 2026: carried as background without fault found. GPT-6 Astra (OpenAI), ChatGPT Pro mode, effort Pro, S6\_09\_PAPERS, 17 November 2026: targeted proof check. Also checked in this writing. Pre-continuation writing events: GPT-6 Astra (OpenAI), ChatGPT Pro mode, effort Pro, S6\_10\_PAPERS, 18 November 2026: targeted proof check; used and rechecked in this run. Continuation metadata events: GPT-6 Astra (OpenAI), ChatGPT Pro mode, effort Pro, S6\_10\_PAPERS continuation, 19 November 2026: used and rechecked in this run; status-vocabulary reconciliation; metadata reconciliation only; no new mathematical check. \emph{Pre-continuation recorded proof scope.} The integral coinvariant calculation, compact-support boundary map, universal coefficient calculation, and component count were re-derived. The rational count is credited to AC26 section 10.5, printed page 91}{the hypotheses, notation and definitions stated here; \cref{cdp:st:D.1.1}; \cref{cdp:eq:D1}}

Let $f:Y\to C$ be a proper holomorphic surjection from a connected
compact complex threefold to a compact Riemann surface, with
connected fibres. Let $\Delta$ be a nonempty finite set containing
the critical values, $s=|\Delta|$, $g=g(C)$,
$D=f^{-1}(\Delta)_{\rm red}$, and let $r$ be the number of
irreducible components of $D$. Write $F$ for a smooth fibre and
\[
t'=\dim_{\Q}H_1(F;\Q)_{\pi_1(C\setminus\Delta)}
   =\dim_{\Q}H^1(F;\Q)^{\pi_1(C\setminus\Delta)}.
\]
For $b_2(Y)=0$ the exact component count is
\[
\boxed{\quad r=2g+s-1+t'-b_1(Y).\quad}
\]
Under CDP20's integral assumptions
$H^1(Y;\Z)=H^2(Y;\Z)=0$, the base is $\PP^1$ and
\[
r=s-1+t',\qquad
H_1(F;\Z)_{\pi_1(C\setminus\Delta)}\text{ is free abelian}.
\]
The last assertion concerns the indicated quotient of $H_1(F;\Z)$;
it does not assert that $H_1(F;\Z)$ itself is torsion-free.
\end{proposition}



The rational component-count formula is the correction in \citet[Section~10.5, p.~91]{AC2026}. The integral torsion-free conclusion under the displayed integral hypotheses is an additional clause; its full argument is included here.

\begin{proof}
\noindent\textbf{Step 1.} Put $U=Y\setminus D$ and $C^\circ=C\setminus\Delta$.
The map $U\to C^\circ$ is a differentiable fibre bundle. Since
$\pi_2(C^\circ)=0$, its homotopy sequence is
\[
1\longrightarrow\pi_1(F)\longrightarrow\pi_1(U)
 \longrightarrow\pi_1(C^\circ)\longrightarrow1.
\]
The last group is free of rank $2g+s-1$. Lifting its free generators
gives a section of this extension. Proposition \ref{cdp:st:D.1.1} therefore gives
\[
H_1(U;\Z)\simeq
\Z^{2g+s-1}\oplus H_1(F;\Z)_{\pi_1(C^\circ)}.       \tag{D1}\label{cdp:eq:D1}
\]
The second equality defining $t'$ follows by duality: the dual of
the coinvariant quotient is precisely the space of invariant
linear functionals on $H_1(F;\Q)$.

\noindent\textbf{Step 2.} The cohomology sequence with compact supports for the open inclusion
$U\subset Y$ contains
\[
H^4(Y;\Q)\longrightarrow H^4(D;\Q)
\longrightarrow H_c^5(U;\Q)
\longrightarrow H^5(Y;\Q)\longrightarrow H^5(D;\Q)=0.
\]
The top-degree cohomology of the compact pure complex surface $D$
is $\Q^r$. One can see this by deleting the singular and intersection
loci, whose real dimensions are at most two, and taking compactly
supported top cohomology of its $r$ oriented smooth components.
Poincar\'e duality identifies $H_c^5(U;\Q)$ with $H_1(U;\Q)$ and
$H^5(Y;\Q)$ with $H_1(Y;\Q)$. Since $H^4(Y;\Q)=0$, taking
dimensions and using \eqref{cdp:eq:D1} proves the formula.

\noindent\textbf{Step 3.} For the integral refinement, the universal coefficient theorem shows
first that $H_1(Y;\Z)$ is finite, using $H^1(Y;\Z)=0$, and then that
it is zero, using
$\operatorname{Ext}^1(H_1(Y;\Z),\Z)\subset H^2(Y;\Z)=0$.
The group $H^4(Y;\Z)$ can be finite and nonzero. Its homomorphism
to $H^4(D;\Z)=\Z^r$ is nevertheless zero. The same sequence now
gives $H_1(U;\Z)\simeq\Z^r$, since $H^5(Y;\Z)=H_1(Y;\Z)=0$.
Equation \eqref{cdp:eq:D1} then proves that the coinvariant summand is free.
Finally $H^1(C;\Q)\hookrightarrow H^1(Y;\Q)$ by the Leray edge
map, so $g=0$.
\end{proof}






\begin{proposition}[Refuted]\label{cdp:st:D.1.5}
\StatementRecord{cdp:st:D.1.5}{proved}{Historical events: Claude Fable 5.2 (Anthropic), effort high, 16 November 2026: carried as background without fault found. Claude Fable 5.2 (Anthropic), effort high, 17 November 2026: carried as background without fault found. GPT-6 Astra (OpenAI), ChatGPT Pro mode, effort Pro, S6\_09\_PAPERS, 17 November 2026: targeted proof check. Also checked in this writing. Pre-continuation writing events: GPT-6 Astra (OpenAI), ChatGPT Pro mode, effort Pro, S6\_10\_PAPERS, 18 November 2026: targeted proof check; used and rechecked in this run. Continuation metadata events: GPT-6 Astra (OpenAI), ChatGPT Pro mode, effort Pro, S6\_10\_PAPERS continuation, 19 November 2026: used and rechecked in this run; status-vocabulary reconciliation; metadata reconciliation only; no new mathematical check. \emph{Pre-continuation recorded proof scope.} The free involution, integral monodromy coinvariants, finite-cover coherent direct images and restriction rank were checked. The rational elliptic product's direct images and Betti numbers follow from its displayed Weierstrass model. Both examples retain their actual Betti hypotheses}{the hypotheses, notation and definitions stated here; \cref{cdp:st:D.1.1}; \citet[Section~10.5, p.~91]{AC2026}}


\textbf{Explicit second examples for the erroneous steps.}
Let $E_0=F_0=E'=\C/(\Z+i\Z)$ and $a=1/2\in E_0$.
The involution
\[
(e,z)\longmapsto(e+a,-z)
\]
acts freely on $E_0\times F_0$. Put
$S=(E_0\times F_0)/\langle(a,-1)\rangle$,
$Y=S\times E'$, and $C=E_0/\langle a\rangle$.
The induced map $Y\to C$ is a smooth torus fibration with fibre
$F_0\times E'$ and monodromy
$\operatorname{diag}(-I_2,I_2)$ around the nontrivial deck character.
Its integral homology coinvariants are
$\Z^2\oplus(\Z/2\Z)^2$. After removing $s\geq1$ fibres,
Proposition \ref{cdp:st:D.1.1} gives
\[
H_1(U;\Z)=\Z^{s+3}\oplus(\Z/2\Z)^2,
\]
whereas the uncorrected direct-sum argument has free rank $s+5$.
This is an explicit compact complex realization of the faulty
group-theoretic step, independent of the retained six-sphere.

Let $\eta$ be the nontrivial order-two degree-zero line bundle on
$C$ associated to the double cover $E_0\to C$. Finite-cover descent
of fibre cohomology gives
\[
R^1f_*\OO_Y=\eta\oplus\OO_C,\qquad R^2f_*\OO_Y=\eta.
\]
The invariant part of the exterior algebra on the three
antiholomorphic one-form directions, with signs $(1,-1,1)$,
gives $h^q(Y,\OO_Y)=(1,2,1,0)$. The restriction from
$H^1(Y,\OO_Y)$ to the two-dimensional $H^1(F_0\times E',\OO)$
has image precisely the $E'$ direction and has rank one.
All these homomorphisms are induced by the displayed finite
\'etale quotient and projection. This example has $b_1(Y)=4$;
it is a counterexample to the unchecked restriction argument, not
a second example satisfying the integral CDP20 hypotheses.

There is also an explicit negative summand over $\PP^1$. Let
$h:R\to\PP^1$ be the relatively minimal elliptic surface with
Weierstrass equation
\[
y^2=x^3+(t^4+1)x+t^6.
\]
Its fundamental line bundle is $\OO(1)$ and its discriminant has
twelve simple zeroes. Relative duality gives
$R^1h_*\OO_R=\OO(-1)$. For $Y'=R\times E'\to\PP^1$,
the product K\"unneth maps give
\[
R^1f_*\OO_{Y'}=\OO\oplus\OO(-1),\qquad
R^2f_*\OO_{Y'}=\OO(-1),
\]
and the restriction in degree one again has rank one. This
projective example has $b_1=2$, $b_2=11$ and $a=3$, so it does not
replace the missing second counterexample under all the standing
CDP20 hypotheses.
\end{proposition}


\Uses{the hypotheses, notation and definitions stated here; \cref{cdp:st:D.1.1}; \citet[Section~10.5, p.~91]{AC2026}}


The bielliptic-product rank defect is already described in \citet[Section~10.5, p.~91]{AC2026}. The full integral torsion calculation, the coherent direct images and the separate rational-elliptic product below retain their additional precise content.





\begin{proposition}[A second count counterexample]\label{cdp:st:D.1.6}
\StatementRecord{cdp:st:D.1.6}{proved}{Historical events: Claude Fable 5.2 (Anthropic), effort high, 16 November 2026: carried as background without fault found. GPT-6 Astra (OpenAI), ChatGPT Pro mode, effort Pro, S6\_09\_PAPERS, 17 November 2026: targeted proof check. Also checked in this writing. Pre-continuation writing events: GPT-6 Astra (OpenAI), ChatGPT Pro mode, effort Pro, S6\_10\_PAPERS, 18 November 2026 04:50:08.441928+00:00: targeted proof check. GPT-6 Astra (OpenAI), ChatGPT Pro mode, effort Pro, S6\_10\_PAPERS, 18 November 2026 04:51:18.698751+00:00: used and rechecked in this run. Continuation metadata events: GPT-6 Astra (OpenAI), ChatGPT Pro mode, effort Pro, S6\_10\_PAPERS continuation, 19 November 2026: used and rechecked in this run; status-vocabulary reconciliation; metadata reconciliation only; no new mathematical check. \emph{Pre-continuation recorded proof scope.} proved counterexample; rejected printed count remains refuted. Full displayed Wang counterexample, exact monodromy algebra and operative integral specialization input. Opened Wang attached Theorem1.1 p2, construction/finite and rank-one/rank-two filling maps pp2–8, actual marked period and cohomology matrices pp9–11, and the complete integral Leray computation pp12–15. Independently computed product identity, order3, common invariant lattice Z\ensuremath{\psi} and homology coinvariant Smith diagonal(1,1,1,0). The I2 period vector(2,1) is sent to(1,0) by (z,v)\ensuremath{\mapsto}(v,v\textsuperscript{2}/z); the inverse is (x,y)\ensuremath{\mapsto}(x\textsuperscript{2}/y,x), giving one component. The other three reduced fillings are irreducible by their explicit unit-square fan or free order3 quotient. Independently read/rederived iv:rank\_two\_specialization, including its full cross-Chern line and actual unit-preserving flow; this closes the integral rank-two source input. Wang’s fixed meridian translation \ensuremath{\ell}/3 gives contraction coefficients(1,2,2,1), with no additional affine defect because the other sections extend. The resulting H2=H3=Z/2 gives H1coho=H2coho=0 by UCT; the rank1 invariant group gives r=4. The rejected count is seven. This event proves the counterexample and does not reinstate the rejected CDP formula. Review provenance: New derivation after reading the supplied proof; not a blind review and not an independent human check}{the hypotheses, notation and definitions stated here; \cite{Wang2026}; \cref{cdp:st:D.1.13}}


\textbf{A second count counterexample under the integral hypotheses.}
Wang's attached construction $X_2(\lambda)$ has integral homology
$H_0=H_6=\Z$, $H_2=H_3=\Z/2\Z$, and the other reduced
homology groups zero. Hence its integral cohomology satisfies
$H^1=H^2=0$. Its fibration has four exceptional values, over which
the reduced fibres are irreducible: the multiplicity-three
bielliptic fibre, two rank-one fillings, and the rank-two filling.
In the rank-one filling over the $I_2$ value, the primitive period
exponent is $(2,1)$; the component count is one, not two.
These are the construction's statements
\texttt{intro:main}, \texttt{fill:rankone},
\texttt{fill:a2}, \texttt{fill:log}, and
\texttt{main:theorem} in the attached primary source
\cite{Wang2026}.

In its cohomology basis the four displayed monodromies are
\[
\begin{aligned}
B_*&=\begin{pmatrix}1&0&-6&-2\\0&-1&-1&1\\0&1&0&-1\\0&0&0&1\end{pmatrix},&\qquad
B_2&=\begin{pmatrix}1&0&6&3\\0&1&2&1\\0&0&1&0\\0&0&0&1\end{pmatrix},\\
B_+&=\begin{pmatrix}1&0&0&0\\0&2&1&0\\0&-1&0&0\\0&0&0&1\end{pmatrix},&\qquad
B_-&=\begin{pmatrix}1&0&0&-1\\0&1&1&0\\0&0&1&0\\0&0&0&1\end{pmatrix}.
\end{aligned}
\]
Their product is one and their common invariant cohomology group
has rank one. On homology the matrices are $B_i^{-t}$; the Smith
diagonal of their difference matrix is $(1,1,1,0)$. Consequently
$t'=1$ and the corrected count is $r=4-1+1=4$, while the printed
count gives seven. The exact matrix calculation is included in the
program. The topological construction and its local fillings are
being used here as cited primary-source results, not inferred from
the matrices alone. Its coherent direct images are derived separately
in Statement \ref{cdp:st:D.1.13} below.
\end{proposition}


\Uses{the hypotheses, notation and definitions stated here; \cite{Wang2026}; \cref{cdp:st:D.1.13}}





\begin{proposition}\label{cdp:st:D.1.7}
\StatementRecord{cdp:st:D.1.7}{proved}{Historical events: Claude Fable 5.2 (Anthropic), effort high, 16 November 2026: carried as background without fault found. GPT-6 Astra (OpenAI), ChatGPT Pro mode, effort Pro, S6\_09\_PAPERS, 17 November 2026: targeted proof check. Also checked in this writing. Pre-continuation writing events: GPT-6 Astra (OpenAI), ChatGPT Pro mode, effort Pro, S6\_10\_PAPERS, 18 November 2026: targeted proof check. Continuation metadata events: GPT-6 Astra (OpenAI), ChatGPT Pro mode, effort Pro, S6\_10\_PAPERS continuation, 19 November 2026: used and rechecked in this run; status-vocabulary reconciliation; metadata reconciliation only; no new mathematical check. \emph{Pre-continuation recorded proof scope.} The injection from real degree-one cohomology to coherent cohomology was re-derived using imaginary pluriharmonic primitives on normalization resolutions and reduction of nilpotents}{the hypotheses, notation and definitions stated here; the explicit constructions and calculations in the proof}

For a connected compact complex space $S$, including a nonreduced
one, the coefficient map
\[
H^1(S;\mathbb R)\longrightarrow H^1(S,\OO_S)
\]
is injective as a map of real vector spaces.
\end{proposition}


\begin{proof}
\noindent\textbf{Step 1.} First take $S$ reduced. Represent a real cohomology class by a
locally constant real \v Cech cocycle $a_{ij}$. If its image is zero,
after refining the open cover there are holomorphic functions $u_i$
with $a_{ij}=u_i-u_j$. The imaginary parts of the $u_i$ agree on
overlaps and give a global continuous pluriharmonic function.
Pull it back to a compact resolution of each irreducible component.
The maximum principle makes each such pullback constant; intersection
of components and connectedness of $S$ make the constants equal.
Thus, after subtracting one imaginary constant from all $u_i$,
each $u_i$ is real-valued and holomorphic, hence locally constant.
They form a real coboundary for $a_{ij}$. For nonreduced $S$, a
class zero in $H^1(S,\OO_S)$ also becomes zero after applying
$\OO_S\to\OO_{S_{\rm red}}$. The underlying topological space
and its real cohomology do not change under reduction, so the
reduced assertion proves the result.
\end{proof}






\begin{proposition}\label{cdp:st:D.1.8}
\StatementRecord{cdp:st:D.1.8}{proved}{Historical events: Claude Fable 5.2 (Anthropic), effort high, 16 November 2026: carried as background without fault found. GPT-6 Astra (OpenAI), ChatGPT Pro mode, effort Pro, S6\_09\_PAPERS, 17 November 2026: targeted proof check. Also checked in this writing. Pre-continuation writing events: GPT-6 Astra (OpenAI), ChatGPT Pro mode, effort Pro, S6\_10\_PAPERS, 18 November 2026: targeted proof check; used and rechecked in this run. Continuation metadata events: GPT-6 Astra (OpenAI), ChatGPT Pro mode, effort Pro, S6\_10\_PAPERS continuation, 19 November 2026: used and rechecked in this run; status-vocabulary reconciliation; metadata reconciliation only; no new mathematical check. \emph{Pre-continuation recorded proof scope.} The generization injection was re-derived from the meridian multiplicity and surjectivity of the complement fundamental group}{the hypotheses, notation and definitions stated here; \cref{cdp:st:D.1.2}}

Let $f:Y\to C$ be as in \ref{cdp:st:D.1.2}, and let $j:C^\circ\hookrightarrow C$
be the smooth-fibre locus. The natural generization homomorphism
\[
R^1f_*\Z_Y\longrightarrow
j_*\bigl((R^1f_*\Z_Y)|_{C^\circ}\bigr)
\]
is injective. Every nonzero global section of $R^1f_*\Z_Y$ is
nonzero in $H^1(f^{-1}(c);\Z)$ at every $c\in C$.
\end{proposition}


\begin{proof}
\noindent\textbf{Step 1.} It suffices to check a critical value $c$. Represent a stalk class
by $a\in H^1(V;\Z)$ for $V=f^{-1}(\Delta_c)$ and a sufficiently
small disc $\Delta_c$. Put $V^*=f^{-1}(\Delta_c\setminus\{c\})$.
Assume that the restriction of $a$ to a smooth fibre is zero.
The homotopy sequence for $V^*\to\Delta_c^*$ shows that
$a|_{V^*}$ is the pullback of an integer multiple $k$ of the
positive meridian class of the punctured disc.

\noindent\textbf{Step 2.} Choose a smooth point of one irreducible component of the central
divisor, away from all its other components. There the map has
the form $t=z^m u$ with $m>0$ and $u$ a holomorphic unit.
A small positive normal circle has base winding $m$ and bounds
its normal disc in $V$. Evaluating $a$ on that circle gives
$km=0$, whence $k=0$.
The inclusion $V^*\hookrightarrow V$ induces a surjection of
fundamental groups: a loop can be perturbed off a complex divisor,
including its higher-codimension singular strata. Therefore
$a|_{V^*}=0$ implies $a=0$.
Taking the direct limit over the discs proves the asserted stalk
injection; proper topological base change identifies its source
with $H^1(f^{-1}(c);\Z)$.

\noindent\textbf{Step 3.} If a global section vanished at one point, it would vanish on a
neighbourhood of that point. Its restriction to the connected
smooth locus is a section of a local system, so it would vanish
everywhere on that locus. The just-proved injection then makes
it zero everywhere, as required.
\end{proof}






\begin{proposition}\label{cdp:st:D.1.9}
\StatementRecord{cdp:st:D.1.9}{proved}{Historical events: Claude Fable 5.2 (Anthropic), effort high, 16 November 2026: carried as background without fault found. GPT-6 Astra (OpenAI), ChatGPT Pro mode, effort Pro, S6\_09\_PAPERS, 17 November 2026: targeted proof check. Also checked in this writing. Pre-continuation writing events: GPT-6 Astra (OpenAI), ChatGPT Pro mode, effort Pro, S6\_10\_PAPERS, 18 November 2026: targeted proof check. Continuation metadata events: GPT-6 Astra (OpenAI), ChatGPT Pro mode, effort Pro, S6\_10\_PAPERS continuation, 19 November 2026: used and rechecked in this run; status-vocabulary reconciliation; metadata reconciliation only; no new mathematical check. \emph{Pre-continuation recorded proof scope.} The transgression, coherent lift, irrational unitary character, and the passage to all Picard components by a finite power were re-derived}{the hypotheses, notation and definitions stated here; \cref{cdp:eq:D5}; \cref{cdp:st:D.1.8}; \cref{cdp:st:D.1.7}}

Let $Y$ be a connected compact complex manifold with
$b_1(Y)=b_2(Y)=0$, and let $f:Y\to\PP^1$ be a proper
holomorphic surjection with connected fibres. There is a class
$\beta\in H^1(Y,\OO_Y)$ such that, for every irrational real number
$r$, the line bundle
\[
L_r=\exp(2\pi i\,r\beta)
\]
has infinite order on every scheme-theoretic fibre of $f$.
Every connected component of $\Pic(Y)$ likewise contains a real
one-parameter family with this property at every irrational real
parameter.
No local-freeness assumption on $R^1f_*\OO_Y$ or constant-dimension
assumption on fibre cohomology is required.
\end{proposition}


\begin{proof}
\noindent\textbf{Step 1.} The low-degree integral Leray sequence and connectedness of the
fibres give an injection with finite cokernel
\[
d_2:H^0(\PP^1,R^1f_*\Z_Y)
 \hookrightarrow H^2(\PP^1;\Z)=\Z.                   \tag{D5}\label{cdp:eq:D5}
\]
Indeed $H^1(Y;\Z)=0$ and $H^2(Y;\Z)$ is finite.
Thus the image in \eqref{cdp:eq:D5} is $d\Z$ for an integer $d>0$; for the
stronger integral hypothesis $H^2(Y;\Z)=0$ one has $d=1$.
Choose $\gamma$ mapped to $d$. By Lemma \ref{cdp:st:D.1.8},
each $\gamma_c\in H^1(f^{-1}(c);\Z)$ is nonzero.
The coefficient morphism $\Z_Y\to\OO_Y$ sends $\gamma$ to a
global section of $R^1f_*\OO_Y$. Coherent Leray and
$H^1(\PP^1,\OO)=H^2(\PP^1,\OO)=0$ identify this space with
$H^1(Y,\OO_Y)$; let $\beta$ be its inverse image.

\noindent\textbf{Step 2.} For every full fibre $Y_c$, naturality of the edge maps and of
restriction gives
\[
\beta|_{Y_c}=\iota_c(\gamma_c)\in H^1(Y_c,\OO_{Y_c}), 
\]
where $\iota_c$ is the coefficient map. This uses the natural
coherent base-change homomorphism, but does not require it to be
an isomorphism. If $L_r|_{Y_c}$ had finite order, the exponential
sequence on the full fibre would give an integer $n>0$ and
$\lambda\in H^1(Y_c;\Z)$ with
\[
nr\,\iota_c(\gamma_c)=\iota_c(\lambda).
\]
Lemma \ref{cdp:st:D.1.7} implies $nr\gamma_c=\lambda$ in real cohomology.
The group $H^1(Y_c;\Z)$ is a finitely generated free abelian group,
and $\gamma_c$ is a nonzero element of it. A nonzero integral
coordinate of $\gamma_c$ now implies that $r$ is rational, a
contradiction. This works for every $c$ and proves the simultaneous
assertion in $\Pic^0(Y)$.

\noindent\textbf{Step 3.} For a component represented by a line bundle $A$, choose $N>0$
with $A^N\in\Pic^0(Y)$ and write $A^N=\exp(2\pi i b)$.
Then
\[
A_r=A\otimes\exp\bigl(2\pi i(r\beta-b)/N\bigr)
\]
belongs to that component and satisfies $A_r^N=L_r$. A group
element has finite order exactly when any fixed positive power of
it has finite order. Therefore each $A_r|_{Y_c}$ is non-torsion
for every irrational real $r$, proving the componentwise statement.
\end{proof}






\begin{proposition}\label{cdp:st:D.1.10}
\StatementRecord{cdp:st:D.1.10}{proved}{Historical events: Claude Fable 5.2 (Anthropic), effort high, 16 November 2026: carried as background without fault found. GPT-6 Astra (OpenAI), ChatGPT Pro mode, effort Pro, S6\_09\_PAPERS, 17 November 2026: targeted proof check. Also checked in this writing. Pre-continuation writing events: GPT-6 Astra (OpenAI), ChatGPT Pro mode, effort Pro, S6\_10\_PAPERS, 18 November 2026 04:42:38.084787+00:00: targeted proof check. Continuation metadata events: GPT-6 Astra (OpenAI), ChatGPT Pro mode, effort Pro, S6\_10\_PAPERS continuation, 19 November 2026: used and rechecked in this run; status-vocabulary reconciliation; metadata reconciliation only; no new mathematical check. \emph{Pre-continuation recorded proof scope.} proved. Full displayed logarithm, punctured-base coefficient and invariant-lattice span argument; no stronger surjectivity claim. Choose N annihilating finite H2(Y,Z); H1(Y,Z)=0 makes the exponential map H1(OY)\ensuremath{\to}Pic0(Y) an isomorphism. The actual Z\textasciicircum{}r\ensuremath{\to}Pic0 divisor map has kernel of rank s−1: every principal vertical divisor descends through the algebraic reduction to a rational function on P1, and full-fibre degree-zero divisors supply the entire rational relation space. D1.2 gives logarithm rank t\ensuremath{{}^{\prime}}. Canonical trivializations on U put those logarithms in the image of H1(U,Z); the E4 locally free coherent Leray bundle makes restriction of the logarithms injective. The punctured-base integral Leray sequence has cokernel H1(F,Z)\textasciicircum{}mon and base kernel. The base coefficient image vanishes by Stein cohomology; D1.7 makes the remaining integral coefficient map injective on a fibre. Thus the logarithms give a full-rank finite-index subgroup of exactly the invariant coefficient lattice, so its complex span lies in the restriction image. The final equivalence is explicitly only the lattice-spanning step of the earlier argument. Review provenance: New derivation after reading the supplied proof; not a blind review and not an independent human check}{the hypotheses, notation and definitions stated here; \cref{cdp:st:D.1.2}; \cref{cdp:st:E4}; \cref{cdp:st:D.1.7}}

Let $f:Y\to\PP^1$ be a holomorphic algebraic reduction of a compact
complex threefold with $b_1(Y)=b_2(Y)=0$ and $a(Y)=1$.
At a general smooth fibre, the image of the restriction
$H^1(Y,\OO_Y)\to H^1(F,\OO_F)$ contains the complex span of
the coefficient homomorphism
\[
\iota_F:H^1(F;\Z)^{\pi_1(C^\circ)}
 \longrightarrow H^1(F,\OO_F),                         \tag{D4}\label{cdp:eq:D4}
\]
which is induced by $\Z_F\hookrightarrow\OO_F$.
The final spanning step of CDP's rank argument reaches the whole
target exactly when this span is $H^1(F,\OO_F)$. Surjectivity of
the full restriction map need not have its image generated only by
this lattice, so that last equivalence describes the stated argument.
\end{proposition}


\begin{proof}
\noindent\textbf{Step 1.} Let $D_1,\ldots,D_r$ be the prime components above $\Delta$, and
choose $N>0$ annihilating the finite group $H^2(Y;\Z)$.
The logarithms, in the exponential identification
$\Pic^0(Y)=H^1(Y,\OO_Y)$, of
$\OO_Y(ND_1),\ldots,\OO_Y(ND_r)$ generate a free abelian
subgroup of rank $r-s+1=t'$.
Here is the rank computation with its maps specified. The map is
$\Z^r\to\Pic^0(Y)$, $(n_i)\mapsto\OO_Y(N\sum n_iD_i)$.
A relation means that $N\sum n_iD_i$ is the divisor of a
meromorphic function. Since $f$ is the algebraic reduction, that
function is $f^*h$ for $h\in\mathcal M(\PP^1)$.
Its divisor is a sum of full fibres with total base degree zero.
Such relations have rank $s-1$, and those given by divisors of
rational functions supported on $\Delta$ already have this rank.
Statement \ref{cdp:st:D.1.2} gives the last equality.

\noindent\textbf{Step 2.} Over $U=f^{-1}(C^\circ)$ each of these divisor line bundles has
its canonical trivialization. Exactness of the exponential sequence
therefore places its logarithm in the coefficient image of
$H^1(U;\Z)$. Coherent Leray identifies restriction
$H^1(Y,\OO_Y)\to H^1(U,\OO_U)$ with restriction of global
sections of $E$; this is injective because \ref{cdp:st:E4} makes $E$ locally free.
Thus the rank-$t'$ group of logarithms remains injective on $U$.

\noindent\textbf{Step 3.} Topological Leray on the punctured curve gives
\[
0\to H^1(C^\circ;\Z)\to H^1(U;\Z)
\to H^1(F;\Z)^{\pi_1(C^\circ)}\to0.
\]
The first group's coefficient image vanishes because
$H^1(C^\circ,\OO)=0$. Conversely a class whose coefficient image
in $H^1(U,\OO)$ vanishes restricts to zero in $H^1(F,\OO_F)$;
Lemma \ref{cdp:st:D.1.7} makes its integral restriction to $F$ zero. Thus the
coefficient image of $H^1(U;\Z)$ identifies exactly with the
invariant lattice in \eqref{cdp:eq:D4}. The logarithm group maps to a subgroup
of the same rank $t'$, hence of finite index, in that lattice.
Its complex span equals the span of the entire lattice, and the
restriction image is a complex vector subspace. This proves the
assertion and the statement about the last spanning step.
\end{proof}






\begin{proposition}\label{cdp:st:D.1.11}
\StatementRecord{cdp:st:D.1.11}{proved}{Historical events: Claude Fable 5.2 (Anthropic), effort high, 16 November 2026: carried as background without fault found. GPT-6 Astra (OpenAI), ChatGPT Pro mode, effort Pro, S6\_09\_PAPERS, 17 November 2026: targeted proof check. Also checked in this writing. Pre-continuation writing events: GPT-6 Astra (OpenAI), ChatGPT Pro mode, effort Pro, S6\_10\_PAPERS, 18 November 2026: targeted proof check. Continuation metadata events: GPT-6 Astra (OpenAI), ChatGPT Pro mode, effort Pro, S6\_10\_PAPERS continuation, 19 November 2026: used and rechecked in this run; status-vocabulary reconciliation; metadata reconciliation only; no new mathematical check. \emph{Pre-continuation recorded proof scope.} The incidence estimate for rank at least two and the rank-one period-ratio argument were re-derived; the resulting exceptional set has measure zero}{the hypotheses, notation and definitions stated here; \cref{cdp:st:E4}; \cref{cdp:st:D.1.9}; \cref{cdp:st:D.1.7}; \cref{cdp:st:D.1.8}}

Let $f:Y\to\PP^1$ be a proper connected-fibre holomorphic map from
a connected compact complex manifold with $b_1(Y)=b_2(Y)=0$.
In every component of $\Pic(Y)$, the line
bundles which are non-torsion on every full fibre form a dense
$G_\delta$ set of full Lebesgue measure in its complex vector-space
coordinates. If evaluation of global sections of $R^1f_*\OO_Y$
has rank one, then $\Pic^0(Y)=\C$ and the exceptional set is
countable. For two-torus fibres this alternative holds whenever
the splitting has a negative summand.
In particular the general-line-bundle conclusion of CDP20
Proposition 5.6 holds with this precise quantifier, without the
nonnegative-summand assertion of Proposition 5.5.
\end{proposition}


\begin{proof}
\noindent\textbf{Step 1.} Let $E=R^1f_*\OO_Y$, a vector bundle by \ref{cdp:st:E4}, and use
$\gamma,\beta$ from \ref{cdp:st:D.1.9}. The section of $E$ corresponding to
$\beta$ has no zero. In fact its value maps naturally to the
nonzero class $\iota_c(\gamma_c)$ in the cohomology of every full
fibre; Lemmas \ref{cdp:st:D.1.7} and \ref{cdp:st:D.1.8} prove nonvanishing. Write
$E=\bigoplus_j\OO(a_j)$ and let $k$ count the nonnegative degrees.
Evaluation has constant rank $k$ onto their direct-summand subbundle
$E_+$. Nonvanishing of $\beta$ gives $k\geq1$. When $k=1$, its
sole nonnegative degree is zero, since a positive-degree line bundle
has no nowhere-zero section. In that case $H^0(E)=\C\beta$.

\noindent\textbf{Step 2.} Suppose first $k\geq2$, and put $V=H^1(Y,\OO_Y)=H^0(E)$ and
$N=\dim_{\C}V$. On a simply connected disc in the smooth-fibre
locus where degree-one base change holds, trivialize the integral
local system. For a flat integral
class $\lambda$, the equation
\[
z|_{F_c}=\iota_c(\lambda),\qquad (z,c)\in V\times\Delta
\]
defines an analytic set. The parameter values where
$\iota_c(\lambda)\in(E_+)_c$ form an analytic subset of the disc;
over it the fibres are affine spaces of dimension $N-k$.
The incidence set thus has complex dimension at most
$N-k+1\leq N-1$. Its projection to $V$ has Lebesgue
measure zero. It is also meagre: stratify and exhaust that set
by countably many compact sets; the images are closed sets with
empty interior, by their real dimension at most $2N-2$.
Take countably many discs, integral classes and positive integer
denominators to include every torsion condition.

\noindent\textbf{Step 3.} There are finitely many excluded fibres, including the exceptional
values of degree-one base change. At each, the natural map
$E\otimes\C(c)\to H^1(Y_c,\OO_{Y_c})$ is injective. This follows
directly from the long exact sequence for multiplication by a local
base parameter: its cokernel is the parameter-torsion in
$R^2f_*\OO_Y$. The restriction $V\to H^1(Y_c,\OO_{Y_c})$
therefore has rank $k\geq2$. The torsion conditions are contained in
countably many affine subspaces of complex codimension at least
two. Combining these with the preceding projections gives a
meagre null set. The complement is dense and $G_\delta$, as claimed.

\noindent\textbf{Step 4.} It remains to prove the case $k=1$.
Now $V=\C\beta$. On a simply connected disc in the same open locus,
the projection of a flat integral class $\lambda$ to
$E/(\OO\beta)$ is holomorphic. Unless it vanishes identically,
its zero set is discrete. Such zeroes give only countably many
parameters $z$ for which $\iota_c(\lambda)=z\beta_c$.
We prove that an identically vanishing projection gives a constant
parameter, including after analytic continuation.

\noindent\textbf{Step 5.} Pass to the universal cover of the smooth locus, on which
$H^1(F;\Z)$ is marked. The identity theorem continues the relation
$\iota_c(\lambda)\in\C\iota_c(\gamma)$ everywhere on this cover.
If $\lambda$ and $\gamma$ are real-dependent, their ratio is a
constant rational number. Otherwise let
$W_{\Q}=\Q\gamma+\Q\lambda$. The coefficient map
$H^1(F;\R)\to H^1(F,\OO_F)$ is real-linearly injective by \ref{cdp:st:D.1.7}.
The images of the two real-independent classes span the entire
real two-plane underlying $\C\beta_c$, so $W_{\R}$ is precisely
its inverse image. Analytic continuation around a base loop
sends $\lambda$ to a class whose coefficient image again belongs
to this line. Hence it preserves $W_{\Q}$ and fixes $\gamma$.

\noindent\textbf{Step 6.} Choose an integral basis $(\gamma_0,\lambda_0)$ of
$W_{\Q}\cap H^1(F;\Z)$ with $\gamma$ a positive integral multiple
of $\gamma_0$. The induced monodromy on this lattice fixes
$\gamma_0$ and preserves the orientation of $W_{\R}$: the complex
structure induced on that plane by the coefficient map varies
continuously along the loop, so its complex orientation cannot change.
It therefore sends $\lambda_0$ to $\lambda_0+k\gamma_0$ for an
integer $k$. Put
\[
\tau(c)=\frac{\iota_c(\lambda_0)}{\iota_c(\gamma_0)}.
\]
This holomorphic function on the universal cover never takes a
real value, by real injectivity of the coefficient map, and its
imaginary part has constant sign. Monodromy changes $\tau$ by
integers. Accordingly either $\exp(2\pi i\tau)$ or
$\exp(-2\pi i\tau)$ descends to a holomorphic function of absolute
value less than one on $\PP^1$ minus finitely many points. The
removable singularity theorem extends it to $\PP^1$, where it is
constant. Its logarithmic derivative gives $d\tau=0$. Thus the
original ratio is constant as asserted.

\noindent\textbf{Step 7.} There are countably many integral classes on a countable disc
cover. The isolated-zero alternatives and the constant alternatives
just proved give a countable exceptional set on smooth fibres.
Each exceptional full fibre contributes another countable set,
because $z\mapsto z\beta_c$ is injective and its intersection with
the rational integral coefficient image is countable. This proves
the assertion in $\Pic^0(Y)$.

\noindent\textbf{Step 8.} Finally for any component $A\Pic^0(Y)$, the map taking the $N$th
power, for an integer $N$ killing its component class, is the affine
isomorphism $z\mapsto b+Nz$ of the complex parameter space used
in \ref{cdp:st:D.1.9}. It preserves the torsion condition in every fibre, and
preserves countability, meagreness and measure-zero sets. This
proves all the componentwise assertions.
\end{proof}



The quantifier in \ref{cdp:st:D.1.11} is explicit. In the globally generated
two-torus case, the torsion parameters at even one fixed smooth fibre are
dense: the rational period lattice is dense in
$H^1(F,\OO_F)$, and restriction onto that space is surjective.
Thus the good set in that case cannot contain a nonempty open
subset of the Picard parameter space. A dense $G_\delta$ set of
full measure is the appropriate conclusion of the incidence
argument, rather than a Zariski-open conclusion.




\begin{proposition}\label{cdp:st:D.1.12}
\StatementRecord{cdp:st:D.1.12}{proved}{Historical events: Claude Fable 5.2 (Anthropic), effort high, 16 November 2026: carried as background without fault found. GPT-6 Astra (OpenAI), ChatGPT Pro mode, effort Pro, S6\_09\_PAPERS, 17 November 2026: targeted proof check. Also checked in this writing. Pre-continuation writing events: GPT-6 Astra (OpenAI), ChatGPT Pro mode, effort Pro, S6\_10\_PAPERS, 18 November 2026: targeted proof check; used and rechecked in this run. Continuation metadata events: GPT-6 Astra (OpenAI), ChatGPT Pro mode, effort Pro, S6\_10\_PAPERS continuation, 19 November 2026: used and rechecked in this run; status-vocabulary reconciliation; metadata reconciliation only; no new mathematical check. \emph{Pre-continuation recorded proof scope.} The Euler degree relation, generically nonzero determinant cup map on tori and primary Kodaira surfaces, torsion elimination, and full-fibre base-change dimensions were re-derived. Borcea's invariant differential forms were opened on printed page376; the direct-image theorem remains its separately cited input}{the hypotheses, notation and definitions stated here; \cref{cdp:st:E4}; \cite[\S3, equations (9)--(12)]{Borcea1984}}

For a proper connected-fibre holomorphic map $f:Y\to C$ from a
compact complex threefold with $b_2(Y)=0$ and general fibre a
two-torus or a primary Kodaira surface, the cup-product homomorphism
is an isomorphism
\[
\bigwedge^2R^1f_*\OO_Y\xrightarrow{\simeq}R^2f_*\OO_Y.
\]
Every full fibre has structure-sheaf cohomology dimensions
$(h^0,h^1,h^2)=(1,2,1)$, and these two direct images satisfy
base change at every fibre. This proves these two cases of CDP20
Proposition 5.1 by an index and cup-product argument.
\end{proposition}



The torus-fibred six-sphere case and the cup-product index method are \citet[Theorem~2.3, pp.~5--6]{Viaclovsky2026Combs}. The statement here has the displayed broader hypotheses: $b_2=0$, an arbitrary base curve, and also primary Kodaira general fibres. The complete proof specifies the additional integral comparison for those hypotheses.

\begin{proof}
\noindent\textbf{Step 1.} Write $E=R^1f_*\OO_Y$, which is locally free of rank two by \ref{cdp:st:E4},
and $G=R^2f_*\OO_Y$, which has rank one. Put
$d=\deg E$, $g=\deg(G/G_{\rm tor})$, and
$\ell=\operatorname{length}(G_{\rm tor})$.
Since $c_1(Y)=0$ rationally, Hirzebruch--Riemann--Roch gives
$\chi(Y,\OO_Y)=\int_Yc_1c_2/24=0$.
Coherent Leray and Riemann--Roch on $C$ give, with the rank terms
cancelled,
\[
0=\chi(\OO_C)-\chi(E)+\chi(G)=g+\ell-d.
\]
The cup product $\det E\to G$ is an isomorphism on a general
fibre. For a two-torus this is its exterior-algebra description.
For a primary Kodaira surface, use invariant $(1,0)$-forms
$\alpha,\eta$ with $d\alpha=0$ and
$d\eta=c\alpha\wedge\overline\alpha$, $c\ne0$.
Their conjugates are $\overline\partial$-closed. The class
$[\overline\alpha\wedge\overline\eta]$ is nonzero because its
Serre pairing with the holomorphic two-form $\alpha\wedge\eta$
is the nonzero integral of a volume form. This also proves that
$[\overline\alpha],[\overline\eta]$ are independent; since
$h^1(\OO)=2$ and $h^2(\OO)=1$, it proves the required cup-product
isomorphism. The invariant forms and their differential are
the forms in \cite[\S3, equations (9)--(12)]{Borcea1984}.
The composite of $\det E\to G$ with the torsion-free quotient is a
nonzero homomorphism of line bundles, so $g\geq d$. The displayed
identity and $\ell\geq0$ force $g=d$, $\ell=0$. A nonzero map
of line bundles of equal degree has no zero divisor, proving the
claimed global isomorphism.

\noindent\textbf{Step 2.} Top-degree base change gives
$G\otimes\C(c)\simeq H^2(Y_c,\OO_{Y_c})$, of dimension one.
The multiplication-parameter exact sequence then identifies
$E\otimes\C(c)$ with $H^1(Y_c,\OO_{Y_c})$, since $G$ has no
parameter-torsion. \ref{cdp:st:E4} gives $h^0=1$. These arguments retain the
scheme-theoretic fibres throughout.
\end{proof}






\begin{proposition}\label{cdp:st:D.1.4}
\StatementRecord{cdp:st:D.1.4}{proved}{Historical events: Claude Fable 5.2 (Anthropic), effort high, 16 November 2026: carried as background without fault found. GPT-6 Astra (OpenAI), ChatGPT Pro mode, effort Pro, S6\_09\_PAPERS, 17 November 2026: targeted proof check. Also checked in this writing. Pre-continuation writing events: GPT-6 Astra (OpenAI), ChatGPT Pro mode, effort Pro, S6\_10\_PAPERS, 18 November 2026: targeted proof check. Continuation metadata events: GPT-6 Astra (OpenAI), ChatGPT Pro mode, effort Pro, S6\_10\_PAPERS continuation, 19 November 2026: used and rechecked in this run; status-vocabulary reconciliation; metadata reconciliation only; no new mathematical check. \emph{Pre-continuation recorded proof scope.} The actual Leray edge map, evaluation of the split direct-image bundle, and fibre base-change map were composed explicitly. A summand O(a) contributes to the evaluation image exactly for a nonnegative degree}{the hypotheses, notation and definitions stated here; \cref{cdp:st:D.1.2}; \cref{cdp:st:E4}}

For any $f:Y\to C$ as in \ref{cdp:st:D.1.2} put $E=R^1f_*\OO_Y$.
The coherent Leray sequence gives
\[
0\longrightarrow H^1(C,\OO_C)\xrightarrow{f^*}H^1(Y,\OO_Y)
\xrightarrow{e}H^0(C,E)\longrightarrow0.               \tag{D2}\label{cdp:eq:D2}
\]
At a smooth fibre where cohomological base change holds, the
restriction homomorphism is exactly the composite
\[
H^1(Y,\OO_Y)\xrightarrow{e}H^0(C,E)
\xrightarrow{\operatorname{ev}_c}E\otimes\C(c)
\xrightarrow{\simeq}H^1(F_c,\OO_{F_c}).               \tag{D3}\label{cdp:eq:D3}
\]
By Statement \ref{cdp:st:E4}, $E$ is locally free. For $C=\PP^1$, write
$E=\bigoplus_{j=1}^q\OO(a_j)$.
Then the rank of \eqref{cdp:eq:D3} at a general point is exactly
$\#\{j:a_j\geq0\}$. In particular, the nonnegative splitting
degrees claimed by CDP20 Proposition 5.5 are equivalent to its
claimed general-fibre surjectivity when $E$ has rank two and is
locally free.
\end{proposition}


\begin{proof}
\noindent\textbf{Step 1.} Stein factorization gives $f_*\OO_Y=\OO_C$. A curve has no coherent
cohomology in degree two, so the low-degree Leray sequence is \eqref{cdp:eq:D2}.
The edge map is induced by restriction to inverse images of open
sets in $C$. Restricting these same representatives to $F_c$ gives
\eqref{cdp:eq:D3}, which proves the asserted identification of the actual maps.
Statement \ref{cdp:st:E4} proves local freeness by resolving the fibre support,
normalizing a cyclic base change with reduced central fibre, and
taking the invariant direct image. For a line bundle
$\OO(a)$ on $\PP^1$, its evaluation has rank one at every point
for $a\geq0$, and its space of global sections is zero for $a<0$.
Applying this separately to the direct summands gives the formula.
\end{proof}






\begin{proposition}[Refuted printed surjectivity and positivity, with a second example]\label{cdp:st:D.1.13}
\StatementRecord{cdp:st:D.1.13}{proved}{Historical events: Claude Fable 5.2 (Anthropic), effort high, 16 November 2026: carried as background without fault found. GPT-6 Astra (OpenAI), ChatGPT Pro mode, effort Pro, S6\_09\_PAPERS, 17 November 2026: targeted proof check. Also checked in this writing. Pre-continuation writing events: GPT-6 Astra (OpenAI), ChatGPT Pro mode, effort Pro, S6\_10\_PAPERS, 18 November 2026 04:50:08.441928+00:00: targeted proof check. GPT-6 Astra (OpenAI), ChatGPT Pro mode, effort Pro, S6\_10\_PAPERS, 18 November 2026 04:51:18.698751+00:00: used and rechecked in this run. Continuation metadata events: GPT-6 Astra (OpenAI), ChatGPT Pro mode, effort Pro, S6\_10\_PAPERS continuation, 19 November 2026: used and rechecked in this run; status-vocabulary reconciliation; metadata reconciliation only; no new mathematical check. \emph{Pre-continuation recorded proof scope.} proved counterexample; rejected printed surjectivity and positivity remain refuted. Full canonical-line, relative-duality and coherent-splitting argument on the actual Wang fillings. Re-derived the M× canonical map \ensuremath{\omega}S\ensuremath{\mapsto}\ensuremath{\omega}S\ensuremath{\wedge}dlogv and its invariance under the actual −6P lift. The unimodular rank-one map (z,v)\ensuremath{\mapsto}(v,v\textsuperscript{2}/z) converts the degenerating vector to a primitive unit vector; dq\ensuremath{\wedge}dlogu=−du\ensuremath{\wedge}dv at uv=q gives the extension. At the rank-two chart, q=z0z1z2 and the unimodular character determinant give q\ensuremath{\prod}dlogzi=\ensuremath{\prod}dzi. At infinity, s=r\textsuperscript{3}, x\_s=r\textsuperscript{4}X, y\_s=r\textsuperscript{6}Y gives ds\ensuremath{\wedge}dx\_s/y\_s=3dr\ensuremath{\wedge}dX/Y; the order3 deck weights multiply to one, and the actual line-coordinate rescaling v\ensuremath{{}^{\prime}}=r\textsuperscript{2}v adds only a dr term to dlogv. The affine logarithmic translation likewise preserves the form, so KY=f*O(−1) holds through every filling. Flat Gorenstein relative duality and the actual D1.12 cup/base-change isomorphism give R\textsuperscript{2}=detR\textsuperscript{1}=O(−1). D1.9’s nonzero full-fibre coefficient section is a subbundle O; Ext\textsuperscript{1}(O(−1),O)=H\textsuperscript{1}O(1)=0 splits R\textsuperscript{1}. Leray gives(1,1,0,0), with restriction rank1 on each smooth torus. The integral source topology was rechecked in D1.6. A dominating divisor would restrict to a nonzero effective topologically trivial divisor on a Kahler fibre, impossible by its positive intersection; this gives the actual meromorphic descent to P1 and a=1. Any biholomorphism to the sphere member would induce an H2 isomorphism, contradicting Z/2 versus0. Review provenance: New derivation after reading the supplied proof; not a blind review and not an independent human check}{the hypotheses, notation and definitions stated here; \cref{cdp:st:D.1.6}; \cref{cdp:st:D.1.12}; \cref{cdp:st:D.1.9}; \cref{cdp:st:D.1.4}; \cref{cdp:st:E1}}


For Wang's threefold $Y=X_2(\lambda)$ in \ref{cdp:st:D.1.6} one has
\[
K_Y\simeq f^*\OO_{\PP^1}(-1),\qquad
R^1f_*\OO_Y\simeq\OO\oplus\OO(-1),\qquad
R^2f_*\OO_Y\simeq\OO(-1).
\]
Consequently $h^q(Y,\OO_Y)=(1,1,0,0)$, and restriction to every
smooth two-torus fibre has rank one. This threefold satisfies
$H^1(Y;\Z)=H^2(Y;\Z)=0$, $a(Y)=1$, and is a second
counterexample to CDP20 Corollary 5.2 and Proposition 5.5.
Its integral homology torsion distinguishes it from the retained
integral homology sphere.
\end{proposition}


\begin{proof}
\noindent\textbf{Step 1.} We compute the canonical line through the actual local fillings.
The rational elliptic surface $\pi:S\to\PP^1$ in this construction
has Weierstrass equation
\[
y^2=x^3-3(1+t)x+(2+3t+2t^2)
\]
and fundamental line bundle $\OO(1)$; hence
$K_S=\pi^*\OO(-1)$. On the complement of the zero section in
$M=\OO_S(P-O)$, the determinant of the relative cotangent sequence
identifies its canonical line with the pullback of $K_S$.
In a local multiplicative coordinate $v$, this identification sends
a local two-form $\omega_S$ to $\omega_S\wedge d\log v$.
Changing $v$ by a nonzero holomorphic function on $S$ adds a
one-form from $S$ to $d\log v$, whose wedge with $\omega_S$ is zero.
This verifies the coordinate independence of the identification.
The contracting generator covers translation by $-6P$ on the
elliptic fibres. Translation preserves the relative elliptic
differential, and the extra base terms disappear after wedging
with the base differential. Its line-bundle lift changes $d\log v$
only by a base one-form. It therefore preserves this canonical
identification. The latter descends to
$K_Y=f^*\OO(-1)$ on the open quotient.

\noindent\textbf{Step 2.} At either rank-one filling, let $q$ be the base coordinate and
$(x,y)$ the two multiplicative coordinates of the construction.
At $t=0$ the explicit change is $x=v$, $y=v^2/z$; its exponent
matrix is unimodular. At $\alpha_+$ no change is required.
Thus the preceding three-form is a nonzero holomorphic unit times
\[
dq\wedge d\log x\wedge d\log y.
\]
The nondegenerating $y$ direction becomes an elliptic curve after
quotienting by its contracting period. Its logarithmic differential
is a nonvanishing relative one-form; changes by a base differential
do not affect the displayed wedge. At a double curve the equation
is $uv=q$, and
$dq\wedge d\log u=-du\wedge dv$. Therefore the three-form extends
without a zero or pole over the full rank-one filling.

\noindent\textbf{Step 3.} At $\alpha_-$ the rank-two filling uses a unimodular period lattice
and the standard toric charts $q=z_0z_1z_2$. A pair of independent
fibre characters, together with $q$, is a unimodular system of
characters in each smooth chart. The same form is, up to a nonzero
constant and a holomorphic unit,
\[
q\,d\log z_0\wedge d\log z_1\wedge d\log z_2
=dz_0\wedge dz_1\wedge dz_2.
\]
It extends without zero or pole, including at the triple points.

\noindent\textbf{Step 4.} It remains to check the order-three logarithmic filling at infinity.
Put $s=1/t$, $x_s=s^2x$ and $y_s=s^3y$. Its Weierstrass
coefficients have orders three and four in $s$, giving type $IV^*$.
On the cubic base change $s=r^3$, put
$x_s=r^4X$, $y_s=r^6Y$. The good-reduction equation is
\[
Y^2=X^3-3r(r^3+1)X+(2+3r^3+2r^6),
\]
whose central elliptic curve is smooth. The canonical form satisfies
\[
ds\wedge\frac{dx_s}{y_s}
=3\,dr\wedge\frac{dX}{Y}.
\]
It has no zero or pole. For $\zeta=e^{2\pi i/3}$ the root deck
action has weights $r\mapsto\zeta r$, $X\mapsto\zeta^2X$,
$Y\mapsto Y$; the weights on $dr$ and $dX/Y$ multiply to one.
The additional invariant translation used in the logarithmic
transform preserves the holomorphic two-form on the torus cover.
Hence the three-form is invariant and descends through the free
cyclic action as a nonvanishing three-form. This proves that the
canonical identification from the open quotient extends everywhere,
giving $K_Y=f^*\OO(-1)$.

\noindent\textbf{Step 5.} The map is flat and Gorenstein, since its source is smooth and its
base is a smooth curve. Relative Serre duality gives
\[
(R^2f_*\OO_Y)^*=f_*\omega_{Y/\PP^1}
=f_*f^*\OO(1)=\OO(1).
\]
Statement \ref{cdp:st:D.1.12} gives local freeness and
$\det(R^1f_*\OO_Y)=R^2f_*\OO_Y=\OO(-1)$.
The nowhere-zero section from \ref{cdp:st:D.1.9} embeds $\OO$ as a subbundle
of $R^1f_*\OO_Y$ with quotient $\OO(-1)$. The extension splits
because $H^1(\PP^1,\OO(1))=0$. Coherent Leray now computes all
four displayed structure-sheaf cohomology groups, and \ref{cdp:st:D.1.4}
identifies the degree-one restriction with rank-one evaluation.

\noindent\textbf{Step 6.} The homology computation in \ref{cdp:st:D.1.6} and the universal coefficient
theorem give $H^1(Y;\Z)=H^2(Y;\Z)=0$; specifically the torsion
in $H_2$ contributes to $H^3$, not to $H^2$.
The general fibre is K\"ahler and $b_2(Y)=0$, so \ref{cdp:st:E1} proves that
this map is the algebraic reduction and $a(Y)=1$.
The retained six-sphere has zero second integral homology, whereas
$H_2(Y;\Z)=\Z/2\Z$. An isomorphism of complex manifolds would
induce an isomorphism of these groups, which is impossible.
This completes the independent counterexample.
\end{proof}






\begin{proposition}[Proved comparison with identical integral monodromy]\label{cdp:st:D.1.15}
\StatementRecord{cdp:st:D.1.15}{proved}{Historical events: Claude Fable 5.2 (Anthropic), effort high, 16 November 2026: carried as background without fault found. GPT-6 Astra (OpenAI), ChatGPT Pro mode, effort Pro, S6\_09\_PAPERS, 17 November 2026: targeted proof check. Also checked in this writing. Pre-continuation writing events: GPT-6 Astra (OpenAI), ChatGPT Pro mode, effort Pro, S6\_10\_PAPERS, 18 November 2026 04:42:38.084787+00:00: targeted proof check. Continuation metadata events: GPT-6 Astra (OpenAI), ChatGPT Pro mode, effort Pro, S6\_10\_PAPERS continuation, 19 November 2026: used and rechecked in this run; status-vocabulary reconciliation; metadata reconciliation only; no new mathematical check. \emph{Pre-continuation recorded proof scope.} proved. Full explicit two-family construction, real period basis, holomorphic quotient maps and coherent splitting. The real-linear period matrix has determinant −(1+|t|\textsuperscript{2})\textsuperscript{2}, and its O(1) frame at infinity has a real basis. These actual global sections define a free properly discontinuous lattice action on O(1)\textasciicircum{}2; their real frame gives a smooth product trivialization, proving properness/compactness and the constant marked integral local system. Fibrewise Hodge decomposition on a complex torus identifies the kernel of the coefficient map with the image of P\textasciicircum{}t:O(−1)\textasciicircum{}2\ensuremath{\to}O\textasciicircum{}4. The independently computed identity BP\textasciicircum{}t=0, gcd1 of the six polynomial minors, and rank2 of the infinity leading coefficients give the exact global sequence with quotient O(1)\textasciicircum{}2. This is the actual Dolbeault coefficient quotient, not merely an equal-rank comparison. The constant family has the same fixed lattice and quotient H1(F,O)\ensuremath{\otimes}O=O\textsuperscript{2}. Both smooth products have b1=4. Exact output is analysis/II\_independent\_batch2\_exact.json. Review provenance: New derivation after reading the supplied proof; not a blind review and not an independent human check}{the hypotheses, notation and definitions stated here; the explicit constructions and calculations in the proof}


Two smooth proper families of two-tori over $\PP^1$ have the same
constant marked $H^1(-;\Z)$ local system but have
$R^1f_*\OO=\OO^2$ and $R^1f_*\OO=\OO(1)^2$, respectively.
Thus the integral representation alone does not determine the
coherent splitting for general compact torus fibrations. Both
examples have $b_1=4$; this comparison does not assert a new
counterexample under $b_1=b_2=0$.
\end{proposition}


\begin{proof}
\noindent\textbf{Step 1.} Let $V=\OO_{\PP^1}(1)^{\oplus2}$. On the affine coordinate $t$,
four global sections of $V$ are the columns of
\[
P(t)=\begin{pmatrix}1&i&t&it\\-t&it&1&-i\end{pmatrix}.
\]
They are a real basis in every fibre. Indeed, writing
$z_1=a+ib$, $z_2=c+id$, their real linear combination is
$(z_1+t z_2,\overline z_2-t\overline z_1)$.
Its vanishing gives $(1+|t|^2)\overline z_2=0$ and then $z_1=0$.
At infinity, the $\OO(1)$ frame changes $P(t)$ to $t^{-1}P(t)$,
whose leading columns again form a real basis. The lattice
generated by these sections therefore acts freely and properly by
holomorphic fibre translations on $V$. Its quotient $Z$ is a
compact smooth holomorphic torus bundle over $\PP^1$.
The displayed real basis identifies it differentiably with
$\PP^1\times(\R^4/\Z^4)$ and marks its constant integral local
system.

\noindent\textbf{Step 2.} The fibrewise period map includes the bundle of holomorphic
one-forms by $P(t)^t$. The Dolbeault coefficient map gives the
exact sequence
\[
0\longrightarrow\OO(-1)^2
 \xrightarrow{P(t)^t}\OO^4
 \longrightarrow R^1f_*\OO_Z\longrightarrow0.
\]
The last map can be represented explicitly by
\[
B(t)=\begin{pmatrix}-it&-t&i&1\\-i&1&-it&t\end{pmatrix}.
\]
The identity $B(t)P(t)^t=0$ holds polynomially. Its two rows have
rank two at every finite $t$, as is seen from their two-by-two
minors, and their leading coefficients have rank two at infinity.
Thus this is the exact complex of vector bundles
\[
0\longrightarrow\OO(-1)^2\longrightarrow\OO^4
 \xrightarrow{B}\OO(1)^2\longrightarrow0,
\]
proving $R^1f_*\OO_Z\simeq\OO(1)^2$. The polynomial identity is
also checked in the attached exact program.

\noindent\textbf{Step 3.} For the constant family $Z_0=Z_{t=0}\times\PP^1$, the same
marked fibre lattice gives the identical constant integral local
system, whereas coherent K\"unneth gives
$R^1f_*\OO_{Z_0}=H^1(Z_{t=0},\OO)\otimes\OO=\OO^2$.
The different splittings are therefore compared by an explicit
identity of the integral representations and explicit coherent
quotient maps.
\end{proof}






\section{Complete smooth-fibre classification and actual Hopf occurrence}\label{sec:II:03}


\begin{proposition}[Refuted surface assertion, with explicit example]\label{cdp:st:C.3.2}
\StatementRecord{cdp:st:C.3.2}{proved}{Historical events: Claude Fable 5.2 (Anthropic), effort high, 16 November 2026: carried as background without fault found. GPT-6 Astra (OpenAI), ChatGPT Pro mode, effort Pro, S6\_09\_PAPERS, 17 November 2026: targeted proof check. Also checked in this writing. Pre-continuation writing events: GPT-6 Astra (OpenAI), ChatGPT Pro mode, effort Pro, S6\_10\_PAPERS, 18 November 2026 04:23:40Z: targeted proof check. Continuation metadata events: GPT-6 Astra (OpenAI), ChatGPT Pro mode, effort Pro, S6\_10\_PAPERS continuation, 19 November 2026: used and rechecked in this run; status-vocabulary reconciliation; metadata reconciliation only; no new mathematical check. \emph{Pre-continuation recorded proof scope.} proved as stated. Full secondary Kodaira example, complex structure, free cyclic quotient, integral cohomology and canonical order. Re-derived the Heisenberg lattice law and sigma-preservation polynomial, sigma\textasciicircum{}6=1, sigma*alpha=tau·alpha and sigma*beta=beta. The type-(1,1) differential proves integrability, and alpha wedge beta gives the primary canonical trivialization. The t/6 shift makes every nonidentity quotient power free. Abelianization of Gamma semidirect Z uses det(A−I)=1 and gives H1=Z; Euler, duality and UCT give H2=0. The primitive sixth-root action on the canonical frame proves exact holomorphic order6. The exact matrix/lattice replay passed in II\_independent\_D\_exact\_checks.out. Review provenance: New derivation after reading the supplied proof; not a blind review and not an independent human check}{the hypotheses, notation and definitions stated here; the explicit constructions and calculations in the proof}


The assertion that every secondary Kodaira surface has torsion in
$H_1(-;\Z)$ is false. There is a secondary Kodaira surface $Q$ with
\[
H_1(Q;\Z)=\Z,\qquad H^2(Q;\Z)=0,
\]
whose canonical bundle has holomorphic order six.
\end{proposition}


\begin{proof}
\noindent\textbf{Step 1.} Let $H$ be the real Heisenberg group in coordinates $(x,y,z)$ with law
\[
(x,y,z)(x',y',z')
=\left(x+x',y+y',z+z'+\frac{xy'-yx'}2\right).
\]
Let $\Gamma\subset H$ be generated by
$a=(1,0,0)$, $b=(0,1,0)$ and $c=(0,0,1)$. Thus
\[
\Gamma=\{(x,y,z):x,y\in\Z,\ z-xy/2\in\Z\},
\qquad [a,b]=c.
\]
It is a discrete cocompact subgroup. On $H\times\mathbb R_t$ write
\[
\tau=e^{\pi i/3},\quad
\eta=dz-\frac{x\,dy-y\,dx}{2},\quad
\alpha=dx+\tau\,dy,\quad
\beta=\eta-\frac{dx+dy}{2}+i\,dt.
\]
These two complex one-forms and their conjugates are linearly
independent over $\C$. Declaring $\alpha,\beta$ to have type $(1,0)$
defines a left-invariant almost complex structure. It is integrable:
$d\alpha=0$ and $d\beta=-dx\wedge dy$ has type $(1,1)$, so the
ideal generated by the $(1,0)$ forms has no $(0,2)$ obstruction.
The structure descends to
$P=\Gamma\backslash H\times(\mathbb R/\Z)$.
The nonzero form $\alpha\wedge\beta$ is $\bar\partial$-closed and
trivializes $K_P$. The projection induced by $x+\tau y$ is a
holomorphic elliptic bundle over $\C/(\Z+\tau\Z)$, with nonzero
central Chern class $[a,b]=c$. Its first Betti number is three, so
$P$ is a primary Kodaira surface.

\noindent\textbf{Step 2.} Define the Lie-group automorphism
\[
\sigma(x,y,z)=(-y,x+y,z-y/2).
\]
It preserves $\Gamma$: the difference between the new central
coordinate and half the product of the new horizontal coordinates
is
\[
(z-xy/2)+xy+y(y-1)/2.
\]
This is integral on $\Gamma$. The inverse preserves the lattice too,
and direct substitution gives $\sigma^6=1$. Moreover
$\sigma^*\alpha=\tau\alpha$ and $\sigma^*\beta=\beta$.
Consequently
\[
g:P\longrightarrow P,\qquad
g([x,y,z],t)=([\sigma(x,y,z)],t+1/6)
\]
is a holomorphic automorphism of order six. Every nonidentity power
is fixed-point-free because of its translation of the $t$-circle.
The quotient $Q=P/\langle g\rangle$ is therefore a compact complex
surface and $P\to Q$ is finite \emph{\'etale} of degree six.

\noindent\textbf{Step 3.} The contractible manifold $H\times\mathbb R$ is its universal cover.
The deck group is $\Gamma\rtimes_\sigma\Z$, where the infinite
generator is the lift $(\sigma,t\mapsto t+1/6)$; its sixth power is
the original integral translation in $t$. Since
$\Gamma^{\rm ab}=\Z^2$ and the induced matrix is
\[
A=\begin{pmatrix}0&-1\\1&1\end{pmatrix},\qquad
\det(A-I)=1,
\]
abelianization gives $H_1(Q;\Z)=\Z$.
The finite cover and $e(P)=0$ give $e(Q)=0$. Poincar\'e duality then
gives $b_2(Q)=0$, and the universal coefficient theorem gives
$H^2(Q;\Z)=0$ because $H_1(Q;\Z)$ has no torsion.
Finally $g^*(\alpha\wedge\beta)=\tau\alpha\wedge\beta$.
Every holomorphic section of $K_P^m$ is a constant multiple of
$(\alpha\wedge\beta)^m$, by compactness. Such a section descends
precisely when $6$ divides $m$. Thus $K_Q$ has holomorphic order six.
The surface has $b_1=1$ and a primary Kodaira surface as a finite
\'etale cover, so it is secondary Kodaira. The matrix and lattice
calculations have been checked exactly by the attached program.
\end{proof}






\begin{proposition}[Refuted surface assertion, with explicit example]\label{cdp:st:C.3.3}
\StatementRecord{cdp:st:C.3.3}{proved}{Historical events: Claude Fable 5.2 (Anthropic), effort high, 16 November 2026: carried as background without fault found. GPT-6 Astra (OpenAI), ChatGPT Pro mode, effort Pro, S6\_09\_PAPERS, 17 November 2026: targeted proof check. Also checked in this writing. Pre-continuation writing events: GPT-6 Astra (OpenAI), ChatGPT Pro mode, effort Pro, S6\_10\_PAPERS, 18 November 2026 04:23:40Z: targeted proof check. Continuation metadata events: GPT-6 Astra (OpenAI), ChatGPT Pro mode, effort Pro, S6\_10\_PAPERS continuation, 19 November 2026: used and rechecked in this run; status-vocabulary reconciliation; metadata reconciliation only; no new mathematical check. \emph{Pre-continuation recorded proof scope.} proved as stated. Full secondary Hopf example, with the exact finite group certificate actually rerun. Inspected the rational-pair Q(sqrt5) quaternion implementation and reran all14400 products and the commutator computation: 120 unit elements, closure, and commutator set size120. The displayed SU2 representation respects Hamilton multiplication and has no nonidentity element fixing a nonzero vector. Radial scaling gives the free proper product deck action and the explicit quotient S1×S3/G. Its universal cover and perfect finite factor yield H1=Z, while Euler, duality and UCT yield H2=0. The finite primary cover has degree120; the noncyclic deck group makes the quotient secondary. Review provenance: New derivation after reading the supplied proof; not a blind review and not an independent human check}{the hypotheses, notation and definitions stated here; the explicit constructions and calculations in the proof}


The assertion that a secondary Hopf surface has torsion in
$H_1(-;\Z)$ is false. There is a secondary Hopf surface $Q_H$ with
$H_1(Q_H;\Z)=\Z$ and $H^2(Q_H;\Z)=0$.
\end{proposition}


\begin{proof}
\noindent\textbf{Step 1.} Write a quaternion as $(q_0,q_1,q_2,q_3)$ and put
$\varphi=(1+\sqrt5)/2$. Let $G$ be the following set of 120 unit
quaternions: the eight signed coordinate units; the sixteen points
$(\pm1,\pm1,\pm1,\pm1)/2$; and the 96 points obtained by even
coordinate permutations of
$(0,\pm1,\pm\varphi,\pm\varphi^{-1})/2$, with independent signs.
Quaternion multiplication makes this the binary icosahedral group.
For an entirely finite verification, the attached program enumerates
these points in $\Q(\sqrt5)^4$, verifies unit norm and all 14,400
multiplication-table entries, and computes its commutator subgroup.
Every element is a commutator, in particular $G^{\rm ab}=0$.
This is an exact computation over rational pairs representing
$\Q(\sqrt5)$.

\noindent\textbf{Step 2.} Use the faithful holomorphic representation
\[
(a,b,c,d)\longmapsto
\begin{pmatrix}a+ib&c+id\\-c+id&a-ib\end{pmatrix}\in SU(2)
\]
and let the generator of $\Z$ act on $\C^2\setminus\{0\}$ by
$v\mapsto2v$. The action of $\Z\times G$ is free. An equality
$2^ngv=v$ first gives $n=0$ by norms, and a unit quaternion whose
matrix has eigenvalue one is the identity. Proper discontinuity
follows from the radial scaling and finiteness of $G$.
Polar coordinates identify the quotient smoothly with
\[
Q_H=(\C^2\setminus\{0\})/(2^{\Z}\times G)
\simeq (\mathbb R/(\log2)\Z)\times(S^3/G).
\]
It is compact, and its universal cover is $\C^2\setminus\{0\}$.
The fundamental group is $\Z\times G$, whose abelianization is
$\Z$. The quotient by $2^{\Z}$ alone is a primary Hopf surface
covering $Q_H$ with degree 120. The fundamental group of $Q_H$ is
noncyclic, so $Q_H$ is secondary Hopf. The product description gives
$e(Q_H)=0$, and the same duality and universal-coefficient argument
as before gives $H^2(Q_H;\Z)=0$.
\end{proof}



These two explicit finite-cover maps locate two additional failures
in the surface-level exclusions used in the proof of CDP20,
Lemma 4.3. They do not construct these surfaces as fibres of a
threefold satisfying the standing hypotheses. Such a realization or
a separate exclusion remains part of C.3.




\begin{proposition}\label{cdp:st:C.3.4}
\StatementRecord{cdp:st:C.3.4}{proved}{Historical events: Claude Fable 5.2 (Anthropic), effort high, 16 November 2026: carried as background without fault found. GPT-6 Astra (OpenAI), ChatGPT Pro mode, effort Pro, S6\_09\_PAPERS, 17 November 2026: targeted proof check. Also checked in this writing. Pre-continuation writing events: GPT-6 Astra (OpenAI), ChatGPT Pro mode, effort Pro, S6\_10\_PAPERS, 18 November 2026: targeted proof check. Continuation metadata events: GPT-6 Astra (OpenAI), ChatGPT Pro mode, effort Pro, S6\_10\_PAPERS continuation, 19 November 2026: used and rechecked in this run; status-vocabulary reconciliation; metadata reconciliation only; no new mathematical check. \emph{Pre-continuation recorded proof scope.} An adjacent effective divisor with rationally trivial first Chern class has positive integral against a Kahler form on a proposed resolution, giving the contradiction}{the hypotheses, notation and definitions stated here; the explicit constructions and calculations in the proof}

Let $Y$ be a compact complex threefold with $b_2(Y)=0$, and let
$f:Y\to C$ be proper, holomorphic, surjective, with connected fibres.
Every irreducible component of a reducible reduced fibre has no
compact K\"ahler resolution. In particular, a fibre all of whose
components admit compact K\"ahler resolutions is irreducible.
\end{proposition}


\begin{proof}
\noindent\textbf{Step 1.} Let $D$ be a component of a reducible fibre. The incidence graph of
the components is connected, so some other component $D'$ meets $D$.
Both are Cartier divisors on the smooth threefold. Let
$\mu:\widehat D\to D\hookrightarrow Y$ be a resolution by a compact
smooth surface. Pullback of the canonical section of $\OO_Y(D')$
is a nonzero holomorphic section of $\mu^*\OO_Y(D')$. Its zero
divisor $E$ is effective and nonzero: the intersection $D\cap D'$
is a nonempty divisor on the irreducible surface $D$, and pullback
to the resolution preserves its nonzero vanishing order along a
component dominating that intersection. On the other hand
$c_1(\mu^*\OO_Y(D'))=0$ in $H^2(\widehat D;\R)$.
A K\"ahler form $\omega$ on $\widehat D$ would give
\[
0=\int_{\widehat D}c_1(\OO_{\widehat D}(E))\wedge\omega
 =\int_E\omega>0,
\]
where positive multiplicities are retained in the last integral.
This contradiction proves the assertion for every component.
\end{proof}






\begin{proposition}[Proved exclusion of bielliptic fibres]\label{cdp:st:C.3.7}
\StatementRecord{cdp:st:C.3.7}{proved}{Historical events: Claude Fable 5.2 (Anthropic), effort high, 16 November 2026: carried as background without fault found. GPT-6 Astra (OpenAI), ChatGPT Pro mode, effort Pro, S6\_09\_PAPERS, 17 November 2026: targeted proof check. Also checked in this writing. Pre-continuation writing events: GPT-6 Astra (OpenAI), ChatGPT Pro mode, effort Pro, S6\_10\_PAPERS, 18 November 2026 04:33:41.407298+00:00: targeted proof check; used and rechecked in this run. Continuation metadata events: GPT-6 Astra (OpenAI), ChatGPT Pro mode, effort Pro, S6\_10\_PAPERS continuation, 19 November 2026: used and rechecked in this run; status-vocabulary reconciliation; metadata reconciliation only; no new mathematical check. \emph{Pre-continuation recorded proof scope.} proved. Full relative-Albanese contradiction re-derived; exact primary construction passage still being opened. The finite orbit sum preserves the Albanese image and supplies an equivariant line/divisor construction. Resolving only exceptional base values leaves the general fibre unchanged. Connected general Albanese fibres force degree-one Stein factor, hence H1(base surface) injects into the blown-up total space. Even-b1 surface Kahlerness gives omega\_B. Every H2 class upstairs is a sum of exceptional divisors supported over omitted values, so g*omega\_B restricts to0 on a general bielliptic fibre, contradicting its strictly positive pairing with a fibre Kahler form. Peters Theorem1.1 p6 is currently opened; Kawai Proposition2 full local construction passage remains a separate opening gate. Review provenance: New derivation after reading the supplied proof; not a blind review and not an independent human check}{the hypotheses, notation and definitions stated here; \cite[Proposition~2, pp.~605--608]{Kawai1969}; \cite[Theorem~1.1]{Peters1990}}


A compact complex threefold $Y$ with $b_1(Y)=b_2(Y)=0$ has no
proper connected-fibre holomorphic map to a compact curve whose
general fibre is bielliptic.
\end{proposition}


\begin{proof}
\noindent\textbf{Step 1.} We use precisely the relative Albanese construction of
\cite[Proposition~2, pp.~605--608]{Kawai1969}, and not its subsequent
corollary. Applied to bielliptic fibres, which have
$p_g=0$ and $h^0(\Omega^1)=1$, it gives a compact complex surface
$B$, a holomorphic map $h:B\to C$, and a meromorphic map
$\Phi:Y\dashrightarrow B$ over $C$. On the complement $C^\circ$
of finitely many values, $B^\circ\to C^\circ$ is a smooth family
of elliptic curves and
\[
 \Phi|_{F_c}:F_c\longrightarrow B_c
\]
is the Albanese morphism of the bielliptic surface, with connected
elliptic fibres. To specify the comparison being imported, the
construction forms the relative period-lattice quotient, glues its
local Albanese maps by fibre translations, and compactifies at each
missing value by a normalized finite base change and the projective
space associated with the direct images of a divisor which is the
inverse image of a local Albanese section. In making this last
step equivariant, use the sum of the finite orbit of that divisor;
a selected section need not itself be invariant under the covering
group. Resolve the normalized base change first if needed, so the
divisor defines a line bundle. The orbit sum restricts on every
general fibre to the pullback of a positive-degree divisor on its
Albanese elliptic curve; sufficiently large powers give the same
Albanese image. Its naturally linearized line bundle and direct
images carry the finite covering-group action, and their quotient
glues the construction to $B^\circ$. These are the maps asserted by
Proposition~2, with that equivariant choice made explicit. The
later assertion in Kawai's corollary about multiple fibres is not
used here.

\noindent\textbf{Step 2.} Resolve $B$ over $C\setminus C^\circ$ and resolve the induced
meromorphic map by a sequence of blowups with smooth centres
supported over those same values. This produces a diagram
\[
 \widetilde Y\xrightarrow{\ g\ }\widetilde B,
 \qquad \pi:\widetilde Y\longrightarrow Y,
 \qquad (h\circ\rho)\circ g=f\circ\pi,
\]
where $\rho:\widetilde B\to B$ is a resolution and both spaces
are smooth and compact. The map $g$ has connected fibres:
its general fibres are the connected Albanese fibres just described,
and its Stein factorization has finite degree one over the normal
surface $\widetilde B$. Hence
$H^1(\widetilde B;\Q)\hookrightarrow H^1(\widetilde Y;\Q)$.
Blowing up smooth centres does not change first cohomology, so
$b_1(\widetilde B)=0$. The compact-surface K\"ahler criterion
\cite[Theorem~1.1]{Peters1990} therefore gives a K\"ahler form
$\omega_B$ on $\widetilde B$.

\noindent\textbf{Step 3.} The blowup cohomology formula gives
\[
 H^2(\widetilde Y;\R)
 =\pi^*H^2(Y;\R)\oplus\bigoplus_j\R[E_j]
 =\bigoplus_j\R[E_j],
\]
where the $E_j$ are the exceptional prime-divisor classes (with
the corresponding total transforms for successive blowups).
Every $E_j$ is supported above one of the removed values in $C$.
Consequently its restriction to a general smooth fibre $F_c$,
identified with its unchanged inverse image in $\widetilde Y$,
is zero. In particular
\[
 [g^*\omega_B]|_{F_c}=0\quad\hbox{in }H^2(F_c;\R).
\]
Choose a K\"ahler form $\omega_F$ on the bielliptic surface $F_c$.
The Albanese morphism $g|_{F_c}$ is a submersion to the elliptic
curve $\widetilde B_c$. Thus its pullback of $\omega_B$ is a
nonzero semipositive $(1,1)$-form and
\[
 \int_{F_c}g^*\omega_B\wedge\omega_F>0.
\]
The vanishing cohomology class makes the same integral zero by
Stokes' theorem, a contradiction. All maps in this argument are
the displayed relative Albanese map and its modifications; no
comparison between the two elliptic fibrations of a bielliptic
surface is needed.
\end{proof}






\begin{proposition}[Proved exclusion of elliptic Hopf general fibres]\label{cdp:st:C.3.8}
\StatementRecord{cdp:st:C.3.8}{proved}{Historical events: Claude Fable 5.2 (Anthropic), effort high, 16 November 2026: carried as background without fault found. GPT-6 Astra (OpenAI), ChatGPT Pro mode, effort Pro, S6\_09\_PAPERS, 17 November 2026: targeted proof check. Also checked in this writing. Pre-continuation writing events: GPT-6 Astra (OpenAI), ChatGPT Pro mode, effort Pro, S6\_10\_PAPERS, 18 November 2026 04:33:41.407298+00:00: targeted proof check; used and rechecked in this run. Continuation metadata events: GPT-6 Astra (OpenAI), ChatGPT Pro mode, effort Pro, S6\_10\_PAPERS continuation, 19 November 2026: used and rechecked in this run; status-vocabulary reconciliation; metadata reconciliation only; no new mathematical check. \emph{Pre-continuation recorded proof scope.} proved. Full anticanonical and countable-locus exclusion argument re-derived. Checked invariance of both explicit bivectors under the actual contraction derivative. On an elliptic Hopf surface all curves are vertical by zero rational H2, so a common fibre-multiplicity multiple of the effective anticanonical divisor is a positive base pullback. The finite-cover invariant section ring is integral by the actual orbit polynomial, preserving Iitaka dimension. On the smooth base the loci h0(K\textasciicircum{}−m)\ensuremath{\geq}2 are analytic; countability forces one whole-base locus when a very general fibre is elliptic Hopf. Relative linear systems then give a projective surface image and algebraic dimension\ensuremath{\geq}2. Review provenance: New derivation after reading the supplied proof; not a blind review and not an independent human check}{the hypotheses, notation and definitions stated here; \cref{cdp:st:C.3.1}; \cite[Proposition~3.1 and its proof]{SS2002}}


The very general Hopf fibre of the algebraic reduction in \ref{cdp:st:C.3.1}
has algebraic dimension zero. In particular neither primary nor
secondary Hopf surfaces of algebraic dimension one can be its
general fibre. This exclusion uses no restriction-surjectivity
assertion from CDP98 or CDP20.
\end{proposition}


\begin{proof}
\noindent\textbf{Step 1.} We first verify that a Hopf surface $H$ with $a(H)=1$ satisfies
$\kappa(H,K_H^{-1})=1$. On a primary Hopf cover the contraction
normal forms are diagonal, or
\[
A(z,w)=(\alpha z+\lambda w^n,\beta w),\qquad
\lambda\ne0,\quad \alpha=\beta^n.
\]
The invariant bivectors $zw\,\partial_z\wedge\partial_w$ in the
diagonal case and $w^{n+1}\,\partial_z\wedge\partial_w$ in the
second case give a nonzero effective anticanonical divisor.
The invariance is an equality under the derivative of the
displayed contraction: its determinant is $\alpha\beta$, and
the displayed coefficient transforms by the same scalar.
These normal forms and the elliptic-fibration description are
also used in \cite[Proposition~3.1 and its proof]{SS2002}.

\noindent\textbf{Step 2.} The algebraic reduction of an elliptic Hopf surface is a
holomorphic map to $\PP^1$, with smooth reductions of all fibres.
Every curve is vertical: its intersection with a general fibre
is zero since $H^2(H;\R)=0$, whereas a horizontal curve would
have positive intersection equal to its degree over the base.
Thus the nonzero anticanonical divisor is a positive sum of
reduced fibres. If $m_iD_i$ is the full fibre supported on
$D_i$, each $\OO(m_iD_i)$ is the pullback of $\OO_{\PP^1}(1)$.
A common multiple of the $m_i$ therefore makes a positive power
of $K_H^{-1}$ the pullback of a positive-degree line bundle on
$\PP^1$. Its Iitaka dimension is one.

\noindent\textbf{Step 3.} For a secondary Hopf surface take a normal finite primary cover
$p:H'\to H$. The cover is unramified, so
$K_{H'}^{-1}=p^*K_H^{-1}$, and $a(H')=a(H)=1$ by invariance of
transcendence degree under a finite map. The anticanonical
section ring of $H$ is the invariant subring of that of $H'$.
The latter is integral over the invariant subring: a homogeneous
section $s$ satisfies the monic orbit polynomial
$\prod_{g\in\operatorname{Gal}(H'/H)}(T-gs)$. Their fraction
fields have the same transcendence degree, proving
$\kappa(H,K_H^{-1})=1$ also in this case.

\noindent\textbf{Step 4.} Now, for each positive integer $m$, the locus
\[
Z_m=\{c\in C^\circ:
 h^0(Y_c,K_Y^{-m}|_{Y_c})\ge2\}
\]
is a closed analytic subset of the smooth base $C^\circ$.
If the very general Hopf fibre had algebraic dimension one,
the preceding calculation would make the union of these
countably many loci contain the complement of a countable
subset of $C^\circ$. A proper closed analytic subset of a
curve is discrete and countable, so some $Z_m$ is the whole
connected smooth base. Coherent base change then gives
$\operatorname{rank}f_*K_Y^{-m}\ge2$.
The relative linear system defines a meromorphic map of $Y$
to $\PP(f_*K_Y^{-m})$ over $C$ with positive-dimensional image
on a general fibre. Its image has dimension at least two.
The projective bundle over the compact curve is projective,
so pulling back its meromorphic functions gives $a(Y)\ge2$,
contrary to the hypothesis. In fact the same argument shows
that, when the smooth fibres are Hopf surfaces, the values with
$a(Y_c)=1$ form at most a countable set.
\end{proof}






\begin{proposition}[Proved exclusion of secondary Kodaira general fibres]\label{cdp:st:C.3.9}
\StatementRecord{cdp:st:C.3.9}{proved}{Historical events: Claude Fable 5.2 (Anthropic), effort high, 16 November 2026: carried as background without fault found. GPT-6 Astra (OpenAI), ChatGPT Pro mode, effort Pro, S6\_09\_PAPERS, 17 November 2026: targeted proof check. Also checked in this writing. Pre-continuation writing events: GPT-6 Astra (OpenAI), ChatGPT Pro mode, effort Pro, S6\_10\_PAPERS, 18 November 2026: used and rechecked in this run. GPT-6 Astra (OpenAI), ChatGPT Pro mode, effort Pro, S6\_10\_PAPERS, 18 November 2026 04:33:41.407298+00:00: targeted proof check; used and rechecked in this run. Continuation metadata events: GPT-6 Astra (OpenAI), ChatGPT Pro mode, effort Pro, S6\_10\_PAPERS continuation, 19 November 2026: used and rechecked in this run; status-vocabulary reconciliation; metadata reconciliation only; no new mathematical check. \emph{Pre-continuation recorded proof scope.} proved. Complete secondary-Kodaira Picard-quotient and finite-cover norm argument re-derived. The integral transgression makes Pic0(S) a C* source with integral period. The canonical cover invariants give a line transverse to the pulled-back elliptic Picard tangent, hence a nonzero integer-power map to Q=Pic0(P)/Pic0(E). The oriented degree-d class alpha is the bounded punctured-disc moduli coordinate and must be constant over the punctured projective line. Matching one fixed ambient exponential with its degree class leaves a finite-order invariant elliptic Picard point; order2,3,4,6 implies the12th power is trivial. The degree-m norm gives the exact M\textasciicircum{}(12m)=rho*O(12mN) identity, forcing a rank\ensuremath{\geq}2 relative linear system. The only remaining gate is the exact current primary moduli/cover locus, not the displayed norm computation. Review provenance: New derivation after reading the supplied proof; not a blind review and not an independent human check}{the hypotheses, notation and definitions stated here; \cref{cdp:st:D.1.9}; \cref{cdp:st:E4}; \cref{cdp:st:D.1.7}; \cref{cdp:eq:C.SK.1}; \cite[Theorem~4 and Remark~1]{Borcea1984}; \cite[\S1, pp.~2--3]{SS2002}; \cref{cdp:eq:C.SK.2}; \cref{cdp:eq:C.SK.3}; \cref{cdp:eq:C.SK.4}; \cref{cdp:eq:C.SK.5}}


A compact complex threefold with $b_1(Y)=b_2(Y)=0$ and
$a(Y)=1$ cannot have a holomorphic algebraic reduction
$f:Y\to\PP^1$ whose general fibre is a secondary Kodaira
surface. Rational Betti-number vanishing suffices; no integral
torsion-freeness of the fibre is assumed.
\end{proposition}


\begin{proof}
\noindent\textbf{Step 1.} Suppose that the smooth general fibres $S_c$ are secondary
Kodaira surfaces. They have $h^1(\OO_{S_c})=b_1(S_c)=1$.
The nonzero integral transgression section of \ref{cdp:st:D.1.9} has
coefficient image $\beta\in H^1(Y,\OO_Y)$ and gives a
nowhere-zero section of $\mathscr E=R^1f_*\OO_Y$.
This coherent sheaf is a line bundle by \ref{cdp:st:E4}. Thus
$\mathscr E=\OO_{\PP^1}$ and
$H^1(Y,\OO_Y)=\C\beta$. For every $z\in\C$ set
\[
 M_z=\exp(2\pi\mathrm i z\beta)\in\Pic^0(Y).
\tag{C.SK.1}\label{cdp:eq:C.SK.1}
\]
Its restriction satisfies $M_{z+1}|_{S_c}\simeq M_z|_{S_c}$,
because $\beta|_{S_c}$ is the coefficient image of an
integral class. It is nonzero as a tangent vector of
$\Pic^0(S_c)$ by \ref{cdp:st:D.1.7}. The exponential sequence therefore
identifies $\Pic^0(S_c)$ with $\C^*$ and makes \eqref{cdp:eq:C.SK.1}, on
each smooth fibre, a nonconstant homomorphism from
$\C/\Z$ to this group.

\noindent\textbf{Step 2.} We specify the primary cover and the invariant used below.
Let $m>1$ be the order of $K_{S_c}$, and let
\[
 p_c:P_c\longrightarrow S_c
\]
be its connected cyclic canonical cover, with deck group
$G=\mu_m$. It is finite \emph{\'etale}. Its canonical bundle
is trivial, and it is non-K\"ahler, because a finite
\emph{\'etale} quotient of a K\"ahler surface is K\"ahler
by averaging a K\"ahler form. The classification of
surfaces with trivial canonical bundle consequently makes
$P_c$ a primary Kodaira surface. Write
$\pi_c:P_c\to E_c$ for its Albanese elliptic bundle, and
$d>0$ for its topological degree. These covers exist in
holomorphic families locally on the smooth base: the
evaluation of a local frame of $f_*K_{Y/\PP^1}^{m}$
trivializes $K_{Y/\PP^1}^{m}$, since it is nowhere zero on
each fibre, and taking its cyclic root gives the required
finite \emph{\'etale} covering family. The integers $m,d$
are locally constant and hence constant on the connected
smooth base.

\noindent\textbf{Step 3.} The exact primary-surface invariant we use is
\[
 Q_c=\Pic^0(P_c)/\pi_c^*\Pic^0(E_c)\simeq\C^*.
\tag{C.SK.2}\label{cdp:eq:C.SK.2}
\]
The Picard pullback in the denominator is injective:
$\pi_{c*}\OO_{P_c}=\OO_{E_c}$ and the projection formula
show that a line with trivial pullback is trivial. For
any line $L_c$ of degree $d$ on $E_c$, the pullback
$\pi_c^*L_c$ belongs to $\Pic^0(P_c)$, and its class in
$Q_c$ is independent of $L_c$. Choose the group
isomorphism $q_c:Q_c\to\C^*$ that places this
positive-degree class in the punctured unit disc, and
denote its value by $\alpha(c)$.
This fixes the possible inversion ambiguity of the
group coordinate. The descriptions in
\cite[Theorem~4 and Remark~1]{Borcea1984} and
\cite[\S1, pp.~2--3]{SS2002} identify $\alpha$ as the
intrinsic holomorphic punctured-disc coordinate of the
moduli of primary Kodaira surfaces. More explicitly,
$P_c$ is the associated elliptic bundle of a line of
degree $d$ with fibre $\C^*/\langle\alpha(c)\rangle$;
quotienting by $\pi_c^*\Pic^0(E_c)$ removes the choice
of that line within its degree. In a holomorphic
family the local integral marking and the coefficient
map give this quotient a holomorphic $\C^*$ coordinate,
and the chosen positive-degree orientation is locally
constant. This is also immediate in Borcea's complete
period family. Thus $\alpha(c)$ is holomorphic.

\noindent\textbf{Step 4.} The locally constructed canonical covers have the same
isomorphism class on overlaps, so this intrinsic
function glues on the smooth base. It has
$0<|\alpha(c)|<1$. It extends across the finitely many
omitted values by the removable singularity theorem
and is constant on $\PP^1$. Write its nonzero constant
value as $\alpha$.

\noindent\textbf{Step 5.} Consider the exact homomorphism of connected Picard
groups
\[
 \Pic^0(S_c)\xrightarrow{\ p_c^*\ }\Pic^0(P_c)
          \longrightarrow Q_c\xrightarrow{\ q_c\ }\C^*.
\tag{C.SK.3}\label{cdp:eq:C.SK.3}
\]
Its differential is nonzero. Indeed finite-cover trace
identifies $H^1(S_c,\OO_{S_c})$ with the $G$-invariant
line in $H^1(P_c,\OO_{P_c})$. The other line,
$\pi_c^*H^1(E_c,\OO_{E_c})$, has nontrivial $G$ action:
otherwise the unique holomorphic one-form of $P_c$
would descend to $S_c$, whereas a secondary Kodaira
surface has no holomorphic one-form. Hence the invariant
line has zero intersection with the denominator's
tangent line in \eqref{cdp:eq:C.SK.2}. This proves the claim.
Consequently \eqref{cdp:eq:C.SK.3} is a nonzero isogeny of $\C^*$
groups. Composing it with \eqref{cdp:eq:C.SK.1} gives
\[
 q_c\bigl([p_c^*(M_z|_{S_c})]\bigr)
       =\exp(2\pi\mathrm i k z)
       \qquad(k\in\Z\setminus\{0\}).
\tag{C.SK.4}\label{cdp:eq:C.SK.4}
\]
Here integrality follows from the period $z\mapsto z+1$;
every holomorphic group homomorphism $\C^*\to\C^*$ is
an integral power. The integer $k$ is locally constant
in the holomorphic families just described, and is
therefore a single integer on the connected smooth base.

\noindent\textbf{Step 6.} Let $\rho_c:S_c\to\PP^1$ be its elliptic algebraic
reduction. The induced finite map $r_c:E_c\to\PP^1$
satisfies $\rho_c p_c=r_c\pi_c$. Write $h=\deg r_c$;
this positive integer is constant. Choose one positive
integer $N$ killing the torsion group $H^2(S_c;\Z)$
and satisfying $d\mid hN$. This choice is uniform
because the smooth fibres are differentiably locally
trivial. Then
$A_c=\rho_c^*\OO_{\PP^1}(N)$ belongs to $\Pic^0(S_c)$,
and \eqref{cdp:eq:C.SK.2} gives
\[
 q_c([p_c^*A_c])=\alpha^{hN/d}.
\tag{C.SK.5}\label{cdp:eq:C.SK.5}
\]
This identity is fibrewise and uses only degrees; it
does not require the divisors $\rho_c^{-1}(x)$ to
extend across a singular fibre of $f$.

\noindent\textbf{Step 7.} Choose $z_0\in\C$ with
$\exp(2\pi\mathrm i k z_0)=\alpha^{hN/d}$ and put
$M=M_{z_0}$. By \eqref{cdp:eq:C.SK.4}--\eqref{cdp:eq:C.SK.5}, the difference of
$p_c^*(M|_{S_c})$ and $p_c^*A_c$ is
$\pi_c^*R_c$ for some $R_c\in\Pic^0(E_c)$.
Both of the first two lines come from $S_c$ and are
$G$-invariant. Injectivity of $\pi_c^*$ implies that
$R_c$ is fixed by the induced nontrivial linear action
on $\Pic^0(E_c)$. An elliptic-curve linear automorphism
of finite order other than one has order
$r\in\{2,3,4,6\}$. If it fixes a point $R_c$, then
\[
 rR_c=(1+\zeta+\cdots+\zeta^{r-1})R_c=0
\]
in the elliptic group. Thus $R_c^{\otimes12}$ is
trivial. The norm of the finite cover satisfies
$\operatorname{Nm}_{p_c}(p_c^*L)=L^{\otimes m}$.
Applying it to the difference now proves the exact
line-bundle identity
\[
 M^{\otimes12m}|_{S_c}
       \simeq \rho_c^*\OO_{\PP^1}(12mN)
\tag{C.SK.6}\label{cdp:eq:C.SK.6}
\]
for every smooth general fibre.

\noindent\textbf{Step 8.} The right side has at least two linearly independent
sections, by pullback from $\PP^1$. Coherent generic
base change gives
$\operatorname{rank}f_*M^{\otimes12m}\ge2$.
Its relative linear system has a positive-dimensional
image on a general fibre and therefore a surface
image in the projective bundle over $\PP^1$.
Pullback of meromorphic functions from this image
gives $a(Y)\ge2$, the desired contradiction.
This proves the exclusion without the arbitrary-twist
cohomological-flatness assertion used in the printed
Kodaira exclusion.
\end{proof}






\begin{proposition}[Proved exclusion of primary Kodaira general fibres]\label{cdp:st:C.3.10}
\StatementRecord{cdp:st:C.3.10}{proved}{Historical events: Claude Fable 5.2 (Anthropic), effort high, 16 November 2026: carried as background without fault found. GPT-6 Astra (OpenAI), ChatGPT Pro mode, effort Pro, S6\_09\_PAPERS, 17 November 2026: targeted proof check. Also checked in this writing. Pre-continuation writing events: GPT-6 Astra (OpenAI), ChatGPT Pro mode, effort Pro, S6\_10\_PAPERS, 18 November 2026 04:33:41.407298+00:00: used and rechecked in this run. GPT-6 Astra (OpenAI), ChatGPT Pro mode, effort Pro, S6\_10\_PAPERS, 18 November 2026 04:34:17.629108+00:00: targeted proof check; used and rechecked in this run. Continuation metadata events: GPT-6 Astra (OpenAI), ChatGPT Pro mode, effort Pro, S6\_10\_PAPERS continuation, 19 November 2026: used and rechecked in this run; status-vocabulary reconciliation; metadata reconciliation only; no new mathematical check. \emph{Pre-continuation recorded proof scope.} proved. Complete primary-Kodaira proof re-derived, including both integral transgression cases and all central shears. For k\ensuremath{\neq}0 the quotient exponential matches the constant degree class directly. For k=0, real coefficient injectivity and saturation put gamma in base H1; its primitive period q=exp(2pi i tau) is bounded and constant, forcing identity base monodromy. Constant (j,alpha,d) and the countable integral orbit in the complete marked family give local products. The group action fixes both central generators and adds only central shears to base generators; each invariant H2 basis class changes by an exact a wedge b term. The logarithmic de Rham degree-zero projection, the canonical identity-monodromy extension, and blowup O-pushforward produce a global nonzero section of G. The full-fibre wedge isomorphism lifts it to v, and the exact two linear maps on the quotient have one fixed nonzero scalar, yielding the ambient positive-degree line. Input D.2.3 was re-read for its analytic comparison role; no separate new proof event is assigned to it. Review provenance: New derivation after reading the supplied proof; not a blind review and not an independent human check}{the hypotheses, notation and definitions stated here; \cref{cdp:st:C.3.9}; \cite[pp.~373--374]{Borcea1984}; \cref{cdp:st:D.1.12}; \cref{cdp:st:D.1.9}; \cref{cdp:st:D.1.7}; \cite[\S1]{SS2002}; \cite[Remark~1]{Borcea1984}; \cref{cdp:eq:C.PK.4}; \cite[\S3]{Borcea1984}; \cref{cdp:st:D.2.3}; \cite[Theorem~2.18 and \S\S2.19--2.21]{Steenbrink1976}; \cref{cdp:eq:C.PK.5}; \cref{cdp:eq:C.PK.1}; \cref{cdp:eq:C.PK.2}; \cref{cdp:eq:C.PK.7}}


A compact complex threefold with $b_1(Y)=b_2(Y)=0$ and
$a(Y)=1$ cannot have a holomorphic algebraic reduction
$f:Y\to\PP^1$ whose general fibre is a primary Kodaira
surface. Together with \ref{cdp:st:C.3.9} this excludes both Kodaira
types using only rational Betti-number vanishing.
\end{proposition}


\begin{proof}
\noindent\textbf{Step 1.} Let $B=\PP^1\setminus\Sigma$ be the smooth base and write
$P_c=f^{-1}(c)$ and $\pi_c:P_c\to E_c$ for the canonical
elliptic algebraic reduction of a smooth fibre. The degree
$d>0$ is constant. All the smooth fibres have the fixed
primary Kodaira differentiable type; every complex structure
on that type is again a primary Kodaira structure
\cite[pp.~373--374]{Borcea1984}.
Use the intrinsic quotient and its orientation from (C.SK.2):
\[
 Q_c=\Pic^0(P_c)/\pi_c^*\Pic^0(E_c),\qquad
 q_c:Q_c\xrightarrow{\simeq}\C^*.
\]
The class of the pullback of a degree-$d$ line has coordinate
$\alpha(c)\in\Delta^*$. The holomorphic moduli comparison
in the proof of \ref{cdp:st:C.3.9} applies directly here, so this bounded
function extends to $\PP^1$ and has a single nonzero
constant value $\alpha$.

\noindent\textbf{Step 2.} Put $\mathscr E=R^1f_*\OO_Y$ and
$\mathscr G=R^2f_*\OO_Y$. The primary Kodaira case of
\ref{cdp:st:D.1.12} proves, on the original threefold and every full fibre,
\[
 \operatorname{rank}\mathscr E=2,\qquad
 \operatorname{rank}\mathscr G=1,\qquad
 \bigwedge^2\mathscr E\xrightarrow{\simeq}\mathscr G.
\tag{C.PK.1}\label{cdp:eq:C.PK.1}
\]
Both sheaves are locally free and satisfy base change.
Let $\gamma$ be the integral transgression section of \ref{cdp:st:D.1.9},
and let $\beta\in H^0(\PP^1,\mathscr E)=H^1(Y,\OO_Y)$
be its coefficient image. This section is nowhere zero.
For $u\in H^1(P_c,\OO_{P_c})$ define the complex-linear
functional $\lambda_c$ by the exact identity
\[
 q_c\bigl([\exp(2\pi\mathrm i z u)]\bigr)
     =\exp(2\pi\mathrm i z\lambda_c(u))
     \quad(z\in\C).
\tag{C.PK.2}\label{cdp:eq:C.PK.2}
\]
Its kernel is the line
$A_c=\pi_c^*H^1(E_c,\OO_{E_c})$.
The function $k(c)=\lambda_c(\beta_c)$ takes integral
values, because $\beta_c$ is the coefficient image of
an integral class. It is holomorphic, by the actual relative
Picard quotient in (C.SK.2), hence is a constant integer $k$.

\noindent\textbf{Step 3.} When $k\ne0$, choose $n>0$ with $nd\ge2$ and choose
$z\in\C$ with $\exp(2\pi\mathrm i kz)=\alpha^n$.
The ambient line $M=\exp(2\pi\mathrm i z\beta)$ then
has, on every smooth fibre, the same $Q_c$ class as the
pullback of a degree-$nd$ line on $E_c$. Its difference
from that pullback belongs to $\pi_c^*\Pic^0(E_c)$.
Thus
$M|_{P_c}=\pi_c^*L_c$ for a line $L_c$ of degree $nd$,
and $h^0(P_c,M|_{P_c})=nd\ge2$. The relative linear
system gives $a(Y)\ge2$, as at the end of \ref{cdp:st:C.3.9}.
It remains to obtain the same contradiction when $k=0$.

\noindent\textbf{Step 4.} In that case $\beta_c\in A_c$. The coefficient map
$H^1(E_c;\R)\to H^1(E_c,\OO_{E_c})$ is an isomorphism
of real vector spaces. Its pullback, together with the real
injectivity of \ref{cdp:st:D.1.7} on $P_c$, proves the exact implication
\[
 \beta_c\in A_c
 \quad\Longrightarrow\quad
 \gamma_c\in\pi_c^*H^1(E_c;\R).
\tag{C.PK.3}\label{cdp:eq:C.PK.3}
\]
The lattice $\pi_c^*H^1(E_c;\Z)$ is saturated in
$H^1(P_c;\Z)$: the elliptic bundle induces a surjection
on fundamental groups, so the cokernel of the induced
injection on integral $H^1$ is torsion free. Thus $\gamma$
lies in this rank-two integral local system. Divide it by
its fixed positive divisibility to obtain a primitive fixed
section $\gamma_0$. On a simply connected open set choose
a complementary integral section $\delta$ with oriented
intersection number one. The ratio of their nonzero
coefficient images is a holomorphic elliptic period
\[
 \tau=\frac{\iota(\delta)}{\iota(\gamma_0)},
 \qquad \operatorname{Im}\tau>0,
\]
after choosing the sign convention of the ordered pair.
Monodromy fixes $\gamma_0$ and replaces $\delta$ by
$\delta+N\gamma_0$, $N\in\Z$. Therefore
$\exp(2\pi\mathrm i\tau)$ is a single-valued bounded
holomorphic function on $B$. It extends to $\PP^1$ and
is constant. Its logarithmic derivative gives $d\tau=0$,
and a nonzero integral translation cannot fix this constant
period. Consequently the Albanese integral $H^1$ local
system has identity monodromy, and $j(E_c)$ is constant.

\noindent\textbf{Step 5.} Both moduli $(j,\alpha)$ and the degree $d$ are now
constant. Borcea's classification, in the precise form
\cite[\S1]{SS2002}, makes all smooth fibres biholomorphic
to one primary Kodaira surface $P$. Here is the actual
local product comparison we use. Borcea's complete
holomorphic marked period family \cite[Remark~1]{Borcea1984}
expresses our family, near any smooth base point, as the
pullback by a holomorphic map into his period domain
$\mathcal D\subset\PP(H^2(P;\C))$. By his Theorem~3,
constant isomorphism class means that the image of this
map lies in one orbit of the integral cohomology action.
That orbit is countable, since the action factors through
$\operatorname{GL}(H^2(P;\Z)/\operatorname{torsion})$.
A connected subset of a countable subset of a complex
manifold is one point: a real coordinate separating two
points would have a nontrivial interval as its connected
image. Therefore the period parameter is constant and
its pullback is a product. This proves the needed
holomorphic local triviality from the specified complete
family.

\noindent\textbf{Step 6.} We compute its entire action on $H^2(P;\C)$, including
vertical shears. Write the fundamental group as
\[
 \Gamma=\langle x,y,w,u\mid
       w,u\text{ central},\ [x,y]=w^d\rangle.
\tag{C.PK.4}\label{cdp:eq:C.PK.4}
\]
Its centre $\langle w,u\rangle$ is the fundamental group
of the elliptic fibre, and its quotient is the fundamental
group of the elliptic base. A biholomorphism inducing the
identity on the base $H^1$ induces the identity on this
quotient. The relation in \eqref{cdp:eq:C.PK.4} then forces it to fix
$w$. On an elliptic fibre its linear part is a holomorphic
group automorphism. Such an automorphism acts by a complex
scalar on the universal cover; fixing the nonzero period
$w$ forces that scalar to be one. It therefore fixes both
$w$ and $u$. The remaining possible actions on $x,y$ are
exactly multiplication by central elements.

\noindent\textbf{Step 7.} For clarity, their induced cohomology action can be read
off without discarding any of those elements. In the real
nilmanifold presentation choose invariant one-forms
$a,b,c,e$ with
\[
 da=db=de=0,\qquad dc=d\,a\wedge b.
\]
The four classes $[a\wedge c],[b\wedge c],
[a\wedge e],[b\wedge e]$ form a basis of
$H^2(P;\R)$, as in \cite[\S3]{Borcea1984}.
An automorphism of \eqref{cdp:eq:C.PK.4} with the stated action extends
to the real nilpotent group: in logarithmic coordinates it
fixes the two central generators and adds central vectors
to the two base generators. Its pullback fixes $a,b$ and
replaces each of $c,e$ by itself plus a real linear
combination of $a,b$. Each of the four displayed classes
therefore changes only by a multiple of the exact form
$a\wedge b$. Since the universal cover of $P$ is
contractible, a map of $P$ and this group-model map with
the same induced fundamental-group map are homotopic.
This computes the action of the original biholomorphism
and proves it is the identity on $H^2$.
In particular every monodromy of our isotrivial family,
whose action on the base $H^1$ was just proved to be the
identity, acts identically on $H^2(P;\C)$.

\noindent\textbf{Step 8.} We next construct a global section of $\mathscr G$ by an
actual morphism of complexes. Resolve all exceptional
fibres by successive blowups with smooth centres contained
in those fibres, writing
$\rho:\widetilde Y\to Y$ and $h=f\rho$. For each such
blowup, $\rho_*\OO_{\widetilde Y}=\OO_Y$ and
$R^j\rho_*\OO_{\widetilde Y}=0$ for $j>0$.
One local verification uses the formal neighbourhoods of
the exceptional projective bundle: their successive
conormal layers restrict to nonnegative twists
$\OO_{\PP^{r-1}}(l)$, $l\ge0$, which have no higher
cohomology. The analytic theorem on formal functions
gives the asserted vanishing; normality gives the
degree-zero equality. The Leray edge is consequently the
canonical isomorphism
\[
 R^2h_*\OO_{\widetilde Y}\xrightarrow{\simeq}
 R^2f_*\OO_Y=\mathscr G.
\tag{C.PK.5}\label{cdp:eq:C.PK.5}
\]
The same identities hold for a composite by derived
pushforward.

\noindent\textbf{Step 9.} Let $D=(h^{-1}(\Sigma))_{\mathrm{red}}$ and use the
relative logarithmic de Rham complex
\[
 \mathscr K^\bullet=
 \Omega_{\widetilde Y}^\bullet(\log D)\big/
 \bigl(h^*\Omega_{\PP^1}^1(\log\Sigma)
       \wedge\Omega_{\widetilde Y}^{\bullet-1}(\log D)\bigr).
\]
Projection onto degree zero is a morphism of complexes
$\mathscr K^\bullet\to\OO_{\widetilde Y}[0]$:
it is the identity in degree zero and zero in every
positive degree. The exact analytic comparison of \ref{cdp:st:D.2.3},
namely \cite[Theorem~2.18 and \S\S2.19--2.21]{Steenbrink1976},
applies to this resolution with all its fibre multiplicities
and without a K\"ahler hypothesis. It gives the logarithmic
de Rham bundle
$\mathscr H^2=\mathbb R^2h_*\mathscr K^\bullet$.
Identity monodromy makes every residue zero, by the
residue interval $[0,1)$ and the exponential comparison.
The global flat frame on $B$ therefore extends across
every point of $\Sigma$, proving
$\mathscr H^2\simeq H^2(P;\C)\otimes\OO_{\PP^1}$.
Taking degree-two direct hypercohomology of the displayed
projection and then \eqref{cdp:eq:C.PK.5} gives a holomorphic map
\[
 p:H^2(P;\C)\otimes\OO_{\PP^1}
          \longrightarrow\mathscr G.
\tag{C.PK.6}\label{cdp:eq:C.PK.6}
\]
Over $B$ it is precisely the coefficient projection
$H^2(P_c;\C)\to H^2(P_c,\OO_{P_c})$.

\noindent\textbf{Step 10.} In an isotrivial chart this latter map is constant. It is
surjective: the class of
$\overline\alpha\wedge\overline\eta$ used in the proof
of \ref{cdp:st:D.1.12} is represented by a closed invariant two-form
and has nonzero Dolbeault class. The charts also show
that its kernel is a flat subspace. Choose a constant flat
class $v_0\in H^2(P;\C)$ whose projection is nonzero,
and put $s=p(v_0)\in H^0(\PP^1,\mathscr G)$.
This section is nowhere zero on $B$; zeros at points
of $\Sigma$ cause no problem.

\noindent\textbf{Step 11.} By \eqref{cdp:eq:C.PK.1} and the nowhere-zero section $\beta$, the
following is an exact sequence of vector bundles on the
entire projective line:
\[
 0\longrightarrow\OO_{\PP^1}\beta
  \longrightarrow\mathscr E
  \xrightarrow{\ \beta\wedge-\ }\mathscr G
  \longrightarrow0.
\tag{C.PK.7}\label{cdp:eq:C.PK.7}
\]
Since $H^1(\PP^1,\OO)=0$, choose
$v\in H^0(\PP^1,\mathscr E)=H^1(Y,\OO_Y)$ with
$\beta\wedge v=s$.
We identify exactly the scalar that this lift induces in
the quotient \eqref{cdp:eq:C.PK.2}. In each isotrivial chart,
$\beta$ is the coefficient image of a fixed integral
class, $v_0$ is a fixed de Rham class, and $s=p(v_0)$
is a fixed nonzero class in $H^2(P,\OO_P)$.
The intrinsically oriented quotient $q_c$ is also fixed
by the chart: every automorphism preserves its
positive-degree class $\alpha\in\Delta^*$, so it cannot
induce inversion of $\C^*$.
Both linear maps $u\mapsto\beta\wedge u$ and
$u\mapsto\lambda(u)s$ have kernel $A$ and are nonzero.
Thus they differ by a constant $\kappa\in\C^*$ on
that chart. The constants agree on overlaps, because
all the maps and sections used in their definition are
the actual global ones. The connectedness of $B$ gives
one constant $\kappa$, and \eqref{cdp:eq:C.PK.7} yields
\[
 \lambda_c(v_c)=\kappa^{-1}\ne0
 \quad\hbox{for every }c\in B.
\tag{C.PK.8}\label{cdp:eq:C.PK.8}
\]
This proves constancy of the quotient differential,
including its normalization; it is not a boundedness
assertion about an unspecified function on $B$.

\noindent\textbf{Step 12.} Choose $n>0$ with $nd\ge2$ and $z\in\C$ with
$\exp(2\pi\mathrm i z\kappa^{-1})=\alpha^n$.
The single ambient line $M=\exp(2\pi\mathrm i z v)$
has this quotient coordinate on every smooth fibre.
Exactly as in the $k\ne0$ case,
$M|_{P_c}=\pi_c^*L_c$ with $\deg L_c=nd\ge2$.
Its direct image has rank at least two and its relative
linear system gives a surface image over $\PP^1$,
contradicting $a(Y)=1$. This completes the proof.
\end{proof}






\begin{proposition}[Proved exclusion of Inoue general fibres]\label{cdp:st:C.3.12}
\StatementRecord{cdp:st:C.3.12}{proved}{Historical events: Claude Fable 5.2 (Anthropic), effort high, 16 November 2026: carried as background without fault found. GPT-6 Astra (OpenAI), ChatGPT Pro mode, effort Pro, S6\_09\_PAPERS, 17 November 2026: targeted proof check. Also checked in this writing. Pre-continuation writing events: GPT-6 Astra (OpenAI), ChatGPT Pro mode, effort Pro, S6\_10\_PAPERS, 18 November 2026 04:35:54.089557+00:00: targeted proof check; used and rechecked in this run. Continuation metadata events: GPT-6 Astra (OpenAI), ChatGPT Pro mode, effort Pro, S6\_10\_PAPERS continuation, 19 November 2026: used and rechecked in this run; status-vocabulary reconciliation; metadata reconciliation only; no new mathematical check. \emph{Pre-continuation recorded proof scope.} proved. Complete Inoue closed-form and meridian-extension argument re-derived. Local completeness keeps first-coordinate real-affine transitions constant in the parameter. Liouville forces a lifted family isomorphism to be independent of the C coordinate in its H component; the unique common boundary fixed point forces that component to be aw+b, whose real holomorphic parameters are constant. Thus dlog Im(w) glues on the entire smooth total space, including finite quotients. For a transverse punctured disc, the developing map on a horodisc obeys W(tau+1)=aW+b. Schwarz–Pick bounds |log a| by1/(Im(tau)−c), so every reduced-divisor meridian has zero period. General-position two-chains identify the H1 kernel with those meridians and yield the actual global real cohomology class. The final all-smooth-fibre argument excludes virtually nilpotent Hopf/Kodaira groups by the expanding lattice eigenvalue, and excludes properly elliptic groups through their non-solvable hyperbolic orbifold quotient. Review provenance: New derivation after reading the supplied proof; not a blind review and not an independent human check}{the hypotheses, notation and definitions stated here; \cref{cdp:st:C.3.1}; \cite[pp.~274, 278, 280]{Inoue1974}; \cite[Sections 2.1--2.3]{KhaledTeleman2026}; \cref{cdp:eq:C.IN.2}; \cref{cdp:eq:C.IN.3}; \cite[Theorem 3.2]{BMS2012}; \cite[Section 3.3]{BMS2012}; \cref{cdp:eq:C.IN.6}}


Under the hypotheses of \ref{cdp:st:C.3.1}, the very general smooth fibre is
not an Inoue surface. More precisely, a proper holomorphic map from
a compact connected complex manifold to a compact Riemann surface,
whose fibres over the complement of a finite set are Inoue surfaces,
has a class in $H^1(Y;\mathbb R)$ restricting nontrivially to each
such fibre. Thus these families cannot have $b_1(Y)=0$.
This assertion permits arbitrary multiplicities and arbitrary
singularities of the omitted fibres.
\end{proposition}


\begin{proof}
\noindent\textbf{Step 1.} We first establish the precise family construction needed for the
assertion about $H^1$. Write $B^\circ$ for the complement of the
finite set and $Y^\circ=f^{-1}(B^\circ)$. Put
$\mathbb H=\{w\in\mathbb C:\operatorname{Im}w>0\}$.
Each Inoue surface is a quotient of $\mathbb H\times\mathbb C$.
In the three standard presentations the first coordinates of all
deck transformations have the form
\[
 w\longmapsto a w+b,\qquad a>0,\quad b\in\mathbb R.
\tag{C.IN.1}\label{cdp:eq:C.IN.1}
\]
They include a nonzero real translation and a transformation
$w\mapsto\alpha w$, with $\alpha>1$. Inoue's original
Propositions 2 and 5 establish infinitesimal rigidity for types
$S_M$ and $S^{-}$; his Proposition 4 gives a locally complete
one-parameter family for type $S^{+}$
\cite[pp.~274, 278, 280]{Inoue1974}.
In that family the parameter changes only the second coordinate:
local coordinate transitions have the form
\[
\begin{aligned}
 w_i&=\alpha^{m_{ij}}w_j+a_{ij},\\
 z_i&=z_j+b_{ij}w_j+c_{ij}+t\,m_{ij},
\end{aligned}
\tag{C.IN.2}\label{cdp:eq:C.IN.2}
\]
where $m_{ij}\in\mathbb Z$ and the other coefficients are
independent of $t$. In particular the real affine transformations
in the first coordinate are constant throughout this complete
family. The standard deck presentations are also recorded in
\cite[Sections 2.1--2.3]{KhaledTeleman2026}.

\noindent\textbf{Step 2.} After shrinking a simply connected base disc $V\subset B^\circ$,
local completeness identifies $f^{-1}(V)$ with the pullback of one
of these families. Its universal cover over $V$ is therefore
$\mathbb H\times\mathbb C\times V$, with holomorphic covering
map to $f^{-1}(V)$. In the $S^{+}$ case the parameter $t$ in
\eqref{cdp:eq:C.IN.2} is a holomorphic function on $V$; the other two cases
are products. The smooth real one-form
\[
 \theta_V=d\log\operatorname{Im}w
\tag{C.IN.3}\label{cdp:eq:C.IN.3}
\]
is invariant under every deck transformation, where $d$ is the
exterior derivative on the \emph{whole} covering space. It thus
descends to a closed real one-form on $f^{-1}(V)$.

\noindent\textbf{Step 3.} We prove that these forms agree on overlaps, including base
directions. Over a simply connected smaller overlap, a transition
between the two family descriptions lifts holomorphically to the
universal covers. On each fibre its first coordinate is independent
of $z$: composing with a Cayley transform gives a bounded entire
function of $z$, which is constant. Applying the same argument to
the inverse proves that the resulting first coordinate is an
automorphism $h$ of $\mathbb H$. This automorphism conjugates the
first-coordinate deck groups. Their common fixed-point set on
$\partial\mathbb H=\mathbb R\cup\{\infty\}$ is exactly
$\{\infty\}$: the nonzero translation has this as its sole fixed
point, and all their elements fix it. Consequently $h(\infty)=
\infty$ and $h(w)=a w+b$ with $a>0$ and $b\in\mathbb R$.
In a holomorphic family of these transitions, $a$ and $b$ are
holomorphic functions of the base parameter; this follows, for
example, by taking the $w$ derivative and then evaluating at a
fixed $w$. Their real values force them to be locally constant.
Therefore
\[
 d\log\operatorname{Im}(a w+b)
 =d\log(a\operatorname{Im}w)
 =d\log\operatorname{Im}w
\tag{C.IN.4}\label{cdp:eq:C.IN.4}
\]
as forms on the total family. Thus the $\theta_V$ glue to an
actual closed form $\theta$ on $Y^\circ$.
The same argument includes a finite \'etale quotient of an Inoue
surface. Locally over a simply connected base disc its finite
cover extends to the family: lift a differentiable marking by
Ehresmann's theorem and equip the resulting covering space with
the pulled-back complex structure. Passing to the Galois closure
when necessary gives a family of finite covers of Inoue surfaces;
these are again Inoue surfaces, as is also seen from their
finite-index lattice presentations. Every deck transformation of
the quotient normalizes a finite-index Inoue deck group on
$\mathbb H\times\mathbb C$. Its first-coordinate group still
contains a nonzero translation and has common boundary fixed
point precisely $\infty$. The preceding conjugation argument
therefore makes each quotient deck transformation real affine in
$w$, and its coefficients are constant over the base disc.
Consequently \eqref{cdp:eq:C.IN.3} descends through the finite quotient as
well. Its pullback to the Inoue cover has the nonzero period
$\log\alpha$, so the descended fibre class is nonzero.
For a fibre loop represented on its universal cover by the deck
transformation $w\mapsto\alpha w$, its integral is
$\log\alpha\ne0$. Hence $[\theta]$ restricts nontrivially to
the fibre.

\noindent\textbf{Step 4.} We now check every meridian of the divisor
$D=Y\setminus Y^\circ$. Let $\iota:\Delta\to Y$ be a small
holomorphic transverse disc at a smooth point of a reduced
irreducible component of $D$, avoiding all other components.
The pullback of the preceding transverse coordinate atlas on
$\Delta^*$ has constant real affine transitions. On the universal
cover
\[
 \mathbb H_c=\{\operatorname{Im}\tau>c\}
 \longrightarrow\Delta^*,\qquad
 \tau\longmapsto\exp(2\pi\mathrm i\tau),
\]
analytic continuation consequently gives a holomorphic developing
map $W:\mathbb H_c\to\mathbb H$ with
$W(\tau+1)=aW(\tau)+b$ for some $a>0$, $b\in\mathbb R$.
This follows by trivializing the pulled-back constant transition
cocycle on the simply connected covering space; it requires no
extension of a family trivialization across $D$.
The period of $\iota^*\theta$ on the positive puncture loop is
$\log a$.

\noindent\textbf{Step 5.} Use the Poincar\'e metrics
$|dw|/\operatorname{Im}w$ and
$|d\tau|/(\operatorname{Im}\tau-c)$.
The function $\log\operatorname{Im}w$ has differential of norm
one for the first metric. Schwarz--Pick and the horizontal segment
from $\tau$ to $\tau+1$ therefore give
\[
\begin{aligned}
 |\log a|
 &\le d_{\mathbb H}\bigl(W(\tau),W(\tau+1)\bigr)\\
 &\le d_{\mathbb H_c}(\tau,\tau+1)
 \le\frac{1}{\operatorname{Im}\tau-c}.
\end{aligned}
\tag{C.IN.5}\label{cdp:eq:C.IN.5}
\]
Letting $\operatorname{Im}\tau\to+\infty$ proves $a=1$.
Every meridian period of $\theta$ is thus zero. This calculation
applies to the reduced divisor and is unaffected by the
scheme-theoretic multiplicities of its components.

\noindent\textbf{Step 6.} For clarity, the exact topological extension is as follows.
The map $H_1(Y^\circ;\mathbb R)\to H_1(Y;\mathbb R)$ is
surjective: loops can be moved away from a complex divisor.
Its kernel is generated by the meridians of the irreducible
components of $D$. Indeed, a bounding real two-chain can be put
in general position, transverse to the smooth divisor locus and
disjoint from its singular locus, which has real codimension at
least four. Removing small discs about its intersection points
leaves a chain in $Y^\circ$ whose additional boundary circles
are meridians. This proves the kernel assertion, including
relations and arbitrary singularities of $D$.
The period homomorphism of $\theta$ therefore factors uniquely
through $H_1(Y;\mathbb R)$ and defines a class
\[
 \widetilde\theta\in H^1(Y;\mathbb R),\qquad
 \widetilde\theta|_{Y^\circ}=[\theta].
\tag{C.IN.6}\label{cdp:eq:C.IN.6}
\]
Its restriction to an Inoue fibre is nonzero by \eqref{cdp:eq:C.IN.3}.
This proves the asserted extension and the obstruction $b_1(Y)>0$.

\noindent\textbf{Step 7.} We finish by explaining why a very general Inoue fibre in \ref{cdp:st:C.3.1}
places the entire smooth locus in this situation. Ehresmann's
theorem identifies the diffeomorphism type, and in particular
the fundamental group, of all smooth fibres. The first two
Betti numbers are then $(1,0)$. These surfaces are minimal by
\ref{cdp:st:C.3.1} and non-K\"ahler. The compact-surface classification
gives class VII surfaces or non-K\"ahler elliptic surfaces
\cite[Theorem 3.2]{BMS2012}.
The Inoue group is solvable and is not virtually nilpotent.
For the latter assertion, its normal lattice is $\mathbb Z^3$
in type $S_M$, and a three-dimensional integral Heisenberg group
in types $S^{+}$ and $S^{-}$. Conjugation by a generator of the
cyclic quotient has an eigenvalue $\alpha>1$ on the rational
abelianization of this lattice. Every finite-index subgroup
contains a positive power of that generator and a finite-index
sublattice, so the eigenvalue becomes $\alpha^k>1$. In a
nilpotent group repeated commutators make every such conjugation
unipotent on this rational abelianization. This proves the
claimed failure of virtual nilpotence.

\noindent\textbf{Step 8.} A Hopf group is virtually cyclic, and a Kodaira group is virtually
nilpotent; neither can equal this group. For a properly elliptic
surface of Euler number zero the canonical bundle formula makes
the base orbifold hyperbolic: its orbifold canonical degree is
positive exactly in Kodaira dimension one. The elliptic
fibration induces a surjection of its fundamental group onto the
orbifold group. That group has a finite-index closed surface
group of genus at least two, and hence is not solvable. This
also cannot be a quotient of an Inoue group. The orbifold exact
sequence is the one used in \cite[Section 3.3]{BMS2012}.
Thus every smooth fibre is an Inoue surface. Applying \eqref{cdp:eq:C.IN.6}
contradicts $b_1(Y)=0$ and completes the exclusion.
\end{proof}






\begin{proposition}\label{cdp:st:C.3.1}
\StatementRecord{cdp:st:C.3.1}{proved}{Historical events: Claude Fable 5.2 (Anthropic), effort high, 15 November 2026: carried as background without fault found. Claude Fable 5.2 (Anthropic), effort high, 16 November 2026: carried as background without fault found. GPT-6 Astra (OpenAI), ChatGPT Pro mode, effort Pro, S6\_09\_PAPERS, 17 November 2026: targeted proof check. Also checked in this writing. Pre-continuation writing events: GPT-6 Astra (OpenAI), ChatGPT Pro mode, effort Pro, S6\_10\_PAPERS, 18 November 2026 04:33:41.407298+00:00: used and rechecked in this run. GPT-6 Astra (OpenAI), ChatGPT Pro mode, effort Pro, S6\_10\_PAPERS, 18 November 2026 04:35:54.089557+00:00: targeted proof check. GPT-6 Astra (OpenAI), ChatGPT Pro mode, effort Pro, S6\_10\_PAPERS, 18 November 2026 05:35:21+00:00: used and rechecked in this run. Continuation metadata events: GPT-6 Astra (OpenAI), ChatGPT Pro mode, effort Pro, S6\_10\_PAPERS continuation, 19 November 2026: used and rechecked in this run; status-vocabulary reconciliation; metadata reconciliation only; no new mathematical check. \emph{Pre-continuation recorded proof scope.} proved. Full Chern/minimality/relative-pluricanonical proof and the very-general classification scope. Connected-fibre Leray kills the base genus. The actual tangent sequence restricts KY and c2(Y) to the fibre, giving rational c1=0 and integral c2=0. Adjunction forbids a(−1)-curve. Each coherent f*KY\textasciicircum{}m is torsion-free on P1; rank\ensuremath{\geq}2 gives a projective relative-linear-system image of dimension\ensuremath{\geq}2, contradicting a(Y)=1. Removing the countable exceptional base-change loci yields kappa\ensuremath{\leq}0 on a very general fibre. Peters Theorem3.1 p10, Teleman p2 and BBK Theorem1 p1 were opened for the exact classification and finite-cover scopes. The phrase referring to the type should explicitly remain very general. The final exact-list clause retains the current gates on C.3.7–12. Review provenance: New derivation after reading the supplied proof; not a blind review and not an independent human check. The receiving wording now explicitly states that the surface-type list concerns the very general fibre. The exact external-locus readings on C.3.7--12 are separately recorded}{the hypotheses, notation and definitions stated here; \cref{cdp:st:C.3.7}; \cref{cdp:st:C.3.8}; \cref{cdp:st:C.3.9}; \cref{cdp:st:C.3.10}; \cref{cdp:st:C.3.6}; \cref{cdp:st:D.2.7}; \cref{cdp:st:C.3.12}; \cite[Theorem 3.1]{Peters1990}; \cite[p.~2]{Teleman2008}; \cite[Theorem 1]{BBK2019}}

Let $f:Y\to C$ be a holomorphic algebraic reduction of a connected
compact complex threefold with $b_1(Y)=b_2(Y)=0$ and $a(Y)=1$.
The base is $C\simeq\PP^1$. Every smooth fibre $F$ has
\[
 c_1(K_F)=0\text{ in }H^2(F;\Q),\qquad
 c_2(F)=0\text{ in }H^4(F;\Z),\qquad
 \chi(F,\OO_F)=0,
\]
and contains no exceptional curve of the first kind. A very general
smooth fibre has $\kappa(F)\leq0$. The very general fibre's type belongs to the following
list: a complex two-torus, a bielliptic surface, a primary or secondary
Kodaira surface, a primary or secondary Hopf surface, or an Inoue
surface. Every listed type has Euler characteristic zero. This is a
preliminary necessary list. Statement \ref{cdp:st:C.3.7} below excludes the
bielliptic entry, and \ref{cdp:st:C.3.8} makes the very general Hopf fibre have
algebraic dimension zero. \ref{cdp:st:C.3.9} and \ref{cdp:st:C.3.10} exclude both Kodaira
entries. Statement \ref{cdp:st:C.3.6} realizes both primary and secondary Hopf
entries, and \ref{cdp:st:D.2.7} realizes a fixed abelian surface as general
fibre. Statement \ref{cdp:st:C.3.12} excludes the Inoue entry. Thus the
final very general list consists exactly of two-tori and
primary or secondary Hopf surfaces of algebraic dimension zero.
\end{proposition}


\begin{proof}
\noindent\textbf{Step 1.} Connected fibres give $f_*\Q_Y=\Q_C$ and the degree-one edge injection
$H^1(C;\Q)\hookrightarrow H^1(Y;\Q)$. Therefore the genus of $C$ is
zero. The holomorphic tangent sequence on a smooth fibre, with the
last map induced by the differential of $f$, is
\[
0\longrightarrow T_F\longrightarrow T_Y|_F
 \xrightarrow{df} \OO_F\otimes T_{C,f(F)}\longrightarrow0.
\]
The normal factor is a trivial holomorphic line bundle. Thus
$K_F=K_Y|_F$ and $c_2(F)=c_2(Y)|_F$. Rational $H^2(Y)$ and $H^4(Y)$
vanish, and $H^4(F;\Z)=\Z$ has no torsion. These facts give the two
Chern-class assertions. Noether's formula gives $\chi(\OO_F)=0$.
An exceptional curve $E\simeq\PP^1$ with $E^2=-1$ would have
$K_F\cdot E=-1$ by adjunction, contradicting the rational vanishing of
$c_1(K_F)$.

\noindent\textbf{Step 2.} For completeness we derive the Kodaira-dimension restriction directly.
Fix $m>0$ and put $G_m=f_*(K_Y^m)$. This is a coherent torsion-free
sheaf on the smooth curve $C$: a section killed by a nonzero base
function vanishes on a dense open subset of $Y$ and hence everywhere.
It is consequently a vector bundle. On the nonempty open set of
cohomological base change its fibre is $H^0(F,K_F^m)$. A rank at least
two would give a meromorphic map
\[
 Y\dashrightarrow \PP_C(G_m^*)
\]
whose restriction to a general fibre is defined by at least two
independent sections of a line bundle and has positive-dimensional
image. Its image has dimension at least two. The target is projective:
a vector bundle on a compact Riemann surface is algebraic, and its
projectivization is a projective variety. The image is therefore
projective and its rational functions pull back to a subfield of
$\mathcal M(Y)$ of transcendence degree at least two. This contradicts
$a(Y)=1$. Thus $\rk G_m\leq1$ for every $m$. Removing the countable
union of the exceptional base-change sets gives the claimed very
general-fibre assertion.

\noindent\textbf{Step 3.} We now apply the Enriques--Kodaira classification in the form of
\cite[Theorem 3.1]{Peters1990}. The surface is minimal by the
adjunction argument. Rational and ruled surfaces have a rational
curve on which the canonical class has nonzero degree, so they are
excluded. K3 and Enriques surfaces have Euler numbers $24$ and $12$,
respectively. Properly elliptic and general-type surfaces have
positive Kodaira dimension. The remaining K\"ahler possibilities are
tori and bielliptic surfaces; the remaining Kodaira-dimension-zero
non-K\"ahler possibilities are the two types of Kodaira surfaces.
For class VII, $b_1(F)=1$, and $e(F)=0$ gives $b_2(F)=0$. The
classification of this case gives Hopf and Inoue surfaces; its
precise scope, including finite covers, is stated in
\cite[p.~2]{Teleman2008} and \cite[Theorem 1]{BBK2019}.
No assertion about torsion in $H_1(F;\Z)$ was used.
\end{proof}






\begin{proposition}[Proved complete cohomological dichotomy]\label{cdp:st:C.3.11}
\StatementRecord{cdp:st:C.3.11}{proved}{Historical events: Claude Fable 5.2 (Anthropic), effort high, 16 November 2026: carried as background without fault found. GPT-6 Astra (OpenAI), ChatGPT Pro mode, effort Pro, S6\_09\_PAPERS, 17 November 2026: targeted proof check. Also checked in this writing. Pre-continuation writing events: GPT-6 Astra (OpenAI), ChatGPT Pro mode, effort Pro, S6\_10\_PAPERS, 18 November 2026 04:30:14.617321+00:00: used and rechecked in this run. GPT-6 Astra (OpenAI), ChatGPT Pro mode, effort Pro, S6\_10\_PAPERS, 18 November 2026 04:35:54.089557+00:00: targeted proof check. Continuation metadata events: GPT-6 Astra (OpenAI), ChatGPT Pro mode, effort Pro, S6\_10\_PAPERS continuation, 19 November 2026: used and rechecked in this run; status-vocabulary reconciliation; metadata reconciliation only; no new mathematical check. \emph{Pre-continuation recorded proof scope.} proved. Complete coherent full-fibre dichotomy argument and its exact dependency scopes. Smooth Ehresmann transport fixes Betti numbers. In the b1=4 branch the surface Kahler criterion and zero rational c1 eliminate every non-torus minimal Kahler type; D.1.12 supplies the actual full-fibre wedge/base-change maps. In the b1=1,b2=0 branch positivity of a nonzero holomorphic2-form would give a nonzero deRham class, hence pg=0 and q=1. E4 and the nonzero transgression trivialize R1. HRR and Leray force the rank-zero R2 sheaf to have length0. The parameter exact sequence then gives full-fibre dimensions(1,1,0), and exponential/Leray give Pic0(Y)=C beta. Final type exclusions retain the explicit source-locus gates on the independently reviewed C.3.7–12 inputs. Review provenance: New derivation after reading the supplied proof; not a blind review and not an independent human check}{the hypotheses, notation and definitions stated here; \cref{cdp:st:C.3.1}; \cref{cdp:st:C.3.12}; \cref{cdp:st:C.3.6}; \cref{cdp:st:D.2.7}; \cref{cdp:st:D.1.9}; \cref{cdp:st:C.3.7}; \cref{cdp:st:C.3.9}; \cref{cdp:st:C.3.10}; \cref{cdp:st:C.3.8}; \cref{cdp:st:D.1.12}; \cref{cdp:st:E4}}


Under the hypotheses of \ref{cdp:st:C.3.1}, a very general smooth fibre
is a two-torus, a primary or secondary Hopf surface of
algebraic dimension zero. The Inoue exclusion is \ref{cdp:st:C.3.12}.
Both Hopf
entries and the torus entry occur, by \ref{cdp:st:C.3.6} and \ref{cdp:st:D.2.7}.
For the entire smooth locus there are exactly the following
two cohomological alternatives:
\[
\begin{array}{c|c|c}
\text{smooth fibres}&
 (h^0(\OO),h^1(\OO),h^2(\OO))&
 \text{coherent direct images}\\ \hline
\text{two-tori}&(1,2,1)&
 R^2f_*\OO_Y\simeq\det R^1f_*\OO_Y\\
b_1=1,\ b_2=0&(1,1,0)&
 R^1f_*\OO_Y\simeq\OO_{\PP^1},\
 R^2f_*\OO_Y=0 .
\end{array}
\tag{C.PK.9}\label{cdp:eq:C.PK.9}
\]
The indicated cohomology dimensions hold on every
scheme-theoretic full fibre in the respective branch,
and both direct images satisfy base change there.
In the second branch
$H^q(Y,\OO_Y)=(\C,\C,0,0)$ and
$\Pic^0(Y)=\C\beta$ with the exponential convention
of \ref{cdp:st:D.1.9}. The first line records a cohomological
alternative rather than a positivity assertion about
the splitting of $R^1f_*\OO_Y$.
\end{proposition}


\begin{proof}
\noindent\textbf{Step 1.} \ref{cdp:st:C.3.7}, \ref{cdp:st:C.3.9}, \ref{cdp:st:C.3.10} and \ref{cdp:st:C.3.12} remove the bielliptic,
Kodaira and Inoue types from \ref{cdp:st:C.3.1}; \ref{cdp:st:C.3.8} gives the Hopf algebraic
dimension. All smooth fibres are diffeomorphic over the
connected smooth base. In the torus branch they have
$b_1=4$, $b_2=6$, and are K\"ahler by the compact-surface
criterion. Their rational first Chern class and second
Chern number vanish by \ref{cdp:st:C.3.1}. The surface classification
then makes them two-tori: an effective nonzero canonical
divisor has strictly positive degree against a K\"ahler
form, whereas this rational first Chern class has zero
degree. The minimal K\"ahler surfaces of Kodaira dimension
zero with these Betti numbers are exactly the two-tori.
\ref{cdp:st:D.1.12} supplies this branch's full-fibre assertions.

\noindent\textbf{Step 2.} In the Hopf branch every smooth fibre has
$b_1=1$ and $b_2=0$. A nonzero holomorphic two-form
is closed and has strictly positive integral against
its conjugate, so it defines a nonzero de Rham
cohomology class. Thus $b_2=0$ forces $p_g=0$.
Since $\chi(\OO)=0$ by \ref{cdp:st:C.3.1}, one has $q=1$.
By \ref{cdp:st:E4}, $\mathscr E=R^1f_*\OO_Y$ is a line bundle,
and its nowhere-zero transgression section $\beta$
from \ref{cdp:st:D.1.9} trivializes it. The coherent sheaf
$\mathscr G=R^2f_*\OO_Y$ has rank zero, hence finite
length. Hirzebruch--Riemann--Roch on $Y$ and coherent
Leray give the exact equality
\[
 0=\chi(Y,\OO_Y)
   =\chi(\OO_{\PP^1})-\chi(\mathscr E)+\chi(\mathscr G)
   =1-1+\operatorname{length}(\mathscr G).
\]
Therefore $\mathscr G=0$. Top-degree base change and
the parameter-multiplication exact sequence give
$h^2(\OO_{Y_c})=0$ and $h^1(\OO_{Y_c})=1$ on each
full fibre; \ref{cdp:st:E4} gives $h^0(\OO_{Y_c})=1$.
Leray now gives the stated global structure-sheaf
cohomology. Finally $H^1(Y;\Z)=0$ and the exponential
sequence identify $\Pic^0(Y)$ with
$H^1(Y,\OO_Y)=\C\beta$, as asserted.
\end{proof}






\begin{proposition}\label{cdp:st:C.3.5}
\StatementRecord{cdp:st:C.3.5}{proved}{Historical events: Claude Fable 5.2 (Anthropic), effort high, 16 November 2026: carried as background without fault found. GPT-6 Astra (OpenAI), ChatGPT Pro mode, effort Pro, S6\_09\_PAPERS, 17 November 2026: targeted proof check. Also checked in this writing. Pre-continuation writing events: GPT-6 Astra (OpenAI), ChatGPT Pro mode, effort Pro, S6\_10\_PAPERS, 18 November 2026 04:23:40Z: targeted proof check. Continuation metadata events: GPT-6 Astra (OpenAI), ChatGPT Pro mode, effort Pro, S6\_10\_PAPERS continuation, 19 November 2026: used and rechecked in this run; status-vocabulary reconciliation; metadata reconciliation only; no new mathematical check. \emph{Pre-continuation recorded proof scope.} proved as stated. Full Euler-number identity with no sign assigned to an individual singular fibre. Ehresmann transport makes the smooth fibres a finite-type bundle and their Euler number is0 by the Chern restriction calculation, without requiring every complex type to be classified. Additivity of compactly supported Euler characteristic therefore leaves precisely the finite exceptional-fibre sum. Reduction preserves the underlying topological fibre; Chern–Gauss–Bonnet gives the integral c3 equality. Review provenance: New derivation after reading the supplied proof; not a blind review and not an independent human check}{the hypotheses, notation and definitions stated here; \cref{cdp:st:C.3.1}}

For every threefold and fibration in \ref{cdp:st:C.3.1}, including every type on
its necessary list, the Euler-number identity is
\[
\int_Y c_3(Y)=e(Y)=\sum_{c\in\Delta}e\bigl((Y_c)_{\rm red}\bigr),
\]
where $\Delta$ contains the finitely many critical values.
The identity alone assigns no sign to the summands.
\end{proposition}


\begin{proof}
\noindent\textbf{Step 1.} All the listed smooth surfaces have Euler number zero by \ref{cdp:st:C.3.1}.
Additivity of compactly supported Euler characteristic for the
stratification by the smooth-fibre locus and the finitely many
exceptional fibres gives
$e(Y)=e_c(C\setminus\Delta)e(F)+\sum_{c\in\Delta}e(Y_c)$.
The first term vanishes. A scheme-theoretic fibre and its reduction
have the same underlying topological space, so their topological
Euler numbers agree. Chern--Gauss--Bonnet gives the first equality.
\end{proof}






\begin{proposition}[Proved actual occurrence of both Hopf types]\label{cdp:st:C.3.6}
\StatementRecord{cdp:st:C.3.6}{proved}{Historical events: Claude Fable 5.2 (Anthropic), effort high, 16 November 2026: carried as background without fault found. GPT-6 Astra (OpenAI), ChatGPT Pro mode, effort Pro, S6\_09\_PAPERS, 17 November 2026: targeted proof check. Also checked in this writing. Pre-continuation writing events: GPT-6 Astra (OpenAI), ChatGPT Pro mode, effort Pro, S6\_10\_PAPERS, 18 November 2026 04:23:40Z: targeted proof check. Continuation metadata events: GPT-6 Astra (OpenAI), ChatGPT Pro mode, effort Pro, S6\_10\_PAPERS continuation, 19 November 2026: used and rechecked in this run; status-vocabulary reconciliation; metadata reconciliation only; no new mathematical check. \emph{Pre-continuation recorded proof scope.} proved as stated. Full primary and secondary Hopf occurrence constructions and meromorphic-field calculation. Norm estimates prove freeness/properness of the C-action, and the two strictly monotone norm equations give its actual smooth global slice S3×S3. Local tautological-line frames give the diagonal Hopf fibre. After meromorphic Hartogs extension, the normally convergent sum of the unit cocycle produces the explicit gauge v(Qw)/v(w)=u(0)/u(w). Rational independence of1,sqrt2 makes every joint eigenfunction one monomial, so fibre meromorphic functions are constant. Local sections descend the field exactly to P1. The commuting scalar mu\_m action is free; transfer and UCT give b1=b2=0 and H2=Z/m in the secondary examples, and the actual fibre deck group is Z×mu\_m. Review provenance: New derivation after reading the supplied proof; not a blind review and not an independent human check}{the hypotheses, notation and definitions stated here; \cref{cdp:eq:C.H.1}}


There are compact complex threefolds with $b_1=b_2=0$ and $a=1$
whose holomorphic algebraic reduction is a smooth bundle over
$\PP^1$ with primary Hopf fibre. There are also such threefolds
with secondary Hopf fibre. In these examples every fibre has Euler
number zero and $c_3=0$.
\end{proposition}


\begin{proof}
\noindent\textbf{Step 1.} Put $a_1=1$, $a_2=\sqrt2$ and
$W=(\C^2\setminus\{0\})_z\times(\C^2\setminus\{0\})_w$.
The additive complex Lie group $\C$ acts holomorphically by
\[
t\cdot(z,w_1,w_2)
 =\bigl(e^t z,e^{ia_1t}w_1,e^{ia_2t}w_2\bigr).
\tag{C.H.1}\label{cdp:eq:C.H.1}
\]
This action is free. Indeed $e^t z=z$ implies $t=2\pi in$ with
$n\in\Z$; a nonzero $w_j$ then gives $e^{-2\pi a_jn}=1$ and
therefore $n=0$. It is proper as well. If a point and its translate
range in fixed compact subsets of $W$, the $z$-norm bounds
$\operatorname{Re}t$. Writing $s=\operatorname{Im}t$, the identity
\[
\|t\cdot w\|^2
 =e^{-2a_1s}|w_1|^2+e^{-2a_2s}|w_2|^2
\]
and the positive constants $a_1,a_2$ bound $s$ both above and
below, since both $w$-norms are bounded away from zero. The inverse
image of a compact subset under the action graph is consequently
closed and bounded in the group variable and compact in the other
variable. This proves properness.

\noindent\textbf{Step 2.} Every orbit meets the real submanifold
$\{\|z\|=\|w\|=1\}=S^3\times S^3$ exactly once.
The first equation fixes $\operatorname{Re}t=-\log\|z\|$.
The displayed squared $w$-norm is strictly decreasing in $s$,
with limits $+\infty$ and $0$, so the second equation has one
solution. Its derivative is strictly negative, and the implicit
function theorem makes the resulting cross-section smooth.
Thus the complex quotient $Y=W/\C$ is a smooth compact complex
threefold diffeomorphic to $S^3\times S^3$. In particular
$b_1(Y)=b_2(Y)=0$ and $e(Y)=0$.

\noindent\textbf{Step 3.} The map
\[
f:Y\longrightarrow\PP^1,\qquad [z,w]\longmapsto[z_1:z_2]
\]
is well-defined and holomorphic. Choose a holomorphic nonzero
frame of the tautological line on a small base disc. Scaling the
$z$-coordinate to this frame leaves exactly the transformations
$t=2\pi in$, which act on $w$ by powers of
\[
Q=\operatorname{diag}(q_1,q_2),\qquad
q_j=e^{-2\pi a_j}.
\]
This identifies $f$ over that disc with the product having fibre
$H=(\C^2\setminus\{0\})/\langle Q\rangle$.
The quotient is a primary Hopf surface, and these local products
prove that $f$ is a proper holomorphic submersion with connected
fibres.

\noindent\textbf{Step 4.} We verify $a(H)=0$ directly. A meromorphic function on $H$ lifts
to a $Q$-invariant meromorphic function on $\C^2\setminus\{0\}$.
It extends meromorphically across the origin: the closure of its
polar divisor is analytic by the codimension-two extension theorem
for divisors, is locally principal on the smooth ball, and
multiplying by a local defining function reduces extension to
Hartogs' theorem for a holomorphic function. Write the germ as
$g/h$ with relatively prime convergent power series. The identity
$(g\circ Q)h=g(h\circ Q)$ gives
\[
g\circ Q=ug,\qquad h\circ Q=uh
\]
for one holomorphic unit $u$: divisibility follows from relative
primality, and preservation of the vanishing order makes the
quotient a unit. Write $u/u(0)=e^\ell$ with $\ell(0)=0$ and set
\[
v(w)=\exp\left(\sum_{n\ge0}\ell(Q^nw)\right).
\]
The sum converges normally on a sufficiently small ball, because
$Q$ is a strict contraction. It gives
$v(Qw)/v(w)=u(0)/u(w)$. Consequently both $vg$ and $vh$ are
eigenfunctions of pullback by $Q$ with the same eigenvalue $u(0)$.
The monomial $w_1^rw_2^s$ has eigenvalue
$e^{-2\pi(r+s\sqrt2)}$. These numbers are distinct for distinct
$(r,s)\in\Z_{\ge0}^2$. Comparing power-series coefficients shows
that each nonzero eigenfunction is a single monomial and that the
two eigenfunctions are proportional. The meromorphic function is
therefore constant, first near the origin and then everywhere by
analytic continuation.

\noindent\textbf{Step 5.} Every meromorphic function on $Y$ is constant on each general
fibre of $f$. On a product chart, restricting to a local section
gives the descended meromorphic function on the base; equality on
the fibres and the identity theorem identify its pullback with
the original function. These local functions agree on overlaps
and extend at the finitely many possible poles on the base.
Thus $\mathcal M(Y)=f^*\mathcal M(\PP^1)$, proving $a(Y)=1$
and proving that $f$ is the holomorphic algebraic reduction.

\noindent\textbf{Step 6.} For any integer $m>1$, let $\mu_m$ act on $W$ by
$(z,w)\mapsto(z,\zeta w)$. It commutes with \eqref{cdp:eq:C.H.1} and induces
a free action on $Y$. A fixed orbit would give $t=2\pi in$ from
the $z$-coordinate; equality of norms in a nonzero $w_j$ then
gives $n=0$ and $\zeta=1$. The quotient
$Y_m=Y/\mu_m$ is a smooth compact complex threefold. Transfer
injects $H^j(Y_m;\Q)$ into $H^j(Y;\Q)$, so $b_1(Y_m)=b_2(Y_m)=0$.
The simply connected cover $Y\simeq S^3\times S^3$ gives
$\pi_1(Y_m)=\mu_m$, and the universal coefficient theorem gives
$H^2(Y_m;\Z)\simeq\Z/m\Z$: its free part is zero by the
rational calculation, and its torsion is
$\operatorname{Ext}(H_1(Y_m;\Z),\Z)$. Thus these secondary
examples meet exactly the rational hypotheses of C.3; the primary
example also has $H^1(Y;\Z)=H^2(Y;\Z)=0$.
The induced map to $\PP^1$ has fibre
\[
H_m=(\C^2\setminus\{0\})/
       \bigl(\langle Q\rangle\times\mu_m\bigr).
\]
Its universal cover is $\C^2\setminus\{0\}$ and its fundamental
group is $\Z\times\mu_m$, which is noncyclic. It is therefore a
secondary Hopf surface, with the displayed primary cover of degree
$m$. Any meromorphic function on $H_m$ pulls back to one on $H$,
so $a(H_m)=0$. The same local descent proves
$\mathcal M(Y_m)=f_m^*\mathcal M(\PP^1)$ and $a(Y_m)=1$.
Finally transfer of Euler characteristic under the finite covering,
or the smooth-fibration Euler formula, gives $e(Y_m)=0$.
\end{proof}






\section{Cotangent torsion, conductor and the printed endpoint error}\label{sec:II:04}


Throughout this section $i:S\hookrightarrow Y$ is the inclusion of a reduced effective Cartier divisor into a complex manifold of dimension $n+1$, $j:S_{\rm reg}\hookrightarrow S$ is the inclusion of its regular locus, and $A$ is an invertible analytic $\OO_S$-module. We write
\[
 N^*=\OO_Y(-S)|_S,\qquad V=\Omega_Y^1|_S,\qquad
 T=\ker(\Omega_S^1\longrightarrow j_*\Omega^1_{S_{\rm reg}}),\qquad Q=\Omega_S^1/T.
\]
All maps between sheaves below are $\OO_S$-linear maps of coherent analytic sheaves, except for the displayed restriction into $j_*\Omega^1_{S_{\rm reg}}$. The latter is also $\OO_S$-linear; its kernel and image have the coherent descriptions proved below.

\begin{proposition}[The conormal filtration and its exact section criterion]\label{cdp:st:A.1.1}
\StatementRecord{cdp:st:A.1.1}{proved}{Historical events: Claude Fable 5.2 (Anthropic), effort high, 16 November 2026: carried as background without fault found. Claude Fable 5.2 (Anthropic), effort high, 17 November 2026: carried as background without fault found. GPT-6 Astra (OpenAI), ChatGPT Pro mode, effort Pro, S6\_09\_PAPERS, 17 November 2026: targeted proof check. Also checked in this writing. Pre-continuation writing events: GPT-6 Astra (OpenAI), ChatGPT Pro mode, effort Pro, S6\_10\_PAPERS, 18 November 2026: targeted proof check; used and rechecked in this run. Continuation mathematical and source checks: GPT-6 Astra (OpenAI), ChatGPT Pro mode, effort Pro, S6\_10\_PAPERS continuation, 19 November 2026: computation or formal certificate; independent re-proof; targeted proof check; used and rechecked in this run. Continuation metadata events: GPT-6 Astra (OpenAI), ChatGPT Pro mode, effort Pro, S6\_10\_PAPERS continuation, 19 November 2026: used and rechecked in this run; status-vocabulary reconciliation; metadata reconciliation only; no new mathematical check. \emph{Continuation recorded proof scope.} Complete current proof re-derived before reading: the reduced-hypersurface conormal injection, coherent torsion kernel, both exact filtration sequences and the two canonical connecting-map criteria. New independent node and cusp module computations passed28 exact checks. Renewed source comparison identifies the classical conormal and long-exact-cohomology mechanisms and the exact AC26 normal-crossing obstruction through the same N*, V, T, Q and kernel-of-delta maps; source construction dependencies receive no additional proof credit.}{the hypotheses, notation and definitions stated here; the explicit constructions and calculations in the proof}
\label{wbII:prop:criterion}\label{wbII:prop:filtration}\label{wbII:lem:conormal}\label{wbII:lem:conormal-injective}\label{wbII:prop:missing-global-section-criterion}\label{wbII:prop:missing-cotangent-filtration}
The conormal map $d:N^*\to V$, defined by $[F]\mapsto dF|_S$ for a local defining equation $F$, is injective. With $K=\ker(V\to Q)$, there are canonical exact sequences
\begin{align}
0&\longrightarrow N^*\xrightarrow{d}K\longrightarrow T\longrightarrow0,
\label{cdp:A:filtration1}\\
0&\longrightarrow K\longrightarrow V\longrightarrow Q\longrightarrow0.
\label{cdp:A:filtration2}
\end{align}
Let $\partial_A:H^0(S,Q\otimes A)\to H^1(S,K\otimes A)$ and $\delta_A:H^0(S,T\otimes A)\to H^1(S,N^*\otimes A)$ be their connecting maps. Then
\[
 H^0(S,V\otimes A)=0
 \quad\Longleftrightarrow\quad
 H^0(S,K\otimes A)=0\ \hbox{and}\ \partial_A\hbox{ is injective}.
\tag{A.1}\label{cdp:eq:A.1}
\]
Under the explicit hypotheses $H^0(S,N^*\otimes A)=0$ and $H^0(S,Q\otimes A)=0$, $H^0(S,V\otimes A)$ is canonically $\ker\delta_A$.
\end{proposition}


\begin{proof}
\noindent\textbf{Step 1.} Locally $R=\C\{x_0,\ldots,x_n\}/(F)$ is reduced. Its total quotient ring is the product of the fraction fields of its irreducible components. At the generic point of every component, $S$ is smooth and $dF$ is nonzero. Consequently $a\,dF=0$ in $R^{n+1}$ implies that $a$ vanishes in each of those fields. The map from the reduced ring $R$ into their product is injective, so $a=0$. This proves injectivity in the usual right-exact conormal sequence. Taking the inverse image of $T\subset\Omega_S^1$ gives \eqref{cdp:A:filtration1} and \eqref{cdp:A:filtration2}. Tensoring by $A$ is exact because $A$ is locally free. The cohomology sequence of \eqref{cdp:A:filtration2} identifies the image of $H^0(VA)$ in $H^0(QA)$ with $\ker\partial_A$ and its kernel with $H^0(KA)$, proving \eqref{cdp:eq:A.1}. When $H^0(QA)=0$, it identifies $H^0(VA)$ with $H^0(KA)$. The first sequence then identifies this group with $\ker\delta_A$ when $H^0(N^*A)=0$.
\end{proof}





\begin{proposition}[The Jacobian fractional ideal for every reduced hypersurface]\label{cdp:st:A.2.1}
\StatementRecord{cdp:st:A.2.1}{proved}{Historical events: Claude Fable 5.2 (Anthropic), effort high, 16 November 2026: carried as background without fault found. GPT-6 Astra (OpenAI), ChatGPT Pro mode, effort Pro, S6\_09\_PAPERS, 17 November 2026: targeted proof check. Also checked in this writing. Pre-continuation writing events: GPT-6 Astra (OpenAI), ChatGPT Pro mode, effort Pro, S6\_10\_PAPERS, 18 November 2026 04:12:55Z: targeted proof check. Continuation mathematical and source checks: GPT-6 Astra (OpenAI), ChatGPT Pro mode, effort Pro, S6\_10\_PAPERS continuation, 19 November 2026: computation or formal certificate; independent re-proof; used and rechecked in this run. Continuation metadata events: GPT-6 Astra (OpenAI), ChatGPT Pro mode, effort Pro, S6\_10\_PAPERS continuation, 19 November 2026: used and rechecked in this run; status-vocabulary reconciliation; metadata reconciliation only; no new mathematical check. \emph{Continuation recorded proof scope.} Continuation complete current proof derived before reading at08:02:36 UTC: generic kernel q dF, inverse Jacobian/Hom map, coherent ideal, full conormal line transition and torsion quotient. New exact polynomial/module program passed173checks at its recorded finite scope; universal proof written separately. Granger–Schulze2014 Proposition3.4 pp9–10 and Corollary3.5 p10 of the downloaded author original reopened; exact multiplier q maps through [F] tensor q to q dF and its torsion quotient. Known residue/inverse-Jacobian credit retained. No broader dependency proof credit.}{the hypotheses, notation and definitions stated here; \cref{cdp:A:filtration1}}

Let $\mathcal J\subset\OO_S$ be the Jacobian ideal, generated locally by $F_{x_0},\ldots,F_{x_n}$ modulo $F$. In the total quotient sheaf $\mathcal M_S$ define
\[
 \mathcal J^{-1}=\{q\in\mathcal M_S:q\mathcal J\subset\OO_S\}.
\]
There are canonical identifications
\[
 K\simeq N^*\otimes\mathcal J^{-1},\qquad
 T\simeq N^*\otimes(\mathcal J^{-1}/\OO_S),
\tag{A.2}\label{cdp:eq:A.2}
\]
under which the first map in \eqref{cdp:A:filtration1} is induced by $\OO_S\hookrightarrow\mathcal J^{-1}$. Thus (A.1), with this explicitly defined fractional ideal, is an exact necessary and sufficient criterion for every reduced hypersurface.
\end{proposition}


\begin{proof}
\noindent\textbf{Step 1.} The ideal $J$ contains a nonzerodivisor: it is nonzero at every minimal prime, and prime avoidance supplies an element outside all minimal primes. Every $R$-linear map $\varphi:J\to R$ is multiplication by a unique element of the total quotient ring. Indeed, for a nonzerodivisor $a\in J$, the identity $a\varphi(b)=b\varphi(a)$ for $b\in J$ gives the multiplier $q=\varphi(a)/a$. Hence $J^{-1}=\Hom_R(J,R)$ is coherent.

\noindent\textbf{Step 2.} At the generic points, an element $\alpha\in K$ is uniquely $q\,dF$ for a tuple of meromorphic functions $q$. The coefficients of $\alpha$ belong to $R$ exactly when $qF_{x_i}\in R$ for every $i$, which is the condition $q\in J^{-1}$. Conversely such a vector restricts to zero in the cotangent space of every regular component, hence belongs to $K$. The map $[F]\otimes q\mapsto q\,dF|_S$ proves the first identification. Its quotient by $R\,dF$ gives the second. Replacing $F$ by $uF$, with $u$ a holomorphic unit, replaces $dF|_S$ by $u\,dF|_S$ and replaces the conormal generator by $u[F]$. Coordinate changes transform the derivative vector by an invertible Jacobian matrix. These are exactly the transitions in \eqref{cdp:eq:A.2}, so the identifications are canonical.
\end{proof}






\begin{proposition}[Normalization, conductor, and the excess kernel]\label{cdp:st:A.2.2}
\StatementRecord{cdp:st:A.2.2}{proved}{Historical events: Claude Fable 5.2 (Anthropic), effort high, 16 November 2026: carried as background without fault found. GPT-6 Astra (OpenAI), ChatGPT Pro mode, effort Pro, S6\_09\_PAPERS, 17 November 2026: targeted proof check. Also checked in this writing. Pre-continuation writing events: GPT-6 Astra (OpenAI), ChatGPT Pro mode, effort Pro, S6\_10\_PAPERS, 18 November 2026 04:12:55Z: targeted proof check. GPT-6 Astra (OpenAI), ChatGPT Pro mode, effort Pro, S6\_10\_PAPERS, 18 November 2026 05:35:21+00:00: used and rechecked in this run. Continuation metadata events: GPT-6 Astra (OpenAI), ChatGPT Pro mode, effort Pro, S6\_10\_PAPERS continuation, 19 November 2026: used and rechecked in this run; status-vocabulary reconciliation; metadata reconciliation only; no new mathematical check. \emph{Pre-continuation recorded proof scope.} proved as stated. Complete displayed conductor, excess-kernel and global-section argument. Re-derived the monic Weierstrass trace pairing: the coefficient-of-z\textasciicircum{}(m−1) pairing is unimodular and equals Tr(r/P\ensuremath{{}^{\prime}}), so integral normalization elements belong to P\ensuremath{{}^{\prime}}\textasciicircum{}−1 R. Choosing a basis of permissible projection directions gives JB\ensuremath{\subset}R. Multipliers q with qJ\ensuremath{\subset}B descend exactly when their J-products vanish in B/R. The exact ring morphisms r\ensuremath{\to}(r,r mod c) and (b,rbar)\ensuremath{\to}b−rbar yield the conductor square; locally free tensor and H0 preserve its kernel, including nonreduced conductor schemes. Review provenance: New derivation after reading the supplied proof; not a blind review and not an independent human check}{the hypotheses, notation and definitions stated here; the explicit constructions and calculations in the proof}
\label{wbII:lem:conductor}\label{wbII:lem:conductor-square}
Let $\eta:\widetilde S\to S$ be the finite analytic normalization, $\mathcal B=\eta_*\OO_{\widetilde S}$, and $\mathfrak c=\Ann_{\OO_S}(\mathcal B/\OO_S)$ its conductor. There are canonical inclusions
\[
 \mathcal J\subset\mathfrak c,\qquad
 \mathcal B\subset\mathcal J^{-1},\qquad
 \eta_*\eta^*N^*\hookrightarrow K.
\tag{A.3}\label{cdp:eq:A.3}
\]
In a local ring $R$ with normalization $B$, put $D=\Hom_B(JB,B)$, viewed as a fractional ideal. Then
\[
 J^{-1}=\ker\left[D\longrightarrow
 \Hom_R\bigl(J,(B/\mathfrak c)/(R/\mathfrak c)\bigr)\right],
\quad q\longmapsto\bigl(b\mapsto qb\bmod R\bigr).
\tag{A.4}\label{cdp:eq:A.4}
\]
For any locally free $\OO_S$-module $E$, the conductor gives a second, independent exact section criterion:
\begin{multline}
H^0(S,E)=
\ker\left[H^0(\widetilde S,\eta^*E)\oplus H^0(D_c,E|_{D_c})\right.\\
\left.\longrightarrow H^0(\widetilde D_c,\eta^*E|_{\widetilde D_c})\right],
\qquad (a,b)\longmapsto a|_{\widetilde D_c}-\eta^*b,
\tag{A.5}\label{cdp:eq:A.5}
\end{multline}
where $D_c$ and $\widetilde D_c$ are defined by $\mathfrak c$ in $S$ and $\widetilde S$. In particular, the ambient group vanishes exactly when the map in \eqref{cdp:eq:A.5}, for $E=V\otimes A$, is injective.
\end{proposition}


\begin{proof}
\noindent\textbf{Step 1.} We prove $JB\subset R$ rather than assume a conductor formula. Choose a generic linear projection for which Weierstrass preparation writes the equation as a unit times a monic polynomial $P(z)$ over $A_0=\C\{x_1,\ldots,x_n\}$. Thus $R=A_0[z]/(P)$ is finite free over $A_0$ and generically separable. The functional taking the coefficient of $z^{m-1}$ in a remainder modulo $P$, where $m=\deg P$, generates $\Hom_{A_0}(R,A_0)$ as an $R$-module: the matrix of the pairing on $1,z,\ldots,z^{m-1}$ has zeros below the reverse diagonal and ones on that diagonal. Over the fraction field, Lagrange interpolation identifies that functional with $r\mapsto\operatorname{Tr}(r/P'(z))$. Therefore the trace dual is $P'(z)^{-1}R$.

\noindent\textbf{Step 2.} For $b\in B$ and $r\in R$, the element $br$ is integral over the integrally closed ring $A_0$. Its trace belongs to $A_0$. It follows that $b$ lies in this trace dual, so $P'(z)b\in R$. Modulo $F$, the derivative in the chosen projection direction differs from $P'(z)$ by a unit. The permissible projection directions form a nonempty open set; choosing a linear basis of such directions proves $F_{x_i}B\subset R$ for every $i$. This proves $J\subset\mathfrak c$ and hence $B\subset J^{-1}$. The last inclusion of \eqref{cdp:eq:A.3} is (A.2) and the projection formula.

\noindent\textbf{Step 3.} The module $D$ consists exactly of the meromorphic $q$ satisfying $qJ\subset B$. Requiring that these products descend to $R$ is precisely the vanishing of their classes in $B/R=(B/\mathfrak c)/(R/\mathfrak c)$, proving \eqref{cdp:eq:A.4}. Finally the ring map
\[
 R\longrightarrow B\oplus R/\mathfrak c,
 \qquad r\longmapsto(r,r\bmod\mathfrak c)
\]
has image the kernel of $(b,\bar r)\mapsto b-\bar r\in B/\mathfrak c$. Surjectivity onto $B/\mathfrak c$ comes from the first summand. Tensoring this exact sequence by the free module corresponding to $E$ and taking global sections gives \eqref{cdp:eq:A.5}. All conductor restrictions are scheme-theoretic, so this argument includes nonreduced conductors.
\end{proof}



The fractional ideal appearing in (A.2) has an established interpretation. Granger--Schulze, Proposition~3.4 and Corollary~3.5 of \cite{GS2014}, identify $J^{-1}$ with the module of logarithmic residues and prove the inclusions in \eqref{cdp:eq:A.3}. The preceding proof gives the maps needed here directly. Their Theorems~1.1 and~1.3 also give the following exact extension of the normal-crossings case.




\begin{proposition}[The canonical normal-crossings formula]\label{cdp:st:A.1.2}
\StatementRecord{cdp:st:A.1.2}{proved}{Historical events: Claude Fable 5.2 (Anthropic), effort high, 16 November 2026: carried as background without fault found. GPT-6 Astra (OpenAI), ChatGPT Pro mode, effort Pro, S6\_09\_PAPERS, 17 November 2026: targeted proof check. Also checked in this writing. Pre-continuation writing events: GPT-6 Astra (OpenAI), ChatGPT Pro mode, effort Pro, S6\_10\_PAPERS, 18 November 2026: targeted proof check. Continuation mathematical and source checks: GPT-6 Astra (OpenAI), ChatGPT Pro mode, effort Pro, S6\_10\_PAPERS continuation, 19 November 2026: computation or formal certificate; independent re-proof; used and rechecked in this run. Continuation metadata events: GPT-6 Astra (OpenAI), ChatGPT Pro mode, effort Pro, S6\_10\_PAPERS continuation, 19 November 2026: used and rechecked in this run; status-vocabulary reconciliation; metadata reconciliation only; no new mathematical check. \emph{Continuation recorded proof scope.} Continuation complete current proof derived before reading at08:02:36 UTC: all-dimensional branch divisibility, exact inverse, fractional normalization/idempotent map, double/triple quotients and global gluing. New exact polynomial/module program passed173checks at its recorded finite scope; universal proof written separately. Original AC26 pp86–87 Definition10.1/Lemmas10.2–10.3 reopened; the same branchwise coefficient $u a_i G_i$ and conductor descent/gluing maps matched, retaining the conormal unit. Known source construction credit retained. No broader dependency proof credit.}{the hypotheses, notation and definitions stated here; \cref{cdp:st:A.2.1}}

For a reduced normal-crossings divisor, the map in (A.3) is an isomorphism. In a chart with $F=x_1\cdots x_r$, it is explicitly
\[
 \bigoplus_{i=1}^r R/(x_i)\longrightarrow K,
 \qquad (a_i)_i\longmapsto
 \sum_{i=1}^r a_i\prod_{j\ne i}x_j\,dx_i.
\tag{A.6}\label{cdp:eq:A.6}
\]
The tensor factor $N^*$ is trivialized here by $[F]$. At a double crossing $T\simeq R/(x,y)$; at a triple crossing,
\[
 T\simeq R^3/\bigl(R(1,1,1)+R(x,0,0)+R(0,y,0)+R(0,0,z)\bigr).
\tag{A.7}\label{cdp:eq:A.7}
\]
\end{proposition}


\begin{proof}
\noindent\textbf{Step 1.} On the normalization branch $x_i=0$, an ambient form maps to zero in the cotangent bundle of that branch precisely when all coefficients except that of $dx_i$ vanish. Thus a coefficient of $dx_i$ must vanish on every branch $x_j=0$ for $j\ne i$. In the convergent power-series ring it is divisible by $\prod_{j\ne i}x_j$. Its remaining factor is defined modulo $x_i$, which gives exactly \eqref{cdp:eq:A.6}. The coefficients of different $dx_i$ show injectivity. Coefficients of the coordinates not among $x_1,\ldots,x_r$ vanish on every branch and hence vanish in $R$. This proves surjectivity in every dimension.

\noindent\textbf{Step 2.} For $r=1$ the map is $a\mapsto a\,dx$. For $r=2$ it is $(a,b)\mapsto ay\,dx+bx\,dy$, and quotienting by $dF=y\,dx+x\,dy$ leaves the generator $y\,dx=-x\,dy$, annihilated by $(x,y)$. For $r=3$ the branch generators are $yz\,dx,xz\,dy,xy\,dz$. Their sum is $dF$ and their respective annihilators are $(x),(y),(z)$, proving \eqref{cdp:eq:A.7}. Multiplication of $F$ by a unit and changes of coordinates were checked in \ref{cdp:st:A.2.1}. Formula \eqref{cdp:eq:A.6} is the restriction of the canonical fractional-ideal map, so it glues without a chosen branch ordering, equation, or coordinate system.
\end{proof}






\begin{proposition}[Pinch points and cuspidal edges]\label{cdp:st:A.2.4}
\StatementRecord{cdp:st:A.2.4}{proved}{Historical events: Claude Fable 5.2 (Anthropic), effort high, 16 November 2026: carried as background without fault found. GPT-6 Astra (OpenAI), ChatGPT Pro mode, effort Pro, S6\_09\_PAPERS, 17 November 2026: targeted proof check. Also checked in this writing. Pre-continuation writing events: GPT-6 Astra (OpenAI), ChatGPT Pro mode, effort Pro, S6\_10\_PAPERS, 18 November 2026 04:12:55Z: targeted proof check. Continuation mathematical and source checks: GPT-6 Astra (OpenAI), ChatGPT Pro mode, effort Pro, S6\_10\_PAPERS continuation, 19 November 2026: computation or formal certificate; independent re-proof; targeted proof check; used and rechecked in this run. Continuation metadata events: GPT-6 Astra (OpenAI), ChatGPT Pro mode, effort Pro, S6\_10\_PAPERS continuation, 19 November 2026: used and rechecked in this run; status-vocabulary reconciliation; metadata reconciliation only; no new mathematical check. \emph{Continuation recorded proof scope.} Complete current proof coverage is the ledger reviewer\textquotesingle s universal local derivation frozen at 09:16:58 UTC together with VI\textquotesingle s isolated module and Hartogs route frozen at 09:18:47 UTC before target reading. Both pinching and cusp inverse-Jacobian maps, conductors, cyclic quotients and global conormal transitions match. The fresh program passes 920 exact checks within its stated bounds; universal proofs are written separately. Granger--Schulze\textquotesingle s exact source maps were reopened. The Hartogs extension input was also opened in Simoni\v c, arXiv:1608.00950v1, Theorem~1, p.~1; its actual restriction/extension map on the punctured bidisc matches the glued multiplier. The supporting-row source gate is complete.}{\cref{cdp:st:A.2.1,cdp:st:A.1.2}; \cite[Proposition 3.4, pp.~9--10, and Corollary 3.5, p.~10, author version]{GS2014}; \cite[Theorem 1, p.~1]{Simonic2016Hartogs}; the displayed normalization and multiplier maps}

At the pinch point $F=x^2-y^2z$, with normalization $x=uv,y=u,z=v^2$, write $B=\C\{u,v\}$. Then
\[
 K=B\,dF,\qquad T\simeq B/R\simeq R/(x,y),
\quad \tau=2yz\,dx-2xz\,dy-xy\,dz
\tag{A.8}\label{cdp:eq:A.8}
\]
is a torsion generator. At the cuspidal edge $F=x^2-y^3$, with normalization $x=t^3,y=t^2$ and unchanged coordinate $z$, write $B=\C\{t,z\}$. Then
\[
 K=t^{-1}B\,dF,\qquad
 T\simeq t^{-1}B/R\simeq R/(x,y^2),
 \quad\tau=2y\,dx-3x\,dy.
\tag{A.9}\label{cdp:eq:A.9}
\]
The conductor is $(x,y)$ for the pinch point and $(x,y)$ for the cusp; in the cusp normalization it is $t^2B$. An ordinary surface, meaning a surface whose singularities are double normal crossings, triple normal crossings, and pinch points, satisfies canonically $K=\eta_*\eta^*N^*$ and $T=N^*\otimes(\eta_*\OO_{\widetilde S}/\OO_S)$. For the cuspidal edge the additional quotient $J^{-1}/B$ is generated by $t^{-1}$ over $\C\{z\}$.
\end{proposition}


\begin{proof}
\noindent\textbf{Step 1.} For the pinch point,
\[
 R=\C\{u,v^2,uv\}=\C\{v^2\}+uB,
 \quad J=(uv,uv^2,u^2)=u(u,v)B.
\]
The equality of the two descriptions of $J$ follows by separating even and odd powers of $v$. An element $q\in J^{-1}$ satisfies $(qu)(u,v)\subset B$. Since $B$ is a regular local ring, $\Hom_B((u,v),B)=B$: a meromorphic function whose products with $u$ and $v$ are holomorphic has no pole along any divisor, and its possible isolated singularity is removable by the holomorphic extension theorem. Hence $qu=h\in B$. The requirements $hv,hv^2\in R$, restricted to $u=0$, say respectively that $h(0,v)v$ and $h(0,v)v^2$ are even power series. The first requires $h(0,v)$ to be odd and the second to be even. Therefore $h(0,v)=0$, so $h\in uB$ and $q\in B$. The reverse inclusion follows either directly or from (A.3). Also $B/R$ is generated by $v$ and its annihilator is $(u,uv)=(y,x)$. Multiplying $dF$ by $v=x/y$ gives the displayed $\tau$.

\noindent\textbf{Step 2.} For the cusp, $R=\C\{t^2,t^3,z\}$ and $J=(t^3,t^4)$. A meromorphic $q$ with $qt^3,qt^4\in R$ has no poles away from $t=0$; writing its Laurent expansion in $t$ shows that the lowest allowed exponent is $-1$. Indeed an exponent $-3$ produces a forbidden $t$ term after multiplication by $t^4$, an exponent $-2$ produces one after multiplication by $t^3$, and smaller exponents produce a pole. Conversely every $t^{-1}B$ has both products in $R$. Thus $J^{-1}=t^{-1}B$. Modulo $R$, the quotient has the $\C\{z\}$-basis $t^{-1},t$; multiplication by $y$ sends the first basis vector to the second, while $x$ and $y^2$ kill both. This proves the cyclic presentation and \eqref{cdp:eq:A.9}. The conductor calculations follow directly from the displayed descriptions of $R\subset B$. Formula (A.2) and the three local computations prove the ordinary-surface assertion by gluing.
\end{proof}






\begin{proposition}[Vertical residue poles at a tacnode--cusp intersection]\label{cdp:st:A.2.5}
\StatementRecord{cdp:st:A.2.5}{proved}{Historical events: Claude Fable 5.2 (Anthropic), effort high, 16 November 2026: carried as background without fault found. GPT-6 Astra (OpenAI), ChatGPT Pro mode, effort Pro, S6\_09\_PAPERS, 17 November 2026: targeted proof check. Also checked in this writing. Pre-continuation writing events: GPT-6 Astra (OpenAI), ChatGPT Pro mode, effort Pro, S6\_10\_PAPERS, 18 November 2026 04:12:55Z: targeted proof check. Continuation metadata events: GPT-6 Astra (OpenAI), ChatGPT Pro mode, effort Pro, S6\_10\_PAPERS continuation, 19 November 2026: used and rechecked in this run; status-vocabulary reconciliation; metadata reconciliation only; no new mathematical check. \emph{Pre-continuation recorded proof scope.} proved as stated. Complete local normalization, conductor, Laurent-module and vertical-residue counterexample. Adjoining v=x/(y\textsuperscript{2}z) gives finite normal B. Parity shows c=u\textsuperscript{2}v\textsuperscript{2}B. The three Jacobian products exclude v-poles below −1 and the u\textasciicircum{}−2 coefficient; writing q=(a+vb)/(uv) forces b divisible by u. Thus J\textasciicircum{}−1=A0+(uv)\textasciicircum{}−1 A0, J\textasciicircum{}−1/B=A0/(uv\textsuperscript{2}), and the R-annihilator modulo R is (x,y\textsuperscript{3}z\textsuperscript{2}). Multiplication gives the displayed regular gradient vectors, while the multiplier v\textasciicircum{}−1 remains a genuine vertical pole. No compact realization is inferred. Review provenance: New derivation after reading the supplied proof; not a blind review and not an independent human check}{the hypotheses, notation and definitions stated here; the explicit constructions and calculations in the proof}

For $F=x^2-y^4z^3$, normalization is $x=u^2v^3$, $y=u$, $z=v^2$. Put $A_0=\C\{u,v^2\}$ and $B=\C\{u,v\}$. Then
\[
 \begin{aligned}
 R&=A_0\oplus u^2v^3A_0,&\mathfrak c&=u^2v^2B,\\
 J^{-1}&=A_0\oplus(uv)^{-1}A_0,&
 J^{-1}/B&\simeq A_0/(uv^2).
 \end{aligned}
 \tag{A.9a}\label{cdp:eq:A.9a}
\]
In a trivialization of $N^*$ the torsion module is
\[
 T\simeq R/(x,y^3z^2),\qquad
 \tau=2yz\,dx-4xz\,dy-3xy\,dz.
 \tag{A.9b}\label{cdp:eq:A.9b}
\]
In particular $v^{-1}\in J^{-1}\setminus B$. For the ruling with base coordinate $v$, this is a residue constant on its generic fibres with a pole on the vertical divisor $v=0$, although the horizontal conductor $u=0$ has coefficient two.
\end{proposition}


\begin{proof}
\noindent\textbf{Step 1.} The element $v=x/(y^2z)$ belongs to the fraction field and satisfies $v^2=z$. Adjoining it gives the regular ring $B$, which is therefore the finite normalization. Separating even and odd powers of $v$ gives the displayed description of $R$. An even monomial $u^av^b$ belongs to the conductor exactly when $a\ge2$ and $b\ge2$; an odd monomial already in $R$ has $a\ge2$, $b\ge3$ and is then in the conductor. Thus $\mathfrak c=u^2v^2B$.

\noindent\textbf{Step 2.} The normalized Jacobian generators, with irrelevant nonzero constants removed, are $u^2v^3,u^3v^6,u^4v^4$. If $qJ\subset R$, the first product shows that $q$ has no other polar divisor, with possible orders at most two along $u=0$ and three along $v=0$. A term of order $v^{-3}$ would produce a forbidden $v^1$ term in $qu^4v^4$; a term of order $v^{-2}$ would produce one in $qu^2v^3$. Both coefficients must vanish, so $q$ has order at least $-1$ in $v$. A coefficient $a(v)u^{-2}$ would require both $a(v)v^3$ and $a(v)v^6$ to be even Laurent series, from the $u^0$ and $u^1$ coefficients of the first two products. The first requirement makes $a$ odd and the second makes it even, so $a=0$. Thus $q=t/(uv)$ with $t\in B$.

\noindent\textbf{Step 3.} Write $t=a(u,v^2)+v b(u,v^2)$. The first Jacobian product equals $uv^2t$ and belongs to $R$ exactly when $b$ is divisible by $u$. The other two products automatically belong to $R$ under this condition. This gives $J^{-1}=A_0\oplus(uv)^{-1}A_0$. The odd part of $B$ is $(uv)^{-1}(uv^2)A_0$, giving $J^{-1}/B$ in \eqref{cdp:eq:A.9a}. The odd part of $R$ is $(uv)^{-1}u^3v^4A_0$, so $J^{-1}/R$ is generated by $(uv)^{-1}$ with annihilator $(x,y^3z^2)$. Multiplication of $dF$ by $(uv)^{-1}$ gives \eqref{cdp:eq:A.9b}.

\noindent\textbf{Step 4.} The particular vertical pole has the completely regular gradient products
\[
 v^{-1}dF=2y^2z\,dx-4xyz\,dy-3xy^2\,dz.
\]
Pulling this expression to $B$ gives zero. It is a nonzero ambient-kernel vector, with multiplier $v^{-1}$, rather than a section of the normalization submodule. The exact identities and finite monomial checks are saved in \path{code/A_vertical_pole.py}; the Laurent argument above establishes the full convergent module. This local example does not assert the existence of a compact CDP fibration realizing that configuration.
\end{proof}



For ordinary surfaces the exact ambient-vanishing criterion is (A.1) with $K\otimes A=\eta_*\eta^*(N^*\otimes A)$. For surfaces with cuspidal-edge charts it is (A.1) with the locally specified sheaf $N^*\otimes t^{-1}B$. These specify the first group and the target of the connecting homomorphism, including the extra polar term at the cusp. Under $H^0(QA)=0$ they reduce respectively to the vanishing of these explicitly computed kernel sections. The general criterion (A.4)--(A.5) retains the full conductor and applies without a singularity-type restriction. The exact polynomial identities and exponent computations are recorded in \path{code/A_local_algebra.py} and its output.




\begin{proposition}[A transverse cusp whose critical-curve equation is only a symbolic square]\label{cdp:st:A.2.6}
\StatementRecord{cdp:st:A.2.6}{proved}{Historical events: Claude Fable 5.2 (Anthropic), effort high, 16 November 2026: carried as background without fault found. GPT-6 Astra (OpenAI), ChatGPT Pro mode, effort Pro, S6\_09\_PAPERS, 17 November 2026: targeted proof check. Also checked in this writing. Pre-continuation writing events: GPT-6 Astra (OpenAI), ChatGPT Pro mode, effort Pro, S6\_10\_PAPERS, 18 November 2026 04:12:55Z: targeted proof check. Continuation metadata events: GPT-6 Astra (OpenAI), ChatGPT Pro mode, effort Pro, S6\_10\_PAPERS continuation, 19 November 2026: used and rechecked in this run; status-vocabulary reconciliation; metadata reconciliation only; no new mathematical check. \emph{Pre-continuation recorded proof scope.} proved as stated. Complete non-lci transverse-cusp example, symbolic/ordinary square comparison and exact conormal poles. Re-derived the rank-five C\{x\} algebra via the 1,y,z,yz,z\textsuperscript{2} multiplication table and distinct residues mod 5. The three minimal generators give height two and non-lci origin. Exact differentiation gives the critical-curve identities and reducedness. On x=t\textsuperscript{5}, the Hessian rank-one direction and nonzero cubic give the cusp by solving the transverse critical coordinate, completing its square analytically and taking a local cube root of the residual unit. The weight-35 part of I\textsuperscript{2} is precisely C·xy·g1\textsuperscript{2}, excluding F; generic conormal injectivity gives symbolic-square membership. The normalized conormal wedge is t\textasciicircum{}15(5t\textsuperscript{4},6t\textsuperscript{5},7t\textsuperscript{6}), and q=t\textasciicircum{}−5(5t\textsuperscript{8}a−2b)\textsuperscript{2}, k=−20t\textasciicircum{}−1 give the stated poles. Independent exact certificate analysis/II\_independent\_A26\_certificate.json passed. This is not an all-curve bound. Review provenance: New derivation after reading the supplied proof; not a blind review and not an independent human check}{the hypotheses, notation and definitions stated here; \cref{cdp:st:A.5.5}}

Put
\[
 \begin{aligned}
 I&=(y^2-xz,\ yz^2-x^4,\ z^3-x^3y)\subset\C\{x,y,z\},\\
 F&=4x^7-15x^3yz^2+10x^2y^3z-3xy^5+4z^5.
 \end{aligned}                                                   \tag{A.NL.1}\label{cdp:eq:A.NL.1}
\]
The reduced critical locus of the reduced hypersurface $F=0$ is the
irreducible curve $D$ parametrized by $(x,y,z)=(t^5,t^6,t^7)$.
Its transverse type at every point with $t\ne0$ is the ordinary cusp
$A_2$. The curve is not a local complete intersection at the origin,
and
\[
             F\in I^{(2)}\setminus I^2.                         \tag{A.NL.2}\label{cdp:eq:A.NL.2}
\]
Thus the ordinary-square membership used to define the conormal
quadratic jet in \ref{cdp:st:A.5.5} does not extend to an arbitrary reduced
critical curve, even with generic transverse type $A_2$.
This is a local example, without an asserted compact CDP realization.
\end{proposition}


\begin{proof}
\noindent\textbf{Step 1.} Write the three displayed generators of $I$ as $g_1,g_2,g_3$.
The quotient is a finite $\C\{x\}$-module generated by
$1,y,z,yz,z^2$: their products by $y$ and $z$ reduce to this list
using $y^2=xz$, $yz^2=x^4$, and $z^3=x^3y$.
For example these identities also give $y^5=x^6$ and $z^5=x^7$,
which make the analytic quotient finite over $\C\{x\}$.
The parametrization sends the five generators to
$1,t^6,t^7,t^{13},t^{14}$. Their exponents have distinct residues
modulo five, so they are linearly independent over $\C\{t^5\}$.
Consequently this actual ring homomorphism is injective and has
image $\C\{t^5,t^6,t^7\}$; it proves the stated normalization and
the primeness of $I$. The three generators are minimal. Indeed their
initial terms have degrees two, three, three; after the degree-two
coefficient vanishes, the cubic terms $yz^2$ and $z^3$ are independent
modulo the span of $xg_1,yg_1,zg_1$. A height-two ideal with three
minimal generators is not a complete intersection.

\noindent\textbf{Step 2.} The following polynomial identities determine the critical locus:
\[
 \begin{aligned}
 F_y&=-15xg_1^2,\\
 F_x&=(17xyz-3y^3)g_1-28x^2g_2,\\
 F_z&=10x^2yg_1+20zg_3.
 \end{aligned}                                                   \tag{A.NL.3}\label{cdp:eq:A.NL.3}
\]
At $x=0$, the equations $F_x=F_z=0$ give $y=z=0$.
At $x\ne0$, $F_y=0$ gives $g_1=0$, and then $F_x=0$
gives $g_2=0$. The value $z=0$ is impossible there, so $F_z=0$
also gives $g_3=0$. Conversely all three partials vanish on $D$.
Weighted Euler's identity, with weights $(5,6,7)$ and degree $35$,
puts $F$ in its Jacobian ideal. The critical locus is therefore
exactly $D$ as a reduced set. A repeated factor of $F$ would put
a two-dimensional component in this critical locus, which proves
that $F$ is reduced.

\noindent\textbf{Step 3.} The Hessian along the normalization has the exact factorization
\[
 \operatorname{Hess}(F)|_D
 =2t^{21}
 \begin{pmatrix}7t^2\\0\\-5\end{pmatrix}
 \begin{pmatrix}7t^2&0&-5\end{pmatrix},
 \qquad \frac{1}{6}F_{yyy}|_D=-20t^{17}.                       \tag{A.NL.4}\label{cdp:eq:A.NL.4}
\]
For $t\ne0$, the plane $x=t^5$ is transverse to $D$.
The restricted Hessian has rank one and its radical is the
$y$ direction; its cubic coefficient in that direction is nonzero.
The holomorphic splitting lemma and a local cube root of a unit
give the transverse cusp equation. At the generic point of $D$,
its conormal is locally free and injects into ambient differentials.
Since $F$ and $dF$ vanish there, $F\in I^2\C\{x,y,z\}_I$,
which is precisely $F\in I^{(2)}$.

\noindent\textbf{Step 4.} To decide ordinary-square membership, use the positive weights
$(5,6,7)$. The three generators have weights $12,20,21$.
Among products of two generators, only $g_1^2$ can contribute
to weight $35$: it needs a coefficient of weight $11$, whose
space is spanned by $xy$. Products $g_1g_2$ and $g_1g_3$ would
need weights $3$ and $2$, which do not occur; the other products
already have greater weight. Thus the weight-$35$ part of $I^2$
is exactly $\C\,xyg_1^2$. Its $x^7$ coefficient is zero, whereas
that of $F$ is four. Taking the weight-$35$ component of a
convergent series identity is legitimate because all weights are
positive. This proves $F\notin I^2$ in the analytic local ring.

\noindent\textbf{Step 5.} For comparison of the curve invariants, the normalization conductor
is $t^{10}\C\{t\}$: the semigroup gaps are $1,2,3,4,8,9$.
The $B$-module generated by the differential image, where $B=\C\{t\}$, is $t^4B\,dt$.
The smallest valuation of a two-by-two minor of the displayed
generators' Jacobian matrix is $19$. Thus the normalized conormal
determinant correction is $19-4=15$, while the conductor exponent
is $10$; the complete-intersection adjunction comparison cannot
identify these two quantities here.

\noindent\textbf{Step 6.} The symbolic-square correction is visible directly on the free
normalized conormal. Set $a=dg_1|_D$ and
$b=dg_2|_D+2t^8a$. Then
\[
 \begin{aligned}
 a&=(-t^7,2t^6,-t^5),& b&=(-6t^{15},5t^{14},0),\\
 dg_3|_D&=-3t^9a+tb,&
 q&=t^{-5}(5t^8a-2b)^2.
 \end{aligned}                                                   \tag{A.NL.5}\label{cdp:eq:A.NL.5}
\]
The first two rows generate the entire normalized conormal image,
and $q$ is the quadratic term, namely half the Hessian.
If $A,B$ are the dual basis, its radical frame $2A+5t^8B$
is represented in the ambient normal quotient by $t^{-6}\partial_y$.
Its intrinsic cubic is therefore $-20t^{-1}$. These are actual
quadratic and cubic poles of orders five and one. The corresponding
combined invariant $q^3k^2$ has pole order $17$; no bound for such
poles on all non-lci curves is deduced from this one computation.
All substitutions, polynomial identities, weighted coefficient
spaces, conormal comparisons, and Jacobian minors are checked
exactly in \path{code/A_non_lci_cusp.py}.
\end{proof}




\begin{proposition}[A complete nodal-family computation]\label{cdp:st:B1}
\StatementRecord{cdp:st:B1}{proved}{Historical events: Claude Fable 5.2 (Anthropic), effort high, 16 November 2026: carried as background without fault found. GPT-6 Astra (OpenAI), ChatGPT Pro mode, effort Pro, S6\_09\_PAPERS, 17 November 2026: targeted proof check. Also checked in this writing. Pre-continuation writing events: GPT-6 Astra (OpenAI), ChatGPT Pro mode, effort Pro, S6\_10\_PAPERS, 18 November 2026: targeted proof check. Continuation metadata events: GPT-6 Astra (OpenAI), ChatGPT Pro mode, effort Pro, S6\_10\_PAPERS continuation, 19 November 2026: used and rechecked in this run; status-vocabulary reconciliation; metadata reconciliation only; no new mathematical check. \emph{Pre-continuation recorded proof scope.} The nodal normalization u maps to (u\textasciicircum{}2-1,u(u\textasciicircum{}2-1)); the descent scalar, torsion length, inverse-image module, product cohomology, and extension of the gluing line were re-derived}{the hypotheses, notation and definitions stated here; the explicit constructions and calculations in the proof}
\label{wbII:thm:nodal-test}\label{wbII:prop:family-geometry}\label{wbII:prop:nodal-counterexample}
Let $0<\varepsilon<4/27$, let $\Delta=\{|t|<\varepsilon\}$, and define
\[
 \mathcal C=\{y^2z=x^2(x+z)+tz^3\}\subset\PP^2\times\Delta,
 \qquad X=\mathcal C\times\PP^1.
\]
The projection $X\to\Delta$ is proper, its total space is smooth, and its central fibre is $S=C_0\times\PP^1$, where $C_0$ is an irreducible nodal cubic. For each $\lambda\in\C^*$ let $A_\lambda$ be the line bundle whose pullback to the normalization $\nu:\PP^1_u\to C_0$ is trivial and whose sections in this trivialization satisfy $h(1)=\lambda h(-1)$. Put $A=\operatorname{pr}_1^*A_\lambda$ and $d_\lambda=1$ for $\lambda=1$, and $d_\lambda=0$ otherwise. In the order $q=0,1,2$, the coherent cohomology dimensions are
\[
\begin{array}{c|c}
\text{sheaf on }S & (h^0,h^1,h^2)\\\hline
 N^*_{S/X}\otimes A &(d_\lambda,d_\lambda,0)\\
 K_S\otimes A &(1,0,0)\\
 T_S\otimes A &(1,0,0)\\
 \widetilde\Omega^1_S\otimes A &(0,1+d_\lambda,d_\lambda)\\
 \Omega^1_S\otimes A &(1,1+d_\lambda,d_\lambda)\\
 \Omega^1_X|_S\otimes A &(1,1+d_\lambda,d_\lambda).
\end{array}                                                    \tag{B.1}\label{cdp:eq:B.1}
\]
Here $T_S$ denotes the torsion of the cotangent sheaf, and $K_S$ its inverse image in the ambient cotangent sheaf. In particular the two proposed vanishings hold exactly for $\lambda\ne1$, and the ambient space of sections has dimension one for every $\lambda\in\C^*$. After shrinking $\Delta$, every $A_\lambda$ extends to a line bundle on $\mathcal C$ of the form $\OO_{\mathcal C}(P-Q)$ for holomorphic sections $P,Q$.
\end{proposition}



\begin{proof}
\noindent\textbf{Step 1.} In the affine chart $z=1$, the defining equation is $F=y^2-x^3-x^2-t=0$. Its partial derivative in $t$ is $-1$, so the total space is smooth there. Its only point with $z=0$ is $[0:1:0]$. In the chart $y=1$ the equation is $z-x^3-x^2z-tz^3=0$, whose $z$ derivative is one at this point. Fibre singularities in the affine chart satisfy $y=0$ and $x(3x+2)=0$; they occur at $t=0,-4/27$. Our disc contains only the first. The quadratic term $y^2-x^2$ at the origin has two distinct linear factors, and its smoothing parameter has nonzero differential. The holomorphic Morse lemma with parameter therefore supplies local coordinates with $t=vw$. Properness follows from the closed inclusion into $\PP^2\times\Delta$.

\noindent\textbf{Step 2.} The normalization is
\[
 u\longmapsto(x,y)=(u^2-1,u(u^2-1)),
\]
completed by the point at infinity. The two points above the node are $1,-1$. The descent sequence, an exact sequence of coherent $\OO_{C_0}$-modules, is
\[
0\longrightarrow A_\lambda\longrightarrow\nu_*\OO_{\PP^1}
 \xrightarrow{\ h\mapsto h(1)-\lambda h(-1)\ }\C_{\rm node}
 \longrightarrow0.                                      \tag{B.2}\label{cdp:eq:B.2}
\]
Its cohomology is the two-term complex $\C\xrightarrow{1-\lambda}\C$. Thus $h^0(A_\lambda)=h^1(A_\lambda)=d_\lambda$.

\noindent\textbf{Step 3.} Write $V=\Omega^1_{\mathcal C}|_{C_0}$. The normalization map induces a surjective $\OO_{C_0}$-linear morphism $\Omega^1_{C_0}\to\nu_*\Omega^1_{\PP^1}$; at a node $vw=0$, it sends $dv$ to $(dv,0)$ and $dw$ to $(0,dw)$. Its kernel is the length-one module generated by $v\,dw=-w\,dv$. Consequently
\[
 T_{C_0}=\C_{\rm node},\qquad
 \widetilde\Omega^1_{C_0}=\nu_*\OO_{\PP^1}(-2).
\]
The kernel of $V\to\nu_*\Omega^1_{\PP^1}$ is $\nu_*\OO_{\PP^1}$. Locally its generators are $w\,dv$ and $v\,dw$ on the two branches; the identification is the normal-crossing conormal identification of Part A with the globally fixed conormal $dt$. Tensoring by $A_\lambda$ preserves exactness, and finite pushforward and the projection formula identify the two resulting normalization terms with $\nu_*\OO$ and $\nu_*\OO(-2)$. Hence
\[
 (h^0,h^1)(C_0,V\otimes A_\lambda)=(1,1),\qquad
 (h^0,h^1)(C_0,\Omega^1_{C_0}\otimes A_\lambda)=(1,1).
\]
The product morphisms give the direct decompositions
\[
 \Omega^1_X|_S\otimes A
 =(V\otimes A_\lambda)\boxtimes\OO
 \ \oplus\ A_\lambda\boxtimes\OO(-2),
\]
and the analogous decomposition for $\Omega^1_S$. The torsion-free quotient is obtained by replacing the first factor by $\nu_*\OO(-2)$. The torsion is $\C_{\rm node}\boxtimes\OO$, its ambient inverse image is $\nu_*\OO\boxtimes\OO$, and $N^*_{S/X}=\OO_S$. The coherent Kunneth isomorphism and $H^0(\PP^1,\OO)=\C$, $H^1(\PP^1,\OO(-2))=\C$ now give every entry of \eqref{cdp:eq:B.1}. All other cohomology of these two line bundles on $\PP^1$ vanishes.

\noindent\textbf{Step 4.} We give the extension explicitly. Fix $Q(0)$ at $u=0$, and take $P(0)$ at $u=a=(1-\lambda)/(1+\lambda)$; for $\lambda=-1$ take $a=\infty$. For $\lambda\ne0$ both points are in the smooth locus. With $b=0$, the rational function $r(u)=(u-a)/(u-b)$ has divisor $P(0)-Q(0)$, and a local section in its pullback trivialization descends precisely with scalar
\[
 \frac{r(1)}{r(-1)}
 =\frac{(1-a)(1+b)}{(1-b)(1+a)}=\lambda.
\]
For $a=\infty$ the same expression is its limit, namely $-1$. The projection $\mathcal C\to\Delta$ is a submersion at these points, so the holomorphic implicit-function theorem extends them to sections on a smaller disc. The point at infinity itself is a constant section. For $\lambda=1$ take $P=Q$ identically. Sections through the smooth fibre locus are effective Cartier divisors in the smooth surface $\mathcal C$; their difference therefore defines the claimed extension. Pullback along $X\to\mathcal C$ extends $A$. This proves the claim also for all roots of unity, including the identity.
\end{proof}



The calculation recovers the local example in the supplied correction note, Lemma 4.3 and Theorem 4.4, with the gluing convention fixed by \eqref{cdp:eq:B.2}. The computation applies to roots of unity as well: the only exceptional scalar for the conormal vanishing is $1$. The symbolic computation and the exact gluing matrices are recorded in \texttt{code/B\_examples.py} and its output.




\begin{proposition}[A projective example with twelve length-one contributions]\label{cdp:st:B2}
\StatementRecord{cdp:st:B2}{proved}{Historical events: Claude Fable 5.2 (Anthropic), effort high, 16 November 2026: carried as background without fault found. GPT-6 Astra (OpenAI), ChatGPT Pro mode, effort Pro, S6\_09\_PAPERS, 17 November 2026: targeted proof check. Also checked in this writing. Pre-continuation writing events: GPT-6 Astra (OpenAI), ChatGPT Pro mode, effort Pro, S6\_10\_PAPERS, 18 November 2026 04:23:40Z: targeted proof check. Continuation metadata events: GPT-6 Astra (OpenAI), ChatGPT Pro mode, effort Pro, S6\_10\_PAPERS continuation, 19 November 2026: used and rechecked in this run; status-vocabulary reconciliation; metadata reconciliation only; no new mathematical check. \emph{Pre-continuation recorded proof scope.} proved as stated. Full projective twelve-node example, length-one coherent direct images and all fibre/global cohomology dimensions. Re-derived the Weierstrass discriminant and smoothness charts, the degree-zero nontriviality of P−O on each integral genus-one fibre, and the acyclicity of its two signs. The local map dv\ensuremath{\to}v·omega, dw\ensuremath{\to}−w·omega has kernel dt and image the reduced node ideal; derived pushforward gives exactly one skyscraper in degree1 at each node. Contraction T\_E\ensuremath{\otimes}K\_E\ensuremath{\to}Omega\_E and product Kunneth shift these to R2g, with all other degrees zero. The complete B1 normalization calculation was re-read for its calling role; Cartier Serre duality gives fibre dimensions (0,1,1), and the degree1 fibre term is precisely Tor1 of the degree2 skyscraper. Review provenance: New derivation after reading the supplied proof; not a blind review and not an independent human check}{the hypotheses, notation and definitions stated here; \cref{cdp:eq:B.6}; \cref{cdp:eq:B.5}; \cref{cdp:eq:B.7}; \cref{cdp:st:B1}}

Let $h:E\to\PP^1_t$ be the Weierstrass surface defined on the finite chart by
\[
 y^2z=x^3+(t^4+1)xz^2+t^6z^3,                          \tag{B.3}\label{cdp:eq:B.3}
\]
and at $s=1/t$ by $x_\infty=s^2x$, $y_\infty=s^3y$, $z_\infty=z$. Let $O=[0:1:0]$ and $P=(0,t^3,1)$ be its two sections and put
\[
 M_E=\OO_E(P-O),\quad Y=E\times\PP^1,\quad
 g=h\circ\operatorname{pr}_E,\quad
 L=K_Y\otimes\operatorname{pr}_E^*M_E^{-1},\quad
 \mathcal E=T_Y\otimes L.
\]
The surface $E$ is smooth and projective. Its singular fibres are twelve irreducible nodal cubics, at the distinct roots $t_i$ of
\[
 D(t)=31t^{12}+12t^8+12t^4+4.
\]
For every such fibre $S_i=g^{-1}(t_i)$, with $A=(K_Y\otimes L^{-1})|_{S_i}$, we have
\[
H^0(S_i,N^*_{S_i/Y}\otimes A)=0,
\qquad H^0(S_i,\widetilde\Omega^1_{S_i}\otimes A)=0,
\]
\[
 (h^0,h^1,h^2)(S_i,\mathcal E|_{S_i})=(0,1,1).
\]
As coherent sheaves on the base, with noncanonical identifications of their one-dimensional fibres,
\[
 R^qg_*\mathcal E=0\ (q\ne2),\qquad
 R^2g_*\mathcal E=\bigoplus_{i=1}^{12}\C_{t_i}.        \tag{B.4}\label{cdp:eq:B.4}
\]
In particular $(h^0,h^1,h^2,h^3)(Y,\mathcal E)=(0,0,12,0)$.
\end{proposition}



\begin{proof}
\noindent\textbf{Step 1.} The transition rule realizes \eqref{cdp:eq:B.3} as a cubic in a projective plane bundle over $\PP^1$, with Weierstrass coefficients in $\OO(4)$ and $\OO(6)$; hence its total space is projective. Its discriminant is $-16D(t)$. Exact polynomial division gives
\[
 \gcd(D,D')=1,\qquad \gcd(D,t^4+1)=1.
\]
The leading coefficient $31$ is its nonzero value in the chart at infinity. Each discriminant root therefore has nonzero $c_4=-48(t^4+1)$ and is a nodal cubic. A simple discriminant zero is transverse to the nodal discriminant and makes the surface smooth at that node. Alternatively the Groebner basis of the affine total-space singularity ideal $(F,F_x,F_y,F_t)$ is $(1)$; the exact computation is attached. At the zero section smoothness follows from $F_z=1$ in the chart $y=1$. At infinity the coefficients are $a=1+s^4,b=1$, with discriminant $-16\cdot31$ at $s=0$. These checks cover the surface.

\noindent\textbf{Step 2.} The sections $P,O$ are disjoint and meet the smooth locus of every fibre. For $P$, a singular point would require both $a(t)=0$ and $t=0$, which is impossible. Each fibre is an integral arithmetic-genus-one curve with trivial dualizing sheaf. The line bundle $M_E$ has degree zero on every fibre and is nontrivial there. On a smooth fibre, a trivialization would give a meromorphic function with one simple pole and one simple zero, hence a degree-one map to $\PP^1$, contradicting genus one. On a nodal fibre the normalization calculation in (B.2) says that two distinct smooth points determine a nontrivial degree-zero line bundle. Thus $H^q(E_t,M_E^{\pm1}|_{E_t})=0$ for both $q=0,1$: a nonzero section of a degree-zero line bundle on an integral curve is a trivialization, and Serre duality gives the assertion in degree one. Grauert's theorem, applied to the proper flat map $h$, yields
\[
 Rh_*M_E=Rh_*M_E^{-1}=0.                               \tag{B.5}\label{cdp:eq:B.5}
\]

\noindent\textbf{Step 3.} We compute the cotangent direct image, including its scheme length. Let $Z\subset E$ be the reduced set of the twelve critical points. The relative dualizing sheaf is $\omega_{E/\PP^1}=h^*\OO(1)$. Indeed the residue frame $dx/(2y)$ has this transition under $x_\infty=s^2x,y_\infty=s^3y$, and extends as a dualizing frame at a node. It follows that $K_E=h^*\OO(-1)$. There is an exact sequence of coherent sheaves on $E$
\[
0\longrightarrow h^*\Omega^1_{\PP^1}
 \xrightarrow{dh}\Omega^1_E
 \longrightarrow\mathcal I_Z\otimes\omega_{E/\PP^1}
 \longrightarrow0.                                   \tag{B.6}\label{cdp:eq:B.6}
\]
Away from $Z$ this is the ordinary relative cotangent sequence. At a node choose coordinates $vw=t$. The map from $\Omega^1_E$ to the relative dualizing line sends $dv\mapsto v\omega$, $dw\mapsto-w\omega$; its image is $(v,w)\omega$ and its kernel is generated by $w\,dv+v\,dw=dt$. This proves \eqref{cdp:eq:B.6} on every stalk, with the critical scheme reduced.

\noindent\textbf{Step 4.} Tensor \eqref{cdp:eq:B.6} by $M_E^{-1}$. By \eqref{cdp:eq:B.5}, projection formula, and the sequence for the ideal of $Z$, its derived direct image has only degree one, namely
\[
 R^1h_*(\Omega^1_E\otimes M_E^{-1})
 =h_*(\omega_{E/\PP^1}\otimes M_E^{-1}|_Z),
 \qquad h_*(\Omega^1_E\otimes M_E^{-1})=0.             \tag{B.7}\label{cdp:eq:B.7}
\]
The right side is the sum of twelve reduced skyscrapers. In particular \eqref{cdp:eq:B.7} calculates the length rather than merely the residue dimension at each point.

\noindent\textbf{Step 5.} Contraction on the smooth surface gives the canonical isomorphism $T_E\otimes K_E\simeq\Omega^1_E$. Consequently the product tangent decomposition identifies
\[
\mathcal E=(\Omega^1_E\otimes M_E^{-1})\boxtimes\OO(-2)
 \ \oplus\ (K_E\otimes M_E^{-1})\boxtimes\OO.
\]
The second summand has zero derived direct image by \eqref{cdp:eq:B.5}. The first has only the tensor product of degree-one terms: \eqref{cdp:eq:B.7} and $H^1(\PP^1,\OO(-2))=\C$. The relative coherent Kunneth comparison gives \eqref{cdp:eq:B.4}. The Leray spectral sequence over a curve then gives the four stated global dimensions.

\noindent\textbf{Step 6.} On a singular fibre $S_i=E_{t_i}\times\PP^1$, the restriction of $M_E$ is a nontrivial $A_\lambda$. Statement \ref{cdp:st:B1} gives the two vanishings and ambient dimensions $(1,1,0)$. Adjunction for the Cartier fibre gives $\omega_{S_i}=K_Y|_{S_i}$, because its normal line is trivial. Serre duality for this compact Gorenstein surface identifies, for $q=0,1,2$,
\[
 H^q(S_i,\mathcal E|_{S_i})^\vee
 =H^{2-q}(S_i,\Omega_Y^1|_{S_i}\otimes A).
\]
This proves the fibre cohomology. It also explains the fibre value $h^1=1$ while $R^1g_*\mathcal E=0$: tensoring a reduced skyscraper in degree two with the residue field has a degree-one Tor term.
\end{proof}






\begin{proposition}[The hypotheses satisfied by the compact example]\label{cdp:st:B3}
\StatementRecord{cdp:st:B3}{proved}{Historical events: Claude Fable 5.2 (Anthropic), effort high, 16 November 2026: carried as background without fault found. GPT-6 Astra (OpenAI), ChatGPT Pro mode, effort Pro, S6\_09\_PAPERS, 17 November 2026: targeted proof check. Also checked in this writing. Pre-continuation writing events: GPT-6 Astra (OpenAI), ChatGPT Pro mode, effort Pro, S6\_10\_PAPERS, 18 November 2026 04:23:40Z: targeted proof check. Continuation metadata events: GPT-6 Astra (OpenAI), ChatGPT Pro mode, effort Pro, S6\_10\_PAPERS continuation, 19 November 2026: used and rechecked in this run; status-vocabulary reconciliation; metadata reconciliation only; no new mathematical check. \emph{Pre-continuation recorded proof scope.} proved as stated. Full topological and algebraic scope of the projective twelve-node example. The exact residue transition gives K\_E=h*O(−1); relative duality and Leray give q=pg=P2=0. Castelnuovo makes E rational. The blowup decomposition of rational surfaces preserves H1=0 and freeness of H2; Euler additivity over the twelve nodes gives e(E)=12 and b2(E)=10. Integral Kunneth then gives H2(Y)=Z11 and e(Y)=24. Projectivity yields algebraic dimension3. Peters1990 current opening includes the precise blowup topology pp8–9 and the stated Castelnuovo criterion p27; the input is credited as classical. Review provenance: New derivation after reading the supplied proof; not a blind review and not an independent human check}{the hypotheses, notation and definitions stated here; \cref{cdp:st:B2}; \cite{CDP20}}

The threefold $Y$ of Statement \ref{cdp:st:B2} has $H^1(Y,\Z)=0$, $H^2(Y,\Z)\simeq\Z^{11}$, algebraic dimension $3$, and $e(Y)=24$. Its map $g$ is holomorphic, proper, surjective, and has connected fibres. Thus it tests the fibrewise inference in Section 7 of \cite{CDP20}; it does not satisfy the vanishing of $H^2(Y,\Z)$ or algebraic dimension one required by that paper's Theorem 2.2.
\end{proposition}


\begin{proof}
\noindent\textbf{Step 1.} The residue calculation above gives $K_E=h^*\OO(-1)$. Relative duality gives $R^1h_*\OO_E=\OO(-1)$, since every fibre is an integral genus-one curve, and $h_*\OO_E=\OO$. Thus $q(E)=p_g(E)=0$ and $H^0(E,K_E^{\otimes m})=0$ for every positive $m$. Castelnuovo's rationality criterion, in the form $q=P_2=0$ for a smooth projective surface, makes $E$ rational. A rational smooth projective surface is obtained from a minimal rational surface by point blowups; its first integral homology vanishes and its second integral cohomology is free. Euler additivity gives $e(E)=12$, since the twelve nodal fibres each have Euler characteristic one and the smooth elliptic fibre has Euler characteristic zero. Poincare duality gives $b_2(E)=10$. The integral Kunneth isomorphism for $Y=E\times\PP^1$ gives the asserted groups and Euler number. Projectivity and dimension give $a(Y)=3$. Connectedness and surjectivity follow directly from the integral cubic fibres of $h$ and the second product factor.
\end{proof}



The construction in Statement \ref{cdp:st:B2} supplies the compact local counterexample, including its bundles, cohomology groups and direct-image lengths. The universal assertion under every hypothesis of CDP20 Theorem 2.2 is proved in \ref{cdp:st:C4e} below: positive third Chern number forces one fixed irreducible nonnormal reduced support whose normalization has the required nonzero group for every ambient Picard twist. Its proof uses the full-fibre comparisons of A.3 and the conductor arguments of A.5. The named-family tests are \ref{cdp:st:C4g}.




\begin{proposition}[All thickenings and their lifting maps]\label{cdp:st:A.3.1}
\StatementRecord{cdp:st:A.3.1}{proved}{Historical events: Claude Fable 5.2 (Anthropic), effort high, 16 November 2026: carried as background without fault found. GPT-6 Astra (OpenAI), ChatGPT Pro mode, effort Pro, S6\_09\_PAPERS, 17 November 2026: targeted proof check. Also checked in this writing. Pre-continuation writing events: GPT-6 Astra (OpenAI), ChatGPT Pro mode, effort Pro, S6\_10\_PAPERS, 18 November 2026 04:12:55Z: targeted proof check. Continuation mathematical and source checks: GPT-6 Astra (OpenAI), ChatGPT Pro mode, effort Pro, S6\_10\_PAPERS continuation, 19 November 2026: computation or formal certificate; independent re-proof; used and rechecked in this run. Continuation metadata events: GPT-6 Astra (OpenAI), ChatGPT Pro mode, effort Pro, S6\_10\_PAPERS continuation, 19 November 2026: used and rechecked in this run; status-vocabulary reconciliation; metadata reconciliation only; no new mathematical check. \emph{Continuation recorded proof scope.} The whole current proof was independently derived before target reading at 09:16:58 UTC: actual Cartier ideal-power multiplication and transition maps, locally free tensor exactness, connecting-map cocycles, lifting fibres and induction for all thickenings. Fresh exact module and transition checks form part of the 920-check program; no universal conclusion is inferred from finite samples. Original Stacks Tags 0B3P and 0C4T, Divisors p.~26, and Tag 01DZ, Cohomology pp.~2--3, were opened and matched through the stated analytic local maps. These are classical inputs; the supporting-row source gate is complete.}{\cite[Tags 0B3P and 0C4T, Divisors chapter p.~26; Tag 01DZ, Cohomology chapter pp.~2--3]{StacksProject}; the explicit local multiplication maps for the analytic Cartier ideal}

Let $S$ be any effective Cartier divisor in a complex manifold, let $E$ be any locally free sheaf on that manifold, and let $E_k=E|_{kS}$. For each $k\ge2$ there is a canonical exact sequence
\[
0\longrightarrow (E|_S)\otimes(N^*)^{k-1}
 \longrightarrow E_k\longrightarrow E_{k-1}\longrightarrow0.
\tag{A.10}\label{cdp:eq:A.10}
\]
Write $G_{k-1}=(E|_S)\otimes(N^*)^{k-1}$ and let
$\epsilon_k:H^0((k-1)S,E_{k-1})\to H^1(S,G_{k-1})$
be the connecting map. Then
\[
0\longrightarrow H^0(S,G_{k-1})\longrightarrow H^0(kS,E_k)
 \longrightarrow\ker\epsilon_k\longrightarrow0.
\tag{A.11}\label{cdp:eq:A.11}
\]
In particular $H^0(kS,E_k)=0$ exactly when $H^0(S,G_{k-1})=0$ and $\epsilon_k$ is injective. Starting with $H^0(S,E|_S)$, equations \eqref{cdp:eq:A.10}--\eqref{cdp:eq:A.11} compute the group on $mS$ through all its extension maps. The vanishing of every $H^0(S,G_j)$, $0\le j<m$, is a sufficient criterion. These statements apply to $E=\Omega_Y^1\otimes B$ for every line bundle $B$, and include smooth, normal-crossings, and all other reduced supports.
\end{proposition}


\begin{proof}
\noindent\textbf{Step 1.} With $I=\OO_Y(-S)$, the sequence is the tensor product of $E$ with
$0\to I^{k-1}/I^k\to\OO_Y/I^k\to\OO_Y/I^{k-1}\to0$.
The Cartier property identifies $I^{k-1}/I^k$ with $(N^*)^{k-1}$ on $S$; local freeness makes the tensor product exact. Taking global sections proves \eqref{cdp:eq:A.11} and its equivalence. Induction on $k$ proves the sufficient criterion. No splitting of the thickening, cyclic base change, or assumption about the normal bundle has been used.
\end{proof}






\begin{proposition}[Counterexample: Vanishing on the support alone does not imply vanishing on a multiple fibre]\label{cdp:st:A.3.2}
\StatementRecord{cdp:st:A.3.2}{proved}{Historical events: Claude Fable 5.2 (Anthropic), effort high, 16 November 2026: carried as background without fault found. GPT-6 Astra (OpenAI), ChatGPT Pro mode, effort Pro, S6\_09\_PAPERS, 17 November 2026: targeted proof check. Also checked in this writing. Pre-continuation writing events: GPT-6 Astra (OpenAI), ChatGPT Pro mode, effort Pro, S6\_10\_PAPERS, 18 November 2026 04:12:55Z: targeted proof check. Continuation metadata events: GPT-6 Astra (OpenAI), ChatGPT Pro mode, effort Pro, S6\_10\_PAPERS continuation, 19 November 2026: used and rechecked in this run; status-vocabulary reconciliation; metadata reconciliation only; no new mathematical check. \emph{Pre-continuation recorded proof scope.} proved as stated. Complete compact quotient counterexample, finite-thickening basis and all additional cohomology dimensions. The diagonal translation/rotation action is free and its quotient map étale. The character line has exact order m by constancy of holomorphic functions on connected F. The tangent and normal factors have weights 1,1,2; on C[z]/z\textasciicircum{}m the invariant monomials are exactly the three listed. Averaging commutes with cohomology. Globally the two O(−1) tangent summands are acyclic; \ensuremath{\Omega}P1(−0) has Cech classes dz/z and dz of weights 0 and 1, hence exactly one invariant class. Kunneth gives (0,1,2,1). Review provenance: New derivation after reading the supplied proof; not a blind review and not an independent human check}{the hypotheses, notation and definitions stated here; the explicit constructions and calculations in the proof}

There is a compact smooth complex threefold $g:Y\to\PP^1$ with connected general fibre an abelian surface and a smooth fibre support $S$ of multiplicity $m$, for every $m\ge3$, and a line bundle $B$ on $Y$, such that
\[
H^0(S,(\Omega_Y^1\otimes B)|_S)=0,
\qquad h^0(mS,(\Omega_Y^1\otimes B)|_{mS})=3.
\tag{A.12}\label{cdp:eq:A.12}
\]
Its conormal has exact order $m$. The replacement is (A.10)--(A.11), with the graded twists retained.
\end{proposition}


\begin{proof}
\noindent\textbf{Step 1.} Take an abelian surface $F$ and a translation $a:F\to F$ of exact order $m$. Let $\zeta$ be a primitive $m$th root of unity. The action generated by $(p,z)\mapsto(a(p),\zeta z)$ on $F\times\PP^1$ is free, so the quotient $Y$ is a smooth compact threefold and the quotient map is finite etale. The invariant map $t=z^m$ gives $g:Y\to\PP^1$. Away from $0$ and $\infty$ its fibre is isomorphic to $F$: the $m$ lifts of a base point are permuted transitively and freely. At $0$ the fibre is $mS$, where $S=F/\langle a\rangle$. Set $B=\OO_Y(-S)$.

\noindent\textbf{Step 2.} The connected cyclic cover $F\to S$ identifies $N^*=B|_S$ with the character line of weight one. Its order is exactly $m$: otherwise a nontrivial proper character power would have a nonzero invariant trivialization on the connected compact cover, whose holomorphic functions are constant. Pullback to $F$ splits the ambient cotangent bundle into two invariant tangent forms and the normal form $dz$ of weight one. Consequently
\[
 (\Omega_Y^1\otimes B)|_S\simeq (N^*)^{\oplus2}\oplus(N^*)^2.
\]
The nontrivial flat lines on the abelian surface $S$ have no holomorphic sections. For completeness, their pullbacks to $F$ are trivial; a section would be a constant function transforming by a nontrivial character, which is zero. Since $m\ge3$, both characters in the preceding splitting are nontrivial, proving the first equality.

\noindent\textbf{Step 3.} The inverse image of $mS$ is the ringed space with structure sheaf $\OO_F\otimes\C[z]/(z^m)$. Write $e$ for the local generator of the pulled-back line $B$; it has weight one. All global sections on this inverse image are linear combinations of $z^k\alpha_1e,z^k\alpha_2e,z^kdz\,e$ for $0\le k<m$, where $\alpha_1,\alpha_2$ are a basis of invariant one-forms on $F$. The first two have weight $k+1$ and the third weight $k+2$. The invariant monomials are exactly
$z^{m-1}\alpha_1e,z^{m-1}\alpha_2e,z^{m-2}dz\,e$.
Taking invariants is exact over $\C$, by averaging over the finite group. These three monomials therefore give the claimed basis and prove the second equality. Translations act identically on $H^q(F,\OO_F)$, whose dimensions are $\binom{2}{q}$. The same calculation gives all fibre groups: $h^q(S,E|_S)=(0,0,0)$ and $h^q(mS,E|_{mS})=(3,6,3)$ for $q=0,1,2$. For completeness, on the compact threefold itself the two tangent summands pull back to $\OO_F\boxtimes\OO_{\PP^1}(-1)$ and have no cohomology. The normal summand is $\OO_F\boxtimes\Omega^1_{\PP^1}(-\{0\})$. Its first cohomology on $\PP^1$ has the Cech basis $dz/z,dz$, of weights zero and one, respectively. Indeed sections near zero are multiples of $z\,dz$, and sections near infinity are multiples of $z^{-2}dz$, leaving exactly those two Laurent classes. Its invariant part is one-dimensional. Kunneth and finite-cover descent therefore give $h^q(Y,E)=(0,1,2,1)$ for $q=0,1,2,3$.
\end{proof}



For this example the complete dimensions on $kS$, $1\le k\le m$, are $0$ for $k\le m-2$, $1$ for $k=m-1$, and $3$ for $k=m$; the same character calculation proves this. More generally the character count with $B=\OO_Y(-bS)$ is
\[
\begin{aligned}
h^0(kS,(\Omega_Y^1\otimes B)|_{kS})
&=2\#\{0\le j<k:j+b\equiv0\pmod m\}\\
&\quad+\#\{0\le j<k:j+b+1\equiv0\pmod m\}.
\end{aligned}
\tag{A.13}\label{cdp:eq:A.13}
\]
The finite counts for $3\le m\le12$ and all characters are checked in \path{code/A_multiple_fibres.py}.




\begin{proposition}[The precise finite-translation step in Remarks 7.2 and 7.4]\label{cdp:st:A.3.3}
\StatementRecord{cdp:st:A.3.3}{proved}{Historical events: Claude Fable 5.2 (Anthropic), effort high, 16 November 2026: carried as background without fault found. GPT-6 Astra (OpenAI), ChatGPT Pro mode, effort Pro, S6\_09\_PAPERS, 17 November 2026: targeted proof check. Also checked in this writing. Pre-continuation writing events: GPT-6 Astra (OpenAI), ChatGPT Pro mode, effort Pro, S6\_10\_PAPERS, 18 November 2026 04:12:55Z: targeted proof check. Continuation mathematical and source checks: GPT-6 Astra (OpenAI), ChatGPT Pro mode, effort Pro, S6\_10\_PAPERS continuation, 19 November 2026: computation or formal certificate; independent re-proof; used and rechecked in this run. Continuation metadata events: GPT-6 Astra (OpenAI), ChatGPT Pro mode, effort Pro, S6\_10\_PAPERS continuation, 19 November 2026: used and rechecked in this run; status-vocabulary reconciliation; metadata reconciliation only; no new mathematical check. \emph{Continuation recorded proof scope.} The whole current proof was independently derived at 09:44:23 UTC before target reading: actual Picard cocycle translations and inverses, finite dense-open intersections in the specified components, all Cartier graded-factor maps and the unsplit full-fibre induction. The new exact rational program passed 459 checks within its stated finite bounds; the universal proof is written separately. The original Stacks Cartier and cohomology inputs were reopened at Tags 0B3P and 0C4T, Divisors p.~26, and Tag 01DZ, Cohomology pp.~2--3. These inputs are classical; the supporting-row source gate is complete.}{the hypotheses, notation and definitions stated here; \cref{cdp:st:A.3.1}}

For a fixed multiple fibre $mS$, suppose that the set of line bundles $B$ for which $H^0(S,(\Omega_Y^1\otimes B)|_S)=0$ contains a dense open subset $U$ of the relevant Picard component. The intersection
\[
\bigcap_{j=0}^{m-1}\{B:B\otimes\OO_Y(-jS)\in U\}
\tag{A.14}\label{cdp:eq:A.14}
\]
is a dense open subset on which the full-fibre ambient group vanishes, whenever the displayed translations preserve that component. The same assertion for a finite union of components follows by using the appropriate dense open subset in each translated component.
\end{proposition}


\begin{proof}
\noindent\textbf{Step 1.} Tensoring by the fixed line bundle $\OO_Y(-jS)$ is a biholomorphic translation between Picard components. Its inverse image of a dense open set is dense and open. The finite intersection retains these properties. Its defining condition is exactly the vanishing of every graded group in \ref{cdp:st:A.3.1}, which proves the result.
\end{proof}



This finite-translation argument supplies the thickening step intended in Remark~7.2 once ambient vanishing on the smooth support is established. On a smooth support, the conormal sequence has $Q=\Omega_S^1$ and $K=N^*$, so the two requisite support vanishings do establish that ambient vanishing. Remark~7.4 requires the same argument on a singular support; a statement solely about $Q=\widetilde\Omega_S^1$ leaves the kernel in \ref{cdp:st:A.1.1} to be checked. Thus the thickening operation itself has a complete exact treatment, while the cited singular-support assertion cannot supply its missing kernel hypothesis. Example \ref{cdp:st:A.3.2} disproves the inference from a single support vanishing even with torsion conormal; it does not refute a generic-vanishing assertion already established for all the finitely translated twists.




\begin{proposition}[The smooth-multiple assertion of Remark 7.2]\label{cdp:st:A.3.4}
\StatementRecord{cdp:st:A.3.4}{proved}{Historical events: Claude Fable 5.2 (Anthropic), effort high, 16 November 2026: carried as background without fault found. GPT-6 Astra (OpenAI), ChatGPT Pro mode, effort Pro, S6\_09\_PAPERS, 17 November 2026: targeted proof check. Also checked in this writing. Pre-continuation writing events: GPT-6 Astra (OpenAI), ChatGPT Pro mode, effort Pro, S6\_10\_PAPERS, 18 November 2026 04:12:55Z: targeted proof check. Continuation mathematical and source checks: GPT-6 Astra (OpenAI), ChatGPT Pro mode, effort Pro, S6\_10\_PAPERS continuation, 19 November 2026: computation or formal certificate; independent re-proof; targeted proof check; used and rechecked in this run. Continuation metadata events: GPT-6 Astra (OpenAI), ChatGPT Pro mode, effort Pro, S6\_10\_PAPERS continuation, 19 November 2026: used and rechecked in this run; status-vocabulary reconciliation; metadata reconciliation only; no new mathematical check. \emph{Continuation recorded proof scope.} A complete alternative assertion proof and 1,712 fresh exact checks were frozen at 10:00:42 UTC before target reading. They retain the actual Picard/transgression maps, all rational normal-line shifts and every full-fibre layer. The new Step~5 finite Dolbeault kernel and analytic-minor argument matches that before-reading derivation; the exact classical Green/Hodge source was opened in Demailly, author-PDF body pp.~15 and 32. The retained root-cover and finite-flat descent route in Steps~1--4, including the final singular-support scope sentence, is awaiting its separate fresh residual comparison. No whole-current-proof completion is inferred from the alternative proof.}{the hypotheses, notation and definitions stated here; \cref{cdp:st:D.1.9}; \cref{cdp:st:A.4.9}; \cref{cdp:st:A.4.7}; \cref{cdp:st:A.3.1}}

Under the hypotheses of CDP20 Proposition~7.3, let $mS$, with $m\ge2$, be a full fibre whose reduced support $S$ is smooth. For every real irrational transgression parameter $r$ of \ref{cdp:st:D.1.9},
\[
 H^0\bigl(mS,(\Omega_Y^1\otimes L_r)|_{mS}\bigr)=0.
 \tag{A.15a}\label{cdp:eq:A.15a}
\]
The same simultaneous vanishing on all smooth multiple fibres $(m\ge2)$ holds on the complement of a proper complex analytic subset of $\Pic(Y)$. The assertion includes the entire scheme-theoretic fibre, with every nilpotent layer.
\end{proposition}


\begin{proof}
\noindent\textbf{Step 1.} Choose a disc $\Delta$ about the fibre value and write $U=g^{-1}(\Delta)$. Form the base change $t=s^m$ and normalize it, obtaining $p:U'\to U$ and $g':U'\to\Delta'$. This map $p$ is finite étale of degree $m$. Indeed, near a point of the smooth support write $g=u z^m$ with $u$ a holomorphic unit. The normalization is given by
\[
 w^m=u,\qquad s=wz.
\]
The derivative $m w^{m-1}$ is invertible, so these equations give an unramified cover of the original smooth chart. They cover the normalized base change, since they coincide away from $z=0$ and are normal. The local constructions glue by uniqueness of normalization. The divisor $S'=(g')^{-1}(0)$ is smooth and reduced, $g'$ is a submersion near it, and $p^{-1}(mS)=mS'$ as Cartier divisors.

\noindent\textbf{Step 2.} Pull back the unitary local system (A.22) on $U$ to $U'$; its associated holomorphic line is $A'=p^*L_r$. Every nearby smooth fibre of $g'$ maps isomorphically onto the corresponding nearby fibre of $g$. Its character is therefore the irrational exponential of a nonzero integral class. The smooth-fibre calculation in \ref{cdp:st:A.4.9} gives zero cohomology in degrees zero and one for that local system. Proper smooth local triviality at $S'$ transports the same cohomology statement to $S'$. If a component presentation of the cover is used, the argument applies on each component and the groups add.

\noindent\textbf{Step 3.} A holomorphic section of a unitary-flat line on a compact smooth complex manifold is parallel: its squared norm is plurisubharmonic and hence constant by the maximum principle. Thus $H^0(S',A')=0$ follows from the zero degree-zero local-system cohomology. \ref{cdp:st:A.4.7} gives $H^0(S',\Omega_{S'}^1\otimes A')=0$ from degree-one vanishing. The conormal of the reduced full fibre $S'$ is trivial, with global generator $dg'$. Its conormal sequence, tensored by $A'$, now gives
$H^0(S',(\Omega_{U'}^1\otimes A')|_{S'})=0$.

\noindent\textbf{Step 4.} All the graded factors for $mS'$ in \ref{cdp:st:A.3.1} equal this same restricted bundle, because the conormal is trivial. Induction in (A.10) proves $H^0(mS',(\Omega_{U'}^1\otimes A')|_{mS'})=0$. Since $p$ is étale, $p^*\Omega_U^1=\Omega_{U'}^1$, and finite faithfully flat pullback injects sections on $mS$ into sections on $mS'$. One can verify that injection locally by tensoring the inclusion of a module into its faithfully flat scalar extension, or by the trace divided by $m$. This proves \eqref{cdp:eq:A.15a}.

\noindent\textbf{Step 5.} There are finitely many non-submersion values: the critical locus is a compact analytic set with finitely many irreducible components, and $dg=0$ makes $g$ constant on the smooth locus, hence on the closure, of each component. Fix the root cover above each smooth multiple fibre and write its smooth central surface as $S'_c$. The exponential sequence identifies $\Pic(Y)$ with the complex vector space $H^1(Y,\OO_Y)$. Exponentiating linear combinations of a basis of holomorphic Cech cocycles gives the actual Picard family. On $S'_c$ consider the holomorphic family of bundles
\[
 E_B=(\Omega^1_{U'_c}\otimes p_c^*B)|_{S'_c}.
\]
We give the analytic nonvanishing locus by an explicit finite kernel map. Near a parameter $B_0$, smooth gauges obtained from a partition of unity and holomorphic parameter-dependent logarithms of the transition ratios identify the underlying smooth bundles. Their Dolbeault operators have the form $D_b=D_0+A_b$, where $A_b$ has order zero, depends holomorphically on $b$, vanishes at $B_0$, and satisfies
\[
 D_0A_b+A_bD_0+A_b^2=0.
\]
Choose Hermitian metrics on $S'_c$ and the central bundle. Let $i$ include the finite-dimensional harmonic spaces, let $\pi$ be harmonic projection, let $G$ be the Green operator for the Dolbeault Laplacian, and put $h=D_0^*G$. The classical Green and Hodge identities give
\[
 D_0h+hD_0=1-i\pi,\qquad \pi i=1,
 \qquad \pi h=0,\qquad hi=0;
\]
see \cite[Part~I, Theorem~3.15, author-PDF body p.~15, and Theorems~7.1--7.2, body p.~32]{DemaillyHodgeAuthor2002}. These results apply to compact Hermitian manifolds without a K\"ahler hypothesis. On fixed sufficiently high Sobolev spaces, shrink the parameter neighbourhood so that the Neumann inverses $T_b^q=(1+hA_b)^{-1}$ exist in the degrees and adjacent regularities used below. Define the finite holomorphic matrix
\[
 M_b=\pi^1 A_b T_b^0 i^0:
 \mathcal H^0(E_{B_0})\longrightarrow\mathcal H^1(E_{B_0}).
\]
The two inverse kernel maps are $v\mapsto T_b^0i^0v$ and $u\mapsto\pi^0u$. Indeed, $D_bu=0$ implies
\[
 u-i\pi u=hD_0u=-hA_bu,
 \qquad u=T_b^0i^0\pi^0u.
\]
Conversely, for $u=T_b^0i^0v$, the contraction and integrability identities give
\[
 \begin{aligned}
 D_bu&=i\pi A_bu+hD_0A_bu
       =i\pi A_bu-hA_bD_bu,\\
 D_bu&=T_b^1i^1M_bv,\qquad
 \pi^0T_b^0i^0v=v.
 \end{aligned}
\]
Thus $\ker M_b$ is isomorphic by these actual maps to the section space of $E_B$; elliptic regularity makes every resulting Sobolev solution smooth. The nonzero-section locus is the zero set of the maximal minors expressing failure of full column rank of $M_b$, with the zero-domain and insufficient-row cases interpreted directly. It is therefore a closed complex analytic subset $Z_c$ of the parameter space. The conclusion follows from this kernel calculation, not from numerical upper semicontinuity alone. Step~3 shows that $Z_c$ excludes every real irrational $L_r$. Outside $Z_c$, all the graded factors on $mS'_c$ vanish, and Step~4 gives the full-fibre vanishing after descent. The finite union of the $Z_c$ is proper because the same $L_r$ lies outside every member. This proves the simultaneous generic assertion and the smooth-multiple conclusion of Remark~7.2. For a singular support the identical root cover retains its singularities and the kernel $N^*\otimes\mathcal J^{-1}$; the present proof does not erase that term in Remark~7.4.
\end{proof}






\begin{proposition}[Every reducible full fibre has generic ambient vanishing]\label{cdp:st:A.3.5}
\StatementRecord{cdp:st:A.3.5}{proved}{Historical events: Claude Fable 5.2 (Anthropic), effort high, 16 November 2026: carried as background without fault found. GPT-6 Astra (OpenAI), ChatGPT Pro mode, effort Pro, S6\_09\_PAPERS, 17 November 2026: targeted proof check. Also checked in this writing. Pre-continuation writing events: GPT-6 Astra (OpenAI), ChatGPT Pro mode, effort Pro, S6\_10\_PAPERS, 18 November 2026: targeted proof check. Continuation metadata events: GPT-6 Astra (OpenAI), ChatGPT Pro mode, effort Pro, S6\_10\_PAPERS continuation, 19 November 2026: used and rechecked in this run; status-vocabulary reconciliation; metadata reconciliation only; no new mathematical check. \emph{Pre-continuation recorded proof scope.} The adjacent-component divisor, negative twists, Krull intersection, and the full Cartier multiplicity filtration were re-derived}{the hypotheses, notation and definitions stated here; the explicit constructions and calculations in the proof}

Let $Y$ be a smooth compact complex threefold with $H^2(Y,\Q)=0$, and let $F=\sum_i m_iS_i$ be an effective Cartier divisor with connected reducible support. Every connected component of $\Pic(Y)$ contains a dense open subset, with proper closed analytic complement, on which
\[
 H^0(F,(\Omega_Y^1\otimes L)|_F)=0.
 \tag{A.15b}\label{cdp:eq:A.15b}
\]
The multiplicities may be unequal. For a proper connected fibration over a compact curve, the vanishing is simultaneous on all reducible full fibres after intersecting finitely many such open sets. Consequently a full fibre with ambient sections for every Picard twist has irreducible support.
\end{proposition}


\begin{proof}
\noindent\textbf{Step 1.} Every component $S_i$ meets a different component $S_j$, because the support is connected and reducible. The canonical section of $\OO_Y(S_j)$ restricts to a nonzero section on $S_i$. It is a nonzerodivisor in every local ring of $S_i$: otherwise one local branch of $S_i$ would be contained in $S_j$, and analytic continuation on the global irreducible component would give $S_i=S_j$. Its zero scheme is therefore a nonzero effective Cartier divisor $D$ on $S_i$.

\noindent\textbf{Step 2.} For any fixed locally free sheaf $E$ on the reduced irreducible compact surface $S_i$, the injections
\[
 H^0(S_i,E(-nD))\hookrightarrow H^0(S_i,E)
\]
form a decreasing sequence of finite-dimensional subspaces. Their intersection is zero. At a point of $D$, a section in the intersection has its germ in every power of the Cartier ideal; the Krull intersection theorem for the Noetherian analytic local ring makes this germ zero. The identity theorem on the reduced irreducible surface, applied in local trivializations of $E$, then makes the global section zero. A descending sequence of finite-dimensional subspaces stabilizes, so $H^0(S_i,E(-nD))=0$ for all sufficiently large $n$.

\noindent\textbf{Step 3.} Fix a Picard component and a representative $L_0$ in it. The integral class $c_1(\OO_Y(S_j))$ is torsion because $H^2(Y,\Q)=0$; choose a positive integer $q$ killing it. Then $\OO_Y(-qnS_j)$ belongs to $\Pic^0(Y)$ for every $n$. Apply the preceding argument to $E=(\Omega_Y^1\otimes L_0)|_{S_i}$. It supplies, in the chosen Picard component, a twist $L_0\otimes\OO_Y(-qnS_j)$ with zero ambient group on $S_i$. Local holomorphic Picard families and proper semicontinuity show that the nonvanishing locus on that component is a proper closed analytic subset. Its complement is dense and open. The same argument works in every Picard component.

\noindent\textbf{Step 4.} Choose an ordered list $A_1,\ldots,A_N$ of the prime divisors $S_i$, with $S_i$ repeated $m_i$ times, and put $F_0=0$, $F_k=A_1+\cdots+A_k$. The exact Cartier sequences, tensored with the locally free ambient bundle, are
\[
\begin{split}
0\longrightarrow
 (\Omega_Y^1\otimes L\otimes\OO_Y(-F_{k-1}))|_{A_k}
 \longrightarrow (\Omega_Y^1\otimes L)|_{F_k}\\
 \longrightarrow(\Omega_Y^1\otimes L)|_{F_{k-1}}
 \longrightarrow0.
\end{split}
\]
This follows locally from multiplication by an equation of $F_{k-1}$, and remains exact when a component occurs repeatedly. For each $k$ take the dense open set just proved on $A_k$ and translate it by the fixed global line $\OO_Y(-F_{k-1})$, using the appropriate Picard component. Their finite intersection is dense open. Every displayed first term has no sections there, so induction from $F_0=0$ proves \eqref{cdp:eq:A.15b} on the full scheme $F_N=F$.

\noindent\textbf{Step 5.} A proper holomorphic map from a smooth compact manifold to a smooth curve has finitely many critical values, since its proper analytic critical-value set does not contain a general point. Its smooth connected fibres are irreducible. Thus there are finitely many reducible fibres, and a further finite intersection gives the asserted simultaneous vanishing. This proves the final consequence and resolves the unequal-multiplicity transfer in this case without any assertion about torsion-free differentials on individual components.
\end{proof}






\begin{proposition}[Both fibre groups on the specified $W$]\label{cdp:st:A.4.1}
\StatementRecord{cdp:st:A.4.1}{proved}{Historical events: Claude Fable 5.2 (Anthropic), effort high, 16 November 2026: carried as background without fault found. GPT-6 Astra (OpenAI), ChatGPT Pro mode, effort Pro, S6\_09\_PAPERS, 17 November 2026: targeted proof check. Also checked in this writing. Pre-continuation writing events: GPT-6 Astra (OpenAI), ChatGPT Pro mode, effort Pro, S6\_10\_PAPERS, 18 November 2026: targeted proof check. Continuation metadata events: GPT-6 Astra (OpenAI), ChatGPT Pro mode, effort Pro, S6\_10\_PAPERS continuation, 19 November 2026: used and rechecked in this run; status-vocabulary reconciliation; metadata reconciliation only; no new mathematical check. \emph{Pre-continuation recorded proof scope.} The topologically trivial line pulls back trivially to dP6 by the exponential sequence. Cotangent forms pull back through the normalization and extend across the blowup points, where the projective-plane calculation kills them. The exact conormal kernel gives the claimed ambient dimension}{the hypotheses, notation and definitions stated here; \cref{cdp:st:A.1.2}; \cref{cdp:st:A.1.1}}

For the reduced normal-crossings fibre $W$ with normalization $\eta:\operatorname{dP}_6\to W$ in the specified threefold $X$, and every $L\in\Pic(X)$, put $A=(K_X\otimes L^{-1})|_W$. Then
\[
H^0(W,Q\otimes A)=0,
\qquad H^0(W,\Omega_X^1|_W\otimes A)\simeq\C.
\tag{A.15}\label{cdp:eq:A.15}
\]
The same two conclusions hold with $A=B|_W$ for every $B\in\Pic(X)$.
\end{proposition}



The nonvanishing on the compact example and the torsion-free-cotangent vanishing are contained in \citet[Theorem~10.5 and Section~10.4]{AC2026}. The present calculation sharpens the ambient section dimension to exactly one; that sharper equality and its proof are distinguished from the source nonvanishing.

\begin{proof}
\noindent\textbf{Step 1.} The only global topological input is $H^2(X,\Z)=0$, and the only identification of the fibre used is its specified normalization as the blowup of $\PP^2$ at three points. A blowup at a smooth point has $\pi_*\OO=\OO$ and $R^1\pi_*\OO=0$. The latter follows from the theorem on formal functions: successive exceptional-neighborhood quotients are $\OO_{\PP^1}(k)$, $k\ge0$, whose first cohomology is zero. Thus $H^1(\operatorname{dP}_6,\OO)=H^1(\PP^2,\OO)=0$. The exponential sequence makes $c_1:\Pic(\operatorname{dP}_6)\to H^2(\operatorname{dP}_6,\Z)$ injective. The pullback of every line bundle on $X$ has zero first Chern class and is therefore trivial.

\noindent\textbf{Step 2.} The differential of normalization gives $\Omega_W^1\to\eta_*\Omega^1_{\operatorname{dP}_6}$. Its target is torsion-free as an $\OO_W$-module, so it factors through $Q$. The resulting map is injective because it is an isomorphism on the dense regular locus and its source is torsion-free. A holomorphic one-form on the blowup restricts away from the exceptional curves to a one-form on $\PP^2$ minus three points. Its coefficients extend over these points by Hartogs' theorem, and $H^0(\PP^2,\Omega^1_{\PP^2})=0$ by the Euler sequence. Hence $H^0(\operatorname{dP}_6,\Omega^1)=0$. The projection formula and the triviality of $\eta^*A$ prove the first equality in \eqref{cdp:eq:A.15}.

\noindent\textbf{Step 3.} Because $W$ is the full reduced fibre, its conormal is trivial. \ref{cdp:st:A.1.2} identifies $K\otimes A$ with $\eta_*\eta^*A=\eta_*\OO_{\operatorname{dP}_6}$. The normalization is connected and compact, so its global holomorphic functions are $\C$. The cohomology sequence of \ref{cdp:st:A.1.1} and the first equality in \eqref{cdp:eq:A.15} identify this group with the ambient group. The same argument applies to every $B$.
\end{proof}






\begin{proposition}[Counterexample: A specific implication in the printed proof of Proposition 7.3]\label{cdp:st:A.4.2}
\StatementRecord{cdp:st:A.4.2}{proved}{Historical events: Claude Fable 5.2 (Anthropic), effort high, 16 November 2026: carried as background without fault found. GPT-6 Astra (OpenAI), ChatGPT Pro mode, effort Pro, S6\_09\_PAPERS, 17 November 2026: targeted proof check. Also checked in this writing. Pre-continuation writing events: GPT-6 Astra (OpenAI), ChatGPT Pro mode, effort Pro, S6\_10\_PAPERS, 18 November 2026: targeted proof check. Continuation metadata events: GPT-6 Astra (OpenAI), ChatGPT Pro mode, effort Pro, S6\_10\_PAPERS continuation, 19 November 2026: used and rechecked in this run; status-vocabulary reconciliation; metadata reconciliation only; no new mathematical check. \emph{Pre-continuation recorded proof scope.} The product E times P1 has the explicitly constant restricted twists and cotangent cohomology tuple (1,2,1), which disproves exactly the stated implication}{the hypotheses, notation and definitions stated here; the explicit constructions and calculations in the proof}

The implication that a ruled minimal surface $S_0$ with a one-parameter family of numerically trivial line bundles $M_t$ and nonzero sections of $\Omega^1_{S_0}\otimes M_t$ must be ruled over a curve of genus at least two is false when the family $M_t$ is allowed to be constant, as the printed proof explicitly allows.
\end{proposition}


\begin{proof}
\noindent\textbf{Step 1.} Take $S_0=E\times\PP^1$ for an elliptic curve $E$ and take $M_t=\OO_{S_0}$ for every parameter. The product cotangent decomposition gives
\[
 h^q(S_0,\Omega^1_{S_0})=(1,2,1),\qquad q=0,1,2.
\]
Indeed the $E$ cotangent summand contributes $(1,1,0)$, while the $\PP^1$ cotangent summand contributes $(0,1,1)$ by the Kunneth formula and $h^1(\PP^1,\OO(-2))=1$. A nonzero invariant form on $E$ provides the asserted sections for every parameter, whereas the ruling has genus one. This refutes the stated implication. It is a counterexample to that implication with its explicit data; it supplies no threefold satisfying all of Proposition~7.3's ambient hypotheses.
\end{proof}






\begin{proposition}[The omitted genus-one case for the standard rank-one fibre]\label{cdp:st:A.4.3}
\StatementRecord{cdp:st:A.4.3}{proved}{Historical events: Claude Fable 5.2 (Anthropic), effort high, 16 November 2026: carried as background without fault found. GPT-6 Astra (OpenAI), ChatGPT Pro mode, effort Pro, S6\_09\_PAPERS, 17 November 2026: targeted proof check. Also checked in this writing. Pre-continuation writing events: GPT-6 Astra (OpenAI), ChatGPT Pro mode, effort Pro, S6\_10\_PAPERS, 18 November 2026: targeted proof check. Continuation metadata events: GPT-6 Astra (OpenAI), ChatGPT Pro mode, effort Pro, S6\_10\_PAPERS continuation, 19 November 2026: used and rechecked in this run; status-vocabulary reconciliation; metadata reconciliation only; no new mathematical check. \emph{Pre-continuation recorded proof scope.} The line pulls back from the elliptic base, the relative cotangent has fibre degree minus two, and the remaining base form descends precisely when the actual scalar gluing is one}{the hypotheses, notation and definitions stated here; the explicit constructions and calculations in the proof}

Let $E$ be an elliptic curve, let $P\in\Pic(E)$, and let $T=\PP_E(\OO_E\oplus P)$ with its zero and infinity sections $D_0,D_\infty$. Glue $D_0$ to $D_\infty$ by a translation of $E$ to form the compact normal-crossings surface $S$, and denote its normalization by $\eta:T\to S$. Every line bundle $A$ on $S$ whose pullback $\eta^*A$ is topologically trivial satisfies
\[
H^0(S,\widetilde\Omega_S^1\otimes A)=
\begin{cases}\C,&A\simeq\OO_S,\\0,&A\not\simeq\OO_S.\end{cases}
\tag{A.16}\label{cdp:eq:A.16}
\]
\end{proposition}


\begin{proof}
\noindent\textbf{Step 1.} For the ruling $p:T\to E$, every topologically trivial line bundle on $T$ is $p^*P_A$ for a unique $P_A\in\Pic^0(E)$. Indeed its degree on every ruling fibre is zero, so the adjunction map from the pullback of its rank-one direct image is an isomorphism, checked on each fibre by $H^0(\PP^1,\OO)=\C$ and $H^1(\PP^1,\OO)=0$. The line on $E$ has degree zero because its first Chern class pulls back to zero and $p$ has a section.

\noindent\textbf{Step 2.} The relative cotangent sequence and restriction to ruling fibres give
\[
 H^0(T,\Omega_T^1\otimes p^*P_A)
 =p^*H^0(E,\Omega_E^1\otimes P_A).
\]
The relative quotient has no sections because its restriction to every fibre is $\OO_{\PP^1}(-2)$. Since $\Omega_E^1\simeq\OO_E$, the displayed group vanishes for nontrivial $P_A$. This proves \eqref{cdp:eq:A.16} in that case by the injection into normalization differentials.

\noindent\textbf{Step 3.} For $P_A=\OO_E$, choose a trivialization of $\eta^*A$. The descent datum for $A$ is a nonzero constant $\lambda$ identifying the two restrictions to $D_0,D_\infty$, because these are trivial lines on the compact connected curve $E$. Multiplying the chosen trivialization by a constant changes neither $\lambda$ nor its meaning. The local calculation for $xy=0$ identifies the image of $\widetilde\Omega_S^1$ in $\eta_*\Omega_T^1$ with pairs of forms whose tangential restrictions to the double curve agree. The unique normalization form is a scalar multiple of $p^*\alpha$, where $\alpha$ is a nonzero invariant form on $E$. Translation pulls $\alpha$ back to itself. Twisted descent is therefore the equation $\alpha=\lambda\alpha$, so a nonzero section exists exactly when $\lambda=1$. That equality also says precisely that the trivialization descends to $A\simeq\OO_S$. This proves both dimensions.
\end{proof}






\begin{proposition}[Canonical triviality on an irreducible reduced torus fibre]\label{cdp:st:A.4.4}
\StatementRecord{cdp:st:A.4.4}{proved}{Historical events: Claude Fable 5.2 (Anthropic), effort high, 16 November 2026: carried as background without fault found. GPT-6 Astra (OpenAI), ChatGPT Pro mode, effort Pro, S6\_09\_PAPERS, 17 November 2026: targeted proof check. Also checked in this writing. Pre-continuation writing events: GPT-6 Astra (OpenAI), ChatGPT Pro mode, effort Pro, S6\_10\_PAPERS, 18 November 2026: targeted proof check. Continuation metadata events: GPT-6 Astra (OpenAI), ChatGPT Pro mode, effort Pro, S6\_10\_PAPERS continuation, 19 November 2026: used and rechecked in this run; status-vocabulary reconciliation; metadata reconciliation only; no new mathematical check. \emph{Pre-continuation recorded proof scope.} The rank-one canonical direct image, vertical evaluation divisor, torsion canonical restriction, and coherent Euler characteristic calculation were re-derived}{the hypotheses, notation and definitions stated here; \cref{cdp:st:A.5.2}}

Let $g:Y\to C$ be a proper holomorphic map from a connected smooth compact threefold onto a smooth compact curve, with connected fibres and general fibre a complex two-torus. Every irreducible reduced full fibre $S=g^{-1}(c)$ has $\omega_S\simeq\OO_S$. When $H^2(Y,\Q)=0$, it therefore satisfies $h^q(S,\OO_S)=(1,2,1)$.
\end{proposition}


\begin{proof}
\noindent\textbf{Step 1.} The coherent sheaf $g_*\omega_Y$ is torsion-free: a nonzero function on an open subset of the base cannot annihilate a holomorphic section of a line bundle on the dominating smooth total space. It has generic rank one because the smooth torus fibre has one holomorphic top form; after restricting to a base open set where Grauert base change holds, this identifies its rank. On the smooth curve it is consequently a line bundle $M$. The evaluation morphism $g^*M\to\omega_Y$ is generically nonzero, so its zero scheme is an effective Cartier divisor $D$ and $\omega_Y\simeq g^*M\otimes\OO_Y(D)$. No component of $D$ dominates the base: on a general smooth torus fibre its defining section is a nonzero holomorphic top form and is nowhere zero. Thus $D$ is vertical.

\noindent\textbf{Step 2.} All vertical components meeting the specified irreducible full fibre are that fibre itself. Consequently $\OO_Y(D)|_S$ is a power of $\OO_Y(S)|_S$, and the latter is trivial because a coordinate on $C$ defines the full reduced fibre. The pullback $g^*M|_S$ is also trivial. Adjunction gives $\omega_S=(\omega_Y\otimes\OO_Y(S))|_S\simeq\OO_S$. The fibre is reduced and connected, so $h^0(\OO_S)=1$; it is a Cartier divisor, hence Cohen--Macaulay and Gorenstein, so Serre duality gives $h^2(\OO_S)=1$. The divisor Riemann--Roch calculation in \ref{cdp:st:A.5.2} gives $\chi(\OO_S)=0$ under the stated rational cohomology hypothesis, yielding $h^1=2$.
\end{proof}






\begin{proposition}[The full torsion-free cotangent assertion for irreducible normal-crossings torus fibres]\label{cdp:st:A.4.5}
\StatementRecord{cdp:st:A.4.5}{proved}{Historical events: Claude Fable 5.2 (Anthropic), effort high, 16 November 2026: carried as background without fault found. GPT-6 Astra (OpenAI), ChatGPT Pro mode, effort Pro, S6\_09\_PAPERS, 17 November 2026: targeted proof check. Also checked in this writing. Pre-continuation writing events: GPT-6 Astra (OpenAI), ChatGPT Pro mode, effort Pro, S6\_10\_PAPERS, 18 November 2026: targeted proof check; used and rechecked in this run. Continuation metadata events: GPT-6 Astra (OpenAI), ChatGPT Pro mode, effort Pro, S6\_10\_PAPERS continuation, 19 November 2026: used and rechecked in this run; status-vocabulary reconciliation; metadata reconciliation only; no new mathematical check. \emph{Pre-continuation recorded proof scope.} The Jacobian cokernel gives the reduced singular-locus structure sheaf, so admissibility does not assume a prior Kodaira smoothing. Schroer-Siebert Propositions2.4-2.6,3.1 and4.1 were opened. Rational, ruled and Hopf form calculations and the parallel-section calculation give the stated ambient kernel}{the hypotheses, notation and definitions stated here; \cref{cdp:st:A.4.4}; \cite[Definition~2.1]{SS2002}; \cref{cdp:st:A.4.1}; \cref{cdp:st:A.4.3}; \cref{cdp:st:A.1.2}}

Let $g:Y\to C$ have the hypotheses of \ref{cdp:st:A.4.4} and let $S$ be a singular irreducible reduced full fibre with normal crossings. For every nontrivial unitary-flat line bundle $A$ on $S$,
\[
H^0(S,\widetilde\Omega_S^1\otimes A)=0,
\qquad
H^0(S,\Omega_Y^1|_S\otimes A)
 \simeq H^0(\widetilde S,\eta^*A).
\tag{A.17}\label{cdp:eq:A.17}
\]
The latter group has dimension one when $\eta^*A$ is holomorphically trivial, and zero otherwise. The normalization is exactly one of the following: a diagonal Hopf surface with two disjoint elliptic curves identified; a ruled surface over an elliptic curve with two disjoint sections identified by translation; or a rational surface with an anticanonical hexagon as in Schr\"oer--Siebert Proposition~2.6.
\end{proposition}


\begin{proof}
\noindent\textbf{Step 1.} We check the hypotheses of the classification, including its deformation sheaf. \ref{cdp:st:A.4.4} gives $\omega_S\simeq\OO_S$. Dualizing the locally free conormal resolution of $\Omega_S^1$ yields
\[
 \mathcal T^1_S=\mathcal Ext^1_{\OO_S}(\Omega_S^1,\OO_S)
 \simeq\OO_S/\mathcal J=\OO_{S_{\rm sing,red}}.
\]
Here the conormal is trivial because $S$ is the full reduced fibre; the local identities $\mathcal J=(x,y)$ at a double point and $\mathcal J=(xy,xz,yz)$ at a triple point show that this is the reduced singular locus. Hence $S$ is d-semistable in the exact sense of \cite[Definition~2.1]{SS2002}. That definition imposes irreducibility, normal crossings, d-semistability and trivial dualizing sheaf; it does not require a prior smoothing to Kodaira surfaces. Propositions~2.4, 2.5 and 2.6 there give the three normalization cases. In the ruled case their Proposition~4.1 says precisely that trivial $\omega_S$ forces the section identification to be a translation. The comparison with their notation is $X_{\mathrm{SS}}=S$ and $S_{\mathrm{SS}}=\widetilde S$.

\noindent\textbf{Step 2.} In the rational case $\eta^*A$ is trivial: its first Chern class is torsion because the line is unitary-flat, whereas the integral second cohomology of a smooth rational surface is torsion-free and its Picard group injects into that group by the exponential sequence. Such a surface has no holomorphic one-forms, by birational invariance as proved in \ref{cdp:st:A.4.1}. The injection $\widetilde\Omega_S^1\to\eta_*\Omega^1_{\widetilde S}$ therefore proves the first vanishing.

\noindent\textbf{Step 3.} In the elliptic ruled case the proof of \ref{cdp:st:A.4.3} applies to a decomposable ruled surface of any degree: the two disjoint sections express it as $\PP_E(\OO_E\oplus P)$, with $P\in\Pic(E)$. The integral second cohomology of this ruled surface is torsion-free, by the projective-bundle formula. A unitary-flat line has torsion first Chern class, so its pullback is topologically trivial and comes from $\Pic^0(E)$. No degree assumption on $P$ occurs in the cotangent calculation, and \ref{cdp:st:A.4.3} proves the required vanishing for every nontrivial $A$.

\noindent\textbf{Step 4.} In the Hopf case, pull back to a finite primary diagonal Hopf cover. Its universal covering is $\C^2\setminus\{0\}$ with contraction $(z_1,z_2)\mapsto(a_1z_1,a_2z_2)$, where $|a_i|<1$. A holomorphic one-form valued in a unitary-flat line lifts to $f_1(z)dz_1+f_2(z)dz_2$, with each $f_i$ holomorphic on the punctured space and hence on $\C^2$ by Hartogs' theorem. For a contraction character $\rho$ of modulus one, equivariance says that the coefficient of $z_1^{k_1}z_2^{k_2}dz_i$ can be nonzero only when
\[
 a_1^{k_1}a_2^{k_2}a_i=\rho
\]
(or its inverse, according to the convention for the associated line). The modulus on the left is strictly less than one and the character has modulus one. Thus every coefficient is zero. Finite etale pullback is injective on sections, so the twisted one-form group on the Hopf normalization also vanishes, and again the normalization injection proves the first assertion on $S$.

\noindent\textbf{Step 5.} Finally \ref{cdp:st:A.1.2} and the first vanishing identify the ambient group with $H^0(\widetilde S,\eta^*A)$. A section of a unitary-flat line on a compact connected complex manifold is either zero or parallel and nowhere zero: its squared norm is locally a sum of squares of holomorphic functions, hence plurisubharmonic, and the maximum principle makes it constant. It therefore trivializes the line, and in the trivial case the section space is $\C$. This proves the dimension statement.
\end{proof}






\begin{proposition}[Euler numbers of these irreducible normal-crossings fibres]\label{cdp:st:A.4.6}
\StatementRecord{cdp:st:A.4.6}{proved}{Historical events: Claude Fable 5.2 (Anthropic), effort high, 16 November 2026: carried as background without fault found. GPT-6 Astra (OpenAI), ChatGPT Pro mode, effort Pro, S6\_09\_PAPERS, 17 November 2026: targeted proof check. Also checked in this writing. Pre-continuation writing events: GPT-6 Astra (OpenAI), ChatGPT Pro mode, effort Pro, S6\_10\_PAPERS, 18 November 2026: targeted proof check; used and rechecked in this run. Continuation metadata events: GPT-6 Astra (OpenAI), ChatGPT Pro mode, effort Pro, S6\_10\_PAPERS continuation, 19 November 2026: used and rechecked in this run; status-vocabulary reconciliation; metadata reconciliation only; no new mathematical check. \emph{Pre-continuation recorded proof scope.} The normalization-conductor Euler formula gives zero for the Hopf and ruled cases, and six minus six plus two for the rational hexagon. The global Euler sum and middle-degree pairing give the stated two possible values}{the hypotheses, notation and definitions stated here; \cref{cdp:st:A.4.5}}

For $S$ in \ref{cdp:st:A.4.5},
\[
e(S)=\begin{cases}
2,&\widetilde S\text{ is rational},\\
0,&\widetilde S\text{ is elliptic ruled or Hopf}.
\end{cases}
\tag{A.18}\label{cdp:eq:A.18}
\]
Thus an arbitrary compact torus fibration whose fibres are irreducible reduced normal-crossings surfaces or multiples of smooth surfaces of Euler number zero has
\[
c_3(Y)=e(Y)=2r,
\tag{A.19}\label{cdp:eq:A.19}
\]
where $r$ is the number of rational-normalization fibres. For $b_1(Y)=b_2(Y)=0$, one has $r\in\{0,1\}$ and $b_3(Y)=2-2r$.
\end{proposition}


\begin{proof}
\noindent\textbf{Step 1.} Write $D\subset\widetilde S$ for the conductor support and $N\subset S$ for the singular curve. Normalization is a homeomorphism off these sets, so additivity of the topological Euler characteristic gives $e(S)=e(\widetilde S)-e(D)+e(N)$. In the Hopf case $e(\widetilde S)=0$, $D$ is a disjoint union of two elliptic curves, and $N$ is one elliptic curve; all three Euler numbers vanish. The ruled case has the same conductor curves and $e(\widetilde S)=e(\PP^1)e(E)=0$.

\noindent\textbf{Step 2.} In the rational case the normalization is a Hirzebruch surface blown up twice, so its Euler number is $4+2=6$. The conductor is a cycle of six rational curves, with six nodes, so $e(D)=12-6=6$. Its image consists of three rational curves with exactly two triple points; counting the three preimages of each triple point once gives $e(N)=6-2(3-1)=2$. Hence $e(S)=2$.

\noindent\textbf{Step 3.} Away from finitely many critical values the proper holomorphic map is a differentiable fibre bundle, and a two-torus has Euler number zero. Additivity for a stratification by this open base and the exceptional points therefore gives $e(Y)=\sum_c e((Y_c)_{\rm red})$. Nilpotents do not change the underlying topological space. The assumed fibre classes and \eqref{cdp:eq:A.18} give \eqref{cdp:eq:A.19}. The equality with $c_3$ is the complex-oriented Gauss--Bonnet theorem. Poincare duality for the connected oriented six-manifold gives $e(Y)=2-2b_1+2b_2-b_3$. Under the two Betti vanishings this is $2-b_3$, so $b_3=2-2r\ge0$ and $r$ is zero or one.
\end{proof}






\begin{proposition}[Unitary-flat differentials on every compact reduced complex surface]\label{cdp:st:A.4.7}
\StatementRecord{cdp:st:A.4.7}{proved}{Historical events: Claude Fable 5.2 (Anthropic), effort high, 16 November 2026: carried as background without fault found. GPT-6 Astra (OpenAI), ChatGPT Pro mode, effort Pro, S6\_09\_PAPERS, 17 November 2026: targeted proof check. Also checked in this writing. Pre-continuation writing events: GPT-6 Astra (OpenAI), ChatGPT Pro mode, effort Pro, S6\_10\_PAPERS, 18 November 2026: targeted proof check. Continuation metadata events: GPT-6 Astra (OpenAI), ChatGPT Pro mode, effort Pro, S6\_10\_PAPERS continuation, 19 November 2026: used and rechecked in this run; status-vocabulary reconciliation; metadata reconciliation only; no new mathematical check. \emph{Pre-continuation recorded proof scope.} Stokes closure of the pulled-back unitary form, constancy of primitives on exceptional fibres, finite descent of continuous primitives, and the local-system Cech injection were re-derived. No holomorphic descent through a non-seminormal point was used}{the hypotheses, notation and definitions stated here; the explicit constructions and calculations in the proof}

Let $S$ be a compact reduced complex surface, let $\mathbb A$ be a rank-one unitary complex local system on its underlying topological space, and let $A=\mathbb A\otimes_{\C}\OO_S$ be its associated holomorphic line bundle. There is a canonical complex-linear injection
\[
 H^0(S,\widetilde\Omega_S^1\otimes A)
       \longrightarrow H^1(S,\mathbb A).
 \tag{A.20}\label{cdp:eq:A.20}
\]
The construction uses continuous local primitives that are holomorphic on the normalization; it does not assert that those primitives are holomorphic on a non-seminormal surface.
\end{proposition}


\begin{proof}
\noindent\textbf{Step 1.} Let $\eta:\widetilde S\to S$ be normalization and choose a resolution $\pi:Z\to\widetilde S$ such that the exceptional curves have normal crossings; the components of $Z$ are compact smooth surfaces. Pullback of Kähler differentials along $h=\eta\circ\pi:Z\to S$ kills the torsion of $\Omega_S^1$, because the resulting form vanishes on the dense inverse image of $S_{\rm reg}$ and the target is locally free. Thus a section $\alpha$ on the left of \eqref{cdp:eq:A.20} gives a holomorphic one-form $\alpha_Z$ with values in the unitary-flat line $h^*A$. This map on sections is injective for the same dense-open reason.

\noindent\textbf{Step 2.} Write $\nabla$ for the flat unitary connection. Holomorphicity gives $\nabla\alpha_Z$ type $(2,0)$. The Hermitian pairing identifies the tensor product of the flat line with its conjugate with the trivial local system. Stokes' theorem gives, on each connected component of $Z$,
\[
 \int_Z \nabla\alpha_Z\wedge\overline{\nabla\alpha_Z}
 =\int_Z d\bigl(\alpha_Z\wedge\overline{\nabla\alpha_Z}\bigr)=0.
\]
After multiplication by the fixed dimensional sign, the integrand is the nonnegative squared norm of a $(2,0)$-form times a positive volume form. Hence $\nabla\alpha_Z=0$.

\noindent\textbf{Step 3.} We prove the required local primitive assertion at a point $x\in S$. Trivialize $\mathbb A$ on a sufficiently small connected neighbourhood $V$ of $x$; complex analytic spaces are locally contractible, so such a flat trivialization exists. The section $\alpha$ has a local lift to $\Omega_S^1\otimes A$, since the latter surjects onto the torsion-free quotient. Its pullback to each reduced fibre of $h$ over $x$ is zero: the morphism of that fibre to $S$ is constant, so the induced morphism on Kähler differentials is zero.

\noindent\textbf{Step 4.} For each point $\widetilde x\in\eta^{-1}(x)$, a sufficiently small regular neighbourhood of the exceptional fibre $\pi^{-1}(\widetilde x)$ in the good resolution retracts onto that fibre. This neighbourhood can be constructed by plumbing the normal disc bundles of the exceptional curves at their transverse intersection points; it is a neighbourhood of the entire fibre. A sufficiently small inverse image of a neighbourhood of $\widetilde x$ lies inside it by properness. In the smooth case one instead uses a coordinate ball. Every period of $\alpha_Z$ on the regular neighbourhood is zero, because the form is closed and its restriction to the retracting fibre is zero. Integration from a fixed point therefore produces a single-valued holomorphic primitive: a smooth primitive has differential of type $(1,0)$, so its $\bar\partial$ is zero. Normality and $\pi_*\OO_Z=\OO_{\widetilde S}$ descend that primitive to a holomorphic function $b$ on a neighbourhood of $\widetilde x$. Choose its constant so that $b(\widetilde x)=0$. Perform this independently at every point over $x$ and shrink $V$ so that all its inverse image is covered by these neighbourhoods.

\noindent\textbf{Step 5.} These primitives have equal values on every fibre of $\eta$ over $V$. To see this, form the finite analytic fibre product
$W=\widetilde S\times_S\widetilde S$ with projections $p_1,p_2$. On the normalization of every positive-dimensional irreducible component of $W_{\rm red}$, the differential of $p_1^*b-p_2^*b$ is zero on a dense open subset: both differentials are pullbacks of the same local lift of $\alpha$ from $S$. The isolated singularities of $\widetilde S$ can be omitted for this comparison; holomorphicity then extends the equality. Thus the difference is constant on each such component. The same conclusion is automatic on a zero-dimensional component. Shrink $V$ once more so that every component of the finite space $W|_V$ meets the fibre over $x$; components disjoint from that fibre can be removed by properness. Our normalization $b(\widetilde x)=0$ makes all these constants zero. Hence $p_1^*b=p_2^*b$ as continuous functions, proving equality on the fibres of $\eta$.

\noindent\textbf{Step 6.} A finite proper surjection is a topological quotient map. The functions $b$ therefore descend to continuous functions on $V$, holomorphic after pullback to normalization, whose differential on the regular locus is $\alpha$. On overlaps, any two such primitives differ by a locally constant function: their difference is constant on each local normalized branch, and continuity identifies the constants whenever branches meet. With the flat transition functions of $\mathbb A$, these differences define a Cech class in $H^1(S,\mathbb A)$. Changing the chosen primitives changes it by a locally constant coboundary. This defines the canonical linear map \eqref{cdp:eq:A.20}.

\noindent\textbf{Step 7.} If the resulting class is zero, adjust the local primitives by locally constant sections so that they glue. Their pullback is then a global holomorphic section of the unitary-flat line on each compact component of $Z$. Its squared norm is plurisubharmonic, so the maximum principle makes the section parallel; its covariant differential is zero. Consequently $\alpha_Z=0$, and injectivity of pullback gives $\alpha=0$. This proves \eqref{cdp:eq:A.20}.
\end{proof}






\begin{proposition}[The torsion-free cotangent assertion for arbitrary singular supports in a torus fibration]\label{cdp:st:A.4.8}
\StatementRecord{cdp:st:A.4.8}{proved}{Historical events: Claude Fable 5.2 (Anthropic), effort high, 16 November 2026: carried as background without fault found. GPT-6 Astra (OpenAI), ChatGPT Pro mode, effort Pro, S6\_09\_PAPERS, 17 November 2026: targeted proof check. Also checked in this writing. Pre-continuation writing events: GPT-6 Astra (OpenAI), ChatGPT Pro mode, effort Pro, S6\_10\_PAPERS, 18 November 2026: targeted proof check. Continuation metadata events: GPT-6 Astra (OpenAI), ChatGPT Pro mode, effort Pro, S6\_10\_PAPERS continuation, 19 November 2026: used and rechecked in this run; status-vocabulary reconciliation; metadata reconciliation only; no new mathematical check. \emph{Pre-continuation recorded proof scope.} The Koszul acyclicity of a nontrivial torus character, injection from the tube into its complement in degree one, and proper base change were re-derived}{the hypotheses, notation and definitions stated here; \cref{cdp:st:D.1.9}; \cref{cdp:st:D.1.8}; \cref{cdp:eq:A.22}; \cref{cdp:st:A.4.7}}

Let $g:Y\to\PP^1$ be a proper surjective holomorphic map with connected fibres from a smooth compact complex threefold with $b_1(Y)=b_2(Y)=0$, and suppose its general fibre is a complex two-torus. Let $L_r=\exp(2\pi\mathrm i r\beta)\in\Pic^0(Y)$ be the transgression family constructed in \ref{cdp:st:D.1.9}. For every real irrational $r$ and every reduced fibre support $S=(Y_c)_{\rm red}$,
\[
 H^1(S,\mathbb A_{r,S})=0,
 \qquad
 H^0(S,\widetilde\Omega_S^1\otimes L_r|_S)=0,
 \tag{A.21}\label{cdp:eq:A.21}
\]
where $\mathbb A_{r,S}$ is the unitary local system whose associated holomorphic line is $L_r|_S$. In particular this proves Equation~(11) of CDP20 for all singular reduced fibre supports in the torus case, with no assumption on their singularities or their scheme-theoretic multiplicities.
\end{proposition}


\begin{proof}
\noindent\textbf{Step 1.} We specify all comparison maps used from \ref{cdp:st:D.1.9}. On a sufficiently small disc $\Delta$ about $c$, put $U=g^{-1}(\Delta)$ and $U^*=U\setminus S$. The integral section $\gamma$ of $R^1g_*\Z$ has a lift $\gamma_U\in H^1(U,\Z)$ by Leray over the contractible disc. Its image in $H^1(U,\OO_U)$ equals $\beta|_U$, by the coherent Leray sequence. Thus the restriction of the holomorphic line $L_r$ to $U$ is the line associated to the flat character
\[
 \pi_1(U)\longrightarrow S^1,\qquad
 a\longmapsto\exp\bigl(2\pi\mathrm i r\langle\gamma_U,a\rangle\bigr).
 \tag{A.22}\label{cdp:eq:A.22}
\]
Denote this local system by $\mathbb A_U$. \ref{cdp:st:D.1.8}--\ref{cdp:st:D.1.9} prove that $\gamma_U$ restricts nontrivially to every nearby smooth fibre. Irrationality of $r$ therefore makes \eqref{cdp:eq:A.22} nontrivial on such a fibre $F$.

\noindent\textbf{Step 2.} A nontrivial rank-one local system on the real torus underlying $F$ has zero cohomology in every degree. Explicitly, use the standard cellular cochain complex of the product of four circles with coefficients in the character. It is the Koszul complex for the four scalars $\chi(e_j)-1$. At least one scalar is nonzero; contraction in that factor, divided by this scalar, is a contracting homotopy. Therefore $H^q(F,\mathbb A_U|_F)=0$ for every $q$. The map $U^*\to\Delta^*$ is a differentiable fibre bundle. Its local-system Leray spectral sequence has zero fibre-cohomology sheaves, hence
$H^1(U^*,\mathbb A_U|_{U^*})=0$.

\noindent\textbf{Step 3.} Inclusion induces a surjection $\pi_1(U^*)\to\pi_1(U)$: a loop in the smooth real six-manifold $U$ can be perturbed off the divisor, whose real codimension is at least two, relative to a base point in $U^*$. The induced map on first cohomology with these local coefficients is injective. For an explicit verification, identify first cohomology with crossed homomorphisms modulo coboundaries. If a crossed homomorphism $a:G\to M$ becomes $a(q(h))=(\chi(q(h))-1)v$ after pullback by a surjection $q:H\to G$, surjectivity gives $a(g)=(\chi(g)-1)v$ for every $g\in G$. Thus its original class was already zero. Applied to the displayed fundamental-group map, this proves $H^1(U,\mathbb A_U)=0$.

\noindent\textbf{Step 4.} The argument applies after every sufficiently small shrinking of $\Delta$, with the restrictions of the same local system. Taking the direct limit and using proper base change identifies the resulting stalk with $H^1(S,\mathbb A_U|_S)$; scheme-theoretic nilpotents do not affect this topological local system. This proves the first vanishing in \eqref{cdp:eq:A.21}. The second is exactly \ref{cdp:st:A.4.7} applied to the restriction of \eqref{cdp:eq:A.22}. All singular values are treated by the same global parameter $r$, so the assertion is simultaneous.
\end{proof}






\begin{proposition}[CDP20 Proposition 7.3 with a complete replacement proof]\label{cdp:st:A.4.9}
\StatementRecord{cdp:st:A.4.9}{proved}{Historical events: Claude Fable 5.2 (Anthropic), effort max, 16 November 2026: targeted proof check. GPT-6 Astra (OpenAI), ChatGPT Pro mode, effort Pro, S6\_09\_PAPERS, 17 November 2026: targeted proof check. Also checked in this writing. Pre-continuation writing events: GPT-6 Astra (OpenAI), ChatGPT Pro mode, effort Pro, S6\_10\_PAPERS, 18 November 2026: targeted proof check; used and rechecked in this run. Continuation metadata events: GPT-6 Astra (OpenAI), ChatGPT Pro mode, effort Pro, S6\_10\_PAPERS continuation, 19 November 2026: used and rechecked in this run; status-vocabulary reconciliation; metadata reconciliation only; no new mathematical check. \emph{Pre-continuation recorded proof scope.} The torus Koszul complex, finite-cover transfer, Kodaira central extension, Hopf cyclic cover and Inoue mapping-torus matrices give degree-zero and degree-one vanishing for every irrational nonzero integral character on a chosen very general smooth fibre. Proper smooth transport along a path in the connected smooth locus gives an actual diffeomorphism d and the equality of character local systems under d-pullback, hence the same vanishings on every smooth fibre. The tube fundamental-group map, proper base change and finite critical-value Picard argument then prove the stated conclusion. SS2002 section1 pp2-4 and Albanese2020 sections2.1.1-2.1.2 pp3-6 were opened at the exact model loci}{the hypotheses, notation and definitions stated here; \cref{cdp:st:D.1.9}; \cref{cdp:st:C.3.1}; \cref{cdp:st:A.4.8}; \cite[\S1 and the central-extension diagram in the proof of Lemma~1.1]{SS2002}; \cite[\S2.1.1 and Theorem~2.2]{Albanese2020}; \cite[\S2.1.2 and Proposition~2.3]{Albanese2020}; \cref{cdp:st:D.1.8}; \cref{cdp:st:A.4.7}; \cref{cdp:st:A.4.2}; \cref{cdp:eq:A.24}}

Let $Y$ be a connected smooth compact complex threefold with $H^1(Y,\Z)=H^2(Y,\Z)=0$, algebraic dimension one and holomorphic algebraic reduction $g:Y\to C$. There is a dense open subset of $\Pic(Y)$, whose complement is a proper complex analytic subset, such that
\[
 H^0\bigl((Y_c)_{\rm red},
       \widetilde\Omega^1_{(Y_c)_{\rm red}}\otimes L|_{(Y_c)_{\rm red}}\bigr)=0
 \tag{A.23}\label{cdp:eq:A.23}
\]
for every singular reduced fibre support and every $L$ in that subset. In fact the transgression twists $L_r$ of \ref{cdp:st:D.1.9} satisfy \eqref{cdp:eq:A.23} on every reduced fibre support, including smooth and reducible supports, for every real irrational $r$.
\end{proposition}


\begin{proof}
\noindent\textbf{Step 1.} Choose one very general smooth fibre $F_0$. The classification \ref{cdp:st:C.3.1} places this fibre among complex tori, bielliptic surfaces, primary or secondary Kodaira surfaces, and primary or secondary Hopf surfaces or Inoue surfaces. We first prove that for each surface $F$ in this list and each nonzero $\gamma_F\in H^1(F,\Z)$ the local system
\[
 \mathbb A_{r,F}:\quad a\longmapsto
       \exp(2\pi\mathrm i r\langle\gamma_F,a\rangle)
 \quad(r\in\mathbb R\setminus\Q)
\]
satisfies $H^0(F,\mathbb A_{r,F})=H^1(F,\mathbb A_{r,F})=0$. Its character has infinite image, so the assertion in degree zero is immediate.

\noindent\textbf{Step 2.} For a torus the Koszul argument in \ref{cdp:st:A.4.8} applies. A bielliptic surface has a finite étale torus cover. Pullback of the integral class is nonzero: transfer composed with pullback is multiplication by the degree. The character stays nontrivial after pullback, and transfer with local coefficients makes pullback on cohomology injective over $\C$. Thus degree-one vanishing descends from the torus cover.

\noindent\textbf{Step 3.} A primary Kodaira surface is a principal elliptic bundle over an elliptic curve, as described in \cite[\S1 and the central-extension diagram in the proof of Lemma~1.1]{SS2002}. Its fundamental group has a central extension
\[
 1\longrightarrow\Z^2\longrightarrow\Gamma\longrightarrow\Z^2
       \longrightarrow1.
 \tag{A.24}\label{cdp:eq:A.24}
\]
For any nontrivial one-dimensional character $\chi$ of $\Gamma$, if its restriction to the central subgroup is nontrivial, the latter's cohomology with these coefficients is zero by the two-variable Koszul contraction. The Hochschild--Serre spectral sequence then gives $H^1(\Gamma,\C_\chi)=0$. If its central restriction is trivial, $\chi$ descends to a nontrivial character of the quotient $\Z^2$. Centrality makes its action on the cohomology of the central subgroup equal to that scalar character, and the same Koszul contraction kills every term on the second page. Again $H^1(\Gamma,\C_\chi)=0$. First cohomology of a connected space with local coefficients is first cohomology of its fundamental group, so this proves the required assertion. A secondary Kodaira surface has a finite primary cover; the same transfer argument as above applies. Hopf surfaces have finite primary covers, and a primary Hopf surface is diffeomorphic to $S^1\times S^3$; the cover statement and topology are recorded in \cite[\S2.1.1 and Theorem~2.2]{Albanese2020}. The nontrivial circle character has an invertible cellular differential, giving the required vanishing and its finite-cover descent.

\noindent\textbf{Step 4.} For completeness we also give the Inoue calculation, rather than assuming a generic-vanishing theorem for its cotangent bundle. The explicit constructions in \cite[\S2.1.2 and Proposition~2.3]{Albanese2020} identify every Inoue surface with a mapping torus $N_\varphi$ of a compact three-manifold $N$. In the $S_M$ cases $N$ is a three-torus and the eigenvalues of $\varphi^*$ on $H^1(N,\C)$ are $\alpha,\beta,\bar\beta$, with $\alpha>1$ and $|\beta|=\alpha^{-1/2}<1$, after reversing the circle orientation when necessary. In the $S^+$ and $S^-$ cases $N$ is the nontrivial circle bundle over a two-torus; the nonzero Euler class makes $H^1(N,\C)$ the pullback of $H^1(T^2,\C)$ by its Gysin sequence. The monodromy on it has eigenvalues $\alpha,\alpha^{-1}$ or $\alpha,-\alpha^{-1}$, respectively. These assertions follow as well directly from the conjugation relations displayed in that construction: on the abelian rank-three kernel the matrix is $M$, and on the quotient of the Heisenberg kernel by its centre the matrix is $N\in\mathrm{GL}_2(\Z)$ with precisely those eigenvalues.

\noindent\textbf{Step 5.} None of these degree-one eigenvalues has modulus one. In particular there are no invariant degree-one classes on $N$. The ordinary mapping-torus sequence therefore identifies $H^1(F,\Z)$ with the copy of $\Z$ from the circle, since the invariant subgroup of the torsion-free group $H^1(N,\Z)$ is zero. Thus $\mathbb A_{r,F}$ is pulled back from the circle, with monodromy $\lambda\in S^1$ of infinite order. The local-coefficient Leray sequence for $N\to N_\varphi\to S^1$ gives
\[
\begin{split}
0\longrightarrow
 \coker(\lambda-1:H^0(N,\C)\to H^0(N,\C))
 \longrightarrow H^1(F,\mathbb A_{r,F})\\
 \longrightarrow
 \ker(\lambda\varphi^*-1:H^1(N,\C)\to H^1(N,\C))
 \longrightarrow0.
\end{split}
\]
The first term is zero because $\lambda\ne1$, and the last term is zero by the eigenvalue moduli. Using the reciprocal convention for the circle local system replaces $\lambda$ by $\lambda^{-1}$ and gives the same conclusion. This completes the degree-one calculation for every surface type in \ref{cdp:st:C.3.1}.

\noindent\textbf{Step 6.} Let $F_1$ be any smooth fibre. The smooth locus in the connected base curve is the complement of finitely many points and is connected. A path from the value of $F_0$ to that of $F_1$, together with proper smooth local triviality, gives an actual diffeomorphism $d:F_0\to F_1$. For any nonzero $\gamma_1\in H^1(F_1,\Z)$, the class $d^*\gamma_1$ is nonzero and the character formula gives the actual local-system isomorphism
\[
 d^*\mathbb A_{r,F_1,\gamma_1}
       \simeq\mathbb A_{r,F_0,d^*\gamma_1}.
\]
The induced cohomology isomorphism identifies degrees zero and one on $F_1$ with the groups just proved zero on $F_0$. This proves both vanishings on every smooth fibre. No constancy of its complex surface type is needed for these maps.

\noindent\textbf{Step 7.} Now repeat the neighbourhood argument of \ref{cdp:st:A.4.8} with these two fibre-cohomology vanishings. \ref{cdp:st:D.1.8}--\ref{cdp:st:D.1.9} supply the same nonzero integral class on every nearby smooth fibre and the unitary local system on every fibre neighbourhood. The only terms that can contribute to first cohomology of the punctured neighbourhood are its degree-zero and degree-one fibre cohomology, just proved zero. The fundamental-group surjection and proper base change therefore give $H^1((Y_c)_{\rm red},\mathbb A_{r,c})=0$ for every $c$. Applying \ref{cdp:st:A.4.7} proves \eqref{cdp:eq:A.23} for the single global twist $L_r$, simultaneously on all reduced supports.

\noindent\textbf{Step 8.} Finally, the exponential sequence and the integral vanishings identify $\Pic(Y)$ with the complex vector space $H^1(Y,\OO_Y)$. Choose holomorphic Cech cocycles representing a basis; exponentiating their linear combinations gives a holomorphic family of line bundles on $Y\times H^1(Y,\OO_Y)$ whose fibre is the indicated line bundle. There are finitely many singular fibre supports. For each of them, tensor its fixed coherent torsion-free cotangent sheaf by this family and apply proper semicontinuity: the locus where its space of sections is nonzero is a closed complex analytic subset of the parameter space. Each is proper because it excludes $L_r$. Their finite union is proper and has empty interior, giving the asserted dense open set. This proves the statement of Proposition~7.3. The invalid genus-two deduction in its printed proof remains refuted by \ref{cdp:st:A.4.2}; the proof above replaces that argument through the explicit maps (A.20)--\eqref{cdp:eq:A.24}.
\end{proof}






\section{Normalization, conductor genus and local torus exclusions}\label{sec:II:05}


\begin{proposition}[What trivial normalization pullback forces]\label{cdp:st:A.5.1}
\StatementRecord{cdp:st:A.5.1}{proved}{Historical events: Claude Fable 5.2 (Anthropic), effort high, 16 November 2026: carried as background without fault found. GPT-6 Astra (OpenAI), ChatGPT Pro mode, effort Pro, S6\_09\_PAPERS, 17 November 2026: targeted proof check. Also checked in this writing. Pre-continuation writing events: GPT-6 Astra (OpenAI), ChatGPT Pro mode, effort Pro, S6\_10\_PAPERS, 18 November 2026 04:23:40Z: targeted proof check. Continuation metadata events: GPT-6 Astra (OpenAI), ChatGPT Pro mode, effort Pro, S6\_10\_PAPERS continuation, 19 November 2026: used and rechecked in this run; status-vocabulary reconciliation; metadata reconciliation only; no new mathematical check. \emph{Pre-continuation recorded proof scope.} proved. Full normalization-section nonvanishing with canonical-subspace wording repaired. The intrinsic injection is H0(normalization,eta*(N*\ensuremath{\otimes}A))\ensuremath{\to}H0(S,V\ensuremath{\otimes}A). A base coordinate and a trivialization of eta*A identify its domain with H0(O\_normalization); the image subspace is independent of those choices. Every compact normalized component contributes a nonzero constant. The owning agent has implemented this exact canonical-source formulation. Review provenance: New derivation after reading the supplied proof; not a blind review and not an independent human check}{the hypotheses, notation and definitions stated here; \cref{cdp:st:A.1.1}; \cref{cdp:st:A.4.1}}

Let $S$ be a connected reduced full fibre of a proper map from a complex manifold to a smooth curve. For every line bundle $A$ on $S$ with $\eta^*A\simeq\OO_{\widetilde S}$, the ambient group contains the canonical image of $H^0(\widetilde S,\eta^*(N^*\otimes A))$. A local base coordinate and a choice of trivialization of $\eta^*A$ identify this subspace with $H^0(\widetilde S,\OO_{\widetilde S})$. In particular it is nonzero. This applies to the subcase of Proposition~7.3 in which all $\eta^*(L_t|_S)$ are trivial. It is compatible with the vanishing of $H^0(S,Q\otimes A)$ through the exact maps of \ref{cdp:st:A.1.1}, as \ref{cdp:st:A.4.1} demonstrates.
\end{proposition}


\begin{proof}
\noindent\textbf{Step 1.} Tensoring the canonical injection in (A.3) by $A$ and applying the projection formula gives
$\eta_*\eta^*(N^*\otimes A)\hookrightarrow K\otimes A\hookrightarrow V\otimes A$.
Taking global sections preserves injectivity and defines the asserted canonical subspace. A local coordinate of the base trivializes $N^*$; together with a choice of trivialization of $\eta^*A$, it identifies the source with $H^0(\widetilde S,\OO_{\widetilde S})$. Every connected component of the compact normalization has the nonzero constant function $1$, proving nonvanishing. No normal-crossings assumption is needed for this inclusion.
\end{proof}



The line in the printed proof asserting $H^2(\widetilde S,\omega_{\widetilde S})=0$ also cannot hold for a connected compact normal Cohen--Macaulay surface with its dualizing sheaf: Serre duality identifies its dual with $H^0(\widetilde S,\OO)=\C$. Replacing $\omega_{\widetilde S}$ there by $\OO_{\widetilde S}$ matches the vanishing stated at the beginning of that proof, but does not settle its omitted genus-one case. We have established (11) for $W$ and the exact obstruction to deducing ambient vanishing from it. The full statement of Proposition~7.3 is proved below in \ref{cdp:st:A.4.9} by a different argument through unitary local systems. The following argument establishes Proposition~4.11 independently of the monodromy-dependent restrictions.




\begin{proposition}[Normal Gorenstein surfaces with vanishing first Chern class and Euler characteristic]\label{cdp:st:A.5.2}
\StatementRecord{cdp:st:A.5.2}{proved}{Historical events: Claude Fable 5.2 (Anthropic), effort high, 16 November 2026: carried as background without fault found. Claude Fable 5.2 (Anthropic), effort max, 16 November 2026: targeted proof check. GPT-6 Astra (OpenAI), ChatGPT Pro mode, effort Pro, S6\_09\_PAPERS, 17 November 2026: targeted proof check. Also checked in this writing. Pre-continuation writing events: GPT-6 Astra (OpenAI), ChatGPT Pro mode, effort Pro, S6\_10\_PAPERS, 18 November 2026: targeted proof check. Continuation metadata events: GPT-6 Astra (OpenAI), ChatGPT Pro mode, effort Pro, S6\_10\_PAPERS continuation, 19 November 2026: used and rechecked in this run; status-vocabulary reconciliation; metadata reconciliation only; no new mathematical check. \emph{Pre-continuation recorded proof scope.} The ruled-resolution Euler contradiction, rational singularity case, crepant negative-definite discrepancy calculation, and the c1=c2=0 classification consequences were re-derived}{the hypotheses, notation and definitions stated here; \cite[Theorem~2.2]{Peters2021}}

A compact connected normal Gorenstein complex surface $S$ with $c_1(\omega_S)=0$ in $H^2(S,\Q)$ and $\chi(S,\OO_S)=0$ is smooth. Consequently every singular irreducible reduced divisor in a compact complex threefold $Y$ with $H^2(Y,\Q)=0$ is nonnormal. In particular the assertion of CDP20 Proposition~4.11 holds without its monodromy-dependent restrictions on the smooth fibre.
\end{proposition}


\begin{proof}
\noindent\textbf{Step 1.} Let $\pi:Z\to S$ be the minimal resolution and set $\ell=h^0(S,R^1\pi_*\OO_Z)$. Normality gives $\pi_*\OO_Z=\OO_S$, and the fibres have dimension at most one, so Leray gives $\chi(Z,\OO_Z)=-\ell$. We first prove $\ell=0$. The classification table in \cite[Theorem~2.2]{Peters2021}, together with Noether's formula and invariance of $\chi(\OO)$ under blowing up, shows that a smooth compact surface of negative holomorphic Euler characteristic is a blown-up ruled surface over a smooth curve $B$ of genus $g\ge2$. Indeed the sums $c_1^2+c_2$ in every other row of that table are nonnegative. Thus, when $\ell>0$, there is a ruling $p:Z\to B$ and $\chi(Z,\OO_Z)=1-g$, giving $\ell=g-1$.

\noindent\textbf{Step 2.} An exceptional irreducible curve $E$ of $\pi$ dominating $B$ would have $h^1(E,\OO_E)\ge g$: normalize $E$, apply Riemann--Hurwitz to its finite map to $B$, and use the normalization exact sequence on $E$. The exact sequence $0\to\mathcal I_E\to\OO_Z\to\OO_E\to0$, pushed forward by $\pi$, gives a surjection $R^1\pi_*\OO_Z\to R^1\pi_*\OO_E$ because $R^2\pi_*\mathcal I_E=0$. It follows that $\ell\ge h^1(E,\OO_E)\ge g$, a contradiction. Hence every exceptional curve is vertical for $p$. The ruling is constant on each connected fibre of $\pi$ and descends to a holomorphic map $p_S:S\to B$: local holomorphic coordinate functions descend by $\pi_*\OO_Z=\OO_S$, and properness supplies the required neighborhoods. The map $p$ is the composite of blowups and a $\PP^1$-bundle, so $R^1p_*\OO_Z=0$. The Leray sequence for $p=p_S\circ\pi$ then gives
\[
0=R^1p_*\OO_Z\longrightarrow (p_S)_*R^1\pi_*\OO_Z
 \longrightarrow R^2(p_S)_*\OO_S=0.
\]
A nonzero sheaf supported at finitely many points cannot have zero direct image, so again $\ell=0$, a contradiction. This proves rationality of all singularities.

\noindent\textbf{Step 3.} We give the crepancy argument to specify the comparison of canonical bundles. Proper duality applied to $R\pi_*\OO_Z=\OO_S$ identifies $\pi_*\omega_Z$ with the invertible dualizing sheaf $\omega_S$. A local generator of $\omega_S$ therefore pulls back to a holomorphic two-form on $Z$, agreeing with the usual pullback off the exceptional set. Its zero divisor is effective and exceptional. These divisors glue to $D\ge0$ with $K_Z=\pi^*K_S+D$. Every exceptional irreducible curve has arithmetic genus zero by the same surjection to its first cohomology, hence is a smooth rational curve. Minimality of the resolution gives $E^2\le-2$, and adjunction gives $D\cdot E=K_Z\cdot E=-2-E^2\ge0$. The exceptional intersection matrix is negative definite. Thus a nonzero $D$ would satisfy both $D^2<0$ and $D^2=\sum_E d_E(D\cdot E)\ge0$. Therefore $D=0$.

\noindent\textbf{Step 4.} Consequently $c_1(K_Z)=0$ in rational cohomology. It follows by adjunction that $Z$ contains no $(-1)$-curve, so $Z$ is a minimal compact surface; Noether's formula gives $c_2(Z)=0$. We exclude rational curves on each surviving row of the classification. A ruled surface is excluded by $K_Z\cdot F=-2$ on a ruling fibre. General-type surfaces have $K_Z^2>0$ and K3 and Enriques surfaces have positive $c_2$, so none occurs. A torus has no rational curve because a map from $\PP^1$ lifts to its vector-space universal cover and is constant; the same holds for a finite etale quotient of a torus. A primary Kodaira surface is an elliptic bundle over an elliptic curve, so a rational curve would map constantly to the base and then constantly to a fibre; finite etale quotients have the same property. A class VII surface here has $b_2=c_2=0$, and the Hopf--Inoue classification gives no rational curves. Equivalently one can use their standard universal covers: $\C^2\setminus\{0\}$ for a primary Hopf cover and $\mathbb H\times\C$ for an Inoue surface. A map from $\PP^1$ lifts and has constant coordinate functions.

\noindent\textbf{Step 5.} It remains to check the properly elliptic row. For a minimal elliptic surface, the Kodaira fibre list expresses $c_2$ as a sum of nonnegative fibre Euler numbers, with zero contributions exactly at multiples of smooth elliptic curves. Since $c_2(Z)=0$, all fibre reductions are smooth elliptic curves. The canonical bundle formula on a Kähler elliptic surface with $\chi(\OO_Z)=0$ has rational first Chern class
\[
 c_1(K_Z)=\left(2g(B)-2+\sum_i(1-m_i^{-1})\right)[F].
\]
In Kodaira dimension one the coefficient is positive, and the Kähler class has positive integral over $F$, contradicting $c_1(K_Z)=0$. In the non-Kähler case, a horizontal rational curve would be a multisection $C$. For $D=C+nF$ with $n$ sufficiently large, $D^2>0$ and $(K_Z-mD)\cdot F<0$ for $m>0$. Riemann--Roch therefore gives quadratic growth of $h^0(Z,\OO_Z(mD))$, making $Z$ Moishezon and hence projective as a smooth compact surface. This contradicts the non-Kähler case. Vertical rational curves have already been excluded by the fibre list. Thus $Z$ has no rational curves in any case. The exceptional curves of $\pi$ were rational, so there are none and $S$ is smooth.

\noindent\textbf{Step 6.} Finally take an irreducible reduced divisor $S\subset Y$ with $H^2(Y,\Q)=0$. It is a Cartier divisor in the smooth manifold $Y$, hence Gorenstein, and adjunction gives $\omega_S=(K_Y\otimes\OO_Y(S))|_S$, with rational first Chern class zero. The divisor sequence and Hirzebruch--Riemann--Roch give
$\chi(\OO_S)=\chi(\OO_Y)-\chi(\OO_Y(-S))=0$, since $c_1(\OO_Y(S))=0$ rationally. A normal $S$ is therefore smooth by the result just proved.
\end{proof}






\begin{proposition}[The effective conductor divisor on a minimal resolution]\label{cdp:st:A.5.3}
\StatementRecord{cdp:st:A.5.3}{proved}{Historical events: Claude Fable 5.2 (Anthropic), effort max, 16 November 2026: targeted proof check. GPT-6 Astra (OpenAI), ChatGPT Pro mode, effort Pro, S6\_09\_PAPERS, 17 November 2026: targeted proof check. Also checked in this writing. Pre-continuation writing events: GPT-6 Astra (OpenAI), ChatGPT Pro mode, effort Pro, S6\_10\_PAPERS, 18 November 2026: targeted proof check. GPT-6 Astra (OpenAI), ChatGPT Pro mode, effort Pro, S6\_10\_PAPERS, 18 November 2026 05:35:21+00:00: used and rechecked in this run. Continuation metadata events: GPT-6 Astra (OpenAI), ChatGPT Pro mode, effort Pro, S6\_10\_PAPERS continuation, 19 November 2026: used and rechecked in this run; status-vocabulary reconciliation; metadata reconciliation only; no new mathematical check. \emph{Pre-continuation recorded proof scope.} Finite-duality conductor effectiveness, the exceptional intersection-matrix equation, and the horizontal degree-two calculation were re-derived}{the hypotheses, notation and definitions stated here; \cref{cdp:st:A.2.1}; \cite[Lemma~1.1 and Corollary~1.4]{Furushima1993}}

Let $S$ be a compact reduced hypersurface in a smooth threefold, let $\eta:\coprod_i T_i\to S$ be its normalization, and let $\pi_i:Z_i\to T_i$ be the minimal resolutions. Put $h_i=\eta\pi_i$. There are canonical effective integral divisors $D_i$ on $Z_i$ with
\[
 \omega_{Z_i}(D_i)\simeq h_i^*\omega_S.
 \tag{A.25}\label{cdp:eq:A.25}
\]
The nonexceptional part of $D_i$ is exactly the divisorial conductor, with its scheme-theoretic multiplicities. When $S\subset Y$ and $H^2(Y,\Q)=0$, every $K_{Z_i}+D_i$ has rational first Chern class zero. A ruled $Z_i\to B$ with $g(B)\ge2$ then has
\[
 D_i=2C+V,
 \tag{A.26}\label{cdp:eq:A.26}
\]
where $C$ is a section and $V$ is vertical. In the case where $C$ is $\pi_i$-exceptional, $\eta^*(L|_S)|_{T_i}$ is trivial for every $L\in\Pic(Y)$.
\end{proposition}


\begin{proof}
\noindent\textbf{Step 1.} Finite duality gives $\omega_{T_i}=\mathfrak c_i\otimes\eta^*\omega_S$ as rank-one reflexive sheaves, where $\mathfrak c_i$ is the conductor viewed as an ideal on $T_i$. Locally this is the equality
$\Hom_R(B,R)=\{q:qB\subset R\}=\mathfrak c$ for a reduced Gorenstein ring $R$ and its finite normalization $B$, tensored with the invertible dualizing module. The equality uses the meromorphic multiplication identification as in \ref{cdp:st:A.2.1}. The conductor is a maximal Cohen--Macaulay module because it is this dual of the maximal Cohen--Macaulay normalization over the Gorenstein ring. On a normal surface a rank-one maximal Cohen--Macaulay module is reflexive, so the displayed equality is determined by its codimension-one orders. Those orders define an effective conductor divisor $D_i^0$ on $T_i$.

\noindent\textbf{Step 2.} Off the exceptional set the canonical identification pulls back to $\omega_{Z_i}(D_i^{\rm str})\simeq h_i^*\omega_S$, where $D_i^{\rm str}$ denotes strict transform. The meromorphic extension of this line-bundle identification gives integral exceptional coefficients $a_j$ with
$K_{Z_i}+D_i^{\rm str}+\sum_j a_jE_j=h_i^*K_S$.
We verify their signs. For each exceptional irreducible curve $E_j$, minimality excludes a smooth rational $(-1)$-curve. Adjunction therefore gives
$K_{Z_i}\cdot E_j=2p_a(E_j)-2-E_j^2\ge0$;
also $D_i^{\rm str}\cdot E_j\ge0$. The exceptional intersection matrix $M=(E_j\cdot E_k)$ is negative definite with nonnegative off-diagonal entries, and
\[
 M a=-v,\qquad v_j=(K_{Z_i}+D_i^{\rm str})\cdot E_j\ge0.
\]
To see that $a\ge0$ without an unproved sign convention, write $a=a_+-a_-$ with disjoint nonnegative coordinate supports. If $a_-\ne0$, then
$a_-^tMa=a_-^tMa_+-a_-^tMa_->0$,
whereas $Ma=-v$ gives $a_-^tMa\le0$. This contradiction proves $a_j\ge0$. Thus $D_i=D_i^{\rm str}+\sum_j a_jE_j$ proves \eqref{cdp:eq:A.25}. The conductor duality and effective exceptional comparison also appear, for projective surfaces with trivial dualizing sheaf, in \cite[Lemma~1.1 and Corollary~1.4]{Furushima1993}; the argument here supplies the maps and does not use that paper's additional K3-type hypothesis $h^1(\OO_S)=0$.

\noindent\textbf{Step 3.} For a divisor in $Y$, adjunction gives $\omega_S=(K_Y\otimes\OO_Y(S))|_S$, proving the numerical assertion. On a general ruling fibre $f$, $K_{Z_i}\cdot f=-2$, so $D_i\cdot f=2$. Every horizontal irreducible curve has normalization of genus at least $g(B)\ge2$. Such a curve $C_0$ cannot occur in $D_i$ with coefficient one: otherwise
\[
 0=(K_{Z_i}+D_i)\cdot C_0
   =2p_a(C_0)-2+(D_i-C_0)\cdot C_0>0.
\]
Every horizontal coefficient is consequently at least two. The degree-two intersection with $f$ forces exactly one horizontal component, coefficient two and degree one over $B$, giving \eqref{cdp:eq:A.26}. The degree-one finite morphism $C\to B$ is an isomorphism because $B$ is normal; thus $C$ is a smooth section.

\noindent\textbf{Step 4.} Finally, integral $H^2$ of a smooth ruled surface and of its blowups is torsion-free, and its topologically trivial line bundles are precisely the pullbacks of $\Pic^0(B)$. These facts follow from the projective-bundle and blowup cohomology formulas and from $p_*\OO_{Z_i}=\OO_B$, $R^1p_*\OO_{Z_i}=0$. Restriction to the section $C$ is inverse to this Picard pullback. The line $h_i^*L$ has torsion, hence zero, integral first Chern class. When $C$ is exceptional, its restriction to $C$ is trivial because $h_i$ maps $C$ to a point. Therefore $h_i^*L$ is trivial on $Z_i$. Normality gives $(\pi_i)_*\OO_{Z_i}=\OO_{T_i}$, and the projection formula descends this triviality to $T_i$. This proves the last assertion.
\end{proof}






\begin{proposition}[The remaining normalization class for a persistent ambient obstruction]\label{cdp:st:A.5.4}
\StatementRecord{cdp:st:A.5.4}{proved}{Historical events: Claude Fable 5.2 (Anthropic), effort max, 16 November 2026: targeted proof check. GPT-6 Astra (OpenAI), ChatGPT Pro mode, effort Pro, S6\_09\_PAPERS, 17 November 2026: targeted proof check. Also checked in this writing. Pre-continuation writing events: GPT-6 Astra (OpenAI), ChatGPT Pro mode, effort Pro, S6\_10\_PAPERS, 18 November 2026: targeted proof check; used and rechecked in this run. Continuation metadata events: GPT-6 Astra (OpenAI), ChatGPT Pro mode, effort Pro, S6\_10\_PAPERS continuation, 19 November 2026: used and rechecked in this run; status-vocabulary reconciliation; metadata reconciliation only; no new mathematical check. \emph{Pre-continuation recorded proof scope.} Selection of a persistent component, the Gauduchon-degree argument, the primitive-curvature unitary-flat conclusion, and the Picard-image exclusions were re-derived}{the hypotheses, notation and definitions stated here; \cref{cdp:st:A.5.3}; \cref{cdp:st:A.4.9}; \cref{cdp:st:A.1.1}; \cite[Lemma~2.3(1)--(2)]{Teleman2006}; \cite[Lemma~2.3(2)]{Teleman2006}}

Let $Y$ satisfy the CDP hypotheses, and let $S$ be a connected nonnormal reduced fibre support for which
\[
 H^0(S,\Omega_Y^1|_S\otimes L|_S)\ne0
 \quad\hbox{for every }L\in\Pic(Y).
 \tag{A.27}\label{cdp:eq:A.27}
\]
There is a normalization component $T_i$ whose minimal resolution $Z_i$ has Kodaira dimension $-\infty$. Every line from $\Pic(Y)$ pulls back to a unitary-flat line on that $Z_i$. A rational or class VII $Z_i$ has trivial Picard pullback from $Y$, and then
\[
 H^2\!\left(T_i,\omega_{T_i}\otimes
             \eta^*((K_Y^{-1}\otimes L)|_S)|_{T_i}\right)\simeq\C
 \quad\hbox{for every }L\in\Pic(Y).
 \tag{A.28}\label{cdp:eq:A.28}
\]
The same conclusion holds for the exceptional-section case of \ref{cdp:st:A.5.3}. Thus the only normalization type left by this argument, when \eqref{cdp:eq:A.28} has not been obtained, is an irrational ruled surface with nonconstant Picard pullback; in base genus at least two its divisor (A.26) has its horizontal section in the strict conductor.
\end{proposition}


\begin{proof}
\noindent\textbf{Step 1.} \ref{cdp:st:A.4.9} gives $H^0(S,\widetilde\Omega_S^1\otimes L|_S)=0$ on a dense open subset of $\Pic(Y)$. By \ref{cdp:st:A.1.1} and \eqref{cdp:eq:A.27}, $H^0(S,K\otimes L|_S)$ is nonzero there. Its nonvanishing locus is closed analytic in the holomorphic Picard family, hence is all of $\Pic(Y)$.

\noindent\textbf{Step 2.} Pull $K$ to each $Z_i$, quotient the resulting sheaf by its torsion, and take its double dual. The result is a fixed holomorphic line bundle $M_i$, since its generic rank on each component is one and $Z_i$ is a smooth surface. The induced maps give an injection
\[
 H^0(S,K\otimes L|_S)\hookrightarrow
       \bigoplus_i H^0(Z_i,M_i\otimes h_i^*L).
 \tag{A.29}\label{cdp:eq:A.29}
\]
Indeed a section in the kernel vanishes generically on every component, and $K\subset\Omega_Y^1|_S$ is torsion-free over the reduced space. Each summand's nonvanishing locus is closed analytic in the connected complex vector space $\Pic(Y)=H^1(Y,\OO_Y)$. A finite union of proper analytic subsets cannot cover that space. Consequently one component $i$ has
$H^0(Z_i,M_i\otimes h_i^*L)\ne0$ for every $L$.

\noindent\textbf{Step 3.} Fix a Gauduchon metric on $Z_i$, with fundamental form $\omega$ satisfying $\partial\bar\partial\omega=0$. The degree of a line bundle is the integral of its real Chern curvature against $\omega$; it is independent of the chosen Hermitian metric and additive under tensor product. A nonzero section has an effective zero divisor, so its line has nonnegative degree. For any $G\in\Pic(Y)$ the preceding nonvanishing applied to $G^{\otimes n}$ therefore gives
\[
 \deg_\omega(M_i)+n\deg_\omega(h_i^*G)\ge0
 \quad\hbox{for every }n\in\Z.
\]
It follows that $\deg_\omega(h_i^*G)=0$ for every $G$ and every Gauduchon metric.

\noindent\textbf{Step 4.} We record the precise flatness implication. For a topologically trivial line, the real Chern curvature $F$ is exact. If $\int F\wedge\omega=0$, the Gauduchon scalar elliptic equation supplies a real $u$ with $\Lambda_\omega(F-i\partial\bar\partial u)=0$. The remaining form is closed, exact, real and primitive of type $(1,1)$, hence anti-selfdual. Its square integrates both to zero by exactness and to minus its squared norm. It is therefore zero. Changing the line metric by $e^{-u}$ gives zero Chern curvature, hence a flat unitary connection. This is exactly the curvature argument of \cite[Lemma~2.3(1)--(2)]{Teleman2006}. Since $H^2(Y,\Z)=0$, it applies to every $h_i^*G$.

\noindent\textbf{Step 5.} The divisor $D_i$ of \ref{cdp:st:A.5.3} is nonzero. For an irreducible nonnormal support this follows from its nonempty conductor; for a connected reducible support every component meets another component, again giving a nonempty conductor. If $Z_i$ had nonnegative Kodaira dimension, some positive pluricanonical power would have a nonzero section, making $\deg_\omega(K_{Z_i})\ge0$. But $D_i$ is a nonzero effective divisor and therefore has strictly positive degree. Equation (A.25), with $\omega_S=(K_Y\otimes\OO_Y(S))|_S$, would give
\[
 0=\deg_\omega h_i^*(K_Y\otimes\OO_Y(S))
   =\deg_\omega(K_{Z_i})+\deg_\omega(D_i)>0,
\]
a contradiction. Thus $\kappa(Z_i)=-\infty$, whose classification consists of rational, irrational ruled, and class VII surfaces.

\noindent\textbf{Step 6.} For a rational $Z_i$, topologically trivial lines are trivial. For a class VII $Z_i$, $h^1(\OO_{Z_i})=1$ and $b_1(Z_i)=1$. With our exponential convention $u\mapsto\exp(2\pi\mathrm i u)$, the unitary-flat subgroup of $\Pic^0(Z_i)$ has real tangent space equal to the injective image of $H^1(Z_i,\mathbb R)$ in $H^1(Z_i,\OO_{Z_i})$. This is \cite[Lemma~2.3(2)]{Teleman2006} after multiplication by the fixed scalar $\mathrm i$, or follows from the pluriharmonic maximum principle. That real tangent space has dimension one. The holomorphic Picard pullback has complex-linear differential and, by the proved flatness, takes values in this real one-dimensional subspace. Its differential is zero, so the homomorphism is constant and equals zero at the trivial line. Every $h_i^*L$ is therefore trivial. Projection formula and $(\pi_i)_*\OO_{Z_i}=\OO_{T_i}$ again descend triviality to $T_i$. The exceptional-section case has the same conclusion by \ref{cdp:st:A.5.3}.

\noindent\textbf{Step 7.} In each of these cases, the line multiplying $\omega_{T_i}$ in \eqref{cdp:eq:A.28} is trivial, since it too comes from $Y$. A normal complex surface is Cohen--Macaulay, so its dualizing sheaf and Serre duality give $H^2(T_i,\omega_{T_i})=H^0(T_i,\OO_{T_i})^\vee=\C$. This proves \eqref{cdp:eq:A.28} and the stated exact reduction of the remaining normalization problem.
\end{proof}






\begin{proposition}[Uniform negative-degree vanishing for a fixed vector bundle]\label{cdp:st:A.5.6}
\StatementRecord{cdp:st:A.5.6}{proved}{Historical events: Claude Fable 5.2 (Anthropic), effort max, 16 November 2026: targeted proof check. GPT-6 Astra (OpenAI), ChatGPT Pro mode, effort Pro, S6\_09\_PAPERS, 17 November 2026: targeted proof check. Also checked in this writing. Pre-continuation writing events: GPT-6 Astra (OpenAI), ChatGPT Pro mode, effort Pro, S6\_10\_PAPERS, 18 November 2026: targeted proof check; used and rechecked in this run. Continuation metadata events: GPT-6 Astra (OpenAI), ChatGPT Pro mode, effort Pro, S6\_10\_PAPERS continuation, 19 November 2026: used and rechecked in this run; status-vocabulary reconciliation; metadata reconciliation only; no new mathematical check. \emph{Pre-continuation recorded proof scope.} The constant contracted-curvature normalization, maximum-principle inequality, all-positive-and-negative-powers degree argument, exact-curvature flatness and complex-linear Picard differential were re-derived. Teleman2006 v3 Lemma2.3(1)-(2) and proof pp5-6 were opened for the exact real Bott-Chern sequence and solvable contracted-curvature equation}{the hypotheses, notation and definitions stated here; \cite[Lemma~2.3]{Teleman2006}; \cref{cdp:st:A.5.4}}

Let $Z$ be a compact smooth complex surface, $\omega$ a Gauduchon metric, and $E$ a fixed holomorphic vector bundle. There is a real constant $b(E,\omega)$ such that
\[
 \deg_\omega P<b(E,\omega)
 \quad\Longrightarrow\quad H^0(Z,E\otimes P)=0
 \tag{A.30}\label{cdp:eq:A.30}
\]
for every holomorphic line bundle $P$. Thus persistence of sections of $E\otimes h^*L$ for every $L\in\Pic(Y)$ forces $\deg_\omega h^*L=0$ for every $L$, for any holomorphic map $h:Z\to Y$. Persistence on one Picard component already forces degree zero on the pullback of $\Pic^0(Y)$. Under $H^2(Y,\Q)=0$, it forces degree zero and unitary flatness for every ambient Picard pullback. On a class VII surface $Z$ with $H^2(Y,\Z)=0$, the entire holomorphic Picard pullback is the zero map.
\end{proposition}


\begin{proof}
\noindent\textbf{Step 1.} Choose a Hermitian metric on $E$ and let $C_E$ bound the largest eigenvalue of the Hermitian curvature endomorphism $\mathrm i\Lambda_\omega F_E$ on the compact surface. Write $V_\omega=\int_Z\omega^2/2$ and normalize degree by
$\deg_\omega P=(2\pi)^{-1}\int_Z\mathrm iF_P\wedge\omega$.
The scalar Gauduchon elliptic equation changes a Hermitian line metric so that its contracted curvature is the constant
\[
 \mathrm i\Lambda_\omega F_P
       =c_P:=2\pi\deg_\omega(P)/V_\omega.
\]
Indeed the difference from this constant has integral zero, and the image of $u\mapsto\mathrm i\Lambda_\omega\partial\bar\partial u$ is precisely the mean-zero functions for a Gauduchon metric; this is the scalar equation used in \cite[Lemma~2.3]{Teleman2006}. The Hermitian line metric may therefore be rescaled by the corresponding $e^{-u}$.

\noindent\textbf{Step 2.} For a holomorphic section $s$ of $E\otimes P$, the Chern-connection identity is
\[
 \mathrm i\Lambda_\omega\partial\bar\partial|s|^2
 =|\nabla^{1,0}s|^2
  -\langle\mathrm i\Lambda_\omega F_{E\otimes P}s,s\rangle.
\]
It follows directly in a holomorphic frame by differentiating the Hermitian norm; thus no Kähler identity is being assumed. At a positive maximum of $|s|^2$, the left side is nonpositive, whereas the right side is at least
$-(C_E+c_P)|s|^2$.
If $c_P<-C_E$, this is strictly positive, a contradiction. We may take $b(E,\omega)=-C_EV_\omega/(2\pi)$, proving \eqref{cdp:eq:A.30}.

\noindent\textbf{Step 3.} Apply this bound to $P=h^*(L^{\otimes n})$ for every positive and negative integer $n$. Nonzero sections for all $n$ are incompatible with a nonzero degree, proving the persistence assertion. For persistence on a single component, absorb one representative line into the fixed bundle $E$ and apply the same argument to all powers of every line in $\Pic^0(Y)$. If $H^2(Y,\Q)=0$, a positive power of every ambient line lies in $\Pic^0(Y)$, so additivity of degree gives degree zero for every pullback. Its real first Chern class is zero, even if its integral class is torsion. The Chern curvature is therefore exact, and the exact-curvature argument in \ref{cdp:st:A.5.4} again produces a flat unitary connection. On a class VII surface the differential of the holomorphic Picard pullback then lies in a real one-dimensional unitary tangent space. Complex linearity makes that differential zero. The exponential sequence makes $\Pic(Y)$ a connected quotient of $H^1(Y,\OO_Y)$ when $H^2(Y,\Z)=0$, so under that integral hypothesis the entire map is identically zero.
\end{proof}






\begin{proposition}[The ruled obstruction with a reduced conductor image]\label{cdp:st:A.5.9}
\StatementRecord{cdp:st:A.5.9}{proved}{Historical events: Claude Fable 5.2 (Anthropic), effort max, 16 November 2026: targeted proof check. GPT-6 Astra (OpenAI), ChatGPT Pro mode, effort Pro, S6\_09\_PAPERS, 17 November 2026: targeted proof check. Also checked in this writing. Pre-continuation writing events: GPT-6 Astra (OpenAI), ChatGPT Pro mode, effort Pro, S6\_10\_PAPERS, 18 November 2026: targeted proof check. Continuation metadata events: GPT-6 Astra (OpenAI), ChatGPT Pro mode, effort Pro, S6\_10\_PAPERS continuation, 19 November 2026: used and rechecked in this run; status-vocabulary reconciliation; metadata reconciliation only; no new mathematical check. \emph{Pre-continuation recorded proof scope.} Finite Gorenstein duality, the exact conductor Euler sequence, connected boundary under each blowup, and the generic length-two conductor algebra were re-derived. Degree and arithmetic genus force the degree-one smooth elliptic cusp image; the first possible blowup contradicts the canonical-conductor identity, yielding the stated minimality and Picard restriction}{the hypotheses, notation and definitions stated here; \cref{cdp:st:A.5.3}; \cref{cdp:st:A.1.1}; \cref{cdp:st:A.5.8}; \cref{cdp:st:A.5.2}; \cref{cdp:eq:A.RC.3}; \cref{cdp:eq:A.RC.4}; \cref{cdp:st:C2d}; \cref{cdp:eq:A.RC.5}; \cref{cdp:st:A.2.4}}

Let $S\subset Y$ be an irreducible nonnormal reduced Cartier
divisor in a compact smooth complex threefold with
$H^2(Y,\Z)=0$ and $\chi(\OO_S)=0$.
Let $\eta:T\to S$ be its normalization and
$\pi:Z\to T$ its minimal resolution, and put $h=\eta\pi$.
Suppose that $Z$ is irrational ruled, with ruling $p:Z\to B$,
that every exceptional curve of $\pi$ is vertical, and that
the effective divisor of \ref{cdp:st:A.5.3} has the form
\[
 K_Z+D_Z=h^*\omega_S,\qquad D_Z=2C+V,                 \tag{A.RC.1}\label{cdp:eq:A.RC.1}
\]
where $C$ is a section in the strict conductor and $V$ is
vertical. Suppose also that the conductor image scheme
$\Gamma=\operatorname{Spec}_S(\OO_S/\mathfrak c)$ is reduced,
and that $H^0(S,K\otimes L|_S)\ne0$ for every $L\in\Pic(Y)$,
with $K$ as in \ref{cdp:st:A.1.1}.
Then every ambient line bundle pulls back trivially to $T$.
In particular
\[
 H^2\bigl(T,\omega_T\otimes
       \eta^*((K_Y^{-1}\otimes L)|_S)\bigr)\simeq\C
       \quad\hbox{for every }L\in\Pic(Y).            \tag{A.RC.2}\label{cdp:eq:A.RC.2}
\]
The proof forces $B$ and $\Gamma$ to be smooth elliptic curves,
$Z=T$ to be minimal ruled, and $D_Z=2C$ with $C^2=0$.
Thus no local-complete-intersection assumption on the original
conductor image curve is needed in this reduced-scheme case.
\end{proposition}


\begin{proof}
\noindent\textbf{Step 1.} Write $\mathcal B=\eta_*\OO_T$ and
$\mathcal Q=\mathcal B/\OO_S$. We first specify conductor
duality, including its twists. Since $S$ is Gorenstein of
dimension two and $\mathcal B$ is maximal Cohen--Macaulay,
its dual is $\mathfrak c=\mathcal Hom_S(\mathcal B,\OO_S)$,
and maximal Cohen--Macaulay biduality gives
$\mathcal Hom_S(\mathfrak c,\OO_S)=\mathcal B$.
The quotient $\OO_\Gamma$ is Cohen--Macaulay of dimension one:
this follows from the depth lemma in
$0\to\mathfrak c\to\OO_S\to\OO_\Gamma\to0$.
Applying $\mathcal Hom_S(-,\OO_S)$ to that sequence, its
initial terms identify the actual normalization inclusion and
its cokernel as
\[
 0\longrightarrow\OO_S\longrightarrow\mathcal B
 \longrightarrow
 \mathcal Ext^1_S(\OO_\Gamma,\OO_S)\longrightarrow0,
 \qquad
 \mathcal Q\simeq\omega_\Gamma\otimes\omega_S^{-1}|_\Gamma.
                                                               \tag{A.RC.3}\label{cdp:eq:A.RC.3}
\]
Adjunction makes $\omega_S$ the restriction of
$K_Y\otimes\OO_Y(S)$. Its degree is zero on every compact
curve, because $H^2(Y,\Z)=0$. Curve Riemann--Roch and Serre
duality therefore give
$\chi(\mathcal Q)=\chi(\omega_\Gamma)=-\chi(\OO_\Gamma)$.
On the other hand, \ref{cdp:st:A.5.8} gives rational singularities of $T$
and $\chi(\OO_T)=1-g(B)$. The normalization sequence and
$\chi(\OO_S)=0$ now give the exact numerical identity
\[
                   \chi(\OO_\Gamma)=g(B)-1.          \tag{A.RC.4}\label{cdp:eq:A.RC.4}
\]

\noindent\textbf{Step 2.} We next prove that $\Gamma$ is connected. Contract ruling
$(-1)$-curves from $Z$ to a minimal ruled bundle $Z_0\to B$.
Under a single contraction $\sigma:Z_1\to Z_2$ with exceptional
curve $E$, let $D_2=\sigma_*D_1$, let $a$ be the coefficient
of $E$ in $D_1$, and let $m$ be the multiplicity of $D_2$ at
the contraction point. The canonical divisor formula gives
\[
 K_{Z_1}+D_1=\sigma^*(K_{Z_2}+D_2)+(1+a-m)E.
                                                               \tag{A.RC.5}\label{cdp:eq:A.RC.5}
\]
Numerical triviality and intersection with $E$ force $a=m-1$.
The same numerical triviality descends to $Z_2$. On $Z_0$
the pushed divisor is $2C_0+V_0$, with $C_0$ a section and
$V_0$ a nonnegative sum of fibres, so its support is connected.
Reverse the contractions. A centre outside the existing support
would produce coefficient $-1$ on its exceptional curve; that
coefficient persists on strict transforms and contradicts the
effectiveness of the final $D_Z$. A centre of multiplicity one
is a smooth point of one coefficient-one component and its
exceptional coefficient is zero, so it does not disconnect the
support. At a centre of multiplicity at least two the exceptional
curve has positive coefficient and connects all strict-transform
branches through that point. Induction proves that
$\operatorname{Supp}D_Z$ is connected.

\noindent\textbf{Step 3.} Every positive exceptional component of $D_Z$ maps into
$\Gamma$. Indeed away from $\Gamma$ the surface $T$ equals
the Gorenstein surface $S$ and has rational singularities by
\ref{cdp:st:A.5.8}. Such a singularity is crepant on its minimal resolution:
a generator of its dualizing sheaf pulls back holomorphically
by rational duality; its effective exceptional zero divisor
has nonnegative intersection with each exceptional curve by
adjunction and minimality, so negative definiteness makes it
zero, as in \ref{cdp:st:A.5.2}. Thus \eqref{cdp:eq:A.RC.1} has no positive exceptional
coefficient there. The strict conductor maps onto the whole
conductor support. It follows that
$h(\operatorname{Supp}D_Z)=\operatorname{Supp}\Gamma$,
which proves connectedness.

\noindent\textbf{Step 4.} Let $D\subset\Gamma$ be the image of $C$. Localize \eqref{cdp:eq:A.RC.3}
at the generic point of $D$. Because $\Gamma$ is reduced,
$\OO_\Gamma$ there is the function field $k(D)$, and its
dualizing module has dimension one over that field. The
elementary conductor quotient sequence is
\[
 0\longrightarrow\OO_\Gamma\longrightarrow
          \mathcal B/\mathfrak c
   \longrightarrow\mathcal Q\longrightarrow0.
\]
It follows that $\mathcal B/\mathfrak c$ has dimension two
over $k(D)$. The conductor coefficient two on $C$ contributes
$2\deg(C\to D)$ to this dimension. Therefore this degree is
one and there is no other normalization branch over $D$.
The generic transverse ring has conductor $t^2\C\{t\}$ and
is exactly $\C+t^2\C\{t\}=\C\{t^2,t^3\}$, the cusp ring.
Thus $D^\nu\simeq C\simeq B$. This deduction also rules out
any additional vertical branch over the same image curve.

\noindent\textbf{Step 5.} A connected compact reduced curve has only constant global
holomorphic functions. From \eqref{cdp:eq:A.RC.4} we therefore obtain
$h^1(\Gamma,\OO_\Gamma)=2-g(B)$.
Restriction to the reduced component $D$ gives a surjection
on $H^1$, because its kernel is a coherent sheaf on a curve
and has no $H^2$. Normalization of $D$ gives
$h^1(D,\OO_D)\ge g(B)$. Hence
\[
             2-g(B)\ge h^1(D,\OO_D)\ge g(B).         \tag{A.RC.6}\label{cdp:eq:A.RC.6}
\]
As $g(B)\ge1$, all quantities equal one. The normalization
quotient of $D$ has length zero, so $D$ itself is smooth
elliptic. Statement \ref{cdp:st:C2d} applies to this actual smooth curve
in $Y$ and makes the transversal cusp constant along it.
No other singular curve meets $D$. Connectedness of $\Gamma$
therefore gives $\Gamma=D$. Its normalization inverse image
is the single smooth curve $C$, the normalization is smooth
near it, and \ref{cdp:st:C2d} gives $C^2=0$.

\noindent\textbf{Step 6.} Adjunction in \eqref{cdp:eq:A.RC.1} yields $V\cdot C=0$. A ruling
$(-1)$-curve outside $V$ cannot meet $C$, since its
intersection with $K_Z+2C+V$ would be at least $-1+2$.
A component of $V$ cannot meet $C$ either. Thus all
contractions to $Z_0$ avoid $C$ and preserve its square zero.
On the minimal bundle, $K_{Z_0}+2C_0$ is numerically trivial,
so its nonnegative vertical divisor $V_0$ is zero.
The first reverse blowup would be away from $C_0$ and have
coefficient $-1$ by \eqref{cdp:eq:A.RC.5}. There are no such blowups.
Thus $Z$ is minimal ruled and $D_Z=2C$.
An exceptional curve for $Z\to T$ would have negative square;
adjunction with $K_Z\equiv-2C$ would make it rational and
therefore vertical. A minimal ruled bundle has only fibres
of square zero, so no such exceptional curve exists and $Z=T$.

\noindent\textbf{Step 7.} The constant-cusp computation \ref{cdp:st:A.2.4} now gives the injection
$K\hookrightarrow\eta_*(h^*\OO_Y(-S)\otimes\OO_Z(C))$.
Every ambient line is topologically trivial on $Z$ and is
the pullback of a degree-zero line on the elliptic curve $B$.
The section exact sequence, together with $R^1p_*\OO_Z=0$,
gives
\[
 0\longrightarrow\OO_B\longrightarrow p_*\OO_Z(C)
 \longrightarrow N_{C/Z}\longrightarrow0,
 \qquad\deg N_{C/Z}=0.                              \tag{A.RC.7}\label{cdp:eq:A.RC.7}
\]
After a degree-zero twist $P$, both end terms have zero
sections unless $P=\OO_B$ or $P=N_{C/Z}^{-1}$.
The image of $\Pic(Y)$ in $\Pic^0(B)$ is connected:
$H^2(Y,\Z)=0$ makes $\Pic(Y)$ the connected image of
$H^1(Y,\OO_Y)$ under the exponential map.
All-twist persistence of $K$, and translation by the ambient
line $\OO_Y(-S)$, put this entire image in those at most two
points. It is therefore the identity. This proves triviality
of every ambient pullback. Serre duality on $T$ proves
\eqref{cdp:eq:A.RC.2}.
\end{proof}






\begin{proposition}[The genus of a horizontal conductor image from a local torus group]\label{cdp:st:A.5.11}
\StatementRecord{cdp:st:A.5.11}{proved}{Historical events: Claude Fable 5.2 (Anthropic), effort max, 16 November 2026: targeted proof check. GPT-6 Astra (OpenAI), ChatGPT Pro mode, effort Pro, S6\_09\_PAPERS, 17 November 2026: targeted proof check. Also checked in this writing. Pre-continuation writing events: GPT-6 Astra (OpenAI), ChatGPT Pro mode, effort Pro, S6\_10\_PAPERS, 18 November 2026: targeted proof check; used and rechecked in this run. Continuation metadata events: GPT-6 Astra (OpenAI), ChatGPT Pro mode, effort Pro, S6\_10\_PAPERS continuation, 19 November 2026: used and rechecked in this run; status-vocabulary reconciliation; metadata reconciliation only; no new mathematical check. \emph{Pre-continuation recorded proof scope.} Collapsing the finite exceptional-value tree gives the actual continuous map onto the curve quotient and its fundamental-group surjection. Lojasiewicz section3 Theorem2, printed page464, applies to the smooth countable real analytic total space and its closed analytic support. The explicit barycentric retraction and metabelian torus-bundle group then force image genus at most one}{the hypotheses, notation and definitions stated here; \cref{cdp:st:A.5.8}; \cref{cdp:eq:A.Pi.1}; \cref{cdp:eq:A.Pi.2}; \cite[\S3, Theorem~2, p.~464]{Lojasiewicz1964}; \cref{cdp:eq:A.Pi.4}}

Let $f:Y\to C_0$ be a proper holomorphic map from a smooth
compact complex threefold to a smooth curve, with connected
fibres and complex two-torus general fibre. Let $S$ be an
irreducible reduced fibre support, $\eta:T\to S$ its
normalization, and $\pi:Z\to T$ a resolution which is
irrational ruled, with ruling $p:Z\to B$.
Suppose every exceptional curve of $\pi$ is vertical, and
the strict conductor has just one horizontal irreducible
component $C$, which is a section of $p$.
Let $D=\eta\pi(C)$ and let $\Sigma=D^\nu$ be its normalization.
Then $g(\Sigma)\le1$.
In particular, when $C\to D$ has degree one, $B$ is elliptic.
This applies to the horizontal-cusp case of $D_Z=2C+V$;
it allows arbitrary multiplicity of the full fibre and
arbitrary singularities or nonreduced conductor image on
the finitely many vertical curves.
\end{proposition}


\begin{proof}
\noindent\textbf{Step 1.} There is nothing to prove when $\Sigma$ is rational, so
assume $g(\Sigma)\ge1$.
Statement \ref{cdp:st:A.5.8} gives a holomorphic ruling
$\widetilde p:T\to B$ with $p=\widetilde p\pi$.
Identifying $C$ with $B$ by its section map, normalization
of its nonconstant map to $D$ gives a finite surjective
holomorphic map $\phi:B\to\Sigma$.
Put $\psi=\phi\widetilde p:T\to\Sigma$.

\noindent\textbf{Step 2.} We describe precisely the finite identifications needed to
make this map descend to $S$. Every vertical conductor
component of $T$ lies over one of finitely many points of
$B$, so $\psi$ is constant on it. Its image under $\eta$
has rational normalization: a vertical component is the
image of a rational curve on $Z$, and a nonconstant map
from that curve to its image normalization is finite.
Such an image cannot equal $D$, because $g(\Sigma)\ge1$.
Consequently these vertical images meet $D$ at finitely
many points. Away from those points and from the finitely
many singular points of $D$, every normalization preimage
of a point of $D$ lies on $C$, and its $\psi$ value is the
unique corresponding point of $\Sigma$.
At a point outside $D$ with more than one normalization
preimage, all those preimages lie on vertical conductor
components. There are no other horizontal conductor
components by hypothesis. Thus there is a finite set
$A\subset\Sigma$ such that, for every $s\in S$, either
$\psi(\eta^{-1}(s))$ is a singleton or it is contained
in $A$. One can take $A$ to contain the constant values
on all vertical conductor components and all points of
$\Sigma$ above the singularities of $D$ and its
intersections with the vertical image curves.

\noindent\textbf{Step 3.} Choose an embedded finite tree $J$ in the underlying
oriented surface $\Sigma$ containing $A$, and let
$q:\Sigma\to\Sigma/J$ collapse that tree to one point.
Such a tree is obtained by joining the finite set by
arcs and retaining a spanning tree after subdividing
their finitely many crossings. Triangulate the surface
with $J$ a subcomplex. A contractible subcomplex is a
cofibration; contracting it gives a homotopy equivalence
of this CW pair. Concretely the homotopy extension
property extends a contraction of $J$ and supplies a
homotopy inverse to the quotient map. In particular
\[
              q_*:\pi_1(\Sigma)\xrightarrow{\sim}
                         \pi_1(\Sigma/J).          \tag{A.Pi.1}\label{cdp:eq:A.Pi.1}
\]
The map $q\psi$ is constant on every fibre of the finite
proper surjection $\eta$. A proper continuous surjection
between these compact Hausdorff spaces is a quotient
map, so there is a unique continuous map
\[
 a:S\longrightarrow\Sigma/J,
 \qquad a\eta=q\psi.                              \tag{A.Pi.2}\label{cdp:eq:A.Pi.2}
\]
Let $\delta:\Sigma\to D\subset S$ be the curve normalization
map. On its dense smooth locus, lifting through $C$ gives
$a\delta=q$. Continuity gives this equality everywhere.
Equation \eqref{cdp:eq:A.Pi.1} therefore makes $a_*$ surjective on
fundamental groups, without any assertion that a finite
point-identification quotient has the same fundamental
group as $\Sigma$.

\noindent\textbf{Step 4.} We next extend the actual map \eqref{cdp:eq:A.Pi.2} to a neighbourhood
of $S$ in $Y$. The simultaneous analytic triangulation
theorem of \cite[\S3, Theorem~2, p.~464]{Lojasiewicz1964}
triangulates the pair $(Y,S)$, since a complex analytic
subset is real semianalytic. Compactness makes the
subcomplex representing $S$ finite. After barycentric
subdivision it is full: vertices lying in it span a
face in that same subcomplex. Write $\lambda_v$ for
barycentric coordinates and let $L$ be its vertex set.
On the open neighbourhood where $\sum_{v\in L}\lambda_v>0$,
the formula
\[
 r(x)=\frac{\sum_{v\in L}\lambda_v(x)v}
                  {\sum_{v\in L}\lambda_v(x)}       \tag{A.Pi.3}\label{cdp:eq:A.Pi.3}
\]
defines a continuous retraction to $S$ in triangulation
coordinates. Its value lies in the face spanned by the
$L$ vertices of the same simplex, so these formulas
agree across faces. This supplies an actual open
neighbourhood $N$ and a map $r:N\to S$ fixed on $S$.

\noindent\textbf{Step 5.} Let $c=f(S)$. Properness allows a small disc $\Delta$
about $c$ with $U=f^{-1}(\Delta)\subset N$ and no other
critical value. The map $(ar)|_U:U\to\Sigma/J$ is
surjective on fundamental groups: its restriction to
$\delta(\Sigma)\subset S\subset U$ is $q$, whose
fundamental-group map is an isomorphism.
Set $U^*=U\setminus S=f^{-1}(\Delta^*)$.
The inclusion induces a surjection
$\pi_1(U^*)\twoheadrightarrow\pi_1(U)$.
Indeed a loop in the smooth six-manifold $U$ can be
perturbed relative to its base point to be transverse
to an analytic stratification of $S$; every stratum
has real codimension at least two, so the resulting
one-dimensional loop misses every stratum.

\noindent\textbf{Step 6.} Over $\Delta^*$, the proper submersion is a smooth
torus bundle. The fundamental-group exact sequence
and $\pi_2(\Delta^*)=0$ give
\[
 1\longrightarrow\Z^4\longrightarrow\pi_1(U^*)
       \longrightarrow\Z\longrightarrow1.
\]
A lift of the generator of $\Z$ splits this group
extension, hence
\[
 \pi_1(U^*)\simeq\Z^4\rtimes_T\Z
 \twoheadrightarrow\pi_1(U)
 \twoheadrightarrow\pi_1(\Sigma/J)
 \simeq\pi_1(\Sigma).                             \tag{A.Pi.4}\label{cdp:eq:A.Pi.4}
\]
The first group is metabelian: its commutator subgroup
is contained in the abelian normal subgroup $\Z^4$.
Every quotient in \eqref{cdp:eq:A.Pi.4} is therefore metabelian.
A closed orientable surface of genus at least two
has a fundamental group surjecting onto the free
group on two generators: in its usual surface
presentation, retain two $a$ generators and send all
$b$ generators and the other $a$ generators to the
identity. That free group surjects onto $S_4$ by
sending its generators to $(1234)$ and $(12)$;
their conjugates supply the adjacent transpositions.
The commutator subgroup of $S_4$ is the nonabelian
group $A_4$, so this quotient is not metabelian.
This contradicts \eqref{cdp:eq:A.Pi.4}. Thus $g(\Sigma)=1$ in the
positive-genus case. When $C\to D$ has degree one,
$\Sigma\simeq C\simeq B$, which proves the final
assertion.
\end{proof}






\begin{proposition}[The arithmetic genus of conductor images from coherent duality]\label{cdp:st:A.5.12}
\StatementRecord{cdp:st:A.5.12}{proved}{Historical events: Claude Fable 5.2 (Anthropic), effort max, 16 November 2026: targeted proof check. GPT-6 Astra (OpenAI), ChatGPT Pro mode, effort Pro, S6\_09\_PAPERS, 17 November 2026: targeted proof check. Also checked in this writing. Pre-continuation writing events: GPT-6 Astra (OpenAI), ChatGPT Pro mode, effort Pro, S6\_10\_PAPERS, 18 November 2026: targeted proof check. Continuation metadata events: GPT-6 Astra (OpenAI), ChatGPT Pro mode, effort Pro, S6\_10\_PAPERS continuation, 19 November 2026: used and rechecked in this run; status-vocabulary reconciliation; metadata reconciliation only; no new mathematical check. \emph{Pre-continuation recorded proof scope.} Canonical torsion on the full multiple fibre, coherent Euler and Hodge dimensions, conductor duality, and the genus bound on the reduced image were re-derived}{the hypotheses, notation and definitions stated here; \cref{cdp:eq:A.CD.2}; \cref{cdp:eq:A.CD.3}; \cref{cdp:eq:A.CD.1}}

Let $f:Y\to C_0$ be a proper connected-fibre map from a
smooth compact complex threefold with $H^2(Y,\Z)=0$ to
a smooth curve, and let its smooth general fibre have
trivial canonical bundle. This includes complex
two-tori and primary Kodaira surfaces.
Let $mS$ be a full fibre with irreducible
reduced nonnormal support $S$. Write $\eta:T\to S$ for
normalization and $\pi:Z\to T$ for a resolution.
If the pullback homomorphism
$\Pic(Y)\to\Pic(Z)$ is nonconstant, then every reduced
irreducible component $D$ of the conductor image scheme
has arithmetic genus at most one.
In particular every such component whose normalization
has positive genus is a smooth elliptic curve.
The conductor image scheme itself is allowed to be
nonreduced, and no local-complete-intersection assumption
is imposed on its reduced components.
\end{proposition}


\begin{proof}
\noindent\textbf{Step 1.} First we determine exactly the needed canonical line on
the reduced support, without a multiplicity formula for
the direct image. The sheaf $f_*K_Y$ is torsion-free of
generic rank one on the smooth curve and hence is a
line bundle. Choose a frame near $c=f(S)$ and evaluate
it in $K_Y$. A coordinate $t$ on the base disc supplies
the conormal frame $dt$, and adjunction identifies
$K_Y|_F$ with $K_F$ after dividing by this frame.
Under that line-bundle isomorphism, the evaluated
section restricts on a general smooth fibre to a
nonzero canonical two-form, which is nowhere zero.
Its zero divisor is therefore
vertical. Shrinking the base disc to remove the other
vertical zero divisors, irreducibility of the support
gives an effective divisor $kS$, with $k\ge0$, and an
actual line-bundle comparison on $U=f^{-1}(\Delta)$:
\[
 K_Y|_U\simeq\OO_U(kS),\qquad
 \omega_S\simeq\OO_Y((k+1)S)|_S,\qquad
 \omega_S^{\otimes m}\simeq\OO_S.                 \tag{A.CD.1}\label{cdp:eq:A.CD.1}
\]
The last isomorphism uses the local base equation whose
zero divisor is exactly the full fibre $mS$.
In particular no assertion $k=m-1$ is used.

\noindent\textbf{Step 2.} Hirzebruch--Riemann--Roch in
$0\to\OO_Y(-S)\to\OO_Y\to\OO_S\to0$ gives
$\chi(\OO_S)=0$: the rational Chern character of
$\OO_Y(-S)$ equals that of $\OO_Y$, since its first
Chern class is zero. A nonzero section of the torsion
line $\omega_S$ trivializes it. Indeed its $m$th power
is a nonzero holomorphic function on the compact
connected reduced space and hence is a nonzero constant.
Thus, with $e=h^0(S,\omega_S)$, we have
\[
 e\in\{0,1\},\qquad
 e=1\Longleftrightarrow\omega_S\simeq\OO_S,
 \qquad h^1(S,\OO_S)=1+e.                         \tag{A.CD.2}\label{cdp:eq:A.CD.2}
\]
Here $h^0(\OO_S)=1$ and Serre duality identifies
$h^2(\OO_S)$ with $e$.

\noindent\textbf{Step 3.} Let $\mathcal Q=\eta_*\OO_T/\OO_S$ and let
$\Gamma\subset S$ be the conductor image scheme.
Both $S$ and $T$ are connected and compact, so their
global holomorphic functions are the same constants.
The normalization sequence consequently gives an
injection and the indicated kernel:
\[
 H^0(S,\mathcal Q)\hookrightarrow H^1(S,\OO_S),
 \qquad
 \operatorname{im}H^0(S,\mathcal Q)
 =\ker\bigl(H^1(S,\OO_S)\to H^1(T,\OO_T)\bigr).
                                                               \tag{A.CD.3}\label{cdp:eq:A.CD.3}
\]
Because $H^2(Y,\Z)=0$, $\Pic(Y)$ is the connected
exponential image of $H^1(Y,\OO_Y)$. A nonconstant
holomorphic group homomorphism out of it has nonzero
differential. The differential of the stated Picard
pullback is the actual composite
\[
 H^1(Y,\OO_Y)\longrightarrow H^1(S,\OO_S)
 \longrightarrow H^1(T,\OO_T)
 \longrightarrow H^1(Z,\OO_Z).
\]
Thus its middle arrow has positive rank. Equations
\eqref{cdp:eq:A.CD.2}--\eqref{cdp:eq:A.CD.3} imply $h^0(\mathcal Q)\le e$.

\noindent\textbf{Step 4.} The conductor is Cohen--Macaulay of dimension one,
and the exact duality comparison (A.RC.3) applies
without a reducedness assumption on $\Gamma$:
$\mathcal Q=\omega_\Gamma\otimes\omega_S^{-1}|_\Gamma$.
Serre duality on that curve therefore gives
\[
 h^1(\Gamma,\omega_S|_\Gamma)
       =h^0(S,\mathcal Q)\le e.                   \tag{A.CD.4}\label{cdp:eq:A.CD.4}
\]
For any reduced irreducible component $D$ of $\Gamma$,
restriction of the line gives an exact quotient
$\omega_S|_\Gamma\to\omega_S|_D$. Its kernel is a
coherent sheaf on a one-dimensional compact space and
has no $H^2$, so the corresponding map on $H^1$ is
surjective. Hence
$h^1(D,\omega_S|_D)\le e$.

\noindent\textbf{Step 5.} For $e=1$, \eqref{cdp:eq:A.CD.2} makes this precisely
$p_a(D)=h^1(D,\OO_D)\le1$.
For $e=0$, the same inequality makes
$h^1(D,\omega_S|_D)=0$.
The restriction $\omega_S|_D$ is torsion by \eqref{cdp:eq:A.CD.1}
and has degree zero. Curve Riemann--Roch now gives
\[
 0\le h^0(D,\omega_S|_D)
       =\chi(D,\omega_S|_D)=1-p_a(D),
\]
again proving $p_a(D)\le1$.
Finally, for a reduced irreducible curve the
normalization exact sequence gives
$p_a(D)=g(D^\nu)+\operatorname{length}
(\nu_*\OO_{D^\nu}/\OO_D)$.
Positive normalization genus therefore forces genus
one and zero normalization quotient. A normal complex
curve is smooth, which proves the final assertion.
\end{proof}






\begin{proposition}[The conductor bound when the general fibre has $q=1$ and $p_g=0$]\label{cdp:st:A.5.13}
\StatementRecord{cdp:st:A.5.13}{proved}{Historical events: Claude Fable 5.2 (Anthropic), effort max, 16 November 2026: targeted proof check. GPT-6 Astra (OpenAI), ChatGPT Pro mode, effort Pro, S6\_09\_PAPERS, 17 November 2026: targeted proof check. Also checked in this writing. Pre-continuation writing events: GPT-6 Astra (OpenAI), ChatGPT Pro mode, effort Pro, S6\_10\_PAPERS, 18 November 2026: targeted proof check; used and rechecked in this run. Continuation metadata events: GPT-6 Astra (OpenAI), ChatGPT Pro mode, effort Pro, S6\_10\_PAPERS continuation, 19 November 2026: used and rechecked in this run; status-vocabulary reconciliation; metadata reconciliation only; no new mathematical check. \emph{Pre-continuation recorded proof scope.} The rank-one R1 bundle, elimination of R2 torsion by the Euler characteristic, the Picard injection, and conductor genus bound were re-derived using the stated known direct-image input}{the hypotheses, notation and definitions stated here; \cref{cdp:st:D.1.9}; \cref{cdp:st:E4}; \cref{cdp:st:A.5.12}}

Let $f:Y\to\PP^1$ be a proper connected-fibre map from
a smooth compact complex threefold with
$H^1(Y,\Z)=H^2(Y,\Z)=0$. Suppose its smooth general
fibre has $h^1(\OO)=1$ and $h^2(\OO)=0$.
Then
\[
 R^1f_*\OO_Y=\OO_{\PP^1},\qquad R^2f_*\OO_Y=0,
 \qquad\Pic(Y)=\C\beta,                            \tag{A.Q1.1}\label{cdp:eq:A.Q1.1}
\]
where the last identification uses the exponential
coordinates of \ref{cdp:st:D.1.9} and its nonzero transgression
class $\beta$. Every full fibre has structure-sheaf
cohomology dimensions $(1,1,0)$.
For an irreducible reduced support $S$, every
nonconstant ambient Picard pullback to a normalization
resolution forces every reduced conductor component
to have arithmetic genus at most one. In particular
each positive-genus component is smooth elliptic,
without any lci or conductor-reducedness hypothesis.
\end{proposition}


\begin{proof}
\noindent\textbf{Step 1.} By \ref{cdp:st:E4}, $E=R^1f_*\OO_Y$ is a line bundle. The section
induced by \ref{cdp:st:D.1.9}'s integral transgression class is
nowhere zero: its image under the natural base-change
map at each point is the nonzero coefficient class on
that full fibre. A zero of the section in $E$ would
have zero image under that map. Thus $E\simeq\OO$.
Coherent Leray now gives $H^1(Y,\OO_Y)=\C\beta$.
The exponential sequence and the two stated integral
cohomology vanishings give the Picard identification
in \eqref{cdp:eq:A.Q1.1}.

\noindent\textbf{Step 2.} The coherent sheaf $G=R^2f_*\OO_Y$ has generic rank
zero, hence is a torsion sheaf. Riemann--Roch on $Y$
gives $\chi(\OO_Y)=\int_Yc_1c_2/24=0$, since
$c_1(Y)$ is zero in rational cohomology. The Euler
characteristic of coherent Leray is therefore
\[
 0=\chi(\OO_{\PP^1})-\chi(E)+\chi(G)
      =\operatorname{length}G.
\]
Thus $G=0$. Top-degree base change gives
$H^2(F,\OO_F)=0$ for every full Cartier fibre $F$.
The multiplication-by-parameter sequence then gives
$H^1(F,\OO_F)=E\otimes\C(c)$, of dimension one;
\ref{cdp:st:E4} gives $h^0(F,\OO_F)=1$.

\noindent\textbf{Step 3.} For a full fibre $F=mS$ with irreducible reduced
support, the quotient $\OO_F\to\OO_S$ induces a
surjection on $H^2$, because its kernel has support
of dimension two and has no $H^3$.
Hence $h^2(S,\OO_S)=0$. Riemann--Roch in the Cartier
sequence on $Y$ gives $\chi(\OO_S)=0$ as in \ref{cdp:st:A.5.12},
so $h^1(S,\OO_S)=1$.

\noindent\textbf{Step 4.} Put $\mathcal Q=\eta_*\OO_T/\OO_S$ for the
normalization $T$ and write $\Gamma$ for the conductor
image scheme. Nonconstant Picard pullback has nonzero
differential factoring through
$H^1(S,\OO_S)\to H^1(T,\OO_T)$.
The source has dimension one, so this map has zero
kernel. The normalization sequence (A.CD.3) gives
$H^0(\mathcal Q)=0$. Conductor duality and Serre
duality give
\[
 H^1(\Gamma,\omega_S|_\Gamma)=0,
 \qquad H^1(D,\omega_S|_D)=0                      \tag{A.Q1.2}\label{cdp:eq:A.Q1.2}
\]
for each reduced irreducible component $D$, where
the second map is the restriction surjection in the
proof of \ref{cdp:st:A.5.12}. Adjunction makes $\omega_S$ ambient,
so $\deg(\omega_S|_D)=0$ by $H^2(Y,\Z)=0$.
Curve Riemann--Roch yields
$0\le h^0(D,\omega_S|_D)=1-p_a(D)$.
The normalization genus and its nonnegative delta
length then give the asserted smooth elliptic
conclusion whenever that genus is positive.
\end{proof}






\begin{proposition}[Conductor-cleared jets on a complete-intersection critical curve]\label{cdp:st:A.5.5}
\StatementRecord{cdp:st:A.5.5}{proved}{Historical events: Claude Fable 5.2 (Anthropic), effort high, 16 November 2026: carried as background without fault found. GPT-6 Astra (OpenAI), ChatGPT Pro mode, effort Pro, S6\_09\_PAPERS, 17 November 2026: targeted proof check. Also checked in this writing. Pre-continuation writing events: GPT-6 Astra (OpenAI), ChatGPT Pro mode, effort Pro, S6\_10\_PAPERS, 18 November 2026: targeted proof check. GPT-6 Astra (OpenAI), ChatGPT Pro mode, effort Pro, S6\_10\_PAPERS, 18 November 2026 05:35:21+00:00: used and rechecked in this run. Continuation metadata events: GPT-6 Astra (OpenAI), ChatGPT Pro mode, effort Pro, S6\_10\_PAPERS continuation, 19 November 2026: used and rechecked in this run; status-vocabulary reconciliation; metadata reconciliation only; no new mathematical check. \emph{Pre-continuation recorded proof scope.} The lci conormal injection, rank-one Hessian, exact conductor-denominator cubic and quartic identities, and degree inequalities were re-derived. The quartic identity forces degree-zero quotient and conductor length zero; the cubic discriminant does the same for the triple case. The stated local example has an actual double conductor pole}{the hypotheses, notation and definitions stated here; \cref{cdp:eq:A.J.4}; \cref{cdp:eq:A.J.5}; \cref{cdp:eq:A.J.2}; \cref{cdp:eq:A.J.3}; \cref{cdp:eq:A.J.7}; \cref{cdp:eq:A.J.9}; \cref{cdp:eq:A.J.8}; \cref{cdp:eq:A.J.10}}

Let $Y$ be a smooth compact complex threefold with
$H^2(Y,\mathbb Q)=0$, let $S\subset Y$ be a reduced divisor,
and let $D\subset\operatorname{Sing}(S)$ be a reduced
irreducible compact curve which is a local complete intersection
in $Y$. Assume that the generic plane-curve transversal to $D$
in $S$ is an ordinary cusp, an ordinary tacnode, or an ordinary
triple point.
If the normalization $C\to D$ has positive genus, then $D$ is
a smooth elliptic curve. The transversal type is the indicated
cusp, tacnode, or triple point at every point of $D$;
in particular no degeneration of that type occurs along $D$.

In the cusp case the proof constructs, on $C$, the actual
meromorphic cubic on the radical of the Hessian and bounds its
pole divisor by twice the conductor. In the tacnode case the
corresponding fourth-order invariant has poles bounded by four
times the conductor. These intermediate jets need not extend
holomorphically before the conductor correction.
\end{proposition}


\begin{proof}
\noindent\textbf{Step 1.} Let $\nu:C\to D$ be the normalization and $j:C\to Y$
its composite with the inclusion,
$I=\mathcal I_{D/Y}$, and
\[
 N=\bigl(\nu^*(I/I^2)\bigr)^*,\qquad
 M=j^*\mathcal O_Y(S).
\]
These are a rank-two vector bundle and a line bundle on $C$.
Let $\Xi$ be the conductor divisor on $C$, so that
$\nu_*\mathcal O_C(-\Xi)$ is the conductor ideal in
$\mathcal O_D$, and put $c=\deg\Xi\geq0$.
Adjunction for the complete-intersection curve and finite
duality give exact line-bundle comparisons
\[
 \det N=\nu^*\omega_D\otimes j^*K_Y^{-1},\qquad
 \omega_C=\nu^*\omega_D\otimes\mathcal O_C(-\Xi).
                                                        \tag{A.J.1}\label{cdp:eq:A.J.1}
\]
For the second comparison one can work directly in the total
quotient ring: $\operatorname{Hom}_{\mathcal O_D}
(\nu_*\mathcal O_C,\mathcal O_D)$ consists precisely of
elements multiplying the normalization into $\mathcal O_D$,
which is its conductor. Since $D$ is Gorenstein, tensoring
this identity by its invertible dualizing sheaf is finite
duality. Both $M$ and $j^*K_Y$ have degree zero, by the
stated rational cohomology hypothesis. Consequently
\[
 \deg\det N=2g(C)-2+c.                              \tag{A.J.2}\label{cdp:eq:A.J.2}
\]
This is the normal bundle of the complete-intersection ideal,
with its specified normalization pullback. It is the bundle
needed in the following jet calculation.

\noindent\textbf{Step 2.} First treat an ordinary cusp. Choose a local defining section
$f$ of $S$, using a local frame of $\mathcal O_Y(S)$.
Since $D\subset\operatorname{Sing}(S)$, its conormal class
maps to zero under
$I/I^2\to\Omega_Y^1|_D$. This homomorphism is injective:
its source is locally free on the reduced curve and the map
is injective at every generic point. Therefore $f\in I^2$.
The associated-graded identity for a regular sequence gives
$I^2/I^3=\operatorname{Sym}^2(I/I^2)$. The class of $f$
is thus a quadratic form
\[
 q\in H^0\bigl(C,\operatorname{Sym}^2N^*\otimes M\bigr).
\]
It has rank one generically. Let $R\subset N$ be its
saturated radical and $Q=N/R$. Both are line bundles.
The quadratic form annihilates $R$ everywhere by holomorphic
continuation, and gives a nonzero section of $M\otimes Q^{-2}$.
In particular
\[
 \deg Q\leq0.                                      \tag{A.J.3}\label{cdp:eq:A.J.3}
\]

\noindent\textbf{Step 3.} On $D_{\rm reg}$, the third-order term on the radical is
intrinsic. Indeed, after writing the quadratic form as
$a\ell^2$, subtract a holomorphic lift of $a\ell^2$ from
$f$ and restrict its cubic class to $R^3$. Changing that lift
adds terms divisible by $\ell$, which vanish on $R^3$.
Changing normal generators has the same effect. Its nonzero
value characterizes the ordinary-cusp transversal among the
rank-one quadratic cases. Denote this cubic by $k$.

\noindent\textbf{Step 4.} We now extend precisely this section across $C\setminus
\nu^{-1}(D_{\rm reg})$. Work at a point of $D$ and choose
two generators $X,Y$ of $I$. After a constant invertible linear
change of these generators, the normalized linear factor of
$q$ has the form $X+hY$ on every normalization branch over
that point. Such a choice exists: the saturated quotient
$N\to Q$ is surjective at each of the finitely many points,
so a constant vector can be chosen outside their finitely
many kernels. Thus $h$ is holomorphic on the normalization,
and
\[
 q=a(X+hY)^2,
\]
where $a=q_{XX}$ is the pullback of an actual coefficient
in $\mathcal O_D$. Choose a local conductor generator $d$
on the normalization. Both $d$ and $dh$ belong to
$\mathcal O_D$, since the conductor is an ideal in the
normalization contained in that ring. Choose lifts
$\widehat d,\widehat{dh},\widehat a$ to $\mathcal O_Y$,
and set
\[
 G=\widehat d^{\,2} f-
    \widehat a\bigl(\widehat d X+\widehat{dh}Y\bigr)^2.
                                                        \tag{A.J.4}\label{cdp:eq:A.J.4}
\]
The quadratic class of $G$ is zero after normalization.
It is already zero on $D$: the map from the locally free
sheaf $\operatorname{Sym}^2(I/I^2)$ to its normalization
pullback is injective because $D$ is reduced. Hence
$G\in I^3$.

\noindent\textbf{Step 5.} Its class in $I^3/I^4=\operatorname{Sym}^3(I/I^2)$,
pulled back to $C$ and evaluated on $R^3$, is holomorphic.
Dividing that class by $d^2$ gives $k$ on the dense regular
locus. To check this equality, there $d$ is a unit and
\eqref{cdp:eq:A.J.4} is exactly the subtraction of the square used in the
definition of the radical cubic. Different choices of lifts
add cubic terms containing its linear factor, and therefore
have zero restriction to $R^3$. Thus these local expressions
are the same meromorphic section wherever they overlap.
Since $d$ has divisor $\Xi$, the resulting bound is
\[
 0\ne k\in H^0\bigl(C,M\otimes R^{-3}
                               \otimes\mathcal O_C(2\Xi)\bigr).
                                                        \tag{A.J.5}\label{cdp:eq:A.J.5}
\]
This proves the claimed pole control without a splitting of
the formal neighbourhood of the singular curve.

\noindent\textbf{Step 6.} Taking degrees in \eqref{cdp:eq:A.J.5} gives $\deg R\leq2c/3$.
Together with \eqref{cdp:eq:A.J.2}--\eqref{cdp:eq:A.J.3}, this yields
\[
 2g(C)-2+c=\deg R+\deg Q\leq\frac{2c}{3},
 \qquad 2g(C)-2+\frac c3\leq0.                     \tag{A.J.6}\label{cdp:eq:A.J.6}
\]
For $g(C)\geq1$, this forces $g(C)=1$ and $c=0$.
Conductor zero means $\mathcal O_D=\nu_*\mathcal O_C$;
a normal complex curve is smooth. Hence $D=C$ is elliptic.
Now \eqref{cdp:eq:A.J.2}--\eqref{cdp:eq:A.J.5} give $\deg R=\deg Q=0$.
The Hessian coefficient and radical cubic are nonzero
sections of degree-zero line bundles on this elliptic curve,
so neither has a zero. Locally the holomorphic Morse lemma
in the nondegenerate normal direction, with the curve
coordinate as parameter, writes $f=y^2+g(x,z)$.
The conditions just proved give $g=x^3u(x,z)$ with $u$ a
holomorphic unit. Replacing $x$ by a local cube root of $u$
times $x$ gives the actual local equation $y^2+x^3=0$.
This proves the asserted constant transversal type.

\noindent\textbf{Step 7.} Next treat an ordinary tacnode. Retain $N,M,R,Q,q,a,h,d,P$
from the cusp calculation, with
$P=\widehat d X+\widehat{dh}Y$ and
$G=\widehat d^{\,2}f-\widehat aP^2\in I^3$.
The generic radical cubic is now zero, so its continuation
$[G]_3|_{R^3}$ is zero everywhere on the normalization.
Since $\ell=X+hY$ is monic in $X$, polynomial division gives
an actual holomorphic quadratic $C_2$ on the normalization
such that
\[
 [G]_3=\ell C_2.
\]
Each coefficient of $dC_2$ is in the conductor and hence in
$\mathcal O_D$. Lift those coefficients and let $H$ be the
resulting homogeneous quadratic in $X,Y$. Set
\[
 K=\widehat d^{\,2}G-PH.
\]
Its cubic class pulls back to
$d^2\ell C_2-d\ell\cdot dC_2=0$.
The same locally free graded-piece injection used above gives
$K\in I^4$. There is an exact identity in the local ring of
the ambient smooth threefold:
\[
 4\widehat a\widehat d^{\,6}f
 -(2\widehat a\widehat d^{\,2}P+H)^2
 =4\widehat a\widehat d^{\,2}K-H^2.               \tag{A.J.7}\label{cdp:eq:A.J.7}
\]
On the dense smooth locus, complete the square in the
nondegenerate normal direction. If the relevant terms there
are $a\ell^2+b\ell r^2+er^4$, the fourth-order invariant is
$\Delta=ae-b^2/4$. A substitution
$\ell\mapsto\ell+ur^2$ leaves this expression unchanged;
terms divisible by $\ell^2$ and higher changes in the radical
coordinate do not alter it. Thus it is the intrinsic section
of $M^2\otimes Q^{-2}\otimes R^{-4}$ on that dense locus.
It is nonzero precisely for the ordinary-tacnode transversal
among the rank-one Hessian cases with zero radical cubic.

\noindent\textbf{Step 8.} Divide \eqref{cdp:eq:A.J.7} by $4\widehat d^{\,6}$ and take the
fourth-order class on $R^4$. Since
$[H]_2=dC_2$, this expresses exactly that same invariant as
\[
 \Delta=
 \frac{a[K]_4|_{R^4}-(C_2|_{R^2})^2/4}{d^4}.
                                                        \tag{A.J.8}\label{cdp:eq:A.J.8}
\]
Its numerator is holomorphic on the normalization.
Different choices agree on the dense regular locus and hence
give one meromorphic section. Therefore
\[
 0\ne\Delta\in H^0\bigl(C,M^2\otimes Q^{-2}
                    \otimes R^{-4}\otimes\mathcal O_C(4\Xi)\bigr).
                                                        \tag{A.J.9}\label{cdp:eq:A.J.9}
\]
Write $q_0=\deg Q\leq0$. Using
$\deg R=2g(C)-2+c-q_0$, the degree of the line bundle
in \eqref{cdp:eq:A.J.9} is
\[
 -2q_0-4\bigl(2g(C)-2+c-q_0\bigr)+4c
       =-8g(C)+8+2q_0.
\]
Its nonzero section and $g(C)\geq1$ force
$g(C)=1$ and $q_0=0$.

\noindent\textbf{Step 9.} The section $q$ of $M\otimes Q^{-2}$ is now nowhere zero.
In the local presentation above, $a=q_{XX}$ is therefore a
unit on every normalization branch. Since it belongs to
$\mathcal O_D$, it is a unit in that local ring as well.
The coefficient $h=q_{XY}/a$ consequently belongs to
$\mathcal O_D$; here $q_{XY}$ denotes the coefficient in
the symmetric matrix, so the cross term is $2q_{XY}XY$.
We can thus repeat the entire construction with $d=1$.
Polynomial division by the monic $X+hY$ then takes place
over $\mathcal O_D$ itself; the zero radical cubic descends
by the same injectivity argument. Formula \eqref{cdp:eq:A.J.8} now has
no poles. The resulting nonzero section of
$M^2\otimes Q^{-2}\otimes R^{-4}$ has degree $-4c$,
which forces $c=0$.

\noindent\textbf{Step 10.} Thus again $D=C$ is smooth elliptic. Both the nonzero
quadratic coefficient and $\Delta$ are sections of
degree-zero line bundles and are nowhere zero. The parametric
Morse lemma first gives $f=y^2+g(x,z)$, and the vanished
radical cubic and nonzero fourth invariant give
$g=x^4u(x,z)$ with $u$ a unit. Taking a local fourth root
produces the actual equation $y^2+x^4=0$ at every point of
$D$. This proves the tacnode assertion.

\noindent\textbf{Step 11.} For an ordinary triple transversal, the quadratic class
vanishes generically, hence vanishes in the locally free
sheaf $I^2/I^3$. Thus $f\in I^3$, and its initial cubic
is an honest global section
$b\in H^0(C,\operatorname{Sym}^3N^*\otimes M)$.
The discriminant of a binary cubic is homogeneous of degree
four in its coefficients and transforms by the sixth power
of the determinant of a change of variables. Its value is
therefore a section
\[
 0\ne\operatorname{disc}(b)\in
 H^0\bigl(C,M^4\otimes(\det N)^{-6}\bigr).
                                                        \tag{A.J.10}\label{cdp:eq:A.J.10}
\]
It is nonzero because the three generic tangent lines are
distinct. Equations \eqref{cdp:eq:A.J.2} and \eqref{cdp:eq:A.J.10} give
$2g(C)-2+c\leq0$, again forcing $g(C)=1$ and $c=0$
under the genus assumption. The discriminant has degree
zero and is nowhere zero. The three simple tangent roots
extend holomorphically on each small curve disc. The implicit
function theorem gives three smooth local branches of $S$
with those pairwise distinct normal tangents. Choose the
first two branches as the coordinate axes; the third is
$y=xu(x,z)$ with $u$ a unit, and replacing $y$ by $y/u(x,z)$
straightens it to $y=x$. This gives the local equation
$xy(x-y)=0$ up to a unit at every point of $D$.
\end{proof}




The pole allowance in \eqref{cdp:eq:A.J.5} has a concrete local use. Put
$h=z^2-x^3$, $D=(h,w)=0$ in $\mathbb C^3$, and
\[
 f=xw^2+2zwh+x^2h^2+h^3.
\]
On the normalization $x=s^2,z=s^3$ of $D$, its quadratic
class is $s^2(w+sh)^2$. Completing the square on $D_{\rm reg}$
gives
\[
 f=x\left(w+\frac zx h\right)^2+
                  \left(1-\frac1x\right)h^3.
\]
The radical cubic has a pole of order two at $s=0$.
The conductor generator is $d=x=s^2$, and the exact
conductor-cleared identity is
\[
 x^2f-x(xw+zh)^2=x(x-1)h^3\in I_D^3.
\]
Thus the corrected jet comparison applies even in this local
model where holomorphic extension of the uncorrected cubic
would fail. The equalities are checked exactly in
\texttt{code/E\_singular\_curve.py}.

The local calculation below is related to the conductor
derivations and characteristic-zero pole calculation in
Reid~\cite[\S3.9, \S4.4--4.8]{Reid1994Nonnormal}.
We give the canonical-module comparison directly, without
using an anticanonical polarization or a classification of
the surface.




\begin{proposition}[A pure double conductor forces a smooth cusp image]\label{cdp:st:A.5.7}
\StatementRecord{cdp:st:A.5.7}{proved}{Historical events: Claude Fable 5.2 (Anthropic), effort max, 16 November 2026: targeted proof check. GPT-6 Astra (OpenAI), ChatGPT Pro mode, effort Pro, S6\_09\_PAPERS, 17 November 2026: targeted proof check. Also checked in this writing. Pre-continuation writing events: GPT-6 Astra (OpenAI), ChatGPT Pro mode, effort Pro, S6\_10\_PAPERS, 18 November 2026: targeted proof check. Continuation metadata events: GPT-6 Astra (OpenAI), ChatGPT Pro mode, effort Pro, S6\_10\_PAPERS continuation, 19 November 2026: used and rechecked in this run; status-vocabulary reconciliation; metadata reconciliation only; no new mathematical check. \emph{Pre-continuation recorded proof scope.} Finite Gorenstein duality identifies the conductor quotient with its dualizing module. The derivative-graph calculation, fractional ideal endomorphism ring, and the exact coordinate change s'=s+u(s)z force a smooth conductor and the cusp product ring}{the hypotheses, notation and definitions stated here; \cref{cdp:eq:A.J.13}; \cref{cdp:eq:A.J.14}; \cref{cdp:eq:A.J.12}}

Let $S$ be a reduced Gorenstein complex surface and
$\pi:\overline S\to S$ its finite normalization.
Suppose that, over a neighbourhood of a point $p\in S$,
$\overline S$ is smooth and its conductor ideal is exactly
$\mathcal O_{\overline S}(-2C)$ for a smooth reduced
divisor $C$; its local branches over $p$ are disjoint.
Suppose that the generic transversal of $S$ along the
reduced image $D=\pi(C)$ is an ordinary cusp.
Then $D$ is smooth at $p$, $C\to D$ is an isomorphism
near the preimage of $p$, and the germ of $S$ at $p$ is
analytically a product of an ordinary cusp and a smooth disc.
No complete-intersection hypothesis on $D$ is required.
\end{proposition}


\begin{proof}
\noindent\textbf{Step 1.} Write $R=\mathcal O_{S,p}$ and let $\overline R$ be its
finite semilocal normalization. It is a product of regular
two-dimensional local analytic rings. Write $\mathfrak c$
for the conductor, and set
\[
 A=R/\mathfrak c,\qquad
 \overline A=\overline R/\mathfrak c,\qquad
 Q=\overline R/R=\overline A/A.
\]
We first give the exact duality comparison needed below.
The $R$-module $\overline R$ is maximal Cohen--Macaulay.
Since $R$ is Gorenstein,
$\operatorname{Ext}_R^j(\overline R,R)=0$ for $j>0$.
Its dual $\operatorname{Hom}_R(\overline R,R)$ is
precisely the conductor, viewed as a fractional ideal:
a homomorphism is multiplication in the total quotient ring,
and multiplication takes $\overline R$ into $R$ exactly
for conductor elements. Applying $\operatorname{Hom}_R(-,R)$
to $0\to R\to\overline R\to Q\to0$ therefore gives
\[
 \operatorname{Ext}_R^1(Q,R)=A.                    \tag{A.J.11}\label{cdp:eq:A.J.11}
\]
The depth lemma shows that $Q$ is a Cohen--Macaulay module
of dimension one. Gorenstein duality on such modules is the
involutive functor $M\mapsto\operatorname{Ext}_R^1(M,R)$.
Consequently $A$ is Cohen--Macaulay and
\[
 Q\simeq\operatorname{Ext}_R^1(A,R)=\omega_A,
 \qquad \operatorname{End}_A(Q)=A.                \tag{A.J.12}\label{cdp:eq:A.J.12}
\]
The final equality is the canonical multiplication map:
duality identifies its endomorphism ring with
$\operatorname{End}_A(A)=A$.

\noindent\textbf{Step 2.} The ordinary-cusp hypothesis says that $A$ is generically
reduced. Its nilradical is therefore supported at finitely
many points; a one-dimensional Cohen--Macaulay ring has
no nonzero submodule with that support. Hence $A$ is reduced.
Choose local coordinates $(s_i,z_i)$ on the normalization
branches with $C=(z_i=0)$. Put
\[
 B=\prod_i\mathbb C\{s_i\},\qquad
 K=\prod_i\operatorname{Frac}\bigl(\mathbb C\{s_i\}\bigr).
\]
The specified conductor gives the actual ring comparison
\[
 \overline A=B[\epsilon]/(\epsilon^2),             \tag{A.J.13}\label{cdp:eq:A.J.13}
\]
component by component. The projection
$\overline A\to B$ is injective on $A$, since its kernel
is nilpotent and $A$ is reduced. It identifies $B$ with
the finite normalization of the curve ring $A$.

\noindent\textbf{Step 3.} Under \eqref{cdp:eq:A.J.13} the inclusion of $A$ takes the form
\[
 a\longmapsto a+\delta(a)\epsilon,
\]
where $\delta:A\to B$ is a complex-linear analytic
derivation. It extends to the total quotient ring, and on
each normalization branch has the form
$\delta(a)=u_i\,da/ds_i$ for a meromorphic germ $u_i$.
Let $u=(u_i)$ and write derivatives componentwise.
The map
\[
 \Phi:B[\epsilon]/(\epsilon^2)\longrightarrow K,
 \qquad b+c\epsilon\longmapsto c-u b'
                                                        \tag{A.J.14}\label{cdp:eq:A.J.14}
\]
is $A$-linear for this specified action: multiplication by
$a+u a'\epsilon$ changes its value to
$ac+u a'b-u(ab)'=a(c-u b')$. Its kernel at every generic
point is exactly the graph of the extended derivation on
$K$. It therefore induces a generic isomorphism from $Q$
to its image in $K$. The module $Q\simeq\omega_A$ is
torsion-free on the reduced curve, so this induced map is
injective.

\noindent\textbf{Step 4.} Differentiation of convergent power-series germs onto
convergent power-series germs is surjective, by termwise
integration. Hence the image in \eqref{cdp:eq:A.J.14} is exactly
\[
 I=B+uB\subset K.                                \tag{A.J.15}\label{cdp:eq:A.J.15}
\]
It is an invertible fractional ideal on the semilocal
normalization $B$. Every $A$-linear endomorphism of this
rank-one module is multiplication by an element of $K$;
preserving $I$ is equivalent, branch by branch, to belonging
to $B$. Thus $\operatorname{End}_A(I)=B$ as subrings of
$K$. Combining this with the actual $A$-linear isomorphism
$Q\xrightarrow{\Phi}I$ and \eqref{cdp:eq:A.J.12} gives $A=B$.
Because $A$ is local, there is a single branch over $p$,
and its ring is regular. This proves that $D$ is smooth
and that $C\to D$ is an isomorphism there.

\noindent\textbf{Step 5.} Finally $A=B=\mathbb C\{s\}$ makes $u=\delta(s)$
holomorphic. The analytic coordinate change
$s'=s+u(s)z$ has invertible differential at $z=0$ and
identifies the image of $A$ modulo $z^2$ with the functions
of $s'$ alone. Since $R$ is the inverse image of $A$ under
$\overline R\to\overline A$, it follows that
\[
 R=\mathbb C\{s'\}+z^2\mathbb C\{s',z\}
   =\mathbb C\{s',z^2,z^3\}
   \simeq\mathbb C\{s',x,y\}/(y^2-x^3).
                                                        \tag{A.J.16}\label{cdp:eq:A.J.16}
\]
This is the asserted product germ.
\end{proof}






\begin{proposition}[Rational singularities for a vertically resolved ruled normalization]\label{cdp:st:A.5.8}
\StatementRecord{cdp:st:A.5.8}{proved}{Historical events: Claude Fable 5.2 (Anthropic), effort max, 16 November 2026: targeted proof check. GPT-6 Astra (OpenAI), ChatGPT Pro mode, effort Pro, S6\_09\_PAPERS, 17 November 2026: targeted proof check. Also checked in this writing. Pre-continuation writing events: GPT-6 Astra (OpenAI), ChatGPT Pro mode, effort Pro, S6\_10\_PAPERS, 18 November 2026: targeted proof check. Continuation metadata events: GPT-6 Astra (OpenAI), ChatGPT Pro mode, effort Pro, S6\_10\_PAPERS continuation, 19 November 2026: used and rechecked in this run; status-vocabulary reconciliation; metadata reconciliation only; no new mathematical check. \emph{Pre-continuation recorded proof scope.} The ruling is constant on every connected exceptional fibre and descends through the normal pushforward. Its Leray sequence and the faithful pushforward of the finite-support R1 sheaf force rational singularities and the stated coherent cohomology}{the hypotheses, notation and definitions stated here; \cref{cdp:eq:A.J.19}}

Let $\pi:Z\to T$ be a resolution of a compact connected
normal complex surface, where $Z$ is an irrational ruled
surface with ruling $p:Z\to B$. Suppose that every
exceptional curve of $\pi$ is vertical for $p$.
Then there is a unique holomorphic map
$\widetilde p:T\to B$ satisfying
$p=\widetilde p\circ\pi$, and
\[
 R\pi_*\mathcal O_Z=\mathcal O_T,
 \qquad R^1\widetilde p_*\mathcal O_T=0.           \tag{A.J.17}\label{cdp:eq:A.J.17}
\]
In particular all singularities of $T$ are rational. For
$g=g(B)$, its complete structure-sheaf cohomology dimensions
are
\[
 (h^0(T,\mathcal O_T),h^1(T,\mathcal O_T),
                         h^2(T,\mathcal O_T))=(1,g,0).
                                                        \tag{A.J.18}\label{cdp:eq:A.J.18}
\]
\end{proposition}


\begin{proof}
\noindent\textbf{Step 1.} The exceptional fibre of the proper birational map $\pi$
over a point is connected, because $T$ is normal and
$\pi_*\mathcal O_Z=\mathcal O_T$. Its curve components
are vertical by assumption, so their intersection
connectivity makes $p$ constant on that whole fibre.
It is also constant on the one-point fibres. We verify
holomorphic descent, rather than only descending the
underlying continuous map.

\noindent\textbf{Step 2.} Fix $t\in T$ and let $U\subset B$ be a coordinate disc
containing $p(\pi^{-1}(t))$. Properness makes
$\pi(Z\setminus p^{-1}(U))$ closed and disjoint from
$t$. There is therefore an open neighbourhood $W$ of $t$
with $\pi^{-1}(W)\subset p^{-1}(U)$. The coordinate
function of $p$ on $\pi^{-1}(W)$ is a section of
$\pi_*\mathcal O_Z(W)=\mathcal O_T(W)$, so descends
to a holomorphic function on $W$. Its values lie in the
coordinate image of $U$; the inverse coordinate gives a
holomorphic map $W\to U$. These local maps agree after
pullback to the surjective map $\pi$, hence agree on
overlaps. They construct the required unique
$\widetilde p$.

\noindent\textbf{Step 3.} Contracting vertical ruling $(-1)$-curves factors $p$ as
\[
 Z\xrightarrow{b} Z_0\xrightarrow{p_0}B,
\]
where $Z_0\to B$ is a $\mathbb P^1$-bundle and $b$
is a finite sequence of point blowups. Every such blowup
has direct image $\mathcal O$ and no higher direct image
of $\mathcal O$: the successive layers on its
infinitesimal exceptional curves are
$\mathcal O_{\mathbb P^1}(k)$ for $k\geq0$, whose
degree-one cohomology is zero, and formal functions gives
the assertion. The bundle $p_0$ has the same direct-image
property by the cohomology of $\mathbb P^1$ and local
triviality. Composition consequently gives
\[
 p_*\mathcal O_Z=\mathcal O_B,
 \qquad R^j p_*\mathcal O_Z=0\quad(j>0).          \tag{A.J.19}\label{cdp:eq:A.J.19}
\]

\noindent\textbf{Step 4.} The sheaf $\mathcal T=R^1\pi_*\mathcal O_Z$ has
finite support at the singular points of $T$. The
low-degree Leray sequence for $p=\widetilde p\circ\pi$
contains the actual arrows
\[
 0\longrightarrow R^1\widetilde p_*\mathcal O_T
 \longrightarrow R^1p_*\mathcal O_Z
 \longrightarrow\widetilde p_*\mathcal T
 \longrightarrow R^2\widetilde p_*\mathcal O_T.    \tag{A.J.20}\label{cdp:eq:A.J.20}
\]
The last term is zero: the nonconstant proper map
$\widetilde p$ has fibres of complex dimension one.
The middle term is zero by \eqref{cdp:eq:A.J.19}. Thus
$R^1\widetilde p_*\mathcal O_T=0$ and
$\widetilde p_*\mathcal T=0$. Pushforward is faithful
on a sheaf with finite support, since its stalk sections
at finitely many points cannot disappear under a proper
map. Hence $\mathcal T=0$. Higher $R^j\pi_*$ vanish
for $j>1$ by the dimension of its fibres. This proves
\eqref{cdp:eq:A.J.17} and the defining rational-singularity property.

\noindent\textbf{Step 5.} Finally \eqref{cdp:eq:A.J.17} identifies every
$H^q(T,\mathcal O_T)$ with $H^q(Z,\mathcal O_Z)$.
Equation \eqref{cdp:eq:A.J.19} identifies the latter with
$H^q(B,\mathcal O_B)$. A compact connected smooth curve
has dimensions $(1,g)$ in degrees zero and one and no
higher cohomology, proving \eqref{cdp:eq:A.J.18}.
\end{proof}






\begin{proposition}[Arithmetic genus two and the embedded delta-one comparison]\label{cdp:st:A.5.14}
\StatementRecord{cdp:st:A.5.14}{proved}{Historical events: Claude Fable 5.2 (Anthropic), effort high, 16 November 2026: carried as background without fault found. GPT-6 Astra (OpenAI), ChatGPT Pro mode, effort Pro, S6\_09\_PAPERS, 17 November 2026: targeted proof check. Also checked in this writing. Pre-continuation writing events: GPT-6 Astra (OpenAI), ChatGPT Pro mode, effort Pro, S6\_10\_PAPERS, 18 November 2026: targeted proof check. Continuation metadata events: GPT-6 Astra (OpenAI), ChatGPT Pro mode, effort Pro, S6\_10\_PAPERS continuation, 19 November 2026: used and rechecked in this run; status-vocabulary reconciliation; metadata reconciliation only; no new mathematical check. \emph{Pre-continuation recorded proof scope.} Canonical torsion and the full coherent Euler calculation bound the conductor image arithmetic genus. The colength-one local subring calculation gives an ordinary node or cusp, hence the actual lci ideal; the previously derived normal-jet theorem then forces the asserted constant elliptic types}{the hypotheses, notation and definitions stated here; \cref{cdp:eq:A.J.31}; \cref{cdp:st:A.5.5}}

Let $f:Y\to C_0$ be a proper connected-fibre map from a
smooth compact complex threefold with $H^2(Y,\Q)=0$ to
a smooth curve. Suppose its smooth general fibre has
trivial canonical bundle. Let $mS$ be a full fibre with
irreducible reduced nonnormal support, and let $D$ be
a reduced irreducible component of its conductor image.
Then
\[
                         p_a(D)\le2.                \tag{A.J.30}\label{cdp:eq:A.J.30}
\]
If the normalization of $D$ has positive genus, then
$D$ is a local complete intersection in $Y$; every
singular germ of $D$ is an ordinary node or an ordinary
cusp, and the sum of its delta invariants is at most
one. In particular, when the generic transverse type
along such a positive-genus $D$ is $A_2$, $A_3$, or an
ordinary triple point, $D$ is smooth elliptic and the
transverse analytic type is constant along all of $D$.
These conclusions do not require a nonconstant ambient
Picard pullback.
\end{proposition}


\begin{proof}
\noindent\textbf{Step 1.} The evaluation, adjunction and fibre-equation comparisons
in (A.CD.1) use no Picard assumption and apply with
rational second-cohomology vanishing. They give a torsion
dualizing line $\omega_S$. The divisor Riemann--Roch
calculation gives $\chi(\OO_S)=0$. Thus, for
$e=h^0(S,\omega_S)\in\{0,1\}$, the same Serre-duality
calculation as in (A.CD.2) gives
$h^1(\OO_S)=1+e$; if $e=1$, then
$\omega_S\simeq\OO_S$.

\noindent\textbf{Step 2.} Write $\eta:T\to S$ for normalization,
$\mathcal Q=\eta_*\OO_T/\OO_S$, and $\Gamma$ for
the full conductor image scheme. Equality of the
constant global functions on $S$ and $T$ gives the
actual injection $H^0(\mathcal Q)\hookrightarrow
H^1(\OO_S)$. Conductor duality and curve Serre duality
therefore give
\[
 h^1(\Gamma,\omega_S|_\Gamma)
       =h^0(S,\mathcal Q)\le1+e.                  \tag{A.J.31}\label{cdp:eq:A.J.31}
\]
The restriction quotient
$\omega_S|_\Gamma\twoheadrightarrow\omega_S|_D$
has a kernel supported in dimension at most one, so
its $H^1$ map is surjective. For $e=1$ it follows that
$p_a(D)=h^1(\OO_D)\le2$. For $e=0$, put
$L=\omega_S|_D$. This torsion line has degree zero,
and \eqref{cdp:eq:A.J.31} gives $h^1(D,L)\le1$. Curve
Riemann--Roch gives
\[
 p_a(D)=1+h^1(D,L)-h^0(D,L)\le2,
\]
proving \eqref{cdp:eq:A.J.30} in both cases.

\noindent\textbf{Step 3.} Let $\nu:\widehat D\to D$ be normalization. The
normalization sequence gives the exact numerical
comparison
\[
 p_a(D)=g(\widehat D)+
        \sum_{x\in D}\dim_{\C}
          (\nu_*\OO_{\widehat D}/\OO_D)_x.       \tag{A.J.32}\label{cdp:eq:A.J.32}
\]
If $g(\widehat D)>0$, every local delta is consequently
at most one and their sum is at most one. We give the
analytic local rings at a point with delta one.
Let $A=\OO_{D,x}$ and let
$B=\bigoplus_{i=1}^r\C\{t_i\}$ be its finite
normalization. Here $\dim_{\C}(B/A)=1$. The residue map
of the local ring $A$ has the same value on every
normalization branch. The resulting quotient of
$\C^r$ by its diagonal has dimension $r-1$, so
$r\le2$.

\noindent\textbf{Step 4.} For $r=2$, the subring
\[
 A_0=\{(b_1,b_2)\in\C\{t_1\}\oplus\C\{t_2\}:
                                  b_1(0)=b_2(0)\}
\]
contains $A$ and has colength one in $B$. Hence
$A=A_0$. The explicit homomorphism
$\C\{u,v\}/(uv)\to A_0$ sends
$u$ to $(t_1,0)$, $v$ to $(0,t_2)$, and a constant
to the corresponding equal pair. Expanding a
convergent series shows that it is an isomorphism.
This is the ordinary node germ.

\noindent\textbf{Step 5.} For $r=1$, an element of $A$ with valuation one would
be a local coordinate on the normalization disc by
the holomorphic inverse function theorem. Convergent
composition in the analytic algebra $A$ would then
give $\C\{t\}\subset A$, contrary to its colength
one. Therefore every nonconstant element of $A$ has
valuation at least two. We have
$A\subset\C+t^2\C\{t\}$; both rings have
colength one in $\C\{t\}$, so equality holds.
The map
\[
 \C\{u,v\}/(v^2-u^3)\longrightarrow A,
                 \qquad (u,v)\longmapsto(t^2,t^3)
\]
is an isomorphism: every series with no linear term
is uniquely the sum of an even series in $t^2$ and
$t^3$ times a series in $t^2$. This is the ordinary
cusp germ. A delta-zero germ is its smooth
normalization disc.

\noindent\textbf{Step 6.} We finally check the embedding into the specified
threefold. These analytic rings have embedding
dimension at most two. In local coordinates of the
smooth threefold, the tangent space of $D$ therefore
has codimension at least one. There is an element
$z$ of its ideal with $dz(x)\ne0$. The inverse
function theorem makes $z$ an ambient coordinate,
and the ideal of $D$ is the inverse image of a
height-one ideal in the regular analytic ring
$\C\{u,v\}=\OO_{Y,x}/(z)$. This ideal is principal:
a pure reduced curve in a smooth surface is a Cartier
divisor, or equivalently its reduced height-one ideal
is a product of the prime factors in this factorial
ring. Thus $I_{D,x}=(z,h)$, a regular sequence.
This proves the asserted lci embedding, including
at the node and cusp. Statement \ref{cdp:st:A.5.5} now applies
to the three specified generic transverse types
and proves the final smooth elliptic and constant-type
conclusion.
\end{proof}






\begin{proposition}[Eliminating vertical conductor on a smooth ruled cusp normalization]\label{cdp:st:A.5.10}
\StatementRecord{cdp:st:A.5.10}{proved}{Historical events: Claude Fable 5.2 (Anthropic), effort max, 16 November 2026: targeted proof check. GPT-6 Astra (OpenAI), ChatGPT Pro mode, effort Pro, S6\_09\_PAPERS, 17 November 2026: targeted proof check. Also checked in this writing. Pre-continuation writing events: GPT-6 Astra (OpenAI), ChatGPT Pro mode, effort Pro, S6\_10\_PAPERS, 18 November 2026: targeted proof check. Continuation metadata events: GPT-6 Astra (OpenAI), ChatGPT Pro mode, effort Pro, S6\_10\_PAPERS continuation, 19 November 2026: used and rechecked in this run; status-vocabulary reconciliation; metadata reconciliation only; no new mathematical check. \emph{Pre-continuation recorded proof scope.} The actual saturated line and rank-two quotient, the bound twice the branch divisor at most the conductor, and both epsilon cases of the quadratic Chern identity were re-derived. Equality forces genus one, no vertical divisor or critical ruling divisor, and smooth minimal normalization}{the hypotheses, notation and definitions stated here; \cref{cdp:eq:A.J.26}; \cref{cdp:eq:A.J.23}; \cref{cdp:eq:A.J.21}; \cref{cdp:eq:A.J.27}; \cref{cdp:eq:A.J.22}; \cref{cdp:st:A.5.7}}

Let $Y$ be a smooth compact complex threefold with
$H^2(Y,\mathbb Q)=0$, let $S\subset Y$ be a reduced
irreducible divisor, and suppose its normalization
$h:Z\to S\subset Y$ is a smooth irrational
ruled surface $p:Z\to B$.
Suppose its conductor divisor is $D_Z=2C+V$, where
$C$ is a ruling section and $V$ is vertical, and
the generic transversal along $h(C)$ is an ordinary cusp.
Then $B$ is elliptic, $V=0$, $C^2=0$, and $Z$ is
already minimal.
The image $h(C)$ is smooth elliptic and $S$ is locally
an ordinary-cusp product along it. No local
complete-intersection hypothesis on the image curve is
required. The conductor-image scheme along the vertical
components may initially be nonreduced.
\end{proposition}


\begin{proof}
\noindent\textbf{Step 1.} Let $g=g(B)\geq1$, let $F$ be the full ruling fibre
class, and write $V=\sum_i c_i\Gamma_i$ for its
irreducible vertical components. Put
$v=V\cdot C$, $n=C^2$ and $k=2g-2$.
Conductor duality and Cartier adjunction give the exact
line-bundle comparison
\[
 K_Z\otimes\mathcal O_Z(2C+V)
       =h^*(K_Y\otimes\mathcal O_Y(S)).
                                                        \tag{A.J.21}\label{cdp:eq:A.J.21}
\]
Its right side has rational first Chern class zero.
Intersecting with $C$ and using adjunction on the section
gives
\[
 n=2-2g-v=-k-v.                                  \tag{A.J.22}\label{cdp:eq:A.J.22}
\]
Let $D_p$ be the divisorial common zero of
$dp:T_Z\to p^*T_B$. Its remaining zero scheme is
zero-dimensional, say $\Xi$, so the image is
$I_\Xi\otimes p^*T_B(-D_p)$.
The kernel $T_p$ is an invertible sheaf, by the depth
lemma. Taking determinants gives the actual comparison
\[
 T_p=K_Z^{-1}\otimes p^*K_B\otimes\mathcal O_Z(D_p).
\]
The section makes $dp$ surjective at every point of $C$.
Consequently $D_p$ and $\Xi$ are disjoint from $C$,
and $T_p|_C=N_{C/Z}$.

\noindent\textbf{Step 2.} Put $\mathcal F=h^*T_Y$. Poincare duality on the compact
oriented six-manifold $Y$ also gives
$H^4(Y,\mathbb Q)=0$. Therefore
\[
 c_1(\mathcal F)=0,\qquad c_2(\mathcal F)=0
        \quad\hbox{in rational cohomology}.       \tag{A.J.23}\label{cdp:eq:A.J.23}
\]
The differential is generically injective. Let
$L\subset\mathcal F$ be the saturation of the image
of $T_p\xrightarrow{dh}\mathcal F$.
It is invertible: a saturated rank-one subsheaf of a
locally free sheaf on a smooth surface is reflexive.
The map $T_p\to L$ has no vertical divisorial zero.
Indeed $h$ is finite, so no ruling fibre is contracted;
in characteristic zero its derivative along that fibre
is nonzero generically.

\noindent\textbf{Step 3.} At the generic point of $C$, an ordinary cusp has
normalization differential equivalent to that of
$(u,s)\mapsto(u^2,u^3,s)$.
The ruling tangent either lies in its kernel along $C$
or does not. The vanishing order in the first case is
exactly one, and in the second case it is zero.
Write $\epsilon=1$ or $\epsilon=0$ for these respective
alternatives. There are no other horizontal conductor
components, and the normalization is an immersion away
from the conductor. Hence the actual line comparison is
\[
 L=T_p\otimes\mathcal O_Z(\epsilon C).
                                                        \tag{A.J.24}\label{cdp:eq:A.J.24}
\]
This checks both positions of the derivative kernel
relative to the ruling tangent.

\noindent\textbf{Step 4.} The quotient $E=\mathcal F/L$ is torsion-free of rank
two. The differential induces a generically nonzero map
\[
 T_Z/T_p=I_\Xi\otimes p^*T_B(-D_p)
                    \longrightarrow E\longrightarrow E^{**}.
\]
It extends uniquely to a map from $p^*T_B(-D_p)$ into
$E^{**}$, since both are locally free and the omitted
set has codimension two; locally this is Hartogs
extension of its coefficient functions.
Let $M$ be its saturated image line in the rank-two
bundle $E^{**}$. At the generic point of $C$, the
two-by-two differential minors have common vanishing
order one. Removing the order $\epsilon$ already
assigned to $L$ leaves order $1-\epsilon$ for this map.
At the generic point of $\Gamma_i$, take a plane transversal
to its image and a normalization parameter $u$ for the
corresponding plane-curve branch. If $m_i$ is its
multiplicity, its derivative modulo the nonzero tangent
along $\Gamma_i$ vanishes to order $r_i=m_i-1$.
These are all codimension-one zeros. Put
$R=\sum_i r_i\Gamma_i$ and $r=R\cdot C$. We obtain
\[
 M=p^*T_B(-D_p)\otimes
           \mathcal O_Z\bigl((1-\epsilon)C+R\bigr).
                                                        \tag{A.J.25}\label{cdp:eq:A.J.25}
\]

\noindent\textbf{Step 5.} The plane-curve conductor gives the inequalities
\[
 2R\leq V,\qquad 2r\leq v.                      \tag{A.J.26}\label{cdp:eq:A.J.26}
\]
To spell out their local length content, on a single
plane branch of multiplicity $m_i$ the classes
$u,u^2,\ldots,u^{m_i-1}$ are linearly independent in
the normalization quotient. Its delta invariant is thus
at least $m_i-1$. The branch is Gorenstein. Finite
duality identifies the length of its ring modulo its
conductor with the length of its normalization quotient;
the intrinsic conductor exponent is consequently twice
that delta invariant. Other branches add their
nonnegative intersection multiplicities to this exponent.
For the last comparison, writing the plane equation
as $\prod_j f_j$ multiplies the residue denominator on
branch $i$ by $\prod_{j\ne i}f_j$, whose valuation is
exactly the sum of those intersection multiplicities.
This proves \eqref{cdp:eq:A.J.26} for every branch, including cases
with nonreduced conductor-image scheme.

\noindent\textbf{Step 6.} Let $\ell_1$ be the length of $E^{**}/E$. Saturation
of $M$ gives a zero-dimensional subscheme $\Xi_M$ and a
line bundle $N$ with
\[
 0\longrightarrow M\longrightarrow E^{**}
       \longrightarrow I_{\Xi_M}\otimes N\longrightarrow0.
\]
Write $\ell_2=\operatorname{length}\Xi_M$.
First Chern classes give $c_1(N)=-c_1(L)-c_1(M)$.
The second Chern classes in these typed quotient
sequences and \eqref{cdp:eq:A.J.23} give the intersection identity
\[
 L^2+L\cdot M+M^2=\ell_1+\ell_2\geq0.            \tag{A.J.27}\label{cdp:eq:A.J.27}
\]
Indeed $c_2(E)=L^2$,
$c_2(E^{**})=L^2-\ell_1$, and
$c_2(E^{**})=-L\cdot M-M^2+\ell_2$.

\noindent\textbf{Step 7.} Set $W=V+D_p$ and $B_v=R-D_p$, both vertical
divisors, the first effective. The exact comparisons
above and \eqref{cdp:eq:A.J.21} give the numerical classes
\[
 L\equiv(2+\epsilon)C+W+kF,\qquad
 M\equiv(1-\epsilon)C+B_v-kF.
\]
The intersection form on vertical divisors is negative
semidefinite. One can check this directly by contracting
the ruling blowups: every vertical class is a multiple
of the pullback fibre class plus exceptional classes,
on which the form is negative definite. Hence
\[
 q_v=W^2+W\cdot B_v+B_v^2
      =(W+B_v/2)^2+3B_v^2/4\leq0.
\]
Its vanishing implies $W^2=B_v^2=0$; these classes
are then numerical multiples of the fibre class.

\noindent\textbf{Step 8.} For $\epsilon=1$, using $D_p\cdot C=0$ in the
left side of \eqref{cdp:eq:A.J.27} gives
\[
 \ell_1+\ell_2=-6k-3v+3r+q_v
                  \leq-6k-3v/2.
\]
Nonnegativity, with $k\geq0$ and $v\geq0$,
forces $g=1$, $v=0$, $r=0$ and $q_v=0$.

\noindent\textbf{Step 9.} For $\epsilon=0$, the same intersection calculation reads
\[
 0\leq\ell_1+\ell_2
       =-6k-2v+4r+q_v\leq-6k.                   \tag{A.J.28}\label{cdp:eq:A.J.28}
\]
It forces $g=1$, $\ell_1=\ell_2=0$, $2r=v$,
and $q_v=0$.
Thus $E=E^{**}$ is locally free, and
$L\to\mathcal F$ is a subbundle inclusion.
Along $C$, the composite
$T_C\to\mathcal F|_C\to E|_C$ vanishes generically,
because $dh$ has rank one there and its nonzero ruling
tangent image is $L|_C$. It therefore vanishes
everywhere. The exact subbundle sequence gives an
actual nonzero holomorphic line-bundle map
\[
 T_C\longrightarrow L|_C=T_p|_C=N_{C/Z}.
                                                        \tag{A.J.29}\label{cdp:eq:A.J.29}
\]
It is nonzero since the finite map $C\to h(C)$ has
nonzero derivative generically. Its Hom line has degree
$n-(2-2g)=-v$ by \eqref{cdp:eq:A.J.22}. A nonzero section on a
compact curve has nonnegative degree, so $v=0$ in this
alternative too.

\noindent\textbf{Step 10.} In both alternatives, $v=r=0$ and $q_v=0$.
Thus $W$ is numerically a fibre multiple and
$W\cdot C=V\cdot C+D_p\cdot C=0$, so it is
numerically zero. An effective nonzero divisor on the
projective ruled surface has positive intersection
with an ample class. It follows that $W=0$ and hence
$V=D_p=0$. Now $B_v=R$ is effective, has square zero
and has zero intersection with $C$, so the same
argument gives $R=0$. Equation \eqref{cdp:eq:A.J.22} gives $C^2=0$.
Finally $K_Z+2C$ is numerically trivial. A ruling
$(-1)$-curve $E_0$ would satisfy
$0=(K_Z+2C)\cdot E_0=-1+2C\cdot E_0$, impossible
because $C\cdot E_0$ is an integer. There are therefore
no ruling blowups and $Z$ is minimal.

\noindent\textbf{Step 11.} Statement \ref{cdp:st:A.5.7} now applies everywhere along
the conductor: the normalization is smooth and its
conductor is exactly $2C$. It makes the image smooth
and the surface locally the cusp product. The image is
elliptic because $C\to h(C)$ is an isomorphism.
\end{proof}






\begin{proposition}[A noncontracted horizontal cusp has zero Euler number]\label{cdp:st:A.5.20}
\StatementRecord{cdp:st:A.5.20}{proved}{Historical events: Claude Fable 5.2 (Anthropic), effort high, 16 November 2026: carried as background without fault found. GPT-6 Astra (OpenAI), ChatGPT Pro mode, effort Pro, S6\_09\_PAPERS, 17 November 2026: targeted proof check. Also checked in this writing. Pre-continuation writing events: GPT-6 Astra (OpenAI), ChatGPT Pro mode, effort Pro, S6\_10\_PAPERS, 18 November 2026: targeted proof check. Continuation metadata events: GPT-6 Astra (OpenAI), ChatGPT Pro mode, effort Pro, S6\_10\_PAPERS continuation, 19 November 2026: used and rechecked in this run; status-vocabulary reconciliation; metadata reconciliation only; no new mathematical check. \emph{Pre-continuation recorded proof scope.} The degree-one image, conductor-genus bound and jet calculation give a smooth elliptic constant cusp. The normalized cusp powers have degree zero; the first possible point blowup contradicts the actual canonical-conductor divisor identity, forcing smooth minimal normalization}{the hypotheses, notation and definitions stated here; \cref{cdp:st:A.5.14}; \cref{cdp:st:A.5.5}; \cref{cdp:st:C2d}; \cref{cdp:st:A.5.18}}

Let $f:Y\to C_0$ be a proper holomorphic map from a smooth
compact complex threefold with $b_2(Y)=0$, with connected
fibres and smooth general fibre a complex two-torus.
Let $mS$ be a full fibre with irreducible reduced support,
$\eta:T\to S$ its normalization and $\pi:Z\to T$ an
irrational ruled resolution with ruling $p:Z\to B$.
Suppose every exceptional curve of $\pi$ is vertical
and the resolution conductor is $2C+V$, where $C$
is a noncontracted section and $V$ is vertical.
If the generic transversal at $D=\eta\pi(C)$ is an
ordinary cusp, then $D$ and $C$ are smooth elliptic
curves and $C\to D$ is an isomorphism. The cusp type
is constant along $D$. The normalization $T$ is a
smooth minimal ruled surface over this elliptic curve,
its conductor is twice a section, and $e(S)=0$.
This permits an initially singular normalization;
smoothness is part of the conclusion.
\end{proposition}


\begin{proof}
\noindent\textbf{Step 1.} A transverse cusp has exactly one normalization
branch. The map from the horizontal section to the
normalization of its image consequently has degree
one. Both curves are smooth, so this finite map
identifies $C\simeq B$ with $D^\nu$. Its genus is
positive. Statement \ref{cdp:st:A.5.14}, including its actual
delta-one lci comparison, and \ref{cdp:st:A.5.5} therefore make
$D$ itself smooth elliptic and its cusp transversal
constant at every point. Thus $C\to D$ is an
isomorphism, without any assumption about the
ambient Picard pullback.

\noindent\textbf{Step 2.} For the exact local normalization $x=u^3,y=u^2$,
let $C_T\subset T$ be the image of $C$. The
normalization is smooth near this curve. Its normal
line $P=N_{C_T/T}$ has $P^2=R$ and $P^3=Q$
in the degree-zero radical and Hessian quotient
lines of \ref{cdp:st:C2d}. Hence $\deg P=0$. The normalization
is smooth near $C$, and no other conductor curve
meets it. Factor $\pi$ through its minimal
resolution, descending the ruling across its
vertical exceptional blowdowns as in \ref{cdp:st:A.5.18}.
On that minimal resolution the section is isomorphic
to $C_T$ in its unchanged smooth neighbourhood. It
therefore has self-intersection zero and is disjoint from
the effective vertical conductor. The canonical
comparison is the actual line-bundle identity
$K+2C+V=(\eta\pi)^*\omega_S$, of zero rational
Chern class.

\noindent\textbf{Step 3.} Apply (A.T.22) to each ruling blowdown. Its
exceptional coefficient is exactly the multiplicity
at its centre minus one. No ruling $(-1)$-curve
can meet $C$, since the intersection with $K+2C+V$
would then be positive. On the minimal ruled
bundle the vertical conductor is zero because it
is disjoint from a section. A first reverse blowup
away from that section would have coefficient
$-1$, so there were no ruling blowups. All
resolution-exceptional curves are vertical, and
the square-zero fibres of the resulting bundle
cannot be exceptional. This identifies the minimal
resolution with $T$ itself and proves the asserted
smoothness and minimality.

\noindent\textbf{Step 4.} Finally the normalization is an isomorphism off
$C$ and $D$. The exact Euler comparison is
$e(S)=e(T)-e(C)+e(D)$; all three terms vanish,
since $T$ is ruled over an elliptic curve and
$C,D$ are elliptic. Thus $e(S)=0$, also for the
underlying topological space of the full fibre $mS$.
\end{proof}






\begin{proposition}[Hirzebruch germs cannot occur in a torus degeneration]\label{cdp:st:A.5.15}
\StatementRecord{cdp:st:A.5.15}{proved}{Historical events: Claude Fable 5.2 (Anthropic), effort high, 16 November 2026: carried as background without fault found. GPT-6 Astra (OpenAI), ChatGPT Pro mode, effort Pro, S6\_09\_PAPERS, 17 November 2026: targeted proof check. Also checked in this writing. Pre-continuation writing events: GPT-6 Astra (OpenAI), ChatGPT Pro mode, effort Pro, S6\_10\_PAPERS, 18 November 2026: targeted proof check; used and rechecked in this run. Continuation metadata events: GPT-6 Astra (OpenAI), ChatGPT Pro mode, effort Pro, S6\_10\_PAPERS continuation, 19 November 2026: used and rechecked in this run; status-vocabulary reconciliation; metadata reconciliation only; no new mathematical check. \emph{Pre-continuation recorded proof scope.} The explicit two-endpoint circle construction, torus-link linking number, nonzero self-intersection and weighted scaling into the actual torus fibre were re-derived. Nemethi1991 TheoremL1, Remark2.2 and Example2.4(a) pp2-3 were opened. Michel-Pichon-Weber2005 Theorems4.1-4.2 pp6-7 and Lemma4.6 with proof p11 were opened for the independent Seifert check, retaining that source's gcd(m,k,l)=1 hypothesis. The direct cycle construction treats the full stated range}{the hypotheses, notation and definitions stated here; \cref{cdp:eq:A.J.33a}; \cref{cdp:eq:A.J.33b}; \cite[Theorem L1 and Example 2.4(a)]{Nemethi1991Join}; \cite[Theorem 4.1 and Lemma 4.6]{MPW2005Hirzebruch}; \cref{cdp:eq:A.J.35}; \cite[Theorem 4.2(2)]{MPW2005Hirzebruch}; \cref{cdp:eq:A.J.34}; \cref{cdp:eq:A.J.37}}

Let $f:Y\to C_0$ be a proper holomorphic map from a
smooth complex threefold to a smooth curve, with smooth
nearby fibres complex two-tori. There is no point of
a fibre at which, in holomorphic ambient coordinates
$(x,y,z)$ and a coordinate on the base, the map has
the equation
\[
 f(x,y,z)=u(x,y,z)\bigl(x^a-y^bz^c\bigr)^m,
 \quad a\ge2,\quad b,c,m\ge1,                                    \tag{A.J.33}\label{cdp:eq:A.J.33}
\]
where $u$ is a holomorphic unit. No Betti-number
assumption on $Y$ is needed. This includes the pinch
point, every germ $x^2-y^{2k}z^{2r+1}$ with
$k\ge1$, $r\ge0$, and the reducible germ
$x^2-y^2z^2$.
\end{proposition}


\begin{proof}
\noindent\textbf{Step 1.} Put $h=x^a-y^bz^c$. We first compare the actual map
with its reduced equation. This polynomial is weighted
homogeneous of degree $a(b+c)$ for weights
$(b+c,a,a)$. On a small simply connected coordinate
neighbourhood choose a holomorphic logarithm of $u$.
Define the holomorphic coordinate change
\[
 \Phi(x,y,z)=\left(
 x u^{1/(ma)},\;
 y u^{1/(m(b+c))},\;
 z u^{1/(m(b+c))}\right).
                                                        \tag{A.J.34}\label{cdp:eq:A.J.34}
\]
Its derivative at the origin is an invertible diagonal
matrix. The holomorphic inverse function theorem and
weighted homogeneity give
$h\circ\Phi=u^{1/m}h$ and therefore
$f=(h\circ\Phi)^m$ exactly. For a sufficiently small
nonzero regular value $t$, a chosen Milnor fibre of
$h$ at one of the $m$ roots of $t$ is consequently a
compact oriented submanifold with boundary of the
smooth global fibre $Y_t$. Denote this Milnor fibre
by $M$.

\noindent\textbf{Step 2.} We first give a direct spherical cycle, without an
irreducibility assumption. Put $d=\gcd(b,c)$,
$b'=b/d$ and $c'=c/d$, and consider the smooth affine
curve
\[
 B=\{(x,v)\in\C^2:x^a-v^d=1\}.
\]
The affine surface $H=\{h=1\}$ maps holomorphically
to $B$ by $v=y^{b'}z^{c'}$. The curve $B$ is
connected: projection to the $v$-line has $a$ sheets,
and a small loop about any root of $v^d+1$ cyclically
permutes them. It is smooth because its two partial
derivatives cannot vanish on $x^a-v^d=1$.
The function $v$ has $a$ distinct simple zeros
$P_j=(\zeta_a^j,0)$ on $B$.

\noindent\textbf{Step 3.} Choose two of these zeros, $P$ and $Q$, and an
embedded arc $\gamma:[0,1]\to B$ joining them and
avoiding the other zeros. Choose smooth complex
functions $A(t),B_0(t)$ on its interior such that
\[
 A(t)^{b'}B_0(t)^{c'}=v(\gamma(t)),\qquad
 |A(t)|=|v(\gamma(t))|^{1/(2b')},\quad
 |B_0(t)|=|v(\gamma(t))|^{1/(2c')}.
\]
They exist by choosing a continuous argument along
the arc. The continuous map
\[
 (t,\theta)\longmapsto
 \bigl(x(\gamma(t)),A(t)e^{\mathrm i c'\theta},
                 B_0(t)e^{-\mathrm i b'\theta}\bigr)
                                                        \tag{A.J.33a}\label{cdp:eq:A.J.33a}
\]
from $[0,1]\times S^1$ to $H$ collapses each boundary
circle to the corresponding point $P$ or $Q$.
It therefore defines an oriented spherical cycle
$\sigma_\gamma:S^2\to H$. Coprimality of $b'$ and
$c'$ makes each interior circle embedded. The whole
map is injective away from the collapsed ends,
because the base arc is embedded.

\noindent\textbf{Step 4.} Let $\gamma'$ be a nearby arc with the same endpoints
and disjoint interior, obtained by an isotopy in a
thin disc neighbourhood of $\gamma$. The isotopy
keeps all other zeros of $v$ outside that disc.
Choose the arguments in \eqref{cdp:eq:A.J.33a} continuously during
the isotopy; both coordinates tend uniformly to zero
at the ends. This gives a continuous homotopy of
maps from $S^2$, so the resulting spherical cycles
are homologous. Their supports meet only at $P$ and
$Q$. Near either endpoint, $(y,z)$ are actual
holomorphic coordinates on $H$, since $a x^{a-1}$
is nonzero there. Arrange the two arcs to be radial
with different arguments in the local $v$-coordinate.
On a sufficiently small sphere in these coordinates,
the links of the two cycles are disjoint parallel
curves of slope $(c',-b')$ on a Clifford torus.
Their oriented linking number is $-b'c'$.
Indeed the torus framing of a $(p,q)$ curve has
linking number $pq$ with its parallel copy: push the
two curves to the opposite sides of the Heegaard
splitting of $S^3$ into its two solid tori and count
meridian discs, obtaining $p q$ signed intersections.
Here $p=c'$ and $q=-b'$.

\noindent\textbf{Step 5.} The local intersection of the two coned cycles is
that linking number. This equality follows directly
by replacing the two cones inside the small ball by
transverse relative two-chains with the same boundary
links; their intersection is the defining linking
number in its boundary three-sphere. At $Q$ both
cycle orientations reverse relative to the outgoing
radial coordinate, so the two reversals leave the
local intersection sign unchanged. Consequently
\[
 [\sigma_\gamma]^2
   =[\sigma_\gamma]\cdot[\sigma_{\gamma'}]
   =-2b'c'\ne0.                                  \tag{A.J.33b}\label{cdp:eq:A.J.33b}
\]
Both cycles and their homotopy have compact support.
The weighted homothety
$(x,y,z)\mapsto(s^{b+c}x,s^a y,s^a z)$ maps
$H$ biholomorphically onto $\{h=s^{a(b+c)}\}$.
For sufficiently small positive $s$, it puts those
compact supports into any prescribed Milnor
neighbourhood of the origin. Thus \eqref{cdp:eq:A.J.33b} is a
spherical class of nonzero square in an actual
local Milnor fibre $M$.

\noindent\textbf{Step 6.} For an independent check when $\gcd(a,b,c)=1$,
we compute the full homotopy type and the rank of
the intersection form from two primary results.
The join theorem
\cite[Theorem L1 and Example 2.4(a)]{Nemethi1991Join}
applies to arbitrary analytic germs, including
non-isolated ones. Set $d=\gcd(b,c)$. The Milnor
fibre of $y^bz^c$ is a disjoint union of $d$ annuli:
its equations in a Milnor polydisc are
$y^{b/d}z^{c/d}=\zeta\epsilon^{1/d}$, and each
component is parametrized by
$y=A w^{c/d}$, $z=B w^{-b/d}$. Coprimality makes
this parametrization injective, and the two radius
bounds give a nonempty annulus in $w$.
The Milnor fibre of $x^a$ consists of $a$ points.
The join theorem thus identifies $M$, by a homotopy
equivalence, with the union of $a$ cones on a
disjoint union of $d$ circles, all glued along their
common base. Contracting one cone and then choosing
arcs between the $d$ base components gives
\[
 M\simeq
 \bigvee^{(a-1)d}S^2\ \vee\
 \bigvee^{(a-1)(d-1)}S^1.                          \tag{A.J.35}\label{cdp:eq:A.J.35}
\]
In particular, every class in $H_2(M,\Z)$ is a sum
of images of spherical fundamental classes.

\noindent\textbf{Step 7.} For the same irreducible germ, the boundary calculation
of \cite[Theorem 4.1 and Lemma 4.6]{MPW2005Hirzebruch}
gives a Seifert three-manifold with base genus
\[
 g_\partial=(a-1)(d-1),\qquad
 e_0=\frac{a}{(b/d)(c/d)}>0.                     \tag{A.J.36}\label{cdp:eq:A.J.36}
\]
Here $e_0$ is its rational Seifert Euler number in
the orientation convention of that paper. Its
nonzero value gives
$b_1(\partial M)=2g_\partial$. To make this last
deduction explicit, the rational abelianized Seifert
presentation has a regular-fibre generator $v$,
exceptional-fibre generators $q_i$, the relations
$\alpha_iq_i+\beta_iv=0$ and
$\sum_iq_i+ev=0$, and $2g_\partial$ free base
generators. Eliminating the $q_i$ gives $e_0v=0$.
Thus all fibre generators vanish over $\Q$, leaving
exactly the $2g_\partial$ base generators.

\noindent\textbf{Step 8.} Let $I_M:H_2(M,\Q)\to H_2(M,\partial M;\Q)$ be
the intersection morphism. Poincare--Lefschetz duality
and the homology sequence of the pair give
\[
 \operatorname{rank} I_M
   =b_2(M)+b_1(M)-b_1(\partial M)=a-1>0.          \tag{A.J.37}\label{cdp:eq:A.J.37}
\]
Indeed $H_3(M)=0$ by \eqref{cdp:eq:A.J.35}, so
$H_3(M,\partial M)\simeq H^1(M)$ injects into
$H_2(\partial M)$. Its cokernel is the kernel of
$I_M$, which has dimension
$b_1(\partial M)-b_1(M)$. This proves the displayed
rank formula. For $a=2$, $b=1$, $c=\ell$, the
more precise boundary is $L(2\ell,1)$ by
\cite[Theorem 4.2(2)]{MPW2005Hirzebruch}; the same
pair sequence over $\Z$ makes the rank-one
intersection form have absolute determinant
$2\ell$. The special case $\ell=1$ is the ordinary
quadratic singularity and has determinant two.

\noindent\textbf{Step 9.} Finally let $j:M\hookrightarrow Y_t$ be the actual
embedding obtained from \eqref{cdp:eq:A.J.34}. A complex two-torus
has contractible universal cover $\C^2$, so
$\pi_2(Y_t)=0$. In particular the spherical class
constructed in \eqref{cdp:eq:A.J.33a} maps to zero in
$H_2(Y_t,\Z)$. Intersection numbers of compact
two-cycles in the interior of $M$ are preserved by
this oriented codimension-zero embedding: the same
small transverse perturbations compute them in
both manifolds. Its square in $Y_t$ is therefore
both zero and $-2b'c'$, contradicting \eqref{cdp:eq:A.J.33b}.
In the irreducible case, \eqref{cdp:eq:A.J.35}--\eqref{cdp:eq:A.J.37} provide
an additional contradiction on the whole intersection
space.
This proves the exclusion for the actual map,
including every multiplicity in \eqref{cdp:eq:A.J.33}.
\end{proof}



The homotopy and boundary theorems just used are
established results; our application is the displayed
embedding and intersection comparison. The explicit
cycle also treats the reducible monomials. The exact
Seifert arithmetic, homology dimensions and ranks
are checked in \path{code/E_Hirzebruch.py}. No
reduction of an arbitrary nonordinary surface germ
to \eqref{cdp:eq:A.J.33} is asserted.




\begin{proposition}[Isolated singularities are excluded and suspension germs are transverse products]\label{cdp:st:A.5.17}
\StatementRecord{cdp:st:A.5.17}{proved}{Historical events: Claude Fable 5.2 (Anthropic), effort high, 16 November 2026: carried as background without fault found. GPT-6 Astra (OpenAI), ChatGPT Pro mode, effort Pro, S6\_09\_PAPERS, 17 November 2026: targeted proof check. Also checked in this writing. Pre-continuation writing events: GPT-6 Astra (OpenAI), ChatGPT Pro mode, effort Pro, S6\_10\_PAPERS, 18 November 2026: targeted proof check; used and rechecked in this run. Continuation metadata events: GPT-6 Astra (OpenAI), ChatGPT Pro mode, effort Pro, S6\_10\_PAPERS continuation, 19 November 2026: used and rechecked in this run; status-vocabulary reconciliation; metadata reconciliation only; no new mathematical check. \emph{Pre-continuation recorded proof scope.} The isolated critical-point Morsification, product-domain boundary transversality for all factor multiplicities, transport through the connected discriminant complement, two-factor monomial model and singular one-factor Morse model were re-derived. The remaining smooth factor yields the exact product A-type germ. Milnor Lemma6.1 with proof pp55-56 and AppendixB LemmasB1-B3 and positivity proof pp111-114 were read in a reproduction of the original monograph; the receiving proof supplies the analytic boundary transport beyond the polynomial statement of Lemma6.1}{the hypotheses, notation and definitions stated here; \cref{cdp:st:A.5.15}; \cite[Lemma 6.1 and Appendix B]{Milnor1968}; \cref{cdp:eq:A.J.40}; \cref{cdp:eq:A.J.39}}

Let $f:Y\to C_0$ have the hypotheses of \ref{cdp:st:A.5.15}.
No local equation $f=u h^m$, with $u$ a holomorphic
unit and $m\ge1$, can have a reduced equation $h$
with an isolated critical point at the origin.
Suppose its local equation at a point of a fibre is
\[
             f=(x^a-g(y,z))^m,
             \qquad a\ge2,\quad m\ge1,             \tag{A.J.38}\label{cdp:eq:A.J.38}
\]
where $g\ne0$ and $g(0,0)=0$. Then the reduced
plane germ $\{g=0\}_{\rm red}$ is smooth. If
$b$ is its multiplicity in $\operatorname{div}(g)$,
there are local holomorphic coordinates in which
$f=(x^a-y^b)^m$.
Consequently, at every point where the full fibre
is a uniform multiple of its reduced support and
the reduced equation has a nonzero Hessian, its
critical surface germ is a constant
transverse $A_{b-1}$ product. A nonordinary change
of transverse type at a point with nonzero Hessian
is excluded.
\end{proposition}


\begin{proof}
\noindent\textbf{Step 1.} We first prove the isolated assertion, including its
comparison with the actual global fibre. Replace $h$
by $H=u^{1/m}h$, taking a holomorphic root on a small
coordinate neighbourhood. Then $f=H^m$ exactly.
This replacement preserves an isolated critical point
as a germ. Indeed every local irreducible component
of the critical locus of $H$ has constant $H$-value,
necessarily zero when it contains the origin. On
$H=0$ one has $dH=u^{1/m}dh$. Thus its critical
locus has the same zero-dimensional germ as that
of $h$.

\noindent\textbf{Step 2.} Choose a small closed ball $K\subset\C^3$ such that
$H^{-1}(0)$ is smooth along its boundary intersection
and transverse to $\partial K$.
The standard isolated-singularity Milnor-ball
construction supplies this choice; see
\cite[Lemma 6.1 and Appendix B]{Milnor1968}.
The restriction $H|_{\partial K}$ is a submersion
at its zero set. Compactness therefore supplies
$\eta>0$ for which it is a submersion wherever
$|H|<2\eta$ on the boundary. For the linear
deformation
\[
 H_\lambda(w)=H(w)-\lambda_1w_1
                         -\lambda_2w_2-\lambda_3w_3,
                                                        \tag{A.J.38a}\label{cdp:eq:A.J.38a}
\]
this boundary submersion persists for all parameters
in a sufficiently small polydisc $P$ and all values
of modulus less than $\eta$. The gradient map of
$H$ has an isolated zero and consequently is a finite
holomorphic map germ of positive local degree. Choose
a small regular value $\lambda$ of that map. Its
nonempty finite inverse image consists of critical
points of $H_\lambda$ with invertible Hessian. All
these points and their critical values tend to the
origin as $\lambda\to0$.

\noindent\textbf{Step 3.} Remove the critical-value discriminant from
$P\times\Delta_\eta$. The relative critical locus
over this domain misses the boundary, so its
projection is proper and its image is a proper
complex analytic subset. Its complement is path
connected. The projection of the corresponding
smooth level sets in $K$ is a proper submersion
with boundary, and hence an Ehresmann bundle.
At any one of the nondegenerate critical points,
the holomorphic Morse lemma gives
$H_\lambda-H_\lambda(p)=X^2+Y^2+Z^2$.
A local Milnor fibre at a sufficiently small regular
value embeds in the corresponding level set in $K$.
The bundle transports it, by an oriented embedding,
into $K\cap H^{-1}(\epsilon_0)$ for a nonzero small
$\epsilon_0$, which is a submanifold of the actual
global torus $Y_{\epsilon_0^m}$.

\noindent\textbf{Step 4.} For a positive real local Milnor value, the real
locus $X,Y,Z\in\mathbb R$ is a two-sphere. It is
Lagrangian for the standard K\"ahler form; multiplication
by $\mathrm i$ identifies its tangent bundle with
its normal bundle with the opposite oriented Euler
class. Its self-intersection is therefore
$-\chi(S^2)=-2$. Equivalently this is the case
$a=2$, $b=c=1$ of the explicit cycle in \ref{cdp:st:A.5.15},
after the linear change $Y^2+Z^2=(Y+\mathrm iZ)(Y-\mathrm iZ)$.
The transported sphere has the same square in the
actual torus, whose $\pi_2$ is zero. This is the
required contradiction. No suspension representation
of the original isolated singularity has been used.

\noindent\textbf{Step 5.} We prove the comparison of Milnor fibres needed for
the argument. Factor the analytic plane function as
\[
                   g=u\prod_{i=1}^r h_i^{b_i},
\]
where the $h_i$ are distinct reduced irreducible
germs and $u$ is a unit. Deform each $h_i$
holomorphically to $h_{i,s}$ in a finite-dimensional
parameter polydisc $P$, and retain its exponent:
\[
 g_s=u\prod_i h_{i,s}^{b_i},\qquad
 H_s=x^a-g_s.                                    \tag{A.J.39}\label{cdp:eq:A.J.39}
\]
We use only sufficiently small parameters.

\noindent\textbf{Step 6.} Choose a small closed ball $K\subset\C^2_{y,z}$
such that each $h_i=0$ meets $\partial K$
transversely in its smooth part, and their boundary
links are disjoint. Choose its representative so
that every critical value of $g$ on $K$ is zero.
This is possible because the critical locus of a
holomorphic function has finitely many local
irreducible components, the function is constant
on each of them, and each component of the germ
contains the origin. The boundary smoothness,
transversality and disjointness persist for
$h_{i,s}$ by compactness and the implicit function
theorem. Near any of those links, the other factors
are units, so there are parameter-dependent local
coordinates $(v,w)$ with $g_s=v^{b_i}$, obtained
by taking a root of that unit. The boundary
$\partial K$ is transverse to the $w$ direction.
Consequently the restriction of $g_s$ to
$\partial K$ is a submersion at every point where
$0<|g_s|<\eta_0$, for a uniform sufficiently small
$\eta_0>0$. Away from these link neighbourhoods
$|g_s|$ has a uniform positive lower bound on the
boundary. This proves the asserted uniform
transversality, including at parameters for which
the germs inside $K$ change their singularities.

\noindent\textbf{Step 7.} Take a small disc $D_x$ of radius $\rho$ with
$\rho^a<\eta_0/2$, then shrink $P$ so that every
critical value of $g_s$ on $K$ has modulus less than
$\rho^a/4$. This shrinking follows by compactness:
a sequence of critical points with parameters
tending to zero would otherwise limit to a
nonzero critical value of $g$ in the chosen small
representative. The latter has no such value.
Fix $0<\eta<\rho^a/4$ and put
$W=\overline D_x\times K$, a compact manifold
with corners contained in the original coordinate
chart in $Y$.

\noindent\textbf{Step 8.} Every $H_s^{-1}(\epsilon)$ with
$0<|\epsilon|<\eta$ is transverse to this boundary.
On the face $\partial D_x\times K$, the equation
makes $|g_s|>3\rho^a/4$, so $dg_s\ne0$; its plane
directions give the submersion. On the face
$\overline D_x\times\partial K$, the $x$
derivative gives it when $x\ne0$. At $x=0$ we
have $g_s=-\epsilon\ne0$, and the preceding
boundary calculation gives it. At the corner
$|x|=\rho$, that same calculation for $g_s$
supplies both real target directions. The corner
stratum is therefore also transverse. We use
Ehresmann's theorem with corners: local lifts of
base tangent vectors can be chosen tangent to
each boundary face by these submersion statements,
then patched by a partition of unity and
integrated on the compact fibres. Thus no change
of the chosen domain $W$ is required.

\noindent\textbf{Step 9.} Inside $P\times\{0<|\epsilon|<\eta\}$ remove the
critical-value discriminant $\mathcal D$ of the
maps $H_s$ on $W$. This is a proper complex
analytic subset. The relative critical locus has
no point on the boundary at such a value, so its
projection is proper over every compact subset
of this parameter domain, and the proper mapping
theorem applies. For fixed $s$ its critical
values are discrete, since a holomorphic function
is constant on each connected critical component.
The complement of a proper complex analytic
subset of a connected complex manifold is path
connected. The transversality just proved and
the implicit function theorem give a proper
smooth submersion, including its boundary,
\[
 \{(s,p):p\in W,\ (s,H_s(p))\in
                    (P\times\Delta_\eta^*)\setminus\mathcal D\}
 \longrightarrow
       (P\times\Delta_\eta^*)\setminus\mathcal D,
 \quad (s,p)\longmapsto(s,H_s(p)).                 \tag{A.J.40}\label{cdp:eq:A.J.40}
\]
Ehresmann's theorem identifies all its compact
fibres by orientation-preserving diffeomorphisms
of manifolds with corners.
In particular a sufficiently small local Milnor
fibre of an interior singularity of some $H_s$,
at a value near zero, embeds into a smooth fibre
in \eqref{cdp:eq:A.J.40}, hence into
$W\cap H_0^{-1}(\epsilon_0)$ for a sufficiently
small nonzero $\epsilon_0$. This last space is
an actual submanifold of the global torus
$Y_{\epsilon_0^m}$ by \eqref{cdp:eq:A.J.38}; no identification
of it with a standard Milnor fibre is needed.
Choose the value to avoid every other
critical value, and choose its local Milnor ball
in the interior of $W$. This is the precise
comparison we use. Arbitrary deformations which
change the exponents along the boundary links
are not used in \eqref{cdp:eq:A.J.40}.

\noindent\textbf{Step 10.} Suppose first that $r\ge2$. Choose two distinct
factors $h_i,h_j$. Their map
$(h_i,h_j):(\C^2,0)\to(\C^2,0)$ is finite,
since their common zero germ is zero dimensional.
Its positive local degree is their intersection
multiplicity. A sufficiently small regular value
$(s_i,s_j)$ consequently has nonempty finite
inverse image. At each inverse-image point,
$h_i-s_i$ and $h_j-s_j$ are transverse coordinates.
Choose the remaining small shifts to avoid those
finitely many points. The deformation \eqref{cdp:eq:A.J.39}
then has an interior point at which
\[
            H_s=x^a-U v^{b_i}w^{b_j},
\]
with $U$ a unit independent of $x$. Replacing
$v$ by $vU^{1/b_i}$ makes this exactly
$x^a-v^{b_i}w^{b_j}$ by an invertible holomorphic
coordinate change. \ref{cdp:st:A.5.15} constructs a spherical
class of nonzero square in its local Milnor fibre.
The embedding \eqref{cdp:eq:A.J.40} puts that class in the
actual smooth torus fibre of $f=H_0^m$.
Preservation of intersection
numbers and $\pi_2$-vanishing give the same
contradiction as in \ref{cdp:st:A.5.15}. Thus $r=1$.

\noindent\textbf{Step 11.} Suppose its reduced factor $h=h_1$ is singular.
Write $g=u h^b$ and absorb a local $b$th root of
$u$ into $h$. This leaves its reduced zero germ
unchanged. Since $h$ is reduced, its gradient
has an isolated zero at the origin: a critical
curve would have constant value zero and would
be a repeated component of $h$. The finite
gradient map has a positive local degree.
Choose a small regular value
$(\lambda_1,\lambda_2)$ of this map. Then
\[
        h_s=h-\lambda_1y-\lambda_2z-\lambda_0
\]
has nondegenerate critical points near zero;
regularity of the gradient value says exactly
that their Hessians are invertible. Choose
$\lambda_0$ so that their critical values are
nonzero. At one such point $p$ put
$v=h_s(p)\ne0$. The Hessian of $h_s^b$ is
$b v^{b-1}$ times the Hessian of $h_s$, and
is invertible. The holomorphic Morse lemma on
the two plane variables identifies the germ of
$H_s-H_s(p)$ with
\[
                         x^a-Y^2-Z^2.
\]
A linear change of $(Y,Z)$ gives the Hirzebruch
germ $x^a-YZ$, already excluded by \ref{cdp:st:A.5.15}.
Its Milnor fibre embeds by \eqref{cdp:eq:A.J.40}, producing
the same contradiction. Hence $h$ is smooth.
Taking $y=h u^{1/b}$ as a coordinate proves the
asserted transverse product.

\noindent\textbf{Step 12.} For the Hessian consequence, take a local $m$th
root of the actual map, absorbing its unit.
Its reduced equation has the same nonzero
Hessian rank as every other reduced equation
at the critical point. The holomorphic splitting
lemma splits off a nondegenerate square and
gives $f=(x^2-g(y,z))^m$ exactly. If $g=0$,
the equation would not be reduced, so $g\ne0$.
At a critical point $dg(0)=0$. The result just
proved gives $g=y^b$ with $b\ge2$ in suitable
holomorphic coordinates, which is the constant
transverse $A_{b-1}$ product.
\end{proof}






\begin{proposition}[Every horizontal tacnode image is elliptic]\label{cdp:st:A.5.18}
\StatementRecord{cdp:st:A.5.18}{proved}{Historical events: Claude Fable 5.2 (Anthropic), effort high, 16 November 2026: carried as background without fault found. GPT-6 Astra (OpenAI), ChatGPT Pro mode, effort Pro, S6\_09\_PAPERS, 17 November 2026: targeted proof check. Also checked in this writing. Pre-continuation writing events: GPT-6 Astra (OpenAI), ChatGPT Pro mode, effort Pro, S6\_10\_PAPERS, 18 November 2026: targeted proof check. Continuation metadata events: GPT-6 Astra (OpenAI), ChatGPT Pro mode, effort Pro, S6\_10\_PAPERS continuation, 19 November 2026: used and rechecked in this run; status-vocabulary reconciliation; metadata reconciliation only; no new mathematical check. \emph{Pre-continuation recorded proof scope.} The actual quotient map to the conductor image, branch and discriminant budgets, simple-branch sphere obstruction, and minimality were re-derived. Both delta-at-most-two non-lci rings were computed with their symbolic squares, exact conormal frames, conductor poles and cleared quartic degrees; the rational-cover addendum was included}{the hypotheses, notation and definitions stated here; \cref{cdp:st:A.5.14}; \cref{cdp:st:A.5.5}; \cref{cdp:st:A.5.8}; \cref{cdp:eq:A.T.3}; \cref{cdp:st:A.5.17}; \cref{cdp:eq:A.T.4}; \cref{cdp:eq:A.T.5}; \cref{cdp:eq:A.T.7}; \cref{cdp:eq:A.T.8}; \cref{cdp:st:A.5.15}; \cref{cdp:eq:A.T.9}; \cref{cdp:eq:A.T.12}; \cref{cdp:eq:A.T.13}; \cref{cdp:eq:A.T.15}; \cref{cdp:eq:A.T.18}; \cref{cdp:eq:A.T.20}; \cref{cdp:eq:A.T.19}; \cref{cdp:st:C2d}; \cref{cdp:eq:A.T.22}}

Let $f:Y\to C_0$ be a proper holomorphic map from a smooth
compact complex threefold with $b_2(Y)=0$, with connected
fibres and smooth general fibre a complex two-torus.
Let $mS$ be a full fibre with irreducible reduced support,
$\eta:T\to S$ its normalization, and $\pi:Z\to T$ a
resolution with an irrational ruling $p:Z\to B$.
Suppose every exceptional curve of $\pi$ is vertical and
the conductor divisor on $Z$ is $2C+V$, where $C$ is a
section of $p$ and $V$ is vertical. Let $D=\eta\pi(C)$.
Suppose the generic plane-curve transversal to $D$ in $S$ is an
ordinary tacnode.

Then $D$ is a smooth elliptic curve, $B=C$ is elliptic,
and $C\to D$ is an unramified double cover. The transversal
type along $D$ is everywhere the same ordinary tacnode.
The normalization $T$ is a smooth minimal ruled surface
over this elliptic curve. Its conductor is twice a
section, and $e(S)=0$. In particular its minimal
resolution is the identity.
In particular a rational image of
the horizontal tacnodal conductor cannot occur. The proof
allows vertical conductor and rational singularities of
$T$; it does not assume that the normalization map is
immersive at the special points under examination.
\end{proposition}


\begin{proof}
\noindent\textbf{Step 1.} If the normalization of $D$ has positive genus, \ref{cdp:st:A.5.14}
and \ref{cdp:st:A.5.5} make $D$ smooth elliptic and its transversal
type constant. There are two generic normalization
branches of a tacnode. A vertical component of $Z$
is rational, so it cannot dominate this elliptic curve.
Thus both branches lie on $C$, and $C\to D$ is an
unramified double cover: the constant local equation
$x^2-y^4=0$ has two separately smooth normalization
branches on every sufficiently small disc in $D$.
Riemann--Hurwitz gives $g(C)=g(B)=1$.
It remains to exclude the case $D^\nu=\Sigma\simeq\PP^1$.

\noindent\textbf{Step 2.} The finite map $C\to\Sigma$ has degree at least two,
because $C\simeq B$ has positive genus. Since a generic
tacnode has exactly two normalization branches, this
degree is exactly two. Consequently no vertical
conductor component can also dominate $D$: its generic
points would supply an additional normalization branch.
Write $g=g(C)\ge1$ and let $\mathcal B\subset\Sigma$
be the reduced branch divisor of this double cover.
Riemann--Hurwitz gives
\[
                   \deg\mathcal B=2g+2.             \tag{A.T.1}\label{cdp:eq:A.T.1}
\]

\noindent\textbf{Step 3.} We first bound the number of branches of the singular
image curve itself. \ref{cdp:st:A.5.8} supplies the descended ruling
$\widetilde p:T\to B$. Identify $C$ with $B$ using the
section, let $\phi:B\to D$ be its actual map to the
image curve, and put
\[
                    \psi=\phi\widetilde p:T\to D.
\]
Every vertical conductor component has constant $\psi$
value. Its image under $\eta$ differs from $D$, as just
proved, and meets $D$ in finitely many points. Thus
there is a finite set $A\subset D$ such that, for each
$s\in S$, either $\psi(\eta^{-1}(s))$ is a singleton
or all its values belong to $A$. Include in $A$ the
constant values on the vertical conductor components
and the values of $\phi$ over their finitely many
intersections with $D$. Fibres over the remaining
points of $D$ have their preimages on $C$, where all
the $\psi$ values equal that same point of $D$.

\noindent\textbf{Step 4.} Choose an embedded finite tree $J\subset D$ containing
$A$, in a triangulation of the compact analytic curve.
It can be constructed by joining these finitely many
points by arcs, subdividing their intersections and
retaining a spanning tree. The quotient
$q:D\to D/J$ is a homotopy equivalence: a contractible
CW subcomplex is a cofibration, and the homotopy
extension property extends its contraction. Since
$q\psi$ is constant on every fibre of the finite
proper quotient $\eta$, there is a continuous map
\[
 a:S\to D/J,\qquad a\eta=q\psi,\qquad a|_D=q.     \tag{A.T.2}\label{cdp:eq:A.T.2}
\]
The last equality holds first on the dense regular
locus of $D$ and then everywhere by continuity.

\noindent\textbf{Step 5.} Use the explicit triangulated neighbourhood retraction
$r:N\to S$ of (A.Pi.3). For a sufficiently small disc
$\Delta$ around the fibre value, properness gives
$U=f^{-1}(\Delta)\subset N$. The map $(ar)|_U$ is
surjective on fundamental groups because its restriction
to $D$ is the homotopy equivalence $q$. Perturbing loops
off the real-codimension-two analytic support gives
$\pi_1(U\setminus S)\twoheadrightarrow\pi_1(U)$.
The torus-bundle exact sequence over $\Delta^*$ therefore
gives the actual quotient maps
\[
 \Z^4\rtimes\Z\twoheadrightarrow\pi_1(U)
       \twoheadrightarrow\pi_1(D/J)\simeq\pi_1(D).
                                                        \tag{A.T.3}\label{cdp:eq:A.T.3}
\]
The first group is metabelian, so the last is metabelian.

\noindent\textbf{Step 6.} For each singular point $x$ of $D$, let $r_x$ be the
number of its normalization branches and put
$\beta=\sum_x(r_x-1)$. The finite normalization map
identifies exactly those finite sets of points on
$\PP^1$; since it is a compact Hausdorff quotient,
this is also the actual underlying topological quotient.
Adjoining arcs between the points to be identified
gives a CW presentation homotopy equivalent to
$S^2\vee\bigvee^\beta S^1$. Hence
$\pi_1(D)$ is the free group $F_\beta$.
For $\beta\ge2$ this group surjects onto $S_4$ by
sending two free generators to $(1234)$ and $(12)$;
the commutator $A_4$ of that quotient is nonabelian.
Equation \eqref{cdp:eq:A.T.3} therefore forces $\beta\le1$.
By \ref{cdp:st:A.5.14} the total delta invariant $\delta=p_a(D)$
is at most two. If $s$ denotes the number of points
of $\Sigma$ over $\operatorname{Sing}D$, then
\[
 s=\#\operatorname{Sing}D+\beta\le\delta+\beta\le3.
                                                        \tag{A.T.4}\label{cdp:eq:A.T.4}
\]

\noindent\textbf{Step 7.} First suppose that $D$ is a local complete intersection.
We retain the exact bundles and jets of \ref{cdp:st:A.5.5}.
Let $\nu:\Sigma\to D$ be normalization, $I=\mathcal I_D$,
\[
 N=(\nu^*(I/I^2))^*,\qquad
 M=\nu^*(\mathcal O_Y(S)|_D).
\]
Thus $M$ is a degree-zero line bundle on $\Sigma$.
Let $\Xi$ be the curve-normalization
conductor divisor, $c=\deg\Xi=2\delta$, and let
$R\subset N$ and $Q=N/R$ be the saturated radical and
quotient of the generic rank-one Hessian. Write
$q_0=\deg Q$. The adjunction and conductor comparisons
(A.J.1) give
\[
 \deg N=-2+c,\quad \deg R=-2+c-q_0,\quad
 0\ne a\in H^0(\Sigma,M\otimes Q^{-2}),\quad q_0\le0.
                                                        \tag{A.T.5}\label{cdp:eq:A.T.5}
\]
Here $a$ denotes the intrinsic quadratic coefficient,
whose zero divisor has degree $-2q_0$.
The actual quartic invariant $\Delta=ae-b^2/4$ of
\ref{cdp:st:A.5.5} is meromorphic in
$L^2$, where $L=M\otimes Q^{-1}\otimes R^{-2}$.
Its pole bound (A.J.9) gives a nonzero holomorphic
section $\sigma$ of the actual square line bundle
$L^2(4\Xi)=(L(2\Xi))^2$. Its zero divisor has degree
\[
               d_\Delta=8+2q_0.                    \tag{A.T.6}\label{cdp:eq:A.T.6}
\]

\noindent\textbf{Step 8.} The double cover $C\to\Sigma$ is obtained by adjoining
$\sqrt{-\Delta}$ to the function field, with the line
trivializations specified by $L$. Indeed over the
generic point a tacnode has normal expansion
$a\ell^2+b\ell r^2+er^4$, and the two normalization
branches are distinguished by the two roots of
$a(\ell/r^2)^2+b(\ell/r^2)+e$. Their discriminant
is $-4\Delta$. Changing a local frame of $L$ multiplies
the representing function of $\Delta$ by a square,
so this describes the same degree-two extension on
every overlap. At a point with local parameter $z$,
write that representing function as $z^n u(z)$,
$u(0)\ne0$. A root of the unit and removal of the
even power of $z$ show that its normalization is
ramified exactly when $n$ is odd. Adding $4\Xi$ changes
each order by an even number. Thus $\mathcal B$ is
exactly the odd-multiplicity support of
$\operatorname{div}(\sigma)$, and
\[
                  2g+2\le 8+2q_0.                 \tag{A.T.7}\label{cdp:eq:A.T.7}
\]

\noindent\textbf{Step 9.} At a point of $D_{\rm reg}$ where $a\ne0$, \ref{cdp:st:A.5.17}
applies to the actual local $m$th root of the full-fibre
map. Its nonzero Hessian makes the local equation a
constant transverse product, necessarily of type
$A_3$ because this is the generic type along $D$.
The conductor cover is therefore unramified there.
Every branch point outside $\nu^{-1}(\operatorname{Sing}D)$
is consequently a zero of $a$. Equations \eqref{cdp:eq:A.T.4}--\eqref{cdp:eq:A.T.5}
give the second bound
\[
                    2g+2\le -2q_0+s.              \tag{A.T.8}\label{cdp:eq:A.T.8}
\]
Adding \eqref{cdp:eq:A.T.7} and \eqref{cdp:eq:A.T.8} yields
$4g+4\le8+s\le11$. Since $g\ge1$, we have
\[
           g=1,\qquad \deg\mathcal B=4,
                    q_0\in\{-2,-1,0\}.            \tag{A.T.9}\label{cdp:eq:A.T.9}
\]

\noindent\textbf{Step 10.} We next prove the exact local obstruction required in
the equality cases. Let a smooth point of $D$ be a
branch point at which $\Delta$ has a simple zero.
The preceding nonzero-Hessian argument gives $a(0)=0$.
In coordinates $D=(x,y)=0$ with curve parameter $z$,
choose the holomorphic saturated Hessian factor as the
first normal coordinate. The quadratic jet is then
$a(z)x^2$. The radical cubic vanishes identically,
so a local reduced equation has the exact expansion
\[
 f_0=a(z)x^2+x^2(xA+yB)+xy^2b(z)+xy^3C+y^4E,
                                                        \tag{A.T.10}\label{cdp:eq:A.T.10}
\]
with holomorphic coefficients. The identity
$\Delta=ae-b^2/4$, $a(0)=0$ and the simple zero of
$\Delta$ imply $b(0)=0$, $a'(0)\ne0$ and $e(0)\ne0$.
Indeed an order at least two for $a$, together with
$b(0)=0$, would give an order at least two for $\Delta$;
and then $\Delta'(0)=a'(0)e(0)$ is nonzero.
Thus $b/a$ is holomorphic and $E(0)\ne0$.

\noindent\textbf{Step 11.} Set $X=x+b(z)y^2/(2a(z))$. In this actual analytic
coordinate change the coefficient of $Xy^2$ cancels.
All terms created from $x^2(xA+yB)$ lie in
$(X^2)+(Xy^3)+(y^4)$. Therefore, collecting the
terms of $X$-degree zero and one, gives exactly
\[
 f_0=X^2A_1(X,y,z)+Xy^3B_1(y,z)+y^4C_1(y,z),
 \qquad C_1(0)\ne0.
\]
Choose a fourth root of $C_1$ and define
\[
 Y=yC_1^{1/4}+
                  \frac{XB_1}{4C_1^{3/4}}.
                                                        \tag{A.T.11}\label{cdp:eq:A.T.11}
\]
The map $(X,y,z)\mapsto(X,Y,z)$ is locally
biholomorphic, and direct expansion gives
$f_0-Y^4\in(X^2)$. Hence $f_0=Y^4+X^2H(X,Y,z)$.
Its quadratic jet along $D$ is still $a(z)X^2$,
so $H(0,0,z)=a(z)$ and $\partial H/\partial z(0)\ne0$.
The further coordinate $Z=H$ gives exactly
\[
                         f_0=Y^4+X^2Z.             \tag{A.T.12}\label{cdp:eq:A.T.12}
\]
To apply this to the actual map, first take a local
$m$th root of its unit. This multiplies the reduced
equation by a unit, which changes neither zero order
of $a$ nor of $\Delta$. The same coordinate construction
then makes the full map $(Y^4+X^2Z)^m$.
Statement \ref{cdp:st:A.5.15} excludes it. Consequently every
smooth branch point has order at least three in
$\operatorname{div}(\sigma)$.

\noindent\textbf{Step 12.} We also record precisely when the conductor denominator
can be removed at a singular point $x\in D$.
If the intrinsic section $a$ is nonzero at every
normalization point above $x$, choose an actual
constant linear combination of the two lci conormal
generators which is nonzero under the normalized
quotient at all those finitely many points. The
corresponding diagonal entry $a_0$ of the symmetric
Hessian matrix belongs to $\mathcal O_{D,x}$ and
has nonzero value on every normalization branch.
It is therefore a unit in $\mathcal O_{D,x}$.
The other normalized factor coefficient is the
actual quotient $h=q_{XY}/a_0\in\mathcal O_{D,x}$.
Thus square completion and monic polynomial division
in the proof of (A.J.7)--(A.J.8) both take place
over $\mathcal O_{D,x}$ itself. The regular-sequence
graded pieces inject into their normalization
pullbacks, so their cancellations already descend
on $D$. We may repeat that proof with $d=1$.
It gives a holomorphic section $\Delta$ of $L^2$
above $x$, without a conductor pole. At a branch
point above $x$, this proves the sharper inequality
\[
 \operatorname{ord}\sigma\ge
                        4\operatorname{ord}\Xi+1.
                                                        \tag{A.T.13}\label{cdp:eq:A.T.13}
\]

\noindent\textbf{Step 13.} We now exhaust the three integer cases in \eqref{cdp:eq:A.T.9}.
If $q_0=-2$, the divisor of $\sigma$ has degree four
and has four odd-multiplicity points, so each is
simple. At most three lie above singularities of
$D$, by \eqref{cdp:eq:A.T.4}. The remaining smooth branch point
contradicts \eqref{cdp:eq:A.T.12}.
If $q_0=0$, the section $a$ has no zeros and all
four branch points would have to lie above the
singular locus; this contradicts $s\le3$.

\noindent\textbf{Step 14.} Finally let $q_0=-1$, so that $\deg\operatorname{div}(a)=2$
and $d_\Delta=6$. If $k$ of the four branch points
lie above $\operatorname{Sing}D$, then the local
obstruction and parity give
\[
               6\ge k+3(4-k)=12-2k.
\]
Hence $k=s=3$, and equality holds: the three singular
preimages have order one in $\operatorname{div}(\sigma)$
and the one smooth branch point has order three.
The relations $s\le\delta+\beta\le3$ now force
$\delta=2$, $\beta=1$ and two singular points of
delta invariant one. The exact delta-one
classification in \ref{cdp:st:A.5.14} identifies them as one
ordinary node and one cusp. Their normalization
conductor orders are respectively $(1,1)$ and $2$;
these follow directly from the displayed node and
cusp rings in that proof.

\noindent\textbf{Step 15.} At the cusp preimage, if $a$ were nonzero,
\eqref{cdp:eq:A.T.13} would make the order of $\sigma$ at least
$4\cdot2+1=9$, exceeding its total degree six.
So $a$ vanishes there. It also vanishes at the
smooth branch point. These two distinct points
exhaust the degree-two divisor of $a$. The section
is consequently nonzero at both node preimages.
Apply \eqref{cdp:eq:A.T.13} at the node: each of its two branch
points has order at least $4\cdot1+1=5$ in
$\operatorname{div}(\sigma)$, again exceeding the
total degree six. This excludes every rational lci image.

\noindent\textbf{Step 16.} We finish by treating all non-lci image germs permitted
by $\delta\le2$ and $\beta\le1$, with explicit analytic
rings and transverse coordinate maps. For a unibranch
reduced curve ring $A\subset\C\{t\}$ of colength at
most two, an element of order one makes $A=\C\{t\}$.
Otherwise its least positive valuation is two or three:
a least valuation at least four leaves three linearly
independent missing coefficients. If it is three and
the colength is two, the inclusion
$A\subset\C+t^3\C\{t\}$ is equality, since the latter
already has colength two. If it is two, choose a
parameter with $t^2\in A$ by taking a square root of
the leading unit. The even part of every convergent
series belongs to $\C\{t^2\}\subset A$. Subtracting
it shows that the least odd element can be written
$t^{2k+1}v(t^2)$; division by its unit puts
$t^{2k+1}$ in $A$. Consequently
$A=\C\{t^2\}\oplus t^{2k+1}\C\{t^2\}$, of colength
$k\le2$. These are the planar cusp and the planar
$A_4$ curve, besides the smooth case.

\noindent\textbf{Step 17.} For two branches, write their image rings as $A_1,A_2$.
The two projections of $A\subset A_1\oplus A_2$ are
surjective. Their kernels define a common finite
quotient $Q_0$ and an exact fibre-product description
$A=A_1\times_{Q_0}A_2$. Thus
$\delta=\delta_1+\delta_2+\operatorname{length}Q_0$.
If both branches are smooth, the quotient is
$\C\{t\}/(t^k)$ with $k\le2$; after changing the
second parameter, the two maps identify the same
$k$-jet. The map
\[
 x\longmapsto(t,s),\qquad
 y\longmapsto(t^k,-s^k)
\]
identifies this fibre product with
$\C\{x,y\}/(y^2-x^{2k})$: write a matching pair as
the sum of its average and $y$ times its difference
divided by $2x^k$. These are the node and tacnode.
If one branch is singular, the delta budget forces
one cusp, one smooth branch, and $Q_0=\C$.
Three branches are excluded by $\beta\le1$.

\noindent\textbf{Step 18.} An abstract plane curve ring in the given smooth
threefold is lci by the actual transverse-coordinate
construction in \ref{cdp:st:A.5.14}. Thus the only remaining
non-lci analytic rings, necessarily consuming the
whole delta budget at a single point, are
\[
\begin{array}{ll}
 A=\C\{t^3,t^4,t^5\},
 &I=(y^2-xz,\ yz-x^3,\ z^2-x^2y),\\[2pt]
 A=\C\{t^2,t^3\}\times_{\C}\C\{s\},
 &I=(xz,\ yz,\ y^2-x^3).
\end{array}                                             \tag{A.T.14}\label{cdp:eq:A.T.14}
\]
These are normal forms of the actual embeddings.
Their three displayed coordinate functions form a
basis of the cotangent space of each ring. Lift them
to $\mathcal O_{Y,x}$; their differentials are linearly
independent because the embedding dimension is three.
They are therefore ambient coordinates by the inverse
function theorem. The displayed ideals are their
kernels, obtained respectively by the parametrizations
$(t^3,t^4,t^5)$ and the union of $(t^2,t^3,0)$ with
$(0,0,s)$.

We show that the local reduced equation $f_0$ actually
lies in $I^2$ in each case. In the first case put
\[
 g_1=y^2-xz,\quad g_2=yz-x^3,\quad
 g_3=z^2-x^2y,\quad
 F=x^5+xy^3+z^3-3x^2yz.
\]
The exact identities
\[
 xF=g_2^2-g_1g_3,\qquad
 yF=xg_1^2+g_2g_3,\qquad
 zF=g_3^2+xg_1g_2                              \tag{A.T.15}\label{cdp:eq:A.T.15}
\]
show that the maximal ideal annihilates the class of
$F$ modulo $I^2$. Moreover
$I^{(2)}=I^2+(F)$ in the convergent local ring.
Here is a finite-module verification of that equality.
Let $J=I^2+(F)$. The quotient $\C\{x,y,z\}/J$ is
finite over the discrete valuation ring $\C\{x\}$,
since its reduction is the finite parametrized curve.
Its generic rank is nine: the curve has degree three
over the $x$ coordinate, and the square of its smooth
codimension-two ideal has length three along each
generic normal slice. Modulo $x$, the quotient has
the nine-element basis
\[
 1,y,y^2,y^3,\ z,yz,y^2z,\ z^2,yz^2,
\]
because its ideal is $(y^4,y^3z,y^2z^2,z^3)$.
A finite module over a discrete valuation ring is
the sum of its free and torsion parts; equality of
generic rank and residue-space dimension here makes
the torsion part zero. Thus $J$ is $x$-saturated.
By \eqref{cdp:eq:A.T.15}, $J$ and $I^2$ agree after inverting $x$.
The curve is lci away from the origin, so its symbolic
square differs from its ordinary square only by a
module supported there. Hence
$I^{(2)}=I^2:x^\infty=J$.

On its normalization let
$a_1=dg_1$ and $b_1=dg_2+t\,dg_1$.
They form a basis of the normalized conormal image;
$dg_3=-2t^2a_1+t b_1$.
The quadratic term of $F$ in this basis is
\[
 q_F=\frac3t a_1^2-\frac3{t^2}a_1b_1
                           +\frac1{t^3}b_1^2,
 \qquad \det(q_F)=\frac3{4t^4}.                  \tag{A.T.16}\label{cdp:eq:A.T.16}
\]
Every equation in $I^2$ has a holomorphic quadratic
form in this same basis. If $f_0=cF+I^2$ had
$c(0)\ne0$, the determinant of its quadratic form
would have the noncancellable leading pole
$3c(0)^2/(4t^4)$: the cross terms with a holomorphic
form have pole at most three. It could not have
generic rank one. Thus $c(0)=0$, and \eqref{cdp:eq:A.T.15} gives
$f_0\in I^2$.

For the second ring put $G=y^2-x^3$ and
$P=(z,G)$, $Q_1=(x,y)$. Then $I=P\cap Q_1$ and
\[
 I^{(2)}=P^2\cap Q_1^2=I^2+(zG).                \tag{A.T.17}\label{cdp:eq:A.T.17}
\]
Indeed $P^2=(z^2,zG,G^2)$, while $zG,G^2$ already
belong to $Q_1^2$. Multiplication by $z^2$ is
injective modulo $Q_1^2$, so membership of
$z^2A+zGB+G^2C$ in $Q_1^2$ forces $A\in Q_1^2$.
This gives the asserted generating ideal directly.
Furthermore
$x(zG)=(xz)G$, $y(zG)=(yz)G$, and
$z(zG)=(yz)^2-x(xz)^2$ all lie in $I^2$.
On the cusp branch use $a_1=d(xz)=t^2dz$ and
$b_1=dG=-3t^4dx+2t^3dy$.
The quadratic form of $zG$ is $t^{-2}a_1b_1$,
of determinant $-1/(4t^4)$. The same rank-one
argument forces its coefficient to vanish at the
origin, and the three identities again give
$f_0\in I^2$.

For clarity, the bundle used at these non-lci points
is now defined by its actual morphism. Put
\[
 K=\operatorname{im}\bigl(\nu^*(I/I^2)
                         \longrightarrow\nu^*\Omega_Y^1\bigr),
 \qquad N=K^*.
\]
The coherent image $K$ is locally free of rank two
on the smooth curve $\Sigma$. Let $R_{\rm ram}$ be
the zero divisor of $d\nu$, and let $K_{\rm sat}$
be the kernel of the surjection
$\nu^*\Omega_Y^1\to\Omega_\Sigma^1(-R_{\rm ram})$.
The torsion quotient $K_{\rm sat}/K$ defines a divisor
$J_{\rm jac}$. Taking determinants in these actual
maps gives
\[
 \deg N=-2+\deg(J_{\rm jac}-R_{\rm ram}).         \tag{A.T.18}\label{cdp:eq:A.T.18}
\]
For $(t^3,t^4,t^5)$, the wedge $a_1\wedge b_1$
has coefficients $t^8(3,4t,5t^2)$, while $d\nu$
has common factor $t^2$. Thus the correction is
$8-2=6$. For the cusp branch in the second case,
the wedge is $-t^5(2,3t,0)$ and the tangent common
factor is $t$, giving correction four. On its
smooth branch the wedge is $s^2(dx\wedge dy)$,
giving correction two. In both cases \eqref{cdp:eq:A.T.18}
therefore gives $\deg N=4$.

Because $f_0\in I^2$, its quadratic form is a
holomorphic section of $\operatorname{Sym}^2K\otimes M$,
independent of its expression as a sum of products
of ideal generators. This follows either from its
actual Hessian or from injectivity of
$\operatorname{Sym}^2K\to\nu^*\operatorname{Sym}^2\Omega_Y^1$.
Its saturated radical and quotient again give line
bundles $R,Q$, with $q_0=\deg Q\le0$ and
$\deg R=4-q_0$. The generic quartic invariant
is meromorphic in $L^2$, where
$L=M\otimes Q^{-1}\otimes R^{-2}$; now
\[
                        \deg L^2=-16+2q_0.        \tag{A.T.19}\label{cdp:eq:A.T.19}
\]

We bound its poles using actual transverse affine
coordinate maps, with the curve parameter held
fixed during every normal differentiation. For the
space cusp use, for $t\ne0$,
\[
 (u,v,t)\longmapsto
 \left(t^3,\ t^4+\frac{v}{3t^5},\
                   t^5-\frac{u}{t^3}+\frac{2v}{3t^4}\right).
                                                        \tag{A.T.20}\label{cdp:eq:A.T.20}
\]
This is locally biholomorphic and its two transverse
derivatives are the dual frame to $a_1,b_1$.
Each coefficient of a derivative has pole at most
five. The first normal derivatives of all three
$g_i$ are holomorphic, since
$dg_1=a_1$, $dg_2=b_1-ta_1$ and
$dg_3=-2t^2a_1+t b_1$.
Second normal derivatives have pole at most ten.
For $f_0=\sum c_{ij}g_i g_j$, the cubic Taylor
coefficients consequently have pole at most ten:
the only surviving terms use a first and a second
derivative of the $g_i$, or a first derivative of
$c_{ij}$ and two first derivatives of the $g_i$.
Every fourth derivative of a holomorphic ambient
function in \eqref{cdp:eq:A.T.20} has pole at most twenty.

Choose a holomorphic change of the $N$ frame adapted
to its saturated radical and quotient. It does not
increase these bounds. The radical cubic is zero
because the generic transversal is $A_3$.
Writing the relevant Taylor terms as
$a\ell^2+b\ell r^2+er^4$ therefore gives a
holomorphic $a$, a mixed cubic $b$ of pole at most
ten, and $e$ of pole at most twenty. The invariant
$\Delta=ae-b^2/4$ has pole at most twenty.
These are the same intrinsic jets on overlaps:
to weighted order four, where $\ell$ has weight
two and $r$ has weight one, an adapted change is
$\ell'=\alpha\ell+u r^2$, $r'=\beta r$.
Direct substitution multiplies $\Delta$ by
$\alpha^{-2}\beta^{-4}$; changes of the defining
section supply the factor $M^2$. Higher weighted
terms and changes of the curve parameter do not
alter these coefficients. This verifies the
transition law for the line $L^2$ in \eqref{cdp:eq:A.T.19}.

There is only one singular preimage $p$ in this case.
Thus $\Delta$ is a holomorphic section of the square
line $L^2(20p)$, whose degree is $4+2q_0$.
As in \eqref{cdp:eq:A.T.7}, the even pole correction preserves
branch parity, so
\[
 \deg\mathcal B\le4+2q_0,
 \qquad \deg\mathcal B\le-2q_0+1.
\]
The second bound again uses \ref{cdp:st:A.5.17} away from that
single singular image preimage. Adding them gives
$2\deg\mathcal B\le5$, contradicting
$\deg\mathcal B=2g+2\ge4$.

For the cusp-plus-line germ, the corresponding
actual affine coordinate maps are
\[
 (u,v,t)\longmapsto
       \left(t^2,t^3+\frac{v}{2t^3},\frac{u}{t^2}\right),
 \qquad
 (u,v,s)\longmapsto\left(\frac us,\frac vs,s\right).
                                                        \tag{A.T.21}\label{cdp:eq:A.T.21}
\]
Their first generator derivatives are holomorphic
in the respective normalized conormal frames.
The same product differentiation gives cubic pole
bounds six and two, and the fourth derivative bounds
are twelve and four. Thus $\Delta$ has poles bounded
by $12p+4q$ at the two normalization preimages.
The line $L^2(12p+4q)$ is a square and has degree
$2q_0\le0$ by \eqref{cdp:eq:A.T.19}. A nonzero section has no
positive-degree odd divisor, whereas this double
cover has at least four branch points. This is the
last contradiction. All non-lci cases have now
been excluded, and the positive-genus case at the
start gives the asserted elliptic curves and constant
transversal type.

We finally prove the normalization and Euler assertions,
keeping track of any unnecessary blowups in the original
choice of resolution. Factor $\pi$ through the minimal
resolution $\pi_{\min}:Z_{\min}\to T$, with the resulting
sequence of exceptional-curve blowdowns
$q:Z\to Z_{\min}$. Every curve contracted by $q$ is
vertical for $p$, by the stated assumption on $\pi$.
The equality of structure-sheaf direct images for each
blowdown, in local coordinates on $B$, gives a descended
ruling $p_{\min}:Z_{\min}\to B$ with
$p=p_{\min}q$. The image $C_{\min}$ is still a section,
and the pushed-forward effective conductor is
$D_{\min}=2C_{\min}+V_{\min}$ with $V_{\min}$ vertical.
The normalization-canonical comparison gives
$K_{Z_{\min}}+D_{\min}=\pi_{\min}^*\eta^*\omega_S$
as line bundles, and its rational Chern class is zero.

Statement \ref{cdp:st:C2d} gives a smooth normalization near the
inverse image of the elliptic curve $D$, with normal
line of degree zero on its unramified double cover.
The minimal resolution is an isomorphism there, so
$C_{\min}^2=0$. No other conductor or exceptional
curve meets it. Hence $V_{\min}\cdot C_{\min}=0$.
No ruling $(-1)$-curve can meet $C_{\min}$: one in
$V_{\min}$ is disjoint from it, while one outside
$V_{\min}$ would have
\[
 (K_{Z_{\min}}+2C_{\min}+V_{\min})\cdot E
       =-1+2C_{\min}\cdot E+V_{\min}\cdot E\ge1,
\]
contrary to the vanishing Chern class.

Contract ruling $(-1)$-curves to a minimal ruled bundle.
Every contraction therefore avoids the section. For
one blowdown $\sigma:W\to W'$ let $D'=\sigma_*D_W$,
let $a$ be the coefficient of its exceptional curve,
and let $d$ be the multiplicity of $D'$ at the centre.
The actual divisor comparison is
\[
 K_W+D_W=\sigma^*(K_{W'}+D')+(1+a-d)E.
                                                        \tag{A.T.22}\label{cdp:eq:A.T.22}
\]
Its intersection with $E$ gives $a=d-1$, so the
log-canonical line and its zero rational Chern class
descend. The section remains disjoint from the
effective vertical part at every step. On a minimal
ruled bundle every nonzero effective vertical divisor
is a positive sum of fibres and meets the section.
The final vertical part is consequently zero. If
there had been a blowup in the reverse sequence,
its first centre would be away from this section,
and \eqref{cdp:eq:A.T.22} would give coefficient $-1$ on its
exceptional curve. Its strict transform retains that
negative coefficient under subsequent blowups, a
contradiction to effectiveness. Thus no ruling
blowdowns were needed: $Z_{\min}$ is already a
minimal ruled bundle and $D_{\min}=2C_{\min}$.

All exceptional curves of $\pi_{\min}$ would be
vertical. A minimal ruled bundle has irreducible
fibres of square zero, which cannot form the
negative-definite exceptional divisor of a surface
resolution. Hence $\pi_{\min}$ has no exceptional
curve and is an isomorphism. This proves the assertion
about $T$. Finally $\eta$ is an isomorphism away from
its conductor section $C_T$ and the elliptic image
$D$. Additivity of the topological Euler number on
these actual decompositions gives
\[
 e(S)=e(T)-e(C_T)+e(D)
      =e(\PP^1)e(B)-0+0=0.                       \tag{A.T.23}\label{cdp:eq:A.T.23}
\]
The full multiple fibre has the same underlying
topological space and therefore the same Euler number.
\end{proof}



\paragraph{Connected double covers in the two non-lci cases.}
The non-lci calculation also excludes a connected
degree-two conductor cover when its covering curve
is rational, for either of the two image rings in
\eqref{cdp:eq:A.T.14}. In the cusp-plus-line case the cleared
quartic has degree $2q_0\le0$, so it has no branch
divisor at all. In the space-cusp case a rational
double cover has branch degree two. The already
proved inequalities give
\[
                2\le4+2q_0,\qquad 2\le-2q_0+1,
\]
so the integer $q_0$ is $-1$. The cleared quartic
then has degree two and exactly two simple branch
points. Only one point lies above the singular
image germ. The other is a smooth simple branch
point, excluded by the actual coordinate reduction
\eqref{cdp:eq:A.T.10}--\eqref{cdp:eq:A.T.12}. Thus these two non-lci images
cannot support a connected double cover of any
nonnegative genus in a torus degeneration.

The integer branch budgets, the discriminant square
weights and the coordinate divisibility used in
\eqref{cdp:eq:A.T.10}--\eqref{cdp:eq:A.T.12} are checked in
\path{code/E_tacnode_bounds.py}; the non-lci charts,
symbolic generators and cubic/quartic pole orders
are checked in \path{code/E_low_delta_tacnode.py}.
The topological
quotient \eqref{cdp:eq:A.T.2} targets the actual singular curve
$D/J$; replacing it by its normalization would
lose the free-group obstruction in \eqref{cdp:eq:A.T.4}.




\begin{proposition}[Ordinary-triple images have arithmetic genus at most one]\label{cdp:st:A.5.19}
\StatementRecord{cdp:st:A.5.19}{proved}{Historical events: Claude Fable 5.2 (Anthropic), effort high, 16 November 2026: carried as background without fault found. GPT-6 Astra (OpenAI), ChatGPT Pro mode, effort Pro, S6\_09\_PAPERS, 17 November 2026: targeted proof check. Also checked in this writing. Pre-continuation writing events: GPT-6 Astra (OpenAI), ChatGPT Pro mode, effort Pro, S6\_10\_PAPERS, 18 November 2026: targeted proof check. Continuation metadata events: GPT-6 Astra (OpenAI), ChatGPT Pro mode, effort Pro, S6\_10\_PAPERS continuation, 19 November 2026: used and rechecked in this run; status-vocabulary reconciliation; metadata reconciliation only; no new mathematical check. \emph{Pre-continuation recorded proof scope.} The lci cubic discriminant and each non-lci symbolic-cube presentation were re-derived, including the space cusp, cusp-plus-line and three-axis rings. Their exact frame poles give negative cleared discriminant degree. The remaining etale covers, elliptic minimality, and the nonspherical local model were checked}{the hypotheses, notation and definitions stated here; \cref{cdp:st:A.5.18}; \cref{cdp:st:A.5.14}; \cref{cdp:st:A.5.5}; \cref{cdp:eq:A.U.2}; \cref{cdp:eq:A.U.4}; \cref{cdp:eq:A.U.5}; \cref{cdp:eq:A.U.1}; \cref{cdp:st:C2d}}

Let $f:Y\to C_0$ be a proper holomorphic map from a smooth
compact complex threefold with $b_2(Y)=0$, with connected
fibres and smooth general fibre a complex two-torus.
Let $mS$ be a full fibre with irreducible reduced support,
and let $D\subset\operatorname{Sing}S$ be an irreducible
curve with generic ordinary-triple plane transversal.
Then $D$ is a local complete intersection in $Y$ and
$p_a(D)\le1$. The possibilities for its underlying curve
are a smooth rational curve, a rational curve with one
ordinary node or cusp, and a smooth elliptic curve.
In the elliptic case the transversal type is constant.

When $D$ is singular rational, every component of the
normalized conductor curve dominating $D$ is rational,
and its map to $D^\nu\simeq\PP^1$ has degree one.
Consequently, in the irrational ruled setup of \ref{cdp:st:A.5.18}
with the generic transversal replaced by an ordinary
triple point, the image of the noncontracted horizontal
section is either smooth rational or smooth elliptic.
In the elliptic case the horizontal cover is unramified
of degree three, the normalization is smooth minimal
elliptic ruled, and the fibre has Euler number zero.
In the rational case the horizontal cover has degree
two or three; these cases are not excluded by this
statement.
\end{proposition}


\begin{proof}
\noindent\textbf{Step 1.} By \ref{cdp:st:A.5.14}, $p_a(D)\le2$. If its normalization has
positive genus, that statement and \ref{cdp:st:A.5.5} already give
a smooth elliptic curve of constant transverse type.
Suppose its normalization $\nu:\Sigma\to D$ is
$\PP^1$, with total delta invariant $\delta\le2$.

\noindent\textbf{Step 2.} First let $D$ be lci. Put
$N=(\nu^*(I_D/I_D^2))^*$ and
$M=\nu^*(\mathcal O_Y(S)|_D)$.
The quadratic normal jet vanishes generically, so
its vanishing descends on the reduced curve through
the locally free graded piece $I_D^2/I_D^3$.
Thus the initial cubic is an actual section of
$\operatorname{Sym}^3N^*\otimes M$.
Adjunction and the curve conductor give
$\deg N=-2+2\delta$, while $\deg M=0$.
Its nonzero binary-cubic discriminant is a section
of $M^4\otimes(\det N)^{-6}$, of degree
\[
                         12-12\delta.             \tag{A.U.1}\label{cdp:eq:A.U.1}
\]
Hence $\delta\le1$.

\noindent\textbf{Step 3.} We exclude every non-lci germ of delta at most two
by its symbolic cube. The one- and two-branch analytic
classification is proved in (A.T.14). There is one
additional possibility when three branches occur:
delta at least two is forced by identifying their
constants, so equality makes the ring exactly
\[
 \{(a,b,c)\in\C\{t_1\}\oplus\C\{t_2\}\oplus\C\{t_3\}:
                      a(0)=b(0)=c(0)\}.
\]
Its actual embedding in the smooth threefold is the
union of the three coordinate axes, with ideal
$I=(xy,xz,yz)$. Indeed the three branch parameters
give a basis of the cotangent space, so their lifts
are ambient coordinates as in \ref{cdp:st:A.5.18}.
Thus there are exactly three non-lci cases to consider;
each consumes the whole delta budget at one point.

\noindent\textbf{Step 4.} For the space cusp retain $I=(g_1,g_2,g_3)$ and
$F$ from (A.T.15). Its exact symbolic-cube identity is
\[
                         I^{(3)}=I^3+FI.           \tag{A.U.2}\label{cdp:eq:A.U.2}
\]
Here is a verification in the convergent local ring.
The right-hand side agrees with $I^3$ after inverting
$x$, by (A.T.15). Its quotient is finite over
$\C\{x\}$, of generic rank $3\cdot6=18$: a square
normal plane has six monomials of total degree less
than three. Modulo $x$, its ideal in $\C\{y,z\}$ is
\[
 (y,z)^6+(y^2z^3,yz^4,z^5).
\]
The fifteen monomials of total degree at most four,
together with $y^5,y^4z,y^3z^2$, give residue dimension
eighteen. The same discrete-valuation-module argument
used for (A.T.15) makes the quotient $x$-torsion-free.
It is therefore the saturation of $I^3$, which equals
the symbolic cube since the curve is lci off the origin.
This proves \eqref{cdp:eq:A.U.2}.

\noindent\textbf{Step 5.} For the cusp-plus-line ideal $I=(xz,yz,G)$ with
$G=y^2-x^3$, one has
\[
                   I^{(3)}=I^3+(zG)I.             \tag{A.U.3}\label{cdp:eq:A.U.3}
\]
To verify this directly, write $P=(z,G)$ and
$Q=(x,y)$. Their powers are primary because these
prime ideals are complete intersections. The
symbolic cube is $P^3\cap Q^3$.
Write an element of $P^3$ as
$z^3A+z^2GB+zG^2C+G^3E$.
The last two terms are already in $Q^3$.
Reducing the other terms modulo $Q^3$ gives
$zA+GB=0$ there. The constant-in-$(x,y)$ coefficient
of $B$ must therefore be divisible by $z$; write
$B=zB'+B_Q$, $B_Q\in Q$.
It follows that $A+GB'\in Q^3$.
Consequently
\[
 P^3\cap Q^3=(z^3Q^3,z^2GQ,zG^2,G^3)
                        =I^3+(zG)I,
\]
as claimed.

\noindent\textbf{Step 6.} For the three axes, membership of
$x^a y^b z^c$ in the symbolic cube means
$a+b,a+c,b+c\ge3$.
If one exponent is zero, the other two are at
least three and the monomial belongs to $I^3$.
If all are positive, either at least two are at
least two, in which case it belongs to $(xyz)I$,
or two are one, which violates one of the three
inequalities. Hence
\[
                         I^{(3)}=I^3+(xyz)I.       \tag{A.U.4}\label{cdp:eq:A.U.4}
\]

\noindent\textbf{Step 7.} In all three cases the generic ordinary-triple
equation $f_0$ belongs to $I^{(3)}$: on the smooth
locus of each curve branch its order in the two
normal variables is at least three. Define the
bundle by its actual conormal image,
\[
 K=\operatorname{im}\bigl(\nu^*(I/I^2)
                           \to\nu^*\Omega_Y^1\bigr),
 \qquad N=K^*.
\]
For the first two cases (A.T.18) gives $\deg N=4$.
For the three axes, on the branch $(t,0,0)$ the
frame is $t\,dy,t\,dz$, with conormal determinant
order two and no tangent ramification. Each of the
three branches has the same correction, for total
six, so again $\deg N=-2+6=4$.

\noindent\textbf{Step 8.} Use the actual affine maps (A.T.20)--(A.T.21).
For the three axes add
\[
              (u,v,t)\longmapsto(t,u/t,v/t)
                                                        \tag{A.U.5}\label{cdp:eq:A.U.5}
\]
on the first branch, and the two analogous permuted
maps on the other branches. In these charts the
first normal derivatives of the ideal generators
are holomorphic in the specified conormal frame.
An $I^3$ term consequently has holomorphic initial
cubic. The extra term in each of \eqref{cdp:eq:A.U.2}--\eqref{cdp:eq:A.U.4}
has initial cubic equal to the quadratic term of
its displayed symbolic-square generator, multiplied
by the first derivative of an ideal generator.

\noindent\textbf{Step 9.} For the space cusp that quadratic term is (A.T.16),
so the cubic pole is at most three. For the
cusp-plus-line germ it is $t^{-2}uv$ on the cusp
and $s^{-1}v^2$ on the smooth branch, so the pole
bounds are two and one. For the three axes it is
$t^{-1}uv$ in \eqref{cdp:eq:A.U.5}, giving pole bound one on
each branch. Coefficients multiplying these terms
are holomorphic; their derivatives do not enter
the initial cubic because the lower normal jets
of the products vanish.

\noindent\textbf{Step 10.} The initial cubic is intrinsic on the dense smooth
part of $D$, as all lower jets are zero. Its
transition under a normal-coordinate change is
therefore the actual cubic tensor in
$\operatorname{Sym}^3N^*\otimes M$; higher-order
coordinate terms make no contribution. The binary
discriminant is homogeneous of degree four in
these coefficients. It is a meromorphic section
of $M^4\otimes(\det N)^{-6}$, of degree $-24$,
with pole divisor of total degree at most twelve
in each of the three cases. After this exact pole
correction its line has negative degree, at most
$-12$. It cannot have the required nonzero
holomorphic section. This excludes all non-lci
cases and proves that $D$ is lci with $p_a(D)\le1$.
The list of underlying curves follows from the
delta-one classification in \ref{cdp:st:A.5.14}.

\noindent\textbf{Step 11.} We give the finite-cover comparison at a singular
rational image. Here $\delta=1$, so \eqref{cdp:eq:A.U.1} has
degree zero and its nonzero section is nowhere zero
on $\Sigma=\PP^1$. The cubic defines a subscheme
$R_3\subset\PP(N)$ of tangent-root lines. Every
fibre consists of three distinct points. In a
projective coordinate chart, its derivative with
respect to the root coordinate is nonzero, so the
implicit function theorem makes
$R_3\to\Sigma$ a finite etale map of degree three.
There is no component in a whole projective fibre,
since its nonzero discriminant forbids an identically
zero cubic there.

\noindent\textbf{Step 12.} Over the generic point of $D$, normalization of an
ordinary triple transversal separates exactly its
three tangent lines. This supplies an isomorphism
between the function-field algebra of $R_3$ and
that of the normalized conductor curves dominating
$D$, after their maps are factored through $\Sigma$.
Both are normal finite curves over $\Sigma$.
Taking the integral closure of $\mathcal O_\Sigma$
in this same algebra therefore identifies the two
finite covers everywhere. Every connected finite
etale cover of the simply connected Riemann sphere
has degree one. Thus each of these conductor
components is rational and maps isomorphically to
$\Sigma$, proving the second assertion.

\noindent\textbf{Step 13.} In the irrational ruled setting the horizontal
section has positive genus. Its finite map to a
rational image normalization must have degree at
least two. The preceding comparison excludes a
singular rational image, and the generic triple
transversal bounds its degree above by three.
The remaining rational image is therefore smooth,
with degree two or three. If the image is elliptic,
no vertical rational curve can dominate it, so all
three branches belong to the horizontal section.
Statement \ref{cdp:st:C2d} gives an unramified degree-three
cover, normal degree zero and constant type.
The exact minimal-resolution factorization and
blowdown argument (A.T.22) applies without change
to this coefficient-two section: its proof uses
only those properties. It makes the normalization
smooth minimal elliptic ruled with conductor twice
the section. The actual normalization Euler formula
(A.T.23) then gives zero.
\end{proof}



The singular rational possibilities cannot be
removed by a blanket local spherical-cycle argument.
For example
\[
                     f=x(yz)(x-yz)
                                                        \tag{A.U.6}\label{cdp:eq:A.U.6}
\]
has ordinary-triple transversals along the nodal curve
$D=(x,yz)=0$ and an additional ordinary-double
critical line $y=z=0$. For $t\ne0$, put
$B_t=\{(x,v):xv(x-v)=t\}$. The exact maps
\[
 B_t\times\C^*\longrightarrow f^{-1}(t),
 \quad (x,v,y)\longmapsto(x,y,v/y),\qquad
 (x,y,z)\longmapsto(x,yz,y)
                                                        \tag{A.U.7}\label{cdp:eq:A.U.7}
\]
are inverse holomorphic maps. Nonzero $t$ forces
$v,y$ and $z$ to be nonzero. The smooth affine
curve $B_t$ and $\C^*$ are aspherical, so this
regular affine fibre has no spherical homotopy
class. This is a local model, not a compact
$b_2=0$ example. A deformation smoothing $yz$
changes the additional critical line's transverse
type on the boundary and does not satisfy the
comparison used in \ref{cdp:st:A.5.17}.

Exact ideal, normal-coordinate and discriminant
calculations are retained in
\path{code/E_triple_images.py} and
\path{code/E_triple_images.out}; the independently
computed symbolic saturations are also in
\path{code/A_low_delta_cusp.py}.




\begin{proposition}[Rational normalization bounds with the exact constant-sheaf map]\label{cdp:st:A.5.16}
\StatementRecord{cdp:st:A.5.16}{proved}{Historical events: Claude Fable 5.2 (Anthropic), effort high, 16 November 2026: carried as background without fault found. GPT-6 Astra (OpenAI), ChatGPT Pro mode, effort Pro, S6\_09\_PAPERS, 17 November 2026: targeted proof check. Also checked in this writing. Pre-continuation writing events: GPT-6 Astra (OpenAI), ChatGPT Pro mode, effort Pro, S6\_10\_PAPERS, 18 November 2026: targeted proof check. Continuation mathematical and source checks: GPT-6 Astra (OpenAI), ChatGPT Pro mode, effort Pro, S6\_10\_PAPERS continuation, 19 November 2026: targeted proof check; used and rechecked in this run. Continuation metadata events: GPT-6 Astra (OpenAI), ChatGPT Pro mode, effort Pro, S6\_10\_PAPERS continuation, 19 November 2026: used and rechecked in this run; status-vocabulary reconciliation; metadata reconciliation only; no new mathematical check. \emph{Pre-continuation recorded proof scope.} The conductor duality, rational-resolution Leray maps, injection from topological to coherent H1, anticanonical boundary genus, and exact normalization-conductor Euler formula were re-derived. A negative curve-image Euler value would identify two points on the sole elliptic component and force boundary arithmetic genus at least two, contradicting the proved bound}{the hypotheses, notation and definitions stated here; \cref{cdp:st:A.5.17}; \cref{cdp:st:A.5.12}; \cref{cdp:eq:A.RT.5}; \cref{cdp:st:A.5.9}; \cref{cdp:eq:A.RT.9}; \cref{cdp:st:A.5.3}; \cref{cdp:eq:A.RT.8}}

Let $f:Y\to B$ be a proper holomorphic map from a smooth compact
complex threefold with $b_2(Y)=0$, with connected fibres and smooth
general fibre a two-torus. Let $mS$ be a full fibre with irreducible
nonnormal reduced support. Write $\eta:T\to S$ for normalization,
$\pi:Z\to T$ for its minimal resolution, and $\Gamma\subset S$
for the conductor-image scheme. Suppose that $Z$ is rational.
Put $Q=\eta_*\OO_T/\OO_S$ and $e=h^0(S,\omega_S)$.
Then $e\in\{0,1\}$, and
\[
\begin{gathered}
 R^1\pi_*\OO_Z=0,\qquad
 H^1(T;\mathbb C)=H^1(T,\OO_T)=H^2(T,\OO_T)=0,\\
 H^0(S,Q)\xrightarrow{\ \simeq\ }H^1(S,\OO_S),\qquad
 h^0(S,Q)=1+e,\quad h^1(S,Q)=e,\quad
 \chi(\Gamma,\OO_\Gamma)=-1.
\end{gathered}
\tag{A.RT.1}\label{cdp:eq:A.RT.1}
\]
In particular $T$ has rational singularities. The coefficient map
\[
 H^1(S;\mathbb C)\longrightarrow H^1(S,\OO_S)
\tag{A.RT.2}\label{cdp:eq:A.RT.2}
\]
is injective, so $b_1(S)\le1+e\le2$.

There is an effective integral divisor $D$ on $Z$ with
$\OO_Z(D)\simeq K_Z^{-1}$, whose strict-transform part is the
normalization conductor. It satisfies
\[
 h^0(D,\OO_D)=h^1(D,\OO_D)=1,
 \quad D_{\rm red}\text{ connected},
 \quad p_a(D_{\rm red})\le1.
\tag{A.RT.3}\label{cdp:eq:A.RT.3}
\]
The conductor image is connected,
$p_a(\Gamma_{\rm red})\le2$, and
$e(\Gamma_{\rm red})\ge0$. Set-theoretically
$D_{\rm red}=\pi^{-1}(\eta^{-1}(\Gamma_{\rm red}))$.
Exceptional curves over points outside the conductor have
coefficient zero in $D$; they are not to be counted as conductor
branches. With $E_p=(\pi^{-1}(p))_{\rm red}$ the exact Euler
formula is
\[
 e(S)=e(Z)-e(D_{\rm red})+e(\Gamma_{\rm red})
       -\sum_{p\in\operatorname{Sing}T\setminus\eta^{-1}(\Gamma)}
                  \bigl(e(E_p)-1\bigr).
\tag{A.RT.4}\label{cdp:eq:A.RT.4}
\]
In the torus smoothing of the statement, \ref{cdp:st:A.5.17} excludes
every point in this sum. Hence $T$ is smooth outside the
normalization conductor, and the sum is zero.
\end{proposition}


\begin{proof}
\noindent\textbf{Step 1.} The canonical evaluation argument in \ref{cdp:st:A.5.12} gives
$\omega_S^m\simeq\OO_S$ and $\chi(S,\OO_S)=0$.
Thus $e=0$ or $1$, with $e=1$ precisely when
$\omega_S\simeq\OO_S$, and
$h^1(S,\OO_S)=1+e$.

\noindent\textbf{Step 2.} We first prove a useful fact without the rationality assumption.
Finite duality identifies the normal dualizing sheaf with
\[
 \omega_T=\mathfrak c_T\otimes\eta^*\omega_S
       \ \subset\ \eta^*\omega_S,
\tag{A.RT.5}\label{cdp:eq:A.RT.5}
\]
where $\mathfrak c_T$ is the conductor ideal on $T$.
A nonzero global section of $\omega_T$ is therefore a nonzero
section of the torsion line $\eta^*\omega_S$ vanishing on its
nonempty conductor. Its $m$th tensor power would be a nonzero
holomorphic function on compact connected $T$ vanishing there,
which is impossible. Consequently
$H^2(T,\OO_T)=H^0(T,\omega_T)^\vee=0$.

\noindent\textbf{Step 3.} For the proper resolution $\pi$, normality gives
$\pi_*\OO_Z=\OO_T$, and its fibres have dimension at most one.
The low-degree Leray sequence is
\[
\begin{split}
0\longrightarrow H^1(T,\OO_T)\longrightarrow H^1(Z,\OO_Z)
 \longrightarrow H^0(T,R^1\pi_*\OO_Z)\\
 \longrightarrow H^2(T,\OO_T)\longrightarrow H^2(Z,\OO_Z)
 \longrightarrow0.
\end{split}
\tag{A.RT.6}\label{cdp:eq:A.RT.6}
\]
Rationality of $Z$ makes both its positive-degree groups zero.
Together with \eqref{cdp:eq:A.RT.5}, this gives $H^1(T,\OO_T)=0$ and
$R^1\pi_*\OO_Z=0$, since this last sheaf has finite support.
The fibres of $\pi$ are connected. Proper base change for the
constant sheaf gives $R^0\pi_*\mathbb C_Z=\mathbb C_T$.
The degree-one topological Leray edge map is consequently an
injection $H^1(T;\mathbb C)\hookrightarrow H^1(Z;\mathbb C)=0$.

\noindent\textbf{Step 4.} In the normalization exact sequence
\[
 0\longrightarrow\OO_S\longrightarrow\eta_*\OO_T
       \longrightarrow Q\longrightarrow0,
\tag{A.RT.7}\label{cdp:eq:A.RT.7}
\]
the map on $H^0$ is the identity on constants. The preceding
vanishings therefore make its connecting map
$H^0(Q)\to H^1(\OO_S)$ an isomorphism and its next connecting
map $H^1(Q)\to H^2(\OO_S)$ an isomorphism. These are the
maps and dimensions in \eqref{cdp:eq:A.RT.1}. In particular $\chi(Q)=1$.
The exact conductor-duality isomorphism proved in \ref{cdp:st:A.5.9} is
\[
 Q\simeq\omega_\Gamma\otimes\omega_S^{-1}|_\Gamma.
\tag{A.RT.8}\label{cdp:eq:A.RT.8}
\]
The conductor image is a proper Cohen--Macaulay curve, including
any nilpotent structure. Serre duality on it gives
$\chi(Q)=-\chi(\omega_S|_\Gamma)$. A torsion line bundle has
degree zero on every reduced component, and therefore also has
total degree zero on this possibly nonreduced curve. Curve
Riemann--Roch gives
$\chi(\omega_S|_\Gamma)=\chi(\OO_\Gamma)$.
This proves $\chi(\OO_\Gamma)=-1$.

\noindent\textbf{Step 5.} Here is the asserted injectivity of the complex coefficient map.
In the category of sheaves of complex vector spaces set
\[
 \mathcal V=\operatorname{coker}
     (\mathbb C_S\longrightarrow\eta_*\mathbb C_T).
\]
The diagram from constants to holomorphic functions induces the
specific sheaf map
\[
 \mathcal V\longrightarrow Q,
 \quad[(a_1,\ldots,a_r)]\longmapsto[(a_1,\ldots,a_r)] .
\tag{A.RT.9}\label{cdp:eq:A.RT.9}
\]
At a point $s\in S$, the $r$ entries refer to the connected
normalization germs above $s$. A tuple of branchwise constants
that comes from an element of $\OO_{S,s}$ has equal values on
all these germs, because evaluation at their common image point
is one scalar. It is therefore a diagonal tuple. This proves
that \eqref{cdp:eq:A.RT.9} is injective on every stalk. No $\OO_S$-module
structure on $\mathcal V$ is asserted or needed.

\noindent\textbf{Step 6.} A finite map has zero positive-degree direct images of the
constant sheaf, by proper base change on its finite fibres.
Since $H^1(T;\mathbb C)=0$ and the map on constants is an
isomorphism, the connecting map from
$0\to\mathbb C_S\to\eta_*\mathbb C_T\to\mathcal V\to0$
identifies $H^0(S,\mathcal V)$ with $H^1(S;\mathbb C)$.
Naturality of the connecting maps gives the commutative square
\[
\begin{array}{ccc}
 H^0(S,\mathcal V)&\xrightarrow{\ \simeq\ }&H^1(S;\mathbb C)\\
 \big\downarrow&&\big\downarrow\\
 H^0(S,Q)&\xrightarrow{\ \simeq\ }&H^1(S,\OO_S).
\end{array}
\tag{A.RT.10}\label{cdp:eq:A.RT.10}
\]
Its left arrow is injective by \eqref{cdp:eq:A.RT.9}, proving \eqref{cdp:eq:A.RT.2}.

\ref{cdp:st:A.5.3} gives the effective integral divisor $D$ and the exact
line-bundle relation
$K_Z\otimes\OO_Z(D)\simeq(\eta\pi)^*\omega_S$.
The torsion line on the right is trivial on the smooth rational
surface: its Picard group is torsion-free and has zero identity
component. Hence $\OO_Z(-D)\simeq K_Z$. The divisor sequence
$0\to K_Z\to\OO_Z\to\OO_D\to0$, rationality of $Z$ and
Serre duality prove \eqref{cdp:eq:A.RT.3}. In detail the first positive
cohomology of $K_Z$ is zero, and $H^2(K_Z)=\mathbb C$;
the positive cohomology of $\OO_Z$ is zero. Thus $D$ is
connected. Restriction $\OO_D\to\OO_{D_{\rm red}}$ is
surjective on $H^1$, because the kernel is supported in
dimension one. Consequently $p_a(D_{\rm red})\le1$.

\noindent\textbf{Step 7.} We identify its support, rather than discarding exceptional
components. On a minimal resolution every exceptional
irreducible curve $E$ has $K_Z\cdot E\ge0$. Indeed
$E^2<0$, and adjunction gives
$K_Z\cdot E=2p_a(E)-2-E^2$; a smooth rational
$(-1)$-curve is excluded by minimality. If $E$ has coefficient
zero in $D$, then
$0=(K_Z+D)\cdot E$ forces $K_Z\cdot E=D\cdot E=0$.
Thus $E$ is disjoint from $D_{\rm red}$.
Every exceptional fibre over the normalization conductor meets
its strict transform. Connectedness of that fibre then forces
all its exceptional curves to have positive coefficient in $D$.
Conversely, a positive-coefficient exceptional curve over a point
outside the conductor would lie in a component of $D_{\rm red}$
disjoint from every strict conductor curve. The connectedness
just proved, and the nonempty strict conductor, exclude this.
This proves the stated support equality and the zero-coefficient
assertion. It also proves connectedness of $\Gamma_{\rm red}$.

\noindent\textbf{Step 8.} For completeness the bound on its arithmetic genus follows
directly from \eqref{cdp:eq:A.RT.8}. One has
$h^1(\Gamma,\omega_S|_\Gamma)=h^0(Q)=1+e$,
and restriction to $\Gamma_{\rm red}$ is surjective on $H^1$.
For $e=1$ the line is trivial, giving
$p_a(\Gamma_{\rm red})\le2$. For $e=0$ its $H^1$ has
dimension at most one, and degree-zero Riemann--Roch gives
$p_a(\Gamma_{\rm red})=1-h^0(\omega_S|_{\Gamma_{\rm red}})
+h^1(\omega_S|_{\Gamma_{\rm red}})\le2$.
Finally normalization and resolution are isomorphisms off the
conductor and the finitely many exceptional fibres.
Stratifying by these disjoint closed sets gives \eqref{cdp:eq:A.RT.4} by
additivity of the compactly supported Euler characteristic.
At a point outside the conductor, $T\to S$ is an
isomorphism and $S$ is normal. A normal complex surface
has isolated singularities. The reduced local equation
therefore has an isolated critical point whenever that
point is singular. Since the full fibre has equation
$u h^m$, the isolated clause of \ref{cdp:st:A.5.17} excludes it.
This proves the final assertion and eliminates the displayed
sum by an actual local comparison with the nearby torus.

\noindent\textbf{Step 9.} We prove the asserted Euler bound for the conductor image,
including its possible positive-genus component. Put
$C=\Gamma_{\rm red}$. The normalization exact sequence for
the connected reduced curve $D_{\rm red}$ shows that the sum
of its normalized component genera is at most
$p_a(D_{\rm red})\le1$. Any component with genus one is
smooth, is unique, and meets the remaining rational curves
only in trees; a second attachment path joining distinct
points of that elliptic component would increase the
arithmetic genus to at least two. Each component of $C$ is
the finite image of a strict conductor component of
$D_{\rm red}$. Riemann--Hurwitz therefore shows that its
normalization is rational or elliptic, with at most one
elliptic-normalized component.

\noindent\textbf{Step 10.} Write $n_0$ for the number of rational-normalized components
of $C$, $n_1\in\{0,1\}$ for the number of its
elliptic-normalized components, and $r_p,\delta_p$ for the
number of branches and the delta invariant at a singular
point. The exact coherent and topological normalization maps
give
\[
 \chi(\OO_C)=n_0-\sum_p\delta_p,
 \qquad e(C)=2n_0-\sum_p(r_p-1).
\]
For each reduced curve germ, independent constants on its
normalized branches give $\delta_p\ge r_p-1$.
The already proved arithmetic-genus bound and connectedness
give $\chi(\OO_C)\ge-1$, hence $e(C)\ge n_0-1$.
If $n_0\ge1$ this is nonnegative. If $n_0=0$, there is
one elliptic-normalized component and
$\sum\delta_p\le1$. A negative value could therefore only
be $-1$, from one ordinary node. Indeed its curve ring has
colength one in the direct sum of its two normalization
rings, and is contained in the subring of pairs with equal
constant value; equality of colengths makes these rings
equal. The two points over that node are distinct on the
elliptic normalization.

\noindent\textbf{Step 11.} Here the hypersurface embedding excludes this last case.
At least one point $t\in T$ over the node receives both
local conductor branches. Otherwise partition the local
irreducible hypersurface branches according to which of
the two conductor branches they receive. Both parts are
nonempty. Two surface branches from opposite parts meet
in a curve, since their distinct equations in the regular
three-dimensional ambient local ring generate an ideal
of height two. That intersection curve belongs to the
conductor image, so is one of its two branches. Both
surface branches receive it, contradicting the partition.
For a locally irreducible hypersurface there is only one
normalization point and the assertion is immediate.

\noindent\textbf{Step 12.} The connected resolution fibre $\pi^{-1}(t)$ lies in
$D_{\rm red}$ by the support equality proved above.
The strict conductor curves over the two branches are
elliptic, because a rational curve cannot map
nonconstantly to an elliptic curve. There is at most one
such curve on $D_{\rm red}$. If the two strict curves
were different, they would already give total genus two.
If they coincide, they meet $\pi^{-1}(t)$ at distinct
points: the map of that normal elliptic curve factors
through the normalization of its image, so one point
cannot map to both points over the node. The connected
resolution fibre then supplies an additional path joining
those points, outside the elliptic curve. Its components
are rational, since $T$ has rational singularities.
This produces a graph cycle in addition to the elliptic
genus, again forcing $p_a(D_{\rm red})\ge2$.
Both alternatives are impossible. Thus $e(C)\ge0$.
\end{proof}



We record the exact minimal-model calculation that follows from
\ref{cdp:st:A.5.16}. This is a reduction of the outstanding Euler calculation;
the configurations with negative complement Euler number are
retained in the calculation.
Choose any sequence of contractions of smooth rational
$(-1)$-curves from $Z$ to a minimal rational surface $Z_0$,
and let $D_0$ be the pushforward of $D$. At one contraction
$\sigma:Z_i\to Z_{i-1}$, let $E$ be the exceptional curve,
$a$ its coefficient in $D_i$, and $r$ the multiplicity of
$D_{i-1}$ at the centre. The actual divisor comparison is
\[
 K_{Z_i}+D_i
   =\sigma^*(K_{Z_{i-1}}+D_{i-1})+(1+a-r)E.
 \tag{A.RT.11}\label{cdp:eq:A.RT.11}
\]
Both canonical-plus-boundary line bundles are trivial, so
intersection with $E$ gives $a=r-1$. In particular the centre
lies on $D_{i-1}$, and effectivity is retained at every stage.
If $r=1$, the centre is a smooth point of a reduced boundary
component, $a=0$, and the strict boundary is isomorphic to the
old boundary. The blowup increases its complement Euler number
by one. If $r\ge2$, the new reduced boundary is exactly the
full inverse image of the old reduced boundary. Both the
surface and the boundary replace the centre by $\mathbf P^1$,
so the complement Euler number is unchanged. Thus
\[
 e(Z\setminus D_{\rm red})
 =e(Z_0\setminus (D_0)_{\rm red})+N_0,
 \qquad N_0\ge0,
 \tag{A.RT.12}\label{cdp:eq:A.RT.12}
\]
where $N_0$ is the number of multiplicity-one centres in this
specified sequence of reverse blowups. Formula \eqref{cdp:eq:A.RT.12} concerns
the resolution complement. For the actual torus smoothing,
the isolated clause of \ref{cdp:st:A.5.17} has already removed the
off-conductor exceptional correction in \eqref{cdp:eq:A.RT.4}. This local
argument is needed: on an abstract minimal $\mathbb F_2$,
its exceptional $(-2)$-section can be disjoint from an
anticanonical divisor without having arisen from one of the
$N_0$ blowups.

On $\mathbf P^2$, the divisor $D_0$ is an effective cubic.
Its possible reduced supports and their complement Euler
numbers are
\[
\begin{array}{c|c}
 \text{reduced support}&e(\mathbf P^2\setminus(D_0)_{\rm red})\\ \hline
 \text{smooth cubic, nodal cubic, cuspidal cubic}&3,\ 2,\ 1\\
 \text{line and conic, transverse or tangent}&1,\ 0\\
 \text{three lines forming a triangle, or concurrent}&0,\ -1\\
 \text{two distinct lines, or one line}&0,\ 1.
\end{array}
\tag{A.RT.13}\label{cdp:eq:A.RT.13}
\]
For example the union of three concurrent lines has Euler
number $6-3+1=4$, whereas $e(\mathbf P^2)=3$.
The last row retains respectively a double line plus a line,
and a triple line. These are all cubic possibilities: an
irreducible singular cubic has arithmetic genus one and
therefore exactly a node or a cusp; the other degrees partition
three as $2+1$ or $1+1+1$.

On $\mathbb F_n$, with $n=0$ or $n\ge2$, write $C_0^2=-n$
for a minimal section and $F$ for a ruling fibre. Then
\[
 D_0\sim 2C_0+(n+2)F.
 \tag{A.RT.14}\label{cdp:eq:A.RT.14}
\]
Its horizontal part has degree two over the ruling base.
If that part consists of two distinct sections $C_1+C_2$,
write $V$ for the vertical part, $d$ for its total degree and
$k=C_1\cdot C_2$. Adjunction on either section gives
$k+d=2$. Let $q$ be the number of distinct intersection
points of $C_1$ and $C_2$, and let $r$ count those points
whose ruling fibre occurs in $V_{\rm red}$. A vertical fibre
meets the two sections in two points unless they coincide at
that fibre. Additivity therefore gives the exact formula
\[
 e((D_0)_{\rm red})=4-q+r,
 \qquad e(\mathbb F_n\setminus(D_0)_{\rm red})=q-r\ge0.
 \tag{A.RT.15}\label{cdp:eq:A.RT.15}
\]
This includes two sections and a fibre through their common
point: that union has Euler number four and complement Euler
number zero.

If the horizontal part is an irreducible bisection $C$, the
same adjunction calculation gives $p_a(C)=1-d$, so $d=0$
or $1$. For $d=0$, its Euler number is at most two, and the
complement Euler number is at least two. For $d=1$, the
bisection is a smooth rational curve, the vertical part is one
reduced fibre, and their set-theoretic intersection has one or
two points. The complement Euler number is respectively one
or two.

The remaining horizontal part is $2C$, where $C$ is a section.
Write $C\sim C_0+kF$. The vertical divisor has total degree
$d=n+2-2k$, and, with $v$ its number of distinct fibres,
\[
 e((D_0)_{\rm red})=2+v,
 \qquad e(\mathbb F_n\setminus(D_0)_{\rm red})=2-v.
 \tag{A.RT.16}\label{cdp:eq:A.RT.16}
\]
If $C\ne C_0$, positivity of intersection with $C_0$ gives
$k\ge n$; hence $d\le2-n$ and $v\le2$.
For $n=0$ one also has $d\le2$. Consequently the only
negative values on a minimal ruled surface occur in the
explicit configurations
\[
 Z_0=\mathbb F_n,\quad n\ge2,\qquad
 D_0=2C_0+\sum_{j=1}^{v}a_jF_j,\qquad
 a_j\ge1,\quad \sum_j a_j=n+2,\quad v\ge3.
 \tag{A.RT.17}\label{cdp:eq:A.RT.17}
\]
The fibres $F_j$ here are distinct. Equations
\eqref{cdp:eq:A.RT.11}--\eqref{cdp:eq:A.RT.17}, together with \eqref{cdp:eq:A.RT.4}, isolate the
remaining contributions exactly; an arbitrary identification
of their conductor branches has not been substituted for an
actual hypersurface normalization map.

There is a further fundamental-group restriction on precisely
these negative minimal-model configurations. Their reduced
supports are unions of rational curves with simply connected
underlying spaces: three concurrent lines on $\mathbf P^2$,
or the section with its attached ruling fibres in \eqref{cdp:eq:A.RT.17}.
Every reverse blowup in \eqref{cdp:eq:A.RT.11} preserves this property.
A multiplicity-one blowup leaves the support isomorphic to
the preceding one. A positive-coefficient blowup at a smooth
point adds a rational leaf, at an ordinary double point
subdivides the incidence edge, and at the initial concurrent
point replaces that vertex by an exceptional $\mathbf P^1$
with three attached branches. Subsequent supports have only
the resulting ordinary double points. None of these explicit
operations introduces a topological loop. Consequently
$D_{\rm red}$ is simply connected throughout these cases.

The actual map $h=\eta\pi$ expresses $S$ as the topological
pushout of
\[
 D_{\rm red}\hookrightarrow Z,
 \qquad D_{\rm red}\xrightarrow{\ h|_{D_{\rm red}}\ }
                           \Gamma_{\rm red}.
 \tag{A.RT.18}\label{cdp:eq:A.RT.18}
\]
Indeed $h$ is an isomorphism off these closed sets and the
whole preimage of $\Gamma_{\rm red}$ is $D_{\rm red}$.
The continuous map from this pushout to $S$ is a bijection
from a compact space to a Hausdorff space, and hence is a
homeomorphism. A triangulation of the pair
$(Z,D_{\rm red})$ makes its inclusion a cofibration.
The double mapping-cylinder model for \eqref{cdp:eq:A.RT.18} therefore
has the same homotopy type as the pushout, and van Kampen
gives
\[
 \pi_1(S)\simeq
 \pi_1(Z)*_{\pi_1(D_{\rm red})}\pi_1(\Gamma_{\rm red})
 \simeq\pi_1(\Gamma_{\rm red}).
 \tag{A.RT.19}\label{cdp:eq:A.RT.19}
\]
All normalizations of the image components are $\mathbf P^1$,
by Riemann--Hurwitz. Gluing their finitely many points over
singularities identifies the topological curve with a graph
of simply connected spheres. Its fundamental group is a
free group, of rank $b_1(\Gamma_{\rm red})$.

That free group is a quotient of a metabelian group. To see
the exact comparison, triangulate $(Y,S)$ and retract a small
regular neighbourhood $N$ onto $S$, as in \ref{cdp:st:A.5.11}. Properness
puts a sufficiently small full fibre tube $U=f^{-1}(\Delta)$
inside $N$. Its retraction is surjective on fundamental
groups because it restricts to the identity on $S$.
General position makes
$\pi_1(U\setminus S)\to\pi_1(U)$ surjective, and the
torus bundle over $\Delta^*$ has
$\pi_1(U\setminus S)=\mathbb Z^4\rtimes\mathbb Z$.
This last group has abelian commutator subgroup, so its
quotient \eqref{cdp:eq:A.RT.19} is metabelian. A free group of rank at
least two is not metabelian: it maps onto $S_4$, whose
commutator group $A_4$ is not abelian. We obtain
\[
 \pi_1(S)\simeq\pi_1(\Gamma_{\rm red})\in\{1,\mathbb Z\},
 \qquad b_1(\Gamma_{\rm red})\le1.
 \tag{A.RT.20}\label{cdp:eq:A.RT.20}
\]
Under the full CDP Betti hypotheses, the nonzero class on
every support from \ref{cdp:st:D.1.9} excludes the trivial group in
\eqref{cdp:eq:A.RT.20}. This last use of the global transgression is
separate from the preceding $b_2(Y)=0$ calculations.

For these two boundary classes one can further exclude a
one-component image with first Betti number one. Its normalization
is $\mathbf P^1$, and exactly one point $q$ identifies two distinct
normalization points $\alpha,\beta$. Other singular points are
unibranch and do not change this topological quotient. The universal
cover of the image is an infinite chain of normalized spheres $C_j$,
with $\beta_j$ identified with $\alpha_{j+1}$.

The actual map $h|_{D_{\rm red}}$ lifts to this cover because
$D_{\rm red}$ is simply connected. Each noncontracted rational
component factors through the normalization of the image and maps
surjectively onto it; its lift therefore covers one whole $C_j$.
Each contracted component lifts to a point. There are finitely many
components, so their connected lifted image is a finite chain
$C_m,\ldots,C_M$.

Choose a preimage in $D_{\rm red}$ of the left endpoint node and
let $t_-$ be its image under $\pi$ in the normalization $T$ of $S$.
The resolution fibre $\pi^{-1}(t_-)$ is connected and lies in
$D_{\rm red}=h^{-1}(\Gamma_{\rm red})$. Its lift is a single point
of the discrete fibre over $q$. Every local conductor branch of $T$
through $t_-$ has a strict transform meeting this resolution fibre.
It can receive only the branch directed into $C_m$: receiving the
other branch would force the lifted image to contain $C_{m-1}$.
At the right endpoint the same argument produces a point $t_+$
receiving only the opposite branch. The two points are distinct,
because their connected resolution fibres lift to different endpoint
nodes, including when $m=M$.

Distinct normalization points over $q$ correspond to distinct local
irreducible hypersurface branches of $S$ in the smooth threefold.
Their two distinct equations have an intersection curve: the ideal
they generate has height two. Every such curve is an actual
conductor-image branch received by both normalization points.
There are only the two image branches at $q$, whereas $t_-$ and
$t_+$ receive disjoint ones. This contradiction excludes the
one-component, first-Betti-number-one image.

Since every image component is rational-normalized, write
$n_\Gamma$ for their number and $\beta_\Gamma\le1$ for the
first Betti number. Their exact topological normalization
formula now yields, on these negative boundary classes,
\[
 e(\Gamma_{\rm red})=n_\Gamma+1-\beta_\Gamma\ge2.
 \tag{A.RT.21}\label{cdp:eq:A.RT.21}
\]
For one component the excluded case is the only value below
two; for at least two components the inequality follows
directly from $\beta_\Gamma\le1$.




\begin{proposition}[A homogeneous binary-product grid is excluded]\label{cdp:st:A.5.21}
\StatementRecord{cdp:st:A.5.21}{proved}{Historical events: Claude Fable 5.2 (Anthropic), effort high, 16 November 2026: carried as background without fault found. GPT-6 Astra (OpenAI), ChatGPT Pro mode, effort Pro, S6\_09\_PAPERS, 17 November 2026: targeted proof check. Also checked in this writing. Pre-continuation writing events: GPT-6 Astra (OpenAI), ChatGPT Pro mode, effort Pro, S6\_10\_PAPERS, 18 November 2026 04:23:40Z: targeted proof check. Continuation metadata events: GPT-6 Astra (OpenAI), ChatGPT Pro mode, effort Pro, S6\_10\_PAPERS continuation, 19 November 2026: used and rechecked in this run; status-vocabulary reconciliation; metadata reconciliation only; no new mathematical check. \emph{Pre-continuation recorded proof scope.} proved as stated. Full constant- and variable-coefficient binary-grid exclusion, including the actual fibre comparisons. Re-derived the unit-absorbing biholomorphism and cyclic degree-d cover of the complement of the projective root set. A path between two simple v-zeros lifts to a sphere with local discs z=c·conjugate(y); two distinct phases have intersection −1 at each endpoint, yielding square −2. Weighted homothety embeds the cycles in an actual nearby fibre. For variable coefficients, distinct analytic roots factor h\_s, its critical zero locus is the same two axes, and the normal-gradient lower bound dominates the tangential boundary correction. The proper critical discriminant complement gives orientation-preserving transport; the exact scaling identity embeds the transported sphere in the global torus. Its nullhomotopic sphere image contradicts its retained intersection square. Review provenance: New derivation after reading the supplied proof; not a blind review and not an independent human check}{the hypotheses, notation and definitions stated here; \cref{cdp:st:A.5.15}; \cref{cdp:eq:A.BP.2}; \cref{cdp:eq:A.BP.8}; \cref{cdp:eq:A.BP.6}}

Let $f:Y\to C_0$ be a proper holomorphic map from a smooth
complex threefold with smooth nearby fibres complex
two-tori. Let $P(X,V)$ be a homogeneous polynomial of
degree $d\ge2$, a product of distinct complex linear
forms, and suppose $P(1,0)\ne0$.
There is no point at which the actual fibre map has
the local equation
\[
             f=u(x,y,z)\,P(x,yz)^m,
             \qquad m\ge1,                         \tag{A.BP.1}\label{cdp:eq:A.BP.1}
\]
with $u$ a holomorphic unit. In particular this excludes
$\prod_{i=1}^d(x-a_i yz)$ for distinct constants $a_i$,
and every positive full-fibre multiplicity of that germ.
No Betti-number assumption on $Y$ is required.

The same exclusion holds for the variable-coefficient
equation $f=u h^m$ with
\[
 h=\sum_{i=0}^d a_i(x,y,z)x^{d-i}(yz)^i,
 \qquad P_0(X,V)=\sum_{i=0}^d a_i(0)X^{d-i}V^i,
                                                        \tag{A.BP.4}\label{cdp:eq:A.BP.4}
\]
provided $P_0$ has $d$ distinct projective roots and
$P_0(1,0)\ne0$. The comparison with its constant
initial form is the explicit analytic family below.
\end{proposition}


\begin{proof}
\noindent\textbf{Step 1.} Put $h=P(x,yz)$. It is weighted homogeneous of degree
$2d$ for the positive weights $(2,1,1)$.
On a small simply connected coordinate neighbourhood
take a holomorphic logarithm of the unit and set
\[
 \Phi(x,y,z)=\bigl(xu^{1/(md)},
                  yu^{1/(2md)},zu^{1/(2md)}\bigr).
                                                        \tag{A.BP.2}\label{cdp:eq:A.BP.2}
\]
Its derivative at the origin is invertible, and direct
homogeneity gives $h\Phi=u^{1/m}h$. Thus
$f=(h\Phi)^m$ exactly. It suffices to construct a
compact spherical cycle of nonzero square in $h=1$,
because weighted homotheties then put it inside an
arbitrarily small actual fibre of \eqref{cdp:eq:A.BP.1}.

\noindent\textbf{Step 2.} The curve $B=\{P(x,v)=1\}\subset\C^2$ is smooth:
Euler's identity $xP_x+vP_v=dP$ precludes a common
zero of its two partial derivatives on that level.
It is also connected. The map
\[
 B\longrightarrow\PP^1\setminus\{P=0\},
                       \quad(x,v)\longmapsto[x:v]
\]
is an etale cover of degree $d$. Locally a projective
representative has nonzero $P$-value, and its $d$
rescalings with $P=1$ give the sheets. Around any
one of the simple projective zeros of $P$, that
value has a simple zero, so its inverse $d$th root
has monodromy a full $d$-cycle. This monodromy is
transitive on the sheets, proving connectedness.

\noindent\textbf{Step 3.} The function $v$ on $B$ has exactly $d$ distinct
simple zeros, since
$P(x,0)=P(1,0)x^d$ and $P_x(x,0)=d/x$ at those
points. Choose an embedded arc $\gamma$ in $B$
joining two of its zeros and avoiding all other
zeros. On the interior choose a continuous argument
of $v(\gamma(t))$ and two functions $A(t),B_0(t)$
with $A B_0=v(\gamma(t))$ and
$|A|=|B_0|=|v(\gamma(t))|^{1/2}$.
The map
\[
 (t,\theta)\longmapsto
   \bigl(x(\gamma(t)),A(t)e^{\mathrm i\theta},
                        B_0(t)e^{-\mathrm i\theta}\bigr)
                                                        \tag{A.BP.3}\label{cdp:eq:A.BP.3}
\]
lies in $h=1$ and collapses both boundary circles
of $[0,1]\times S^1$ to points. It is therefore a
spherical cycle, with embedded interior circles.

\noindent\textbf{Step 4.} A parallel-arc isotopy with the same endpoints gives
a homologous spherical cycle whose support meets
the first only at those endpoints. At each endpoint,
$P_x\ne0$, so $(y,z)$ are actual holomorphic
coordinates on $h=1$. The links of the two cycles
are parallel curves of slope $(1,-1)$ on a Clifford
torus in the local three-sphere. Their linking
number is $-1$. As in the explicit relative-chain
comparison in \ref{cdp:st:A.5.15}, that number is the local
intersection of the two coned cycles. Both
orientations reverse at the second endpoint,
so its intersection sign is unchanged. The
spherical class consequently has square $-2$.

\noindent\textbf{Step 5.} The weighted homothety
$(x,y,z)\mapsto(\lambda^2x,\lambda y,\lambda z)$
maps $h=1$ biholomorphically to $h=\lambda^{2d}$
and shrinks the compact cycles and their isotopy
to the origin. Together with \eqref{cdp:eq:A.BP.2}, this puts
them inside an actual smooth global fibre of $f$.
That fibre is a complex two-torus and has zero
$\pi_2$, so the spherical class has zero homology
image. The oriented codimension-zero embedding
preserves the intersection of the compact cycles,
giving both zero and $-2$ for its square. This
contradiction proves the statement.

\noindent\textbf{Step 6.} We now prove the variable-coefficient assertion,
including the boundary comparison. Absorb a local
$m$th root of the unit into $h$; this preserves the
form \eqref{cdp:eq:A.BP.4} and its two initial-form conditions.
For a complex parameter $s$ set
\[
 h_s=\sum_{i=0}^d a_i(s^2x,sy,sz)x^{d-i}(yz)^i.
                                                        \tag{A.BP.5}\label{cdp:eq:A.BP.5}
\]
This is holomorphic also at $s=0$, with
$h_0=P_0(x,yz)$. For $s\ne0$ there is the exact
identity
\[
       h(s^2x,sy,sz)=s^{2d}h_s(x,y,z).           \tag{A.BP.6}\label{cdp:eq:A.BP.6}
\]
The leading coefficient $a_0$ is a unit. The implicit
function theorem applied to its degree-$d$ root
polynomial gives distinct holomorphic functions
$b_i(s,x,y,z)$ and the exact factorization
\[
       h_s=a_0(s^2x,sy,sz)
                     \prod_{i=1}^d(x-b_i(s,x,y,z)yz)
                                                        \tag{A.BP.7}\label{cdp:eq:A.BP.7}
\]
on one small parameter and coordinate neighbourhood.
Each factor defines a smooth surface, since its
$x$ derivative is a unit. For $i\ne j$ the ideal
generated by the two factors equals $(x,yz)$:
their difference is the unit $b_j-b_i$ times $yz$.
Thus the zero level has precisely the same two
critical curves $D=(x,yz)=0$ for every $s$.

\noindent\textbf{Step 7.} Choose a small closed ball $K$ in the coordinate
chart such that the two axes of $D$ meet its boundary
transversely and the zero-level surfaces away from
those axes are transverse to the boundary. At
$s=0$ this follows directly for the smooth graph
surfaces $x=b_i(0)yz$ and sufficiently small radii.
It persists for small $s$ by compactness away from
neighbourhoods of the two boundary links.
Near either link, use normal coordinates $(U,V)$
for the smooth critical curve and a parameter $w$
along it. The normal initial term of $h_s$ is a
degree-$d$ binary form with distinct roots.
On the compact parameterized link its normal
gradient therefore satisfies a uniform bound
\[
 \|\partial_{U,V}h_s\|\ge c\|(U,V)\|^{d-1},
 \qquad |\partial_w h_s|\le C\|(U,V)\|^d
                                                        \tag{A.BP.8}\label{cdp:eq:A.BP.8}
\]
after the neighbourhood is shrunk. The first bound
follows because the gradient of a homogeneous
reduced binary form has no zero on the unit sphere;
higher normal terms have one additional order.
The second follows from $h_s\in(U,V)^d$.

\noindent\textbf{Step 8.} Transversality of the critical curve to $\partial K$
allows a boundary coordinate with one real component
of $w$ solved as a smooth function of $(U,V)$ and
the remaining component. The correction to normal
differentiation tangent to this boundary is a
bounded multiple of $\partial_w h_s$.
By \eqref{cdp:eq:A.BP.8} it is smaller by one normal order than
the surjective complex normal derivative. Thus
$h_s|_{\partial K}$ is a real submersion at every
nearby point with $h_s\ne0$, uniformly in $s$.
Away from the two link neighbourhoods, compactness
and the preceding zero-level transversality give
the same conclusion for all sufficiently small
nonzero values.

\noindent\textbf{Step 9.} Remove the proper critical-value discriminant from
the product of the parameter disc and a sufficiently
small punctured value disc. The critical locus over
that domain cannot meet the boundary, so its
projection is proper and its image is complex
analytic. Its complement is connected. The level
sets in $K$ form a proper submersion with boundary
there, and Ehresmann transport identifies their
compact fibres by oriented diffeomorphisms.
Choose a small sphere cycle of the already proved
model $h_0=P_0(x,yz)$, and transport it to a fibre
of $h_{s_0}$ for $s_0\ne0$. Equation \eqref{cdp:eq:A.BP.6}
then embeds that fibre into the actual fibre
of $h^m$ with value $s_0^{2dm}\epsilon^m$.
It carries the spherical class of square $-2$
into a global torus. The same intersection and
asphericity contradiction proves \eqref{cdp:eq:A.BP.4}.
\end{proof}



The condition $P(1,0)\ne0$ is used to produce the
distinct collapsing circles. The exact product
model (A.U.6), for which $P(X,V)=XV(X-V)$, has
$P(1,0)=0$ and is outside this statement.
The variable-coefficient extension uses only the
specific comparison \eqref{cdp:eq:A.BP.5}--\eqref{cdp:eq:A.BP.8}, with its
unchanged boundary critical curves. The weighted scaling,
simple-zero and endpoint-square identities are
checked in \path{code/E_binary_products.py}.

\paragraph{Two exact scalar-zero comparisons.}
The homogeneous germ
\[
                  h=xy\{z(x-y)+xy\}
\tag{A.BP.9}\label{cdp:eq:A.BP.9}
\]
has a cubic normal tensor along $(x,y)=0$ which
vanishes identically at $z=0$. On $xy\ne0$ put
\[
 v=x-y,\qquad Z=y(y+v)z+y^2(2y+v).
\]
These are actual coordinates with $y$, with inverse
$x=y+v$, $z=(Z-y^2(2y+v))/(y(y+v))$, and
\[
                              h=vZ+y^4.
\tag{A.BP.10}\label{cdp:eq:A.BP.10}
\]
Choose an arc between two roots of $y^4=1$ avoiding
$y=0$ and the other roots. The matching circles
$v=\rho A e^{\mathrm i\theta}$,
$Z=\rho^{-1}B e^{-\mathrm i\theta}$, where
$AB=1-y^4$ and $|A|=|B|=|1-y^4|^{1/2}$,
collapse at its endpoints. Choose the positive
constant $\rho$ small enough that $|v|<|y|/2$
on the arc. The entire spherical cycle then lies
in the indicated coordinate chart. Its two endpoint
linkings are $-1$, as in \eqref{cdp:eq:A.BP.3}, so its square
is $-2$. Ordinary homothety in $(x,y,z)$ shrinks
it because \eqref{cdp:eq:A.BP.9} has degree four. The same
weighted unit map as \eqref{cdp:eq:A.BP.2}, with all three
weights equal to one, includes every local unit
and every positive full-fibre multiplicity.
Thus this particular scalar-zero germ is excluded.

The following different scalar-zero model records
the limit of that argument:
\[
                     h=xz(y^2-x^2z).
\tag{A.BP.11}\label{cdp:eq:A.BP.11}
\]
For $h=t\ne0$, the substitution $u=y/x$ is
invertible and gives $h=x^3 z(u^2-z)$.
The map forgetting $x$ is a finite etale triple
cover of $E=\{z(u^2-z)\ne0\}$.
There is an explicit finite etale double cover
\[
 \C^*\times(\C\setminus\{1,-1\})\longrightarrow E,
                  \qquad(s,v)\longmapsto(s^2,sv)
\tag{A.BP.12}\label{cdp:eq:A.BP.12}
\]
in coordinates $(z,u)$, with deck involution
$(s,v)\mapsto(-s,-v)$. Its source is aspherical,
so both $E$ and the regular affine fibre of
\eqref{cdp:eq:A.BP.11} are aspherical by these actual covering
maps. The quadratic surface branch has normalization
$(v,u)\mapsto(v,vu,u^2)$ and conductor
$2(v=0)+2(u=0)$ in the full reduced surface.
The other two normalization branches have conductors
$2(y=0)+(z=0)$ on $x=0$, and
$(x=0)+2(y=0)$ on $z=0$.
Thus their reduced conductor supports have only
normal-crossing pairs. This is an exact local
model, not a construction of a compact $b_2=0$
fibre. It shows why a vanishing cubic coefficient
alone does not justify a local spherical-cycle
exclusion. The identities and the two finite-cover
formulas are checked in \path{code/E_cubic_double.py}.



\section{Normal-crossing topology and canonical-flat component comparisons}\label{sec:II:08}


\begin{proposition}[The normalized double curves and the coherent Euler formula]\label{cdp:st:C.2.NC.2}
\StatementRecord{cdp:st:C.2.NC.2}{proved}{Historical events: Claude Fable 5.2 (Anthropic), effort high, 16 November 2026: carried as background without fault found. GPT-6 Astra (OpenAI), ChatGPT Pro mode, effort Pro, S6\_09\_PAPERS, 17 November 2026: targeted proof check. Also checked in this writing. Pre-continuation writing events: GPT-6 Astra (OpenAI), ChatGPT Pro mode, effort Pro, S6\_10\_PAPERS, 18 November 2026: targeted proof check. Continuation metadata events: GPT-6 Astra (OpenAI), ChatGPT Pro mode, effort Pro, S6\_10\_PAPERS continuation, 19 November 2026: used and rechecked in this run; status-vocabulary reconciliation; metadata reconciliation only; no new mathematical check. \emph{Pre-continuation recorded proof scope.} Adjunction on normalized double curves, triple-point incidence, the alternating sign-line normalization resolution, and both topological Euler identities were re-derived}{the hypotheses, notation and definitions stated here; \cref{cdp:eq:C.NC.6}; \cref{cdp:eq:C.NC.7}}

Let $S$ be a connected compact normal-crossings divisor in a smooth compact complex threefold $Y$ with $H^2(Y,\Q)=0$. Write $\nu:\coprod_iT_i\to S$ for normalization, $D_i\subset T_i$ for its conductor, $C_\alpha$ for the connected components of the normalization of the double locus, and $\tau$ for the number of triple points. Let $R_\alpha$ be the reduced divisor on $C_\alpha$ of preimages of triple points, counted separately for the three double-curve branches at each triple point. Then
\[
 e(C_\alpha)=\deg R_\alpha,\qquad
 \sum_\alpha e(C_\alpha)=3\tau,
 \qquad
 \sum_i\chi(T_i,\OO_{T_i})=\frac{\tau}{2}.
 \tag{C.NC.4}\label{cdp:eq:C.NC.4}
\]
In particular every normalized double curve has genus at most one, and $\tau$ is even. The remaining topological identities are
\[
\begin{aligned}
 \sum_i e(D_i)&=2\sum_\alpha e(C_\alpha)-3\tau,\\
 e(S)&=\sum_i e(T_i)-\sum_\alpha e(C_\alpha)+\tau.
\end{aligned}
 \tag{C.NC.5}\label{cdp:eq:C.NC.5}
\]
\end{proposition}


\begin{proof}
\noindent\textbf{Step 1.} Adjunction and the local residue comparison give a canonical holomorphic line-bundle isomorphism
\[
 \nu^*\omega_S|_{T_i}\simeq\omega_{T_i}(D_i),
 \qquad \omega_S=(\omega_Y\otimes\OO_Y(S))|_S.
 \tag{C.NC.6}\label{cdp:eq:C.NC.6}
\]
For example, the residue of $dx\wedge dy\wedge dz/(xy)$ on the branch $x=0$ is $\pm dy\wedge dz/y$, and at $xyz=0$ it is $\pm dy\wedge dz/(yz)$. These expressions agree under changes of local equations by the usual determinant and unit factors of the displayed line bundles, proving the comparison. Since both line bundles on the right in the expression for $\omega_S$ come from $Y$, its first Chern class and that of $\omega_{T_i}(D_i)$ vanish rationally.

\noindent\textbf{Step 2.} Normalize every irreducible curve in every $D_i$. The resulting disjoint union $\widehat D$ maps to $\coprod_\alpha C_\alpha$ by an étale map of total degree two over each $C_\alpha$. Away from a triple point this follows from the two branches of $xy=0$. At a triple point, each of the three normalized double branches has exactly two preimages, and their maps in the remaining coordinate are local isomorphisms. This proves étaleness at the missing points as well. The covering over a given $C_\alpha$ may be connected; no splitting is assumed.

\noindent\textbf{Step 3.} Adjunction for a nodal divisor on a smooth surface, pulled back to $\widehat D$, identifies the restriction of $\omega_{T_i}(D_i)$ with the canonical line of that normalized conductor curve twisted by all its preimages of nodes of $D_i$. Such nodes are exactly the preimages of triple points of $S$. This is the elementary local dualizing-form identity $dw/w$ on each branch of a node. The resulting degree is zero by \eqref{cdp:eq:C.NC.6}. Since the covering of $C_\alpha$ is étale and its marked divisor is the pullback of $R_\alpha$, it follows that
$\deg\omega_{C_\alpha}+\deg R_\alpha=0$.
Thus $e(C_\alpha)=\deg R_\alpha$. There are exactly three marks over every triple point, even when two or three belong to one $C_\alpha$, proving their summed formula.

\noindent\textbf{Step 4.} For the coherent calculation, the alternating branch-resolution sequence is
\[
 0\longrightarrow\OO_S\longrightarrow\nu_*\OO_{\coprod T_i}
 \longrightarrow\bigoplus_\alpha(j_\alpha)_*\varepsilon_\alpha
 \longrightarrow\bigoplus_{p\in S^{[3]}}\C_p
 \longrightarrow0.
 \tag{C.NC.7}\label{cdp:eq:C.NC.7}
\]
Here $j_\alpha:C_\alpha\to S$ is the finite normalization map of that double curve, and $\varepsilon_\alpha$ is the sign line of its two-sheeted conductor covering. Locally order the branches and send $(f_i)$ to $(f_j-f_i)_{i<j}$; the next map is the alternating sum of the three values at a triple point. The sequence is exact at $xy=0$ because a pair of functions glues precisely when its restrictions agree, and at $xyz=0$ because pairwise differences satisfy exactly the single alternating relation. Equivalently, compare coefficients of convergent power series on the three coordinate planes; equality on their coordinate axes is exactly the condition for these series to come from one series modulo $xyz$. Changing a local branch ordering multiplies the appropriate difference by $-1$, producing $\varepsilon_\alpha$. At a triple point this sign convention only chooses a basis of the one-dimensional last term. Hence the sequence is globally well-defined with the indicated local systems, including self-intersections.

\noindent\textbf{Step 5.} The sign line satisfies $\varepsilon_\alpha^{\otimes2}\simeq\OO_{C_\alpha}$, so it has degree zero and $\chi(C_\alpha,\varepsilon_\alpha)=1-g(C_\alpha)=e(C_\alpha)/2$. The divisor sequence on $Y$ and Hirzebruch--Riemann--Roch give $\chi(S,\OO_S)=0$, since $c_1(\OO_Y(S))=0$ rationally. Taking Euler characteristics in \eqref{cdp:eq:C.NC.7} now gives
\[
 0=\sum_i\chi(T_i,\OO_{T_i})-\frac12\sum_\alpha e(C_\alpha)+\tau,
\]
which proves the last equality in \eqref{cdp:eq:C.NC.4}.

\noindent\textbf{Step 6.} Finally $\widehat D$ has Euler number $2\sum_\alpha e(C_\alpha)$. Each triple point produces three nodes in $\coprod D_i$, and identifying the two preimages of each node subtracts one. This proves the first identity in \eqref{cdp:eq:C.NC.5}. The double locus in $S$ has normalization $\coprod C_\alpha$ and three preimages over each triple point, so its Euler number is $\sum_\alpha e(C_\alpha)-2\tau$. Normalization is an isomorphism off the conductor. Additivity of Euler characteristic therefore gives the second identity in \eqref{cdp:eq:C.NC.5}.
\end{proof}






\begin{proposition}[Uniform multiple normal-crossings fibres]\label{cdp:st:C.2.NC.3}
\StatementRecord{cdp:st:C.2.NC.3}{proved}{Historical events: Claude Fable 5.2 (Anthropic), effort high, 16 November 2026: carried as background without fault found. GPT-6 Astra (OpenAI), ChatGPT Pro mode, effort Pro, S6\_09\_PAPERS, 17 November 2026: targeted proof check. Also checked in this writing. Pre-continuation writing events: GPT-6 Astra (OpenAI), ChatGPT Pro mode, effort Pro, S6\_10\_PAPERS, 18 November 2026: targeted proof check. Continuation metadata events: GPT-6 Astra (OpenAI), ChatGPT Pro mode, effort Pro, S6\_10\_PAPERS continuation, 19 November 2026: used and rechecked in this run; status-vocabulary reconciliation; metadata reconciliation only; no new mathematical check. \emph{Pre-continuation recorded proof scope.} The uniform-multiplicity nearby-cycle formula and the normalized root cover w\textasciicircum{}m=u were re-derived, allowing a disconnected cover}{the hypotheses, notation and definitions stated here; \cref{cdp:st:C.2.NC.1}}

Let $g:Y\to C$ be a proper holomorphic map from a connected smooth compact complex threefold with $H^2(Y,\Q)=0$ to a smooth compact curve, with connected fibres and smooth general fibre of Euler characteristic zero. A fibre $mS$ with reduced normal-crossings support has
\[
 e(S)=\tau(S)\ge0.
 \tag{C.NC.8}\label{cdp:eq:C.NC.8}
\]
The normalized base change $t=s^m$ is finite étale of degree $m$ over a neighbourhood of $S$ and has reduced normal-crossings central fibre $S'$. Its possibly disconnected components do not affect the comparisons
$e(S')=m e(S)$ and $\tau(S')=m\tau(S)$.
\end{proposition}


\begin{proof}
\noindent\textbf{Step 1.} The local formula \ref{cdp:st:C.2.NC.1} gives
\[
 0=e(Y_t)=m\sum_i e(T_i\setminus D_i).
\]
By (C.NC.5),
\[
 e(S)-\sum_i e(T_i\setminus D_i)
 =\sum_\alpha e(C_\alpha)-2\tau=\tau.
\]
This proves \eqref{cdp:eq:C.NC.8} directly, with no base-change assumption.

\noindent\textbf{Step 2.} To verify the additional comparison, locally write $t=u(z_1\cdots z_r)^m$ with $u$ nowhere zero. The normalization of $s^m=u(z_1\cdots z_r)^m$ is given by
$w^m=u$ and $s=wz_1\cdots z_r$.
The derivative of $w^m-u$ with respect to $w$ is nonzero because $w$ is a unit. Thus its map to the original chart is étale of degree $m$, and the lifted central fibre has equation $z_1\cdots z_r=0$. These local normalizations glue by the universal property of normalization; their degree and étaleness are local properties. Restriction is therefore an étale degree-$m$ covering $S'\to S$, and each triple point has exactly $m$ preimages. A lifted triangulation proves the Euler-number comparison. If the cover has several connected components, the degrees and Euler numbers add over them. The two displayed comparisons remain valid without any connectedness assertion about that cover.
\end{proof}






\begin{proposition}[Reducible normal-crossings fibres with unequal multiplicities]\label{cdp:st:C.2.NC.4}
\StatementRecord{cdp:st:C.2.NC.4}{proved}{Historical events: Claude Fable 5.2 (Anthropic), effort high, 16 November 2026: carried as background without fault found. GPT-6 Astra (OpenAI), ChatGPT Pro mode, effort Pro, S6\_09\_PAPERS, 17 November 2026: targeted proof check. Also checked in this writing. Pre-continuation writing events: GPT-6 Astra (OpenAI), ChatGPT Pro mode, effort Pro, S6\_10\_PAPERS, 18 November 2026: targeted proof check. Continuation metadata events: GPT-6 Astra (OpenAI), ChatGPT Pro mode, effort Pro, S6\_10\_PAPERS continuation, 19 November 2026: used and rechecked in this run; status-vocabulary reconciliation; metadata reconciliation only; no new mathematical check. \emph{Pre-continuation recorded proof scope.} The classification and adjacent-divisor exclusion, coherent triple-point count, weighted nearby Euler relation, and the irreducible alternatives were re-derived}{the hypotheses, notation and definitions stated here; \cref{cdp:st:C.2.NC.3}; \cref{cdp:st:C.2.NC.2}; \cite[Theorem~2.2]{Peters2021}; \cite[Theorem~3.2]{BMS2012}}

Under the hypotheses on $g:Y\to C$ in \ref{cdp:st:C.2.NC.3}, let a fibre be $\sum_i m_iS_i$ with connected reducible normal-crossings support $S$. Its multiplicities may be unequal. Then
\[
 \tau(S)=0,\qquad e(S)=0.
 \tag{C.NC.9}\label{cdp:eq:C.NC.9}
\]
Every normalized conductor is a disjoint union of smooth elliptic curves, and every normalization component has Euler number zero. More generally, any singular normal-crossings fibre support under these hypotheses has either no triple point or exactly two triple points; in the latter case it is irreducible and its normalization is rational.
\end{proposition}


\begin{proof}
\noindent\textbf{Step 1.} First consider a normalization component $T_i$ whose conductor $D_i$ is nonempty. If $T_i$ is Kähler and has nonnegative Kodaira dimension, choose a nonzero section of $\omega_{T_i}^{\otimes n}$ and denote its effective zero divisor by $E$. The effective divisor $E+nD_i$ has rational Chern class zero by (C.NC.6). Its integral against a Kähler form is positive because $D_i$ is nonempty, whereas its cohomology class gives integral zero. This is impossible. A Kähler $T_i$ is consequently rational or a blown-up ruled surface. For a ruling $T_i\to B$ and a general ruling fibre $F$, one has $K_{T_i}\cdot F=-2$, hence $D_i\cdot F=2$. Some conductor curve therefore dominates $B$. Its normalization has genus at most one by \ref{cdp:st:C.2.NC.2}, so Riemann--Hurwitz forces $g(B)\le1$. The Kähler possibilities have
$\chi(\OO_{T_i})=1$ in the rational case and $0$ in the elliptic ruled case; their Euler numbers are nonnegative.

\noindent\textbf{Step 2.} For a non-Kähler $T_i$, the Enriques--Kodaira table \cite[Theorem~2.2]{Peters2021} leaves class VII, Kodaira and properly elliptic minimal models. Class VII has $c_1^2=-b_2$ and $c_2=b_2$, so Noether's formula gives $\chi(\OO)=0$. Kodaira surfaces have both Chern numbers zero. A minimal non-Kähler elliptic surface has topological Euler number zero by Kodaira's theorem in the precise form \cite[Theorem~3.2]{BMS2012}; its canonical square is zero by the same classification table, so Noether's formula again gives $\chi(\OO)=0$. Blowing up preserves the holomorphic Euler characteristic and adds one to the topological Euler characteristic. Thus every non-Kähler $T_i$ has $\chi(\OO_{T_i})=0$ and $e(T_i)\ge0$. This derives all numerical classification facts used here from the cited statements.

\noindent\textbf{Step 3.} A rational normalization component cannot meet a different global component of $S$. Indeed, for $j\ne i$, pull back the canonical section of $\OO_Y(S_j)$ to $T_i$. It is a nonzero section whose zeros are precisely the inverse image of $S_i\cap S_j$. Its line bundle has zero rational first Chern class. For a smooth rational surface $T_i$, $H^1(T_i,\OO)=0$ and $H^2(T_i,\Z)$ is torsion-free, by the corresponding facts for $\PP^2$ or a Hirzebruch surface and the blowup formula. The exponential sequence therefore makes that pulled-back line bundle trivial. A nonzero section of a trivial line bundle on a connected compact surface is a nonzero constant and has no zeros. Hence $S_i\cap S_j$ is empty. This proof places no restriction on intersections between two local branches belonging to the same $S_i$.

\noindent\textbf{Step 4.} The intersection graph of the global components of a connected reducible fibre is connected. Each component has nonempty conductor, and the preceding paragraph excludes rational normalizations. Thus every $\chi(\OO_{T_i})$ is zero, and (C.NC.4) gives $\tau=0$. Each $C_\alpha$ is now elliptic and has no marked point. The conductor covering is étale, so the normalized conductor components are elliptic; there are no nodes of $D_i$, since these would be preimages of triple points. Consequently $e(D_i)=0$. The weighted Euler formula now reads
\[
 0=e(Y_t)=\sum_i m_i e(T_i).
\]
All summands are nonnegative and every $m_i$ is positive, so each $e(T_i)=0$. Formula (C.NC.5) gives $e(S)=0$.

\noindent\textbf{Step 5.} For a singular irreducible support the conductor is nonempty and the same classification argument gives $\chi(\OO_{T_1})\in\{0,1\}$, with value one exactly in the rational case. Formula (C.NC.4) then gives $\tau=0$ or $2$, respectively. This proves the final assertion and also proves that any triple point lies on an irreducible support with rational normalization.
\end{proof}






\begin{proposition}[Global positivity and the all-twist normalization obstruction in the normal-crossings class]\label{cdp:st:C.2.NC.5}
\StatementRecord{cdp:st:C.2.NC.5}{proved}{Historical events: Claude Fable 5.2 (Anthropic), effort high, 16 November 2026: carried as background without fault found. GPT-6 Astra (OpenAI), ChatGPT Pro mode, effort Pro, S6\_09\_PAPERS, 17 November 2026: targeted proof check. Also checked in this writing. Pre-continuation writing events: GPT-6 Astra (OpenAI), ChatGPT Pro mode, effort Pro, S6\_10\_PAPERS, 18 November 2026: targeted proof check. Continuation metadata events: GPT-6 Astra (OpenAI), ChatGPT Pro mode, effort Pro, S6\_10\_PAPERS continuation, 19 November 2026: used and rechecked in this run; status-vocabulary reconciliation; metadata reconciliation only; no new mathematical check. \emph{Pre-continuation recorded proof scope.} The global Euler sum, parity of b3, rational-normalization Picard pullback, and the exact ambient cotangent dimension were re-derived}{the hypotheses, notation and definitions stated here; \cref{cdp:st:C.2.NC.3}; \cref{cdp:st:C.3.5}; \cref{cdp:st:C.3.1}; \cref{cdp:st:C.2.NC.4}; \cref{cdp:st:A.1.2}}

Let $g:Y\to C$ have the hypotheses of \ref{cdp:st:C.2.NC.3}, and assume every fibre reduction has normal crossings; smooth reductions are included. Then
\[
 c_3(Y)=e(Y)=\sum_{c\in C}\tau((Y_c)_{\rm red})=2r\ge0,
 \tag{C.NC.10}\label{cdp:eq:C.NC.10}
\]
where $r$ is the number of fibre supports with rational normalization. Every such support is irreducible, nonnormal, and has exactly two triple points. Equality $c_3(Y)=0$ holds exactly when every fibre support is free of triple points. If also $b_1(Y)=0$, then $r\in\{0,1\}$ and $b_3(Y)=2-2r$.

In the positive case there is an irreducible nonnormal reduced support $S$ with rational normalization $\eta:T\to S$ such that, for every $L\in\Pic(Y)$,
\[
 H^2\!\left(T,K_T\otimes\eta^*((K_Y^{-1}\otimes L)|_S)\right)\simeq\C,
 \qquad
 H^0\!\left(S,\Omega_Y^1|_S\otimes(K_Y\otimes L^{-1})|_S\right)\simeq\C.
 \tag{C.NC.11}\label{cdp:eq:C.NC.11}
\]
The scheme-theoretic fibre may be a multiple of $S$; this conclusion does not assert that its multiplicity is one. By \ref{cdp:st:C.3.5}, all these conclusions apply to every fibration under the CDP hypotheses whose fibre reductions have normal crossings, for every smooth-fibre type in \ref{cdp:st:C.3.1}.
\end{proposition}


\begin{proof}
\noindent\textbf{Step 1.} There are finitely many critical values: the critical locus is analytic, its image is analytic by properness, and it does not contain a general point of the curve. Away from these values $g$ is a differentiable fibre bundle whose fibre has Euler number zero. Euler characteristic is additive over this open base and the finite exceptional set, so $e(Y)=\sum_c e((Y_c)_{\rm red})$. Nilpotents leave the underlying space unchanged. An irreducible support has a single multiplicity, so \ref{cdp:st:C.2.NC.3} applies; a reducible support is covered by \ref{cdp:st:C.2.NC.4}. A smooth support has Euler number zero directly from (C.NC.1). Thus every summand equals its triple-point count. \ref{cdp:st:C.2.NC.4} identifies the positive summands and their value two. Gauss--Bonnet gives $c_3(Y)=e(Y)$. Poincaré duality gives $e(Y)=2-2b_1(Y)+2b_2(Y)-b_3(Y)$, proving the assertion about $r$ and $b_3$. \ref{cdp:st:C.3.5} proves the required zero Euler number for every smooth fibre under the CDP hypotheses, so that application uses no assumption that the general fibre is a torus.

\noindent\textbf{Step 2.} For a positive support, every line bundle on $Y$ pulls back trivially to its smooth rational normalization: its rational first Chern class is zero, and the injectivity and torsion-freeness used in \ref{cdp:st:C.2.NC.4} apply. In particular $\eta^*((K_Y^{-1}\otimes L)|_S)$ is trivial. Serre duality on the connected smooth compact surface $T$ identifies the first group in \eqref{cdp:eq:C.NC.11} with $H^0(T,\OO_T)^\vee\simeq\C$.

\noindent\textbf{Step 3.} For the ambient group set $A=(K_Y\otimes L^{-1})|_S$ and $N^*=\OO_Y(-S)|_S$. Both $\eta^*A$ and $\eta^*N^*$ are trivial, even when the full fibre is $mS$. The injection of torsion-free differentials into normalization differentials gives
$H^0(S,\widetilde\Omega_S^1\otimes A)\hookrightarrow H^0(T,\Omega_T^1)=0$.
The canonical normal-crossings kernel isomorphism \ref{cdp:st:A.1.2} therefore gives
\[
 H^0(S,\Omega_Y^1|_S\otimes A)
 =H^0(T,\eta^*(N^*\otimes A))=\C.
\]
This proves the second equality. It establishes the universal normalization target requested in B.3 for the full normal-crossings class, including unequal multiplicities. It leaves the unrestricted nonordinary-fibre assertion to be decided separately.
\end{proof}






\begin{proposition}[The normal jets along cusp, tacnode and triple curves]\label{cdp:st:C2d}
\StatementRecord{cdp:st:C2d}{proved}{Historical events: Claude Fable 5.2 (Anthropic), effort high, 16 November 2026: carried as background without fault found. GPT-6 Astra (OpenAI), ChatGPT Pro mode, effort Pro, S6\_09\_PAPERS, 17 November 2026: targeted proof check. Also checked in this writing. Pre-continuation writing events: GPT-6 Astra (OpenAI), ChatGPT Pro mode, effort Pro, S6\_10\_PAPERS, 18 November 2026: targeted proof check. Continuation metadata events: GPT-6 Astra (OpenAI), ChatGPT Pro mode, effort Pro, S6\_10\_PAPERS continuation, 19 November 2026: used and rechecked in this run; status-vocabulary reconciliation; metadata reconciliation only; no new mathematical check. \emph{Pre-continuation recorded proof scope.} The normal-jet line weights for cusp, tacnode, and triple germs, the degree 2g-2 of the normal line, the equality case, and the inverse-Jacobian pole calculation were re-derived}{the hypotheses, notation and definitions stated here; \cref{cdp:eq:C.J.2}; \cref{cdp:eq:C.J.3}; \cref{cdp:st:A.2.4}}

Let $Y$ be a smooth complex threefold with $b_2(Y)=0$, let $S$ be a reduced Cartier divisor, and let $D\subset\operatorname{Sing}S$ be a smooth compact connected curve. Suppose the general plane section transverse to $D$ has one of the equations
\[
 x^2-y^3=0,\qquad x^2-y^4=0,\qquad xy(x-y)=0.
 \tag{C.J.1}\label{cdp:eq:C.J.1}
\]
Then $g(D)\le1$. For $g(D)=1$, the indicated transverse analytic type is constant along all of $D$. The inverse image of $D$ in the normalization is respectively an isomorphic copy, a finite etale cover of degree two, or a finite etale cover of degree three. Each connected component $C$ of that cover has genus one and normal line of degree zero in the normalization, which is smooth near $C$. At these three transverse types, the inverse Jacobian ideal pulls back into the line allowing at most a simple pole on each normalization branch.
\end{proposition}


\begin{proof}
\noindent\textbf{Step 1.} Put $N=N_{D/Y}$ and $B=\OO_Y(S)|_D$. Since every ambient line bundle has zero rational first Chern class, $\deg B=\deg K_Y|_D=0$. Adjunction for the smooth curve gives
\[
 \det N=K_D\otimes K_Y^{-1}|_D,\qquad
 \deg N=2g(D)-2.                                      \tag{C.J.2}\label{cdp:eq:C.J.2}
\]
A local defining section of $S$ belongs to $I_D^2$. Its quadratic normal jet is a well-defined section of $\operatorname{Sym}^2N^*\otimes B$. In the first two cases its generic rank is one. Its saturated radical $R\subset N$ is a line subbundle on the smooth curve, and $Q=N/R$ is a line bundle. The quadratic jet gives a nonzero section of $B\otimes Q^{-2}$, so
$\deg Q\le0$.

\noindent\textbf{Step 2.} In the cusp case, the cubic jet on the radical is an intrinsic nonzero section of $B\otimes R^{-3}$. To check its independence, choose coordinates $(x,y,z)$ with $D=(x=y=0)$ and $R$ the $y$ direction. A change of lift of a radical vector changes the cubic by a Hessian term with one radical argument, which is zero. Changing the local retraction to $D$ differentiates the identity that the quadratic jet vanishes on $R$, again giving zero. Thus the local cubics glue across all points, including possible zeros of the quadratic coefficient. Consequently $\deg R\le0$. Equation \eqref{cdp:eq:C.J.2} proves $g(D)\le1$.

\noindent\textbf{Step 3.} In the tacnode case the radical cubic vanishes identically. Write the relevant normal terms as
\[
 a x^2+bxy^2+cy^4.
\]
Completing the square has invariant numerator
\[
 \Delta=ac-b^2/4\in
 H^0(D,B^2\otimes Q^{-2}\otimes R^{-4}).               \tag{C.J.3}\label{cdp:eq:C.J.3}
\]
This numerator is holomorphic even at a zero of $a$. Under $x\mapsto x+uy^2$, the coefficients become $b+2au$ and $c+bu+au^2$, and \eqref{cdp:eq:C.J.3} is unchanged. Linear changes preserving the radical give precisely the displayed line-bundle weights. Changes of the retraction to $D$ add terms containing the radical cubic or a radical Hessian argument, so also leave this quartic numerator unchanged. These checks give its global definition. It is nonzero by the generic tacnode assumption. Hence
\[
 \deg Q\le0,\qquad \deg Q+2\deg R\le0,
 \qquad \deg N\le\tfrac12\deg Q\le0.
\]
This proves the same genus bound.

\noindent\textbf{Step 4.} For the ordinary triple case the quadratic jet is identically zero. The binary cubic normal jet is a section of $\operatorname{Sym}^3N^*\otimes B$. Its discriminant is a nonzero section of
$B^4\otimes(\det N^*)^6$. The discriminant has degree four in the cubic coefficients and weight six under a linear change of the two normal coordinates; these weights also follow by writing it as the square of the product of the three pairwise determinants of the linear factors. Its nonnegative divisor degree gives $\deg N\le0$ and the genus bound again.

\noindent\textbf{Step 5.} When $g(D)=1$, all the inequalities are equalities. In the cusp case both the quadratic and radical cubic sections have no zeros. In the tacnode case $a$ and $\Delta$ have no zeros. In the triple case the cubic discriminant has no zeros. The holomorphic splitting lemma with parameter $z$ completes the square in the first two cases. The remaining one-variable term is respectively a unit times $y^3$ or $y^4$; taking a local root of the unit gives the corresponding equation in \eqref{cdp:eq:C.J.1}. In the triple case its three distinct tangent lines split after an etale ordering cover; the holomorphic implicit function theorem gives three smooth branches meeting transversely in each normal plane. Straightening two branches to the axes and rescaling the third branch's unit gives the third equation. Thus these are constant transverse types throughout $D$. In particular no other singular curve meets $D$.

\noindent\textbf{Step 6.} In the cusp normalization, $x=u^3,y=u^2$, up to the choice of which coordinate is the radical. Its normal line $P=N_{C/\widetilde S}$ satisfies $P^2=R$ and $P^3=Q$, so has degree zero. For the tacnode, the two branches have $x=\pm y^2$ and normal line the pullback of $R$, again of degree zero. Their permutation cover is etale because $\Delta$ is nowhere zero. For the ordinary triple, order the three tangent lines by a finite etale cover. Three pairwise transverse line subbundles in a rank-two bundle make it a sum $P\oplus P$: the third is a graph of an isomorphism between the first two. Since the determinant has degree zero, $\deg P=0$. Each normalization branch has this tangent-root normal line. Etale descent gives the claimed degrees and covers before ordering.

\noindent\textbf{Step 7.} Finally the inverse-Jacobian assertion can be checked on these exact transverse equations. For the cusp, $R=\C\{u^2,u^3\}$ and $\mathcal J^{-1}=u^{-1}\C\{u\}$, as in \ref{cdp:st:A.2.4}. For the tacnode, normalization functions are pairs whose constant and linear terms agree. A putative pole of order two in a fractional multiplier has leading coefficients $(a_1,a_2)$. Multiplication by $2x$ forces $a_1=-a_2$ at order zero, while multiplication by $-4y^3$ forces $a_1=a_2$ at order one. Thus both vanish. For the ordinary triple, use branches $(0,u),(u,0),(u,u)$. The partial derivatives have leading coefficient vectors $(-1,0,1)$ and $(0,1,-1)$ at order two. A pole of order two would make both products constant on all three branches, forcing all its leading coefficients to vanish. A pole of larger order is already excluded by regularity of these products on normalization. These computations prove the simple-pole bound. The covariance and leading-coefficient equations are checked exactly in \path{code/C_jets.py}.
\end{proof}







\section{Irreducible supports: topology, conductor and equality}
\label{sec:II:night-sign}
The following proofs retain every Part~4 result of the night run. The final classification is Theorem~\ref{night:P4.58}; its exact normalization maps and its local exclusions are proved below.
\paragraph{Standing torus-fibration setting.} The ambient space is a smooth compact complex threefold $Y$ with $b_1(Y)=b_2(Y)=0$, and $f:Y\to B$ is proper and surjective with connected fibres and smooth general fibre a complex two-torus. Every statement below concerning an irreducible support uses its actual full fibre $mS$, normalization and resolution. The global assertions expressly retaining the hypothesis that every support is irreducible or has normal crossings are used with that hypothesis. Constructions and comparison examples with different Betti numbers state their own assumptions.
\subsection{The explicit threefold with identity monodromy}

Throughout this part we orient complex spaces by their complex structures.
All homology groups in the example have coefficients in $\Z$. A meridian is
counterclockwise in its indicated holomorphic coordinate; at infinity this
coordinate is $1/t$. We write
\[
 E=\C/(\Z+i\Z)=\C^*/q^\Z,\qquad q=e^{-2\pi},\qquad
 A=E^2,\qquad \Lambda=\Z e_1\oplus\Z e_2\oplus\Z e_3\oplus\Z e_4,
\]
where $(e_1,e_2,e_3,e_4)=((1,0),(i,0),(0,1),(0,i))$. The comparison
$E\longrightarrow\C^*/q^\Z$ is $z\mapsto e^{2\pi iz}$.

\begin{proposition}[P4.1]\label{night:P4.1}\label{stmt:P4.1}
\StatementRecord{night:P4.1}{proved}{Historical events: GPT-6 Astra (OpenAI), ChatGPT Pro mode, effort Pro, S6\_06\_BLIND, 16--18 November 2026: targeted proof check. Also checked in this writing. Pre-continuation writing events: GPT-6 Astra (OpenAI), ChatGPT Pro mode, effort Pro, S6\_10\_PAPERS, 18 November 2026: targeted proof check; used and rechecked in this run. Continuation metadata events: GPT-6 Astra (OpenAI), ChatGPT Pro mode, effort Pro, S6\_10\_PAPERS continuation, 19 November 2026: used and rechecked in this run; status-vocabulary reconciliation; metadata reconciliation only; no new mathematical check. \emph{Pre-continuation recorded proof scope.} The night-run re-derivation was by GPT-6 Astra (OpenAI, ChatGPT Pro mode, effort Pro), 16--18 November 2026. Also checked in this writing by GPT-6 Astra (OpenAI), ChatGPT Pro mode, effort Pro, 18 November 2026. The two actual blowup charts and line frames give lambda\_B=v lambda\_A, and the canonical section gives the punctured product map lambda\_A/u=lambda\_B/s. The scalar quotient is free, proper and Hausdorff. Tautological evaluation identifies the exceptional nonzero line with C2 minus zero equivariantly. The other component has primitive circle Chern class, b1=3 and trivial canonical line, matching the currently opened BHPV classification table. The map uv gives the exact reduced SNC central divisor}{\ref{night:P4.3}}

Let $V=\operatorname{Bl}_{(0,0)}(B\times\Delta)$, where $B=E$, let $R$ be
the exceptional curve, and put $L=\OO_V(R)$. The holomorphic quotient
$U=L^\times/q^\Z$ is a smooth Hausdorff threefold, proper over $\Delta$.
Its noncentral fibres are $E\times B$. Its central fibre is the reduced
union of a primary Kodaira surface and the scalar Hopf surface
$(\C^2\setminus\{0\})/q^\Z$, meeting transversally in one elliptic curve.
\end{proposition}
\begin{proof}
\noindent\textbf{Step 1.} We use coordinates $(z,t)$ at the blown-up point. The blowup has charts
\[
 (z,t)=(u,uv),\qquad (z,t)=(rs,s),\qquad r=v^{-1},\quad s=uv.
\]
The frames $u^{-1}$ and $s^{-1}$ for $L$ give fibre coordinates related by
$\lambda_B=v\lambda_A$. Each nonzero fibre coordinate, modulo multiplication
by $q$, gives a product chart with fibre $E$. The transitions are holomorphic
translations on $E$. The scalar action is free and properly discontinuous:
the logarithm of the norm of a nonzero fibre coordinate changes by the
nonzero constant $\log q$, and compact subsets meet only finitely many
translates. Thus the quotient is Hausdorff and smooth. These product charts
have compact fibre $E$, so $U\to V$ is proper. The blowup and $B\times\Delta
\to\Delta$ are proper, proving properness of their composite.

\noindent\textbf{Step 2.} The canonical section of $\OO(R)$ has coefficients $u$ and $s$. Over
$t\ne0$ it never vanishes. The ratios $\lambda_A/u=\lambda_B/s$ therefore
give the specified holomorphic product trivialization over $\Delta^*$.

\noindent\textbf{Step 3.} The central divisor in $V$ is $B'\cup R$, with $R^2=-1$ and $R\cdot B'=1$.
Thus $L|_R=\OO_{\Pp^1}(-1)$ and $L|_{B'}$ has degree one. Over $B'$ the
quotient is a principal elliptic bundle with primitive nonzero circle Chern
class. Its canonical bundle is trivial: the base differential wedged with
$d\lambda/\lambda$ is invariant under the line-bundle transition, because
the additional term is pulled back from the one-dimensional base. The
circle-bundle calculation in P4.3 gives $b_1=3$. The surface classification
The classification \cite[Chapter~VI, Theorem~1.1 and Table~10, pp.~243--244]{BHPV2004E} consequently identifies this surface as
a primary Kodaira surface.

\noindent\textbf{Step 4.} The tautological evaluation morphism
$\OO_{\Pp^1}(-1)^\times\to\C^2\setminus\{0\}$ is a biholomorphism: a
nonzero vector uniquely determines its line. It intertwines both scalar
actions, and hence identifies the other component with the stated Hopf
surface. Over $B'\cap R$ the components meet in $E$. In local product
coordinates the map is $(u,v,c)\mapsto uv$; this proves transversality,
multiplicity one, and the absence of other critical points.
\end{proof}

\begin{proposition}[P4.2]\label{night:P4.2}\label{stmt:P4.2}
\StatementRecord{night:P4.2}{proved}{Historical events: GPT-6 Astra (OpenAI), ChatGPT Pro mode, effort Pro, S6\_06\_BLIND, 16--18 November 2026: targeted proof check. Also checked in this writing. Pre-continuation writing events: GPT-6 Astra (OpenAI), ChatGPT Pro mode, effort Pro, S6\_10\_PAPERS, 18 November 2026: targeted proof check. Continuation metadata events: GPT-6 Astra (OpenAI), ChatGPT Pro mode, effort Pro, S6\_10\_PAPERS continuation, 19 November 2026: used and rechecked in this run; status-vocabulary reconciliation; metadata reconciliation only; no new mathematical check. \emph{Pre-continuation recorded proof scope.} The night-run re-derivation was by GPT-6 Astra (OpenAI, ChatGPT Pro mode, effort Pro), 16--18 November 2026. Also checked in this writing by GPT-6 Astra (OpenAI), ChatGPT Pro mode, effort Pro, 18 November 2026. The logarithm divided by 2pi i defines a single-valued elliptic-valued transition because its increment is the integral lattice generator. The three displayed holomorphic group isomorphisms, their determinant-one real lattice matrices, local properness and point separation give the smooth compact Hausdorff total space. Their punctured trivializations make the remaining return an actual translation, which acts identically on integral H1}{\ref{night:P4.1}}

The three markings and the principal-bundle clutch in the prompt define a
smooth compact complex threefold $f:Y\to\Pp^1$ with connected fibres,
exactly three singular fibres of the form in P4.1, and identity monodromy
on $H_1$ of the regular fibre.
\end{proposition}
\begin{proof}
\noindent\textbf{Step 1.} The map $\ell(t)=(\log t)/(2\pi i)$ with values in $E$ is single-valued and
holomorphic on an annulus: changing the logarithm adds an integer.
Consequently $a_{\rm out}=a_{\rm in}+(\ell(t),0)$ is a holomorphic transition
for a principal $A$-bundle on $\Pp^1$. Product charts with compact fibre
prove that this bundle is Hausdorff and compact.

\noindent\textbf{Step 2.} We choose disjoint discs about $0,1,\infty$ and replace their product models
by P4.1 using the holomorphic group isomorphisms
\[
 C_i\times B_i\longrightarrow A:\qquad
 (z,w)\mapsto(w,z),\quad(iz,iz+w),\quad(z,iz+w),
\]
respectively. Each lattice matrix $S_i$ displayed below has determinant
one. These maps, $\lambda_B=v\lambda_A$, the canonical punctured-disc
trivialization in P4.1, and the clutch $\ell$ specify all transitions.
Different base points separate in the base. Points over the same base
point lie in one Hausdorff local model, which separates them. The gluing is
therefore Hausdorff. Properness holds over each base chart and hence over
the base; compactness follows.

\noindent\textbf{Step 3.} The two central components have connected elliptic fibres over connected
curves and meet each other. All remaining fibres equal $A$. The equations
$uv=t$ give exactly the indicated critical values. The punctured-disc
trivializations identify all fibre lattices with $\Lambda$, and the remaining
transition is a translation on $A$. Translation induces the identity on
$H_1(A;\Z)$, giving the asserted monodromy representation.
\end{proof}

\begin{proposition}[P4.3]\label{night:P4.3}\label{stmt:P4.3}
\StatementRecord{night:P4.3}{proved}{Historical events: GPT-6 Astra (OpenAI), ChatGPT Pro mode, effort Pro, S6\_06\_BLIND, 16--18 November 2026: targeted proof check. Also checked in this writing. Pre-continuation writing events: GPT-6 Astra (OpenAI), ChatGPT Pro mode, effort Pro, S6\_10\_PAPERS, 18 November 2026: targeted proof check. Continuation metadata events: GPT-6 Astra (OpenAI), ChatGPT Pro mode, effort Pro, S6\_10\_PAPERS continuation, 19 November 2026: used and rechecked in this run; status-vocabulary reconciliation; metadata reconciliation only; no new mathematical check. \emph{Pre-continuation recorded proof scope.} The night-run re-derivation was by GPT-6 Astra (OpenAI, ChatGPT Pro mode, effort Pro), 16--18 November 2026. Also checked in this writing by GPT-6 Astra (OpenAI), ChatGPT Pro mode, effort Pro, 18 November 2026. The blowup is the explicit Hopf-boundary homotopy pushout B wedge S2. The circle Chern class is zero on a regular B and minus one on the exceptional sphere, yielding the pushout T3 from S1 to S3. After its product with the norm circle, the only cellular boundaries are r to a and rc to a wedge c. This proves the two precise inclusion quotients and free ranks. The canonical section trivializes the stated meridian bounding chains, including every mixed product}{the hypotheses and explicit local calculation stated here}

Order the local fibre lattice as $(a,c,b,d)$, where $(a,c)$ are the cycles
$(1,i)$ in $C_i$. The regular-fibre inclusions induce the quotient maps
\[
 H_1(A)=\Lambda\longrightarrow\Lambda/\Z a=H_1(U),\qquad
 H_2(A)=\bigwedge^2\Lambda\longrightarrow
 \bigwedge^2\Lambda/\Z(a\wedge c)=H_2(U).
\]
Both target groups are free. The meridian and every product of the
meridian with a fibre one-cycle map to zero in these degrees.
\end{proposition}
\begin{proof}
\noindent\textbf{Step 1.} A Hermitian metric gives a diffeomorphism
$L^\times/q^\Z\cong S(L)\times S^1_c$, with factors the circle bundle and
the logarithm of the norm modulo $\log q$. We first specify a homotopy
model for $V$. Blowing up a four-ball replaces it by the disc bundle of
$\OO(-1)$ along its Hopf boundary. The complement of the ball in
$B\times\Delta$ has homotopy type $B\vee S^3$; the boundary sphere is
the $S^3$ summand, since its map to $B$ is nullhomotopic. The resulting
homotopy pushout is $B\vee S^2$. Equivalently, on simply connected
covers the replacement adds one sphere at each lift of the centre and
no other homology; the comparison is an isomorphism on $\pi_1$ and on
the homology of those covers, and hence a homotopy equivalence. We may
choose the $B$ summand to be a regular $B_t$ and the sphere to be $R$,
joining their base points by a path.

\noindent\textbf{Step 2.} The circle Chern class is zero on $B_t$ and $-1$ on $R$. Therefore
$S(L)$ is represented by the homotopy pushout
\[
 T^3_{(a,b,d)}\longleftarrow S^1_a\longrightarrow S^3.
\]
The second map is a Hopf fibre and is nullhomotopic. Thus this pushout
has the homotopy type $(T^3\cup_aD^2)\vee S^3$. We multiply by $S^1_c$.
The resulting integral cellular complex starts with the exterior-cell
complex of $T^4$, ordered by $(a,c,b,d)$, and adds cells $r$ in degree
two, $rc,h$ in degree three, and $hc$ in degree four. Its only nonzero
differentials are
\[
 \partial r=a,\qquad \partial(rc)=a\wedge c.
\]
We choose the orientation of $r$ to give the first equality and use the
product orientation for $rc$. The complete differential matrices are
printed below. The complex has free homology of ranks $(1,3,5,5,2)$;
the inclusion of the torus subcomplex gives exactly the two quotient maps.

\noindent\textbf{Step 3.} We represent a base point, and both base generators $b,d$, away from
$0\in B$. Their products with the disc miss $R$, so the canonical section
trivializes $L$ over the resulting chains. The meridian and the products
with $b,d$ consequently bound. The products with $a,c$ bound over one
such base point. This proves all the remaining assertions on the maps.
\end{proof}

The lattice matrices, with columns $(a_i,c_i,b_i,d_i)$ in the common basis,
are
\[
S_1=\begin{pmatrix}0&0&1&0\\0&0&0&1\\1&0&0&0\\0&1&0&0\end{pmatrix},\quad
S_2=\begin{pmatrix}0&-1&0&0\\1&0&0&0\\0&-1&1&0\\1&0&0&1\end{pmatrix},\quad
S_3=\begin{pmatrix}1&0&0&0\\0&1&0&0\\0&-1&1&0\\1&0&0&1\end{pmatrix}.
\]
Put $Q_i=S_i^{-1}$. Deleting the first row of $Q_i$ gives $R_{1,i}$.
The row indices of $\bigwedge^2Q_i$ are $(ac,ab,ad,cb,cd,bd)$ and its
column indices are $(12,13,14,23,24,34)$; deleting its first row gives
$R_{2,i}$. Thus these are the integer matrices of the local quotient maps,
with target bases $(c_i,b_i,d_i)$ and $(a_ib_i,a_id_i,c_ib_i,c_id_i,b_id_i)$.

\begin{proposition}[P4.4]\label{night:P4.4}\label{stmt:P4.4}
\StatementRecord{night:P4.4}{proved}{Historical events: GPT-6 Astra (OpenAI), ChatGPT Pro mode, effort Pro, S6\_06\_BLIND, 16--18 November 2026: targeted proof check. Also checked in this writing. Pre-continuation writing events: GPT-6 Astra (OpenAI), ChatGPT Pro mode, effort Pro, S6\_10\_PAPERS, 18 November 2026: targeted proof check. GPT-6 Astra (OpenAI), ChatGPT Pro mode, effort Pro, S6\_10\_PAPERS, 18 November 2026 05:30:45.837959+00:00: computation or formal certificate. Continuation metadata events: GPT-6 Astra (OpenAI), ChatGPT Pro mode, effort Pro, S6\_10\_PAPERS continuation, 19 November 2026: used and rechecked in this run; status-vocabulary reconciliation; metadata reconciliation only; no new mathematical check. \emph{Pre-continuation recorded proof scope.} The night-run re-derivation was by GPT-6 Astra (OpenAI, ChatGPT Pro mode, effort Pro), 16--18 November 2026. Also checked in this writing by GPT-6 Astra (OpenAI), ChatGPT Pro mode, effort Pro, 18 November 2026. Reconstructed F1 and F2 directly from the three holomorphic marking matrices and logarithmic clutch, including every e\_j wedge g\_i mixed term and the two base meridian rows. The current independent replay gives determinant F1=-1, fifteen unit Smith factors for F1, twenty-nine for F2, and reduced relation determinants -1 and 1. The Mayer-Vietoris exact sequence and explicit van Kampen relations then give integral H1=H2=0 and trivial pi1. The computation scope is recorded separately in analysis/sign\_P44\_integral\_replay.json}{the hypotheses and explicit local calculation stated here}

The threefold $Y$ has $H_1(Y;\Z)=H_2(Y;\Z)=0$ and $\pi_1(Y)=1$.
\end{proposition}
\begin{proof}
\noindent\textbf{Step 1.} Let $P$ be the complement of three small open base discs. The remaining
bundle $W=f^{-1}(P)$ is trivial. We use $W=A\times P$, let $s_1,s_2$
generate $H_1(P)$, and set $s_3=-s_1-s_2$. The overlaps are
$N_i=A\times S^1$. Positive continuation of $\ell$ around zero adds $e_1$.
With
\[
 g_1=e_1,\quad g_2=g_3=0,\qquad
 v_1=\binom10,\quad v_2=\binom01,\quad v_3=\binom{-1}{-1},
\]
the overlap maps into $W$ are
\[
 x\mapsto x,\quad s_i\mapsto g_i+v_i,\qquad
 x\wedge y\mapsto x\wedge y,\quad
 x\wedge s_i\mapsto x\wedge g_i+x\otimes v_i.
\]
These are the maps represented by the products of the cycles, or by the
induced maps into the aspherical space $A\times P$; in particular they
retain the degree-two contribution of the clutch.

\noindent\textbf{Step 2.} We use the ordered basis $(e_1,e_2,e_3,e_4,s_1,s_2)$ for $H_1(W)$ and
append the three bases $(c_i,b_i,d_i)$ for the local groups. On the source
we concatenate the three bases $(e_1,e_2,e_3,e_4,s_i)$. The resulting
$15\times15$ Mayer--Vietoris matrix $F_1$ has, in its $i$th five-column
block, upper block
\[
 \begin{pmatrix}I_4&g_i\\0_{2\times4}&v_i\end{pmatrix}
\]
and local block $(-R_{1,i},0)$ in the $i$th summand; its other local
blocks vanish. The sign is the usual difference of the two inclusions.

\noindent\textbf{Step 3.} For $H_2(W)$ we use
\[
 (12,13,14,23,24,34,e_1s_1,\ldots,e_4s_1,e_1s_2,\ldots,e_4s_2)
\]
and append the three local five-element bases. The source is the
concatenation of $(12,13,14,23,24,34,e_1s_i,\ldots,e_4s_i)$.
The $29\times30$ matrix $F_2$ has top six rows
$[I_6,(e_j\wedge g_i)_{j=1}^4]$ in its $i$th block; its next eight rows
have the coefficients of $e_j\otimes v_i$ in the displayed order, and
its $i$th local block is $[-R_{2,i},0]$. All other blocks vanish.
Finally
\[
 F_0=\begin{pmatrix}1&1&1\\-1&0&0\\0&-1&0\\0&0&-1\end{pmatrix}
\]
has zero kernel. Thus the required exact sequence is
\[
 \Z^{30}\xrightarrow{F_2}\Z^{29}\longrightarrow H_2(Y)
 \longrightarrow\Z^{15}\xrightarrow{F_1}\Z^{15}\longrightarrow H_1(Y)
 \longrightarrow\Z^3\xrightarrow{F_0}\Z^4.
\]

\noindent\textbf{Step 4.} We verify integral surjectivity by hand. Eliminating the local $H_1$
groups and two independent base meridians leaves $\Lambda$ modulo
$e_1,a_1,a_2,a_3$. The relation matrix is
\[
 R_{H_1}=\begin{pmatrix}1&0&0&1\\0&0&1&0\\0&1&0&0\\0&0&1&1\end{pmatrix},
 \qquad\det R_{H_1}=-1.
\]
Hence $F_1$ is an integral isomorphism. In degree two the fibre
$\bigwedge^2\Lambda$ columns eliminate the fifteen local generators,
leaving the three relations $\eta_i=a_i\wedge c_i$. The second meridian
columns kill $\Lambda\otimes s_2$, the third then kill
$\Lambda\otimes s_1$, and the first leave $e_1\wedge\Lambda$. The six
relations, in the order $e_{12},e_{13},e_{14},\eta_1,\eta_2,\eta_3$, have
matrix
\[
R_{H_2}=\begin{pmatrix}
1&0&0&0&1&1\\0&1&0&0&0&-1\\0&0&1&0&1&0\\
0&0&0&0&-1&0\\0&0&0&0&0&-1\\0&0&0&1&1&1
\end{pmatrix},\qquad \det R_{H_2}=1.
\]
This proves surjectivity of $F_2$ over $\Z$ and excludes torsion. The
executed independent Smith calculation gives fifteen unit factors for
$F_1$ and twenty-nine unit factors for $F_2$, and computes $\det F_1=-1$.

\noindent\textbf{Step 5.} For the fundamental group the local pushout gives
$\pi_1(U_i)=\Lambda/\Z a_i$, with trivial meridian. Van Kampen starts
from $\pi_1(W)=\Lambda\times F(s_1,s_2)$ and imposes
\[
 a_1=a_2=a_3=0,\qquad s_1+e_1=0,\quad s_2=0,\quad s_1+s_2=0.
\]
Both base generators and $e_1$ vanish. The same unimodular matrix
$R_{H_1}$ then kills the fibre lattice, proving $\pi_1(Y)=1$.
\end{proof}

The program \texttt{code/part4/example\_mayer\_vietoris.py} recomputes all
entries, differentials, determinants and Smith forms exactly. Its output
and machine-readable matrices are supplied alongside it. For completeness
we print every integer matrix, including the zero differentials:
% Generated by code/part4/matrices_to_tex.py from example_matrices.json.
\begingroup
\setlength{\arraycolsep}{1.7pt}
\renewcommand{\arraystretch}{0.88}
\par\smallskip\noindent
$S_1$ has size $4\times 4$; its full integer matrix is
\[
\resizebox{!}{11.2mm}{$\displaystyle S_1=\begin{pmatrix}
0 & 0 & 1 & 0\\
0 & 0 & 0 & 1\\
1 & 0 & 0 & 0\\
0 & 1 & 0 & 0\\
\end{pmatrix}$}
\]
\par\smallskip\noindent
$S_2$ has size $4\times 4$; its full integer matrix is
\[
\resizebox{!}{11.2mm}{$\displaystyle S_2=\begin{pmatrix}
0 & -1 & 0 & 0\\
1 & 0 & 0 & 0\\
0 & -1 & 1 & 0\\
1 & 0 & 0 & 1\\
\end{pmatrix}$}
\]
\par\smallskip\noindent
$S_3$ has size $4\times 4$; its full integer matrix is
\[
\resizebox{!}{11.2mm}{$\displaystyle S_3=\begin{pmatrix}
1 & 0 & 0 & 0\\
0 & 1 & 0 & 0\\
0 & -1 & 1 & 0\\
1 & 0 & 0 & 1\\
\end{pmatrix}$}
\]
\par\smallskip\noindent
$R_{1,1}$ has size $3\times 4$; its full integer matrix is
\[
\resizebox{!}{8.399999999999999mm}{$\displaystyle R_{1,1}=\begin{pmatrix}
0 & 0 & 0 & 1\\
1 & 0 & 0 & 0\\
0 & 1 & 0 & 0\\
\end{pmatrix}$}
\]
\par\smallskip\noindent
$R_{1,2}$ has size $3\times 4$; its full integer matrix is
\[
\resizebox{!}{8.399999999999999mm}{$\displaystyle R_{1,2}=\begin{pmatrix}
-1 & 0 & 0 & 0\\
-1 & 0 & 1 & 0\\
0 & -1 & 0 & 1\\
\end{pmatrix}$}
\]
\par\smallskip\noindent
$R_{1,3}$ has size $3\times 4$; its full integer matrix is
\[
\resizebox{!}{8.399999999999999mm}{$\displaystyle R_{1,3}=\begin{pmatrix}
0 & 1 & 0 & 0\\
0 & 1 & 1 & 0\\
-1 & 0 & 0 & 1\\
\end{pmatrix}$}
\]
\par\smallskip\noindent
$R_{2,1}$ has size $5\times 6$; its full integer matrix is
\[
\resizebox{!}{14.0mm}{$\displaystyle R_{2,1}=\begin{pmatrix}
0 & -1 & 0 & 0 & 0 & 0\\
0 & 0 & 0 & -1 & 0 & 0\\
0 & 0 & -1 & 0 & 0 & 0\\
0 & 0 & 0 & 0 & -1 & 0\\
1 & 0 & 0 & 0 & 0 & 0\\
\end{pmatrix}$}
\]
\par\smallskip\noindent
$R_{2,2}$ has size $5\times 6$; its full integer matrix is
\[
\resizebox{!}{14.0mm}{$\displaystyle R_{2,2}=\begin{pmatrix}
1 & 0 & 0 & 1 & 0 & 0\\
0 & 0 & 0 & 0 & 1 & 0\\
0 & -1 & 0 & 0 & 0 & 0\\
1 & 0 & -1 & 0 & 0 & 0\\
1 & 0 & -1 & 1 & 0 & 1\\
\end{pmatrix}$}
\]
\par\smallskip\noindent
$R_{2,3}$ has size $5\times 6$; its full integer matrix is
\[
\resizebox{!}{14.0mm}{$\displaystyle R_{2,3}=\begin{pmatrix}
1 & 1 & 0 & 0 & 0 & 0\\
0 & 0 & 1 & 0 & 0 & 0\\
0 & 0 & 0 & 1 & 0 & 0\\
1 & 0 & 0 & 0 & 1 & 0\\
1 & 1 & 0 & 0 & 1 & 1\\
\end{pmatrix}$}
\]
\par\smallskip\noindent
$\partial_1$ has size $1\times 4$; its full integer matrix is
\[
\resizebox{!}{2.8mm}{$\displaystyle \partial_1=\begin{pmatrix}
0 & 0 & 0 & 0\\
\end{pmatrix}$}
\]
\par\smallskip\noindent
$\partial_2$ has size $4\times 7$; its full integer matrix is
\[
\resizebox{!}{11.2mm}{$\displaystyle \partial_2=\begin{pmatrix}
0 & 0 & 0 & 0 & 0 & 0 & 1\\
0 & 0 & 0 & 0 & 0 & 0 & 0\\
0 & 0 & 0 & 0 & 0 & 0 & 0\\
0 & 0 & 0 & 0 & 0 & 0 & 0\\
\end{pmatrix}$}
\]
\par\smallskip\noindent
$\partial_3$ has size $7\times 6$; its full integer matrix is
\[
\resizebox{!}{19.599999999999998mm}{$\displaystyle \partial_3=\begin{pmatrix}
0 & 0 & 0 & 0 & 1 & 0\\
0 & 0 & 0 & 0 & 0 & 0\\
0 & 0 & 0 & 0 & 0 & 0\\
0 & 0 & 0 & 0 & 0 & 0\\
0 & 0 & 0 & 0 & 0 & 0\\
0 & 0 & 0 & 0 & 0 & 0\\
0 & 0 & 0 & 0 & 0 & 0\\
\end{pmatrix}$}
\]
\par\smallskip\noindent
$\partial_4$ has size $6\times 2$; its full integer matrix is
\[
\resizebox{!}{16.799999999999997mm}{$\displaystyle \partial_4=\begin{pmatrix}
0 & 0\\
0 & 0\\
0 & 0\\
0 & 0\\
0 & 0\\
0 & 0\\
\end{pmatrix}$}
\]
\par\smallskip\noindent
$F_0$ has size $4\times 3$; its full integer matrix is
\[
\resizebox{!}{11.2mm}{$\displaystyle F_0=\begin{pmatrix}
1 & 1 & 1\\
-1 & 0 & 0\\
0 & -1 & 0\\
0 & 0 & -1\\
\end{pmatrix}$}
\]
\par\smallskip\noindent
$F_1$ has size $15\times 15$; its full integer matrix is
\[
\resizebox{0.94\linewidth}{!}{$\displaystyle F_1=\begin{pmatrix}
1 & 0 & 0 & 0 & 1 & 1 & 0 & 0 & 0 & 0 & 1 & 0 & 0 & 0 & 0\\
0 & 1 & 0 & 0 & 0 & 0 & 1 & 0 & 0 & 0 & 0 & 1 & 0 & 0 & 0\\
0 & 0 & 1 & 0 & 0 & 0 & 0 & 1 & 0 & 0 & 0 & 0 & 1 & 0 & 0\\
0 & 0 & 0 & 1 & 0 & 0 & 0 & 0 & 1 & 0 & 0 & 0 & 0 & 1 & 0\\
0 & 0 & 0 & 0 & 1 & 0 & 0 & 0 & 0 & 0 & 0 & 0 & 0 & 0 & -1\\
0 & 0 & 0 & 0 & 0 & 0 & 0 & 0 & 0 & 1 & 0 & 0 & 0 & 0 & -1\\
0 & 0 & 0 & -1 & 0 & 0 & 0 & 0 & 0 & 0 & 0 & 0 & 0 & 0 & 0\\
-1 & 0 & 0 & 0 & 0 & 0 & 0 & 0 & 0 & 0 & 0 & 0 & 0 & 0 & 0\\
0 & -1 & 0 & 0 & 0 & 0 & 0 & 0 & 0 & 0 & 0 & 0 & 0 & 0 & 0\\
0 & 0 & 0 & 0 & 0 & 1 & 0 & 0 & 0 & 0 & 0 & 0 & 0 & 0 & 0\\
0 & 0 & 0 & 0 & 0 & 1 & 0 & -1 & 0 & 0 & 0 & 0 & 0 & 0 & 0\\
0 & 0 & 0 & 0 & 0 & 0 & 1 & 0 & -1 & 0 & 0 & 0 & 0 & 0 & 0\\
0 & 0 & 0 & 0 & 0 & 0 & 0 & 0 & 0 & 0 & 0 & -1 & 0 & 0 & 0\\
0 & 0 & 0 & 0 & 0 & 0 & 0 & 0 & 0 & 0 & 0 & -1 & -1 & 0 & 0\\
0 & 0 & 0 & 0 & 0 & 0 & 0 & 0 & 0 & 0 & 1 & 0 & 0 & -1 & 0\\
\end{pmatrix}$}
\]
\par\smallskip\noindent
$F_2$ has size $29\times 30$; its full integer matrix is
\[
\resizebox{0.94\linewidth}{!}{$\displaystyle F_2=\begin{pmatrix}
1 & 0 & 0 & 0 & 0 & 0 & 0 & -1 & 0 & 0 & 1 & 0 & 0 & 0 & 0 & 0 & 0 & 0 & 0 & 0 & 1 & 0 & 0 & 0 & 0 & 0 & 0 & 0 & 0 & 0\\
0 & 1 & 0 & 0 & 0 & 0 & 0 & 0 & -1 & 0 & 0 & 1 & 0 & 0 & 0 & 0 & 0 & 0 & 0 & 0 & 0 & 1 & 0 & 0 & 0 & 0 & 0 & 0 & 0 & 0\\
0 & 0 & 1 & 0 & 0 & 0 & 0 & 0 & 0 & -1 & 0 & 0 & 1 & 0 & 0 & 0 & 0 & 0 & 0 & 0 & 0 & 0 & 1 & 0 & 0 & 0 & 0 & 0 & 0 & 0\\
0 & 0 & 0 & 1 & 0 & 0 & 0 & 0 & 0 & 0 & 0 & 0 & 0 & 1 & 0 & 0 & 0 & 0 & 0 & 0 & 0 & 0 & 0 & 1 & 0 & 0 & 0 & 0 & 0 & 0\\
0 & 0 & 0 & 0 & 1 & 0 & 0 & 0 & 0 & 0 & 0 & 0 & 0 & 0 & 1 & 0 & 0 & 0 & 0 & 0 & 0 & 0 & 0 & 0 & 1 & 0 & 0 & 0 & 0 & 0\\
0 & 0 & 0 & 0 & 0 & 1 & 0 & 0 & 0 & 0 & 0 & 0 & 0 & 0 & 0 & 1 & 0 & 0 & 0 & 0 & 0 & 0 & 0 & 0 & 0 & 1 & 0 & 0 & 0 & 0\\
0 & 0 & 0 & 0 & 0 & 0 & 1 & 0 & 0 & 0 & 0 & 0 & 0 & 0 & 0 & 0 & 0 & 0 & 0 & 0 & 0 & 0 & 0 & 0 & 0 & 0 & -1 & 0 & 0 & 0\\
0 & 0 & 0 & 0 & 0 & 0 & 0 & 1 & 0 & 0 & 0 & 0 & 0 & 0 & 0 & 0 & 0 & 0 & 0 & 0 & 0 & 0 & 0 & 0 & 0 & 0 & 0 & -1 & 0 & 0\\
0 & 0 & 0 & 0 & 0 & 0 & 0 & 0 & 1 & 0 & 0 & 0 & 0 & 0 & 0 & 0 & 0 & 0 & 0 & 0 & 0 & 0 & 0 & 0 & 0 & 0 & 0 & 0 & -1 & 0\\
0 & 0 & 0 & 0 & 0 & 0 & 0 & 0 & 0 & 1 & 0 & 0 & 0 & 0 & 0 & 0 & 0 & 0 & 0 & 0 & 0 & 0 & 0 & 0 & 0 & 0 & 0 & 0 & 0 & -1\\
0 & 0 & 0 & 0 & 0 & 0 & 0 & 0 & 0 & 0 & 0 & 0 & 0 & 0 & 0 & 0 & 1 & 0 & 0 & 0 & 0 & 0 & 0 & 0 & 0 & 0 & -1 & 0 & 0 & 0\\
0 & 0 & 0 & 0 & 0 & 0 & 0 & 0 & 0 & 0 & 0 & 0 & 0 & 0 & 0 & 0 & 0 & 1 & 0 & 0 & 0 & 0 & 0 & 0 & 0 & 0 & 0 & -1 & 0 & 0\\
0 & 0 & 0 & 0 & 0 & 0 & 0 & 0 & 0 & 0 & 0 & 0 & 0 & 0 & 0 & 0 & 0 & 0 & 1 & 0 & 0 & 0 & 0 & 0 & 0 & 0 & 0 & 0 & -1 & 0\\
0 & 0 & 0 & 0 & 0 & 0 & 0 & 0 & 0 & 0 & 0 & 0 & 0 & 0 & 0 & 0 & 0 & 0 & 0 & 1 & 0 & 0 & 0 & 0 & 0 & 0 & 0 & 0 & 0 & -1\\
0 & 1 & 0 & 0 & 0 & 0 & 0 & 0 & 0 & 0 & 0 & 0 & 0 & 0 & 0 & 0 & 0 & 0 & 0 & 0 & 0 & 0 & 0 & 0 & 0 & 0 & 0 & 0 & 0 & 0\\
0 & 0 & 0 & 1 & 0 & 0 & 0 & 0 & 0 & 0 & 0 & 0 & 0 & 0 & 0 & 0 & 0 & 0 & 0 & 0 & 0 & 0 & 0 & 0 & 0 & 0 & 0 & 0 & 0 & 0\\
0 & 0 & 1 & 0 & 0 & 0 & 0 & 0 & 0 & 0 & 0 & 0 & 0 & 0 & 0 & 0 & 0 & 0 & 0 & 0 & 0 & 0 & 0 & 0 & 0 & 0 & 0 & 0 & 0 & 0\\
0 & 0 & 0 & 0 & 1 & 0 & 0 & 0 & 0 & 0 & 0 & 0 & 0 & 0 & 0 & 0 & 0 & 0 & 0 & 0 & 0 & 0 & 0 & 0 & 0 & 0 & 0 & 0 & 0 & 0\\
-1 & 0 & 0 & 0 & 0 & 0 & 0 & 0 & 0 & 0 & 0 & 0 & 0 & 0 & 0 & 0 & 0 & 0 & 0 & 0 & 0 & 0 & 0 & 0 & 0 & 0 & 0 & 0 & 0 & 0\\
0 & 0 & 0 & 0 & 0 & 0 & 0 & 0 & 0 & 0 & -1 & 0 & 0 & -1 & 0 & 0 & 0 & 0 & 0 & 0 & 0 & 0 & 0 & 0 & 0 & 0 & 0 & 0 & 0 & 0\\
0 & 0 & 0 & 0 & 0 & 0 & 0 & 0 & 0 & 0 & 0 & 0 & 0 & 0 & -1 & 0 & 0 & 0 & 0 & 0 & 0 & 0 & 0 & 0 & 0 & 0 & 0 & 0 & 0 & 0\\
0 & 0 & 0 & 0 & 0 & 0 & 0 & 0 & 0 & 0 & 0 & 1 & 0 & 0 & 0 & 0 & 0 & 0 & 0 & 0 & 0 & 0 & 0 & 0 & 0 & 0 & 0 & 0 & 0 & 0\\
0 & 0 & 0 & 0 & 0 & 0 & 0 & 0 & 0 & 0 & -1 & 0 & 1 & 0 & 0 & 0 & 0 & 0 & 0 & 0 & 0 & 0 & 0 & 0 & 0 & 0 & 0 & 0 & 0 & 0\\
0 & 0 & 0 & 0 & 0 & 0 & 0 & 0 & 0 & 0 & -1 & 0 & 1 & -1 & 0 & -1 & 0 & 0 & 0 & 0 & 0 & 0 & 0 & 0 & 0 & 0 & 0 & 0 & 0 & 0\\
0 & 0 & 0 & 0 & 0 & 0 & 0 & 0 & 0 & 0 & 0 & 0 & 0 & 0 & 0 & 0 & 0 & 0 & 0 & 0 & -1 & -1 & 0 & 0 & 0 & 0 & 0 & 0 & 0 & 0\\
0 & 0 & 0 & 0 & 0 & 0 & 0 & 0 & 0 & 0 & 0 & 0 & 0 & 0 & 0 & 0 & 0 & 0 & 0 & 0 & 0 & 0 & -1 & 0 & 0 & 0 & 0 & 0 & 0 & 0\\
0 & 0 & 0 & 0 & 0 & 0 & 0 & 0 & 0 & 0 & 0 & 0 & 0 & 0 & 0 & 0 & 0 & 0 & 0 & 0 & 0 & 0 & 0 & -1 & 0 & 0 & 0 & 0 & 0 & 0\\
0 & 0 & 0 & 0 & 0 & 0 & 0 & 0 & 0 & 0 & 0 & 0 & 0 & 0 & 0 & 0 & 0 & 0 & 0 & 0 & -1 & 0 & 0 & 0 & -1 & 0 & 0 & 0 & 0 & 0\\
0 & 0 & 0 & 0 & 0 & 0 & 0 & 0 & 0 & 0 & 0 & 0 & 0 & 0 & 0 & 0 & 0 & 0 & 0 & 0 & -1 & -1 & 0 & 0 & -1 & -1 & 0 & 0 & 0 & 0\\
\end{pmatrix}$}
\]
\par\smallskip\noindent
$R_{H_1}$ has size $4\times 4$; its full integer matrix is
\[
\resizebox{!}{11.2mm}{$\displaystyle R_{H_1}=\begin{pmatrix}
1 & 0 & 0 & 1\\
0 & 0 & 1 & 0\\
0 & 1 & 0 & 0\\
0 & 0 & 1 & 1\\
\end{pmatrix}$}
\]
\par\smallskip\noindent
$R_{H_2}$ has size $6\times 6$; its full integer matrix is
\[
\resizebox{!}{16.799999999999997mm}{$\displaystyle R_{H_2}=\begin{pmatrix}
1 & 0 & 0 & 0 & 1 & 1\\
0 & 1 & 0 & 0 & 0 & -1\\
0 & 0 & 1 & 0 & 1 & 0\\
0 & 0 & 0 & 0 & -1 & 0\\
0 & 0 & 0 & 0 & 0 & -1\\
0 & 0 & 0 & 1 & 1 & 1\\
\end{pmatrix}$}
\]
\endgroup


\begin{proposition}[P4.5]\label{night:P4.5}\label{stmt:P4.5}
\StatementRecord{night:P4.5}{proved}{Historical events: GPT-6 Astra (OpenAI), ChatGPT Pro mode, effort Pro, S6\_06\_BLIND, 16--18 November 2026: targeted proof check. Also checked in this writing. Pre-continuation writing events: GPT-6 Astra (OpenAI), ChatGPT Pro mode, effort Pro, S6\_10\_PAPERS, 18 November 2026: targeted proof check. Continuation metadata events: GPT-6 Astra (OpenAI), ChatGPT Pro mode, effort Pro, S6\_10\_PAPERS continuation, 19 November 2026: used and rechecked in this run; status-vocabulary reconciliation; metadata reconciliation only; no new mathematical check. \emph{Pre-continuation recorded proof scope.} The night-run re-derivation was by GPT-6 Astra (OpenAI, ChatGPT Pro mode, effort Pro), 16--18 November 2026. Also checked in this writing by GPT-6 Astra (OpenAI), ChatGPT Pro mode, effort Pro, 18 November 2026. Stratified Euler additivity gives zero for the actual smooth torus fibres and Kodaira-Hopf central supports, and Poincare duality with the computed integral H1 and H2 gives the asserted torsion-free cohomology and b3=2. An effective horizontal polar divisor would restrict to a nonzero effective divisor of zero real Chern class on a Kahler torus. Excluding it, removing vertical poles by local powers of the base coordinate, and connected proper descent identify the meromorphic function field with that of P1}{\ref{night:P4.4}}

For this example,
\[
 a(Y)=1,\qquad c_3[Y]=0,\qquad b_3(Y)=2,
\]
and the integral homology is free of ranks $(1,0,0,2,0,0,1)$.
\end{proposition}
\begin{proof}
\noindent\textbf{Step 1.} The regular fibre has Euler number zero. Both components of each special
fibre are elliptic bundles over compact curves, and their intersection is
elliptic, so additivity gives zero for each special support. Stratifying
the base and applying multiplicativity to the smooth bundle gives
$e(Y)=0$. The Euler class of the real bundle underlying a complex
rank-three tangent bundle is its top Chern class, with the complex
orientation. Hence $c_3[Y]=0$. Duality and P4.4 give $b_4=b_5=0$ and
$e(Y)=2-b_3$, so $b_3=2$. Integral duality and the universal-coefficient
theorem identify the torsion of $H_3$ with that of $H_2$, which is zero.
All other groups have already been determined.

\noindent\textbf{Step 2.} We also have $H^2(Y;\Z)=0$. A horizontal effective divisor would restrict
to a nonzero effective divisor on the Kähler torus $A$ with zero first
Chern class. Its pairing with a Kähler class would be positive by
integration over the divisor and zero by the Chern-class computation.
Thus every divisor on $Y$ is vertical. The polar divisor of a meromorphic
function is consequently vertical; on a general fibre the function is
holomorphic and therefore constant. Local holomorphic sections give a
descended meromorphic function on the regular part of the base. Near
each exceptional base value a sufficiently high power of the base
coordinate removes the vertical poles. Connected fibres then give
holomorphic descent there, and dividing by that power gives meromorphic
descent. Therefore $\C(Y)=f^*\C(\Pp^1)$ and $a(Y)=1$.
\end{proof}

\subsection{The normal-crossings sign theorem}

\begin{proposition}[P4.6]\label{night:P4.6}\label{stmt:P4.6}
\StatementRecord{night:P4.6}{proved}{Historical events: GPT-6 Astra (OpenAI), ChatGPT Pro mode, effort Pro, S6\_06\_BLIND, 16--18 November 2026: targeted proof check. Also checked in this writing. Pre-continuation writing events: GPT-6 Astra (OpenAI), ChatGPT Pro mode, effort Pro, S6\_10\_PAPERS, 18 November 2026: targeted proof check. Continuation metadata events: GPT-6 Astra (OpenAI), ChatGPT Pro mode, effort Pro, S6\_10\_PAPERS continuation, 19 November 2026: used and rechecked in this run; status-vocabulary reconciliation; metadata reconciliation only; no new mathematical check. \emph{Pre-continuation recorded proof scope.} The night-run re-derivation was by GPT-6 Astra (OpenAI, ChatGPT Pro mode, effort Pro), 16--18 November 2026. Also checked in this writing by GPT-6 Astra (OpenAI), ChatGPT Pro mode, effort Pro, 18 November 2026. The branchwise normal derivative gives the actual line isomorphism N\_i=nu\_i*O(F) tensor O(-D\_i), also for self-intersecting global components. Its tangent Chern sequence gives c2=K2 and K numerically minus D. The unramified branch-choice double cover and three conductor nodes at each triple point give e(Gamma)=3T. The alternating normalization complex carries the degree-zero sign line on its double term; coherent HRR for the Cartier sequence gives T=2 sum chi(O\_Si). Topological inclusion-exclusion then gives e(F)=T and parity, without a multiplicity assumption}{the hypotheses and explicit local calculation stated here}

Let $X$ be a compact complex threefold with $b_2(X)=0$, and let $F$ be a
reduced compact normal-crossings divisor. We allow different local
branches of one global irreducible component to meet. If $T(F)$ is the
number of ordinary triple points, then
\[
 e(F)=T(F),\qquad T(F)\text{ is even}.
\]
This assertion concerns the reduced support, regardless of the
multiplicities with which it occurs in a fibre.
\end{proposition}
\begin{proof}
\noindent\textbf{Step 1.} Write $\nu:S=\coprod S_i\to X$ for the normalization immersion, and $D_i$
for the conductor divisor on $S_i$. The surfaces $S_i$ are smooth and
the curves $D_i$ are reduced nodal divisors. The quotient
$N_i=\nu_i^*T_X/T_{S_i}$ is a line bundle. A local equation of $F$ gives
\[
 N_i\simeq \nu_i^*\OO_X(F)\otimes\OO_{S_i}(-D_i).
\]
Indeed, for the coordinate equation $x_1\cdots x_k=0$ the normal
derivative on $x_1=0$ is $x_2\cdots x_k$, whose zero divisor is $D_i$.
These are morphisms of holomorphic vector bundles, including when
$S_i$ is not embedded globally.

\noindent\textbf{Step 2.} Since $b_2(X)=0$, both $c_1(T_X)$ and $c_1(\OO_X(F))$ vanish in real
cohomology. Duality gives $b_4(X)=0$, so $c_2(T_X)$ also vanishes. The
tangent exact sequence yields
\[
 K_{S_i}\equiv-D_i,\qquad c_2(S_i)=K_{S_i}^2.
\]
If $a=\sum_i\chi(\OO_{S_i})$, the surface Riemann--Roch formula
$12\chi(\OO_{S_i})=K_{S_i}^2+e(S_i)$ consequently gives
$\sum_i e(S_i)=6a$.

\noindent\textbf{Step 3.} Let $\Gamma$ be the normalization of the double-curve locus, with one
sheet for each unordered pair of local branches. The normalization of
$\coprod D_i$ is an unramified degree-two cover of $\Gamma$, recording
the choice of branch. Each triple point contributes three nodes to
$\coprod D_i$. Adjunction gives
$\chi(\OO_{D_i})=-D_i\cdot(D_i+K_{S_i})/2=0$. The normalization
sequence for these nodal curves therefore gives
\[
 0=2\chi(\OO_\Gamma)-3T(F)=e(\Gamma)-3T(F).
\]
The normalization complex of $F$ is the exact alternating restriction
complex of the local coordinate hyperplanes. Permuting the branches
globally replaces its double-curve term by $\OO_\Gamma$ tensored with
the sign line of the two-sheeted branch cover. That line has order at
most two, and hence degree zero on every component; it has the same
coherent Euler characteristic as $\OO_\Gamma$. The triple term has
length one at each triple point. It follows that
\[
 \chi(\OO_F)=a-\chi(\OO_\Gamma)+T(F)=a-T(F)/2.
\]
The exact sequence for the Cartier divisor $F$ and analytic
Hirzebruch--Riemann--Roch give $\chi(\OO_F)=0$, since
$\operatorname{ch}(\OO_X(-F))=\operatorname{ch}(\OO_X)$ in rational
cohomology. Thus $T(F)=2a$. Finally ordinary Euler additivity gives
\[
 e(F)=\sum_i e(S_i)-e(\Gamma)+T(F)=6a-2T(F)=T(F).
\]
The alternating count is also verified directly at a point with one,
two or three branches. This proves the identity and parity, without
using fibre multiplicities or a base change.
\end{proof}

\begin{proposition}[P4.7]\label{night:P4.7}\label{stmt:P4.7}
\StatementRecord{night:P4.7}{proved}{Historical events: GPT-6 Astra (OpenAI), ChatGPT Pro mode, effort Pro, S6\_06\_BLIND, 16--18 November 2026: targeted proof check. Also checked in this writing. Pre-continuation writing events: GPT-6 Astra (OpenAI), ChatGPT Pro mode, effort Pro, S6\_10\_PAPERS, 18 November 2026: targeted proof check. Continuation metadata events: GPT-6 Astra (OpenAI), ChatGPT Pro mode, effort Pro, S6\_10\_PAPERS continuation, 19 November 2026: used and rechecked in this run; status-vocabulary reconciliation; metadata reconciliation only; no new mathematical check. \emph{Pre-continuation recorded proof scope.} The night-run re-derivation was by GPT-6 Astra (OpenAI, ChatGPT Pro mode, effort Pro), 16--18 November 2026. Also checked in this writing by GPT-6 Astra (OpenAI), ChatGPT Pro mode, effort Pro, 18 November 2026. The connected-fibre constant-sheaf edge injects H1 of the base. Smooth-fibre adjunction gives c2=0 from the ambient H4 vanishing. Properness makes the critical-value set finite, and Ehresmann plus Euler additivity gives the exact sum over supports. Poincare duality and the alternating middle pairing force the nonnegative sum to lie in \{0,2\}. No torus hypothesis is used in this numerical argument}{\ref{night:P4.6}}

Let $f:X\to B$ be a proper surjective holomorphic map with connected
fibres from a compact complex threefold satisfying $b_1(X)=b_2(X)=0$
to a smooth compact curve. Suppose all reduced fibre supports have
normal crossings. Then
\[
 B=\Pp^1,\qquad
 c_3[X]=\sum_{b\in B}T((f^{-1}(b))_{\mathrm{red}})\in\{0,2\},\qquad
 b_3(X)=2-c_3[X].
\]
The numerical conclusion does not require the general fibre to be a
torus.
\end{proposition}
\begin{proof}
\noindent\textbf{Step 1.} Connected fibres give $f_*\R=\R$. The beginning of the constant-sheaf
Leray sequence injects $H^1(B;\R)$ into $H^1(X;\R)$, so $B$ has genus
zero. A smooth fibre $G$ has trivial normal line bundle. Its tangent
sequence gives $c_2(G)=c_2(X)|_G=0$, hence $e(G)=0$ for any such fibre.
The critical values are finite. Ehresmann's theorem gives a smooth
bundle over their complement. Euler additivity and P4.6 give the
displayed sum. Duality gives $e(X)=2-b_3(X)$, and the nondegenerate
alternating cup product on $H^3(X;\R)$ makes $b_3$ even. The
nonnegative sum is therefore zero or two.
\end{proof}

\begin{proposition}[P4.8]\label{night:P4.8}\label{stmt:P4.8}
\StatementRecord{night:P4.8}{proved}{Historical events: GPT-6 Astra (OpenAI), ChatGPT Pro mode, effort Pro, S6\_06\_BLIND, 16--18 November 2026: targeted proof check. Also checked in this writing. Pre-continuation writing events: GPT-6 Astra (OpenAI), ChatGPT Pro mode, effort Pro, S6\_10\_PAPERS, 18 November 2026: targeted proof check; used and rechecked in this run. Continuation metadata events: GPT-6 Astra (OpenAI), ChatGPT Pro mode, effort Pro, S6\_10\_PAPERS continuation, 19 November 2026: used and rechecked in this run; status-vocabulary reconciliation; metadata reconciliation only; no new mathematical check. \emph{Pre-continuation recorded proof scope.} The night-run re-derivation was by GPT-6 Astra (OpenAI, ChatGPT Pro mode, effort Pro), 16--18 November 2026. Also checked in this writing by GPT-6 Astra (OpenAI), ChatGPT Pro mode, effort Pro, 18 November 2026. The equality e=K2 gives k=2(1-g) in the ruled case and minimality with b2=0 in the non-Kahler classVII case. The actual effective conductor excludes Inoue covers and, for irreducible supports, its torsion adjunction kills every positive plurigenus. Pullback of an adjacent component excludes Kahler normalization in a reducible support. The zero triple count makes the conductor curves disjoint elliptic curves with the unramified branch maps. The normalized root cover of a smooth multiple support is a proper submersion and supplies the finite unramified torus cover. BHPV, BBK, OVV, Khaled-Teleman and the earlier Catanese exact source loci are now matched; the printed Remark3.4 pointer is corrected to Remark3.5}{\ref{night:P4.12}; \ref{night:P4.6}; \ref{night:P4.7}}

For a torus fibration in P4.7 whose reduced supports have normal
crossings, an irreducible singular support with
rational normalization has Euler number two; every other irreducible
support has Euler number zero. In the zero case the supports have the
forms stated in the prompt: a smooth support has a finite unramified
torus cover; a singular irreducible support has smooth minimal
elliptic-ruled or Hopf normalization; a reducible support has smooth
minimal Hopf or non-Kähler elliptic normalization components, with
smooth elliptic double curves and finite unramified branch maps.
Thus $c_3[X]=0$ precisely when every fibre belongs to these zero cases.
\end{proposition}
\begin{proof}
\noindent\textbf{Step 1.} We apply \cite[Chapter~VI, Theorem~1.1 and Table~10, pp.~243--244]{BHPV2004E} to the smooth compact
normalization components, using the identity $e(S_i)=K_{S_i}^2$ proved
in P4.6. The relevant non-Kähler classification is stated in
\cite[Theorem~4.2 and proof, pp.~17--21; Theorem~5.1 and proof, pp.~27--29]{OVV2020}: a minimal non-Kähler
surface is elliptic or of class VII, and the reduced fibres of a
non-Kähler elliptic surface are smooth elliptic curves.

\noindent\textbf{Step 2.} For a component with nonempty conductor, a Kähler form $\omega$ would
give $K_{S_i}\cdot\omega=-D_i\cdot\omega<0$. No positive power of
the canonical bundle can then have a nonzero section. The classification
makes a Kähler component rational or ruled. If it is obtained by $k$
blowups of a ruled surface over a genus-$g$ curve, then
\[
 e=4(1-g)+k,\qquad K^2=8(1-g)-k.
\]
Equality gives $k=2(1-g)$, so $g\le1$ and the genus-one case is minimal.
A rational component has $\chi(\OO)=1$, hence $e=6$ here, including
the case with minimal model $\Pp^2$.

\noindent\textbf{Step 3.} A minimal non-Kähler elliptic surface has $e=K^2=0$, by its smooth
elliptic reduced fibres. A minimal class-VII surface has $b_1=1$ and
$p_g=0$, so $\chi(\OO)=0$ and $K^2=-e$, with $e=b_2\ge0$.
Blowups preserve $\chi(\OO)$ and change $(e,K^2)$ by $(1,-1)$.
Consequently $e=K^2$ forces minimality and $e=0$ in all these
non-Kähler cases. The class-VII case has $b_2=0$. Theorem 1 of
Buonerba--Bogomolov--Kurnosov \cite[Theorem~1, p.~1]{BBK2019} supplies a finite unramified Hopf or
Inoue cover. The conductor excludes an Inoue cover: Inoue surfaces
have no compact curves, by Inoue's theorem as stated and credited in
\cite[Theorem~3.1, p.~303]{KhaledTeleman2026Journal}.
The type-III family has the explicit type-II unramified double cover
of \cite[proof of Remark~3.5, p.~305]{KhaledTeleman2026Journal}. A finite unramified quotient of a Hopf surface has universal
cover $\C^2\setminus\{0\}$ and is therefore a Hopf surface.

\noindent\textbf{Step 4.} For an irreducible singular support there is one normalization
component. P4.6 gives $T=2\chi(\OO_{S_1})$, which is two in the
rational case and zero in the other cases. P4.12 proves that the
dualizing sheaf of an irreducible support is torsion. Adjunction on
the normalization therefore makes $K_{S_1}+D_1$ torsion. Since
$D_1\ne0$, every positive plurigenus of $S_1$ vanishes: multiplying
a pluricanonical section by an appropriate positive power of the
conductor section would give a section of a torsion line bundle with
a zero, whose sufficiently high power is a nonzero constant. Thus a
non-Kähler normalization in this irreducible case is of class VII,
and hence is the Hopf case already obtained.

\noindent\textbf{Step 5.} For a reducible connected fibre, each normalization component meets
another global irreducible component. Pulling back the divisor of
that component gives a nonzero effective divisor on $S_i$ with zero
real first Chern class. A Kähler form would pair positively with it,
contradicting this zero class. Hence every component is non-Kähler,
has $\chi(\OO)=0$, and $T=0$. The conductor components are disjoint
and smooth, because all their nodes would come from triple points.
For any conductor component $C$, adjunction and $K+D\equiv0$ give
$2g(C)-2=(K+C)\cdot C=0$, since the other conductor components
are disjoint. Thus the double curves are smooth elliptic, and their
branch maps are the unramified maps already described in P4.6.

\noindent\textbf{Step 6.} Finally let the support $G$ of a multiple fibre $mG$ be smooth. The
normalized cyclic base change $t=u^m$ is locally a disjoint union of
smooth charts: the original equation is $t=z^m$ times a holomorphic
unit, whose root can be extracted locally. These charts show that
the new central fibre maps finitely and unramifiedly to $G$, and
that the new map to the $u$-disc is a proper submersion. Its nearby
fibres are the original tori. Catanese's theorem
\cite[Theorem~2.1]{Catanese2004} states that every
deformation of a complex torus is a complex torus; its definition
uses precisely a proper holomorphic submersion over a connected
disc, with no Kähler hypothesis on the total space. We have verified
all these hypotheses. The new central fibre is therefore a torus;
Ehresmann's theorem also shows that it is connected. This gives
the required finite unramified cover of $G$.
\end{proof}

\begin{proposition}[P4.9]\label{night:P4.9}\label{stmt:P4.9}
\StatementRecord{night:P4.9}{proved}{Historical events: GPT-6 Astra (OpenAI), ChatGPT Pro mode, effort Pro, S6\_06\_BLIND, 16--18 November 2026: targeted proof check. Also checked in this writing. Pre-continuation writing events: GPT-6 Astra (OpenAI), ChatGPT Pro mode, effort Pro, S6\_10\_PAPERS, 18 November 2026: targeted proof check. Continuation metadata events: GPT-6 Astra (OpenAI), ChatGPT Pro mode, effort Pro, S6\_10\_PAPERS continuation, 19 November 2026: used and rechecked in this run; status-vocabulary reconciliation; metadata reconciliation only; no new mathematical check. \emph{Pre-continuation recorded proof scope.} The night-run re-derivation was by GPT-6 Astra (OpenAI, ChatGPT Pro mode, effort Pro), 16--18 November 2026. Also checked in this writing by GPT-6 Astra (OpenAI), ChatGPT Pro mode, effort Pro, 18 November 2026. Followed the actual local irreducible sign, rational and classVII equality statements, zero-Euler normalization maps, and positive-ruled exclusions through P4.26, P4.32, P4.35-37, P4.54-58 and P4R.18-20. The global conclusion retains the hypothesis that all supports are irreducible or NC. This assembly review adds no new individual proof credit to those separately reviewed inputs}{\ref{night:P4.26}; \ref{night:P4.32}; \ref{night:P4.35}; \ref{night:P4.37}; \ref{night:P4.54}; \ref{night:P4.57}; \ref{night:P4.58}; \ref{night:P4R.18}; \ref{night:P4R.20}}

In the requested irreducible and normal-crossing fibre classes,
the complete sign, exact Euler values and normalization equality
classification hold as stated in P4.58. In particular every
irreducible singular support with rationally resolved normalization
has Euler number two, and every other irreducible support has
Euler number zero.
\end{proposition}
\begin{proof}
\noindent\textbf{Step 1.} P4.26 proves the unrestricted irreducible numerical sign;
P4.35--P4.37 give the rational and class-VII exact values.
P4.32 proves the zero-Euler classification. The positive ruled
alternative is excluded by the complete sequence P4.54--P4.57
and P4R.18--P4R.20, synthesized in P4.58.
\end{proof}
We checked CDP 2020, Lemma~4.6, Propositions~4.7,
4.11 and~7.3, deriving the sign and classification independently
from its local conductor identities. P4.14 records the precise
incorrect curve-classification inference in Proposition~7.3.
Earlier proof obligations and unsuccessful reductions are retained
in the working notes; the final classification is P4.58.
The original stated-value row also claimed a sign for arbitrary
reducible supports without normal crossings. Such supports are
outside the two classes explicitly requested in Part~4.2; the
proved global assertion retains the stated irreducible-or-NC
support hypothesis.

\subsection{Explicit checks on the hypotheses}

\begin{proposition}[P4.10]\label{night:P4.10}\label{stmt:P4.10}
\StatementRecord{night:P4.10}{proved}{Historical events: GPT-6 Astra (OpenAI), ChatGPT Pro mode, effort Pro, S6\_06\_BLIND, 16--18 November 2026: targeted proof check. Also checked in this writing. Pre-continuation writing events: GPT-6 Astra (OpenAI), ChatGPT Pro mode, effort Pro, S6\_10\_PAPERS, 18 November 2026: targeted proof check; used and rechecked in this run. Continuation metadata events: GPT-6 Astra (OpenAI), ChatGPT Pro mode, effort Pro, S6\_10\_PAPERS continuation, 19 November 2026: used and rechecked in this run; status-vocabulary reconciliation; metadata reconciliation only; no new mathematical check. \emph{Pre-continuation recorded proof scope.} The night-run re-derivation was by GPT-6 Astra (OpenAI, ChatGPT Pro mode, effort Pro), 16--18 November 2026. Also checked in this writing by GPT-6 Astra (OpenAI), ChatGPT Pro mode, effort Pro, 18 November 2026. The two Weierstrass charts have twelve simple nodal discriminant values, smooth total surface and relative dualizing line O(1), so the displayed direct-image computation makes the surface rational. Blowing up the actual diagonal ideal (x-u,y-v) in xy=uv gives both displayed smooth charts and reduced normal-crossing central equations. No new surface component appears. Each exceptional P1 changes the fibre Euler number from one to two. The extended section and two factor surface maps kill the actual complement generators of pi1. The result has c3=24 and a nonzero Kahler class in H2}{the hypotheses and explicit local calculation stated here}

There is a smooth projective threefold $Z$ with $b_1(Z)=0$ and a
connected torus fibration $Z\to\Pp^1$ whose reduced special fibres
are irreducible normal-crossings surfaces, but $c_3[Z]=24$.
Consequently the hypothesis $b_2=0$ cannot be omitted, even in the
intersection of the two classes of fibres in the question.
\end{proposition}
\begin{proof}
\noindent\textbf{Step 1.} On the two affine charts of $\Pp^1$, with $s=t^{-1}$, we take
\[
 y^2z=x^3-3xz^2+t^6z^3,\qquad
 y_\infty^2z=x_\infty^3-3s^4x_\infty z^2+z^3,
\]
glued by $x_\infty=s^2x$ and $y_\infty=s^3y$. These are the
Weierstrass cubics in $\Pp(\OO\oplus\OO(2)\oplus\OO(3))$ for
the line bundle $\OO_{\Pp^1}(1)$. They define an elliptic surface
$\pi:S\to\Pp^1$ with zero section $[0:1:0]$. The discriminant
is a nonzero constant times $t^{12}-4$, or $1-4s^{12}$ on the
other chart. There are twelve simple zeros and a smooth fibre at
infinity. At a discriminant zero the cubic has a double root
$x=1$ or $x=-1$ and a distinct third root, so the fibre has an
ordinary node. Since $6t^5\ne0$ there, the total surface is smooth.
At the zero section the derivative in $z$ is nonzero. The map near
each node is therefore $xy=t$ after holomorphic coordinate changes
and translation of the base.

\noindent\textbf{Step 2.} The relative dualizing differential $dx/y$ has transitions
$\pi^*\OO(1)$ and generates the dualizing sheaf also at the nodes.
Thus $\omega_S=\pi^*\OO(-1)$. Relative duality gives
$R^1\pi_*\OO_S=\OO(-1)$, and connected fibres give
$\pi_*\OO_S=\OO$. Leray yields $h^1(\OO_S)=0$, while all
positive plurigenera vanish. The classification \cite[Chapter~VI, Theorem~1.1, p.~243]{BHPV2004E} makes $S$
rational; in the ruled case the base genus equals $h^1(\OO_S)=0$.
In particular $\pi_1(S)=1$, since the minimal rational models
$\Pp^2$ and the Hirzebruch surfaces are simply connected and a
point blowup preserves the fundamental group.

\noindent\textbf{Step 3.} The projective fibre product $W=S\times_{\Pp^1}S$ has exactly
twelve singularities, at the pair of nodes in each nodal fibre.
Their equations are $xy=uv$, with base function $xy=uv$. We blow
up the coherent ideal of the diagonal $\Delta_S\subset W$. This
divisor is Cartier away from the twelve points, so the blowup is
the identity there. Put $r=x-u$, $h=y-v$ at a singular point.
The equation becomes $ry+uh=0$ and the diagonal ideal is $(r,h)$.
The blowup has charts
\[
 h=ar,\quad y=-au,\quad x=u+r,\quad f=-aux;
 \qquad
 r=bh,\quad u=-by,\quad x=-bv,\quad f=-byv.
\]
The independent coordinates are $(a,u,x)$ and $(b,y,v)$.
Both charts are smooth, with $b=a^{-1}$ on the overlap, and the
exceptional fibre is $\Pp^1$. Thus the resulting $Z$ is smooth
and projective. The displayed equations make every special fibre
reduced and normal crossings. Its support remains irreducible:
the original product of two irreducible nodal cubics is irreducible,
and the exceptional locus has dimension one, so introduces no
surface component. Its normalization is rational, by the
birational map from $\Pp^1\times\Pp^1$.

\noindent\textbf{Step 4.} A nodal cubic has Euler number $2-1=1$, so its self-product has
Euler number one. The resolution replaces one point by $\Pp^1$;
each resolved special support therefore has Euler number two.
All other fibres are $E_t\times E_t$ with Euler number zero.
Hence $e(Z)=12\cdot2=24=c_3[Z]$.

\noindent\textbf{Step 5.} For completeness we compute $\pi_1(Z)$. Every loop can be moved
off the fibre divisors, so the fundamental group of their complement
surjects onto $\pi_1(Z)$. The complement is a bundle with fibre
$E\times E$ and a section over a twelve-punctured sphere. The
section extends over $\Pp^1$, killing the base generators. The
two maps
\[
 S\longrightarrow Z:\quad p\mapsto(p,O(\pi(p))),\qquad
 p\mapsto(O(\pi(p)),p)
\]
are defined because their images miss every conifold point, and
extend the two factor inclusions of $E$. Since $S$ is simply
connected they kill both fibre fundamental groups. Thus
$\pi_1(Z)=1$ and $b_1(Z)=0$. A projective Kähler class is nonzero
in $H^2(Z;\R)$, because its cube has positive integral. This
identifies the omitted hypothesis: $b_2(Z)>0$.
\end{proof}

\begin{proposition}[P4.11]\label{night:P4.11}\label{stmt:P4.11}
\StatementRecord{night:P4.11}{proved}{Historical events: GPT-6 Astra (OpenAI), ChatGPT Pro mode, effort Pro, S6\_06\_BLIND, 16--18 November 2026: targeted proof check. Also checked in this writing. Pre-continuation writing events: GPT-6 Astra (OpenAI), ChatGPT Pro mode, effort Pro, S6\_10\_PAPERS, 18 November 2026: targeted proof check. Continuation metadata events: GPT-6 Astra (OpenAI), ChatGPT Pro mode, effort Pro, S6\_10\_PAPERS continuation, 19 November 2026: used and rechecked in this run; status-vocabulary reconciliation; metadata reconciliation only; no new mathematical check. \emph{Pre-continuation recorded proof scope.} The night-run re-derivation was by GPT-6 Astra (OpenAI, ChatGPT Pro mode, effort Pro), 16--18 November 2026. Also checked in this writing by GPT-6 Astra (OpenAI), ChatGPT Pro mode, effort Pro, 18 November 2026. Freeness and properness of the exact complex action, its norm-level S3 times S3 model, and the induced elliptic pencil with Hopf fibres were re-derived. The fibre fundamental group is Z, whereas an etale quotient of a two-torus has a finite-index Z4 subgroup. The P4.7 dependency is used only for the final numerical comparison}{\ref{night:P4.7}}

There is a smooth compact complex threefold $C$ with
$b_1(C)=b_2(C)=0$, $e(C)=0$, and a holomorphic submersion
$C\to\Pp^1$ whose fibres are primary Hopf surfaces. No fibre has
a finite unramified torus cover. Thus the torus hypothesis is
necessary for the equality classification as stated in the prompt.
P4.7 proves it unnecessary for the numerical normal-crossings
sign theorem.
\end{proposition}
\begin{proof}
\noindent\textbf{Step 1.} The additive group $\C$ acts on
$V=(\C^2\setminus\{0\})\times(\C^2\setminus\{0\})$ by
\[
 \tau\cdot(z,w)=(e^\tau z,e^{i\tau}w).
\]
For $\tau=a+ib$, the logarithms of the two Euclidean norms change
by $a$ and $-b$. The action is free. It is proper: when both a
point and its translate lie in prescribed compact subsets of $V$,
the two norm ratios bound $a$ and $b$. Each orbit has a unique
representative with both norms one, obtained with
$a=-\log\|z\|$ and $b=\log\|w\|$. This gives an explicit
diffeomorphism $C=V/\C\cong S^3\times S^3$ and proves compactness.
Holomorphic transverse slices, for example $z_j=1$ after choosing
a local logarithm, give quotient charts with holomorphic
transitions. Hence $C$ is a smooth compact complex threefold.
The product cell complex gives free integral homology of ranks
$(1,0,0,2,0,0,1)$ and therefore the stated Betti and Euler numbers.

\noindent\textbf{Step 2.} The invariant map $[z,w]\mapsto[z_1:z_2]$ is a holomorphic
submersion to $\Pp^1$. In an affine base chart write
$z=\lambda(1,\xi)$ and set $\lambda=1$ using the action.
The choices of logarithm differ by $2\pi ik$, changing $w$ by
$e^{-2\pi k}$. Thus we obtain the explicit local fibre bundle
with fibre
$H=(\C^2\setminus\{0\})/\langle w\mapsto e^{-2\pi}w\rangle$.
Its universal cover retracts onto the simply connected $S^3$, so
$\pi_1(H)=\Z$. A connected finite unramified torus cover would
inject $\Z^4$ as a finite-index subgroup of $\Z$. Every subgroup
of $\Z$ is cyclic, whereas $\Z^4$ has rank four, which is
impossible. This proves the failure of the stated smooth-support
classification using the exact covering morphism.
\end{proof}

The executed program \texttt{code/part4/necessity\_examples.py}
checks the degree and squarefreeness of the discriminant, both
resolution substitutions, the Euler count $24$, and the product
cell ranks used above. Its complete output is supplied.

\subsection{Irreducible supports: canonical torsion and the normal case}

\begin{proposition}[P4.12]\label{night:P4.12}\label{stmt:P4.12}
\StatementRecord{night:P4.12}{proved}{Historical events: GPT-6 Astra (OpenAI), ChatGPT Pro mode, effort Pro, S6\_06\_BLIND, 16--18 November 2026: targeted proof check. Also checked in this writing. Pre-continuation writing events: GPT-6 Astra (OpenAI), ChatGPT Pro mode, effort Pro, S6\_10\_PAPERS, 18 November 2026: targeted proof check; used and rechecked in this run. Continuation metadata events: GPT-6 Astra (OpenAI), ChatGPT Pro mode, effort Pro, S6\_10\_PAPERS continuation, 19 November 2026: used and rechecked in this run; status-vocabulary reconciliation; metadata reconciliation only; no new mathematical check. \emph{Pre-continuation recorded proof scope.} The night-run re-derivation was by GPT-6 Astra (OpenAI, ChatGPT Pro mode, effort Pro), 16--18 November 2026. Also checked in this writing by GPT-6 Astra (OpenAI), ChatGPT Pro mode, effort Pro, 18 November 2026. The rank-one direct-image generator has divisor aS with 0<=a<m, so adjunction and O(mS)=O give torsion omega\_S. On the minimal resolution the principal-function intersection equation and the scaled quadratic identity prove negative definiteness; minimal-resolution adjunction gives discrepancies at most zero. Formal-functions comparison gives p=h0(R1 pi\_*O)>=h1(O\_Z) for every exceptional divisor Z. The canonical-trivial, rational, class-VII and positive-genus ruled alternatives each give the stated contradiction unless the normal support is smooth. For a nonnormal support the finite-duality conductor and the same matrix calculation give an effective nonzero pole divisor, hence Kodaira dimension minus infinity. Mumford's original section I p.6 and the CDP20 conductor locus are opened; no unrestricted class-VII classification is invoked}{\ref{night:P4.6}}

In the standing torus-fibration setting of this section, let $f^{-1}(b)=mS$
have irreducible reduced support. Then
\[
 \omega_S^{\otimes m}\simeq\OO_S,\qquad
 \chi(\OO_S)=0,\qquad h^1(\OO_S)\in\{1,2\}.
\]
Every normal such support is smooth. For a singular, necessarily
nonnormal support, a minimal resolved model of its normalization
has Kodaira dimension $-\infty$.
\end{proposition}
\begin{proof}
\noindent\textbf{Step 1. Canonical torsion.}
Choose a disc $\Delta$ with no other critical value, with coordinate
$t$ vanishing at $b$, and write $U=f^{-1}(\Delta)$. The coherent
direct image $f_*\omega_U$ is torsion-free of rank one: nonzero
base functions act injectively on sections, and a nearby torus has
one-dimensional canonical space. It is therefore a line bundle
on the disc. A generator gives a section $v$ of $\omega_U$,
nonzero and nowhere vanishing on every noncentral fibre. Its
divisor is $aS$ for an integer $a\ge0$. If $a\ge m$, then $v/t$
would be holomorphic, contradicting primitivity of the generator
at zero. We have $0\le a<m$ and
\[
 \omega_U\simeq\OO_U(aS),\qquad\OO_U(mS)\simeq\OO_U.
\]
Adjunction gives $\omega_S\simeq\OO_U((a+1)S)|_S$, proving
$\omega_S^{\otimes m}\simeq\OO_S$. The ambient
Riemann--Roch computation of P4.6, applied to $S$, gives
$\chi(\OO_S)=0$. Reduced irreducibility gives $h^0(\OO_S)=1$.
A nonzero section of a torsion line bundle has a power equal to
a nonzero constant, and hence is nowhere zero; its space of
sections consequently has dimension at most one. Serre duality
on the Gorenstein surface $S$ gives
$h^2(\OO_S)=h^0(\omega_S)\in\{0,1\}$, proving the assertion
about $h^1$.

\noindent\textbf{Step 2. Exceptional curves and discrepancies.}
For a minimal resolution $\pi:T\to S$ of a normal Gorenstein
surface write $K_T=\pi^*K_S+\sum a_jE_j$. Its exceptional
intersection matrix $M$ has nonnegative off-diagonal entries
and is negative definite. We recall the local analytic proof.
Pulling back a generic holomorphic function vanishing at a
contracted point gives $H+\sum m_jE_j$, with $m_j>0$ and $H$
the strict transform. Intersecting gives $Mm\le0$, strictly in
at least one entry of each connected block. After scaling by
$\operatorname{diag}(m_j)$, the negative quadratic form is a
sum of nonnegative terms $c_{ij}(z_i-z_j)^2$ and $d_i z_i^2$,
with connected support and some $d_i>0$ in every block. It is
positive definite. This is the quadratic identity of
\cite[\S I, unnumbered assertion and proof, p.~6]{Mumford1961}.
The displayed local holomorphic function and exceptional blocks give
all its inputs for these analytic germs.

Adjunction gives
$K_T\cdot E_j=2p_a(E_j)-2-E_j^2\ge0$: the sole negative
possibility is a smooth rational $(-1)$-curve, excluded by
minimality of the resolution. The inverse of a negative definite
matrix with nonnegative off-diagonal entries has nonpositive
entries. Indeed, put $Q=-M$ and choose $c$ larger than all its
eigenvalues and diagonal entries; then
$Q^{-1}=c^{-1}\sum_{k\ge0}(I-Q/c)^k$ is entrywise nonnegative.
The discrepancy equations $Ma=(K_T\cdot E_j)_j$ imply
$a_j\le0$. Hence $K_T=\pi^*K_S-E$ with $E$ an effective
exceptional integral divisor.

Put $p=\dim H^0(S,R^1\pi_*\OO_T)$. Leray, with this sheaf
supported at finitely many points, gives
\begin{equation}\label{eq:P4chi-resolution}
 \chi(\OO_T)=\chi(\OO_S)-p=-p.
\end{equation}
For any effective exceptional cycle $Z$, the higher direct image
$R^2\pi_*I_Z$ vanishes because all fibres of $\pi$ have dimension
at most one. The exact sequence of $Z$ gives a surjection
$R^1\pi_*\OO_T\to R^1\pi_*\OO_Z$, and therefore
\begin{equation}\label{eq:P4exceptional-genus}
 p\ge h^1(Z,\OO_Z).
\end{equation}

\noindent\textbf{Step 3. Exclusion of normal singular supports.}
Suppose first that $E=0$. Then $K_T$ is torsion and the
classification gives Kodaira dimension zero and nonnegative
holomorphic Euler characteristic. Equation
\eqref{eq:P4chi-resolution} forces $p=0$ and $\chi(\OO_T)=0$.
The remaining surfaces are tori, bielliptic surfaces and Kodaira
surfaces. They have no rational curves: a map from $\Pp^1$ to
a torus is constant, bielliptic surfaces have finite unramified
torus covers, and a primary Kodaira surface has an elliptic
fibration with smooth elliptic reduced fibres over an elliptic
curve. A map from $\Pp^1$ to a secondary Kodaira surface lifts
to its finite unramified primary Kodaira cover. On a torsion
canonical surface a negative curve is a smooth rational
$(-2)$-curve by adjunction. Thus there are no exceptional
curves, and $S$ is smooth.

Suppose next that $E\ne0$. Since $K_T+E$ is torsion, all
positive plurigenera vanish, so $T$ is rational, ruled, or a
blowup of a class-VII surface. Rationality contradicts
\eqref{eq:P4chi-resolution}, since $\chi(\OO_T)=1$. In class
VII that equation gives $p=0$, whereas
\[
 \chi(\OO_E)=-E\cdot(E+K_T)/2=0.
\]
The nonempty effective compact divisor $E$ has $h^0(\OO_E)\ge1$,
so $h^1(\OO_E)\ge1$, contradicting
\eqref{eq:P4exceptional-genus}.

Finally suppose $T$ is ruled, after point blowups, over a curve
$C$ of genus $g\ge1$. Equation \eqref{eq:P4chi-resolution}
gives $p=g-1$. A general ruling fibre $F$ has $K_T\cdot F=-2$,
so $E\cdot F=2$. An exceptional component $E_j$ therefore
maps onto $C$. Its normalization has genus at least $g$ by
Riemann--Hurwitz, so $h^1(\OO_{E_j})\ge g$. Taking $Z=E_j$
in \eqref{eq:P4exceptional-genus} gives $p\ge g$, a contradiction.
This exhausts the normal case.

\noindent\textbf{Step 4. The resolved normalization of a nonnormal support.}
Write $\nu:N\to S$ for the finite normalization and $\pi:T\to N$
for its minimal resolution. The residue comparison of dualizing
sheaves gives a pole divisor along the conductor; its strict
transform $D$ on $T$ is nonzero and effective. For completeness,
pull back a local generator of $\omega_S$ as a meromorphic
two-form. Its canonical divisor has the form $-D+\sum a_jE_j$.
The equations for the exceptional coefficients have right side
$(K_T+D)\cdot E_j\ge0$, because both $K_T\cdot E_j$ and
$D\cdot E_j$ are nonnegative. The same inverse-matrix argument
gives $a_j\le0$. Taking the $m$th power to trivialize the pulled
back $\omega_S$ consequently gives
\[
 mK_T+D'\sim0
\]
for a nonzero effective integral divisor $D'$. This is also the
conductor formula of \cite[Lemma~4.6, p.~7 of the attached-source edition]{CDP20}; its local proof does
not use the disputed global restriction conclusion. Any
positive pluricanonical section on $T$, multiplied by the
appropriate effective section and raised to a power, would give
a nonzero holomorphic function on compact $T$ with a zero.
It follows that $\kappa(T)=-\infty$. The classification \cite[Chapter~VI, Theorem~1.1, p.~243]{BHPV2004E}
gives precisely the stated alternatives for the minimal model.
\end{proof}

The ruled-genus and nonnormal-conductor arguments needed after
this reduction are proved below and completed in P4.58.

\begin{proposition}[P4.13]\label{night:P4.13}\label{stmt:P4.13}
\StatementRecord{night:P4.13}{proved}{Historical events: GPT-6 Astra (OpenAI), ChatGPT Pro mode, effort Pro, S6\_06\_BLIND, 16--18 November 2026: targeted proof check. Also checked in this writing. Pre-continuation writing events: GPT-6 Astra (OpenAI), ChatGPT Pro mode, effort Pro, S6\_10\_PAPERS, 18 November 2026: targeted proof check; used and rechecked in this run. Continuation metadata events: GPT-6 Astra (OpenAI), ChatGPT Pro mode, effort Pro, S6\_10\_PAPERS continuation, 19 November 2026: used and rechecked in this run; status-vocabulary reconciliation; metadata reconciliation only; no new mathematical check. \emph{Pre-continuation recorded proof scope.} The night-run re-derivation was by GPT-6 Astra (OpenAI, ChatGPT Pro mode, effort Pro), 16--18 November 2026. Also checked in this writing by GPT-6 Astra (OpenAI), ChatGPT Pro mode, effort Pro, 18 November 2026. The exceptional multiplicity bound, explicit local monodromy traces on every normal-crossing stratum, spectral-sequence alternating trace, quasi-unipotence and determinant sign and bound were re-derived. A'Campo1975 section2 pp237-243 and section3 pp243-245 were subsequently opened in primary-host indexed text, including the exact analytic hypotheses and stalk-trace proof. Direct visual-scan retrieval remains separately unavailable}{the hypotheses and explicit local calculation stated here}

Let a proper holomorphic degeneration of two-dimensional complex
tori have smooth total threefold and central fibre $mS$ with
irreducible reduced support. If $A$ is the monodromy on $H_1$ of
a nearby torus, then
\[
 m e(S_{\mathrm{reg}})=\det(I-A^m),\qquad
 0\le e(S_{\mathrm{reg}})\le16/m.
\]
This assertion requires no Betti-number hypothesis on a compactification.
\end{proposition}
\begin{proof}
\noindent\textbf{Step 1.} We resolve the pair by a proper holomorphic modification
$\pi:U'\to U$ which is an isomorphism away from $S$ and over
$S_{\mathrm{reg}}$. The total central divisor has normal crossings.
The strict transform has multiplicity $m$, and each exceptional
component has multiplicity at least $2m$. Indeed, its first centre
lies over a singular point of the reduced divisor, whose order is
at least two; subsequent centres either retain such an exceptional
multiplicity or add further positive multiplicities.

\noindent\textbf{Step 2.} We spell out the trace comparison. At a point on $k$ branches,
with local equation $z_1^{n_1}\cdots z_k^{n_k}=t$, the nearby
fibre retracts onto a disjoint union of real $(k-1)$-tori.
Monodromy permutes these components and translates any component
that it preserves. For $k\ge2$ the alternating cohomological
trace is zero. On a single branch of multiplicity $n$ there are
$n$ contractible nearby components in a cyclic permutation;
the $d$th iterate has trace $n$ when $n\mid d$, and zero otherwise.
Alternating traces in the nearby-cycle spectral sequence, with a
finite stratification of the compact central divisor, therefore give
\[
 L(A^d;H^*(T^4;\C))=\sum_{n_j\mid d}n_j e(E_j^\circ).
\]
This is the global form of A'Campo 1975, Theorem 1; his proper
analytic construction in Section 2 and spectral-sequence proof in
Section 3 give precisely this comparison. The explicit local
calculation verifies the hypotheses without a Kähler assumption.
For $d=m$ the strict transform is the sole contribution, and its
open stratum is $S_{\mathrm{reg}}$.

\noindent\textbf{Step 3.} These same charts prove quasi-unipotence. A common multiple $N$
of the $n_j$ makes monodromy act trivially on all local
nearby-cohomology sheaves, hence on the $E_2$ page. Its action
on the filtered abutment is consequently unipotent. Every
eigenvalue of $A$ is a root of unity. Since torus cohomology is
the exterior algebra on $H_1$, its Lefschetz number is
$\det(I-A^m)$. Complex conjugate eigenvalues contribute
$|1-\lambda^m|^2$; real eigenvalues are $1$ or $-1$ and
contribute zero or two. Thus this determinant is nonnegative
and at most $2^4$, as required.
\end{proof}

\begin{proposition}[P4.14]\label{night:P4.14}\label{stmt:P4.14}
\StatementRecord{night:P4.14}{proved}{Historical events: GPT-6 Astra (OpenAI), ChatGPT Pro mode, effort Pro, S6\_06\_BLIND, 16--18 November 2026: targeted proof check. Also checked in this writing. Pre-continuation writing events: GPT-6 Astra (OpenAI), ChatGPT Pro mode, effort Pro, S6\_10\_PAPERS, 18 November 2026: targeted proof check; used and rechecked in this run. Continuation metadata events: GPT-6 Astra (OpenAI), ChatGPT Pro mode, effort Pro, S6\_10\_PAPERS continuation, 19 November 2026: used and rechecked in this run; status-vocabulary reconciliation; metadata reconciliation only; no new mathematical check. \emph{Pre-continuation recorded proof scope.} The night-run re-derivation was by GPT-6 Astra (OpenAI, ChatGPT Pro mode, effort Pro), 16--18 November 2026. Also checked in this writing by GPT-6 Astra (OpenAI), ChatGPT Pro mode, effort Pro, 18 November 2026. For y(y-x\textasciicircum{}2)=0, the normalization map is (a,b) with b-a divisible by t\textasciicircum{}2; its inverse expression a+y(b-a)/x\textasciicircum{}2 gives delta two and conductor order two on both branches. For xy(x-y)=0, equality of constants gives two conditions and the linear relation gives one more; the displayed degree-k monomials separate every coefficient for k>=2, giving delta three and order-two conductor on all three branches. Plane-cubic adjunction proves trivial dualizing sheaf and arithmetic genus one. These are actual reducible curves contradicting the sentence in CDP20 Proposition7.3, attached-source p.16, without asserting a counterexample to the global threefold sign theorem}{the hypotheses and explicit local calculation stated here}

The reducible-curve inference in \cite[proof of Proposition~7.3,
p.~16 of the attached-source edition]{CDP20} omits the two plane cubics
\[
 Z(Y^2-XZ)=0,\qquad XY(X-Y)=0.
\]
They are reduced Gorenstein curves of arithmetic genus one with
trivial dualizing sheaf and rational normalization, but have
respectively two tangent components and three concurrent
components. Their conductor order is two on every normalization
branch. These curves refute the asserted nodal-cycle
classification at that step; they are not counterexamples to
the threefold sign assertion.
\end{proposition}
\begin{proof}
\noindent\textbf{Step 1.} Plane adjunction gives $\omega_C=\OO_C(3-3)=\OO_C$ and
$p_a(C)=1$. The first cubic has a line tangent to a smooth
conic at $[1:0:0]$. Its local equation is $y(y-x^2)=0$.
The normalization restrictions are $(p(t,0),p(t,t^2))$.
Their image is exactly the pairs $(a,b)$ with $b-a$ divisible
by $t^2$: the inverse expression is
$p(x,y)=a(x)+y(b(x)-a(x))/x^2$. Thus the normalization
quotient has length two and the conductor exponent is two on
each branch.

\noindent\textbf{Step 2.} For the second cubic the branches are $(t,0),(0,t),(t,t)$.
All constant coefficients must agree, and the third linear
coefficient is the sum of the first two. These three conditions
are sufficient. In every degree $k\ge2$, the monomials
$x^k-x^{k-1}y$, $y^k-xy^{k-1}$ and $x^{k-1}y$ independently
realize the coefficients on the three branches. Hence the
normalization quotient has length three and each conductor
exponent is again two. The normalization Euler characteristics
are two and three; subtracting these quotient lengths gives
$\chi(\OO_C)=0$ in both cases.

\noindent\textbf{Step 3.} The program \texttt{code/part4/genus\_one\_conductor.py}
executes the exact jet restriction matrices through orders
$2,3,4,6,8$, computing the ranks and conductor exponents.
When these branches vary in a surface, their permutation may
assemble them into a single normalization curve covering the
image conductor with degree two or three. Consequently a
unique horizontal conductor component with multiplicity two
does not alone force a degree-one map or generically reduced
image conductor. An additional argument is needed at that
precise point of the general irreducible case.
\end{proof}

The branch permutations in P4.14 have explicit local realizations.
Take an elliptic curve $C$ with a free translation of order two or
three. The map $(x,y)\mapsto(x,x^2-y)$ exchanges the two tacnode
branches and preserves $y(y-x^2)$. The order-three map
$(x,y)\mapsto(-y,x-y)$ cycles the three branches and preserves
$xy(x-y)$. The diagonal action on $C$ times a small invariant
neighbourhood of the origin is free. Its quotient is a smooth
threefold; the divisor is irreducible because the action permutes
its branches transitively. Its normalization conductor is one
elliptic curve mapping to $C/\langle\tau\rangle$ with degree two
or three. These specify the finite morphisms in P4.14 explicitly;
they make no assertion of a compact torus smoothing.

\begin{proposition}[P4.15]\label{night:P4.15}\label{stmt:P4.15}
\StatementRecord{night:P4.15}{proved}{Historical events: GPT-6 Astra (OpenAI), ChatGPT Pro mode, effort Pro, S6\_06\_BLIND, 16--18 November 2026: targeted proof check. Also checked in this writing. Pre-continuation writing events: GPT-6 Astra (OpenAI), ChatGPT Pro mode, effort Pro, S6\_10\_PAPERS, 18 November 2026: targeted proof check. Continuation metadata events: GPT-6 Astra (OpenAI), ChatGPT Pro mode, effort Pro, S6\_10\_PAPERS continuation, 19 November 2026: used and rechecked in this run; status-vocabulary reconciliation; metadata reconciliation only; no new mathematical check. \emph{Pre-continuation recorded proof scope.} The night-run re-derivation was by GPT-6 Astra (OpenAI, ChatGPT Pro mode, effort Pro), 16--18 November 2026. Also checked in this writing by GPT-6 Astra (OpenAI), ChatGPT Pro mode, effort Pro, 18 November 2026. The actual derivative is a sheaf injection of rank two into rank three. The depth lemma excludes a zero-dimensional torsion subsheaf in its cokernel, yielding I\_Z tensor L and e(T)=K\_T\textasciicircum{}2-length(Z). Substitution in each ruled and class-VII Chern formula gives the stated bounds and minimal zero cases. Adjunction on a minimal elliptic ruled surface makes every coefficient-one horizontal conductor smooth elliptic and excludes vertical poles; the doubled-section case uses the independently proved nonnegative Euler number of the open support. On a Hopf surface rational curves are excluded by their lift to C\textasciicircum{}2 minus zero. At a finite immersion every local pair of surface branches intersects along the same smooth conductor curve, so the actual branch maps are finite unramified elliptic maps and the support Euler number is zero}{\ref{night:P4.12}; \ref{night:P4.13}; \ref{night:P4.8}}

Let $mS$ be an irreducible singular fibre in the setting of 4.2.
Suppose its normalization $T$ is smooth and $d\nu:T_T\to\nu^*T_X$
has rank two away from finitely many points. Let $\ell$ be the
length of its maximal-minor scheme. Then
\[
 e(T)=K_T^2-\ell.
\]
The alternatives for $T$ are a rational surface with
$e(T)=6-\ell/2$ and $0\leq\ell\leq6$, a minimal ruled surface
over an elliptic curve with $\ell=0$, or a Hopf surface with
$\ell=0$. In either of the last two alternatives $e(S)=0$.
The reduced conductor and singular locus are disjoint unions of
smooth elliptic curves, and their normalization maps are finite
and unramified.
\end{proposition}
\begin{proof}
\noindent\textbf{Step 1.} The differential is injective as a sheaf map: its kernel vanishes
on a dense open set, and a locally free sheaf is torsion-free.
Its rank-one cokernel $Q$ is locally free away from finitely many
points. The depth lemma applied to the length-one locally free
resolution gives depth at least one at those points, excluding
a finite-length subsheaf. All torsion would be supported there,
so $Q$ is torsion-free. On the smooth surface we obtain
\[
 Q=\mathcal I_Z\otimes L,\qquad
 L=K_T\otimes\nu^*(\det T_X),\qquad \operatorname{length}(Z)=\ell.
\]
The two-column presentation by $d\nu$ identifies the local ideal
of $Z$ with its three maximal minors. The classes $c_1(T_X)$ and
$c_2(T_X)$ vanish in real cohomology, since $b_2(X)=b_4(X)=0$.
The second Chern class of the exact sequence therefore gives
\[
 0=e(T)-K_T^2+\ell.
\]
P4.12 gives $\kappa(T)=-\infty$. For a ruled surface over a
genus-$g$ curve, blown up $k$ times, we have
\[
 e(T)=4(1-g)+k,\qquad K_T^2=8(1-g)-k.
\]
Consequently $4(g-1)+2k+\ell=0$. Thus $g\leq1$, and $g=1$
forces $k=\ell=0$. In the class-VII case, including blowups,
$K_T^2=-e(T)$ and $e(T)=b_2(T)\geq0$. Hence
$2e(T)+\ell=0$, and the surface is minimal with $b_2=0$.
Its nonempty conductor, the finite-cover classification, and the
absence of curves on Inoue surfaces proved in P4.8 force it to
be Hopf. For rational $T$, Noether's formula gives
$2e(T)+\ell=12$. The inequality $e(T)\geq3$ gives the asserted
range. When $\ell=0$, the cokernel is a line bundle and $d\nu$
has rank two everywhere.

\noindent\textbf{Step 2.} We now compute the image Euler number in the two latter cases.
The conductor divisor $D$ satisfies $K_T+D\equiv0$. First let
$T$ be a minimal ruled surface over an elliptic curve $E$;
it is a $\mathbb P^1$-bundle. For a ruling fibre $F$, $D\cdot F=2$.
If a horizontal component $H$ occurs with multiplicity one,
adjunction gives
\[
 2p_a(H)-2=(K_T+H)H=-(D-H)H\leq0.
\]
The map to $E$ gives geometric genus at least one. Equality
therefore holds, $H$ is smooth elliptic, and $H$ is disjoint
from every other component of $D$. There is no vertical
component because every ruling fibre meets $H$. The possibilities
are one elliptic double-section or two disjoint elliptic sections.

\noindent\textbf{Step 3.} The other possibility is $D=2H+\sum_j a_jF_j$, where $H$ is a
section, $F_j$ are distinct ruling fibres and $a_j>0$.
The reduced conductor has Euler number equal to the number $k$
of these fibres: $e(H)=0$, $e(F_j)=2$, and each meets $H$ once.
The immersive normalization identifies $T\setminus D_{\mathrm{red}}$
with $S_{\mathrm{reg}}$. Hence
\[
 e(S_{\mathrm{reg}})=e(T)-e(D_{\mathrm{red}})=-k.
\]
P4.13 forces $k=0$. Thus in every case the reduced conductor is
a disjoint union of smooth elliptic curves.

\noindent\textbf{Step 4.} On a Hopf surface every irreducible curve $H$ satisfies
$H^2=K_TH=0$, so $p_a(H)=1$. A rational curve would lift from
$\mathbb P^1$ to the universal cover $\mathbb C^2\setminus\{0\}$,
where both coordinate functions must be constant. Thus $H$
is smooth elliptic. Distinct curves have intersection zero and
are disjoint, proving the same conductor description.

\noindent\textbf{Step 5.} To check the image, at a singular point the finite immersive
normalization presents the germ as distinct smooth surface
branches $S_i$. The reduced conductor on each branch is a single
smooth curve germ $C_i$. Any two branches intersect in a curve
in the smooth ambient threefold, and its support lies in both
$C_i$ and $C_j$. Their irreducibility forces those curve germs
to agree. Thus the image singular locus is one smooth curve
germ, and the conductor maps locally isomorphically to it.
Globally this map is finite and unramified. Riemann--Hurwitz
then makes every image component elliptic. Finally
\[
 e(S)=e(T)-e(D_{\mathrm{red}})+e(\operatorname{Sing}S)=0.
\]
\end{proof}

\begin{proposition}[P4.16]\label{night:P4.16}\label{stmt:P4.16}
\StatementRecord{night:P4.16}{proved}{Historical events: GPT-6 Astra (OpenAI), ChatGPT Pro mode, effort Pro, S6\_06\_BLIND, 16--18 November 2026: targeted proof check. Also checked in this writing. Pre-continuation writing events: GPT-6 Astra (OpenAI), ChatGPT Pro mode, effort Pro, S6\_10\_PAPERS, 18 November 2026: targeted proof check. Continuation metadata events: GPT-6 Astra (OpenAI), ChatGPT Pro mode, effort Pro, S6\_10\_PAPERS continuation, 19 November 2026: used and rechecked in this run; status-vocabulary reconciliation; metadata reconciliation only; no new mathematical check. \emph{Pre-continuation recorded proof scope.} The night-run re-derivation was by GPT-6 Astra (OpenAI, ChatGPT Pro mode, effort Pro), 16--18 November 2026. Also checked in this writing by GPT-6 Astra (OpenAI), ChatGPT Pro mode, effort Pro, 18 November 2026. The explicit cover of x\textasciicircum{}2-zy\textasciicircum{}n=1 has U=C* times C and two contractible graph neighborhoods, with each overlap C* times a disc. Van Kampen and Mayer-Vietoris give pi1=0 and H2=Z. The compactification performs n blowups over each of (1,0) and (-1,0); its rational-tree boundary makes the compact class primitive. Orthogonality to that boundary gives the class sum E\_plus minus sum E\_minus, of square -2n. Positive weights (n+2,2,4) carry compact representatives and transverse companions into an actual local fibre. The weighted holomorphic unit-root map removes a unit exactly, and the full-fibre root comparison embeds the resulting spherical class into a nearby aspherical torus, where its nonzero square is impossible. Archived numerical runs are not counted as current computations}{the hypotheses and explicit local calculation stated here}

For every integer $n\geq1$, the reduced hypersurface germ
$x^2-zy^n=0$ cannot occur in a fibre of a proper holomorphic
degeneration of complex two-tori with smooth total threefold,
even with any positive common multiplicity on that germ.
In particular the Whitney umbrella ($n=2$) is excluded.
No hypothesis on $b_2$ is needed.
\end{proposition}
\begin{proof}
\noindent\textbf{Step 1.} Put $F_n=\{x^2-zy^n=1\}\subset\mathbb C^3$. We exhibit its
compact spherical homology class and compute its square.
The open subset $U=\{y\ne0\}$ is
$\mathbb C^*\times\mathbb C$ in coordinates $(y,x)$.
For $0<\delta<1$, take the two disjoint graph neighbourhoods
\[
 N_\pm:\quad x=\pm\sqrt{1+zy^n},\quad |zy^n|<\delta.
\]
Their $(y,z)$ domains contract by simultaneous real scaling.
Each $N_\pm\cap U$ has coordinates $(y,zy^n)$ in
$\mathbb C^*\times\{|w|<\delta\}$. Its map to $U$ takes
the generating $y$-loop to the generating $y$-loop. These three
sets cover $F_n$. Van Kampen gives $\pi_1(F_n)=0$.
Mayer--Vietoris identifies $H_2(F_n;\mathbb Z)$ with the
kernel of $(a,b)\mapsto a+b$ on the two intersection
$H_1$ groups, hence with $\mathbb Z$. Hurewicz represents its
generator by a map $S^2\to F_n$.

\noindent\textbf{Step 2.} We compute the intersection form by compactification. Blow up
$\mathbb P^1_x\times\mathbb P^1_y$ successively $n$ times above
each point $(1,0)$ and $(-1,0)$. Near either point let
$a=x^2-1$, $b=y$. The successive charts at the centres are
\[
 a=b^j a_j,\qquad z=\frac{a_j}{b^{n-j}}.
\]
For $j<n$, blow up the next centre $a_j=b=0$.
At $j=n$, $z=a_n$ is regular. The final exceptional curve
minus its point on the preceding boundary is precisely the
line $y=0,x=\pm1$ in $F_n$. Every earlier exceptional curve
is a pole divisor of $z$. Consequently $F_n=T_n\setminus B$,
where $T_n$ is this smooth projective blown-up surface and
$B$ consists of the two infinity divisors, the strict transform
of $y=0$, and the first $n-1$ exceptional curves on each chain.
Both infinity divisors are removed, so no extension of the
rational function at their intersection is asserted or needed.

\noindent\textbf{Step 3.} In the total-transform basis $H_x,H_y,E_{\pm,j}$ the nonzero
products are $H_xH_y=1$ and $E_{\pm,j}^2=-1$. The boundary
classes are
\[
 H_x,\ H_y,\ H_y-E_{+,1}-E_{-,1},\quad
 E_{+,j}-E_{+,j+1},\ E_{-,j}-E_{-,j+1}\quad(1\leq j<n).
\]
They form a connected tree of rational curves; thus $H_1(B)=0$.
Compact-support duality and the sequence of $(T_n,B)$ identify
the image of $H_2(F_n)$ in $H_2(T_n)$ with the kernel of
intersection against these boundary classes. Orthogonality to
$H_x,H_y$ removes both ruling coefficients. Orthogonality to
the exceptional differences makes the coefficients constant
along each chain. The strict transform of $y=0$ makes the two
constants opposite. Thus the kernel is generated by the
primitive integral vector
\[
 v=\sum_{j=1}^n E_{+,j}-\sum_{j=1}^n E_{-,j},\qquad v^2=-2n.
\]
The sequence introduces no multiple: $H_1(B)=0$ injects
$H^2(T_n,B)$ into $H^2(T_n)$ with precisely the stated kernel.
Therefore the compact spherical generator has square $-2n$.
The exact matrices, primitive kernels, boundary graphs and
blowup identities for $n=1,\ldots,12$ are independently checked
by \texttt{code/part4/milnor\_intersection.py}; its output and
JSON data are included. The proof is uniform in $n$.

\noindent\textbf{Step 4.} The defining polynomial is weighted homogeneous of degree
$d=2n+4$ with positive weights $(n+2,2,4)$. Scaling by
\[
 (x,y,z)\longmapsto
 (\varepsilon^{n+2}x,\varepsilon^2y,\varepsilon^4z)
\]
identifies $F_n$ with the fibre of value $\varepsilon^{2n+4}$.
It carries any compact representative of $v$, together with a
compact transverse representative for its square, into an
arbitrarily small neighbourhood of the origin.

\noindent\textbf{Step 5.} We also check the unit and multiplicity in an actual defining
map. Take an analytic $m$-th root of its unit. For a unit $h$,
the local analytic isomorphism
\[
 \Phi(x,y,z)=
 \bigl(h^{(n+2)/d}x,h^{2/d}y,h^{4/d}z\bigr)
\]
satisfies $g\circ\Phi=hg$, where $g=x^2-zy^n$; its derivative
at the origin is invertible. Thus the actual local map to the
base is $(g\circ\Phi)^m$. The small local fibre of $g$ embeds
in the corresponding smooth compact torus fibre.

\noindent\textbf{Step 6.} The spherical representative is null-homotopic in that torus,
whose universal cover is the contractible space $\mathbb C^2$.
It is therefore zero in $H_2$, and its self-intersection there
is zero. The two compact representatives already chosen compute
the same intersection in the open local fibre and in the torus.
Their intersection is $-2n$, which is a contradiction.
\end{proof}

P4.16 excludes the displayed pinching and ramification germs.
The subsequent arguments retain singular normalizations and
nonimmersive normalization maps and complete their classification
in P4.58.

\begin{proposition}[P4.17]\label{night:P4.17}\label{stmt:P4.17}
\StatementRecord{night:P4.17}{proved}{Historical events: GPT-6 Astra (OpenAI), ChatGPT Pro mode, effort Pro, S6\_06\_BLIND, 16--18 November 2026: targeted proof check. Also checked in this writing. Pre-continuation writing events: GPT-6 Astra (OpenAI), ChatGPT Pro mode, effort Pro, S6\_10\_PAPERS, 18 November 2026: targeted proof check. Continuation metadata events: GPT-6 Astra (OpenAI), ChatGPT Pro mode, effort Pro, S6\_10\_PAPERS continuation, 19 November 2026: used and rechecked in this run; status-vocabulary reconciliation; metadata reconciliation only; no new mathematical check. \emph{Pre-continuation recorded proof scope.} The night-run re-derivation was by GPT-6 Astra (OpenAI, ChatGPT Pro mode, effort Pro), 16--18 November 2026. Also checked in this writing by GPT-6 Astra (OpenAI), ChatGPT Pro mode, effort Pro, 18 November 2026. With a graph neighborhoods the same exact cover gives H2=ker(Z\textasciicircum{}a to Z) and a simply connected fibre. The a exceptional chains in the compactification yield the primitive differences v\_i=sum E\_i,j-sum E\_i+1,j; their intersection form is minus n times the Cartan matrix A\_(a-1). The positive-weight and unit-root maps give the local torus contradiction. The finite normalization map (u,v) to (u\textasciicircum{}n v,u\textasciicircum{}a,v\textasciicircum{}a) has its exact mu\_a action and degree-a function-field extension, so the normal invariant ring is the normalization. For (a,n)=(2,3) it is UW=V\textasciicircum{}2 with x=UV,y=U,z=W; for n=a the conductor map is v to v\textasciicircum{}a. This proves the exclusions through the displayed morphisms}{\ref{night:P4.16}}

For all $a\geq2$ and $n\geq1$, the germ $x^a-zy^n=0$ is
excluded in the setting of P4.16. Its affine nearby fibre has
$\pi_1=0$ and compact second-homology intersection lattice
$A_{a-1}(-n)$. This includes the order-three cyclic-ramification
germ $x^3-zy^3$ and examples with singular normalization.
\end{proposition}
\begin{proof}
\noindent\textbf{Step 1.} Replace the two roots of $x^2=1$ in P4.16 by the $a$ distinct
roots of $x^a=1$. The open set $y\ne0$ is still
$\mathbb C^*\times\mathbb C$. There are $a$ contractible graph
neighbourhoods, each meeting it in $\mathbb C^*\times\Delta$
with the same generating $y$-loop. Thus van Kampen gives
$\pi_1=0$, and Mayer--Vietoris gives
\[
 H_2(F_{a,n};\mathbb Z)=\ker\bigl(\mathbb Z^a
   \xrightarrow{(c_i)\mapsto\sum_i c_i}\mathbb Z\bigr).
\]
At each root, $x^a-1$ is a local coordinate. The same $n$
successive blowup charts give a compactification with $a$
exceptional chains and boundary classes
\[
 H_x,\ H_y,\ H_y-\sum_i E_{i,1},\qquad
 E_{i,j}-E_{i,j+1}\quad(1\leq i\leq a,\ 1\leq j<n).
\]
This boundary is a tree of rational curves. The exact sequence
of the pair identifies compact $H_2$ integrally with the kernel
of intersection against these classes. Its primitive basis is
\[
 v_i=\sum_{j=1}^n E_{i,j}-\sum_{j=1}^n E_{i+1,j}
       \quad(1\leq i<a).
\]
The products are $v_i^2=-2n$, $v_iv_{i+1}=n$, with all other
off-diagonal products zero. This is $A_{a-1}(-n)$.
Hurewicz makes every class spherical. Positive weights
$(n+1,a,a)$, of degree $a(n+1)$, scale compact representatives
into arbitrarily small nearby fibres. The explicit unit removal
and torus argument of P4.16 now exclude $v_1$, whose square
is $-2n$. The program independently checks the complete
matrices and primitive bases for $2\leq a\leq6$, $1\leq n\leq6$.

\noindent\textbf{Step 2.} We specify the normalization morphism. Let $\mu_a$ act on
$\mathbb C^2_{u,v}$ by $(u,v)\mapsto(\zeta u,\zeta^{-n}v)$.
Then
\[
 x=u^nv,\qquad y=u^a,\qquad z=v^a
\]
gives a finite map to $x^a=zy^n$: both $u$ and $v$ satisfy
monic equations over its coordinate ring. The function-field
degree is $a$, with the displayed faithful $\mu_a$ action.
Consequently the induced quotient map is finite and birational.
The invariant ring is integrally closed, so this is the
normalization. For $(a,n)=(2,3)$ it has coordinates
$U=u^2,V=uv,W=v^2$ with $UW=V^2$, and its map is
$(x,y,z)=(UV,U,W)$. Thus it has an $A_1$ singularity.
For $n=a$ the quotient is smooth with coordinates $(y,v)$,
and its conductor map is $v\mapsto v^a$. This verifies
explicitly the cyclic branch-permutation description.
\end{proof}

\begin{proposition}[P4.18]\label{night:P4.18}\label{stmt:P4.18}
\StatementRecord{night:P4.18}{proved}{Historical events: GPT-6 Astra (OpenAI), ChatGPT Pro mode, effort Pro, S6\_06\_BLIND, 16--18 November 2026: targeted proof check. Also checked in this writing. Pre-continuation writing events: GPT-6 Astra (OpenAI), ChatGPT Pro mode, effort Pro, S6\_10\_PAPERS, 18 November 2026: targeted proof check. Continuation metadata events: GPT-6 Astra (OpenAI), ChatGPT Pro mode, effort Pro, S6\_10\_PAPERS continuation, 19 November 2026: used and rechecked in this run; status-vocabulary reconciliation; metadata reconciliation only; no new mathematical check. \emph{Pre-continuation recorded proof scope.} The night-run re-derivation was by GPT-6 Astra (OpenAI, ChatGPT Pro mode, effort Pro), 16--18 November 2026. Also checked in this writing by GPT-6 Astra (OpenAI), ChatGPT Pro mode, effort Pro, 18 November 2026. Writing b=dB,c=dC and Bp+Cq=1 gives the literal punctured-fibre coordinates y=h\textasciicircum{}p t\textasciicircum{}C,z=h\textasciicircum{}q t\textasciicircum{}-B, with t=y\textasciicircum{}q z\textasciicircum{}-p. Over a disc joining two simple zeros of x\textasciicircum{}a-1, the two graph-neighborhood loop maps have matrix [[B,C],[q,-p]], determinant -1, and kill both generators. The finite map (X,y,z) to (X,y,z\textasciicircum{}C) has degree C; pullback of the compact Poincare-dual class of square -2B has square -2BC. Hurewicz, positive weights and the exact unit-root comparison place a nonzero spherical square in an actual nearby torus, giving the contradiction. The separate mixed example normalizes by (u,v) to (u\textasciicircum{}3 v\textasciicircum{}2,v,u\textasciicircum{}2); its monomial ring gives conductor u\textasciicircum{}2 v\textasciicircum{}2 and its derivative minors have common factor u, verifying the asserted positive-dimensional rank drop}{\ref{night:P4.15}; \ref{night:P4.16}}

For $a\geq2$ and $b,c\geq1$, the reduced germ
$x^a-y^bz^c=0$ cannot occur with a common positive multiplicity
in a torus degeneration with smooth total threefold. Its affine
nearby fibre contains a spherical class of nonzero square.
In particular this excludes $x^2-z^3y^4$, whose smooth
normalization has a divisorial differential rank drop.
\end{proposition}
\begin{proof}
\noindent\textbf{Step 1.} Write $d=\gcd(b,c)$, $b=dB$, $c=dC$, so $\gcd(B,C)=1$.
The nearby fibre $F=\{x^a-y^bz^c=1\}$ maps to the smooth
affine curve $A=\{x^a-w^d=1\}$ by $w=y^Bz^C$.
The function $w$ on $A$ has $a$ simple zeros. Choose an
embedded arc between two zeros avoiding the others, and a
small open disc $D$ around it. In a holomorphic coordinate
$\xi$ on $D$, write $w=u(\xi)(\xi-\alpha)(\xi-\beta)$
with $u$ a unit. Taking a $B$-th root of $u$, rescaling $y$,
and making an affine coordinate change identifies the preimage
of $D$ with
\[
 W=\{y^Bz^C=X^2-1,\ X\in D'\},
\]
where $D'$ is a disc containing both $-1$ and $1$.

\noindent\textbf{Step 2.} Choose $Bp+Cq=1$. With $h=X^2-1$, the explicit trivialization
away from the two special fibres is
\[
 y=h^pt^C,\quad z=h^qt^{-B},\quad t=y^qz^{-p}.
\]
Its fundamental group is $F(\beta_-,\beta_+)\times\mathbb Z\tau$.
The neighbourhood of each special fibre is the graph
$X=\pm\sqrt{1+y^Bz^C}$ on $|y^Bz^C|<\delta$, a contractible
domain by real scaling. Its punctured part has two generating
loops in $y,z$, whose images in $\langle\beta_\pm,\tau\rangle$
are the columns of
\[
 \begin{pmatrix}B&C\\q&-p\end{pmatrix},\qquad \det=-1.
\]
Thus adjoining each neighbourhood kills $\beta_\pm$ and $\tau$.
Van Kampen gives $\pi_1(W)=0$.

\noindent\textbf{Step 3.} The map
\[
 \pi:W\longrightarrow W'=\{y^BZ=X^2-1,\ X\in D'\},
 \qquad (X,y,z)\longmapsto(X,y,z^C)
\]
is proper finite of degree $C$. The compact class $v$ in
P4.16, of square $-2B$, is already supported over $D'$.
Indeed, the cover just given with $C=1$ has the identical
Mayer--Vietoris maps before and after enlarging $D'$ to the
whole plane: both punctured bases have the same two meridians,
and the same two graph neighbourhoods are attached. The
inclusion is therefore an isomorphism on $H_2$ and preserves
the primitive class and its square.
For its compact-support dual $\alpha\in H_c^2(W';\mathbb Z)$,
properness and degree give
\[
 \int_W(\pi^*\alpha)^2=C\int_{W'}\alpha^2=-2BC.
\]
One may verify this by smooth compactly supported forms:
outside the measure-zero branch divisor the map is an oriented
$C$-sheeted covering. The dual of $\pi^*\alpha$ is spherical
by $\pi_1(W)=0$ and Hurewicz. Its compact representatives
embed in $F$ with the same nonzero square.

\noindent\textbf{Step 4.} The polynomial has positive weights $(b+c,a,a)$. Weighted
scaling places these representatives in arbitrarily small
nearby fibres. P4.16's explicit unit removal handles the unit
and common multiplicity. The spherical class becomes zero
in a torus, contradicting its square $-2BC$.

\noindent\textbf{Step 5.} For the mixed example the normalization is the finite
birational map
\[
 (u,v)\longmapsto(u^3v^2,v,u^2).
\]
Finiteness follows from $u^2=z$, and $u=x/(zy^2)$ gives
birationality. The image ring contains all even powers of $u$
and the odd monomials $u^iv^j$ with $i\geq3,j\geq2$.
Its conductor is $u^2v^2\mathbb C\{u,v\}$, giving divisor
$2\{u=0\}+2\{v=0\}$. Its differential minors are
$3u^2v^2,-4u^4v,-2u$, with common divisorial factor $u$.
Thus this example explicitly lies outside P4.15's finite
rank-drop assumption. The 64 exponent matrices for
$1\leq b,c\leq8$, transfer squares and normalization
identities are executed in
\texttt{code/part4/binomial\_milnor.py}; its output and
JSON data are saved alongside it.
\end{proof}

\begin{proposition}[P4.19]\label{night:P4.19}\label{stmt:P4.19}
\StatementRecord{night:P4.19}{proved}{Historical events: GPT-6 Astra (OpenAI), ChatGPT Pro mode, effort Pro, S6\_06\_BLIND, 16--18 November 2026: targeted proof check. Also checked in this writing. Pre-continuation writing events: GPT-6 Astra (OpenAI), ChatGPT Pro mode, effort Pro, S6\_10\_PAPERS, 18 November 2026: targeted proof check. Continuation metadata events: GPT-6 Astra (OpenAI), ChatGPT Pro mode, effort Pro, S6\_10\_PAPERS continuation, 19 November 2026: used and rechecked in this run; status-vocabulary reconciliation; metadata reconciliation only; no new mathematical check. \emph{Pre-continuation recorded proof scope.} The night-run re-derivation was by GPT-6 Astra (OpenAI, ChatGPT Pro mode, effort Pro), 16--18 November 2026. Also checked in this writing by GPT-6 Astra (OpenAI), ChatGPT Pro mode, effort Pro, 18 November 2026. The plane and Whitney umbrella normalization maps, conductor ideals x squared and v squared, ramification at the origin, and the actual nonzero-fibre isomorphism with C-star times C were re-derived. The comparison cites the separately stated sphere obstructions and does not assert a compact counterexample}{\ref{night:P4.16}; \ref{night:P4.18}}

The germ $y(x^2-zy^2)=0$ has ramified conductor, of multiplicity
two on each normalization component, but its affine nearby
fibre is $\mathbb C^*\times\mathbb C$. Ramification alone
therefore does not give the spherical intersection obstruction
of P4.16--P4.18. This is a local statement, not a compact
counterexample to 4.2.
\end{proposition}
\begin{proof}
\noindent\textbf{Step 1.} One normalization component is the plane $y=0$; the other is
the Whitney normalization $(u,v)\mapsto(uv,v,u^2)$.
The singular locus is the smooth $z$-axis. The plane conductor
is $x^2=0$, obtained by restricting the other component's
equation. On the Whitney component the intrinsic conductor
is $v\mathbb C\{u,v\}$; multiplication by the added plane
equation $y=v$ gives $v^2\mathbb C\{u,v\}$. Directly, the
union ring is the fibre product of the component rings over
$\mathbb C\{x,z\}/(x^2)$, giving exactly these conductor ideals.
The reduced conductor cover is the identity sheet from the
plane together with $u\mapsto u^2$ from the Whitney component.
The latter is ramified.

\noindent\textbf{Step 2.} On $y(x^2-zy^2)=1$ we have $y\ne0$, and the inverse maps
\[
 (x,y,z)\longmapsto(y,x),\qquad
 (y,x)\longmapsto\left(x,y,\frac{yx^2-1}{y^3}\right)
\]
identify this fibre with $\mathbb C^*\times\mathbb C$.
Its $\pi_2$ and $H_2$ vanish. Hence every compact spherical
class in an open subset has zero square after inclusion,
and therefore has zero square already in that open subset.
\end{proof}

\begin{proposition}[P4.20]\label{night:P4.20}\label{stmt:P4.20}
\StatementRecord{night:P4.20}{proved}{Historical events: GPT-6 Astra (OpenAI), ChatGPT Pro mode, effort Pro, S6\_06\_BLIND, 16--18 November 2026: targeted proof check. Also checked in this writing. Pre-continuation writing events: GPT-6 Astra (OpenAI), ChatGPT Pro mode, effort Pro, S6\_10\_PAPERS, 18 November 2026: targeted proof check. Continuation metadata events: GPT-6 Astra (OpenAI), ChatGPT Pro mode, effort Pro, S6\_10\_PAPERS continuation, 19 November 2026: used and rechecked in this run; status-vocabulary reconciliation; metadata reconciliation only; no new mathematical check. \emph{Pre-continuation recorded proof scope.} The night-run re-derivation was by GPT-6 Astra (OpenAI, ChatGPT Pro mode, effort Pro), 16--18 November 2026. Also checked in this writing by GPT-6 Astra (OpenAI), ChatGPT Pro mode, effort Pro, 18 November 2026. Boundary-preserving deformations, the finite flat critical-rank scheme, the exact crosscap normal form, and transport into the actual torus were re-derived. For immersive sheets, dependent gradients give an explicit nonzero critical value and an arbitrarily small quadratic Morse perturbation, excluding non-NC intersections. The final global applications retain P4.6/P4.8 hypotheses}{\ref{night:P4.16}; \ref{night:P4.6}; \ref{night:P4.8}}

Let $mS$ be an irreducible fibre of a proper holomorphic
degeneration of complex two-tori with smooth total threefold.
Suppose the normalization of $S$ is smooth and, outside
finitely many points, $S$ is smooth or has the ordinary-double
germ $xy=0$. Then $S$ has normal crossings everywhere.
In the global setting of 4.2, P4.6--P4.8 apply to this subclass.
\end{proposition}
\begin{proof}
\noindent\textbf{Step 1.} We give the boundary and transport arguments explicitly.
Take a small ball $B$ around an exceptional point, with
boundary transverse to the smooth and double strata away
from that point. Let $g$ be a reduced equation. An analytic
$m$-th root of the unit in the actual defining function lets
us write that function as $g^m$.

\noindent\textbf{Step 2.} Consider a small holomorphic deformation $G_s$ of $g$ whose
zero set near $\partial B$ retains its smoothly moving smooth
and ordinary-double strata, transverse to $\partial B$.
There is a uniform $\delta>0$ such that
$G_s|_{\partial B}$ is submersive over $0<|t|<\delta$.
At a boundary point on a double curve, choose coordinates
$a,b,c$ with $G_s=ab$, absorbing a unit in $a$.
Transversality of the double curve to $\partial B$ allows one
real $c$-coordinate to be solved from the boundary equation.
The complex coordinates $a,b$ remain free there, so $ab$ is
submersive when nonzero. At a smooth point use $G_s=a$.
Compactness gives uniform bounds; points outside these
neighbourhoods stay away from zero.

\noindent\textbf{Step 3.} For each fixed small $s$ there are finitely many nonzero
critical values in this disc. After shrinking, the analytic
critical set in a neighbourhood of the compact closed ball
has finitely many relevant local components. The function is constant on
each analytic critical component. A component at a nonzero
small value cannot meet $\partial B$, so it is compact;
the maximum principle excludes a positive-dimensional
compact analytic subset of a ball. Thus the relevant critical
points are isolated and finite. Paths avoiding these values
join a small regular value to one fixed regular value $t_0$.
Proper Ehresmann transport on the closed ball, with the
verified boundary submersion, identifies the fibres along
such paths. For sufficiently small deformation the fibres
$G_s=t_0$ also form a smooth proper family to $s=0$, by
compactness. Hence every compact spherical class of nonzero
square in a local nearby fibre of $G_s$ transports into
$g=t_0$. Its inclusion in the actual torus fibre is impossible.

\noindent\textbf{Step 4.} At a normalization preimage write the map germ as
$\phi:(\mathbb C^2,0)\to(\mathbb C^3,0)$. Away from the
exceptional point it is immersive with transverse multiple
images. If $d\phi$ drops rank, its maximal minors define a
nonempty finite scheme. Small generic polynomial
perturbations have transverse rank-one points near it.
We check both genericity and conservation of their positive
number. Varying linear coefficients makes the universal
derivative map submersive onto the $3$-by-$2$ matrix space.
The rank-one locus is smooth of codimension two; the rank-zero
locus has codimension six. Generic perturbations are transverse
to the first and avoid the second. The height-two maximal-minor
ideal has its Hilbert--Burch resolution. Its one-parameter
zero scheme is Cohen--Macaulay of dimension one; the finite
special fibre excludes a vertical component. Thus it is flat
and preserves the positive length. Further generic polynomial
coefficients keep the finitely many rank-one images off other
sheets: incidence of a zero-dimensional critical set with a
two-dimensional image in a three-dimensional target has
negative expected dimension.

\noindent\textbf{Step 5.} Each transverse rank-one point is a Whitney crosscap, as we
now verify. Choose its first target coordinate as a source
coordinate $u$, and write $\phi=(u,F(u,v),H(u,v))$.
Transversality of the equations $F_v=H_v=0$ gives a linear
combination with nonzero second $v$-derivative. The
holomorphic Morse lemma with parameter $u$, followed by
subtracting the critical value, makes the second coordinate
$v^2$. Subtract the even part of the third coordinate to
write it as $vA(u,v^2)$. Transversality gives $A_u(0,0)\ne0$.
The target change $u\mapsto A(u,Y)$ and the matching source
change give exactly $(u,v)\mapsto(u,v^2,uv)$.
Its image is the crosscap of P4.16, with a compact spherical
class of square $-4$ in its nearby fibre.

\noindent\textbf{Step 6.} The deformed images have holomorphic equations $G_s$ with
$G_0=g$. To see this, a generic projection of the original
finite normalization to two target coordinates is finite;
Weierstrass preparation preserves this after shrinking the
deformation. The total image is an analytic hypersurface
in $\mathbb C^3\times\Delta$. Its central equation has
multiplicity one, as seen at a smooth point of $S$; a unit
adjustment gives $G_0=g$. Outside a smaller ball the original
map is a transverse immersion, and this compact boundary
data persist under small deformation. We therefore have
the boundary property already proved. Transporting the
crosscap class contradicts its torus inclusion. The
normalization is consequently immersive everywhere.

\noindent\textbf{Step 7.} Write its image germ as a unit times $g_1\cdots g_r$, where
each branch is smooth and $v_i=dg_i(0)\ne0$. Absorb the
unit in one factor. Independent gradients give $r\leq3$
and normal-crossing coordinates by the inverse function
theorem. Suppose instead $\sum_i c_iv_i=0$ is a nonzero
relation, and let $I$ be its nonzero support. Set
$a_i=\varepsilon/c_i$ for $i\in I$ and $a_j=\eta$ for
$j\notin I$, where $0<|\varepsilon|\ll|\eta|\ll1$.
Perturb constants so that $h_i(0)=a_i$. Keep all gradients
except one $i_0\in I$, and change that gradient by
\[
 \Delta v_{i_0}=-a_{i_0}\sum_{j\notin I}\frac{v_j}{a_j}.
\]
This is arbitrarily small and makes $\sum_i dh_i(0)/a_i=0$.
Thus $H=\prod_i h_i$ is critical at the origin and has
nonzero critical value $\prod_i a_i$. Its Hessian divided
by that value is
\[
 \sum_i\frac{\operatorname{Hess}(h_i)}{a_i}
 -\sum_i\frac{dh_i\otimes dh_i}{a_i^2}.
\]
An arbitrarily small generic quadratic change to $h_{i_0}$
makes this matrix invertible: its determinant is a nonzero
polynomial in the freely variable symmetric Hessian.
The branch functions remain smooth and retain their
ordinary-double/smooth boundary strata. We have constructed
a holomorphic Morse point off the divisor.

\noindent\textbf{Step 8.} Its nearby fibre contains the real vanishing sphere in
$w_1^2+w_2^2+w_3^2=\epsilon$. Its normal bundle is
$iTS^2$, with the opposite orientation to $TS^2$ under the
complex surface orientation. Its self-intersection is
therefore $-e(S^2)=-2$. Boundary transport again puts this
class in the original torus fibre, a contradiction. Hence
the original branch gradients are independent everywhere,
which proves the normal-crossing assertion.
\end{proof}

The rank-one count is also stated in Gim\'enez Conejero--Pint\'er,
arXiv:2401.14951v2, Lemma 2.3; Section 5.3, p.~33,
discusses the stable-image/nearby-fibre comparison.
The proof above verifies the boundary transport required here.
It does not extend that argument to cuspidal, tacnodal or
triple-curve transverse types.

\paragraph{Exact source locations used here.}
The classification inputs are \cite[Chapter~VI, Theorem~1.1 and
Table~10, pp.~243--244]{BHPV2004E},
\cite[Theorems~4.2 and~5.1, with proofs, pp.~17--21 and~27--29]{OVV2020},
\cite[Theorem~1, p.~1]{BBK2019}, and
\cite[Theorem~2.1]{Catanese2004}.
The no-curve theorem is credited to Inoue by
\cite[Theorem~3.1, p.~303]{KhaledTeleman2026Journal}; the actual
unramified double cover is in its proof of Remark~3.5, p.~305.
The exceptional intersection identity is
\cite[\S I, p.~6]{Mumford1961}; the proof above supplies its
local analytic function and its exact exceptional-divisor map.
The current source records distinguish theorem-statement openings
from complete source-proof readings. The sign and equality assertions
are proved through the displayed normalization, conductor, filling
and intersection maps, with the stated irreducible-or-NC scope for
the global conclusion.

\subsection{The remaining smooth ruled normalizations}

\begin{proposition}[P4.21]\label{night:P4.21}\label{stmt:P4.21}
\StatementRecord{night:P4.21}{proved}{Historical events: GPT-6 Astra (OpenAI), ChatGPT Pro mode, effort Pro, S6\_06\_BLIND, 16--18 November 2026: targeted proof check. Also checked in this writing. Pre-continuation writing events: GPT-6 Astra (OpenAI), ChatGPT Pro mode, effort Pro, S6\_10\_PAPERS, 18 November 2026: targeted proof check. Continuation metadata events: GPT-6 Astra (OpenAI), ChatGPT Pro mode, effort Pro, S6\_10\_PAPERS continuation, 19 November 2026: used and rechecked in this run; status-vocabulary reconciliation; metadata reconciliation only; no new mathematical check. \emph{Pre-continuation recorded proof scope.} The night-run re-derivation was by GPT-6 Astra (OpenAI, ChatGPT Pro mode, effort Pro), 16--18 November 2026. Also checked in this writing by GPT-6 Astra (OpenAI), ChatGPT Pro mode, effort Pro, 18 November 2026. On a smooth rational normalization, torsion adjunction gives an actual anticanonical divisor. Its complete cubic and ruled minimal-model cases have the stated elliptic, cycle or simply connected rational-support topology. Under each allowed blowup the pole coefficient is multiplicity minus one; collapsing the inserted rational tree gives the explicit fundamental-group isomorphism. Every initially nontrivial reduced nodal case remains reduced nodal and gives the NC conclusion by P4.20. In the other cases both normalization and conductor are simply connected, so the conductor pushout identifies pi1(S) with the curve graph free group. Virtual abelianness bounds its rank by one and gives positive curve Euler number; the open-support bound completes the positive sign}{\ref{night:P4.12}; \ref{night:P4.13}; \ref{night:P4.20}; \ref{night:P4.6}; \ref{night:P4.8}; \ref{night:P4R.1}; \ref{night:P4R.3}}

In the standing torus-fibration setting of this section, an irreducible singular support $S$
with smooth rational normalization $T$ has $e(S)>0$.
\end{proposition}
\begin{proof}
\noindent\textbf{Step 1.} Let $D$ be the conductor divisor and
$C=\operatorname{Sing}(S)_{\mathrm{red}}$. By P4.12,
$K_T+D\equiv0$. Contract $T$ to a minimal rational surface
$T_0$ and push $D$ forward to $D_0$. At every reversed
blowup the exceptional coefficient is $\mu-1$, where
$\mu$ is the multiplicity of the current conductor at the
centre. In particular every centre lies on the conductor.

\noindent\textbf{Step 2.} We give the complete minimal-surface case split. On
$\mathbb P^2$, the divisor $D_0$ is a plane cubic. An
irreducible cubic is smooth elliptic, rational nodal, or
rational cuspidal. A reducible cubic is a smooth conic
and a line, three lines, a double line and another line,
or a triple line. The conic and line meet at two
transverse points or at one tangency; three distinct
lines form a triangle or are concurrent. Thus either
$D_0$ is reduced nodal and has fundamental group
$\mathbb Z$ or $\mathbb Z^2$, or its reduced support is
simply connected and all normalized components are
rational.

\noindent\textbf{Step 3.} On a Hirzebruch surface, write $F$ for a ruling fibre.
Since $D_0\cdot F=2$, its horizontal part is a double
section, two distinct reduced sections, or a reduced
irreducible bisection. A double section together with
vertical fibres is a rational tree. For the other cases
write
\[
 D_0=H_1+H_2+\sum_j a_jF_j
 \quad\hbox{or}\quad
 D_0=H+\sum_j a_jF_j.
\]
In the first case adjunction on either section gives
\[
 H_1\cdot H_2+\sum_j a_j=2.
\]
The support is therefore a rational tree or a single
rational cycle. A cycle occurs only with all coefficients
one and ordinary double intersection points. This
includes the cases of two section-intersection points,
one section-intersection point and one vertical fibre,
or two vertical fibres. Coincident intersections or a
vertical fibre through the section-intersection point
instead give trees.

\noindent\textbf{Step 4.} For a bisection, adjunction gives
\[
 p_a(H)=1-\sum_j a_j.
\]
There is either no vertical component and $H$ has
arithmetic genus one, or one coefficient-one vertical
fibre and a smooth rational bisection. In the former
case $H$ is smooth elliptic, rational nodal, or rational
cuspidal. In the latter it meets the fibre at two distinct
points or at one tangency. Again the only nontrivial
fundamental groups are $\mathbb Z$ and $\mathbb Z^2$,
and those cases have reduced nodal conductor.

\noindent\textbf{Step 5.} We follow these alternatives through the blowups. At a
point of conductor multiplicity at least two, the new
conductor contains the exceptional $\mathbb P^1$.
Its contraction is the quotient which collapses that
connected simply connected subcomplex to the old
point. Van Kampen therefore makes the induced map on
fundamental groups an isomorphism. At a point of
multiplicity one the conductor is smooth and reduced,
and its strict transform is isomorphic to it. Thus
$\pi_1(D_{\mathrm{red}})$ is unchanged throughout.
If the initial group was nontrivial, the conductor was
reduced nodal. A smooth-point blowup adds no conductor
component, and a node blowup inserts one component of
coefficient one. It remains reduced nodal.

\noindent\textbf{Step 6.} In this latter case the generic transverse conductor
exponent of every normalization branch is one. The
exponent is its intrinsic conductor exponent plus the
intersection multiplicities with the other branches.
A singular branch has intrinsic exponent at least two.
Exponent one therefore means two smooth branches
meeting transversely. Consequently $S$ has ordinary
double curves outside finitely many points. P4.20 makes
$S$ normal crossing everywhere, and P4.6--P4.8 give
$e(S)=2$ for this rational normalization.

\noindent\textbf{Step 7.} In the remaining case $D_{\mathrm{red}}$ is connected
and simply connected, and all its normalized components
are rational. The conductor pushout verified in P4R.3
gives
\[
 \pi_1(S)=\pi_1(T)*_{\pi_1(D_{\mathrm{red}})}\pi_1(C)
          =\pi_1(C).
\]
Every component of $C$ has rational normalization,
because it is the finite nonconstant image of a
rational component of $D_{\mathrm{red}}$. P4R.1 makes
$\pi_1(C)$ virtually abelian. For a connected curve
with rational normalized components, this is the free
group $F_b$ of its point-identification graph, so
$b\le1$. If $c\ge1$ is its number of components, then
\[
 e(C)=c+1-b\ge c\ge1.
\]
P4.13 gives $e(S_{\mathrm{reg}})\ge0$. Additivity proves
$e(S)=e(S_{\mathrm{reg}})+e(C)>0$.
\end{proof}

\begin{proposition}[P4.22]\label{night:P4.22}\label{stmt:P4.22}
\StatementRecord{night:P4.22}{proved}{Historical events: GPT-6 Astra (OpenAI), ChatGPT Pro mode, effort Pro, S6\_06\_BLIND, 16--18 November 2026: targeted proof check. Also checked in this writing. Pre-continuation writing events: GPT-6 Astra (OpenAI), ChatGPT Pro mode, effort Pro, S6\_10\_PAPERS, 18 November 2026: targeted proof check. Continuation metadata events: GPT-6 Astra (OpenAI), ChatGPT Pro mode, effort Pro, S6\_10\_PAPERS continuation, 19 November 2026: used and rechecked in this run; status-vocabulary reconciliation; metadata reconciliation only; no new mathematical check. \emph{Pre-continuation recorded proof scope.} The night-run re-derivation was by GPT-6 Astra (OpenAI, ChatGPT Pro mode, effort Pro), 16--18 November 2026. Also checked in this writing by GPT-6 Astra (OpenAI), ChatGPT Pro mode, effort Pro, 18 November 2026. Adjunction and Riemann-Hurwitz force coefficient-one horizontal curves on a minimal elliptic ruled surface to be smooth elliptic, unramified and disjoint. Their subsequent blowups have zero exceptional pole coefficient, so P4.15 forces the minimal zero case. In the doubled-section alternative the actual conductor inclusion induces the normalization fundamental-group isomorphism, and the pushout identifies the support group with that of the conductor image. The graph Euler formula gives nonnegativity. Equality leaves a single elliptic image; finiteness and Riemann-Hurwitz exclude rational poles and therefore all blowups. The transverse conductor order two has exactly the cusp, tangent-double, and ordinary-triple possibilities, each giving the claimed zero case through its actual normalization map}{\ref{night:P4.13}; \ref{night:P4.15}; \ref{night:P4R.1}; \ref{night:P4R.3}}

In the standing torus-fibration setting of this section, an irreducible singular support
$S$ whose smooth normalization $T$ is birationally ruled
over an elliptic curve has $e(S)\ge0$. Equality holds
exactly when $T$ is minimal.
\end{proposition}
\begin{proof}
\noindent\textbf{Step 1.} Push the effective conductor to the minimal ruled
surface $T_0\to B$, where $B$ is elliptic. We have
$K_{T_0}+D_0\equiv0$. If a horizontal component $H$
has coefficient one, adjunction gives $p_a(H)\le1$,
whereas its nonconstant map to $B$ gives geometric
genus at least one. Thus $H$ is smooth elliptic and
$(D_0-H)\cdot H=0$. There is no vertical component.
The horizontal part is one elliptic bisection or two
disjoint elliptic sections, all with coefficient one.
All subsequent blowups have smooth coefficient-one
conductor centres; their exceptional curves are absent
from the conductor. The normalization differential
has no divisorial rank drop. P4.15 makes $T$ minimal,
makes the differential immersive everywhere, and
gives $e(S)=0$.

\noindent\textbf{Step 2.} The other possibility is
\[
 D_0=2H+\sum_j a_jF_j,
\]
where $H$ is a smooth elliptic section. As in P4R.3,
$D_{\mathrm{red}}$ is $H$ with rational trees attached,
and its inclusion in $T$ induces an isomorphism on
fundamental groups. The conductor pushout and P4R.1
make $\pi_1(C)$ virtually abelian. The reduced-curve
calculation in P4R.3 gives $e(C)\ge0$, and P4.13 then
gives $e(S)\ge0$.

\noindent\textbf{Step 3.} If $e(S)=0$, the curve calculation forces $C$ to have
one component, genus-one smooth normalization, and
no point-identification graph cycles. Each conductor
component maps finitely and nonconstantly onto $C$.
Riemann--Hurwitz excludes rational conductor
components. Thus $D_{\mathrm{red}}=H$ and no blowup
occurred: the first blowup on $2H$ would create a
vertical conductor component of coefficient one.
Therefore $T$ is minimal.

\noindent\textbf{Step 4.} Conversely suppose $T$ is minimal. In the double-section
case let $v$ be the number of distinct vertical fibres
in the conductor. Since $e(T)=e(H)=0$ and each vertical
$\mathbb P^1$ meets $H$ once,
\[
 e(S_{\mathrm{reg}})=e(T\setminus D_{\mathrm{red}})=-v.
\]
P4.13 forces $v=0$, hence $D=2H$. At a general point
of its image curve, every normalization branch has
conductor exponent two. The branch exponent equals
its intrinsic conductor exponent plus the sum of its
intersections with the other branches. A singular
branch together with another branch would have
exponent greater than two. Thus either there is one
ordinary cusp branch, or all branches are smooth:
two with intersection multiplicity two, or three
with pairwise intersection multiplicity one. Four
smooth branches would have exponent at least three,
so the list is exhaustive.

\noindent\textbf{Step 5.} In the cusp case $H\to C$ has generic degree one.
The normalization of $C$ is therefore $H$, an elliptic
curve. Virtual abelianness of $\pi_1(C)$ excludes any
additional point-identification graph cycle, so
$e(C)=0$. Additivity gives $e(S)=0$. In the other two
cases the normalization differential is immersive at
the generic point of $H$. There is no other conductor
component, so its rank-drop set is finite. P4.15
applies and again gives $e(S)=0$.
\end{proof}


\subsection{Class-VII minimal models with positive second Betti number}

\begin{proposition}[P4.23]\label{night:P4.23}\label{stmt:P4.23}
\StatementRecord{night:P4.23}{proved}{Historical events: GPT-6 Astra (OpenAI), ChatGPT Pro mode, effort Pro, S6\_06\_BLIND, 16--18 November 2026: targeted proof check. Also checked in this writing. Pre-continuation writing events: GPT-6 Astra (OpenAI), ChatGPT Pro mode, effort Pro, S6\_10\_PAPERS, 18 November 2026: targeted proof check. Continuation metadata events: GPT-6 Astra (OpenAI), ChatGPT Pro mode, effort Pro, S6\_10\_PAPERS continuation, 19 November 2026: used and rechecked in this run; status-vocabulary reconciliation; metadata reconciliation only; no new mathematical check. \emph{Pre-continuation recorded proof scope.} The night-run re-derivation was by GPT-6 Astra (OpenAI, ChatGPT Pro mode, effort Pro), 16--18 November 2026. Also checked in this writing by GPT-6 Astra (OpenAI), ChatGPT Pro mode, effort Pro, 18 November 2026. The effective numerically anticanonical divisor supplies the exact hypothesis of the opened Dloussky theorem. On an Inoue-Hirzebruch surface the negative-definite cycle matrix forces the reduced divisor; on a parabolic Inoue surface the numerically trivial cycle coefficient is one because an effective nonzero divisor cannot be torsion. Their reduced nodal transforms give the NC Chern contradiction at positive b2. In the remaining cycle-with-trees case the full conductor pushout is <t,c\_1,...,c\_b | t\textasciicircum{}k=w>. The injective cyclic-amalgam case embeds F\_b; the zero-word cases are checked as the corresponding actual free products. When k=0 and w is nontrivial, any nontrivial quotient factor produces a free subgroup on t and g t g\textasciicircum{}-1, so virtual abelianness forces b<=1. The exact curve Euler count and open-support bound then give strict positivity}{\ref{night:P4.12}; \ref{night:P4.13}; \ref{night:P4.20}; \ref{night:P4.21}; \ref{night:P4.6}; \ref{night:P4R.1}}

In the standing torus-fibration setting of this section, let an irreducible singular support
$S$ have smooth normalization $T$. If its minimal model $T_0$
belongs to class VII and $b_2(T_0)>0$, then $e(S)>0$.
\end{proposition}
\begin{proof}
\noindent\textbf{Step 1.} Let $D$ be the conductor. P4.12 makes $K_T+D$ torsion;
pushing down the point blowups gives $K_{T_0}+D_0$ torsion,
with $D_0$ effective and nonzero. This is a numerically
anticanonical divisor in the sense of Dloussky,
Definition~2.20 of \cite{Dloussky2006}. His Theorem~0.2
therefore gives a global spherical shell. Proposition~2.27
and the square-zero-divisor paragraph after Lemma~2.22
leave the parabolic Inoue or Inoue--Hirzebruch cases,
or a connected rational cycle with attached rational
trees. The latter configuration and $\pi_1(T_0)=\mathbb Z$
are stated explicitly in \cite[\S1, pp.~284--285]{DOT2003}.
We verify the consequences needed
for our conductor.

\noindent\textbf{Step 2.} For an Inoue--Hirzebruch surface, the reduced cycle or
union of two cycles is numerically anticanonical. The
curve intersection matrix is negative definite. Subtracting
that divisor from $D_0$ and intersecting with every curve
therefore makes every coefficient difference zero.
Thus $D_0$ is reduced nodal. In the parabolic Inoue case
there is an elliptic curve $E$ and a disjoint rational
cycle $\Gamma$, with
\[
 K_{T_0}+E+\Gamma\sim0,\qquad \Gamma\equiv0.
\]
These relations are stated after Dloussky's Lemma~2.22.
Intersections with $E$ and the rational curves give
$D_0=E+k\Gamma$: the cycle matrix has kernel generated
by the sum of its components. Torsion of $K_{T_0}+D_0$
then makes $\mathcal O((k-1)\Gamma)$ torsion. A nonzero
effective divisor cannot have torsion associated line
bundle on a compact connected reduced complex space:
a power of its defining section would be a nonzero
constant and could have no zero. Hence $k=1$ and
$D_0$ is again reduced nodal.

\noindent\textbf{Step 3.} The blowup coefficient $\mu-1$ preserves the reduced
nodal property, as in P4.21. Thus $S$ has ordinary
double curves outside finitely many points. P4.20
makes it normal crossing everywhere. Its normalization
is then immersive and P4.6 gives $e(T)=K_T^2$.
For a class-VII surface, however,
\[
 e(T)=b_2(T)>0,\qquad K_T^2=-b_2(T),
\]
a contradiction. These two exceptional minimal-model
cases are excluded.

\noindent\textbf{Step 4.} In the remaining case the support of the conductor is
a connected rational cycle with attached rational trees.
It is the whole connected curve configuration: Dloussky's
Lemma~2.22(2) puts every curve meeting the anticanonical
support in that support, and Proposition~2.27 gives its
connectedness. Its fundamental group is $\mathbb Z$.
The blowups preserve this group by the quotient argument
of P4.21, and preserve the surface fundamental group.
Consequently
\[
 \pi_1(D_{\mathrm{red}})=\mathbb Z,\qquad
 \pi_1(T)=\mathbb Z.
\]
Every normalized component of
$C=\operatorname{Sing}(S)_{\mathrm{red}}$ is rational,
being the finite nonconstant image of a rational conductor
component. The curve is connected. Write $\pi_1(C)=F_b$
for the free group of its point-identification graph.
The conductor pushout gives
\[
 \pi_1(S)=
 \langle t,c_1,\ldots,c_b\mid t^k=w(c_1,\ldots,c_b)\rangle.
\]
Here $k$ records the conductor-cycle map to $\pi_1(T)$,
and $w$ records its map to $\pi_1(C)$. Neither map is
assumed to be an isomorphism.

\noindent\textbf{Step 5.} Virtual abelianness from P4R.1 forces $b\le1$.
If $k\ne0$ and $w\ne1$, both cyclic amalgam maps are
injective; alternating nonidentity coset representatives
give the reduced normal form and show that $F_b$ injects
in the amalgam. If $w=1$, the group is
$(\mathbb Z/k\mathbb Z)*F_b$, with
$\mathbb Z/0\mathbb Z=\mathbb Z$, and $F_b$ again
injects. A nonabelian free subgroup is impossible in a
virtually abelian group. Finally, if $k=0$ and $w\ne1$,
the group is
\[
 \mathbb Z*\bigl(F_b/\langle\!\langle w\rangle\!\rangle\bigr).
\]
A free product of an infinite cyclic group and a nontrivial
group contains a nonabelian free subgroup: for a nontrivial
$g$ in the second factor, $t$ and $gtg^{-1}$ have no
nonempty reduced relation. Thus the second factor must
be trivial. Its abelianization has rank at least $b-1$,
which again gives $b\le1$.

\noindent\textbf{Step 6.} For the number $c\ge1$ of components of $C$, we obtain
\[
 e(C)=c+1-b\ge c\ge1.
\]
By P4.13, $e(S_{\mathrm{reg}})\ge0$, so
$e(S)=e(S_{\mathrm{reg}})+e(C)>0$.
\end{proof}

\subsection{Minimality of a smooth Hopf normalization}

\begin{proposition}[P4.24]\label{night:P4.24}\label{stmt:P4.24}
\StatementRecord{night:P4.24}{proved}{Historical events: GPT-6 Astra (OpenAI), ChatGPT Pro mode, effort Pro, S6\_06\_BLIND, 16--18 November 2026: targeted proof check. Also checked in this writing. Pre-continuation writing events: GPT-6 Astra (OpenAI), ChatGPT Pro mode, effort Pro, S6\_10\_PAPERS, 18 November 2026: targeted proof check; used and rechecked in this run. Continuation metadata events: GPT-6 Astra (OpenAI), ChatGPT Pro mode, effort Pro, S6\_10\_PAPERS continuation, 19 November 2026: used and rechecked in this run; status-vocabulary reconciliation; metadata reconciliation only; no new mathematical check. \emph{Pre-continuation recorded proof scope.} The night-run re-derivation was by GPT-6 Astra (OpenAI, ChatGPT Pro mode, effort Pro), 16--18 November 2026. Also checked in this writing by GPT-6 Astra (OpenAI), ChatGPT Pro mode, effort Pro, 18 November 2026. The primary Hopf normal form, compact-curve monomial calculation, nowhere-zero foliation, blowup saddle charts, saturated derivative Chern identity, integral quadratic bound, and actual smooth elliptic conductor maps were re-derived. Khaled Theorem 2.1(1) and formula (1) were opened in the primary preprint}{the hypotheses and explicit local calculation stated here}

In the standing torus-fibration setting of this section, let $S$ be an irreducible
singular fibre support with smooth normalization $T$.
If the minimal model of $T$ is a Hopf surface, then
$T$ is already minimal. The reduced conductor on $T$
and its reduced image in $S$ are disjoint unions of
smooth elliptic curves. Each map between their
components is a finite unramified covering, and
$e(S)=0$.
\end{proposition}
\begin{proof}
\noindent\textbf{Step 1.} Write $\nu:T\to S\hookrightarrow Y$ for the finite
holomorphic normalization map followed by the closed
embedding. We use the effective integral conductor
divisor $D$ with $K_T+D$ torsion. Let $T\to T_0$
be the sequence of $k$ point blowups giving the
minimal Hopf model. At each step the conductor
coefficient on the exceptional curve is $\mu-1$,
where $\mu$ is the multiplicity of the conductor
at the centre. Effectiveness therefore puts every
centre on the current conductor.

\noindent\textbf{Step 2.} There is a finite unramified primary Hopf cover
$T'_0\to T_0$. Pulling back every blowup gives a
finite unramified cover $T'\to T$ and a blowup
presentation of $T'$ over $T'_0$. A centre has
exactly the degree of the covering many lifts,
so the new number of blowups is that degree times
$k$. The composite $T'\to Y$ is finite and is
immersive away from the pulled-back conductor.
The rational Chern classes of its pulled-back
ambient tangent bundle vanish. Indeed $b_2(Y)=0$
gives $c_1(T_Y)=c_2(T_Y)=0$ in rational cohomology,
using Poincar\'e duality in degree four. The
argument below applies to this finite map, without
requiring it to be a normalization. Proving
minimality of $T'$ proves $k=0$. We therefore use
$T,T_0,\nu,k$ for the covering objects during the
Chern calculation.

\noindent\textbf{Step 3.} We first construct on $T_0$ a nonsingular
holomorphic foliation tangent to every compact
curve. Khaled, \emph{Real structures on primary
Hopf surfaces}, Theorem~2.1(1) and formula~(1),
gives
\[
 T_0=(\mathbb C^2\setminus\{0\})/\langle F\rangle,
 \qquad F(z,w)=(\alpha z+\lambda w^n,\beta w),
\]
where $0<|\alpha|\le|\beta|<1$ and
$\lambda(\alpha-\beta^n)=0$. We verify the needed
curve statement from these coordinates. Every
compact irreducible curve is smooth elliptic:
adjunction and $b_2(T_0)=0$ give arithmetic genus
one, and a rational normalization would lift to
$\mathbb C^2\setminus\{0\}$ and have constant
coordinate functions.

\noindent\textbf{Step 4.} The map from the fundamental group of such an
elliptic curve to $\pi_1(T_0)=\mathbb Z$ has
nonzero image, since zero image would give the
same impossible compact lift. Its full inverse
image has finitely many connected components.
The Remmert--Stein extension across the origin
is an analytic curve germ; a power $F^h$ fixes
one of its branches. On the normalization of
that branch the induced contraction can be
written $t\mapsto\eta t$, $0<|\eta|<1$.
For completeness, if the initial coordinate
has return map $g(t)=\eta t+O(t^2)$, choose
$|\eta|<\rho<\sqrt{|\eta|}$. On a sufficiently
small disc, $|g(t)|\le\rho|t|$, and the
successive differences in $g^j(t)/\eta^j$
are bounded by a constant times
$(\rho^2/|\eta|)^j|t|^2$. Their convergent sum
gives the asserted linearizing coordinate.

\noindent\textbf{Step 5.} When $\lambda=0$, comparison of coefficients in
$z(\eta t)=\alpha^h z(t)$ and
$w(\eta t)=\beta^h w(t)$ shows that each
nonzero coordinate is a single monomial.
Thus every branch away from the coordinate axes
has parametrization $(a t^p,b t^q)$. Its ratio
$p/q=\log|\alpha|/\log|\beta|$ is independent
of the branch. If such a branch exists, take
coprime positive integers $p,q$ in this ratio
and the vector field
$p z\partial_z+q w\partial_w$.
If no such branch exists, use
$z\partial_z+w\partial_w$.
Both are invariant under $F$, have no zero on
$\mathbb C^2\setminus\{0\}$, and are tangent
to every compact curve.

\noindent\textbf{Step 6.} When $\lambda\ne0$, we have $\alpha=\beta^n$ and
\[
 F^h(z,w)=
 (\alpha^h z+h\lambda\alpha^{h-1}w^n,\beta^h w).
\]
If a normalized invariant branch had
$w(t)=b t^q\ne0$, the coefficient of $t^{nq}$
in the first coordinate identity would give
$0=h\lambda\alpha^{h-1}b^n$, a contradiction.
Thus the only curve is the image of $w=0$.
The vector field $n z\partial_z+w\partial_w$
is invariant, nowhere zero, and tangent to it.
This proves the foliation assertion.

\noindent\textbf{Step 7.} Pull this foliation through the point blowups
and saturate its tangent sheaf. It is a line
bundle $L\subset T_T$. Every strict transform
of a conductor component and every exceptional
curve is invariant. The only singularities of
the foliation are simple saddles. To verify
this assertion and its multiplicities, at a
regular point use coordinates $(x,y)$ with
leaves $y=\text{constant}$. In the chart
$x=u,y=uv$ the saturated generator is
$u\partial_u-v\partial_v$; there is one
saddle, and the other blowup chart is regular
away from it. At an existing saddle use
$b x\partial_x-a y\partial_y$, with positive
integers $a,b$. The two charts have generators
\[
 b u\partial_u-(a+b)v\partial_v,\qquad
 (a+b)u\partial_u-a v\partial_v.
\]
One saddle is replaced by two. Each has
zero ideal $(u,v)$ and length one.
Consequently there are exactly $k$ saddles.

\noindent\textbf{Step 8.} Let $V$ be the divisor measuring the
divisorial zero of the pulled-back transverse
one-form before saturation. Numerically
\[
 c_1(L)=-K_T+V.
\]
The omitted pullback line bundle is topologically
trivial because it comes from $T_0$.
The preceding charts give, at every blowup $\mu$,
\[
 K'=\mu^*K+E,\qquad
 V'=\mu^*V+\varepsilon E,\qquad
 \varepsilon=
 \begin{cases}0&\text{at a regular point},\\
 1&\text{at a saddle}.
 \end{cases}
\]
Starting with $V=0$ and $K_{T_0}^2=0$, we obtain
\[
 K_T^2=-k,\qquad K_TV-V^2=0.
\]
The second identity is preserved because its
change is $\varepsilon^2-\varepsilon=0$.

\noindent\textbf{Step 9.} Put $B=\nu^*T_Y$ and compose the holomorphic
maps $L\to T_T\to B$. This map has no
divisorial zero. Away from the conductor this
follows from immersivity. At the generic point
of a conductor component the foliation is
tangent to that curve, and finiteness of
$\nu$ makes its differential along the curve
nonzero. A nonconstant holomorphic map of
curves has nonzero derivative at the generic
point. Thus the cokernel $Q=B/L$ is a
torsionfree sheaf of rank two. Its double
dual $Q^{**}$ is locally free, and
\[
 z_0=\operatorname{length}(Q^{**}/Q)
\]
is finite. At a point, write the map $L\to B$
as a column $(f_1,f_2,f_3)$ in trivializations.
If this column vanishes, the ideal
$(f_1,f_2,f_3)$ has finite colength. Dualizing
$0\to\mathcal O\to\mathcal O^3\to Q\to0$
gives $\mathcal E xt^1(Q,\mathcal O)
=\mathcal O/(f_1,f_2,f_3)$.
Dualizing $0\to Q\to Q^{**}\to Q^{**}/Q\to0$
and using the two-dimensional regular local
ring identifies this with
$\mathcal E xt^2(Q^{**}/Q,\mathcal O)$,
which has the same length. At each foliation
saddle all three coefficients vanish.
It follows that
\[
 z_0\ge k.
\]

\noindent\textbf{Step 10.} Let $E_B\subset B$ be the saturation of the
generic rank-two image of $d\nu:T_T\to B$.
Its determinant differs from that of $T_T$
by an effective integral divisor $R$, the
divisorial rank-drop divisor:
\[
 c_1(E_B)=-K_T+R.
\]
The support of $R$ is contained in the
conductor. Saturate $E_B/L$ in $Q^{**}$.
The resulting rank-one reflexive subsheaf
is a line bundle $J$, with
$c_1(J)=R-V$. Its torsionfree rank-one
quotient has the form
$\mathcal I_{Z_1}\otimes J_1$ with
$c_1(J_1)=K_T-R$ and
$z_1=\operatorname{length}(Z_1)\ge0$.
All these assertions concern sheaves on the
smooth surface $T$; determinants are computed
off finite sets and therefore determine their
line-bundle classes.

\noindent\textbf{Step 11.} The identities $c_1(B)=c_2(B)=0$ give
\[
 c_2(Q)=(-K_T+V)^2,\qquad
 c_2(Q^{**})=(-K_T+V)^2-z_0.
\]
The exact sequence with $J$ gives another
expression for the same number:
\[
 (-K_T+V)^2-z_0=(R-V)(K_T-R)+z_1.
\]
Using $K_T^2=-k$ and $K_TV-V^2=0$, this becomes
\[
 G(R):=R^2-K_TR-VR=k+z_0+z_1\ge2k.
\]

\noindent\textbf{Step 12.} We compare this expression with the actual
blowup sequence. Push $R$ and $V$ down at
each step; $R$ has integral coefficients.
Writing $R'=\mu^*R+cE$ gives $c\in\mathbb Z$,
without any positivity requirement on $c$.
The intersection calculation is
\[
 G(R')-G(R)=-c^2+(1+\varepsilon)c.
\]
For $\varepsilon=0$ this is at most zero
for every integer $c$; for $\varepsilon=1$
it is $1-(c-1)^2\le1$. On $T_0$ every
divisor is numerically zero, so $G(R_0)=0$.
Consequently
\[
 G(R)\le\#\{\text{blowups at saddles}\}\le k.
\]
The lower bound $G(R)\ge2k$ forces
$k=z_0=z_1=0$. Returning through the finite
unramified cover proves that the original
normalization $T$ is minimal.

\noindent\textbf{Step 13.} It remains to identify the conductor maps,
rather than only their Euler number.
The reduced conductor on a minimal Hopf
surface consists of disjoint smooth elliptic
curves. Since $z_0=0$ on the primary cover,
the differential of $\nu$ along each such
curve is nowhere zero, also before that
cover. Thus every branch of its image is
smooth. At an image point, two different
curve germs would have to come from two
different normalization points. Their
surface branches in the smooth threefold
intersect in a curve germ. This germ is
contained in the conductor image of both
branches; the unique conductor germ at
each normalization point forces the two
image germs to be equal. Therefore the
image conductor is locally a single smooth
curve. Its components are disjoint, and
each conductor map is finite unramified.
Riemann--Hurwitz makes every image component
elliptic. Finally, normalization is an
isomorphism off these curves, so
\[
 e(S)=e(T)-e(D_{\mathrm{red}})
                 +e(\operatorname{Sing}S)=0.
\]
\end{proof}


\subsection{The reduced resolved-conductor case}

\begin{proposition}[P4.25]\label{night:P4.25}\label{stmt:P4.25}
\StatementRecord{night:P4.25}{proved}{Historical events: GPT-6 Astra (OpenAI), ChatGPT Pro mode, effort Pro, S6\_06\_BLIND, 16--18 November 2026: targeted proof check. Also checked in this writing. Pre-continuation writing events: GPT-6 Astra (OpenAI), ChatGPT Pro mode, effort Pro, S6\_10\_PAPERS, 18 November 2026: targeted proof check; used and rechecked in this run. Continuation metadata events: GPT-6 Astra (OpenAI), ChatGPT Pro mode, effort Pro, S6\_10\_PAPERS continuation, 19 November 2026: used and rechecked in this run; status-vocabulary reconciliation; metadata reconciliation only; no new mathematical check. \emph{Pre-continuation recorded proof scope.} The night-run re-derivation was by GPT-6 Astra (OpenAI, ChatGPT Pro mode, effort Pro), 16--18 November 2026. Also checked in this writing by GPT-6 Astra (OpenAI), ChatGPT Pro mode, effort Pro, 18 November 2026. The actual normalized boundary gives the log-canonical resolution inequalities. Liu-Rollenske2012's main theorem, Table1, Definition2.3, and Propositions3.1/4.1 were opened at their exact author-preprint loci. Every listed non-NC germ was deformed with its full multiplicity and boundary Hessian retained, reducing to the proved crosscap or tangent-sheet obstruction. The final global applications retain P4.6/P4.8 hypotheses}{\ref{night:P4.16}; \ref{night:P4.20}; \ref{night:P4.6}; \ref{night:P4.8}}

Every semi-log-canonical reduced fibre support in
Part~4.2 has normal crossings. In particular, let
$\rho:T\to N\to S$ be a resolution of the normalization
of an irreducible support, and let $D$ be the integral
divisor defined by
$\omega_T(D)\simeq\rho^*\omega_S$. When this actual
divisor is reduced and has normal crossings, $S$
has normal crossings, its normalization is smooth,
and P4.6--P4.8 apply.
\end{proposition}
\begin{proof}
\noindent\textbf{Step 1.} We first explain why the stated resolved-conductor
case is semi-log-canonical. At a generic conductor
point, an exponent one on every normalization
branch forces the transverse plane curve to be an
ordinary node. Indeed the conductor exponent on
branch $i$ is
$2\delta_i+\sum_{j\ne i}I_{ij}$; positivity of
intersection multiplicities leaves exactly two
smooth transverse branches. The support $S$ is a
Cartier hypersurface, hence is Cohen--Macaulay and
has invertible dualizing sheaf. On its normal
normalization $N$, the pair consisting of $N$
and its reduced conductor is log canonical.
To check this last assertion, pull its canonical
divisor to $T$. The resulting boundary is exactly
$D$, by the displayed isomorphism and the
divisorial orders of a local generator. A smooth
surface with reduced normal-crossing boundary is
log canonical: blowing up a point gives new
boundary coefficient $\mu-1\le1$, and this
bound persists under successive point blowups.
Discrepancy comparison therefore gives the
assertion on $N$. This is precisely the
semi-log-canonical definition of
Liu--Rollenske, Definition~2.3 and Lemma~2.2,
which explicitly applies to complex analytic
spaces.

\noindent\textbf{Step 2.} We now treat all local semi-log-canonical
hypersurface germs. We use the complete analytic
classification in Liu--Rollenske,
\emph{Two-dimensional semi-log-canonical
hypersurfaces}, main Theorem, Table~1,
Propositions~3.1 and~4.1. The normal singular
germs are isolated. Any such germ is excluded
in our family: a small generic linear
perturbation of its reduced equation has
Morse critical points, the number counted
with multiplicity being the positive
colength of its gradient ideal. The equality
of this number before and after perturbation
follows by conservation of the local
intersection multiplicity of the three
gradient equations. A constant perturbation
puts their critical values away from zero.
The original zero divisor is smooth near
a small transverse boundary sphere, so it
retains that boundary property under these
perturbations. The proper boundary
transport proved in P4.20 carries a
vanishing sphere of square $-2$ into a
nearby torus fibre, which is impossible.

\noindent\textbf{Step 3.} For the nonnormal germs we give the exact
deformations and check their boundary strata.
The ordinary double point and $xyz=0$
already have normal crossings. A pinch point
$x^2+y^2z=0$ is excluded by P4.16.
The remaining double-point equations are
\[
 x^2+y^2z^2,\qquad
 x^2+y^2(z^2+y^{q-2}),\quad q\ge3.
\]
For the first equation the two normalization
branches are the smooth graphs $x=\pm i yz$.
Their tangent planes coincide at the origin,
whereas they meet transversely off that point
along the two axes. The local proof of P4.20
therefore supplies a vanishing sphere after
a deformation retaining its boundary
ordinary-double strata.

\noindent\textbf{Step 4.} For finite $q$ put
\[
 G_s=x^2+y^2(z^2+y^{q-2}+s).
\]
For $s\ne0$ small, at each point
$(0,0,z_0)$ with $z_0^2=-s$ the function
$Z=z^2+y^{q-2}+s$ is a local coordinate:
its $z$ derivative is $2z_0\ne0$.
The actual equation there is exactly
$x^2+y^2Z$, a pinch point. At the original
boundary double axis $x=y=0$, $z\ne0$,
the transverse Hessian determinant is
$4(z^2+s)$, which stays nonzero on the
compact boundary when $s$ is small.
All boundary smooth points remain smooth.
Thus the zero divisor retains the same
boundary smooth and double strata.
The spherical class of square $-4$
from the pinch-point nearby fibre transports,
by the detailed P4.20 construction, to
the original nearby torus fibre.
This excludes every finite $q$.

\noindent\textbf{Step 5.} The remaining triple-point equations are
\[
 xyz+x^p,\quad p\ge3,\qquad
 xyz+x^p+y^q,\quad q\ge p\ge3.
\]
For the first put
$G_s=xyz+x^p+s x^2
=x(yz+x^{p-1}+s x)$.
When $s\ne0$, its two local normalization
branches are smooth: the derivative in $x$
of the second factor at the origin is $s$.
Their tangent planes coincide, so the
origin is not a normal crossing. Along
the boundary $y$ axis the transverse
Hessian determinant is $-y^2$, and along
the boundary $z$ axis it is $-z^2$.
Both axes persist as ordinary-double
strata. The local argument of P4.20 on
the deformed germ supplies a compact
spherical class of nonzero square, and
the same boundary transport gives the
contradiction.

\noindent\textbf{Step 6.} For the second equation put
\[
 G_s=xyz+x^p+y^q+s x^2+s y^3,
 \qquad 0<|s|\ll1.
\]
Its transverse Hessian along the boundary
$z$ axis has determinant $-z^2$,
so its sole boundary double stratum
persists. We check the smooth normalization
of its germ at the origin explicitly.
Since $\partial^2G_s/\partial x^2(0)=2s$,
the critical equation in $x$ has a unique
analytic solution $x=\xi(y,z)$ with
$\xi$ divisible by $yz$. The parametric
holomorphic Morse change in $x$ writes
\[
 G_s=X^2+y^2H(y,z).
\]
Divisibility follows directly on substituting
$\xi$: each of $\xi yz,\xi^p,s\xi^2$
is divisible by $y^2$, as are the two
pure $y$ terms. Moreover
\[
 H_y(0,0)=
 \begin{cases}1+s&q=3,\\s&q>3,\end{cases}
 \qquad H(0,z)=-z^2/(4s).
\]
The finite
birational map
\[
 \{v^2+H(y,z)=0\}\longrightarrow\{G_s=0\},
 \qquad (v,y,z)\longmapsto(X=yv,y,z),
\]
followed by the inverse Morse coordinate
change, is the normalization. Its source
is smooth because $H_y(0,0)\ne0$,
and the map is finite because $v$
satisfies the displayed monic equation.
It is birational on $y\ne0$ and is
therefore the normalization by normality
of its source. At the origin its
differential has rank one, so this germ
is not a normal crossing. P4.20 again
produces a spherical class, with the
already verified boundary transport.

\noindent\textbf{Step 7.} These calculations exhaust the analytic
classification. Thus only the smooth,
ordinary-double, and triple-normal-crossing
germs remain. Their local equations give
the claimed normal-crossing support and
smooth normalization.
\end{proof}


\subsection{Numerical sign for arbitrary irreducible supports}

\begin{proposition}[P4.26]\label{night:P4.26}\label{stmt:P4.26}
\StatementRecord{night:P4.26}{proved}{Historical events: GPT-6 Astra (OpenAI), ChatGPT Pro mode, effort Pro, S6\_06\_BLIND, 16--18 November 2026: targeted proof check. Also checked in this writing. Pre-continuation writing events: GPT-6 Astra (OpenAI), ChatGPT Pro mode, effort Pro, S6\_10\_PAPERS, 18 November 2026: targeted proof check. Continuation metadata events: GPT-6 Astra (OpenAI), ChatGPT Pro mode, effort Pro, S6\_10\_PAPERS continuation, 19 November 2026: used and rechecked in this run; status-vocabulary reconciliation; metadata reconciliation only; no new mathematical check. \emph{Pre-continuation recorded proof scope.} The night-run re-derivation was by GPT-6 Astra (OpenAI, ChatGPT Pro mode, effort Pro), 16--18 November 2026. Also checked in this writing by GPT-6 Astra (OpenAI), ChatGPT Pro mode, effort Pro, 18 November 2026. The rational pole-support transfers, conductor pushout and group presentations, class-VII alternatives, local ruled and Hopf comparisons, normal cover case, Euler sum and parity were re-derived with the source's exact separation between the single-support conclusion and the global all-support hypothesis}{\ref{night:P4.12}; \ref{night:P4.13}; \ref{night:P4.21}; \ref{night:P4.23}; \ref{night:P4.25}; \ref{night:P4.6}; \ref{night:P4.7}; \ref{night:P4.8}; \ref{night:P4R.1}; \ref{night:P4R.6}; \ref{night:P4R.7}; \ref{night:P4R.8}; \ref{night:P4R.9}}

In Part~4.2 every irreducible reduced fibre support
has nonnegative topological Euler number. For a
singular irreducible support, rational resolution
of the normalization gives strictly positive
Euler number; so does a resolution whose minimal
model is of class VII with positive second
Betti number. These assertions permit singular
normalization.

Every normal-crossing support has the
nonnegative even Euler number of P4.6.
Consequently, for a fibration all of whose
reduced supports are irreducible or have
normal crossings,
\[
 B=\mathbb P^1,\qquad c_3(Y)\in\{0,2\},
 \qquad b_3(Y)=2-c_3(Y).
\]
This last conclusion does not assume a
classification of the equality cases for
irreducible nonnormal supports.
\end{proposition}
\begin{proof}
\noindent\textbf{Step 1.} For an irreducible singular support use the
minimal resolution of its normalization
$\rho:T\to N\to S$, and put
$A=\rho^{-1}(\operatorname{Sing}S)_{\rm red}$,
$C=\operatorname{Sing}(S)_{\rm red}$.
P4R.6 gives an effective integral divisor
$D$ with support $A$ and $K_T+D$ torsion.
The resolution $T$ has Kodaira dimension
$-\infty$, by P4.12. We explain the two
remaining transfers from smooth
normalization to this resolution.

\noindent\textbf{Step 2.} Suppose first that $T$ is rational.
Push $D$ to a minimal rational surface.
The complete anticanonical-support calculation
in P4.21 has the following exact alternatives.
The reduced support is connected; either
its fundamental group is trivial and all
its normalized components are rational,
or the divisor itself is reduced with
normal crossings. On $\mathbb P^2$ these
are the plane cubic configurations
enumerated there. On a ruled minimal
surface they are the degree-two horizontal
configurations there. Those arguments use
only the effective integral divisor and
$K+D\equiv0$, so apply to the present
resolved divisor.

\noindent\textbf{Step 3.} We check their preservation through the
actual blowup morphisms. A blowup at
conductor multiplicity $\mu$ gives
exceptional coefficient $\mu-1$.
For $\mu=1$, the centre is a smooth
point of the reduced support and its
strict transform is isomorphic to that
support. For $\mu\ge2$, the whole
exceptional $\mathbb P^1$ belongs to
the new support. Its contraction is
the blowup map on supports. Collapsing
this connected simply connected
subcomplex does not change the
fundamental group, by van Kampen.
The reduced normal-crossing case stays
reduced normal crossing: a smooth
point gives coefficient zero and a
double point gives coefficient one.
Thus the same alternatives hold for
the actual divisor $D$ on $T$.

\noindent\textbf{Step 4.} In the reduced normal-crossing alternative,
P4.25 makes $S$ normal crossing. Its
normalization is smooth, so this minimal
normalization resolution is the
normalization itself. The rational case
of P4.6--P4.8 gives $e(S)=2$.
In the other alternative $\pi_1(A)=1$
and $\pi_1(T)=1$. The proper-map pushout
$T\cup_A C\to S$ is a homeomorphism:
it is a continuous compact-to-Hausdorff
bijection, since $\rho$ is an isomorphism
off $A$ and $A$ is the full inverse
image of $C$. Compatible triangulation
makes $A\to T$ a cofibration.
Van Kampen therefore gives
$\pi_1(S)\simeq\pi_1(C)$.
Every normalized component of $C$ is
rational: a nonconstant map from a
normalized component of $A$ onto it
is a finite map from $\mathbb P^1$.
Contracted components of $A$ contribute
no exception to this statement.
The curve $C$ is connected, and its
fundamental group is a free group
$F_b$ on the cycles of its
point-identification graph. P4R.1
implies $b\le1$. If $c\ge1$ is the
number of its components, then
\[
 e(C)=c+1-b\ge1.
\]
P4.13 gives $e(S_{\rm reg})\ge0$,
so $e(S)>0$.

\noindent\textbf{Step 5.} Now suppose that the minimal model
$T_0$ is of class VII with
$b_2(T_0)>0$. The divisor pushed to
$T_0$ satisfies the hypotheses of
Dloussky's numerically anticanonical
theorem exactly as in P4.23.
The verified sources and their
normalization-independent hypotheses
are specified there. The
Inoue--Hirzebruch and parabolic Inoue
alternatives give a reduced
normal-crossing divisor, preserved
under the blowups just checked.
P4.25 then makes $S$ normal crossing
and $T=N$ smooth. The immersive
normalization Chern formula would
give $e(T)=K_T^2$, whereas
$e(T)=b_2(T)>0$ and
$K_T^2=-b_2(T)$. These alternatives
cannot occur.

\noindent\textbf{Step 6.} In the intermediate alternative
$A$ is a connected rational cycle
with attached rational trees, and
\[
 \pi_1(A)=\mathbb Z,\qquad
 \pi_1(T)=\mathbb Z.
\]
Both statements survive the actual
blowups; they concern the resolution
and do not require $N$ to be smooth.
Every normalized component of $C$
is again rational and $C$ is connected.
With $\pi_1(C)=F_b$, the pushout
presentation is exactly
\[
 \pi_1(S)=
 \langle t,c_1,\ldots,c_b
 \mid t^k=w(c_1,\ldots,c_b)\rangle
\]
for the actual two conductor-group
homomorphisms. A component contracted
to a point merely induces the zero
homomorphism on its group. The
amalgam and free-product normal-form
argument in P4.23, which allowed
arbitrary $k$ and $w$, gives $b\le1$
from virtual abelianness.
Thus $e(C)=c+1-b\ge1$, and
$e(S)>0$ follows again from P4.13.

\noindent\textbf{Step 7.} This exhausts these two resolution
classes. Ruled resolutions of positive
base genus are covered by P4R.7
and P4R.9, and Hopf minimal models
by P4R.8. The class-VII minimal
model with $b_2=0$ is Hopf or Inoue,
by the verified classification in
P4.8. An Inoue model has no curves,
so it cannot carry the nonempty
effective divisor pushed down from
$D$. Therefore it does not occur
here. All irreducible singular
supports have now been covered.

\noindent\textbf{Step 8.} A normal irreducible support is
smooth by P4.12. The cyclic
base-change comparison and torus
deformation theorem in P4.8 show
that it has a finite unramified
torus cover, hence Euler number
zero. This proves the full
irreducible assertion.

\noindent\textbf{Step 9.} For the final global statement,
$B=\mathbb P^1$ follows from
$b_1(Y)=0$ and connected fibres,
as proved in P4.7. Euler
additivity over the finitely many
exceptional base values gives
\[
 e(Y)=\sum_{b\in B}
 e\bigl((f^{-1}(b))_{\rm red}\bigr),
\]
because a smooth torus has Euler
number zero and the complement
is a proper smooth fibre bundle.
Every summand is nonnegative
under the stated support
hypothesis. Poincar\'e duality
and $b_1=b_2=0$ give
$e(Y)=2-b_3(Y)$; the nondegenerate
skew-symmetric cup-product
pairing on $H^3(Y;\mathbb Q)$
makes $b_3(Y)$ even. Thus
$0\le e(Y)\le2$ and $e(Y)$ is
even, which proves the asserted
values and their Chern-number
interpretation.
\end{proof}


\subsection{Equality in the Hopf-resolution case}

\begin{proposition}[P4.27]\label{night:P4.27}\label{stmt:P4.27}
\StatementRecord{night:P4.27}{proved}{Historical events: GPT-6 Astra (OpenAI), ChatGPT Pro mode, effort Pro, S6\_06\_BLIND, 16--18 November 2026: targeted proof check. Also checked in this writing. Pre-continuation writing events: GPT-6 Astra (OpenAI), ChatGPT Pro mode, effort Pro, S6\_10\_PAPERS, 18 November 2026: targeted proof check. Continuation metadata events: GPT-6 Astra (OpenAI), ChatGPT Pro mode, effort Pro, S6\_10\_PAPERS continuation, 19 November 2026: used and rechecked in this run; status-vocabulary reconciliation; metadata reconciliation only; no new mathematical check. \emph{Pre-continuation recorded proof scope.} The night-run re-derivation was by GPT-6 Astra (OpenAI, ChatGPT Pro mode, effort Pro), 16--18 November 2026. Also checked in this writing by GPT-6 Astra (OpenAI), ChatGPT Pro mode, effort Pro, 18 November 2026. The exact Hopf pole graph has r elliptic roots and v rational vertices after k blowups, so equality forces k=v and zero conductor-image Euler number. In the all-elliptic image case every rational pole contracts and no elliptic root does. The resolution-to-minimal-model and resolution-to-normalization maps then have exactly the same connected fibres; normalizing their graph image gives proper finite birational projections to normal surfaces, hence a holomorphic isomorphism of the normalization with the minimal Hopf surface. In the sole remaining graph count, the homomorphism from the Hopf group to Z has finite kernel and each elliptic-root map has infinite image by the actual universal-cover lift argument. Quotienting each elliptic vertex by the saturated image of its rank-one kernel gives a surjection to a graph of cyclic groups. Removing its unique trivial-edge cycle leaves a tree amalgam containing Z; restoring it yields a free product with Z containing F2, contradicting virtual abelianness. Thus the equality normalization is smooth minimal Hopf}{\ref{night:P4.24}; \ref{night:P4R.1}; \ref{night:P4R.6}; \ref{night:P4R.8}}

Let \(S\) be an irreducible singular reduced fibre support in Part 4.2.
Suppose that the minimal resolution \(T\to N\) of its normalization has a
Hopf minimal model \(T_0\). Then
\[
 e(S)=0\quad\Longleftrightarrow\quad
 N\text{ is a smooth minimal Hopf surface}.
\]
In this case the reduced conductor and its image are disjoint unions of
smooth elliptic curves, and all component maps are finite unramified covers.
\end{proposition}
\begin{proof}
\noindent\textbf{Step 1.} Write \(\rho:T\to N\to S\), \(A=\rho^{-1}(\operatorname{Sing}S)_{\rm red}\),
and \(C=\operatorname{Sing}(S)_{\rm red}\). Let \(k\) be the number of
blowups in \(\pi:T\to T_0\). P4R.6 and P4R.8 give an effective integral pole
divisor whose reduced support \(A\) has \(r\) connected components. Each is
a smooth elliptic root with a rational tree attached. Let \(v\) be the
number of its rational components. Every such component is one of the
\(k\) blowup exceptional curves. We have
\[
 e(S)=k-v+e(C),\qquad k-v\ge0,\qquad e(C)\ge0.
\]
Thus equality forces \(k=v\) and \(e(C)=0\).

\noindent\textbf{Step 2.} Let \(c,s,g,b\) denote the numbers of irreducible components of \(C\),
connected components of \(C\), elliptic normalized components of \(C\), and
the sum of its point-identification graph cycle ranks. P4R.8 gives
\[
 e(C)=c+s-2g-b,\qquad
 b+r-s\le1,\qquad c\ge g,\quad r\ge g.
\]
Consequently \(e(C)=0\) implies \(c+r-2g\le1\).

\noindent\textbf{Step 3.} First suppose \(c=g\). Every normalized component of \(C\) is elliptic.
Every rational component of \(A\) is contracted by \(T\to N\), since a
finite nonconstant map from \(\mathbb P^1\) to an elliptic normalization
does not exist. A root cannot be contracted. Indeed all rational curves
attached to it are already contracted, so its whole rooted connected
component would be a connected exceptional block disjoint from the strict
conductor. P4R.6 excludes precisely such a block. Hence all roots are
noncontracted.

\noindent\textbf{Step 4.} The equality \(k=v\) now says that every blowup exceptional curve of
\(\pi:T\to T_0\) belongs to \(A\) and is contracted by \(T\to N\).
Conversely every exceptional curve of \(T\to N\) belongs to \(A\), by
P4R.6, and cannot be a root. The two proper birational maps therefore have
exactly the same exceptional curves. Their fibres are identical: each
nontrivial fibre is a connected component of the exceptional locus, and
their fibres are connected by Stein factorization and normality of the
targets.

\noindent\textbf{Step 5.} Here is the holomorphic comparison, rather than only a topological
identification. Take the image of
\[
 (\pi,\rho_N):T\longrightarrow T_0\times N
\]
and then its normalization \(G\). The two projections \(G\to T_0,N\)
are proper birational maps with zero-dimensional fibres, because the
original two maps have identical fibres. They are finite. A finite
birational morphism to a normal complex space is an isomorphism, as follows
on local rings from integral closedness. Thus \(N\simeq T_0\)
biholomorphically. It is a smooth minimal Hopf surface.

\noindent\textbf{Step 6.} The only other possibility is \(c=g+1,r=g\). There is one rational
normalized component of \(C\), and all \(g=r\) roots are noncontracted and
map to distinct elliptic normalized components of \(C\). Each such map is
an isogeny. Also
\[
 b+r-s=1.
\]
We exclude this possibility by a quotient of the exact fundamental-group
pushout; the quotient is constructed explicitly below.

\noindent\textbf{Step 7.} Represent \(C\) as its disjoint normalized smooth curves with a finite
graph of point identifications. This is a graph of spaces with elliptic
vertex groups \(\mathbb Z^2\), trivial rational vertex groups, and trivial
edge groups. Its graph has \(s\) connected components and total cycle rank
\(b\). Add a vertex for \(T\), and add one root edge for each connected
component of \(A\). The root edge group is \(\mathbb Z^2\), since its
rational tree is simply connected. Its map to the \(C\) side factors
through the elliptic normalization receiving that root; its map to the
\(T\) side is the actual inclusion. This graph-of-spaces presentation is
the ordinary van Kampen presentation of the homotopy pushout
\(T\cup_A C=S\). It has total graph cycle rank \(b+r-s=1\).

\noindent\textbf{Step 8.} Let \(\Gamma=\pi_1(T)=\pi_1(T_0)\). It is virtually infinite cyclic, and
\(b_1(T)=1\). The primitive map from its abelianization to its free part
therefore gives a surjection
\[
 h:\Gamma\longrightarrow\mathbb Z
\]
with finite kernel. Indeed a cyclic subgroup of finite index maps
nontrivially to \(\mathbb Z\), so the kernel meets that subgroup trivially
and is finite.

\noindent\textbf{Step 9.} For root \(H_i\), write
\[
 a_i:\pi_1(H_i)=\mathbb Z^2\longrightarrow\Gamma,\qquad
 q_i:\pi_1(H_i)=\mathbb Z^2\longrightarrow
       \pi_1(\widetilde C_i)=\mathbb Z^2.
\]
The second map is injective of finite index because it is induced by an
isogeny. The composite \(h a_i\) is nonzero. To see this, a finite image
of \(a_i\) would become trivial after a finite cover of \(H_i\); that
compact elliptic cover would lift holomorphically to the universal Hopf
cover \(\mathbb C^2\setminus\{0\}\), where its coordinate functions are
constant. Thus \(a_i\) has infinite image; since \(\ker h\) is finite,
\(h a_i\ne0\).

\noindent\textbf{Step 10.} Put \(K_i=\ker(h a_i)\), a primitive rank-one subgroup of
\(\mathbb Z^2\). Quotient the \(i\)-th elliptic vertex group by the
saturation of \(q_i(K_i)\). The resulting vertex group is \(\mathbb Z\).
Quotient the root edge group by \(K_i\); it too becomes \(\mathbb Z\).
Its two induced maps into the \(T\) vertex \(\mathbb Z\) and the elliptic
vertex \(\mathbb Z\) are nonzero multiplications and hence injective.
The latter assertion follows because an element outside \(K_i\) cannot
map under the injective rational map \(q_i\otimes\mathbb Q\) into the
rational span of \(q_i(K_i)\). This verifies compatibility of the chosen
quotients with both edge homomorphisms.

\noindent\textbf{Step 11.} We have constructed a surjective homomorphism from \(\pi_1(S)\) onto the
fundamental group \(G\) of a graph of groups. Its nontrivial vertex groups
are infinite cyclic, its root edges are infinite cyclic with injective
maps, and all point-identification edges are trivial. The root edges
form a star: the roots map to distinct elliptic normalizations. They
therefore contain no graph cycle. The unique graph cycle of the complete
graph contains a point-identification edge, whose edge group is trivial.

\noindent\textbf{Step 12.} Delete one such edge on the cycle. The remaining graph is a connected
tree, and the normal form for iterated amalgamated products along
injective edge maps embeds its \(T\) vertex group \(\mathbb Z\) into its
fundamental group \(G_0\). Restoring the removed edge is an HNN extension
over the trivial group, hence
\[
 G\simeq G_0*\mathbb Z.
\]
Choosing a nontrivial element of the embedded \(\mathbb Z\subset G_0\)
and the new stable letter gives an embedded free group \(F_2\), by reduced
free-product normal form. But every quotient of a virtually abelian group
is virtually abelian, and a virtually abelian group has no \(F_2\)
subgroup. This contradicts P4R.1. The case \(c=g+1,r=g\) is impossible.

\noindent\textbf{Step 13.} We have proved that equality makes \(N\) smooth minimal Hopf. Conversely
P4.24 applies to such an actual normalization and proves the unramified
elliptic-conductor description and \(e(S)=0\). This completes both
directions.
\end{proof}


\subsection{Cuspidal transverse type and the ruled equality case}

\begin{proposition}[P4.28]\label{night:P4.28}\label{stmt:P4.28}
\StatementRecord{night:P4.28}{proved}{Historical events: GPT-6 Astra (OpenAI), ChatGPT Pro mode, effort Pro, S6\_06\_BLIND, 16--18 November 2026: targeted proof check. Also checked in this writing. Pre-continuation writing events: GPT-6 Astra (OpenAI), ChatGPT Pro mode, effort Pro, S6\_10\_PAPERS, 18 November 2026: targeted proof check. Continuation metadata events: GPT-6 Astra (OpenAI), ChatGPT Pro mode, effort Pro, S6\_10\_PAPERS continuation, 19 November 2026: used and rechecked in this run; status-vocabulary reconciliation; metadata reconciliation only; no new mathematical check. \emph{Pre-continuation recorded proof scope.} The night-run re-derivation was by GPT-6 Astra (OpenAI, ChatGPT Pro mode, effort Pro), 16--18 November 2026. Also checked in this writing by GPT-6 Astra (OpenAI), ChatGPT Pro mode, effort Pro, 18 November 2026. The transverse Hessian along the smooth critical curve is treated over the one-variable DVR: a minimal diagonal entry divides the off-diagonal entry, so the square completion is holomorphic. The rank-one case reduces to X\textasciicircum{}2+y\textasciicircum{}3 h; a simple-zero perturbation of h gives the excluded X\textasciicircum{}2+y\textasciicircum{}3 Z through an actual coordinate change. In the rank-zero case a small cubic perturbation, Weierstrass division and a Tschirnhaus change give Y\textasciicircum{}3+x\textasciicircum{}2 H; a simple-zero perturbation of H gives the excluded Y\textasciicircum{}3+x\textasciicircum{}2 Z. On the two nested boundaries the A2 normal form gives the required real submersion, and proper boundary transport carries the compact nonzero-square cycles to the original local fibre. The two coordinate perturbations are used on their specified nested balls. The remaining unit case is the actual product cusp, whose normalization is (u,z) to (u\textasciicircum{}3,-u\textasciicircum{}2,z), with conductor coordinate z unchanged}{\ref{night:P4.16}; \ref{night:P4.18}; \ref{night:P4.20}}

Let $p$ lie on a reduced fibre support of Part~4.2.
Suppose its reduced singular locus is a smooth curve
near $p$, and its generic transverse plane-curve germ
is the cusp $a^2+b^3=0$. Then the actual germ at $p$
is analytically isomorphic to
$\{a^2+b^3=0\}\times(\mathbb C,0)$.
In particular its normalization is smooth.
\end{proposition}
\begin{proof}
\noindent\textbf{Step 1.} We extend the boundary argument of P4.20 in the
precise form used here. Suppose a deformation retains
the smooth critical curve near a fixed transverse
boundary sphere and retains transverse type $A_2$.
Near a boundary point, the holomorphic Morse lemma
in the nondegenerate transverse variable followed
by extraction of a cube root of a unit gives
coordinates with defining function $a^2+b^3$.
The third coordinate parametrizes the critical
curve. Its transversality to the boundary lets
one real third-coordinate variable be solved
from the boundary equation; the complex
coordinates $a,b$ remain free. Thus the
restriction to the boundary is submersive
at every nonzero level. Compactness and
the smooth part of the boundary give the
same uniform nonzero-level submersion as
in P4.20. Its finite-critical-value and
proper Ehresmann transport proof applies
unchanged. A compact spherical class of
nonzero square in a deformed nearby
fibre would therefore lie in an actual
nearby torus fibre, which is impossible.
All deformations below retain this
boundary transverse type.

\noindent\textbf{Step 2.} Choose holomorphic coordinates $(x,y,z)$
so the smooth critical curve is $x=y=0$.
The transverse Hessian along it is a
symmetric rank-one matrix over the
discrete valuation ring $\mathbb C\{z\}$.
It is not the zero matrix generically.
It can be diagonalized holomorphically
to $h(z)x^2$. Here is the elementary
ring calculation. The determinant is
zero. Choose a diagonal entry of
minimal order; one exists since the
square of the off-diagonal entry is
the product of the two diagonal entries.
The ratio of the off-diagonal entry
to the chosen entry is holomorphic.
Completing its square gives the desired
invertible coordinate change. The
coefficient $h$ is not identically zero.

\noindent\textbf{Step 3.} First suppose that the Hessian has
rank one at $p$. The Morse lemma with
parameters $(y,z)$ writes the reduced
equation as
\[
 g=X^2+y^3h_1(y,z).
\]
Indeed after removing the $X$ direction,
the remaining function and its first
two $y$ derivatives vanish on the
critical curve; the third derivative
is not identically zero there. If
$h_1(0,0)\ne0$, extracting its cube
root in the $y$ coordinate proves the
asserted product form.

\noindent\textbf{Step 4.} Otherwise $h_1(0,z)$ has a zero of
finite positive order. Use
$g_\eta=X^2+y^3(h_1+\eta)$.
For arbitrarily small generic $\eta$,
the function $h_1(0,z)+\eta$ has a
simple zero $z_\eta$ near zero.
This follows directly from one-variable
Weierstrass preparation, avoiding the
finite critical values of $h_1(0,z)$.
At $(0,0,z_\eta)$ the function
$Z=h_1(y,z)+\eta$ is a coordinate,
since its $z$ derivative is nonzero.
The actual germ is
$X^2+y^3Z$, which is excluded by
P4.16: its nearby fibre has a
spherical class of square $-6$.
On a fixed sufficiently small boundary,
$h_1(0,z)$ is nonzero at the critical
curve, and this persists for small
$\eta$. The deformed transverse type
is therefore still $A_2$ there.
The preceding boundary transport gives
the contradiction.

\noindent\textbf{Step 5.} It remains to exclude Hessian rank
zero at $p$, so $h(0)=0$. In the
diagonalizing coordinates add
$\varepsilon y^3$ to $g$. The
critical curve and its transverse
Hessian do not change. The cubic
coefficient in the transverse kernel
remains nonzero near the original
compact boundary for small
$\varepsilon$, so this deformation
has the boundary property just proved.
Choose $\varepsilon$ so that the
coefficient of $y^3$ at the origin
is nonzero.

\noindent\textbf{Step 6.} On a smaller neighbourhood, Weierstrass
preparation in $y$ gives a unit times
\[
 P=y^3+a(x,z)y^2+b(x,z)y+c(x,z).
\]
Restriction to $x=0$ is a unit times
$y^3$, so all three coefficients
vanish at $x=0$. The critical-curve
equation $g_x(0,0,z)=0$ first gives
$c_x(0,z)=0$. In the identity
$g=U P$ the mixed derivative on
the critical curve is then
$g_{xy}=U b_x$, since the terms
containing $c,c_x,b$ vanish there.
The zero $xy$ Hessian entry therefore
gives $b_x(0,z)=0$. Thus
$x^2\mid b$, $x^2\mid c$, and
$x\mid a$.
The Tschirnhaus change $y=Y-a/3$
therefore gives exactly
\[
 P=Y^3+x^2H(x,Y,z).
\]
The coefficient $H(0,0,z)$ is a unit
multiple of the original Hessian
coefficient $h(z)$. It is not
identically zero but vanishes at
zero. Perturb it to $H+\eta$.
At a simple root of
$H(0,0,z)+\eta$, the coordinates
$(x,Y,Z=H+\eta)$ give
$Y^3+x^2Z$. P4.18 excludes this
binomial germ by its compact
spherical class of nonzero square.

\noindent\textbf{Step 7.} There are two nested transport
neighbourhoods in this argument.
First fix the original ball for
$g\to g+\varepsilon y^3$.
Then choose a smaller ball in the
Weierstrass coordinate neighbourhood,
where $g+\varepsilon y^3$ has only
the original critical curve and its
boundary has transverse type $A_2$.
The second perturbation transports
its spherical class into a nearby
fibre on this smaller ball. Inclusion
puts it in a fibre on the original
ball; the first boundary transport
puts it in a fibre of $g$.
This does not require extending the
second perturbation outside its
coordinate neighbourhood.
Units in the reduced equations are
retained during the perturbations.
At a binomial special point a unit
can be absorbed by the positive
weighted coordinate scaling of
P4.16--P4.18.

\noindent\textbf{Step 8.} We have ruled out both a zero
Hessian and a vanishing cubic
coefficient. The product form
follows. Its normalization is
the finite map
\[
 (u,z)\longmapsto (a=u^3,b=-u^2,z),
\]
whose source is smooth, whose
inverse rational parameter off
$b=0$ is $u=-a/b$, and whose
restriction to the conductor is
the identity on the $z$ coordinate.
\end{proof}

\begin{proposition}[P4.29]\label{night:P4.29}\label{stmt:P4.29}
\StatementRecord{night:P4.29}{proved}{Historical events: GPT-6 Astra (OpenAI), ChatGPT Pro mode, effort Pro, S6\_06\_BLIND, 16--18 November 2026: targeted proof check. Also checked in this writing. Pre-continuation writing events: GPT-6 Astra (OpenAI), ChatGPT Pro mode, effort Pro, S6\_10\_PAPERS, 18 November 2026: targeted proof check. Continuation metadata events: GPT-6 Astra (OpenAI), ChatGPT Pro mode, effort Pro, S6\_10\_PAPERS continuation, 19 November 2026: used and rechecked in this run; status-vocabulary reconciliation; metadata reconciliation only; no new mathematical check. \emph{Pre-continuation recorded proof scope.} The night-run re-derivation was by GPT-6 Astra (OpenAI, ChatGPT Pro mode, effort Pro), 16--18 November 2026. Also checked in this writing by GPT-6 Astra (OpenAI), ChatGPT Pro mode, effort Pro, 18 November 2026. The zero-Euler cusp graph, descent of the ruling, rationality of normalization singularities, Cohen-Macaulay reduced conductor, and vanishing conductor delta were re-derived. The product cusp normal form then gives the actual smooth minimal normalization}{\ref{night:P4.12}; \ref{night:P4.13}; \ref{night:P4.22}; \ref{night:P4.28}; \ref{night:P4R.7}; \ref{night:P4R.9}}

Let $S$ be an irreducible singular support
in Part~4.2 whose normalization has ruled
resolution, and suppose $e(S)=0$.
In the double-horizontal-divisor case,
when the generic transverse type is
cuspidal $A_2$, the normalization is
a smooth minimal elliptic ruled surface.
The image conductor is a smooth elliptic
curve and every point along it has the
product cusp germ of P4.28.
\end{proposition}
\begin{proof}
\noindent\textbf{Step 1.} Use the notation of P4R.7 and P4R.9.
The actual resolved pole support is
a section $H$ with rational trees,
and its inclusion in the ruled
resolution $T$ induces an isomorphism
on fundamental groups. Thus
$\pi_1(S)=\pi_1(C)$, where
$C=\operatorname{Sing}(S)_{\rm red}$.
The virtually abelian curve-group
calculation gives $e(C)\ge0$, and
P4.13 gives $e(S_{\rm reg})\ge0$.
Equality makes both numbers zero.
The same connected-curve calculation
shows that $C$ has exactly one
normalized component, of genus one,
and has no point-identification
cycles.

\noindent\textbf{Step 2.} All rational components of the
resolved pole support are contracted
to points: a nonconstant proper map
from a rational curve to the
elliptic normalization of $C$ is
impossible. Hence the strict
conductor consists only of $H$.
The generic cusp has one
normalization branch, so
$H\to C$ has generic degree one
and is the normalization of $C$.
In particular $H$ and the ruling
base both have genus one.

\noindent\textbf{Step 3.} Every exceptional curve of
$T\to N$ is now rational and
vertical for the ruling, since
$H$ is not contracted. The ruling
$r:T\to H$ descends to a holomorphic
map $\bar r:N\to H$. For an explicit
descent, the image of
$(\rho_N,r):T\to N\times H$
projects properly and bijectively
to $N$, because $r$ is constant
on each exceptional fibre. Its
normalized graph projects finitely
and birationally to the normal
space $N$, hence is isomorphic
to $N$.

\noindent\textbf{Step 4.} The normalization singularities
are rational. Indeed a sequence
of point blowups of a
$\mathbb P^1$-bundle has
$R^1r_*\mathcal O_T=0$.
For a blowup this follows from
the two affine charts, or from
the exact sequence of the
exceptional $\mathbb P^1$ with
$H^1(\mathcal O_{\mathbb P^1})=0$;
the assertion for the whole
sequence follows by Leray.
Since $\rho_{N*}\mathcal O_T
=\mathcal O_N$, the low-degree
Leray sequence for
$r=\bar r\circ\rho_N$ gives
\[
 \bar r_*R^1\rho_{N*}\mathcal O_T=0.
\]
Here $R^2\bar r_*\mathcal O_N=0$
because $\bar r$ has one-dimensional
fibres. The sheaf
$R^1\rho_{N*}\mathcal O_T$ is
supported at finitely many
normalization singular points,
and direct image of a nonzero
such sheaf to $H$ is nonzero.
Thus it vanishes. In particular
$\chi(\mathcal O_N)
=\chi(\mathcal O_T)=0$.

\noindent\textbf{Step 5.} We next check the scheme structure
of the conductor, to prove that
its image curve is smooth.
Write $I$ for the conductor ideal
on $S$ and on $N$, and let
$C_{\rm sch}\subset S$ and
$D_{\rm sch}\subset N$ be its
two zero schemes. Locally
$\mathcal O_S$ is Gorenstein
of dimension two and
$\nu_*\mathcal O_N$ is a
maximal Cohen--Macaulay module.
Finite duality identifies
its dual module with $I$,
which is again maximal
Cohen--Macaulay. The depth
lemma in
$0\to I\to\mathcal O_S
\to\mathcal O_{C_{\rm sch}}\to0$
therefore gives depth one
for the conductor curve.
At a generic cusp its local
length in the transverse
direction is one. A
one-dimensional Cohen--Macaulay
scheme which is generically
reduced is reduced: any
nilpotent ideal would have
zero-dimensional support,
contradicting the absence of
embedded associated points.
Consequently $C_{\rm sch}=C$.

\noindent\textbf{Step 6.} Put $L=\nu^*\omega_S$, a
torsion line bundle by P4.12.
Finite duality on $N$ gives
\[
 I=\omega_N\otimes L^{-1}.
\]
Serre duality on the compact
Cohen--Macaulay surface $N$
gives $\chi(I)=\chi(L)$.
The rational-singularity
comparison just proved and
the projection formula give
$\chi(L)=\chi(\rho_N^*L)$.
Riemann--Roch on the smooth
surface $T$, with the rational
Chern class of this torsion
line bundle equal to zero,
gives $\chi(\rho_N^*L)
=\chi(\mathcal O_T)=0$.
Therefore
$\chi(\mathcal O_{D_{\rm sch}})
=\chi(\mathcal O_N)-\chi(I)=0$.

\noindent\textbf{Step 7.} The exact conductor square is
\[
 0\longrightarrow\mathcal O_S
 \longrightarrow
 \nu_*\mathcal O_N\oplus\mathcal O_C
 \longrightarrow\nu_*\mathcal O_{D_{\rm sch}}
 \longrightarrow0.
\]
Taking Euler characteristics and
using $\chi(\mathcal O_S)=0$
from P4.12 yields
$\chi(\mathcal O_C)=0$.
Its normalization is the
smooth elliptic curve $H$,
so the normalization sequence
gives $\chi(\mathcal O_C)
 =-\sum_p\delta(C,p)$.
All delta invariants vanish;
thus $C$ is smooth.

\noindent\textbf{Step 8.} Every point of $C$ now satisfies
P4.28, so $S$ has the product
cusp germ there and $N$ is
smooth along the conductor.
It was already smooth away
from the conductor. P4.22
therefore applies and makes
this elliptic ruled
normalization minimal.
\end{proof}


\subsection{The conormal calculation and ordinary-triple equality}

\begin{proposition}[P4.30]\label{night:P4.30}\label{stmt:P4.30}
\StatementRecord{night:P4.30}{proved}{Historical events: GPT-6 Astra (OpenAI), ChatGPT Pro mode, effort Pro, S6\_06\_BLIND, 16--18 November 2026: targeted proof check. Also checked in this writing. Pre-continuation writing events: GPT-6 Astra (OpenAI), ChatGPT Pro mode, effort Pro, S6\_10\_PAPERS, 18 November 2026: targeted proof check. Continuation metadata events: GPT-6 Astra (OpenAI), ChatGPT Pro mode, effort Pro, S6\_10\_PAPERS continuation, 19 November 2026: used and rechecked in this run; status-vocabulary reconciliation; metadata reconciliation only; no new mathematical check. \emph{Pre-continuation recorded proof scope.} The night-run re-derivation was by GPT-6 Astra (OpenAI, ChatGPT Pro mode, effort Pro), 16--18 November 2026. Also checked in this writing by GPT-6 Astra (OpenAI), ChatGPT Pro mode, effort Pro, 18 November 2026. Differentiation on the elliptic normalization has image Omega\_E(-R) and a primitive rank-two kernel K of degree deg R. The symbolic conormal Q=P/P\textasciicircum{}(2) injects into j\_*K with finite quotient of length tau. At a unibranch germ of multiplicity m the derivative common order is m-1. The Pluecker minors of the gradient submodule bound the colength below by 2m in embedding dimension three and by m in the plane case; a smooth germ contributes zero. Hence tau>=deg R and chi(Q)<=0. The exact symbolic-square sequence gives chi(O/P\textasciicircum{}(2))<=-sum delta. Equality forces every delta to vanish, so the conductor curve is smooth. The bound is derived from the explicit gradient and normalization maps}{the hypotheses and explicit local calculation stated here}

Let $C$ be a compact irreducible reduced curve in the smooth
threefold $Y$ of Part~4.2. Suppose its normalization is an elliptic
curve $E$, and every point of $C$ is unibranch. Write $P$ for its
ideal in $Y$ and $P^{(2)}$ for the ideal of functions vanishing to
order at least two at the generic point of $C$. Then
\[
 \chi(\mathcal O_Y/P^{(2)})\le0.
\]
Equality forces $C$ to be smooth. The assertion uses
$c_1(TY)\cdot C=0$, which follows from $b_2(Y)=0$.
\end{proposition}
\begin{proof}
\noindent\textbf{Step 1.} Let $j:E\to Y$ be the normalization map followed by inclusion.
The derivative map
\[
 j^*\Omega_Y\longrightarrow\Omega_E
\]
has image $\Omega_E(-R)$ for an effective divisor $R$.
Its kernel $K$ is locally free of rank two: dividing the derivative
row by its common zero gives a surjection onto a line bundle.
Since $\deg j^*\Omega_Y=0$ and $\deg\Omega_E=0$, we have
$\deg K=\deg R$.

\noindent\textbf{Step 2.} Differentiation induces an injection
\[
 Q:=P/P^{(2)}\longrightarrow j_*K.
\]
Indeed its kernel consists precisely of functions with zero first
normal derivative at a generic smooth point of $C$, which is the
defining order-two condition for $P^{(2)}$. Both sides agree at the
generic point. Their quotient has finite length, denoted by $\tau$.
Riemann--Roch on the smooth elliptic curve gives
\[
 \chi(Q)=\deg R-\tau.
\]
We prove $\tau\ge\deg R$ by local equations, including at a
singular point of $C$.

\noindent\textbf{Step 3.} Parametrize its sole branch by $t$. Choose ambient coordinates
vanishing at the point, and let
$m=\min_i\operatorname{ord}_t x_i(t)$ be the multiplicity.
The common order of the three derivative entries is $m-1$,
so the coefficient of $R$ is $m-1$. The saturated kernel $K$
is a direct summand of the free rank-three module over
$\mathbb C\{t\}$. Consequently a generator of $\bigwedge^2K$
has at least one unit Pl\"ucker coordinate. If $L$ is the
$\mathbb C\{t\}$-span of the differentiated defining equations,
then the colength of $L$ in $K$ is the minimum order of their
two-by-two minors. This is the diagonal elementary-divisor
calculation for a rank-two submodule over a discrete valuation
ring. The local length of $j_*K/Q$ is at least this colength,
since $Q\subset L\subset K$.

\noindent\textbf{Step 4.} At a smooth point $m=1$ and there is nothing to prove.
At a singular point of embedding dimension three, every defining
equation belongs to the square of the ambient maximal ideal.
Every entry of its gradient therefore has order at least $m$
on the parametrization, and every two-by-two minor has order
at least $2m$. Hence the local quotient length is at least
$2m\ge m-1$.

\noindent\textbf{Step 5.} At a singular point of embedding dimension two, an invertible
ambient coordinate change writes the ideal as $(z,h(x,y))$,
where $h$ is a reduced irreducible plane equation without a
linear term. The conormal image has a free $dz$ summand.
Its other generator is $dh$; both its plane coefficients have
order at least $m$. The primitive rank-one kernel of the plane
derivative row has a unit coefficient, so $dh$ is its generator
multiplied by a function of order at least $m$. The quotient
length is therefore at least $m\ge m-1$. Embedding dimension
one is the smooth case. These cases prove the local bound and
therefore $\chi(Q)\le0$.

\noindent\textbf{Step 6.} The normalization exact sequence gives
$\chi(\mathcal O_C)=-\sum_p\delta(C,p)$.
Finally the exact sequence
\[
 0\longrightarrow P/P^{(2)}\longrightarrow
 \mathcal O_Y/P^{(2)}\longrightarrow\mathcal O_C
 \longrightarrow0
\]
gives
\[
 \chi(\mathcal O_Y/P^{(2)})
 =-\sum_p\delta(C,p)+\deg R-\tau\le0.
\]
Equality forces all delta invariants to vanish, so the
normalization is an isomorphism and $C$ is smooth.
\end{proof}

\begin{proposition}[P4.31]\label{night:P4.31}\label{stmt:P4.31}
\StatementRecord{night:P4.31}{proved}{Historical events: GPT-6 Astra (OpenAI), ChatGPT Pro mode, effort Pro, S6\_06\_BLIND, 16--18 November 2026: targeted proof check. Also checked in this writing. Pre-continuation writing events: GPT-6 Astra (OpenAI), ChatGPT Pro mode, effort Pro, S6\_10\_PAPERS, 18 November 2026: targeted proof check. Continuation metadata events: GPT-6 Astra (OpenAI), ChatGPT Pro mode, effort Pro, S6\_10\_PAPERS continuation, 19 November 2026: used and rechecked in this run; status-vocabulary reconciliation; metadata reconciliation only; no new mathematical check. \emph{Pre-continuation recorded proof scope.} The night-run re-derivation was by GPT-6 Astra (OpenAI, ChatGPT Pro mode, effort Pro), 16--18 November 2026. Also checked in this writing by GPT-6 Astra (OpenAI), ChatGPT Pro mode, effort Pro, 18 November 2026. The graph and ruled-pole bounds give a single unibranch conductor image with elliptic normalization and a degree-three horizontal section of genus g. The rational-exceptional contraction descends the ruling to the normalization and gives chi(O\_N)=1-g. The conductor image is pure Cohen-Macaulay with one associated prime; its generic ordinary-triple calculation is P\textasciicircum{}2, so primary localization and saturation identify its global ideal with P\textasciicircum{}(2). Finite duality, torsion adjunction and Serre duality give chi(D\_conductor)=0 and chi(C\_conductor)=g-1. P4.30 forces g=1 and a smooth image; the degree-three section map is unramified. Distinct section points cannot be joined by a vertical exceptional block. Each of the three hypersurface branches is normal, and its pairwise intersection with another branch is the reduced smooth Cartier curve; this forces each branch regular and pairwise transverse. The explicit coordinate changes take the third equation a x+b y to X-Y, yielding XY(X-Y), and the normalization is consequently smooth minimal elliptic ruled}{\ref{night:P4.22}; \ref{night:P4.29}; \ref{night:P4.30}; \ref{night:P4R.6}; \ref{night:P4R.7}; \ref{night:P4R.9}}

Let $S$ be an irreducible singular support in Part~4.2 whose
normalization has ruled resolution, and suppose $e(S)=0$.
In the double-horizontal-divisor case, when the generic
transverse type is the ordinary triple point $xy(x-y)=0$,
the normalization is a smooth minimal elliptic ruled surface.
Its conductor image is a smooth elliptic curve, and its
reduced conductor is an unramified cover of that curve of
degree three. Every germ along the image is analytically
the product of an ordinary plane triple point with a disc.
\end{proposition}
\begin{proof}
\noindent\textbf{Step 1.} Use the minimal resolution $T\to N\to S$ and the section $H$
of P4R.7 and P4R.9. Put $g=g(H)$ and
$C=\operatorname{Sing}(S)_{\rm red}$.
The same fundamental-group and Euler argument as in P4.29
shows that $C$ is connected, has just one normalized
component, that this normalization $E$ has genus one, and
that its point-identification graph has no cycles. Thus
every point of $C$ is unibranch. All rational components
of the resolved pole support are contracted, and the
only strict conductor component is $H$. Its finite map
to $C$ has degree three, since an ordinary transverse
triple point has three normalization branches.
In particular $g\ge1$ by Riemann--Hurwitz for $H\to E$.

\noindent\textbf{Step 2.} All exceptional curves of $T\to N$ are rational and
vertical for the ruling. Exactly the normalized-graph
descent and Leray calculation of P4.29 give a morphism
$N\to H$ and
\[
 R^1\rho_{N*}\mathcal O_T=0,\qquad
 \chi(\mathcal O_N)=\chi(\mathcal O_T)=1-g.
\]
Thus the normalization singularities are rational; no
smoothness assertion is used here.

\noindent\textbf{Step 3.} Let $I$ be the conductor ideal, and let $C_{\rm sch}$
and $D_{\rm sch}$ be its zero schemes in $S$ and $N$.
Finite duality and the depth argument of P4.29 show that
$\mathcal O_{C_{\rm sch}}$ is Cohen--Macaulay of pure
dimension one. If $P\subset\mathcal O_Y$ is the reduced
ideal of $C$, the inverse image $J$ of $I$ in
$\mathcal O_Y$ therefore has no embedded associated
prime. Its only associated prime locally is $P$.
The generic ordinary triple point has conductor equal
to the square of its maximal ideal in the transverse
plane. To check this directly, normalize its three
lines. Constant functions on them have a common value;
their linear terms are the evaluations of one linear
form on three distinct directions, a two-dimensional
space. Quadratic forms evaluate surjectively on those
three directions. All terms of degree at least two
can therefore be chosen independently on the three
branches, and a nonzero constant or linear term cannot
belong to the conductor. This proves the stated square
ideal, including its scheme structure.

\noindent\textbf{Step 4.} Consequently $J_P=P_P^2$, and in fact
\[
 J=P^{(2)}.
\]
For completeness, if $a\in P^{(2)}$, some $s\notin P$
satisfies $sa\in J$. Multiplication by $s$ is injective
on $\mathcal O_Y/J$, because its only associated prime
is $P$, so $a\in J$. The reverse inclusion follows by
localizing at $P$. Thus the conductor image is the
actual symbolic double curve to which P4.30 applies.

\noindent\textbf{Step 5.} Put $L=\nu^*\omega_S$, which is torsion. Finite duality
gives $I=\omega_N\otimes L^{-1}$. Serre duality, the
rational-singularity comparison and Riemann--Roch on
$T$ give
\[
 \chi(I)=\chi(L)=\chi(\rho_N^*L)
   =\chi(\mathcal O_T)=1-g.
\]
Hence $\chi(\mathcal O_{D_{\rm sch}})=0$.
Taking Euler characteristics in the exact conductor
square, and using $\chi(\mathcal O_S)=0$, yields
\[
 \chi(\mathcal O_{C_{\rm sch}})=g-1.
\]
P4.30 now gives $g-1\le0$. Since $g\ge1$, equality
holds, $g=1$, and $C$ is smooth. Riemann--Hurwitz then
makes the degree-three map $H\to C$ unramified.

\noindent\textbf{Step 6.} We finish by proving normalization smoothness from
the actual local morphisms. Over a point of $C$ the
three distinct points of $H$ remain distinct in $N$:
every exceptional block is vertical for the ruling
$T\to H$ and therefore cannot join different points
of $H$. The finite normalization consequently gives
three local irreducible branches of $S$. Each branch
is a hypersurface in $Y$ and is smooth at the generic
point of $C$. Its singularities are isolated, so the
hypersurface Cohen--Macaulay property and regularity
in codimension one make it normal. The corresponding
component of $N$ is that branch itself.

\noindent\textbf{Step 7.} Write its three local equations as $f_1,f_2,f_3$.
The intersection of any two branches is supported on
$C$ and has generic multiplicity one, because the
generic transverse triple point is ordinary. On a
normal branch the other equation is a nonzero divisor.
Its quotient is a one-dimensional Cohen--Macaulay
ring, generically reduced with support $C$, so it is
reduced and is exactly $\mathcal O_C$. Thus $C$ is a
smooth Cartier curve in each branch. A local ring of
dimension two with a nonzero divisor whose quotient
is a regular one-dimensional local ring has embedding
dimension at most two: lift its one maximal-ideal
generator and add the nonzero divisor. It is regular.
All three branches are therefore smooth.

\noindent\textbf{Step 8.} Any two branches meet transversely along $C$.
Choose their equations as coordinates $x,y$, with
$z$ a coordinate on $C$. The third equation is
$a(x,y,z)x+b(x,y,z)y$, where both $a$ and $b$ are
units by the pairwise transversality just proved.
The invertible coordinate change $X=ax$, $Y=-by$
writes the reduced union exactly as
$XY(X-Y)=0$. This is the asserted product triple
point, and its normalization is smooth. Away from
the conductor, normalization singularities were
already excluded by P4R.6. Finally P4.22 applies to
this smooth elliptic ruled normalization and proves
its minimality.
\end{proof}


\subsection{Complete zero-Euler classification in the requested classes}

\begin{proposition}[P4.32]\label{night:P4.32}\label{stmt:P4.32}
\StatementRecord{night:P4.32}{proved}{Historical events: GPT-6 Astra (OpenAI), ChatGPT Pro mode, effort Pro, S6\_06\_BLIND, 16--18 November 2026: targeted proof check. Also checked in this writing. Pre-continuation writing events: GPT-6 Astra (OpenAI), ChatGPT Pro mode, effort Pro, S6\_10\_PAPERS, 18 November 2026: targeted proof check. Continuation metadata events: GPT-6 Astra (OpenAI), ChatGPT Pro mode, effort Pro, S6\_10\_PAPERS continuation, 19 November 2026: used and rechecked in this run; status-vocabulary reconciliation; metadata reconciliation only; no new mathematical check. \emph{Pre-continuation recorded proof scope.} The night-run re-derivation was by GPT-6 Astra (OpenAI, ChatGPT Pro mode, effort Pro), 16--18 November 2026. Also checked in this writing by GPT-6 Astra (OpenAI), ChatGPT Pro mode, effort Pro, 18 November 2026. The zero-Euler normalization and conductor alternatives were checked against their stated local inputs. The local conclusion uses the single support, ambient b1=b2=0, and nearby tori; the final global assertion retains the source hypothesis that every support is irreducible or NC}{\ref{night:P4.12}; \ref{night:P4.22}; \ref{night:P4.24}; \ref{night:P4.26}; \ref{night:P4.27}; \ref{night:P4.29}; \ref{night:P4.31}; \ref{night:P4.8}; \ref{night:P4R.11}; \ref{night:P4R.7}; \ref{night:P4R.9}}

An irreducible support in Part~4.2 has Euler number zero exactly
in the following cases: it is smooth with a finite unramified
torus cover, or it is singular with smooth minimal elliptic ruled
normalization, or it is singular with smooth minimal Hopf
normalization. In the ruled double-horizontal case the conductor
image is a smooth elliptic curve; the local germs are product
cusps, product tacnodes, or product ordinary plane triple points,
and the reduced conductor maps have degrees one, two, or three
respectively and are unramified.

For a fibration whose supports are all irreducible or normal
crossing, $c_3(Y)=0$ exactly when every irreducible support has
one of these forms and every reducible support has the
zero-Euler normal-crossing form of P4.8.
\end{proposition}
\begin{proof}
\noindent\textbf{Step 1.} Normal irreducible supports are smooth with the stated finite
unramified cover by P4.8 and P4.12. For a singular irreducible
support, use its minimal normalization resolution $T$.
P4.26 excludes rational resolutions and class-VII minimal
models with positive second Betti number from equality.
The Inoue minimal model has no conductor curves and is
excluded there as well. A Hopf minimal model is covered
completely by P4.27: equality makes the actual normalization
a smooth minimal Hopf surface, and the converse is P4.24.

\noindent\textbf{Step 2.} It remains to use the positive-genus ruled resolution.
The reduced-horizontal alternatives of P4R.9 already give
smooth minimal elliptic ruled normalization. In the remaining
case the pole support is a section $H$ with rational trees,
and the coefficient of $H$ is two. Its inclusion in $T$
induces an isomorphism on fundamental groups. The conductor
pushout and the curve-group calculation of P4R.7--P4R.9 give
$\pi_1(S)=\pi_1(C)$ and $e(C)\ge0$ for the reduced image
conductor $C$. Equality in
$e(S)=e(S_{\rm reg})+e(C)$ makes $C$ a connected curve
with elliptic normalization and no point-identification
cycles. Every rational pole component is contracted, since
there is no finite nonconstant map from a rational curve
to an elliptic curve. Thus $H$ is the sole strict conductor
component, and every branch of a generic transverse
plane-curve singularity has conductor exponent two.

\noindent\textbf{Step 3.} We enumerate those transverse germs explicitly. For a branch
$B_i$ of a reduced plane curve, its conductor exponent is
\[
 c_i=2\delta(B_i)+\sum_{j\ne i}I(B_i,B_j).
\]
If a branch is singular, $c_i=2$ forces
$\delta(B_i)=1$ and no other branch. The one-gap
parametrization of a delta-one irreducible plane curve
has orders two and three and gives the ordinary cusp.
If all branches are smooth, two branches have intersection
multiplicity two and give a tacnode, or three branches
have all pairwise intersection multiplicities one and
give an ordinary triple point. Four branches would give
each conductor exponent at least three. One smooth branch
is nonsingular and has exponent zero. These exhaust the
possibilities.

\noindent\textbf{Step 4.} P4.29, P4R.11, and P4.31 respectively treat the cusp,
tacnode, and ordinary-triple cases, without assuming that
the initial normalization or image conductor was smooth.
They give the asserted smooth minimal elliptic ruled
normalization and the precise local conductor maps.
The converse for a smooth minimal elliptic ruled
normalization is P4.22. This proves the entire irreducible
zero-Euler classification.

\noindent\textbf{Step 5.} Finally P4.26 expresses $c_3(Y)$ as a finite sum of
nonnegative fibre Euler numbers. Its vanishing is therefore
equivalent to vanishing of every summand. The irreducible
summands have just been classified, and P4.8 gives the
complete reducible normal-crossing zero case. This proves
the final assertion.
\end{proof}


\subsection{Normalization cohomology for the remaining positive cases}

\begin{proposition}[P4.33]\label{night:P4.33}\label{stmt:P4.33}
\StatementRecord{night:P4.33}{proved}{Historical events: GPT-6 Astra (OpenAI), ChatGPT Pro mode, effort Pro, S6\_06\_BLIND, 16--18 November 2026: targeted proof check. Also checked in this writing. Pre-continuation writing events: GPT-6 Astra (OpenAI), ChatGPT Pro mode, effort Pro, S6\_10\_PAPERS, 18 November 2026: targeted proof check. Continuation metadata events: GPT-6 Astra (OpenAI), ChatGPT Pro mode, effort Pro, S6\_10\_PAPERS continuation, 19 November 2026: used and rechecked in this run; status-vocabulary reconciliation; metadata reconciliation only; no new mathematical check. \emph{Pre-continuation recorded proof scope.} The night-run re-derivation was by GPT-6 Astra (OpenAI, ChatGPT Pro mode, effort Pro), 16--18 November 2026. Also checked in this writing by GPT-6 Astra (OpenAI), ChatGPT Pro mode, effort Pro, 18 November 2026. Finite duality, the coherent conductor Euler relation, the two possibilities for the horizontal component, descent of the ruling, and the Leray length calculation were re-derived}{\ref{night:P4.12}; \ref{night:P4.26}; \ref{night:P4.29}}

Let $S$ be an irreducible singular support in Part~4.2, and let
$\rho:T\to N\to S$ be its minimal normalization resolution.
Then $H^2(N,\mathcal O_N)=0$. If
$p=\dim H^0(N,R^1\rho_{N*}\mathcal O_T)$ is the sum of the
geometric genera of the normalization singularities, then
\[
 p=h^1(T,\mathcal O_T)-h^1(N,\mathcal O_N).
\]
In particular a rational resolution forces every normalization
singularity to be rational and gives $\chi(\mathcal O_N)=1$.
For the scheme-theoretic image conductor $C_{\rm sch}$ we have
\[
 \chi(\mathcal O_{C_{\rm sch}})
   =-\chi(\mathcal O_N)=h^1(N,\mathcal O_N)-1.
\]

In the ruled double-horizontal case, with section $H$ of genus
$g\ge1$, a contracted $H$ gives $h^1(N,\mathcal O_N)=0$ and
$p=g$. A noncontracted $H$ gives $p=0$ and
$h^1(N,\mathcal O_N)=g$.
\end{proposition}
\begin{proof}
\noindent\textbf{Step 1.} Write $I_S\subset\mathcal O_S$ for the conductor ideal.
Finite duality gives
\[
 \nu_*\omega_N=I_S\otimes\omega_S.
\]
The dualizing line bundle $\omega_S$ is torsion by P4.12.
A nonzero section of a torsion line bundle on a compact
irreducible reduced complex space has no zero: raising it
to a trivializing positive power gives a nonzero holomorphic
function, hence a nonzero constant. Such a section cannot
belong to the proper conductor ideal along the nonempty
conductor. Therefore $H^0(N,\omega_N)=0$. Serre duality
on the compact Cohen--Macaulay normal surface gives
$H^2(N,\mathcal O_N)=0$.

\noindent\textbf{Step 2.} The resolution is birationally ruled or of class VII and
has $H^2(T,\mathcal O_T)=0$, by P4.12 and the classification
already used in P4.26. The low-degree Leray sequence,
using $\rho_{N*}\mathcal O_T=\mathcal O_N$ and the
finite support of $R^1\rho_{N*}\mathcal O_T$, is therefore
\[
 0\longrightarrow H^1(N,\mathcal O_N)
 \longrightarrow H^1(T,\mathcal O_T)
 \longrightarrow H^0(N,R^1\rho_{N*}\mathcal O_T)
 \longrightarrow0.
\]
This proves the formula for $p$. If $T$ is rational,
$h^1(\mathcal O_T)=0$, so $p=0$. Each nonnegative local
geometric genus is zero, giving the asserted rational
singularities and Euler characteristic.

\noindent\textbf{Step 3.} For completeness the conductor Euler identity does not
require rational singularities. Put $L=\nu^*\omega_S$,
and write $I_N=\omega_N\otimes L^{-1}$ for the conductor
on $N$. Serre duality gives $\chi(I_N)=\chi(L)$.
Riemann--Roch on $T$, applied to the torsion line bundle
$\rho_N^*L$, gives
$\chi(\rho_N^*L)=\chi(\mathcal O_T)$.
The projection formula makes
$R^1\rho_{N*}\rho_N^*L=L\otimes R^1\rho_{N*}\mathcal O_T$;
its finite length is still $p$. Leray therefore gives
$\chi(L)=\chi(\mathcal O_N)$.
Consequently the conductor scheme $D_{\rm sch}\subset N$
has $\chi(\mathcal O_{D_{\rm sch}})=0$.
The exact conductor square, together with
$\chi(\mathcal O_S)=0$, proves
$\chi(\mathcal O_{C_{\rm sch}})=-\chi(\mathcal O_N)$.
Since $N$ is compact connected and $H^2(\mathcal O_N)=0$,
the remaining equality follows.

\noindent\textbf{Step 4.} Finally suppose the ruled pole support is $H$ with
rational vertical trees. If $H$ is contracted, the
composition $H\to T\to N$ is constant. Pullback from
$H^1(N,\mathcal O_N)$ to $H^1(H,\mathcal O_H)$ is zero.
But restriction
$H^1(T,\mathcal O_T)\to H^1(H,\mathcal O_H)$ is an
isomorphism: the ruling has section $H$, a
$\mathbb P^1$-bundle has zero first higher direct image
of its structure sheaf, and the actual point blowups
preserve these groups. The Leray injection above
therefore forces $H^1(N,\mathcal O_N)=0$, so $p=g$.
If $H$ is not contracted, every exceptional curve is
vertical. The ruling descends by the normalized-graph
comparison of P4.29. Its Leray calculation gives $p=0$,
and the same exact sequence gives $h^1(\mathcal O_N)=g$.
\end{proof}


\subsection{The value one in the rational ordinary-triple sector}

\begin{proposition}[P4.34]\label{night:P4.34}\label{stmt:P4.34}
\StatementRecord{night:P4.34}{proved}{Historical events: GPT-6 Astra (OpenAI), ChatGPT Pro mode, effort Pro, S6\_06\_BLIND, 16--18 November 2026: targeted proof check. Also checked in this writing. Pre-continuation writing events: GPT-6 Astra (OpenAI), ChatGPT Pro mode, effort Pro, S6\_10\_PAPERS, 18 November 2026: targeted proof check. Continuation metadata events: GPT-6 Astra (OpenAI), ChatGPT Pro mode, effort Pro, S6\_10\_PAPERS continuation, 19 November 2026: used and rechecked in this run; status-vocabulary reconciliation; metadata reconciliation only; no new mathematical check. \emph{Pre-continuation recorded proof scope.} The night-run re-derivation was by GPT-6 Astra (OpenAI, ChatGPT Pro mode, effort Pro), 16--18 November 2026. Also checked in this writing by GPT-6 Astra (OpenAI), ChatGPT Pro mode, effort Pro, 18 November 2026. The symbolic double-curve Euler formula, the planar defect 5 delta, the nonplanar valuation bounds, and the exclusion of the apparent equality-six unibranch case were re-derived. The final exclusion of Euler one uses exactly the stated global upper bound}{\ref{night:P4.26}; \ref{night:P4.30}; \ref{night:P4.31}; \ref{night:P4.33}; \ref{night:P4.8}}

An irreducible singular support with rationally resolved
normalization and generic transverse ordinary triple point
cannot have Euler number one. For a fibration all of whose
supports are irreducible or normal crossing, its Euler number
is therefore two.
\end{proposition}
\begin{proof}
\noindent\textbf{Step 1.} Suppose its Euler number were one. The rational
normal-crossing case has Euler number two by P4.8, so we
are in the simply connected resolved-conductor alternative
of P4.26. That proof gives $\pi_1(S)=\pi_1(C)=F_b$,
$b\le1$, and
$e(C)=c+1-b$ for the number $c$ of normalized rational
components. Since $e(S_{\rm reg})\ge0$, the value one
forces $c=b=1$ and $e(S_{\rm reg})=0$.
Thus $C$ has normalization $E=\mathbb P^1$, exactly one
point with two analytic branches, and only unibranch
singular points elsewhere.

\noindent\textbf{Step 2.} P4.33 proves that $N$ has rational singularities and
that its image conductor scheme has
$\chi(\mathcal O_{C_{\rm sch}})=-1$.
The generic ordinary-triple conductor and pure
Cohen--Macaulay argument of P4.31 identify this scheme
with the symbolic double curve $\mathcal O_Y/P^{(2)}$.
At a point with two analytic branches, the symbolic ideal
means order at least two along each branch. The same
argument applies: a nonzero finite-support defect in the
conductor quotient would violate depth one.

\noindent\textbf{Step 3.} We compute its Euler characteristic and the possible
local defects. On $E$ use the bundles $K$ and the
ramification divisor $R$ of P4.30. This time
$\deg K=2+\deg R$, since $E=\mathbb P^1$.
Write $\tau_p$ for the local length of the conormal
quotient $j_*K/(P/P^{(2)})$ and $r_p$ for the sum of
the coefficients of $R$ over $p$. Normalization and
Riemann--Roch give
\[
 \chi(\mathcal O_Y/P^{(2)})
 =5-\sum_p d_p,
 \qquad d_p=\delta(C,p)+\tau_p-r_p.
\]
We prove the precise gap needed here: at an ordinary
node $d_p=5$; at every planar singularity
$d_p=5\delta(C,p)$; and at every nonplanar singularity
having at most two branches $d_p\ge7$.

\noindent\textbf{Step 4.} First consider a planar curve, with ideal $(z,h(x,y))$.
Its conormal is free over $\mathcal O_C$, generated by
$dz,dh$. On each normalized branch, write
$dh=a_i u_i$ for a primitive plane kernel generator,
and put $v_i=\operatorname{ord}a_i$.
The inclusion of each of these two free summands into
its normalized module has length $\delta(C,p)$.
The second summand has the additional vanishing orders
$v_i$. Hence
\[
 \tau_p=2\delta(C,p)+\sum_i v_i.
\]
We verify $\sum_i(v_i-r_i)=2\delta(C,p)$ directly by
curve duality. Let $A$ be the reduced plane-curve ring,
$B$ its normalization, and $I=\operatorname{Hom}_A(B,A)$
its conductor. Applying $\operatorname{Hom}_A(-,A)$ to
$0\to A\to B\to B/A\to0$ gives
$A/I\simeq\operatorname{Ext}^1_A(B/A,A)$.
Here $\operatorname{Ext}^1_A(B,A)=0$ because $B$ is
maximal Cohen--Macaulay over the one-dimensional
Gorenstein ring $A$. The dual of a finite-length module
has the same length, so $\operatorname{length}(A/I)=\delta$,
and $\operatorname{length}(B/I)=2\delta$.
If the conductor exponent on branch $i$ is $c_i$,
this says $\sum_i c_i=2\delta$.
Choose a generic plane coordinate $x$ with minimal
valuation on every normalized branch. The dualizing
generator $dx/h_y$ then pulls back with
order $r_i-v_i$. The identity
$\omega_B=I\omega_A$ therefore gives
$c_i=v_i-r_i$. Combining these equalities proves
$\tau_p-r_p=4\delta$ and the stated planar formula.

\noindent\textbf{Step 5.} Now consider a nonplanar unibranch curve of multiplicity
$m$. Its embedding dimension is three. As in P4.30,
all gradient entries have valuation at least $m$, so
$\tau_p\ge2m$, while $r_p=m-1$ and
$\delta\ge m-1$. A nonplanar unibranch curve has
$m\ge3$. To see the last assertion, if $x=t^2$, write
each other coordinate as $a(x)+t b(x)$. After subtracting
the even parts, choose the nonzero $b$ of least order.
The other odd coordinate is an analytic multiple of it
over $\mathbb C\{x\}$. This is a smooth ambient equation
containing the curve, so its embedding dimension is at
most two. The case where both odd parts vanish contradicts
the birational normalization parametrization.

\noindent\textbf{Step 6.} The estimates give $d_p\ge2m\ge6$. Equality at six
would require $m=3$, $\delta=2$, and $\tau_p=6$.
The valuation semigroup of a unibranch curve has exactly
$\delta$ gaps, as follows by filtering its local ring
and its normalization by order. Thus the two gaps here
are one and two, and we can choose generators with
initial orders three, four, and five. They form an
ambient coordinate system: their classes generate the
three-dimensional cotangent space, since the square
of the maximal ideal has order at least six.
In a defining relation, the quadratic monomials
$x^2$ and $xy$ have respectively the unique least
orders six and seven, so their coefficients vanish.
Every gradient with valuation three consequently has
its initial vector in the single $z$ direction
(from a possible $xz$ term). Two such initial
gradient vectors are dependent. Every two-by-two
minor therefore has valuation at least seven,
contradicting $\tau_p=6$. Hence $d_p\ge7$ in this case.

\noindent\textbf{Step 7.} Finally consider a nonplanar point with two branches
of multiplicities $m_1,m_2$. Two smooth curve branches
in a smooth threefold are contained in a smooth
surface. Indeed choose a common immersive coordinate
$x$ and write them as graphs
$(y,z)=(a_i(x),b_i(x))$. Divide their difference pair
by its common power of $x$; the equation
\[
 \frac{b_1-b_2}{x^k}(y-a_1)
 -\frac{a_1-a_2}{x^k}(z-b_1)=0
\]
has a nonzero linear derivative and contains both.
Thus in the nonplanar case $m_1+m_2\ge3$.
The absence of a linear defining equation gives,
on the two branches separately,
$\tau_p\ge2(m_1+m_2)$.
The reduced-curve normalization sequence gives
$\delta(C,p)=\delta_1+\delta_2+
\operatorname{length}\mathcal O_Y/(P_1+P_2)
\ge m_1+m_2-1$.
Also $r_p=m_1+m_2-2$.
It follows that
$d_p\ge2(m_1+m_2)+1\ge7$.
This proves the gap in all cases.

\noindent\textbf{Step 8.} At the unique two-branch point, a planar singularity
has defect five if its delta invariant is one, and
at least ten otherwise. A nonplanar point has defect
at least seven. Every other singular point is
unibranch and contributes at least five if planar,
and at least seven if nonplanar. Smooth points
contribute zero. Thus the sum of the defects can
be five or at least seven, but never six.
The required value $\chi(\mathcal O_{C_{\rm sch}})=-1$
would give exactly six. This contradiction excludes
$e(S)=1$.

\noindent\textbf{Step 9.} The final statement uses P4.26: the sum of all the
nonnegative fibre Euler numbers is at most two.
The rational-resolution case is strictly positive,
and we have just excluded one, leaving exactly two.
\end{proof}


\subsection{The conductor graph excludes Euler number one in the rational case}

\begin{proposition}[P4.35]\label{night:P4.35}\label{stmt:P4.35}
\StatementRecord{night:P4.35}{proved}{Historical events: GPT-6 Astra (OpenAI), ChatGPT Pro mode, effort Pro, S6\_06\_BLIND, 16--18 November 2026: targeted proof check. Also checked in this writing. Pre-continuation writing events: GPT-6 Astra (OpenAI), ChatGPT Pro mode, effort Pro, S6\_10\_PAPERS, 18 November 2026: targeted proof check. Continuation metadata events: GPT-6 Astra (OpenAI), ChatGPT Pro mode, effort Pro, S6\_10\_PAPERS continuation, 19 November 2026: used and rechecked in this run; status-vocabulary reconciliation; metadata reconciliation only; no new mathematical check. \emph{Pre-continuation recorded proof scope.} The night-run re-derivation was by GPT-6 Astra (OpenAI, ChatGPT Pro mode, effort Pro), 16--18 November 2026. Also checked in this writing by GPT-6 Astra (OpenAI), ChatGPT Pro mode, effort Pro, 18 November 2026. The pairwise-intersection two-label argument forces a conductor graph cycle. Rational exceptional fibres were checked by the ideal sequence and coherent relative dimension one, and contraction of their simply connected trees preserves the conductor fundamental group. The Euler-one case is excluded without a transverse-type assumption}{\ref{night:P4.26}; \ref{night:P4.33}; \ref{night:P4.8}; \ref{night:P4R.6}}

No irreducible support with rationally resolved normalization
has Euler number one. Consequently, in a fibration whose
supports are all irreducible or normal crossing, every
irreducible singular support with rationally resolved
normalization has Euler number exactly two.

More generally, an irreducible support cannot have reduced
image conductor equal to a rational curve with exactly one
point-identification cycle when its reduced normalization
conductor has no graph cycles. This excludes that conductor
case for every ruled-resolution double-horizontal pole tree.
\end{proposition}
\begin{proof}
\noindent\textbf{Step 1.} We first prove the graph assertion independently of
the transverse singularity type and of ramification.
Let $D\subset N$ be the reduced normalization conductor,
and suppose its point-identification graph is a tree.
Let $c$ be its number of normalized compact components.
Suppose its image $C$ is irreducible with rational
normalization and exactly one point-identification
cycle. There is then a unique point $p\in C$ with
two analytic branches; all other points are unibranch.
Denote the two preimages of $p$ on $C^\nu$ by
$p_+,p_-$. Every normalized component of $D$ maps
finitely and surjectively to $C^\nu$.

\noindent\textbf{Step 2.} The finite set $\nu^{-1}(p)$ has $n$ points and
parametrizes the local irreducible hypersurface
branches of $S$ at $p$. Every pair of such branches
intersects along a curve: two nonunit local
hypersurface equations in a regular local ring of
dimension three generate an ideal of height at most
two. Their common curve lies in $C$. Consequently
the nonempty subsets of $\{+,-\}$ recording which
branches of $C$ are contained in each hypersurface
branch pairwise intersect. Such a family of subsets
of a two-element set has a common element. If a
singleton occurs its element belongs to every set;
if no singleton occurs every set is the whole set.
Interchanging the labels, every point of
$\nu^{-1}(p)$ therefore has at least one conductor
branch mapping to the plus branch of $C$.

\noindent\textbf{Step 3.} Let $r_+$ and $r_-$ count the points on $D^\nu$
over $p_+$ and $p_-$ respectively. The preceding
paragraph gives $r_+\ge n$. Surjectivity of the
map from each of the $c$ normalized components
gives $r_-\ge c$. The contribution of the points
over $p$ to the point-identification graph of $D$
is exactly
\[
 \sum_{q\in\nu^{-1}(p)}(r_q-1)
       =r_++r_--n\ge c.
\]
For a connected reduced compact curve with $c$
normalized components, the cycle rank of this graph
is $1-c+\sum_q(r_q-1)$. All other contributions
are nonnegative. Its rank is therefore at least
one, contradicting the tree hypothesis.
This proves the graph assertion.

\noindent\textbf{Step 4.} Now let the normalization have rational resolution
and suppose $e(S)=1$. The normal-crossing alternative
has Euler number two by P4.8. In the other alternative
P4.26 gives a connected simply connected resolved
pole support $A\subset T$, and
\[
 \pi_1(S)=\pi_1(C)=F_b,\qquad b\le1,
 \qquad e(C)=c_C+1-b.
\]
Its normalized components are rational.
Since $e(S_{\rm reg})\ge0$, Euler number one forces
$c_C=b=1$ and $e(S_{\rm reg})=0$.
It remains to verify the tree property for the
actual conductor $D\subset N$, allowing singular $N$.

\noindent\textbf{Step 5.} P4.33 makes every singularity of $N$ rational.
Every connected exceptional fibre of $T\to N$
is simply connected. Here is the precise curve
calculation. Over a small Stein neighbourhood,
$R^1\rho_{N*}\mathcal O_T=0$ gives
$H^1(\mathcal O_T)=0$ on its inverse image.
This proper map has every fibre of dimension at
most one over that Stein neighbourhood. Its coherent
cohomological dimension is consequently at most
one, so the second higher direct image of every
coherent sheaf vanishes. Applied to the ideal of the reduced
exceptional curve $E$, Leray and the Stein property
give $H^2(\mathcal O_T(-E))=0$ there. Its ideal
sequence consequently gives $H^1(E,\mathcal O_E)=0$.
The curve normalization exact sequence now implies
that all normalized components are rational and
that their point-identification graph is a tree.
Thus $E$ is simply connected.

\noindent\textbf{Step 6.} By P4R.6, $A$ is the full inverse image of $D$.
The proper map $A\to D$ is the quotient which
collapses these connected exceptional curves to
points and is an isomorphism elsewhere.
Compatible triangulation makes these curve
inclusions cofibrations. Collapsing a connected
simply connected subcomplex preserves fundamental
groups by van Kampen. Collapsing the finitely many
exceptional fibres therefore gives $\pi_1(D)=0$.
Each normalized component of $D$ is rational,
being birational to a strict pole component in
$T$. Hence its point-identification graph is a
tree. The already proved graph assertion contradicts
the description of $C$. This excludes $e(S)=1$.

\noindent\textbf{Step 7.} For the final more general application, the
resolved pole support in the ruled double-section
case is $H$ with rational trees attached.
Every exceptional fibre is connected and is a
union of its components. If it does not contain
$H$, it is a connected rational subtree. If it
contains $H$, all the remaining components form
rational trees attached to that component once.
Contracting such a connected subcurve cannot
create a cycle in the graph: contraction of a
connected subgraph of a tree again has a tree
as quotient. The only positive-genus vertex,
if contracted, is replaced by a point, and this
also introduces no graph edge. Thus the image
conductor $D$ has no graph cycles, and the same
contradiction applies.

\noindent\textbf{Step 8.} Finally, for a fibration in the two requested
support classes, P4.26 makes the total sum of
nonnegative fibre Euler numbers at most two.
The rational-resolution case is strictly positive;
excluding one leaves exactly two.
\end{proof}


\subsection{Eliminating all positive-Euler Hopf resolutions}

\begin{proposition}[P4.36]\label{night:P4.36}\label{stmt:P4.36}
\StatementRecord{night:P4.36}{proved}{Historical events: GPT-6 Astra (OpenAI), ChatGPT Pro mode, effort Pro, S6\_06\_BLIND, 16--18 November 2026: targeted proof check. Also checked in this writing. Pre-continuation writing events: GPT-6 Astra (OpenAI), ChatGPT Pro mode, effort Pro, S6\_10\_PAPERS, 18 November 2026: targeted proof check; used and rechecked in this run. Continuation metadata events: GPT-6 Astra (OpenAI), ChatGPT Pro mode, effort Pro, S6\_10\_PAPERS continuation, 19 November 2026: used and rechecked in this run; status-vocabulary reconciliation; metadata reconciliation only; no new mathematical check. \emph{Pre-continuation recorded proof scope.} The night-run re-derivation was by GPT-6 Astra (OpenAI, ChatGPT Pro mode, effort Pro), 16--18 November 2026. Also checked in this writing by GPT-6 Astra (OpenAI), ChatGPT Pro mode, effort Pro, 18 November 2026. The Hopf finite-cover reduction, saturated derivative line and rank-two image, Chern identity, integral quadratic blowup increment, and equality propagation through every exceptional curve were re-derived. Minimality forces the actual normalization to be smooth and minimal; the Hopf foliation normal-form input remains a separately cited dependency}{\ref{night:P4.24}; \ref{night:P4R.6}}

Every irreducible singular support in Part~4.2 whose
normalization has a Hopf minimal model has smooth minimal
Hopf normalization and Euler number zero. The assertion
does not assume initially that the normalization is smooth
or that its Euler number vanishes.
\end{proposition}
\begin{proof}
\noindent\textbf{Step 1.} Let $T\to N\to S\subset Y$ be the minimal normalization
resolution, and let $\pi:T\to T_0$ contract its $k$ blowups
to a minimal Hopf surface. We first take a finite unramified
primary Hopf cover $T'_0\to T_0$ and pull back the actual
sequence of point blowups. This gives a finite unramified
map $T'\to T$ and a blowup morphism $T'\to T'_0$.
Compose with $T\to Y$. All Chern calculations below take
place on this explicit smooth cover; no cover of the
singular normalization $N$ is asserted. To avoid a second
set of symbols we write $T,T_0,k,\rho$ for these covered
objects until the last paragraph.

\noindent\textbf{Step 2.} The primary Hopf normal form and proof of P4.24 give a
nowhere-zero holomorphic foliation on $T_0$ tangent to
every compact curve. Its saturated pullback to $T$ is a
line subbundle away from finitely many simple saddle
zeros. It is tangent to every strict conductor component
and every created exceptional curve. The actual conductor
centres lie on these invariant curves. Every blowup centre
is either a regular point of this foliation or one of its
saddles. Let $\varepsilon=0$ in the first case and
$\varepsilon=1$ in the second. The same explicit two-chart
calculation as in P4.24 gives
\[
 c_1(L)=-K_T+V,\qquad
 V'=\mu^*V+\varepsilon E,
 \qquad K_T^2=-k,\qquad K_TV-V^2=0.
\]
Topologically trivial line-bundle factors from $T_0$ do
not change these integral intersection calculations.

\noindent\textbf{Step 3.} Put $B=\rho^*TY$. Its first and second rational Chern
classes vanish: $b_2(Y)=0$ gives $c_1(TY)=0$ rationally,
and Poincar\'e duality gives $H^4(Y;\mathbb Q)=0$ and
hence $c_2(TY)=0$. The derivative $L\to B$ can now have
a divisorial zero, since $\rho$ contracts exceptional
curves. Denote that effective integral divisor by $Z$.
Dividing by it gives the saturated line injection in
codimension one
\[
 L'=L(Z)\longrightarrow B,
 \qquad c_1(L')=-K_T+V+Z.
\]
The quotient $Q=B/L'$ is torsionfree of rank two.
Write $z_0\ge0$ for the length of $Q^{**}/Q$.
The double dual is locally free on the smooth surface,
and the Chern sequence gives
\[
 c_2(Q^{**})=(-K_T+V+Z)^2-z_0.
\]

\noindent\textbf{Step 4.} Let $E_B\subset B$ be the saturated rank-two image
of $d\rho:TT\to B$. Its determinant is
$\det E_B=\omega_T^{-1}(R)$, where the effective
integral divisor $R$ is the divisorial zero of the
determinant map. Since $E_B$ is saturated in $B$,
the divided foliation line $L'$ is contained in it.
The image of $E_B/L'$ in $Q^{**}$ is saturated in
codimension one. Its rank-one double dual
$J=(E_B/L')^{**}$ is a line bundle. The map into
$Q^{**}$ extends across the finitely many omitted
points by Hartogs, since both sheaves are locally
free there. Its saturation is therefore a line bundle
of first Chern class $R-V-Z$; the rank-one quotient
is an ideal sheaf of a zero-dimensional scheme
times a line bundle of first Chern class $K_T-R$.
Let its length be $z_1\ge0$. Thus
\[
 (-K_T+V+Z)^2-z_0
       =(R-V-Z)(K_T-R)+z_1.
\]
Using $K_T^2=-k$ and $K_TV-V^2=0$, this is precisely
\[
 \mathcal B(R,Z)
 :=R^2-K_TR-VR-ZR-K_TZ+2VZ+Z^2
 =k+z_0+z_1.                                      \tag{P4.36.1}
\]

\noindent\textbf{Step 5.} We now evaluate the left side through the actual
point blowups, retaining all integral coefficients.
Push $R$ and $Z$ down along the chosen sequence.
At one reversed blowup write
\[
 R'=\mu^*R+cE,\quad Z'=\mu^*Z+dE,
 \quad K'=\mu^*K+E,\quad V'=\mu^*V+\varepsilon E.
\]
The integers $c,d$ are not assumed positive.
Since $E^2=-1$ and a pulled-back divisor has zero
intersection with $E$, direct expansion gives
\[
 \mathcal B(R',Z')-\mathcal B(R,Z)
 =-c^2+(1+\varepsilon)c+cd+(1-2\varepsilon)d-d^2.
\]
Set $a=c-1$ and $b=d-(1-\varepsilon)$. For each
$\varepsilon\in\{0,1\}$ this is exactly
\[
 1-(a^2+b^2-ab)\le1.                            \tag{P4.36.2}
\]
The quadratic form is
$\frac12(a^2+b^2+(a-b)^2)$; it is nonnegative
and is zero exactly when $a=b=0$.
All divisor intersections on the initial Hopf
surface are zero because $b_2(T_0)=0$.
Summing (P4.36.2) therefore gives
$\mathcal B(R,Z)\le k$.
Equation (P4.36.1) forces equality, $z_0=z_1=0$,
and at every single blowup
\[
 c=1,\qquad d=1-\varepsilon.                    \tag{P4.36.3}
\]

\noindent\textbf{Step 6.} Every created exceptional curve has positive
coefficient in $Z$. At a regular centre, (P4.36.3)
adds one to its nonnegative pulled-back coefficient.
At a saddle centre it adds zero, but at least one
previously created exceptional curve passes through
that saddle. The original Hopf foliation has no
zero; the two-chart blowup calculation places each
new saddle on the new exceptional curve. Induction
therefore gives a positive pulled-back coefficient
at every later saddle blowup as well.

\noindent\textbf{Step 7.} The foliation is tangent to each exceptional curve.
If its image under $\rho$ were a curve, the derivative
along its tangent line would be nonzero at a generic
point: a nonconstant holomorphic map of smooth
normalized curves in characteristic zero is
generically immersive. Positive divisorial order
of $d\rho|_L$ contradicts that fact. Thus every
created exceptional curve is contracted by $\rho$,
and hence by the normalization resolution, since
the normalization map is finite.

\noindent\textbf{Step 8.} Return to the original, possibly secondary Hopf
surface. Every exceptional curve downstairs has
an exceptional lift in the finite unramified cover;
the preceding argument makes its image a point.
Thus all exceptional curves of the original
$\pi:T\to T_0$ would be contracted by the minimal
resolution $T\to N$. If $k>0$, the exceptional curve
of its last blowup is a smooth rational curve of
self-intersection $-1$, which a minimal resolution
cannot contract. Hence $k=0$.
There is no further exceptional contraction from
a minimal Hopf surface: every compact curve has
self-intersection zero by $b_2=0$, whereas a
nonempty exceptional configuration has negative
definite intersection matrix, as proved in P4R.6.
Therefore $T=N=T_0$ is a smooth minimal Hopf
surface. P4.24 gives $e(S)=0$ and its complete
elliptic-conductor description.
\end{proof}


\subsection{Excluding the remaining class-VII minimal models}

\begin{proposition}[P4.37]\label{night:P4.37}\label{stmt:P4.37}
\StatementRecord{night:P4.37}{proved}{Historical events: GPT-6 Astra (OpenAI), ChatGPT Pro mode, effort Pro, S6\_06\_BLIND, 16--18 November 2026: targeted proof check. Also checked in this writing. Pre-continuation writing events: GPT-6 Astra (OpenAI), ChatGPT Pro mode, effort Pro, S6\_10\_PAPERS, 18 November 2026: targeted proof check; used and rechecked in this run. Continuation metadata events: GPT-6 Astra (OpenAI), ChatGPT Pro mode, effort Pro, S6\_10\_PAPERS continuation, 19 November 2026: used and rechecked in this run; status-vocabulary reconciliation; metadata reconciliation only; no new mathematical check. \emph{Pre-continuation recorded proof scope.} The night-run re-derivation was by GPT-6 Astra (OpenAI, ChatGPT Pro mode, effort Pro), 16--18 November 2026. Also checked in this writing by GPT-6 Astra (OpenAI), ChatGPT Pro mode, effort Pro, 18 November 2026. The class-VII Chern identity, negative-definite initial quadratic form, and at-most-one contribution of every blowup were re-derived. The resulting strict inequality contradicts the ambient derivative identity. The exact Favre and Dloussky-Oeljeklaus source loci for the initial foliation and curve basis are tracked separately}{\ref{night:P4.23}; \ref{night:P4.26}; \ref{night:P4.36}}

No irreducible singular support in Part~4.2 has normalization
whose minimal resolution has a class-VII minimal model with
positive second Betti number. Consequently every class-VII
normalization occurring here is a smooth minimal Hopf surface
and has Euler number zero.
\end{proposition}
\begin{proof}
\noindent\textbf{Step 1.} The numerically anticanonical classification used in P4.23
and P4.26 reduces this question to an intermediate Kato
minimal model $T_0$: the Inoue--Hirzebruch and parabolic
alternatives were excluded there by their reduced
normal-crossing conductor and the Chern identity.
Write $n=b_2(T_0)>0$, let $A_0$ be its maximal reduced
curve divisor, and let $D_0$ be the effective integral
pole divisor pushed down from the normalization resolution.
Thus $K_{T_0}+D_0$ is numerically zero. The support of
$D_0$ is $A_0$.

\noindent\textbf{Step 2.} We state the precise additional primary input.
Favre, \emph{Dynamique des applications rationnelles},
Theorem~1.2.24(1)--(2), printed p.~35, says that the
$n$ rational curve classes span $H^2(T_0;\mathbb Q)$
and that their intersection form is negative definite.
Theorem~1.2.30, printed p.~42, and its proof on
pp.~44 and~46 give a holomorphic foliation with conormal
line
\[
 N_{\mathcal F}^*=\mathcal O_{T_0}(-A_0)
                       \otimes\text{a flat line}.
\]
Theorem~1.2.32 and its proof on p.~46 show that every
component of $A_0$ is invariant and that its singularities
are precisely the intersection points, with local
first integrals $z^\alpha w^\beta$, $\alpha,\beta>0$.
The invariant-curve and order-one assertions are also
Dloussky--Oeljeklaus, Theorem~3.1, pp.~1528--1529.
The actual source URLs and exact theorem locators are
recorded in the bibliography and literature log.

\noindent\textbf{Step 3.} It follows from the determinant of tangent and normal
lines that the foliation tangent line has numerical
first Chern class $-K_{T_0}-A_0$. The local generator
at an intersection is
$\beta z\partial_z-\alpha w\partial_w$.
Blowing up a regular point or one of these saddles
again gives only regular points and saddles with
opposite-sign nonzero eigenvalues. Thus every actual
blowup centre has the same two possibilities
$\varepsilon=0,1$ as in P4.36, with the same formulas.

\noindent\textbf{Step 4.} Let $T\to T_0$ contain $k$ blowups. Write
$c_1(L)=-K_T+V$ for the saturated foliation line;
initially $V_0=-A_0$, and
$V'=\mu^*V+\varepsilon E$ at a blowup.
The reduced divisor $A_0$ is a rational cycle with
rational trees attached. Its normalization has
$n$ rational components and $n$ nodes, including
self-nodes, so $\chi(\mathcal O_{A_0})=0$.
Adjunction, or the divisor exact sequence on $T_0$,
therefore gives
\[
 (K_{T_0}+A_0)A_0=0.
\]
Hence $K_{T_0}V_0-V_0^2=0$.
Under each blowup this quantity changes by
$\varepsilon^2-\varepsilon=0$.
Since $K_{T_0}^2=-n$, we obtain
\[
 K_TV-V^2=0,\qquad K_T^2=-n-k.
\]

\noindent\textbf{Step 5.} Define $B=\rho^*TY$, the divisorial determinant zero
$R$, the divisorial foliation zero $Z$, and the
nonnegative lengths $z_0,z_1$ exactly as in P4.36.
Its saturated sheaf calculation, including the
Hartogs extension of the rank-one double dual,
now gives
\[
 \mathcal B(R,Z)
 =R^2-K_TR-VR-ZR-K_TZ+2VZ+Z^2
 =n+k+z_0+z_1.                                  \tag{P4.37.1}
\]
Every reversed blowup changes $\mathcal B$ by at
most one, by the unrestricted integer identity
(P4.36.2). It remains to evaluate the initial term.

\noindent\textbf{Step 6.} Push $R,Z$ down to effective divisors $R_0,Z_0$
on $T_0$. Their supports are contained in $A_0$:
the derivative is generically injective away from
the resolved conductor, and in any event all
compact curves of the intermediate surface are
components of its maximal divisor. Put $D=D_0$,
$A=A_0$. Using $K=-D$ numerically and $V=-A$ gives
\[
 \mathcal B_0
 =R_0^2+DR_0+AR_0-Z_0R_0+DZ_0-2AZ_0+Z_0^2.
\]
For divisor classes define
$q(X,Y)=X^2-XY+Y^2$.
Adjunction above gives $DA=A^2$, while $D^2=-n$.
Expansion consequently gives the exact completion
\[
 \mathcal B_0
       =n+q(R_0+D,Z_0+D-A).                     \tag{P4.37.2}
\]
The form on divisor classes is negative definite.
Moreover
\[
 q(X,Y)=\tfrac12\bigl(X^2+Y^2+(X-Y)^2\bigr),
\]
so it is nonpositive, and it vanishes only for
$X=Y=0$. The first argument $R_0+D$ is nonzero:
it is a nonzero effective combination of the
$n$ independent curve classes. Thus
$\mathcal B_0<n$.
After the $k$ blowups we still have
$\mathcal B(R,Z)<n+k$, contradicting (P4.37.1).
This excludes the intermediate minimal model.

\noindent\textbf{Step 7.} All class-VII cases with positive second Betti
number have now been excluded. The zero-Betti
Inoue case has no conductor curves and was
excluded in P4.26. P4.36 treats the remaining
Hopf case and proves the final assertion.
\end{proof}


\subsection{Two exact local exclusions for ordinary-double curves}

\begin{proposition}[P4.38]\label{night:P4.38}\label{stmt:P4.38}
\StatementRecord{night:P4.38}{proved}{Historical events: GPT-6 Astra (OpenAI), ChatGPT Pro mode, effort Pro, S6\_06\_BLIND, 16--18 November 2026: targeted proof check. Also checked in this writing. Pre-continuation writing events: GPT-6 Astra (OpenAI), ChatGPT Pro mode, effort Pro, S6\_10\_PAPERS, 18 November 2026 05:10:56.410177+00:00: computation or formal certificate; targeted proof check. Continuation metadata events: GPT-6 Astra (OpenAI), ChatGPT Pro mode, effort Pro, S6\_10\_PAPERS continuation, 19 November 2026: used and rechecked in this run; status-vocabulary reconciliation; metadata reconciliation only; no new mathematical check. \emph{Pre-continuation recorded proof scope.} The night-run re-derivation was by GPT-6 Astra (OpenAI, ChatGPT Pro mode, effort Pro), 16--18 November 2026. Also checked in this writing by GPT-6 Astra (OpenAI), ChatGPT Pro mode, effort Pro, 18 November 2026. proved. Full planar-singular and cusp/tacnode-plus-axis ordinary-double exclusion, including the actual boundary comparison. Re-derived P\textasciicircum{}(2)=P\textsuperscript{2} for the reduced complete-intersection plane curve and the exact Morse splitting f=Z\textsuperscript{2}+h\textsuperscript{2}u. A morsification and nonzero critical-value shift give Hessian block 2h\_s Hess(h\_s), hence a Morse sphere. For P=(xz,yz,y\textsuperscript{2}−x\textasciicircum{}k), k=3,4, directly intersected (z,h)\textsuperscript{2} with (x,y)\textsuperscript{2}, obtaining the five displayed generators. The axis quadratic determinant after h\ensuremath{\to}h−sx\textsuperscript{2} has divided factor F=(-sA+Bz)(A+Dz)−C\textsuperscript{2}z\textsuperscript{2}/4, with F\_z=AB and F\_s=−A\textsuperscript{2}. Its simple nonzero root gives the actual coordinate Y\textsuperscript{2}+x\textsuperscript{2}Z. Followed the compact sphere of square −4 in P4.16 and both nested proper boundary transports; the contradiction uses the null-homotopy of every sphere in the actual nearby T4. The unchanged critical\_curve\_node.py replay passed. Review provenance: New derivation after reading the supplied proof; not a blind review and not an independent human check}{\ref{night:P4.16}; \ref{night:P4.20}}

An actual irreducible support in Part~4.2 cannot have a singular
planar reduced critical curve at a point at which every generic
transverse germ is an ordinary double point. Nor can its reduced
critical curve there be the transverse union of a plane cusp
and a smooth branch, again with ordinary-double transverse germs
on both components. The latter assertion also holds with the
cusp replaced by a plane tacnode.
These exclusions allow singular normalization.
\end{proposition}
\begin{proof}
\noindent\textbf{Step 1.} We retain the boundary transport of P4.20 throughout. At a fixed
compact boundary, ordinary-double transverse germs form an open
condition. A sufficiently small deformation retaining that
condition gives the same nearby nonzero-level boundary
submersion. The finite-critical-value argument in that proof
then transports any compact vanishing sphere from the deformed
germ into the preceding nearby fibre. When a coordinate change
or a unit requires a smaller ball, we perform this construction
on that ball first, and then transport through the earlier
larger ball. Thus every sphere constructed below reaches the
original nearby four-torus.

\noindent\textbf{Step 2.} First suppose that the reduced critical curve has ideal
$P=(z,h(x,y))$, where $h$ is reduced and singular at the
origin. The equation $f$ lies in $P^{(2)}$: its generic
transverse ordinary double point vanishes to order two on
each component. The complete-intersection square $P^2$ is
unmixed, so $P^{(2)}=P^2$. We therefore have the exact expression
\[
 f=A z^2+Bzh+Ch^2.
\]
Adding a small constant multiple of $z^2$ preserves the critical
curve and the ordinary-double boundary condition. We choose it
so that the $z^2$ coefficient at the origin is nonzero. The
Morse lemma in $z$, with parameters $x,y$, then gives
\[
 f=Z^2+h(x,y)^2u(x,y).
\]
Indeed the critical $z$ coordinate vanishes on $h=0$ and
is divisible by $h$, so its substitution has the displayed
square factor. Add a small constant to $u$ and shrink the
ball so that $u$ is a unit. Absorbing its square root in $h$
gives $Z^2+H(x,y)^2$, with the same singular reduced zero
curve. A reduced singular plane curve has an isolated critical
point. A generic small morsification $H_s$ has a Morse critical
point, and a small additive constant makes all these critical
values nonzero. At any such point $Z^2+H_s^2$ has a
nondegenerate isolated critical point: its plane Hessian is
$2H_s\operatorname{Hess}H_s$, and its $Z$ Hessian is $2$.
Its vanishing sphere has self-intersection $-2$.
The moving curve $H_s=0$ remains transverse to the chosen
boundary, with nondegenerate normal quadratic form there.
The preceding boundary transport therefore contradicts
$\pi_2(T^4)=0$.

\noindent\textbf{Step 3.} For the second assertion write
\[
 h=y^2-x^k,\qquad P=(xz,yz,h),\qquad k=3\text{ or }4.
\]
This is the union of the plane curve $z=h=0$ and the
transverse line $x=y=0$. Its symbolic square is
\[
 P^{(2)}=(z,h)^2\cap(x,y)^2
  =(x^2z^2,xyz^2,y^2z^2,zh,h^2).
\]
For $k=4$, the first equality uses the two smooth branches of
$h=0$ separately; their complete-intersection square is unmixed
and hence is still $(z,h)^2$. To verify the intersection
formula directly, write an element of $(z,h)^2$ as
$az^2+bzh+ch^2$. The last two terms belong to $(x,y)^2$,
and $z$ is a nonzero divisor modulo this ideal. The first term
belongs to it exactly when $a\in(x,y)^2$.
Consequently
\[
 f=Azh+z^2(Bx^2+Cxy+Dy^2)+Gh^2.
\]
Small constant changes in $A,B$ retain the boundary condition
and make $A(0)B(0)\ne0$. Replace $h$ in this expression by
$h_s=y^2-x^k-sx^2$. The new critical curve contains
$z=h_s=0$ and $x=y=0$, and its boundary ordinary-double
germs persist. On the latter axis the transverse quadratic
matrix has coefficients
\[
 a=-sAz+Bz^2,\qquad b=Cz^2,\qquad c=Az+Dz^2,
\]
with off-diagonal entry $b/2$. Here the four coefficient
functions are restricted to that axis. Its determinant
divided by $z^2$ is
\[
 F(s,z)=(-sA+Bz)(A+Dz)-C^2z^2/4.
\]
At $(s,z)=(0,0)$ we have $F_z=AB\ne0$ and
$F_s=-A^2$. The implicit-function theorem gives a simple root
$z(s)=sA(0)/B(0)+O(s^2)$, nonzero for small nonzero $s$.
There the $y^2$ coefficient is nonzero, the Hessian has rank
one, and its determinant has a simple zero in the axis
direction. The parametrized Morse lemma gives
$Y^2+x^2H(x,z)$, with $H_z\ne0$. The exact coordinate
change $Z=H(x,z)$ yields $Y^2+x^2Z$.
P4.16 gives its compact spherical class of square $-4$.
Nested boundary transport once more places this class in
the original nearby four-torus, a contradiction.
\end{proof}

The exact quadratic matrix, determinant, and the two nonzero
derivatives in this proof are verified by
\texttt{code/part4/critical\_curve\_node.py}; its saved output
reports PASS.


\subsection{A coefficient-two exceptional curve over an ordinary-triple germ}

\begin{proposition}[P4.39]\label{night:P4.39}\label{stmt:P4.39}
\StatementRecord{night:P4.39}{proved}{Historical events: GPT-6 Astra (OpenAI), ChatGPT Pro mode, effort Pro, S6\_06\_BLIND, 16--18 November 2026: targeted proof check. Also checked in this writing. Pre-continuation writing events: GPT-6 Astra (OpenAI), ChatGPT Pro mode, effort Pro, S6\_10\_PAPERS, 18 November 2026 05:10:56.410177+00:00: computation or formal certificate; targeted proof check. Continuation metadata events: GPT-6 Astra (OpenAI), ChatGPT Pro mode, effort Pro, S6\_10\_PAPERS continuation, 19 November 2026: used and rechecked in this run; status-vocabulary reconciliation; metadata reconciliation only; no new mathematical check. \emph{Pre-continuation recorded proof scope.} The night-run re-derivation was by GPT-6 Astra (OpenAI, ChatGPT Pro mode, effort Pro), 16--18 November 2026. Also checked in this writing by GPT-6 Astra (OpenAI), ChatGPT Pro mode, effort Pro, 18 November 2026. proved. Full coefficient-two exceptional-curve exclusion for uniform ordinary-triple transverse germs; arithmetic-genus wording repair recorded. The ambient blowup residue comparison gives D1+(d−2)e and 2=(d−2)n+a. Re-derived positivity using the negative exceptional intersection matrix and arithmetic adjunction. The two received expressions 2g(H)−2 must be 2p\_a(H)−2 until smoothness is deduced: positive geometric genus gives p\_a\ensuremath{\geq}1, and coefficient one forces p\_a=1, hence g=1 and zero delta, so H is smooth elliptic. For d=4 the quartic Hessian vanishes along an actual critical arc, giving multiplicity at least3 and excluding every positive-genus reduced tangent component. For d=3, depressed Weierstrass preparation and P\textasciicircum{}(3)=P\textsuperscript{3} give u\textsuperscript{3}+h\textsuperscript{2}Bu+h\textsuperscript{3}C. Re-derived both the double-line quadratic normalization, whose discriminant has degree\ensuremath{\leq}2, and the isolated-point chart v\textasciicircum{}e+xG. At a simple zero the actual coordinate W=G gives v\textasciicircum{}e+xW; its cyclic-quotient resolution is rational. The three pole cases and all first-blowup degree bounds passed the unchanged exact programs. This proves the statement with the explicitly supplied adjunction correction. Review provenance: New derivation after reading the supplied proof; not a blind review and not an independent human check. The receiving adjunction calculation now uses arithmetic genus until its equality proves that the exceptional curve is smooth elliptic}{\ref{night:P4R.6}}

Let $S=(f=0)$ be a reduced hypersurface germ in a smooth
threefold, with one-dimensional reduced critical locus $C$.
Suppose that every generic transverse germ along $C$ is an
ordinary triple point. On the minimal resolution $T$ of its
normalization, let $D$ be the integral total pole divisor
defined by the residue of an ambient volume form divided by
$f$. No exceptional curve of positive genus has coefficient
two in $D$.

In particular, the contracted horizontal section of a
positive-genus ruled resolution in Part~4.2 cannot lie over
a point where all local critical-curve branches have generic
type $D_4$. This assertion does not replace that local
hypothesis by a mixture of transverse types.
\end{proposition}
\begin{proof}
\noindent\textbf{Step 1.} We blow up the ambient origin, write $S_1$ for the strict
transform and $N_1$ for its normalization, and take the
minimal resolution $T_1$ of $N_1$. The proper map $N_1\to S$
lifts to the normalization $N$ by its universal property.
Thus $T_1\to N$ is a resolution and dominates $T$ by a
sequence of point blowups. The strict transform of a proposed
positive-genus exceptional curve $H$ remains present and has
the same pole coefficient two.

\noindent\textbf{Step 2.} Put $d=\operatorname{mult}_0f$, and denote the ambient
exceptional divisor by $E$. Actual adjunction gives
\[
 \omega_{S_1}=\pi^*\omega_S\otimes
        \mathcal O((2-d)E)|_{S_1}.
\]
Let $D_1$ be the total conductor divisor on $T_1$ relative
to $\omega_{S_1}$, and let $e$ be the Cartier pullback of
$E$. With the specified residue comparison, the total pole
of the original germ on $T_1$ is
\[
                    D_1+(d-2)e.                       \tag{4.39a}
\]
The divisor $D_1$ is effective. Its strict conductor
coefficients are positive integers. For exceptional curves
of the minimal resolution of $N_1$, the negative-matrix
argument of P4R.6 applies locally: $K_{T_1}F\ge0$, since
there is no smooth rational exceptional $(-1)$-curve;
adding the strict conductor retains nonnegative intersection.
The inverse exceptional intersection matrix has nonpositive
entries, so all exceptional coefficients are nonnegative.

\noindent\textbf{Step 3.} If $H$ is exceptional over $N_1$, its coefficient in $D_1$
is strictly positive. A zero coefficient would give
\[
 0=(K_{T_1}+D_1)H
   =2p_a(H)-2-H^2+D_1H>0.
\]
If that coefficient is one, adjunction gives
\[
 0=2p_a(H)-2+(D_1-H)H.
\]
The normalization of $H$ has positive genus, so $p_a(H)\ge g(H^\nu)\ge1$. Both terms in the last equality are nonnegative, hence $p_a(H)=1$ and $(D_1-H)H=0$. The normalization formula $p_a(H)=g(H^\nu)+\sum\delta(H,p)$ now gives $g(H^\nu)=1$ and zero delta invariant at every point. Thus $H$ is smooth elliptic and is disjoint from every other positive component of $D_1$. A zero-coefficient exceptional
curve cannot meet a positive one, by its own intersection
equation and $K_{T_1}F\ge0$. The connected exceptional
fibre is consequently $H$ alone and meets no strict conductor.
The entire germ of $S_1$ at its image is normal with an
isolated simple elliptic singularity. Indeed, any additional
hypersurface branch would intersect this branch in a curve
and would give strict conductor meeting $H$.

\noindent\textbf{Step 4.} Set $n=\operatorname{ord}_H(e)>0$ and
$a=\operatorname{coeff}_H(D_1)\ge0$. Formula (4.39a)
gives the exact identity
\[
                       2=(d-2)n+a.                    \tag{4.39b}
\]
Thus $d\ge5$ is impossible. For $d=4$ we have
$n=1,a=0$, and $H$ is not exceptional over $N_1$.
At its generic point the normalization is an isomorphism,
so $H$ is birational to a reduced component of the projective
quartic tangent cone. Every second derivative of $f$ vanishes
on $C$. Its leading term along any analytic arc of $C$
therefore makes that quartic have a point of multiplicity
at least three. An irreducible quartic has arithmetic genus
three and loses at least three there. A smooth elliptic cubic
component together with the remaining line could have
multiplicity at most two at the point. All components of
degree at most two are rational. Thus every reduced tangent
component is rational, excluding $H$ for $d=4$.

\noindent\textbf{Step 5.} We are left with $d=3$ and
\[
                   (n,a)=(2,0)\quad\hbox{or}\quad(1,1).
\]
Choose a coordinate $u$ with nonzero pure cubic initial
coefficient, apply Weierstrass preparation, and depress the
monic cubic. Multiplication by the preparation unit does not
affect Hessian vanishing on $C$, because $f=df=0$ there.
We obtain
\[
 f=\text{unit}\cdot\bigl(u^3+b(x,z)u+c(x,z)\bigr).
\]
The second $u$ derivative is $6u$, so $C$ lies in the smooth
surface $u=0$. Write its reduced ideal as $(u,h(x,z))$.
The complete-intersection cube is unmixed and equals its
symbolic cube. Generic order three along $C$ and coefficient
comparison therefore give the exact expression
\[
 f=\text{unit}\cdot
       \bigl(u^3+h^2B(x,z)u+h^3C(x,z)\bigr).           \tag{4.39c}
\]
If $C$ is singular, $\operatorname{ord}h\ge2$ and the
tangent cone is the triple plane $u^3$. The degree-four
term restricts to zero on this plane. After the ambient
blowup the entire reduced exceptional line is critical,
so none of its points is an isolated normal singularity.
If $C$ is smooth, take $x=h$. Every homogeneous term $f_j$
then has degree at least three in $u,x$. The cubic tangent
cone is a binary cubic in these two variables, hence three
lines, counted with multiplicities, through the tangent
direction of $C$.

\noindent\textbf{Step 6.} First suppose that $H$ dominates a tangent-cone line.
For $(n,a)=(2,0)$ the normalization is generically an
isomorphism, and $H$ is birational to that line. For
$(n,a)=(1,1)$, a transverse normalization branch has conductor
exponent one. It is smooth and meets exactly one other smooth
branch transversely: the intrinsic conductor exponent of a
singular plane branch is at least two, and each other branch
contributes a positive intersection multiplicity. Thus the
generic transverse germ of $S_1$ is a node. A multiplicity-one
tangent line is generically smooth and cannot have $a=1$.
For a multiplicity-three line, the two geometric branches of
that node have exceptional orders whose weighted sum is three.
A connected double cover would give equal orders and an even
sum. Thus the order-one branch has degree one over the line
and is rational.

\noindent\textbf{Step 7.} The remaining possibility is a double line together with a
distinct simple line. Here $C$ is smooth by (4.39c). Make a
linear change so that $f_3=xu^2$. On the $x$ chart of the
blowup put $u=xv,z=xw$. The strict equation is
\[
 v^2+x f_4(1,v,w)+x^2f_5(1,v,w)+\cdots.
\]
An ordinary double curve on $x=v=0$ requires
$f_4(1,0,w)=0$ identically. Its transverse quadratic form is
\[
 v^2+xvB(w)+x^2C(w),
 \qquad \deg B\le1,\quad\deg C\le2.                  \tag{4.39d}
\]
The degree bounds follow from the degree-at-least-three
condition in $u,x$ on every $f_j$: here
$B=\partial_v f_4(1,0,w)$ and $C=f_5(1,0,w)$.
The normalization over the exceptional line adjoins a root
of $q^2+B(w)q+C(w)$. Its discriminant has degree at most
two. The resulting quadratic extension of $\mathbb C(w)$
is branched at at most two points, including infinity.
Riemann--Hurwitz makes its compact normalization rational;
the components of a split cover are rational as well.

\noindent\textbf{Step 8.} Finally suppose that $H$ is exceptional over a point of $N_1$.
We must have $n=a=1$, and the first part of the proof would
make $S_1$ normal with an isolated simple elliptic singularity
there. This cannot happen. For singular $C$ the whole
exceptional line is critical, as already proved. For smooth
$C$ its strict transform passes through the common point of
the tangent lines, so that point is not an isolated normal
singularity. Away from it, at a line of multiplicity
$e=2$ or $3$, divide by the nonzero factor from any other
line. The strict equation has the form
\[
                         v^e+xG(v,w,x),
\]
where $G(0,w,0)$ is affine linear in $w$, by the same
normal-degree bound. If this function is nonzero at a point,
the surface is smooth. If it is identically zero, the entire
line is critical. At a simple zero, $G_w$ is a unit and the
exact coordinate change $W=G(v,w,x)$ gives $v^e+xW$, a
rational double point of type $A_{e-1}$. A multiplicity-one
line is smooth away from intersections. There is therefore
no isolated simple elliptic singularity. All possibilities
for $H$ have been excluded.
\end{proof}

The positive-integer solutions of (4.39b), the monomial degree
bounds in (4.39d), and the forced fourth-order factor of the
homogeneous quadratic discriminant are checked in
\texttt{code/part4/d4\_pole.py}. Its saved output reports PASS.


\subsection{The excess-one curve calculation and the ordinary-double sector}

\begin{proposition}[P4.40]\label{night:P4.40}\label{stmt:P4.40}
\StatementRecord{night:P4.40}{proved}{Historical events: GPT-6 Astra (OpenAI), ChatGPT Pro mode, effort Pro, S6\_06\_BLIND, 16--18 November 2026: targeted proof check. Also checked in this writing. Pre-continuation writing events: GPT-6 Astra (OpenAI), ChatGPT Pro mode, effort Pro, S6\_10\_PAPERS, 18 November 2026 05:10:56.410177+00:00: computation or formal certificate; targeted proof check. Continuation metadata events: GPT-6 Astra (OpenAI), ChatGPT Pro mode, effort Pro, S6\_10\_PAPERS continuation, 19 November 2026: used and rechecked in this run; status-vocabulary reconciliation; metadata reconciliation only; no new mathematical check. \emph{Pre-continuation recorded proof scope.} The night-run re-derivation was by GPT-6 Astra (OpenAI, ChatGPT Pro mode, effort Pro), 16--18 November 2026. Also checked in this writing by GPT-6 Astra (OpenAI), ChatGPT Pro mode, effort Pro, 18 November 2026. proved. Full excess-one analytic curve-ring classification and its global ordinary-double application. For A\ensuremath{\subset}R=Cdiag+\ensuremath{\oplus}t\_i C\{t\_i\} with length(R/A)=1, the quotient is killed by m\_A. Writing a missing element v with v\textsuperscript{2}=a+cv, c\ensuremath{\neq}0 would let the implicit-function theorem solve a root in A and force v into A; therefore m\_R\textsuperscript{2}\ensuremath{\subset}A. Classified the exact first-jet hyperplanes: one nonzero coefficient gives a cusp plus independent lines; two give the actual tacnode generators (t,t),(t\textsuperscript{2},−t\textsuperscript{2}); three give a plane ordinary triple because the second-jet product determinant is nonzero. Minimal generator lifts give the asserted ambient coordinate isomorphisms. In the global pole-tree case H must contract because its coefficient2 cannot dominate a generic conductor of exponent1. The current conductor duality gives \ensuremath{\chi}OC=−1; the two rational normalized components and one graph cycle give total delta3, identification length2 and excess1. Every resulting embedded singular germ is one of the actual exclusions in P4.38. The unchanged curve\_excess.py replay passed. Review provenance: New derivation after reading the supplied proof; not a blind review and not an independent human check}{\ref{night:P4.33}; \ref{night:P4.35}; \ref{night:P4.38}; \ref{night:P4R.4}}

In the scope of Part~4.2 where every fibre support is
irreducible or normal crossing, no irreducible support with
positive-genus ruled resolution has generic transverse type
$A_1$ on every component of its critical curve.
\end{proposition}
\begin{proof}
\noindent\textbf{Step 1.} We first give the exact local curve calculation needed in
the proof. Let $A$ be a reduced complex curve local ring with
$r$ branches, let $B=\bigoplus_{i=1}^r\mathbb C\{t_i\}$
be its normalization, and let $R\subset B$ be the ring whose
constant terms agree. Then
\[
 \operatorname{length}(R/A)=\delta(A)-(r-1).
\]
Suppose this length is one. The simple $A$-module $R/A$
is annihilated by $\mathfrak m_A$, so
$\mathfrak m_A R\subset A$. Choose
$v\in\mathfrak m_R$ spanning $R/A$ and write
$v^2=a+cv$, where $a\in\mathfrak m_A$ and $c\in\mathbb C$.
If $c\ne0$, the convergent square-root series solves
$b^2-cb-a=0$ with $b\in\mathfrak m_A$. Then
$(v-b)(v+b-c)=0$ and the second factor is a unit, forcing
$v=b$, a contradiction. Thus $v^2\in A$ and
\[
                         \mathfrak m_R^2\subset A.
\]
Consequently $A$ is exactly the inverse image of a hyperplane
\[
                         \sum_i a_i f_i'(0)=0
\]
in $\mathfrak m_R/\mathfrak m_R^2=\mathbb C^r$.
Rescaling the parameters with nonzero $a_i$ gives actual
analytic ring isomorphisms with the following models.

\noindent\textbf{Step 2.} With one nonzero coefficient, that branch has ring
$\mathbb C\{t^2,t^3\}$ and the others are independent smooth
branches joined at their constant terms. The embedding
dimension is $r+1$: for $r=1$ this is a plane cusp, for
$r=2$ a cusp and transverse line, and for $r=3$ it does
not embed in a smooth threefold. With two nonzero
coefficients, the pair has equal first derivatives. Its
ring is generated by $(t,t)$ and $(t^2,-t^2)$, giving the
plane tacnode $y^2=x^4$. The other $r-2$ branches are
independent transverse lines. For $r=3$ and three nonzero
coefficients, the first jets are the restrictions of linear
forms to three distinct lines in a plane. Their products
span all three second-jet coordinates, so the ring is
exactly the ordinary plane triple point. Equivalently,
the latter ring has precisely the same constant and
first-derivative conditions and colength three in $B$.
These cases exhaust $r\le3$. The analytic ring
isomorphisms lift to ambient coordinate changes: choose
minimal algebra generators, lift their independent linear
parts, and apply the inverse-function theorem. In embedding
dimension two we first straighten the smooth containing
surface. Every realizable model is therefore one of the
local classes excluded in P4.38 for ordinary-double
transverse type.

\noindent\textbf{Step 3.} Now suppose the stated ruled support exists. Its horizontal
pole section has coefficient two. It cannot map
nonconstantly to the conductor, since an ordinary double
point has conductor exponent one on each normalization
branch. The section is contracted. P4.33 gives
$\chi(\mathcal O_{C_{\rm sch}})=-1$.
The conductor image is Cohen--Macaulay and generically
reduced in type $A_1$. It is therefore reduced, so
$C_{\rm sch}=C$ and $\chi(\mathcal O_C)=-1$.
All normalized components of $C$ are rational.

\noindent\textbf{Step 4.} The already proved group and Euler calculations leave
exactly two normalized rational components and one
point-identification cycle. Indeed the rational-tree case
has Euler number at least three by P4R.4; the
one-component cycle is excluded by P4.35; and the global
nonnegative Euler sum is at most two. The normalization
sequence and the graph formula consequently give
\[
 \sum_p\delta_p=3,\qquad
 \sum_p(r_p-1)=2,
 \qquad
 \sum_p\bigl(\delta_p-(r_p-1)\bigr)=1.
\]
Every summand in the last sum is nonnegative. There is
exactly one excess-one point and every point has at most
three local branches. The local calculation classifies that
point by an actual analytic isomorphism, and P4.38 excludes
each possible model. This is the required contradiction.
\end{proof}

The first-jet cases, the second-jet product determinant, and
the global excess computation are verified in
\texttt{code/part4/curve\_excess.py}; its saved output reports
PASS. The proof retains the ordinary-double boundary
condition and makes no substitution of transverse type
$A_2$ for $A_1$.


\subsection{The ordinary-node defect in transverse type \texorpdfstring{$A_3$}{A3}}

\begin{proposition}[P4.41]\label{night:P4.41}\label{stmt:P4.41}
\StatementRecord{night:P4.41}{proved}{Historical events: GPT-6 Astra (OpenAI), ChatGPT Pro mode, effort Pro, S6\_06\_BLIND, 16--18 November 2026: targeted proof check. Also checked in this writing. Pre-continuation writing events: GPT-6 Astra (OpenAI), ChatGPT Pro mode, effort Pro, S6\_10\_PAPERS, 18 November 2026: used and rechecked in this run. GPT-6 Astra (OpenAI), ChatGPT Pro mode, effort Pro, S6\_10\_PAPERS, 18 November 2026 05:10:56.410177+00:00: computation or formal certificate; targeted proof check; used and rechecked in this run. Continuation metadata events: GPT-6 Astra (OpenAI), ChatGPT Pro mode, effort Pro, S6\_10\_PAPERS continuation, 19 November 2026: used and rechecked in this run; status-vocabulary reconciliation; metadata reconciliation only; no new mathematical check. \emph{Pre-continuation recorded proof scope.} The night-run re-derivation was by GPT-6 Astra (OpenAI, ChatGPT Pro mode, effort Pro), 16--18 November 2026. Also checked in this writing by GPT-6 Astra (OpenAI), ChatGPT Pro mode, effort Pro, 18 November 2026. proved. Full ordinary-node conductor-module calculation, exact ring-condition matrices, defect bound and equality normalization. Classified the normalized rank-one module at A=C\{x,y\}/(xy) as the diagonal line or the full two-coordinate space, giving J\ensuremath{\cong}A,q=1 or J\ensuremath{\cong}B,q=0 after the recorded branch factors. Opened and re-derived de Jong–van Straten Theorem1.12 with its codimension-two Cohen–Macaulay and Gorenstein hypotheses. In the free case I=(g,k\textsuperscript{2}), its adjugate forces f\ensuremath{\in}I\textsuperscript{2}. A unit g\textsuperscript{2} coefficient yields a boundary-preserving nonzero-value Morse point, so its vanishing gives defect\ensuremath{\geq}5. In the nonfree case, solving xy=\ensuremath{\alpha}a+\ensuremath{\beta}b identifies the exact pole term V(0)\ensuremath{\alpha}(x,0)/x in the null cubic; its vanishing forces the displayed branch orders. Equality is d=(4,1), b=(0,1), q=0. The invertible coordinate Y=y−ell(x)a gives I=(z\textsuperscript{2},Yz,x\textsuperscript{2}Y). Computed every new 2×2 minor: they force a=yA,c=zC,b=−x\textsuperscript{2}C+zB modulo I, and f\ensuremath{\in}(yz\textsuperscript{2})+I\textsuperscript{2}. The smooth-factor substitution y=z\textsuperscript{2}v and the Morse centre divisible by x\textsuperscript{2}y produce Z\textsuperscript{2}+x4Y. Its actual finite birational normalization (x,u)\ensuremath{\to}(x,−u\textsuperscript{2},x\textsuperscript{2}u) is smooth. The unchanged a3\_node.py replay passed. Review provenance: New derivation after reading the supplied proof; not a blind review and not an independent human check}{\ref{night:P4.28}; \ref{night:P4R.10}; \ref{night:P4R.14}}

At a point of an actual support in Part~4.2, suppose the
entire reduced critical curve is the ordinary node
$C=\{z=xy=0\}$ and both branches have generic transverse
type $A_3$. Let $d_i,b_i,q$ be the Hessian zero orders,
projected conormal colengths, and normalization descent
defect of P4R.10. Then
\[
 \delta(C)+\sum_{i=1}^2\left(\frac{d_i}{2}+b_i\right)+q
                         \ge\frac92.                 \tag{4.41a}
\]
Here $R=0$ and $\delta(C)=1$. Equality makes the normalization
of the hypersurface smooth over the point.
\end{proposition}
\begin{proof}
\noindent\textbf{Step 1.} Put $P=(z,xy)$, and let $I$ be the ambient ideal of the
scheme-theoretic conductor image. We have $P^2=P^{(2)}\subset I$,
and $J=P/I$ is torsionfree of rank one on both branches.
After dividing its normalized image by the two branchwise
vanishing orders, it is an $A=\mathcal O_C$ submodule of
$B=\mathbb C\{x\}\oplus\mathbb C\{y\}$ generating both
summands. The conductor of $A$ in $B$ is $\mathfrak m_A$,
so $\mathfrak m_AB\subset J$. Thus $J$ is determined by
a subspace of $\mathbb C^2$ projecting onto both coordinates.
If that subspace is a line, rescaling gives $J\simeq A$ and
$q=1$. If it is the whole space, $J\simeq B$ and $q=0$.
These are actual module isomorphisms. The divided orders are
the $b_i$ already in the formula.

\noindent\textbf{Step 2.} We use the ring condition of the actual conductor:
\[
 \operatorname{Hom}_{\mathcal O_S}(I,I)=
 \operatorname{Hom}_{\mathcal O_S}(I,\mathcal O_S).
                                                               \tag{4.41b}
\]
The right side is the normalization algebra by finite duality,
and $I$ is an ideal of that algebra. For the determinant
presentation of a Cohen--Macaulay codimension-two conductor
inside a Gorenstein hypersurface, this condition is equivalent
to membership of the adjugate entries in $I$. This is precisely
de Jong--van Straten, Theorem~1.12, printed p.~532
\cite{DeJongVanStraten1990Normalization}. Its hypotheses hold
here: the ambient space is regular, the reduced support is a
hypersurface, and its conductor quotient is Cohen--Macaulay.
We give all the matrices used below.

\noindent\textbf{Step 3.} Suppose first that $q=1$. The surjection
$P/P^2\to J$ is a surjection $A^2\to A$, with primitive
kernel generator $g$. Taking a complementary generator $k$
of $P$ gives $I=(g,k^2)$. A determinant presentation
$f=ag+bk^2$ has two-by-two adjugate containing $a,b,g,k^2$.
Condition (4.41b) says $a,b\in I$, so the exact equation is
\[
                   f=Ag^2+Bgk^2+Ck^4.                \tag{4.41c}
\]
We claim that $A(0)$ cannot be a unit. If $g_z(0)\ne0$,
straighten $g$, take $k$ to be a nodal equation in $g=0$,
and apply the parameter Morse lemma in $g$. The residual
is $k^4$ times a unit after a small change in its fourth-order
coefficient. A morsification of the node and a shift of its
critical value away from zero give an isolated Morse point
of nonzero critical value.

\noindent\textbf{Step 4.} If $g_z(0)=0$, its $xy$ coefficient is a unit and we take
$k=z$. Completing the square in (4.41c), including the
square root of $A$ in its first factor, gives exactly
\[
                         f=G^2+z^4U,
\]
with $G(x,y,0)=xy$ times a unit and $G_z(0)=0$.
Add a small constant to $U$, then take the actual coordinate
$Z=zU^{1/4}$. The parameter Morse lemma for $G$ in $x,y$
gives $G=XY+a(Z)$, where $a\in(Z^2)$. A small change
$a\mapsto a+\eta Z^2$ makes $a''(0)\ne0$. Replacing
$G$ by $G-\epsilon$ produces at the origin a critical point
of value $\epsilon^2$. The $XY$ Hessian entry is
$-2\epsilon$, and the $ZZ$ entry is
$-2\epsilon a''(0)$, so the point is Morse.
These deformations retain type $A_3$ along the moving
critical curve at each fixed boundary. The nested boundary
transport of P4.28 and P4R.14 places its square-$-2$ sphere
in the actual nearby torus, a contradiction. Hence $A(0)=0$.

\noindent\textbf{Step 5.} If $g_z(0)$ is a unit, both saturated Hessian lines approach
$dz$, so $b_1=b_2=1$ and $d_i\ge1$. If $g_z(0)=0$,
the differential of $g$ on each branch has exactly the
common factor $t_i$ and a remaining nonzero coefficient
in the other normal direction. Thus $b_i=0,d_i\ge3$.
In either case the left side of (4.41a) is at least five,
including $q=1$.

\noindent\textbf{Step 6.} Now let $q=0$. Choose actual generators $a,b$ of $P$ mapping
to $(1,0),(0,1)$ in $J\simeq B$. Their coefficient matrix
is invertible modulo the maximal ideal, hence invertible
locally. We obtain
\[
 I/P^2=(ya,xb),\qquad
 f=xU(x)b^2+yV(y)a^2\pmod {P^3}.                      \tag{4.41d}
\]
The second identity records the Hessian kernel $b$ on the
$x$ branch and $a$ on the $y$ branch. Write the exact
relation $xy=\alpha a+\beta b$. Near a general point of
the $x$ axis, solve this for $y$ using coordinates $x,a,b$.
The null cubic coefficient contains
$V(0)\alpha(x,0)/x$; the remaining contribution from
$P^3$ is holomorphic at the node. Its identical vanishing
in type $A_3$ forces $V(0)\alpha(0)=0$. The other branch
gives $U(0)\beta(0)=0$.

\noindent\textbf{Step 7.} If both $a_z(0),b_z(0)$ are nonzero, the relation for $xy$
has $\alpha(0)\beta(0)\ne0$. Hence $U(0)=V(0)=0$,
both $b_i=1$, and both $d_i\ge2$. The defect is at least
five. Otherwise interchange the branches so that $a_z(0)$
is a unit and $b_z(0)=0$. Then $\alpha(0)=0$ and
$\beta(0)$ is a unit, forcing $U(0)=0$. On the $x$ axis,
$db$ has exactly the common factor $x$ and remaining
$dy$ coefficient a unit. On the $y$ axis, $da$ is primitive
with a unit $dz$ coefficient. Thus
\[
 b_x=0,\quad d_x=3+\operatorname{ord}_xU\ge4,
 \qquad
 b_y=1,\quad d_y=1+\operatorname{ord}_yV\ge1.
\]
This proves (4.41a). Equality requires
$\operatorname{ord}_xU=1$ and $V(0)\ne0$.

\noindent\textbf{Step 8.} We identify the normalization in this exact equality case.
Use $a$ as the ambient coordinate $z$. The generator relation
with $\beta$ a unit gives
\[
                    I=(a^2,ya,x^2y-x\alpha a).
\]
Modulo its first two generators, replace $\alpha$ by
$\alpha(x,0,0)=x\ell(x)$. The actual coordinate
$Y=y-\ell(x)a$ puts the ideal into
$I=(z^2,Yz,x^2Y)$. Renaming $Y$ as $y$, its Hilbert--Burch
matrix and determinant extension are
\[
 \begin{pmatrix}y&0\\-z&x^2\\0&-z\end{pmatrix},
 \qquad
 M=\begin{pmatrix}y&0&a\\-z&x^2&b\\0&-z&c\end{pmatrix},
 \qquad
 \det M=az^2+byz+cx^2y.
\]
The new two-by-two minors are
$yb+za,yc,-zc,-x^2a,-za,x^2c+zb$, up to signs.
Their membership in $I$ gives exactly
\[
             a=yA,\qquad c=zC,\qquad b=-x^2C+zB
                          \pmod I.
\]
Therefore $f\in(yz^2)+I^2$, so
\[
           f=Uyz^2+Vz^4+Wx^2y^2z+Tx^4y^2.            \tag{4.41e}
\]
The equality orders $d_y=1,d_x=4$ say precisely
$U(0)T(0)\ne0$.

\noindent\textbf{Step 9.} Substitute $y=z^2v$ and divide (4.41e) by $z^4$.
At the finite root its $v$ derivative is $U(0)$.
The implicit-function theorem supplies the actual smooth
factor $g_1=y-z^2v(x,z)$. Division by it gives
\[
                  g_2=z^2A+x^2yzB+x^4yC,
                  \qquad A(0)C(0)\ne0.
\]
For clarity, this division retains the factors $x^2,x^4$:
subtracting the expression at $y=z^2v$ supplies one factor
$y-z^2v$; every extra substituted $y$ supplies $z^2$ and
is absorbed into the first term. The parameter Morse
centre in $z$ is divisible by $x^2y$. Substitution leaves
$x^4y$ times a unit. An actual coordinate change therefore
gives $g_2=Z^2+x^4Y$. Its finite birational normalization is
\[
                 (x,u)\longmapsto(x,-u^2,x^2u),
\]
with target coordinates $(x,Y,Z)$. It is smooth; finiteness
follows from $u^2=-Y$ and birationality from $u=Z/x^2$
away from $x=0$. The other branch $g_1$ is smooth already.
Thus the full normalization is smooth over the node.
\end{proof}

All determinant minors, the four minimum-defect cases,
the factor substitution, and the normalization equation
are verified by \texttt{code/part4/a3\_node.py}.
The output reports PASS. Its initial failed assertion
compared an expanded polynomial with a factored expression;
that output is retained, and the corrected assertion tests
their exact zero difference.

\begin{proposition}[P4.42]\label{night:P4.42}\label{stmt:P4.42}
\StatementRecord{night:P4.42}{proved}{Historical events: GPT-6 Astra (OpenAI), ChatGPT Pro mode, effort Pro, S6\_06\_BLIND, 16--18 November 2026: targeted proof check. Also checked in this writing. Pre-continuation writing events: GPT-6 Astra (OpenAI), ChatGPT Pro mode, effort Pro, S6\_10\_PAPERS, 18 November 2026 05:10:56.410177+00:00: computation or formal certificate; targeted proof check. Continuation metadata events: GPT-6 Astra (OpenAI), ChatGPT Pro mode, effort Pro, S6\_10\_PAPERS continuation, 19 November 2026: used and rechecked in this run; status-vocabulary reconciliation; metadata reconciliation only; no new mathematical check. \emph{Pre-continuation recorded proof scope.} The night-run re-derivation was by GPT-6 Astra (OpenAI, ChatGPT Pro mode, effort Pro), 16--18 November 2026. Also checked in this writing by GPT-6 Astra (OpenAI), ChatGPT Pro mode, effort Pro, 18 November 2026. proved. Full two-two-branch cycle exclusion, including planar different, nonplanar small-delta model and equality rigidity. For E=t\textasciicircum{}a e1\ensuremath{\oplus}t\textasciicircum{}b e2 and primitive ell=u e1+v e2, independently derived d/2\ensuremath{\geq}max(a−ord u,b−ord v) and projected length min(a+ord v,b+ord u); the sum is at least a+b. The exact plane-curve dualizing pullback dx/h\_y gives sum(a+b)=2delta+degR and local defect\ensuremath{\geq}3delta+q. Nonplanar branches have gradient and Hessian order at least their multiplicity, giving the stated positive contribution after ramification. The sole small delta2 nonplanar two-branch model is cusp plus transverse line; its symbolic-square extra class zh has maximal-ideal annihilator, and rank-one Hessian forces its constant coefficient to vanish, so f\ensuremath{\in}P\textsuperscript{2}. Its conormal orders are (2,3),(1,1), giving defect\ensuremath{\geq}8. Thus each two-branch junction contributes\ensuremath{\geq}9/2, with equality only at the ordinary node of P4.41 and smooth normalization. Re-derived \ensuremath{\chi}OCsch=8−sum(defects) on two normalized P1s. The uncontracted case contradicts \ensuremath{\chi}=g−1\ensuremath{\geq}0; the contracted case forces equality everywhere, and P4R.14 plus absence of any additional positive defect makes the normalization smooth, contradicting a contracted positive-genus section. Review provenance: New derivation after reading the supplied proof; not a blind review and not an independent human check}{\ref{night:P4.33}; \ref{night:P4.34}; \ref{night:P4.40}; \ref{night:P4.41}; \ref{night:P4R.10}; \ref{night:P4R.13}; \ref{night:P4R.14}; \ref{night:P4R.6}}

In the requested global support classes, a positive-genus
ruled-resolution support whose critical curve has generic
type $A_3$ everywhere cannot have two two-branch points as
its point-identification cycle. The two points may be
nonordinary. This includes additional
unibranch singular points of the critical curve.
\end{proposition}
\begin{proof}
\noindent\textbf{Step 1.} After P4R.13, the remaining positive ruled case has two
rational normalized conductor components and one graph
cycle. In the stated instance the cycle is provided by
two two-branch points. We first extend the local lower
bound of P4.41 to these points.

\noindent\textbf{Step 2.} At a planar curve point, write its ideal as $(z,h)$.
Its symbolic square equals its ordinary square, so the
rank-one Hessian belongs to $\operatorname{Sym}^2 E_0$,
where $E_0$ is the normalized conormal image. On a branch
put $E_0=\langle t^a e_1,t^b e_2\rangle$ and write the
primitive Hessian direction as $\ell=u e_1+v e_2$.
One of $u,v$ is a unit. Membership of
$t^d\ell^2$ in $\operatorname{Sym}^2E_0$ and projection
onto the quotient by $\ell$ give
\[
 \frac d2\ge\max\{a-\operatorname{ord}u,
                         b-\operatorname{ord}v\},\qquad
 b_{\rm proj}=\min\{a+\operatorname{ord}v,
                         b+\operatorname{ord}u\}.
\]
The sum of the maximum and minimum on the right is
$a+b$. Thus $d/2+b_{\rm proj}\ge a+b$.
The plane different formula proved in P4.34 is
$\sum(a+b)=2\delta+\deg R$. The whole local defect
is consequently at least $3\delta+q$.
A nonordinary planar two-branch point has $\delta\ge2$
and contributes at least six.

\noindent\textbf{Step 3.} For a nonplanar point, the ambient Hessian vanishes: a
nonzero rank-one Hessian would give $f=X^2+g(Y,Z)$ by
the parameter Morse lemma and would put the entire
critical curve in the smooth plane $X=0$.
On a branch of multiplicity $m$, all conormal generators
have order at least $m$, and every Hessian entry has
order at least $m$. Thus $b_{\rm proj}\ge m$, $d\ge m$,
and, since $r=m-1$,
\[
                  d/2+b_{\rm proj}-r\ge1+m/2\ge1.
\]
When $\delta\ge3$, a nonplanar two-branch point
therefore contributes at least five.
The only remaining possibility has $\delta=2$.
The exact excess-one classification in P4.40 identifies
it with a cusp and transverse line,
\[
                  P=(xz,yz,h),\qquad h=y^2-x^3.
\]
Here $P^{(2)}/P^2$ is generated by $zh$ and is
annihilated by the maximal ideal. Write $f=Azh+P^2$.
Along the cusp, the normal quadratic determinant is
$4KG-A^2$, with $K(0)=0$. Its rank is one generically,
so $A(0)=0$ and in fact $f\in P^2$.
The conormal lattices on the two normalized branches
are respectively $\operatorname{diag}(t^2,t^3)$ and
$\operatorname{diag}(t,t)$. Their total colength is seven,
while $\deg R=1$ and $\delta=2$.
Applying the same symmetric-square inequality gives
defect at least $2+7-1=8$.
Every two-branch point thus contributes at least $9/2$,
and equality is possible only at the ordinary node of
P4.41, whose normalization is smooth.

\noindent\textbf{Step 4.} Apply the rank-one Hessian calculation
of P4R.10 to the two normalized projective lines. On each
line the rank-two conormal kernel has degree $2+\deg R$.
The rank-one quotient has Euler characteristic
$3+\deg R-\deg L$. Combining both lines with the
normalization sequence of $C$ gives the exact formula
\[
 \chi(\mathcal O_{C_{\rm sch}})
   =8-\sum_p\left[
       \delta_p+\sum_{i\text{ over }p}
          \left(\frac{d_i}{2}+b_i-r_i\right)+q_p
                 \right].                            \tag{4.42a}
\]
Every extra unibranch singular point has a strictly positive
bracket by the local conormal inequalities of P4R.10.
At a smooth point the bracket is nonnegative. The two
two-branch points contribute at least nine together by the
local calculation just given.

\noindent\textbf{Step 5.} If the horizontal section is not contracted, P4.33 gives
$\chi(\mathcal O_{C_{\rm sch}})=g-1\ge0$, contradicting
(4.42a), whose right side is at most $-1$.
If the section is contracted, that Euler characteristic is
$-1$. Equality must then hold at both two-branch points and every
other bracket must vanish. There are no additional
unibranch singular points. Both cycle points are ordinary
nodes and the normalization is smooth over them by P4.41.
At every smooth point of the
critical curve, P4R.14 gives either the product tacnode
or the displayed $A_1$ normalization degeneration. The
latter has Hessian zero order one and contributes $1/2$,
so it is excluded by the equality in (4.42a). Thus the
normalization is smooth along all of the conductor, and
it is smooth away from the conductor by P4R.6. A contracted
positive-genus horizontal section in its minimal resolution
is impossible. This proves the assertion.
\end{proof}


\subsection{The uniform cuspidal sector}

\begin{proposition}[P4.43]\label{night:P4.43}\label{stmt:P4.43}
\StatementRecord{night:P4.43}{proved}{Historical events: GPT-6 Astra (OpenAI), ChatGPT Pro mode, effort Pro, S6\_06\_BLIND, 16--18 November 2026: targeted proof check. Also checked in this writing. Pre-continuation writing events: GPT-6 Astra (OpenAI), ChatGPT Pro mode, effort Pro, S6\_10\_PAPERS, 18 November 2026 05:10:56.410177+00:00: computation or formal certificate; targeted proof check. Continuation metadata events: GPT-6 Astra (OpenAI), ChatGPT Pro mode, effort Pro, S6\_10\_PAPERS continuation, 19 November 2026: used and rechecked in this run; status-vocabulary reconciliation; metadata reconciliation only; no new mathematical check. \emph{Pre-continuation recorded proof scope.} The night-run re-derivation was by GPT-6 Astra (OpenAI, ChatGPT Pro mode, effort Pro), 16--18 November 2026. Also checked in this writing by GPT-6 Astra (OpenAI), ChatGPT Pro mode, effort Pro, 18 November 2026. proved. Full uniform-cusp positive-Euler exclusion via the actual normalization and conductor pushout maps. Two local hypersurface branches intersect in a curve in the regular threefold. That would make the generic transverse germ reducible, so uniform A2 makes the finite normalization bijective; properness gives an actual homeomorphism. The resolution contracts vertical rational trees without changing \ensuremath{\pi}1, and contracting the horizontal section kills the entire ruling surface group because its inclusion induces an isomorphism. Thus \ensuremath{\pi}1(N) is either a surface group of genus g\ensuremath{\geq}1 or trivial. The explicit conductor pushout and the positive-Euler graph reduction give \ensuremath{\pi}1(S)=\ensuremath{\pi}1(C)=Z for two rational components and one cycle. Abelianization has rank2g or0 for N, so the homeomorphism cannot induce the required group isomorphism. Followed the exact roles of R7/R9/R13, including exclusion of an elliptic conductor component. Review provenance: New derivation after reading the supplied proof; not a blind review and not an independent human check}{\ref{night:P4R.13}; \ref{night:P4R.7}}

In the requested full-fibration classes, a support with
positive-genus ruled resolution cannot have positive
Euler number and generic transverse type $A_2$ along
every component of its critical curve.
\end{proposition}
\begin{proof}
\noindent\textbf{Step 1.} We use the exact normalization map
$\nu:N\longrightarrow S$ and its minimal resolution
$\rho:T\longrightarrow N$. A germ of $S$ cannot have
two distinct analytic hypersurface branches. Indeed,
their intersection has dimension one in the smooth
ambient threefold. A general point of an irreducible
component of that intersection has at least two local
hypersurface branches, so its transverse curve germ is
reducible. The transverse cusp $A_2$ is irreducible.
This contradicts the stated uniform transverse type.
The points of the finite normalization over a germ are
its analytic branches. Consequently $\nu$ is bijective.
It is proper, so its continuous map from a compact space
to the Hausdorff space $S$ is a homeomorphism.

\noindent\textbf{Step 2.} Let $B$ be the base curve of the minimal ruled surface,
of genus $g\ge1$, and let $H$ be the horizontal section
of coefficient two in the total pole divisor. Point
blowups preserve the fundamental group, and the ruling
gives $\pi_1(T)=\pi_1(B)$. By the fibre-tree calculation
in P4R.7, every vertical exceptional configuration is a
tree of rational curves. Collapsing any such connected
tree preserves the fundamental group: apply van Kampen
to its tubular neighbourhood and the complement, or
equivalently to the cofibration quotient by that
simply connected tree. These quotients are the actual
topological contractions of the resolution map.

\noindent\textbf{Step 3.} If $H$ is not contracted by $\rho$, all resolution
exceptional configurations are of this vertical kind.
It follows that $\pi_1(N)=\pi_1(B)$. If $H$ is
contracted, its connected exceptional configuration
contains $H$. The inclusion $H\hookrightarrow T$
induces an isomorphism on fundamental groups, since
the ruling restricts to the identity section and all
point blowups preserve that isomorphism. The quotient
which collapses this configuration therefore kills
the whole group. The other contracted configurations
are rational trees, giving $\pi_1(N)=1$.

\noindent\textbf{Step 4.} On the other hand, the proved positive ruled reduction,
including P4R.13, gives a conductor curve $C$ with two
rational normalized components and one graph cycle.
The conductor cofibration pushout in P4R.7 gives
$\pi_1(S)=\pi_1(C)=\mathbb Z$. This contradicts the
homeomorphism $N\simeq S$: neither the trivial group
nor the genus-$g$ surface group is $\mathbb Z$, as their
abelianizations have ranks zero and $2g$, respectively.
\end{proof}


\subsection{Further uniform transverse sectors}

\begin{proposition}[P4.44]\label{night:P4.44}\label{stmt:P4.44}
\StatementRecord{night:P4.44}{proved}{Historical events: GPT-6 Astra (OpenAI), ChatGPT Pro mode, effort Pro, S6\_06\_BLIND, 16--18 November 2026: targeted proof check. Also checked in this writing. Pre-continuation writing events: GPT-6 Astra (OpenAI), ChatGPT Pro mode, effort Pro, S6\_10\_PAPERS, 18 November 2026 05:10:56.410177+00:00: computation or formal certificate; targeted proof check. Continuation metadata events: GPT-6 Astra (OpenAI), ChatGPT Pro mode, effort Pro, S6\_10\_PAPERS continuation, 19 November 2026: used and rechecked in this run; status-vocabulary reconciliation; metadata reconciliation only; no new mathematical check. \emph{Pre-continuation recorded proof scope.} The night-run re-derivation was by GPT-6 Astra (OpenAI, ChatGPT Pro mode, effort Pro), 16--18 November 2026. Also checked in this writing by GPT-6 Astra (OpenAI), ChatGPT Pro mode, effort Pro, 18 November 2026. proved. Full contraction of the horizontal section in the remaining uniform-A3 three-branch case. After the two-node-cycle exclusion, the conductor has two rational normalized components and a single three-branch junction. If H survives, its map to one normalized P1 has degree2, using both generic tacnode preimages. The explicit P4R.14 normalizations each have two distinct points over a smooth critical point, so this double cover is unramified there. The exact Euler identity bounds the total defect by9−g\ensuremath{\leq}8. A nonplanar triple consumes\ensuremath{\geq}5, a planar triple\ensuremath{\geq}9, and any extra unibranch singularity\ensuremath{\geq}2. Hence H has at most two possible branch values over the triple plus one elsewhere. A degree2 map from genus g\ensuremath{\geq}1 requires2g+2\ensuremath{\geq}4 distinct branch values. The mismatch proves contraction using the actual finite map, not only the incidence graph. Review provenance: New derivation after reading the supplied proof; not a blind review and not an independent human check}{\ref{night:P4.33}; \ref{night:P4.42}; \ref{night:P4R.14}; \ref{night:P4R.6}}

In the remaining positive ruled case with generic transverse
type $A_3$ on every conductor component, the horizontal
section of the pole divisor is contracted by the resolution
map. The point-identification cycle occurs at a single
three-branch point of the critical curve.
\end{proposition}
\begin{proof}
\noindent\textbf{Step 1.} P4.42 excludes a cycle at two two-branch points. The
remaining graph has two rational normalized components
and one cycle, so its cycle occurs at one three-branch
point $p$. Suppose that the horizontal section $H$ is
not contracted. It maps finitely to one normalized
conductor component $C_H^\nu=\mathbb P^1$.
The generic tacnode has two conductor preimages.
The degree of $H\to\mathbb P^1$ is greater than one,
since $H$ has genus $g\ge1$, and hence is exactly two.
No other, vertical conductor component can map onto
$C_H$: the horizontal map already supplies both generic
preimages.

\noindent\textbf{Step 2.} At a point where the entire critical curve is smooth,
P4R.14 gives either the product tacnode or the germ
$X(Y^2+XZ)$. Each has two distinct normalization points
over the critical-curve point. Each point lies on the
strict conductor, by P4R.6. Thus the two preimages must
both lie on $H$, and the degree-two map is unramified
there. All ramification is over singular points of $C$.

\noindent\textbf{Step 3.} By (4.42a) and P4.33 the sum of the local defects is
$9-g\le8$. A nonplanar three-branch point has
$\delta\ge2$, and the branchwise estimate in P4.42
adds at least one per branch; its defect is at least
five. A planar three-branch point has $\delta\ge3$
and contributes at least $3\delta\ge9$, which is
already impossible. Every additional unibranch singular
point contributes at least two: in the planar case
the bound is $3\delta$, and in the nonplanar case it is
at least $\delta+1$. There is consequently at most one
additional unibranch singular point.

\noindent\textbf{Step 4.} The normalization $C_H^\nu$ has at most two points
above $p$, since the other global conductor component
supplies at least one of its three branches. Including
the possible additional unibranch point leaves at most
three possible branch values for $H\to\mathbb P^1$.
Every ramification point of a degree-two map has index
two, and its branch value has only that ramified
preimage. Riemann--Hurwitz requires $2g+2\ge4$ distinct
branch values. This contradiction proves that $H$ is
contracted.
\end{proof}

\begin{proposition}[P4.45]\label{night:P4.45}\label{stmt:P4.45}
\StatementRecord{night:P4.45}{proved}{Historical events: GPT-6 Astra (OpenAI), ChatGPT Pro mode, effort Pro, S6\_06\_BLIND, 16--18 November 2026: targeted proof check. Also checked in this writing. Pre-continuation writing events: GPT-6 Astra (OpenAI), ChatGPT Pro mode, effort Pro, S6\_10\_PAPERS, 18 November 2026 05:10:56.410177+00:00: computation or formal certificate; targeted proof check. Continuation metadata events: GPT-6 Astra (OpenAI), ChatGPT Pro mode, effort Pro, S6\_10\_PAPERS continuation, 19 November 2026: used and rechecked in this run; status-vocabulary reconciliation; metadata reconciliation only; no new mathematical check. \emph{Pre-continuation recorded proof scope.} The night-run re-derivation was by GPT-6 Astra (OpenAI, ChatGPT Pro mode, effort Pro), 16--18 November 2026. Also checked in this writing by GPT-6 Astra (OpenAI), ChatGPT Pro mode, effort Pro, 18 November 2026. proved. Full planar-singular ordinary-triple exclusion with explicit critical-value deformation. The complete-intersection symbolic cube gives f=Au\textsuperscript{3}+Bu\textsuperscript{2}h+Cuh\textsuperscript{2}+Dh\textsuperscript{3}. Replaced h by an independent H in the actual analytic F and made arbitrarily small coefficient changes so its binary cubic has a simple critical root with nonzero critical value. The implicit-function root u=Hv gives F=H\textsuperscript{3}beta with beta a unit. Inverted the analytic map H\ensuremath{\to}H beta\textasciicircum{}(1/3), then chose H\_s so the critical-value function is exactly a morsification hhat\_s\textsuperscript{3}. At a nonzero-value plane Morse point the Schur complement is3hhat\_s\textsuperscript{2} Hess(hhat\_s), while F\_uu is nonzero. The resulting full Morse point retains ordinary-triple transverse boundary germs. The exact P4.20/P4.28 boundary transport embeds its square−2 sphere in the actual nearby torus, giving the exclusion. Review provenance: New derivation after reading the supplied proof; not a blind review and not an independent human check}{\ref{night:P4.28}}

In an actual torus degeneration, a point of a planar
singular critical curve cannot have generic transverse
type $D_4$ along all its local branches.
\end{proposition}
\begin{proof}
\noindent\textbf{Step 1.} Write the full reduced critical-curve ideal as
$P=(u,h(x,y))$, where $h$ is reduced and singular.
Generic multiplicity three gives $f\in P^{(3)}$.
Powers of this complete-intersection prime-by-components
ideal are unmixed, so $P^{(3)}=P^3$. Thus there are
holomorphic coefficients such that
\[
 f=A u^3+B u^2h+Cuh^2+Dh^3.
\]
Introduce an independent variable $H$ and the actual
holomorphic function
\[
 F(u,x,y,H)=A(u,x,y)u^3+B(u,x,y)u^2H
               +C(u,x,y)uH^2+D(u,x,y)H^3.
\]
Then $f=F(u,x,y,h)$. Arbitrarily small constant changes
of the four coefficients preserve the $D_4$ transverse
germs on a fixed boundary and make the binary cubic
at the origin have distinct roots and a simple critical
point as a polynomial in $u/H$. Its critical value is
nonzero: a zero critical value would be a repeated root.
We also choose its leading coefficient nonzero.

\noindent\textbf{Step 2.} After substituting $u=Hv$, the quotient $F_u/H^2$
is holomorphic. The implicit-function theorem gives
a holomorphic critical root $v=v(x,y,H)$ and
\[
          F(Hv(x,y,H),x,y,H)=H^3\beta(x,y,H),
                       \qquad \beta(0)\ne0.
\]
Choose the local cube root of the unit $\beta$.
The map $H\mapsto H\beta(x,y,H)^{1/3}$ has nonzero
$H$ derivative at the origin and has a holomorphic
inverse with parameters $x,y$. Put
\[
             \widehat h(x,y)=h(x,y)
                           \beta(x,y,h(x,y))^{1/3}.
\]
It defines the same singular reduced plane curve as $h$.
Choose a small morsification $\widehat h_s$ with a
Morse critical point of nonzero critical value. This
exists by the positive local degree of the gradient
of the singular reduced plane function; a small constant
shift avoids zero critical values. Use the displayed
inverse to solve exactly
\[
         H_s\beta(x,y,H_s)^{1/3}=\widehat h_s,
             \qquad f_s=F(u,x,y,H_s(x,y)).
\]
On its $u$-critical branch the critical-value function
is exactly $\widehat h_s^3$. At the selected plane
Morse point, $H_s\ne0$ and $F_{uu}\ne0$. The Schur
complement of the full Hessian is
$3\widehat h_s^2\operatorname{Hess}(\widehat h_s)$.
Consequently $f_s$ has a nondegenerate three-variable
critical point with nonzero critical value.

\noindent\textbf{Step 3.} All deformations above can be chosen successively small
on the fixed boundary. There the moving curve $H_s=0$
has the same ordinary-triple transverse germ: its binary
cubic has three distinct linear factors. The boundary
transport used in P4.28 therefore identifies the local
nearby fibres through nested representatives. The Morse
vanishing sphere embeds in a nearby global torus fibre
and has self-intersection $-2$, contradicting its
asphericity. This proves the assertion.
\end{proof}

\begin{proposition}[P4.46]\label{night:P4.46}\label{stmt:P4.46}
\StatementRecord{night:P4.46}{proved}{Historical events: GPT-6 Astra (OpenAI), ChatGPT Pro mode, effort Pro, S6\_06\_BLIND, 16--18 November 2026: targeted proof check. Also checked in this writing. Pre-continuation writing events: GPT-6 Astra (OpenAI), ChatGPT Pro mode, effort Pro, S6\_10\_PAPERS, 18 November 2026 05:10:56.410177+00:00: computation or formal certificate; targeted proof check. Continuation metadata events: GPT-6 Astra (OpenAI), ChatGPT Pro mode, effort Pro, S6\_10\_PAPERS continuation, 19 November 2026: used and rechecked in this run; status-vocabulary reconciliation; metadata reconciliation only; no new mathematical check. \emph{Pre-continuation recorded proof scope.} The night-run re-derivation was by GPT-6 Astra (OpenAI, ChatGPT Pro mode, effort Pro), 16--18 November 2026. Also checked in this writing by GPT-6 Astra (OpenAI), ChatGPT Pro mode, effort Pro, 18 November 2026. proved. Full uniform-D4 positive-Euler exclusion using the complete conormal/descent defect. The generic conductor is the actual symbolic double ideal P\textasciicircum{}(2), by the conductor square and depth-one argument; Riemann–Roch on the two normalized P1s gives \ensuremath{\chi}=10−sum(delta+tau−r), so the uncontracted total is11−g\ensuremath{\leq}10. At a nonplanar triple of smooth branches, a generic linear section has length3 and embedding dimension2, hence ideal(v,w)\textsuperscript{2}. The nonzerodivisor Tor sequence and Nakayama lift exactly3 defining generators. The normalized conormal needs6 generator values and its branch lattices already have colength\ensuremath{\geq}6; the descent therefore adds\ensuremath{\geq}3, producing defect\ensuremath{\geq}11. A singular branch gives total multiplicity\ensuremath{\geq}4; the only apparent small case m=4,delta=3 is the cusp plus two independent lines, of embedding dimension4, excluded by the explicit P4.40 ring model. The remaining two-junction alternative forces two planar ordinary nodes and g=1; each node is excluded by P4.45. Contracted H is separately excluded by the corrected P4.39 proof. Review provenance: New derivation after reading the supplied proof; not a blind review and not an independent human check}{\ref{night:P4.31}; \ref{night:P4.33}; \ref{night:P4.34}; \ref{night:P4.39}; \ref{night:P4.40}; \ref{night:P4.45}}

In the requested full-fibration classes, a support with
positive-genus ruled resolution cannot have positive
Euler number and generic transverse type $D_4$ along
every component of its critical curve.
\end{proposition}
\begin{proof}
\noindent\textbf{Step 1.} If the horizontal section is contracted, P4.39 already
gives a contradiction. Suppose it is not contracted.
The remaining critical curve has two rational normalized
components and one graph cycle. For the actual
conductor scheme, ordinary-triple transverse type gives
the ideal $P^{(2)}$, as in P4.31. Let $Q=P/P^{(2)}$,
and let $\tau_p$ be the full colength of its differentiated
image in the direct image of the normalized rank-two
kernel $K$. This includes both the branchwise lattice
colength and the further descent colength. Riemann--Roch
on the two normalized projective lines gives
\[
 \chi(\mathcal O_{C_{\rm sch}})
       =10-\sum_p(\delta_p+\tau_p-r_p).
\]
P4.33 identifies the left side with $g-1$, so the total
defect is $11-g\le10$.

\noindent\textbf{Step 2.} We prove that a three-branch point contributes at least
eleven. A planar point contributes $5\delta\ge15$ by
the exact different calculation in P4.34. Consider a
nonplanar point. If its three branches are smooth, the
curve has multiplicity three and embedding dimension
three. A generic linear parameter has an Artinian
section of length three and embedding dimension two,
with ideal $(v,w)^2$. The defining curve ideal therefore
has exactly three minimal generators. To justify the
lifting, the parameter is a nonzerodivisor on the
Cohen--Macaulay curve ring, so the Tor term vanishes in
the ideal sequence modulo that parameter; Nakayama
then lifts precisely the three generators.
Each branchwise conormal lattice has rank two and
colength at least two, since every defining gradient
has positive order. Their sum has colength at least
six. Modulo the three branch maximal ideals its
normalization has six generator values. The image of
$Q$, generated by three elements over the local curve
ring, supplies a vector space of dimension at most
three. The additional descent colength is therefore
at least three. Since $\delta\ge2$ and $R=0$, the
full defect is at least $2+6+3=11$.

\noindent\textbf{Step 3.} If one of the three branches is singular, let $m\ge4$
be total multiplicity. The branchwise gradient bound
gives $\tau\ge2m$, while $\deg R=m-3$.
The general inequality $\delta\ge m-1$ leaves only
$m=4,\delta=3$ as a possible value below the required
bound. That value has excess $\delta-(3-1)=1$.
The exact classification in P4.40 shows that its only
singular-branch model is a cusp and two independent
lines, of embedding dimension four. It does not embed
in the ambient smooth threefold. Thus $\delta\ge4$
in the present case, and the full defect is at least
$\delta+2m-(m-3)=\delta+m+3\ge11$.
A cycle at one three-branch point is impossible.

\noindent\textbf{Step 4.} The other possible cycle has two two-branch points.
Each contributes at least five by P4.34; equality is
possible only for a planar ordinary node. A nonplanar
two-branch point contributes at least seven, and a
nonordinary planar point at least ten. The total bound
of ten therefore forces two ordinary planar nodes
and $g=1$. Each such node is forbidden by P4.45.
This proves the assertion.
\end{proof}

The local Hessian identities, rank-one lattice identity,
and exact integer defect bounds are checked in
\texttt{code/part4/uniform\_sectors.py}; the saved output
reports PASS.


\subsection{The full quadratic conductor at three axes}

\begin{proposition}[P4.47]\label{night:P4.47}\label{stmt:P4.47}
\StatementRecord{night:P4.47}{proved}{Historical events: GPT-6 Astra (OpenAI), ChatGPT Pro mode, effort Pro, S6\_06\_BLIND, 16--18 November 2026: targeted proof check. Also checked in this writing. Pre-continuation writing events: GPT-6 Astra (OpenAI), ChatGPT Pro mode, effort Pro, S6\_10\_PAPERS, 18 November 2026 05:10:56.410177+00:00: computation or formal certificate; targeted proof check; used and rechecked in this run. Continuation metadata events: GPT-6 Astra (OpenAI), ChatGPT Pro mode, effort Pro, S6\_10\_PAPERS continuation, 19 November 2026: used and rechecked in this run; status-vocabulary reconciliation; metadata reconciliation only; no new mathematical check. \emph{Pre-continuation recorded proof scope.} The night-run re-derivation was by GPT-6 Astra (OpenAI, ChatGPT Pro mode, effort Pro), 16--18 November 2026. Also checked in this writing by GPT-6 Astra (OpenAI), ChatGPT Pro mode, effort Pro, 18 November 2026. proved. Full all-three-coefficient ordinary-axis equality normalization and boundary-preserving cone-polar exclusion. The maximal-ideal annihilator of xyz modulo P\textsuperscript{2} and the leading rank-one Hessian force f\ensuremath{\in}P\textsuperscript{2}. Defect9 gives descent length1; the rank-two quotient-value matrix and actual branch-coordinate multiplications lift its kernel to g\ensuremath{\in}I with quadratic part axy+bxz+cyz. The complete intersection(g,xyz) contains P\textsuperscript{2}, agrees generically with I, and its Cohen–Macaulay torsionfree quotient rules out a finite kernel, proving exact ideal equality. The adjugate condition gives f\ensuremath{\in}I\textsuperscript{2}, hence f=G\textsuperscript{2}+(xyz)\textsuperscript{2}U with a nondegenerate quadratic G. A unit U gives two ordinary quadratic surface branches. For U(0)=0, the fixed-axis scaling yields the explicit tangential function r6 H(t)\textsuperscript{2}(r ell(t)−1). Independently recomputed its derivative and all3 nonzero-value, nondegenerate critical points, including their common rational value and Hessian determinant. The implicit-function theorem lifts these to actual nonzero-value Morse points with critical values of order epsilon7. Boundary A3 is preserved because only U changes and its outer transverse coefficients stay units. The unchanged a3\_three\_axes.py replay passed. Review provenance: New derivation after reading the supplied proof; not a blind review and not an independent human check}{\ref{night:P4.41}}

Consider an actual torus degeneration at a point where the
entire reduced critical curve is the three coordinate axes,
with generic transverse type $A_3$ on each axis. Suppose
the local Hessian--conductor defect is nine, the three
Hessian orders are two, and the quartic equation is the
square of $a xy+b xz+c yz$ with $abc\ne0$.
Then the normalization has only rational singularities
over that point. In particular it cannot contain the
contraction of a positive-genus horizontal section.
\end{proposition}
\begin{proof}
\noindent\textbf{Step 1.} Put $P=(xy,xz,yz)$ and $h=xyz$. The symbolic square is
$P^{(2)}=P^2+(h)$, and the actual conductor image ideal
$I$ contains $P^{(2)}$. Also $f\in P^{(2)}$.
Write $f=Ah+P^2$. On the $x$-axis, a nonzero constant
term of $A$ would give a leading rank-two transverse
Hessian, since the quadratic term is then a nonzero
multiple of $xyz$. Generic type $A_3$ has rank one.
Thus $A(0)=0$. Since $\mathfrak m h\subset P^2$,
we have $f\in P^2$.

\noindent\textbf{Step 2.} The normalized conormal lattice is $tK$ on each axis,
so its projected colength is one. The curve has
$\delta=2$ and $R=0$. The stated defect therefore
forces the remaining normalization descent colength
to be one. The normalized quotient values form a
two-dimensional subspace of $\mathbb C^3$ with one
relation involving all three coordinates. Multiplying
a source generator that is a unit on one branch by
that branch's axis coordinate generates its positive
orders alone. Hence all positive-order branch functions
occur in the image, and the leading kernel vector
lifts to an actual element $g\in I$ with quadratic
part $a xy+b xz+c yz$.

\noindent\textbf{Step 3.} We prove the exact ideal equality $I=(g,h)$.
On each coordinate plane, $g$ is the product of its
two coordinate functions times a unit. Thus
$I_0=(g,h)$ has radical $P$. Furthermore $P^2\subset I_0$:
modulo $h$, multiplication of $g$ by $xy$ gives a unit
times $(xy)^2$, and the other two squares follow in
the same way; all mixed products are multiples of $h$.
At the generic point of every axis, $I_0$ is the
tacnode conductor ideal, generated by a normal coordinate
and the square of the null coordinate. It agrees
generically with $I$. The ideal $I_0$ is a
codimension-two complete intersection, so
$\mathcal O/I_0$ is Cohen--Macaulay and $P/I_0$ is
torsionfree on the reduced curve. The map
$P/I_0\to P/I$ has finite-length kernel, which must
vanish. This proves $I=I_0$.

\noindent\textbf{Step 4.} The actual conductor ring condition, with the hypotheses
verified in P4.41, gives $f\in I^2$ by the two-by-two
adjugate calculation of de Jong--van Straten,
Theorem~1.12. Write $f=A g^2+Bgh+Ch^2$.
Its quartic part makes $A(0)\ne0$. Completing the square
and absorbing $\sqrt A$ into the first factor gives
the exact equation
\[
                         f=G^2+h^2U.                 \tag{4.47a}
\]
The quadratic part of $G$ is nondegenerate, and
$(G,h)=(g,h)$. If $U(0)\ne0$, the two factors
$G\pm h\sqrt{-U}$ have nondegenerate quadratic
Hessians. Each is an isolated $A_1$ surface singularity,
so the full normalization has rational singularities.

\noindent\textbf{Step 5.} We exclude $U(0)=0$ by a deformation that keeps the
three axes fixed. Add an arbitrarily small generic
linear form to $U$, and then replace $U$ by
$U-\epsilon$. On a fixed boundary the transverse
$A_3$ coefficient is a unit, so these changes retain
the boundary transverse type. A diagonal linear change
puts the quadratic part of $G$ in the form
$G_2=xy+xz+yz$, changing $h$ only by a nonzero constant.
Let $L$ be the resulting linear part of $U$.
Under $(x,y,z)=\epsilon(X,Y,Z)$, equation (4.47a) becomes
\[
 \frac{f_\epsilon}{\epsilon^4}
   =G_\epsilon(X,Y,Z)^2+\epsilon^3(XYZ)^2
                          \bigl(L(X,Y,Z)-1+O(\epsilon)\bigr),
 \qquad G_\epsilon=G_2+O(\epsilon).
\]
Near a nonzero point of the smooth cone $G_2=0$,
use $G_\epsilon$ as the normal coordinate. The
parameter Morse centre is $O(\epsilon^3)$.
Substitution leaves the tangential function
$\epsilon^3(\Phi+O(\epsilon))$, where
\[
                 \Phi=(XYZ)^2(L-1)|_{G_2=0}.
\]

\noindent\textbf{Step 6.} We verify a nonempty open set of choices of $L$ for
which $\Phi$ has a nondegenerate critical point of
nonzero value. Parametrize the projective conic by
\[
              [X:Y:Z]=[t(t+1):t+1:-t].
\]
For $L=X+Y+Z$, set
$H(t)=-t^2(t+1)^2$ and $\ell(t)=t^2+t+1$.
In radial coordinate $r$,
$\Phi=r^6H(t)^2(r\ell(t)-1)$.
Radial criticality gives $r=6/(7\ell(t))$, and
\[
 \frac{d}{dt}\left(\frac{H}{\ell^3}\right)
  =\frac{t(t-1)(t+1)(t+2)(2t+1)}{(t^2+t+1)^4}.
\]
The three critical points $t=1,-2,-1/2$ are simple
and avoid zeros and poles. Their radial coordinates
are $2/7,2/7,8/7$. At each point,
\[
       \Phi=-\frac{1024}{823543},\qquad
       \det\operatorname{Hess}_{r,t}\Phi
                     =\frac{4718592}{1977326743}.
\]
These numbers are nonzero. Simplicity and avoidance
of zeros and poles are nonvanishing algebraic conditions
on the coefficients of $L$. The explicit example proves
that they hold on a nonempty Zariski open set, so an
arbitrarily small generic linear perturbation of any
initial linear part is sufficient.

\noindent\textbf{Step 7.} After eliminating the normal Morse coordinate, divide
the two tangential critical equations by $\epsilon^3$.
They are analytic in $\epsilon$ and have invertible
Jacobian at the displayed points. The implicit-function
theorem lifts them to nondegenerate three-variable
critical points of $f_\epsilon$. Their critical values
have leading term $\epsilon^7$ times the displayed
nonzero value. They tend to the original point.
Since the deformation retains boundary $A_3$, the
nested boundary transport embeds a Morse vanishing
sphere of square $-2$ into a nearby global torus fibre.
This contradicts asphericity. Therefore $U(0)=0$
cannot occur, proving the assertion.
\end{proof}

The conic derivative, all three exact critical values
and Hessian determinants, and the related conductor
matrices are executed in
\texttt{code/part4/a3\_three\_axes.py}.
The output reports PASS. The working record explicitly
rejects the earlier proposed change $G\mapsto G-\epsilon$:
with $h$ fixed that change replaces the boundary tacnode
by two ordinary double points and does not supply the
required boundary transport.


\subsection{The three-branch tacnode defect budget}

\begin{proposition}[P4.48]\label{night:P4.48}\label{stmt:P4.48}
\StatementRecord{night:P4.48}{proved}{Historical events: GPT-6 Astra (OpenAI), ChatGPT Pro mode, effort Pro, S6\_06\_BLIND, 16--18 November 2026: targeted proof check. Also checked in this writing. Pre-continuation writing events: GPT-6 Astra (OpenAI), ChatGPT Pro mode, effort Pro, S6\_10\_PAPERS, 18 November 2026 05:10:56.410177+00:00: computation or formal certificate; targeted proof check. Continuation metadata events: GPT-6 Astra (OpenAI), ChatGPT Pro mode, effort Pro, S6\_10\_PAPERS continuation, 19 November 2026: used and rechecked in this run; status-vocabulary reconciliation; metadata reconciliation only; no new mathematical check. \emph{Pre-continuation recorded proof scope.} The night-run re-derivation was by GPT-6 Astra (OpenAI, ChatGPT Pro mode, effort Pro), 16--18 November 2026. Also checked in this writing by GPT-6 Astra (OpenAI), ChatGPT Pro mode, effort Pro, 18 November 2026. proved. Full planar/ordinary-axis triple defect bound and exact quartic-square classification. A planar conductor quotient is generated by2 elements while its three normalized branches need3 values, so q\ensuremath{\geq}1 and 3delta+q\ensuremath{\geq}10. At the three axes, P\textasciicircum{}(2)=P\textsuperscript{2}+(xyz) and the leading transverse determinant force f\ensuremath{\in}P\textsuperscript{2}; all normalized conormal lattices are exactly tK. Re-derived the three principal-minor equations for the quartic tensor. When all diagonal coefficients are nonzero, substituting the quadratic null direction on the x-axis makes its cubic coefficient linear in x proportional to F−BD/A; higher total-degree terms cannot change that coefficient. Hence the full quartic is a square. A zero square gives defect\ensuremath{\geq}19/2; a single nonzero coefficient leaves one Hessian order\ensuremath{\geq}3 and a quotient-value matrix of rank\ensuremath{\leq}2, again giving\ensuremath{\geq}19/2. At least two coefficients give Hessian orders(2,2,2), projected orders(1,1,1), and the displayed rank-two three-row matrix forces q\ensuremath{\geq}1. Equality is exactly the retained data. The finite matrix and rational-cone checks passed in the unchanged replay. Review provenance: New derivation after reading the supplied proof; not a blind review and not an independent human check}{\ref{night:P4.42}; \ref{night:P4.47}}

At a planar three-branch critical-curve point with uniform
generic transverse type $A_3$, the local conductor defect
is at least ten. At an ordinary nonplanar three-branch
point, the defect is at least nine. Equality in the
second assertion forces Hessian orders $(2,2,2)$,
projected colengths $(1,1,1)$, descent colength one,
and a quartic square $(a xy+b xz+c yz)^2$ with at
least two nonzero coefficients. The case $abc\ne0$
has rational normalization by P4.47.
\end{proposition}
\begin{proof}
\noindent\textbf{Step 1.} For a planar critical curve the ideal is $P=(u,h)$.
The rank-one nilradical $J=P/I$ of the conductor image
is generated by at most two elements over $\mathcal O_C$.
Its normalized image has at least one generator value
on each of the three branches. Modulo the three branch
maximal ideals, the source supplies a vector space of
dimension at most two, while the normalized image
supplies dimension three. Hence its descent colength
$q$ is at least one. The planar symmetric-square and
different calculation in P4.42 gives local defect
$\ge3\delta+q$. Since a planar three-branch point has
$\delta\ge3$, its defect is at least ten.

\noindent\textbf{Step 2.} For an ordinary nonplanar point use actual coordinates
with $P=(xy,xz,yz)$. The argument in P4.47 gives
$f\in P^2$. On each axis, its normalized conormal
lattice is exactly $tK$, so the projected colength
is exactly one. Every Hessian order is at least two.
Write the quartic form in the basis $(xy,xz,yz)$ as
\[
 A(xy)^2+C(xz)^2+E(yz)^2
      +2B(xy)(xz)+2D(xy)(yz)+2F(xz)(yz).
\]
The rank-one Hessians give
$AC=B^2$, $AE=D^2$, $CE=F^2$.
If a diagonal coefficient vanishes, its adjacent
off-diagonal coefficients vanish and the remaining
two-by-two block has rank at most one. If all three
diagonal coefficients are nonzero, the vanishing null
cubic on the $x$-axis gives $F=BD/A$. Higher-degree
terms cannot alter the coefficient linear in $x$
of that normal cubic. Thus the whole quartic is a
square $(a xy+b xz+c yz)^2$.

\noindent\textbf{Step 3.} If the quartic is zero, all three Hessian orders are
at least three, and the local defect is at least
$2+3+9/2=19/2$. If the square has just one nonzero
coefficient, say $a xy$, the $z$-axis Hessian order
is at least three. In the source-generator basis
$(xy,xz,yz)$, the three leading quotient-value rows
have the forms
\[
                      (0,1,0),\qquad(0,0,1),
                      \qquad(0,*,*).
\]
Their rank is at most two, so $q\ge1$ and again
the defect is at least $19/2$.

\noindent\textbf{Step 4.} If at least two coefficients are nonzero, all three
Hessian orders are two. The leading quotient-value
matrix has rank two. For example, up to rescaling
its rows it is
\[
 \begin{pmatrix}-b&a&0\\-c&0&a\\0&-c&b\end{pmatrix}.
\]
It has zero determinant and rank two whenever at
least two of $a,b,c$ are nonzero. Therefore $q\ge1$,
and the local defect is at least $2+3+3+1=9$.
Equality forces $q=1$ and the displayed Hessian and
projection orders. P4.47 applies when all three
quadratic coefficients are nonzero.
\end{proof}


\subsection{The remaining ordinary-axis conductor}

\begin{proposition}[P4.49]\label{night:P4.49}\label{stmt:P4.49}
\StatementRecord{night:P4.49}{proved}{Historical events: GPT-6 Astra (OpenAI), ChatGPT Pro mode, effort Pro, S6\_06\_BLIND, 16--18 November 2026: targeted proof check. Also checked in this writing. Pre-continuation writing events: GPT-6 Astra (OpenAI), ChatGPT Pro mode, effort Pro, S6\_10\_PAPERS, 18 November 2026 05:10:56.410177+00:00: computation or formal certificate; targeted proof check; used and rechecked in this run. Continuation metadata events: GPT-6 Astra (OpenAI), ChatGPT Pro mode, effort Pro, S6\_10\_PAPERS continuation, 19 November 2026: used and rechecked in this run; status-vocabulary reconciliation; metadata reconciliation only; no new mathematical check. \emph{Pre-continuation recorded proof scope.} The night-run re-derivation was by GPT-6 Astra (OpenAI, ChatGPT Pro mode, effort Pro), 16--18 November 2026. Also checked in this writing by GPT-6 Astra (OpenAI), ChatGPT Pro mode, effort Pro, 18 November 2026. proved. Full two-form conductor classification, both normalization algebras, all polar deformations and exact rational-singularity deductions. Lifted the missing-coefficient quadratic kernel to g=x(y+z)−y z c. The Hilbert–Burch minors and the torsionfree generic comparison give I=(g,xyz,(yz)\textsuperscript{2}). Analytic division modulo(x,y+z,yz), followed by the invertible map x\ensuremath{\to}x−yzB, gives exactly the two lambda=0,1 ideals. For lambda0 independently solved every adjugate condition using I:h,I:x,I:k, obtaining f=A x\textsuperscript{2}k\textsuperscript{2}+B x\textsuperscript{2}kh+Cxh\textsuperscript{2}+Dh4 and the stated square/coordinate completion. A unit C gives a smooth factor and the finite birational normal quadric u(y+z)=yz with X=−u\textsuperscript{2}. C(0)=0 gives the explicit weighted polar; recomputed its reduced critical polynomial, nonzero critical values and Hessian resultants. For lambda1 independently followed the three colon ideals, the coefficient module and the exact identities reducing to Ag\textsuperscript{2}+Bgxh+Cgh\textsuperscript{2}+D(xh)\textsuperscript{2}+kappa p. Re-derived the four multiplication relations from the actual adjugate fractions v=−g/h and w=−(ak+bh)/h\textsuperscript{2}. For kappa a unit, the last relation eliminates k, and the displayed invertible linear map identifies the tangent-cone ideal with the prime rational cubic cone ideal; same dimension prevents an additional initial equation. Its blowup is smooth along the degree3 rational exceptional curve, whose O(3n) thickening layers have H1=0, giving rationality by formal functions. For kappa(0)=0, completed the normal square and used XH\textsuperscript{3}=−Yp,H4=−Kp on G=0 to obtain the stated tangential polar. Its exact nonzero value and Hessian at(t,r)=(1/2,−7/6) give actual Morse points of order epsilon8. Both deformations preserve the fixed-axis A3 boundary data. The unchanged a3\_missing\_coefficient.py passed all colon, multiplication, resultant and tangent-cone checks. Review provenance: New derivation after reading the supplied proof; not a blind review and not an independent human check}{\ref{night:P4.41}; \ref{night:P4.47}; \ref{night:P4.48}}

At an ordinary three-axis critical-curve point with uniform
generic transverse type $A_3$ and local defect nine, the
normalization has rational singularities. This includes the
quartic square with exactly one missing coefficient left
by P4.47--P4.48.
\end{proposition}
\begin{proof}
\noindent\textbf{Step 1.} Only the one-missing-coefficient case remains. Rescale
its leading conductor form to $x(y+z)$.
As in P4.47, the descent colength one lifts the leading
kernel vector to $g\in I$ with this quadratic part.
Write $g=x(\alpha y+\beta z)+yz\gamma$, where
$\alpha,\beta$ are units and $\gamma(0)=0$.
Replacing $y,z$ by $\alpha y,\beta z$ is an analytic
coordinate change preserving the axes. Put $k=y+z$,
$h=yz$. Then $g=xk-hc$, with $c(0)=0$.
The ideal $I_0=(g,xh,h^2)$ lies in the actual conductor
$I$ and has Hilbert--Burch matrix
\[
                 \begin{pmatrix}h&0\\-k&h\\c&-x\end{pmatrix}.
\]
It has radical $(xy,xz,yz)$, is Cohen--Macaulay, and
agrees with $I$ generically on each axis. The
torsionfree finite-kernel comparison used in P4.47
therefore gives the exact equality $I=I_0$.

\noindent\textbf{Step 2.} There are exactly two required ideal forms. Since
$\mathcal O/(x,k,h)=\mathbb C\{y\}/(y^2)$, analytic
division writes $c=\lambda y+xA+kB+hC$ for a scalar
$\lambda$. Replace $g$ by $g+xhA+h^2C$, and then
set $x'=x-hB$. These generator and coordinate changes
give
\[
                    I=(x'k-\lambda yh,x'h,h^2).
\]
If $\lambda=0$ this is $(xk,xh,h^2)$. If
$\lambda\ne0$, rescaling $x'=\lambda X$ gives
$(Xk-yh,Xh,h^2)$. Every change is invertible near
the origin. We treat the two forms separately.

\paragraph{The ideal $(xk,xh,h^2)$.}
Appending $(a,b,c)$ to its Hilbert--Burch matrix gives
determinant $axk+bxh+ch^2$. The adjugate ring condition
of de Jong--van Straten, Theorem~1.12, applies to the
actual conductor as in P4.41. The colon ideals
\[
 I:h=(x,h),\qquad I:x=(k,h),\qquad I:k=(x,h^2)
\]
give exactly
$f=A x^2k^2+B x^2kh+Cxh^2+Dh^4$, with $A(0)\ne0$.
Complete the square in $k$, absorb $\sqrt A$ into $x$,
and replace $y$ by $y(1+\gamma z)$ to make
$k+\gamma h$ the new $y+z$. These are analytic
coordinate maps; $h$ changes only by a unit.
Absorbing that unit in the coefficients gives
\[
                       f=x^2k^2+xh^2C+h^4D.          \tag{4.49a}
\]
If $C(0)\ne0$, substitute $x=h^2v$ and divide by
$h^4$. The finite $v$-root is simple, giving a smooth
factor $x-h^2v(y,z)$. Its quotient is $xk^2+h^2U$
with $U$ a unit. The coordinate $X=x/U$ turns it
into $Xk^2+h^2$. Its finite birational normalization is
\[
                 u=h/k,\qquad X=-u^2,
                 \qquad u(y+z)=yz.
\]
The last equation is an $A_1$ quadric, so the
normalization has rational singularities.

\noindent\textbf{Step 3.} Suppose $C(0)=0$. Perturb $D$ by a small constant
and the linear part of $C$ generically, retaining
$C(0)=0$, and then replace $C$ by $C-\epsilon$.
On a fixed boundary the transverse tacnode coefficients
are units, so these changes retain boundary $A_3$.
Under $x=\epsilon^3X$, $y=\epsilon Y$,
$z=\epsilon Z$, the leading scaled equation is
\[
 \frac f{\epsilon^8}=F+O(\epsilon),\qquad
 F=X^2K^2+XH^2(L-1)+D_0H^4,
 \quad K=Y+Z,\ H=YZ.
\]
For $L=K$, $D_0=1$, put $W=Y-Z$.
At $W=0$ the $X$-critical coordinate is
$X=-K^2(K-1)/32$, and the remaining derivative is
$K^5(12K^2+7K-3)/512$.
At either root of $12K^2+7K-3$, the full critical
value and Hessian determinant are respectively
\[
 \frac{K^6(K+1)(3K-1)}{1024},\qquad
 -\frac{3K^{10}(K+1)(3K-1)(28K^2+14K-5)}{32768}.
\]
They are nonzero: their numerator resultants with the
critical polynomial are $8748$ and $4923269424$.
Nondegeneracy, nonzero value and avoidance of the axes
hold on a nonempty algebraic open set of coefficients.
An arbitrarily small generic perturbation is therefore
available. The implicit-function theorem lifts these
Morse points through the higher weighted terms, with
nonzero critical values of order $\epsilon^8$.
The boundary comparison embeds their square-$-2$
vanishing spheres in a nearby torus fibre, a
contradiction. Thus $C(0)=0$ is excluded.

\paragraph{The ideal $(g,xh,h^2)$, where $g=xk-yh$.}
Append $(a,b,d)$ to
\[
                         \begin{pmatrix}h&0\\-k&h\\y&-x\end{pmatrix}.
\]
Its determinant is $ag+bxh+dh^2$. The new adjugate
entries give the conditions
\[
       ax,\ ah,\ bx+dh,\ kd+by,\ dh-ay,\ ak+bh\in I.
\]
The exact colon identities are
\[
 I:h=(x,h),\qquad I:x=(xk,h,z^2),\qquad I:x^2=(k,h).
\]
The first two conditions give $a=hA$ modulo $I$;
then $dh-ay$ gives $d=Ay+xD+hE$.
Initially $bx+dh$ permits
$b=-Ak+hB+z^2B_2$ modulo $I$.
The condition $ak+bh$ forces $z^2B_2\in(x,h)$,
so $B_2\in(x,y)$. Its $x$-part lies in $I$, since
$xz^2=h^2-xh$ modulo $I$, and its $y$-part is
absorbed in $hB$. Thus $b=-Ak+hB$ modulo $I$.
The remaining condition is $y(D+B)+kE\in(x,h)$.

\noindent\textbf{Step 4.} Write $T=D+B$. Modulo $(x,h,k)$, the equality
$yT=0$ in $\mathbb C\{y\}/(y^2)$ forces $T(0)=0$.
Write $T=xW+yU+kV$. Since $k$ is a nonzerodivisor
modulo $(x,h)$, the general solution is
$E=-yU-yV+xE_1+hE_2$. Substitution in the determinant
gives $f\in(p)+I^2$, where
\[
                            p=yh^2(x-h).
\]
The exact identities
$xh^3=ygh^2-yp$ and $h^4=ygh^2-kp$ then put the
equation in the form
\[
 f=A g^2+B gxh+Cgh^2+D(xh)^2+\kappa p,
                       \qquad A(0)\ne0.              \tag{4.49b}
\]

\noindent\textbf{Step 5.} We first treat $\kappa(0)\ne0$.
Use the appended coefficients
$a=Ag+Bxh+Ch^2$, $b=Dxh+\kappa yh$,
$d=-\kappa yh$. The actual conductor identifies
$\operatorname{Hom}(I,\mathcal O_S)$ with the
normalization, generated by $1,n_1,n_2$, where the
adjugate fractions are
\[
          n_1=-a/h,\qquad n_2=-(ak+bh)/h^2.
\]
Since $A$ is a unit, the same actual normalization
algebra is generated by $1,v,w$, where
$v=(n_1+Bx+Ch)/A=-g/h$ and $w=n_2$.
Direct multiplication gives the exact relations
\begin{align*}
 Av^2-Bxv-Chv+Dx^2+\kappa xy-\kappa yh&=0,\\
 xw&=Ayv-Bxy-(\kappa+C)yh,\\
 vw&=\kappa h-Dxy,\\
 w^2+[(\kappa+B)y+Ck]w-A\kappa k+ADy^2
             -D(\kappa+C)yh-B\kappa h&=0.
\end{align*}
At the origin they give $v=w=0$, so there is one
normalization point. The last equation eliminates $k$
analytically, with $k$ in the maximal-ideal square.
Thus $h=y(k-y)$ has quadratic initial term $-y^2$.
Write $a=A_0$, $b=B_0$, $d=D_0$, $u=\kappa_0$.
The three quadratic equations become
\[
 av^2-bxv+dx^2+uxy=0,\qquad
 xw-ayv+bxy=0,\qquad vw+uy^2+dxy=0.
\]
Set $V=av-(b/2)x$, $\delta=ad-b^2/4$,
$Y=\delta x+auy$, $W_0=w+(b/2)y$, and
$W=auW_0+\delta V$. This invertible linear change
turns them into
\[
                       V^2+xY=0,\qquad
                       VW+Y^2=0,\qquad xW-YV=0.
\]
They define the irreducible rational normal cubic cone.
The full tangent-cone ideal contains this prime ideal
and has the same dimension two, so it equals it.
Blowing up the normalization point is smooth along its
exceptional rational normal cubic: the two independent
initial equations have independent derivatives along
that smooth projective curve in every blowup chart.
The exceptional curve is $\mathbb P^1$ with normal
bundle $\mathcal O(-3)$. The exact sequences of its
thickenings have successive terms $\mathcal O(3n)$,
all with $H^1=0$. Formal functions gives
$R^1\rho_*\mathcal O=0$. The singularity is rational.

\noindent\textbf{Step 6.} Finally suppose $\kappa(0)=0$. Perturb $D$ and the
linear part of $\kappa$ generically, and replace
$\kappa$ by $\kappa-\epsilon$. The fixed axes retain
boundary $A_3$: equation (4.49b) stays in the same
ring-condition module, and its boundary transverse
quartic coefficients remain units.
Scale $x=\epsilon^2X$, $y=\epsilon Y$, $z=\epsilon Z$.
Put $K=Y+Z$, $H=YZ$, $G=XK-YH$.
After division by $\epsilon^6$, the first terms are
\[
 A_0G^2+\epsilon(B_0GXH+C_0GH^2)
   +\epsilon^2\bigl[D_0(XH)^2+p(L-1)\bigr].
\]
Near a nonzero smooth point of $G=0$, eliminate the
normal Morse coordinate. On that cone the identities
$XH^3=-Yp$, $H^4=-Kp$ absorb the square-completion
corrections into the coefficients, leaving the leading
tangential function
\[
                    \Phi=D'(XH)^2+p(L'-1)|_{G=0}.
\]
Choose $D'=1$, $L'=-Z$ to exhibit a nonempty open
set of suitable coefficients. The parametrization
\[
 K=r,\quad Y=tr,\quad Z=(1-t)r,\quad
 X=t^2(1-t)r^2
\]
gives $\Phi=t^4(1-t)^4r^7[1+r(t^2-t+1)]$.
At $t=1/2$, $r=-7/6$, its value and Hessian are
\[
 -\frac{823543}{573308928},\qquad
 \operatorname{diag}\left(\frac{117649}{1990656},
                          \frac{15647317}{214990848}\right).
\]
The Hessian determinant is
$1840891197733/427972821516288$, nonzero.
The angular critical numerator factors as
$(2t-1)(3t^2-3t-4)$, so this point is simple and
avoids the axes and $K=0$. The generic coefficient
condition is algebraic and nonempty. After normal
elimination, divide the tangential critical equations
by $\epsilon^2$ and apply the implicit-function theorem.
The resulting full Morse points have nonzero critical
values of order $\epsilon^8$ and tend to the origin.
Boundary transport gives the forbidden square-$-2$
vanishing sphere in a nearby torus fibre. Thus
$\kappa(0)=0$ is excluded, completing both ideal forms.
\end{proof}

All adjugates, colon ideals, normalization relations,
linear coordinate identities, and both polar Hessian
calculations are executed in
\texttt{code/part4/a3\_missing\_coefficient.py}.
The saved output reports PASS.

\begin{proposition}[P4.50]\label{night:P4.50}\label{stmt:P4.50}
\StatementRecord{night:P4.50}{proved}{Historical events: GPT-6 Astra (OpenAI), ChatGPT Pro mode, effort Pro, S6\_06\_BLIND, 16--18 November 2026: targeted proof check. Also checked in this writing. Pre-continuation writing events: GPT-6 Astra (OpenAI), ChatGPT Pro mode, effort Pro, S6\_10\_PAPERS, 18 November 2026 05:10:56.410177+00:00: computation or formal certificate; targeted proof check. Continuation metadata events: GPT-6 Astra (OpenAI), ChatGPT Pro mode, effort Pro, S6\_10\_PAPERS continuation, 19 November 2026: used and rechecked in this run; status-vocabulary reconciliation; metadata reconciliation only; no new mathematical check. \emph{Pre-continuation recorded proof scope.} The night-run re-derivation was by GPT-6 Astra (OpenAI, ChatGPT Pro mode, effort Pro), 16--18 November 2026. Also checked in this writing by GPT-6 Astra (OpenAI), ChatGPT Pro mode, effort Pro, 18 November 2026. proved. Full final uniform-A3 positive-Euler exclusion, with every local budget alternative followed. The verified P4.42/P4.44 maps leave a contracted positive-genus section, two rational normalized conductor components, one triple point and exact total defect9. P4.48 excludes planar triples; P4R.16 excludes the tacnode-plus-line excess-one model; P4R.17 exhausts delta4 nonplanar models and proves each bound>9. The uniform branch estimate delta+3+m/2 excludes delta\ensuremath{\geq}5. At delta2 the exact ring is the three axes, which consumes the full budget, and P4.47/P4.49 make its normalization rational. Any additional singular critical-curve point contributes positively and is absent. At each smooth critical point the actual forms in P4R.14 give smooth or A1 normalization, and R6 excludes exceptional blocks away from the conductor. Rational normalization singularities cannot contract H: the ruled restriction H1(T,O)\ensuremath{\to}H1(H,O) is an isomorphism, whereas a contraction would kill that map, as quantified in P4.33. This gives the stated contradiction. Review provenance: New derivation after reading the supplied proof; not a blind review and not an independent human check}{\ref{night:P4.33}; \ref{night:P4.42}; \ref{night:P4.44}; \ref{night:P4.47}; \ref{night:P4.48}; \ref{night:P4.49}; \ref{night:P4R.14}; \ref{night:P4R.16}; \ref{night:P4R.17}; \ref{night:P4R.6}}

In the requested full-fibration classes, no support
with positive-genus ruled resolution has positive
Euler number and generic transverse type $A_3$
on every critical-curve component.
\end{proposition}
\begin{proof}
\noindent\textbf{Step 1.} P4.42 and P4.44 leave only a contracted horizontal
positive-genus section and a cycle at one three-branch
point of a conductor with two rational normalized
components. The total defect is nine by P4.33 and
(4.42a). P4.48 excludes planar triples; P4R.16
excludes the nonplanar excess-one tacnode-plus-line
case; the exact excess-two classification and conormal
calculation in P4R.17 exclude
every delta-four nonplanar triple. The branchwise bound
$\delta+3+m/2$ excludes all larger delta invariants,
where $m\ge3$ is total multiplicity.
The only remaining triple is the ordinary three-axis
point. P4.48 gives local defect at least nine, so it
uses the entire budget, and P4.47 and P4.49 give
rational normalization there.

\noindent\textbf{Step 2.} Every additional singular critical-curve point has
strictly positive defect and therefore is absent.
At a smooth point of the entire critical curve,
P4R.14 gives the product tacnode or its $A_1$
normalization degeneration; both have rational
normalization. Away from the conductor, P4R.6
excludes exceptional configurations. The normalization
therefore has rational singularities everywhere.
Its minimal resolution cannot contract the horizontal
positive-genus section. This contradiction proves
the assertion.
\end{proof}


\subsection{Local intersection forms and the quartic tangent obstruction}

\begin{proposition}[P4.51]\label{night:P4.51}\label{stmt:P4.51}
\StatementRecord{night:P4.51}{proved}{Historical events: GPT-6 Astra (OpenAI), ChatGPT Pro mode, effort Pro, S6\_06\_BLIND, 16--18 November 2026: targeted proof check. Also checked in this writing. Pre-continuation writing events: GPT-6 Astra (OpenAI), ChatGPT Pro mode, effort Pro, S6\_10\_PAPERS, 18 November 2026: targeted proof check. Continuation metadata events: GPT-6 Astra (OpenAI), ChatGPT Pro mode, effort Pro, S6\_10\_PAPERS continuation, 19 November 2026: used and rechecked in this run; status-vocabulary reconciliation; metadata reconciliation only; no new mathematical check. \emph{Pre-continuation recorded proof scope.} The night-run re-derivation was by GPT-6 Astra (OpenAI, ChatGPT Pro mode, effort Pro), 16--18 November 2026. Also checked in this writing by GPT-6 Astra (OpenAI), ChatGPT Pro mode, effort Pro, 18 November 2026. The local Milnor-cycle lift to the torus cover associated with the specialization kernel W was re-derived; dim W<=3 implies the pairing on Lambda\textasciicircum{}2 W is zero}{\ref{night:P4R.4}}

Let $S$ be an irreducible reduced support in the torus
fibration of Part~4.2 with $b_1(S)>0$. Every local
Milnor fibre of a reduced local equation of $S$ has
identically zero rational intersection form. This applies
in particular to the remaining positive ruled support,
whose fundamental group is $\mathbb Z$.
\end{proposition}
\begin{proof}
\noindent\textbf{Step 1.} Let $U$ be a sufficiently small full tube and $F_t$ a
nearby compact smooth torus fibre. P4R.4 identifies the
homomorphism
\[
 \pi_1(F_t)=\mathbb Z^4\longrightarrow
 \pi_1(U)=\pi_1(S)
\]
and gives a finite cyclic cokernel, of order equal to
the fibre multiplicity. Its map on rational first
homology is therefore onto. Write
\[
 W=\ker\bigl(H_1(F_t;\mathbb Q)\longrightarrow
                  H_1(S;\mathbb Q)\bigr).
\]
Then $\dim W=4-b_1(S)\le3$.

\noindent\textbf{Step 2.} Choose a coordinate ball about the point, with its
closed ball contained in $U$. The local full fibre
equation is $u g^m$, with $g$ reduced and $u$ a
holomorphic unit. We take a holomorphic $m$-th root of
$u$ on the ball and absorb it into $g$. The actual
nearby fibre in the ball is consequently the disjoint
union of the local Milnor fibres
$g=\zeta t^{1/m}$. Each such fibre $M$ is an oriented
codimension-zero submanifold with boundary of $F_t$.
Its inclusion into $U$ factors through the contractible
ball. Thus every loop of $M$ maps to zero in
$H_1(U;\mathbb Q)$, and the image subgroup
\[
 H=\operatorname{im}\bigl(\pi_1(M)\longrightarrow
                                 \mathbb Z^4\bigr)
\]
satisfies $H\otimes\mathbb Q\subset W$.

\noindent\textbf{Step 3.} The covering criterion lifts $M\to F_t$ to the
connected cover $F_H\to F_t$ corresponding to $H$.
This cover is a $K(H,1)$, with $H$ free abelian. The
map on second rational homology is therefore exactly
\[
 \bigwedge^2(H\otimes\mathbb Q)\longrightarrow
              \bigwedge^2 H_1(F_t;\mathbb Q).
\]
Every class from $H_2(M;\mathbb Q)$ consequently maps
into $\bigwedge^2W$. The torus intersection pairing is
exterior multiplication followed by the orientation
isomorphism $\bigwedge^4H_1(F_t;\mathbb Q)=\mathbb Q$.
It vanishes on $\bigwedge^2W$ because $\dim W\le3$.

\noindent\textbf{Step 4.} This ambient pairing is the local intersection
pairing on $M$. We represent finitely many cycles
inside its interior and make them transverse there;
the orientation-preserving inclusion into $F_t$
retains all intersection points and signs. Thus the
local pairing is zero. The symbolic exterior-square
matrix and its zero Gram matrix are recorded in
\texttt{code/part4/local\_intersection.py} and its
executed output.
\end{proof}

\begin{proposition}[P4.52]\label{night:P4.52}\label{stmt:P4.52}
\StatementRecord{night:P4.52}{proved}{Historical events: GPT-6 Astra (OpenAI), ChatGPT Pro mode, effort Pro, S6\_06\_BLIND, 16--18 November 2026: targeted proof check. Also checked in this writing. Pre-continuation writing events: GPT-6 Astra (OpenAI), ChatGPT Pro mode, effort Pro, S6\_10\_PAPERS, 18 November 2026: targeted proof check; used and rechecked in this run. Continuation metadata events: GPT-6 Astra (OpenAI), ChatGPT Pro mode, effort Pro, S6\_10\_PAPERS continuation, 19 November 2026: used and rechecked in this run; status-vocabulary reconciliation; metadata reconciliation only; no new mathematical check. \emph{Pre-continuation recorded proof scope.} The night-run re-derivation was by GPT-6 Astra (OpenAI, ChatGPT Pro mode, effort Pro), 16--18 November 2026. Also checked in this writing by GPT-6 Astra (OpenAI), ChatGPT Pro mode, effort Pro, 18 November 2026. The compact-cycle, conductor-pole, and multiplicity reduction were followed. Siersma 1991 section 7.1, printed page 462, gives (j2)\_*(T-I)=VAR'' j; sections 6.2-6.5, printed pages 460-461, preserve the required ordinary degree-two groups. The boundary Wang sequence yields the claimed rank bound; no degree-two variation isomorphism was assumed}{\ref{night:P4.20}; \ref{night:P4.39}; \ref{night:P4.51}; \ref{night:P4R.15}; \ref{night:P4R.6}}

In the remaining positive ruled support of Part~4.2,
the image of a contracted horizontal pole section has
ambient multiplicity exactly three. This restriction
imposes no generic-transverse-type hypothesis on
any component of the critical curve.
\end{proposition}
\begin{proof}
\noindent\textbf{Step 1.} We first give the comparison with a leading
homogeneous polynomial. Write the reduced equation
as $g=g_d+g_{d+1}+\cdots$, and suppose that $g_d$ is
reduced. Every finite collection of cycles in
$M_d=\{g_d=1\}$, together with its intersections,
transports into a local Milnor fibre of $g$.
Indeed, choose a compact neighbourhood of the cycle
representatives. Euler's identity shows that
$dg_d\ne0$ on $M_d$. The functions
\[
 G_\epsilon(x)=\epsilon^{-d}g(\epsilon x)
\]
converge with their derivatives to $g_d$ on that
neighbourhood. We solve
\[
 dG_\epsilon(V_\epsilon)
                     =-\partial_\epsilon G_\epsilon
\]
using the Hermitian gradient divided by its squared
norm. A cutoff supported in a slightly larger
compact neighbourhood, equal to one along the
trajectories under consideration, gives a flow for
small positive real $\epsilon$. It transports those
cycles to $G_\epsilon=1$. Holomorphic scaling by
$\epsilon$ then embeds them in $g=\epsilon^d$ inside
the fixed original ball, preserving orientations
and intersections. This compares smooth nonzero
level sets on compact neighbourhoods; no assertion
about equisingularity of the zero set is needed.
By P4.51 the pairing on $M_d$ is consequently zero.

\noindent\textbf{Step 2.} For a reduced homogeneous polynomial its global
fibre and a sufficiently large compact Milnor fibre
have the same homology and scalar monodromy. One
can see this directly from the squared radius on
the affine fibre. This is a proper semialgebraic
function with finitely many critical values; above
the largest, its gradient supplies a product end.
A large invariant ball therefore gives a compact
core. Rescaling identifies that core with a small
local Milnor fibre. The Milnor-fibre CW dimension
bound gives zero homology above degree two, and
reducedness gives connectedness.

\noindent\textbf{Step 3.} We need one exact variation identity. For such a
homogeneous polynomial, let $\mu_i$ be the
transverse plane-curve Milnor numbers at the
singular points of its reduced projective curve,
and put $M=\sum_i\mu_i$. Write $F$ for its compact
Milnor fibre and $\partial_2F$ for the boundary
part over the singular-link circles. Let
\[
 \begin{split}
 j&:H_2(F;\mathbb Q)\longrightarrow
                         H_2(F,\partial F;\mathbb Q),\\
 j_2&:H_2(F;\mathbb Q)\longrightarrow
                         H_2(F,\partial_2F;\mathbb Q).
 \end{split}
\]
Siersma's chain-level variation identity
\cite[\S3.6, pp.~451--452; \S7.1, p.~462]
{Siersma1991Variation} is
\[
                 j_2(T-I)=\operatorname{VAR}^{II}j.
                                                            \tag{4.52a}
\]
Here $T$ is the Milnor monodromy on $H_2(F;\mathbb Q)$.
The identity does not require the variation map to
be an isomorphism in dimension two. The dimension-two
Wang sequence is explicitly treated in that source,
\S6.3. Poincar\'e--Lefschetz duality identifies $j$
with the intersection form. Its vanishing makes
$\operatorname{im}(T-I)$ lie in $\ker j_2$.
The exact sequence of $(F,\partial_2F)$ identifies
that kernel with the image of $H_2(\partial_2F)$.
Each boundary piece is the mapping torus of a
transverse Milnor fibre. Its Wang sequence gives
\[
 H_2(\partial_2F;\mathbb Q)
       =\bigoplus_i\ker(A_i-I)\subset
                        \bigoplus_iH_1(F_i;\mathbb Q),
\]
where $A_i$ is the vertical monodromy. Thus zero
intersection form implies
\[
                 \operatorname{rank}(T-I)\le M.             \tag{4.52b}
\]
An empty $\partial_2F$ gives the same assertion with
$M=0$.

\noindent\textbf{Step 4.} Suppose now that the reduced projective quartic
$D=\{g_4=0\}$ has a normalization component of
positive genus. By degree, either $D$ is an
irreducible quartic with normalization genus
$g\ge1$, or it is a smooth elliptic cubic and a
line. In the first case
$\sum\delta_i=3-g\le2$, so
$M=\sum(2\delta_i-r_i+1)\le4$.
In the second case B\'ezout gives total intersection
length three and at least one intersection point.
Hence $\sum\delta_i=3$ and
$\sum(r_i-1)\ge1$, giving $M\le5$.

\noindent\textbf{Step 5.} The map $\{g_4=1\}\to\mathbb P^2\setminus D$ is
an unramified fourfold cover. The normalization
sequence and plane-curve genus formula give
\[
              \chi(F)=4\chi(\mathbb P^2\setminus D)
                      =4(7-M).
\]
The three nonidentity scalar monodromies act freely
on the compact invariant core, so their Lefschetz
numbers are zero. It follows by the character
formula for the cyclic group of order four that
\[
 H_0(F;\mathbb C)-H_1(F;\mathbb C)+H_2(F;\mathbb C)
       =(7-M)\,\mathbb C[\mathbb Z/4]
\]
as virtual representations. Every nontrivial
character therefore has multiplicity at least
$7-M$ in $H_2$. The finite-order monodromy is
semisimple, and hence
\[
             \operatorname{rank}(T-I)
                   \ge3(7-M)\ge6>M.                       \tag{4.52c}
\]
This contradicts (4.52b).

\noindent\textbf{Step 6.} We apply this obstruction through the actual pole
comparison. Blow up the ambient origin, normalize
the strict transform, and take its minimal
resolution $T_1$. As in P4.39, this resolution maps
to the original normalization and dominates its
minimal resolution by point blowups. The strict
transform of the positive-genus horizontal section
$H$ retains pole coefficient two. If $E$ is the
ambient exceptional divisor, put
$n=\operatorname{ord}_H(E)>0$, and let $a\ge0$
be the coefficient of $H$ in the total conductor
divisor relative to the strict-transform dualizing
sheaf. Adjunction gives exactly
\[
                          2=(d-2)n+a.                     \tag{4.52d}
\]
Effectiveness follows from the exceptional
intersection matrix argument of P4R.6. An
exceptional positive-genus curve over the new
normalization has strictly positive coefficient,
by its adjunction equation. Thus $d\ge5$ is
impossible. If $d=4$, then $n=1,a=0$; the curve
$H$ is neither exceptional over the new
normalization nor over its conductor. The
normalization is an isomorphism at its generic
point. Its image is a multiplicity-one component
of the quartic exceptional curve, birational to
$H$. It has positive genus. Such a component has
degree three or four, leaving at most one linear
factor; the entire quartic polynomial is reduced.
The preceding contradiction excludes this case.

\noindent\textbf{Step 7.} For clarity, an isolated critical point cannot
hide a contracted $H$. A generic small linear
perturbation of an isolated hypersurface critical
point has a positive finite number of Morse
critical points, the local degree of its gradient.
A small constant shift makes a chosen critical
value nonzero. On a fixed sufficiently small ball
the original zero-level boundary is smooth and
transverse; compactness retains the boundary
submersion over a uniform punctured value disc.
The finite-critical-value and path-transport
argument of P4.20 then carries the Morse
vanishing sphere, of square $-2$, into the
original local Milnor fibre. P4.51 excludes it.
This argument does not assume smoothness of the
normalization.

\noindent\textbf{Step 8.} At a nonisolated point with $d=2$, the Hessian is
nonzero. P4R.15 gives the exact product equation
$X^2+Y^m$. For even $m$ its two analytic branches
are smooth. For odd $m$ its normalization is the
explicit smooth parametrization
$(Y,X,Z)=(u^2,\mathrm i u^m,Z)$.
Neither case has an exceptional curve in its
minimal normalization resolution. A point with
$d=1$ is smooth as well. Therefore the image of
$H$ has multiplicity exactly three.

\noindent\textbf{Step 9.} The cyclic character calculation and all integer
bounds are independently recomputed in
\texttt{code/part4/quartic\_tangent.py}; no numerical
approximation enters the argument.
\end{proof}


\subsection{The cubic double line and its weighted leading fibre}

\begin{proposition}[P4.53]\label{night:P4.53}\label{stmt:P4.53}
\StatementRecord{night:P4.53}{proved}{Historical events: GPT-6 Astra (OpenAI), ChatGPT Pro mode, effort Pro, S6\_06\_BLIND, 16--18 November 2026: targeted proof check. Also checked in this writing. Pre-continuation writing events: GPT-6 Astra (OpenAI), ChatGPT Pro mode, effort Pro, S6\_10\_PAPERS, 18 November 2026: targeted proof check. Continuation metadata events: GPT-6 Astra (OpenAI), ChatGPT Pro mode, effort Pro, S6\_10\_PAPERS continuation, 19 November 2026: used and rechecked in this run; status-vocabulary reconciliation; metadata reconciliation only; no new mathematical check. \emph{Pre-continuation recorded proof scope.} The night-run re-derivation was by GPT-6 Astra (OpenAI, ChatGPT Pro mode, effort Pro), 16--18 November 2026. Also checked in this writing by GPT-6 Astra (OpenAI), ChatGPT Pro mode, effort Pro, 18 November 2026. The pole identity, weighted degree-five discriminant, the connected double-cover branch bound, both coprime and common-factor cases, and the variation-rank contradiction were re-derived. The exceptional horizontal case was checked by adjunction and the minimal-resolution intersection inequality}{\ref{night:P4.51}; \ref{night:P4.52}}

In the remaining positive ruled support, blow up the
multiplicity-three image of the contracted horizontal
section $H$, normalize the strict transform $S_1$, and
take its minimal resolution $T_1$. The strict transform
of $H$ is exceptional over a point of $S_1$. The entire
germ of $S_1$ there is normal and has an isolated simple
elliptic singularity. In particular $H$ cannot dominate
a curve in the ambient exceptional plane. No generic
transverse type is imposed on the original critical curve.
\end{proposition}
\begin{proof}
\noindent\textbf{Step 1.} Use the comparison maps of P4.52, and let $D_1$ be the
effective total conductor relative to $\omega_{S_1}$.
Write $n=\operatorname{ord}_H(E)>0$ and
$a=\operatorname{coeff}_H(D_1)$. The pole identity is
$2=n+a$, so $(n,a)=(2,0)$ or $(1,1)$.

\noindent\textbf{Step 2.} Suppose $H$ dominates a reduced component $L$ of the
projective cubic tangent curve. At its generic point
a transverse Weierstrass presentation has degree
$m$, the multiplicity of that component in the cubic.
Factoring over the completed valuation field of the
exceptional coordinate $r$ gives the exact degree
identity
\[
 \sum_{H'}\operatorname{ord}_{H'}(r)
                 [\mathbb C(H'):\mathbb C(L)]=m.             \tag{4.53a}
\]
The factors group the roots by valuation branches;
the two factors in each summand are the ramification
index and residue-field degree. Residue extensions
are separable in characteristic zero.

\noindent\textbf{Step 3.} A positive-genus component of a cubic is a smooth
elliptic cubic of multiplicity one. The strict
transform is smooth along it, with $n=1,a=0$,
contrary to $2=n+a$. Every other reduced component
is rational. The finite map from $H$ then has degree
$h\ge2$. Equation (4.53a) excludes $n=2$, since
$m\le3$. Thus $n=a=1$.

\noindent\textbf{Step 4.} At a generic point of $L$, the $h$ points on $H$
give distinct branches of the transverse plane germ,
each with conductor exponent one. The branch formula
\[
             c_i=2\delta_i+\sum_{j\ne i}I_{ij}
\]
forces $h=2$ and exactly two smooth branches meeting
once, with no additional branch. Each has $r$-order
one, so their total Weierstrass degree is two. Hence
$m=2$, and the tangent polynomial is $y^2\ell$,
where $y=0$ is the double line and $\ell$ a different
linear form.

\noindent\textbf{Step 5.} Choose complementary coordinates $x,z$. The
restriction $\ell(x,z)$ is nonzero; any $y$ term in
$\ell$ has higher weight for the weights used below.
In a generic blowup chart along the double line,
write $v$ for the transverse coordinate. The
strict-transform equation begins
\[
                 \ell v^2+r(f_4|_L)+\cdots.
\]
A nonzero $f_4|_L$ would give ramification order two
and residue degree one. Therefore $f_4$ is divisible
by $y$. The two branches meet once, so after
$v=ru$ their distinct leading slopes solve
\[
                         \ell u^2+G_3u+G_5=0,
\]
where $G_3,G_5$ are forms of degrees three and five
in $x,z$. The discriminant
$\Delta=G_3^2-4\ell G_5$ is nonzero and nonsquare
in $\mathbb C(\mathbb P^1)$. The actual residue-field
extension of $H\to L$ is
$\mathbb C(\mathbb P^1)(\sqrt\Delta)$.

\noindent\textbf{Step 6.} Consequently the original equation has weighted
leading polynomial
\[
          P_5=\ell(x,z)y^2+G_3(x,z)y+G_5(x,z),               \tag{4.53b}
\]
of degree five for weights $(1,2,1)$ on $(x,y,z)$,
with no smaller-weight terms. This identifies the
normalized projective component birationally with
$H$ through its function field.

\noindent\textbf{Step 7.} Let $D=\{P_5=0\}\subset\mathbb P(1,2,1)$.
The quadratic is irreducible over $\mathbb C(x/z)$,
since the degree-two residue cover $H$ is connected.
The only possible further factor is a common
vertical factor of $\ell,G_3,G_5$. Since $\ell$
is linear, either $P_5$ is irreducible or
$P_5=\ell Q_4$, with $Q_4$ monic quadratic in $y$.
Both alternatives are reduced.

\noindent\textbf{Step 8.} In the irreducible case, let $k_i$ be the
multiplicities of the distinct zeros of the
homogeneous degree-six discriminant. The normalized
double cover has genus
$(\#\{i:k_i\text{ odd}\}-2)/2$. Positive genus
therefore requires at least four odd multiplicities,
and gives
\[
              \sum_i(k_i-1)=6-\#\{i\}\le2.                  \tag{4.53c}
\]
Away from $\ell=0$ and the weighted vertex,
completing the quadratic identifies the projective
curve germ with $w^2-s^{k_i}$, up to a unit and
holomorphic coordinates. Its transverse Milnor
number is $k_i-1$. At $\ell=0$, if $G_3\ne0$ the
finite root is simple and the other point is at
the weighted vertex. If $G_3=0$, absence of a
common factor gives $G_5\ne0$, and $\Delta$ has
a simple zero. There is no additional singularity.
In the chart $y=1$ the vertex is the quotient of
$\mathbb C^2$ by $(x,z)\mapsto(-x,-z)$; the
lifted equation $\ell+G_3+G_5$ has nonzero
differential. The original polynomial also has
nonzero derivative in the $\ell$ direction on
the $y$-axis, so that axis is not critical.
The total transverse Milnor-number sum is thus
$M\le2$.

\noindent\textbf{Step 9.} In the reducible case write
$G_3=\ell A_2$, $G_5=\ell A_4$. Then
$Q_4=y^2+A_2y+A_4$ has discriminant of degree four.
A connected positive-genus double cover requires
four distinct simple roots. Its normalization is
elliptic and $Q_4$ is smooth on the weighted plane;
it avoids the weighted vertex. The rational
component $\ell=0$ meets it in total length two,
read directly from the restricted monic quadratic.
Either there are two transverse intersections,
with total Milnor number two, or one tangency of
order two, with Milnor number three. There are no
other critical rays. Hence in all cases
\[
                              M\le3.                       \tag{4.53d}
\]

\noindent\textbf{Step 10.} The weighted plane has Euler number three: its
strata $x\ne0$, $x=0,z\ne0$, and $x=z=0$ are
$\mathbb C^2$, $\mathbb C$, and a point.
In the irreducible case $e(D)\le e(H)\le0$.
In the reducible case
$e(D)=0+2-\#(Q_4\cap\{\ell=0\})\le1$.
Therefore
\[
               \chi\bigl(\mathbb P(1,2,1)\setminus D\bigr)
                         \ge2.                            \tag{4.53e}
\]

\noindent\textbf{Step 11.} The weighted scalar action of a primitive fifth
root, with weights $(1,2,1)$, acts freely on
$F=\{P_5=1\}$. Its quotient is
$\mathbb P(1,2,1)\setminus D$: the only possible
nontrivial stabilizer of the weighted
$\mathbb C^*$ action is at the weighted vertex,
which belongs to $D$. This is an unramified
fivefold cover. The compact-core and character
calculation of P4.52 gives
\[
 \operatorname{rank}(T-I)
 \ge4\chi\bigl(\mathbb P(1,2,1)\setminus D\bigr)
 \ge8.
\]
The variation identity (4.52a) applies to every
one-dimensional hypersurface critical locus. Its
boundary mapping-torus estimate would give
$\operatorname{rank}(T-I)\le M\le3$ if the
intersection form were zero.

\noindent\textbf{Step 12.} It is zero here by the actual anisotropic scaling
comparison. Set
\[
 D_\epsilon(x,y,z)=(\epsilon x,\epsilon^2y,\epsilon z).
\]
Then $\epsilon^{-5}f(D_\epsilon(x,y,z))$ converges
with derivatives to $P_5$. The compact-cycle flow
of P4.52 transports its pairings into the actual
local Milnor fibre, whose pairing vanishes by
P4.51. This contradiction excludes every dominating
$H$.

\noindent\textbf{Step 13.} Thus $H$ is exceptional over the new normalization.
Its coefficient $a$ is positive, so $n=a=1$.
Adjunction gives
\[
                    0=2g(H)-2+(D_1-H)\cdot H.
\]
It follows that $H$ is elliptic and meets no other
positive pole component. A zero-coefficient
exceptional curve cannot meet a positive one:
its own intersection equation has nonnegative
canonical term on a minimal resolution and
nonnegative remaining terms, with a strictly
positive term at such a meeting. The connected
exceptional fibre is therefore $H$ alone and is
disjoint from the strict conductor. The
normalization is an isomorphism at its image;
another analytic hypersurface branch would
intersect this branch in a curve and create
conductor meeting $H$. Hence the entire
strict-transform germ is normal, with the single
smooth elliptic exceptional curve $H$. This is
precisely an isolated simple elliptic singularity.

\noindent\textbf{Step 14.} The discriminant multiplicity partitions and the
exact rank bounds are recomputed in
\texttt{code/part4/cubic\_weighted.py} and its
executed output.
\end{proof}


\subsection{Exclusion of the contracted horizontal section}

\begin{proposition}[P4.54]\label{night:P4.54}\label{stmt:P4.54}
\StatementRecord{night:P4.54}{proved}{Historical events: GPT-6 Astra (OpenAI), ChatGPT Pro mode, effort Pro, S6\_06\_BLIND, 16--18 November 2026: targeted proof check. Also checked in this writing. Pre-continuation writing events: GPT-6 Astra (OpenAI), ChatGPT Pro mode, effort Pro, S6\_10\_PAPERS, 18 November 2026: targeted proof check; used and rechecked in this run. Continuation metadata events: GPT-6 Astra (OpenAI), ChatGPT Pro mode, effort Pro, S6\_10\_PAPERS continuation, 19 November 2026: used and rechecked in this run; status-vocabulary reconciliation; metadata reconciliation only; no new mathematical check. \emph{Pre-continuation recorded proof scope.} The night-run re-derivation was by GPT-6 Astra (OpenAI, ChatGPT Pro mode, effort Pro), 16--18 November 2026. Also checked in this writing by GPT-6 Astra (OpenAI), ChatGPT Pro mode, effort Pro, 18 November 2026. The cofinal ideal filtrations, exact associated graded section ring, analytic coordinate changes preserving r, original-coordinate blowdown map, and all three weighted Milnor-fibre contradictions were re-derived. Grauert 1960 section 6 Hauptsatze I and IIa, printed page 59, was opened. The prime and order-four quotient-curve arguments were checked separately}{\ref{night:P4.51}; \ref{night:P4.52}; \ref{night:P4.53}; \ref{night:P4R.18}}

In the remaining positive ruled support of Part~4.2,
the horizontal pole section is not contracted by the
normalization resolution. It consequently maps with
degree two or three onto a rational normalized conductor
component whose generic transverse germ is $D_4$.
No transverse-type restriction is imposed on the other
conductor component.
\end{proposition}
\begin{proof}
\noindent\textbf{Step 1.} By P4.53, a contracted section would be the sole smooth
elliptic exceptional curve $H$ of an isolated normal
hypersurface singularity on the first ambient strict
transform. Its self-intersection is $-d<0$, and the
actual exceptional coordinate $r$ has order one along
$H$. Write $\pi:\widetilde V\to(V,0)$ for this local
resolution. We establish the following three coordinate
weights and leading degrees, retaining the actual $r$:
\[
 (1,1,1;3),\qquad(1,1,2;4),\qquad(1,2,3;6).
                                                        \tag{4.54w.1}
\]

\paragraph{1. Analytic formal functions and the exact associated graded ring}

\noindent\textbf{Step 2.} Put \(L=\mathcal O_H(-H)\), of degree \(d\), and
\[
 I_m=(\pi_*\mathcal O_{\widetilde V}(-mH))_0
     =\{a\in\mathcal O_{V,0}:\operatorname{ord}_H(a)\ge m\}.
\]
For \(m\ge1\), the restrictions of
\(\mathcal O(-mH)\) to successive infinitesimal neighbourhoods
of \(H\) have graded factors \(L^{m+j}\), \(j\ge0\).
Each has zero first cohomology, by Serre duality on the elliptic
curve and its strictly positive degree. Their extensions also
have zero first cohomology.

\noindent\textbf{Step 3.} The analytic formal-functions theorem is Grauert1960,
Section6, Hauptsatz IIa, printed p.59; Hauptsatz I on the same
page supplies coherence \cite[\S6, Hauptsätze I and IIa]{Grauert60}. The maximal-ideal pullback and the
ideal of \(H\) have the same radical. On a finite cover of
the compact exceptional curve, Noetherianity gives one uniform
power of the latter contained in the former. Their powers are
therefore cofinal. Applying the theorem with
\(\mathcal O(-mH)\) gives a zero completion for
\(R^1\pi_*\mathcal O(-mH)\). This coherent module is
supported at the origin, so it is zero. Consequently
\[
 I_m/I_{m+1}\simeq H^0(H,L^m)\qquad(m\ge1),             \tag{4.54w.2}
\]
through the actual restriction maps in
\[
 0\longrightarrow\mathcal O(-(m+1)H)
 \longrightarrow\mathcal O(-mH)
 \longrightarrow L^m\longrightarrow0.
\]
These maps respect multiplication. Since normality gives
\(\pi_*\mathcal O_{\widetilde V}=\mathcal O_V\), and a
function vanishes at the contracted point exactly when its
restriction to \(H\) vanishes, \(I_1=\mathfrak m\).
The degree-zero piece is \(\mathbb C\). Thus
\[
 \operatorname{gr}_{I}\mathcal O_{V,0}
   =\bigoplus_{m\ge0}H^0(H,L^m)                         \tag{4.54w.3}
\]
as graded algebras, not only as vector spaces.

\noindent\textbf{Step 4.} The same formal-functions theorem in degree zero identifies
\[
 \widehat{\mathcal O}_{V,0}
      \simeq\varprojlim_m H^0(mH,\mathcal O_{mH})
      \simeq\varprojlim_m\mathcal O_{V,0}/I_m.          \tag{4.54w.4}
\]
The second equality uses the just-proved first-direct-image
vanishing. This supplies the exact completeness and separation
needed in the leading-equation argument below; no unproved equivalence between a
divisorial and an ordinary degree filtration is being used.

\noindent\textbf{Step 5.} Because \(\mathfrak m^2\subset I_2\), (4.54w.2) gives a
surjection
\(\mathfrak m/\mathfrak m^2\to H^0(H,L)\).
The source has dimension three and the target dimension \(d\).
Hence \(d\le3\).

\paragraph{2. The three section rings and their actual elliptic curves}

\noindent\textbf{Step 6.} For \(d=3\), Riemann--Roch shows that \(L\) separates
points and tangent vectors: subtracting any length-two
subscheme leaves a degree-one line bundle with one section.
The complete linear series therefore embeds \(H\) as a smooth
plane cubic. The restriction sequence on \(\mathbb P^2\),
using \(H^1(\mathbb P^2,\mathcal O(k))=0\), gives its whole
section ring as
\[
 \mathbb C[x,y,z]/(F_3),\qquad \deg(x,y,z)=(1,1,1),
\]
where \(F_3=0\) is that smooth cubic.

\noindent\textbf{Step 7.} For \(d=2\), the complete pencil is base-point free and gives
a degree-two map \(p:H\to\mathbb P^1\), with
\(L=p^*\mathcal O(1)\). Trace splits
\(p_*\mathcal O_H=\mathcal O\oplus B^{-1}\).
Its Euler characteristic is zero, so \(\deg B=2\).
The algebra structure is therefore a double-cover equation
\(z^2=F_4(x,y)\), where the binary quartic has four distinct
zeros because the double cover is smooth. Projection formula
gives the entire section ring with generator weights
\((1,1,2)\) and that single degree-four relation.

\noindent\textbf{Step 8.} For \(d=1\), write \(L=\mathcal O_H(p)\).
Riemann--Roch gives a degree-one section \(x\), one new
degree-two section \(y\), and one new degree-three section
\(z\). Dividing by powers of \(x\) gives functions of pole
orders two and three at \(p\). Their Weierstrass relation
has weighted degree six. Equivalently, after graded changes
only for describing this section ring, it is
\[
 z^2=y^3+A x^4y+B x^6,
 \qquad 4A^3+27B^2\ne0,
 \qquad\deg(x,y,z)=(1,2,3).
\]
The pole-order count shows that the monomials with at most one
factor \(z\), reduced by this relation, give every space
\(H^0(H,L^m)\). Thus this is the whole section ring. Its
weighted projective curve is the actual smooth \(H\).
The coordinate choice below does not require completing the
square or applying any of these optional graded simplifications.

\paragraph{3. Retaining the given coordinate using only permitted changes}

\noindent\textbf{Step 9.} Start with analytic ambient coordinates \((r,u,v)\), retaining
the given \(r\). Their images form a basis of
\(\mathfrak m/\mathfrak m^2\). The map to
\(I_1/I_2=H^0(L)\) has the rank just computed.

\noindent\textbf{Step 10.} When \(d=3\), it is an isomorphism, so any linear complement
to \(r\) supplies the other two valuation-one coordinates.

\noindent\textbf{Step 11.} When \(d=2\), choose a linear coordinate \(s\) whose
leading section complements that of \(r\), and choose the
remaining independent linear coordinate \(v_0\) in the
kernel of the degree-one restriction. Thus
\(\operatorname{ord}_H(v_0)\ge2\). Its degree-two leading
section must be independent of \(r^2,rs,s^2\). Otherwise
all degree-two initial sections of analytic functions would
lie in that three-dimensional span: a Taylor expansion shows
that every term of ordinary order at least three has valuation
at least three, and the linear terms of valuation two are
accounted for by \(v_0\). This contradicts
\(h^0(L^2)=4\). Therefore \(v_0\) has valuation exactly
two and is the required indecomposable degree-two generator.
The coordinates \((r,s,v_0)\) were obtained by an invertible
linear change.

\noindent\textbf{Step 12.} When \(d=1\), first choose two independent linear coordinates
\(u_0,v_0\) in the kernel of the degree-one restriction.
At least one linear combination, denoted \(s\), has a
degree-two initial section independent of \(r^2\), since
\(h^0(L^2)=2\). Choose the remaining independent linear
coordinate \(v_1\), and subtract a scalar multiple of \(s\)
so that its degree-two initial term is \(\alpha r^2\).
Set
\[
                         v_2=v_1-\alpha r^2.
\]
This is a triangular analytic coordinate change, and
\(\operatorname{ord}_H(v_2)\ge3\). Its degree-three
initial section must be independent of \(r^3,rs\).
Indeed, if it were absent or in that span, any Taylor series
of valuation three would have degree-three initial section
in that same span: the possible terms are precisely the
linear \(v_2\) term and the weighted-degree-three monomials
\(r^3,rs\); its valuation-one and valuation-two terms
must first vanish. This contradicts \(h^0(L^3)=3\).
Thus \(v_2\) has valuation exactly three and supplies the
indecomposable degree-three generator. The three ambient
coordinates remain analytically invertible, because only
an invertible linear change and the displayed quadratic
translation were performed.

\noindent\textbf{Step 13.} This argument proves directly that the chosen analytic
coordinates have the prescribed orders and generate the
graded ring. It avoids a circular appeal to filtered
Nakayama or to a coordinate normalization that might move
the fixed \(r\).

\paragraph{4. The actual weighted leading equation}

\noindent\textbf{Step 14.} Write the defining analytic equation in the coordinates just
constructed, with the weights in (4.54w.1). Its weighted initial
form is the single defining relation of the graded section
ring, up to a nonzero scalar. Here is an explicit proof of that
last assertion.

\noindent\textbf{Step 15.} Let \(F_D\) be the graded relation from the section-ring calculation above, expressed
in the actual chosen initial generators. The analytic function
\(F_D(r,s,v)\) has valuation at least \(D+1\).
At each successive valuation, (4.54w.3) writes its initial
section as a weighted homogeneous polynomial in the three
generators. Subtract that polynomial. This produces a formal
series \(G=F_D+F_{D+1}+\cdots\) whose image in every
\(\mathcal O_{V,0}/I_m\) is zero. Since all generator
weights are at most three, the correction polynomials tend
to zero in ordinary power-series degree; their sum is a
well-defined formal power series. By (4.54w.4), its image in
\(\widehat{\mathcal O}_{V,0}\) is zero.

\noindent\textbf{Step 16.} The completed hypersurface ideal is principal, generated by
the actual analytic equation \(f\). Thus \(G=h f\) for
a formal series \(h\). Taking weighted initial forms gives
\(F_D=\operatorname{in}(h)\operatorname{in}(f)\).
The polynomial \(F_D\) is irreducible because its weighted
projective curve is the smooth connected elliptic curve
\(H\). Also \(f\) is nonconstant and vanishes at the
origin. Hence \(\operatorname{in}(h)\) is a nonzero
constant, and \(\operatorname{in}(f)\) is a scalar multiple
of \(F_D\). This proves (4.54w.1) for the actual defining
equation, with all higher terms retained.

\paragraph{5. The exact descent to the original blowdown coordinates}

\noindent\textbf{Step 17.} In the chart with original coordinates
\(x=r u\), \(y=r v\), every change above has the form
\[
 s=a u+b v+c r,\qquad
 v'=a'u+b'v+c'r-\alpha r^2,
 \qquad ab'-a'b\ne0.
\]
It descends to the actual holomorphic coordinate map
\[
 \begin{aligned}
 x'&=r s=a x+b y+c r^2,\\
 y'&=r v'=a'x+b'y+c'r^2-\alpha r^3,\\
 r'&=r.
 \end{aligned}                                        \tag{4.54w.5}
\]
Its Jacobian determinant at the origin is \(ab'-a'b\),
so the inverse-function theorem supplies an analytic inverse.
This is the required morphism preserving the original
parameter and blowdown structure.

\paragraph{Leading polynomials in the original coordinates.}
Apply this calculation to the simple elliptic point supplied
by P4.53. The original equation is $r^3f_1$, and its original
coordinates are $(x,y,z)=(rs,rv,r)$. The three pairs of
original weights and degrees are
\[
 (2,2,1;6),\qquad (2,3,1;7),\qquad (3,4,1;9).
\]
For a strict-transform monomial $r^as^bv^c$, multiplication
by $r^3$ gives $x^by^cz^{3+a-b-c}$. Its original weighted
degree is three plus its strict-transform weighted degree.
Analyticity of the original equation makes every
negative-power coefficient zero. Thus the actual weighted
leading term is a polynomial of the indicated degree.
On $z\ne0$, the ratios
\[
     \left(x/z^{1+w_s},\ y/z^{1+w_v}\right)
\]
identify its projective component with the strict-transform
elliptic relation on $r\ne0$. This is a birational map
of the actual residue-field curves; their smooth
normalizations are the same elliptic curve $H$.

\noindent\textbf{Step 18.} Anisotropic scaling with these weights and the compact-cycle
flow in P4.52 transports every pairing of the leading
Milnor fibre into the actual local torus fibre. By P4.51,
all these pairings must vanish. We now contradict this
in each of the three cases.

\paragraph{Degree six.}
The leading polynomial is exactly
\[
                    P_6(x,y,z)=G_3(x,y,z^2),
\]
where $G_3(X,Y,Z)$ defines a smooth projective cubic.
The map $(x,y,z)\mapsto(x,y,z^2)$ gives a finite double
cover from $F=\{P_6=1\}$ to $\{G_3=1\}$, branched over
\[
                       B=\{G_3(x,y,0)=1\}.
\]
The homogeneous cubic fibre has Euler number nine: it
is a threefold unramified cover of the complement of
an elliptic cubic in $\mathbb P^2$, whose Euler number
is three. Let $r\in\{1,2,3\}$ be the number of distinct
zeros of the nonzero binary cubic $G_3(x,y,0)$ on
$\mathbb P^1$. The curve $B$ is a threefold cover of
that punctured projective line. Hence
\[
          \chi(B)=3(2-r),\qquad \chi(F)=18-\chi(B)=12+3r.
\]

\noindent\textbf{Step 19.} Weighted monodromy has order six. Its only
nonidentity element with a fixed set is its cube,
which negates $z$ and fixes $x,y$; its fixed set is
$B$. The other four nonidentity elements have no
fixed point on $F$. We choose an invariant
semialgebraic compact core large enough for $F$
and $B$. The finite-order Lefschetz character of
$H_0-H_1+H_2$ is therefore $\chi(F)$ at the identity,
$\chi(B)$ at the cube, and zero elsewhere. Its
trivial-character multiplicity is
\[
                       \frac{\chi(F)+\chi(B)}6=3.
\]
Since $H_0$ has only the trivial character, the
nontrivial-character dimension in $H_2$ is at least
$\chi(F)-3=9+3r\ge12$.

\noindent\textbf{Step 20.} Every nonzero critical orbit of $P_6$ lies on $z=0$
at a tangency of the cubic with $Z=0$. If the
intersection orders are $e_i$, then $\sum e_i=3$.
At a tangency, smoothness of $G_3$ makes its
$Z$-derivative a unit, and the exact transverse
germ is $u^{e_i}+z^2$, after holomorphic coordinates
and a unit. Consequently
\[
                       M=\sum(e_i-1)=3-r\le2.
\]
This also shows reducedness of $P_6$: a repeated
factor would give a two-dimensional critical set,
whereas the derivative calculation gives only
these critical orbits. The variation inequality
(4.52b) contradicts the lower bound twelve for
the nontrivial monodromy rank.

\paragraph{Weighted curve cohomology.}
We record the two curve facts needed for the last
cases, with their quotient comparisons. For
$\mathbb P(a,b,c)=\operatorname{Proj}\mathbb C[u,v,w]$
and a reduced equation of weighted degree $e$,
the exact graded sequence sheafifies to
\[
 0\longrightarrow\mathcal O_{\mathbb P}(-e)
 \longrightarrow\mathcal O_{\mathbb P}
 \longrightarrow\mathcal O_D\longrightarrow0.
\]
The three standard affine charts give its \v Cech
complex. The middle cohomology of every
$\mathcal O_{\mathbb P}(k)$ is zero. The top
cohomology of $\mathcal O_{\mathbb P}(-e)$ has
basis the Laurent monomials $u^{-i}v^{-j}w^{-k}$,
$i,j,k\ge1$, of weighted degree $-e$. Thus
\[
 h^1(\mathcal O_D)=
 \#\{\hbox{nonnegative weighted monomials of degree }
                         e-(a+b+c)\},
 \qquad h^0(\mathcal O_D)=1.                         \tag{4.54a}
\]
In particular the arithmetic genus is zero below
degree $a+b+c$, and is one in degrees $a+b+c$
and $a+b+c+1$ when there is exactly one variable
of weight one.

\noindent\textbf{Step 21.} At a cyclic quotient vertex of prime order two
or three, with no invariant linear ambient
coordinate, a smooth coarse curve has a smooth
plane pullback. To prove this, its normalized
branches form one orbit. A single singular
branch of multiplicity $m\ge2$ has a faithful
normalization action of order $p$; every invariant
ambient function has order at least $2m>p$,
whereas a regular quotient requires a function
of normalization order $p$. For $p$ transitively
permuted branches the stabilizer is trivial,
and every invariant function has order at least
two, whereas a regular quotient requires order
one. Both alternatives are impossible. The
remaining single branch is smooth, and its
reduced equation has the character of a nonzero
linear ambient coordinate. These arguments
use the actual analytic invariant rings:
averaging makes the invariants functor exact,
so every quotient-curve function lifts to an
invariant ambient function.

\paragraph{Degree seven.}
Consider weights $(1,2,3)$ and degree seven.
Formula (4.54a) gives arithmetic genus zero for
curves of degree at most five and arithmetic
genus one in degrees six and seven. Thus a
positive-genus component has degree six or
seven, genus one, and multiplicity one. There
is at most one further degree-one factor;
the polynomial is reduced.

\noindent\textbf{Step 22.} In degree seven the whole curve $D$ is
irreducible with arithmetic and normalization
genus one. Its total delta invariant is zero,
so it is a smooth coarse elliptic curve.
At the quotient vertices of orders two and
three, the preceding argument makes its plane
pullback smooth. Elsewhere the ambient plane
is smooth. Thus the affine cone has no
nonzero critical orbit and $M=0$.

\noindent\textbf{Step 23.} In degree six write $D=E\cup\{z=0\}$, with
$z$ the weight-one coordinate. The curve $E$
is smooth coarse elliptic. It cannot contain
either quotient vertex: its degree-six
equation has trivial character there, while
both local-coordinate characters are
nontrivial, so a smooth lift would require
an impossible invariant linear equation.
The genus formula for $D$ gives total delta
one. Its two smooth components therefore
meet once transversely in the smooth part
of the ambient weighted plane. There is
one ordinary-double critical orbit and
$M=1$.

\noindent\textbf{Step 24.} The weighted plane has Euler number three,
from its affine-cell stratification. In the
two cases $e(D)=0$ or one, so its complement
has Euler number three or two. Both quotient
vertices belong to $D$, since degree seven
is divisible by neither two nor three.
The weighted $\mathbb C^*$ action is free
off $D$, so $\{P_7=1\}$ is an unramified
sevenfold cover of that complement. The
free order-seven monodromy gives nontrivial
$H_2$ rank at least eighteen or twelve,
respectively. This contradicts (4.52b),
whose upper bound is zero or one.

\paragraph{Degree nine.}
For weights $(1,3,4)$, (4.54a) gives
arithmetic genus zero below degree eight
and one in degrees eight and nine.
A positive-genus component must have one
of these two degrees, and multiplicity one.

\noindent\textbf{Step 25.} An elliptic degree-eight component is
impossible. It necessarily contains the
weight-three vertex, since eight is not
divisible by three. The local plane
characters are $1,1$ modulo three, whereas
its equation has character two. The
prime-order quotient argument would give
a smooth lift, whose equation must have
character one. This is a contradiction.

\noindent\textbf{Step 26.} Thus the entire degree-nine curve $D$ is
irreducible and smooth coarse elliptic.
It avoids the weight-three vertex: its
equation there has character zero,
incompatible with a smooth plane lift.
At the weight-four vertex the plane
characters are $1,3$ and its equation
has character one. We check the lift
without assuming that four is prime.
The transitive orbit of normalized
branches has size one, two, or four.

\noindent\textbf{Step 27.} For four branches the stabilizer is
trivial and every invariant maximal-ideal
function has order at least two on a
branch, contradicting regularity of the
quotient. For two branches the stabilizer
has order two. A singular branch gives
invariant orders at least four, exceeding
two. Two smooth branches would have a
quadratic initial equation, but the
quadratic characters of the ambient
action are only zero and two, never one.
For one branch the normalization action
is faithful of order four. Both ambient
coordinate orders are odd, since their
characters are one and three. A singular
branch therefore has multiplicity at
least three, and all invariant functions
have order at least six, exceeding four.
These contradictions leave a single
smooth branch.

\noindent\textbf{Step 28.} The degree-nine cone is therefore smooth
away from its vertex. It is singular
at the vertex, because no variable has
weight nine. Its hypersurface critical
point is isolated, so the isolated-point
Morse transport of P4.52 gives a spherical
class of square $-2$ in its Milnor fibre.
This again contradicts the zero pairing
transported from the actual torus.

\noindent\textbf{Step 29.} All three possibilities for the sole
simple elliptic exceptional curve have
now been excluded. Together with P4.53
this proves that the horizontal pole
section cannot be contracted. P4R.18
then gives the asserted generic $D_4$
image and degree two or three. The
other conductor component retains its
unrestricted transverse type.

\noindent\textbf{Step 30.} The exact weighted cohomology counts,
order-six characters, quotient quadratic
characters and an explicit smooth cubic
are recomputed in
\path{code/part4/}\allowbreak\path{simple_elliptic_weighted.py}
and its executed outputs.
\end{proof}


\subsection{The conductor graph and the horizontal spectral cover}

\begin{proposition}[P4.55]\label{night:P4.55}\label{stmt:P4.55}
\StatementRecord{night:P4.55}{proved}{Historical events: GPT-6 Astra (OpenAI), ChatGPT Pro mode, effort Pro, S6\_06\_BLIND, 16--18 November 2026: targeted proof check. Also checked in this writing. Pre-continuation writing events: GPT-6 Astra (OpenAI), ChatGPT Pro mode, effort Pro, S6\_10\_PAPERS, 18 November 2026: targeted proof check. Continuation metadata events: GPT-6 Astra (OpenAI), ChatGPT Pro mode, effort Pro, S6\_10\_PAPERS continuation, 19 November 2026: used and rechecked in this run; status-vocabulary reconciliation; metadata reconciliation only; no new mathematical check. \emph{Pre-continuation recorded proof scope.} The night-run re-derivation was by GPT-6 Astra (OpenAI, ChatGPT Pro mode, effort Pro), 16--18 November 2026. Also checked in this writing by GPT-6 Astra (OpenAI), ChatGPT Pro mode, effort Pro, 18 November 2026. The two-sheet graph budgets, symbolic double-conductor defect bound, ruled cubic spectral genus formula, degree-three exclusion, and existence of the unique remaining vertical sheet were re-derived, including both junction and cube exclusions}{\ref{night:P4.31}; \ref{night:P4.34}; \ref{night:P4.54}; \ref{night:P4R.13}; \ref{night:P4R.18}; \ref{night:P4R.19}}

In the remaining positive ruled case, the image $A$ of the
horizontal pole section is a smooth rational curve. The map
$H\to A$ has degree exactly two. The third generic
normalization preimage along $A$ is supplied by one rational
vertical conductor component, mapping isomorphically to $A$.
The other image component retains its unrestricted transverse
type. A separated self-identification point on either image
component is impossible.
\end{proposition}
\begin{proof}
\noindent\textbf{Step 1.} We write $T\to N\to S$ for the minimal normalization
resolution. By P4.54 and P4R.18, the pole section $H$ is
noncontracted, has genus $g\ge1$, and maps with degree two
or three to $A^\nu=\mathbb P^1$. The full generic transverse
germ along $A$ is $D_4$. The reduced image conductor is
$C=A\cup B$, with two rational normalized components and
one graph cycle.

\noindent\textbf{Step 2.} The reduced normalization conductor $D$ has a tree as its
point-identification graph. Its positive-genus component
$H$ is smoothly embedded: the ruling descends to $N$, and
restriction of that ruling is an inverse to $H\to N$.
In particular no two points of $H$ are identified by the
normalization resolution. Every other normalized conductor
component is rational. These are the actual conductor
curves; vertical components and their identifications are
retained throughout the argument.

\noindent\textbf{Step 3.} At any point of $C$, points of $N$ over it correspond to
local irreducible hypersurface branches of $S$. Two such
branches intersect along a curve, by the height bound for
two equations in a regular local ring of dimension three.
Their nonempty supports among the local analytic branches
of $C$ therefore pairwise intersect.

\paragraph{Separated self-identifications.}
Suppose $p\in B\setminus A$ has two analytic branches.
Let $D_B$ be the union of the conductor components mapping
to $B$, with $c_B$ normalized components and $s_B\ge1$
connected components. Its graph is a forest. Every component
maps finitely and surjectively to $B^\nu$, so the numbers
$r_+,r_-$ of its normalized points over the two preimages
of $p$ satisfy $r_+,r_-\ge c_B$.

\noindent\textbf{Step 4.} There are $n$ normalization points of $S$ over $p$.
Their nonempty supports in $\{+,-\}$ pairwise intersect.
They have a common element: a singleton forces its element
into every support, and without a singleton all supports
are the entire two-element set. After interchanging signs,
$r_+\ge n$. The graph contribution at $p$ is consequently
\[
                  r_++r_--n\ge c_B.
\]
The whole forest has contribution $c_B-s_B<c_B$, a
contradiction. The same proof applies to a self-identification
of $A$ away from $B$, using all components mapping to $A$.
Thus the image cycle consists either of two distinct
intersections of $A$ and $B$, or of one three-branch point
belonging to both components.

\paragraph{The actual primary conductor and smoothness of $A$.}
Let $C_{A,\mathrm{sch}}$ be the primary component of the
actual image conductor supported on $A$. The primary
restriction maps of P4R.13, used in P4R.18, give
\[
                    g-1\le\chi(\mathcal O_{C_{A,\mathrm{sch}}}).
                                                            \tag{4.55a}
\]
The generic ordinary-triple conductor is the transverse
maximal ideal squared. Purity of the actual conductor
quotient identifies this primary ideal with $P_A^{(2)}$,
as in P4.31. The exact normalized conormal calculation of
P4.34 therefore gives
\[
 \chi(\mathcal O_{C_{A,\mathrm{sch}}})
       =5-\sum_p\bigl(\delta(A,p)+\tau_p-r_p\bigr).
                                                            \tag{4.55b}
\]
A planar point contributes exactly $5\delta(A,p)$, and a
nonplanar singular point with at most two branches contributes
at least seven. There cannot be three branches of $A$ at one
point: the identity
$\sum_p(r_p(C)-1)=2$ would then leave no identification for
the necessary connection with $B$; placing that connection
at the same point gives at least four branches and also
exceeds the identity.

\noindent\textbf{Step 5.} Since $g\ge1$, (4.55a)--(4.55b) bound the total defect by
five. If $A$ is singular, it is planar everywhere, has total
delta invariant one, and $g=1$. It is then a locally complete
intersection curve in $Y$. Its conormal
$Q=P_A/P_A^2$ is locally free of rank two. Let
$j:\mathbb P^1\to A$ be its normalization and put $E_0=j^*Q$.
Adjunction gives
\[
             \omega_A=(\omega_Y|_A)\otimes(\det Q)^{-1}.
\]
The degree of $j^*\omega_Y$ is zero because $b_2(Y)=0$.
Also $p_a(A)=1$, so $\deg j^*\omega_A=0$. The degree of a
line bundle is preserved by finite birational normalization:
the normalization sequences for that line bundle and for
$\mathcal O_A$ have the same finite-length correction.
Consequently
\[
                              \deg E_0=0.                 \tag{4.55c}
\]

\noindent\textbf{Step 6.} The actual reduced equation has generic order three along
$A$. Powers of a complete-intersection curve ideal are
unmixed, so $P_A^{(3)}=P_A^3$. Its third normal term is an
actual global section
\[
 q_3\in H^0\!\left(A,
          \operatorname{Sym}^3 Q\otimes\mathcal O_Y(S)|_A\right).
\]
The line twist has degree zero and becomes trivial on
$\mathbb P^1$. On the quotient projective bundle
$\mathbb P(E_0)$, with tautological class $\xi$ and fibre
class $F$, the pulled-back cubic has zero divisor of class
$3\xi$. Remove its common zero fibres, of total degree
$b\ge0$. The residual spectral curve $\Gamma$ has class
$3\xi-bF$. Its horizontal components are reduced, because
the generic $D_4$ cubic has three distinct roots.

\noindent\textbf{Step 7.} We specify the spectral morphisms. At a generic point of
$H$, the first normal derivative of $T\to S\subset Y$
gives the direction of its smooth transverse branch and
hence a point of $\mathbb P(E_0)$. On the smooth compact
curve $H$, its two local coordinates extend after cancelling
their common power. This gives a holomorphic map over the
actual map $H\to\mathbb P^1$. Distinct generic preimages
on $H$ give distinct cubic roots. Hence its map onto the
spectral component $\Gamma_H$ is finite and birational,
and $H$ is that component's normalization. In degree three,
$\Gamma=\Gamma_H$. In degree two, precisely one remaining
rational vertical conductor component supplies the third
root, with degree one over $A$. Its direction curve
$\Gamma_R$ is a section and is isomorphic to $\mathbb P^1$.
Thus $\Gamma=\Gamma_H\cup\Gamma_R$. These exhaust the
three generic roots.

\noindent\textbf{Step 8.} For a rank-two bundle of degree $e$ on $\mathbb P^1$,
\[
 \xi^2=e,\qquad \xi F=1,\qquad F^2=0,\qquad
 K_{\mathbb P(E_0)}=-2\xi+(e-2)F.
\]
Adjunction gives the exact arithmetic genus
\[
                    p_a(3\xi-bF)=3e-2-2b.               \tag{4.55d}
\]
With (4.55c) this is $-2-2b$. In degree three it is at
least $g\ge1$. In degree two the exact union sequence gives
\[
 p_a(\Gamma)=p_a(\Gamma_H)+p_a(\Gamma_R)
           +\operatorname{length}(\Gamma_H\cap\Gamma_R)-1
       \ge g-1\ge0.
\]
Both alternatives are impossible. Therefore $A$ is smooth.

\paragraph{The three-branch image point in degree three.}
Suppose now that $H\to A$ has degree three. Then $H$ is
the only conductor component mapping onto $A$. Since $A$
is smooth, a single three-branch image point consists of
one $A$ branch and two $B$ branches, denoted $+$ and $-$.

\noindent\textbf{Step 9.} Let $c_B$ be the number of normalized conductor components
mapping to $B$. The entire conductor has $1+c_B$
normalized components and a tree graph, so its total
point-identification contribution is $c_B$. At the proposed
point write $r_A\ge1$, $r_+,r_-\ge c_B$ for the numbers
of conductor branches, and $n$ for the number of normalization
points. Each normalization point has at most one $A$
conductor branch, because $H$ is embedded.

\noindent\textbf{Step 10.} A singleton support $\{A\}$ forces every point to carry
$A$. Then $n=r_A$, and the graph contribution is
$r_++r_-\ge2c_B>c_B$. A singleton $\{+\}$ forces
$r_+\ge n$, giving contribution at least
$r_A+r_-\ge1+c_B$; the same applies to $\{-\}$.
Thus every support has at least two elements. Consequently
$r_A+r_++r_-\ge2n$, and its graph contribution is at least
\[
       \frac{r_A+r_++r_-}{2}
                         \ge c_B+\frac{r_A}{2}>c_B.
\]
This is again impossible. Only two distinct $A$--$B$
junctions remain in degree three.

\paragraph{Both degree-three junctions are common cubic zeros.}
Let these junctions be $p,q$. If a local normalization
branch over $p$ has conductor only over $A$, pairwise
hypersurface intersections force every normalization point
over $p$ to carry $A$. Thus $n=r_A$, and the graph
contribution at $p$ is $r_B\ge c_B$, using the entire
budget. At $q$ some normalization point must carry both
$A$ and $B$: otherwise an $A$-only point and a $B$-only
point have disjoint supports. That point contributes
positively to the graph, contradicting the budget. There
is therefore no $A$-only local branch at either junction.

\noindent\textbf{Step 11.} We retain all factors in the following local calculation.
Write the smooth curve as $A=\{x=y=0\}$. If its nonzero
cubic normal form has a simple root, use the normal blowup
chart $x=yu$. The implicit-function theorem on the strict
equation gives that root as $u=\alpha(y,z)$. The actual
function has the analytic factor
\[
                         f_1=x-y\alpha(y,z).
\]
This branch is smooth. Exact division writes $f=f_1f_2$,
and restriction of $f_2$ to this branch is $y^2$ times a
unit: the other cubic directions differ from the simple
root at the point. Its intersections with all other
branches are therefore supported only on $A$. It is the
forbidden $A$-only normalization branch.

\noindent\textbf{Step 12.} Every nonzero binary cubic has a simple root unless it is
a cube. P4R.19 excludes a nonzero cube by its exact finite
weighted maps and compact intersection class. Hence the
whole cubic tensor vanishes at both $p$ and $q$.

\paragraph{The degree-three genus contradiction.}
For the now smooth curve $A=\mathbb P^1$, adjunction gives
$\deg N_{A/Y}=-2$, so $E=N_{A/Y}^*$ has degree two.
The actual cubic has at least two common zeros. Its residual
spectral divisor is $3\xi-bF$ with $b\ge2$, and is
irreducible with normalization $H$. Formula (4.55d) gives
\[
                    p_a(\Gamma)=4-2b\le0,
\]
contradicting $p_a(\Gamma)\ge g\ge1$. Degree three is
excluded. The degree-two alternative and its unique rational
third sheet are precisely the possibilities already identified
by the spectral morphism above.

\noindent\textbf{Step 13.} The exact intersection polynomial and the graph inequalities
are recomputed in
\path{code/part4/d4_graph_spectral.py}; its saved output
reports PASS. The inequalities in this proof establish the
unrestricted graph counts.
\end{proof}


\subsection{The other image curve with at most two generic branches}

\begin{proposition}[P4.56]\label{night:P4.56}\label{stmt:P4.56}
\StatementRecord{night:P4.56}{proved}{Historical events: GPT-6 Astra (OpenAI), ChatGPT Pro mode, effort Pro, S6\_06\_BLIND, 16--18 November 2026: targeted proof check. Also checked in this writing. Pre-continuation writing events: GPT-6 Astra (OpenAI), ChatGPT Pro mode, effort Pro, S6\_10\_PAPERS, 18 November 2026: targeted proof check. Continuation metadata events: GPT-6 Astra (OpenAI), ChatGPT Pro mode, effort Pro, S6\_10\_PAPERS continuation, 19 November 2026: used and rechecked in this run; status-vocabulary reconciliation; metadata reconciliation only; no new mathematical check. \emph{Pre-continuation recorded proof scope.} The night-run re-derivation was by GPT-6 Astra (OpenAI, ChatGPT Pro mode, effort Pro), 16--18 November 2026. Also checked in this writing by GPT-6 Astra (OpenAI), ChatGPT Pro mode, effort Pro, 18 November 2026. The rank-two finite-flat ribbon calculation, the graph alternatives, exact completed-square normalization, and the equality 2b=m versus b>=m were re-derived. The SNC conductor support used here follows by tracking the minimal ruled pole divisor through each point blowup in P4R.7 and P4R.9}{\ref{night:P4.55}; \ref{night:P4R.7}; \ref{night:P4R.9}; \ref{night:P4R.19}}

In the remaining case of P4.55, the other conductor image
component $B$ has at least three generic normalization
branches. This conclusion imposes no preliminary bound on
its transverse multiplicity, conductor exponent, or tangency
order.
\end{proposition}
\begin{proof}
\noindent\textbf{Step 1.} We retain the smooth curve $A$, its degree-two horizontal
cover $H$, and its degree-one vertical cover $R$. All other
normalized conductor components map to $B$. Let $r_B$ be
the generic normalization branch number along $B$, and
$c_B$ the number of global conductor components over it.
Their finite degrees sum to $r_B$. We exclude $r_B=1,2$.
The normalized conductor has $2+c_B$ components and a tree
graph, with total identification contribution $1+c_B$.
Distinct points on one normalized component cannot be
identified, since that would create a graph loop.

\paragraph{A length-two intersection calculation.}
Let $F,G$ be two local hypersurface branches whose
scheme-theoretic intersection is supported on the smooth
curve $A$ and has generic length $m\in\{1,2\}$ along it.
The ring
\[
                 Q=\mathbb C\{x,y,z\}/(F,G)
\]
is a one-dimensional complete intersection. Its only
associated prime is the ideal of $A$. Choose $z$ as a
coordinate on $A$. Then $Q$ is finite and torsion-free
over $\mathbb C\{z\}$, hence free of rank $m$. The
reduction map to $\mathbb C\{z\}$ splits through the
chosen coordinate.

\noindent\textbf{Step 2.} For $m=1$, the ring is $\mathbb C\{z\}$, and the linear
parts of $F,G$ have rank two. Both branches are smooth.
For $m=2$, the nilradical is a free rank-one module over
$\mathbb C\{z\}$. Its square vanishes generically and
therefore vanishes, by torsion-freeness. A generator
$\epsilon$ gives the exact analytic algebra isomorphism
\[
                    Q\simeq\mathbb C\{z,\epsilon\}/(\epsilon^2).
                                                            \tag{4.56a}
\]
Thus $Q$ has embedding dimension two, so one of $dF,dG$
is nonzero. If $G$ has generic normal order two along $A$,
it belongs to the ordinary square of its smooth ideal and
$dG$ vanishes on all of $A$. Then $dF\ne0$ and $F$ is
smooth. We use this calculation on the original intersection
scheme, retaining its multiplicity.

\paragraph{Two distinct image junctions.}
At a two-branch image point, let $r_A,r'_B$ count the
normalized conductor points over its two branches, and
let $n$ count normalization points of $S$. Pairwise
intersection of the hypersurface branches gives a common
element in their supports in $\{A,B\}$. Hence
$n\le\max(r_A,r'_B)$, and the graph contribution is at
least $\min(r_A,r'_B)$. Here $r_A=2$ or three, from the
one or two $H$ points and the unique $R$ point.

\noindent\textbf{Step 3.} If $c_B=2$, both degrees are one, and $r'_B=2$ at either
junction. The two contributions total at least four,
exceeding the budget three. If $c_B=1$ and its degree is
two, an unramified junction costs at least two and the
other at least one, exceeding the budget two. Thus one
component $V$ remains, with exactly one normalized point
over each junction.

\noindent\textbf{Step 4.} If $B$ belongs to every branch support at a junction,
there is only one normalization point there and its
contribution is $r_A\ge2$. The other junction makes this
impossible. Therefore $A$ belongs to every support. The
budget forces contribution exactly one at each junction:
every normalization point has one $A$ branch, and the
sole $V$ point is paired with one of them. Connectedness
forces the graph path $H$--$V$--$R$, with its edges at
the two different junctions.

\noindent\textbf{Step 5.} At the $R$--$V$ junction, let $F$ be the corresponding
hypersurface branch. The other branches come from $H$
and carry only $A$. If its two preimages are distinct,
intersection with either such branch has support $A$ and
generic length one, because the full generic $A$ germ
is an ordinary triple point. The preceding calculation
makes $F$ smooth. If the two $H$ preimages coalesce,
their one hypersurface branch $G$ has generic normal
order two along $A$, and the intersection with $F$ has
length two. Formula (4.56a) again makes $F$ smooth.

\noindent\textbf{Step 6.} This is impossible. The branch $F$ carries all the nearby
generic $B$ preimages: either the single branch of a
singular unibranch transverse germ, or both branches
when $V\to B^\nu$ has degree two and ramifies at this
point. In either case $F$ is intrinsically nonnormal
along the nearby $B$ germ. It cannot be smooth in a
neighbourhood of the junction.

\paragraph{A three-branch image point with one $B$ component.}
The image branches are $A,B_+,B_-$, since $A$ is smooth.
If $c_B=1$, the normalized points on $V$ over the two
$B$ signs are distinct and cannot be identified in $N$.
No normalization point can therefore carry both signs.
Pairwise hypersurface intersections force every point
carrying either sign to carry $A$ as well.

\noindent\textbf{Step 7.} The graph budget two now forces exactly one point over
each sign, each paired with exactly one $A$ branch, and
no other identification. The pairs are $H$--$V$ and
$R$--$V$. An extra $H$ preimage is $A$-only. At the
$R$--$V$ point, the other $H$ branch or branches intersect
its hypersurface branch only along $A$, because their
$B$ supports have the opposite sign or are empty. The
same length-one or length-two calculation makes this
branch smooth, contradicting its intrinsic nonnormality
along the $B$ germ whose generic preimages it contains.

\paragraph{A three-branch point with two $B$ components.}
Write the two degree-one components as $V,W$. If $h=1$
or two is the number of $H$ points over the $A$ value,
the total normalized branch count is
\[
                         T=(h+1)+2+2=h+5.
\]
The graph budget three gives $n\ge h+2$ normalization
points. No singleton image support is possible: a
singleton $\{A\}$ forces $n\le h+1$, and a singleton
$\{B_+\}$ or $\{B_-\}$ forces $n\le2$. Hence every
point has at least two image labels and $n\le T/2$.
This excludes $h=2$. For $h=1$ there are precisely
three points, each with exactly two conductor branches
of different image labels. Up to signs and interchange
of $V,W$, their supports are
\[
 \begin{array}{c|c}
 \text{hypersurface branch}&\text{conductor branches}\\ \hline
 F_R&R,\ B_-\text{ on }V\\
 F_H&H,\ B_+\text{ on }W\\
 F_0&B_+\text{ on }V,\ B_-\text{ on }W.
 \end{array}
                                                            \tag{4.56b}
\]
The tree is the path $H$--$W$--$V$--$R$, and $H\to A$
is ramified with local degree two at its point.

\noindent\textbf{Step 8.} The intersection of $F_R$ and $F_H$ is supported on $A$
and has generic length two. The equation of $F_H$ has
normal order two along $A$, so (4.56a) makes $F_R$
smooth. Choose its equation as $X$, a coordinate $Y$
on that surface cutting $A$, and a coordinate $z$ on
$A$. Preparation and depression in $Y$, preserving the
ideal $(X,Y)$, give the actual equation
\[
                       F_H=\text{unit}\,[Y^2+X^2k(X,z)].
                                                            \tag{4.56c}
\]
The second $X$ follows from normal multiplicity two
along every nearby point of $A$.

\noindent\textbf{Step 9.} The monic finite partial normalization is
\[
      N'=\{u^2+k(X,z)=0\},\qquad Y=Xu.
                                                            \tag{4.56d}
\]
It is an isomorphism away from $A$. The branch $F_H$
is generically smooth along $B_+$: the two generic
$B$ sheets belong separately to $F_H$ and $F_0$.
Thus $N'$ has at most an isolated singularity here.
Its hypersurface Cohen--Macaulay property and regularity
in codimension one make it normal, so it is the actual
normalization.

\noindent\textbf{Step 10.} The reduced divisor $X=0$ is supported on $H$ and has
equation $u^2+k(0,z)=0$. The actual $H$ is smoothly
embedded and has degree two with ramification over $A$.
An even zero of $k(0,z)$ gives two branches; an odd
zero of order at least three gives a singular curve.
Therefore its zero is simple, $k_z(0,0)\ne0$, and
$N'$ is smooth. The minimal normalization resolution
is the identity in this neighbourhood.

\noindent\textbf{Step 11.} Let $m\ge1$ be the generic intersection multiplicity of
the two branches along $B$. It is the same on both
local signs, because they belong to the same irreducible
image component. Let $G$ define $F_0$, and set
\[
                      b=\operatorname{ord}_{p,A}(G|_A)>0.
\]
On the smooth branch $F_R$, its divisor is $mB_-$.
Restricting to $A$ gives
\[
                 b=m I_{F_R}(A,B_-)\ge m.                \tag{4.56e}
\]
On the smooth normalization (4.56d), the divisor of $G$
is $mW$. The actual pole support is normal crossing,
so $H$ and $W$ meet transversely. Its restriction to
$H$ consequently has order $m$. But it is the pullback
of the same ambient restriction $G|_A$ under the local
degree-two map. Hence
\[
                              2b=m,
\]
contradicting (4.56e). Every case with $r_B\le2$ is
excluded.

\noindent\textbf{Step 12.} The exact incidence partitions and the normalization
substitution are checked in
\texttt{code/part4/d4\_other\_two\_branches.py}.
Its corrected run reports PASS; the initial structural
symbolic-equality assertion failure is preserved separately.
\end{proof}


\subsection{Exclusion of the two-junction spectral configuration}

\begin{proposition}[P4.57]\label{night:P4.57}\label{stmt:P4.57}
\StatementRecord{night:P4.57}{proved}{Historical events: GPT-6 Astra (OpenAI), ChatGPT Pro mode, effort Pro, S6\_06\_BLIND, 16--18 November 2026: targeted proof check. Also checked in this writing. Pre-continuation writing events: GPT-6 Astra (OpenAI), ChatGPT Pro mode, effort Pro, S6\_10\_PAPERS, 18 November 2026: targeted proof check. Continuation metadata events: GPT-6 Astra (OpenAI), ChatGPT Pro mode, effort Pro, S6\_10\_PAPERS continuation, 19 November 2026: used and rechecked in this run; status-vocabulary reconciliation; metadata reconciliation only; no new mathematical check. \emph{Pre-continuation recorded proof scope.} The night-run re-derivation was by GPT-6 Astra (OpenAI, ChatGPT Pro mode, effort Pro), 16--18 November 2026. Also checked in this writing by GPT-6 Astra (OpenAI), ChatGPT Pro mode, effort Pro, 18 November 2026. The Hessian-zero consequence of at least three branches, cubic spectral identity, two-junction ruling, and every allocation of the simple common cubic zero between linear and quadratic factors were re-derived. Each case contradicts the required number of branches}{\ref{night:P4.55}; \ref{night:P4.56}; \ref{night:P4R.19}}

The remaining degree-two configuration of P4.55--P4.56
cannot have its image-conductor graph cycle realized at
two distinct intersections of its two components.
\end{proposition}
\begin{proof}
\noindent\textbf{Step 1.} The component $B$ has at least three generic normalization
branches, hence transverse multiplicity at least three.
Along the smooth $A$, choose coordinates $A=\{x=y=0\}$,
with $z$ on $A$. The actual equation belongs to $(x,y)^3$;
write $q_3(x,y;z)$ for its normal cubic.

\noindent\textbf{Step 2.} A local branch of $B$ has first nonzero normal displacement
$(a,b)t^q$, with $(a,b)\ne0$. All ambient Hessian entries
of the actual equation vanish identically along $B$, by
its generic multiplicity. The first normal-order term of
its normal Hessian is
\[
       t^q\operatorname{Hess}_{x,y}
                    (q_3(\,\cdot\,,\,\cdot\,;0))(a,b).
\]
Terms of higher normal order and terms from varying $z$
have larger order in $t$. This matrix is therefore zero.
For a homogeneous binary cubic $Q$, a nonzero vector
annihilating this Hessian makes the directional derivative
of $Q$ the zero quadratic. In coordinates with that vector
as first axis, $Q$ is a cube of the other variable, or zero.
P4R.19 excludes the nonzero cube. Thus the cubic vanishes
at every $A$--$B$ junction, including tangent junctions.

\paragraph{The global spectral equality.}
Let $p,q$ be the two distinct junctions. The common cubic
zero divisor has degree $b\ge2$. On $\mathbb P(N_{A/Y}^*)$
its residual spectral curve is $\Gamma_H\cup\Gamma_R$,
of class $3\xi-bF$. The bundle has degree two, $\Gamma_R$
is a section, and $H$ normalizes $\Gamma_H$. Adjunction
and the exact union sequence give
\[
       g(H)+\delta(\Gamma_H)+\Gamma_H\cdot\Gamma_R=5-2b.
\]
More precisely, the displayed formula is an equality of
integers with $\delta(\Gamma_H)$ the total normalization
length. Since $g(H)\ge1$, it forces
\[
 b=2,\qquad g(H)=1,\qquad \delta(\Gamma_H)=0,
                  \qquad\Gamma_H\cdot\Gamma_R=0.          \tag{4.57a}
\]
The two common cubic zeros are simple, the double spectral
cover is smooth, and it is disjoint from the third section.

\paragraph{The graph support at a junction.}
Let $c_B$ count the global conductor components over $B$.
The graph budget is $c_B+1$. At either junction the
contribution is at least $\min(r_A,r'_B)$, where
$r_A\ge2$ and $r'_B\ge c_B$ count normalized points
over its two image branches.

\noindent\textbf{Step 3.} If $c_B=2$, the two contributions total at least four,
exceeding the budget three. If $c_B=1$, a junction with
at least two $B$ preimages costs at least two, and the
other at least one, exceeding the budget two. Thus both
have one $B$ preimage. The argument of P4.56 gives the
path $H$--$V$--$R$, with one edge at each junction.
At the $R$--$V$ junction the local branch carries all
nearby generic $B$ sheets of $V$, regardless of its
finite degree. Its intersection with the $H$ branch or
branches is a length-one or length-two thickening of
smooth $A$. P4.56 makes it smooth, contradicting its
intrinsic nonnormality from the at least three $B$ sheets.

\noindent\textbf{Step 4.} Therefore $c_B\ge3$. If $A$ belongs to every branch
support at a junction, that point costs at least
$r'_B\ge c_B$, and the other costs at least two.
This exceeds $c_B+1$. The two-element support argument
consequently forces $B$ to belong to every local
hypersurface branch at each junction. Branches carrying
only $B$ have been retained in this conclusion.

\paragraph{Choosing distinct normalization points.}
The ruling descends to $r_N:N\to H$. Its restriction
to $H$ is the identity, and $R$ maps to a fixed point
$b_0\in H$. If a point of $H$ and a point of $R$ have
the same image in $N$, that $H$ point is $b_0$.
This includes identifications through a contracted
exceptional chain. All such coincidences therefore lie
over the single value $a_0\in A$ which is the image of
$b_0$. Choose a junction different from $a_0$, and call
it $p$. Its $H$ preimages and its $R$ preimage are
distinct normalization points. Different $H$ points
are already distinct because $H$ is embedded.

\noindent\textbf{Step 5.} The actual local equation factors into the corresponding
irreducible hypersurface factors and any factors carrying
only $B$. For an $A$-containing factor, let $b_i\ge0$
be the common $z$-order of its first nonzero normal term.
For a factor not containing $A$, let $c_j\ge1$ be the
order of its nonzero restriction to $A$ at $p$.
Gauss's lemma over $\mathbb C\{z\}$ gives
\[
                         \sum_i b_i+\sum_j c_j=1.         \tag{4.57b}
\]
Indeed primitive homogeneous polynomials have nonzero
product modulo $z$, so their common coefficient orders
add exactly. The right side is the simple cubic zero
of (4.57a).

\noindent\textbf{Step 6.} If $H$ is unramified here, there are three
$A$-containing factors of normal order one. At least
two have nonzero linear normal terms, by (4.57b).
Their directions differ by (4.57a). The implicit-function
theorem makes their intersection exactly $A$, whereas
both must contain $B$. This is a contradiction.
Thus $H$ is ramified at $p$.

\noindent\textbf{Step 7.} There are now two $A$-containing factors, denoted by
$F_H$ and $F_R$, of normal orders two and one. Their
residual spectral directions differ. Equality (4.57b)
has precisely the following three possibilities.

\paragraph{A factor carrying only $B$ uses the zero.}
Here $b_H=b_R=0$ and one such factor restricts to $A$
with order one. The branch $F_R$ is smooth. In coordinates
$X=F_R$, $A=\{X=Y=0\}$, restriction of $F_H$ to $X=0$
is $Y^2$ times a unit, since its residual direction
differs from the $R$ direction. Thus these two branches
intersect only along $A$. They both contain $B$, which
is impossible.

\paragraph{The quadratic factor uses the zero.}
Here $b_H=1$, $b_R=0$, and there is no other factor.
The branch $F_R$ is again smooth, and in its coordinates
exact division gives
\[
                         F_H|_{X=0}=Y^2h(Y,z).
\]
The value $h(0,z)$ has a simple zero: the residual
quadratic is nonzero on the $R$ direction, so its
common scalar factor survives with order one on this
restriction. Thus $h_z(0,0)\ne0$. Its zero curve is the
extra image branch $B$, with intersection multiplicity
one between $F_H$ and $F_R$.

\noindent\textbf{Step 8.} The smooth $F_R$ branch supplies one generic $B$
normalization branch. Each branch supplied by $F_H$
has positive intersection multiplicity with it. There
can therefore be at most one such branch. The full
$B$ germ would have at most two branches, contradicting
P4.56.

\paragraph{The linear factor uses the zero.}
Finally $b_R=1$, $b_H=0$, with no other factor.
The nonzero quadratic form of $F_H$ has a double root.
Preparation and depression preserving $A$ give
\[
                   F_H=\text{unit}\,[Y^2+X^2k(X,z)].
\]
Its spectral curve is $u^2+k(0,z)=0$. By (4.57a) this
is smooth and its double cover is ramified at $p$.
Hence $k_z(0,0)\ne0$. The exact finite map
\[
 \{u^2+k(X,z)=0\}\longrightarrow\{F_H=0\},
                         \qquad Y=Xu
\]
has smooth source, is birational, and is an isomorphism
off $A$. It is the normalization. Thus $F_H$ supplies
one generic $B$ branch.

\noindent\textbf{Step 9.} For $F_R$, the simple common zero of its linear normal
term makes its quadratic Taylor part
\[
                       zL(X,Y)+Q_2(X,Y),\qquad L\ne0.
\]
There is no $z^2$ term because $F_R$ vanishes on $A$.
Its Hessian has rank at least two. Rank three gives an
isolated critical point. In rank two the holomorphic
Morse lemma with a parameter gives the exact equation
$U^2+V^2+h(W)$. A nonisolated critical curve forces
$h'\equiv0$, hence $h=0$ because the factor vanishes
at the point. The resulting factorization
$(U+iV)(U-iV)$ contradicts local irreducibility.
Therefore $F_R$ has at most isolated singularities and
also supplies one generic $B$ branch. Again the full
branch number is two, contradicting P4.56.

\noindent\textbf{Step 10.} All allocations in (4.57b) have been excluded, proving
the assertion. The exact genus budget and nonzero
Hessian minor are recomputed in
\texttt{code/part4/d4\_two\_junction.py}; its saved
output reports PASS.
\end{proof}


\subsection{The completed irreducible classification and sign theorem}

\begin{proposition}[P4.58]\label{night:P4.58}\label{stmt:P4.58}
\StatementRecord{night:P4.58}{proved}{Historical events: GPT-6 Astra (OpenAI), ChatGPT Pro mode, effort Pro, S6\_06\_BLIND, 16--18 November 2026: targeted proof check. Also checked in this writing. Pre-continuation writing events: GPT-6 Astra (OpenAI), ChatGPT Pro mode, effort Pro, S6\_10\_PAPERS, 18 November 2026: targeted proof check. Continuation metadata events: GPT-6 Astra (OpenAI), ChatGPT Pro mode, effort Pro, S6\_10\_PAPERS continuation, 19 November 2026: used and rechecked in this run; status-vocabulary reconciliation; metadata reconciliation only; no new mathematical check. \emph{Pre-continuation recorded proof scope.} The night-run re-derivation was by GPT-6 Astra (OpenAI, ChatGPT Pro mode, effort Pro), 16--18 November 2026. Also checked in this writing by GPT-6 Astra (OpenAI), ChatGPT Pro mode, effort Pro, 18 November 2026. Every displayed assembly step was checked against its stated local input. The positive ruled branch uses the actual P4.54 contraction exclusion, P4R.18 reduction, P4.55-57 graph and spectral exclusions, and P4R.20 endpoint calculation. The theorem retains the hypothesis that every support is irreducible or NC; individual dependency checks remain separately recorded}{\ref{night:P4.10}; \ref{night:P4.11}; \ref{night:P4.26}; \ref{night:P4.32}; \ref{night:P4.35}; \ref{night:P4.36}; \ref{night:P4.37}; \ref{night:P4.54}; \ref{night:P4.55}; \ref{night:P4.56}; \ref{night:P4.57}; \ref{night:P4.6}; \ref{night:P4.7}; \ref{night:P4.8}; \ref{night:P4R.13}; \ref{night:P4R.18}; \ref{night:P4R.20}}

Let $f:Y\to B$ be the fibration of Part~4.2, with
$b_1(Y)=b_2(Y)=0$, connected fibres and general fibre a
complex two-torus. Suppose every reduced fibre support
is irreducible or has normal crossings. Then
\[
 B=\mathbb P^1,\qquad c_3(Y)\in\{0,2\},
                       \qquad b_3(Y)=2-c_3(Y).
\]
Every reduced support has nonnegative Euler number.
An irreducible singular support whose normalization has
a rational resolution has Euler number two. Every other
irreducible support has Euler number zero. Such a support
is either smooth with a finite unramified torus cover,
or singular with smooth minimal elliptic ruled
normalization, or singular with smooth minimal Hopf
normalization. In particular the class-VII normalization
alternative is realized only by a smooth minimal Hopf
surface.

For a normal-crossing support the Euler number equals
its number of triple points, and is even. Its zero case
is exactly P4.8: in the reducible case the normalization
is a disjoint union of smooth minimal Hopf or non-K\"ahler
elliptic surfaces, with smooth elliptic conductor curves
and finite unramified branch maps. The equality
$c_3(Y)=0$ holds exactly when every support has one of
the corresponding zero-Euler forms.
\end{proposition}
\begin{proof}
\noindent\textbf{Step 1.} P4.26 proves the sign for every irreducible support,
including singular normalizations, and P4.6--P4.8 give
the full normal-crossing calculation. P4.35 proves the
exact value two for the rationally resolved singular
case. P4.36--P4.37 show that every class-VII case has
smooth minimal Hopf normalization and Euler number
zero. P4.32 gives the complete irreducible zero-Euler
classification and its converses.

\noindent\textbf{Step 2.} We finish the positive-genus ruled alternative without
a transverse-type assumption on its vertical conductor
branches. P4R.13 and the preceding reductions show that
a positive-Euler ruled support would have a doubled
horizontal pole section and an image conductor with two
rational normalized components and one graph cycle.
P4.54 excludes contraction of that section, using the
actual weighted normalization and blowdown maps.
P4R.18 reduces its generic image to $D_4$ with degree
two or three. P4.55 makes that image smooth and excludes
degree three, leaving the degree-two section and one
rational vertical third sheet. P4.56 forces at least
three generic branches on the other image component.

\noindent\textbf{Step 3.} The graph alternatives are exhaustive by P4.55. P4.57
excludes two distinct junctions. P4R.20 excludes a sole
three-branch image point, with the exact conductor-path
and endpoint normalization calculation. Thus no
positive-Euler ruled support remains. This proves the
asserted value zero for every nonrational irreducible
support, and P4.32 supplies its actual normalization.

\noindent\textbf{Step 4.} Finally Euler additivity gives $c_3(Y)$ as the sum of
these nonnegative support Euler numbers. The base,
Poincar\'e-duality calculation $e(Y)=2-b_3(Y)$, and
parity are proved in P4.26. They yield the displayed
global values. A sum of nonnegative integers vanishes
exactly when every summand vanishes, giving the equality
classification through P4.32 and P4.8.

\noindent\textbf{Step 5.} The hypothesis checks are the explicit constructions
P4.10 and P4.11. Removing $b_2=0$ permits the proved
Euler value $24$. Removing the torus-fibre hypothesis
permits the Hopf-fibre example, which fails the stated
normalization equality classification. The numerical
normal-crossing sign itself follows from the stronger
calculation in P4.7 and does not require that hypothesis.
\end{proof}


\subsection{The specialization group and the unibranch case}

\begin{proposition}[P4R.1]\label{night:P4R.1}\label{stmt:P4R.1}
\StatementRecord{night:P4R.1}{proved}{Historical events: GPT-6 Astra (OpenAI), ChatGPT Pro mode, effort Pro, S6\_06\_BLIND, 16--18 November 2026: targeted proof check. Also checked in this writing. Pre-continuation writing events: GPT-6 Astra (OpenAI), ChatGPT Pro mode, effort Pro, S6\_10\_PAPERS, 18 November 2026: targeted proof check; used and rechecked in this run. Continuation metadata events: GPT-6 Astra (OpenAI), ChatGPT Pro mode, effort Pro, S6\_10\_PAPERS continuation, 19 November 2026: used and rechecked in this run; status-vocabulary reconciliation; metadata reconciliation only; no new mathematical check. \emph{Pre-continuation recorded proof scope.} The night-run re-derivation was by GPT-6 Astra (OpenAI, ChatGPT Pro mode, effort Pro), 16--18 November 2026. Also checked in this writing by GPT-6 Astra (OpenAI), ChatGPT Pro mode, effort Pro, 18 November 2026. Compatible triangulation is applied in the smooth real analytic total space. Its open-star retraction and a proper small tube give an actual retraction to the irreducible support. General-position discs identify the complement-group kernel with the normal closure of one positive meridian, whose base winding is the full-fibre multiplicity. The image of the torus lattice is abelian and normal with finite cyclic quotient, and the retraction surjects this virtually abelian group onto pi1 of the support}{the hypotheses and explicit local calculation stated here}

For a proper holomorphic degeneration of complex two-tori with
smooth total threefold and central fibre $mS$, where $S$ is
irreducible and reduced, the group $\pi_1(S)$ has an abelian
normal subgroup with cyclic quotient of order dividing $m$.
\end{proposition}
\begin{proof}
We give the neighbourhood and group comparison explicitly.
The compatible triangulation of Lojasiewicz, Theorem~2,
\noindent\textbf{Step 1. Triangulation of semi-analytic sets}, p.~464 \cite{Lojasiewicz1964}, applies to
the smooth real analytic total space and its closed analytic
subset $S$. Its hypotheses hold because the total space is
countable and finite dimensional, and the collection consisting
of $S$ is locally finite and semianalytic. After barycentric
subdivision $S$ is a full subcomplex. In an open star of this
subcomplex, erase the barycentric coordinates at vertices
outside $S$ and normalize the remaining coordinates to sum
to one. Linear interpolation gives a strong deformation
retraction $r:V\to S$. The star can be chosen inside a
relatively compact neighbourhood of $S$.

Properness gives a small disc $\Delta_\epsilon$ with
$U=f^{-1}(\Delta_\epsilon)\subset V$: on the compact
complement of $V$ in a larger closed tube, $|f|$ has a
positive minimum. Thus the inclusion $i:S\to U$ and the
continuous restriction $r:U\to S$ satisfy $r\circ i=\id$.
This retraction comparison also follows from the statement
immediately after Definition~1 of Fernandez de Bobadilla--Kollar,
\noindent\textbf{Step 2. Homotopically trivial deformations}, p.~85 \cite{BobadillaKollar2012}; the preceding
construction supplies the particular comparison needed here.

Write $U^*=U\setminus S$. The proper smooth fibration
$U^*\to\Delta_\epsilon^*$ has fibre $T^4$. Its homotopy
sequence and $\pi_2(\Delta_\epsilon^*)=0$ give
\[
 \pi_1(U^*)=\Lambda\rtimes\langle t\rangle,
 \qquad \Lambda=\mathbb Z^4.
\]
The extension splits because its quotient is infinite cyclic.
The inclusion $U^*\to U$ induces a surjection on fundamental
groups. A loop can be perturbed off the real-codimension-two
smooth stratum of $S$ and off its higher-codimension singular
strata. The same general-position argument for a null-homotopy
disc avoids the singular strata and meets $S_{\rm reg}$ in
finitely many transverse points. Deleting little discs about
these intersections expresses the kernel as the normal closure
of the corresponding meridians. The smooth locus of the
irreducible $S$ is connected, so those meridians are conjugate,
with their positive orientation fixed by a complex transverse
coordinate.

At a smooth point of $S$, the defining map is a unit times
$z^m$. A positive meridian therefore maps to $m$ under
$\pi_1(U^*)\to\mathbb Z$. Let $A$ be the image of $\Lambda$
in $\pi_1(U)$. It is abelian and normal, and
$\pi_1(U)/A$ is cyclic with order dividing $m$.
The homomorphism $r_*:\pi_1(U)\to\pi_1(S)$ is surjective
because $r_*i_*=\id$. Its image $r_*(A)$ is abelian and
normal, and the induced quotient remains cyclic of order
dividing $m$. This proves the assertion for the actual support,
without replacing it by its normalization.
\end{proof}



\begin{proposition}[P4R.2]\label{night:P4R.2}\label{stmt:P4R.2}
\StatementRecord{night:P4R.2}{proved}{Historical events: GPT-6 Astra (OpenAI), ChatGPT Pro mode, effort Pro, S6\_06\_BLIND, 16--18 November 2026: targeted proof check. Also checked in this writing. Pre-continuation writing events: GPT-6 Astra (OpenAI), ChatGPT Pro mode, effort Pro, S6\_10\_PAPERS, 18 November 2026: targeted proof check; used and rechecked in this run. Continuation metadata events: GPT-6 Astra (OpenAI), ChatGPT Pro mode, effort Pro, S6\_10\_PAPERS continuation, 19 November 2026: used and rechecked in this run; status-vocabulary reconciliation; metadata reconciliation only; no new mathematical check. \emph{Pre-continuation recorded proof scope.} The night-run re-derivation was by GPT-6 Astra (OpenAI, ChatGPT Pro mode, effort Pro), 16--18 November 2026. Also checked in this writing by GPT-6 Astra (OpenAI), ChatGPT Pro mode, effort Pro, 18 November 2026. The smooth bijective normalization is a homeomorphism by compactness and the Hausdorff property, so its Euler number and fundamental group are the actual support invariants. Torsion adjunction with a nonzero conductor kills plurigenera. A ruled base of genus at least two has no virtually abelian finite-index surface group. The resulting rational, elliptic-ruled and classVII Chern formulas prove nonnegativity and the precise smooth minimal zero cases; the nonempty conductor excludes the Inoue alternative}{\ref{night:P4.12}; \ref{night:P4.8}; \ref{night:P4R.1}}

In the standing torus-fibration setting of this section, every irreducible singular support
whose normalization is smooth and whose normalization morphism
is bijective has nonnegative topological Euler number.
Its Euler number is zero precisely when that normalization is
a minimal ruled surface over an elliptic curve or a Hopf surface.
\end{proposition}
\begin{proof}
\noindent\textbf{Step 1.} Let $\nu:T\to S$ denote this finite normalization morphism.
A continuous bijection from the compact $T$ to the Hausdorff
$S$ is a homeomorphism. It therefore induces an isomorphism
of fundamental groups and equality of Euler numbers. P4R.1
makes $\pi_1(T)$ virtually abelian.

\noindent\textbf{Step 2.} The effective-conductor and canonical-torsion argument of
P4.12 applies to this same $T$ and proves $\kappa(T)=-\infty$.
The opened classification theorem used there leaves rational,
ruled, and class-VII minimal models, followed by point blowups.
Blowing up a smooth point preserves the fundamental group:
van Kampen applies to the common punctured four-ball with
simply connected boundary $S^3$. A ruled surface over a
genus-$g$ curve has the same fundamental group as that curve,
by the homotopy sequence of its $\mathbb P^1$-bundle.
For $g\ge2$ this group is not virtually abelian. Indeed a
finite-index subgroup corresponds to an unramified cover
of genus $1+d(g-1)>1$, and a genus-$h$ surface group maps
onto the nonabelian free group on $h$ generators by sending
the $b$ generators to the identity in its usual presentation.
Consequently $g\le1$.

\noindent\textbf{Step 3.} A rational surface has positive Euler number. A ruled surface
over an elliptic curve, blown up $k$ times, has Euler number
$k$. A class-VII surface has Euler number $b_2\ge0$, and
each blowup increases it by one. These formulas give
$e(S)=e(T)\ge0$. Equality in the ruled case is exactly
$k=0$. Equality in class VII gives a minimal surface with
$b_2=0$. The conductor is nonempty because $S$ is singular
and nonnormal; P4.12 gives an effective conductor curve on
$T$. The precise class-VII classification and absence of
curves on Inoue surfaces verified in P4.8 therefore force
$T$ to be Hopf. The converse follows from the same Euler
formulas and the homeomorphism $\nu$.
\end{proof}

\subsection{The conductor pushout for ruled normalizations of higher genus}

\begin{proposition}[P4R.3]\label{night:P4R.3}\label{stmt:P4R.3}
\StatementRecord{night:P4R.3}{proved}{Historical events: GPT-6 Astra (OpenAI), ChatGPT Pro mode, effort Pro, S6\_06\_BLIND, 16--18 November 2026: targeted proof check. Also checked in this writing. Pre-continuation writing events: GPT-6 Astra (OpenAI), ChatGPT Pro mode, effort Pro, S6\_10\_PAPERS, 18 November 2026: targeted proof check. Continuation metadata events: GPT-6 Astra (OpenAI), ChatGPT Pro mode, effort Pro, S6\_10\_PAPERS continuation, 19 November 2026: used and rechecked in this run; status-vocabulary reconciliation; metadata reconciliation only; no new mathematical check. \emph{Pre-continuation recorded proof scope.} The night-run re-derivation was by GPT-6 Astra (OpenAI, ChatGPT Pro mode, effort Pro), 16--18 November 2026. Also checked in this writing by GPT-6 Astra (OpenAI), ChatGPT Pro mode, effort Pro, 18 November 2026. Adjunction and Riemann-Hurwitz force the minimal ruled pole divisor over genus at least two to have horizontal part twice a section. The blowup coefficient rule gives only attached rational trees and no coefficient-one exceptional component. The actual finite normalization quotient is the conductor pushout; the conductor inclusion induces the isomorphism of fundamental groups needed for the van Kampen calculation. The curve graph free-product presentation and virtual abelianness bound its Euler number. In the only zero case all rational poles disappear, forcing the original minimal section and contradicting the independently proved nonnegative open-support Euler number}{\ref{night:P4.12}; \ref{night:P4.13}; \ref{night:P4R.1}; \ref{night:P4R.2}}

In the standing torus-fibration setting of this section, let an irreducible singular fibre
support $S$ have smooth normalization $\nu:T\to S$ whose
minimal model is ruled over a curve of genus $g\ge2$.
Write $D$ for the conductor divisor on $T$ and
$C=\operatorname{Sing}(S)_{\rm red}$.
The inclusion $D_{\rm red}\to T$ induces an isomorphism on
fundamental groups, and the normalization pushout induces
$\pi_1(S)\simeq\pi_1(C)$. Moreover,
\[
 e(C)\ge0,\qquad e(S)>0.
\]
The assertion includes conductor components on different
ruling fibres that map onto the same curve in $C$.
\end{proposition}
\begin{proof}
\noindent\textbf{Step 1.} The canonical comparison of P4.12 gives
$K_T+D\equiv0$, with $D$ effective. We first determine the
reduced conductor as an embedded topological space; this
step does not require the ruling to descend to $S$.
Let $\pi:T\to T_0$ be the sequence of point blowdowns to
the minimal ruled surface $r_0:T_0\to B$, and put
$D_0=\pi_*D$. Then $K_{T_0}+D_0\equiv0$ and
$D_0\cdot F=2$ for a ruling fibre $F$.

\noindent\textbf{Step 2.} There are three possibilities for the horizontal part of
$D_0$: a reduced bisection, two distinct reduced sections,
or twice one section. The first two are excluded directly.
For an irreducible horizontal component $H$ of coefficient
one, adjunction and the effectiveness of the other
components give
\[
 2p_a(H)-2=(K_{T_0}+H)\cdot H
       =-(D_0-H)\cdot H\le0.
\]
On the normalization of $H$, Riemann--Hurwitz for its
finite map of degree $d\ge1$ to $B$ gives
$2g(\widetilde H)-2\ge d(2g-2)>0$.
Since $p_a(H)\ge g(\widetilde H)$, these inequalities
are incompatible. Thus
\[
 D_0=2H+\sum_{j=1}^{v}a_jF_j,\qquad a_j\ge1,
\]
where $H$ is a section and the $F_j$ are distinct ruling
fibres. A horizontal curve of degree one is a section:
its normalization maps isomorphically to the smooth base,
and the resulting map to the projective bundle is a
section with closed image. In particular $H$ is smooth.

\noindent\textbf{Step 3.} At one step $\beta:T'\to T''$ of the blowup sequence,
let $\mu$ be the multiplicity of the current conductor
$D''$ at the centre and $E$ the exceptional curve.
Write $D'=D''_{\rm strict}+aE$. The canonical blowup
formula gives
\[
 K_{T'}+D'
   =\beta^*(K_{T''}+D'')+(1+a-\mu)E.
\]
Intersection with $E$ and numerical triviality give
$a=\mu-1$. Effectiveness therefore places every centre
on the conductor. At a smooth conductor point of
coefficient one, the exceptional curve is absent from
the new conductor. At a smooth point of coefficient at
least two, it is a rational leaf attached at one point.
At an intersection of two conductor components it
subdivides their edge. Starting with $H$ and the
transverse fibres $F_j$, these operations preserve
normal crossings and make $D_{\rm red}$ a copy of $H$
with rational trees attached. No cycle is introduced.

\noindent\textbf{Step 4.} The projection of $H$ to $B$ induces an isomorphism on
fundamental groups, and every attached rational tree
is simply connected. A point blowup leaves
$\pi_1(T)$ unchanged, as shown in P4R.2. The preceding
description also proves, at every blowup, that the
inclusion of the reduced conductor induces the same
isomorphism to the fundamental group of the surface.
Hence $\pi_1(D_{\rm red})\to\pi_1(T)$ is an
isomorphism.

\noindent\textbf{Step 5.} We next verify the pushout for the actual singular
support. The morphism $\nu$ is finite, is an
isomorphism away from the conductor, and has smooth
source. Its nonisomorphism locus equals $C$: a normal
point has a neighbourhood identified with its
normalization and is consequently smooth. Conversely
every smooth point is normal. Thus the natural map
\[
 T\underset{D_{\rm red}}{\cup}C\longrightarrow S
\]
is a continuous bijection. Both source and target are
compact Hausdorff spaces, so it is a homeomorphism.
The compatible triangulation used in P4R.1 makes
$D_{\rm red}\subset T$ a subcomplex and hence a
cofibration. Replacing the map $D_{\rm red}\to C$ by
its mapping cylinder computes the same homotopy
pushout: the homotopy extension property for this
cofibration contracts the inserted cylinder without
altering the homotopy type. Van Kampen on that double
mapping cylinder therefore gives
\[
 \pi_1(S)=\pi_1(T)
       *_{\pi_1(D_{\rm red})}\pi_1(C)
       \simeq\pi_1(C).
\]
All three spaces in the amalgam are connected. This
argument retains every identification made by the
finite conductor map, including identifications
between horizontal and vertical branches.

\noindent\textbf{Step 6.} P4R.1 makes this group virtually abelian. Let $c$ be
the number of irreducible components of $C$, let
$g_i$ be the genera of their smooth normalizations,
and let $b$ be the first Betti number of the incidence
graph of the finite point identifications. Replacing
each identification by a small tree joining its
preimages gives a homotopy model consisting of these
closed oriented surfaces joined along a connected
graph. Choosing a spanning tree gives
\[
 \pi_1(C)\simeq
   \Bigl(*_{i=1}^{c}\pi_1(C_i^{\rm norm})\Bigr)*F_b,
 \qquad
 e(C)=c+1-2\sum_i g_i-b.
\]
The formula also covers unibranch singular points:
their normalization is a local topological
homeomorphism and they contribute no graph edge.

\noindent\textbf{Step 7.} A surface factor of genus at least two cannot occur
in a virtually abelian group, by the finite-cover
argument of P4R.2. Two nontrivial surface factors,
two free generators, or one elliptic factor together
with a free generator contain a nonabelian free
subgroup. Such a subgroup cannot lie in a virtually
abelian group: its intersection with a finite-index
abelian subgroup would have finite index and would
be abelian, whereas a finite-index subgroup of a
nonabelian free group has rank at least two. The
possibilities are therefore all $g_i=0$ with
$b\le1$, or exactly one $g_i=1$ with $b=0$.
The displayed Euler formula gives $e(C)\ge0$.
Additivity of compactly supported Euler numbers and
P4.13 now give
\[
 e(S)=e(S_{\rm reg})+e(C)\ge0.
\]

\noindent\textbf{Step 8.} Finally equality is impossible here. It would force
$e(C)=0$, so the same formula gives $c=1$, $g_1=1$
and $b=0$. Every irreducible component of
$D_{\rm red}$ maps finitely and nonconstantly onto
this curve. A rational vertical component would
therefore give a nonconstant map from $\mathbb P^1$
to its genus-one normalization, contrary to
Riemann--Hurwitz. There are no vertical conductor
components. The initial conductor is consequently
$2H$, and there can have been no blowup: the first
centre would lie on $2H$ and would create a vertical
conductor component of coefficient one. Thus
$T=T_0$ and $D_{\rm red}=H$. It follows that
\[
 e(S_{\rm reg})=e(T\setminus H)
       =(4-4g)-(2-2g)=2-2g<0,
\]
contradicting P4.13. This proves strict positivity.
\end{proof}

P4R.3 proves strict positivity for its stated subclass.
The later classification P4.58 excludes every positive-Euler
ruled case, including this subclass, by the complete local
and global conductor arguments P4.54--P4.57 and P4R.18--P4R.20.
The earlier live status is retained in
\path{records/part4_root_superseded_status.md}.

\begin{proposition}[P4R.4]\label{night:P4R.4}\label{stmt:P4R.4}
\StatementRecord{night:P4R.4}{proved}{Historical events: GPT-6 Astra (OpenAI), ChatGPT Pro mode, effort Pro, S6\_06\_BLIND, 16--18 November 2026: targeted proof check. Also checked in this writing. Pre-continuation writing events: GPT-6 Astra (OpenAI), ChatGPT Pro mode, effort Pro, S6\_10\_PAPERS, 18 November 2026: targeted proof check. Continuation metadata events: GPT-6 Astra (OpenAI), ChatGPT Pro mode, effort Pro, S6\_10\_PAPERS continuation, 19 November 2026: used and rechecked in this run; status-vocabulary reconciliation; metadata reconciliation only; no new mathematical check. \emph{Pre-continuation recorded proof scope.} The night-run re-derivation was by GPT-6 Astra (OpenAI, ChatGPT Pro mode, effort Pro), 16--18 November 2026. Also checked in this writing by GPT-6 Astra (OpenAI), ChatGPT Pro mode, effort Pro, 18 November 2026. The nested open-star and tube maps, punctured radial bundle equivalence and unchanged meridian normal closures give the explicit pi1 isomorphism between support and tube. The complement is Lambda semidirect\_A Z and its filling relation is lambda t\textasciicircum{}m. Rational abelianization has dimension ker(A-I); quotienting by Lambda has order m. A rational-tree conductor makes the support simply connected, hence m=1 and det(I-A) is a positive integer. The trace formula and the tree Euler count then give e(S) at least three}{\ref{night:P4.13}; \ref{night:P4R.1}; \ref{night:P4R.3}}

In the setting of P4R.1, choose a positive small-loop
monodromy $A$ on $\Lambda=H_1(T^4;\mathbb Z)$.
For a suitable $\lambda\in\Lambda$, a positive
meridian of the reduced support is represented by
$\mu=\lambda t^m$ in $\Lambda\rtimes_A\langle t\rangle$.
The natural small-tube comparison gives
\[
 \pi_1(S)\simeq
 (\Lambda\rtimes_A\langle t\rangle)/\langle\!\langle\mu\rangle\!\rangle,
 \qquad
 \pi_1(S)/\operatorname{im}\Lambda\simeq\mathbb Z/m,
 \qquad
 b_1(S)=\dim_{\mathbb Q}\ker(A-I).
\]
For the support in P4R.3, a singular-locus incidence
graph that is a tree of rational normalization
components gives $e(S)\ge3$.
\end{proposition}
\begin{proof}
\noindent\textbf{Step 1.} We strengthen the neighbourhood comparison of P4R.1
without assuming that a preassigned tube is itself
one of the chosen open stars. Choose an open-star
neighbourhood $V_0$ retracting onto $S$, a small tube
$U_0\subset V_0$, a second open-star neighbourhood
$V_1$ with $S\subset V_1\subset U_0$, and a smaller
tube $U_1\subset V_1$. Properness supplies both tube
inclusions. The punctured-disc fibrations give an
isomorphism
\[
 \pi_1(U_1\setminus S)\longrightarrow
 \pi_1(U_0\setminus S):
\]
the fibres can be identified by radial parallel
transport and both base fundamental groups are
generated by the same positive loop. As in P4R.1,
the kernels of the maps to $\pi_1(U_i)$ are the
normal closures of the positive smooth-support
meridian. These meridians correspond under the
displayed isomorphism. Hence
$\pi_1(U_1)\to\pi_1(U_0)$ is an isomorphism.

\noindent\textbf{Step 2.} This isomorphism factors through $\pi_1(V_1)$,
which is identified with $\pi_1(S)$ by its
retraction. Thus $\pi_1(S)\to\pi_1(U_0)$ is
surjective. The restriction to $U_0$ of the
retraction of $V_0$ proves injectivity. This
identifies $\pi_1(S)$ with the exact meridian
quotient, not just with a further quotient.

\noindent\textbf{Step 3.} Projection to the cyclic factor sends $\mu$ to
$m$. After dividing by the image of $\Lambda$,
the remaining presentation is therefore exactly
$\langle t\mid t^m=1\rangle$. In the rational
abelianization the presentation becomes
\[
 \frac{(\Lambda/(A-I)\Lambda)\otimes\mathbb Q
       \ \oplus\ \mathbb Qt}
      {\mathbb Q(\overline\lambda+mt)}.
\]
The single displayed relation has nonzero
$t$ coefficient and reduces the dimension by
one. This proves the asserted formula for $b_1$.

\noindent\textbf{Step 4.} In the rational-tree case of P4R.3 the same
pushout calculation gives $\pi_1(S)=1$.
The cyclic quotient then forces $m=1$, and the
$b_1$ formula gives $\det(I-A)\ne0$.
P4.13 gives
$e(S_{\rm reg})=\det(I-A)\ge1$.
A nonempty connected tree of $c\ge1$ rational
normalization components has Euler number
$c+1\ge2$, including any unibranch singular
points. Additivity gives $e(S)\ge3$.
\end{proof}

\subsection{Hopf normalizations and the conductor graph}

\begin{proposition}[P4R.5]\label{night:P4R.5}\label{stmt:P4R.5}
\StatementRecord{night:P4R.5}{proved}{Historical events: GPT-6 Astra (OpenAI), ChatGPT Pro mode, effort Pro, S6\_06\_BLIND, 16--18 November 2026: targeted proof check. Also checked in this writing. Pre-continuation writing events: GPT-6 Astra (OpenAI), ChatGPT Pro mode, effort Pro, S6\_10\_PAPERS, 18 November 2026: targeted proof check. Continuation metadata events: GPT-6 Astra (OpenAI), ChatGPT Pro mode, effort Pro, S6\_10\_PAPERS continuation, 19 November 2026: used and rechecked in this run; status-vocabulary reconciliation; metadata reconciliation only; no new mathematical check. \emph{Pre-continuation recorded proof scope.} The night-run re-derivation was by GPT-6 Astra (OpenAI, ChatGPT Pro mode, effort Pro), 16--18 November 2026. Also checked in this writing by GPT-6 Astra (OpenAI), ChatGPT Pro mode, effort Pro, 18 November 2026. Adjunction on a minimal Hopf surface and the lift of any rational normalization to C2 minus zero show that every compact curve is smooth elliptic and that distinct curves are disjoint. Tracking the complete pole divisor through blowups gives e(Sreg)=k-v. The disconnected conductor pushout has free quotient of rank b+r-s. Its virtually abelian constraint and the exact curve Euler formula leave only one putative negative configuration; the two-branch hypersurface intersection argument makes that image locally unibranch, which forces b=0 and contradicts the negative count}{\ref{night:P4.12}; \ref{night:P4R.1}; \ref{night:P4R.3}}

In the standing torus-fibration setting of this section, an irreducible singular
support whose normalization is smooth and has a Hopf
surface as its minimal model has nonnegative Euler
number. More precisely, for its reduced conductor
image $C$ one has $e(C)\ge0$. If the normalization is
obtained by $k$ point blowups and $v$ of the exceptional
curves belong to the conductor, then
\[
 e(S_{\rm reg})=k-v,\qquad
 e(S)=k-v+e(C)\ge0.
\]
\end{proposition}
\begin{proof}
\noindent\textbf{Step 1.} Let $T_0$ be the minimal Hopf surface. Every
irreducible curve $E$ on $T_0$ is smooth elliptic.
Indeed $b_2(T_0)=0$ gives $E^2=K_{T_0}\cdot E=0$,
so adjunction gives $p_a(E)=1$. A rational
normalization would define a nonconstant map
$\mathbb P^1\to T_0$. It lifts to the universal
cover $\mathbb C^2\setminus\{0\}$, where its two
coordinate functions are constant. This is
impossible. Thus the normalization genus is one
and the sum of the local delta invariants is zero.
Distinct curves are disjoint, since their
intersection number is zero and every local
intersection of distinct complex curves is
positive.

\noindent\textbf{Step 2.} Push the effective conductor $D$ to $D_0$ on
$T_0$. The canonical comparison of P4.12 makes
$K_{T_0}+D_0$ torsion. The divisor $D_0$ is
nonempty: otherwise a positive power of $K_{T_0}$
would be trivial, contradicting the Kodaira
dimension $-\infty$ of a Hopf surface. Its
reduced support consists of $r\ge1$ disjoint
smooth elliptic curves. The blowup coefficient
$\mu-1$ calculation of P4R.3 applies without a
ruling. At a smooth conductor point it adds a
rational leaf precisely when $\mu\ge2$; at a
node it subdivides the corresponding edge.
Consequently $D_{\rm red}$ has $r$ connected
components, each a smooth elliptic root with a
rational tree attached. With $v$ rational
components it has exactly $v$ double points,
and hence $e(D_{\rm red})=v$. Since
$e(T)=k$, normalization away from the conductor
gives $e(S_{\rm reg})=k-v$. Each rational
component was created by one of the $k$ blowups,
so $k-v\ge0$.

\noindent\textbf{Step 3.} Write $s$ for the number of connected components
of $C$, $c$ for its number of irreducible
components, and $g$ for the number of its
normalizations of genus one. Every normalized
component has genus zero or one: it receives a
finite nonconstant map from a normalized
conductor component. Each elliptic component
must receive an elliptic root, and distinct
image components require distinct roots.
Thus $c\ge g$ and $r\ge g$. Let $b$ be the sum
of the cycle ranks of the point-identification
graphs of the $s$ connected components. The
normalization formula gives
\[
 e(C)=c+s-2g-b.
\]

\noindent\textbf{Step 4.} We retain the disconnected conductor in the
fundamental-group comparison. Its pushout has
one vertex for $T$ and one for each connected
component of $C$, with $r$ edges supplied by
the connected components of $D_{\rm red}$.
This graph is connected, because the finite
map $D_{\rm red}\to C$ is surjective.
The van Kampen presentation has $r-s$ stable
letters, in addition to the vertex groups.
The group of each conductor component is the
group of its elliptic root. Its map into $C$
factors, on that root, through the smooth
normalization of one image component; the
attached rational tree adds no fundamental
group generator. Kill the vertex group
$\pi_1(T)$ and the normal subgroups generated
by all normalized component groups of $C$.
Every conductor-group relation then disappears.
The resulting quotient is exactly
\[
 \pi_1(S)\twoheadrightarrow F_{b+r-s}.
\]
This is also the group of the homotopy-pushout
graph obtained by making the same quotients in
the normalization graph model of $C$. P4R.1
therefore gives $b+r-s\le1$. Substitution yields
\[
 e(C)\ge c+r-2g-1\ge-1.
\]

\noindent\textbf{Step 5.} We exclude the only negative possibility by
the actual local hypersurface geometry.
The value $e(C)<0$ in these inequalities
forces $c=r=g$ and $b+r-s=1$. Every image
component then has elliptic normalization.
No rational conductor component can map
finitely onto one of them, by
Riemann--Hurwitz. Thus $D_{\rm red}$ consists
of disjoint smooth elliptic curves. This
conclusion permits blowups at coefficient-one
points; their exceptional curves are absent
from the conductor and do not affect the
following argument.

\noindent\textbf{Step 6.} At a point $p\in C$, suppose there were two
distinct reduced curve germs $C_1,C_2$ of
$C$. Choose normalization preimages over $p$
whose conductor germs map onto $C_1,C_2$.
They must be distinct preimages, since the
reduced conductor has just one smooth germ
at each of its points. They correspond to
distinct irreducible analytic surface
branches $S_1,S_2$ of $S$ at $p$.
The support $S$ is a hypersurface in the
smooth total threefold. Its distinct
surface branches have a one-dimensional
intersection germ through $p$: two local
nonassociate irreducible equations in the
regular three-dimensional local ring have
height two. That intersection lies in the
conductor image of both branches. The
unique conductor germ on each chosen
normalization branch therefore has the
same image germ, giving $C_1=C_2$, a
contradiction.

\noindent\textbf{Step 7.} Hence $C$ is locally unibranch and its
irreducible components are disjoint. Its
point-identification graphs have $b=0$
and $s=c=g$. The Euler formula gives
$e(C)=0$, contrary to the negative
possibility. We have proved $e(C)\ge0$,
and the asserted formulas finish the proof.
\end{proof}

\subsection{The integral pole divisor on a normalization resolution}

\begin{proposition}[P4R.6]\label{night:P4R.6}\label{stmt:P4R.6}
\StatementRecord{night:P4R.6}{proved}{Historical events: GPT-6 Astra (OpenAI), ChatGPT Pro mode, effort Pro, S6\_06\_BLIND, 16--18 November 2026: targeted proof check. Also checked in this writing. Pre-continuation writing events: GPT-6 Astra (OpenAI), ChatGPT Pro mode, effort Pro, S6\_10\_PAPERS, 18 November 2026: targeted proof check; used and rechecked in this run. Continuation metadata events: GPT-6 Astra (OpenAI), ChatGPT Pro mode, effort Pro, S6\_10\_PAPERS continuation, 19 November 2026: used and rechecked in this run; status-vocabulary reconciliation; metadata reconciliation only; no new mathematical check. \emph{Pre-continuation recorded proof scope.} The night-run re-derivation was by GPT-6 Astra (OpenAI, ChatGPT Pro mode, effort Pro), 16--18 November 2026. Also checked in this writing by GPT-6 Astra (OpenAI), ChatGPT Pro mode, effort Pro, 18 November 2026. Finite Gorenstein duality produces the effective horizontal conductor in the pulled meromorphic canonical divisor. The actual principal-function exceptional intersection equation proves negative definiteness and strict positivity of minus the inverse matrix on each block. Minimal-resolution adjunction gives nonnegative right-hand coefficients, so all exceptional pole coefficients are nonnegative. A zero block is a disjoint rational minus-two configuration outside the conductor. Its isolated hypersurface morsification supplies an embedded vanishing sphere of square minus two in a nearby torus via the normalized root cover, contradicting asphericity. Hence every exceptional block meets the conductor and has positive coefficients}{\ref{night:P4.12}}

Let $S$ be an irreducible singular support in
Part~4.2, let $\nu:N\to S$ be its normalization,
and let $\rho:T\to N\to S$ use the minimal
resolution of $N$. There is an effective integral
divisor $D$ on $T$ such that
\[
 \omega_T(D)\simeq\rho^*\omega_S,
 \qquad \operatorname{Supp}D
       =\rho^{-1}(\operatorname{Sing}S)_{\rm red}.
\]
In particular $K_T+D$ is torsion, by P4.12.
The strict-transform conductor has its usual
positive integral coefficients. Every connected
exceptional configuration over a singular point
of $N$ has strictly positive coefficients in $D$.
Every such configuration meets the strict-transform
conductor.
\end{proposition}
\begin{proof}
\noindent\textbf{Step 1.} The reduced support is a Cartier hypersurface in
the smooth threefold, so $\omega_S$ is invertible.
Choose a local generator $\eta$ of this line
bundle. On $S_{\rm reg}$ it is an ordinary
holomorphic two-form; its pullback is a
meromorphic two-form on $T$. Its divisorial
orders are integers. Two generators differ by
a holomorphic unit on $S$, whose pullback is
still a unit. Consequently their negative
divisors define one integral divisor $D$ and
the displayed line-bundle isomorphism.
Finite duality, in the form
\[
 \nu_*\omega_N
  =\mathcal Hom_{\mathcal O_S}
       (\nu_*\mathcal O_N,\omega_S),
\]
identifies the generic strict-transform pole
ideal with the conductor. This is the same
normalization adjunction used in P4.12. Thus
the coefficients away from the exceptional
curves are the positive integral conductor
exponents; away from conductor and exceptional
curves the pullback has no zero or pole.
No assertion that $K_N$ alone is
$\mathbb Q$-Cartier is needed.

\noindent\textbf{Step 2.} We prove positivity along an exceptional
configuration. Its intersection matrix $M$
is negative definite, and we give the local
comparison establishing the matrix property
used here. Choose a nonzero local function
on the normal surface $N$ vanishing at the
image point. Its pullback divisor is
$H+\sum_i n_iE_i$, with $n_i>0$ and $H$
effective without exceptional components.
The strict transform meets the exceptional
configuration, because a curve through the
image point has closure meeting its proper
inverse image. The divisor is principal
near the configuration, so
\[
 Mn=-h,\qquad h_i=H\cdot E_i\ge0,
 \qquad h\ne0.
\]
The off-diagonal entries of $M$ are
nonnegative. On a connected configuration,
writing a real vector as $v_i=n_iw_i$
gives the identity
\[
 v^tMv
 =-\sum_{i<j}M_{ij}n_in_j(w_i-w_j)^2
   -\sum_i n_ih_iw_i^2.
\]
It is strictly negative for every nonzero
$v$: equality forces all $w_i$ equal along
the connected graph and then forces that
common value to be zero. The same maximum
principle, applied to $M(n_iw_i)=-b$ with
$b_i\ge0$, shows that $-M^{-1}$ has
nonnegative entries, strictly positive on
each connected block. This also follows
directly by taking a negative minimum of
the $w_i$ in the displayed weighted
Laplacian identity.

\noindent\textbf{Step 3.} Write $D=D_h+\sum_i d_iE_i$, with $D_h$
the strict conductor. Minimality of the
resolution gives
\[
 K_T\cdot E_i=2p_a(E_i)-2-E_i^2\ge0.
\]
Indeed the only exception for a negative
curve of arithmetic genus zero would be
a smooth rational exceptional curve of
self-intersection $-1$, which is absent
from a minimal resolution. Also
$D_h\cdot E_i\ge0$. Intersecting
$K_T+D=\rho^*K_S$ with the contracted
curves gives
\[
 Md=-b,
 \qquad b_i=(K_T+D_h)\cdot E_i\ge0.
\]
Hence $d_i\ge0$. A zero coefficient on
a connected block forces $b=0$ there.
Every curve in that block is then a
smooth rational $(-2)$-curve, and the
block is disjoint from the strict
conductor.

\noindent\textbf{Step 4.} In fact, every exceptional block disjoint from
the strict conductor is forbidden, independently
of its coefficients. Finite duality identifies
the conductor on the normal $N$, up to the
invertible pullback of $\omega_S$, with the
rank-one reflexive sheaf $\omega_N$. Thus its
support has no isolated point. Any conductor
curve through the image point has strict
transform meeting the exceptional block:
properness makes the closure of its punctured
inverse image meet the fibre. Therefore the
image of a block disjoint from the strict
conductor is away from the conductor, and
$\nu$ is an isomorphism near it. The
corresponding point of $S$ is a normal isolated
hypersurface singularity. The local
Milnor-fibre argument of P4.12 applies
at this point itself. A local reduced
equation $g$ has an isolated critical
point and positive Milnor number. A
Morse perturbation supplies a vanishing
two-sphere of square $-2$ in its Milnor
fibre. The original map is a unit times
$g^m$; taking a local root of the unit,
each of its nearby local fibres is a
disjoint union of copies of a Milnor
fibre of $g$. Such a copy embeds into
the nearby four-torus. Every map
$S^2\to T^4$ is null-homotopic, so that
embedded sphere would have zero
homology class and square zero, a
contradiction. Thus every exceptional
block has positive pole coefficients.
This proves both effectiveness and the
asserted support equality.
\end{proof}

\begin{proposition}[P4R.7]\label{night:P4R.7}\label{stmt:P4R.7}
\StatementRecord{night:P4R.7}{proved}{Historical events: GPT-6 Astra (OpenAI), ChatGPT Pro mode, effort Pro, S6\_06\_BLIND, 16--18 November 2026: targeted proof check. Also checked in this writing. Pre-continuation writing events: GPT-6 Astra (OpenAI), ChatGPT Pro mode, effort Pro, S6\_10\_PAPERS, 18 November 2026: targeted proof check. Continuation metadata events: GPT-6 Astra (OpenAI), ChatGPT Pro mode, effort Pro, S6\_10\_PAPERS continuation, 19 November 2026: used and rechecked in this run; status-vocabulary reconciliation; metadata reconciliation only; no new mathematical check. \emph{Pre-continuation recorded proof scope.} The night-run re-derivation was by GPT-6 Astra (OpenAI, ChatGPT Pro mode, effort Pro), 16--18 November 2026. Also checked in this writing by GPT-6 Astra (OpenAI), ChatGPT Pro mode, effort Pro, 18 November 2026. The degree-two minimal ruled pole divisor, preservation of its SNC tree under blowup, the pushout fundamental group, and the normalization-trace Euler comparison were re-derived}{\ref{night:P4.13}; \ref{night:P4R.1}; \ref{night:P4R.3}; \ref{night:P4R.6}}

In Part~4.2, every irreducible singular
support whose normalization has a
resolution birationally ruled over a
curve of genus at least two has
nonnegative Euler number. The
normalization may be singular.
\end{proposition}
\begin{proof}
\noindent\textbf{Step 1.} Use the minimal resolution and effective
integral divisor of P4R.6. Push this
divisor to the minimal ruled model
$T_0$. Its intersection with a ruling
fibre is two. The adjunction and
Riemann--Hurwitz calculation in P4R.3
therefore gives exactly
\[
 D_0=2H+\sum_j a_jF_j,
 \qquad a_j\in\mathbb Z_{>0},
\]
with $H$ a smooth section. The reverse
blowup coefficient is the integral
number $\mu-1$. At a conductor crossing
one has $\mu\ge2$, so its exceptional
curve remains in the support and
subdivides the edge. At a smooth
conductor point the operation either
adds a rational leaf or leaves the
reduced curve unchanged. Thus
$A=\operatorname{Supp}D$ is connected,
is $H$ with rational trees attached,
and $\pi_1(A)\to\pi_1(T)$ is an
isomorphism.

\noindent\textbf{Step 2.} Put $C=\operatorname{Sing}(S)_{\rm red}$.
The map $\rho$ is proper and is an
isomorphism outside $A$, and P4R.6
identifies $A$ with its full inverse
image of $C$. The continuous map
$T\cup_A C\to S$ is therefore a
compact-to-Hausdorff bijection and a
homeomorphism. The same cofibration
and van Kampen argument as P4R.3 gives
$\pi_1(S)\simeq\pi_1(C)$. This retains
the exceptional components that map
to points, as well as the strict
conductor components that map to
curves. P4R.1 and the normalization
graph calculation in P4R.3 give
$e(C)\ge0$. Finally P4.13 gives
$e(S_{\rm reg})\ge0$, so
$e(S)=e(S_{\rm reg})+e(C)\ge0$.
\end{proof}

\begin{proposition}[P4R.8]\label{night:P4R.8}\label{stmt:P4R.8}
\StatementRecord{night:P4R.8}{proved}{Historical events: GPT-6 Astra (OpenAI), ChatGPT Pro mode, effort Pro, S6\_06\_BLIND, 16--18 November 2026: targeted proof check. Also checked in this writing. Pre-continuation writing events: GPT-6 Astra (OpenAI), ChatGPT Pro mode, effort Pro, S6\_10\_PAPERS, 18 November 2026: targeted proof check. Continuation metadata events: GPT-6 Astra (OpenAI), ChatGPT Pro mode, effort Pro, S6\_10\_PAPERS continuation, 19 November 2026: used and rechecked in this run; status-vocabulary reconciliation; metadata reconciliation only; no new mathematical check. \emph{Pre-continuation recorded proof scope.} The night-run re-derivation was by GPT-6 Astra (OpenAI, ChatGPT Pro mode, effort Pro), 16--18 November 2026. Also checked in this writing by GPT-6 Astra (OpenAI), ChatGPT Pro mode, effort Pro, 18 November 2026. The Hopf rooted-tree count, disconnected conductor pushout, free quotient rank b+r-s, and the equality case with contracted rational subtrees were re-derived. The hypersurface-branch intersection argument forces the image curve to be unibranch in the only potentially negative case}{\ref{night:P4R.5}; \ref{night:P4R.6}; \ref{night:P4R.7}}

In Part~4.2, every irreducible singular support
whose normalization has a resolution with Hopf
minimal model has nonnegative Euler number.
The normalization may be singular. For the
minimal normalization resolution, let $k$ be
the number of blowups from its Hopf minimal
model and $v$ the number of rational components
of the pole divisor. Then
\[
 e(S_{\rm reg})=k-v\ge0,\qquad
 e(\operatorname{Sing}S)\ge0.
\]
\end{proposition}
\begin{proof}
\noindent\textbf{Step 1.} Use $\rho:T\to N\to S$ and the effective
integral divisor $D$ of P4R.6. Put
$A=\operatorname{Supp}D$ and
$C=\operatorname{Sing}(S)_{\rm red}$.
The Hopf-curve and blowup calculations of
P4R.5 use only integrality, effectiveness,
and $K_T+D\equiv0$. They therefore apply
to this divisor: $A$ has $r\ge1$ connected
components, each a smooth elliptic root
with rational trees attached. Each of its
$v$ rational components came from one of
the $k$ blowups; $e(A)=v$, $e(T)=k$, and
$\rho$ is an isomorphism off $A$. This
proves the displayed formula for the
regular locus.

\noindent\textbf{Step 2.} Every one-dimensional component of $C$
receives a nonconstant proper map from a
component of $A$, so its normalization
has genus zero or one. There are no
isolated points of $C$: an isolated
singular point of the hypersurface is
normal by the hypersurface $S_2$ property
and regularity in codimension one, and
P4R.6's vanishing-sphere argument excludes
it. Let $s,c,g,b$ respectively count the
connected components of $C$, its
irreducible components, its genus-one
normalizations, and the cycle ranks of
its point-identification graphs.
Then $e(C)=c+s-2g-b$, $r\ge g$, and
$c\ge g$.

\noindent\textbf{Step 3.} The proper-map pushout $T\cup_A C=S$
and its cofibration justification are
exactly those of P4R.7. A root which
maps to a point gives the zero
homomorphism on its fundamental group;
a noncontracted root maps through the
normalization of one component of $C$.
Every rational tree is simply connected.
Consequently killing $\pi_1(T)$ and
the normalized component groups of $C$
in the disconnected van Kampen
presentation leaves the free quotient
$F_{b+r-s}$, as in P4R.5. Virtual
abelianness gives $b+r-s\le1$, hence
\[
 e(C)\ge c+r-2g-1\ge-1.
\]

\noindent\textbf{Step 4.} A negative value forces $c=r=g$.
Every root is then noncontracted, maps
onto an elliptic normalization, and
supplies one of the $c$ image components.
Every rational component of $A$ is
contracted: a nonconstant map from its
normalization $\mathbb P^1$ to an
elliptic curve contradicts
Riemann--Hurwitz. Thus every exceptional
block is a rational subtree of one
rooted component of $A$. It meets that
root in exactly one point. At least one
intersection is forced by the final
assertion of P4R.6, since the roots are
the only strict conductor components;
two intersections would make a cycle
in the rooted tree. The exceptional
fibre is connected. One can see this
directly by Stein factorization of
$T\to N$: its finite birational factor
is an isomorphism because $N$ is normal.

\noindent\textbf{Step 5.} The conductor on $N$ consequently has
one local analytic branch at each of
its points. At a singular point of $N$,
each local branch has a strict transform
meeting the exceptional fibre, whereas
there is exactly one such intersection
with the smooth root. Away from those
fibres the assertion follows from the
smooth disjoint roots themselves.

\noindent\textbf{Step 6.} The local hypersurface argument of
P4R.5 now applies with $N$ in place of
its smooth resolution. Two different
curve germs of $C$ would require two
different normalization points, because
each point has just one conductor germ.
Their corresponding hypersurface
branches intersect in a curve germ.
That germ belongs to both conductor
images and forces the two chosen germs
to coincide. Thus $C$ is locally
unibranch with disjoint irreducible
components. This gives $b=0$, $s=c=g$,
and $e(C)=0$, a contradiction.
Therefore $e(C)\ge0$, proving the
assertion.
\end{proof}

\begin{proposition}[P4R.9]\label{night:P4R.9}\label{stmt:P4R.9}
\StatementRecord{night:P4R.9}{proved}{Historical events: GPT-6 Astra (OpenAI), ChatGPT Pro mode, effort Pro, S6\_06\_BLIND, 16--18 November 2026: targeted proof check. Also checked in this writing. Pre-continuation writing events: GPT-6 Astra (OpenAI), ChatGPT Pro mode, effort Pro, S6\_10\_PAPERS, 18 November 2026: targeted proof check. Continuation metadata events: GPT-6 Astra (OpenAI), ChatGPT Pro mode, effort Pro, S6\_10\_PAPERS continuation, 19 November 2026: used and rechecked in this run; status-vocabulary reconciliation; metadata reconciliation only; no new mathematical check. \emph{Pre-continuation recorded proof scope.} The night-run re-derivation was by GPT-6 Astra (OpenAI, ChatGPT Pro mode, effort Pro), 16--18 November 2026. Also checked in this writing by GPT-6 Astra (OpenAI), ChatGPT Pro mode, effort Pro, 18 November 2026. Adjunction on the reduced horizontal pole divisor forces equality in every degree inequality; the etale bisection or two disjoint sections, absence of vertical poles, and absence of resolution or nonminimal exceptional curves were re-derived}{\ref{night:P4.13}; \ref{night:P4.22}; \ref{night:P4R.6}; \ref{night:P4R.7}}

In Part~4.2, every irreducible singular
support whose normalization has a
resolution birationally ruled over an
elliptic curve has nonnegative Euler
number. When the horizontal part of
the pole divisor on the minimal ruled
model is reduced, the normalization
is smooth and minimal and the support
has Euler number zero.
\end{proposition}
\begin{proof}
\noindent\textbf{Step 1.} Push the integral divisor of P4R.6 to
the minimal ruled surface $T_0\to B$.
Its horizontal degree is two. A reduced
horizontal component $H$ satisfies
\[
 0\le 2g(\widetilde H)-2
 \le 2p_a(H)-2
 =-(D_0-H)H\le0.
\]
The first inequality is
Riemann--Hurwitz over the elliptic
base. All inequalities are equalities:
every such component is smooth elliptic,
unramified over $B$, and disjoint from
the other components of $D_0$. Every
vertical fibre meets the horizontal
part, so no vertical component occurs.
The reduced-horizontal possibilities
are therefore one smooth elliptic
bisection or two disjoint elliptic
sections, each with coefficient one.

\noindent\textbf{Step 2.} The reverse blowup formula $\mu-1$
then keeps $A=\operatorname{Supp}D$
a disjoint union of their smooth
strict transforms; each new exceptional
curve is absent from $A$. Every
exceptional curve of the normalization
resolution must belong to $A$, by
P4R.6. An exceptional component would
therefore be one of these disjoint
elliptic curves and would be disjoint
from the strict conductor, contrary
to P4R.6. Thus the normalization is
smooth. P4.22, with its hypotheses now
verified, proves minimality and
$e(S)=0$.

\noindent\textbf{Step 3.} In the remaining case
$D_0=2H+\sum a_jF_j$ with $H$ a
smooth section. The same reverse
blowups make $A$ the elliptic root $H$
with rational trees attached. Its
inclusion in $T$ induces an
isomorphism on fundamental groups,
by the ruling and invariance under
point blowups. The resolution pushout
gives $\pi_1(S)=\pi_1(C)$ as in
P4R.7. The virtually abelian curve
group calculation gives $e(C)\ge0$,
and P4.13 gives $e(S_{\rm reg})\ge0$.
Their sum proves $e(S)\ge0$ also
in this remaining case.
\end{proof}
\subsection{The Hessian line and the tacnodal equality case}

\begin{proposition}[P4R.10]\label{night:P4R.10}\label{stmt:P4R.10}
\StatementRecord{night:P4R.10}{proved}{Historical events: GPT-6 Astra (OpenAI), ChatGPT Pro mode, effort Pro, S6\_06\_BLIND, 16--18 November 2026: targeted proof check. Also checked in this writing. Pre-continuation writing events: GPT-6 Astra (OpenAI), ChatGPT Pro mode, effort Pro, S6\_10\_PAPERS, 18 November 2026 05:10:56.410177+00:00: computation or formal certificate; targeted proof check. Continuation metadata events: GPT-6 Astra (OpenAI), ChatGPT Pro mode, effort Pro, S6\_10\_PAPERS continuation, 19 November 2026: used and rechecked in this run; status-vocabulary reconciliation; metadata reconciliation only; no new mathematical check. \emph{Pre-continuation recorded proof scope.} The night-run re-derivation was by GPT-6 Astra (OpenAI, ChatGPT Pro mode, effort Pro), 16--18 November 2026. Also checked in this writing by GPT-6 Astra (OpenAI), ChatGPT Pro mode, effort Pro, 18 November 2026. proved. Full elliptic-primary A3 conductor Euler inequality, exact differentiated morphisms and equality case. Finite duality makes the conductor curve Cohen–Macaulay; generic containment and absence of a zero-dimensional submodule give P\textasciicircum{}(2)\ensuremath{\subset}I. Its nilradical J=P/I is torsionfree rank1. Differentiation embeds P/P\textasciicircum{}(2) into the normalized kernel K of j*OmegaY\ensuremath{\to}OmegaE, with degK=degR. The actual Hessian is a section of L\textsuperscript{2}\ensuremath{\otimes}j*O(S); the latter factor has degree0, so 2degL=sum d. Projection to K/L has kernel exactly I/P\textasciicircum{}(2), determined generically and extended by torsionfreeness. Its branchwise and descent colengths give the displayed exact Euler formula. For a nonplanar branch, all gradient minors have order\ensuremath{\geq}2m; writing a+b=tau and extracting a unit column of t\textasciicircum{}d ell\textsuperscript{2} gives d\ensuremath{\geq}a, hence d/2+b\ensuremath{\geq}m>m−1. For a planar branch, the exact generators dz,dh give E=<dz,t\textasciicircum{}a u>, a\ensuremath{\geq}m; the symmetric-square condition yields d/2+b\ensuremath{\geq}a. Equality of the global Euler number forces delta=0, then R=0 and all remaining nonnegative lengths vanish, so the curve is smooth and the Hessian nowhere zero. The bounded tacnode\_hessian replay is retained only as a supplement to this uniform proof. Review provenance: New derivation after reading the supplied proof; not a blind review and not an independent human check}{\ref{night:P4.29}; \ref{night:P4.30}}

Let $S$ be an irreducible reduced fibre support in Part~4.2.
Suppose that its reduced singular curve $C$ is irreducible and
unibranch at every point, that its normalization $j:E\to C\hookrightarrow Y$
is elliptic, and that its generic transverse germ is a tacnode
$x^2-y^4=0$. Let $C_{\mathrm{sch}}$ be the image conductor scheme.
Then
\[
 \chi(\mathcal O_{C_{\mathrm{sch}}})\leq0.
\]
Equality makes $C$ smooth and makes the transverse Hessian of
the reduced defining equation have rank one at every point of $C$.
\end{proposition}
\begin{proof}
\noindent\textbf{Step 1.} Let $P\subset\mathcal O_Y$ be the reduced ideal of $C$, and let
$I\subset\mathcal O_Y$ be the inverse image of the conductor ideal
of $S$. The conductor depth argument of P4.29 gives a
Cohen--Macaulay quotient $\mathcal O_Y/I$ of pure dimension one.
At the generic tacnode the normalization has two smooth branches;
its conductor is $(x,y^2)$ in the displayed coordinates.
Consequently $I$ contains $P^{(2)}$. Indeed a section of $P^{(2)}$
has zero image in $\mathcal O_Y/I$ at the generic point, and any
nonzero image would generate a finite-length submodule, contrary
to depth one. We obtain
\[
 0\longrightarrow J:=P/I\longrightarrow
 \mathcal O_{C_{\mathrm{sch}}}\longrightarrow\mathcal O_C
 \longrightarrow0.
\]
Here $J^2=0$, because $P^2\subset P^{(2)}\subset I$.
The module $J$ is therefore an $\mathcal O_C$-module. It is
torsionfree of rank one: a finite-support torsion submodule would
again be a finite-length submodule of $\mathcal O_Y/I$.

\noindent\textbf{Step 2.} On $E$, put
\[
 K=\ker\bigl(j^*\Omega_Y\longrightarrow\Omega_E\bigr),
 \qquad\operatorname{im}(dj)=\Omega_E(-R).
\]
These are a rank-two bundle and a line bundle respectively.
The coefficient of $R$ at a point with local parametrization
$x_i=x_i(t)$ is $r=\min_i\operatorname{ord}_t x_i'(t)$.
We have $\deg K=\deg R$, since $\deg j^*\Omega_Y=0$
and $E$ is elliptic. Differentiation identifies
$G=P/P^{(2)}$ with a subsheaf of $j_*K$, as in P4.30.
Write $E_0\subset K$ for the image of its pullback to $E$.
This is a rank-two locally free subsheaf with finite-length
quotient.

\noindent\textbf{Step 3.} The Hessian of a local reduced equation $f$ of $S$ annihilates
the tangent derivative of $C$, since all first derivatives of
$f$ vanish on $C$. It thus defines a quadratic tensor in
$\operatorname{Sym}^2K$. Under a change $f\mapsto u f$ its
restriction to $C$ multiplies by $u$, so globally it is a section
\[
 q\in H^0\bigl(E,\operatorname{Sym}^2K\otimes j^*\mathcal O_Y(S)\bigr).
\]
It has generic rank one. Let $L\subset K$ be the saturation of
its generic image line. On the smooth curve $E$ this is a line
subbundle. The tensor lies in $L^{\otimes2}\otimes
j^*\mathcal O_Y(S)$ everywhere, since that is a subbundle and
the assertion holds on a dense open set. Write its zero orders
as $d_p\geq0$. The normal line of $S$ has degree zero on $E$:
locally over the base disc the divisor $mS$ is principal.
It follows that
\[
 2\deg L=\sum_p d_p.
\]

\noindent\textbf{Step 4.} Generically $I/P^{(2)}$ is exactly the line of the linear
conductor equation $x$, which is the Hessian line. Projection
to $K/L$ therefore gives an injection
\[
 J\longrightarrow j_*(K/L).
\]
To check the kernel at the exceptional points, a section whose
projection vanishes generically has torsion image in $J$; the
torsionfreeness just proved makes that image zero. Let $b_p$ be
the colength of the projected lattice $E_0\to K/L$ at $p$.
Let $q_p$ be the additional nonnegative length between $J$ and
the direct image of this normalized image lattice. Riemann--Roch
on $E$ now gives the exact formula
\[
 \chi(J)=\deg R-\frac12\sum_p d_p-\sum_p b_p-\sum_p q_p.
\]
We prove locally that $d_p/2+b_p\geq r_p$.

\noindent\textbf{Step 5.} At a smooth point of $C$, $r_p=0$, so the claim follows from
nonnegativity. At a singular point put
$m=\min_i\operatorname{ord}_t x_i(t)$, so $r_p=m-1$.
First suppose the embedding dimension is three. Every defining
equation of $C$ has no linear term. Its gradient along the
parametrization has all entries of order at least $m$.
Every two-by-two minor of the conormal image consequently has
order at least $2m$. Since $K$ is saturated in the ambient free
rank-three module, a generator of $\bigwedge^2K$ has a unit
Pl\"ucker coordinate. Thus
\[
 \tau_p:=\operatorname{length}(K/E_0)_p\geq2m.
\]
Let $a_p=\operatorname{length}(L/(E_0\cap L))_p$.
The exact lattice sequence gives $\tau_p=a_p+b_p$.
Choose a primitive local generator $\ell$ of $L$ and write
$q=h\ell^2$ after trivializing the normal line, so
$\operatorname{ord}_t h=d_p$.
Every column of the Hessian is the differential of one of
the first derivatives of $f$. Those first derivatives belong
to $P$, so the columns lie in $E_0$. At least one ambient
coordinate of the primitive vector $\ell$ is a unit; its
corresponding column shows that $h\ell\in E_0$.
Therefore $d_p\geq a_p$, and
\[
 \frac{d_p}{2}+b_p\geq\frac{a_p}{2}+b_p
 \geq\frac{\tau_p}{2}\geq m>m-1=r_p.
\]

\noindent\textbf{Step 6.} Now suppose the embedding dimension is two. Choose coordinates
with $P=(z,H(x,y))$, where $H$ is an irreducible reduced plane
equation with no linear term. Here $P^{(2)}=P^2$.
For an explicit verification, $z,H$ is a regular sequence, and
$P/P^2$ is free over the domain $\mathcal O_Y/P$.
If $s\notin P$ and $sv\in P^2$, reduction modulo $P$ gives
$v\in P$; reduction in this free conormal module then gives
$v\in P^2$. Thus $P^2$ is $P$-primary and equals its symbolic
square. The equation $f$ belongs to this square, so
\[
 f=A z^2+B zH+C H^2.
\]
Let $u$ be a primitive generator of the plane kernel of the
derivative row. We have $dH=w u$, with
$a=\operatorname{ord}_t w\geq m$, because both plane entries
of $dH$ have order at least $m$. In the basis $(dz,u)$ of $K$
we consequently have
\[
 E_0=\langle dz,w u\rangle,
 \qquad q=2A\,dz^2+2Bw\,dz\,u+2Cw^2u^2.
\]
Write the primitive generator of $L$ as
$\ell=\ell_z dz+\ell_u u$, and again $q=h\ell^2$ with
$\operatorname{ord}_t h=d_p$.
When $\ell_z$ is a unit, the projected conormal lattice has
colength $b_p=\min(a,\operatorname{ord}_t\ell_u)$.
The coefficient of $u^2$ in the displayed tensor gives
\[
 d_p+2\operatorname{ord}_t\ell_u\geq2a,
 \qquad d_p/2+b_p\geq a\geq m.
\]
The same inequality holds when $\ell_u=0$, since then $b_p=a$.
When $\ell_z$ is not a unit, $\ell_u$ is a unit. Projection of
$dz$ generates $K/L$, so $b_p=0$; the $u^2$ coefficient gives
$d_p\geq2a$, proving the same bound. These cases cover all
singular curve points and establish the claimed local inequality.

\noindent\textbf{Step 7.} We conclude that $\chi(J)\leq0$. Since normalization gives
$\chi(\mathcal O_C)=-\sum\delta(C,p)$, the conductor exact
sequence yields
\[
 \chi(\mathcal O_{C_{\mathrm{sch}}})
 =-\sum_p\delta(C,p)+\chi(J)\leq0.
\]
Equality forces every delta invariant to vanish, so $C=E$ is
smooth. Then $R=0$, $E_0=K$, and all normalization descent
defects vanish. The exact formula reduces to
$\chi(J)=-\tfrac12\sum d_p$, so equality also forces
$d_p=0$ everywhere. The Hessian has rank one at every point.
The program \texttt{code/part4/tacnode\_hessian.py} verifies the
two lattice inequalities for the stated bounds, the exact
rank-one Hessian on the singular parametrization
$(t^2,t^3,0)$, and the simple Hessian zero of the permitted
local model $zx^2+xy^2$. Its JSON and text outputs retain
the bounds; the preceding inequalities prove the unrestricted
statement.
\end{proof}

\begin{proposition}[P4R.11]\label{night:P4R.11}\label{stmt:P4R.11}
\StatementRecord{night:P4R.11}{proved}{Historical events: GPT-6 Astra (OpenAI), ChatGPT Pro mode, effort Pro, S6\_06\_BLIND, 16--18 November 2026: targeted proof check. Also checked in this writing. Pre-continuation writing events: GPT-6 Astra (OpenAI), ChatGPT Pro mode, effort Pro, S6\_10\_PAPERS, 18 November 2026 05:10:56.410177+00:00: computation or formal certificate; targeted proof check. Continuation metadata events: GPT-6 Astra (OpenAI), ChatGPT Pro mode, effort Pro, S6\_10\_PAPERS continuation, 19 November 2026: used and rechecked in this run; status-vocabulary reconciliation; metadata reconciliation only; no new mathematical check. \emph{Pre-continuation recorded proof scope.} The night-run re-derivation was by GPT-6 Astra (OpenAI, ChatGPT Pro mode, effort Pro), 16--18 November 2026. Also checked in this writing by GPT-6 Astra (OpenAI), ChatGPT Pro mode, effort Pro, 18 November 2026. proved. Full zero-Euler double-horizontal A3 normalization, degree-two conductor and product-tacnode conclusion. Followed the pole-tree pushout to identify the reduced conductor normalization as one elliptic curve. Every rational pole component contracts, so H is its only strict preimage and has degree2. The actual ruling descends through the normalized graph because exceptional fibres are vertical; its Leray sequence gives R1rho O=0. Finite conductor duality and torsion-line Riemann–Roch give \ensuremath{\chi}OCsch=g−1. The independently rederived R10 bound forces g=1, smooth conductor and a nowhere-zero rank-one Hessian; Riemann–Hurwitz makes H\ensuremath{\to}C an unramified double cover. Distinct H points cannot merge under a vertical exceptional fibre. Thus the two local hypersurface factors are smooth graphs in the nondegenerate Morse variable. Their difference is a unit times the square of a coordinate defining the smooth conductor, since the generic intersection order is2. The exact unit-root map gives X\textsuperscript{2}−Y4; its two finite normalization maps are smooth. The remaining normalization and minimality inputs retain their stated R6/P4.22 scopes. Review provenance: New derivation after reading the supplied proof; not a blind review and not an independent human check}{\ref{night:P4.12}; \ref{night:P4.22}; \ref{night:P4.29}; \ref{night:P4R.10}; \ref{night:P4R.6}; \ref{night:P4R.9}}

Let $S$ be an irreducible singular support in Part~4.2 whose
normalization has ruled resolution, and suppose $e(S)=0$.
In the double-horizontal-divisor case with generic transverse
tacnode, the normalization is a smooth minimal elliptic ruled
surface. The conductor image is a smooth elliptic curve, its
reduced preimage is an unramified covering of degree two,
and every local germ is a product tacnode.
\end{proposition}
\begin{proof}
\noindent\textbf{Step 1.} Use $T\to N\to S$, the horizontal section $H$, and
$C=\operatorname{Sing}(S)_{\mathrm{red}}$ as in P4R.9.
The fundamental-group and Euler argument of P4.29 gives an
irreducible curve $C$ with elliptic normalization, unibranch
at every point. All rational curves in the resolved conductor
are contracted, since none maps nonconstantly to an elliptic
curve. Thus $H$ is the only strict conductor component.
Its finite map to $C$ has degree two, and its genus $g$ is at
least one by Riemann--Hurwitz.

\noindent\textbf{Step 2.} Every exceptional curve of $T\to N$ is vertical for the ruling
$T\to H$. The normalized-graph construction of P4.29 descends
this ruling to $N\to H$. Its Leray calculation proves
$R^1\rho_{N*}\mathcal O_T=0$, and hence
$\chi(\mathcal O_N)=1-g$. Finite duality identifies the
conductor on $N$ with $\omega_N\otimes\nu^*\omega_S^{-1}$.
The latter line twist is torsion by P4.12. Serre duality,
the rational-singularity comparison and Riemann--Roch on $T$
therefore give $\chi(I)=1-g$ and
$\chi(\mathcal O_{D_{\mathrm{sch}}})=0$.
The exact conductor square, with $\chi(\mathcal O_S)=0$,
then gives
\[
 \chi(\mathcal O_{C_{\mathrm{sch}}})=g-1.
\]
P4R.10 makes this at most zero. Thus $g=1$, $C$ is smooth,
and the reduced defining equation has rank-one Hessian along
all of $C$. Riemann--Hurwitz makes $H\to C$ unramified.

\noindent\textbf{Step 3.} We verify the local analytic conclusion. Over any point of
$C$, its two preimages on $H$ cannot be identified in $N$:
each exceptional block is vertical for the descended ruling
and hence meets only one point of the section $H$.
There are therefore two local irreducible branches of $S$.
The holomorphic Morse lemma with parameters writes its reduced
equation as $X^2+g(y,z)$. Its two analytic factors are monic
linear polynomials in $X$, by Weierstrass factorization, and
their coefficients sum to zero. Each branch is a graph and
is smooth. Their intersection has support exactly the smooth
curve $C$. Consequently, in coordinates with $C=\{X=y=0\}$,
the factor coefficient is a unit times a power of $y$.
The generic tacnode fixes that power to two. Extracting a square
root of the unit in the $y$ coordinate gives
$X^2-y^4=0$ times the $z$ disc. This also verifies the actual
normalization and conductor morphisms locally.
The normalization was already smooth away from the conductor
by P4R.6. It is therefore smooth everywhere, and P4.22 makes
this elliptic ruled surface minimal.
\end{proof}

\subsection{The components outside the pole divisor in a ruled fibre}

\begin{proposition}[P4R.12]\label{night:P4R.12}\label{stmt:P4R.12}
\StatementRecord{night:P4R.12}{proved}{Historical events: GPT-6 Astra (OpenAI), ChatGPT Pro mode, effort Pro, S6\_06\_BLIND, 16--18 November 2026: targeted proof check. Also checked in this writing. Pre-continuation writing events: GPT-6 Astra (OpenAI), ChatGPT Pro mode, effort Pro, S6\_10\_PAPERS, 18 November 2026 05:10:56.410177+00:00: computation or formal certificate; targeted proof check. Continuation metadata events: GPT-6 Astra (OpenAI), ChatGPT Pro mode, effort Pro, S6\_10\_PAPERS continuation, 19 November 2026: used and rechecked in this run; status-vocabulary reconciliation; metadata reconciliation only; no new mathematical check. \emph{Pre-continuation recorded proof scope.} The night-run re-derivation was by GPT-6 Astra (OpenAI, ChatGPT Pro mode, effort Pro), 16--18 November 2026. Also checked in this writing by GPT-6 Astra (OpenAI), ChatGPT Pro mode, effort Pro, 18 November 2026. proved. Full disjoint-leaf description of fibre components outside the pole support on the actual ruled blowup model. On the actual point-blown-up ruled model, every vertical component E has negative self-intersection by the full multiplicity equation m\_E E\textsuperscript{2}+sum m\_j(E\_j E)=0. If E lies outside D, then DE=2+E\textsuperscript{2}\ensuremath{\geq}0, so E\textsuperscript{2} is−1 or−2. The latter gives DE=0 and is disjoint from D; the resolution is an isomorphism there by R6. Removing the fibre multiplicity is unramified on E, and its finite etale pullback consists of P1s. Actual smooth-family transport places its normal square−2 sphere in a nearby T4, impossible. Thus E\textsuperscript{2}=−1 and DE=1, giving one transverse intersection with a coefficient-one pole component. Two outside components cannot be adjacent: the positive full-fibre multiplicity relations would give both m\_E>m\_Eprime and m\_Eprime>m\_E. Therefore all such components are pairwise disjoint leaves. No multiplicity-one assumption was used. Review provenance: New derivation after reading the supplied proof; not a blind review and not an independent human check}{\ref{night:P4R.3}; \ref{night:P4R.6}; \ref{night:P4R.9}}

Let $T\to N\to S$ be the minimal normalization resolution of an
irreducible support in Part~4.2 with ruled minimal model of positive
base genus. Let $r:T\to B_0$ be its ruling, and let $D$ be the
effective integral pole divisor of P4R.6. Every component $E$ of a
reducible ruling fibre which is absent from $\operatorname{Supp}D$
is a smooth rational curve of self-intersection $-1$. It meets
exactly one pole component, transversely once, and that component
has pole coefficient one. In the double-horizontal-section case,
$E$ is disjoint from the horizontal section. Distinct such
components are disjoint; each is a leaf of the reduced ruling fibre.
These assertions retain the multiplicities of the ruling fibre
and do not require the multiplicity of $E$ to be one.
\end{proposition}
\begin{proof}
\noindent\textbf{Step 1.} A reduced horizontal pole divisor occurs here only over an
elliptic base, by P4R.3, and P4R.9 then makes the normalization
smooth and minimal. In that case there is no reducible ruling
fibre. We may therefore work in the double-horizontal-section
case throughout the proof.
A ruling fibre on $T$ is obtained from a smooth $\mathbb P^1$
by the actual sequence of point blowups. Its reduced support is
a normal-crossing tree of smooth rational curves. Each component
of a reducible fibre has negative self-intersection. One can
verify the latter assertion without a contraction criterion:
if $F_b=\sum m_iE_i$ is the full fibre, then
$0=F_bE_i=m_iE_i^2+\sum_{j\ne i}m_jE_jE_i$;
connectedness and reducibility make the last sum positive.

\noindent\textbf{Step 2.} For $E$ outside the pole support, effectiveness and adjunction give
\[
 0\le DE=-K_TE=2+E^2.
\]
Thus $E^2$ is $-1$ or $-2$. In the second case $DE=0$,
so $E$ is disjoint from the full pole support. The map to $S$
then identifies a neighbourhood of this compact curve with a
neighbourhood in $S_{\rm reg}$, by P4R.6. The local cyclic
base change which removes the multiplicity of the support is
unramified over this regular locus. Its restriction to $E$ is
a disjoint union of copies of $\mathbb P^1$, because
$\pi_1(\mathbb P^1)=0$. A chosen copy has the same normal
degree $-2$. A smooth local trivialization over a sufficiently
small parameter disc transports it into a nearby four-torus,
still with self-intersection $-2$. Every map from $S^2$ to that
torus is null-homotopic, since its universal cover is
$\mathbb R^4$; its homology class and self-intersection must
therefore be zero. This contradiction leaves $E^2=-1$ and
$DE=1$.

\noindent\textbf{Step 3.} All pole coefficients and all intersection multiplicities are
positive integers at an intersection. Hence $DE=1$ says exactly
that there is one intersection, transverse and on a coefficient-one
component. A horizontal pole section of coefficient two cannot
be that component.

\noindent\textbf{Step 4.} It remains to exclude an intersection of two components $E,E'$
outside the pole divisor. Both have square $-1$, and both meet
the pole support. In the full ruling fibre, let their positive
multiplicities be $m,m'$. The two fibre-intersection equations
would give
\[
 m=m'+a,\qquad m'=m+b,
\]
where $a,b$ are positive sums of the multiplicities of their
other vertical neighbours; their respective pole neighbours make both
sums positive. Adding these equations gives $a+b=0$, a
contradiction. Thus no two outside components meet. Each has
exactly the one pole neighbour already found and is a leaf.
\end{proof}

\subsection{Restriction to an elliptic conductor component}

\begin{proposition}[P4R.13]\label{night:P4R.13}\label{stmt:P4R.13}
\StatementRecord{night:P4R.13}{proved}{Historical events: GPT-6 Astra (OpenAI), ChatGPT Pro mode, effort Pro, S6\_06\_BLIND, 16--18 November 2026: targeted proof check. Also checked in this writing. Pre-continuation writing events: GPT-6 Astra (OpenAI), ChatGPT Pro mode, effort Pro, S6\_10\_PAPERS, 18 November 2026: targeted proof check; used and rechecked in this run. Continuation metadata events: GPT-6 Astra (OpenAI), ChatGPT Pro mode, effort Pro, S6\_10\_PAPERS continuation, 19 November 2026: used and rechecked in this run; status-vocabulary reconciliation; metadata reconciliation only; no new mathematical check. \emph{Pre-continuation recorded proof scope.} The night-run re-derivation was by GPT-6 Astra (OpenAI, ChatGPT Pro mode, effort Pro), 16--18 November 2026. Also checked in this writing by GPT-6 Astra (OpenAI), ChatGPT Pro mode, effort Pro, 18 November 2026. The conductor restriction and canonical-trace maps, primary restriction Euler bound, degree-zero jet and discriminant sections, product analytic types, and impossibility of a first blowup were re-derived. Greb 2011 Remark 2.1, printed page 415, was opened for the surface rational-singularity trace criterion; the earlier local jet bounds remain separately recorded dependencies}{\ref{night:P4.30}; \ref{night:P4.31}; \ref{night:P4.33}; \ref{night:P4R.10}; \ref{night:P4R.6}; \ref{night:P4R.7}; \ref{night:P4R.9}}

In Part~4.2, let an irreducible singular support $S$ have a
normalization with ruled resolution. If its reduced image
conductor has a component with elliptic normalization, then
the normalization of $S$ is a smooth minimal elliptic ruled
surface, the image conductor is one smooth elliptic curve,
and $e(S)=0$. Every local germ along this curve is the product
of its generic transverse germ with a disc. In particular no
rational conductor component is attached to that curve.
\end{proposition}
\begin{proof}
\noindent\textbf{Step 1.} Let $\rho_N:T\to N$ be the minimal normalization resolution,
let $\nu:N\to S$ be the finite normalization, and put
$\rho=\nu\rho_N$. The reduced-horizontal case is already
proved by P4R.9. We therefore use the actual double-horizontal
pole divisor $D$ of P4R.6. Its support is a smooth section $H$
with rational trees attached. The pushout and curve-group
calculation of P4R.7--P4R.9 show that the reduced image
conductor $C$ is connected, that its point-identification
graph has no cycles, and that it has exactly one normalized
component of positive genus, of genus one. Denote the
corresponding reduced component by $C_E$ and its normalization
by $j:E\to C_E\hookrightarrow Y$. In particular every point
of $C_E$ is unibranch. Every other normalized component of
$C$ is rational.

\noindent\textbf{Step 2.} A rational component of $D$ cannot map nonconstantly to $E$.
Thus $H$ maps nonconstantly to $E$. Write $g=g(H)\ge1$.
At a general point of $C_E$, every normalization branch has
conductor exponent two, because the coefficient of $H$ in
$D$ is two. The intrinsic conductor exponent of a singular
plane branch is at least two, and each additional branch
increases it by its positive intersection multiplicity.
The possibilities are therefore a single ordinary cusp,
two smooth branches with intersection multiplicity two, or
three smooth branches with pairwise intersection multiplicity
one. These are $A_2,A_3,D_4$, respectively, and $H\to E$
has degree $1,2,3$ in the three cases.

\noindent\textbf{Step 3.} We first construct a map which detects global functions on
the entire conductor, including its rational components.
Since $H$ is not contracted, P4.33 proves that $N$ has
rational singularities. Put $L_N=\nu^*\omega_S$.
Finite duality identifies its conductor ideal with
\[
 I_N=\omega_N\otimes L_N^{-1}.
\]
The canonical trace is an isomorphism
$\rho_{N*}\omega_T=\omega_N$ for rational analytic
singularities. We use Greb, Remark~2.1, p.~415
\cite{Greb2011}. The hypotheses hold here: $N$ is normal
and Cohen--Macaulay, $\rho_N$ is a proper resolution,
P4.33 gives $R^1\rho_{N*}\mathcal O_T=0$, and its
one-dimensional fibres give vanishing of the higher
structure-sheaf direct images. The integral pole identity
of P4R.6 and projection formula consequently give
\[
 \rho_{N*}\mathcal O_T(-D)
 =\rho_{N*}(\omega_T\otimes\rho_N^*L_N^{-1})
 =I_N.
\]
These are equalities of ideal subsheaves of $\mathcal O_N$:
the pole identity uses the pullback of the same residue
section which defines finite duality, and the trace agrees
with it on the birational isomorphism locus. Equality there
determines the maps between these torsionfree sheaves.
The divisor ideal sequence therefore gives an injection
\[
 \mathcal O_N/I_N\longrightarrow\rho_{N*}\mathcal O_D.
\]
The finite conductor square gives an injection from the
image conductor scheme. Composing them yields the actual
morphism of coherent $\mathcal O_S$-algebras
\[
 \mathcal O_{C_{\rm sch}}\longrightarrow
 \nu_*(\mathcal O_N/I_N)\longrightarrow\rho_*\mathcal O_D,
 \qquad\hbox{both arrows injective}.                         \tag{P4R.13.1}
\]

\noindent\textbf{Step 4.} We show that a global function on $D$ is determined by
its value at the generic doubled horizontal section.
At every point blowup $\beta:T'\to T''$ in the actual
map to the minimal ruled surface, the pole divisors obey
$D'=\beta^*D''-E_\beta$. The elementary direct images are
\[
 \beta_*\mathcal O(E_\beta)=\mathcal O,
 \quad R^1\beta_*\mathcal O(E_\beta)=0,
 \quad\beta_*\mathcal O=\mathcal O,
 \quad R^1\beta_*\mathcal O=0.
\]
For the first equality, a meromorphic function allowed
a pole only on $E_\beta$ descends off the centre and
extends across that point by Hartogs. For the vanishing,
the two blowup charts, or their exceptional-curve
filtration, have successive terms
$\mathcal O_{\mathbb P^1}(j)$ with $j\ge-1$;
their first cohomology vanishes. Equivalently one applies
$0\to\mathcal O\to\mathcal O(E_\beta)
\to\mathcal O_{E_\beta}(-1)\to0$ to the usual
blowup structure-sheaf calculation. Projection formula,
the divisor ideal sequence and the absence of second
direct images then give
\[
 \beta_*\mathcal O_{D'}=\mathcal O_{D''},
 \qquad R^1\beta_*\mathcal O_{D'}=0.                         \tag{P4R.13.2}
\]
These particular maps identify $H^0(D,\mathcal O_D)$
with $H^0(D_0,\mathcal O_{D_0})$, where
\[
 D_0=2H_0+V_0,
 \qquad V_0=\sum_j a_jF_j
\]
on the minimal ruled surface. The kernel of restriction
from $D_0$ to $2H_0$ is $\mathcal O_{V_0}(-2H_0)$.
Its successive layers on each thick fibre $a_jF_j$ are
$\mathcal O_{\mathbb P^1}(-2)$, since $F_j^2=0$ and
$H_0F_j=1$. They have no global sections. Restriction
to $2H_0$ is therefore injective on global functions.
This double section is Cohen--Macaulay and has no
zero-dimensional associated point, so a section which
vanishes generically on it vanishes everywhere.

\noindent\textbf{Step 5.} Let $C_{E,\rm sch}$ be the primary component of
$C_{\rm sch}$ supported on $C_E$: locally its ideal
is obtained by localizing the conductor ideal at the
generic point of $C_E$ and contracting. Primary
decomposition of the pure one-dimensional conductor
quotient makes these ideals coherent and compatible
on overlaps. Restriction is a surjection
\[
 \mathcal O_{C_{\rm sch}}\longrightarrow
 \mathcal O_{C_{E,\rm sch}}.
\]
Its map on $H^0$ is injective. Indeed, a global element
of its kernel maps by (P4R.13.1) to a function on $D$
which vanishes generically on $2H$; (P4R.13.2) and the
preceding paragraph make it zero. Its map on $H^1$
is surjective, because the kernel is supported in
dimension at most one and has no second cohomology.
We have proved
\[
 \chi(\mathcal O_{C_{\rm sch}})
 \le\chi(\mathcal O_{C_{E,\rm sch}}).                       \tag{P4R.13.3}
\]

\noindent\textbf{Step 6.} We evaluate the right side using the actual primary
ideal, without removing the rational components from
the left side. In type $A_2$ the primary conductor
is reduced: a pure one-dimensional quotient which
is generically reduced has no nilpotent submodule
of finite support. Thus its Euler characteristic
is $-\sum_p\delta(C_E,p)\le0$. In type $D_4$ the
primary ideal is the symbolic square of the ideal
of $C_E$, by exactly the generic-conductor and
depth-one argument of P4.31. P4.30 gives the same
nonpositive bound, with equality only for smooth
$C_E$.

\noindent\textbf{Step 7.} In type $A_3$, the primary ideal contains that
symbolic square and its nilpotent module is
torsionfree of rank one on $C_E$. The proof of
P4R.10 applies to this primary quotient: its
differentiation map, saturated Hessian line and
projection to the quotient conormal line depend
only on this ideal and on the actual equation of
$S$. The presence of other conductor components
changes none of those maps. More explicitly its
Euler characteristic is
\[
 -\sum_p\delta(C_E,p)+\deg R
 -\frac12\sum_p d_p-\sum_p b_p-\sum_p q_p\le0,
\]
with the same ramification divisor $R$, Hessian
zero orders $d_p$, projected lattice colengths
$b_p$ and descent defects $q_p$ as in P4R.10.
The local inequalities there prove this bound
at every singular point of $C_E$. Equality
makes $C_E$ smooth and its Hessian nowhere zero.

\noindent\textbf{Step 8.} On the other hand P4.33 gives
$\chi(\mathcal O_{C_{\rm sch}})=g-1$. Thus
\[
 0\le g-1
 =\chi(\mathcal O_{C_{\rm sch}})
 \le\chi(\mathcal O_{C_{E,\rm sch}})\le0.                    \tag{P4R.13.4}
\]
All these quantities vanish. We obtain $g=1$,
$C_E=E$ smooth, and an unramified map $H\to E$.
In type $A_3$ its transverse Hessian has rank
one everywhere on $E$.

\noindent\textbf{Step 9.} It remains to exclude attachments to this smooth
elliptic component. Put
\[
 K=\mathcal I_E/\mathcal I_E^2,
 \qquad N_E=\mathcal O_Y(S)|_E.
\]
Both $\deg K$ and $\deg N_E$ are zero. The first
follows from the conormal exact sequence,
$c_1(TY)E=0$, and $\deg\Omega_E=0$; the second
follows because a positive multiple of $S$ is
principal in a tube of this fibre.

\noindent\textbf{Step 10.} In type $A_2$ let $L\subset K$ be the saturated
quadratic-image line and put $Q=K/L$. The nonzero
quadratic tensor is a section of
$L^2\otimes N_E$, so $2\deg L\ge0$.
The cubic jet on the null direction is an
intrinsic section of $Q^3\otimes N_E$.
Here is the coordinate check, including points
where the quadratic tensor vanishes. A change
of a second-order lift changes the cubic
coefficient by the polarization of the
quadratic tensor. That polarization vanishes
on the null direction. A change of tangential
lift contributes a tangential derivative of
the rank-one quadratic tensor; its evaluation
twice on the null direction is again zero.
Changing a reduced defining equation multiplies
the tensor by its transition unit. Thus the
stated line bundle and its holomorphic section
are defined on all of $E$. The section is
nonzero generically by type $A_2$, so
$-3\deg L\ge0$. Both degrees are zero,
and both sections are nowhere zero. The
parametrized Morse lemma and a cube root
of the residual unit now give the product
equation $X^2+Y^3=0$ at every point of $E$.

\noindent\textbf{Step 11.} In type $A_3$ (P4R.13.4) already makes the
quadratic tensor nowhere zero, whence
$\deg L=\deg Q=0$. Its cubic null jet is
identically zero. Removing the quadratic
direction gives an intrinsic residual fourth
jet in $Q^4\otimes N_E$. In local coordinates
its relevant terms and coefficient are
\[
 a x^2+bxy^2+cy^4,
 \qquad q_4=c-\frac{b^2}{4a},\qquad a\ne0.                 \tag{P4R.13.5}
\]
Substitution of the critical coordinate
$x=-by^2/(2a)+O(y^3)$ proves the formula.
The change $x\mapsto x+vy^2$ sends $b$ to
$b+2av$ and $c$ to $c+bv+av^2$, leaving
$q_4$ unchanged. Changes of normal frames
and of equation give the claimed
$Q^4\otimes N_E$ transition. Tangential
lift changes do not affect a function
whose residual order is four. This section
is nonzero generically and has degree zero;
it never vanishes. The residual function
is $y^4$ times a unit, giving a product
tacnode at every point of $E$.

\noindent\textbf{Step 12.} In type $D_4$, the equation belongs to
$\mathcal I_E^3$. This ideal is primary
because $E$ is smooth, so generic order
three proves membership everywhere.
The cubic normal jet belongs to
$\operatorname{Sym}^3K\otimes N_E$.
Its discriminant is a section of
$(\det K)^6\otimes N_E^4$, of degree
zero. It is nonzero generically, hence
nowhere zero. The three distinct normal
tangent directions give three smooth
holomorphic branches locally along $E$.
After straightening the first two to
$x=0,y=0$, the third is $ax+by=0$
with both coefficients units. Setting
$X=ax,Y=-by$ gives the product equation
$XY(X-Y)=0$.

\noindent\textbf{Step 13.} In all three cases the entire singular
locus near every point of $E$ is $E$
itself. A rational conductor tree could
not meet it. Connectedness of $C$ then
forces $C=E$. The displayed local
product forms make $N$ smooth along
its conductor, and P4R.6 already makes
it smooth elsewhere. Hence $T=N$.
A finite normalization cannot contract
a curve, and a rational curve cannot
map nonconstantly to $E$. Thus $D=2H$
has no vertical component. The first
blowup of $2H_0$ in a nonempty actual
blowup sequence would create a vertical
pole component of coefficient one;
later blowups never remove its strict
transform from the pole support.
Consequently no blowup occurs and $N$
is minimal elliptic ruled. Finally
\[
 e(S)=e(N)-e(H)+e(E)=0.
\]
The exact coefficient changes, discriminant
weights and degree identities are executed
in \texttt{code/part4/elliptic\_primary\_restriction.py};
its JSON and text outputs are retained.
\end{proof}

\subsection{The smooth critical curve of transverse type \texorpdfstring{$A_3$}{A3}}

\begin{proposition}[P4R.14]\label{night:P4R.14}\label{stmt:P4R.14}
\StatementRecord{night:P4R.14}{proved}{Historical events: GPT-6 Astra (OpenAI), ChatGPT Pro mode, effort Pro, S6\_06\_BLIND, 16--18 November 2026: targeted proof check. Also checked in this writing. Pre-continuation writing events: GPT-6 Astra (OpenAI), ChatGPT Pro mode, effort Pro, S6\_10\_PAPERS, 18 November 2026 05:10:56.410177+00:00: computation or formal certificate; targeted proof check. Continuation metadata events: GPT-6 Astra (OpenAI), ChatGPT Pro mode, effort Pro, S6\_10\_PAPERS continuation, 19 November 2026: used and rechecked in this run; status-vocabulary reconciliation; metadata reconciliation only; no new mathematical check. \emph{Pre-continuation recorded proof scope.} The night-run re-derivation was by GPT-6 Astra (OpenAI, ChatGPT Pro mode, effort Pro), 16--18 November 2026. Also checked in this writing by GPT-6 Astra (OpenAI), ChatGPT Pro mode, effort Pro, 18 November 2026. proved. Full smooth-critical-axis A3 classification, finite-degree polar argument and explicit simple-elliptic local test. Diagonalized the rank-one Hessian over C\{z\}; the weight-four tensor is hx\textsuperscript{2}+axy\textsuperscript{2}+by4 with discriminant4hb−a\textsuperscript{2}. Independent constant changes make any zero simple away from h=0, producing the actual binomial X\textsuperscript{2}+y4Z and its forbidden compact spherical class. Thus the discriminant is a unit. At h(0)=0, the simple finite root after x=y\textsuperscript{2}v yields f=X(Y\textsuperscript{2}+XH). The off-axis critical equations are H\_z=0 and2H+XH\_X=0; the entire-axis hypothesis makes their intersection isolated. Re-derived the finite positive-degree map Psi=(H\_z,(G−2zH\_z)/2), the measure-zero exclusions of X=0 and H\_X=0, and the exact nonzero critical value−X\textsuperscript{3}H\_X/2. The full Hessian determinant is2X4J, so regular parameter values give actual Morse points and the boundary transport forces H\_z to be a unit. The explicit coordinate Z=H then gives X(Y\textsuperscript{2}+XZ). Re-derived the retained double-plane counterexample: its blowup double cover of P1 has4 branch values, exceptional square−2, residue discrepancy−1 and extra order1 from x. The unchanged polar program verifies all symbolic identities and the6 simple exact test solutions. Review provenance: New derivation after reading the supplied proof; not a blind review and not an independent human check}{\ref{night:P4.16}; \ref{night:P4.28}}

Let $S$ be an actual reduced support in Part~4.2. At a point
where the entire reduced critical locus of its local defining
function is one smooth curve, with generic transverse type
$A_3$, holomorphic coordinates put that function in one of
the forms
\[
 X^2+Y^4,\qquad X(Y^2+XZ).                              \tag{P4R.14.1}
\]
The normalization germs are respectively two smooth branches,
or a smooth plane and an $A_1$ quadric. No positive-genus
exceptional curve occurs over such a point.
\end{proposition}
\begin{proof}
\noindent\textbf{Step 1.} Choose a coordinate $z$ on the critical curve and normal
coordinates $x,y$. Its quadratic transverse Hessian has generic
rank one. Diagonalization of the symmetric matrix over
$\mathbb C\{z\}$, as in P4.28, gives the quadratic term
$h(z)x^2$. The cubic coefficient on the null direction
vanishes identically, by the generic transverse type. With
weights two for $x$ and one for $y$, the weight-four terms are
\[
 h(z)x^2+a(z)xy^2+b(z)y^4,\qquad
 \Delta(z)=4h(z)b(z)-a(z)^2.                             \tag{P4R.14.2}
\]
Every remaining term has larger weight and $\Delta$ is
nonzero generically.

\noindent\textbf{Step 2.} We first prove that $\Delta$ is a unit. A zero persists,
counted with multiplicity, under small independent constant
perturbations of $h,a,b$, by the argument principle on a
small parameter disc. Generic such perturbations make the
zeros simple and disjoint from $h=0$. Indeed $h=\Delta=0$
also requires $a=0$, which the independent parameters avoid;
where $h\ne0$, the derivative of $\Delta$ in the $b$
parameter is $4h\ne0$. These perturbations preserve the
critical axis and its transverse $A_3$ germ near the fixed
boundary. At a simple zero with $h\ne0$, the parametrized
Morse lemma gives
\[
 X^2+y^4U(y,z),\qquad U(0,0)=0,\quad U_z(0,0)\ne0.
\]
The actual coordinate $Z=U(y,z)$ gives $X^2+y^4Z$.
P4.16 supplies a compact spherical class with nonzero square
in its smoothing. The boundary comparison of P4.28 applies
here: near the boundary critical curve, the parametrized
Morse lemma and a fourth root of a unit give $X^2+Y^4$;
the tangential boundary coordinate makes the nonzero levels
submersive. Proper boundary transport therefore carries the
sphere to the actual nearby torus. Its universal cover is
$\mathbb R^4$, so a sphere has zero homology class and zero
square, a contradiction. Thus $\Delta$ is a unit.

\noindent\textbf{Step 3.} For $h(0)\ne0$, Morse splitting and a fourth root of the
residual unit give the first form in (P4R.14.1). Suppose
$h(0)=0$. Then $a(0)\ne0$. The holomorphic function
\[
 \frac{f(y^2v,y,z)}{y^4}
\]
restricts at $y=0$ to $h(z)v^2+a(z)v+b(z)$. Its finite
root $-b(0)/a(0)$ at the origin is simple. The holomorphic
implicit-function theorem gives a root $v=v(y,z)$, and
division by the monic equation of its graph gives
\[
 X=x-y^2v(y,z),\qquad f=Xg.
\]
On $X=0$, the function $g$ is $y^2$ times a unit. Its
parameter Morse centre in $y$ is therefore zero there and
its critical value is zero. The parametrized Morse lemma,
with parameters $X,z$, gives the exact form
\[
 f=X\bigl(Y^2+XH(X,z)\bigr).                            \tag{P4R.14.3}
\]
The remaining critical value is divisible by $X$ by its
vanishing on $X=0$. The function $H(0,z)$ is not identically
zero, since the generic transverse Hessian has rank one;
at the point under consideration $H(0,0)=0$.

\noindent\textbf{Step 4.} Put $G=2H+XH_X$. Off $X=0$, the critical equations of
(P4R.14.3) are $Y=0,H_z=0,G=0$. We prove $H_z(0,0)\ne0$.
Otherwise $H_z$ is a nonzero plane germ not divisible by
$X$, because $H(0,z)$ has finite order at least two.
No component of $H_z=0$ is contained in $G=0$: it would
give an additional critical-curve branch off $X=0$, whereas
the entire critical locus is the stated smooth axis. Hence
$(H_z,G)$ is an isolated complete intersection of positive
intersection multiplicity.

\noindent\textbf{Step 5.} Replace $H$ by $H-\eta z-\varepsilon$. Its plane critical
equations are
\[
 H_z-\eta=0,\qquad G-2\eta z-2\varepsilon=0.              \tag{P4R.14.4}
\]
Their solutions are the fibres of the holomorphic map
\[
 \Psi(X,z)=\left(H_z,\frac{G-2zH_z}{2}\right).
\]
Its zero fibre has the positive local degree just proved.
Choose a small compact representative with no boundary zero.
That degree is unchanged on a sufficiently small open
parameter ball. The universal equations in (P4R.14.4) have
parameter Jacobian determinant two, so regular parameter
values give a nonempty finite set of transverse solutions.
We also choose such a value outside the images under $\Psi$
of $X=0$ and $H_X=0$. These exceptional images have real
four-dimensional measure zero: each source curve has
countably many smooth real two-dimensional strata, and the
smooth image of each has measure zero. The germ $H_X$ is
not identically zero. Indeed in that case $H$ would depend
only on $z$ and both $H_z$ and $G=2H$ would vanish on the
whole line $z=0$, contrary to the isolated intersection.
The two exceptional images cannot exhaust the open ball
on which the positive degree is retained. This argument
uses no analyticity assertion for the image of a nonproper
map.

\noindent\textbf{Step 6.} At each chosen solution, the critical value is
\[
 f_{\eta,\varepsilon}=-\frac{X^3H_X}{2}\ne0.              \tag{P4R.14.5}
\]
For the indicated order of the two equations, put
\[
 J=\det\frac{\partial(G-2\eta z-2\varepsilon,\ H_z-\eta)}
                  {\partial(X,z)}.
\]
The $Y$-Hessian entry is $2X$, and the remaining Hessian
block at a critical point is
\[
 \begin{pmatrix}
 X(3H_X+XH_{XX})&X^2H_{Xz}\\
 X^2H_{Xz}&X^2H_{zz}
 \end{pmatrix}.
\]
Its determinant is $X^3J$, and the full Hessian determinant
is $2X^4J\ne0$. Thus these are isolated holomorphic Morse
points with nonzero critical values. Near the fixed boundary,
$H(0,z)$ remains nonzero under the small perturbation, so the
critical axis retains its transverse $A_3$ type. The same
proper boundary comparison transports a Morse vanishing
sphere, of square $-2$, into an actual nearby torus. This is
impossible. We have proved $H_z(0,0)\ne0$.

\noindent\textbf{Step 7.} The holomorphic coordinate $Z=H(X,z)$, keeping $X,Y$, now
gives the second form of (P4R.14.1). Its normalization has
the plane $X=0$ and the normal quadric $Y^2+XZ=0$. The
quadric's minimal resolution has one rational $(-2)$-curve.
The first form factors as $(X+iY^2)(X-iY^2)$ and its two
normalization branches are smooth. This proves the stated
normalization conclusions.

\noindent\textbf{Step 8.} The distinction between a pole coefficient and this
torus-smoothing conclusion can be checked by an exact
retained local test. Take
\[
 f=x\bigl[y^2+xz(x-z)(x-2z)\bigr].
\]
It has only the smooth $A_3$ critical axis $x=y=0$. The
second normalization branch is the double plane branched
over four distinct lines. Blowing up their common point
and normalizing gives a double cover of the exceptional
$\mathbb P^1$ branched in four points, hence an elliptic
curve of square $-2$. Its discrepancy is $-1$ and the
order of $x$ on it is one, giving pole coefficient two
for the displayed total reduced hypersurface. It is excluded
here by the polar deformation, not by the coefficient alone.
For $H=z(x-z)(x-2z)$, the deformation $H\mapsto H-1$ has
the exact Groebner basis
\[
 3x-5z^4-14z,\qquad 25z^6+50z^3-2.
\]
It gives six simple solutions; adjoining $x$, $H_x$ or
$J$ to the critical ideal gives the unit ideal in each
case. For $H\mapsto H-\varepsilon$ all six scale by
$\varepsilon^{1/3}$ and converge to the origin. The
identities and these exact rational calculations are
executed in \texttt{code/part4/smooth\_a3\_polar.py},
with JSON and text output retained.
\end{proof}

\subsection{A nonzero Hessian and the entire critical curve}

\begin{proposition}[P4R.15]\label{night:P4R.15}\label{stmt:P4R.15}
\StatementRecord{night:P4R.15}{proved}{Historical events: GPT-6 Astra (OpenAI), ChatGPT Pro mode, effort Pro, S6\_06\_BLIND, 16--18 November 2026: targeted proof check. Also checked in this writing. Pre-continuation writing events: GPT-6 Astra (OpenAI), ChatGPT Pro mode, effort Pro, S6\_10\_PAPERS, 18 November 2026 05:10:56.410177+00:00: computation or formal certificate; targeted proof check. Continuation metadata events: GPT-6 Astra (OpenAI), ChatGPT Pro mode, effort Pro, S6\_10\_PAPERS continuation, 19 November 2026: used and rechecked in this run; status-vocabulary reconciliation; metadata reconciliation only; no new mathematical check. \emph{Pre-continuation recorded proof scope.} The night-run re-derivation was by GPT-6 Astra (OpenAI, ChatGPT Pro mode, effort Pro), 16--18 November 2026. Also checked in this writing by GPT-6 Astra (OpenAI), ChatGPT Pro mode, effort Pro, 18 November 2026. proved. Full nonzero-Hessian product theorem allowing initially different transverse branch types. A nonzero complex symmetric Hessian supplies an actual Morse coordinate X and f=X\textsuperscript{2}+g(Y,Z). Factored the nonzero analytic plane function into its distinct irreducible factors with multiplicities. Two distinct factors define a finite positive-degree map; small independent target values give a transverse intersection avoiding all other factors. The unit-root coordinate gives X\textsuperscript{2}+Y\textasciicircum{}a Z\textasciicircum{}b, and P4.18 supplies its nonzero-square compact spherical class. Verified the boundary identification for every multiplicity: exponent1 remains smooth, higher exponents retain the product A-type; proper transport contradicts the torus. With a single factor g=h\textasciicircum{}m, a singular reduced h has finite positive gradient degree; subtracting a small linear form and shifting its value yields a Morse point with nonzero h-value. The full Hessian is diag(2,m v\textasciicircum{}(m−1)Q), of determinant2m\textsuperscript{2}v\textasciicircum{}(2m−2)detQ, and the same boundary transport excludes it. Therefore h is smooth and the actual coordinate Y=h gives the product. The unchanged arbitrary-m symbolic program passed. Review provenance: New derivation after reading the supplied proof; not a blind review and not an independent human check}{\ref{night:P4.18}; \ref{night:P4.28}}

At every critical-curve point of an actual reduced support in
Part~4.2 with nonzero ambient Hessian, holomorphic coordinates
give the exact product germ
\[
                       X^2+Y^m,\qquad m\ge2.
\]
In particular its entire reduced critical curve is smooth.
At a singular point of the reduced critical curve, the ambient
Hessian vanishes. This assertion allows different generic
transverse types on the different branches before the exclusion.
\end{proposition}
\begin{proof}
\noindent\textbf{Step 1.} A nonzero complex symmetric Hessian has a vector on which its
quadratic form is nonzero. The parametrized holomorphic Morse
lemma in that direction gives the exact equation
\[
                        f=X^2+g(Y,Z).
\]
The plane function $g$ is nonzero because the support is
reduced. Factor it in the convergent power-series ring as
\[
             g=u\prod_{i=1}^{r}h_i^{m_i},
             \qquad u(0)\ne0,\quad m_i\ge1,
\]
with distinct reduced irreducible factors $h_i$ vanishing at
the point. At least one exponent is at least two: a reduced
plane function has an isolated critical point, while the
critical locus here contains a curve.

\noindent\textbf{Step 2.} We exclude $r\ge2$ by an actual deformation of these factors.
For two distinct factors $h_i,h_j$, the holomorphic map
\[
 (h_i,h_j):(\mathbb C^2,0)\longrightarrow(\mathbb C^2,0)
\]
has isolated zero fibre and positive finite local degree.
On a sufficiently small compact representative that degree
is retained for small target values. Regular such values
give a nonempty transverse intersection of
$h_i-s_i=0$ and $h_j-s_j=0$. Independently choose the
other constants $s_k$ so that at least one of these
intersection points avoids every other zero factor. This
is possible because the chosen intersection set is finite
and each forbidden $s_k$ belongs to a finite set of values.
All parameters can be arbitrarily small.

\noindent\textbf{Step 3.} At the selected intersection, the two transverse factors
are holomorphic coordinates and every other factor is a
unit. Taking an actual local $m_i$th root of their unit
product and multiplying the first coordinate by it gives
the germ
\[
                        X^2+Y^{m_i}Z^{m_j}.
\]
P4.18 excludes every such germ by its compact spherical
class of nonzero square. The comparison here is to the
original torus degeneration, as follows. Choose a small
Milnor ball before moving the factors. Near its boundary,
the distinct reduced factors are smooth and disjoint.
The moved factor with exponent $m_i\ge2$ retains exactly
the transverse product type $A_{m_i-1}$ there; a factor
of exponent one remains a regular point of the surface.
Parametrized Morse coordinates and the appropriate root
of the remaining unit identify these boundary germs.
Solving the tangential boundary coordinate makes the
nonzero levels submersive. The proper boundary transport
used in P4.28 therefore identifies the needed nearby
levels and transports the spherical class into the actual
nearby torus. Its square cannot be nonzero there. Hence
there is only one factor.

\noindent\textbf{Step 4.} Absorb the unit by writing $h=u^{1/m}h_1$, so that
$g=h^m$. The reduced plane function $h$ is irreducible.
For a singular $h$, its gradient has an isolated zero,
and the finite positive local degree of
$(h_Y,h_Z)$ gives nonempty finite regular fibres under
small constant changes of that gradient. Equivalently,
subtracting a sufficiently small generic linear form
from $h$ gives isolated Morse critical points. Add a
small constant to make their critical values nonzero.
Call the resulting function $h_s$, and choose one such
critical point with value $v\ne0$ and nonsingular Hessian
$Q$. For $f_s=X^2+h_s^m$, the critical value and Hessian
are exactly
\[
 f_s=v^m\ne0,\qquad
 \operatorname{Hess}(f_s)
   =\begin{pmatrix}2&0\\0&m v^{m-1}Q\end{pmatrix}.
\]
Thus the full Hessian determinant is
$2m^2v^{2m-2}\det Q\ne0$. The moving curve $h_s=0$
retains its transverse $A_{m-1}$ type on the fixed
boundary. The same proper comparison transports the
Morse vanishing sphere, of square $-2$, to a nearby
torus, a contradiction. Consequently $h$ is smooth,
and using it as the coordinate $Y$ gives the asserted
product equation.

\noindent\textbf{Step 5.} The Hessian identity is executed symbolically for the
arbitrary positive integer $m$ in
\texttt{code/part4/hessian\_product.py}. Its retained
output also checks the moving-node example
$X^2+(YZ-\varepsilon)^m$ at a nonzero critical value.
\end{proof}


\subsection{The tacnode and transverse-line critical curve}

\begin{proposition}[P4R.16]\label{night:P4R.16}\label{stmt:P4R.16}
\StatementRecord{night:P4R.16}{proved}{Historical events: GPT-6 Astra (OpenAI), ChatGPT Pro mode, effort Pro, S6\_06\_BLIND, 16--18 November 2026: targeted proof check. Also checked in this writing. Pre-continuation writing events: GPT-6 Astra (OpenAI), ChatGPT Pro mode, effort Pro, S6\_10\_PAPERS, 18 November 2026 05:10:56.410177+00:00: computation or formal certificate; targeted proof check. Continuation metadata events: GPT-6 Astra (OpenAI), ChatGPT Pro mode, effort Pro, S6\_10\_PAPERS continuation, 19 November 2026: used and rechecked in this run; status-vocabulary reconciliation; metadata reconciliation only; no new mathematical check. \emph{Pre-continuation recorded proof scope.} The night-run re-derivation was by GPT-6 Astra (OpenAI, ChatGPT Pro mode, effort Pro), 16--18 November 2026. Also checked in this writing by GPT-6 Astra (OpenAI), ChatGPT Pro mode, effort Pro, 18 November 2026. proved. Full tacnode-plus-transverse-line symbolic square, rank-one obstruction and defect11 bound. Used the exact embedded ideal(xz,yz,y\textsuperscript{2}−x4) and its three parametrizations. The normalization sequence gives delta3 and zero ramification. Direct intersection of (z,h)\textsuperscript{2} and(x,y)\textsuperscript{2} gives P\textasciicircum{}(2)=P\textsuperscript{2}+(zh); the three displayed product identities kill zh by the maximal ideal, while its cubic term proves it is nonzero modulo P\textsuperscript{2}. In f=kappa zh+P\textsuperscript{2} the transverse determinant on(t,t\textsuperscript{2},0) has constant term−kappa\textsuperscript{2}, so rank1 forces kappa0 and the exact inclusion f\ensuremath{\in}P\textsuperscript{2}. Recomputed all three primitive gradient matrices, with elementary orders(1,2),(1,2),(1,1). The symmetric-square lattice inequality then gives3+3+3+2=11 before adding the nonnegative descent length, contradicting the exact total9. The unchanged a3\_tacnode\_line.py certificate passed. Review provenance: New derivation after reading the supplied proof; not a blind review and not an independent human check}{\ref{night:P4.33}; \ref{night:P4.40}; \ref{night:P4.42}; \ref{night:P4.44}}

In the remaining positive ruled case with uniform transverse
type $A_3$, the three-branch point of the critical curve cannot
be the union of a plane tacnode and an independent transverse
line. Its local conductor defect is at least eleven, whereas
the total defect in the contracted-section case is nine.
\end{proposition}
\begin{proof}
\noindent\textbf{Step 1.} The analytic ring calculation in P4.40 supplies actual ambient
coordinates for this curve. Its reduced ideal is
\[
 P=(xz,yz,h),\qquad h=y^2-x^4.
\]
Its three parametrizations are $(t,t^2,0)$, $(t,-t^2,0)$,
and $(0,0,t)$. The two tangent branches have intersection
multiplicity two, and their union meets the independent line
with intersection length one. The normalization sequence
therefore gives $\delta=3$. Every branch is immersive, so
its ramification divisor has coefficient zero.

\noindent\textbf{Step 2.} We first determine the symbolic square and the actual
equation. The plane-tacnode ideal $(z,h)$ is a reduced
complete intersection. Its square is unmixed and agrees
with the square of each prime at its generic point; hence
its symbolic square is its ordinary square. Intersecting
with the square of the line ideal gives
\[
 P^{(2)}=(z,h)^2\cap(x,y)^2=P^2+(zh).
\]
Indeed an element of $(z,h)^2$ is $Az^2+Bzh+Ch^2$.
Its last two terms belong to $(x,y)^2$. The first belongs
to that ideal precisely when $A$ does, because $z$ is a
nonzerodivisor modulo $(x,y)^2$. The resulting three
generators $x^2z^2,xyz^2,y^2z^2$ are already in $P^2$.
Furthermore
\[
 xzh=(xz)h,\qquad yzh=(yz)h,\qquad
 z^2h=(yz)^2-x^2(xz)^2.
\]
The class of $zh$ modulo $P^2$ is therefore annihilated
by the ambient maximal ideal. It is nonzero, since it has
a nonzero cubic term and $P^2$ has order at least four.
Thus $P^{(2)}/P^2$ is exactly one copy of the residue field.

\noindent\textbf{Step 3.} Let $f$ be the reduced defining function of the support.
Its generic transverse multiplicity is two, so
$f\in P^{(2)}$ and
\[
                       f=\kappa zh+f_0,\quad f_0\in P^2,
                       \quad\kappa\in\mathbb C.
\]
At $(t,t^2,0)$ with $t\ne0$, use $z,h$ as transverse
coordinates. The quadratic transverse part of $f$ has the
form
\[
 t^2a(t)z^2+(\kappa+t b(t))zh+c(t)h^2,
\]
where $a,b,c$ are holomorphic at $t=0$. The transverse
Hessian has rank one, so its determinant vanishes for
all sufficiently small nonzero $t$. Its constant term at
zero is $-\kappa^2$ (up to the harmless common convention
for the quadratic-form matrix). Hence $\kappa=0$, and
we have proved the exact inclusion $f\in P^2$.

\noindent\textbf{Step 4.} Let $K_i$ be the saturated rank-two kernel of the derivative
of the $i$th parametrization, and let $E_i\subset K_i$ be
the span of the differentiated generators of $P$. On the
two tangent branches choose the bases
\[
                dz,\quad dy\mp2t\,dx.
\]
The generator matrices have columns
\[
 \begin{pmatrix}t&\pm t^2&0\\0&0&\pm2t^2\end{pmatrix}.
\]
Their elementary-divisor orders are $(1,2)$. On the
independent line use $dx,dy$; the matrix is
\[
                   \begin{pmatrix}t&0&0\\0&t&0\end{pmatrix},
\]
with orders $(1,1)$.

\noindent\textbf{Step 5.} For completeness, the rank-one quadratic estimate of
P4.42 applies to these exact lattices as follows. Write
$E=t^a\mathbb C\{t\}e_1\oplus
t^b\mathbb C\{t\}e_2$, and write the saturated Hessian
line as $\ell=u e_1+v e_2$, with $(u,v)$ primitive.
Since $f\in P^2$, the tensor $t^d\ell^2$ belongs to
$\operatorname{Sym}^2E$. Therefore
\[
 \frac d2\ge\max(a-\operatorname{ord}u,
                         b-\operatorname{ord}v).
\]
Projection of $E$ to $K/\ell$ has colength
$b_{\rm proj}=\min(a+\operatorname{ord}v,
b+\operatorname{ord}u)$. Adding these two expressions
gives $d/2+b_{\rm proj}\ge a+b$; the identity follows
by subtracting the same two candidate quantities, with
the zero-component cases read directly. The local defect
formula of P4.42 consequently gives
\[
 \delta+\sum_i(d_i/2+b_i-r_i)+q
       \ge3+(1+2)+(1+2)+(1+1)=11.
\]
Here $q\ge0$ is the further descent colength, retained
rather than discarded from the definition. P4.44 contracts
the horizontal section, and P4.33 and (4.42a) give total
defect nine. The displayed lower bound contradicts that
total. This excludes the stated analytic curve germ.
\end{proof}

The symbolic-square identities, all three normalized
conormal matrices, and the rank-one determinant obstruction
are checked exactly in
\texttt{code/part4/a3\_tacnode\_line.py}.


\subsection{The nonordinary three-branch tacnode sector}

\begin{proposition}[P4R.17]\label{night:P4R.17}\label{stmt:P4R.17}
\StatementRecord{night:P4R.17}{proved}{Historical events: GPT-6 Astra (OpenAI), ChatGPT Pro mode, effort Pro, S6\_06\_BLIND, 16--18 November 2026: targeted proof check. Also checked in this writing. Pre-continuation writing events: GPT-6 Astra (OpenAI), ChatGPT Pro mode, effort Pro, S6\_10\_PAPERS, 18 November 2026 05:10:56.410177+00:00: computation or formal certificate; targeted proof check. Continuation metadata events: GPT-6 Astra (OpenAI), ChatGPT Pro mode, effort Pro, S6\_10\_PAPERS continuation, 19 November 2026: used and rechecked in this run; status-vocabulary reconciliation; metadata reconciliation only; no new mathematical check. \emph{Pre-continuation recorded proof scope.} The night-run re-derivation was by GPT-6 Astra (OpenAI, ChatGPT Pro mode, effort Pro), 16--18 November 2026. Also checked in this writing by GPT-6 Astra (OpenAI), ChatGPT Pro mode, effort Pro, 18 November 2026. proved. Full unrestricted delta-four nonplanar triple classification, exact ambient morphisms and all local bounds. For any normalized conormal lattice E, the exact sequence along the saturated Hessian line gives a+b=e1+e2. A unit Hessian column gives d\ensuremath{\geq}a and E\ensuremath{\subset}t\textasciicircum{}e1K gives b\ensuremath{\geq}e1, hence d/2+b\ensuremath{\geq}e1+e2/2 without assuming f\ensuremath{\in}P\textsuperscript{2}. The nonplanar branch estimate gives delta+3+m/2. For three smooth branches and delta4, classified the entire finite-colength analytic ring from its first2 jets. One first-jet direction gives the exact common-tangent model with E=t\textsuperscript{2}K on each branch, defect13. Two first-jet directions force two proportional branch evaluations; after actual parameter changes, the residual coupling mu\ensuremath{\neq}0 gives two analytic generators and hence a planar ring. For mu=0 the exact contact-three-plus-line ideal has conormal orders(1,3),(1,3),(1,1), giving21/2. For multiplicity4, branch delta2 would leave the minimally glued ring of embedding dimension4; hence the singular branch is the ordinary cusp. The colength-one module makes m0\textsuperscript{2}\ensuremath{\subset}A, and embedding dimension3 forces mA\textsuperscript{2}=m0\textsuperscript{2}. Classified its cotangent hyperplane and product map. A nonzero cusp-cubic coefficient gives independent tangent axes; the only allowed initial cubic xyz has rank-two leading Hessian and therefore vanishes, giving f\ensuremath{\in}m4 and defect11. With zero cubic coefficient, the three other coefficients must all be nonzero; the actual model has ideal(z(x−y),xy(x−y),z\textsuperscript{2}−x\textsuperscript{2}y). Its Hilbert–Burch minors, generic reducedness and depth give exact equality with the curve ideal. Primitive conormal matrices have orders(3,3),(1,2),(1,2), so defect23/2. The delta3 and delta2 cases then follow from the explicitly rederived excess-one ring classification. The unchanged a3\_delta4.py verifies all9 lattices and the4 exact jet identities; its received historical checker wording is not counted as a new check. Review provenance: New derivation after reading the supplied proof; not a blind review and not an independent human check}{\ref{night:P4.40}; \ref{night:P4.42}; \ref{night:P4.48}; \ref{night:P4R.10}; \ref{night:P4R.15}; \ref{night:P4R.16}}

In the remaining positive ruled case with uniform generic transverse
type $A_3$, every nonordinary three-branch critical-curve point has
local conductor defect strictly greater than nine. In particular,
the delta-four nonplanar cases have the explicit lower bounds
$13$, $21/2$, $11$, and $23/2$ below. The only three-branch germ
compatible with the total defect budget nine is the ordinary
nonplanar union of three coordinate axes.
\end{proposition}
\begin{proof}
Use the notation of P4R.10 and P4.42: \[
 D(C,f)=\delta(C)+
        \sum_i\left(\frac{d_i}{2}+b_i-r_i\right)+q,\qquad q\ge0.
\] Here \(d_i\) is the order of the saturated rank-one Hessian, \(b_i\) is the projected conormal colength, and \(r_i=m_i-1\) for a branch of multiplicity \(m_i\).

\par\medskip\noindent\textbf{1. Bounds without an ordinary-square assumption}\par\label{bounds-without-an-ordinary-square-assumption}

Let \(E_i\subset K_i\) be the normalized conormal lattice, with elementary divisors \((t^{e_1},t^{e_2})\), \(e_1\le e_2\). For the saturated Hessian line \(L_i\), put \[
 a_i=\operatorname{length}\bigl(L_i/(L_i\cap E_i)\bigr).
\] The exact lattice sequence gives \(a_i+b_i=e_1+e_2\). Every first derivative of \(f\) belongs to the radical ideal \(P\) of the critical curve. Thus every Hessian column belongs to \(E_i\). Write the Hessian as \(ut^{d_i}\ell_i^2\), with \(u\) a unit and \(\ell_i\) primitive. One coordinate of \(\ell_i\) is a unit; its column gives \(t^{d_i}\ell_i\in E_i\), so \(d_i\ge a_i\). Also \(b_i\ge e_1\), since \(E_i\subset t^{e_1}K_i\). Consequently \[
                  \frac{d_i}{2}+b_i
                    \ge e_1+\frac{e_2}{2}.           \tag{A}
\] This remains valid when a symbolic-square term occurs in \(f\).

Since the curve is nonplanar, \(P\) has no linear equation. Every defining gradient therefore has order at least \(m_i\) on branch \(i\), and \(b_i\ge m_i\). P4R.15 gives zero ambient Hessian, so \(f\) has order at least three and \(d_i\ge m_i\). For total multiplicity \(m=\sum m_i\), \[
                  D(C,f)\ge\delta(C)+3+\frac m2.       \tag{B}
\] Thus \(m\ge5\) gives \(D\ge19/2>9\). Only multiplicities three and four need inspection.

\par\medskip\noindent\textbf{2. Three smooth branches: the complete excess-two classification}\par\label{three-smooth-branches-the-complete-excess-two-classification}

Put \[
 B=\bigoplus_{i=1}^3\mathbb C\{t_i\},\quad
 F^j=\bigoplus_{i=1}^3t_i^j\mathbb C\{t_i\},\quad
 A_0=\mathbb C_{\rm diag}+F^1.
\] The actual ring \(A\subset A_0\) has \(\operatorname{length}(A_0/A)=2\), since \(A_0\) has delta invariant two. Each branch projection of \(A\) is onto.

Let \(W_1\subset F^1/F^2=\mathbb C^3\) be its first-jet image. Each coordinate projection is nonzero. If \(W_1=\mathbb C^3\), Nakayama first gives \(\mathfrak m_AA_0=F^1\) and then, on the finite \(A\)-module \(A_0/A\), gives \(A=A_0\). This contradicts colength two. Thus \(\dim W_1=1\) or \(2\).

\par\medskip\noindent\textbf{2a. One-dimensional first jets}\par\label{a.-one-dimensional-first-jets}

Scale the branch parameters so \(W_1=\mathbb C(1,1,1)\). The two missing first jets use the entire colength, hence \(F^2\subset A\). The exact algebra is \[
 A=\mathbb C_{\rm diag}+\mathbb C(t_1,t_2,t_3)+F^2.
\] It is generated by \[
 x=(t_1,t_2,t_3),\quad y=(t_1^2,0,0),\quad z=(0,t_2^2,0).
\] Indeed \(y,z,x^2-y-z\) give all three quadratic branch ideals, and multiplication by \(x\) gives every higher order. The embedded model is exactly \[
       P=\bigl(y(y-x^2),yz,z(z-x^2)\bigr).             \tag{C}
\] The branches are \((t,t^2,0),(t,0,t^2),(t,0,0)\). Its differentiated generators give \(E_i=t_i^2K_i\) on all three branches. For example, on \((t,t^2,0)\) use the primitive basis \(dy-2t\,dx,dz\). Formula (A) gives \[
                         D\ge4+3\cdot3=13.            \tag{D}
\]

\par\medskip\noindent\textbf{2b. Two-dimensional first jets}\par\label{b.-two-dimensional-first-jets}

The three branch evaluations are nonzero linear forms on \(W_1\). If no two are proportional, their squares are independent: choose the first two as \(s,t\); the third is \(as+bt\), \(ab\ne0\), whose square supplies the nonzero \(st\) coefficient. Thus products give the whole second-jet space. Multiplication by one first jet nonzero on all branches then gives every later jet. This would leave only one missing positive-jet dimension, contrary to colength two.

Two evaluations are therefore proportional. Scale and label the parameters to obtain \(W_1=\{(a,a,b)\}\). Its products give the same two-dimensional second-jet space. The one remaining missing dimension occurs there: if the second-jet image were full, multiplication by a nonvanishing first jet would make every later one full. Hence \(F^3\subset A\), and the pure second-jet image is exactly \(\{(c,c,d)\}\).

A lift of \(W_1\) has quadratic difference \(\lambda a+\mu b\). The analytic changes \(t_i\mapsto t_i+c_it_i^2\), on the first two branches, remove \(\lambda a\). Thus the exact ring has common constants, with first and second coefficients satisfying \[
                     a_1=a_2,\qquad c_1-c_2=\mu a_3, \tag{E}
\] and arbitrary independent terms of order at least three.

For \(\mu\ne0\), the two elements \[
 x=(t_1,t_2,0),\qquad y=(\mu t_1^2,0,t_3)
\] generate this whole ring. The exact identities \[
 \begin{aligned}
 xy/\mu&=(t_1^3,0,0),\\
 x^3-xy/\mu&=(0,t_2^3,0),\\
 y^3-\mu(xy)^2&=(0,0,t_3^3),\\
 y^2-\mu x^2y&=(0,0,t_3^2)
 \end{aligned}
\] give the independent cubic ideals and every permitted quadratic jet. Multiplication by \(x\) or \(y\) gives all the higher branch series. This ring is planar, contrary to the present hypothesis.

Thus \(\mu=0\). The ring is generated by \((t_1,t_2,0),(0,0,t_3),(t_1^3,0,0)\). Naming the independent transverse coordinate \(z\) gives the exact embedded model \[
                       P=(xz,yz,y(y-x^3)).            \tag{F}
\] Its branches are \((t,0,0),(t,t^3,0),(0,0,t)\), and its conormal elementary divisors are \[
                         (1,3),\quad(1,3),\quad(1,1).
\] On the second branch use \(dy-3t^2dx,dz\); the gradients give \(t^3(dy-3t^2dx)\) and \(t\,dz\). Without restricting a symbolic-square coefficient of \(f\), (A) gives \[
                  D\ge4+\frac52+\frac52+\frac32
                       =\frac{21}{2}>9.               \tag{G}
\]

These are actual analytic ring isomorphisms. In each nonplanar case the three displayed generators form a basis of the curve cotangent space. Lift them to the ambient local ring. Their linear Jacobian is invertible, so the inverse-function theorem gives an ambient coordinate isomorphism to the stated model.

\par\medskip\noindent\textbf{3. Multiplicity four: one cusp and two smooth branches}\par\label{multiplicity-four-one-cusp-and-two-smooth-branches}

The singular branch has multiplicity two and delta at least one. Joining three branches costs at least two further delta units. If the branch delta is two, all remaining gluing is minimal, and the local ring is \[
 \mathbb C_{\rm diag}
 +\mathfrak m_{\rm singular}
 +t_2\mathbb C\{t_2\}+t_3\mathbb C\{t_3\}.
\] Its embedding dimension is at least \(2+1+1=4\), impossible in a smooth threefold. A larger branch delta is impossible. Thus this branch is the ordinary cusp \(\mathbb C\{t^2,t^3\}\).

Set \[
 A_0=\mathbb C_{\rm diag}
 +\mathfrak m_{\mathbb C\{t^2,t^3\}}
 +t_2\mathbb C\{t_2\}+t_3\mathbb C\{t_3\}.
\] Now \(A_0\) has delta three, and \(A\subset A_0\) has colength one. The one-dimensional \(A\)-module \(A_0/A\) is killed by \(\mathfrak m_A\), so \(\mathfrak m_AA_0\subset A\). Its elements have zero constant term, giving \(\mathfrak m_AA_0=\mathfrak m_A\). Hence \(\mathfrak m_A\) is an ideal of \(A_0\). The local algebra \(A_0/\mathfrak m_A\) has length two, so its maximal ideal squares to zero. Therefore \[
                         \mathfrak m_0^2\subset A.     \tag{H}
\]

The four cotangent coordinates of \(A_0\), denoted \(a,b,c,d\), are represented by the cusp functions \(t^2,t^3\) and the two smooth parameters. The three-dimensional subspace \(W=\mathfrak m_A/\mathfrak m_0^2\) is a hyperplane \[
                 \lambda a+\mu b+\nu c+\xi d=0.       \tag{I}
\] Its projection to \(a,b\) is onto, since the actual branch image is the full cusp ring, and the two smooth projections are nonzero. Ambient embedding dimension three gives \(\mathfrak m_A^2=\mathfrak m_0^2\). Thus products of \(W\) must generate \(\mathfrak m_0^2/\mathfrak m_0^3\), whose coordinates are \[
                              a^2,\quad ab,\quad c^2,\quad d^2.
                                                               \tag{J}
\]

\par\medskip\noindent\textbf{3a. A nonzero coefficient of the cusp cubic}\par\label{a.-a-nonzero-coefficient-of-the-cusp-cubic}

For \(\mu\ne0\), solve (I) for \(b\). Surjectivity of the cusp projection forces \((\nu,\xi)\ne(0,0)\). The three minimal ambient coordinate functions have parametrizations \[
 \begin{array}{c|c}
 \text{cusp}&(t^2-\lambda't^3,-\nu't^3,-\xi't^3)\\
 \text{second branch}&(0,t,0)\\
 \text{third branch}&(0,0,t).
 \end{array}
\] Their tangent directions are independent.

By P4R.15, \(f\in\mathfrak m^3\). Let \(G_3\) be its homogeneous cubic part. Since \(df\) vanishes on all three branches, \(dG_3\) vanishes on each independent tangent axis. Every cubic monomial other than \(xyz\) has a derivative nonzero on an axis. Hence \(G_3=cxyz\). On the \(x\)-axis its normal Hessian is \[
                         \begin{pmatrix}0&c\\c&0\end{pmatrix}.
\] A nonzero \(c\) gives a nonzero leading minor on the actual branch, contradicting generic rank one. Thus \(c=0\), and \(f\in\mathfrak m^4\).

The Hessian entries have order at least \(2m_i\) on branch \(i\), and \(b_i\ge m_i\). Its defect contribution is at least \(m_i+1\). Therefore \[
                            D\ge4+(4+3)=11>9.         \tag{K}
\]

\par\medskip\noindent\textbf{3b. A zero coefficient of the cusp cubic}\par\label{b.-a-zero-coefficient-of-the-cusp-cubic}

For \(\mu=0\), surjectivity in (J) forces \(\lambda,\nu,\xi\) all to be nonzero. Otherwise two of \(a,c,d\) are proportional on \(W\), so their squares are dependent, and (J) has rank at most three. That creates an extra minimal generator, contrary to embedding dimension three. Scale the branch parameters so that \(a+c+d=0\).

The pure hyperplane generators belong to \(A\) exactly by (H). They give the parametrizations \[
                       (t^2,t^2,t^3),\quad(t,0,0),\quad(0,t,0).
\] Their defining ideal is \[
                 P=(z(x-y),xy(x-y),z^2-x^2y).        \tag{L}
\] A Hilbert--Burch matrix is \[
 \begin{pmatrix}
 xy&z\\
 -z&-x\\
 0&-(x-y)
 \end{pmatrix}.
\] Its maximal minors give (L), with alternating signs. The ideal has codimension two and is Cohen--Macaulay. Its zero set is exactly the cusp and the two axes, and it is generically reduced on each branch. A nonzero nilpotent quotient supported at the origin would violate depth one, so it is the actual reduced intersection ideal.

On the cusp use the primitive conormal basis \[
                     k_1=dx-dy,\qquad k_2=3t\,dx-2dz.
\] The differentiated generators are \(t^3k_1,t^4k_1,t^3(tk_1-k_2)\). Thus its elementary divisors are \((3,3)\). On each smooth axis they are \((1,2)\). Here \(r=1\) on the cusp and zero on the axes. By (A), \[
                  D\ge4+\left(\frac92-1\right)+2+2
                       =\frac{23}{2}>9.               \tag{M}
\]

This exhausts multiplicity four. Together with Sections 1--2 it proves the strict local bound for every nonplanar three-branch curve with delta four. Only actual coordinate isomorphisms, conormal maps, and the rank-one Hessian were used; no boundary deformation has been introduced.

\par\medskip\noindent\textbf{4. Every nonordinary triple exceeds the budget}\par\label{every-nonordinary-triple-exceeds-the-budget}

For a planar three-branch point, P4.48 gives defect at least ten, including the ordinary plane triple point. At a nonplanar point with delta at least five, (B) and \(m\ge3\) give \[
                  D\ge5+3+\frac32=\frac{19}{2}>9.
\] Delta four was exhausted above.

At delta three the excess \(\delta-(r-1)\) is one. The actual analytic ring calculation in P4.40 gives precisely an ordinary plane triple, a plane tacnode with an independent transverse line, or a cusp with two independent lines. The first is excluded by P4.48. The second has defect at least eleven by P4R.16. The third has embedding dimension four and cannot occur in the ambient smooth threefold.

Finally a three-branch point with delta two has zero excess. Its ring equals \(\mathbb C_{\rm diag}+F^1\), generated by the independent branch parameters. It is exactly the three coordinate axes, with ideal \((xy,xz,yz)\). Thus every nonordinary three-branch point is excluded from the total defect budget nine. The only remaining local possibility is the ordinary nonplanar triple treated by P4.48 and the separate full-quartic conductor calculation.

\end{proof}

The analytic model identities, primitive conormal lattices,
Hilbert--Burch minors, cubic obstruction, and jet generators are
checked exactly in \path{code/part4/a3_delta4.py}.


\subsection{The horizontal image in the unrestricted ruled remainder}

\begin{proposition}[P4R.18]\label{night:P4R.18}\label{stmt:P4R.18}
\StatementRecord{night:P4R.18}{proved}{Historical events: GPT-6 Astra (OpenAI), ChatGPT Pro mode, effort Pro, S6\_06\_BLIND, 16--18 November 2026: targeted proof check. Also checked in this writing. Pre-continuation writing events: GPT-6 Astra (OpenAI), ChatGPT Pro mode, effort Pro, S6\_10\_PAPERS, 18 November 2026: targeted proof check; used and rechecked in this run. Continuation metadata events: GPT-6 Astra (OpenAI), ChatGPT Pro mode, effort Pro, S6\_10\_PAPERS continuation, 19 November 2026: used and rechecked in this run; status-vocabulary reconciliation; metadata reconciliation only; no new mathematical check. \emph{Pre-continuation recorded proof scope.} The night-run re-derivation was by GPT-6 Astra (OpenAI, ChatGPT Pro mode, effort Pro), 16--18 November 2026. Also checked in this writing by GPT-6 Astra (OpenAI), ChatGPT Pro mode, effort Pro, 18 November 2026. The A3/D4 conductor alternatives, primary defect budget, at-most-three branch-value graph bound, and degree-two Riemann-Hurwitz contradiction were re-derived, using the previously stated local product models and defect bounds}{\ref{night:P4.33}; \ref{night:P4.42}; \ref{night:P4R.10}; \ref{night:P4R.13}; \ref{night:P4R.14}; \ref{night:P4R.15}; \ref{night:P4R.6}}

In the remaining positive-Euler ruled case of Part~4.2, a
noncontracted horizontal pole section maps with degree two or
three to a rational normalized conductor component whose generic
transverse germ is $D_4$. This assertion imposes no transverse
type or conductor-exponent restriction on the other image component.
\end{proposition}
\begin{proof}
\noindent\textbf{Step 1.} Write $H$ for the horizontal section, $g=g(H)\ge1$, and
$C_H$ for its image in the reduced conductor $C$. The preceding
ruled reduction gives two rational normalized components of $C$
and one point-identification cycle. The nonconstant map
\[
                         H\longrightarrow C_H^\nu=\mathbb P^1
\]
has degree at least two: a degree-one finite map between smooth
curves is an isomorphism, whereas their genera differ. Each
generic point of $H$ has conductor exponent two, by the actual
pole divisor of P4R.6. Thus the generic transverse plane curve
along $C_H$ has at least two normalization branches of exponent two.

\noindent\textbf{Step 2.} For its branches $B_i$ the exact conductor formula is
\[
             c_i=2\delta(B_i)+\sum_{j\ne i}I(B_i,B_j).
\]
A singular branch with $c_i=2$ permits no other branch,
so the two specified branches are smooth. With two total
branches their intersection multiplicity is two, giving $A_3$.
With three total branches, the first two exponents force all
three pairwise intersection multiplicities to equal one;
the third branch is then smooth as well, giving $D_4$.
Four branches would give exponent at least three on every
smooth branch. These possibilities are exhaustive. In type
$A_3$, $H$ has degree exactly two and supplies both generic
normalization preimages. In type $D_4$ its degree is two or
three. We exclude the former transverse type next.

\noindent\textbf{Step 3.} Assume that the transverse germ along $C_H$ is $A_3$.
Let $C_{H,\rm sch}$ be the primary component of the actual
image conductor scheme supported on $C_H$. The primary
restriction construction of P4R.13 applies before its
elliptic-genus calculation, without change. More explicitly,
P4.33 gives rational singularities of the normalization when
$H$ is not contracted, so finite duality and the canonical
trace give the injections
\[
 \mathcal O_{C_{\rm sch}}\longrightarrow
 \nu_*(\mathcal O_N/I_N)\longrightarrow\rho_*\mathcal O_D.
\]
The exact blowup relation $D'=\beta^*D''-E$ identifies
global functions on $D$ with those on $2H_0+V_0$ of the
minimal ruled model. The kernel of restriction to $2H_0$
has successive vertical-fibre layers
$\mathcal O_{\mathbb P^1}(-2)$, with no global sections.
Hence a global function is determined at the generic doubled
horizontal section. It follows that restriction from
$C_{\rm sch}$ to $C_{H,\rm sch}$ is injective on $H^0$
and surjective on $H^1$: the first assertion follows from
those injections and generic determination, and the second
from the dimension-one kernel of the restriction sheaf map.
Therefore
\[
                 g-1=\chi(\mathcal O_{C_{\rm sch}})
                    \le\chi(\mathcal O_{C_{H,\rm sch}}).
                                                            \tag{P4R.18.1}
\]
Every arrow here is the primary restriction of the same
conductor algebra; no other conductor component is deleted
from the left side.

\noindent\textbf{Step 4.} We evaluate the right side on $C_H^\nu=\mathbb P^1$.
Use the actual differentiated conormal, the saturated Hessian
line and the rank-one nilradical of P4R.10. The kernel $K$
of the normalized derivative has degree $2+\deg R$, since
$c_1(Y)\cdot C_H=0$. The Hessian line has degree
$\tfrac12\sum d_p$. Riemann--Roch and the normalization
sequence give exactly
\[
 \chi(\mathcal O_{C_{H,\rm sch}})
 =4-\sum_p\left[\delta(C_H,p)
      +\sum_{i\mapsto p}(d_i/2+b_i-r_i)+q_p\right].
                                                            \tag{P4R.18.2}
\]
All terms are those of this primary curve and its primary
conductor. They remain defined at its intersections with
the other component, because the pure conductor quotient
extends its generic containment of the symbolic square.

\noindent\textbf{Step 5.} Every singular point of $C_H$ contributes at least three
to the sum in (P4R.18.2). At a planar point the ideal is
$(u,h)$, its symbolic square equals its ordinary square,
and the actual equation belongs to that square. The
rank-one symmetric-square estimate of P4.42 and the exact
plane-curve different calculation give a contribution at
least $3\delta(C_H,p)+q_p\ge3$.
At a nonplanar point every defining equation has order at
least two. On a branch of multiplicity $m_i$ its projected
conormal order is at least $m_i$. The point is singular
on the entire critical curve, so P4R.15 makes the ambient
Hessian zero and gives $d_i\ge m_i$. Since $r_i=m_i-1$,
the contribution is at least
\[
                      \delta(C_H,p)+\sum_i(1+m_i/2),
\]
which is at least three at a singular curve point. The
same arguments give nonnegative contributions everywhere;
at a smooth primary point its conormal is the full kernel
and its only possible contribution is the nonnegative
Hessian zero order divided by two.
Combining (P4R.18.1)--(P4R.18.2) gives total primary
defect at most $5-g\le4$. Thus $C_H$ has at most one
singular point, and two different singular points are
impossible even in the presence of arbitrary transverse
types on the other global component.

\noindent\textbf{Step 6.} We count the possible branch values of the degree-two map
$H\to\mathbb P^1$ with the identifications retained.
The normalization graph formula is
\[
                         \sum_{p\in C}(r_p-1)=2.
\]
Connectedness leaves either two two-branch points or one
three-branch point. In the two-point case, two intersections
of the global components give two points on $C_H^\nu$.
There can be at most one additional unibranch singularity
of $C_H$, so these supply at most three possible branch
values. The other connected possibility is a self-identification
on one component and one intersection of the two components.
A self-identification on $C_H$ supplies two normalization
points and already costs at least three in the primary
defect. It permits no additional singular point, so with
the intersection there are again at most three possible
branch values. A self-identification on the other component
supplies no additional point on $C_H^\nu$.

\noindent\textbf{Step 7.} At a three-branch point, the two global components must
supply two and one branches. Thus $C_H^\nu$ has at most
two points over it. Two points already make $C_H$ singular
there and exclude any further singular point by the defect
bound. One point permits at most one further unibranch
singularity. This case also has at most three possible
branch values.

\noindent\textbf{Step 8.} There is no ramification elsewhere. At a point where the
entire critical curve is smooth, P4R.14 gives precisely
the product tacnode or $X(Y^2+XZ)$. Both have two distinct
normalization points over that point of the conductor.
Their two generic preimages lie on $H$, so properness
places both special preimages on $H$ as well. The
degree-two map is therefore unramified there. We have
accounted for every possible ramification value.
Riemann--Hurwitz for a degree-two map from a genus-$g$
curve requires $2g+2\ge4$ distinct branch values; each
ramified fibre consists of its single point of index two.
This contradicts the count of at most three. Thus the
generic germ along $C_H$ must be $D_4$, with degree two
or three as proved at the start.
\end{proof}

The exact connected identification patterns and their primary
defect/branch-value counts are enumerated in
\texttt{code/part4/horizontal\_primary.py}.


\subsection{A triple spectral root with nonzero cubic tensor}

\begin{proposition}[P4R.19]\label{night:P4R.19}\label{stmt:P4R.19}
\StatementRecord{night:P4R.19}{proved}{Historical events: GPT-6 Astra (OpenAI), ChatGPT Pro mode, effort Pro, S6\_06\_BLIND, 16--18 November 2026: targeted proof check. Also checked in this writing. Pre-continuation writing events: GPT-6 Astra (OpenAI), ChatGPT Pro mode, effort Pro, S6\_10\_PAPERS, 18 November 2026: targeted proof check. Continuation metadata events: GPT-6 Astra (OpenAI), ChatGPT Pro mode, effort Pro, S6\_10\_PAPERS continuation, 19 November 2026: used and rechecked in this run; status-vocabulary reconciliation; metadata reconciliation only; no new mathematical check. \emph{Pre-continuation recorded proof scope.} The night-run re-derivation was by GPT-6 Astra (OpenAI, ChatGPT Pro mode, effort Pro), 16--18 November 2026. Also checked in this writing by GPT-6 Astra (OpenAI), ChatGPT Pro mode, effort Pro, 18 November 2026. The exact cubic preparation and all three Newton slopes were checked. The explicit finite proper maps to two-node conic fibrations carry compact cycle classes of nonzero square back with their degrees; the repeated-factor balanced cubic remains connected. Anisotropic scaling transfers these classes to an actual Milnor fibre, contradicting P4.51}{\ref{night:P4.51}; \ref{night:P4.52}}

Let $f$ be the actual reduced local equation of the remaining
positive ruled support, and let $C_0$ be a smooth local branch
of its critical curve with generic transverse type $D_4$.
At a point where the cubic normal tensor along $C_0$ is
nonzero, that tensor is not the cube of a linear form.
Other critical-curve branches at the point are retained,
without a restriction on their generic transverse types.
\end{proposition}
\begin{proof}
\noindent\textbf{Step 1.} Choose coordinates $(x,y,z)$ with $C_0=(x=y=0)$.
Generic normal multiplicity three gives
$f\in(x,y)^{(3)}$. Powers of this smooth
complete-intersection ideal are primary, so
$f\in(x,y)^3$ at the specified point.

\noindent\textbf{Step 2.} Suppose its nonzero cubic normal tensor is a cube.
A linear normal-coordinate change makes it $x^3$, up
to a nonzero scalar. Weierstrass preparation in $x$
gives the exact equality
\[
 f=U(x,y,z)\bigl[x^3+a_2(y,z)x^2+a_1(y,z)x+a_0(y,z)\bigr],
\]
where $U$ is a holomorphic unit. The bracket belongs to
$(x,y)^3$, so $a_2,a_1,a_0$ are divisible by $y,y^2,y^3$,
respectively. Depressing the monic cubic is a holomorphic
coordinate change preserving $(x,y)$ and yields
\[
 f=U(x,y,z)\bigl[x^3+y^2B(y,z)x+y^3C(y,z)\bigr],
 \qquad B(0,0)=C(0,0)=0.                         \tag{R19.1}
\]
The unit remains in this equality and in the scaling
comparison below.

\noindent\textbf{Step 3.} Set
\[
 k=\operatorname{ord}_z B(0,z),\qquad
 l=\operatorname{ord}_z C(0,z),\qquad
 \rho=\min(k/2,l/3),
\]
where the order of an identically zero series is infinity.
At least one order is finite: otherwise all generic normal
cubics would be $x^3$, contrary to generic $D_4$.
Thus $0<\rho<\infty$. Choose $L>3\rho$ and positive
rational weights
\[
       \operatorname{wt}(x)=L+\rho,\qquad
       \operatorname{wt}(y)=L,\qquad
       \operatorname{wt}(z)=1.
\]
Multiplication by a common denominator gives positive
integer weights. The weighted degree of the first face
is $3L+3\rho$. Terms of $B$ containing $y$ have weight
gap at least $L-2\rho>0$, and those of $C$ containing
$y$ have gap at least $L-3\rho>0$. Consequently the
leading polynomial, apart from the nonzero factor $U(0)$,
is exactly one of
\[
\begin{array}{ll}
 P=x(x^2+a y^2z^k),\quad a\ne0,
       & k/2<l/3,\\[2pt]
 P=x^3+b y^3z^l,\quad b\ne0,
       & l/3<k/2,\\[2pt]
 P=Q(x,yz^r),\quad
 Q(X,V)=X^3+aXV^2+bV^3,\quad a,b\ne0,
       & k=2r, l=3r, r\ge1.
\end{array}                                                   \tag{R19.2}
\]
In the last case $r$ is an integer. The depressed cubic
$Q$ is not a cube: a cube with $X^3$ coefficient one and
$X^2V$ coefficient zero would have $a=b=0$. A double
factor of $Q$ is allowed throughout the argument.

\par\medskip\noindent\textbf{The compact class used in the comparison.}
Let $B$ be a connected smooth complex curve and let
$h:B\to\mathbb C$ have two distinct simple zeros $p,q$.
Join them by a smooth embedded arc $\gamma$ whose
interior contains no zero. On the smooth surface
\[
             Z_h=\{(b,u,v)\in B\times\mathbb C^2:
                                      uv=h(b)\},
\]
consider the locus over $\gamma$ with $|u|=|v|$.
It is a circle over each interior point and collapses
to a point at each end. Near an endpoint choose the
arc so that $h$ is a fixed nonzero complex multiple
of a positive real parameter $s$. Then
$u=\sqrt{s}\,e^{\mathrm i\phi}$ and
$v=c\sqrt{s}\,e^{-\mathrm i\phi}$ are the smooth
polar coordinates of a real disc in the local $uv$
chart. The locus is thus a smoothly embedded sphere.

It is Lagrangian for a product K\"ahler form.
Writing $u=R e^{\mathrm i\phi}$,
$v=R e^{\mathrm i(\psi-\phi)}$, the sum of the two
fibre area forms restricts to
$d(R^2/2)\wedge d\psi=0$, since $R,\psi$ depend
only on the arc parameter. The base area form also
restricts to zero. Multiplication by $\mathrm i$
identifies tangent and normal bundles, with the
opposite normal orientation relative to the complex
ambient orientation. Its self-intersection is
therefore $-\chi(S^2)=-2$.
For the compactly supported Poincar\'e dual $\alpha$,
\[
                         \int_{Z_h}\alpha^2=-2.             \tag{R19.3}
\]

If $\pi:W\to Z_h$ is a finite proper holomorphic map
of degree $d>0$ between smooth complex surfaces, then
$\pi^*\alpha$ is compactly supported and
\[
              \int_W(\pi^*\alpha)^2
                       =d\int_{Z_h}\alpha^2=-2d.            \tag{R19.4}
\]
The equality follows from the proper oriented degree;
off the analytic branch locus the map has $d$ positive
sheets, and the branch locus has measure zero for
integration of compactly supported top forms.
Poincar\'e duality gives a compact homology two-class
with that square. Its representative need not itself
be a sphere.

\par\medskip\noindent\textbf{The exact finite maps.}
In the first case of (R19.2), $P=1$ forces $x\ne0$
and gives
\[
           y^2z^k=h(x),\qquad
           h(x)=\frac{1-x^3}{ax},\qquad x\in\mathbb C^*.
\]
The map
\[
 (x,y,z)\longmapsto(x,u=y^2,v=z^k)
\]
is finite proper of degree $2k$ onto $uv=h(x)$.
The target is smooth: at $u=v=0$ the zeros of $h$
are simple, and elsewhere a derivative in $u$ or
$v$ is nonzero. The source is smooth by the weighted
Euler identity on $P=1$. Join two of the three cube
roots of one by an arc in $\mathbb C^*$ avoiding the
third. Equations (R19.3)--(R19.4) give a compact
two-class of square $-4k$ on the leading fibre.

In the second case the exact finite proper map is
\[
 \begin{split}
 \{x^3+b y^3z^l=1\}&\longrightarrow
                 \{uv=(1-x^3)/b\},\\
 (x,y,z)&\longmapsto(x,u=y^3,v=z^l).
 \end{split}
\]
It has degree $3l$, and the base is $\mathbb C$.
The same three simple zeros supply a pulled-back
compact class of square $-6l$.

For the balanced case let
\[
                 B_Q=\{(X,V)\in\mathbb C^2:Q(X,V)=1\}.
\]
It is smooth by the ordinary Euler identity and is
connected even when $Q$ has a double factor. Its
projection to $\mathbb P^1$ minus the distinct zero
directions of $Q$ is an unramified cyclic cover of
degree three. Around a zero direction of multiplicity
one or two its monodromy generates $\mathbb Z/3$.
Such a direction exists because $Q$ is not a cube,
so the cover is connected.

The function $V$ on $B_Q$ has precisely the three
zeros $(\zeta,0)$ with $\zeta^3=1$, and they are
simple because $Q_X(\zeta,0)=3\zeta^2\ne0$.
The surface
\[
             Z_Q=\{(b,u,w):b\in B_Q,\ uw=V(b)\}
\]
is smooth. The exact map
\[
 \begin{split}
       \{Q(x,yz^r)=1\}&\longrightarrow Z_Q,\\
       (x,y,z)&\longmapsto
                    \bigl(b=(x,yz^r),u=y,w=z^r\bigr)
 \end{split}                                                   \tag{R19.5}
\]
is finite proper of degree $r$: $z$ satisfies the
monic equation $z^r=w$ and the other coordinates are
retained. An arc in $B_Q$ joining two $V$-zeros and
avoiding the third supplies a compact class of square
$-2r$ on the source. In particular this argument
retains $Q=(X-\lambda V)^2(X+2\lambda V)$,
$\lambda\ne0$, without a perturbation of its
repeated factor.

\par\medskip\noindent\textbf{Transport to the actual local fibre.}
Let $(w_x,w_y,w_z)$ be the positive integer weights
and let $d$ be the resulting degree. Put
\[
 D_\epsilon(x,y,z)=
       (\epsilon^{w_x}x,\epsilon^{w_y}y,\epsilon^{w_z}z),
 \qquad
 G_\epsilon=\epsilon^{-d}f\circ D_\epsilon.
\]
These actual functions converge with derivatives on
compact sets to $U(0)P$. The weighted Euler identity
makes $dP$ nonzero on $P=1$, including the repeated-factor
case on its zero set. Choose compact representatives
of the two-class above and of the transverse
perturbation computing its square. On a compact
neighbourhood of these chains, the gradient-flow
comparison of P4.52 solves
\[
       dG_\epsilon(V_\epsilon)
                         =-\partial_\epsilon G_\epsilon.
\]
The denominator is bounded away from zero on a
slightly larger neighbourhood; a cutoff gives a
flow for sufficiently small positive $\epsilon$.
It transports the chains and their intersections
to $G_\epsilon=U(0)$. Scaling by $D_\epsilon$ then
embeds them in the original fixed Milnor ball on
\[
                         f=U(0)\epsilon^d.
\]
The oriented intersection number is preserved and
equals $-4k$, $-6l$, or $-2r$. Each is nonzero,
contradicting P4.51. This proves the assertion.
\end{proof}

\par\medskip\noindent\textbf{The two-direction exception is retained.}
The assertion cannot be enlarged to every spectral
collision. The reduced equation
\[
                       f=y(x^2+y^2z^2)
\]
has generic $D_4$ along $x=y=0$ and generic $A_1$
along $x=z=0$. Only two spectral directions collide
at their intersection. Its non-plane normalization
sheets $x=\pm\mathrm i yz$ have pole divisor
$2(y=0)+(z=0)$, with simple normal crossings.
Its affine fibre $f=1$ is exactly $(\mathbb C^*)^2$:
\[
 \begin{split}
       a&=x+\mathrm i yz,\qquad b=x-\mathrm i yz,\\
       x&=(a+b)/2,\qquad y=(ab)^{-1},\qquad
       z=(a-b)ab/(2\mathrm i).
 \end{split}
\]
The displayed maps are inverse and substitution gives
$f=1$. Its intersection form is zero: the product
of two unit circles generates its second homology,
and a radial translation gives a disjoint homologous
copy. More generally the fibre of $y(x^2+y^2z^n)$
becomes $u^2+z^n=t^3$, $t\ne0$, through
$t=1/y$, $u=x/y$. Positive radial dilation gives
a product real direction, along which any compact
cycle can be displaced from another compact cycle.
Its intersection form is also zero. Thus the
other generic transverse type remains unrestricted.

The exact weighted gaps, all three finite-map
equations, their degree-square identities, and the
inverse maps of the two-direction exception are
executed in \path{code/part4/d4_spectral_triple.py};
its \path{.json} and \path{.out} outputs report PASS.


\subsection{The three-branch image point with a double horizontal cover}

\begin{proposition}[P4R.20]\label{night:P4R.20}\label{stmt:P4R.20}
\StatementRecord{night:P4R.20}{proved}{Historical events: GPT-6 Astra (OpenAI), ChatGPT Pro mode, effort Pro, S6\_06\_BLIND, 16--18 November 2026: targeted proof check. Also checked in this writing. Pre-continuation writing events: GPT-6 Astra (OpenAI), ChatGPT Pro mode, effort Pro, S6\_10\_PAPERS, 18 November 2026: targeted proof check. Continuation metadata events: GPT-6 Astra (OpenAI), ChatGPT Pro mode, effort Pro, S6\_10\_PAPERS continuation, 19 November 2026: used and rechecked in this run; status-vocabulary reconciliation; metadata reconciliation only; no new mathematical check. \emph{Pre-continuation recorded proof scope.} The night-run re-derivation was by GPT-6 Astra (OpenAI, ChatGPT Pro mode, effort Pro), 16--18 November 2026. Also checked in this writing by GPT-6 Astra (OpenAI), ChatGPT Pro mode, effort Pro, 18 November 2026. The two-component ruled spectral calculation, graph path, endpoint ribbon normalization, degree-one internal sheets, and incompatible conductor exponents on the same vertical sheet were re-derived. The one-component graph alternative was checked separately. A literal quad typo in the source is a typesetting correction only}{\ref{night:P4.55}}

In the remaining degree-two horizontal sector after P4.55,
suppose that the generic transverse curve along the other
image component $B$ has at least three normalization branches.
The unique image-graph cycle cannot be concentrated at one
point with three local critical-curve branches, namely the
smooth branch $A$ and two local branches of $B$.
\end{proposition}
\begin{proof}
\noindent\textbf{Step 1.} Write these two local branches as $B_+,B_-$ and the
point as $p$. Let $W_1,\ldots,W_c$ be the normalized
conductor components mapping onto $B$, with degrees
$d_i>0$ over $B^\nu$. Put
\[
                         r_B=\sum_i d_i\ge3.
\]
This is the number of generic normalization branches
of the full transverse plane curve along $B$.
The normalization conductor has components
$H,R,W_1,\ldots,W_c$ and a tree graph. Its total
point-identification contribution is $c+1$.
The horizontal section $H$ maps with degree two to
$A$, while the rational vertical curve $R$ maps
with degree one. Distinct points of $H$ are not
identified in the normalization $N$. Distinct
points of $R$ are not identified either: a
contracted vertical tree effecting such an
identification would create a cycle in the ruling
fibre. All conductor branches below are counted
on their normalization.

\noindent\textbf{Step 2.} We first record the exact spectral bound used in
the proof. Let $E=N_{A/Y}^*$, so $\deg E=2$ by
adjunction and $c_1(Y)\cdot A=0$. On $\mathbb P(E)$
use the quotient convention, tautological class
$\xi$, and fibre class $F$. Thus
\[
          \xi^2=2,\qquad \xi F=1,\qquad F^2=0,
          \qquad K_{\mathbb P(E)}=-2\xi.
\]
The actual normal derivative maps of P4.55 identify
$H$ with the normalization of the degree-two
spectral component $\Gamma_H$, and identify the
third sheet with a section $\Gamma_R\cong A$.
Remove the common zero divisor of the cubic tensor,
of degree $b\ge0$. For an integer $t$ write
\[
       \Gamma_R=\xi-tF,\qquad
       \Gamma_H=2\xi+(t-b)F.
\]
If $\delta_H$ is the normalization colength of
$\Gamma_H$ and $n=\Gamma_H\cdot\Gamma_R\ge0$,
adjunction and the intersection product give
\[
 \begin{split}
       g(H)+\delta_H&=p_a(\Gamma_H)=1+t-b,\\
       n&=4-t-b,\\
       g(H)+\delta_H+n&=5-2b.
 \end{split}                                                   \tag{R20.1}
\]
These are the actual residual components of the
cubic, with their common zero divisor retained.
Since $g(H)\ge1$, they imply $b\le2$.

\par\medskip\noindent\textbf{The exact graph when $c\ge2$.}
Let $r_A$ be the number of normalized conductor
points over the branch $A$ at $p$. It is two if
$H\to A$ ramifies there, and three otherwise.
Let $r_+,r_-$ count the points over $B_+,B_-$.
Every $W_i$ is surjective to $B^\nu$, so
$r_+,r_-\ge c$. Let $n_p$ be the number of
normalization points of $S$ over $p$.
Their nonempty supports in $\{A,+,-\}$ pairwise
intersect: distinct local hypersurface components
intersect in a curve.

\noindent\textbf{Step 3.} A singleton support $\{A\}$ forces every point to
carry $A$. Then $n_p\le r_A$ and the graph
contribution is at least $r_++r_-\ge2c>c+1$.
A singleton $\{+\}$ forces every point to carry
$+$ and gives contribution at least
$r_A+r_-\ge c+2$; the same applies to $\{-\}$.
Thus every support has at least two elements, and
\[
 \text{contribution at }p
       \ge\frac{r_A+r_++r_-}{2}\ge c+1.
\]
The whole global budget is $c+1$. Equality is
forced throughout:
\[
                 r_A=2,\qquad r_+=r_-=c,\qquad n_p=c+1.
                                                               \tag{R20.2}
\]
Every point over $p$ has exactly two normalized
conductor branches, with different image labels,
and all graph identifications occur at $p$.
The horizontal cover ramifies over $p$.
Each $W_i$ has precisely one point over each of
$B_+,B_-$ and has valence two; $H$ and $R$ have
valence one. The connected tree is therefore
the path
\[
                    H-W_1-W_2-\cdots-W_c-R.             \tag{R20.3}
\]
Choose the label $+$ at the $H$ endpoint. The next
point on $W_1$ has label $-$. Every internal point
has support $\{+,-\}$, so it joins that point to
the $+$ point of $W_2$. Induction gives the $+$
point on the left of every $W_i$ and its $-$
point on the right. The $R$ endpoint carries $-$.

\noindent\textbf{Step 4.} The endpoint hypersurface factors, denoted
$f_H,f_R$, thus have supports $\{A,+\}$ and
$\{A,-\}$. The $c-1$ internal factors $f_j$
have supports $\{+,-\}$. These are all the local
irreducible hypersurface factors.

\par\medskip\noindent\textbf{The actual length-two ribbon.}
The set-theoretic intersection of $f_H$ and $f_R$
is $A$ alone. Along generic $A$, the first factor
supplies two ordinary-triple branches and the
second supplies the third, so their intersection
multiplicity is two. Their one-dimensional
complete-intersection ring is Cohen--Macaulay,
without a zero-dimensional embedded component.
A coordinate $z$ on the smooth $A$ is a
nonzerodivisor. The ring is finite free of rank
two over $\mathbb C\{z\}$, and its quotient by
$z$ has length two.

\noindent\textbf{Step 5.} If both factors had no linear term, their
restrictions modulo $z$ would belong to
$(x,y)^2$. The Artinian quotient would contain
the independent classes $1,x,y$, contrary to
length two. At least one factor is smooth.
It cannot be $f_H$: a smooth surface is its own
normalization and supplies only one generic
preimage of $A$, whereas this factor supplies
the ramified double cover. Thus $f_R$ is smooth.
Take the actual ambient coordinate $u=f_R$.
On $u=0$, the divisor of $f_H$ is exactly twice
the smooth $A$, so in coordinates
$A=(u=v=0)$ it has restriction $v^2$ times a
unit. Weierstrass preparation in $v$, followed
by depression and using $f_H\in(u,v)^2$, gives
the exact expression
\[
                    f_H=\text{unit}\,(v^2+u^2g(u,z)).     \tag{R20.4}
\]
All changes preserve $A$. The full cubic normal
tensor is therefore, up to a nonvanishing scalar,
\[
        \left(\prod_j f_j(0,0,z)\right)
                     u\bigl(v^2+u^2g(0,z)\bigr).          \tag{R20.5}
\]
The coefficient of $uv^2$ in the last factor is
one, so the common zero order is exactly the
sum of the orders of $f_j|_A$.

\par\medskip\noindent\textbf{The pole on the smooth rational branch.}
On the smooth branch $f_R=0$, its conductor is
the restriction of the product of the other
factors. The factor $f_H$ contributes $2A$.
Each $f_j$ cuts out only $B_-$ there. This
normalization branch is already smooth, and
the pole support on its minimal resolution is
SNC. Hence $A$ and its $B_-$ curve are smooth
and transverse. Choose the latter's equation
$z=0$ on this surface, with $z$ still a
coordinate on $A$, and extend it to the ambient
threefold. Then
\[
             f_j|_{f_R=0}=\text{unit}\,z^{e_j},
                  \qquad e_j\ge1.
\]
Their sum $e=\sum e_j$ is the generic conductor
exponent of the global component $W_c$ on
this branch. By (R20.5) it is also the local
common cubic-zero order. The branch on
$f_R=0$ is smooth; the plane-curve conductor
formula therefore gives
\[
          e=\sum_{\text{other branches}}I
                         \ge r_B-1\ge2.
\]
Combining $b\ge e$ with (R20.1) forces
\[
     b=e=2,\quad r_B=3,\quad
     g(H)=1,\quad \delta_H=n=0.                         \tag{R20.6}
\]
Thus $\Gamma_H$ is smooth and disjoint from
$\Gamma_R$, and there is no further common
zero of the cubic.

\par\medskip\noindent\textbf{The exact normalization at the ramified endpoint.}
The cover $H\to A$ ramifies at $p$ by (R20.2).
Equation (R20.4) presents its spectral cover as
$w^2+g(0,z)=0$. After translating its double
root, smoothness from (R20.6) gives
$g_z(0,0)\ne0$. Hence
\[
 \begin{split}
        N_H&=\{w^2+g(u,z)=0\},\\
        (u,w,z)&\longmapsto(u,v=uw,z)
 \end{split}                                                   \tag{R20.7}
\]
is the actual normalization near that point.
It is smooth, finite, and birational onto the
branch (R20.4); off $u=0$ its inverse is
$w=v/u$. Coordinates $(u,w)$ solve
$z=\psi(u,w)$, with $z|_H$ of order two in
$w$. The minimal normalization resolution
has no exceptional curve over this smooth
point, and $H$ is $u=0$.

\noindent\textbf{Step 6.} The other conductor branch $W_1$ meets $H$
transversely by SNC. Thus $u$ is a coordinate
on $W_1$ and has order one on its image
$B_+$. By (R20.2), $W_1$ has only one point
over the normalized $B_+$ point, so its
global degree is $d_1=1$. The same smooth
SNC argument on $f_R=0$ gives $d_c=1$.
Since $\sum d_i=r_B=3$, the case $c=2$
would have sum two, while $c>3$ is impossible.
Therefore
\[
                 c=3,\qquad d_1=d_2=d_3=1.              \tag{R20.8}
\]
There are exactly two internal factors, and
$e_1=e_2=1$.

\par\medskip\noindent\textbf{The conductor mismatch.}
Each $f_j$ intersects $f_H$ only along $B_+$.
Its pullback to the actual smooth normalization
(R20.7) consequently has divisor supported
only on $W_1$. Its restriction to $H$ vanishes
to order two: the order on $A$ is one and
$H\to A$ has ramification index two.
Transversality of $H$ and $W_1$ makes the
divisor exactly $2W_1$. Thus the generic
intersection multiplicity of $f_H$ with
each internal factor along $B_+$ is two.

\noindent\textbf{Step 7.} The path (R20.3) identifies the same global
component $W_3$ on the smooth $R$ branch
over $B_-$ and on an internal factor over
$B_+$. At $B_-$ its conductor exponent is
$e=2$. At $B_+$ it has intersection
multiplicity two with $f_H$ and positive
intersection multiplicity with the other
internal factor. Its conductor exponent
is at least three. These cannot be two
values of the coefficient on one global
irreducible conductor component.

\noindent\textbf{Step 8.} The full marking is also explicit. Over
$B_+$ the factors $(f_H,f_1,f_2)$ carry
$(W_1,W_2,W_3)$; over $B_-$ the factors
$(f_1,f_2,f_R)$ carry the same ordered
global components. If $p_+,p_-\ge1$ are
the intersections of the internal factors,
the conductor vectors are
\[
         (4,2+p_+,2+p_+),\qquad
         (1+p_-,1+p_-,2).
\]
There is no missing intrinsic branch-delta
term: the exponent two of the smooth $R$
branch forces its two intersections to
equal one, which makes both other generic
$B$ branches smooth as well. The third
entries already give the contradiction.

\par\medskip\noindent\textbf{The case $c=1$.}
Now the conductor graph has $H,R,W$ and
budget two. The points of $W$ over $B_+$
and $B_-$ cannot be identified in $N$,
since that creates a loop in the conductor
tree. Points carrying the two labels are
distinct; their supports intersect, so
both carry $A$. Each contributes at least
one, exhausting the graph budget. Each
has exactly one $A$ and one $B$ branch,
and no other identification occurs.

\noindent\textbf{Step 9.} If $H$ ramifies over $p$, these two
$A$ branches are $H$ and $R$. There are
only the factors with supports
$\{A,+\}$ and $\{A,-\}$. The length-two
ribbon argument makes $f_R$ smooth.
No other factor contains $B_-$, and this
smooth branch has no singularity there.
Thus $B_-$ cannot be a critical curve.

\noindent\textbf{Step 10.} If $H$ is unramified, the three normalized
$A$ preimages are distinct. Two occur
at the $B$ attachments and the third
is an $A$-only normalization point.
There are three factors, each supplying
one generic $D_4$ branch. Every pair
meets only along $A$, with multiplicity
one. Its complete-intersection curve
is the reduced smooth $A$: it is pure,
generically reduced, and has no other
component. Modulo $z$, length one forces
independent normal linear terms.
All three factors are smooth near $p$
and their only pairwise intersection is
$A$. Neither $B_+$ nor $B_-$ can be a
critical branch. This final contradiction
completes the proof.
\end{proof}

The exact spectral intersections, length-two
ribbon basis, normalization equation, scalar
cubic coefficient, endpoint-degree list and
marked conductor vectors are executed in
\path{code/part4/degree_two_triple_point.py}.
Its \path{.json} and \path{.out} certificates
report PASS.

\section{Global CDP replacement and normalization resonance}\label{sec:II:12}


\begin{proposition}[A geometric replacement on an explicit fibre class]\label{cdp:st:C4a}
\StatementRecord{cdp:st:C4a}{proved}{Historical events: Claude Fable 5.2 (Anthropic), effort high, 16 November 2026: carried as background without fault found. GPT-6 Astra (OpenAI), ChatGPT Pro mode, effort Pro, S6\_09\_PAPERS, 17 November 2026: targeted proof check. Also checked in this writing. Pre-continuation writing events: GPT-6 Astra (OpenAI), ChatGPT Pro mode, effort Pro, S6\_10\_PAPERS, 18 November 2026: targeted proof check. Continuation metadata events: GPT-6 Astra (OpenAI), ChatGPT Pro mode, effort Pro, S6\_10\_PAPERS continuation, 19 November 2026: used and rechecked in this run; status-vocabulary reconciliation; metadata reconciliation only; no new mathematical check. \emph{Pre-continuation recorded proof scope.} The exact full-fibre vanishing test, positive base twists, and endpoint sequence were re-derived}{the hypotheses, notation and definitions stated here; \cref{cdp:st:A.4.5}; \cref{cdp:st:A.1.2}; \cref{cdp:st:A.4.6}; \cref{cdp:st:C2a}}

Let $f:Y\to\PP^1$ be a connected-fibre holomorphic map from a compact connected complex threefold with $b_1(Y)=b_2(Y)=0$, and smooth general fibre a complex two-torus. Suppose every full fibre has one of the following forms: a reduced irreducible normal-crossing surface; or a smooth torus family modulo a finite free cyclic action over a cyclic base cover, so that its reduction is smooth and its multiplicity is the covering order. Let
\[
 0\ne\gamma\in H^0(\PP^1,R^1f_*\Z)
\]
be a generator, and suppose that, for every singular reduced full fibre $S$, its pullback to $H^1(\widetilde S,\Q)$ under normalization is nonzero. Then a dense open subset of every component of $\Pic(Y)$ satisfies
\[
 H^q(Y,T_Y\otimes L)=0\quad(0\le q\le3),\qquad
 \chi(Y,T_Y\otimes L)=0,\qquad c_3(Y)=0.                \tag{C.15}\label{cdp:eq:C.15}
\]
The generator exists under the stated rational Betti-number assumptions. Its normalization condition is a condition on an explicit integral cohomology class, with no tangent-cohomology vanishing assumed.
\end{proposition}


\begin{proof}
\noindent\textbf{Step 1.} The integral Leray sequence gives an injection
\[
H^0(\PP^1,R^1f_*\Z)\longrightarrow H^2(\PP^1,\Z)=\Z
\]
whose cokernel injects into the finite group $H^2(Y,\Z)$. Thus its source is infinite cyclic. The degree-one specialization argument of Part D identifies its stalk at every fibre with a nonzero integral class: a class in the preimage of a small disc that vanishes on a nearby fibre comes from the meridian of the punctured disc; a normal disc to any irreducible component makes a nonzero multiple of that meridian nullhomotopic, forcing the class to vanish. The real-to-holomorphic cohomology map on a compact connected complex space is injective on real classes, as proved there by the pluriharmonic maximum principle. Coherent Leray therefore supplies $\beta\in H^1(Y,\OO_Y)$ whose restriction to every full fibre is the image of $\gamma$.

\noindent\textbf{Step 2.} Fix a real irrational number $r$ and put $M_r=\exp(r\beta)$, with the exponential normalized by $2\pi i$. On each fibre this line is represented by the unitary character $\exp(2\pi ir\gamma)$. It is nontrivial, and on a singular fibre its normalization pullback is nontrivial by the hypothesis on $\gamma$. To see this directly, triviality of an $n$th tensor power would place $nr\gamma$ in the integral lattice after applying $H^1(\R)\to H^1(\OO)$; real injectivity makes this an equality in $H^1(\R)$, which is impossible for an irrational multiple of a nonzero integral vector.

\noindent\textbf{Step 3.} The same construction works in every Picard component. Each component contains a finite-order line bundle $P$: for a representative $P_0$, a positive power $P_0^m$ lies in the complex vector group $\Pic^0(Y)$, where division by $m$ is possible. Dividing $P_0$ by the resulting line makes its $m$th power trivial. Replacing $M_r$ by $M_r\otimes P$ retains nontrivial unitary restriction on each fibre and on each specified normalization. Indeed an $m$th power eliminates $P$ and reduces a purported trivialization to the preceding irrationality contradiction.

\noindent\textbf{Step 4.} Set $L=K_Y\otimes(M_r\otimes P)^{-1}$. On every full fibre $F$ the Serre-dual bundle for degree two is
\[
 (T_Y\otimes L)|_F^*\otimes\omega_F
       =\Omega_Y^1|_F\otimes(M_r\otimes P)|_F,
\]
because $F$ is a Cartier fibre and $\omega_F=K_Y|_F$. At a smooth torus fibre both the conormal line and the intrinsic cotangent bundle are trivial; nontrivial unitary degree-zero lines on the torus have zero cohomology. The ambient space of sections is therefore zero. At a reduced irreducible normal-crossing fibre, Statement \ref{cdp:st:A.4.5} proves the torsion-free cotangent vanishing for every nontrivial unitary line. The normal-crossing kernel in \ref{cdp:st:A.1.2} is the normalization pushforward, and its sections vanish because the pulled-back unitary line is nontrivial. A holomorphic section of a unitary line on a compact connected complex manifold has constant norm by the maximum principle and is either zero or a trivialization. Thus the ambient group vanishes there as well.

\noindent\textbf{Step 5.} On a cyclic-cover neighbourhood of a smooth multiple fibre, the covering map to $Y$ is finite etale. The pulled-back line restricts nontrivially to the covering central torus: finite etale pullback is injective on rational degree-one cohomology by its trace, and $r$ is irrational. The scheme-theoretic inverse image of the full fibre is the corresponding thickening of this torus. Each layer in its ideal filtration is the same ambient cotangent restriction tensored with this nontrivial line, since the covering base coordinate trivializes the conormal. Every layer has zero cohomology. The exact filtration sequences therefore give zero sections on the full thickening. Taking finite-group invariants, an exact functor, proves the required ambient vanishing on the original full fibre.

\noindent\textbf{Step 6.} Top-degree coherent base change now gives $R^2f_*(T_Y\otimes L)=0$. For sufficiently large integer $n$, Serre vanishing on the base gives
\[
H^1\bigl(\PP^1,R^1f_*(T_Y\otimes L)\otimes\OO(n)\bigr)=0.
\]
Projection formula and the two-column Leray sequence imply
$H^2(Y,T_Y\otimes L\otimes f^*\OO(n))=0$.
This is an actual twist in the same Picard component up to the fixed finite-order class of $f^*\OO(n)$; choosing $n$ divisible by the finite order of its Chern class keeps the component unchanged. Semicontinuity in the explicit Picard family of C.1 makes the degree-two zero locus dense open in that component. Intersect it with the dense open endpoint-vanishing loci from C.1. On this intersection $\chi(T_Y\otimes L)=-h^1\le0$.

\noindent\textbf{Step 7.} Statement \ref{cdp:st:A.4.6} gives Euler number zero or two for each reduced irreducible normal-crossing fibre; the positive case has rational normalization. Our normalization hypothesis excludes that case, since rational smooth surfaces have $H^1(-,\Q)=0$. Smooth torus fibres and the smooth finite etale quotients in the second fibre class have Euler number zero. Therefore \ref{cdp:st:C2a} gives $c_3(Y)=0$. The index formula in C.1 gives $\chi=0$, forcing $h^1=0$ on the same dense open set. This proves all of \eqref{cdp:eq:C.15}.
\end{proof}






\begin{proposition}[An unavoidable failure on a rational homology six-sphere]\label{cdp:st:C4b}
\StatementRecord{cdp:st:C4b}{proved}{Historical events: Claude Fable 5.2 (Anthropic), effort high, 16 November 2026: carried as background without fault found. GPT-6 Astra (OpenAI), ChatGPT Pro mode, effort Pro, S6\_09\_PAPERS, 17 November 2026: targeted proof check. Also checked in this writing. Pre-continuation writing events: GPT-6 Astra (OpenAI), ChatGPT Pro mode, effort Pro, S6\_10\_PAPERS, 18 November 2026: targeted proof check. Continuation metadata events: GPT-6 Astra (OpenAI), ChatGPT Pro mode, effort Pro, S6\_10\_PAPERS continuation, 19 November 2026: used and rechecked in this run; status-vocabulary reconciliation; metadata reconciliation only; no new mathematical check. \emph{Pre-continuation recorded proof scope.} The local root-cover and finite Picard-component reductions in the generic criterion were re-derived}{the hypotheses, notation and definitions stated here; \cref{cdp:st:C4a}; \cref{cdp:st:C2a}; \cref{cdp:st:A.4.6}}

Let $f:Y\to\PP^1$ be a connected-fibre holomorphic map from a compact complex threefold with $b_1=b_2=b_3=0$. For every $L\in\Pic(Y)$ there is a scheme-theoretic fibre $F$ with
\[
 H^2(F,(T_Y\otimes L)|_F)\ne0,\qquad
 H^0(F,\Omega_Y^1|_F\otimes(K_Y\otimes L^{-1})|_F)\ne0. \tag{C.16}\label{cdp:eq:C.16}
\]
For torus general fibre the reduced fibre Euler numbers sum to two. Within the fibre class of \ref{cdp:st:C4a}, exactly one fibre has rational normalization; it is an irreducible reduced normal-crossing fibre of Euler number two. For that normalization,
\[
H^2(\widetilde S,K_{\widetilde S}\otimes
             \eta^*(K_Y^{-1}\otimes L))=\C
\quad\hbox{for every }L\in\Pic(Y).                    \tag{C.17}\label{cdp:eq:C.17}
\]
\end{proposition}


\begin{proof}
\noindent\textbf{Step 1.} The divisor of any fibre is nonzero and effective. The index argument C.1 therefore gives $H^2(Y,T_Y\otimes M)\ne0$ for every line bundle $M$. Suppose $R^2f_*(T_Y\otimes L)=0$. Twisting by a sufficiently large positive pullback from the base kills $H^1(R^1f_*)$, so Leray would give a line bundle $M=L\otimes f^*\OO(n)$ with $H^2(Y,T_Y\otimes M)=0$, a contradiction. Thus $R^2f_*(T_Y\otimes L)$ is nonzero. Its stalk has nonzero residue at some point, and top-degree base change gives the first nonvanishing in \eqref{cdp:eq:C.16}. Cartier-fibre adjunction and Serre duality give the second, with the exact line bundle displayed there.

\noindent\textbf{Step 2.} For torus general fibre \ref{cdp:st:C2a} gives $\sum e(F_{\rm red})=e(Y)=2$. In the class of \ref{cdp:st:C4a}, \ref{cdp:st:A.4.6} makes every singular reduced contribution either zero or two, with two exactly for the rational hexagonal normalization; smooth finite torus quotients have zero Euler number. There is consequently exactly one positive fibre. Its smooth rational normalization has $\Pic^0=0$ and torsion-free integral second cohomology. Every ambient line bundle has torsion first Chern class, whose pullback is therefore zero. In particular both $\eta^*K_Y$ and $\eta^*L$ are trivial. Serre duality on the connected rational normalization gives \eqref{cdp:eq:C.17}.
\end{proof}






\begin{proposition}[The finite singular-fibre test and a fixed unavoidable fibre]\label{cdp:st:C4c}
\StatementRecord{cdp:st:C4c}{proved}{Historical events: Claude Fable 5.2 (Anthropic), effort high, 16 November 2026: carried as background without fault found. GPT-6 Astra (OpenAI), ChatGPT Pro mode, effort Pro, S6\_09\_PAPERS, 17 November 2026: targeted proof check. Also checked in this writing. Pre-continuation writing events: GPT-6 Astra (OpenAI), ChatGPT Pro mode, effort Pro, S6\_10\_PAPERS, 18 November 2026: targeted proof check. Continuation metadata events: GPT-6 Astra (OpenAI), ChatGPT Pro mode, effort Pro, S6\_10\_PAPERS continuation, 19 November 2026: used and rechecked in this run; status-vocabulary reconciliation; metadata reconciliation only; no new mathematical check. \emph{Pre-continuation recorded proof scope.} The distinction between generic vanishing and vanishing of the R2 torsion sheaf, together with Leray, was re-derived}{the hypotheses, notation and definitions stated here; \cref{cdp:st:A.1.1}; \cref{cdp:st:A.3.1}; \cref{cdp:st:A.3.5}; \cref{cdp:st:D.1.11}}

Let $f:Y\to\PP^1$ be a proper connected-fibre map from a compact complex threefold with $b_1(Y)=b_2(Y)=0$ and smooth general fibre a complex two-torus. Let $\Sigma$ be its finite set of critical values, and fix a connected component $P$ of $\Pic(Y)$. For $c\in\Sigma$ retain the entire Cartier fibre $F_c$, including its multiplicities, and put
\[
 Z_c=\{L\in P:H^0(F_c,\Omega_Y^1|_{F_c}\otimes
                  (K_Y\otimes L^{-1})|_{F_c})\ne0\}.
                                                               \tag{C.18}\label{cdp:eq:C.18}
\]
Then the following are equivalent: every $Z_c$ is a proper subset of $P$; and $H^2(Y,T_Y\otimes L)=0$ on a dense open subset of $P$. In that case the generic twist in $P$ satisfies
\[
 H^q(Y,T_Y\otimes L)=0\quad(q\ne1),\qquad
 \chi(Y,T_Y\otimes L)\le0,\qquad c_3(Y)\le0.          \tag{C.19}\label{cdp:eq:C.19}
\]
For a reduced fibre, membership in the complement of $Z_c$ is exactly the two explicitly typed kernel-and-connecting-map conditions of \ref{cdp:st:A.1.1}, with kernel $N^*\otimes\mathcal J^{-1}$ computed in A.2. For $F_c=mS$, the equivalent complete test is the finite lifting filtration of \ref{cdp:st:A.3.1}. These tests use only sheaves and maps on the finitely many singular full fibres.

When $b_3(Y)=0$, at least one fixed critical fibre $F_c$ satisfies $Z_c=P$. Its support is irreducible, by \ref{cdp:st:A.3.5}. Under the integral assumption $H^2(Y,\Z)=0$, the Picard group is connected, so this is a single fixed fibre for which the ambient group in \eqref{cdp:eq:C.18} is nonzero for every line bundle on $Y$.
\end{proposition}


\begin{proof}
\noindent\textbf{Step 1.} The universal Picard family constructed in C.1 restricts to a holomorphic family on the proper space $F_c\times P$. Its underlying sheaf is flat over $P$, because it is a line-bundle twist of the fixed locally free restriction $\Omega_Y^1|_{F_c}$. Proper semicontinuity therefore makes $Z_c$ a closed complex analytic subset of $P$.

\noindent\textbf{Step 2.} Assume that each $Z_c$ is proper. The complement of their finite union is open dense and has full Lebesgue measure in $P\simeq\C^q$. Statement \ref{cdp:st:D.1.11} supplies a dense set of full measure of line bundles non-torsion on every full fibre. Choose $L$ in the intersection. On every smooth torus fibre, $K_Y|_F$ and the conormal are trivial. The intrinsic cotangent bundle is trivial as well, so the conormal exact sequence expresses the ambient cotangent restriction as an extension of trivial bundles. The line $L|_F$ has zero integral Chern class, since $c_1(L)$ is torsion and $H^2(F,\Z)$ is torsion-free; it is nontrivial by the chosen non-torsion condition. Its cohomology in every degree vanishes, as follows by the unitary Fourier resolution of a topologically trivial line on a torus. The ambient section group is consequently zero at every smooth fibre. At each critical fibre it is zero by the definition of the chosen complement.

\noindent\textbf{Step 3.} Serre duality on each Gorenstein Cartier fibre, with $\omega_{F_c}=K_Y|_{F_c}$, and top-degree base change give $R^2f_*(T_Y\otimes L)=0$. Twisting by $f^*\OO(n)$ preserves every fibre restriction. For sufficiently large $n$, it kills $H^1(\PP^1,R^1f_*(T_Y\otimes L)\otimes\OO(n))$. We choose $n$ divisible by the finite order of $c_1(f^*\OO(1))$, so this twist stays in $P$. The two-column Leray sequence gives an actual line bundle in $P$ with zero $H^2$. Proper semicontinuity makes this zero locus dense open. Intersecting it with the dense open endpoint-vanishing loci of C.1 proves \eqref{cdp:eq:C.19}, using $\chi=c_3/2$.

\noindent\textbf{Step 4.} Conversely, suppose $H^2$ vanishes on a dense open subset of $P$, and choose $L$ in that subset and in the full-measure set of \ref{cdp:st:D.1.11}. The preceding smooth-fibre argument makes $R^2f_*(T_Y\otimes L)$ a torsion sheaf. Leray gives a surjection
\[
 H^2(Y,T_Y\otimes L)\longrightarrow
 H^0(\PP^1,R^2f_*(T_Y\otimes L)).
\]
Thus that torsion sheaf is zero. Top-degree base change and fibre duality imply $L\notin Z_c$ for every $c$. In particular none of the $Z_c$ equals $P$, proving the equivalence.

\noindent\textbf{Step 5.} When $b_3=0$, C.1 gives $H^2(Y,T_Y\otimes L)\ne0$ for every $L$. The equivalence therefore forces some $Z_c$ to be all of $P$. There are only finitely many $c$, and $P$ is an irreducible vector-space torsor; a finite union of proper closed analytic subsets cannot exhaust it. The integral hypothesis makes $P=\Pic(Y)$ by the exponential sequence. This proves the assertion about a fixed fibre. Finally the identifications in A.1--A.3 prove the claimed explicit forms of the fibre tests, without replacing a thickened fibre by its support.
\end{proof}



\ref{cdp:st:C4a} is an explicit geometric instance of \ref{cdp:st:C4c}. The two vanishing groups on the support in \ref{cdp:st:B1}--\ref{cdp:st:B2} do not give its fibre test: their ambient space is one-dimensional. The compact quotient \ref{cdp:st:A.3.2} proves that the full multiple fibre cannot be replaced by one unshifted reduced-support test. \ref{cdp:st:C4c} supplies the fixed full-fibre obstruction in arbitrary torus fibrations. \ref{cdp:st:C4e} proves that its normalization target is nonzero for every twist, with arbitrary singularities and multiplicities.




\begin{proposition}[A finite transgression test for every CDP general-fibre type]\label{cdp:st:C4d}
\StatementRecord{cdp:st:C4d}{proved}{Historical events: Claude Fable 5.2 (Anthropic), effort high, 16 November 2026: carried as background without fault found. GPT-6 Astra (OpenAI), ChatGPT Pro mode, effort Pro, S6\_09\_PAPERS, 17 November 2026: targeted proof check. Also checked in this writing. Pre-continuation writing events: GPT-6 Astra (OpenAI), ChatGPT Pro mode, effort Pro, S6\_10\_PAPERS, 18 November 2026: targeted proof check. Continuation metadata events: GPT-6 Astra (OpenAI), ChatGPT Pro mode, effort Pro, S6\_10\_PAPERS continuation, 19 November 2026: used and rechecked in this run; status-vocabulary reconciliation; metadata reconciliation only; no new mathematical check. \emph{Pre-continuation recorded proof scope.} The all-fibre obstruction and its propagation through base twists were re-derived with the exact support hypotheses}{the hypotheses, notation and definitions stated here; \cref{cdp:st:D.1.9}; \cref{cdp:st:A.4.9}; \cref{cdp:st:A.4.7}; \cref{cdp:st:A.3.4}; \cref{cdp:st:A.5.2}}

Let $f:Y\to\PP^1$ satisfy the full hypotheses of CDP20 Theorem~2.2: $Y$ is a compact connected complex threefold, $H^1(Y,\Z)=H^2(Y,\Z)=0$, and $f$ is its holomorphic algebraic reduction with connected fibres. Let $\beta\in H^1(Y,\OO_Y)$ lift the integral transgression generator of \ref{cdp:st:D.1.9}, and put $M_z=\exp(2\pi i z\beta)$, with the same exponential convention as in that statement. For each critical value $c$, define
\[
 B_c=\{z\in\C:
 H^0(F_c,\Omega_Y^1|_{F_c}\otimes M_z|_{F_c})\ne0\},
 \qquad F_c=f^{-1}(c).
 \tag{C.20}\label{cdp:eq:C.20}
\]
When each $B_c$ is a proper subset of $\C$, a dense open subset of $\Pic(Y)$ satisfies the vanishing and the two inequalities in (C.19). The premise is a finite collection of tests on the actual full fibres, given by A.1--A.3; it is satisfied in particular when every reduced support is smooth.

When $c_3(Y)>0$, there is a fixed critical value $c$ with $B_c=\C$. Its reduced support is nonnormal. Thus a nonnormal reduced fibre is necessary for positive Euler characteristic under the full CDP hypotheses, with arbitrary multiplicities and arbitrary singularities allowed at the other fibres.
\end{proposition}


\begin{proof}
\noindent\textbf{Step 1.} The transgression lift exists because the coherent Leray sequence over $\PP^1$ has no degree-one or degree-two base term for $\OO_{\PP^1}$. The integral degree-one Leray sequence identifies the integral invariant group with $\Z$ under the two stipulated integral vanishings. \ref{cdp:st:D.1.9} proves that its restriction to every full fibre, and to every nearby smooth fibre, is nonzero. On a neighbourhood $U$ of any fibre, $M_r$ for real $r$ has the unitary local system
\[
 a\longmapsto\exp(2\pi i r\langle\gamma_U,a\rangle)
\]
as its flat realization. This equality follows from the common integral lift $\gamma_U$ in the exponential sequence; no global flat realization on $Y$ is used.

\noindent\textbf{Step 2.} The holomorphic family of exponentiated Cech cocycles representing $z\beta$ restricts to a flat family of coherent bundles on $F_c\times\C$. Proper semicontinuity makes $B_c$ a closed analytic subset of the complex line. A proper such subset is discrete and countable. There are finitely many critical values. We can therefore choose a real irrational $r$ outside the union of all $B_c$ when all are proper.

\noindent\textbf{Step 3.} On a smooth fibre $F$, the character just displayed is the irrational exponential of a nonzero integral class. The explicit calculations for tori, bielliptic surfaces, Kodaira surfaces, Hopf surfaces and Inoue surfaces in \ref{cdp:st:A.4.9} give
\[
 H^0(F,\mathbb A_r)=H^1(F,\mathbb A_r)=0.
\]
The unitary maximum principle gives $H^0(F,M_r)=0$, and the injection \ref{cdp:st:A.4.7} gives $H^0(F,\Omega_F^1\otimes M_r)=0$. The conormal sequence of a smooth full fibre, whose conormal line is trivial, then gives
$H^0(F,\Omega_Y^1|_F\otimes M_r)=0$.
At every critical full fibre this group vanishes by our choice of $r$. Set $L=K_Y\otimes M_r^{-1}$. Cartier-fibre duality and top-degree base change give $R^2f_*(T_Y\otimes L)=0$. A sufficiently positive base twist kills $H^1(R^1f_*(T_Y\otimes L))$, so
\[
 H^2(Y,T_Y\otimes L\otimes f^*\OO(n))=0
\]
for an actual integer $n$. The exponential sequence identifies $\Pic(Y)$ with a connected complex vector space. Semicontinuity therefore makes degree-two vanishing dense open in this whole group. Intersecting with the endpoint-vanishing loci of C.1 gives (C.19).

\noindent\textbf{Step 4.} For a full fibre with smooth reduction, \ref{cdp:st:A.3.4} proves the ambient vanishing for every real irrational $r$, including all the nilpotent layers. Its set $B_c$ is consequently proper. This proves the asserted geometric sufficient case.

\noindent\textbf{Step 5.} Finally $c_3(Y)=2-b_3(Y)$, and the index formula is integral. Thus $c_3(Y)>0$ implies $b_3(Y)=0$. C.1 then gives $H^2(Y,T_Y\otimes L)\ne0$ for every line bundle $L$. The preceding proof shows that the finitely many $B_c$ cannot all be proper; hence one fixed $B_c$ is the entire complex line. Its reduced support cannot be smooth by \ref{cdp:st:A.3.4}. A connected normal reduced support would be irreducible. Since it is a Cartier divisor in $Y$, its dualizing line has rational first Chern class zero and its structure-sheaf Euler characteristic is zero by divisor Riemann--Roch, exactly as calculated in \ref{cdp:st:A.5.2}. That statement proves it smooth. We have therefore obtained a nonnormal support, as claimed.
\end{proof}



The fixed fibre in \ref{cdp:st:C4d} carries an ambient section along the entire complex transgression line. For torus general fibre, \ref{cdp:st:C4c} strengthens this to all twists in the Picard group. For a general fibre with $h^1(\OO)=1$ and $h^2(\OO)=0$, \ref{cdp:st:A.5.13} identifies that line with the entire Picard group. The following proof treats both cases through their actual conductor schemes.




\begin{proposition}[The unrestricted normalization target under the integral CDP hypotheses]\label{cdp:st:C4e}
\StatementRecord{cdp:st:C4e}{proved}{Historical events: Claude Fable 5.2 (Anthropic), effort high, 16 November 2026: targeted proof check. GPT-6 Astra (OpenAI), ChatGPT Pro mode, effort Pro, S6\_09\_PAPERS, 17 November 2026: targeted proof check. Also checked in this writing. Pre-continuation writing events: GPT-6 Astra (OpenAI), ChatGPT Pro mode, effort Pro, S6\_10\_PAPERS, 18 November 2026: targeted proof check. Continuation metadata events: GPT-6 Astra (OpenAI), ChatGPT Pro mode, effort Pro, S6\_10\_PAPERS continuation, 19 November 2026: used and rechecked in this run; status-vocabulary reconciliation; metadata reconciliation only; no new mathematical check. \emph{Pre-continuation recorded proof scope.} The persistent inverse-Jacobian reduction, conductor genus cases, connected elliptic Picard image, last nilpotent-layer injection, and closed analytic extension were re-derived}{the hypotheses, notation and definitions stated here; \cref{cdp:st:C.3.1}; \cref{cdp:st:C.3.7}; \cref{cdp:st:C.3.9}; \cref{cdp:st:C.3.10}; \cref{cdp:st:C4c}; \cref{cdp:st:C4d}; \cref{cdp:st:A.5.13}; \cref{cdp:st:A.3.5}; \cref{cdp:st:A.3.4}; \cref{cdp:st:A.5.2}; \cref{cdp:st:A.3.1}; \cref{cdp:st:A.4.9}; \cref{cdp:st:A.1.1}; \cref{cdp:st:A.5.4}; \cref{cdp:st:A.5.3}; \cref{cdp:eq:C.22}; \cref{cdp:st:A.2.3}; \cref{cdp:st:A.5.12}; \cref{cdp:st:C2d}}

Let $f:Y\to\PP^1$ satisfy the integral CDP hypotheses and $c_3(Y)>0$. Then a fixed irreducible nonnormal reduced support $S$ has normalization $\eta:\widetilde S\to S$ satisfying
\[
 H^2\bigl(\widetilde S,\omega_{\widetilde S}\otimes
           \eta^*(K_Y^{-1}\otimes L)\bigr)\ne0
 \qquad\hbox{for every }L\in\Pic(Y).                  \tag{C.21}\label{cdp:eq:C.21}
\]
The assertion concerns the reduced support; no multiplicity-one conclusion is included.
\end{proposition}


\begin{proof}
\noindent\textbf{Step 1.} \ref{cdp:st:C.3.1} and the exclusions \ref{cdp:st:C.3.7}, \ref{cdp:st:C.3.9} and \ref{cdp:st:C.3.10} show
that the general fibre is a two-torus, a Hopf surface, or
an Inoue surface. In the last two cases it has
$h^1(\OO)=1$ and $h^2(\OO)=0$.
For torus general fibre, \ref{cdp:st:C4c} gives a fixed full fibre with a nonzero ambient cotangent section for every ambient twist. In the second case, \ref{cdp:st:C4d} gives persistence along $\C\beta$, and \ref{cdp:st:A.5.13} identifies this with $\Pic(Y)$. From this point the proof uses precisely that persistence property and does not use $c_3(Y)>0$ again. \ref{cdp:st:A.3.5} makes its support irreducible, so it is $mS$ for a single positive integer $m$. \ref{cdp:st:A.3.4} excludes a smooth support, and \ref{cdp:st:A.5.2} then excludes a normal support. Its filtration in \ref{cdp:st:A.3.1} has finitely many graded factors
\[
 \Omega_Y^1|_S\otimes M|_S\otimes\OO_Y(-jS)|_S,
 \qquad 0\le j<m.
\]
Nonvanishing on the filtered sheaf requires nonvanishing on at least one factor. Each factor's nonvanishing locus is closed analytic in the connected vector group $\Pic(Y)$. A finite union of proper such loci cannot cover that group. Hence one factor has nonzero sections for every $M$. Translating by $\OO_Y(-jS)$ shows
$H^0(S,\Omega_Y^1|_S\otimes M|_S)\ne0$ for every $M$.
\ref{cdp:st:A.4.9} kills the torsion-free quotient on a dense open set of twists. The exact sequence \ref{cdp:st:A.1.1} and proper semicontinuity then show that the inverse-Jacobian kernel itself has sections for every twist.

\noindent\textbf{Step 2.} We apply the precise persistent-kernel reduction \ref{cdp:st:A.5.4}. It gives a fixed component of the normalization on whose resolution every twist of one fixed rank-one sheaf has a section. The Gauduchon degree argument there proves that the pullback of every ambient line has degree zero for every Gauduchon metric. Its conductor comparison excludes nonnegative Kodaira dimension, and its class VII argument makes every ambient line trivial on that normalization component. Rational components have the same triviality because their Picard group has no connected part. Either case gives \eqref{cdp:eq:C.21} immediately by Serre duality and descent from resolution. It remains to treat an irrational ruled resolution $p:Z\to B$.

\noindent\textbf{Step 3.} Let $h:Z\to S\subset Y$. \ref{cdp:st:A.5.3} supplies an effective conductor divisor $D_Z$ with
\[
 K_Z+D_Z=h^*(K_Y+S),\qquad D_Z\cdot F=2              \tag{C.22}\label{cdp:eq:C.22}
\]
for a general ruling fibre $F$. Its horizontal part is therefore either reduced of total ruling degree two, or $2C$ for a section $C$. We record the comparison through each ruling blowdown $\sigma:Z\to Z'$ used below. Put $D'=\sigma_*D_Z$, let $a$ be the coefficient of its exceptional curve $E$ in $D_Z$, and put $d=\operatorname{mult}_{\sigma(E)}D'$. Then
\[
 K_Z+D_Z=\sigma^*(K_{Z'}+D')+(1+a-d)E.
 \tag{C.22b}\label{cdp:eq:C.22b}
\]
Intersection with $E$, using $c_1(K_Z+D_Z)=0$ and $E^2=-1$, gives $a=d-1$. Thus the log-canonical divisor is the exact pullback of $K_{Z'}+D'$, and the latter again has zero rational Chern class. No factorization of $h$ through the blowdown is asserted or needed.

\noindent\textbf{Step 4.} We first treat the reduced horizontal part $H$. If a component is exceptional for the normalization resolution, every ambient line restricts trivially to it. Its normalization maps finitely onto $B$, and the norm composed with pullback on $\Pic^0(B)$ is multiplication by the covering degree. Hence this restriction has finite kernel; the connected ambient Picard image is trivial, giving \eqref{cdp:eq:C.21}. We may therefore assume that the horizontal components map to conductor curves. Adjunction on each of them gives
\[
 2p_a(H_i)-2+(D_Z-H_i)\cdot H_i=0.
\]
The second term is nonnegative. Since $H_i$ dominates a positive-genus base, its arithmetic genus is at least one. Thus the base has genus one, each $H_i$ is smooth elliptic, and the horizontal curves are disjoint from each other and from the vertical part. Contract ruling $(-1)$-curves to a minimal bundle. A contracted curve outside $D_Z$ has $D_Z\cdot E=1$, so cannot identify two horizontal points or create a horizontal node. A contracted curve in the vertical part is disjoint from $H$. Downstairs the horizontal curves remain smooth and disjoint, and any nonzero vertical divisor on a minimal ruled bundle meets them. Consequently the downstairs vertical part is zero. Reversing the contractions, a blowup away from the horizontal conductor would give coefficient $-1$ on its exceptional curve in \eqref{cdp:eq:C.22}, which persists under further strict transform and contradicts effectiveness. A blowup at a smooth point of a coefficient-one horizontal curve gives exceptional coefficient zero. Hence $D_Z=H$ has no vertical component. Every codimension-one singularity of $S$ is then an ordinary node, and \ref{cdp:st:A.2.3} identifies the inverse-Jacobian kernel with the normalization kernel globally, including its codimension-two points. A section of a topologically trivial line on the Kahler resolution has an effective zero divisor of zero Kahler volume, so is nowhere zero. Persistence for all twists therefore makes all ambient pullbacks trivial and proves \eqref{cdp:eq:C.21}.

\noindent\textbf{Step 5.} We now have $D_Z=2C+V$, where $C$ is a section and $V$ is vertical. If $C$ is exceptional for $Z\to\widetilde S$, every ambient line is trivial on $C$. Restriction $\Pic^0(Z)\to\Pic^0(C)$ is an isomorphism, because $C$ is a ruling section. This again makes all ambient pullbacks trivial. Otherwise let $D\subset S$ be its image and let $D^\nu$ normalize that curve. When $D^\nu$ is rational, every ambient line pulls back trivially to $D^\nu$, hence to $C$, giving the same conclusion.

\noindent\textbf{Step 6.} Suppose instead that $g(D^\nu)\ge1$. A vertical rational curve cannot map finitely onto $D^\nu$, so all the normalization branches over the general point of $D$ belong to $C$. Their conductor coefficient is two. We derive the possible transverse reduced plane curves rather than assuming a list of singularities. Write its equation as $f=\prod_i f_i$ in the smooth-plane local ring and let $R_i$ and $B_i$ be a branch ring and its normalization. The elements of $R=\OO/(f)$ supported only on branch $i$ form
\[
 \left(\prod_{j\ne i}f_j\right)\bigg|_{f_i}\,R_i.
\]
This equality follows from
$\bigcap_{j\ne i}(f_j)=(\prod_{j\ne i}f_j)$ in that unique
factorization domain. Multiplying by arbitrary elements of
$B_i$ therefore shows that the conductor exponent on branch $i$
is its intrinsic exponent plus $\sum_{j\ne i}I(f_i,f_j)$.
A singular branch has intrinsic exponent at least two:
exponent one would imply $\C+t\C\{t\}=\C\{t\}\subset R_i$.
Exponent two forces $R_i=\C+t^2\C\{t\}=\C\{t^2,t^3\}$,
since an element of order one would again make the ring smooth.
Any additional branch has intersection multiplicity at least
two with a singular branch and is therefore excluded. Thus a
singular branch gives exactly the cusp. When all branches are
smooth, their pairwise positive intersection multiplicities sum
to two on each branch. This gives two branches of contact two
(the tacnode), or three pairwise transverse branches (the
ordinary triple point), and exhausts the possibilities.
The only remaining alternative is a nonconstant pullback
$\Pic(Y)\to\Pic(Z)$; a constant pullback is the zero
homomorphism and already gives \eqref{cdp:eq:C.21}. In that alternative,
the coherent conductor argument \ref{cdp:st:A.5.12} gives $p_a(D)\le1$
for torus general fibre, and \ref{cdp:st:A.5.13} gives the same bound
in the second case.
Since $g(D^\nu)\ge1$, it follows that $D$ itself is smooth
elliptic. This conclusion permits nonreduced conductor image
schemes and imposes no assumption on the original singular
curve. We can therefore apply \ref{cdp:st:C2d} directly. It makes the
transverse type constant, gives the etale conductor cover
$C\to D$, and proves $C^2=0$. In particular $C$ and the
ruling base are elliptic.

\noindent\textbf{Step 7.} Adjunction in \eqref{cdp:eq:C.22} gives $V\cdot C=0$. No ruling $(-1)$-curve meets $C$: for one outside $V$, its intersection with $K_Z+2C+V$ would be at least $-1+2=1$; for one in $V$, the equality $V\cdot C=0$ forbids such an intersection. Thus every contraction to the minimal ruled bundle avoids $C$ and preserves its self-intersection zero. On that bundle $K+2C$ is numerically trivial, so the effective vertical part is numerically trivial and hence zero. A first reverse blowup, necessarily away from $C$, would again produce negative conductor coefficient. There are therefore no reverse blowups: $Z$ is minimal ruled and $D_Z=2C$. There is also no exceptional curve for the normalization resolution: a negative curve $E\ne C$ would satisfy $2p_a(E)-2=E^2-2C\cdot E<0$, hence would be rational and contained in a ruling fibre; a minimal ruled bundle has only self-intersection-zero fibres. Thus the normalization itself is smooth here.

\noindent\textbf{Step 8.} By the explicit inverse-Jacobian computations in \ref{cdp:st:C2d}, its pullback allows at most a simple pole along $C$. There are no other divisorial poles. Hartogs extension for the line bundle on the smooth surface gives an injection of the kernel into $h^*\OO_Y(-S)\otimes\OO_Z(C)$. The exact sequence of the section $C$ and $R^1p_*\OO_Z=0$ give
\[
 0\longrightarrow\OO_B\longrightarrow p_*\OO_Z(C)
   \longrightarrow N_{C/Z}\longrightarrow0,
 \qquad \deg N_{C/Z}=0.                              \tag{C.22a}\label{cdp:eq:C.22a}
\]
For a degree-zero line $P$ on the elliptic curve, both end terms after tensoring by $P$ have zero sections unless $P=\OO_B$ or $P=N_{C/Z}^{-1}$. Thus $H^0(Z,\OO_Z(C)\otimes p^*P)$ is supported at at most two points of $\Pic^0(B)$. The ambient pullback image is a connected complex subgroup of this Picard group. Persistence for every ambient twist forces that subgroup to be a point, and it contains the identity. Therefore every ambient line pulls back trivially to $Z$ in this last case as well.

\noindent\textbf{Step 9.} The equality $\pi_*\OO_Z=\OO_{\widetilde S}$ for the normal resolution descends triviality of a pulled-back line: projection formula identifies its space of sections with that of its trivial pullback, and the same statement for its inverse gives mutually inverse trivializations. Consequently one normalization component has trivial $\eta^*(K_Y^{-1}\otimes L)$ for every $L$. Serre duality on that compact normal Cohen--Macaulay surface identifies its top dualizing cohomology with $H^0(\OO)^\vee=\C$. It is a direct summand of the target in \eqref{cdp:eq:C.21}, proving the claim.
\end{proof}



The local example \ref{cdp:st:A.2.5} has actual vertical inverse-Jacobian
poles, and \ref{cdp:st:A.2.6} has a critical curve that is not a local
complete intersection. Their local rings remain part of the
audit. The global conductor-duality arguments \ref{cdp:st:A.5.12}--\ref{cdp:st:A.5.13} prove
that a positive-genus conductor image in the persistent ruled
case is smooth, and the exact contractions above then remove
the vertical conductor. Thus neither local phenomenon leaves
an unproved step in the stated theorem. The assertion
is about the normalization of a reduced support and does not
claim that its full fibre has multiplicity one.




\begin{proposition}[An unrestricted geometric replacement for CDP20 Theorem 2.2]\label{cdp:st:C4f}
\StatementRecord{cdp:st:C4f}{proved}{Historical events: Claude Fable 5.2 (Anthropic), effort high, 16 November 2026: targeted proof check. GPT-6 Astra (OpenAI), ChatGPT Pro mode, effort Pro, S6\_09\_PAPERS, 17 November 2026: targeted proof check. Also checked in this writing. Pre-continuation writing events: GPT-6 Astra (OpenAI), ChatGPT Pro mode, effort Pro, S6\_10\_PAPERS, 18 November 2026: targeted proof check; used and rechecked in this run. Continuation metadata events: GPT-6 Astra (OpenAI), ChatGPT Pro mode, effort Pro, S6\_10\_PAPERS continuation, 19 November 2026: used and rechecked in this run; status-vocabulary reconciliation; metadata reconciliation only; no new mathematical check. \emph{Pre-continuation recorded proof scope.} The exact geometric criterion was re-derived from the printed support cases and Picard maps. This does not assert universal surjectivity onto normalization cohomology}{the hypotheses, notation and definitions stated here; \cref{cdp:st:C4e}; \cref{cdp:st:C4c}; \cref{cdp:st:A.3.4}; \cref{cdp:st:A.3.5}; \cref{cdp:st:A.5.2}; \cref{cdp:st:A.2.2}; \cref{cdp:eq:C.VC.4}; \cref{cdp:st:A.5.13}; \cref{cdp:st:C4d}; \cref{cdp:st:D.1.11}; \cref{cdp:st:A.4.9}}

Let $f:Y\to\PP^1$ satisfy the integral CDP hypotheses.
For each critical full fibre with
irreducible nonnormal reduced support $S$, let
$\eta_S:T_S\to S$ be its normalization. Then a general
line bundle $L\in\Pic(Y)$ satisfies
\[
 H^q(Y,T_Y\otimes L)=0\quad(q\ne1)
 \tag{C.VC.1}\label{cdp:eq:C.VC.1}
\]
exactly when, for every such $S$, the holomorphic pullback
homomorphism
\[
 \Pic(Y)\longrightarrow\Pic(T_S)
 \tag{C.VC.2}\label{cdp:eq:C.VC.2}
\]
is nonzero. Equivalently, its differential
$H^1(Y,\OO_Y)\to H^1(T_S,\OO_{T_S})$ is nonzero.
In this case $\chi(Y,T_Y\otimes L)\le0$ and $c_3(Y)\le0$.
If one map \eqref{cdp:eq:C.VC.2} is zero, then
$H^2(Y,T_Y\otimes L)\ne0$ for every $L\in\Pic(Y)$.
All fibre multiplicities and all reduced-support
singularities are permitted in this criterion.
\end{proposition}


\begin{proof}
\noindent\textbf{Step 1.} The classification and exclusions cited in \ref{cdp:st:C4e} leave
exactly the two cohomological cases used in that proof:
torus general fibre, or $h^1(\OO)=1$, $h^2(\OO)=0$.
The Picard group of $Y$ is the connected exponential image
of $H^1(Y,\OO_Y)$. A holomorphic group homomorphism out of
this group is zero precisely when its differential is zero.
For each fixed critical full fibre, the ambient-section
nonvanishing locus is closed analytic in this group, as
in \ref{cdp:st:C4c}. \ref{cdp:st:A.3.4} makes it proper when the support is smooth,
and \ref{cdp:st:A.3.5} makes it proper when the support is reducible.
A connected normal support is irreducible and is smooth
by \ref{cdp:st:A.5.2}, so these cases leave just the supports in the
statement.

\noindent\textbf{Step 2.} Suppose the nonvanishing locus for $F=mS$ is the whole
Picard group. The proof of \ref{cdp:st:C4e}, beginning with its
fixed-fibre persistence assertion, proves that \eqref{cdp:eq:C.VC.2}
is zero. Its steps use the actual component filtration,
the inverse-Jacobian kernel, the Gauduchon bound, the
effective conductor comparison, the coherent conductor
genus bound, and the explicit constant-transverse-type
calculations. Each is independent of the Euler number
once persistence has been given.

\noindent\textbf{Step 3.} Conversely, suppose \eqref{cdp:eq:C.VC.2} is zero. The last nonzero
ideal layer of the full fibre is an actual subsheaf
\[
 (\Omega_Y^1\otimes M\otimes\OO_Y(-(m-1)S))|_S
 \hookrightarrow
          (\Omega_Y^1\otimes M)|_{mS}.
 \tag{C.VC.3}\label{cdp:eq:C.VC.3}
\]
The canonical normalization submodule in \ref{cdp:st:A.2.2} gives
an injection into the left side from
\[
 \eta_{S*}\bigl(\eta_S^*((M\otimes\OO_Y(-mS))|_S)\bigr).
 \tag{C.VC.4}\label{cdp:eq:C.VC.4}
\]
Here the first conormal factor $\OO_Y(-S)|_S$ combines
with the displayed $(m-1)$st ideal layer; this comparison
retains every nilpotent layer needed for the injection.
Since $mS$ is a full fibre, $\OO_Y(-mS)|_S$ is a trivial
line, with its one-dimensional scalar factor determined
by a local base parameter. The pullback of $M$ to $T_S$
is trivial by \eqref{cdp:eq:C.VC.2}. Hence \eqref{cdp:eq:C.VC.4} has a nonzero global
section for every $M$, and its injections prove that the
full-fibre nonvanishing locus is the entire Picard group.
This also proves the converse directly for arbitrary $m$.

\noindent\textbf{Step 4.} We have proved an exact geometric characterization of
every improper critical-fibre locus. For torus general
fibre, applying \ref{cdp:st:C4c} gives the equivalence with generic
degree-two vanishing. In the second case \ref{cdp:st:A.5.13} identifies
$\Pic(Y)$ with the transgression line, so properness of
all these loci is exactly the premise of \ref{cdp:st:C4d} and gives
generic degree-two vanishing there as well.
The endpoint vanishings of C.1 then give \eqref{cdp:eq:C.VC.1} and the
index inequalities.
For the last assertion in the torus case, generic non-torsion on every
smooth torus fibre makes $R^2f_*(T_Y\otimes L)$ a
torsion sheaf for a dense set of twists, by \ref{cdp:st:D.1.11}.
The fixed full-fibre obstruction makes it nonzero.
Its nonzero global sections are a quotient of
$H^2(Y,T_Y\otimes L)$ by the Leray sequence.
Thus degree-two cohomology is nonzero on a dense set.
Its proper semicontinuity locus is closed analytic in
the connected Picard family, so it is nonzero for every
twist. In the second case use instead the actual twists
$L_r=K_Y\otimes M_r^{-1}$ for real irrational $r$.
\ref{cdp:st:A.4.9} and the smooth-fibre conormal sequence make their
$R^2f_*(T_Y\otimes L_r)$ torsion. The fixed full-fibre
obstruction again makes that torsion sheaf nonzero,
so Leray gives $H^2(Y,T_Y\otimes L_r)\ne0$.
These twists have accumulation points in the complex
line $\Pic(Y)=\C\beta$. A closed complex analytic subset
of that line containing them is the whole line.
This proves all-twist nonvanishing in the second case
and completes both implications of the criterion.
\end{proof}






\begin{proposition}[The necessary normalization defect in the three named test families]\label{cdp:st:C4g}
\StatementRecord{cdp:st:C4g}{proved}{Historical events: Claude Fable 5.2 (Anthropic), effort high, 16 November 2026: carried as background without fault found. GPT-6 Astra (OpenAI), ChatGPT Pro mode, effort Pro, S6\_09\_PAPERS, 17 November 2026: targeted proof check. Also checked in this writing. Pre-continuation writing events: GPT-6 Astra (OpenAI), ChatGPT Pro mode, effort Pro, S6\_10\_PAPERS, 18 November 2026: targeted proof check. Continuation metadata events: GPT-6 Astra (OpenAI), ChatGPT Pro mode, effort Pro, S6\_10\_PAPERS continuation, 19 November 2026: used and rechecked in this run; status-vocabulary reconciliation; metadata reconciliation only; no new mathematical check. \emph{Pre-continuation recorded proof scope.} Full three-family comparison rederived: finite ambient H2 pulls to zero in free H2(dP6), H1(O\_dP6)=0 makes every pulled-back line trivial, and Serre duality gives the one-dimensional normal target. Checked the C1.3 global nonvanishing route without claiming normalization-map surjectivity; toric strata give2, elliptic-ruled conductor subtraction gives0, and the actual unibranch cuspidal normalization is a proper bijection. All full multiple fibres have the stated actual root covers and pass A3.4 including nilpotent layers. E13/E14/E19 remain named reviewed dependencies}{the hypotheses, notation and definitions stated here; \cref{cdp:st:E13}; \cref{cdp:st:E14}; \cref{cdp:st:A.3.4}; \cref{cdp:st:A.4.1}; \cref{cdp:st:E19}; \cref{cdp:st:C.1.3}; \cref{cdp:st:C5b}}

For every admissible member of the original $(3,4)$ family,
every member of the two-parameter family in \ref{cdp:st:E13}, and every
member of the Liu--Rollenske family in \ref{cdp:st:E14}, there is a reduced
full hexagonal fibre $W$ with normalization
$\eta:\mathrm dP_6\to W$ such that
\[
 \eta^*L=\OO_{\mathrm dP_6}\quad(L\in\Pic(Y)),\qquad
 H^2\bigl(\mathrm dP_6,K_{\mathrm dP_6}\otimes
                    \eta^*(K_Y^{-1}\otimes L)\bigr)=\C.
 \tag{C.TEST.1}\label{cdp:eq:C.TEST.1}
\]
In these three families the singular-fibre Euler contributions are
\[
\begin{array}{c|c|c}
\text{family}&\text{reduced singular supports}&
                 \text{their Euler numbers}\\ \hline
(3,4)&W&2\\
\text{two-parameter \ref{cdp:st:E13}}&W,\ W_1,\ W_2&2,0,0\\
\text{Liu--Rollenske \ref{cdp:st:E14}}&W,\ C_{\rm cusp}&2,0 .
\end{array}
\tag{C.TEST.2}\label{cdp:eq:C.TEST.2}
\]
The other critical full fibres are smooth multiple fibres,
with multiplicities $(3,4)$, $4$, and $4$, respectively.
Every one of those full multiple fibres passes the generic
ambient vanishing test \ref{cdp:st:A.3.4}, including all its nilpotent
layers. Consequently positive $c_3$ in these examples does
not require a failure of that test on a nonreduced fibre.
The reduced full fibre $W$ already supplies the fixed
normalization obstruction, and $H^2(Y,T_Y\otimes L)\ne0$
for every ambient line bundle in each family.
\end{proposition}


\begin{proof}
\noindent\textbf{Step 1.} The original local filling is the hexagonal toric quotient
in \ref{cdp:st:A.4.1}. The two-parameter construction in \ref{cdp:st:E13} places this
same filling at the simple nodal value where its specified
linearization vanishes; its other two simple nodal values
have rank-one fillings. The coordinate and line-bundle
comparisons of \ref{cdp:st:E14} place the hexagonal filling at its
separate simple node and the cuspidal quotient at the
remaining critical value. Those proofs construct the actual
local spaces and their identifications over the punctured
base discs. Thus the normalization of $W$ in each case is
the stated smooth rational surface.

\noindent\textbf{Step 2.} For every member in question, $H^2(Y,\Z)$ is a finite
group. For the original affine family this follows from E19:
the universal coefficient sequence gives
\[
 H^2(Y,\Z)\simeq\operatorname{Ext}^1(H_1(Y,\Z),\Z)
 \simeq\Z/d\Z,\qquad d=|p|,
\]
because $H_2(Y,\Z)$ is torsion. In the two other families
the integral homology is that of the six-sphere. In
particular every ambient first Chern class is torsion.
The group $H^2(\mathrm dP_6,\Z)$ is free, so its pullback
to $\mathrm dP_6$ is zero. The exponential sequence there,
together with $H^1(\mathrm dP_6,\OO)=0$, shows that the
pulled-back line bundle is trivial. This applies to both
$L$ and $K_Y^{-1}\otimes L$, and Serre duality on
$\mathrm dP_6$ proves \eqref{cdp:eq:C.TEST.1}.

\noindent\textbf{Step 3.} All these threefolds have $b_1=b_2=b_3=0$: use \ref{cdp:st:E19} for
every original affine member, and the integral homology
of the six-sphere for the other two families. Any full
fibre is a nonzero effective Cartier divisor. Statement
\ref{cdp:st:C.1.3} therefore gives
\[
 h^2(Y,T_Y\otimes L)\geq 1
 \qquad\text{for every }L\in\Pic(Y).
\]
This argument applies equally when $d>1$ and the integral
cohomology has nonzero torsion. It requires no surjectivity
of the global normalization cohomology map. Indeed \ref{cdp:st:C5b}
computes the actual exceptions to that surjectivity in
the original family.

\noindent\textbf{Step 4.} The toric strata of $W$ consist of one $(\C^*)^2$ stratum,
three $\C^*$ strata and two triple points. Their compactly
supported Euler numbers give $e(W)=2$.
Each rank-one normalization is a $\PP^1$-bundle over an
elliptic curve, with its zero and infinity sections glued.
Both the normalization and these conductor curves have
Euler number zero, so normalization additivity gives
$e(W_i)=0$. At the cusp in \ref{cdp:st:E14}, normalization of the
cuspidal rational curve makes its degree-zero line bundle
trivial. The lifted contracting action on this normalized
line complement is a product of a fixed automorphism of
$\PP^1$ and multiplication by a constant of modulus
strictly less than one on $\C^*$. Its quotient is a
$\PP^1$-bundle over the elliptic curve given by that
multiplicative period. The cusp normalization is
bijective; a proper continuous bijection between these
compact Hausdorff spaces is a homeomorphism. Hence
$e(C_{\rm cusp})=0$. This proves \eqref{cdp:eq:C.TEST.2}.

\noindent\textbf{Step 5.} Finally the explicit good-reduction covers in \ref{cdp:st:E13}--\ref{cdp:st:E14}
and the original root charts have the stated orders.
Their reduced central fibres are smooth finite unramified
torus quotients, and their Euler numbers are zero.
\ref{cdp:st:A.3.4} applies to their actual full fibres, proving the
asserted generic ambient vanishing. Euler localization
therefore gives $c_3(Y)=2$ in all three rows, while the
nonreduced fibres satisfy the complete multiple-fibre
test. This is the required test of the proposed additional
nonreduced-fibre necessity.
\end{proof}




\begin{proposition}[The normalization quotient and its cohomology map]\label{cdp:st:C5a}
\StatementRecord{cdp:st:C5a}{proved}{Historical events: Claude Fable 5.2 (Anthropic), effort high, 16 November 2026: carried as background without fault found. GPT-6 Astra (OpenAI), ChatGPT Pro mode, effort Pro, S6\_09\_PAPERS, 17 November 2026: targeted proof check. Also checked in this writing. Pre-continuation writing events: GPT-6 Astra (OpenAI), ChatGPT Pro mode, effort Pro, S6\_10\_PAPERS, 18 November 2026: targeted proof check. Continuation metadata events: GPT-6 Astra (OpenAI), ChatGPT Pro mode, effort Pro, S6\_10\_PAPERS continuation, 19 November 2026: used and rechecked in this run; status-vocabulary reconciliation; metadata reconciliation only; no new mathematical check. \emph{Pre-continuation recorded proof scope.} Full branch-normal quotient, coordinate-change descent, immersion determinant, twisted exact sequence, finite pushforward and compact-surface Serre pairing read and rederived. The target is H0(nu*(K\_Y tensor L\textasciicircum{}-1)) dual and its image is the stated boundary kernel; no surjectivity is inferred without the exact boundary condition. DGKXv1 Lemma6.11 p81 was opened as the specific-family quotient counterpart; the short general twisted application retains that granular credit}{the local branch-normal coefficient map and adjunction; \citet[Lemma~6.11, p.~81]{DGKXv1} for the original-family quotient; finite direct image and Serre duality on the smooth compact normalization}

Let $Y$ be a smooth compact complex threefold, $S\subset Y$ a reduced normal-crossing divisor, and $\nu:\widetilde S\to Y$ the normalization followed by inclusion. This is a finite holomorphic immersion on each local branch. Put $\mathcal T=T_Y(-\log S)$. There is a canonical exact sequence of coherent $\OO_Y$-modules
\[
0\longrightarrow\mathcal T\longrightarrow T_Y
 \longrightarrow\nu_*N_\nu\longrightarrow0,
\qquad N_\nu=K_{\widetilde S}\otimes\nu^*K_Y^{-1}.      \tag{C.3}\label{cdp:eq:C.3}
\]
For every holomorphic line bundle $L$ on $Y$, the image of the natural cohomology map is exactly
\[
\begin{split}
\im\bigl(H^2(Y,T_Y\otimes L)&\longrightarrow
 H^2(\widetilde S,N_\nu\otimes\nu^*L)\bigr)\\
 &=\ker\bigl(H^2(\widetilde S,N_\nu\otimes\nu^*L)
 \longrightarrow H^3(Y,\mathcal T\otimes L)\bigr).
\end{split}                                             \tag{C.4}\label{cdp:eq:C.4}
\]
The target identifies by Serre duality with
$H^0(\widetilde S,\nu^*(K_Y\otimes L^{-1}))^\vee$.
\end{proposition}


The branch-normal quotient specializes to \citet[Lemma~6.11, p.~81]{DGKXv1} in the original family. We retain the short argument for its general ambient line-bundle twist and the exact image of its cohomology map.

\begin{proof}
\noindent\textbf{Step 1.} In coordinates with $S=(z_1\cdots z_r=0)$, $1\le r\le3$, the map sends a vector field $\sum a_i\partial_{z_i}$ to the normal coefficient $a_i|_{z_i=0}$ on each branch $i\le r$. The separate coefficients lift independently to holomorphic functions, making the map surjective. Its kernel consists precisely of fields with $z_i\mid a_i$ for $i\le r$, which is $T_Y(-\log S)$. Coordinate changes act on each coefficient through the normal-line transition and permute the branches; hence the map and its exactness descend canonically to the normalization. The immersion tangent sequence on $\widetilde S$ gives $\det\nu^*T_Y=\det T_{\widetilde S}\otimes N_\nu$, proving the adjunction expression. Tensoring by the locally free line bundle $L$ preserves exactness. Finite pushforward has no higher direct images, so the long exact sequence in coherent cohomology gives \eqref{cdp:eq:C.4}. The smooth compact surface $\widetilde S$ has the usual Serre pairing, which gives the stated dual target. In particular the vanishing of $H^3(Y,\mathcal T\otimes L)$ makes this cohomology map surjective; no such vanishing has been suppressed in \eqref{cdp:eq:C.4}.
\end{proof}






In this subsection $X=X_c$ is the period construction of \cite{AC26}, compared with the quotient construction of \cite{Engel2026}, with admissible parameter $c$, standard clockwise affine vectors
$v_1=(1,2,-4,0)^t$, $v_2=(-1,-3,3,0)^t$, and untwisted cusp collar. Thus $p=12l_0-4l_1-3l_2=-1$. We write $\mathfrak B=\PP(3,4)$, $\rho:\mathfrak B\to B$ for its coarse map, and $\bar f:X\to\mathfrak B$ for the root factorization. Let $\mathcal T=T_X(-\log W)$ and let $\nu:\mathrm dP_6\to X$ normalize $W$. The root generator is fixed by
\[
 M_n=\bar f^*\OO_{\mathfrak B}(n)
     =\OO_X(n(S_1-S_2)),\quad n\in\Z.
\]
The exponential isomorphism $H^1(X,\OO_X)\to\Pic(X)$ identifies $\Pic(X)$ with $\C$. We scale its coordinate so that $M_z=\exp(z\gamma^{0,1})$ and $M_1=\OO_X(S_1-S_2)$. The sign of this equality is fixed by the clockwise transgression $12\gamma\mapsto-p=1$. This is checked by evaluating the two finite lifts on $v_1,v_2$: their integral periods are $4l_1,3l_2$, while the cusp contributes $-12l_0$, whose sum is $-p$. The exponential sequence therefore gives $\exp(12\gamma^{0,1})=f^*\OO_B(1)$. Its injectivity and the root relation give the stated $M_1$ equality.

\begin{proposition}[Integral twists and the normalization-map exception]\label{cdp:st:C5b}
\StatementRecord{cdp:st:C5b}{proved}{Historical events: Claude Fable 5.2 (Anthropic), effort high, 16 November 2026: carried as background without fault found; used and rechecked in this run. GPT-6 Astra (OpenAI), ChatGPT Pro mode, effort Pro, S6\_09\_PAPERS, 17 November 2026: targeted proof check. Also checked in this writing. Pre-continuation writing events: GPT-6 Astra (OpenAI), ChatGPT Pro mode, effort Pro, S6\_10\_PAPERS, 18 November 2026: targeted proof check. Continuation metadata events: GPT-6 Astra (OpenAI), ChatGPT Pro mode, effort Pro, S6\_10\_PAPERS continuation, 19 November 2026: used and rechecked in this run; status-vocabulary reconciliation; metadata reconciliation only; no new mathematical check. \emph{Pre-continuation recorded proof scope.} The root-bundle maps and evaluation semigroup were re-derived using the Paper I input: the normalization map has rank one exactly for n in \{-6,-3,-2\} union Z>=0, and is zero at n=-1 although h2=1. Paper I's root bundles require their own check records}{the hypotheses, notation and definitions stated here; \cref{cdp:eq:C.8}; \cref{cdp:st:C5a}; \cref{cdp:eq:C.9}}

For $n\in\Z$ define
\[
 r_1(n)\equiv n\pmod3,\quad 0\le r_1(n)<3,
 \qquad r_2(n)\equiv3n\pmod4,\quad0\le r_2(n)<4,
\]
\begin{gather*}
 d(n)=\frac{n-4r_1(n)-3r_2(n)}{12},\qquad
 \epsilon_n=\begin{cases}0&3\mid n,\\1&3\nmid n.\end{cases}
\\
 u(n)=\max\{d(n)+1,0\},\qquad
 v(n)=\max\{-d(n)-1,0\}.
\end{gather*}
The coarse direct images of $T_X\otimes M_n$ are
\begin{align*}
 f_*(T_X\otimes M_n)&=\OO(d(n))\oplus\OO(d(n-1)),\\
 R^1f_*(T_X\otimes M_n)&=\OO(d(n))\oplus\OO(d(n-1))\oplus\OO(d(n+2)),\\
 R^2f_*(T_X\otimes M_n)&=\OO(d(n+6))\oplus\C_{p_1}^{\epsilon_n}.\tag{C.5}\label{cdp:eq:C.5}
\end{align*}
Their global cohomology dimensions are
\begin{align*}
 h^0&=u(n)+u(n-1),\\
 h^1&=v(n)+v(n-1)+u(n)+u(n-1)+u(n+2),\\
 h^2&=v(n)+v(n-1)+v(n+2)+u(n+6)+\epsilon_n,\\
 h^3&=v(n+6).                                             \tag{C.6}\label{cdp:eq:C.6}
\end{align*}
The natural map $H^2(X,T_X\otimes M_n)\to H^2(\mathrm dP_6,K_{\mathrm dP_6})=\C$ has rank one exactly for
\[
 n\in\{-6,-3,-2\}\cup\Z_{\ge0};                         \tag{C.7}\label{cdp:eq:C.7}
\]
it is zero for all other integers. In particular $M_{-1}$ has tangent cohomology $(0,0,1,0)$ and a zero map to the normalization target.
\end{proposition}



\begin{proof}
\noindent\textbf{Step 1.} We check the root-curve and logarithmic tangent inputs of \cref{fam:src:A.1.1} and \cref{fam:src:A.5.2}. The root charts at $p_1,p_2$ are the finite etale torus-family covers used in the construction; finite-group invariants over $\C$ are exact. At the cusp the local equations are $xy=q$ and $xyz=q$, and the vertical logarithmic fields are generated by $x\partial_x-y\partial_y$ and the remaining relative direction. Their affine transition derivatives depend only on the base. The invariant vertical field gives a trivial subline of the vertical bundle $E$. The quotient has twelfth power $\OO(-p_0)$: its automorphy is that of the reciprocal elliptic discriminant, whose only divisor is a simple pole at the cusp. As $\Pic(\mathfrak B)=\Z$, the quotient is $\OO(-1)$; $H^1(\mathfrak B,\OO(1))=0$ splits the extension. Thus $E=\OO\oplus\OO(-1)$.

\noindent\textbf{Step 2.} The normalization complex of $W$ has global scalar complex
\[
 \C\xrightarrow{0}\C^3
 \xrightarrow{(a,b,c)\mapsto(a-b+c,a-b+c)}\C^2.
\]
The rational normalization and the three rational double curves have no positive-degree structure-sheaf cohomology. The displayed complex therefore gives $(1,2,1)$ for the fibre cohomology and its exterior cup product; the two oriented triangular faces give the nondegenerate degree-one pairing. The cyclic-cover fibres have the same dimensions. Proper flat base change on these charts gives locally free $G=R^1\bar f_*\OO_X$, with $R^2\bar f_*\OO_X=\det G$. Relative duality and the determinant of the logarithmic sequence give $\det G=\OO(-1)$. The integral class $\gamma$ maps to a nowhere-zero section of $G$: on a marked smooth torus the two identity columns of the period matrix prevent it from belonging to the span of the two holomorphic period rows, and at the cusp it is a nonzero cycle in the displayed complex. The same splitting argument yields
\[
E=G=\OO\oplus\OO(-1),\quad R^2\bar f_*\OO_X=\OO(-1),
\quad0\to\bar f^*E\to\mathcal T\to\bar f^*\OO(-5)\to0. \tag{C.8}\label{cdp:eq:C.8}
\]
These are the actual root-chart and cusp comparisons needed from the earlier returns; no coarse multiple fibre was replaced by its reduction in this argument.

\noindent\textbf{Step 3.} Choose root sections $x$ of weight $3$, vanishing at $p_2$, and $y$ of weight $4$, vanishing at $p_1$. The Kodaira--Spencer morphism $\OO(-5)\to E\otimes G$ has, by its weights, the form
\[
 \kappa=\begin{pmatrix}0&a y\\b y&c_0x\end{pmatrix}.
\]
Its first coefficient vanishes because $H^0(\OO(5))=0$. The logarithmic period derivative at the cusp is invertible, giving $ab\ne0$. At $p_1$, the elliptic coordinate $\tau-\rho$ has a simple zero on the root chart, so its period derivative is nonzero, giving $c_0\ne0$. Rescaling the four line frames makes the three nonzero coefficients one. Cup product with this extension class gives
\[
\delta_0=(0,y,y,x)^t,\qquad
\delta_1=\begin{pmatrix}-y&0\\-x&y\end{pmatrix}:
 \OO(-5)\oplus\OO(-6)\to\OO(-1)\oplus\OO(-2).
\]
The first is a subbundle inclusion since $x,y$ have no common zero, and its cokernel is $\OO\oplus\OO(-1)\oplus\OO(2)$. The second is injective with determinant $-y^2$; its cokernel is
$Q=\OO(-1)|_{2p_1}$, a length-two module on the order-three chart. The direct-image sequence of \eqref{cdp:eq:C.8} yields
\[
 R^0\bar f_*\mathcal T=E,\quad
 R^1\bar f_*\mathcal T=\OO\oplus\OO(-1)\oplus\OO(2),\quad
 R^2\bar f_*\mathcal T=\OO(-6)\oplus Q.
\]
The last splitting follows from $H^1(Q(6))=0$; it is not designated canonical.

\noindent\textbf{Step 4.} Statement \ref{cdp:st:C5a} gives the quotient $\nu_*K_{\mathrm dP_6}$. The local cusp volume form trivializes $\nu^*K_X$, and every $\nu^*M_z$ is trivial because $H^2(X,\Z)=0$ and $\Pic^0(\mathrm dP_6)=0$. Hence this quotient has cohomology only in degree two, where it is $\C$. At $p_0$, Serre duality and the conormal computation in Part A give $h^2(W,T_X|_W)=1$. Top-degree base change identifies this dimension with the residue dimension of $R^2\bar f_*T_X$. That module has generic rank one, so Nakayama's lemma and the structure theorem over $\C\{q\}$ make it free of rank one at $p_0$. Consequently its line summand is the elementary extension of $\OO(-6)$ by $\C_{p_0}$, namely $\OO(6)$, with no new cusp torsion. We have
\[
 R^0\bar f_*T_X=E,\quad
 R^1\bar f_*T_X=\OO\oplus\OO(-1)\oplus\OO(2),\quad
 R^2\bar f_*T_X=\OO(6)\oplus Q.                       \tag{C.9}\label{cdp:eq:C.9}
\]

\noindent\textbf{Step 5.} For a root line, removing its two stabilizer characters gives
$\rho_*\OO(n)=\OO_B(d(n))$. Averaging the finite stabilizers is exact, so there are no higher direct images of $\rho$. Its sections are equivalently monomials $x^ay^b$ with $3a+4b=n$, and Serre duality gives $h^1\OO(n)=h^0\OO(-n-7)$. At $p_1$, the invariant exponent of $Q(n)$ is $n-1$ modulo three. Its two possible coefficients have orders zero and one, so its coarse length is $\epsilon_n$ and it is killed by the coarse parameter. Apply projection formula to \eqref{cdp:eq:C.9} and then $\rho_*$. This proves \eqref{cdp:eq:C.5}, including the support and length of every torsion summand. Leray over $B$ has columns only zero and one, proving \eqref{cdp:eq:C.6}.

\noindent\textbf{Step 6.} The normalization map on $R^2$ is the quotient in
\[
0\longrightarrow\OO(n-6)\longrightarrow\OO(n+6)
 \longrightarrow\C_{p_0}\longrightarrow0.
\]
Its global-section map is nonzero precisely when a section of $\OO(n+6)$ is nonzero at $p_0$. Both $x,y$ are nonzero at that point, so this occurs precisely when $n+6$ lies in the additive semigroup generated by $3,4$. That semigroup is $\{0,3,4,6,7,8,\ldots\}$, giving \eqref{cdp:eq:C.7}. The $H^1(R^1)$ part maps to zero in the cohomology of a sheaf supported at one point, by the natural Leray filtration. At $n=-1$, \eqref{cdp:eq:C.5} has only one nonzero cohomological contribution, the reduced $\C_{p_1}$ summand in degree two, while $H^3(X,\mathcal T\otimes M_{-1})=H^1(\OO(-7))=\C$. The boundary from the normalization target is an isomorphism. This proves the stated counterexample to a surjection for every twist.
\end{proof}






\begin{proposition}[Every nonintegral twist and its resonance lengths]\label{cdp:st:C5c}
\StatementRecord{cdp:st:C5c}{proved}{Historical events: Claude Fable 5.2 (Anthropic), effort high, 16 November 2026: carried as background without fault found. GPT-6 Astra (OpenAI), ChatGPT Pro mode, effort Pro, S6\_09\_PAPERS, 17 November 2026: targeted proof check. Also checked in this writing. Pre-continuation writing events: GPT-6 Astra (OpenAI), ChatGPT Pro mode, effort Pro, S6\_10\_PAPERS, 18 November 2026: targeted proof check. Continuation metadata events: GPT-6 Astra (OpenAI), ChatGPT Pro mode, effort Pro, S6\_10\_PAPERS continuation, 19 November 2026: used and rechecked in this run; status-vocabulary reconciliation; metadata reconciliation only; no new mathematical check. \emph{Pre-continuation recorded proof scope.} Full nonintegral-twist proof read and rederived, including Picard cocycle equations, smooth Fourier contraction, exact local Koszul complexes, regular C\textasciicircum{}2 versus O/(s\textasciicircum{}2) modules, order-three joint coefficient123ik/(64pi), both finite-character invariant counts, cusp character acyclicity and Leray. Verified epsilon\_n=0 for3|n and1 otherwise by required Taylor exponents n+2 modulo3. Added the missing explicit passage from isolated zeros for each integer tuple to existence of a nonresonant fibre, then finite coherent support. The unchanged earlier root Koszul replay is reused at its recorded scope; this review adds no second execution credit}{the hypotheses, notation and definitions stated here; \cref{cdp:eq:C.13}}

Retain the fixed marked functions $\tau,\mu,\beta_c$ and put $\rho=\tau(z_1)$. For $z\notin\Z$, let $\mathcal R_z$ be the set of regular base points at which integers $(n,p,q,r)$ satisfy
\[
 q\tau+r\mu=p,\qquad z=n-6q\mu-r\beta_c.               \tag{C.10}\label{cdp:eq:C.10}
\]
Let
\[
 \Lambda_1=\Z+(3\beta_c(z_1)+4-2\rho)\Z,\quad
 \Lambda_2=\Z+(2\beta_c(z_2)+3-3i)\Z.
\]
Set $\ell_1(z)=1+\epsilon_n$ for $z=n+k(3\beta_c(z_1)+4-2\rho)$ with $k\ne0$, and zero off $\Lambda_1$. Set $\ell_2(z)=1$ for $z=n+k(2\beta_c(z_2)+3-3i)$ with $k\ne0$ and $n\equiv1,2\pmod4$, and zero otherwise. Then $\mathcal R_z$ is finite. With $N_z=|\mathcal R_z|$ and $J_z=2N_z+\ell_1(z)+\ell_2(z)$, we have
\[
 (h^0,h^1,h^2,h^3)(X,T_X\otimes M_z)=(0,J_z,1+J_z,0),  \tag{C.11}\label{cdp:eq:C.11}
\]
and, writing $\mathfrak m_t$ for the maximal ideal of $\OO_{B,t}$,
\begin{align*}
R^1f_*(T_X\otimes M_z)&=
 \bigoplus_{t\in\mathcal R_z}\C_t^2
 \oplus\C_{p_1}^{\ell_1(z)}\oplus\C_{p_2}^{\ell_2(z)},\\
R^2f_*(T_X\otimes M_z)&=
 \C_{p_0}\oplus\bigoplus_{t\in\mathcal R_z}\OO_{B,t}/\mathfrak m_t^2
 \oplus\C_{p_1}^{\ell_1(z)}\oplus\C_{p_2}^{\ell_2(z)}.\tag{C.12}\label{cdp:eq:C.12}
\end{align*}
The natural map to $H^2(\mathrm dP_6,K_{\mathrm dP_6})$ is surjective for all these nonintegral twists.
\end{proposition}



\begin{proof}
\noindent\textbf{Step 1.} On a marked smooth torus, the exponential sequence presents $\Pic^0$ as constant complex lattice cocycles modulo the integral cocycles and the two holomorphic period rows. The equation for $z\gamma$ to belong to this sum is
\[
 z\gamma=(n,p,q,r)+a(6\mu,\tau,1,0)+b(\beta_c,\mu,0,1).
\]
The last two coordinates give $a=-q,b=-r$, and the first two give \eqref{cdp:eq:C.10}. If $r=0$, positivity of $\operatorname{Im}\tau$ forces $p=q=0$ and $z=n$. Otherwise
\[
 \operatorname{Im}z=-r\left(\operatorname{Im}\beta_c-
           6(\operatorname{Im}\mu)^2/\operatorname{Im}\tau\right)\ne0.
\]
In particular real nonintegral twists have no regular resonance. The fixed values $\mu(z_1)=(2-\rho)/3$ and $\mu(z_2)=(1-i)/2$ give the two displayed lattices by solving the first equation. Their nonreal generators make each representation $n+k\lambda_j$ unique.

\noindent\textbf{Step 2.} We give a coherent local complex to calculate every length. On a smooth simply connected torus-family disc use normalized period coordinates $\zeta=Z(s)x+y$, $x,y\in\R^2/\Z^2$, and translation-invariant forms $G_j=(dx_j)^{0,1}$. The invariant horizontal lift is $h=\partial_s+\sum_i(Z'(s)x)_i\partial_{\zeta_i}$. Since $dx_j$ is closed, $\bar\partial G_j=0$ and
$\bar\partial h=\sum Z'_{ij}e_iG_j$. After subtracting the integer cocycle of a resonance, the line multiplier has invariant Dolbeault form $a_1G_1+a_2G_2$, up to the common nonzero factor $2\pi i$. Thus the invariant relative complex has terms
\[
 C^j=(V\oplus\OO h)\otimes\Lambda^jG,\qquad j=0,1,2,
\]
\[
 d(v\eta)=v(a\wedge\eta),\qquad
 d(h\eta)=h(a\wedge\eta)+\kappa\wedge\eta,             \tag{C.13}\label{cdp:eq:C.13}
\]
where
\[
 a=(z-n+6q\mu+r\beta_c,-p+q\tau+r\mu),\qquad
 \kappa=\begin{pmatrix}6\mu'&\tau'\\\beta_c'&\mu'\end{pmatrix}.
\]
This complex computes the derived direct image. To verify the comparison, decompose relative smooth forms into Fourier characters under real torus translations. The zero character gives \eqref{cdp:eq:C.13}. For every other character the relative Dolbeault differential is exterior multiplication by a nonzero covector; contraction by its normalized dual is a homotopy. On a sufficiently small closed subdisc those covectors are bounded away from zero at bounded frequency and grow linearly at large frequency. The homotopies preserve smooth Fourier series and all their derivatives. The vertical-horizontal filtration of the tangent extension, followed by the filtration in the base antiholomorphic differential, gives a finite filtration with acyclic nonzero-character quotients. Hence their full complexes are acyclic. The zero-character formulas have holomorphic coefficients, so the remaining base Dolbeault resolution is precisely the resolution of \eqref{cdp:eq:C.13}. A finite cyclic stabilizer normalizes translations, so this construction and its averaging projection are equivariant on each root chart.

\noindent\textbf{Step 3.} At a regular resonance, $a'=\kappa^t(q,r)^t$ and $(q,r)\ne(0,0)$. The period determinant calculation in Part E shows that $\kappa$ is invertible at every regular point. Thus $a$ has a simple zero; choosing its direction as the second cocycle frame and changing the vertical frame gives
\[
 a=(0,s),\qquad\kappa=\begin{pmatrix}0&1\\1&0\end{pmatrix}.
\]
In degree one the two generators $e_1G_2,e_2G_2$ are killed by $s$. In degree two the relations eliminate the two vertical generators and leave $hG_1G_2$ killed by $s^2$. This gives the regular modules in \eqref{cdp:eq:C.12}.

\noindent\textbf{Step 4.} At the order-three point put $t=s^3$, $\nu=\mu-(2-\tau)/3$ and $b=\beta_c+4\mu+(2/3)\tau-b(0)$. The order-three characters give $\nu=\nu_2s^2+O(s^3)$ and $b=b_3s^3+O(s^4)$. We check the joint coefficient that controls the module. With elliptic scale $f(s)=Fs^2+O(s^3)$, the Weierstrass equation gives
\[
 \nu_2^2=9F^2/512,\quad \tau_1=i\pi F^2/2,\quad
 6\mu'^2-\tau'\beta_c'=(3F^2/4)s^2+O(s^3).
\]
The first equality follows from the exact tripling coordinate $x(3P)(0)=512/81$ and $x=f^2\wp$: $3\nu$ approaches zero, so $x(3P)(0)=F^2/(9\nu_2^2)$. The determinant equality follows from the identity in Part E and $f^4g_2=12t^3(t-1)$, $f^6g_3=-8t^4(t-1)^2$. Comparing coefficients in $6\nu'^2-\tau'b'$ gives $b_3=7i/(32\pi)$. For a nonintegral resonance the Picard section is $3k(\nu,b)$ in the corresponding frames. Eliminating the unit $\tau'$ leaves transverse coefficient
\[
 3k(b-6\nu\nu'/\tau')
   =\frac{123ik}{64\pi}s^3+O(s^4),\qquad k\ne0.        \tag{C.14}\label{cdp:eq:C.14}
\]
It is nonzero. Equivariant holomorphic frame operations therefore reduce \eqref{cdp:eq:C.13} to
\[
 a=(0,s^2),\qquad\kappa=\begin{pmatrix}0&s\\s&1\end{pmatrix}.
\]
Its degree-one and degree-two modules are both $\C\{s\}/(s^2)\oplus\C\{s\}/(s^3)$. Explicit degree-one cycles are $u=e_1G_2$, $v=e_2G_1+shG_2$; the third cycle is removed by the boundary $su+e_2G_2+sv$. The remaining relations are $s^2u=s^3v=0$. In degree two elimination of $e_2G_1G_2$ leaves the generators $e_1G_1G_2,hG_1G_2$ with annihilators $s^2,s^3$.

\noindent\textbf{Step 5.} We retain the actual affine characters. The integral cocycle at the order-three resonance is $(n,-2k,-k,-3k)$; evaluating on $v_1$ gives $n$. At the order-four resonance it is $(n,-k,-k,-2k)$; evaluation on $v_2$ gives $-n\equiv3n\pmod4$. The invariant exponents at the order-three point are $n+2,n+2$ for the two degree-one generators and $n+2,n$ in degree two. Here an invariant exponent is the required Taylor exponent modulo the stabilizer order. Thus the length-two summand retains precisely the monomials $s^j$ with $0\le j<2$ and $j\equiv n+2\pmod3$: its invariant lengths for $n\equiv0,1,2$ are $0,1,1$. A length-three summand retains one coefficient in every residue class. Both coarse lengths are therefore $1+\epsilon_n$, and multiplication by $t=s^3$ kills them. At the order-four point the regular normal form applies. The two length-one degree-one generators have invariant exponents $3n+1,3n+2$, and the length-two degree-two generator exponent $3n+2$. Their invariant length is one exactly for $n\equiv1,2\pmod4$, and the coarse parameter $s^4$ kills them. This proves all noncusp summands.

\noindent\textbf{Step 6.} At the cusp, the restriction of $M_z$ has normalization pullback trivial and dual-complex character $(a,1)$, $a=e^{2\pi iz}$. Its normalization complex is
\[
 \C\xrightarrow{(a-1,0)^t}\C^2
 \xrightarrow{(0,a-1)}\C.
\]
For $z\notin\Z$ it is acyclic. Semicontinuity and base change make $R\bar f_*M_z$ vanish on a cusp neighbourhood. The logarithmic sequence (C.8) then gives $R\bar f_*(\mathcal T\otimes M_z)=0$ there. Sequence (C.3) contributes exactly the reduced skyscraper $\C_{p_0}$ in degree two for the ordinary tangent bundle. Away from resonances a nontrivial topologically trivial line on a torus has zero cohomology: its unitary Fourier presentation has no zero eigenmode; the tangent extension by trivial bundles retains that acyclicity. We verify that this analytic support is proper. Fix a simply connected disc in the regular base. For each fixed integer tuple $(n,p,q,r)$, a resonance has $r\ne0$ and $a^\prime=\kappa^t(q,r)^t\ne0$, so its zeros are isolated. Each such zero set is countable on this disc, and the union over the countably many tuples is countable. The disc has a point outside this union. At that fibre the preceding line-bundle acyclicity and tangent extension give zero cohomology in every degree. Upper semicontinuity gives a neighbourhood with the same vanishing. Each coherent higher direct image therefore has proper analytic support on the compact base curve, hence finite support. Every regular resonance lies in that support by its computed nonzero local modules, proving that $\mathcal R_z$ is finite. Degree zero is zero because a direct image of a bundle is torsion-free and has generic rank zero. These arguments prove \eqref{cdp:eq:C.12}, and Leray proves \eqref{cdp:eq:C.11}. The logarithmic degree-two direct image is a torsion sheaf for these twists, so its first cohomology vanishes. Equation (C.4) proves the final surjectivity assertion.
\end{proof}



Statements \ref{cdp:st:C5b}--\ref{cdp:st:C5c} reproduce, with the displayed comparisons and independent exact module computations, the groups and stalk lengths claimed in S6\_01 D.1--D.4. The extra conclusion (C.7) corrects the assertion that the map is nonzero for every twist: the target is always one-dimensional, but its global map is sometimes zero. All formulas give $\chi(T_X\otimes M_z)=1$. The generic tuple is $(0,0,1,0)$; it directly contradicts the degree-two vanishing in CDP20 Theorem 2.2(a), while its Euler characteristic and $c_3(X)=2$ contradict parts (b) and (c). A holomorphic universal Picard family is defined by $\bar\partial+z\alpha$ for any global closed Dolbeault representative $\alpha$ of the fixed generator. Proper coherent base change makes every cohomology threshold analytic. Real nonintegral parameters attain the generic tuple, so every higher threshold is a proper locally discrete analytic subset of the Picard line. This establishes the asserted jump-locus interpretation as well as the individual dimensions.



\section{Constant-monodromy period comparison and compact examples}\label{sec:II:13}


\begin{proposition}[Proved obstruction]\label{cdp:st:D.2.1}
\StatementRecord{cdp:st:D.2.1}{proved}{Historical events: Claude Fable 5.2 (Anthropic), effort high, 16 November 2026: carried as background without fault found. GPT-6 Astra (OpenAI), ChatGPT Pro mode, effort Pro, S6\_09\_PAPERS, 17 November 2026: targeted proof check. Also checked in this writing. Pre-continuation writing events: GPT-6 Astra (OpenAI), ChatGPT Pro mode, effort Pro, S6\_10\_PAPERS, 18 November 2026: targeted proof check. Continuation metadata events: GPT-6 Astra (OpenAI), ChatGPT Pro mode, effort Pro, S6\_10\_PAPERS continuation, 19 November 2026: used and rechecked in this run; status-vocabulary reconciliation; metadata reconciliation only; no new mathematical check. \emph{Pre-continuation recorded proof scope.} The integral complement formula gives r=s+3 under trivial four-dimensional rational monodromy; the actual adjacent-component divisor excludes every proposed Kahler resolution}{the hypotheses, notation and definitions stated here; \cref{cdp:st:D.1.2}; \cref{cdp:st:C.3.4}}

A compact complex threefold with $b_1=b_2=0$ and a holomorphic
connected-fibre map to a curve, with general fibre a two-torus and
trivial rational monodromy, has reducible fibres. More precisely,
for every nonempty set of $s$ removed values containing all critical
values, their reduced fibres have $r=s+3$ irreducible components in
total. Every component of each reducible fibre has no compact
K\"ahler resolution. Consequently no such fibration can have all
its fibre components admitting compact K\"ahler resolutions.
\end{proposition}


\begin{proof}
\noindent\textbf{Step 1.} The base is $\PP^1$ by \ref{cdp:st:D.1.2}. Trivial monodromy on the
four-dimensional $H_1(F;\Q)$ gives $t'=4$, so the corrected formula
gives $r=s+3$. If all removed fibres were irreducible one would have
$r=s$, which is impossible. The assertion about each component is
Statement \ref{cdp:st:C.3.4}, applied to this threefold. The last assertion follows
by combining these two conclusions.
\end{proof}






\begin{proposition}[Proved logarithmic period obstruction]\label{cdp:st:D.2.3}
\StatementRecord{cdp:st:D.2.3}{proved}{Historical events: Claude Fable 5.2 (Anthropic), effort high, 16 November 2026: carried as background without fault found. GPT-6 Astra (OpenAI), ChatGPT Pro mode, effort Pro, S6\_09\_PAPERS, 17 November 2026: targeted proof check. Also checked in this writing. Pre-continuation writing events: GPT-6 Astra (OpenAI), ChatGPT Pro mode, effort Pro, S6\_10\_PAPERS, 18 November 2026: targeted proof check; used and rechecked in this run. GPT-6 Astra (OpenAI), ChatGPT Pro mode, effort Pro, S6\_10\_PAPERS, 18 November 2026 04:34:17.629108+00:00: used and rechecked in this run. Continuation metadata events: GPT-6 Astra (OpenAI), ChatGPT Pro mode, effort Pro, S6\_10\_PAPERS continuation, 19 November 2026: used and rechecked in this run; status-vocabulary reconciliation; metadata reconciliation only; no new mathematical check. \emph{Pre-continuation recorded proof scope.} The actual relative logarithmic complex, zero-residue flat extension, annular logarithmic-volume divergence, nilpotent-layer monomial Koszul spectrum a/e, equivariant residue map, and primitive-form power divergence were re-derived with arbitrary multiplicities. Steenbrink 1976 section2 and Proposition1.13 were opened in the original journal text; its analytic hypotheses match and impose no Kahler condition}{the hypotheses, notation and definitions stated here; \cite[\S2.1, Theorem~2.18, \S2.19, Proposition~2.20 and
Theorem~2.21]{Steenbrink1976}; \cite[Proposition~1.13 and proof of
Proposition~2.20]{Steenbrink1976}; \cref{cdp:eq:D.18}; \cref{cdp:eq:D.19}}


Let $h:U\to\Delta$ be a proper holomorphic map from a smooth
complex threefold, with connected fibres, smooth two-torus fibres
over $\Delta^*$, and central fibre $F=\sum_i m_iD_i$ with
normal-crossings reduced support $D=\bigcup_iD_i$.
If every normalization component of $D$ has geometric genus zero,
the monodromy on $H^2(U_t;\C)$ is nontrivial. In particular the
monodromy on $H^1(U_t;\Q)$ is nontrivial. No K\"ahler assumption
on $U$ or $D$ is made, and the positive integers $m_i$ are arbitrary.
\end{proposition}


\begin{proof}
\noindent\textbf{Step 1.} We first identify exactly the analytic comparison used in the
argument. Put
\[
\mathscr K^\bullet=
 \Omega_U^\bullet(\log D)\big/
 \bigl(h^*\Omega_\Delta^1(\log0)\wedge
       \Omega_U^{\bullet-1}(\log D)\bigr),\qquad
\mathscr H^2=\mathbb R^2h_*\mathscr K^\bullet.
\]
In \cite[\S2.1, Theorem~2.18, \S2.19, Proposition~2.20 and
Theorem~2.21]{Steenbrink1976}, the assumptions for these assertions
are precisely a proper holomorphic map from a connected complex
manifold to a disc, smooth off zero, with normal-crossings support.
They precede the additional Hodge-theoretic assumptions of the
later sections. They give a locally free bundle $\mathscr H^2$,
its actual Gauss--Manin connection
\[
\nabla:\mathscr H^2\longrightarrow
\Omega_\Delta^1(\log0)\otimes\mathscr H^2,
\]
and the comparison
$\mathscr H^2|_{\Delta^*}\simeq
R^2h_*\C\otimes\OO_{\Delta^*}$.
The residue eigenvalues belong to $[0,1)$ and the extended
monodromy at zero is
$\exp(-2\pi i\operatorname{Res}_0\nabla)$.
Its canonical base-change map identifies the fibre at zero with
$\mathbb H^2(D,\mathscr K^\bullet/t\mathscr K^\bullet)$.
The local calculation used below is the same monomial Koszul
calculation as \cite[Proposition~1.13 and proof of
Proposition~2.20]{Steenbrink1976}; we give it explicitly for the
subcomplex needed here. No local invariant-cycle surjectivity
or polarization is invoked.

\noindent\textbf{Step 2.} Suppose that the monodromy on $H^2$ were the identity. Its
residue has all eigenvalues zero, and its exponential is the
identity. A nilpotent matrix whose exponential is the identity
is zero. Hence the residue vanishes. In a holomorphic frame
$\nabla=d+A(t)\,dt/t$ with $A(0)=0$, so this is a holomorphic
connection across zero. It has a holomorphic horizontal frame
on the disc. Consequently the coordinates of every holomorphic
section of $\mathscr H^2$ in a flat frame, that is its periods,
are holomorphic and bounded near zero.

\noindent\textbf{Step 3.} The determinant of the locally free relative logarithmic complex is
\[
\mathscr K^2\simeq
 \omega_U(D)\otimes h^*\omega_\Delta(0)^{-1}
 \simeq\omega_{U/\Delta}(D-F).
\]
The direct image $h_*\mathscr K^2$ is a torsion-free
coherent sheaf of rank one, because a smooth torus fibre has
one-dimensional space of holomorphic two-forms. Shrink the disc
and choose a generator $\Omega$ of this line bundle.
Thus $\Omega$ determines a holomorphic section $[\Omega]$ of
$\mathscr H^2$ by the inclusion of the last term of the
relative de Rham complex. Its smooth-fibre comparison is the
ordinary class $[\Omega_t]\in H^2(U_t;\C)$.

\noindent\textbf{Step 4.} First suppose that $\Omega$ restricts nontrivially to a normalized
component $T$ of $D$. There is a canonical isomorphism
\[
 \mathscr K^1|_T\simeq\Omega_T^1(\log D_T),\qquad
 \mathscr K^2|_T\simeq\omega_T(D_T),
\]
where $D_T$ is the reduced conductor. To specify the map, on the
branch $z_i=0$ subtract
$\operatorname{res}_i(\eta)(dt/t)/m_i$ from an absolute logarithmic
one-form $\eta$ and restrict its tangential part. The normal
residue of $dt/t$ is $m_i$, so this kills exactly the subbundle
divided out in $\mathscr K^1$ and is independent of the lift.
The exterior square gives the second isomorphism. If every
residue of $\Omega|_T$ vanished, its local logarithmic expression
would be holomorphic on $T$. This contradicts $p_g(T)=0$.
Thus there is a smooth point of a double curve, away from triple
points, where its residue is nonzero.

\noindent\textbf{Step 5.} In a small chart about that point, write $h=x^ay^b=t$ and use $z$
along the double curve. The relative two-form is
\[
 \Omega_t=g(x,y,z)\,\frac{dx}{x}\wedge dz,
 \qquad g(0,0,z_0)\ne0.
\]
Choose a fixed small $z$-disc and $\epsilon>0$ such that
$|g(x,y,z)|\ge c>0$ whenever $|x|,|y|<\epsilon$ there.
The map to the annulus
$|t|^{1/a}\epsilon^{-b/a}<|x|<\epsilon$
has $b$ sheets. Integration on its inverse image and the fixed
$z$-disc gives
\[
 \int_{U_t}\Omega_t\wedge\overline{\Omega_t}
 \ \ge\ c'\log\frac{\epsilon^{a+b}}{|t|}
 \ \longrightarrow\ +\infty.
\]
Here the integrand is a nonnegative real volume form on the
complex surface; its local expression is a positive constant
times $|g|^2|x|^{-2}\,dA(x)dA(z)$.

\noindent\textbf{Step 6.} On the other hand, $\Omega_t$ is closed, and this integral is
the fixed cup-product pairing of $[\Omega_t]$ with its complex
conjugate. All its periods are bounded by the holomorphic
connection argument, so the pairing is bounded. This is a
contradiction.

\noindent\textbf{Step 7.} It remains to treat the case where $\Omega$ vanishes on every
branch of $D$. Put
\[
 \mathscr J=\OO_U(-D),\qquad
 \mathscr B^\bullet=
 \mathscr J\mathscr K^\bullet/t\mathscr K^\bullet.
\]
Both $\mathscr J\Omega_U^\bullet(\log D)$ and
$t\Omega_U^\bullet(\log D)$ are subcomplexes for the absolute
exterior derivative: locally the logarithmic derivatives of
$z_1\cdots z_r$ and $t=\prod z_i^{m_i}$ are holomorphic
linear combinations of the $d\log z_i$. Hence the locally split
absolute/relative logarithmic exact sequence remains an exact
sequence of complexes after tensoring by $\mathscr J/(t)$:
\[
0\longrightarrow\mathscr B^\bullet[-1]
 \xrightarrow{\ (dt/t)\wedge\ }
 (\mathscr J/(t))\Omega_U^\bullet(\log D)
 \longrightarrow\mathscr B^\bullet\longrightarrow0.
\tag{D.18}\label{cdp:eq:D.18}
\]
Signs in the shifted differential are chosen so that the first
arrow is a morphism of complexes. Its connecting morphism
$\rho_{\mathscr B}:\mathscr B^\bullet\to\mathscr B^\bullet$
in the derived category defines degree-preserving endomorphisms
on both cohomology sheaves and hypercohomology. The inclusion
$\mathscr J/(t)\hookrightarrow\OO_U/(t)$ gives a commutative
diagram from \eqref{cdp:eq:D.18} to the analogous exact sequence for
$\mathscr K^\bullet/t\mathscr K^\bullet$. The induced
endomorphism for the latter is precisely the Gauss--Manin
residue. Thus the map on hypercohomology intertwines these
endomorphisms.

\noindent\textbf{Step 8.} We compute the eigenvalues of $\rho_{\mathscr B}$ without
assuming reduced multiplicities. At a point of $r$ branches,
write
\[
 t=\prod_{i=1}^r z_i^{m_i},\qquad
 e=\gcd(m_1,\ldots,m_r),\qquad
 y=\prod_{i=1}^r z_i^{m_i/e},\qquad t=y^e.
\]
The transverse holomorphic Poincar\'e lemma reduces the relative
complex to the Koszul complex with commuting operators
\[
 D_i=z_i\partial_{z_i}
           -\frac{m_i}{m_1}z_1\partial_{z_1},
 \qquad 2\le i\le r,
\]
on the convergent-power-series module
$(z_1\cdots z_r)/(z_1^{m_1}\cdots z_r^{m_r})$.
The monomials diagonalize these operators. A monomial with some
nonzero eigenvalue has a contracting Koszul homotopy, obtained
by division by that eigenvalue. These divisions preserve
convergence: every nonzero eigenvalue has absolute value at
least $1/m_1$. The surviving monomials have exponent vector
$a(m_1/e,\ldots,m_r/e)$. Membership in $\mathscr J$ requires
$a\ge1$, and quotienting by $t$ requires $a<e$. Consequently
each cohomology stalk is spanned by
\[
 y^a\,d\log z_{i_1}\wedge\cdots\wedge d\log z_{i_q},
 \qquad 1\le a<e,
\]
subject to the relative relation $\sum m_i d\log z_i=0$.
Since $dy^a=(a/e)y^a\,dt/t$, the connecting endomorphism
acts on these representatives by $a/e$. In particular zero is
not an eigenvalue; when $e=1$ the complex is acyclic.

\noindent\textbf{Step 9.} The hypercohomology spectral sequence
\[
 H^p\bigl(D,\mathcal H^q(\mathscr B^\bullet)\bigr)
 \ \Longrightarrow\ \mathbb H^{p+q}(D,\mathscr B^\bullet)
\]
is equivariant for this derived endomorphism. There are only
finitely many branch multiplicities, so all eigenvalues on its
terms and abutment lie in a finite subset of $(0,1)$. For
example the product of the corresponding linear factors
annihilates every cohomology sheaf, and a power annihilates the
abutment. Hence $\rho_{\mathscr B}$ is invertible on
$\mathbb H^2(D,\mathscr B^\bullet)$. The residue on
$\mathscr H^2|_0$ is zero. Equivariance therefore proves that
\[
 \mathbb H^2(D,\mathscr B^\bullet)\longrightarrow
 \mathbb H^2(D,\mathscr K^\bullet/t\mathscr K^\bullet)
       =\mathscr H^2|_0
 \quad\hbox{is the zero map}.
\tag{D.19}\label{cdp:eq:D.19}
\]

\noindent\textbf{Step 10.} The central class of $\Omega$ comes from the left side of
\eqref{cdp:eq:D.19}, because $\Omega\in\mathscr J\mathscr K^2$.
Thus $[\Omega]$ vanishes at zero in the locally free bundle
$\mathscr H^2$, and $[\Omega]/t$ is a holomorphic section of
that bundle. Its periods are bounded. On a smooth fibre this
section is the de Rham class of $\Omega_t/t$.

\noindent\textbf{Step 11.} Primitivity of the generator $\Omega$ implies that its order
$b_i$ along at least one $D_i$ is less than $m_i$; otherwise
$\Omega/t$ would be a holomorphic section of $\mathscr K^2$
over $U$, contradicting its choice as a generator of
$h_*\mathscr K^2$. Choose a smooth point of that component
away from all other branches and all further zeros. There
are coordinates with $t=z^{m_i}$ in which
\[
 \Omega=z^{b_i}a(z,\xi,\eta)\,d\xi\wedge d\eta,
 \qquad a(0,0,0)\ne0.
\]
On one sheet above a fixed small $(\xi,\eta)$-polydisc, the
positive integral of $\Omega_t/t$ against its conjugate has
the lower bound
\[
 \int_{U_t}\frac{\Omega_t}{t}\wedge
                    \overline{\frac{\Omega_t}{t}}
 \ \ge c''|t|^{2(b_i/m_i-1)}\longrightarrow+\infty.
\]
Its expression as the fixed cup-product pairing of bounded
periods again gives a contradiction. This proves nontriviality
of the $H^2$ monodromy in the remaining case.

\noindent\textbf{Step 12.} Finally the natural
cup-product isomorphism
$\bigwedge^2H^1(U_t;\Q)\simeq H^2(U_t;\Q)$ for a two-torus
is monodromy-equivariant, proving the final assertion.
\end{proof}






\begin{proposition}[Proved sign in the full normal-crossings class]\label{cdp:st:D.2.4}
\StatementRecord{cdp:st:D.2.4}{proved}{Historical events: Claude Fable 5.2 (Anthropic), effort high, 16 November 2026: carried as background without fault found. GPT-6 Astra (OpenAI), ChatGPT Pro mode, effort Pro, S6\_09\_PAPERS, 17 November 2026: targeted proof check. Also checked in this writing. Pre-continuation writing events: GPT-6 Astra (OpenAI), ChatGPT Pro mode, effort Pro, S6\_10\_PAPERS, 18 November 2026: targeted proof check. Continuation metadata events: GPT-6 Astra (OpenAI), ChatGPT Pro mode, effort Pro, S6\_10\_PAPERS continuation, 19 November 2026: used and rechecked in this run; status-vocabulary reconciliation; metadata reconciliation only; no new mathematical check. \emph{Pre-continuation recorded proof scope.} The normalized root cover is etale with reduced NC fibre, its rational-normalization covers have geometric genus zero, and its identity monodromy contradicts D.2.3. This excludes every positive NC Euler contribution}{the hypotheses, notation and definitions stated here; \cref{cdp:st:C.2.NC.5}; \cref{cdp:st:C.2.NC.3}; \cref{cdp:st:D.2.3}; \cref{cdp:st:D.2.7}}


For the threefolds and maps in D.2 whose reduced fibre supports
all have normal crossings, the Euler number is
\[
c_3(Y)=0,\qquad b_3(Y)=2.
\]
All multiplicities, including unequal multiplicities on a
reducible support, are allowed. In this class there are no
positive-Euler-number examples with trivial monodromy.
\end{proposition}


\begin{proof}
\noindent\textbf{Step 1.} Statement \ref{cdp:st:C.2.NC.5} gives $c_3(Y)=2r$, where $r$ counts the
irreducible fibre supports with rational normalization. A fibre
of that kind is necessarily $mS$ for one positive integer $m$.
Choose a small disc around its value and perform the normalized
base change $t=u^m$. Statement \ref{cdp:st:C.2.NC.3} identifies the resulting
map to the original neighbourhood as finite \emph{\'etale} of
degree $m$, and its central fibre $S'$ as reduced with normal
crossings. The total space is therefore smooth. Its nearby
fibres are the original connected tori. Stein factorization
over the normal disc shows that its central fibre is connected
as well: the finite Stein factor has generic degree one and is
the disc itself.

\noindent\textbf{Step 2.} Normalization commutes with this \emph{\'etale} base change.
Each normalized component of $S'$ is an unramified finite cover
of the smooth rational normalization $T$ of $S$. Since $T$ is
simply connected, these normalized components are copies of
$T$, and all have geometric genus zero. The local monodromy
after base change is the $m$th power of the original monodromy,
so remains the identity on $H^1$. This contradicts \ref{cdp:st:D.2.3}.
Therefore $r=0$, proving $c_3(Y)=0$. Poincar\'e duality and
$b_1(Y)=b_2(Y)=0$ give $c_3(Y)=2-b_3(Y)$, hence $b_3(Y)=2$.
The example of \ref{cdp:st:D.2.7} realizes this normal-crossings class.
\end{proof}






\begin{proposition}[Proved unrestricted sign under trivial monodromy]\label{cdp:st:D.2.5}
\StatementRecord{cdp:st:D.2.5}{proved}{Historical events: Claude Fable 5.2 (Anthropic), effort high, 16 November 2026: carried as background without fault found. Claude Fable 5.2 (Anthropic), effort high, 17 November 2026: carried as background without fault found. GPT-6 Astra (OpenAI), ChatGPT Pro mode, effort Pro, S6\_09\_PAPERS, 17 November 2026: targeted proof check. Also checked in this writing. Pre-continuation writing events: GPT-6 Astra (OpenAI), ChatGPT Pro mode, effort Pro, S6\_10\_PAPERS, 18 November 2026: targeted proof check; used and rechecked in this run. Continuation metadata events: GPT-6 Astra (OpenAI), ChatGPT Pro mode, effort Pro, S6\_10\_PAPERS continuation, 19 November 2026: used and rechecked in this run; status-vocabulary reconciliation; metadata reconciliation only; no new mathematical check. \emph{Pre-continuation recorded proof scope.} The fixed-fibre and fixed-Cartier-layer reductions, rational-Betti smooth-support exclusion, all-powers Gauduchon degree bound, effective conductor negative canonical degree, and embedded resolution to geometric-genus-zero components were re-derived. D.2.3 then contradicts identity monodromy. Bierstone-Milman Theorem1.6 and its analytic-neighborhood scope were opened}{the hypotheses, notation and definitions stated here; \cref{cdp:st:C4c}; \cref{cdp:st:A.3.5}; \cref{cdp:eq:D.S.2}; \cref{cdp:st:D.1.9}; \cref{cdp:st:A.4.7}; \cref{cdp:st:A.5.2}; \cref{cdp:eq:D.S.4}; \cref{cdp:st:A.5.6}; \cref{cdp:st:A.5.3}; \cref{cdp:eq:D.S.5}; \cite[Theorem~1.6]{BierstoneMilman1997}; \cref{cdp:st:D.2.3}}


Let $f:Y\to\PP^1$ be a proper surjective holomorphic map with
connected fibres from a compact connected complex threefold with
$b_1(Y)=b_2(Y)=0$. Suppose its smooth general fibre is a complex
two-torus and its monodromy representation on $H_1(F,\Q)$ is
trivial. Then
\[
 c_3(Y)\le0,\qquad b_3(Y)\ge2.
 \tag{D.S.1}\label{cdp:eq:D.S.1}
\]
The fibres may have arbitrary singularities and arbitrary
multiplicities. In particular this proves the sign assertion of
D.2 for a holomorphic algebraic reduction; no existence assertion
is required for the conclusion.
\end{proposition}


\begin{proof}
\noindent\textbf{Step 1.} Assume $c_3(Y)>0$. Poincar\'e duality gives
$c_3(Y)=2-b_3(Y)$, and the alternating nondegenerate intersection
pairing on $H^3(Y,\Q)$ makes $b_3(Y)$ even. Hence $b_3(Y)=0$.
Apply \ref{cdp:st:C4c} to the Picard component
$P=K_Y\otimes\Pic^0(Y)$. It gives a fixed critical full fibre
$F$ such that
\[
 H^0\bigl(F,\Omega_Y^1|_F\otimes M|_F\bigr)\ne0
 \qquad(M\in\Pic^0(Y)).
 \tag{D.S.2}\label{cdp:eq:D.S.2}
\]
Indeed the Serre-dual parameter $K_Y\otimes L^{-1}$ runs through
$\Pic^0(Y)$ when $L$ runs through $P$.

\noindent\textbf{Step 2.} Statement \ref{cdp:st:A.3.5} excludes reducible support in \eqref{cdp:eq:D.S.2}, including
all unequal-multiplicity possibilities. Thus $F=mS$ for one
irreducible reduced Cartier divisor $S$ and one positive integer
$m$. We specify why a smooth $S$ is excluded under the present
\emph{rational} Betti-number assumptions. \ref{cdp:st:D.1.9} supplies
$\beta\in H^1(Y,\OO_Y)$ whose restriction near every full fibre
comes from a nonzero integral transgression class. The line
$M_r=\exp(2\pi i r\beta)$, for real irrational $r$, has that
unitary local realization. The normalized root change $t=s^m$
near a smooth $S$ is the finite \emph{\'etale} map of A.3.4:
in a chart $t=uz^m$ it has equations $w^m=u$, $s=wz$.
Its central reduction $S'$ is smooth, and its nearby fibres are
the original tori. Proper smooth local triviality transports
the vanishing of degree-zero and degree-one cohomology of the
pulled-back unitary local system from a nearby torus to $S'$.
The injectivity \ref{cdp:st:A.4.7} and the conormal sequence then give zero
ambient cotangent sections on $S'$. Each layer of $mS'$ has the
same twist, since $s$ trivializes its conormal. The finite
filtration and exact invariants functor for the cyclic cover
give $H^0(mS,\Omega_Y^1|_{mS}\otimes M_r)=0$, contradicting
\eqref{cdp:eq:D.S.2}. This proof uses only the rational version of \ref{cdp:st:D.1.9}.
Therefore $S$ is singular. \ref{cdp:st:A.5.2}, applied using
$\omega_S=(K_Y\otimes\OO_Y(S))|_S$ and
$\chi(\OO_S)=0$, makes $S$ nonnormal.

\noindent\textbf{Step 3.} The ideal filtration of $mS$ has the finitely many locally free
graded factors on $S$
\[
 E_j=(\Omega_Y^1\otimes\OO_Y(-jS))|_S,
 \qquad 0\le j<m.
 \tag{D.S.3}\label{cdp:eq:D.S.3}
\]
If all the groups $H^0(S,E_j\otimes M|_S)$ were zero, the exact
filtration sequences would make \eqref{cdp:eq:D.S.2} zero. Thus their
nonvanishing loci cover $\Pic^0(Y)$. Each locus is closed
complex analytic by the actual Picard family and proper
semicontinuity. The group $\Pic^0(Y)$ is a connected complex
vector space, since $H^1(Y,\Z)=0$. A finite union of proper
closed analytic subsets cannot cover it. Some fixed $j$
therefore satisfies
\[
 H^0(S,E_j\otimes M|_S)\ne0
 \qquad(M\in\Pic^0(Y)).
 \tag{D.S.4}\label{cdp:eq:D.S.4}
\]

\noindent\textbf{Step 4.} Let $\eta:T\to S$ be the finite normalization and
$\pi:Z\to T$ its minimal resolution, and put $h=\eta\pi$,
viewed also as a holomorphic map to $Y$.
Pullback of sections in \eqref{cdp:eq:D.S.4} is injective: $S$ is reduced
and irreducible, and a section of the locally free sheaf $E_j$
vanishing on its dense smooth locus is zero. Thus the fixed
vector bundle $h^*(\Omega_Y^1\otimes\OO_Y(-jS))$ has a
section after every twist by $h^*M$, $M\in\Pic^0(Y)$.
For any Gauduchon metric $\omega$ on $Z$, the uniform
negative-degree bound \ref{cdp:st:A.5.6}, applied to all positive and
negative powers of $M$, gives
\[
 \deg_\omega h^*M=0\qquad(M\in\Pic^0(Y)).
\]
Every line bundle $G$ on $Y$ has a positive power in
$\Pic^0(Y)$, because $H^2(Y,\Z)$ is finite. Additivity of
degree therefore gives $\deg_\omega h^*G=0$ for every such
$G$. No classification of the surface $Z$ is used in this step.

\noindent\textbf{Step 5.} The conductor comparison \ref{cdp:st:A.5.3} gives an effective divisor
$D_Z$ with
\[
 \omega_Z(D_Z)\simeq h^*(K_Y\otimes\OO_Y(S)).
 \tag{D.S.5}\label{cdp:eq:D.S.5}
\]
Moreover $D_Z\ne0$. Indeed a reduced hypersurface is
Cohen--Macaulay and hence satisfies Serre's condition $S_2$;
its nonnormality must therefore occur in codimension one.
The divisorial conductor is nonzero, and its strict transform
is part of $D_Z$.
Taking Gauduchon degree in \eqref{cdp:eq:D.S.5} yields
\[
 \deg_\omega K_Z=-\deg_\omega\OO_Z(D_Z)<0.
\]
A nonzero holomorphic section of $K_Z$ would have an effective
zero divisor, whose degree is nonnegative. This contradiction
proves $p_g(Z)=h^0(Z,K_Z)=0$.

\noindent\textbf{Step 6.} Choose a small closed fibre neighbourhood $U=f^{-1}(\Delta)$.
Embedded resolution of the divisor $F$ gives a proper
holomorphic modification
\[
 q:\widehat U\longrightarrow U
\]
which is a finite composite of blowups with smooth centres in
the successive central divisors, is an isomorphism over
$\Delta^*$, has smooth total space, and makes the reduced
support of $(f\circ q)^{-1}(0)$ a simple normal-crossings
divisor. We use the analytic embedded-resolution theorem in
\cite[Theorem~1.6]{BierstoneMilman1997}; properness of the
central fibre makes its locally finite construction finite
after shrinking $\Delta$. The strict transform of $S$ is a
resolution of its normalization and has geometric genus zero
by birational invariance of holomorphic two-forms on smooth
compact surfaces. A new exceptional component is $\PP^2$ over
a point or a $\PP^1$-bundle over a smooth compact curve.
Both have geometric genus zero: restriction of a holomorphic
two-form to a ruling chart has negative fibre degree, and
$H^0(\PP^2,K_{\PP^2})=0$. Further strict transforms preserve
this property. Consequently every component of the final
central reduced divisor has geometric genus zero.

\noindent\textbf{Step 7.} The map $f\circ q$ remains proper with connected fibres:
a blowup has connected fibres and its structure-sheaf direct
image is the original structure sheaf. Its smooth fibres over
$\Delta^*$ are the original two-tori under the isomorphism
$q|_{\Delta^*}$, so their $H^1$ monodromy remains the identity.
The multiplicities of the resolved central fibre are arbitrary
positive integers and are all permitted in \ref{cdp:st:D.2.3}. That
proposition now gives nontrivial monodromy on $H^2$ and hence
on $H^1$, a contradiction. This proves \eqref{cdp:eq:D.S.1}.
\end{proof}



The sign result uses the analytic logarithmic comparison of
Steenbrink. The explicit compact construction in \ref{cdp:st:D.2.7} proves
that the stated class is nonempty and attains $c_3=0$.
Together these results settle both the existence and the sign
questions in D.2.




\begin{proposition}[Proved extension of the marked torus period]\label{cdp:st:D.2.6}
\StatementRecord{cdp:st:D.2.6}{proved}{Historical events: Claude Fable 5.2 (Anthropic), effort high, 16 November 2026: carried as background without fault found. GPT-6 Astra (OpenAI), ChatGPT Pro mode, effort Pro, S6\_09\_PAPERS, 17 November 2026: targeted proof check. Also checked in this writing. Pre-continuation writing events: GPT-6 Astra (OpenAI), ChatGPT Pro mode, effort Pro, S6\_10\_PAPERS, 18 November 2026: targeted proof check; used and rechecked in this run. Continuation metadata events: GPT-6 Astra (OpenAI), ChatGPT Pro mode, effort Pro, S6\_10\_PAPERS continuation, 19 November 2026: used and rechecked in this run; status-vocabulary reconciliation; metadata reconciliation only; no new mathematical check. \emph{Pre-continuation recorded proof scope.} The bounded-period and positive-volume argument forces a nonzero decomposable limiting two-class. The actual kernel alpha wedge v=0 and its conjugate complement give the extended period plane; evaluation embeds the integral lattice into its dual and defines the proper group family. The killed-cycle obstruction to extending the original Albanese map was checked}{the hypotheses, notation and definitions stated here; \cref{cdp:st:D.2.3}; \cref{cdp:eq:D.20}}


Let $h:U\to\Delta$ be a proper holomorphic map from a smooth
complex threefold, smooth with connected two-torus fibres over
$\Delta^*$. Suppose that its integral $H^1$ monodromy is the
identity. The marked period map
\[
 \Delta^*\longrightarrow
 \mathcal T_2
 :=\{W\in\operatorname{Gr}(2,H^1(U_t;\C)):
                   W\oplus\overline W=H^1(U_t;\C)\}
\]
extends holomorphically across zero with value in $\mathcal T_2$.
In particular the relative Albanese group family over $\Delta^*$
extends to a smooth proper group family of complex two-tori over
$\Delta$. This assertion does not assert an extension of the
Albanese map from the given proper model $U$.
\end{proposition}


\begin{proof}
\noindent\textbf{Step 1.} An embedded resolution of the central fibre changes neither
the smooth fibres nor their marking, so we may use the relative
logarithmic complex of \ref{cdp:st:D.2.3} with arbitrary multiplicities.
Identity on $H^1$ gives identity on $H^2=\bigwedge^2H^1$.
Consequently its logarithmic Gauss--Manin bundle has zero
residue and a holomorphic flat frame. Let $\Omega$ be a
generator of $h_*\mathscr K^2$, as in \ref{cdp:st:D.2.3}.

\noindent\textbf{Step 2.} The second case of that proof, including the exact maps
(D.18)--(D.19), did not use any geometric-genus assumption.
It shows that $\Omega$ cannot vanish on the whole reduced
central support: otherwise its divided cohomology class would
have bounded periods while its positive $L^2$ volume diverges.
Thus $\Omega|_T$ is nonzero on some normalized component $T$.
The first case of that proof shows, without a geometric-genus
assumption, that every conductor residue of every component
restriction is zero, because a nonzero residue would give the
same divergent positive volume. The nonzero restriction
$\Omega|_T$ is therefore a holomorphic two-form on $T$.

\noindent\textbf{Step 3.} Choose a point of $T$ outside the conductor and the zeros of
this two-form. In a small chart the map is $t=z^m$ and
\[
 \Omega=a(z,\xi,\eta)\,d\xi\wedge d\eta,
 \qquad a(0,0,0)\ne0.
\]
Integrating on one sheet over a fixed small
$(\xi,\eta)$-polydisc yields a constant $c>0$ such that
\[
 \int_{U_t}\Omega_t\wedge\overline{\Omega_t}\ge c
 \qquad(0<|t|\ll1).
\tag{D.20}\label{cdp:eq:D.20}
\]
The flat cohomology class $v(t)=[\Omega_t]$ is holomorphic
across zero. Its cup pairing with its conjugate equals the
integral in \eqref{cdp:eq:D.20}, hence
\[
 v(0)\wedge\overline{v(0)}>0
 \quad\hbox{in }H^4(U_t;\R)\simeq\R,
\]
using the complex orientation. In particular $v(0)\ne0$.

\noindent\textbf{Step 4.} Write $V=H^1(U_t;\C)$ with its fixed integral marking.
For $t\ne0$, $v(t)$ is a decomposable element of
$\bigwedge^2V$, spanning $\bigwedge^2H^{1,0}(U_t)$.
The decomposable cone is the closed algebraic set defined by
the Pl\"ucker equations, so $v(0)$ is decomposable. For a
nonzero decomposable $v$, the exact inverse to the Pl\"ucker
map is
\[
 W(v)=\{\alpha\in V:\alpha\wedge v=0\}.
\]
This kernel has dimension two and varies holomorphically near
$v(0)$. The nonzero pairing
$v(0)\wedge\overline{v(0)}$ says exactly that
$W(v(0))\cap\overline{W(v(0))}=0$. The extended period
therefore lies in $\mathcal T_2$, as claimed.

\noindent\textbf{Step 5.} For completeness the extending group family is explicit.
Let $\mathcal W\subset V\otimes\OO_\Delta$ be the extended
rank-two subbundle $W(v(t))$, and let
$\Lambda=H_1(U_t;\Z)$ with its fixed marking. Evaluation of
cohomology on integral cycles gives a holomorphic embedding
of the constant lattice $\Lambda$ in $\mathcal W^*$.
The direct-sum relation with the conjugate makes this a full
real lattice in every fibre, uniformly on compact subsets of
the disc. Hence
\[
 \mathcal A=\mathcal W^*/\Lambda\longrightarrow\Delta
\]
is a smooth proper family of complex Lie groups with
two-torus fibres. On $\Delta^*$, integration of holomorphic
one-forms identifies each group fibre with the Albanese torus
of $U_t$, and these identifications vary holomorphically.
This is the asserted exact isomorphism of relative group
families. Extending a map from the original model $U$ to
$\mathcal A$ requires a further argument; it does not follow
merely from this period construction.
\end{proof}



\paragraph{Scope of \ref{cdp:st:D.2.6}.}
For a fibration over $\PP^1$ with globally trivial marked
$H^1$ monodromy, the local constructions glue to a smooth
proper torus group family over $\PP^1$. The period map may
be nonconstant, as the explicit family in \ref{cdp:st:D.1.15} demonstrates.
The local filling in \ref{cdp:st:D.2.7} proves that the marked Albanese
map need not extend to the given model. Its inclusion of
a nearby fibre kills the nonzero cycle $w$ in (D.25).
An extension $U\to A\times\Delta$ of the marked fibre
isomorphisms would induce a map on $H_1$ whose composite
with that inclusion is the identity on $H_1(A;\Z)$,
which is impossible because it would send $w$ both to
zero and to $w$. This exact topological obstruction
explains why the extending group family does not imply
nonexistence in D.2.




\begin{proposition}[Proved existence]\label{cdp:st:D.2.7}
\StatementRecord{cdp:st:D.2.7}{proved}{Historical events: Claude Fable 5.2 (Anthropic), effort high, 15 November 2026: carried as background without fault found. Claude Fable 5.2 (Anthropic), effort high, 16 November 2026: targeted proof check. GPT-6 Astra (OpenAI), ChatGPT Pro mode, effort Pro, S6\_09\_PAPERS, 17 November 2026: targeted proof check. Also checked in this writing. Pre-continuation writing events: GPT-6 Astra (OpenAI), ChatGPT Pro mode, effort Pro, S6\_10\_PAPERS, 18 November 2026: targeted proof check. Continuation metadata events: GPT-6 Astra (OpenAI), ChatGPT Pro mode, effort Pro, S6\_10\_PAPERS continuation, 19 November 2026: used and rechecked in this run; status-vocabulary reconciliation; metadata reconciliation only; no new mathematical check. \emph{Pre-continuation recorded proof scope.} The O(R) nonzero-line quotient by q\textasciicircum{}Z, Kodaira and Hopf central components, marked local homology quotients and bounding mixed cycles, exact logarithmic clutching, and both integral Mayer-Vietoris relation systems were re-derived. Independent determinant computation gives primitive full-rank relations in degrees one and two. Euler additivity and actual meromorphic descent give all remaining assertions}{the hypotheses, notation and definitions stated here; \cref{cdp:st:D.2.5}; \cref{cdp:eq:D.24}; \cref{cdp:eq:D.25}; \cref{cdp:eq:D.26}; \cref{cdp:eq:D.23}; \cref{cdp:eq:D.28}; \cref{cdp:eq:D.30}; \cref{cdp:eq:D.32}; \cref{cdp:eq:D.31}}


There exists a smooth compact connected complex threefold $Y$
with a holomorphic map $f:Y\to\PP^1$ with connected fibres such
that
\[
 H_1(Y;\Z)=H_2(Y;\Z)=0,\qquad
 a(Y)=1,\qquad c_3(Y)=0,\qquad b_3(Y)=2.
\tag{D.21}\label{cdp:eq:D.21}
\]
Every smooth fibre is the fixed abelian surface
$A=E\times E$, where $E=\C/(\Z+\mathrm i\Z)$, and the marked
integral monodromy on $H_1(A;\Z)$ is the identity. There are
exactly three singular fibres. Each is reduced with simple
normal crossings, and consists of a primary Kodaira surface
and a primary Hopf surface meeting along one smooth elliptic
curve. Consequently the trivial-monodromy class in D.2 is
nonempty, and equality in the unrestricted sign theorem
\ref{cdp:st:D.2.5} is attained.
\end{proposition}


\begin{proof}
\noindent\textbf{Step 1.} We give the local holomorphic filling first, including its
boundary homology maps. Put $q=e^{-2\pi}$. The map
\[
 \C/(\Z+\mathrm i\Z)\longrightarrow\C^*/q^{\Z},
 \qquad z\longmapsto e^{2\pi\mathrm i z},
\tag{D.22}\label{cdp:eq:D.22}
\]
is a holomorphic group isomorphism: its first lattice generator
acts trivially and its second multiplies the exponential by
$q$. Denote these oriented integral generators by $(w,u)$,
respectively.

\noindent\textbf{Step 2.} Let $C$ and $B$ be copies of $E$, with a specified product
identification $A=C\times B$. Blow up $B\times\Delta$ at
$(0,0)$ and write
\[
 b:V=\operatorname{Bl}_{(0,0)}(B\times\Delta)\longrightarrow\Delta,
 \qquad R\subset V\quad\hbox{for the exceptional curve},
 \qquad L=\OO_V(R).
\]
The notation $L^\times$ means the total space with its zero
section removed. Define the complex threefold and its map by
\[
 U=L^\times/q^{\Z}\xrightarrow{\ p\ }V\xrightarrow{\ b\ }\Delta,
\tag{D.23}\label{cdp:eq:D.23}
\]
where $q^{\Z}$ acts by fibrewise scalar multiplication. This
action is free and properly discontinuous. A local
trivialization of $L$ identifies $p$ with the product with
$\C^*/q^{\Z}=C$, so $p$ is a smooth proper principal elliptic
bundle and $U$ is smooth. Both $b$ and $p$ are proper, hence
$U\to\Delta$ is proper.

\noindent\textbf{Step 3.} The canonical section of $\OO_V(R)$ is nonzero off $R$.
It gives an exact holomorphic product identification
\[
 U|_{\Delta^*}\simeq C\times B\times\Delta^*
                         =A\times\Delta^*.
\tag{D.24}\label{cdp:eq:D.24}
\]
The central curve of $V$ is the union of the strict transform
$B_0\simeq B$ and $R\simeq\PP^1$, meeting at one point.
Their normal-crossings equation is $xy=t$. The inverse image
of $B_0$ is the principal elliptic bundle
$\OO_B(0)^\times/q^{\Z}$ of Chern degree one. Its canonical
bundle is trivial, since the vertical tangent is trivial and
the base is elliptic. The primitive Chern transgression
kills exactly one of the two fibre one-cycles, so its first
Betti number is three. This identifies it as a primary
Kodaira surface, directly by the usual definition as a
nontrivial principal elliptic bundle over an elliptic curve.
The inverse image of $R$ is
\[
 \OO_{\PP^1}(-1)^\times/q^{\Z}
       \simeq(\C^2\setminus\{0\})/\langle z\mapsto qz\rangle,
\]
a primary Hopf surface. The two surfaces meet in the single
elliptic fibre over $B_0\cap R$. Since $p$ is a submersion,
their union has local equation $xy=t$ in the smooth
threefold $U$, and both multiplicities are one.

\noindent\textbf{Step 4.} Here are the integral homology maps in the marking \eqref{cdp:eq:D.24}.
Let $\Lambda=H_1(A;\Z)$ and let $[C]=w\wedge u$ in
$H_2(A;\Z)=\bigwedge^2\Lambda$. Then inclusion of a nearby
fibre induces the exact quotients
\[
 H_1(U;\Z)=\Lambda/\Z w,\qquad
 H_2(U;\Z)=\bigwedge^2\Lambda/\Z[C].
\tag{D.25}\label{cdp:eq:D.25}
\]
To verify them, $V$ retracts onto $B_0\cup R$, which has
homotopy type $B\vee S^2$. The first Chern class of $L$
pairs to zero on a nearby $B_t$ and to $-1$ on $R$.
The transgression of the principal $C$-bundle is consequently
the primitive vector $-w$ on this sphere. In degree one it
kills exactly $w$. In degree two the integral Serre sequence
has surviving terms
\[
 E_\infty^{2,0}\simeq\Z,
 \qquad E_\infty^{1,1}\simeq\Z^4,
 \qquad E_\infty^{0,2}=0.
\]
Indeed $H^2(V;\Z)$ modulo the primitive Chern class is
$\Z$, the degree-three cohomology of $V$ is zero, and the
transgression on the exterior square of $H^1(C;\Z)$ is
injective and primitive. The analogous homology sequence,
or its integral dual with these free terms, shows that the
nearby fibre maps isomorphically onto the first two graded
terms. Its fibre class $[C]$ is zero. This proves precisely
the surjectivity and the kernel in \eqref{cdp:eq:D.25}, without a torsion
ambiguity.

\noindent\textbf{Step 5.} Let $\mu$ be the positive base meridian in the boundary
$A\times S^1$. In this same trivialization its inclusion
maps satisfy
\[
 \mu\longmapsto0\quad\hbox{in }H_1(U;\Z),
 \qquad
 x\wedge\mu\longmapsto0\quad\hbox{in }H_2(U;\Z)
 \quad(x\in\Lambda).
\tag{D.26}\label{cdp:eq:D.26}
\]
For a cycle in $C$, choose a fixed point of $B$ different
from the blowup centre and take its product with the base
disc. For a generator of $H_1(B;\Z)$, represent it by a loop
avoiding the blowup centre and take its product with the
disc. On these products the canonical section of $L$ is
nonzero, so these are actual bounding chains in the product
trivialization \eqref{cdp:eq:D.24}. The same fixed-point disc bounds
$\mu$. This proves \eqref{cdp:eq:D.26} with its specified marking.

\noindent\textbf{Step 6.} We now specify the three local fillings. Write
\[
 e_1=(1,0),\quad e_2=(\mathrm i,0),\quad
 e_3=(0,1),\quad e_4=(0,\mathrm i)
\]
for the ordered basis of $\Lambda$. The chosen elliptic
embeddings $\phi_i:E\to A$, their first lattice cycles
$w_i=\phi_i(1)$, and second cycles $u_i=\phi_i(\mathrm i)$ are
\[
\begin{array}{c|c|c|c}
 i&\phi_i(z)&w_i&u_i\\ \hline
 1&(0,z)&e_3&e_4\\
 2&(\mathrm i z,\mathrm i z)&e_2+e_4&-e_1-e_3\\
 3&(z,\mathrm i z)&e_1+e_4&e_2-e_3.
\end{array}
\tag{D.27}\label{cdp:eq:D.27}
\]
In row one take the first factor as complementary $B$; in
rows two and three take the second factor. The resulting
complex $2\times2$ change-of-coordinate matrices have
determinants $-1,\mathrm i,1$, respectively, units in
$\Z[\mathrm i]$. They therefore give exact product
isomorphisms $C_i\times B_i\simeq A$, including their
integral lattices. Apply \eqref{cdp:eq:D.23} in these three markings,
obtaining $U_1,U_2,U_3$.

\noindent\textbf{Step 7.} For the initial smooth $A$-bundle over $\PP^1$, use the
coordinate $t$ near zero and glue the trivial bundle there
to the trivial bundle over $\PP^1\setminus\{0\}$ by
\[
 a_{\rm out}=a_{\rm in}+(\ell(t),0),
 \qquad \ell(t)=\frac{\log t}{2\pi\mathrm i}\pmod{\Z+\mathrm i\Z}.
\tag{D.28}\label{cdp:eq:D.28}
\]
This is a single-valued holomorphic translation: changing
the logarithm adds an integer. It defines a smooth proper
holomorphic principal $A$-bundle $Y_0\to\PP^1$.
Its positive clutching loop has class $v_0=e_1$.
Take disjoint discs at $0,1,\infty$. On the latter two
discs the exterior trivialization in \eqref{cdp:eq:D.28} is already a
product trivialization. Remove these three fibres of $Y_0$
and glue in $U_1,U_2,U_3$ using exactly \eqref{cdp:eq:D.24} and the
original local trivializations of $Y_0$. All gluing maps
are holomorphic isomorphisms over punctured discs. The
resulting $f:Y\to\PP^1$ is a smooth complex threefold with
the claimed three central fibres. It is Hausdorff and
proper: each disc neighbourhood is the proper model $U_i$,
and away from their centres the gluing is an isomorphism
to $Y_0$. A finite such cover of the compact base proves
compactness. Connectedness of all its fibres follows
from the product fibres and the described central unions.
All transition maps over the smooth base are translations
of the marked $A$; hence their induced integral $H_1$
monodromy is the identity.

\noindent\textbf{Step 8.} We compute the global integral homology with these exact
clutchings. Let $M$ be the part above the sphere with the
interiors of the three small discs removed. It is the
product with $A$. Positive meridians $\mu_i$ about the
three deleted discs satisfy
$\mu_1+\mu_2+\mu_3=0$ in $H_1$ of that base. In the
specified local coordinates the boundary inclusion maps
to $M$ are
\[
 x\longmapsto x\quad(x\in\Lambda),\qquad
 \mu_i^{\rm loc}\longmapsto\mu_i+c_i,
 \qquad(c_1,c_2,c_3)=(e_1,0,0).
\tag{D.29}\label{cdp:eq:D.29}
\]
This sign follows directly from \eqref{cdp:eq:D.28}: a positive loop
of $t$ increases $\ell(t)$ by $+1$. On mixed
two-cycles the same map is
\[
 x\wedge\mu_i^{\rm loc}
       \longmapsto x\wedge\mu_i+x\wedge c_i.
\tag{D.30}\label{cdp:eq:D.30}
\]
This formula can also be obtained by the induced map on
$H_2$ of the product boundary torus; a translation loop
of class $c_i$ replaces its last degree-one generator
by $\mu_i+c_i$.

\noindent\textbf{Step 9.} Use the signed Mayer--Vietoris map
$(j_{M*},-j_{U_i*})$ for each boundary. Its degree-one
cokernel is
\[
 H_1(Y;\Z)=
 \Lambda/\langle e_1,w_1,w_2,w_3\rangle=0,
\tag{D.31}\label{cdp:eq:D.31}
\]
because the matrix with these four columns is
\[
 \begin{pmatrix}
 1&0&0&1\\0&0&1&0\\0&1&0&0\\0&0&1&1
 \end{pmatrix},\qquad\det=-1.
\]
There is no subsequent contribution from degree zero:
the three boundary components map injectively into the
four components consisting of $M,U_1,U_2,U_3$.
Moreover the degree-one source has rank $3(4+1)=15$
and the target has rank $(4+2)+3\cdot3=15$.
The preceding cokernel computation makes this homomorphism
surjective over $\Z$, hence an isomorphism. In particular
its kernel contributes nothing to $H_2(Y;\Z)$.

\noindent\textbf{Step 10.} By \eqref{cdp:eq:D.25}, the pure fibre two-cycles map onto each
$H_2(U_i;\Z)$ with kernel $\Z[C_i]$. Eliminate those
local targets in the degree-two Mayer--Vietoris cokernel.
By \eqref{cdp:eq:D.26}, all local mixed cycles map to zero in $U_i$.
Their exterior images \eqref{cdp:eq:D.30} eliminate the eight
generators in $\Lambda\otimes H_1$ of the exterior base,
leaving exactly the relations
$x\wedge(c_1+c_2+c_3)=x\wedge e_1$. Thus the precise
remaining quotient is
\[
 H_2(Y;\Z)=
 \frac{\bigwedge^2\Lambda}
 {e_1\wedge\Lambda+\Z[C_1]+\Z[C_2]+\Z[C_3]}.
\tag{D.32}\label{cdp:eq:D.32}
\]
Modulo $e_1\wedge\Lambda$ its three new relations are
\[
 [C_1]=e_3\wedge e_4,\qquad
 [C_2]=-e_2\wedge e_3+e_3\wedge e_4,\qquad
 [C_3]=-e_2\wedge e_4+e_3\wedge e_4.
\]
They form an integral basis of the rank-three quotient.
Consequently \eqref{cdp:eq:D.32} is zero, including its torsion.
The attached program constructs both full signed
Mayer--Vietoris matrices: their sizes are $15\times15$
and $29\times30$, the first has determinant $+1$, and
all nonzero Smith entries of both are $1$. It also checks
the six-dimensional relation matrix in \eqref{cdp:eq:D.32}.

\noindent\textbf{Step 11.} Each central component is an elliptic bundle over a
compact curve and therefore has Euler number zero; their
intersection is elliptic and also has Euler number zero.
The three central fibres have Euler number zero. The
smooth torus fibre has Euler number zero, so additivity
over the base gives $e(Y)=0$. Gauss--Bonnet and Poincar\'e
duality now give $c_3(Y)=0$ and $b_3(Y)=2$.

\noindent\textbf{Step 12.} Finally every effective divisor on $Y$ is vertical. If an
irreducible divisor dominated the base, its restriction
to a general K\"ahler torus fibre would be a nonzero
effective divisor. Its rational Chern class is zero
because $H^2(Y;\Q)=0$; integration against a K\"ahler
form gives the contradiction of a strictly positive
number equal to zero. A meromorphic function therefore
has only vertical zero and pole divisors, and its
restriction to a general compact fibre is a holomorphic
nowhere-zero function, hence constant. It consequently
descends to a meromorphic function of the base. To verify
extension at a special value with parameter $t$, choose
$N$ so that multiplication by $f^*t^N$ removes the pole
order on every component of the full fibre. The product
has no polar divisor and is holomorphic on the smooth
total space, by the local normality criterion for a
meromorphic function. Properness, connected fibres and
normality of the base give $f_*\OO_Y=\OO_{\PP^1}$ by
Stein factorization. The product therefore descends to a
holomorphic function, proving the required meromorphic
extension after division by $t^N$. Thus
$\mathcal M(Y)=f^*\C(t)$ and $a(Y)=1$.
The universal coefficient theorem also gives
$H^1(Y;\Z)=H^2(Y;\Z)=0$ from \eqref{cdp:eq:D.31}--\eqref{cdp:eq:D.32}.
This proves every assertion in \eqref{cdp:eq:D.21}.
\end{proof}



\paragraph{Literature comparison.}
The construction uses explicit principal elliptic bundles,
one ordinary point blowup in each local quotient surface,
and the specified integral clutchings. Searches for
``complex threefold, Kodaira, Hopf, fibre, monodromy'',
``algebraic dimension, $S^3\times S^3$, torus'', and
``complex torus, blow-up, principal elliptic bundle,
degeneration'' did not identify this marked example in
the sources inspected. 
The complete construction and integral calculation above
are independent of a literature classification of such
fillings. They also show why extending the group family
of \ref{cdp:st:D.2.6} does not imply extension of its marked Albanese
map on the original model: in \eqref{cdp:eq:D.25} one integral
one-cycle of the smooth torus is killed although its
linear monodromy is the identity.




\begin{proposition}[Proved coherent and integral comparisons for \ref{cdp:st:D.2.7}]\label{cdp:st:D.2.8}
\StatementRecord{cdp:st:D.2.8}{proved}{Historical events: Claude Fable 5.2 (Anthropic), effort high, 16 November 2026: carried as background without fault found. GPT-6 Astra (OpenAI), ChatGPT Pro mode, effort Pro, S6\_09\_PAPERS, 17 November 2026: targeted proof check. Also checked in this writing. Pre-continuation writing events: GPT-6 Astra (OpenAI), ChatGPT Pro mode, effort Pro, S6\_10\_PAPERS, 18 November 2026: targeted proof check. Continuation metadata events: GPT-6 Astra (OpenAI), ChatGPT Pro mode, effort Pro, S6\_10\_PAPERS continuation, 19 November 2026: used and rechecked in this run; status-vocabulary reconciliation; metadata reconciliation only; no new mathematical check. \emph{Pre-continuation recorded proof scope.} The three real annihilator lattices span the constant coherent period plane at every fibre; their actual coefficient maps extend its frame. The integral intersection is Z(1,1,0,-1), while the exact torsion locus is the image of Q4 under (a,b,c,d) maps to (a+ib,c+id)/2. Countability, density and failure of openness were checked}{the hypotheses, notation and definitions stated here; \cref{cdp:st:D.2.7}; \cref{cdp:st:D.1.11}; \cref{cdp:st:D.1.12}; \cref{cdp:eq:D.36}; \cref{cdp:st:D.1.7}}


For the exact threefold and markings of \ref{cdp:st:D.2.7},
\[
 R^0f_*\OO_Y=\OO_{\PP^1},\qquad
 R^1f_*\OO_Y=\OO_{\PP^1}^{\oplus2},\qquad
 R^2f_*\OO_Y=\OO_{\PP^1},
 \qquad h^q(Y,\OO_Y)=(1,2,1,0).
\tag{D.33}\label{cdp:eq:D.33}
\]
The coherent restriction map to every full fibre is surjective.
The integral comparison is
\[
 H^0(\PP^1,R^1f_*\Z)
   =\Z(e_1^*+e_2^*-e_4^*)
   \ \subsetneq\ H^1(A;\Z)^{\pi_1}=\Z^4.
\tag{D.34}\label{cdp:eq:D.34}
\]
The map in \eqref{cdp:eq:D.34} is restriction to the smooth base, and the
display gives its exact image. The exceptional torsion locus
for restrictions of $\Pic(Y)=H^1(Y,\OO_Y)\simeq\C^2$ to
at least one fibre is exactly
\[
 \iota\bigl(H^1(A;\Q)\bigr)\subset H^1(A,\OO_A),
 \qquad
 \iota(a,b,c,d)=\tfrac12(a+\mathrm i b,c+\mathrm i d).
\tag{D.35}\label{cdp:eq:D.35}
\]
It is countable and dense. Thus the all-fibre non-torsion locus
of \ref{cdp:st:D.1.11} need not be open, even under integral
$H^1(Y;\Z)=H^2(Y;\Z)=0$.
\end{proposition}


\begin{proof}
\noindent\textbf{Step 1.} Write $\mathscr E=R^1f_*\OO_Y$. Statement \ref{cdp:st:D.1.12} gives its
local freeness, full-fibre base change, and rank two for this
threefold. A local filling $U_i$ retracts onto its central
fibre $S_i$: the retraction of the blown-up base onto
$B_0\cup R$ lifts through its principal compact torus
bundle, by the homotopy lifting property. Therefore (D.25)
identifies
\[
 H^1(S_i;\Z)=H^1(U_i;\Z)
          =\{\lambda\in\Lambda^*:\lambda(w_i)=0\}.
\tag{D.36}\label{cdp:eq:D.36}
\]
The coefficient map into $H^1(U_i,\OO_{U_i})$
corresponds by coherent Leray on the Stein disc to
holomorphic sections of $\mathscr E$. On the punctured
disc, these are constant vectors in $H^1(A,\OO_A)$,
because (D.24) has constant marked periods. The explicit
coefficient map is $\iota$ in \eqref{cdp:eq:D.35}: the $(0,1)$ parts
of $dx,dy$ on $E$ are $d\bar z/2$ and
$\mathrm i\,d\bar z/2$.

\noindent\textbf{Step 2.} Each annihilator in \eqref{cdp:eq:D.36} has real dimension three, so
its complex span in the two-dimensional period space is
all of that space. The constant complex relation among
three chosen integral generators holds over the
punctured disc and therefore holds as an identity of
holomorphic sections over the entire disc. At zero,
their three real vectors remain independent: the map
$H^1(S_i;\R)\to H^1(S_i,\OO_{S_i})$ is real-linearly
injective by \ref{cdp:st:D.1.7}, and \ref{cdp:st:D.1.12} identifies this target
with the fibre of $\mathscr E$. Three real independent
vectors cannot lie in one complex line. Their central
complex span thus still has dimension two. Two of the
constant period vectors forming a basis consequently
extend as a frame of $\mathscr E$ across zero.

\noindent\textbf{Step 3.} These local frames agree with the same constant period
frame on the smooth base. They therefore glue to
$\mathscr E\simeq\OO^{\oplus2}$ on $\PP^1$.
The cup-product isomorphism in \ref{cdp:st:D.1.12} gives
$R^2f_*\OO_Y\simeq\det\mathscr E\simeq\OO$.
Connected fibres give $f_*\OO_Y=\OO$. Leray on $\PP^1$
proves all four dimensions in \eqref{cdp:eq:D.33}, and evaluation
of the trivial rank-two bundle, followed by the proved
base-change isomorphism, proves coherent surjectivity.

\noindent\textbf{Step 4.} The global integral sections are the intersection of
the three annihilators \eqref{cdp:eq:D.36} inside the constant
smooth-fibre lattice. Their defining equations are
$\lambda_3=0$, $\lambda_2+\lambda_4=0$,
$\lambda_1+\lambda_4=0$. The primitive solution is
$(1,1,0,-1)$, giving \eqref{cdp:eq:D.34}. In particular all three
extra fibre components are necessary in the corrected
count: $r=6$, $s=3$, and $t'=4$ give
$6=3-1+4$.

\noindent\textbf{Step 5.} Finally $\Pic(Y)=H^1(Y,\OO_Y)$ by the integral vanishing
already proved. Under its coherent restriction
identification with $H^1(A,\OO_A)$, a line bundle is
torsion on a smooth fibre exactly when some positive
integral multiple of its class belongs to
$\iota(\Lambda^*)$, by the exponential sequence.
This is exactly \eqref{cdp:eq:D.35}. At the singular fibre $S_i$
the same exponential sequence and \eqref{cdp:eq:D.36} replace
$\Lambda^*$ by its sublattice $\operatorname{ann}(w_i)$.
That torsion locus is contained in \eqref{cdp:eq:D.35}, so the union
over all fibres is exactly the displayed set. The
real-linear map $\iota:\R^4\to\C^2$ is an isomorphism;
the image of $\Q^4$ is therefore countable and dense.
Its complement is dense and has full Lebesgue measure,
but has no nonempty open subset. This proves the final
assertion and the exact meaning of the generic locus
in this example.
\end{proof}





\begin{proposition}[Proved existence and exact Euler number]\label{cdp:st:D.2.2}
\StatementRecord{cdp:st:D.2.2}{proved}{Historical events: Claude Fable 5.2 (Anthropic), effort high, 16 November 2026: carried as background without fault found. GPT-6 Astra (OpenAI), ChatGPT Pro mode, effort Pro, S6\_09\_PAPERS, 17 November 2026: targeted proof check. Also checked in this writing. Pre-continuation writing events: GPT-6 Astra (OpenAI), ChatGPT Pro mode, effort Pro, S6\_10\_PAPERS, 18 November 2026: targeted proof check. Continuation metadata events: GPT-6 Astra (OpenAI), ChatGPT Pro mode, effort Pro, S6\_10\_PAPERS continuation, 19 November 2026: used and rechecked in this run; status-vocabulary reconciliation; metadata reconciliation only; no new mathematical check. \emph{Pre-continuation recorded proof scope.} The unrestricted local nonnegativity and the independent logarithmic trivial-monodromy upper bound give zero for every individual support. The exact construction D.2.7 supplies integral homology vanishing and algebraic dimension one}{the hypotheses, notation and definitions stated here; \cref{cdp:st:C.2.G.1}; \cref{cdp:st:D.2.5}; \cref{cdp:st:D.2.7}; \cref{cdp:st:D.2.3}}

The class of compact complex threefolds with $b_1=b_2=0$,
a proper connected-fibre map with smooth two-torus general fibre,
and trivial rational monodromy is nonempty. Every member satisfies
\[
                   c_3(Y)=0,\qquad b_3(Y)=2.
\]
Every individual reduced fibre support has Euler number zero and
the geometric description of \ref{cdp:st:C.2.G.1}. Arbitrary singularities and
full-fibre multiplicities are allowed. There is an example with
integral $H_1(Y,\mathbb Z)=H_2(Y,\mathbb Z)=0$ and
algebraic dimension one.
\end{proposition}


\begin{proof}
\noindent\textbf{Step 1.} The unrestricted lower bound of \ref{cdp:st:C.2.G.1} and the independent
trivial-monodromy upper bound of \ref{cdp:st:D.2.5} give $c_3(Y)=0$.
Poincar\'e duality gives $b_3(Y)=2-c_3(Y)=2$.
\ref{cdp:st:C.2.G.1} expresses $c_3$ as the sum of the nonnegative
Euler numbers of the reduced fibre supports, so each summand
vanishes and satisfies that theorem's exact equality criterion.
Statement \ref{cdp:st:D.2.7} constructs the asserted example with integral
homology vanishing and algebraic dimension one. \ref{cdp:st:D.2.3}--4 also
give a separate local proof of the equality when all supports
have normal crossings.
\end{proof}






\section{Fibration stability and explicit deformation calculations}\label{sec:II:14}


\begin{proposition}[The infinitesimal map and the rank-two boundary]\label{cdp:st:E5}
\StatementRecord{cdp:st:E5}{proved}{Historical events: Claude Fable 5.2 (Anthropic), effort high, 16 November 2026: carried as background without fault found. GPT-6 Astra (OpenAI), ChatGPT Pro mode, effort Pro, S6\_09\_PAPERS, 17 November 2026: targeted proof check. Also checked in this writing. Pre-continuation writing events: GPT-6 Astra (OpenAI), ChatGPT Pro mode, effort Pro, S6\_10\_PAPERS, 18 November 2026: targeted proof check. Continuation mathematical and source checks: GPT-6 Astra (OpenAI), ChatGPT Pro mode, effort Pro, S6\_10\_PAPERS continuation, 19 November 2026: targeted proof check. Continuation metadata events: GPT-6 Astra (OpenAI), ChatGPT Pro mode, effort Pro, S6\_10\_PAPERS continuation, 19 November 2026: used and rechecked in this run; status-vocabulary reconciliation; metadata reconciliation only; no new mathematical check. \emph{Pre-continuation recorded proof scope.} Full degree-one Leray edge and smooth-locus tangent boundary rederived. Cup product sends a eta1+b eta2 to (b u1-a u2) eta1 wedge eta2, so its rank equals the Kodaira–Spencer rank. The image section vanishes on the rank-two open set and globally by the locally free R1O theorem E4, whose exact source statement was opened. Critical ideals (z\textasciicircum{}(m-1)),(x,y),(xy,xz,yz) and the onto logarithmic map were checked; no ordinary special-fibre surjectivity was used. Smooth source and nonsingular compact base are now explicit}{the hypotheses, notation and definitions stated here; \cref{cdp:st:E4}}

Let $g:Y\to C$ be a proper surjective holomorphic map with connected fibres
from a smooth connected compact complex threefold to a nonsingular compact
connected complex curve. Its smooth general fibre is a two-dimensional complex
torus. There is a canonical exact sequence of vector spaces
\begin{equation}\label{cdp:eq:E-Leray-map}
0\to H^1(C,T_C)\xrightarrow{g^*}H^1(Y,g^*T_C)
 \to H^0(C,T_C\otimes R^1g_*\mathcal O_Y)\to0.
\end{equation}
On the smooth locus $U\subset C$, the relative tangent sequence
induces a sheaf morphism
\[
 \delta:T_C\otimes R^1g_*\mathcal O_Y|_U
                  \longrightarrow R^2g_*T_{Y/C}|_U.
\]
At a smooth torus fibre $F$, write its Kodaira--Spencer tensor as
$\kappa=u_1\otimes\eta_1+u_2\otimes\eta_2$, where
$\eta_1,\eta_2$ is a basis of $H^1(F,\mathcal O_F)$ and
$u_i\in\operatorname{Lie}(F)$. Then
\[
 \delta(a\eta_1+b\eta_2)
   =(bu_1-au_2)\otimes(\eta_1\wedge\eta_2).
\]
Its rank is the rank of $\kappa$. When this rank is two at one
regular value, the image of $H^1(Y,T_Y)\to H^1(Y,g^*T_C)$ lies
in $g^*H^1(C,T_C)$. For $C=\mathbb P^1$ this image is zero.
\end{proposition}


\begin{proof}
\noindent\textbf{Step 1.} The projection formula gives
$R^qg_*g^*T_C=T_C\otimes R^qg_*\mathcal O_Y$, and the coherent
Leray spectral sequence has only columns zero and one. Its
degree-one edge sequence is exactly \eqref{cdp:eq:E-Leray-map}.
On $U$ the sequence
$0\to T_{Y/C}\to T_Y\to g^*T_C\to0$ is exact. Its extension
class on $F$ is the Kodaira--Spencer tensor. The connecting map is
cup product with that class; applying alternating multiplication to
the two displayed basis vectors gives the formula and its rank.
Exactness of the higher-direct-image sequence says that the image
of a global tangent class restricts to the kernel of $\delta$.
Rank two holds on an open dense subset after it holds at one point.
By \ref{cdp:st:E4} the source $T_C\otimes R^1g_*\mathcal O_Y$ is locally free,
so a section zero on that subset is zero everywhere. This proves
the image assertion and the final statement.

\noindent\textbf{Step 2.} At a critical point the sheaf map $dg$ need not be onto. Its correct
image is $\mathcal I_{\rm crit}\otimes g^*T_C$, giving
\[
0\to T_{Y/C}\to T_Y\to
 \mathcal I_{\rm crit}\otimes g^*T_C\to0.
\]
For $t=z^m$, the image ideal is $(z^{m-1})$; for $t=xy$ it is
$(x,y)$; for $t=xyz$ it is $(xy,xz,yz)$. On a root cover of a
smooth multiple fibre the ordinary tangent sequence for $z$ applies.
At a normal-crossing fibre the logarithmic map sends
$\sum a_i z_i\partial_{z_i}$ to
$(\sum a_i)t\partial_t$ and is onto the logarithmic base tangent
line. These formulas type the comparison at the special fibres.
The proof above only uses the genuine smooth-locus boundary and
the globally locally free source, so it makes no assumption about
local freeness of $R^2g_*T_{Y/C}$ at a special value.
\end{proof}






\begin{proposition}[Stability of a rank-two torus fibration]\label{cdp:st:E6}
\StatementRecord{cdp:st:E6}{proved}{Historical events: Claude Fable 5.2 (Anthropic), effort high, 16 November 2026: carried as background without fault found. Claude Fable 5.2 (Anthropic), effort high, 17 November 2026: carried as background without fault found. GPT-6 Astra (OpenAI), ChatGPT Pro mode, effort Pro, S6\_09\_PAPERS, 17 November 2026: targeted proof check. Also checked in this writing. Pre-continuation writing events: GPT-6 Astra (OpenAI), ChatGPT Pro mode, effort Pro, S6\_10\_PAPERS, 18 November 2026: targeted proof check; used and rechecked in this run. Continuation metadata events: GPT-6 Astra (OpenAI), ChatGPT Pro mode, effort Pro, S6\_10\_PAPERS continuation, 19 November 2026: used and rechecked in this run; status-vocabulary reconciliation; metadata reconciliation only; no new mathematical check. \emph{Pre-continuation recorded proof scope.} The small-Artin-extension embedding lift, restriction of the map obstruction, dense-fibre vanishing, and nonreduced Douady graph construction were re-derived. Douady sections 9-10, Horikawa II section 4 and Theorem 6.1, and Buchdahl Theorem 11 were opened at their exact loci. The analytic source-variation extension is supplied in the displayed proof}{the hypotheses, notation and definitions stated here; \cref{cdp:st:E5}; \cite[\S4, equations (4.3)--(4.5), and \S6]{Horikawa1974}; \cref{cdp:eq:E-surj-embedding}; \cref{cdp:eq:E-Leray-map}; \cite[Section~9, Theorems~1--2, pp.~83--85]{Douady1966}; \cite[Section~10.1, Proposition~1, p.~86]{Douady1966}; \cite[Section~10.2, Theorem~1, p.~87]{Douady1966}; \cite[Theorem 11]{Buchdahl1999}; \cref{cdp:st:E1}}

Let $g:Y\to\mathbb P^1$ be as in \ref{cdp:st:E5} and suppose its
Kodaira--Spencer tensor has rank two at one regular value. For every
proper deformation $\mathcal Y\to(S,0)$ of $Y$ over a complex
analytic germ, including a nonreduced germ, there is, after shrinking
$S$, a holomorphic map
\[
 G:\mathcal Y\longrightarrow\mathbb P^1\times S
\]
over $S$ whose central restriction is $g$. Every $G_s$ is surjective
with connected fibres and smooth general fibre a complex two-torus.
With $b_2(Y)=0$, every sufficiently small deformation has algebraic
dimension one.
\end{proposition}


\begin{proof}
\noindent\textbf{Step 1.} We prove the lifting assertion at every Artin order; this is more
than the tangent-space assertion of \ref{cdp:st:E5}. The obstruction calculations
below use the Cech construction of holomorphic-map deformations
in \cite[\S4, equations (4.3)--(4.5), and \S6]{Horikawa1974};
we spell out the portions needed here. The injectivity hypothesis
on degree-two tangent cohomology in its Theorem~6.1 is replaced
here by the explicitly proved unobstructedness of the torus source.

\noindent\textbf{Step 2.} First fix a rank-two regular fibre and its embedding $i:F\hookrightarrow
Y$. The normal line is canonically $N=\mathcal O_F\otimes T_c\mathbb
P^1$. The tangent sequence and \ref{cdp:st:E5} imply that
\begin{equation}\label{cdp:eq:E-surj-embedding}
 H^1(F,T_F)\longrightarrow H^1(F,i^*T_Y)
 \quad\hbox{is surjective}.
\end{equation}
Let $A'\to A$ be a small extension of local Artin complex algebras,
with kernel $I$ annihilated by the maximal ideal. Suppose an embedding
$i_A:F_A\hookrightarrow Y_A$ deforming $i$ and a deformation
$Y_{A'}$ lifting $Y_A$ have been specified. The deformation functor
of a complex torus is smooth: with a real lattice marking its small
complex structures are graphs of arbitrary small complex linear maps
between the two complementary Hodge subspaces; constant Beltrami
tensors have zero bracket. This gives a smooth four-dimensional
miniversal germ for $F$. In particular $F_A$ has a lift $F_{A'}$.

\noindent\textbf{Step 3.} Choose local lifts of $i_A$ on sufficiently small coordinate charts.
They exist because the target deformation is smooth over $A'$.
Their differences on overlaps form a one-cocycle in
$i^*T_Y\otimes I$. Changing the lifted transition functions of
$F_{A'}$ by a cocycle $v_{ij}\in T_F\otimes I$ changes this
obstruction by $di(v_{ij})$. Such changes are precisely the action
of $H^1(F,T_F)\otimes I$ on the set of lifts of $F_A$; their
cocycle law has no quadratic term because $I\mathfrak m_{A'}=0$.
Surjectivity \eqref{cdp:eq:E-surj-embedding} therefore makes the map
obstruction zero. A zero-coboundary correction glues the local maps.
The resulting proper map is a closed immersion: the map on local
rings is surjective modulo the nilpotent ideal, hence surjective
by Nakayama, and its underlying central map is the embedding $i$.
Thus the given embedded torus lifts in every ambient small extension.

\noindent\textbf{Step 4.} Now suppose $g_A:Y_A\to\mathbb P^1_A$ has been lifted and that
$Y_{A'}$ is prescribed. Choose local lifts $t_i$ of the target
coordinate of $g_A$; on overlaps their discrepancies define the
holomorphic-map obstruction
\[
 \alpha\in H^1(Y,g^*T_{\mathbb P^1})\otimes I.
\]
Coordinate changes on the target give exactly the transition factors
of $g^*T_{\mathbb P^1}$, and the triple-overlap relation gives its
Cech cocycle law. A coboundary makes the local maps agree, so
$\alpha=0$ is sufficient as well as necessary.

\noindent\textbf{Step 5.} For a rank-two regular value $c$, the divisor
$F_A=g_A^{-1}(c)$ is a smooth deformation of the central torus.
The preceding embedding argument provides a Cartier divisor
$F_{A'}\subset Y_{A'}$ lifting it. This forces the restriction
of $\alpha$ to $F$ to vanish, as can be checked without a comparison
of obstruction theories: choose local equations $t_i-\eta_i$ for
$F_{A'}$, where $\eta_i$ lies in $I\otimes\mathcal O_Y$.
Two such equations differ by a unit. Modulo the central equation
$t=0$, that unit difference disappears, and their overlap equality
reads
\[
 (t_i-t_j)|_F=\eta_i|_F-\eta_j|_F.
\]
This is exactly the assertion that the restricted obstruction
cocycle is a coboundary in $N\otimes I$.

\noindent\textbf{Step 6.} By \eqref{cdp:eq:E-Leray-map} and $H^1(\mathbb P^1,T_{\mathbb P^1})=0$,
$\alpha$ is a section of the locally free sheaf
$T_{\mathbb P^1}\otimes R^1g_*\mathcal O_Y$, tensored with $I$.
Its evaluation at every rank-two regular fibre is zero. That set
is open dense; hence the section and $\alpha$ are zero. The map
therefore lifts across every small extension.

\noindent\textbf{Step 7.} We pass from the formal assertion to an analytic map and
specify the relative parameter space, including its nilpotents.
Let $H$ be the germ at the graph of $g$ in the absolute Douady
space of $\mathcal Y\times\mathbb P^1$, and let
$\pi:\mathcal R\to H$ be its universal compact flat family.
Douady's existence and finiteness theorems
\cite[Section~9, Theorems~1--2, pp.~83--85]{Douady1966}
apply to this analytic ambient space; its local rings and those
of the parameter space may have nilpotents.

\noindent\textbf{Step 8.} After shrinking $H$, one has
\[
 \pi_*\mathcal O_{\mathcal R}=\mathcal O_H,
\]
compatibly with base change. To check this without assuming a
reduced parameter space, represent $R\pi_*\mathcal O_{\mathcal R}$
near the central point by a finite complex of finite free analytic
modules. On the central fibre the unit spans
$H^0(Y,\mathcal O_Y)=\mathbb C$. Split off this unit in degree
zero. The differential is injective modulo the maximal ideal on
the complementary degree-zero summand; an invertible minor
cancels that summand with its image in degree one. The resulting
complex begins with $\mathcal O_H$ in degree zero and zero
differential from it. This proves the displayed equality and its
base-change assertion on the full analytic germ.

\noindent\textbf{Step 9.} Embed a representative of $(S,0)$ in a polydisc. The coordinate
functions of the map $\mathcal R\to\mathcal Y\to S$ are now
pullbacks of unique holomorphic functions on $H$. They satisfy
the defining equations of $S$, because their pullbacks do and the
unit map just proved is injective. They therefore define a
holomorphic morphism $H\to S$, and the universal subspace lies
scheme-theoretically in
\[
 (\mathcal Y\times_S H)\times\mathbb P^1.
\]
The projection of this subspace to $\mathcal Y\times_S H$ is a
proper morphism between two proper flat families over $H$ and
is an isomorphism at the central point. By the isomorphism-locus
criterion \cite[Section~10.1, Proposition~1, p.~86]{Douady1966},
it is an isomorphism after restricting to an open neighbourhood
of that point. The proof of that criterion uses the finite
direct image, flatness and Nakayama's lemma, and applies to
nonreduced $H$. On this open locus the universal subspace is
the graph of an actual map to $\mathbb P^1$.

\noindent\textbf{Step 10.} The universal property of the absolute Douady space consequently
identifies this open locus with the germ $\mathcal H\to S$ of
relative maps from $\mathcal Y/S$ to
$\mathbb P^1\times S/S$. This is the relative version of
Douady's graph construction
\cite[Section~10.2, Theorem~1, p.~87]{Douady1966}; the argument
above supplies the variation of the compact source explicitly.
The just-proved lifting property for every
small Artin extension says that this morphism is formally smooth
at $[g]$. The local analytic smoothness criterion identifies a
formally smooth morphism of finite-dimensional analytic germs with
a projection $S\times(\mathbb C^r,0)\to S$ after shrinking. One
can see the criterion on presentations of their convergent local
algebras: lifting across all square-zero quotients makes the
relative conormal obstruction zero, and the analytic implicit
function theorem eliminates the corresponding equations. The
zero section of this product supplies the required holomorphic
family of graphs. This argument applies to nonreduced $S$.

\noindent\textbf{Step 11.} At one regular value, the family $G$ remains a proper submersion
near the central compact fibre. Its nearby fibres are connected
deformations of a torus. Each map $G_s$ is nonconstant, and hence
surjective. Its Stein factorization has finite degree one because
it has a connected regular fibre in this neighbourhood. Thus
every fibre is connected. Over the connected smooth locus of
$G_s$, Ehresmann's theorem identifies the differentiable fibres
with $T^4$. Every compact complex surface with this differentiable
type is a torus: its even first Betti number makes it K\"ahler
by \cite[Theorem 11]{Buchdahl1999};
its Albanese has dimension two; the exterior algebra on its four
integral degree-one classes shows the Albanese map has degree one;
and a nontrivial exceptional curve of this degree-one map would
include a $(-1)$-curve, impossible for the even intersection form
of $T^4$. Hence the general fibre is a complex two-torus.
Finally $b_2$ is constant in the differentiable deformation family,
and \ref{cdp:st:E1} proves $a(Y_s)=1$.
\end{proof}




\begin{proposition}[Application to the named six-sphere families]\label{cdp:st:E7}
\StatementRecord{cdp:st:E7}{proved}{Historical events: Claude Fable 5.2 (Anthropic), effort high, 16 November 2026: carried as background without fault found. GPT-6 Astra (OpenAI), ChatGPT Pro mode, effort Pro, S6\_09\_PAPERS, 17 November 2026: targeted proof check. Also checked in this writing. Pre-continuation writing events: GPT-6 Astra (OpenAI), ChatGPT Pro mode, effort Pro, S6\_10\_PAPERS, 18 November 2026: targeted proof check; used and rechecked in this run. Continuation metadata events: GPT-6 Astra (OpenAI), ChatGPT Pro mode, effort Pro, S6\_10\_PAPERS continuation, 19 November 2026: used and rechecked in this run; status-vocabulary reconciliation; metadata reconciliation only; no new mathematical check. \emph{Pre-continuation recorded proof scope.} The invertible logarithmic period derivative was re-derived. The actual rank-two local models were opened in the comb paper (periodextension and hexfill), two-parameter paper (mum), and Liu-Rollenske paper (compactification X). The application uses E6 and the stated Paper I period inputs}{the hypotheses, notation and definitions stated here; \cref{cdp:st:E6}; \cite[Proposition labelled \texttt{hexfill}]{Viaclovsky2026Combs}; \cite[Proposition labelled \texttt{mum}]{Viaclovsky2026Two}; \cite[Proposition labelled \texttt{prop: compactification X}]{LiuRollenske2026}; \cite[Proposition labelled \texttt{periodextension},
case (i)]{Viaclovsky2026Combs}; \cref{cdp:eq:E-log-period}; \cref{cdp:st:E1}}

Every member of the following torus-fibred families,
including the period family $X_c$, its members with $|p|>1$,
Viaclovsky's $X_{3k,4l}$, his two-parameter family and his combs,
and the Liu--Rollenske family, satisfies the rank hypothesis of \ref{cdp:st:E6}.
Consequently every sufficiently small deformation retains a
holomorphic torus fibration over $\mathbb P^1$. Every such member
with $b_2=0$, and every sufficiently small deformation of it, has
algebraic dimension one. In particular the fibration-retention
question in Section~10(1) of the supplied comb paper has an
affirmative answer.
\end{proposition}


\begin{proof}
\noindent\textbf{Step 1.} We verify a local condition common to these actual constructions.
At their hexagonal fibre, a marking on the punctured disc puts
the Hodge periods in the form
\begin{equation}\label{cdp:eq:E-log-period}
  (I_2,Z(q)),\qquad
 Z(q)=\frac{\log q}{2\pi i}B+H(q),
 \qquad \det B\ne0,
\end{equation}
where $H$ is holomorphic at zero. Therefore
\[
 qZ'(q)=\frac{B}{2\pi i}+qH'(q)
\]
is invertible for all sufficiently small nonzero $q$. The
Kodaira--Spencer tensor of a marked torus with periods $(I,Z)$
is represented by $Z'$ after the invertible identifications of
the Lie algebra and $H^1(\mathcal O)$ with the chosen period
coordinates. Thus it has rank two at those regular fibres.

\noindent\textbf{Step 2.} Here the existence and nonsingularity of $B$ are checked directly
in the supplied local models, including
\cite[Proposition labelled \texttt{hexfill}]{Viaclovsky2026Combs},
\cite[Proposition labelled \texttt{mum}]{Viaclovsky2026Two}, and
\cite[Proposition labelled \texttt{prop: compactification X}]{LiuRollenske2026}.
The AC cusp has logarithmic matrix
$B_0=\left(\begin{smallmatrix}0&1\\-1&0\end{smallmatrix}\right)$.
The affine twists and finite logarithmic transformations change
the finite collars and constant gluing data, while the punctured
cusp matrix still has this same invertible logarithmic coefficient.
Engel's two deck transformations at the $I_1$-fibre have independent
valuation vectors $(1,0)$ and $(0,1)$, again giving an invertible
$B$. In the two-parameter family the designated fibre over $p_3$
is the full-lattice rank-two filling; its displayed two valuation
vectors are independent. Liu--Rollenske's compactification over
$I_1$ uses the same two independent valuation directions. Each
comparison is a change of the integral marking and of the
holomorphic fibre frame, so preserves the rank of the
Kodaira--Spencer tensor.

\noindent\textbf{Step 3.} For the combs, \cite[Proposition labelled \texttt{periodextension},
case (i)]{Viaclovsky2026Combs} writes a flat frame as
$e_i=\alpha_i+\ell(q)\beta_i$, $f_i=\beta_i$, with
$\ell(q)=\log q/(2\pi i)$, and the extended Hodge subbundle as
the rows of $(I,H(q))$ in this frame. The fixed value of the
smoothed period section at $q=0$ keeps it transverse to the
$f_i$-plane. Its rows therefore remain $(I,H_s(q))$ with
$H_s$ holomorphic. In the flat frame they are
$(I,\ell(q)I+H_s(q))$, which is \eqref{cdp:eq:E-log-period} with
$B=I$. This verifies the needed condition for every comb member,
without a formula for its global splitting type.

\noindent\textbf{Step 4.} All the sources construct the indicated members as compact
connected-fibre holomorphic maps with smooth general two-torus
fibre. \ref{cdp:st:E6} applies, and \ref{cdp:st:E1} gives the algebraic-dimension assertion.
On the differentiable six-sphere these neighbourhoods of the
listed structures constitute an open set on which $a=1$.
\end{proof}






\begin{proposition}[The theta determinant and its exact rank locus for $X_c$]\label{cdp:st:E8}
\StatementRecord{cdp:st:E8}{proved}{Historical events: Claude Fable 5.2 (Anthropic), effort high, 16 November 2026: carried as background without fault found. GPT-6 Astra (OpenAI), ChatGPT Pro mode, effort Pro, S6\_09\_PAPERS, 17 November 2026: targeted proof check. Also checked in this writing. Pre-continuation writing events: GPT-6 Astra (OpenAI), ChatGPT Pro mode, effort Pro, S6\_10\_PAPERS, 18 November 2026: targeted proof check; used and rechecked in this run. Continuation mathematical and source checks: GPT-6 Astra (OpenAI), ChatGPT Pro mode, effort Pro, S6\_10\_PAPERS continuation, 19 November 2026: targeted proof check. Continuation metadata events: GPT-6 Astra (OpenAI), ChatGPT Pro mode, effort Pro, S6\_10\_PAPERS continuation, 19 November 2026: used and rechecked in this run; status-vocabulary reconciliation; metadata reconciliation only; no new mathematical check. \emph{Pre-continuation recorded proof scope.} Full theta determinant proof rederived. Adjoining(mu,beta/6) gives an exact index-six lattice enlargement and the identity-induced degree-six isogeny, preserving Kodaira–Spencer rank. Heat equation, eta derivative and h\_mu=-wp-pi\textasciicircum{}2E2/3 cancel the E2 terms; differentiating wp(mu)=0 gives determinant -g2\textasciicircum{}2/(36pi\textasciicircum{}2g3). Checked root orders1 and2, the resulting determinant order2 at the order-three point, cancellation at the order-four point, and the nondegenerate cusp logarithmic diagonal}{the hypotheses, notation and definitions stated here; the explicit constructions and calculations in the proof}

Use the period convention
$\Omega=\left(\begin{smallmatrix}\tau&\mu\\\mu&\beta/6\end{smallmatrix}\right)$,
Jacobi's odd theta function with periods $1,\tau$, and
$\beta/6=(i/\pi)\log(\vartheta_1(\mu,\tau)/\eta(\tau))+\text{constant}$.
Writing $h=\partial_\mu\log\vartheta_1$, one has the identity
\begin{equation}\label{cdp:eq:E-theta-general}
 \det\frac{d\Omega}{d\tau}
 =-\left(\mu'-\frac{i}{2\pi}h\right)^2
       -\frac{\wp(\mu,\tau)}{4\pi^2}.
\end{equation}
For the actual AC period branch $\wp(\mu,\tau)=0$, this becomes
\[
 \det\frac{d\Omega}{d\tau}
       =-\frac{g_2(\tau)^2}{36\pi^2g_3(\tau)}.
\]
On the smooth root-family base, the Kodaira--Spencer tensor has rank
one at the order-three point and rank two at every other finite
point, including the order-four point. Its determinant has order
two at the order-three point. The logarithmic tensor at the cusp
has rank two.
\end{proposition}


\begin{proof}
\noindent\textbf{Step 1.} The symmetric matrix has a concrete comparison with the raw AC
periods. Adjoin $(\mu,\beta/6)$ to the lattice generated by
$(1,0),(0,1),(\tau,\mu),(6\mu,\beta)$. This enlarges the
lattice by index six; identity on $\mathbb C^2$ induces a finite
holomorphic isogeny of degree six to the torus with period matrix
$(I,\Omega)$. Its differential and pullback on complex first
cohomology are isomorphisms, so the Kodaira--Spencer ranks agree.
This is the comparison used in the determinant calculation.

\noindent\textbf{Step 2.} The convention is
$\vartheta_j(z,\tau)=\theta_j(\pi z\mid\tau)$ and
$q=e^{\pi i\tau}$. The normally convergent theta series and
\cite[eqs.~20.2.1 and 20.13.1]{DLMFTheta2026} give
$\partial_\tau\vartheta_1=(4\pi i)^{-1}\partial_\mu^2\vartheta_1$.
Put $Q=q^2$ and
$E_2=1-24\sum_{n\ge1}nQ^n/(1-Q^n)$.
The normally convergent product
\cite[eq.~23.17.8]{DLMFTheta2026} gives
$\eta'/\eta=\pi i E_2/12$ by termwise logarithmic differentiation.
Likewise \cite[eq.~20.5.1]{DLMFTheta2026} gives
\[
 h(\mu)=\frac1\mu-\frac{\pi^2E_2}{3}\mu+O(\mu^3).
\]
Here the coefficient follows from
$\sum_{n\ge1}Q^n/(1-Q^n)^2
=\sum_{n\ge1}nQ^n/(1-Q^n)$, obtained by comparing the absolutely
convergent double series. The quasi-period equation
\cite[eq.~20.2.6]{DLMFTheta2026} makes $\partial_\mu h$ elliptic.
Consequently $-\partial_\mu h-\pi^2E_2/3$ and $\wp$
have identical poles, principal parts and zero Laurent constant.
Their difference is a holomorphic elliptic function with constant
zero, so they agree. Substituting these three
identities into
$\det\Omega'=(\beta/6)'-\mu'^2$ gives
\eqref{cdp:eq:E-theta-general}. Thus the square identity in Lead~9
uses the equation $\wp(\mu)=0$; a general section requires the
displayed additional term unless its other period terms cancel it.

\noindent\textbf{Step 3.} For the zero branch differentiate $\wp(\mu,\tau)=0$. Put
$c=\pi^2E_2/3$ so $h'=-\wp-c$. Differentiating the heat equation
once in $\mu$, using $\wp''=6\wp^2-g_2/2$, and using
$E_2'=(\pi i/6)(E_2^2-E_4)$, yields at $\wp=0$
\[
 \partial_\tau\wp=\frac{h\wp'}{2\pi i}
                       +\frac{i g_2}{6\pi}.
\]
Consequently
\[
 \mu'-\frac{i h}{2\pi}=-\frac{i g_2}{6\pi\wp'},
 \qquad (\wp')^2=-g_3,
\]
which proves the second formula wherever $g_3\ne0$. It is an
identity of meromorphic expressions and therefore extends across
the branch values in a local uniformizing parameter.

\noindent\textbf{Step 4.} The equation $j(\tau)=1728t$ shows that $\tau-\rho$ has order one
at the order-three root parameter, while $\tau-i$ has order two
at the order-four root parameter. The modular form $g_2$ has a
simple zero at $\rho$ and $g_3$ a simple zero at $i$, with no
other zeros modulo the modular group. Multiplication by
$(d\tau/ds)^2$ therefore gives an order-two zero at the first
point, and cancels the order-two zero of $g_3(\tau(s))$ at the
second. The tensor is nonzero at the first point because
$d\tau/ds\ne0$ there, hence has rank exactly one. Away from the
two root points the coordinate change is unramified and the
formula has nonzero determinant. At the cusp
$\tau=(\log q)/(2\pi i)+O(q)$, $\mu$ has a finite limit, and
$\beta=-\tau+O(1)$; hence the logarithmic derivative has limiting
matrix $\operatorname{diag}(1/(2\pi i),-1/(12\pi i))$.
Its determinant is nonzero.
\end{proof}






\begin{proposition}[Counterexample: A fixed positive-genus target in the unrestricted stability assertion]\label{cdp:st:E10}
\StatementRecord{cdp:st:E10}{proved}{Historical events: Claude Fable 5.2 (Anthropic), effort high, 16 November 2026: carried as background without fault found. Claude Fable 5.2 (Anthropic), effort high, 17 November 2026: carried as background without fault found. GPT-6 Astra (OpenAI), ChatGPT Pro mode, effort Pro, S6\_09\_PAPERS, 17 November 2026: targeted proof check. Also checked in this writing. Pre-continuation writing events: GPT-6 Astra (OpenAI), ChatGPT Pro mode, effort Pro, S6\_10\_PAPERS, 18 November 2026: targeted proof check. Continuation metadata events: GPT-6 Astra (OpenAI), ChatGPT Pro mode, effort Pro, S6\_10\_PAPERS continuation, 19 November 2026: used and rechecked in this run; status-vocabulary reconciliation; metadata reconciliation only; no new mathematical check. \emph{Pre-continuation recorded proof scope.} Full explicit counterexample rederived: the two Hesse critical-value sets are disjoint; the fibre product and moving genus-two double-cover base change are smooth and proper. Fixing five branch points forces the infinitesimal PGL2 field to vanish while the sixth moves, so theta\_C is nonzero. Relative duality/Kunneth and genus-two Riemann–Roch give the three printed cohomology rows. The exact class dg(theta\_Y)=g*theta\_C is nonzero by the Leray injection and obstructs fixed-target extension. Status is proved counterexample; the rejected unrestricted assertion remains false}{the hypotheses, notation and definitions stated here; \cref{cdp:eq:E-Leray-map}; \cref{cdp:st:E5}}

The assertion that every small deformation of a fibration extends
over its original fixed positive-genus target curve $C$ is false even for a torus
fibration with Kodaira--Spencer rank two at a regular value.
\end{proposition}


\begin{proof}
\noindent\textbf{Step 1.} Let $S\to\mathbb P^1_t$ be the rational elliptic surface obtained
by blowing up the nine base points of the Hesse pencil
$x^3+y^3+z^3-3txyz=0$. Its critical values are
$\{1,\omega,\omega^2,\infty\}$, where $\omega^3=1$, $\omega\ne1$.
Pull another copy back by $\phi(t)=1/(t-4)$; its critical values
are $\{4,5,4+\omega,4+\omega^2\}$, disjoint from the first set.
Their fibre product $Z\to\mathbb P^1$ is a smooth projective
threefold, since at every point at least one elliptic-surface
map is a submersion. Its general fibre is a product of two
elliptic curves. Both elliptic $j$-functions are nonconstant,
so their two Kodaira--Spencer derivatives are simultaneously
nonzero off a finite set. The tensor of their product is then
diagonal of rank two.

\noindent\textbf{Step 2.} Let $C_s$ be the smooth genus-two double cover
\[
 y^2=(t-10)(t-11)(t-12)(t-13)(t-14)(t-15-s)
\]
of $\mathbb P^1$, for $s$ in a small disc. All its branch values
avoid the critical values of $Z$. Thus
$Y_s=Z\times_{\mathbb P^1}C_s$ is a smooth projective deformation
of a compact threefold, and $g_s:Y_s\to C_s$ has connected torus
general fibre and rank two at a regular value. The Kodaira--Spencer
class $\theta_C\in H^1(C_0,T_{C_0})$ is nonzero: any infinitesimal
isomorphism of the hyperelliptic covers descends through their
unique hyperelliptic involutions to an infinitesimal projective
transformation fixing five distinct branch values; it is zero,
whereas the sixth branch value moves with derivative one.

\noindent\textbf{Step 3.} All the cohomology used for the obstruction is explicit. Put
$L=\pi_C^*\mathcal O_{\mathbb P^1}(1)$ on $C_0$; the
double-cover canonical formula gives $K_{C_0}=L$, of degree two.
For either rational elliptic surface, relative duality gives
$R^0\pi_*\mathcal O=\mathcal O$ and
$R^1\pi_*\mathcal O=\mathcal O(-1)$: its relative canonical
line is $\pi^*\mathcal O(1)$ and every fibre has
$h^\bullet(\mathcal O)=(1,1)$. Relative Kunneth and the flat
base change to $C_0$ consequently give
\[
 R^qg_*\mathcal O_Y=
 \begin{cases}
  \mathcal O_{C_0}&q=0,\\
  L^{-1}\oplus L^{-1}&q=1,\\
  L^{-2}&q=2,\\
  0&q>2.
 \end{cases}
\]
The vertical tangent sheaf is $g^*(L^{-1}\oplus L^{-1})$:
the invariant vector on each elliptic factor generates its
relative tangent sheaf, also at a node, where the kernel of
$d(xy)$ is generated by $x\partial_x-y\partial_y$.
Riemann--Roch on the genus-two curve gives
$h^0(L^{-a})=0$ and $h^1(L^{-a})=2a+1$ for $a>0$.
Projection and Leray now compute all degrees:
\[
\begin{array}{c|rrrr}
 \mathcal F&h^0&h^1&h^2&h^3\\\hline
 \mathcal O_Y&1&2&6&5\\
 g^*T_{C_0}&0&3&10&7\\
 T_{Y/C_0}&0&6&20&14
\end{array}
\]
All higher groups vanish. In particular
$H^0(C_0,T_{C_0}\otimes R^1g_*\mathcal O_Y)=0$ and the
injection $H^1(C_0,T_{C_0})\to H^1(Y,g^*T_{C_0})$ is an
isomorphism between three-dimensional spaces.

\noindent\textbf{Step 4.} Compatibility of the Cech transition equations for the family
$g_s$ gives
\[
 dg(\theta_Y)=g^*\theta_C
       \quad\text{in }H^1(Y_0,g^*T_{C_0}).
\]
The right side is nonzero because the first map of
\eqref{cdp:eq:E-Leray-map} is injective. Hence the first-order
obstruction to extending $g_0$ with fixed target $C_0$ is nonzero.
This proves the counterexample with an actual smooth proper
family and its obstruction class. The obstruction belongs to
the base-deformation summand specified exactly in \ref{cdp:st:E5}.
\end{proof}






\begin{proposition}[All coprime $(3k,4l)$ multiplicities]\label{cdp:st:E11}
\StatementRecord{cdp:st:E11}{proved}{Historical events: Claude Fable 5.2 (Anthropic), effort high, 16 November 2026: carried as background without fault found. GPT-6 Astra (OpenAI), ChatGPT Pro mode, effort Pro, S6\_09\_PAPERS, 17 November 2026: targeted proof check. Also checked in this writing. Pre-continuation writing events: GPT-6 Astra (OpenAI), ChatGPT Pro mode, effort Pro, S6\_10\_PAPERS, 18 November 2026: targeted proof check. Continuation metadata events: GPT-6 Astra (OpenAI), ChatGPT Pro mode, effort Pro, S6\_10\_PAPERS continuation, 19 November 2026: used and rechecked in this run; status-vocabulary reconciliation; metadata reconciliation only; no new mathematical check. \emph{Pre-continuation recorded proof scope.} Full coprime-multiplicity tangent/Kuranishi argument rederived. The actual root morphism has sigma*O(1)=O(kl), its finite invariant pushforward is exact, and fixed affine translations preserve the vertical/H1O representations. Checked horizontal root tangent versus pulled vertical line, all low-degree vanishings, the actual horizontal-one-form extension, and weighted-monomial contradiction3k(a+2)+4l(b+2)=kl. HRR on each verified six-sphere gives h2=1; the original affine case has its direct acyclic log calculation. Fixed three critical sections kill every infinitesimal base projective motion; the nonzero scale class then forces the Kuranishi local-ring ideal to be zero. All stale PaperI Theorem3.9 references in this block are replaced by the exact canonical fam:src:A.5.2}{the hypotheses, notation and definitions stated here; \cite[Theorem labelled \texttt{t:main}]{Viaclovsky2026Log}; \cref{cdp:st:E7}}

Fix admissible logarithmic-transformation data with multiplicities
$r=3k$, $s=4l$, where $k,l>0$ and $\gcd(r,s)=1$, in the family
constructed from the original elliptic-surface data. For every member
which is diffeomorphic to $S^6$, as in
\cite[Theorem labelled \texttt{t:main}]{Viaclovsky2026Log},
\[
 h^\bullet(X,T_X)=(1,1,1,0),
\]
its Kuranishi germ is a smooth reduced line, and varying the nonzero
linearization scale with those data fixed is locally miniversal.
For the original $(3,4)$ affine-twist family, the same statements
hold whenever $p\ne0$, including $|p|>1$.
\end{proposition}


\begin{proof}
\noindent\textbf{Step 1.} Put $\mathfrak B_{r,s}=\mathbb P(r,s)$. The root morphism
$\sigma:\mathfrak B_{r,s}\to\mathfrak B_{3,4}$ over the identical
coarse $\mathbb P^1$ sends the two root sections to their $k$th and
$l$th powers. In line notation
$\sigma^*\mathcal O_{\mathfrak B_{3,4}}(1)
=\mathcal O_{\mathfrak B_{r,s}}(kl)$.
The vertical transition matrices and the action on fibre
$\mathcal O$-cohomology depend only on the linear monodromy.
The extra affine translations act trivially on both. The same
cusp comparison used in \cref{fam:src:A.5.2} therefore gives the actual bundles
\[
 E=G=\mathcal O\oplus\mathcal O(-kl),\qquad
 R^2\bar f_*\mathcal O_X=\mathcal O(-kl),
\]
and
\[
0\to\bar f^*E\to\mathcal T:=T_X(-\log W)
 \to\bar f^*L\to0,
\quad L=\mathcal O(r+s-rs),
\quad K_X=\bar f^*\mathcal O(kl-r-s).
\]
For example the normal-cover coordinate at the first finite point
is $u=v^k$, and differentiation changes the ordinary root tangent
line from the old one to $T_{\mathfrak B_{r,s}}$; the vertical
bundle and its cohomology are pulled back through $u=v^k$.
This specifies the comparison rather than identifying the two
horizontal tangent lines. At the cusp $W$ remains the same
normal-crossing divisor, so its logarithmic base correction is
$-p_0=-rs$ and its normal quotient is $\eta_*K_{\mathrm dP_6}$.

\noindent\textbf{Step 2.} Locally $\sigma_*\mathcal O=\mathcal O$ and its positive direct
images vanish: the additional cyclic group has invariants
$\mathbb C\{v\}^{\mu_k}=\mathbb C\{u\}$, and averaging is
exact. Consequently the cohomology of the pulled-back lines
$\mathcal O(-kl)$ and $\mathcal O(-2kl)$ is the cohomology of
the original $\mathcal O(-1)$ and $\mathcal O(-2)$, hence zero.
Projection and Leray give
$h^\bullet(X,\bar f^*E)=(1,1,0,0)$.
For the horizontal line, its coarse direct image is
$\mathcal O_{\mathbb P^1}(-1)$: a vector field on a root disc
must have its coefficient divisible by the root parameter,
so $\rho_*T_{\mathfrak B_{r,s}}=T_{\mathbb P^1}(-p_1-p_2)$,
and the extra $-p_0$ gives the asserted degree. Moreover
$L$ and $L\otimes\mathcal O(-kl)$ have negative orbifold degree,
so have no global sections. Thus
$H^0(X,\bar f^*L)=H^1(X,\bar f^*L)=0$.
The two exact tangent sequences now give
$h^0(X,T_X)=h^1(X,T_X)=1$.

\noindent\textbf{Step 3.} We compute degree three independently. On the smooth torus locus,
the boundary on a relative holomorphic one-form is the transpose
of the Kodaira--Spencer tensor. The latter is generically invertible
by \ref{cdp:st:E7}. Hence every ambient holomorphic one-form on an inverse-image
base disc is horizontal. It is the pullback of a base one-form.
At a finite root point the root map is a submersion, so that
one-form extends holomorphically in the root coordinate. At a
generic smooth point of $W$ its local equation is $q=z_1$, which
excludes a pole in its base coefficient. Therefore the canonical
pullback comparison gives
\[
 \bar f_*\Omega_X^1=\Omega_{\mathfrak B_{r,s}}^1.
\]
By Serre duality and the projection formula,
\[
 h^3(X,T_X)=h^0\bigl(\mathfrak B_{r,s},
                         \mathcal O(kl-2r-2s)\bigr)=0.
\]
For the last equality, a weighted monomial of that degree would
give nonnegative integers $a,b$ with
$3k(a+2)+4l(b+2)=kl$. Reduction modulo $k$ makes $b+2$ a positive
multiple of $k$, since $\gcd(k,4l)=1$. The second term alone
is then at least $4kl$, a contradiction. On a six-sphere
$\chi(T_X)=1$, giving $h^2=1$. For the original affine-twist
family, $r=3$, $s=4$, and the horizontal Leray terms are the
acyclic lines $\mathcal O(-5)$ and $\mathcal O(-6)$ on
$\mathbb P(3,4)$. Therefore $h^\bullet(\mathcal T)=(1,1,0,0)$,
and the quotient $\eta_*K_{\mathrm dP_6}$ gives
$h^\bullet(T_X)=(1,1,1,0)$ directly for every admissible affine
twist. This computation does not require the integral homology
to be that of a sphere.

\noindent\textbf{Step 4.} The scale derivative changes the fourth torus period by a nonzero
constant multiple of the parameter derivative in \cref{fam:src:A.5.2}, so its
restricted torus class is nonzero. A first-order trivialization of
the compact member would turn the given fibration into a
first-order deformation of the fixed map $f$ with fixed source.
Such deformations are $H^0(X,f^*T_{\mathbb P^1})
=H^0(\mathbb P^1,T_{\mathbb P^1})$, by connected fibres and
the projection formula. Thus they are precisely infinitesimal
projective transformations of the base. The two multiple critical
values and the hexagonal critical value remain the same three
constant sections in this scale family; the local equations
$t=v^r$, $t-1=v^s$, $q=xyz$ show this also over the dual numbers.
The base vector field is therefore zero. The nonzero torus class
cannot be trivialized, proving the nonzero global scale derivative.
The one-dimensional tangent space and the analytic local-algebra
argument at the end of \cref{fam:src:A.5.2} prove smoothness, reducedness and
miniversality.
\end{proof}






\begin{proposition}[Counterexample: The universal local-freeness assertion of CDP98, Lemma 1.5]\label{cdp:st:E12}
\StatementRecord{cdp:st:E12}{proved}{Historical events: Claude Fable 5.2 (Anthropic), effort high, 16 November 2026: carried as background without fault found. GPT-6 Astra (OpenAI), ChatGPT Pro mode, effort Pro, S6\_09\_PAPERS, 17 November 2026: targeted proof check. Also checked in this writing. Pre-continuation writing events: GPT-6 Astra (OpenAI), ChatGPT Pro mode, effort Pro, S6\_10\_PAPERS, 18 November 2026: targeted proof check. Continuation metadata events: GPT-6 Astra (OpenAI), ChatGPT Pro mode, effort Pro, S6\_10\_PAPERS continuation, 19 November 2026: used and rechecked in this run; status-vocabulary reconciliation; metadata reconciliation only; no new mathematical check. \emph{Pre-continuation recorded proof scope.} The derived Poincare-line pushforward C\_0[-1], its product with O\_P1(-2), and the resulting torsion direct image were re-derived}{the hypotheses, notation and definitions stated here; the explicit constructions and calculations in the proof}

Let $A$ be an elliptic curve with origin $0$, let
$P=\mathcal O_{A\times A}(\Delta-A\times\{0\}-\{0\}\times A)$,
and let $Y=A\times A\times\mathbb P^1$, with $g:Y\to A$
projection to the second factor. For
$V=P\boxtimes\mathcal O_{\mathbb P^1}(-2)$ one has
\[
 g_*V=R^1g_*V=0,\qquad R^2g_*V=\mathbb C_0,
 \qquad h^\bullet(Y,V)=(0,0,1,0).
\]
The sheaf in degree two has length one. For $t\ne0$ the fibre
cohomology is zero in every degree; for $t=0$ it is $(0,1,1)$.
\end{proposition}


\begin{proof}
\noindent\textbf{Step 1.} Each nontrivial degree-zero line on an elliptic curve has no
sections: a nonzero section would have an effective divisor of
degree zero and hence trivialize it. Riemann--Roch and Serre
duality then give zero first cohomology. The trivial line has
cohomology $(1,1)$. To compute the direct image with its length,
put $p:A\times A\to A$ and $M=P\otimes\operatorname{pr}_1^*
\mathcal O_A(0)$. Its fibres have degree one, so
$R^1p_*M=0$ and $p_*M$ is a line bundle. The exact sequence for
the section $\{0\}\times A$ gives the two-term complex
\[
 Rp_*P=\bigl[p_*M\longrightarrow\mathcal O_A\bigr].
\]
The arrow evaluates the unique section of the moving degree-one
line at the origin. Its zero is exactly $t=0$ and is simple:
in a local coordinate the moving zero divisor is the diagonal,
with equation $x-t$. Thus the arrow is the inclusion
$\mathcal O_A(-0)\to\mathcal O_A$, and
$Rp_*P=\mathbb C_0[-1]$. Tensoring with
$R\Gamma(\mathbb P^1,\mathcal O(-2))=\mathbb C[-1]$
gives $Rg_*V=\mathbb C_0[-2]$. This proves the direct-image,
length, and global-cohomology assertions. Fibre Kunneth gives
the stated fibre cohomology.

\noindent\textbf{Step 2.} For any proper flat family over a nonsingular curve and a locally
free sheaf $V$, multiplication by a local parameter gives the
actual replacement exact sequence
\[
0\to (R^qg_*V)_c/t(R^qg_*V)_c
 \to H^q(Y_c,V|_{Y_c})
 \to (R^{q+1}g_*V)_c[t]\to0.
\]
It follows directly from the long direct-image sequence of
$0\to V\xrightarrow{t}V\to V|_{Y_c}\to0$. Here $[t]$ denotes
the kernel of multiplication by $t$. Local constancy of all
fibre cohomology dimensions implies local freeness and base
change by Grauert's theorem; generic vanishing alone permits
the torsion explicitly present above.
\end{proof}



In the supplied CDP98 arXiv source, Lemma~1.5 is used in
Proposition~2.5 to pass from a general-fibre vanishing to
$R^2f_*=0$, in the proof of Theorem~3.1 for the asserted
flatness of the two-parameter direct image, in Lemma~5.4 for
$R^1f_*\mathcal O$, in Theorem~5.6 to propagate sections from a
single fibre, and in Proposition~5.7 to remove a torsion
$R^2f_*\mathcal O$. Its first use is replaced in \ref{cdp:st:E2} by vanishing
on every full fibre; the special structure-sheaf assertion in
degree one is independently established in \ref{cdp:st:E4}. The other uses
require their own hypotheses, rather than the refuted lemma.




\begin{proposition}[The Kuranishi germ of the two-parameter family]\label{cdp:st:E13}
\StatementRecord{cdp:st:E13}{proved}{Historical events: Claude Fable 5.2 (Anthropic), effort high, 16 November 2026: carried as background without fault found. GPT-6 Astra (OpenAI), ChatGPT Pro mode, effort Pro, S6\_09\_PAPERS, 17 November 2026: targeted proof check. Also checked in this writing. Pre-continuation writing events: GPT-6 Astra (OpenAI), ChatGPT Pro mode, effort Pro, S6\_10\_PAPERS, 18 November 2026: targeted proof check; used and rechecked in this run. Continuation metadata events: GPT-6 Astra (OpenAI), ChatGPT Pro mode, effort Pro, S6\_10\_PAPERS continuation, 19 November 2026: used and rechecked in this run; status-vocabulary reconciliation; metadata reconciliation only; no new mathematical check. \emph{Pre-continuation recorded proof scope.} The meromorphic connection and its residues, root-curve weights, Hodge quotient comparison, every Leray dimension, two-by-two residue determinant, and independent critical-value and elliptic-modulus derivatives were re-derived. The inverse-function argument kills the full Kuranishi ideal. The named family and Mumford model statements in the attached primary source were opened}{the hypotheses, notation and definitions stated here; \cite[Proposition labelled \texttt{family}]{Viaclovsky2026Two}; \cref{cdp:eq:E13-connection}; \cref{cdp:eq:E13-residue}; \cref{cdp:st:E11}; \cref{cdp:st:E1}; \cref{cdp:st:E6}}

For every member of the family of
\cite[Proposition labelled \texttt{family}]{Viaclovsky2026Two}, on the parameter
bidiscs with $\mu$ sufficiently close to zero and nonzero sufficiently
small linearization scale $\lambda$, one has
\[
 (h^0,h^1,h^2,h^3)(X_{\mu,\lambda},T_X)=(1,2,2,0).
\]
Its Kuranishi germ is a smooth reduced surface, and the displayed
$(\mu,\lambda)$-family is locally miniversal. In particular its
sufficiently small deformations all have algebraic dimension one.
\end{proposition}


\begin{proof}
\noindent\textbf{Step 1.} We work with the equation and section used in the construction,
\[
 S_\mu:\quad y^2=x^3+tx+\mu t+1,\qquad
 P_\mu=(-\mu,b),\qquad b^2=1-\mu^3,
\]
where the branch of $b$ is close to one. Write $p_1,p_2,p_3$ for its
three distinct nodal values, with the zero of the linearization at
$p_3$. The order-four good-reduction cover at infinity gives the
root curve $\mathfrak B=\mathbb P(1,4)$, with coarse projection
$\rho$ and $\rho^*\mathcal O_{\mathbb P^1}(1)=\mathcal O(4)$.
Let $\bar f:X\to\mathfrak B$ denote the root factorization and
$D=W_1+W_2+W_3$ the reduced divisor of Mumford fibres.

\noindent\textbf{Step 2.} We spell out the bundles and the residue map, since their ranks alone
would not determine the answer. The vertical logarithmic Lie bundle
is
\[
 E=\mathcal O\oplus\mathcal O(-1).
\]
The first summand is generated by multiplication in the line-bundle
coordinate. For a concrete splitting of the second summand put
$\omega=dx/(2y)$ and use the relative differential
\begin{equation}\label{cdp:eq:E13-connection}
 \alpha=\frac{y+b}{x+\mu}\,\frac{dx}{2y}.
\end{equation}
Its residues are $1$ at $P_\mu$ and $-1$ at the zero section; there
is no pole at $-P_\mu$, because its numerator vanishes there.
Consequently, in the rational section of
$\mathcal O(P_\mu-O)$ having divisor $P_\mu-O$, the connection
$\nabla(hs)=(dh+\alpha h)s$ is holomorphic along each smooth
elliptic fibre. Here $s$ is that rational section and a point of
the line bundle is written $v s$; the horizontal lift of a vector
$V$ is therefore $\widetilde V=V-\alpha(V)v\partial_v$.
It is a multiplicative connection: the difference between its
pullback under multiplication and the sum of its two pullbacks is a
holomorphic one-form on the product of two elliptic curves, and its
restrictions to both coordinate axes vanish. Every holomorphic
one-form on that product is the sum of forms pulled back from those
axes, so this difference is zero. Its horizontal lift therefore
commutes with the lifted translation defining $X$.

\noindent\textbf{Step 3.} At infinity put $t=r^{-4}$, $x=r^{-2}x'$, $y=r^{-3}y'$. Formula
\eqref{cdp:eq:E13-connection}, in the relative cotangent sheaf, becomes
\[
 \alpha=\frac{y'+br^3}{x'+\mu r^2}\frac{dx'}{2y'},
 \qquad y'^2=x'^3+x'+\mu r^2+r^6.
\]
This is the same holomorphic connection at the good-reduction
fibre; the limiting section $(0,0)$ is a nonzero two-torsion point,
so it remains disjoint from the zero section. Moreover
$\omega=r\,dx'/(2y')$, which gives the summand $\mathcal O(-1)$
for its dual Lie vector. At a nodal fibre the connection is
logarithmic at the two branches and its lift is a section of the
vertical logarithmic tangent bundle. Thus these calculations prove
the claimed splitting on the entire root curve, with the specified
comparison maps.

\noindent\textbf{Step 4.} The inverse pairing associated with the nondegenerate alternating
class $\xi=u\wedge w+2\psi\wedge\delta$ in the construction induces a
holomorphic isomorphism
\[
 G:=R^1\bar f_*\mathcal O_X\simeq E.
\]
Here is the extension check. On a marked smooth fibre the periods
have cross terms in ratio $1:-2$, so the holomorphic one-form
subspace $F^1\subset H^1(F,\mathbb C)$ is Lagrangian for this pairing.
Pairing with $F^1$ therefore gives the isomorphism
$H^1(F,\mathbb C)/F^1\to(F^1)^*=\operatorname{Lie}(F)$.
At a Mumford point, subtract the logarithmic monodromy term from the
periods. In the resulting holomorphic frame the alternating pairing
is still constant and nondegenerate, and the two holomorphic
one-forms are the logarithmic coordinate differentials. Their dual
frame is the logarithmic Lie frame. The quotient frame gives
$R^1\bar f_*\mathcal O_X$: at a rank-one point its first class is
the class of the loop in the normalization diagram and its second
is the class on the elliptic normalization; at the rank-two point
both classes are loops in the normalization diagram. These are
bases, by the normalization exact sequence. Thus the pairing extends
as an isomorphism at these points. The finite good-reduction cover
gives the same conclusion at infinity, and descent uses invariance
of $\xi$. This identifies the comparison, including its value at
the special fibres. Cup product and the normalization sequences
give $R^2\bar f_*\mathcal O_X=\det G=\mathcal O(-1)$.

\noindent\textbf{Step 5.} The logarithmic tangent and canonical sequences are therefore
\[
\begin{split}
 0&\longrightarrow\bar f^*E\longrightarrow T_X(-\log D)
 \longrightarrow\bar f^*\mathcal O(-7)\longrightarrow0,\\
 K_X&=\bar f^*\mathcal O(-4).
\end{split}
\]
Indeed $K_{\mathfrak B}=\mathcal O(-5)$ and
$T_{\mathfrak B}(-p_1-p_2-p_3)=\mathcal O(5-12)$; taking the
determinant of the displayed logarithmic sequence gives the
canonical formula after removing $D$ of degree $12$.
For root line bundles, cohomology is computed by weighted monomials
and Serre duality with $\mathcal O(-5)$. The projection formula and
Leray give
\[
\begin{split}
 h^\bullet(X,\bar f^*E)&=(1,1,0,0),\\
 h^\bullet(X,\bar f^*\mathcal O(-7))&=(0,1,2,1),\\
 h^\bullet(X,T_X(-\log D))&=(1,2,2,1).
\end{split}
\]
There are no uncomputed differentials in these equalities: the base
has dimension one, and the groups of the first displayed row in
degrees two and three vanish.

\noindent\textbf{Step 6.} Let $Q=\bigoplus_{i=1}^3\eta_{i*}K_{\widetilde W_i}$, the normal
quotient of $T_X(-\log D)\to T_X$. Here
$\widetilde W_1,\widetilde W_2$ are ruled over elliptic curves and
$\widetilde W_3=dP_6$, so
\[
 h^\bullet(Q)=(0,2,3).
\]
It remains to compute the boundary
$H^1(Q)\to H^2(T_X(-\log D))$, between two-dimensional spaces.
Under Serre duality its dual is the residue homomorphism
\begin{equation}\label{cdp:eq:E13-residue}
 H^1\bigl(X,\Omega_X^1(\log D)\otimes K_X\bigr)
 \longrightarrow
 H^1(\widetilde W_1,\mathcal O)\oplus
 H^1(\widetilde W_2,\mathcal O).
\end{equation}
The dual logarithmic sequence, tensored with $K_X$, has horizontal
term $\bar f^*\mathcal O(3)$ and vertical quotient
$\bar f^*(\mathcal O(-4)\oplus\mathcal O(-3))$.
The latter has zero cohomology in degrees zero and one: the coarse
direct images of $\mathcal O(-4)$ and $\mathcal O(-3)$ are both
$\mathcal O_{\mathbb P^1}(-1)$, and every line in their tensor
products with $G$ has negative weighted degree, so has no section.
Hence the
source of \eqref{cdp:eq:E13-residue} is canonically
\[
 H^1(X,\bar f^*\mathcal O(3))
 =H^0\bigl(\mathfrak B,\mathcal O(3)\otimes G\bigr)
 =H^0(\mathcal O(3))\oplus H^0(\mathcal O(2))=\mathbb C^2.
\]
On the affine $t$-line the generators of these last two spaces are
constant in the frame just constructed. The residue in each
summand is a nonzero scalar times evaluation followed by
$G_{p_i}\to H^1(\widetilde W_i,\mathcal O)$.

\noindent\textbf{Step 7.} Write $e_i$ for the double root in the nodal cubic at $p_i$.
The line killed by this last projection is the graph-loop class.
Under $G\simeq E$ it is the toric Lie line: in the local periods
$(q\alpha,\beta),(\gamma,\lambda\delta)$ it is the direction
of the degenerating first multiplicative coordinate, whereas the
second coordinate is the elliptic normalization. If $V$ is dual to
$\omega$, its horizontal lift $\widetilde V$ has vertical
line-coordinate component $-\alpha(V)$. At the nodal boundary put
\[
 a_i=\frac{b}{e_i+\mu}.
\]
Thus the toric line is generated by
$\widetilde V+a_i v\partial_v$, and the projection to the elliptic
Lie line sends $(v\partial_v,\widetilde V)$ to $(1,-a_i)$, up to
a common nonzero factor at that point. Consequently the matrix of
\eqref{cdp:eq:E13-residue}, after nonzero row rescalings, is
\[
 \begin{pmatrix}1&-a_1\\1&-a_2\end{pmatrix},\qquad
 a_1-a_2=\frac{b(e_2-e_1)}{(e_1+\mu)(e_2+\mu)}\ne0.
\]
The inequality is exact: $b\ne0$; distinct nodal values have
distinct double roots, since $p_i=-3e_i^2$; and $e_i=-\mu$ would
place $P_\mu$ at the node and force $b=0$. The boundary therefore
has rank two.

\noindent\textbf{Step 8.} The rank-two cusp implies
$\bar f_*\Omega_X^1=\Omega_{\mathfrak B}^1$, by the generic
Kodaira--Spencer boundary and extension across the local models,
as in \ref{cdp:st:E11}. It follows that
$h^3(T_X)=h^0(\Omega_X^1\otimes K_X)
=h^0(\mathcal O(-9))=0$. The long exact sequence with quotient
$Q$ now gives $h^0(T_X)=1$, $h^1(T_X)=2$, and $h^2(T_X)=2$.

\noindent\textbf{Step 9.} Finally the constructed bidisc has injective Kodaira--Spencer map.
At $\mu=0$, writing $a_i=p_i(0)$, differentiation of
$4t^3+27(\mu t+1)^2=0$ gives $p_i'(0)=-9/(2a_i)$.
An infinitesimal automorphism fixing the multiple value at infinity
has the form $At+B$; the three equalities
$-9/(2a_i)=Aa_i+B$ would make a nonzero quadratic polynomial vanish
at all three $a_i$. The cross-ratio variation is therefore nonzero,
and remains so on a smaller disc. For fixed $\mu$, a normalized
double curve is $\mathbb C^*/\langle c(\mu)\lambda\rangle$.
Its elliptic modulus has nonzero $\lambda$-derivative when
$|c(\mu)\lambda|$ is sufficiently small, as follows from
$j(q)=q^{-1}+744+O(q)$. A trivial first-order deformation of $X$
would trivialize its intrinsic anticanonical pencil, its labelled
exceptional-value configuration, and this double curve; it can
therefore kill neither of the two independent tangent directions.
Thus the bidisc maps to the two-dimensional tangent space of the
Kuranishi germ with invertible differential. Present its local
ring as $\mathbb C\{z_1,z_2\}/I$; the induced substitution into
$\mathbb C\{\mu,\lambda-\lambda_0\}$ is a local analytic
isomorphism before quotienting, by the inverse function theorem.
Every element of $I$ vanishes under it, so $I=0$. This proves the
Kuranishi assertion. \ref{cdp:st:E1} and \ref{cdp:st:E6} give the algebraic-dimension assertion.
\end{proof}






\begin{proposition}[The Liu--Rollenske family and its transverse deformation]\label{cdp:st:E14}
\StatementRecord{cdp:st:E14}{proved}{Historical events: Claude Fable 5.2 (Anthropic), effort high, 16 November 2026: carried as background without fault found. GPT-6 Astra (OpenAI), ChatGPT Pro mode, effort Pro, S6\_09\_PAPERS, 17 November 2026: targeted proof check. Also checked in this writing. Pre-continuation writing events: GPT-6 Astra (OpenAI), ChatGPT Pro mode, effort Pro, S6\_10\_PAPERS, 18 November 2026: targeted proof check; used and rechecked in this run. Continuation metadata events: GPT-6 Astra (OpenAI), ChatGPT Pro mode, effort Pro, S6\_10\_PAPERS continuation, 19 November 2026: used and rechecked in this run; status-vocabulary reconciliation; metadata reconciliation only; no new mathematical check. \emph{Pre-continuation recorded proof scope.} The explicit Weierstrass coordinate comparison, simultaneous cusp smoothing, Hartogs extension across the two-dimensional base, free logarithmic tangent module, and the actual dual-number restriction map with determinant i were re-derived. The tangent tuple and the two independent deformation directions follow from those exact maps. The local scale germ is used on the covering parameter specified in P3K.36}{the hypotheses, notation and definitions stated here; \cite[Proposition labelled
\texttt{prop: surface facts}]{LiuRollenske2026}; \cref{cdp:eq:E14-elliptic-family}; \cite[Theorem labelled \texttt{thm: main}]{LiuRollenske2026}; \cref{cdp:st:E13}; \cref{cdp:st:E4}; \cref{cdp:eq:E14-canonical}; \cref{cdp:eq:E14-jet-map}; \cref{cdp:eq:E14-log-coho}; \cref{cdp:st:E6}; \cref{cdp:st:E1}}

Every member $Y(u)$ of the Liu--Rollenske family, with its scale
sufficiently small and nonzero as in its construction, satisfies
\[
 h^\bullet(Y(u),T_{Y(u)})=(1,2,2,0).
\]
Its Kuranishi germ is a smooth reduced surface. One parameter is
the given scale; a second parameter replaces its cuspidal fibre by
two rank-one nodal fibres. All sufficiently small deformations
retain a torus fibration and have algebraic dimension one.
\end{proposition}


\begin{proof}
\noindent\textbf{Step 1.} We derive the needed holomorphic invariants from the local models.
The commented-out holomorphic-invariants notes in the supplied
Liu--Rollenske source are not used as established results.

\noindent\textbf{Step 2.} In the Weierstrass equation of \cite[Proposition labelled
\texttt{prop: surface facts}]{LiuRollenske2026}, make the explicit
coordinate comparison
\[
 t_{\rm new}=\frac{27}{4}(t_{\rm old}^{-1}-1),\quad
 x_{\rm new}=x_{\rm old}/t_{\rm old}^2,\quad
 y_{\rm new}=y_{\rm old}/t_{\rm old}^3.
\]
It gives $y^2=x^3+tx+t$, with section $P=(-1,i)$; the multiple
value is now infinity, the cusp is at zero and the remaining nodal
value is $-27/4$. Consider the explicit deformation
\begin{equation}\label{cdp:eq:E14-elliptic-family}
 y^2=x^3+tx+t+\epsilon,\qquad
 P_\epsilon=(-1,b),\quad b^2=\epsilon-1,
 \quad b(0)=i.
\end{equation}
Its finite discriminant is $4t^3+27(t+\epsilon)^2$.
The nodal value near $-27/4$ remains simple and holomorphically
labelled. The other two nodes coalesce to the type-II cusp at
$\epsilon=0$. Infinity keeps type $III^*$; on its order-four
good-reduction cover the equation is
\[
 y'^2=x'^3+x'+r^2+\epsilon r^6,
 \qquad P'_{\epsilon}=(-r^2,br^3).
\]

\noindent\textbf{Step 3.} These equations give a simultaneous construction of the compact
threefolds, including the central one. On a disc containing the cusp
and its two nearby nodal values, the total elliptic surface in
\eqref{cdp:eq:E14-elliptic-family} is smooth over $\epsilon$:
$\partial(y^2-x^3-tx-t-\epsilon)/\partial t=-(x+1)$ is a unit
near the cusp. Translation by $-2P_\epsilon$ is an automorphism
there. The line bundle $\mathcal O(P_\epsilon-O)$ and its lifted
translation vary holomorphically; its multiplier has a zero only
at the separate simple nodal value chosen for the hexagonal filling.
For a uniformly small nonzero scale its contraction on the
complement of the zero section is uniform on the compact elliptic
fibres. The quotient is consequently free, properly discontinuous,
proper over the disc, and smooth over the parameter. This quotient
itself fills the cusp and the nearby rank-one nodal fibres. At the
separate node the ordinary rank-two Mumford charts vary
holomorphically. At infinity the fixed order-four torsion translation
and good-reduction comparison vary holomorphically in the displayed
equation. Gluing along their marked punctured collars therefore
constructs a proper holomorphic family over $(\epsilon,u)$ whose
central scale slice is $Y(u)$. Every member has the central
differentiable type by Ehresmann. The supplied topology theorem
\cite[Theorem labelled \texttt{thm: main}]{LiuRollenske2026} thus
identifies these differentiable manifolds with $S^6$.

\noindent\textbf{Step 4.} Write $\mathfrak B=\mathbb P(1,4)$ and retain its convention
$\rho^*\mathcal O(1)=\mathcal O(4)$. On the central threefold the
canonical form on the contracting quotient is the elliptic-surface
two-form wedged with $dv/v$. At the cusp it extends without a zero
or pole: from $t=(y^2-x^3)/(1+x)$ one gets
\[
 dt\wedge\frac{dx}{2y}=-\frac{dx\wedge dy}{1+x}.
\]
The finite-cover and hexagonal calculations are the same explicit
coordinate calculations as in \ref{cdp:st:E13}. Since the elliptic surface has
$K_S=\pi^*\mathcal O_{\mathbb P^1}(-1)$, these comparisons give
\begin{equation}\label{cdp:eq:E14-canonical}
 K_Y=f^*\mathcal O_{\mathbb P^1}(-1)
     =\bar f^*\mathcal O_{\mathfrak B}(-4).
\end{equation}
The calculation is simultaneous in $\epsilon$. Adjunction on every
full fibre gives trivial dualizing sheaf. \ref{cdp:st:E4} gives $h^0=1$;
duality gives $h^2=1$; constancy of its Euler characteristic in the
proper flat family gives $h^1=2$. Thus, also on the two-dimensional
base $(t,\epsilon)$, Grauert's theorem gives locally free $R^qf_*\mathcal O$
with base change, for $q=0,1,2$.

\noindent\textbf{Step 5.} The multiplicative connection
\[
 \alpha=\frac{y+b}{x+1}\frac{dx}{2y},\qquad
 \nabla(hs)=(dh+\alpha h)s
\]
has precisely the residues of the rational section $s$ with divisor
$P_\epsilon-O$. The same proof and good-reduction substitution as
in \ref{cdp:st:E13} give the logarithmic Lie bundle
$E=\mathcal O\oplus\mathcal O(-1)$, including at the central
cusp. There the vector dual to $dx/(2y)$ is the derivation
$2y\partial_x+3x^2\partial_y$ on the cuspidal fibre, and its
horizontal lift is well defined because
$(y+b)/(x+1)$ is regular near the singular locus.

\noindent\textbf{Step 6.} We justify the identification with $G=R^1\bar f_*\mathcal O$
at the cusp without assuming a limiting-period rule there.
The inverse pairing associated with the invariant alternating
class gives $G\simeq E$ over the two-dimensional $(t,\epsilon)$
base away from its single cuspidal point: all remaining points
are smooth, nodal or the finite-cover points already computed in
\ref{cdp:st:E13}. Both $G$ and $E$ are locally free across that point, by the
preceding cohomological-flatness and connection calculations.
The matrix entries of the isomorphism and its inverse extend
across that codimension-two set by Hartogs' theorem. Their products
are the identity on a dense open set, hence everywhere. Thus
\[
 G=\mathcal O\oplus\mathcal O(-1),\qquad
 R^2\bar f_*\mathcal O=\mathcal O(-1)
\]
on the central root curve as well; the second equality also follows
directly from relative duality and \eqref{cdp:eq:E14-canonical}.

\noindent\textbf{Step 7.} Let $D=C+W$ be the sum of the cuspidal and hexagonal fibres.
For the local tangent calculation choose coordinates in which
$t=y^2-x^3$. They are obtained from the earlier formula by the
invertible change $\widetilde x=x/(1+x)^{1/3}$,
$\widetilde y=y/(1+x)^{1/2}$, with the branches equal to one at
zero, and then renaming the two new coordinates. The logarithmic
tangent module is freely generated by
\[
 2y\partial_x+3x^2\partial_y,\quad
 2x\partial_x+3y\partial_y,\quad\partial_z,
\]
where $z$ is a coordinate along the elliptic singular curve. The
second generator maps to $6t\partial_t$;
the first and third generate the relative logarithmic kernel.
Consequently the global exact sequence is
\[
 0\longrightarrow\bar f^*E\longrightarrow T_Y(-\log D)
 \longrightarrow\bar f^*\mathcal O(-3)\longrightarrow0,
\]
because $T_{\mathfrak B}(-C-W)=\mathcal O(5-8)$.
The vertical cohomology is $(1,1,0,0)$. Every line in the
horizontal Leray terms is $\mathcal O(-3)$ or $\mathcal O(-4)$,
which has no cohomology on $\mathbb P(1,4)$. Hence
\begin{equation}\label{cdp:eq:E14-log-coho}
 h^\bullet(T_Y(-\log D))=(1,1,0,0).
\end{equation}

\noindent\textbf{Step 8.} The quotient at $W$ is $\eta_*K_{dP_6}$ with cohomology $(0,0,1)$.
At $C$, the map $v\mapsto v(t)|_C$ identifies the quotient with
the Jacobian ideal
\[
 J=(x^2,y)\subset\mathcal O_C.
\]
Indeed its kernel is precisely the logarithmic tangent module and
$N_{C/Y}=\mathcal O_C$. Put $Z=V(J)\subset C$. We now compute
the actual restriction map
\begin{equation}\label{cdp:eq:E14-jet-map}
 H^1(C,\mathcal O_C)\longrightarrow H^1(Z,\mathcal O_Z).
\end{equation}

\noindent\textbf{Step 9.} In the deformation \eqref{cdp:eq:E14-elliptic-family} the critical
base near the cusp is the smooth curve
\[
 \Gamma:\quad t=-3x^2,\qquad
 \epsilon=x^2(3+2x).
\]
Its map to the $\epsilon$-disc is finite flat of degree two; its
central fibre is $\operatorname{Spec}\mathbb C[x]/(x^2)$.
The lifted translation fixes its $x$-coordinate. This follows on
the punctured curve because translation fixes the unique node of
each nodal cubic, and then follows at zero by holomorphic
continuation. On the line coordinate it is
$(x,v)\mapsto(x,c(x)v)$ with $c$ a holomorphic nowhere-zero
contraction. Thus its quotient critical locus is exactly the
elliptic family
\[
 \mathcal Z=(\Gamma\times\mathbb C^*)/
       \langle(x,v)\mapsto(x,c(x)v)\rangle.
\]
In particular $H^0(Z,\mathcal O_Z)$ and
$H^1(Z,\mathcal O_Z)$ are both free rank-one modules over
$A=\mathbb C[x]/(x^2)$, and $H^2(Z,\mathcal O_Z)=0$.

\noindent\textbf{Step 10.} Restriction in this proper family is a holomorphic bundle map
$G|_\Gamma\to R^1(\mathcal Z/\Gamma)_*\mathcal O_{\mathcal Z}$.
On $\Gamma\setminus\{0\}$ the calculation of \ref{cdp:st:E13} identifies it,
after a nonzero scalar change of target frame, with the row
\[
 \left(1,-\frac{b}{1+x}\right).
\]
The same frame can be chosen at zero. To check this last point,
normalize $C$. Its normalization is a $\mathbb P^1$-bundle over
the elliptic singular curve $E_0$, and its normalization quotient
is $\mathcal O_{E_0}$: the first jet at the cusp transforms with
derivative one under translation on the additive smooth locus.
The normalization exact sequence, together with
$h^1(\mathcal O_C)=2$, $h^2(\mathcal O_C)=1$ and
$h^\bullet(\mathcal O_{\widetilde C})=(1,1,0)$, implies
that $H^1(\mathcal O_C)\to H^1(\mathcal O_{\widetilde C})$
is onto. Restriction from the ruled normalization to its singular
section identifies the latter group with $H^1(\mathcal O_{E_0})$.
Thus restriction to $E_0$ is onto. The common scalar in the row
above cannot have a zero at zero, since both entries would then
vanish and this restriction would be zero. It is a unit. This
proves the asserted holomorphic target-frame comparison.

\noindent\textbf{Step 11.} The equality of these holomorphic maps on the punctured curve
extends over all of $\Gamma$, and in particular over its
nonreduced central fibre. Hence \eqref{cdp:eq:E14-jet-map} is,
in the specified bases, the complex-linear map
\[
 \mathbb C^2\longrightarrow A,\qquad
 (a,d)\longmapsto a+d(-i+ix).
\]
Its matrix in $1,x$ has determinant $i$, so it is an isomorphism.
This is the required jet comparison, not an inference from the
rank at nearby reduced fibres.

\noindent\textbf{Step 12.} The exact sequence $0\to J\to\mathcal O_C\to\mathcal O_Z\to0$
now gives
\[
 h^\bullet(C,J)=(0,1,1).
\]
Together with \eqref{cdp:eq:E14-log-coho} and the hexagonal quotient,
the tangent exact sequence gives
$h^\bullet(T_Y)=(1,2,2,0)$.

\noindent\textbf{Step 13.} It remains to verify that the actual two-parameter family has
bijective Kodaira--Spencer map. In the above local trivialization
of its total space the map to the base is
\[
 f_{\epsilon}=\frac{y^2-x^3-\epsilon}{1+x}.
\]
The derivative $\partial_\epsilon f$ has image $-1+x$ in
$\mathcal O_C/J=\mathbb C[x]/(x^2)$ along the singular curve.
If the ambient first-order deformation were trivial, every
first-order variation of its map would be a global section of
$f^*T_{\mathbb P^1}$, hence the pullback of a projective vector
field, because $f_*\mathcal O_Y=\mathcal O_{\mathbb P^1}$.
Its local image in this Jacobian quotient is constant. Changing
the local ambient trivialization adds $df(v)$ and hence adds
zero in that quotient. The nonzero coefficient of $x$ therefore
proves that the $\epsilon$-direction is nontrivial and cannot be
cancelled by the scale direction, which does not change this local
map. With $\epsilon=0$, a normalized singular elliptic curve is
$\mathbb C^*/\langle c_0u\rangle$; its modulus has nonzero
$u$-derivative for sufficiently small $u$. A trivial deformation
of the total space preserves its anticanonical pencil from
\eqref{cdp:eq:E14-canonical}, its three distinguished special values,
and that elliptic singular curve. Hence the scale direction is
nontrivial too. The two directions span $H^1(T_Y)$.
The inverse-function argument of \ref{cdp:st:E13} consequently makes the
Kuranishi ideal zero and proves miniversality. \ref{cdp:st:E6} and \ref{cdp:st:E1} give the
remaining assertions.
\end{proof}






\begin{proposition}[Tangent cohomology and Kuranishi germs of positive combs]\label{cdp:st:E15}
\StatementRecord{cdp:st:E15}{proved}{Historical events: Claude Fable 5.2 (Anthropic), effort high, 15 November 2026: carried as background without fault found. Claude Fable 5.2 (Anthropic), effort high, 16 November 2026: carried as background without fault found. Claude Opus 5.5 (Anthropic), effort max, sweep 1, scope s6, 18 November 2026: used and rechecked in this run. GPT-6 Astra (OpenAI), ChatGPT Pro mode, effort Pro, S6\_09\_PAPERS, 17 November 2026: targeted proof check. Also checked in this writing. Pre-continuation writing events: GPT-6 Astra (OpenAI), ChatGPT Pro mode, effort Pro, S6\_10\_PAPERS, 18 November 2026: targeted proof check; used and rechecked in this run. Continuation metadata events: GPT-6 Astra (OpenAI), ChatGPT Pro mode, effort Pro, S6\_10\_PAPERS continuation, 19 November 2026: used and rechecked in this run; status-vocabulary reconciliation; metadata reconciliation only; no new mathematical check. \emph{Pre-continuation recorded proof scope.} The invariant-generator orders, nonnegative coarse splittings, all three elementary modifications, cup-product injection, tangent tuple, root-section conditions, uniform cusp filling and annular gluing were re-derived. The period family has the full injective Kodaira-Spencer map, so the inverse-function argument kills the Kuranishi ideal. The source's fixed-value parameter space and its nine-dimensional quotient were matched to the actual sheaf inclusion}{the hypotheses, notation and definitions stated here; \cite{Viaclovsky2026Combs}; \cref{cdp:eq:E15-products}; \cref{cdp:eq:E15-positive}; \cref{cdp:st:E5}; \cref{cdp:st:E13}; \cref{cdp:st:E1}; \cref{cdp:st:E6}}

For each positive-comb manifold $X_m$ constructed in
\cite{Viaclovsky2026Combs}, in particular for every $m\geq12$ in
its stated construction, one has
\[
 h^\bullet(X_m,T_{X_m})=(2m+1,\,8m+1,\,6m+1,\,0).
\]
Its Kuranishi germ is smooth and reduced of dimension $8m+1$.
A locally miniversal family is supplied by all sufficiently nearby
admissible period sections on the root curve, with the affine
attachment classes fixed. Its members have algebraic dimension one
and torus fibrations with the same three critical values,
multiplicities $3,4$, and hexagonal fibre type.
\end{proposition}


\begin{proof}
\noindent\textbf{Step 1.} Let $\mathfrak B=\mathbb P(3,4)$, with coarse projection
$\rho:\mathfrak B\to B$. Write $p_1,p_2$ for the finite-isotropy
points and $p_0$ for the hexagonal value. On the root curve let
$\mathcal H$ be the extended cohomology bundle, $F$ the Hodge
subbundle, $G$ its quotient, and $V$ its dual:
\[
 0\longrightarrow F\longrightarrow\mathcal H
 \longrightarrow G\longrightarrow0,\qquad V=F^*.
\]
The finite-cover comparison identifies
$G=R^1\bar f_*\mathcal O_{X_m}$ by invariant pushforward of the
ordinary torus Hodge sequence. At the hexagon the two deck
characters of the infinite toric quilt map to the two generators
of $H^1(\mathcal O_W)$ in its normalization complex, establishing
the same comparison there. The bundle $\bar f^*V$ is the vertical
logarithmic tangent bundle: its translations are dual to the
holomorphic one-forms on smooth fibres, and at the hexagon the
dual logarithmic coordinate fields extend in every toric chart.
Cup product identifies $R^2\bar f_*\mathcal O$ with
$D:=\bigwedge^2G$. On finite covers this is the ordinary torus
exterior algebra; at the hexagon it is the exterior algebra of
the two cycles in the normalization complex.

\noindent\textbf{Step 2.} We specify finite characters by the orders of their invariant
generators. At $p_1$ the nontrivial Hodge line has invariant
generator $w$ times an equivariant frame, while the nontrivial
line of each of $V$ and $G$ has generator $w^2$ times its frame.
At $p_2$ the corresponding orders are $3$ for $F$ and $1$ for
$V,G$. Each bundle also has one trivial-character line. These
characters persist because comb smoothing fixes the limiting
periods. The tautological quotient on each tooth is
$\mathcal O(1)^{\oplus2}$, so each tooth adds $2$ to the coarse
degree of $G$. Therefore
\[
 \deg\rho_*G=2m-1,\qquad D=\mathcal O_{\mathfrak B}(24m-1).
\]
The determinant of $\mathcal H$ has degree zero on the root
curve: its determinant local system is trivial, and its cusp
residue has trace zero. Thus
\[
 \det F=\mathcal O(1-24m),\quad
 \det V=\mathcal O(24m-1),\quad
 \deg\rho_*V=2m-1.
\]
Indeed the finite fractional correction for $V$ and $G$ is
$2/3+1/4=11/12$; subtracting it from $(24m-1)/12$ gives
$2m-1$. For $F$ the correction $1/3+3/4=13/12$ gives
$\deg\rho_*F=-2m-1$.

\noindent\textbf{Step 3.} The coarse bundles $V_c=\rho_*V$ and $G_c=\rho_*G$ have
nonnegative splitting. To prove this, evaluate the unchanged
invariant real homology class $\delta$ on $F$. Its resulting
holomorphic section of $V$ never vanishes on a smooth fibre:
a real class annihilating a Hodge plane also annihilates its
conjugate. The same argument holds on the finite good-reduction
covers. At the hexagon, $\delta$ is a nonzero combination of
angular cycles and the logarithmic forms have identity angular
period matrix. Its evaluation is nonzero there too, and gives
a nowhere-zero section of $V_c$. The invariant real cohomology
class $\gamma$ likewise gives a nowhere-zero section of $G_c$;
at the cusp it belongs to $\ker N$, whose projection to the
Hodge quotient is an isomorphism in the transverse graph chart.
At each finite point these nonzero invariant central values give
nonzero values of the coarse sections. A split rank-two bundle
of positive degree on $\mathbb P^1$ with a negative summand
cannot have a nowhere-zero section: every section lies in its
other summand, which has positive degree and hence has zeros.
Since $2m-1>0$, both splittings are nonnegative. In particular
\begin{equation}\label{cdp:eq:E15-positive}
 h^0(V_c)=h^0(G_c)=2m+1,\qquad H^1(V_c)=H^1(G_c)=0.
\end{equation}

\noindent\textbf{Step 4.} Set $A=\rho_*(V\otimes G)$ and $L=\mathcal O_{\mathfrak B}(-5)$.
Multiplication of invariant sections gives exact comparisons
\begin{equation}\label{cdp:eq:E15-products}
\begin{split}
 0&\longrightarrow V_c\otimes G_c\longrightarrow A
       \longrightarrow\mathbb C_{p_1}\longrightarrow0,\\
 0&\longrightarrow V_c(2m-1)\longrightarrow\rho_*(V\otimes D)
       \longrightarrow\mathbb C_{p_1}\longrightarrow0,\\
 0&\longrightarrow G_c(-1)\longrightarrow\rho_*(L\otimes G)
       \longrightarrow\mathbb C_{p_1}\longrightarrow0.
\end{split}
\end{equation}
Here is the local verification. At $p_1$, multiplying two
nontrivial invariant generators of $V,G$ gives $w^4$, whereas
their tensor character has generator $w$; multiplication is
therefore one factor $t=w^3$. At $p_2$ the product is $w^2$
and the tensor generator also has order two, giving no defect.
Products involving a trivial character have no defect. The
line $D$ has the nontrivial character of $G$ and coarse degree
$2m-1$, proving the second comparison. The invariant-generator
orders of $L$ are $1,1$, and $\rho_*L=\mathcal O_B(-1)$.
Its product with the nontrivial generator of $G$ has the same
single defect at $p_1$, proving the third comparison.

\noindent\textbf{Step 5.} All summands on the left in \eqref{cdp:eq:E15-products} have zero
$H^1$, by \eqref{cdp:eq:E15-positive}; for the last sequence their
degrees are at least $-1$. Consequently
\[
\begin{array}{c|c|c}
\text{bundle on }B&h^0&h^1\\ \hline
 A&8m+1&0\\
 \rho_*(V\otimes D)&6m&0\\
 \rho_*(L\otimes G)&2m&0
\end{array}
\]
by Riemann--Roch: their degrees are $8m-3$, $6m-2$, $2m-2$
and their ranks $4,2,2$. Also
$\rho_*(L\otimes D)=\mathcal O_B(2m-1)$.
Finite-group invariant pushforward is exact, so these calculations
give the cohomology on the root curve.

\noindent\textbf{Step 6.} The logarithmic tangent sequence is
\[
 0\longrightarrow\bar f^*V\longrightarrow T_{X_m}(-\log W)
 \longrightarrow\bar f^*L\longrightarrow0.
\]
Projection and Leray yield
\[
 h^\bullet(\bar f^*V)=(2m+1,8m+1,6m,0),\qquad
 h^\bullet(\bar f^*L)=(0,2m,2m,0).
\]
Its only possible connecting map between nonzero groups is
$H^0(L\otimes G)\to H^0(V\otimes D)$. This is the map of
sections of the fibrewise cup-product boundary in \ref{cdp:st:E5}. The cusp
has rank two, so its sheaf map is generically injective. A map
from a locally free sheaf on a smooth curve with zero generic
kernel is injective; hence this map on sections is injective.
It follows that
\[
 h^\bullet(T_{X_m}(-\log W))=(2m+1,8m+1,6m,0).
\]
The branch-normal quotient is $\eta_*K_{dP_6}$, since the toric
three-form trivializes the canonical bundle near the fibre.
Its cohomology is $(0,0,1)$, and its exact sequence proves the
asserted tangent-cohomology tuple.

\noindent\textbf{Step 7.} We construct a family with this full tangent space. All nearby
root-period sections have normal bundle $V\otimes G$, tangent
space $H^0(A)$, and obstruction space $H^1(A)=0$. This section
germ can be realized using an ordinary coarse Grassmannian bundle
and finite analytic conditions. Split the equivariant bundle as
$F\oplus G$ on a finite-cover chart and write a nearby plane as
a graph matrix. Dividing the invariant-generator factors shows
that a holomorphic lift requires one entry to vanish at $p_1$
and two at $p_2$. At $p_1$ the entry is from the nontrivial Hodge
line to the trivial quotient line. At $p_2$ it is that entry
and the entry between the two nontrivial lines. These are linear
conditions on the values of the local graph matrix. Their normal
sheaf is $A$. The Cech gluing obstruction across a small Artin
extension lies in $H^1(A)=0$; each local graph chart and its
stated conditions is smooth. The finite-dimensional Douady
section germ with these conditions is therefore smooth of
dimension $8m+1$. This description does not require a global
manifold cover of the root curve.

\noindent\textbf{Step 8.} These sections produce compact threefolds even though their
limiting periods are allowed to vary. At a finite point the
equivariant lift is a small deformation on the compact
good-reduction cover. Admissibility is open there. The affine
action has the same real formula $x\mapsto Tx+v/n$; it is
holomorphic for the new equivariant periods and stays free
because its real action is unchanged. At the hexagon the
transverse graph chart is open, and its periods are
\[
 Z(q)=\frac{\log q}{2\pi i}I+H(q)
\]
with $H$ holomorphic and uniformly bounded on a smaller parameter
polydisc. The translations with unit factors
$\exp(2\pi iH_{ij})$ preserve the same toric fan. On logarithmic
radii their action is
$r\mapsto r+(I+O(1/\log|q|))\lambda$; it is uniformly proper
and free for small $q$. A fixed finite collection of toric unit
polydiscs gives compact representatives. Thus the filling is
proper with local equation $q=x_1x_2x_3$, and its total space
is smooth over the parameter. Varying $H(0)$ varies gluing units
and preserves these estimates and local equations. On the
compact complement of the critical discs admissibility is open.
The annular gluings extend holomorphically with their fixed
integral classes: the lattice--Lie--torus exact sequence and
Cartan vanishing on an annular product extend the original
translations. This constructs a proper holomorphic family with
all $8m+1$ period directions and the specified central member.

\noindent\textbf{Step 9.} Its Kodaira--Spencer map is injective. A first-order trivialization
would make its map variation a section of
$H^0(X_m,f^*T_B)=H^0(B,T_B)$. The three critical values remain
constant by the equations $w^3,w^4,x_1x_2x_3$, so that base
vector field vanishes at three points and is zero. The
trivialization would then be fibre preserving. On a marked
smooth torus the period derivative is the corresponding
constant Dolbeault tensor in
$\operatorname{Lie}(F)\otimes H^1(F,\mathcal O_F)$.
A nonzero constant tensor cannot be a Dolbeault coboundary,
since derivatives of periodic coefficients integrate to zero.
Thus a period section giving a trivial total-space deformation
vanishes on every regular fibre, and hence identically. Its
Kodaira--Spencer map is an isomorphism onto the already computed
$H^1(T_{X_m})$. The inverse-function argument of \ref{cdp:st:E13} then makes
the Kuranishi ideal zero, proving smoothness, reducedness and
miniversality. The topology theorem of the comb construction
identifies these manifolds with $S^6$; \ref{cdp:st:E1} and \ref{cdp:st:E6} give algebraic
dimension one throughout their small deformation germs.
\end{proof}



\paragraph{Comparison with the constructed parameter space.}
The germ in \cite[Lemma labelled \texttt{paramspace} and
Proposition labelled \texttt{effective}]{Viaclovsky2026Combs}
fixes the values of the coarse period section at all three
critical points. Its normal sheaf is $N(-p_0-p_1-p_2)$, where
$N=(\rho_*F)^*\otimes G_c$. The actual inclusion of that germ
into our root-period germ has the differential induced by
\[
0\longrightarrow N(-p_0-p_1-p_2)\longrightarrow A
\longrightarrow
 \mathbb C_{p_0}^{4}\oplus\mathbb C_{p_1}^{3}
                     \oplus\mathbb C_{p_2}^{2}
\longrightarrow0.
\]
The local graph conditions above prove this exact sequence:
the old sheaf sets all four entries to zero at each point,
whereas $A$ imposes zero, one, and two conditions, respectively.
For the published combs its left term has zero $H^1$, as proved
in the cited smoothing construction. Its tangent dimension is
$8m-8$, and the quotient contributes exactly nine directions.
This explains the relation of the published lower bound to the
dimension $8m+1$ proved here by an explicit morphism of parameter
germs. The source asks for the full Kuranishi germ in
Section~10(1). Searches on 15 November 2026 for its title together
with ``Kuranishi'', ``smooth'', and ``$8m+1$'', and for the
two-parameter and Liu--Rollenske titles with ``tangent cohomology'',
located the supplied construction papers but no matching complete
computations \ref{cdp:st:E13}--\ref{cdp:st:E15}. This records the literature comparison.




\begin{proposition}[Rank divisors from the root Kodaira--Spencer morphism]\label{cdp:st:E16}
\StatementRecord{cdp:st:E16}{proved}{Historical events: Claude Fable 5.2 (Anthropic), effort high, 16 November 2026: carried as background without fault found. GPT-6 Astra (OpenAI), ChatGPT Pro mode, effort Pro, S6\_09\_PAPERS, 17 November 2026: targeted proof check. Also checked in this writing. Pre-continuation writing events: GPT-6 Astra (OpenAI), ChatGPT Pro mode, effort Pro, S6\_10\_PAPERS, 18 November 2026: targeted proof check. Continuation metadata events: GPT-6 Astra (OpenAI), ChatGPT Pro mode, effort Pro, S6\_10\_PAPERS continuation, 19 November 2026: used and rechecked in this run; status-vocabulary reconciliation; metadata reconciliation only; no new mathematical check. \emph{Pre-continuation recorded proof scope.} The root-cover chain rule, weighted entry divisibilities, exact degree budgets at rank-one nodes and the finite root, and factorization y\textasciicircum{}2 P(x\textasciicircum{}4,y\textasciicircum{}3) were re-derived. The wedge determinant of the actual Gauss-Manin derivatives has the stated transition factors and detects the full rank-zero locus through its four entries}{the hypotheses, notation and definitions stated here; \cref{cdp:eq:E.16.3}; \cref{cdp:st:E8}; \cref{cdp:st:E13}; \cref{cdp:st:E14}; \cref{cdp:st:E15}; \cref{cdp:st:E7}}

We use the logarithmic Kodaira--Spencer morphism at the specified degenerate fibres and the ordinary morphism on the finite smooth torus covers. For the original affine-twist family, including every $|p|>1$, its only rank defect is rank one at the order-three point. The logarithmic cusp has rank two. For the $(3k,4l)$ logarithmic transforms, on their full root covers the ranks at the two multiple values are respectively
\[
 \begin{cases}1&k=1,\\0&k>1,\end{cases}
 \qquad
 \begin{cases}2&l=1,\\0&l>1.\end{cases}                 \tag{E.16.1}\label{cdp:eq:E.16.1}
\]
Every other smooth point and the logarithmic cusp have rank two.

For the two-parameter family the logarithmic rank is one at each of its two rank-one nodal values, zero at the order-four root point, and two everywhere else, including the hexagonal fibre. Its determinant has simple coarse zeros at the two nodal values and order four in the root coordinate at the multiple value. For the Liu--Rollenske central family the logarithmic rank is two at both the cuspidal and hexagonal fibres and throughout the remaining smooth locus; at the order-four root point it is zero, and its determinant has order four there.

For a specified positive comb section $\sigma$ of degree parameter $m$, let $x,y$ be root coordinates of weights $3,4$, with $y=0$ at the order-three point. There is a nonzero homogeneous polynomial $P_\sigma(X,Y)$ of degree $4m$, uniquely determined up to an overall nonzero scalar by the marked period section, such that
\[
 \det\kappa_\sigma=y^2P_\sigma(x^4,y^3).             \tag{E.16.2}\label{cdp:eq:E.16.2}
\]
The rank-defect set is exactly its zero divisor. At any such point the rank is zero exactly when all four entries of the actual Kodaira--Spencer matrix vanish; otherwise it is one. Formula \eqref{cdp:eq:E.16.3} below computes those entries from $\sigma$. No unspecified choice of a comb section is assigned a fixed root set. The polynomial is nonzero at the hexagonal cusp, so that cusp always has rank two.
\end{proposition}


\begin{proof}
\noindent\textbf{Step 1.} Adding a linearization constant or changing only the admissible affine lifts does not change the derivative of the marked period matrix. \ref{cdp:st:E8} therefore proves the first assertion. On a root cover of a higher logarithmic transform the original local good-reduction coordinate is $s=u^k$ at the first multiple value and $s=u^l$ at the second. The lifted period map is its actual holomorphic pullback, not a change of frame with a pole. The chain rule gives $k u^{k-1}\kappa(s)$ or $l u^{l-1}\kappa(s)$. \ref{cdp:st:E8} gives rank one at the original order-three centre and rank two at the original order-four centre. The covering derivative is nonzero off its centre, proving \eqref{cdp:eq:E.16.1} and the absence of further defects. The cusp coordinate is unchanged.

\noindent\textbf{Step 2.} For the two-parameter family \ref{cdp:st:E13} proves, with its extension maps,
\[
 E=G=\OO\oplus\OO(-1),\qquad
 \kappa:\OO(-7)\longrightarrow E\otimes G
 \quad\text{on }\PP(1,4).
\]
Thus the four entries have weights $7,6,6,5$. Write $x$ for the weight-one coordinate vanishing at the root point and $y$ for the weight-four coordinate. Every monomial of these respective weights is divisible by $x^3,x^2,x^2,x$. Consequently the entire matrix vanishes at the root point, and its determinant is divisible by $x^4$. The determinant is a section of $\OO(12)$, whose coarse degree is three. At each of the two rank-one Mumford values, the limiting logarithmic derivative is the nonzero rank-one integral monodromy operator in the local torus frames. Its determinant therefore vanishes at least once, whereas its matrix has rank one. These two coarse zeros, together with the root zero of order at least four, account for all of the degree of $\OO(12)$. The determinant is nonzero because the logarithmic rank at the hexagonal cusp is two. All the stated orders are therefore exact, with no remaining zeros. This proves the two-parameter assertion without assuming a theta identity for its varying section.

\noindent\textbf{Step 3.} For the Liu--Rollenske central member \ref{cdp:st:E14} proves the analogous map from $\OO(-3)$ on $\PP(1,4)$. Its entries have weights $3,2,2,1$; these spaces consist respectively of multiples of $x^3,x^2,x^2,x$. The determinant is a section of $\OO(4)$ and is a constant multiple of $x^4$. The constant is nonzero by the hexagonal local monodromy. Hence the only rank defect is the root point, where all four entries vanish. In particular this calculation determines the value of the extended logarithmic boundary at the cusp, even though its fibre is not normal crossing.

\noindent\textbf{Step 4.} For a positive comb, \ref{cdp:st:E15} gives
\[
 \det V=\det G=\OO(24m-1),\qquad
 \kappa_\sigma:\OO(-5)\longrightarrow V\otimes G.
\]
The determinant is a section of $\det V\otimes\det G\otimes\OO(10)=\OO(48m+8)$. On $\PP(3,4)$ every weighted monomial $x^ay^b$ of this degree satisfies $b\equiv2\pmod3$ and $a\equiv0\pmod4$. Factoring $y^2$ therefore gives exactly the polynomial in \eqref{cdp:eq:E.16.2}, of coarse degree $4m$.

\noindent\textbf{Step 5.} Here is the exact computation of that polynomial for any of the period sections in the construction. On a marked local root disc let the two columns of a holomorphic $4\times2$ matrix $U$ frame the Hodge subbundle of the canonical flat extension. Let $\nabla$ be its explicitly marked Gauss--Manin connection, and let $D$ be the local generator of $T_{\mathfrak B}(-p_0)$: $\partial_s$ away from the cusp and $q\partial_q$ at the cusp. Then
\[
 \kappa_\sigma=(\nabla_DU)\bmod\operatorname{im}U,
 \qquad
 \det\kappa_\sigma\ \text{is represented by}\
 \det(U_1,U_2,\nabla_DU_1,\nabla_DU_2).              \tag{E.16.3}\label{cdp:eq:E.16.3}
\]
The first expression uses the quotient frame of $G$ and the dual frame of $V=F^*$; it specifies every matrix entry and therefore also the rank-zero test. Changing the Hodge frame multiplies the determinant by the square of its determinant, exactly its transition in $\det V\otimes\det G\otimes T_{\mathfrak B}(-p_0)^{-2}$. Changing a root parameter gives the inverse square of the tangent-frame change. Thus the local calculations glue to the stated global section and determine $P_\sigma$ from the actual period data. In the cusp comparison of \ref{cdp:st:E7} and \ref{cdp:st:E15}, the logarithmic derivative has the fixed invertible integral leading matrix; its determinant is nonzero. This proves both the nonzero-polynomial assertion and the exact rank-divisor description.
\end{proof}






\begin{proposition}[Counterexample: The embedded-component omission in the divisor-twisting lemma]\label{cdp:st:E17}
\StatementRecord{cdp:st:E17}{proved}{Historical events: Claude Fable 5.2 (Anthropic), effort high, 16 November 2026: carried as background without fault found. GPT-6 Astra (OpenAI), ChatGPT Pro mode, effort Pro, S6\_09\_PAPERS, 17 November 2026: targeted proof check. Also checked in this writing. Pre-continuation writing events: GPT-6 Astra (OpenAI), ChatGPT Pro mode, effort Pro, S6\_10\_PAPERS, 18 November 2026: targeted proof check. Continuation metadata events: GPT-6 Astra (OpenAI), ChatGPT Pro mode, effort Pro, S6\_10\_PAPERS continuation, 19 November 2026: used and rechecked in this run; status-vocabulary reconciliation; metadata reconciliation only; no new mathematical check. \emph{Pre-continuation recorded proof scope.} The square-zero embedded-point example and the corrected stable support formula by Krull intersection were re-derived}{the hypotheses, notation and definitions stated here; \cref{cdp:st:E2}}

CDP98, Proposition~1.1, in its printed generality of arbitrary
nonreduced compact spaces, is false. The exact replacement is as
follows. For an effective Cartier divisor $D$ on a compact complex
space $Z$ and a locally free sheaf $\mathcal V$, the spaces
$H^0(Z,\mathcal V(-kD))$, viewed inside $H^0(Z,\mathcal V)$,
stabilize for $k\gg0$ at
\[
 \{s\in H^0(Z,\mathcal V):\operatorname{supp}(s)\cap|D|
       =\varnothing\}.
\]
In particular they eventually vanish when $Z$ has no embedded
associated components and $D$ meets every irreducible component.
This last case applies to every full Cartier fibre used in \ref{cdp:st:E2}.
\end{proposition}


\begin{proof}
\noindent\textbf{Step 1.} Let $C=\mathbb P^1$, let $p,q$ be distinct points, and form the
square-zero complex space
\[
 |Z|=C,\qquad\mathcal O_Z=\mathcal O_C\oplus\mathbb C_p,
 \qquad(a,u)(b,v)=(ab,av+bu).
\]
Near $p$ its analytic local ring is
$\mathbb C\{z,e\}/(ze,e^2)$, and away from $p$ it is the
ordinary smooth curve. These charts glue through the quotient
by $e$, so this specifies a compact connected complex space.
It has one irreducible component, of positive dimension, and an
embedded associated point at $p$. The divisor $D=q$ is effective
Cartier: its local equation is a coordinate at $q$ and a unit
near $p$. For every integer $k$ there is the exact, split
description as an $\mathcal O_C$-module
\[
 \mathcal O_Z(-kD)=\mathcal O_C(-kq)\oplus\mathbb C_p.
\]
It gives every cohomology group:
\[
 h^0=\max(1-k,0)+1,\qquad h^1=\max(k-1,0),\qquad
 h^j=0\quad(j\geq2).
\]
In particular the nonzero section of the square-zero summand
survives for every $k\geq1$, refuting the printed proposition.

\noindent\textbf{Step 2.} For the replacement, the descending spaces in the statement are
finite dimensional, hence stabilize. A section in their
intersection belongs to $\bigcap_k I_D^k\mathcal V_z=0$ at
every $z\in|D|$, by Krull intersection in the analytic local
ring. Its coherent support is therefore disjoint from $D$.
Conversely a section with support disjoint from $D$ vanishes
near $D$ and is a section of $\mathcal V(-kD)$ for every $k$.
This proves the exact stable-space formula. When there are no
embedded associated components, every nonzero submodule of a
locally free sheaf has support containing an irreducible
component of $Z$. Such a support cannot avoid $D$ under the
stated intersection hypothesis, proving the vanishing.
\end{proof}



\paragraph{Audit of the printed non-holomorphic proof.}
We now give its entire chain of reductions and identify its
dependencies. CDP20, Section~3, first uses constancy of the index
under Picard twists and Riemann--Roch to deduce its parts (b)
and (c) from generic vanishing. Serre duality reduces generic
vanishing to degrees zero and two. Degree zero invokes CDP98,
Corollary~1.3; that corollary uses only Proposition~1.1 on the
smooth manifold, a multiple of an effective divisor lying in
$\operatorname{Pic}^0$, and semicontinuity. \ref{cdp:st:E17} proves the
version of Proposition~1.1 needed there, and \ref{cdp:st:E2} spells out the
exponential-family semicontinuity argument. Thus this dependency
is valid and involves none of the cotangent or monodromy steps.

For degree two the printed proof resolves the meromorphic pencil
by $\sigma:\widehat Y\to Y$, with a connected map
$h:\widehat Y\to\mathbb P^1$, and writes
\[
 h^*\mathcal O(1)=\sigma^*P\otimes\mathcal O(-E),
\]
where $E$ is the exceptional divisor of the pulled-back base
ideal. The map on second cohomology from $Y$ to $\widehat Y$
is an isomorphism, since a composition of smooth blowups has
$R\sigma_*\mathcal O=\mathcal O_Y$. The divisor sequence
reduces vanishing on $\widehat Y$ to vanishing after $-tE$
and to $H^2(tE,\sigma^*W)=0$. The printed proof asserts
the exceptional higher-image vanishing needed for this step.
For this particular pencil divisor it has the following direct
justification. Pull back the evaluation sequence on $\mathbb P^1$
and tensor by $\sigma^*P^{-t}$ to obtain, for $t\geq1$,
\[
\begin{split}
0\longrightarrow&\bigl(\mathcal O(E)\otimes
              \sigma^*P^{-(t+1)}\bigr)^{\oplus t}
 \longrightarrow(\sigma^*P^{-t})^{\oplus(t+1)}\\
 &\longrightarrow\mathcal O(-tE)\longrightarrow0.
\end{split}
\]
The evaluation sequence remains exact under pullback because it
is a sequence of vector bundles with locally free quotient.
The middle term has zero $R^2\sigma_*$, by the projection
formula and $R\sigma_*\mathcal O=\mathcal O_Y$.
Every coherent sheaf has zero $R^3\sigma_*$ because the
fibres have dimension at most two. Hence
$R^2\sigma_*\mathcal O(-tE)=0$; the divisor sequence now gives
\[
 R^1\sigma_*\mathcal O_{tE}=
 R^2\sigma_*\mathcal O_{tE}=0.
\]
The support of $\sigma_*\mathcal O_{tE}$ has dimension at most
one, so Leray and the projection formula give the required
$H^2(tE,\sigma^*W)=0$. No identification of this pushforward
with a structure sheaf is needed.

This special verification matters. The analogous assertion for
an arbitrary exceptional divisor is false. Blow up a point
$p\in\mathbb P^3$ and then a smooth plane cubic in the exceptional
$\mathbb P^2$. For the second exceptional divisor $S$, projection
to that cubic is a $\mathbb P^1$-bundle; thus
$h^\bullet(S,\mathcal O_S)=(1,1,0)$ and the composite contraction
to $\mathbb P^3$ has
$R^1\sigma_*\mathcal O_S=\mathbb C_p$.
This example, also given in \cite{LRS2019}, concerns arbitrary
exceptional divisors and does not refute the just-proved assertion
for the divisor of a resolved pencil. The graph proof of \ref{cdp:st:E2}
avoids the distinction by computing its own higher direct images
from its hypersurface equation.

The printed argument then replaces its twist $M$ by
$M\otimes P^{t+k}$ and applies Leray to $h$. It must prove
\[
 R^2h_*\sigma^*(W\otimes M\otimes P^k)=0
\]
and then the vanishing of $H^1$ of the corresponding $R^1$
twisted by $\mathcal O(t)$; the latter is ordinary Serre
vanishing on $\mathbb P^1$. For the former it uses top-degree
base change and fibrewise Serre duality. It invokes CDP98,
Proposition~1.1, first on components meeting $E$, and then
successively on components at increasing graph distance from
them. Coefficients in an auxiliary effective divisor are chosen
in reverse order of this distance, and the power of $E$ is
chosen last. \ref{cdp:st:E2} gives this argument with the actual pullback
maps, full multiplicity filtrations, and a finite-support
argument establishing uniformity. These fibres are Cartier in
a smooth threefold and have no embedded components, so \ref{cdp:st:E17}'s
proved version of the twisting lemma applies.

These are the only earlier CDP results invoked in Section~3:
CDP98, Corollary~1.3 and Proposition~1.1. Its opening reference
to CDP98, Theorem~2.1 identifies the proof being rewritten,
rather than using that theorem as a premise. CDP98, Lemma~1.5
is not invoked; its erroneous replacement of fibrewise vanishing
by generic vanishing is precisely what the full component
argument removes. No result from CDP20, Sections~4--7 is used.
In particular the cotangent-filtration and monodromy failures
treated in Parts A and D do not enter \ref{cdp:st:E2}. The supplied LRS
graph proof has the same dependency on the no-embedded-component
twisting lemma, and supplies its own higher-image calculation.
Its displayed adjunction formula has a sign misprint:
for a Cartier fibre $F\subset Y$ it is
$\omega_F=(K_Y+F)|_F$, not $(K_Y-F)|_F$.
\ref{cdp:st:E2} uses the actual dualizing line and includes nonreduced fibres,
so that misprint has no role in the proof here.




\begin{proposition}[Counterexample: The square theta formula for arbitrary sections]\label{cdp:st:E18}
\StatementRecord{cdp:st:E18}{proved}{Historical events: Claude Fable 5.2 (Anthropic), effort high, 16 November 2026: carried as background without fault found. GPT-6 Astra (OpenAI), ChatGPT Pro mode, effort Pro, S6\_09\_PAPERS, 17 November 2026: targeted proof check. Also checked in this writing. Pre-continuation writing events: GPT-6 Astra (OpenAI), ChatGPT Pro mode, effort Pro, S6\_10\_PAPERS, 18 November 2026: targeted proof check; used and rechecked in this run. Continuation mathematical and source checks: GPT-6 Astra (OpenAI), ChatGPT Pro mode, effort Pro, S6\_10\_PAPERS continuation, 19 November 2026: computation or formal certificate; targeted proof check; used and rechecked in this run. Continuation metadata events: GPT-6 Astra (OpenAI), ChatGPT Pro mode, effort Pro, S6\_10\_PAPERS continuation, 19 November 2026: used and rechecked in this run; status-vocabulary reconciliation; metadata reconciliation only; no new mathematical check. \emph{Pre-continuation recorded proof scope.} Matched the primary NIST formulas with the exact convention vartheta\_j(z,tau)=theta\_j(pi*z|tau), nome exp(pi*i*tau), and Weierstrass periods 1,tau. The theta and eta products give the positive logarithm and a full lattice; the theta quotient has exactly the multipliers 1,-1 and divisor P-O. Adjoining (1/2,v) has exact index six because comparison of imaginary parts forces a multiple of six, giving the identity-induced isogeny and the same Kodaira-Spencer rank. The half-period shift makes h=0 while the heat/product calculation and e1 formula give determinant -(theta3\textasciicircum{}4+theta4\textasciicircum{}4)/12, nonzero at i. The general correction follows by differentiating the convergent products and matching the elliptic principal part, including every pi factor. The actual tangent direct-image long exact sequence has boundary ranks one and two, giving ranks 2,3,1. All source and normalization gates in the former partial review are closed}{the hypotheses, notation and definitions stated here; \cref{cdp:eq:E-theta-general}; \cref{cdp:st:E5}}

The square-only formula in Lead~9 is false as a local period
identity for arbitrary elliptic sections and linearizations.
The exact formula for the theta expression specified there is
\eqref{cdp:eq:E-theta-general}. The assertion for the actual original
period family remains true, by its equation $\wp(\mu)=0$.
\end{proposition}


\begin{proof}
\noindent\textbf{Step 1.} \noindent\textbf{Step 1.} We use $\vartheta_j(z,\tau)=\theta_j(\pi z\mid\tau)$ in the
notation of \cite[eqs.~20.2.1--2, 20.2.6 and 20.2.12]{DLMFTheta2026}.
Thus the nome is $q=e^{\pi i\tau}$, and the Weierstrass periods
are $2\omega_1=1$, $2\omega_3=\tau$. The products below are
\cite[eqs.~20.4.3--5 and 23.17.8]{DLMFTheta2026}, with this
specified change of theta argument.

\noindent\textbf{Step 2.} Let $\Delta$ be a sufficiently small disc about $i$ in the upper
half-plane and choose the holomorphic logarithm of the nonvanishing
function $\vartheta_1(1/2,\tau)/\eta(\tau)$ which is real at $i$.
Put
\[
 \mu(\tau)=\frac12,\qquad
 v(\tau)=\frac{i}{\pi}
       \log\frac{\vartheta_1(1/2,\tau)}{\eta(\tau)}+i,
 \qquad \beta(\tau)=6v(\tau).
\]
At $\tau=i$, the product expansions give
\[
 \frac{\vartheta_1(1/2,i)}{\eta(i)}
 =2e^{-\pi/6}\prod_{n\geq1}(1+e^{-2\pi n})^2>1.
\]
Thus $\operatorname{Im}v(i)>1$, and after shrinking $\Delta$
one has $\operatorname{Im}v>0$. Define
\[
 \Lambda_\tau=\mathbb Z(1,0)+\mathbb Z(0,1)
 +\mathbb Z(\tau,1/2)+\mathbb Z(3,\beta(\tau)),
 \qquad T_\tau=\mathbb C^2/\Lambda_\tau.
\]
The real determinant of these four periods is nonzero, since
$\operatorname{Im}\tau>0$ and $\operatorname{Im}\beta>0$.
They form a locally uniformly full lattice, so the quotients form
a proper holomorphic family of compact complex two-tori over
$\Delta$. Each fibre has
$h^\bullet(\mathcal O)=(1,2,1)$ and
$h^\bullet(T)=(2,4,2)$, with zero higher groups, by its
translation-invariant Dolbeault complex.

\noindent\textbf{Step 3.} \noindent\textbf{Step 2.} This family also has the required elliptic line-bundle
description. Quotient first by $(1,0),(0,1),(\tau,1/2)$.
Writing $z$ for the first coordinate and $a=\exp(2\pi iz_2)$
for the second gives the complement of zero in the flat line
bundle on $\mathbb C/(\mathbb Z+\tau\mathbb Z)$ with multipliers
$1$ and $-1$. For $P=1/2$, the actual meromorphic section
\[
 r(z)=\frac{\vartheta_1(z-P,\tau)}{\vartheta_1(z,\tau)}
\]
has divisor $P-O$ and satisfies $r(z+1)=r(z)$ and
$r(z+\tau)=-r(z)$. These equations identify the displayed
quotient line bundle with $\mathcal O(P-O)$. The last lattice
generator acts by
$(z,a)\mapsto(z+3,\exp(2\pi i\beta)a)$.
Its elliptic translation is $-6P=0$ in the elliptic curve;
its scalar is a holomorphic contracting linearization. This
constructs the family from an elliptic family, a section and a
linearization explicitly.

\noindent\textbf{Step 4.} \noindent\textbf{Step 3.} There is a specified relation to the symmetric period matrix
$\Omega=\left(\begin{smallmatrix}\tau&\mu\\\mu&v\end{smallmatrix}\right)$.
Adjoining $(\mu,v)$ enlarges $\Lambda_\tau$ to the lattice
generated by $(1,0),(0,1),(\tau,\mu),(\mu,v)$; its index is six,
because $6(\mu,v)=(3,\beta)$ and no smaller multiple belongs
to the original real lattice. Identity on $\mathbb C^2$ induces
a finite holomorphic isogeny of degree six between these families.
Its differential is the identity on the Lie spaces and its
pullback on complex first cohomology is an isomorphism. The
Kodaira--Spencer tensors therefore have the same rank under
these explicitly given comparisons.

\noindent\textbf{Step 5.} \noindent\textbf{Step 4.} The identity $\vartheta_1(1/2+z,\tau)=\vartheta_2(z,\tau)$
and evenness of $\vartheta_2$ give
$\partial_\mu\log\vartheta_1(1/2,\tau)=0$. Since $\mu'=0$,
the proposed square is zero. The actual determinant, by the
proved heat-equation formula, is
\cite[eqs.~20.7.5, 20.13.1, 23.3.9 and 23.6.2]{DLMFTheta2026}
\[
 \det\Omega'(i)=-\frac{\wp(1/2;i)}{4\pi^2}
 =-\frac{\vartheta_3(0,i)^4+\vartheta_4(0,i)^4}{12}\ne0.
\]
The last inequality follows from the theta products at the real
nome $e^{-\pi}$, whose factors give strictly positive theta
constants. Thus even the zero sets in the asserted formula
differ in this fully specified family. Extra holomorphic terms
in a different linearization contribute their derivatives to
the determinant; they must be retained, just as the $\wp$ term
must be retained when the section is not a zero of $\wp$.

\noindent\textbf{Step 6.} \noindent\textbf{Step 5.} For completeness, the direct images and total-space cohomology
in this local example are also determined. Write
$g:\mathscr T\to\Delta$ for the displayed proper family.
Translation-invariant forms, base change, and triviality of
holomorphic vector bundles on a disc give
\[
 R^qg_*\mathcal O_{\mathscr T}
   =\mathcal O_\Delta^{\binom2q},\qquad
 R^qg_*T_{\mathscr T/\Delta}
   =\mathcal O_\Delta^{2\binom2q}\quad(0\leq q\leq2),
\]
with all higher images zero. Shrink once more so that the
nonzero determinant stays nonzero. In the direct-image sequence
of $0\to T_{\mathscr T/\Delta}\to T_{\mathscr T}
\to g^*T_\Delta\to0$, the first boundary sends $1$ to the
nowhere-zero tensor $\kappa$, and the next boundary is the
rank-two isomorphism of \ref{cdp:st:E5}. Thus
\[
 R^qg_*T_{\mathscr T}=
 \begin{cases}
  \mathcal O_\Delta^2&q=0,\\
  \mathcal O_\Delta^3&q=1,\\
  \mathcal O_\Delta&q=2,\\
  0&q>2.
 \end{cases}
\]
\noindent\textbf{Step 6.} Since $\Delta$ is Stein, every total-space group
$H^q(\mathscr T,\mathcal F)$ for these three sheaves is exactly
the space of sections on $\Delta$ of the corresponding displayed
direct image. This specifies every cohomological degree without
assigning finite dimensions to the holomorphic-function spaces
on a disc.
\end{proof}





\section{Connected deformation chains and two independent periods}\label{sec:II:push-stability}
\begin{corollary}[arbitrary connected deformation chains]\label{push:src:E.2.4}
\StatementRecord{push:src:E.2.4}{proved}{Historical events: Claude Fable 5.2 (Anthropic), effort high, 16 November 2026: carried as background without fault found. GPT-6 Astra (OpenAI), ChatGPT Pro mode, effort Pro, S6\_09\_PAPERS, 17 November 2026: targeted proof check. Also checked in this writing. Pre-continuation writing events: GPT-6 Astra (OpenAI), ChatGPT Pro mode, effort Pro, S6\_10\_PAPERS, 18 November 2026: targeted proof check. Continuation metadata events: GPT-6 Astra (OpenAI), ChatGPT Pro mode, effort Pro, S6\_10\_PAPERS continuation, 19 November 2026: used and rechecked in this run; status-vocabulary reconciliation; metadata reconciliation only; no new mathematical check. \emph{Pre-continuation recorded proof scope.} The canonical-pencil, full-punctured-disc, generic-disc, finite-projection, and Kuranishi closure steps were checked in their stated order. This event concerns the assembly; each substantive proof has a separate record}{the hypotheses, notation and definitions stated here; \cref{push:src:E.2.5}; \cref{push:src:E.2.24}; \cref{push:src:E.2.21}; \cref{push:src:E.2.22}; \cref{push:src:E.2.23}}

The ordered coprime pair $(3k,4l)$ is constant in every
proper smooth holomorphic family over a connected complex
analytic base which contains a constructed sphere member.
Two constructed members are deformation equivalent through
such families, allowing arbitrary intermediate members and
finite chains, if and only if their pairs agree.
Moreover different pairs cannot occur in one continuous
family in the $C^\infty$ space of integrable complex
structures on $S^6$.
\end{corollary}


\begin{proof}
\noindent\textbf{Step 1.} The detailed proof is \ref{push:src:E.2.5}--\ref{push:src:E.2.24} below.  The intrinsic
pencil persists for either canonical slope, the limiting
period stays in the admissible interior, and the two finite
root charts and the actual analytic cusp are recovered.
\ref{push:src:E.2.21} then identifies every full punctured-disc smooth
completion with the prescribed member.  The additional
step \ref{push:src:E.2.22} propagates from any nonempty known open subset
of a disc: finitely many successive discrete analytic
exceptional sets are filled in reverse order.  \ref{push:src:E.2.23} uses
a finite local projection and normalization of a curve
through two specified points to pass to arbitrary analytic
bases and deformation chains.
The classical $C^\infty$ Kuranishi local theorem then
gives the continuous-family conclusion in \ref{push:src:E.2.24}.

\noindent\textbf{Step 2.} Local miniversality alone would prove only openness.
Nor does a boundary point of an arbitrary open locus
necessarily admit a holomorphic disc whose punctured
image lies in that locus.  The generic-disc propagation
and finite-projection steps are the substantive additions
which establish the required closedness.
\end{proof}



\subsubsection*{Smooth specializations of the canonical pencil}

\begin{proposition}[a fractional canonical pencil at a smooth limit]\label{push:src:E.2.5}
\StatementRecord{push:src:E.2.5}{proved}{Historical events: Claude Fable 5.2 (Anthropic), effort high, 16 November 2026: carried as background without fault found. GPT-6 Astra (OpenAI), ChatGPT Pro mode, effort Pro, S6\_09\_PAPERS, 17 November 2026: targeted proof check. Also checked in this writing. Pre-continuation writing events: GPT-6 Astra (OpenAI), ChatGPT Pro mode, effort Pro, S6\_10\_PAPERS, 18 November 2026: targeted proof check. Continuation metadata events: GPT-6 Astra (OpenAI), ChatGPT Pro mode, effort Pro, S6\_10\_PAPERS continuation, 19 November 2026: used and rechecked in this run; status-vocabulary reconciliation; metadata reconciliation only; no new mathematical check. \emph{Pre-continuation recorded proof scope.} The unique fractional canonical line in the divisible Picard group, mod-parameter injection of sections, CDP-LRS meromorphic-map obstruction, connected algebraic reduction, and actual common Cartier divisor of the evaluated sections were re-derived}{the hypotheses, notation and definitions stated here; \citet[Theorem~2.1]{CDP2020D}; \citet[Theorem~3.1 and Corollary~3.2]{LRS2019}}

Let $\pi:\mathcal X\to\Delta$ be a proper holomorphic submersion
whose fibres are diffeomorphic to $S^6$.  Suppose that for every
$t\ne0$ the fibre is one of the members of \cref{push:src:E.2.1}, with the same
coprime multiplicities $(3k,4l)$.  Set
\[
 A=12kl,\qquad D=kl-3k-4l\ne0.
\]
There is a unique holomorphic line bundle $\mathcal L$ on
$\mathcal X$, up to isomorphism, satisfying
\[
 \mathcal L^{D}=K_{\mathcal X/\Delta}^{A}.
\]
Its restriction to $X_t$ for $t\ne0$ is the coarse pencil
$f_t^*\mathcal O_{\mathbb P^1}(1)$.  The coherent sheaf
$\pi_*\mathcal L$ is locally free of rank two, and its central
fibre injects into $H^0(X_0,\mathcal L_0)$.

The central manifold has algebraic dimension one and a
holomorphic algebraic reduction $g_0:X_0\to\mathbb P^1$ with
connected fibres.  If $s_0,s_1$ are the two central sections
obtained from a local frame of $\pi_*\mathcal L$, there are a
nonconstant rational map $R:\mathbb P^1\to\mathbb P^1$ of
degree $e\ge1$ and an effective Cartier divisor $E_0$ such that
\[
 [s_0:s_1]=R\circ g_0,\qquad
 \mathcal L_0\simeq g_0^*\mathcal O(e)\otimes\mathcal O(E_0).
\]
Here $[s_0:s_1]$ denotes the holomorphic extension of the ratio;
the displayed divisor is its common fixed divisor.  The
statement does not assert that $E_0=0$.
\end{proposition}


\begin{proof}
\noindent\textbf{Step 1.} By Ehresmann's theorem the smooth bundle over the contractible
disc is differentiably trivial.  Consequently
$H^1(\mathcal X,\mathbb Z)=H^2(\mathcal X,\mathbb Z)=0$.
The exponential sequence identifies
\[
 \operatorname{Pic}(\mathcal X)=H^1(\mathcal X,\mathcal O_{\mathcal X})
\]
as groups.  This is a complex vector space, so multiplication by
the nonzero integer $D$ is an automorphism.  This proves the
existence and uniqueness of the \emph{relative} root, including
when $D<0$.  The same argument on each fibre makes its Picard
group uniquely divisible.  \cref{push:src:E.2.1} gives
$K_{X_t}^{A}=(f_t^*\mathcal O(1))^{D}$, so uniqueness identifies
$\mathcal L_t$ with this actual pencil, not just with a numerical
class.

\noindent\textbf{Step 2.} Proper direct-image coherence gives $\mathcal V=\pi_*\mathcal L$.
It is torsion-free: a local section killed by a nonzero function
of $t$ vanishes on the inverse image of an open set, hence
vanishes as a section of the invertible sheaf on the smooth
total space.  A coherent torsion-free sheaf on a smooth curve
is locally free.  Its rank is two because the noncentral
algebraic reductions have connected fibres.  Applying $\pi_*$
to
\[
 0\longrightarrow\mathcal L\xrightarrow{\,t\,}\mathcal L
 \longrightarrow i_*\mathcal L_0\longrightarrow0
\]
shows directly that $\mathcal V/t\mathcal V$ injects into
$H^0(X_0,\mathcal L_0)$.  Thus the two central sections are
linearly independent even if central cohomology jumps; no
surjectivity of base change has been assumed.

\noindent\textbf{Step 3.} We use precisely the non-holomorphic-meromorphic-map theorem
of \citet[Theorem~2.1]{CDP2020D}, also written out in
\citet[Theorem~3.1 and Corollary~3.2]{LRS2019}.
For a compact complex threefold with $b_2=0$, a nonconstant
meromorphic map to a curve which fails to be holomorphic
would force $c_3=0$.  Since $c_3(X_0)[X_0]=2$, every
meromorphic function on $X_0$ extends as a holomorphic map
to $\mathbb P^1$.  This argument uses neither the different
holomorphic-fibration assertion of CDP, Theorem~2.2, nor its
Corollary~2.3.  Those assertions cannot be applied here.

\noindent\textbf{Step 4.} In particular the ratio of $s_0,s_1$ gives a nonconstant
holomorphic map $h_0:X_0\to\mathbb P^1$.  Also $a(X_0)\le1$:
otherwise two algebraically independent meromorphic functions
would give a holomorphic map $X_0\to\mathbb P^1\times\mathbb P^1$
with two-dimensional image.  At a point of differential rank
two, after choosing finite affine target coordinates $u,v$
centred at its image, the global meromorphic function $u/v$
would have the local indeterminacy of the quotient of two
independent coordinates.  It cannot extend holomorphically
there.  This contradicts the preceding theorem.  Hence
$a(X_0)=1$.

\noindent\textbf{Step 5.} Take the Stein factorization $h_0=R\circ g_0$.  Its intermediate
space is a smooth compact curve.  Connectedness of the fibres
injects its $H^1(\ ,\mathbb Z)$ into $H^1(X_0,\mathbb Z)=0$
by the degree-one Leray sequence, so it is $\mathbb P^1$.
Every meromorphic function on $X_0$ factors rationally through
$g_0$: the image of its holomorphic map together with $g_0$
is an analytic curve in $\mathbb P^1\times\mathbb P^1$;
normalizing this curve gives a finite factor of $g_0$, whose
degree is one because $g_0$ has connected fibres.  Thus $g_0$
is the algebraic reduction, and $R$ has some degree $e\ge1$.

\noindent\textbf{Step 6.} Finally, on a target chart where the first homogeneous
coordinate of $h_0$ is nonzero, $s_1=h_0s_0$ in a local
trivialization, and the common ideal is generated by $s_0$.
The analogous assertion holds on the other chart.  These local
generators define a nonzero section of
$\mathcal L_0\otimes h_0^*\mathcal O(-1)$ and an effective
Cartier zero divisor $E_0$.  This gives the stated factorization
of $\mathcal L_0$ with no residual codimension-two base locus.
\end{proof}



\begin{proposition}[positive canonical slope preserves the fibration]\label{push:src:E.2.6}
\StatementRecord{push:src:E.2.6}{proved}{Historical events: Claude Fable 5.2 (Anthropic), effort high, 16 November 2026: carried as background without fault found. GPT-6 Astra (OpenAI), ChatGPT Pro mode, effort Pro, S6\_09\_PAPERS, 17 November 2026: targeted proof check. Also checked in this writing. Pre-continuation writing events: GPT-6 Astra (OpenAI), ChatGPT Pro mode, effort Pro, S6\_10\_PAPERS, 18 November 2026: targeted proof check. Continuation metadata events: GPT-6 Astra (OpenAI), ChatGPT Pro mode, effort Pro, S6\_10\_PAPERS continuation, 19 November 2026: used and rechecked in this run; status-vocabulary reconciliation; metadata reconciliation only; no new mathematical check. \emph{Pre-continuation recorded proof scope.} The positive-slope canonical restriction, exclusion of positive Kodaira dimension by actual meromorphic functions, specialization of connected pencil divisors, elimination of the vertical common divisor, and contradiction from opposite canonical slopes were re-derived}{the hypotheses, notation and definitions stated here; \cref{push:src:E.2.5}; \cref{push:src:E.2.22}; \cref{push:src:E.2.24}; \citet[Proposition~3.1.2]{AbashevaDeev2024AB}}

Under the hypotheses of \ref{push:src:E.2.5}, suppose $D>0$.  Then $E_0=0$
and $e=1$.  After shrinking the disc and choosing a frame of
$\pi_*\mathcal L$, its evaluation defines a proper flat
holomorphic map
\[
 g:\mathcal X\longrightarrow\mathbb P^1\times\Delta
\]
which is the intrinsic algebraic reduction on every fibre.
Its smooth general central fibre is a complex two-torus.

In particular, two families of the type in \ref{push:src:E.2.5} with positive
canonical slope cannot have the same smooth central manifold
unless their multiplicity pairs agree.  A positive-slope and
a negative-slope family cannot have a common smooth central
manifold either.  This proposition concerns a common smooth
specialization.  The additional propagation and closure
arguments for arbitrary connected chains are \ref{push:src:E.2.22}--\ref{push:src:E.2.24}.
\end{proposition}


\begin{proof}
\noindent\textbf{Step 1.} Choose a very general smooth fibre $F$ of $g_0$, avoiding the
images of all vertical components of $E_0$ and the jumping
sets for the pluricanonical direct images below.  Its normal
bundle in $X_0$ is trivial.  Adjunction gives
$K_F=K_{X_0}|_F$, and $c_1(K_F)=0$ integrally because
$H^2(X_0,\mathbb Z)=0$.  In particular $F$ is minimal: a
$(-1)$-curve would have canonical intersection $-1$.
Restriction of the root identity in \ref{push:src:E.2.5} gives
\[
 K_F^{A}\simeq\mathcal O_F(D E_0|_F).
\]
The right hand side is effective since $D>0$, so
$\kappa(F)\ge0$.

\noindent\textbf{Step 2.} We next show $\kappa(F)\le0$, without assuming that the fibres
of an algebraic reduction have algebraic dimension zero.
If $\kappa(F)>0$, some $r>0$ has $h^0(F,K_F^r)\ge2$.
Our very-general choice ensures that this is the generic rank
of the coherent sheaf $g_{0*}K_{X_0}^r$.  A sufficiently
positive twist of this sheaf on $\mathbb P^1$ is globally
generated.  Two of its global sections can therefore be
chosen to be linearly independent on a general fibre.  By
projection formula they are sections of
$K_{X_0}^r\otimes g_0^*\mathcal O(N)$; their ratio is a
meromorphic function on $X_0$ which is nonconstant on $F$.
This contradicts the algebraic-reduction assertion of \ref{push:src:E.2.5}.
The choice of a very general fibre is legitimate: for each
$r$ the exceptional set for the generic cohomology rank is
a proper analytic subset of the base curve, and there are
only countably many $r$.

\noindent\textbf{Step 3.} Thus $\kappa(F)=0$.  The classification of minimal compact
complex surfaces makes $K_F$ torsion; the relevant consequence
is also explicitly recorded in the proof of
\citet[Proposition~3.1.2]{AbashevaDeev2024AB}.
An effective divisor whose line bundle is torsion is zero:
raising its canonical section to a torsion-killing power
gives a holomorphic function on a compact connected manifold,
and that function is a nonzero constant.  Therefore
$E_0|_F=0$.  Every horizontal component of $E_0$ would meet
such a general fibre in a nonempty effective divisor.
Consequently $E_0$ has only vertical components, supported
over finitely many points of $\mathbb P^1$.

\noindent\textbf{Step 4.} We now use connectedness of actual divisor limits to remove
both this vertical fixed part and the possible degree $e$.
Choose a linear combination $s=as_0+bs_1$ for which its
moving central divisor is the union of the $e$ distinct
smooth fibres over $R^{-1}([-b:a])$, all disjoint from the
support of $E_0$.  Also choose it to be a regular fibre on
one specified noncentral member.  Both requirements omit
only finitely many choices in the projective line of
combinations.  Extend it by the corresponding combination
of the frame of $\pi_*\mathcal L$.

\noindent\textbf{Step 5.} Its total zero divisor $\mathcal D\subset\mathcal X$ is a
relative effective Cartier divisor: the section is not
identically zero on any fibre, by independence of the two
sections, including at zero.  In local coordinates, $t$ is
a non-zero-divisor modulo its defining equation.  Hence
$\mathcal D\to\Delta$ is proper and flat.  Its general
noncentral fibre is a smooth connected torus.  There are no
vertical irreducible components of $\mathcal D$, since such
a component would have to be the whole irreducible
threefold $X_t$.  It follows that $\mathcal D$ is reduced
and irreducible: its unique generic component has multiplicity
one, and a Cartier divisor in a smooth manifold has no
embedded components.  Normalize $\mathcal D$ and take the
Stein factorization over $\Delta$.  The finite factor has
degree one over the punctured generic locus; normality of
the disc makes it the disc itself.  Its central fibre, and
therefore the image $\mathcal D_0$, is connected.

\noindent\textbf{Step 6.} But the support of
\[
 \mathcal D_0=E_0+g_0^*R^*([-b:a])
\]
is a disjoint union of the $e$ chosen smooth connected fibres
and, if nonempty, the support of $E_0$ over other base points.
Connectedness forces $E_0=0$ and $e=1$.

\noindent\textbf{Step 7.} The central two sections now generate $\mathcal L_0$ at
every point.  The noncentral two sections already generate
their pencils, so the evaluation
$\pi^*\pi_*\mathcal L\to\mathcal L$ is surjective after
shrinking.  It defines the asserted map $g$.  It is proper
because $\pi$ is proper.  On each $X_t$ the nonconstant
map to the smooth curve $\mathbb P^1$ is flat: its local
structure sheaf is torsion-free over the local discrete
valuation ring of the curve.  The total source is flat
over $t$, so the local criterion for flatness applied to
the parameter $t-t_0$ gives flatness over $(z,t)$.
Equivalently, the local fibre equations form a regular
sequence in the smooth source.  At a regular
point of the central base curve, properness makes $g$ a
proper submersion on a neighbourhood of the corresponding
compact fibre.  Ehresmann's theorem identifies that central
surface differentiably with nearby complex two-tori.
The classification of compact complex surfaces, or its
special case for the differentiable four-torus, then makes
the central surface a complex two-torus as well.

\noindent\textbf{Step 8.} Suppose now two positive-slope families have the same
central manifold $X_0$.  Its algebraic reduction is intrinsic
up to $\operatorname{PGL}_2$, so the preceding result makes
both central pencil lines the same non-torsion line
$L=g_0^*\mathcal O(1)$.  With indices $i=1,2$ the identities
$K_{X_0}^{A_i}=L^{D_i}$ give
$L^{D_1A_2-D_2A_1}=\mathcal O$.  The line $L$ is non-torsion
because a positive power has a nonzero section with zeros.
Thus $D_1/A_1=D_2/A_2$, and the exact recovery of the ordered
pair from this slope in \cref{push:src:E.2.3} gives the assertion.

\noindent\textbf{Step 9.} Finally, a common smooth limit from opposite signs would
have a nonzero section of a positive power of $K_{X_0}$
and of a negative power, by upper semicontinuity (or by
the same injective direct-image argument).  Raising them
to a common power, their nonzero product is a constant;
both are nowhere zero.  Hence $K_{X_0}$ is torsion.
The exponential sequence and $H^1(X_0,\mathbb Z)=
H^2(X_0,\mathbb Z)=0$ make $\operatorname{Pic}(X_0)$
torsion-free, so $K_{X_0}$ would be trivial.  Its nonzero
holomorphic three-form $\Omega$ is closed and exact because
$H^3(X_0,\mathbb C)=0$.  Stokes' theorem would give
$\int_{X_0}\Omega\wedge\overline\Omega=0$, whereas its
appropriate constant multiple is a positive volume form.
This contradiction excludes the common limit.
\end{proof}



\paragraph{Scope after \ref{push:src:E.2.6}}

\ref{push:src:E.2.5} proves that the central meromorphic reduction is
holomorphic, and \ref{push:src:E.2.6} removes its fixed divisor and degree
loss for $D>0$.  The positive-slope proof alone does not
exclude the effective antipluricanonical divisor allowed
by the restriction identity when $D<0$.  \ref{push:src:E.2.8}--\ref{push:src:E.2.10}
remove that possibility by the horizontal-divisor vanishing
argument.  \ref{push:src:E.2.13} and \ref{push:src:E.2.16}--\ref{push:src:E.2.21} then recover the finite
charts, cusp, and actual smooth completion.  The separate
generic-disc propagation and closure arguments in
\ref{push:src:E.2.22}--\ref{push:src:E.2.24} prove \ref{push:src:E.2.4} for arbitrary analytic bases,
deformation chains, and continuous complex structures.
The present proposition retains its stated one-sided,
positive-slope scope.

\Uses{the hypotheses, notation and definitions stated here; \cref{push:src:E.2.6}; \cref{push:src:E.2.5}; \cref{push:src:E.2.8}; \cref{push:src:E.2.10}; \cref{push:src:E.2.13}; \cref{push:src:E.2.16}; \cref{push:src:E.2.21}; \cref{push:src:E.2.22}; \cref{push:src:E.2.24}; \cref{push:src:E.2.4}}


\begin{proposition}[an extended fibration cannot reach the zero period]\label{push:src:E.2.7}
\StatementRecord{push:src:E.2.7}{proved}{Historical events: Claude Fable 5.2 (Anthropic), effort high, 16 November 2026: carried as background without fault found; used and rechecked in this run. GPT-6 Astra (OpenAI), ChatGPT Pro mode, effort Pro, S6\_09\_PAPERS, 17 November 2026: targeted proof check. Also checked in this writing. Pre-continuation writing events: GPT-6 Astra (OpenAI), ChatGPT Pro mode, effort Pro, S6\_10\_PAPERS, 18 November 2026: targeted proof check. Continuation metadata events: GPT-6 Astra (OpenAI), ChatGPT Pro mode, effort Pro, S6\_10\_PAPERS continuation, 19 November 2026: used and rechecked in this run; status-vocabulary reconciliation; metadata reconciliation only; no new mathematical check. \emph{Pre-continuation recorded proof scope.} The bounded punctured parameter extends; the exact return S\_-h A\_j\textasciicircum{}n cannot be identity at a finite cap or a cusp because the nilpotent matrices have the displayed distinct images. The holomorphic period lift remains strictly inside the determinant domain}{the hypotheses, notation and definitions stated here; \cref{push:src:E.2.5}; \cref{push:src:E.2.6}; \cref{push:src:E.2.20}; \cref{push:src:E.2.21}; \cref{push:src:E.2.22}; \cref{push:src:E.2.24}}

Under the hypotheses of \ref{push:src:E.2.5}, suppose in addition that the
intrinsic algebraic reductions extend to a holomorphic map
\[
 g:\mathcal X\longrightarrow\mathbb P^1\times\Delta
\]
whose restriction to $X_0$ is nonconstant.  This hypothesis
holds in \ref{push:src:E.2.6}.  Normalize the punctured family to fixed
discrete logarithmic data with multiplicities $(3k,4l)$, and
write its isomorphism parameter as
\[
 z=\exp\!\left(-\frac{\pi i c}{kl}\right).
\]
The resulting holomorphic function $z:\Delta^*\to\mathbb C^*$
extends across zero and satisfies $z(0)\ne0$.  Its central
value lies strictly inside the admissible period disc.

This statement concerns the period function and the extended
relative fibration.  Identification of $X_0$ with the member
at $z(0)$ is proved separately in \ref{push:src:E.2.20}--\ref{push:src:E.2.21}; propagation
to arbitrary deformation bases and continuous families is
proved in \ref{push:src:E.2.22}--\ref{push:src:E.2.24}.
\end{proposition}


\begin{proof}
\noindent\textbf{Step 1.} Local miniversality in \cref{push:src:E.2.1} supplies holomorphic local
classifying functions for the family on $\Delta^*$.
\cref{push:src:E.1.2}--\cref{push:src:E.1.3} normalize the discrete data by holomorphic
family comparisons.  For fixed data the comparison shifts
are precisely $c\mapsto c+2kl\mathbb Z$, so the displayed
exponential is a single-valued holomorphic function on
$\Delta^*$.  An automorphism with the same marked parameter
has identity linear action on the homology of an ordinary
torus: in the necessity argument of \cref{push:src:E.1.2} its parameters
are $\alpha=1$ and $h=0$.  Consequently the homology local
system of the punctured family is also the one obtained
from these local classifying maps; possible translation
automorphisms of the comparisons add no homology monodromy.

\noindent\textbf{Step 2.} The admissible parameter domain is a punctured disc
$0<|z|<R_{k,l}$, with finite positive radius.  The removable
singularity theorem therefore extends $z$ holomorphically
to $\Delta$.  It remains to exclude a zero at the centre.

\noindent\textbf{Step 3.} We first control the moving three critical values without
assuming that they remain distinct at zero.  Since
$\pi:\mathcal X\to\Delta$ is a submersion, the critical
locus of $g$ is the analytic locus where its differential
has rank less than two.  Properness makes its image
\(\mathcal C\subset\mathbb P^1\times\Delta\) a closed
analytic set.  Its intersection with the central target
curve is the finite critical-value set of the nonconstant
map $g_0:X_0\to\mathbb P^1$.  In particular
$\mathcal C$ has no component equal to that whole curve.
After shrinking $\Delta$, its reduced one-dimensional
part is finite over $\Delta$ and over $\Delta^*$ consists
of the three critical-value sections.

\noindent\textbf{Step 4.} Those sections are individually marked by the linear
peripheral monodromy orders $3$, $4$, and infinity.
The normalization of the closure of each section is
therefore finite of degree one over the disc.  It is
isomorphic to the disc, so the three values
$c_1(s),c_2(s),c_0(s)$ extend as holomorphic maps to
$\mathbb P^1$, with collisions allowed at $s=0$.
The unique projective transformation $A_s$ sending this
ordered triple to $(0,1,\infty)$ has meromorphic
coefficients at zero, by its explicit cross-ratio formula.
This is the critical-value argument of \cref{push:src:B.4.4} in a setting
where no vertical discriminant component is present.

\noindent\textbf{Step 5.} Choose a point $v\in\mathbb P^1$ outside all central
critical values.  Properness and closedness of
$\mathcal C$ allow us to shrink the disc so that
$(v,s)\notin\mathcal C$ for every $s$.  Then
\[
 \mathcal T=g^{-1}(\{v\}\times\Delta)\longrightarrow\Delta
\]
is a proper smooth holomorphic family of surfaces.
Its noncentral fibres are ordinary complex two-tori, and
Ehresmann's theorem makes its integral first-homology
monodromy around $s=0$ the identity.

\noindent\textbf{Step 6.} In the original marked base, this same family is obtained
by evaluating the period construction at
\[
 t(s)=A_s(v)\in\mathbb P^1\setminus\{0,1,\infty\}.
\]
The function $t(s)$ is meromorphic at zero.  Its loop on a
sufficiently small positive $s$-circle is therefore either
null-homotopic in the marked base or conjugate to a power
of one of its three peripheral loops.  Indeed, if its
limit is ordinary, its image lies in an ordinary small
disc; if its limit is a marked point, a local coordinate
has a zero or pole of finite order and its image lies
in the corresponding punctured disc.  This observation
also covers a constant $t(s)$.

\noindent\textbf{Step 7.} Let $\Lambda$ have ordered basis
$(\widehat\gamma,\widehat u,\widehat w,\widehat\delta)$,
put $E=\widehat\delta\otimes\gamma$, and write
$S_a=I+aE$.  The two finite peripheral monodromies are
\[
 A_1=\begin{pmatrix}
 1&0&0&0\\6&0&1&0\\-6&-1&-1&0\\-2&1&0&1
 \end{pmatrix},\qquad
 A_2=\begin{pmatrix}
 1&0&0&0\\0&0&-1&0\\-6&1&0&0\\3&0&1&1
 \end{pmatrix}.
\]
They have orders three and four, respectively, even for
the larger multiplicities: the extra finite deck
transformations are translations and do not affect
$H_1$.  In the cusp marking the third monodromy is
\[
 A_0=I+N,\qquad
 N\widehat u=\widehat w,\quad
 N\widehat\gamma=-\widehat\delta,\quad
 N\widehat w=N\widehat\delta=0.
\]
All peripheral matrices preserve $\widehat\delta$ and
$\gamma$, so commute with every $S_a$.  The change to
the cusp marking preserves the same rank-one operator
$E$; conjugating a peripheral loop consequently does
not alter the cancellation test below.

\noindent\textbf{Step 8.} Suppose that $z(s)$ has a zero of order $d>0$ at zero.
Continuation of a logarithm around the positive loop
changes $c$ by $-h$, where
\[
 h=2kl\,d>0.
\]
Since $\Pi_{c+h}=\Pi_cS_h$, a comparison of the identical
complex universal covers from parameter $c-h$ back to
$c$ acts on the marked lattice by $S_{-h}$.
The affine translation part of a closing comparison
acts trivially on $H_1$.  Thus the return map of the
family $\mathcal T$ has homology action, up to a common
conjugation,
\[
 S_{-h},\qquad S_{-h}A_1^n,\qquad
 S_{-h}A_2^n,\qquad\hbox{or}\qquad S_{-h}A_0^n
\]
for an integer $n$.  Reversing the convention for a
peripheral generator only replaces $n$ by $-n$.
The parameter factor and the base factor commute, so
the order in this display has no effect.

\noindent\textbf{Step 9.} None of these matrices is the identity.  In the two
finite cases, taking the third or fourth power would
give $S_{-3h}=I$ or $S_{-4h}=I$, impossible because
$E\ne0$ and $h>0$.  In the cusp case, $N^2=EN=NE=0$
and hence
\[
 S_{-h}A_0^n=I+nN-hE.
\]
Its coefficient of $\widehat w$ on $\widehat u$ is
$n$, so identity would force $n=0$, and then $h=0$.
The ordinary case is $S_{-h}\ne I$.  This contradicts
the identity monodromy of the smooth family over
the full disc.  Therefore $z(0)\ne0$.

\noindent\textbf{Step 10.} Finally, the extended holomorphic function satisfies
$|z(0)|\le R_{k,l}$.  Equality would make it a constant
of maximal modulus, contradicting $|z(s)|<R_{k,l}$
for $s\ne0$.  Thus $0<|z(0)|<R_{k,l}$, as claimed.
The matrix identities and the factor $h=2kld$ are
checked exactly in \path{code/e_push_E_zero_end.py}.
\end{proof}



\paragraph{Scope and subsequent completion}

For either canonical slope, \ref{push:src:E.2.10} supplies the
fibration-persistence hypothesis of this proposition,
so a one-sided smooth specialization has an interior
period value.  Equality of that value with an existing
model does not by itself prove uniqueness of its smooth
completion.  \ref{push:src:E.2.11}--\ref{push:src:E.2.12} exclude critical values
appearing only at the central parameter; \ref{push:src:E.2.13}--\ref{push:src:E.2.20}
recover the actual finite and cusp charts, and \ref{push:src:E.2.21}
proves the full one-sided identification.  The distinct
problem of propagation through arbitrary deformation
branches is settled by \ref{push:src:E.2.22}--\ref{push:src:E.2.24}.

\Uses{the hypotheses, notation and definitions stated here; \cref{push:src:E.2.10}; \cref{push:src:E.2.11}; \cref{push:src:E.2.12}; \cref{push:src:E.2.13}; \cref{push:src:E.2.20}; \cref{push:src:E.2.21}; \cref{push:src:E.2.22}; \cref{push:src:E.2.24}}


\subsubsection*{A restriction on the negative-slope fixed divisor}

\begin{proposition}[horizontal divisors force non-K\"ahler fibre normalizations]\label{push:src:E.2.8}
\StatementRecord{push:src:E.2.8}{proved}{Historical events: Claude Fable 5.2 (Anthropic), effort high, 16 November 2026: carried as background without fault found. GPT-6 Astra (OpenAI), ChatGPT Pro mode, effort Pro, S6\_09\_PAPERS, 17 November 2026: targeted proof check. Also checked in this writing. Pre-continuation writing events: GPT-6 Astra (OpenAI), ChatGPT Pro mode, effort Pro, S6\_10\_PAPERS, 18 November 2026: targeted proof check; used and rechecked in this run. GPT-6 Astra (OpenAI), ChatGPT Pro mode, effort Pro, S6\_10\_PAPERS, 18 November 2026 05:35:21+00:00: used and rechecked in this run. Continuation metadata events: GPT-6 Astra (OpenAI), ChatGPT Pro mode, effort Pro, S6\_10\_PAPERS continuation, 19 November 2026: used and rechecked in this run; status-vocabulary reconciliation; metadata reconciliation only; no new mathematical check. \emph{Pre-continuation recorded proof scope.} The complement count forces all supports irreducible in the b1=1 branch. The restriction of a horizontal effective divisor to any proposed Kahler normalization resolution has zero Chern class and positive integral, giving the stated contradiction}{the hypotheses, notation and definitions stated here; \cref{push:src:E.2.5}; \citet[Lemma~4.2]{CDP2020D}; \citet[Lemma~3.3 and Corollary~3.4]{Fujino2021E}; \cref{push:src:E.2.6}}

Let $X$ be a compact complex threefold diffeomorphic to $S^6$, and
let $f:X\to\mathbb P^1$ be a surjective holomorphic map with
connected fibres.  Suppose that a smooth general fibre $F$ has
$b_1(F,\mathbb Q)=1$.  Then every reduced fibre of $f$ is
irreducible, and the monodromy on $H^1(F,\mathbb Q)$ is trivial.

If $H\subset X$ is a horizontal prime divisor, the normalization
of every reduced fibre lies outside Fujiki's class $\mathcal C$.
Equivalently, no such normalization admits a K\"ahler resolution.
In particular no fibre has Moishezon normalization.

Without the hypothesis $b_1(F)=1$, every irreducible component
of a \emph{reducible} fibre of any such connected map from a
complex $S^6$ also has normalization outside $\mathcal C$.

Consequently, in \ref{push:src:E.2.5} with $D<0$, a horizontal component of
$E_0$ is possible only if every reduced fibre of the central
algebraic reduction $g_0$ is irreducible and has normalization
outside $\mathcal C$.  This necessary condition does not exclude
non-K\"ahler normalizations and does not prove persistence of the
relative pencil in the negative-slope case.
\end{proposition}


\begin{proof}
\noindent\textbf{Step 1.} Choose a nonempty finite set $\Sigma\subset\mathbb P^1$ containing
all critical values; a regular value can be added if necessary.
Write $s=|\Sigma|$, put $A=f^{-1}(\Sigma)_{\mathrm{red}}$,
and let $r$ be the number of irreducible components of $A$.
The compact analytic set $A$ is triangulable.  Each of its
irreducible two-dimensional components has its complex
fundamental class, and their intersections have real dimension
at most two.  Hence
\[
 H^4(A,\mathbb Q)\simeq\mathbb Q^r.
\]
Alexander duality in the oriented six-sphere gives
\[
 b_1(X\setminus A,\mathbb Q)=r.
\]
On the other hand, $f:X\setminus A\to U:=\mathbb P^1\setminus\Sigma$
is a proper smooth fibre bundle.  The base $U$ has the homotopy
type of a graph with $s-1$ loops.  The degree-one Leray--Serre
sequence therefore gives an exact sequence
\[
 0\longrightarrow H^1(U,\mathbb Q)
 \longrightarrow H^1(X\setminus A,\mathbb Q)
 \longrightarrow H^1(F,\mathbb Q)^{\pi_1(U)}
 \longrightarrow0.
\]
In particular, the correct formula retains the invariant
subspace:
\[
 r=s-1+\dim H^1(F,\mathbb Q)^{\pi_1(U)}\le s.
\]
Every nonempty fibre over $\Sigma$ has at least one irreducible
component, so $r\ge s$.  Equality follows throughout.  Thus
every fibre over $\Sigma$ has precisely one reduced irreducible
component and the invariant subspace is all of the
one-dimensional $H^1(F,\mathbb Q)$.  Fibres outside $\Sigma$ are
smooth and connected, hence irreducible.  This proves both
assertions without assuming monodromy triviality at the outset.
It is the rank-one repair of the component-count argument;
the unrestricted formula in \citet[Lemma~4.2]{CDP2020D} is not
being used.

\noindent\textbf{Step 2.} Let $S=f^{-1}(t)_{\mathrm{red}}$.  A horizontal prime $H$ maps
surjectively to the base by properness, so $H\cap S\ne\varnothing$.
Since $S$ is irreducible, $S\subset H$ would imply $S=H$,
contrary to horizontality.  The canonical section of
$\mathcal O_X(H)$ consequently restricts to a nonzero section
on the reduced irreducible space $S$, with a nonempty zero
set.  Its zero divisor is a nonzero effective Cartier divisor
on $S$.  Indeed the section cannot vanish identically on
any local irreducible branch: its pullback would then vanish
on an open subset of the irreducible normalization, hence
everywhere by analytic continuation.  In the reduced local
ring it consequently avoids every minimal prime and is a
non-zero-divisor, including at self-intersections.

\noindent\textbf{Step 3.} Let $\nu:S^\nu\to S$ be the normalization.  If $S^\nu$ were
in $\mathcal C$, it would have a compact K\"ahler resolution
$\mu:\widehat S\to S^\nu$; this precise surface statement is
\citet[Lemma~3.3 and Corollary~3.4]{Fujino2021E}.
The pulled-back line bundle
\[
 L=(\nu\mu)^*(\mathcal O_X(H)|_S)
\]
has $c_1(L)=0$ integrally, since $H^2(X,\mathbb Z)=0$.
Its pulled-back section is nonzero and has zeros: normalization
is finite surjective, and modification is proper surjective.
Thus it defines a nonzero effective divisor $B$ on
$\widehat S$.  For a K\"ahler form $\omega$,
\[
 0=\langle c_1(L)\smile[\omega],[\widehat S]\rangle
   =\int_B\omega>0,
\]
a contradiction.  The argument applies to the reduced support
even when the original fibre is nonreduced.  Moishezon spaces
are in $\mathcal C$, so the stated consequence follows.

\noindent\textbf{Step 4.} For the assertion about a reducible fibre, let $S_i$ be one
of its reduced irreducible components.  The fibre is connected,
so another component $S_j$ meets $S_i$; otherwise $S_i$ would
be an open-and-closed component of its support.  Distinctness
ensures $S_i\not\subset S_j$.  Apply exactly the preceding
canonical-section and positivity argument to
$\mathcal O_X(S_j)|_{S_i}$.  Its first Chern class is zero
because $H^2(X,\mathbb Z)=0$, whereas its effective Cartier
zero divisor is nonempty.  This proves the additional assertion
without any hypothesis on the general fibre's first Betti number.

\noindent\textbf{Step 5.} For the application, choose a very general smooth fibre $F$
of $g_0$ meeting a horizontal component of $E_0$ and avoiding
its vertical components.  As in \ref{push:src:E.2.6}, adjunction gives
$c_1(K_F)=0$ and $F$ is minimal.  The identity of \ref{push:src:E.2.5} now
restricts to
\[
 K_F^{-A}=\mathcal O_F((-D)E_0|_F),
 \qquad -D>0,
\]
with a nonzero effective divisor on the right.  If a positive
power $K_F^m$ had a nonzero section, multiplying a suitable
power of that section by the $m$th power of the displayed
anticanonical section would give a nonzero holomorphic
function with zeros on $F$.  Compactness forbids this.
Thus $\kappa(F)=-\infty$.  Moreover $F$ is not K\"ahler,
since its nonzero effective divisor $(-D)E_0|_F$ has zero
real Chern class.  The Enriques--Kodaira classification then
places $F$ in class VII and gives $b_1(F)=1$; see
\citet[Chapter~VI, Theorem~1.1 and Table~10, pp.~243--244]
{BHPV2004E}.  The preceding argument applies with $f=g_0$
and any horizontal prime component $H$ of $E_0$.
\end{proof}



The proof supplies a necessary restriction on an unknown
smooth specialization.  It neither classifies the singular
non-K\"ahler normalizations nor replaces the unresolved
non-normal-fibre step in the holomorphic-fibration argument
of \citet[\S7]{CDP2020D}.  That step is not used here:
\ref{push:src:E.2.9}--\ref{push:src:E.2.10} prove the needed horizontal-divisor vanishing,
and \ref{push:src:E.2.21}--\ref{push:src:E.2.24} complete the actual specialization and
deformation-chain arguments.

\subsubsection*{Removing the negative-slope horizontal fixed part}

\begin{proposition}[no horizontal divisor when all fibres are irreducible]\label{push:src:E.2.9}
\StatementRecord{push:src:E.2.9}{proved}{Historical events: Claude Fable 5.2 (Anthropic), effort high, 16 November 2026: carried as background without fault found. GPT-6 Astra (OpenAI), ChatGPT Pro mode, effort Pro, S6\_09\_PAPERS, 17 November 2026: targeted proof check. Also checked in this writing. Pre-continuation writing events: GPT-6 Astra (OpenAI), ChatGPT Pro mode, effort Pro, S6\_10\_PAPERS, 18 November 2026: targeted proof check. Continuation metadata events: GPT-6 Astra (OpenAI), ChatGPT Pro mode, effort Pro, S6\_10\_PAPERS continuation, 19 November 2026: used and rechecked in this run; status-vocabulary reconciliation; metadata reconciliation only; no new mathematical check. \emph{Pre-continuation recorded proof scope.} Negative divisor twists, Krull intersection on every Cartier multiplicity layer, full-fibre duality, compact nested jumping sets, and the affine Picard-line avoidance argument were re-derived. The final contradiction uses the explicitly cited analytic cohomology-and-base-change input and chi(B)=1}{the hypotheses, notation and definitions stated here; \cref{push:src:E.2.8}; \citet[\S3, Theorem~1, p.~81, and the complete-intersection
calculation on p.~91]{RR70push}; \cite[Chapter~III, Theorem~4.12(ii)]{BS76push}; \cite[Chapter~III, Theorem~4.1]{BS76push}}

Let $X$ be a compact complex threefold diffeomorphic to $S^6$,
and let $f:X\to\mathbb P^1$ be a surjective holomorphic map
with connected fibres.  If every reduced fibre is irreducible,
then $X$ has no horizontal prime divisor for $f$.
\end{proposition}


\begin{proof}
\noindent\textbf{Step 1.} Suppose that $H$ is such a divisor.  It is Cartier because
$X$ is smooth.  Since $H$ is horizontal and every reduced
fibre $S$ is irreducible, $H$ meets $S$ and does not contain
it.  As in \ref{push:src:E.2.8}, restriction of its canonical section is a
non-zero-divisor on every local branch of $S$.

\noindent\textbf{Step 2.} We first record the elementary vanishing used below.  On a
compact reduced irreducible complex space $S$, let $D$ be a
nonzero effective Cartier divisor and $V$ a vector bundle.
Multiplication by the powers of the canonical section
identifies
\[
 H^0(S,V(-nD))\subset H^0(S,V),\qquad n\ge0,
\]
with a decreasing sequence of finite-dimensional vector
spaces.  A section in their intersection vanishes to every
order along $D$.  Krull intersection in a local ring at a
point of $D$ makes its germ zero.  After pullback to the
irreducible normalization, analytic continuation makes the
section zero everywhere.  The decreasing sequence therefore
stabilizes at zero.  In particular
\[
 H^0(S,V(-nD))=0\quad\hbox{for all sufficiently large }n.
\]
The same conclusion holds for a vector bundle on a thickened
fibre $mS$.  Indeed $S$ is Cartier in $X$, and the finite
filtration by powers of its ideal has successive quotients
\[
 V|_S\otimes\mathcal O_S(-jS),\qquad 0\le j<m.
\]
Apply the preceding reduced-space statement to these finitely
many vector bundles, with the same nonzero divisor $H|_S$,
and then use their exact sequences.  This also proves that
the multiplication maps used here remain injective on $mS$:
its only associated component is $S$ and $H$ contains none
of its local branches.

\noindent\textbf{Step 3.} Let $B$ be any holomorphic vector bundle on $X$.  The map
$f$ is flat, since its local rings are torsion-free over the
discrete valuation rings of the smooth base.  Every analytic
fibre $F_t$ is a Cartier, Cohen--Macaulay, Gorenstein surface,
possibly nonreduced.  Its dualizing line is
\[
 \omega_{F_t}=K_X|_{F_t},
\]
where the constant one-dimensional factor from the target
has been trivialized.  Serre duality on this compact
Gorenstein surface gives
\[
 H^2(F_t,B(nH)|_{F_t})^*
 \simeq H^0\bigl(F_t,(B^*\otimes K_X)(-nH)|_{F_t}\bigr).
\]
The compact-space duality being used includes singular
Cohen--Macaulay spaces; it is the absolute duality theorem
of \citet[\S3, Theorem~1, p.~81, and the complete-intersection
calculation on p.~91]{RR70push}, already specified in \cref{push:src:E.3.1},
not a smooth-fibre replacement.
The preceding filtration argument makes the right hand
side zero for all sufficiently large $n$, for each $t$.

\noindent\textbf{Step 4.} For fixed $n$, the nonvanishing locus
\[
 Z_n=\{t\in\mathbb P^1:
        h^2(F_t,B(nH)|_{F_t})>0\}
\]
is closed by semicontinuity.  Multiplication on the dual
section spaces gives $Z_{n+1}\subset Z_n$.  Their intersection
is empty by fibrewise vanishing.  Compactness of $\mathbb P^1$
therefore supplies a single $N$ with $Z_N=\varnothing$.
Constant-zero cohomology and analytic cohomology-and-base-change
\cite[Chapter~III, Theorem~4.12(ii)]{BS76push} give
\[
 R^2f_*(B(NH))=0.
\]
Higher direct images vanish in degrees greater than two by
the fibre dimension.  The coherent sheaf $R^1f_*(B(NH))$
has vanishing first cohomology after a sufficiently positive
twist on $\mathbb P^1$.  The projection formula and Leray
sequence consequently give, for some integer $a\gg0$,
\[
 H^2(X,B\otimes M_+)=H^3(X,B\otimes M_+)=0,
 \qquad
 M_+=\mathcal O_X(NH)\otimes f^*\mathcal O(a).
\]
All coherence, semicontinuity and constant-zero base-change
uses are the proper analytic theorems audited in \cref{push:src:E.3.1};
the argument does not infer base change from coherence alone.

\noindent\textbf{Step 5.} Independently, the reduced-space vanishing argument applied
to the connected smooth manifold $X$ and its nonzero divisor
$H$ gives
\[
 H^0(X,B\otimes M_-)=0,
 \qquad M_-=\mathcal O_X(-bH),\quad b\gg0.
\]
We must arrange these different vanishings for one line
bundle.  The exponential sequence gives
\[
 \operatorname{Pic}(X)\simeq H^1(X,\mathcal O_X)
\]
because $H^1(X,\mathbb Z)=H^2(X,\mathbb Z)=0$.
Let $\ell_-,\ell_+$ be the corresponding additive classes.
The affine line
\[
 \ell(z)=(1-z)\ell_-+z\ell_+,\qquad z\in\mathbb C,
\]
is realized by a holomorphic family of line bundles on
$X\times\mathbb C$: choose additive holomorphic Cech
cocycles representing the two classes and exponentiate their
linear combination.  Equivalently, use the integrable
Dolbeault operators obtained by a holomorphic linear
combination of fixed closed $(0,1)$-form representatives.
Thus no representability or universal-Picard-space assertion
is needed here.

\noindent\textbf{Step 6.} For the proper projection $X\times\mathbb C\to\mathbb C$,
the cohomology jumping loci
\[
 T_j=\{z:h^j(X,B\otimes M_z)>0\},\qquad j=0,2,3,
\]
are closed analytic subsets of $\mathbb C$ by the proper
analytic finite-complex theorem
\cite[Chapter~III, Theorem~4.1]{BS76push}: the condition is
expressed by the rank minors in a local finite free complex.
In particular no closedness claim from the corrected
Grauert Section~7, Satz~2 is used.  The set $T_0$ omits $z=0$,
while $T_2$ and $T_3$ omit $z=1$.  Each is therefore proper,
and their finite union cannot fill the irreducible complex
line.  Choose $z$ outside that union.  Then
\[
 \chi(X,B\otimes M_z)=-h^1(X,B\otimes M_z)\le0.
\]
Every line bundle on $X$ is topologically trivial, so
Hirzebruch--Riemann--Roch makes this Euler characteristic
equal to $\chi(X,B)$.

\noindent\textbf{Step 7.} Finally take $B=T_X$.  Since $H^2(X,\mathbb Z)=H^4(X,\mathbb Z)=0$,
the only contributing term in its Riemann--Roch formula is
\[
 \chi(X,T_X)=\frac12\langle c_3(X),[X]\rangle
            =\frac12\chi_{\mathrm{top}}(S^6)=1.
\]
This contradicts the preceding inequality.  Hence no
horizontal prime $H$ exists.
\end{proof}



\begin{proposition}[persistence for either canonical slope]\label{push:src:E.2.10}
\StatementRecord{push:src:E.2.10}{proved}{Historical events: Claude Fable 5.2 (Anthropic), effort high, 16 November 2026: carried as background without fault found. GPT-6 Astra (OpenAI), ChatGPT Pro mode, effort Pro, S6\_09\_PAPERS, 17 November 2026: targeted proof check. Also checked in this writing. Pre-continuation writing events: GPT-6 Astra (OpenAI), ChatGPT Pro mode, effort Pro, S6\_10\_PAPERS, 18 November 2026: targeted proof check. Continuation metadata events: GPT-6 Astra (OpenAI), ChatGPT Pro mode, effort Pro, S6\_10\_PAPERS continuation, 19 November 2026: used and rechecked in this run; status-vocabulary reconciliation; metadata reconciliation only; no new mathematical check. \emph{Pre-continuation recorded proof scope.} Both slope signs eliminate the horizontal common divisor through E.2.8-E.2.9, after which connected specialization gives the intrinsic pencil and recovers the same exact ordered multiplicity pair}{the hypotheses, notation and definitions stated here; \cref{push:src:E.2.5}; \cref{push:src:E.2.7}; \cref{push:src:E.2.4}; \cref{push:src:E.2.6}; \cref{push:src:E.2.8}; \cref{push:src:E.2.9}}

Under the hypotheses of \ref{push:src:E.2.5}, for every nonzero
$D=kl-3k-4l$ one has $E_0=0$ and $e=1$.
After shrinking the disc, the evaluation of
$\pi_*\mathcal L$ gives a proper flat holomorphic map
\[
 g:\mathcal X\longrightarrow\mathbb P^1\times\Delta
\]
whose restriction to every fibre is its intrinsic algebraic
reduction.  A smooth general central fibre of $g$ is a
complex two-torus.  Its normalized punctured period has
an interior limit as in \ref{push:src:E.2.7}.

Two families satisfying \ref{push:src:E.2.5} cannot have the same smooth
central manifold unless their ordered multiplicity pairs
agree, regardless of the signs of their canonical slopes.
Uniqueness of a smooth completion at an interior period and
the arbitrary deformation-chain statement \ref{push:src:E.2.4} are not
asserted.
\end{proposition}


\begin{proof}
\noindent\textbf{Step 1.} The case $D>0$ is \ref{push:src:E.2.6}.  If $D<0$ and $E_0$ had a
horizontal prime component, the final part of \ref{push:src:E.2.8} would
give $b_1(F)=1$ for a smooth general fibre of $g_0$ and
would make every reduced fibre of $g_0$ irreducible.
\ref{push:src:E.2.9} would then forbid precisely that horizontal component.
Thus for either sign $E_0$ is vertical.

\noindent\textbf{Step 2.} The connected-divisor-limit argument in the middle of the
proof of \ref{push:src:E.2.6} only uses verticality, not the sign of $D$.
Its generic choice of pencil section has central zero
support equal to the disjoint union of the vertical fixed
part and $e$ distinct regular fibres.  The total divisor
is reduced and irreducible with connected central support
after normalization and Stein factorization.  Consequently
$E_0=0$ and $e=1$.  Surjectivity of evaluation, flatness and
the torus assertion follow from the remainder of that same
proof.  The hypothesis of \ref{push:src:E.2.7} is now satisfied for either
sign, and gives the interior period limit.

\noindent\textbf{Step 3.} Suppose two such one-sided families have the same central
manifold $X_0$.  Its algebraic reduction is intrinsic, so
both pencil lines are the same line
$L=g_0^*\mathcal O(1)$ up to isomorphism.  With indices
$i=1,2$, the two root identities give
\[
 K_{X_0}^{A_i}=L^{D_i},\qquad
 L^{D_1A_2-D_2A_1}=\mathcal O.
\]
The line $L$ is not torsion, because a positive power has
a nonzero section with zeros.  Hence $D_1/A_1=D_2/A_2$.
\cref{push:src:E.2.3} recovers the ordered coprime multiplicity pair from
this rational number.  This proves the common-limit
assertion without identifying the limit with a member of
the given period construction.
\end{proof}



The vanishing argument is an application of standard proper
cohomology, duality and Riemann--Roch machinery; no novelty
claim is intended.  It proves the needed absence of a
horizontal fixed divisor under the rank-one hypothesis.
It does not establish the different claim that every
holomorphic Hopf-surface fibration on a complex six-sphere
is impossible.

\subsubsection*{An extension criterion at an isolated base point}

\begin{proposition}[purity with a smooth torus Jacobian]\label{push:src:E.2.11}
\StatementRecord{push:src:E.2.11}{proved}{Historical events: Claude Fable 5.2 (Anthropic), effort high, 16 November 2026: carried as background without fault found. GPT-6 Astra (OpenAI), ChatGPT Pro mode, effort Pro, S6\_09\_PAPERS, 17 November 2026: targeted proof check. Also checked in this writing. Pre-continuation writing events: GPT-6 Astra (OpenAI), ChatGPT Pro mode, effort Pro, S6\_10\_PAPERS, 18 November 2026: targeted proof check. Continuation metadata events: GPT-6 Astra (OpenAI), ChatGPT Pro mode, effort Pro, S6\_10\_PAPERS continuation, 19 November 2026: used and rechecked in this run; status-vocabulary reconciliation; metadata reconciliation only; no new mathematical check. \emph{Pre-continuation recorded proof scope.} The vertical Lie fields extend across codimension two, integrate by properness, and give a real compact-torus action. Connected stabilizers have trivial tangent representation and vanish; the finite transverse germ cannot have isolated ramification. The resulting free transitive group action supplies the actual torsor extension}{the hypotheses, notation and definitions stated here; the explicit constructions and calculations in the proof}

Let $B\subset\C^2$ be a sufficiently small bidisc about $0$, put
$U=B\setminus\{0\}$, and let $g:X\to B$ be a proper flat
holomorphic map from a connected smooth complex $(n+2)$-fold,
with connected fibres of dimension $n$.  Let $G\to B$ be a
smooth proper holomorphic family of complex $n$-tori with a
holomorphic zero section and its group law.  Suppose that
$X_U\to U$ is a holomorphic $G_U$-torsor.  Then $g$ is a
submersion on all of $X$, and the given action extends uniquely
to make $X\to B$ a holomorphic $G$-torsor.

No polarization, K\"ahler metric on $X$, or section of the
given torsor on $U$ is assumed.  The smooth extension $G$ is
an explicit hypothesis.
\end{proposition}


\begin{proof}
\noindent\textbf{Step 1.} Shrink $B$ so that the holomorphic Lie bundle of $G$ has a
frame $e_1,\ldots,e_n$.  Its period local system has a basis
of holomorphic sections
\[
 \lambda_1(b),\ldots,\lambda_{2n}(b)\in\C^n,
 \qquad
 G_b=\C^n/\sum_{a=1}^{2n}\Z\lambda_a(b).
\]
Indeed the lattice is locally constant over the simply
connected bidisc, and the local inverses of the relative
exponential map make its sections holomorphic.  The periods
are a real basis at every $b$.

\noindent\textbf{Step 2.} The infinitesimal torsor action gives holomorphic vector
fields $V_1,\ldots,V_n$ on $X_U$.  Flatness makes
$Z=g^{-1}(0)$ have complex codimension two in $X$.
Hartogs extension in local coordinates therefore extends
each coefficient of each $V_j$ uniquely across $Z$; uniqueness
glues these local extensions to holomorphic vector fields
on $X$.  This use of Hartogs applies to functions on the
complement of an analytic set of codimension at least two
in a complex manifold, hence also coefficientwise to any
locally free sheaf.  It does not assert extension across a
divisor.  The identities
\[
 dg(V_j)=0,\qquad [V_i,V_j]=0
\]
hold on the dense open subset $X_U$ and consequently on $X$.

\noindent\textbf{Step 3.} Every $V_j$, and every linear combination with coefficients
pulled back from $B$, is complete.  Its integral curve stays
in a compact fibre; a finite endpoint of a maximal solution
would have an accumulation point in that fibre, where the
holomorphic ordinary differential equation extends the
solution.  The flows thus define a holomorphic action of
$\C^n$ on $X$ over $B$.

\noindent\textbf{Step 4.} For $a=(a_1,\ldots,a_{2n})\in\R^{2n}$ take the time-one
flow of the vertical holomorphic vector field whose Lie
coordinate is
\[
 \sum_{j=1}^{2n}a_j\lambda_j(g(x)).
\]
These maps depend smoothly on $a$ and add under composition:
all their coefficients are constant along every vertical
flow.  For $a\in\Z^{2n}$ the map is the identity on $X_U$,
hence everywhere by holomorphic uniqueness.  We obtain an
action of the fixed compact real torus
$T=\R^{2n}/\Z^{2n}$ by holomorphic automorphisms of $X$,
over $B$.  On the fibre over $b$, it is the holomorphic
action of $G_b$ obtained by factoring the $\C^n$ action by
its period lattice.  The action of $T$ is free on $X_U$.
The variable period identities also make the holomorphic
Lie-bundle action descend to a holomorphic action
$G\times_BX\to X$.  Its restriction therefore acts on
each scheme-theoretic fibre, including its nilpotent
structure.

\noindent\textbf{Step 5.} We show that its stabilizers have dimension zero everywhere.
Let $x\in X$, $b=g(x)$, and let $H$ be the identity component
of the stabilizer.  Viewed in $G_b$, it is a compact complex
subgroup: it is the identity component of the inverse image
of $x$ under the holomorphic orbit map $G_b\to X_b$.
Put
\[
 K_x=\ker(dg_x)=T_x(X_b),
\]
where the right-hand side is the Zariski tangent space of
the scheme-theoretic fibre.  Differentiating the holomorphic
$G_b$ action on this fibre gives a holomorphic representation
$H\to\operatorname{GL}(K_x)$.  It is trivial, because every
holomorphic function on the compact connected complex group
$H$ is constant.  The induced representation on
$T_xX/K_x$ is also trivial: this quotient injects into
$T_bB$ by $dg_x$, and every group element acts over $B$.
Averaging a Hermitian inner product over the compact group
$H$ gives an invariant complement to $K_x$.  The full
derivative representation on $T_xX$ is therefore trivial.

\noindent\textbf{Step 6.} Here the required local linearization has a direct proof.
Choose holomorphic coordinates $z$ centred at $x$ and a
sufficiently small $H$-invariant neighbourhood on which all
$z(hy)$ are defined.  Such a neighbourhood exists by
compactness of $H$ and continuity of the action.  With
normalized Haar measure set
\[
 \Psi(y)=\int_H z(hy)\,dh.
\]
This map is holomorphic, invariant under $H$, and has
derivative the identity at $x$.  It is a holomorphic
coordinate system after shrinking.  Its invariance then
forces every $h\in H$ to be the identity near $x$, and hence
on connected $X$.  This contradicts freeness on $X_U$ unless
$H$ is trivial.  Thus every stabilizer is finite, and the
$n$ vector fields $V_j$ are linearly independent everywhere.

\noindent\textbf{Step 7.} Their commuting flows give local holomorphic coordinates
$(z,s)\in\C^n\times\C^2$ in which they span the $z$
directions: choose a transverse complex surface and apply
the holomorphic inverse function theorem to the flow map.
In these coordinates $g$ has the form
\[
 g(z,s)=h(s).
\]
Equidimensionality implies that the germ $h^{-1}(b)$ is
zero-dimensional.  Consequently the germ
$h:(\C^2,0)\to(B,b)$ is finite; this is the local finiteness
criterion obtained by Weierstrass division for an isolated
fibre.  If $dg_x$ failed to have rank two, the determinant
of $dh$ would vanish at $s=0$.  It is not identically zero,
since the map is a submersion away from $Z$.  Its zero set
would therefore contain a complex curve through $0$.
Finiteness sends that curve to a curve of critical values.
This contradicts smoothness over $B\setminus\{0\}$.
Hence $g$ is a submersion everywhere.

\noindent\textbf{Step 8.} The $G_b$ orbits now have the full dimension of every smooth
connected fibre and are compact.  Each fibre is a single
orbit.  Its finite stabilizer acts trivially on its tangent
space, because $G_b$ is abelian; it also acts trivially on
the base tangent quotient.  Averaging as above, now over
that finite group, again makes its full derivative trivial
and then makes the action locally trivial.  Freeness on
$X_U$ forces the stabilizer to be trivial.  Thus the extended
action is free and transitive on every fibre.  Local
holomorphic sections of the submersion identify $X$ with
$G$ over small base opens, which proves the torsor assertion.
The extension is unique because it agrees with the original
one on the dense open subset $X_U$.
\end{proof}



\paragraph{Scope and literature check}

This extension criterion has precisely the displayed hypotheses. The proof does not establish purity for every proper flat map.  In particular it does
not apply to B.4.15: that construction has vertical
threefolds over the bidisc origin and is not flat over the
bidisc.  Searches on 15 November 2026 for purity and
discriminants of complex-torus fibrations, Hartogs extension
of torus actions, and isolated critical values did not
locate the exact statement under these hypotheses without
an added K\"ahler assumption.  The complete argument is given above.

\Uses{the hypotheses, notation and definitions stated here; the argument and qualifications immediately following}


\begin{proposition}[no additional central critical values]\label{push:src:E.2.12}
\StatementRecord{push:src:E.2.12}{proved}{Historical events: Claude Fable 5.2 (Anthropic), effort high, 16 November 2026: carried as background without fault found. GPT-6 Astra (OpenAI), ChatGPT Pro mode, effort Pro, S6\_09\_PAPERS, 17 November 2026: targeted proof check. Also checked in this writing. Pre-continuation writing events: GPT-6 Astra (OpenAI), ChatGPT Pro mode, effort Pro, S6\_10\_PAPERS, 18 November 2026: targeted proof check. Continuation metadata events: GPT-6 Astra (OpenAI), ChatGPT Pro mode, effort Pro, S6\_10\_PAPERS continuation, 19 November 2026: used and rechecked in this run; status-vocabulary reconciliation; metadata reconciliation only; no new mathematical check. \emph{Pre-continuation recorded proof scope.} The ordinary, finite-root and cusp period limits were checked separately. Identity parameter monodromy supplies the integral root divisibility; the cusp nilpotent return excludes cusp collapse. The smooth Jacobian and E.2.11 exclude an isolated extra critical value}{the hypotheses, notation and definitions stated here; \cref{push:src:E.2.10}; \cref{push:src:E.2.7}; \cref{push:src:E.2.11}}

In the smooth-specialization setting of \ref{push:src:E.2.10}, let
\[
 g:\mathcal X\longrightarrow\mathbb P^1\times\Delta
\]
be the extended intrinsic fibration, and let
$c_1(s),c_2(s),c_0(s)$ be the three labelled critical-value
sections of \ref{push:src:E.2.7}.  The critical values of $g_0$ are exactly
the distinct points among $c_1(0),c_2(0),c_0(0)$.
In particular, no additional critical value occurs only
on the central target curve.

If the normalizing projective transformations $A_s$ of
\ref{push:src:E.2.7} have a degenerate limit, the central regular fibres
are isomorphic complex two-tori and $g_0$ has at most two
critical values.  If the constant image of the limiting
rank-one matrix is $0$ or $1$, this torus is the
corresponding finite-root abelian surface.  If that image
is an ordinary source point, it is the ordinary torus at
that point and need not be an abelian surface.  The cusp
cannot be the constant image.
\end{proposition}


\begin{proof}
\noindent\textbf{Step 1.} \ref{push:src:E.2.7} gives $0<|z(0)|<R_{k,l}$.  A local logarithm
therefore extends the normalized parameter $c(s)$
holomorphically to the disc.  The critical image is closed
and contains the three punctured sections, so it contains
their three limiting points.  Suppose that a different
central critical value $b_0=(v_0,0)$ exists.  Take a small
bidisc $B$ about $b_0$ which avoids all three section
closures and all other central critical values.  Put
$U=B\setminus\{b_0\}$.  The restricted family over $U$
is a proper smooth family of complex two-tori.  We verify
that its Jacobian extends to a smooth holomorphic torus
group family over $B$, so that \ref{push:src:E.2.11} applies.

\noindent\textbf{Step 2.} The normalized source coordinate is $t=A_s(v)$.
Represent $A_s$ by a matrix with holomorphic entries,
after clearing poles and removing a common power of
$s$, so its central matrix is nonzero.  If that matrix
is invertible, $t$ extends holomorphically near $b_0$ and
its value there is ordinary: a marked value would put
$b_0$ on the corresponding section closure.  The
ordinary period model with parameter $c(s)$ then gives
the desired smooth Jacobian over $B$.

\noindent\textbf{Step 3.} Otherwise the central matrix has rank one.  Denote its
constant image by $t_*$ and its exceptional target point
by $v_*$.  For each marked source value different from
$t_*$, its preimage tends to $v_*$.  At least two of the
three marked values are different from $t_*$, so $v_*$
belongs to at least two labelled section closures.
Consequently $v_0\ne v_*$.  After shrinking $B$, the
function $t(s,v)$ tends uniformly to $t_*$ there.
For an ordinary $t_*$ the ordinary period model again
extends directly.

\noindent\textbf{Step 4.} Suppose instead that $t_*=0$ or $1$, and use a local
coordinate $\xi$ at that marked point.  The divisor of
$\xi(t(s,v))$ on $B$ is supported on $s=0$, since $B$
avoids the corresponding critical-value graph.
Consequently
\[
 \xi(t(s,v))=s^n u(s,v),\qquad n>0,
\]
where $u$ is a nowhere-zero holomorphic function.
This also explains why no zero of the leading
coefficient may be overlooked: such a zero would put
the point on the labelled graph closure.
Choose a nearby $v_1$ for which $(v_1,0)$ is regular.
The actual family over the full $s$-disc at $v_1$ is
smooth, so its integral degree-one monodromy is the
identity.  The parameter $c(s)$ has no winding.  By
the marked comparison of \ref{push:src:E.2.7}, this monodromy is
$A_1^n$ or $A_2^n$, respectively.  Their orders are
three and four, so $e\mid n$, where $e=3$ or $4$.
A holomorphic $e$th root of $u$ exists on the bidisc,
and
\[
 r(s,v)=s^{n/e}u(s,v)^{1/e}
\]
is an actual holomorphic root coordinate on $B$.
Pull back the smooth finite-root torus period model
by $(c(s),r(s,v))$.  This produces the required
smooth torus group family $G\to B$.  The larger
multiplicities $3k,4l$ do not change this Jacobian
argument: their additional deck transformations
are translations, and the linear Jacobian model
is the original order-three or order-four root model.

\noindent\textbf{Step 5.} For completeness $t_*=\infty$ is impossible.  In the
cusp coordinate the same divisor argument gives
$q=s^n u$ with $n>0$.  The smooth family at $v_1$
would have monodromy $A_0^n$.  But
$A_0=I+N$, $N^2=0$, and
$N\widehat u=\widehat w\ne0$, so $A_0^n\ne I$.

\noindent\textbf{Step 6.} It remains to identify $G|_U$ with the actual Jacobian,
rather than merely comparing individual tori.
The punctured bidisc $U$ is simply connected.  Its
integral first-homology local system therefore has
a global marking.  Over $s\ne0$ that marking agrees
with the reference period marking, after one constant
integral change of basis; the transition comparisons
of \cref{push:src:E.1.3} have identity linear homology action at
fixed marked parameter.  Both holomorphic period
subbundles extend on $U$, and they agree on its dense
open subset $s\ne0$.  They agree everywhere on $U$.
Equivalently, the two quotients of their holomorphic
Lie bundles by the same marked period lattice are
isomorphic torus group families on $U$.  This is a
Jacobian identification and makes no assertion that
the actual affine torsor has a section on $B$.

\noindent\textbf{Step 7.} The total space is smooth and $g$ is proper and flat.
The only possible nonsmooth fibre over $B$ lies above
$b_0$, and its codimension in the total space is two.
All hypotheses of \ref{push:src:E.2.11} now hold.  That lemma makes
$g$ smooth also above $b_0$, a contradiction.  The
central critical-value set is therefore precisely
the set of labelled limits.

\noindent\textbf{Step 8.} For the final assertions, in the rank-one case every
ordinary central target point different from $v_*$
has the same limiting marked torus: it is obtained
from $(c(0),t_*)$, or from the finite-root value
$(c(0),r=0)$.  Equality of the extended period
subbundles proves this assertion without assuming
separatedness for threefold completions.
All marked values except possibly $t_*$ have their
preimages tending to $v_*$.  Hence there are at most
two central critical values; if $t_*$ is ordinary,
there is exactly one.
At $t_*=0$ or $1$, the finite linear automorphism of
the root torus has complex eigenvalues $1$ and a
primitive third or fourth root of unity.  The
rational idempotents for these two distinct
cyclotomic factors split that torus, up to isogeny,
as a product of two elliptic curves.  It is therefore
an abelian surface.  This argument makes no
projectivity assertion for an ordinary limiting
torus, and supplies no polarization on the whole
central threefold.
\end{proof}



\paragraph{Scope of \ref{push:src:E.2.11}--\ref{push:src:E.2.12}}

Purity here is obtained from the explicit good
Jacobian near a hypothetical extra critical value.
Smooth total space and flatness alone were not used
as a general purity theorem.  These results do not
give uniqueness of an interior smooth completion,
do not exclude the remaining isotrivial central
fibrations, and do not prove persistence of the
canonical pencil on an arbitrary new deformation
branch.

\Uses{the hypotheses, notation and definitions stated here; \cref{push:src:E.2.11}; \cref{push:src:E.2.12}}


\begin{proposition}[persistence of an isolated finite multiple fibre]\label{push:src:E.2.13}
\StatementRecord{push:src:E.2.13}{proved}{Historical events: Claude Fable 5.2 (Anthropic), effort high, 16 November 2026: carried as background without fault found. GPT-6 Astra (OpenAI), ChatGPT Pro mode, effort Pro, S6\_09\_PAPERS, 17 November 2026: targeted proof check. Also checked in this writing. Pre-continuation writing events: GPT-6 Astra (OpenAI), ChatGPT Pro mode, effort Pro, S6\_10\_PAPERS, 18 November 2026: targeted proof check. Continuation metadata events: GPT-6 Astra (OpenAI), ChatGPT Pro mode, effort Pro, S6\_10\_PAPERS continuation, 19 November 2026: used and rechecked in this run; status-vocabulary reconciliation; metadata reconciliation only; no new mathematical check. \emph{Pre-continuation recorded proof scope.} The divisor equality g*Gamma=mD, normalization by the explicit graphs r=zeta a u\textasciicircum{}(1/m), finite etaleness over the actual total space, connected Stein factor, holomorphic period roots, and torsor purity were re-derived. The central multiplicity is exactly the marked m}{the hypotheses, notation and definitions stated here; \cref{push:src:E.2.10}; \cref{push:src:E.2.12}; \cref{push:src:E.2.11}}

In the setting of \ref{push:src:E.2.10}, suppose the central value of
one finite labelled critical section is distinct from
the limits of the other two labelled sections.  Write
its multiplicity as $m$, with $m=3k$ or $4l$, and its
linear monodromy order as $e=3$ or $4$, respectively.
Near that section, normalize the base change taking
an $m$th root of its defining equation.

The resulting total space is smooth, and its map to
the root bidisc is a proper smooth family of two-tori.
The deck group acts freely on the total space.
Consequently the actual central reduced finite fibre
is smooth, has exactly multiplicity $m$, and belongs
to the same smooth family of finite free affine
quotients as the punctured reduced finite fibres.

This applies to both finite fibres when the three
labelled critical values have distinct limits.  In
the two-critical-value collapsed-target case, it
applies to the isolated finite labelled fibre.
\end{proposition}


\begin{proof}
\noindent\textbf{Step 1.} Choose target coordinates $(v,s)$ in a bidisc about
the section, and let $\Gamma$ be its graph.  Take a
local defining function $\xi$ for $\Gamma$.
On the smooth fourfold the divisor identity is
\[
 g^*\Gamma=m\mathcal D,
\]
where $\mathcal D$ is the prime closure of the
punctured reduced multiple-fibre family.  Indeed
this equality holds away from $s=0$.  The only
possible extra divisorial component would be
$X_0$, but $g_0$ is nonconstant, so $X_0$ is not
mapped into the point $\Gamma(0)$.  The closure
$\mathcal D$ is Cartier because the total space
is smooth.  Thus locally, even if its central
restriction is initially singular or reducible,
\[
 g^*\xi=a^m u
\]
for a local equation $a$ of $\mathcal D$ and a
holomorphic unit $u$.

\noindent\textbf{Step 2.} Let $B'\to B$ be the smooth root bidisc
$r^m=\xi$, and let $Y$ be the normalization
of $\mathcal X_B\times_B B'$.  A local $m$th
root of $u$ factors its defining equation as
\[
 r^m-a^m u=
 \prod_{\zeta^m=1}(r-\zeta a u^{1/m}).
\]
Its normalized branches are separate copies of
the original smooth local fourfold.  Therefore
$Y\to\mathcal X_B$ is a finite etale cover,
and $Y$ is smooth.  The deck group permutes these
branches freely, including over $a=0$.
This is an assertion about the normalized cover;
the unnormalized fibre product generally has
intersecting branches.

\noindent\textbf{Step 3.} The induced map $g':Y\to B'$ is proper and
equidimensional of relative dimension two.
Since both source and base are smooth, the
two pulled-back base parameters form a regular
sequence in each local ring, giving flatness.
Over $s\ne0$ this is the actual smooth root
torus family.  Its connectedness makes $Y$
connected and its generic fibre connected.
The Stein factor of $g'$ is finite birational
over the normal bidisc and hence is the
bidisc itself, so all fibres are connected.
Equivalently this follows from normality of
$Y$ and the finite birational form of its
Stein factorization.

\noindent\textbf{Step 4.} By \ref{push:src:E.2.12} the original central map has no
critical value away from the labelled limits.
Our choice of the section therefore makes
$g'$ smooth away from the origin of $B'$.
It remains to supply a smooth Jacobian over
the full root bidisc.

\noindent\textbf{Step 5.} If the normalizing projective transformations
$A_s$ have an invertible limit, the source
coordinate at the finite marked value has
the form
\[
 \xi_{\rm src}=\xi\,u_1(v,s)
\]
with $u_1$ a unit.  Thus its order-$e$ root
is $r^{m/e}u_1^{1/e}$, and the exact finite-root
period model with parameter $c(s)$ extends
smoothly to the full root bidisc.

\noindent\textbf{Step 6.} If $A_s$ has rank-one limit, the isolated
section in the hypothesis must be the
preimage of its constant image $t_*$:
all other marked preimages tend to the
exceptional point.  Thus $t_*$ is precisely
the finite marked value under consideration.
Near the isolated section the source
coordinate has divisor equal to that section
plus a positive multiple of $s=0$, so
\[
 \xi_{\rm src}=s^n\xi\,u_1(v,s),\qquad n>0,
\]
with a nowhere-zero unit $u_1$.  Choose a
nearby central regular target point with
$\xi\ne0$.  Its actual smooth torus family
has identity parameter monodromy.  As in
\ref{push:src:E.2.12}, $c(s)$ has no winding and hence
$e\mid n$.  The required holomorphic root
coordinate on $B'$ is now
\[
 s^{n/e}r^{m/e}u_1^{1/e}.
\]
It again gives a smooth reference Jacobian
over $B'$.  On the punctured bidisc its
marking and holomorphic period subbundle
agree with the actual Jacobian by the same
simply-connected argument as \ref{push:src:E.2.12}.

\ref{push:src:E.2.11} makes $g'$ a smooth torsor over the
entire root bidisc.  In particular $r$ is
a submersion coordinate on the central
threefold upstairs.  The unramified deck
cover then shows that the original fibre
has exactly multiplicity $m$.  Restricting
to $r=0$ gives a proper smooth family of
two-tori in the parameter $s$, with free
affine deck action.  Its quotient is a
proper smooth family of the reduced finite
fibres, proving all assertions.
\end{proof}



\paragraph{Scope of the finite-chart argument}

This argument identifies the smooth root charts
at the isolated finite critical values.  The
excluded target-collision cases are treated in
\ref{push:src:E.2.14}--\ref{push:src:E.2.15}.  The actual cusp is recovered
separately in \ref{push:src:E.2.16}--\ref{push:src:E.2.20}, and \ref{push:src:E.2.21} glues
these comparisons into the full one-sided
completion theorem.

\Uses{the hypotheses, notation and definitions stated here; \cref{push:src:E.2.14}; \cref{push:src:E.2.15}; \cref{push:src:E.2.16}; \cref{push:src:E.2.20}; \cref{push:src:E.2.21}}


\begin{proposition}[extension of the circle field and exclusion of finite-root collapse]\label{push:src:E.2.14}
\StatementRecord{push:src:E.2.14}{proved}{Historical events: Claude Fable 5.2 (Anthropic), effort high, 16 November 2026: carried as background without fault found. GPT-6 Astra (OpenAI), ChatGPT Pro mode, effort Pro, S6\_09\_PAPERS, 17 November 2026: targeted proof check. Also checked in this writing. Pre-continuation writing events: GPT-6 Astra (OpenAI), ChatGPT Pro mode, effort Pro, S6\_10\_PAPERS, 18 November 2026: targeted proof check. Continuation metadata events: GPT-6 Astra (OpenAI), ChatGPT Pro mode, effort Pro, S6\_10\_PAPERS continuation, 19 November 2026: used and rechecked in this run; status-vocabulary reconciliation; metadata reconciliation only; no new mathematical check. \emph{Pre-continuation recorded proof scope.} The primitive generator cannot vanish identically on the central fibre, its integral period extends, and compact averaging preserves the fixed curve and its two normal weights. A collapsed finite-root period would give a nontrivial compact complex group action with a fixed point, contradicting its trivial derivative representation}{the hypotheses, notation and definitions stated here; \cref{push:src:E.2.10}; \cref{push:src:E.2.12}; \cref{push:src:E.2.11}}

In the setting of \ref{push:src:E.2.10}, the normalized vector field
$\delta$ of the punctured family extends to a holomorphic
relative vector field on $\mathcal X$ after shrinking
the parameter disc.  Its time-one flow is the identity,
and its real circle action on the central fibre is
equivariantly diffeomorphic to the prescribed circle
action on the punctured fibres.  In particular its
central fixed locus is a smooth rational curve with
normal bundle $\mathcal O(-1)\oplus\mathcal O(-1)$
and the same two opposite circle weights.

The rank-one limit of the normalizing projective
transformations $A_s$ cannot have constant image
$0$ or $1$.  More generally, a collapsed constant
image is impossible whenever the limiting torus's
complex $\delta$ line has two real-linearly
independent periods.  Thus the finite-root
isotrivial possibilities in \ref{push:src:E.2.12} are excluded.
\end{proposition}


\begin{proof}
\noindent\textbf{Step 1.} Let $\pi:\mathcal X\to\Delta$ be the original
proper holomorphic submersion.  The coherent sheaf
\[
 \mathcal H=\pi_*T_{\mathcal X/\Delta}
\]
is torsion-free: multiplication by a nonzero base
function is injective on the vector bundle upstairs
and hence on its sections.  Its generic rank is
one by \cref{push:src:E.2.1}.  On the smooth disc it is therefore
locally free of rank one.  Choose a generator $V$
after shrinking the disc.  On $\Delta^*$ it has
the form
\[
 V|_{X_s}=b(s)\delta_s,
\]
with $b$ holomorphic and nowhere zero there.
The normalized fields $\delta_s$ agree on overlaps
of local classifying maps: at fixed discrete data
the comparison maps of \cref{push:src:E.1.3} have fibre-linear
part the identity and fix the marked period
$\widehat\delta$.  Their translations also
preserve the invariant vector field.
Moreover $V$ is vertical for $g$, because this
holds on the dense punctured family.

\noindent\textbf{Step 2.} Choose a regular central target value and a small
bidisc around it on which $g$ is smooth.  The
ordinary or finite-root comparison used in
\ref{push:src:E.2.12} supplies its smooth marked Jacobian, with
the marked period $\widehat\delta$ extending as
a nonzero holomorphic Lie section.  The resulting
invariant field $\delta_B$ on the actual torsor
is nonvanishing there.  Comparing one nonzero
coefficient of $\delta_B$ with $V$ shows that
$b(s)$ extends holomorphically to $s=0$.
Its independence of the target coordinate follows
from that equality on $s\ne0$ and holomorphic
uniqueness.

\noindent\textbf{Step 3.} If $b(0)=0$, the restriction $V|_{X_0}$ vanishes
on a nonempty open subset of connected $X_0$,
and therefore vanishes everywhere.  In local
coordinates for the submersion $\pi$, every
coefficient of $V$ is then divisible by $s$.
The quotients glue uniquely to a global
holomorphic relative vector field $V/s$.
This contradicts the fact that $V$ generates
the free module $\mathcal H$.  Hence $b(0)\ne0$,
and $V/b(s)$ gives the required global extension
of $\delta$.

\noindent\textbf{Step 4.} The extended field is complete on each compact
fibre of $\pi$.  Its time-one flow is the
identity on the punctured family and hence on
the whole family by holomorphic uniqueness.
Its real flow consequently defines an action
of $\mathbb R/\mathbb Z$ on $\mathcal X$
over the parameter disc.  Average a smooth
metric and an Ehresmann horizontal distribution
over this compact group.  Parallel transport
then gives equivariant diffeomorphisms of
the fibres.  The central action is therefore
the same effective semifree action, with a
single fixed two-sphere, as on the prescribed
fibres.

\noindent\textbf{Step 5.} The fixed set of a holomorphic compact-group
action is a complex submanifold, as follows
from the averaging coordinates used in
\ref{push:src:E.2.11}.  Here it is a proper smooth family
of curves of genus zero, and hence each fibre
is $\mathbb P^1$.  Its two holomorphic normal
weight lines have the same degrees as in
the punctured family.  Those degrees are
$-1,-1$ by C.3.11, or directly by the two
fixed toric edge normals.  The characters
distinguish the lines, so their central
restrictions are exactly
$\mathcal O(-1)\oplus\mathcal O(-1)$.

\noindent\textbf{Step 6.} Suppose now that $A_s$ has rank-one limit
with constant image a finite marked value.
\ref{push:src:E.2.12} identifies every central regular
torus, with its marked $\delta$ direction,
with the corresponding finite-root torus
at parameter $c(0)$.  The invariant lattice
of $A_j$ has rank two and spans its
one-dimensional complex fixed subspace.
It contains $\widehat\delta$; its second
primitive period is independent over
$\mathbb R$.  The complex fixed subspace
is precisely the $\delta$ line, since
the other complex eigenvalue is a
primitive third or fourth root of unity.
Consequently, in the normalization in
which the $\delta$ period is $1$, there
is a constant second period
\[
 \kappa\in\mathbb C\setminus\mathbb R.
\]
This is a period of the invariant field
on every central regular fibre.  The
time-$\kappa$ flow is therefore the
identity on their dense union, and
hence on all of $X_0$.

\noindent\textbf{Step 7.} The flow of the nonzero field $\delta_0$
thus factors through a nontrivial
holomorphic action of the compact
elliptic curve
$E=\mathbb C/(\mathbb Z+\kappa\mathbb Z)$.
Such an action has no fixed point on
a connected complex manifold unless
it is trivial.  Indeed its derivative
at a fixed point is a holomorphic
representation of the compact complex
group into a general linear group,
so is constant.  Averaging local
coordinates over the compact group
then makes the action trivial near
that point, and holomorphic uniqueness
makes it trivial everywhere.

\noindent\textbf{Step 8.} The elliptic action here is nontrivial
because $\delta_0\ne0$.  It therefore
has no fixed point.  This contradicts
the fixed rational curve of the
extended real circle action: a zero
of the generating field is fixed
by its entire complex flow.  Equivalently,
the resulting nowhere-zero real vector
field would give $\chi(X_0)=0$,
whereas $X_0$ is diffeomorphic to $S^6$.
The same argument applies to an ordinary
constant-image torus if its $\delta$
line has a second independent period.
\end{proof}



\paragraph{What the elliptic-action argument does not assume}

The argument uses the global one-dimensional
space of vector fields on the known punctured
members and the actual marked period comparison
on a central regular bidisc.  It does not
assume that all translations of an ordinary
central torus extend over the singular fibre.
For an ordinary torus the $\delta$ line may
have only the one specified real period;
such a case is not excluded by the
elliptic-action argument alone.

\Uses{the hypotheses, notation and definitions stated here; the argument and qualifications immediately following}


\begin{proposition}[an isotrivial volume obstruction]\label{push:src:E.2.15}
\StatementRecord{push:src:E.2.15}{proved}{Historical events: Claude Fable 5.2 (Anthropic), effort high, 16 November 2026: carried as background without fault found. GPT-6 Astra (OpenAI), ChatGPT Pro mode, effort Pro, S6\_09\_PAPERS, 17 November 2026: targeted proof check. Also checked in this writing. Pre-continuation writing events: GPT-6 Astra (OpenAI), ChatGPT Pro mode, effort Pro, S6\_10\_PAPERS, 18 November 2026: targeted proof check. Continuation metadata events: GPT-6 Astra (OpenAI), ChatGPT Pro mode, effort Pro, S6\_10\_PAPERS continuation, 19 November 2026: used and rechecked in this run; status-vocabulary reconciliation; metadata reconciliation only; no new mathematical check. \emph{Pre-continuation recorded proof scope.} The countable marked period orbit forces constant periods. The exact pluricanonical coefficient P=q\textasciicircum{}(-D-2A) phi(1/q), finite volume and subharmonic mean estimate imply P=q\textasciicircum{}(1-A)H. For D>=-A the entire phi extends with zero at infinity, a contradiction}{the hypotheses, notation and definitions stated here; \cref{push:src:E.2.12}; \cref{push:src:E.2.10}}

Let $f:X\to\PP^1$ be a proper holomorphic map with connected
fibres from a compact connected complex threefold.  Suppose
that for some $p\in\PP^1$ the restriction over
$U=\PP^1\setminus\{p\}$ is a smooth family of mutually
isomorphic complex two-tori.  If
\[
 K_X^{\otimes A}\simeq f^*\mathcal O_{\PP^1}(D),
 \qquad A\in\Z_{>0},\quad D\in\Z,
 \tag{E.2.15.1}
\]
then necessarily $D<-A$.  In particular the ordinary-source
collapsed-target case of \ref{push:src:E.2.12} cannot occur for the
six-sphere families under consideration.
\end{proposition}


\begin{proof}
\noindent\textbf{Step 1.} Choose an affine coordinate $t$ on $U\simeq\C$ and put
$q=t^{-1}$ near $p$.  Since $U$ is simply connected, its
integral first-homology local system has a global marking.
The holomorphic period subbundle is constant: locally the
marked periods belong to the countable integral-change-of-basis
orbit of one fixed torus, and a continuous map from a connected
base into such a countable subset of a complex Grassmannian
is constant.  For the last assertion, any two distinct image
points can be separated by a real continuous coordinate;
its connected image would contain an interval and would be
uncountable.  Thus the Jacobian group family is the product
with one fixed torus $T$.

\noindent\textbf{Step 2.} Choose a nonzero translation-invariant holomorphic two-form
$\eta_T$ on $T$.  Local holomorphic origins identify our
family with its Jacobian, and two choices differ by a
translation.  Hence $\eta_T$ defines a global nonvanishing
relative canonical form $\eta$ on $X_U$.  Its fibre volume
is a fixed positive number:
\[
 \int_{X_t}\frac14\eta\wedge\overline\eta=V>0.
 \tag{E.2.15.2}
\]
This construction uses neither a global torsor section nor
a polarization on $X$.  The canonical determinant
isomorphism for a submersion gives the nonvanishing
absolute three-form $\eta\wedge dt$ on $X_U$; this notation
does not require choosing an absolute lift of $\eta$.

\noindent\textbf{Step 3.} Use (E.2.15.1) to pull back a meromorphic section of
$\mathcal O(D)$ whose divisor is $D[p]$.  We obtain a
meromorphic pluricanonical form $\sigma$ with
\[
 \operatorname{div}(\sigma)=D f^*(p).
\]
On $X_U$ it is holomorphic and nowhere zero.  Its quotient
by $(\eta\wedge dt)^{\otimes A}$ is a holomorphic function
constant on each compact connected fibre.  It therefore
descends to an entire nowhere-zero function $\phi$ on $\C$:
\[
 \sigma=\phi(t)(\eta\wedge dt)^{\otimes A}.
 \tag{E.2.15.3}
\]

\noindent\textbf{Step 4.} On a small disc about $p$, the meromorphic section
$q^{-D}\sigma$ is in fact a holomorphic nowhere-zero section
of $K_X^{\otimes A}$ over the whole inverse image.  This
follows directly from (E.2.15.1), or by cancellation of the
actual divisors, with their multiplicities.  Denote it by
$\Omega_A$.  On the punctured inverse image, since
$dt=-q^{-2}dq$, it is
\[
 \Omega_A=(-1)^A P(q)(\eta\wedge dq)^{\otimes A},
 \qquad P(q)=q^{-D-2A}\phi(1/q).
 \tag{E.2.15.4}
\]
The function $P$ is single-valued and holomorphic on the
punctured disc.  No extension of a comparison with another
threefold is asserted here.

\noindent\textbf{Step 5.} A nowhere-zero holomorphic $A$-canonical form determines
a smooth positive density: if it is
$h(z)(dz_1\wedge dz_2\wedge dz_3)^{\otimes A}$, take
$|h|^{2/A}$ times the ordinary complex-coordinate volume
density.  This is independent of the coordinates.  The
density associated with $\Omega_A$ has finite integral
over the inverse image of a smaller closed disc, because
that inverse image is compact.  Fibre integration of
(E.2.15.4), using (E.2.15.2), consequently gives
\[
 \int_{0<|q|<\epsilon}|P(q)|^{2/A}\,d\mathrm{Area}(q)
 <\infty.
 \tag{E.2.15.5}
\]
Only a fixed positive normalization constant is suppressed.
This step concerns a smooth positive density on the actual
total space and requires no K\"ahler metric.

\noindent\textbf{Step 6.} We spell out the isolated-singularity consequence of
(E.2.15.5), including when $2/A<1$.  The function
$|P|^{2/A}$ is subharmonic.  If $|q|=r$, the mean-value
inequality on the disc of radius $r/2$ centred at $q$ gives
\[
 |P(q)|^{2/A}
 \leq \frac{4}{\pi r^2}
       \int_{r/2<|z|<3r/2}|P(z)|^{2/A}\,d\mathrm{Area}(z).
\]
The last integral tends to zero as $r\to0$, by
integrability.  Hence $r^A|P(q)|\to0$, uniformly on
$|q|=r$.  The function $q^A P(q)$ therefore has a removable
singularity and vanishes at zero.  For some holomorphic
$H$ on the disc,
\[
 P(q)=q^{1-A}H(q).
 \tag{E.2.15.6}
\]
Equivalently, the pole order of $P$ is at most $A-1$.

\noindent\textbf{Step 7.} Combining (E.2.15.4) and (E.2.15.6) yields
\[
 \phi(1/q)=q^{D+A+1}H(q).
 \tag{E.2.15.7}
\]
If $D\geq-A$, the right-hand side extends to zero at
$q=0$.  Thus $\phi$ extends holomorphically to $\PP^1$
and vanishes at $p$.  A holomorphic function on $\PP^1$
is constant, so $\phi=0$, contradicting its construction.
This proves $D<-A$.

\noindent\textbf{Step 8.} In \ref{push:src:E.2.10} the actual canonical-root equality has
\[
 A=12kl,\qquad D=kl-3k-4l,
 \qquad
 \frac DA=\frac1{12}-\frac1{4l}-\frac1{3k}
       \geq-\frac12>-1.
\]
If the collapsed constant source point is ordinary,
\ref{push:src:E.2.12} says that the central map has precisely one
critical value and is isotrivial on its complement.
All hypotheses above apply to the actual central
threefold, and give the required contradiction.
\end{proof}



\paragraph{Scope of the volume argument}

The proof uses the equality of line bundles in
(E.2.15.1), not just the topological vanishing of
$c_1(K_X)$.  It also uses constancy of the torus periods,
which makes the fibre volume in (E.2.15.2) constant.
It does not assert a general good-reduction or
comparison theorem for non-K\"ahler torus degenerations.
Together with the separate elliptic-action exclusion
in \ref{push:src:E.2.14}, it removes the two types of collapsed
target permitted by \ref{push:src:E.2.12}.  Identification of a
noncollapsed interior completion, and the behaviour
on arbitrary new deformation branches, remain
distinct questions.

\Uses{the hypotheses, notation and definitions stated here; \cref{push:src:E.2.14}; \cref{push:src:E.2.12}}


\begin{proposition}[the central cusp is reduced and irreducible]\label{push:src:E.2.16}
\StatementRecord{push:src:E.2.16}{proved}{Historical events: Claude Fable 5.2 (Anthropic), effort high, 16 November 2026: carried as background without fault found. GPT-6 Astra (OpenAI), ChatGPT Pro mode, effort Pro, S6\_09\_PAPERS, 17 November 2026: targeted proof check. Also checked in this writing. Pre-continuation writing events: GPT-6 Astra (OpenAI), ChatGPT Pro mode, effort Pro, S6\_10\_PAPERS, 18 November 2026: targeted proof check; used and rechecked in this run. Continuation metadata events: GPT-6 Astra (OpenAI), ChatGPT Pro mode, effort Pro, S6\_10\_PAPERS continuation, 19 November 2026: used and rechecked in this run; status-vocabulary reconciliation; metadata reconciliation only; no new mathematical check. \emph{Pre-continuation recorded proof scope.} The common invariant covectors and complement count give three irreducible supports. The explicit triangle-group representations contradict simple connectedness when the cusp multiplicity exceeds one. Euler additivity gives two, and Honda-Viaclovsky v1 Corollary3.7 p.11 supplies the irreducible H1 invariant-cycle equality}{the hypotheses, notation and definitions stated here; \cref{push:src:E.2.10}; \cref{push:src:E.2.17}; \cref{push:src:E.2.20}; \cref{push:src:E.2.12}; \cref{push:src:E.2.14}; \cref{push:src:E.2.15}; \cref{push:src:E.2.13}; \cref{push:src:E.2.7}; \cite[Corollary 3.7, arXiv version~1]{HondaViaclovsky2025D}}

In the one-sided smooth specialization of \ref{push:src:E.2.10},
the three labelled critical values remain distinct.
The two finite fibres have their original exact
multiplicities $m=3k$ and $n=4l$.  Every reduced
central critical fibre is irreducible, and the
cusp fibre has multiplicity one.  Its underlying
surface $S$ satisfies
\[
 \chi_{\rm top}(S)=2,\qquad b_1(S;\mathbb Q)=2.
\]
This supplies the reduced irreducible Cartier central
surface used in the analytic cusp identification of
\ref{push:src:E.2.17}--\ref{push:src:E.2.20}.
\end{proposition}


\begin{proof}
\noindent\textbf{Step 1.} \ref{push:src:E.2.12}, \ref{push:src:E.2.14} and \ref{push:src:E.2.15} exclude every degenerate
limit of the normalizing projective transformations.
They therefore extend as invertible projective
transformations, and the three labelled critical
values remain distinct.  \ref{push:src:E.2.13} proves the stated
finite-fibre multiplicities.

\noindent\textbf{Step 2.} Normalize the three sections to $0,1,\infty$.
Over $B^\circ=\mathbb P^1\setminus\{0,1,\infty\}$,
the relative map is a proper submersion over
$B^\circ\times\Delta$.  Hence its integral
first-homology monodromies on the central fibre
are the same matrices $A_1,A_2,A_0$ as in
\ref{push:src:E.2.7}.  The common invariant space of the dual
matrices $A_1^t,A_2^t$ has dimension one:
in the displayed marking it is generated by
the covector $\gamma$.  This is obtained by
solving their simultaneous fixed equations;
the exact calculation is recorded in
\path{code/e_push_interior_cusp.py}.

\noindent\textbf{Step 3.} Let $r$ be the total number of irreducible
components of the three reduced critical
fibres, and let $X^\circ$ be their complement
in $X_0$.  Alexander duality on $S^6$ gives
\[
 b_1(X^\circ;\mathbb Q)=r.
\]
The degree-one Leray sequence over the
punctured projective line, whose second
cohomology vanishes, gives instead
\[
 b_1(X^\circ;\mathbb Q)
 =b_1(B^\circ;\mathbb Q)
   +\dim H^1(T;\mathbb Q)^{\rm mon}
 =2+1=3.
\]
There are three nonempty critical fibres,
so each has exactly one reduced component.

\noindent\textbf{Step 4.} Write the cusp fibre as $dS$ with $d\ge1$.
We prove $d=1$ by a direct meridian argument.
The smooth fibration over $B^\circ$ has
connected fibres, hence
$\pi_1(X^\circ)\to\pi_1(B^\circ)$ is
surjective.  The inclusion
$X^\circ\hookrightarrow X_0$ is surjective
on fundamental groups, and its kernel is
normally generated by meridians about the
three reduced divisor components.  To see
the latter assertion, make a null-homotopy
disc transverse to their smooth loci and
disjoint from their singular loci, which
have real codimension at least four.
Removing small discs about the finitely
many intersections expresses its boundary
as a product of conjugate meridians.

\noindent\textbf{Step 5.} A meridian of a multiplicity-$a$ component
maps to the $a$th power of the corresponding
base meridian, since locally the base
function is a unit times the $a$th power
of a reduced local equation.  Consequently
there is a surjection
\[
 \pi_1(X_0)\longrightarrow
 \left\langle a,b,c\ \middle|\
 a^m=b^n=c^d=abc=1\right\rangle.
 \tag{E.2.16.1}
\]
This is the usual orbifold fundamental
group, but the surjection has just been
proved without assuming a stack
factorization.  Its source is trivial.

\noindent\textbf{Step 6.} For clarity the group on the right is
nontrivial for every $d\ge2$ here.
Except for $(m,n,d)=(3,4,2)$, one has
\[
 \frac1m+\frac1n+\frac1d<1.
\]
Indeed for $d\ge3$ the sum is at most
$11/12$; for $d=2$, increasing either
$m=3k$ or $n=4l$ above $3,4$ makes it
at most $23/24$.  Put
$\alpha=\pi/m$, $\beta=\pi/n$,
$\theta=\pi/d$, and choose $\ell>0$ with
\[
 \cosh\ell=
 \frac{\cos\theta+\cos\alpha\cos\beta}
      {\sin\alpha\sin\beta}>1.
\]
Let
\[
 R_\alpha=\begin{pmatrix}
 \cos\alpha&\sin\alpha\\
 -\sin\alpha&\cos\alpha
 \end{pmatrix},\qquad
 D_\ell=\operatorname{diag}(e^{\ell/2},e^{-\ell/2}).
\]
In $\operatorname{PSL}_2(\mathbb R)$ the
elements represented by $R_\alpha$ and
$D_\ell R_\beta D_\ell^{-1}$ have orders
$m,n$, and their product has trace
$-2\cos\theta$ and projective order $d$.
Taking its inverse for $c$ gives a
nontrivial representation of (E.2.16.1).
For the exceptional triple, the
permutations $a=(123)$ and $b=(1324)$
have orders three and four, and
$ab=(24)$ has order two, giving a
nontrivial representation in $S_4$.
Thus $d\ge2$ contradicts (E.2.16.1),
and the cusp is reduced.

\noindent\textbf{Step 7.} The ordinary torus fibres have Euler
characteristic zero.  The finite reduced
fibres are smooth finite free quotients
of two-tori by \ref{push:src:E.2.13}, so also have
Euler characteristic zero.  Additivity
over the three singular values therefore
gives $\chi_{\rm top}(S)=\chi(X_0)=2$.
Finally, the valid irreducible-fibre
degree-one invariant-cycle statement
\cite[Corollary 3.7, arXiv version~1]{HondaViaclovsky2025D}
gives
\[
 b_1(S;\mathbb Q)=
 \dim H^1(T;\mathbb Q)^{A_0}=2.
\]
Its irreducibility hypothesis has just
been proved.  This use is separate from
the false unrestricted equality
converse audited in \cref{push:src:B.4.17}.
\end{proof}



\begin{proposition}[persistence of the normal-crossing cusp]\label{push:src:E.2.17}
\StatementRecord{push:src:E.2.17}{proved}{Historical events: Claude Fable 5.2 (Anthropic), effort high, 16 November 2026: carried as background without fault found. GPT-6 Astra (OpenAI), ChatGPT Pro mode, effort Pro, S6\_09\_PAPERS, 17 November 2026: targeted proof check. Also checked in this writing. Pre-continuation writing events: GPT-6 Astra (OpenAI), ChatGPT Pro mode, effort Pro, S6\_10\_PAPERS, 18 November 2026: targeted proof check; used and rechecked in this run. Continuation metadata events: GPT-6 Astra (OpenAI), ChatGPT Pro mode, effort Pro, S6\_10\_PAPERS continuation, 19 November 2026: used and rechecked in this run; status-vocabulary reconciliation; metadata reconciliation only; no new mathematical check. \emph{Pre-continuation recorded proof scope.} The finite critical scheme conserves Milnor number. Lazzeri Theorem2 and its corollary/remark, printed pp.274-275, exclude multiple singular points with the full original Milnor number. The circle-weight factor q=Tg(T,z) and the two simple zero sections give actual relative xy and xyz charts}{the hypotheses, notation and definitions stated here; \cref{push:src:E.2.10}; \cref{push:src:E.2.14}; \cref{push:src:E.2.16}; \citet[Theorem 2 and its corollary, pp.~274--275]{Lazzeri1973Root}}

In the smooth fixed-pair specialization of \ref{push:src:E.2.10},
the cusp divisor is a relative normal-crossing divisor
over the parameter disc.  More precisely, after
shrinking that disc, at every point of the total cusp
divisor there are holomorphic coordinates
$(z_1,z_2,z_3,s)$, with $s$ the original parameter,
in which the cusp coordinate of the intrinsic map is
one of
\[
 q=z_1,\qquad q=z_1z_2,\qquad q=z_1z_2z_3.
 \tag{E.2.17.1}
\]
There are exactly two relative triple-point sections.
Both lie on the fixed rational curve of \ref{push:src:E.2.14}.
In particular the normalization of the total cusp
divisor is smooth and proper over the parameter
disc, and its central fibre is the normalization
of the central cusp.
\end{proposition}


\begin{proof}
\noindent\textbf{Step 1.} By \ref{push:src:E.2.14}--\ref{push:src:E.2.16} the three labelled critical
values remain distinct, the cusp has multiplicity
one, and its reduced central fibre is irreducible.
We normalize its target section to $q=0$.
On a sufficiently small target disc, zero is the
only critical value on every parameter fibre.
The extended field $\delta$ is zero precisely on
the relative fixed rational curve $\mathcal C$,
and the two normal characters there are $+1,-1$.
For $s\ne0$ the two triple points of the prescribed
cusp belong to $\mathcal C_s$; outside that curve
the cusp is smooth or has an ordinary double crossing.

\noindent\textbf{Step 2.} We first record exactly the local singularity result
we use.  If $h_s(a,b)$ is a holomorphic family and
$h_0=0$ is a reduced plane-curve germ with an
isolated critical point, choose a sufficiently
small representative of that germ.  The sum of
the multiplicities of all nearby critical points
of $h_s$, regardless of their critical values,
is $\mu(h_0)$.  One can see this directly from
the analytic critical scheme
\[
 \mathbb C\{a,b,s\}/
       ((h_s)_a,(h_s)_b).
\]
After taking a sufficiently small representative,
it is finite over the parameter disc.  Its two
equations are a regular sequence, and its
central fibre is zero-dimensional.  It is thus
Cohen--Macaulay of dimension one, with $s$ a
nonzerodivisor, and is finite flat over the
disc.  Its degree is the central Jacobian-algebra
dimension, namely $\mu(h_0)$.

\noindent\textbf{Step 3.} The non-splitting theorem of
\citet[Theorem 2 and its corollary, pp.~274--275]{Lazzeri1973Root}
says that if a deformation of an isolated
hypersurface singularity has more than one
singular point in a single deformed fibre,
their summed Milnor numbers are strictly
smaller than the original Milnor number.
The strict numerical inequality is the
conclusion of the cited corollary's proof.
Consequently, if all the critical points
just counted have critical value zero,
there can be only one such point.  This
application uses a flat deformation of
plane hypersurfaces: the map
$(a,b,s)\mapsto(h_s(a,b),s)$ is flat near
the central germ because its two base
parameters form a regular sequence in
the smooth source.  No assertion of
separatedness of compact threefolds enters
this local argument.

\noindent\textbf{Step 4.} Let $p$ be a central cusp point outside
$\mathcal C_0$.  A relative holomorphic
flow box for the nonzero field $\delta$
gives coordinates $(a,b,v,s)$ in which
$\delta=\partial_v$.  Invariance of the
intrinsic map gives
\[
 q=h_s(a,b).
\]
The central equation is reduced, by
\ref{push:src:E.2.16}.  A reduced plane-curve germ is
either smooth or has an isolated critical
point.  In the smooth case the implicit
function theorem already gives the first
form in (E.2.17.1).

\noindent\textbf{Step 5.} In the singular case every critical point
of $h_s$ in a small representative lifts
to a critical point of the actual map
on $X_s$.  All its critical values are
therefore zero.  For $s\ne0$ each such
point is a node, since the flow box is
disjoint from the fixed curve and hence
from the two triple points.  Conservation
and non-splitting give one node and
$\mu(h_0)=1$.  Thus the central Hessian
is nondegenerate.  The implicit function
theorem gives its unique critical-point
section, whose critical value is zero.
The holomorphic Morse lemma with
parameters then gives $q=ab$ throughout
this neighbourhood, with $s$ unchanged.

\noindent\textbf{Step 6.} It remains to treat the fixed curve.
Holomorphic coordinates obtained by
averaging over the extended compact
circle may be chosen in the form
$(z,x,y,s)$, with $\mathcal C$ defined
by $x=y=0$ and with circle weights
$0,+1,-1,0$.  The original parameter
$s$ is retained by this averaging.
The holomorphic invariant function $q$
vanishes on $\mathcal C$, so its
convergent power series has the form
\[
 q=F_s(T,z)=T g_s(T,z),\qquad T=xy.
 \tag{E.2.17.2}
\]
The coefficient
$Q_s(z)=g_s(0,z)$ is intrinsically a
section of the dual of the product of
the two normal weight lines along
$\mathcal C_s$.  By \ref{push:src:E.2.14} this is
$\mathcal O_{\mathbb P^1}(2)$.

\noindent\textbf{Step 7.} The central section $Q_0$ is not
identically zero.  Otherwise the
central equation in every such chart
would be divisible by $x^2y^2$,
contradicting reducedness of the
central cusp.  Its zero divisor
therefore has degree two.  Suppose
that both zeros coincide at a point
$z=0$.  The plane equation
$F_0(T,z)=Tg_0(T,z)$ is reduced:
a repeated factor of $g_0$ would
remain repeated after $T=xy$, and
$T$ does not divide $g_0$ because
$Q_0$ is not identically zero.
Thus $F_0$ has an isolated plane
critical point there.

\noindent\textbf{Step 8.} For $s\ne0$ the two zeros of $Q_s$
are the two distinct triple points
of the prescribed cusp.  Under the
supposed collision they both lie
in the chosen small plane
representative.  They are distinct
critical points of $F_s$ with value
zero.  Every other critical point
of $F_s$ also has value zero:
this is automatic at $T=0$, and
at $T\ne0$ it follows by lifting
to $xy=T$, where the quotient
map $(x,y,z)\mapsto(T,z)$ is a
submersion.  All sufficiently small
$T$ have such a lift in the original
coordinate chart.  A nonzero
critical value would contradict
the actual critical-value set of
the intrinsic map.

\noindent\textbf{Step 9.} Conservation of all critical
multiplicities and Lazzeri's
non-splitting theorem now give a
contradiction to the two distinct
critical points.  Hence $Q_0$ has
two simple zeros.  They extend
as two disjoint holomorphic
sections over the parameter disc.
Where $Q_0\ne0$, the factor
$g_s(xy,z)$ is a unit and can
be absorbed into $y$, giving
$q=xy$.  At a simple zero,
$\partial g_0/\partial z\ne0$;
replace $z$ by $g_s(xy,z)$ to
obtain $q=xyz$, again keeping
$s$ unchanged.  This proves
all local forms in (E.2.17.1).

\noindent\textbf{Step 10.} Properness makes the central cusp
compact, so finitely many such
neighbourhoods cover it and a
single shrinking of the parameter
disc suffices.  The normalization
of a relative normal-crossing
divisor separates its local
branches.  In the coordinates
above every normalized branch
is smooth over $s$, and these
descriptions agree on overlaps.
Normalization is finite for
complex analytic spaces; its
composition with the proper
cusp map is therefore proper.
Its fibres are precisely the
normalizations of the cusp
fibres in these local charts.
This proves the final assertions.
\end{proof}



\paragraph{Literature and remaining identification}

Non-splitting is a known theorem of
Lazzeri, not a new singularity-theoretic
claim.  We opened the original paper
and use its strict numerical
corollary.  The application combines
that theorem with the actual extended
circle action, reducedness, and the
absence of extra critical values
proved above.  Searches for
non-splitting, Milnor-number
conservation and degenerations of
torus-fibred complex six-spheres
did not locate this application.  The
proposition supplies relative
normal-crossing charts and a
simultaneous normalization.  It
does not by itself identify the entire
analytic cusp neighbourhood.  The
prescribed marked Mumford completion
is recovered by \ref{push:src:E.2.18}--\ref{push:src:E.2.20}.

\Uses{the hypotheses, notation and definitions stated here; \cref{push:src:E.2.18}; \cref{push:src:E.2.20}}


\begin{proposition}[the normalized cusp and its marked hexagon]\label{push:src:E.2.18}
\StatementRecord{push:src:E.2.18}{proved}{Historical events: Claude Fable 5.2 (Anthropic), effort high, 16 November 2026: carried as background without fault found. GPT-6 Astra (OpenAI), ChatGPT Pro mode, effort Pro, S6\_09\_PAPERS, 17 November 2026: targeted proof check. Also checked in this writing. Pre-continuation writing events: GPT-6 Astra (OpenAI), ChatGPT Pro mode, effort Pro, S6\_10\_PAPERS, 18 November 2026: targeted proof check; used and rechecked in this run. Continuation metadata events: GPT-6 Astra (OpenAI), ChatGPT Pro mode, effort Pro, S6\_10\_PAPERS continuation, 19 November 2026: used and rechecked in this run; status-vocabulary reconciliation; metadata reconciliation only; no new mathematical check. \emph{Pre-continuation recorded proof scope.} Topology preserves the six rational conductor curves and degrees. The projective minus-one divisor contraction in Grauert1962 section4 no.8(a), printed pages364-365, was opened in reproduced original journal text, with primary-scan access separately unavailable. Contraction gives the rational b2-one pair with its triangle; exact evaluation and base change extend the logarithmic tangent frame}{the hypotheses, notation and definitions stated here; \cref{push:src:E.2.17}; \citet[\S4, no.~8(a), pp.~364--365]{Grauert1962E}; \cite[Chapter~VI, Theorem~1.1 and Table~10]{BHPV2004E}}

Let $\nu:\mathcal S\to\mathcal W$ be the
simultaneous normalization of the relative cusp
in \ref{push:src:E.2.17}, and let $\mathcal D\subset\mathcal S$
be its reduced conductor.  After shrinking the
parameter disc, each fibre pair is the standard
toric pair
\[
 (\mathcal S_s,\mathcal D_s)
 \simeq(\mathrm{dP}_6,\hbox{its toric hexagon}).
\]
The six conductor components may be labelled over
the disc so that their incidence, self-intersections
$-1$, opposite-pair involution, and endpoint
identifications are those of the punctured family.
The relative logarithmic tangent bundle has a
global holomorphic frame:
\[
 T_{\mathcal S/\Delta}(-\log\mathcal D)
 \simeq\mathcal O_{\mathcal S}^{\oplus2}.
\]
This assertion identifies the normalized pair
and supplies a relative frame; it does not by
itself identify the analytic smoothing neighbourhood.
\end{proposition}


\begin{proof}
\noindent\textbf{Step 1.} In the relative charts of \ref{push:src:E.2.17}, the normalized
branches are smooth over the original parameter.
Their conductor branches and their intersections
are also smooth over that parameter.  The
normalization of the conductor is consequently
a proper smooth family of curves, and the
normalized intersection strata are finite
etale over the disc.  Their connected components
and incidence are constant.  Since the disc
is simply connected, label the six curve
families $\mathcal D_1,\ldots,\mathcal D_6$
and the six vertex sections by continuation
from a punctured fibre.  Each curve fibre
is a smooth curve of genus zero.  Its two
distinct vertex sections remain distinct;
the same relative local charts exclude a
new self-intersection or an additional
incidence at the centre.  Thus the central
conductor is a hexagon of six embedded
rational curves.  The degree of each
relative normal line is constant in its
proper curve family, so all six central
self-intersections equal $-1$.

\noindent\textbf{Step 2.} Write $S=\mathcal S_0$ and $D_i=\mathcal D_{i,0}$.
Ehresmann transport gives $b_1(S)=0$ and
$b_2(S)=4$.  Contract the three disjoint
curves $D_1,D_3,D_5$.  The analytic
Castelnuovo contraction gives a smooth
compact complex surface $Y$; the needed
analytic statement, including smoothness
of the resulting point, is
\citet[\S4, no.~8(a), pp.~364--365]{Grauert1962E}.
Only these three contractions on the
central surface are used; no relative
contraction theorem is assumed.

\noindent\textbf{Step 3.} The images $C_2,C_4,C_6$ of the remaining
curves are smooth rational curves.  Each
meets two contracted curves transversely
once, so has self-intersection
$-1+1+1=1$.  Each pair meets once, at
the image of their intervening contracted
curve; these three intersection points
are distinct and the intersections are
transverse.  Thus they form a triangle.
Also $b_1(Y)=0$ and $b_2(Y)=1$.

\noindent\textbf{Step 4.} Any one of these curves, say $C=C_2$,
has nonnegative intersection with every
effective divisor: distinct irreducible
curves have nonnegative local intersection
multiplicities and $C^2=1$.  Adjunction
gives $K_Y\cdot C=-3$.  A nonzero section
of $K_Y^{\otimes r}$ for $r>0$ would
give an effective divisor of intersection
$-3r$ with $C$, impossible.  Hence
$\kappa(Y)=-\infty$.  The compact-surface
classification
\cite[Chapter~VI, Theorem~1.1 and Table~10]{BHPV2004E}
together with $b_1(Y)=0$ makes $Y$ rational:
class~VII has first Betti number one,
and a ruled surface has the first Betti
number of its base curve.  A rational
surface with $b_2=1$ is $\mathbb P^2$.
The curves $C_2,C_4,C_6$ are lines,
since their self-intersections are one.
Their three pairwise intersection points
are distinct, so they form a projective
triangle.  Blowing up its three vertices
recovers exactly the toric degree-six
del Pezzo surface and its hexagonal
boundary.  This proves the assertion
about the central pair; the punctured
pairs were already the prescribed ones.

\noindent\textbf{Step 5.} The conductor involution is holomorphic
over the disc.  Its permutation of the
six labelled curve families, and its
permutation of the vertex sections,
are locally constant.  They are therefore
the same opposite-pair and endpoint
identifications as on a punctured fibre.
The remaining scalar gluing parameters
are not claimed to be determined by
this discrete observation.

\noindent\textbf{Step 6.} Finally put
$\mathcal E=T_{\mathcal S/\Delta}(-\log\mathcal D)$.
For every fibre pair just identified,
the two toric infinitesimal actions
trivialize its logarithmic tangent
bundle.  Thus
\[
 \mathcal E|_{\mathcal S_s}\simeq
 \mathcal O_{\mathcal S_s}^{\oplus2},
 \qquad h^0(\mathcal S_s,\mathcal E_s)=2.
\]
Constant-cohomology base change, in the
form used in \cref{push:src:E.3.1}, makes the direct
image a rank-two holomorphic vector
bundle on the disc and identifies its
fibres with these section spaces.
Choose a frame after shrinking the
disc.  The evaluation map from its
pullback to $\mathcal E$ is an
isomorphism on every fibre, hence
is an isomorphism everywhere by
Nakayama's lemma.  This gives the
claimed relative frame without
choosing a potentially singular
family of projective coordinate
changes on the punctured disc.
\end{proof}



\begin{proposition}[extension of the marked cusp torus action]\label{push:src:E.2.19}
\StatementRecord{push:src:E.2.19}{proved}{Historical events: Claude Fable 5.2 (Anthropic), effort high, 16 November 2026: carried as background without fault found. GPT-6 Astra (OpenAI), ChatGPT Pro mode, effort Pro, S6\_09\_PAPERS, 17 November 2026: targeted proof check. Also checked in this writing. Pre-continuation writing events: GPT-6 Astra (OpenAI), ChatGPT Pro mode, effort Pro, S6\_10\_PAPERS, 18 November 2026: targeted proof check. Continuation metadata events: GPT-6 Astra (OpenAI), ChatGPT Pro mode, effort Pro, S6\_10\_PAPERS continuation, 19 November 2026: used and rechecked in this run; status-vocabulary reconciliation; metadata reconciliation only; no new mathematical check. \emph{Pre-continuation recorded proof scope.} The branchwise logarithmic tangent identifications, conductor equalizer, Grauert evaluation, complete commuting flows, implicit period extension, primitive invariant lattice equality, and constant integral triple-point characters were re-derived}{the hypotheses, notation and definitions stated here; \cref{push:src:E.2.17}; \cref{push:src:E.2.18}; \cref{push:src:E.2.16}}

Let
\[
 f:\mathcal N\longrightarrow
 B=\Delta_q\times\Delta_s
\]
be the proper cusp neighbourhood of the smooth
specialization in \ref{push:src:E.2.17}, and put
$\mathcal W=f^{-1}(q=0)$.  After shrinking the
bidisc, the rank-two vertical logarithmic tangent
bundle is trivial and is generated by two commuting
complete holomorphic vector fields.  Their action
factors through a holomorphic action of
$(\C^*)^2$ over $B$, with the integral marking
given by the periods $\mathbf w,\boldsymbol\delta$.
On the normalization of every $\mathcal W_s$
this is its usual toric action.

The marked rank-two lattice is primitive in the
period lattice of each ordinary torus fibre and
equals the full invariant lattice of the cusp
monodromy.  At the two relative triple-point
sections, the three nonzero tangent characters
are constant in $s$ and equal to those of the
prescribed Mumford cusp.  In the toric chart
lattice coordinates, transported from that
punctured model by their fixed integral change
of basis from the period marking, they are
respectively
\[
 \{(1,0),(0,1),(-1,-1)\},\qquad
 \{(-1,0),(0,-1),(1,1)\}.
 \tag{E.2.19.1}
\]
The assertion constructs the action on the actual
neighbourhood.  It does not assert that an
arbitrary equivariant quotient is smooth or
that the entire completion has already been
identified with a prescribed toric quotient.
\end{proposition}


\begin{proof}
\noindent\textbf{Step 1. The logarithmic bundle and its restriction.}
The coordinates in \ref{push:src:E.2.17} make $\mathcal W$ a
relative normal-crossing divisor over $s$.
Define
\[
 \mathcal L=
 \ker\left[
 d\log q:
 T_{\mathcal N/\Delta_s}(-\log\mathcal W)
       \longrightarrow\mathcal O_{\mathcal N}
 \right].
 \tag{E.2.19.2}
\]
The map is surjective: if
$q=z_1\cdots z_k$, it sends each
$z_i\partial_{z_i}$, $i\leq k$, to one.
Consequently $\mathcal L$ is locally free of
rank two.  Its sections are actual holomorphic
vector fields, tangent to the branches, and
annihilate both $q$ and $s$.

Let $\nu:\mathcal S\to\mathcal W$ and
$\mathcal D\subset\mathcal S$ be the
simultaneous normalization and conductor.
Restriction to a branch gives a canonical
isomorphism
\[
 \nu^*(\mathcal L|_{\mathcal W})
 \simeq T_{\mathcal S/\Delta_s}(-\log\mathcal D).
 \tag{E.2.19.3}
\]
This can be checked at every stratum.  For
$q=xy$, a basis is
$x\partial_x-y\partial_y,\partial_z$;
on the branch $x=0$ it becomes
$-y\partial_y,\partial_z$, a logarithmic
basis there.  For $q=xyz$, use
$x\partial_x-z\partial_z$ and
$y\partial_y-z\partial_z$.  Their
restrictions to $x=0$ have coefficient
matrix of determinant one in the basis
$y\partial_y,z\partial_z$.
The one-branch case is immediate.
The same computation keeps $s$ fixed.

By \ref{push:src:E.2.18} the right-hand side of
(E.2.19.3) has a global holomorphic
frame.  On the two preimages of each
conductor curve, the identifications
in (E.2.19.3) describe descent by an
invertible two-by-two matrix.  Its
entries are holomorphic on a proper
$\mathbb P^1$ family over $\Delta_s$,
so are constant along each curve
fibre and are holomorphic functions
of $s$.  Denote these matrices by
$G_j(s)$.

For $s\ne0$ both invariant logarithmic
fields in the prescribed Mumford model
descend.  Indeed, its period maps in
multiplicative coordinates multiply
the two coordinates by units depending
only on the base.  Their derivative
on the two vertical logarithmic
fields is the identity.  The toric
small-resolution charts preserve
these same fields.  Thus the entire
two-dimensional space of sections
on the normalized pair descends for
$s\ne0$.  In any chosen frame this
means $G_j(s)=I$, and the identity
theorem gives $G_j(0)=I$ as well.

There is no additional descent
condition hidden at a triple point.
For a normal-crossing local ring,
functions on the normalization
descend precisely when their
restrictions to every pair of
branches agree.  This is the usual
equalizer for
$\C\{x,y,z\}/(xyz)$, and likewise
for one or two branches.  Tensoring
with the locally free bundle
$\mathcal L|_{\mathcal W}$ proves
the same assertion for its sections.
Consequently the chosen normalized
frame descends, and
\[
 \mathcal L|_{\mathcal W}
       \simeq\mathcal O_{\mathcal W}^{\oplus2}.
 \tag{E.2.19.4}
\]

\noindent\textbf{Step 2. A frame on the whole neighbourhood.}
The map $f$ is proper and flat; this
also follows locally from the equations
$(q,s)=(z_1\cdots z_k,s)$.
The bundle $\mathcal L$ is flat over
$B$.  On each ordinary fibre it is
the tangent bundle of a complex
two-torus.  On each $q=0$ fibre it
is trivial by (E.2.19.4).  All these
fibres are compact and connected,
so in every case
\[
 h^0(\mathcal N_{q,s},
       \mathcal L|_{\mathcal N_{q,s}})=2,
\]
and those sections generate the
bundle pointwise.  The proper
constant-cohomology base-change
theorem, in the form established
in \cref{push:src:E.3.1}, makes $f_*\mathcal L$
a rank-two vector bundle on $B$.
Its evaluation map
\[
 f^*f_*\mathcal L\longrightarrow\mathcal L
 \tag{E.2.19.5}
\]
is an isomorphism on every fibre,
and hence everywhere by Nakayama's
lemma.  Shrink the bidisc and
choose a holomorphic frame $V_1,V_2$.

These fields commute.  On every
ordinary compact torus, every
holomorphic vector field is a
translation field, so their
bracket vanishes there.  The
ordinary locus is dense, and
the bracket is holomorphic.
Each field is complete because
it is vertical and every fibre
of the proper map is compact:
the holomorphic local flow can
be continued for every complex
time on that compact invariant
fibre.  Holomorphic dependence
on initial points and parameters
gives a holomorphic action
\[
 \operatorname{Act}:\C^2\times
       \mathcal N\longrightarrow\mathcal N
\]
over $B$.

\noindent\textbf{Step 3. Holomorphic extension of the two toric periods.}
Choose a smooth point $p_0$ of the
central cusp away from its conductor.
The map $f$ is a submersion there,
so there is a holomorphic section
$p(q,s)$ through $p_0$.
On the central normalized toric
pair the two fields are a basis
of the toric logarithmic Lie
algebra.  Its open orbit is
$(\C^*)^2$.  Let
$v_1^0,v_2^0\in\C^2$ be a
basis of its full period lattice
in the coordinates $V_1,V_2$.
They are complex-linearly
independent.

Near each $v_i^0$, write
$\operatorname{Act}(v,p(q,s))=p(q,s)$
using two vertical coordinates
in a submersion chart around
$p_0$.  The derivative in $v$
is invertible at $(v_i^0,0,0)$,
since the infinitesimal action
is an isomorphism on that open
orbit.  The holomorphic implicit
function theorem gives a unique
holomorphic vector $v_i(q,s)$
near $v_i^0$ solving this equation.
For $q\ne0$, fixing $p(q,s)$
means that this translation is
the identity on the entire
torus fibre.  The equality of
holomorphic maps then extends
over $q=0$.  Thus $v_i(q,s)$
is a period on every fibre.
Their complex-linear
independence persists after
shrinking $B$.

Use these two columns as a
change of the frame.  Because
their coefficients are pulled
back from $B$, the new fields
still commute and are complete.
They have period one, and the
action consequently factors
through
$\C^2/\Z^2\simeq(\C^*)^2$,
using the exponential
$z\mapsto e^{2\pi i z}$
in each coordinate.

\noindent\textbf{Step 4. Primitivity and the actual marking.}
On $B^\circ=\Delta_q^*\times\Delta_s$
the period lattices of the
translation action are the
integral first-homology local
system of the torus family.
The continuously varying
periods $v_1,v_2$ represent
flat integral classes: in a
local differentiable
trivialization their flow
loops vary continuously, and
integral homology is discrete.
They generate a rank-two
invariant sublattice $L$.

Suppose $L$ were not primitive
at some ordinary fibre.  There
would be a rational pair
$(a_1,a_2)\notin\Z^2$ such
that $a_1v_1+a_2v_2$ is an
integral period there.  This
is an equality of flat
integral classes after
multiplication by a common
denominator, so it persists
throughout the connected
ordinary base.  In particular
the corresponding fixed
finite-order element of
$(\C^*)^2$ acts identically
there.  By the identity
theorem it acts identically
on the central open orbit.
This contradicts the choice
of $v_1^0,v_2^0$ as the full
central toric period lattice.
Thus $L$ is primitive.

The marked cusp monodromy,
constant in $s$ by \ref{push:src:E.2.16},
has the form
\[
 \widehat\gamma\longmapsto
       \widehat\gamma-\widehat\delta,
 \qquad
 \widehat u\longmapsto\widehat u+\widehat w,
 \qquad
 \widehat w\longmapsto\widehat w,
 \qquad
 \widehat\delta\longmapsto\widehat\delta.
 \tag{E.2.19.6}
\]
Its full integral invariant
lattice is
$\Z\widehat w\oplus\Z\widehat\delta$.
Since $L$ is a primitive
rank-two sublattice of this
rank-two lattice, equality
follows.  An integral
unimodular change of the two
periods therefore gives
exactly the marking
$\mathbf w,\boldsymbol\delta$.
This also proves that no
extra finite kernel has been
silently introduced.

\noindent\textbf{Step 5. Characters at the triple points.}
The two triple-point sections
are fixed by the torus action:
on the normalization they
are toric vertices.  Restrict
the derivative representation
to the compact real torus
$(S^1)^2$ and average a
Hermitian metric along either
section.  Its characters
are integral pairs and
vary continuously in $s$;
hence their multiset is
locally constant.  Equivalently,
Fourier averaging gives the
individual character
subbundles.  They coincide
at $s=0$ with the known
characters for $s\ne0$.

For clarity, the two unit
triangles in the prescribed
cocharacter lattice have
height-one rays
\[
 \{(0,0,1),(1,0,1),(0,1,1)\},
 \qquad
 \{(1,0,1),(0,1,1),(1,1,1)\}.
\]
Inverting their unimodular
ray matrices gives the two
character triples in
(E.2.19.1); the sum in
each triple is zero, as
required by invariance of
$q=z_1z_2z_3$.
The calculation is recorded
in \path{code/d_push_cusp_torus_action.py}.
All action and marking
assertions follow.
\end{proof}



\begin{proposition}[identification of the actual analytic cusp]\label{push:src:E.2.20}
\StatementRecord{push:src:E.2.20}{proved}{Historical events: Claude Fable 5.2 (Anthropic), effort high, 16 November 2026: carried as background without fault found. GPT-6 Astra (OpenAI), ChatGPT Pro mode, effort Pro, S6\_09\_PAPERS, 17 November 2026: targeted proof check. Also checked in this writing. Pre-continuation writing events: GPT-6 Astra (OpenAI), ChatGPT Pro mode, effort Pro, S6\_10\_PAPERS, 18 November 2026: targeted proof check. Continuation metadata events: GPT-6 Astra (OpenAI), ChatGPT Pro mode, effort Pro, S6\_10\_PAPERS continuation, 19 November 2026: used and rechecked in this run; status-vocabulary reconciliation; metadata reconciliation only; no new mathematical check. \emph{Pre-continuation recorded proof scope.} The relative stratum-preserving transport, exact fundamental-group quotient, lifted torus action, adjacent-origin comparison q\textasciicircum{}n times a holomorphic unit, unimodular character charts, and finite fundamental-domain shrinking were re-derived. The inverse coordinate maps and exact period deck action give the actual analytic biholomorphism}{the hypotheses, notation and definitions stated here; \cref{push:src:E.2.10}; \cref{push:src:E.2.7}; \cref{push:src:E.2.16}; \cref{push:src:E.2.18}; \cref{push:src:E.2.19}; \cite[p.~30]{AC26Root}; \cref{push:src:E.2.12}}

In the smooth specialization of \ref{push:src:E.2.10}, normalize
the three distinct critical-value sections to
$0,1,\infty$ and write $q=1/t$ near the cusp.
Let $c(s)$ be the holomorphic interior period
lift supplied by \ref{push:src:E.2.7}.  After shrinking the
$(q,s)$-bidisc, its inverse image in the actual
family is biholomorphic, over that bidisc, to
the prescribed marked Mumford cusp family at
parameter $c(s)$.  The isomorphism induces
the given identification of the regular
torus Jacobians and their integral markings.
Its affine origin may be fixed by choosing
a holomorphic section through a smooth
point of the reduced cusp.
\end{proposition}


\begin{proof}
Write $f:\mathcal N\to\Delta_q\times\Delta_s$
for the proper actual cusp neighbourhood and
$\mathcal W=f^{-1}(q=0)$.  \ref{push:src:E.2.16}--\ref{push:src:E.2.18}
give its reduced relative normal-crossing
divisor, its normalized toric hexagon, and
the unchanged opposite-side identifications.
\ref{push:src:E.2.19} gives an actual holomorphic action
of $H=(\C^*)^2$ over the bidisc.  Its
period marking is the primitive vanishing
lattice
\[
 \Lambda_{\rm tor}
   =\Z\widehat w\oplus\Z\widehat\delta.
 \tag{E.2.20.1}
\]
We reconstruct the neighbourhood from
this action, keeping track of the
extension of its characters.

\noindent\textbf{Step 1. The topology and the exact covering.}
The relative coordinates
$q=x,xy,xyz$ give smooth real lifts
of the two real parameter directions
of $s$ which annihilate $dq$ and
are tangent to every cusp stratum.
Away from the cusp such lifts exist
because $f$ is a submersion.
A smooth partition of unity preserves
both these linear requirements.
On the inverse image of a smaller
closed $q$-disc the lifts integrate
along a sufficiently short parameter
path: $q$ remains constant, and
properness prevents escape.
Thus their flows identify the cusp
neighbourhoods, all strata, and
the regular-fibre inclusion maps
by $q$-preserving diffeomorphisms.

For the known Mumford cusp we
check the covering calculation
directly.  Each maximal toric
chart above $|q|<\epsilon$ is
\[
 \{(x,y,z)\in\C^3:|xyz|<\epsilon\},
\]
which is star-shaped and simply
connected.  Every intersection
contains the same connected dense
open torus $(\C^*)^2\times\Delta_q^*$.
Each connected component of an
intersection is open and must meet
that dense set, so the intersection
is connected.  The same holds for
the intersection with a finite
union of preceding charts.
Van Kampen proves simple
connectedness of each finite
union; a loop lies in such a
union by compactness.  The full
toric cover is therefore simply
connected.  The known free,
properly discontinuous period
action has deck group
\[
 \overline\Lambda=\Lambda/\Lambda_{\rm tor}
                 \simeq\Z^2.
\]
On a regular fibre its covering
is $\C^2/\Lambda_{\rm tor}\to
\C^2/\Lambda$.  The regular-fibre
map to the cusp body's fundamental
group is exactly the projection
to $\overline\Lambda$, and the
$q$-meridian dies through the
section $(x,y,z)=(q,1,1)$.
These statements also agree
with the original manuscript's
Proposition~4.7 and Corollary~4.8,
whose original PDF was checked
\cite[p.~30]{AC26Root}; the
preceding argument verifies
the step we use.  The
$q$-preserving transport transfers
this exact description to the
actual neighbourhood.

Take the universal covering
\[
 \widetilde{\mathcal N}\longrightarrow\mathcal N.
\]
Its deck group is
$\overline\Lambda$, and its
central normalized components
form the planar unfolding of
the hexagon.  The $H$ action
lifts.  Indeed the fundamental
loops of $H$ at a central open
orbit lie in one normalized
del Pezzo component, which is
simply connected, and hence die
in $\pi_1(\mathcal N)$.
The identity lift is unique
and satisfies the action law.
Choose a holomorphic section
$p(q,s)$ through a smooth
central cusp point and lift
it to the cover.  On $q\ne0$
its orbit map is an isomorphism
\[
 H\times(\Delta_q^*\times\Delta_s)
       \xrightarrow{\ \sim\ }
 \widetilde{\mathcal N}|_{q\ne0}.
 \tag{E.2.20.2}
\]
The stabilizer has been removed
exactly by (E.2.20.1), using
the just-verified regular-fibre
inclusion map.  In particular
the two characters of $H$
are globally defined on this
open covering space.

Connectedness of $H$ makes
its lifted action commute
with every deck transformation:
conjugation acts continuously
on the discrete deck group
and is the identity at the
identity of $H$.  On
(E.2.20.2), the two deck
generators therefore act by
\[
 (h,q,s)\longmapsto(P_i(q,s)h,q,s),
 \qquad i=1,2.
 \tag{E.2.20.3}
\]
The functions $P_i$ are the
actual marked period multipliers.
Changing the affine origin
does not change them.
The marked period comparison
of \ref{push:src:E.2.7} and \ref{push:src:E.2.12} identifies
them with the prescribed
$q$-monomials times holomorphic
units at $c(s)$.

\noindent\textbf{Step 2. Character extension across adjacent components.}
For a central component $D_v$,
choose a holomorphic base
section $a_v(q,s)$ through
its dense $H$ orbit.  Its
orbit map is an open
holomorphic embedding
\[
 H\times(\Delta_q\times\Delta_s)
       \longrightarrow\widetilde{\mathcal N}.
\]
It is injective on the
central open orbit and on
each regular fibre, and
its differential is invertible
along those loci.  Its
image contains the whole
regular covering.  Two
such origins differ for
$q\ne0$ by a holomorphic
translation $h_{vw}(q,s)$
in $H$.  We must prove
that this function has no
essential singularity at
$q=0$.

If both origins lie in
the same central component,
their ratio extends as a
holomorphic $H$-valued
unit.  The central open
orbit is free and transitive,
so this follows from the
holomorphic inverse-function
theorem for the orbit map,
with both parameters retained.

For two adjacent components,
use a crossing chart with
$q=ab$ and an orbit coordinate
$z\ne0$.  The stabilizer
of a general conductor
point is a primitive
connected one-dimensional
subtorus of $H$.  Its
two normal characters are
$+1,-1$: on each normalized
toric component these are
the primitive boundary
characters, and invariance
of $q$ makes them opposite.
Compact-circle averaging
gives these exact local
linear coordinates.  A
complementary torus direction
provides the orbit coordinate;
an invariant unit in $q$
can be absorbed while
preserving the weights.

Fix small nonzero constants
$a_0,b_0,z_0$ in this chart
and take the two special
origins
\[
 (a,b,z)=(q/b_0,b_0,z_0),\qquad
 (a,b,z)=(a_0,q/a_0,z_0).
 \tag{E.2.20.4}
\]
The stabilizer cocharacter
evaluated at $a_0b_0/q$
sends the first to the
second.  This equality
is justified by the local
action itself: a path
in that subgroup can
increase the modulus of
$a$ and decrease that
of $b$, with $ab=q$
fixed, staying inside
the crossing chart;
the phase can be varied
without changing either
modulus.  Thus the
possibly large group
parameter causes no
unjustified use of a
local coordinate outside
its domain.

Comparing (E.2.20.4) with
arbitrary chosen origins
adds only holomorphic
$H$-valued units, by the
same-component argument.
The adjacent comparison
therefore has the form
\[
 h_{vw}(q,s)=q^{n_{vw}}u_{vw}(q,s),
 \qquad n_{vw}\in\Z^2,
 \tag{E.2.20.5}
\]
where $u_{vw}$ is holomorphic
with values in $H$ on the
whole bidisc.  Every central
component is reached by
a finite adjacent path.
Multiplying these expressions
proves the same assertion
for its comparison with
the initial origin.
Different paths agree
on the regular cover,
and hence define identical
meromorphic characters.
The exponents are locally
constant in $s$, since
(E.2.20.5) holds with a
nonvanishing unit on the
full bidisc.

\noindent\textbf{Step 3. The fan and the analytic coordinate maps.}
At either triple-point
section the whole $H$
fixes the point.
Compact-torus averaging
gives three eigen-coordinates
with the constant integer
characters of \ref{push:src:E.2.19}.
Their only primitive
positive invariant relation
is their product.
The reduced normal-crossing
equation therefore writes
$q$ as that product times
a holomorphic invariant
unit.  The unit is a
function of the product
and $s$; absorb it into
one eigen-coordinate to
obtain exactly $q=xyz$.

The orders of the two
global characters along
$D_v$, together with
$\operatorname{ord}_{D_v}q=1$,
give a ray $(v,1)$.
At a triple point the
three incident rays
form a unimodular cone.
Indeed all adjacent
orders and tangent
characters are constant
in $s$ and coincide
with the known ones
for $s\ne0$.
The entire resulting
fan is consequently
the prescribed $A_2$
triangulation, with
its fixed labels.
No additional refinement
is being introduced.

In a triple chart a
global character is
a Laurent monomial in
$x,y,z$ times a function
of $q,s$.  Evaluate it
on a section keeping
two nonzero coordinates
fixed and solving the
third from $q$.
The origin comparisons
above show that the
remaining function is
a power of $q$ times
a holomorphic unit.
Thus the dual monomial
coordinates of the
unimodular fan cone
pull back to $x,y,z$
times holomorphic units
in $q,s$.  They extend
holomorphically and
give a local
biholomorphism through
the triple point.
The crossing and
open-orbit charts give
the same conclusion
on the other strata.
The maps agree on
the dense regular
cover and therefore
glue.

We finally check the
common neighbourhood
and the deck action.
We do not take an
infimum of infinitely
many chart radii.
Choose finitely many
charts covering one
compact normalized
central del Pezzo
component and finitely
many overlaps for
the two deck generators.
Shrink the bidisc once
for this finite data.
On the regular cover
the character map
intertwines the actual
deck action with the
reference multipliers
$P_i(q,s)$ in
(E.2.20.3).  Both maps
extend in the charts
just constructed, so
their equality extends
to the central strata.
The reference multipliers
preserve the prescribed
invariant fan.  Deck
translates of the finite
chart collection cover
the entire central
unfolding and use the
same base bidisc.

We have obtained local
biholomorphisms at all
central strata, bijections
on their labelled torus
orbits, and the original
isomorphism on the whole
regular cover.  The
inverse coordinate maps
are the same unimodular
monomial charts with
the inverse units.
Quotienting by the
identical free properly
discontinuous deck
actions gives the
claimed biholomorphism
of actual analytic
cusp neighbourhoods.
\end{proof}



\paragraph{Scope and source comparison}

The toric covering and its
period action are the known
Mumford construction.  The
point here is to recover
that construction from
the actual smooth limit,
using the action proved
in \ref{push:src:E.2.19}.  In particular
the absence of an essential
singularity in the character
comparison is proved in
(E.2.20.4)--(E.2.20.5);
it is not an assumed
meromorphic comparison.
The local identification
does not alone identify
the global threefold or
settle arbitrary deformation
chains.  \ref{push:src:E.2.21} supplies the
finite affine and global collar
comparison; \ref{push:src:E.2.22}--\ref{push:src:E.2.24}
provide the additional propagation
and closure arguments for general
bases and continuous structures.

\Uses{the hypotheses, notation and definitions stated here; \cref{push:src:E.2.19}; \cref{push:src:E.2.21}; \cref{push:src:E.2.22}; \cref{push:src:E.2.24}}


\begin{proposition}[global identification of a smooth completion]\label{push:src:E.2.21}
\StatementRecord{push:src:E.2.21}{proved}{Historical events: Claude Fable 5.2 (Anthropic), effort high, 16 November 2026: carried as background without fault found; targeted proof check; used and rechecked in this run. GPT-6 Astra (OpenAI), ChatGPT Pro mode, effort Pro, S6\_09\_PAPERS, 17 November 2026: targeted proof check. Also checked in this writing. Pre-continuation writing events: GPT-6 Astra (OpenAI), ChatGPT Pro mode, effort Pro, S6\_10\_PAPERS, 18 November 2026: targeted proof check. Continuation metadata events: GPT-6 Astra (OpenAI), ChatGPT Pro mode, effort Pro, S6\_10\_PAPERS continuation, 19 November 2026: used and rechecked in this run; status-vocabulary reconciliation; metadata reconciliation only; no new mathematical check. \emph{Pre-continuation recorded proof scope.} The finite affine norm classes and integral collar classes are locally constant under their actual lattice maps. Finite averaging and the explicit H1(V)=0 Laurent splitting remove the remaining holomorphic cocycle, including its parameter dependence}{the hypotheses, notation and definitions stated here; \cref{push:src:E.2.5}; \cref{push:src:E.2.7}; \cref{push:src:E.2.10}; \cref{push:src:E.2.12}; \cref{push:src:E.2.14}; \cref{push:src:E.2.15}; \cref{push:src:E.2.13}; \cref{push:src:E.2.20}}

Let $\pi:\mathcal X\to\Delta$ satisfy the hypotheses of
\ref{push:src:E.2.5}, with fixed coprime multiplicities $(3k,4l)$ on
every noncentral fibre.  Then $X_0$ is a member of that
same constructed family.  Its isomorphism parameter is
the interior value $z(0)$ of \ref{push:src:E.2.7}.

After choosing a holomorphic lift $c(s)$ of that parameter
and shrinking the disc, the actual family is isomorphic
to the prescribed family at $c(s)$.  The ambiguity in
this comparison is the connected translation automorphism
group already determined in E.1.
\end{proposition}


\begin{proof}
\ref{push:src:E.2.10} extends the intrinsic pencil to a proper flat map
\[
 g:\mathcal X\longrightarrow\mathbb P^1\times\Delta.
\]
\ref{push:src:E.2.7} supplies its interior nonzero period $z(s)$ and hence
a holomorphic logarithmic lift $c(s)$.  \ref{push:src:E.2.12} excludes
additional critical values, and \ref{push:src:E.2.14}--\ref{push:src:E.2.15} exclude
target collapse.  Normalize the three distinct labelled
sections to $0,1,\infty$ by a holomorphic projective
transformation over the disc.

On the ordinary locus the actual marked Jacobian has
period matrix $\Pi_{c(s)}(t)$.  More intrinsically, its
holomorphic period kernel in the complexified marked
lattice equals that of this matrix.  This equality holds
for $s\ne0$ by the given family identification and extends
over $s=0$ by the identity theorem.  The local marking
choices agree: an automorphism of a fixed marked member
has identity action on ordinary $H_1$, by \cref{push:src:E.1.2}.  The
actual smooth family at any fixed ordinary target point
also has identity homology return around $s=0$.

\ref{push:src:E.2.13} gives smooth actual normalized root covers at the
two finite values, of degrees $3k$ and $4l$, with their
free cyclic deck actions.  Their marked Jacobians are
the corresponding reference root Jacobians.  \ref{push:src:E.2.20}
identifies the actual cusp neighbourhood over the full
$(q,s)$-bidisc with the prescribed marked Mumford
neighbourhood, including its period multipliers.
It remains to identify the finite affine actions and
the global translations on the overlaps.

\noindent\textbf{Step 1. The finite affine classes.}
Choose local holomorphic origins on a finite root torus
family.  Every deck map is affine on the compact torus
fibres.  Its linear part is the fixed marked finite
monodromy.  Subtracting this part gives a holomorphic
translation section.  Lift it to the Lie bundle on
simply connected refinements of the root charts.
The finite-order relation says that the sum of its
successive linear translates is an integral period.
This lattice vector is locally constant in $s$:
the lifted translation and period vectors vary
continuously, whereas the integral coefficients are
discrete.  Changing the lift changes it by precisely
the integral coboundary allowed in E.1.

Consequently the finite affine cocycle has the same
integral class at $s=0$ as at $s\ne0$.  After matching
these classes, the remaining difference is a cocycle
in the complex vector sheaf.  The explicit finite
averaging in E.1 removes it, since division by the
finite group order is allowed in that sheaf.
This constructs the finite-root comparisons with
holomorphic dependence on $s$.  Equivalently one
may use a finite chart cover and its deck translates;
the same averaging solves the equivariant comparison.
The reduced finite fibres are the resulting free
quotients; they have not been assumed to be tori.

\noindent\textbf{Step 2. The integral collar classes.}
Use the two finite-root neighbourhoods, the cusp
neighbourhood, ordinary simply connected charts,
and fixed annular collars to cover the target.
Choose local holomorphic origins on the ordinary
torus charts.  A transition between two such
descriptions is fibrewise affine.  The common
integral marking makes its linear part the
prescribed one, leaving a translation section.

Lift these sections to the Lie bundles on simply
connected refinements of the collars.  Differences
on deck translates and sums on triple overlaps
are integral periods.  Every resulting integral
vector is constant in $s$.  This comparison takes
place on fixed collars away from the singular
centres; it does not evaluate a logarithm at a
singular fibre.  The finite and cusp chart
identifications specify exactly which integral
changes of origins are allowed.

For $s\ne0$ these classes are those of the
prescribed logarithmic-transform data, modulo
the integral changes of origins in E.1.
They are therefore the same at $s=0$.
One can normalize them once on the connected
parameter disc.  This includes the cusp
winding class, since \ref{push:src:E.2.20} identifies the
actual marked period action there.  Thus
there is no additional central integral
translation or fan datum.

\noindent\textbf{Step 3. Solving the holomorphic remainder.}
After this discrete normalization, the remaining
comparison is the additive cocycle in the
extension sheaf $\mathcal V$ of E.1.  Its
finite extension conditions are exactly
those of the already identified root
charts.  Its cusp extension condition
is the actual toric collar condition
proved in \ref{push:src:E.2.20}.  The calculation in
E.1 therefore applies to this actual
cover and gives
\[
 0\longrightarrow\mathcal O_{\mathbb P^1}
       \boldsymbol\delta
 \longrightarrow\mathcal V
 \longrightarrow\mathcal O_{\mathbb P^1}(-\infty)
 \longrightarrow0,
 \qquad H^1(\mathbb P^1,\mathcal V)=0.
\]
The explicit finite averaging and Laurent-series
splitting in E.1 kill the additive cocycle.
The resulting translations glue the local
comparisons to a biholomorphism of $X_0$
with the prescribed member at $c(0)$.

The same solution works over a smaller parameter
disc.  On the fixed collars, Cauchy integral
formulas for the Laurent coefficients preserve
holomorphic dependence on $s$, as does finite
averaging.  Thus the solved translations and
their glued comparison are holomorphic in
$s$.  Their ambiguity is the connected
scalar translation group of E.1.  This
proves both the central and family assertions.
\end{proof}



\paragraph{Scope of \ref{push:src:E.2.21}}

The known-pair locus is open by \cref{push:src:E.2.2} and
is now stable under full punctured-disc
smooth specializations.  These two facts
alone do not prove that an arbitrary open
known locus in a reduced analytic base
is topologically closed.  A holomorphic
disc through its boundary need not have
its entire punctured image in that locus.
The needed generic propagation argument
is proved in \ref{push:src:E.2.22}.  \ref{push:src:E.2.23}--\ref{push:src:E.2.24} then
establish the arbitrary-base and continuous-family
conclusions asserted in \ref{push:src:E.2.4}.

\Uses{the hypotheses, notation and definitions stated here; \cref{push:src:E.2.21}; \cref{push:src:E.2.22}; \cref{push:src:E.2.23}; \cref{push:src:E.2.24}; \cref{push:src:E.2.4}}


\begin{proposition}[propagation from a known open subset of a disc]\label{push:src:E.2.22}
\StatementRecord{push:src:E.2.22}{proved}{Historical events: Claude Fable 5.2 (Anthropic), effort high, 16 November 2026: carried as background without fault found; targeted proof check; used and rechecked in this run. GPT-6 Astra (OpenAI), ChatGPT Pro mode, effort Pro, S6\_09\_PAPERS, 17 November 2026: targeted proof check. Also checked in this writing. Pre-continuation writing events: GPT-6 Astra (OpenAI), ChatGPT Pro mode, effort Pro, S6\_10\_PAPERS, 18 November 2026: targeted proof check; used and rechecked in this run. Continuation metadata events: GPT-6 Astra (OpenAI), ChatGPT Pro mode, effort Pro, S6\_10\_PAPERS continuation, 19 November 2026: used and rechecked in this run; status-vocabulary reconciliation; metadata reconciliation only; no new mathematical check. \emph{Pre-continuation recorded proof scope.} The intrinsic line and pencil propagate outside finitely many successive analytic discrete layers. Finite normalization, immersion and conormal rank tests give actual relative NC charts. Equality of the full period kernels, strict determinant threshold and collar comparison classify the remaining fibres; reverse-order filling handles accumulating later layers}{the hypotheses, notation and definitions stated here; \cref{push:src:E.2.21}; \cref{push:src:E.2.5}; \cite[Chapter~II, Theorem~8.8]{Demailly2012E}; \cite[Chapter~VI, Theorem~1.1 and Table~10]{BHPV2004E}; \cite[Chapter~II, Definition~7.9 and Theorem~7.12]{Demailly2012E}; \cref{push:src:E.2.18}; \cref{push:src:E.2.19}; \cref{push:src:E.2.20}}

Let $\pi:\mathcal X\to\Delta$ be a proper holomorphic
submersion with fibres diffeomorphic to $S^6$.  Suppose
there is a nonempty open subset $U\subset\Delta$ such
that every $X_s$, $s\in U$, is a member of the fixed
coprime multiplicity family $(3k,4l)$ of \cref{push:src:E.2.1}.
Then every fibre over $\Delta$ is a member of that
same family.
\end{proposition}


\begin{proof}
Local miniversality in \cref{push:src:E.2.1} supplies actual
holomorphic comparisons with the prescribed family
locally over $U$, with locally constant integral
markings.  Thus the known open set has the relative
root and normal-crossing charts used below; only
pointwise isomorphisms are not being substituted
for these relative comparisons.

The proof first propagates the actual geometric data
outside finitely many successive discrete analytic
exceptional sets.  It then fills these sets in reverse
order by \ref{push:src:E.2.21}.  Later exceptional sets are only
claimed to be discrete on the domain then under
consideration; they may accumulate at points removed
at an earlier stage.

\noindent\textbf{Step 1. The intrinsic line and the relative pencil.}
Put $A=12kl$ and $D=kl-3k-4l\ne0$.
The argument of \ref{push:src:E.2.5} constructs the unique line
bundle $\mathcal L$ on the whole total space with
\[
 \mathcal L^D=K_{\mathcal X/\Delta}^{A}.
\]
Indeed Ehresmann transport over the contractible disc
gives $H^1(\mathcal X,\Z)=H^2(\mathcal X,\Z)=0$,
so the exponential sequence identifies its Picard
group with a complex vector space.  This construction
needs no hypothesis on the fibres outside $U$.

Proper coherence makes $\mathcal V=\pi_*\mathcal L$
a torsion-free, hence locally free, sheaf on the
disc.  Its rank is two because it has that rank
on $U$.  At every parameter its fibre injects
into $H^0(X_s,\mathcal L_s)$ by the multiplication
exact sequence used in \ref{push:src:E.2.5}.  Choose a frame
of $\mathcal V$ on the disc.  The common zero
locus of its two evaluated sections is closed
analytic in $\mathcal X$.  Its image is closed
analytic by properness and Remmert's theorem
\cite[Chapter~II, Theorem~8.8]{Demailly2012E}.
It misses $U$, so is a discrete analytic subset.
Remove it temporarily.

On the remaining connected parameter curve $S$
the sections define a proper holomorphic map
\[
 g:\mathcal X_S\longrightarrow\mathbb P^1\times S,
 \qquad \mathcal L=g^*\mathcal O(1).
 \tag{E.2.22.1}
\]
Each $g_s$ is nonconstant because its two sections
are independent.  Its fibres have pure dimension
two; the smooth total space and smooth target
then give flatness.  The fibres are connected:
the finite Stein factor is generically of degree
one, as it is so above $U$, and a finite
birational map to the normal target is an
isomorphism.  Its source is irreducible, so
there is no additional vertical component
in the Stein factor.  We retain
\[
 K_{\mathcal X_S/S}^{A}=g^*\mathcal O(D).
 \tag{E.2.22.2}
\]

\noindent\textbf{Step 2. Critical values and ordinary tori.}
The relative critical locus of $g$ is analytic
and proper over the target.  Its image has
finite intersection with each target fibre:
the nonconstant proper map $g_s$ to a curve
has finitely many critical values.  Thus the
reduced critical image is finite over $S$.
Its one-dimensional part has degree three,
as can be read over $U$; every such component
dominates $S$.  Remove the images of its
zero-dimensional components and its branch
and singular points.  These are further
proper analytic discrete subsets, absent
over the known open set.

A finite etale parameter cover now labels
the three critical sections.  Pointwise
classification after this cover descends;
images of later discrete exceptional sets
remain discrete because this cover is finite.
One may work on each connected component,
each of which meets the inverse image of $U$.
Send the labelled sections to $0,1,\infty$.
There are exactly these three critical values
on the current parameter domain.

The ordinary family over
$B^\circ\times S$, where
$B^\circ=\mathbb P^1\setminus\{0,1,\infty\}$,
is smooth and proper.  Its compact surface
fibres have the topology of the known two-tori
by Ehresmann's theorem and connectedness.
Adjunction gives integral $c_1(K_F)=0$,
since $F$ is a smooth fibre in a complex
six-sphere and its normal line is trivial.
These surfaces are complex tori.  To spell
out the classification step, $b_1(F)=4$
is even, so $F$ is K\"ahler.  The zero
canonical class excludes exceptional curves
and ruled surfaces.  On a K\"ahler surface,
a nonzero effective divisor cannot have
zero real Chern class; hence a nonzero
pluricanonical section has no zeros and
cannot give positive Kodaira dimension.
The compact-surface classification leaves
the Kodaira-dimension-zero list, in which
$b_1=4$ identifies a two-torus
\cite[Chapter~VI, Theorem~1.1 and Table~10]{BHPV2004E}.

\noindent\textbf{Step 3. The finite root charts.}
For any labelled critical section $\Gamma_j$,
every component of the Cartier divisor $g^*\Gamma_j$
dominates $S$.  A vertical divisor component would
be an entire fibre $X_s$, making $g_s$ constant.
Every dominating component meets $U$, where the
reduced fibre is irreducible.  Thus there is
exactly one such component, and its multiplicity
is the known one.  In particular
\[
 g^*\Gamma_1=(3k)\mathcal D_1,
 \qquad g^*\Gamma_2=(4l)\mathcal D_2,
 \qquad g^*\Gamma_0=\mathcal W,
 \tag{E.2.22.3}
\]
with reduced total hypersurfaces on the right.
For clarity, two distinct dominating components
could not both restrict to the known irreducible
reduced fibre.  Each is a Cartier prime in
the smooth total space, is torsion-free and
hence flat over the smooth parameter curve,
and has pure two-dimensional fibres.  Thus
its fibre over any parameter in $U$ contains
the whole known irreducible reduced surface.
At a general smooth point of
that fibre the known local equation has only
one reduced total branch.  Properness forces
every dominating component to meet that open
part of the known fibre, a contradiction.

For a finite multiplicity $m$, write locally
$q=a^m u$, with $a$ a reduced local equation
of $\mathcal D_j$ and $u$ a unit.  The normalization
of the cyclic root base change $r^m=q$ is
locally the disjoint union of the graphs
\[
 r=\zeta a u^{1/m},\qquad \zeta^m=1.
\]
It is therefore smooth and etale over the
actual total space, with free cyclic deck
action.  On $r\ne0$ its map to the root
disc times $S$ is already smooth.  Its
critical locus is supported over $r=0$,
is proper over $S$, and is empty over $U$.
Delete its discrete analytic image for
each of the two roots.  The resulting
proper smooth root families have connected
fibres by their Stein factors and are
families of two-tori by the same topology,
adjunction and surface argument above.

\noindent\textbf{Step 4. A finite analytic test for relative normal crossings.}
The total cusp $\mathcal W$ in (E.2.22.3)
is a reduced embedded hypersurface in
the smooth fourfold.  Take its finite
analytic normalization
\[
 \nu:\widetilde{\mathcal W}\longrightarrow
       \mathcal W\longrightarrow\mathcal X_S.
\]
Existence and finiteness are the analytic
normalization theorem
\cite[Chapter~II, Definition~7.9 and Theorem~7.12]{Demailly2012E}.
The normalization is proper over $S$.
First remove the proper images of its
singular locus, and then of the critical
locus of its map to $S$.  Each is analytic
and misses $U$.  On the resulting domain
the normalization is smooth over $S$.
Next remove the proper image of the
rank-drop locus of
\[
 T_{\widetilde{\mathcal W}/S}
       \longrightarrow\nu^*T_{\mathcal X_S/S}.
\]
This locus too is analytic and absent over
$U$.  The finite map $\nu$ is now an immersion.

For $r=2,3,4$, let $R_r$ be the locus
of ordered, pairwise distinct preimages
in the $r$-fold fibre product of
$\widetilde{\mathcal W}$ over $\mathcal X_S$.
Its partial diagonals are both open and
closed: a finite immersion is locally
an embedding, so its diagonal is open,
and separatedness makes it closed.
Consequently $R_r$ is closed analytic
in the proper fibre product and is
proper over $S$.

At a point of $R_r$ the $r$ local sheets
have relative conormal lines in the same
rank-three bundle $T^*_{\mathcal X_S/S}$.
On $R_2$ and $R_3$, delete the proper
images of the rank-drop loci of these
two or three lines.  Delete also the
image of $R_4$.  These are closed
analytic images and miss $U$, hence
are discrete.  There are only finitely
many such deletions.

At a remaining point the total cusp
has at most three local sheets, each
embedded smoothly and submersive over
$S$, with independent vertical normal
covectors.  Since the \emph{total}
$\mathcal W$ is a reduced embedded
hypersurface, its local ideal is exactly
the product of the distinct local
equations of those sheets.  No finite
subring or seminormality assumption
is being used.  Independence of the
vertical normals makes these equations,
the parameter on $S$, and the remaining
coordinates a holomorphic coordinate
system.  Thus
\[
 q=x,\qquad q=xy,\qquad\hbox{or}\qquad q=xyz
 \tag{E.2.22.4}
\]
after absorbing a unit.  These are
actual relative normal-crossing charts;
in particular the remaining cusp fibres
are reduced.  This proof did not assume
that normalization commutes with base
change before obtaining the charts.

\noindent\textbf{Step 5. The normalized toric pair and its action.}
The smooth proper normalization and
its conductor strata have the same
topology, incidence and normal degrees
as over $U$.  A further finite etale
cover can label these finitely many
strata.  The argument of \ref{push:src:E.2.18} applies
to every fibre: the six conductor
curves are a hexagon of $(-1)$-curves,
three alternate contractions give a
triangle in $\mathbb P^2$, and the
normalized pair is $(\mathrm{dP}_6,D)$.
Its relative logarithmic tangent bundle
is the evaluation pullback of a rank-two
bundle on the parameter curve.

In this evaluation description the
conductor descent maps are holomorphic
endomorphisms of that rank-two bundle.
They are constant along the compact
conductor curves.  They equal the identity
over $U$, where both toric logarithmic
fields descend.  The holomorphic identity
theorem on the connected parameter curve
makes them the identity everywhere.
This is the descent proof of \ref{push:src:E.2.19};
it requires a nonempty known open set,
not specifically a punctured disc.
Grauert evaluation, commutation on the
ordinary tori, completeness and the
implicit extension of the two primitive
toric periods then give the actual
marked cusp torus action locally on
the current parameter curve.  The
integral monodromy and triple-point
characters are constant by transport
and are the prescribed ones.

\noindent\textbf{Step 6. Continuation of the full period kernel.}
Fix an ordinary target point $t_*$.
The integral local system over
$B^\circ\times S$ has constant marked
peripheral representation, up to a
choice of marking: its matrices are
transported from those over $U$.
Parameter monodromy commutes with the
base representation.  It therefore
preserves the cusp invariant lattice
\[
 \ker(A_0-I)=\Z\widehat w\oplus\Z\widehat\delta.
\]
The wedge of the two corresponding
holomorphic period vectors at $t_*$
has a well-defined zero locus on $S$:
changing its integral basis multiplies
the wedge by a nonzero determinant.
It is nonzero over $U$, so that zero
locus is discrete analytic.  Remove
it as one final exceptional layer.

On a connected universal cover of
the remaining parameter curve, normalize
the Lie frame at $t_*$ by these two
columns.  The $\widehat\gamma$ column
then defines a holomorphic function
$c(s)$ by subtracting the fixed
reference value of $\beta^{\rm part}(t_*)$.
All other column identities at $t_*$
hold over a lift of $U$ and hence
everywhere by the identity theorem.

Now work on the universal cover of
the whole ordinary base times this
parameter cover.  The actual holomorphic
period kernel is a rank-two subbundle
of the complexified marked rank-four
lattice.  Compare it with
\[
 \ker\Pi_{c(s)}(t).
 \tag{E.2.22.5}
\]
The latter matrix always has complex
rank two, since its $\widehat w$ and
$\widehat\delta$ columns are the
standard basis.  These two holomorphic
Grassmannian maps agree over the
product with a lift of $U$.  They
therefore agree on the entire
connected ordinary cover.
This comparison does not assume
that the two actual columns were
independent at every $t$; equality
of the kernels proves it afterwards.
In a local actual Lie frame the
identity says that the full period
matrix is $W(t,s)\Pi_{c(s)}(t)$,
where $W$ is its two marked columns.
If $W$ had rank at most one, every
period column would have that rank,
contradicting the real rank four
of the compact torus lattice.

The actual ordinary and finite-root
fibres are compact tori.  The exact
period determinant is nonzero there.
Its sign is the negative sign on
$U$ and stays negative by connectedness
and nonvanishing.  Equality of the
period kernels also extends to the
smooth finite root charts; their
nonzero determinant preserves this
sign at the root point where the
maximum in \cref{fam:src:C.2.3} is attained.  Thus
$\operatorname{Im}c(s)+M<0$ by
the precise threshold of \cref{fam:src:C.2.3}.
This places $c(s)$ in the
admissible interior domain, including
strictness at its boundary.  The
cusp multipliers are the corresponding
marked reference multipliers.  The
proof of \ref{push:src:E.2.20} now reconstructs
the actual cusp at each parameter:
relative normal-crossing transport
gives its known cover and incidence,
the action gives the same characters,
and adjacent-origin comparison gives
the monomial times unit extensions.
The proof works along any path from
$U$ and locally in the current
parameter; it does not need an
already known punctured neighbourhood
at the point being classified.

The finite affine lattice classes
and integral collar classes are
locally constant on this connected
domain, hence remain those over
$U$.  The additive comparison in
E.1 and \ref{push:src:E.2.21} identifies every
remaining actual threefold with
the reference member at $c(s)$.
The pointwise conclusion descends
from every finite parameter cover
used to label the critical or
conductor strata.

\noindent\textbf{Step 7. Filling the exceptional layers.}
We have made only finitely many
successive deletions.  At each stage
the deleted subset was closed
analytic and discrete on the
parameter domain then present.
Finite covers do not change that
property after taking their proper
images.  On the final domain all
fibres are known members of the
fixed pair.  At a point of the
last deleted set, choose a small
disc containing no other point of
that set and avoiding all earlier
deleted sets.  Its punctured disc
consists of known members.  \ref{push:src:E.2.21}
identifies the central fibre as
well.  Fill the entire last layer
in this way, then the preceding
one, and continue in reverse order.
This is valid even when a later
layer accumulates at an earlier
deleted point.  Finally every
fibre over the original disc is
known.  No analytic curve selection
for an arbitrary open locus has
been assumed.
\end{proof}



\begin{theorem}[separation under arbitrary deformation chains]\label{push:src:E.2.23}
\StatementRecord{push:src:E.2.23}{proved}{Historical events: Claude Fable 5.2 (Anthropic), effort high, 16 November 2026: carried as background without fault found; used and rechecked in this run. Claude Fable 5.2 (Anthropic), effort high, 17 November 2026: carried as background without fault found. GPT-6 Astra (OpenAI), ChatGPT Pro mode, effort Pro, S6\_09\_PAPERS, 17 November 2026: targeted proof check. Also checked in this writing. Pre-continuation writing events: GPT-6 Astra (OpenAI), ChatGPT Pro mode, effort Pro, S6\_10\_PAPERS, 18 November 2026: targeted proof check; used and rechecked in this run. Continuation metadata events: GPT-6 Astra (OpenAI), ChatGPT Pro mode, effort Pro, S6\_10\_PAPERS continuation, 19 November 2026: used and rechecked in this run; status-vocabulary reconciliation; metadata reconciliation only; no new mathematical check. \emph{Pre-continuation recorded proof scope.} A finite local projection with unique zero fibre gives a proper analytic curve through a known point and a specified boundary point. Normalizing that curve and applying generic-disc propagation proves closure on every reduced analytic branch. Reduction preserves fibre classes, and connected period data give the converse}{the hypotheses, notation and definitions stated here; \cref{push:src:E.2.4}; \cref{push:src:E.2.22}; \citet[Chapter~II, Theorem~4.19, pp.~95--96]{Demailly2012E}; \cite[Chapter~II, Theorems~7.6 and~7.12]{Demailly2012E}}

In a proper smooth holomorphic family of compact complex
manifolds over a connected reduced complex analytic base,
if one fibre is a sphere member of \cref{push:src:E.2.1} with coprime
multiplicities $(3k,4l)$, then every fibre is a constructed
member with that same pair.

Consequently two of the constructed sphere members are
deformation equivalent, allowing chains of proper smooth
families over connected complex analytic bases, if and
only if their ordered multiplicity pairs agree.  In
particular the distinct coprime pairs give distinct
deformation classes, and \ref{push:src:E.2.4} is resolved by this theorem.
\end{theorem}


\begin{proof}
\noindent\textbf{Step 1.} Fix the pair and let $U$ be the locus of its constructed
members in the base $B$.  It is open by the actual local
miniversality of \cref{push:src:E.2.1} and \cref{push:src:E.2.2}.  We prove that it is
closed without assuming a punctured analytic disc can
be selected inside an arbitrary open subset.

\noindent\textbf{Step 2.} First consider any proper smooth family over a connected
Riemann surface.  If its known-pair locus is nonempty,
Ehresmann's theorem makes all its fibres diffeomorphic
to $S^6$, and that locus is open.  At any point of its closure choose a
coordinate disc meeting that locus.  The intersection
contains a nonempty open subset of the disc, so \ref{push:src:E.2.22}
makes the entire disc known.  Thus the known-pair locus
is closed and equals the connected Riemann surface.
Only this disc propagation, not a special choice of
punctured disc, is used below.

\noindent\textbf{Step 3.} Let $b\in\overline U\subset B$.  If $b\in U$ there
is nothing to prove.  Otherwise choose a small local
representative and an irreducible branch $T$ through
$b$ which meets $U$ arbitrarily near $b$.  There are
only finitely many such local branches.  Its dimension
is positive.  The local parametrization theorem gives
a finite proper projection
\[
 \rho:T\longrightarrow\Delta^d,
 \qquad \rho^{-1}(0)=\{b\},
 \tag{E.2.23.1}
\]
after shrinking the representative and translating
$b$ to the origin.  This is the finite projection
of \citet[Chapter~II, Theorem~4.19, pp.~95--96]{Demailly2012E}.
The isolated central fibre is part of choosing
the generic projection for the germ; the
Weierstrass representative is then taken small
enough to contain no other point over zero.

\noindent\textbf{Step 4.} Choose $a\in U\cap T$ in this representative,
with $\rho(a)\ne0$.  Let $L$ be the complex
line through $0$ and $\rho(a)$, restricted
to the polydisc.  An irreducible component
$C$ of the reduced analytic set
$\rho^{-1}(L)$ which contains $a$ has
dimension one.  Indeed $d-1$ linear
equations give local dimension at least
one, while finiteness of $\rho$ gives
dimension at most one.  Its finite
proper image in $L$ is one-dimensional,
so is all of $L$.  Therefore $C$ contains
the unique point $b$ over zero.

\noindent\textbf{Step 5.} Normalize $C$.  Its normalization
$\widetilde C$ is a connected Riemann
surface, finite over $C$; analytic
normalization exists and a normal
curve is smooth by
\cite[Chapter~II, Theorems~7.6 and~7.12]{Demailly2012E}.
It has points above both $a$ and $b$.
Pull back the given proper smooth
family to $\widetilde C$.  Its
known-pair locus contains the
inverse image of the open set
$U\cap C$ near $a$.  The preceding
Riemann-surface argument makes
every fibre on $\widetilde C$
known, in particular the fibre
above $b$.  Thus $b\in U$.

\noindent\textbf{Step 6.} We have proved that $U$ is open
and closed.  Since $B$ is
connected and $U$ is nonempty,
$U=B$.  A nonreduced base
does not add fibre isomorphism
classes: restriction to its
reduction preserves both the
fibres and smoothness of the
family.  Thus the same conclusion
also applies to such a base
when it is admitted in the
definition of deformation
equivalence.

\noindent\textbf{Step 7.} Applying this result to each
link in a finite deformation
chain preserves the pair.
Conversely, E.1 identifies
all admissible sphere data
with a fixed pair with the
connected period family,
using its holomorphic
comparisons and its period
lattice $2kl\Z$.
The admissible half-plane
and its punctured-disc
quotient are connected.
Thus all members with
that pair are deformation
equivalent.  This proves
both directions.
\end{proof}



\begin{corollary}[separation for continuous complex structures]\label{push:src:E.2.24}
\StatementRecord{push:src:E.2.24}{proved}{Historical events: Claude Fable 5.2 (Anthropic), effort high, 16 November 2026: carried as background without fault found; targeted proof check; used and rechecked in this run. GPT-6 Astra (OpenAI), ChatGPT Pro mode, effort Pro, S6\_09\_PAPERS, 17 November 2026: targeted proof check. Also checked in this writing. Pre-continuation writing events: GPT-6 Astra (OpenAI), ChatGPT Pro mode, effort Pro, S6\_10\_PAPERS, 18 November 2026: targeted proof check; used and rechecked in this run. Continuation metadata events: GPT-6 Astra (OpenAI), ChatGPT Pro mode, effort Pro, S6\_10\_PAPERS continuation, 19 November 2026: used and rechecked in this run; status-vocabulary reconciliation; metadata reconciliation only; no new mathematical check. \emph{Pre-continuation recorded proof scope.} Meersseman Theorem3.4 and Remark3.6, printed pp.898-899, give a connected neighborhood, smooth-operator retraction and isotopic diffeomorphism comparison. Applying the arbitrary-base result to the connected Kuranishi reduction proves clopenness and its descent through the indicated quotient maps}{the hypotheses, notation and definitions stated here; \citet[Section~3.2, Theorem~3.4 and Remark~3.6,
pp.~898--899]{Meersseman2019E}; \cite[Section~2.1, pp.~883--884]{Meersseman2019E}; \cref{push:src:E.2.23}}

Let $\mathcal I(S^6)$ be the space of smooth integrable
almost-complex structures on the fixed smooth sphere,
with its $C^\infty$ topology.  For a fixed coprime pair
$(3k,4l)$, let $\mathcal I_{k,l}$ consist of the structures
whose complex manifolds are biholomorphic to constructed
members with that pair.  Then $\mathcal I_{k,l}$ is open
and closed in $\mathcal I(S^6)$.

Consequently a continuous family of such structures
over any connected topological parameter space cannot
contain constructed members with two different pairs.
The same separation holds in the quotient spaces by
$\operatorname{Diff}_0(S^6)$ or by all orientation-preserving
diffeomorphisms, with their quotient topologies.
This statement does not identify the number of marked
connected components belonging to a single pair.
\end{corollary}


\begin{proof}
\noindent\textbf{Step 1.} We use the classical Kuranishi local-completeness
theorem at an arbitrary $J_0\in\mathcal I(S^6)$,
including an initially unconstructed structure.
The complex orientation is locally constant, so the
oriented formulation applies on each such neighbourhood.
The precise $C^\infty$ version used here is
\citet[Section~3.2, Theorem~3.4 and Remark~3.6,
pp.~898--899]{Meersseman2019E}.  It supplies a connected
open neighbourhood $\mathcal U$ of $J_0$, a
finite-dimensional Kuranishi analytic space $K$
containing $J_0$, and a continuous retraction
$\Xi:\mathcal U\to K$ such that every $J\in\mathcal U$
is obtained from $\Xi(J)$ by a diffeomorphism isotopic
to the identity.  In that theorem the local product
map is first constructed in Sobolev completions;
elliptic regularity gives precisely this restriction
to smooth operators.  The topology on these operators
is the $C^\infty$ topology of
\cite[Section~2.1, pp.~883--884]{Meersseman2019E}.

\noindent\textbf{Step 2.} The Kuranishi space carries its proper smooth
holomorphic deformation family.  We can take $K$
connected: in the stated local theorem it is the
continuous retract of the connected $\mathcal U$.
If necessary replace $K$ by its reduction, which
changes neither its points nor the complex manifolds
in its fibres.  Suppose $\mathcal U$ contains some
$J\in\mathcal I_{k,l}$.  Its Kuranishi representative
$\Xi(J)$ has the same complex manifold up to
biholomorphism.  \ref{push:src:E.2.23} applied to the connected
Kuranishi family therefore makes every point of
$K$ a member of the same pair.  Since all points
of $\mathcal U$ are biholomorphic to their
Kuranishi representatives, it follows that
$\mathcal U\subset\mathcal I_{k,l}$.

\noindent\textbf{Step 3.} Thus at every $J_0$ the neighbourhood $\mathcal U$
is either contained in $\mathcal I_{k,l}$ or
disjoint from it.  This proves both openness and
closedness.  The inverse image of this set under
a continuous map from a connected parameter space
is then empty or the whole space, proving the
continuous-family assertion.  Finally
$\mathcal I_{k,l}$ is saturated under the indicated
diffeomorphism groups, because its definition uses
only the biholomorphism class.  Its image in either
quotient is open and closed by the definition of
the quotient topology.  No assertion that the
saturated set itself is connected is needed.
\end{proof}





\subsection{E.5. The two-parameter family and its actual singular fibres}
\begin{proposition}[the input configuration]\label{push:src:E.5.1}
\StatementRecord{push:src:E.5.1}{proved}{Historical events: Claude Fable 5.2 (Anthropic), effort high, 16 November 2026: carried as background without fault found. Claude Opus 5.5 (Anthropic), effort max, sweep 1, scope s6, 18 November 2026: used and rechecked in this run. GPT-6 Astra (OpenAI), ChatGPT Pro mode, effort Pro, S6\_09\_PAPERS, 17 November 2026: targeted proof check. Also checked in this writing. Pre-continuation writing events: GPT-6 Astra (OpenAI), ChatGPT Pro mode, effort Pro, S6\_10\_PAPERS, 18 November 2026 04:50:08.441928+00:00: targeted proof check. GPT-6 Astra (OpenAI), ChatGPT Pro mode, effort Pro, S6\_10\_PAPERS, 18 November 2026 04:51:18.698751+00:00: used and rechecked in this run. Continuation metadata events: GPT-6 Astra (OpenAI), ChatGPT Pro mode, effort Pro, S6\_10\_PAPERS continuation, 19 November 2026: used and rechecked in this run; status-vocabulary reconciliation; metadata reconciliation only; no new mathematical check. \emph{Pre-continuation recorded proof scope.} proved. Full comparison of the actual Weierstrass and threefold fibre configurations. Directly changed s=1/t,X=s\textsuperscript{2}x,Y=s\textsuperscript{3}y to obtain Y\textsuperscript{2}=X\textsuperscript{3}+s\textsuperscript{3}X+s\textsuperscript{6} and valuations(3,6,9). Opened ViaTwo Lemma2.1p3 for the exact III* identification. Its Proposition3.1pp6–7 gives the actual multiplicative periods with exponent pairs(1,0),(0,0) at the first two points and(1,0),(0,1) at the third, and its free order4 affine quotient pp7–8 gives4B*. The independent four-matrix calculation has nilpotent ranks1,1,2. Thus the earlier three-rank-one proposal is contradicted by the exact construction maps; the source’s own stated configuration is credited. Review provenance: New derivation after reading the supplied proof; not a blind review and not an independent human check}{the hypotheses, notation and definitions stated here; \cite{ViaTwo26push}; \cite[Lemma 2.1]{ViaTwo26push}; \cref{push:src:E.5.2}}

The proposed configuration with three rank-one reduced singular fibres
does not describe the construction of \cite{ViaTwo26push}. The elliptic surface
$y^2=x^3+tx+1$ has fibre $III^*$ at infinity.  The constructed
threefold has two rank-one fibres at $p_1,p_2$, a rank-two
hexagonal fibre at $p_3$, and the multiple bielliptic fibre
$4B_*$ at infinity.
\end{proposition}


\begin{proof}
\noindent\textbf{Step 1.} Put $s=t^{-1}$, $X=s^2x$, and $Y=s^3y$.  The minimal equation is
$Y^2=X^3+s^3X+s^6$, with coefficient and discriminant orders
$(3,6,9)$; this is the $III^*$ equation, as in
\cite[Lemma 2.1]{ViaTwo26push}.  The discriminant at finite
points is $-16(4t^3+27)$ with three simple zeros.  In the
threefold's multiplicative periods the orders at these points
are respectively $(1,0),(1,0),(1,1)$:
\[
 h(z,v)=(q\alpha(q)z,\beta(q)v),\qquad
 k(z,v)=(\gamma(q)z,\lambda q^\epsilon\delta(q)v),
\]
where all four coefficients are units, $\epsilon=0$ at $p_1,p_2$,
and $\epsilon=1$ at $p_3$.  Thus the first two torus ranks
drop by one and the last by two.  The explicit matrices of
\ref{push:src:E.5.2} independently have nilpotent ranks $1,1,2$; their exact
calculation rules out the assertion that all three are rank one.
These are the types stated in Proposition 3.1 and Theorem 1.1
of the cited source, not a correction to that source.
\end{proof}



\begin{proposition}[monodromy, centralizer, and integral lattices]\label{push:src:E.5.2}
\StatementRecord{push:src:E.5.2}{proved}{Historical events: Claude Fable 5.2 (Anthropic), effort high, 16 November 2026: carried as background without fault found. GPT-6 Astra (OpenAI), ChatGPT Pro mode, effort Pro, S6\_09\_PAPERS, 17 November 2026: targeted proof check. Also checked in this writing. Pre-continuation writing events: GPT-6 Astra (OpenAI), ChatGPT Pro mode, effort Pro, S6\_10\_PAPERS, 18 November 2026 04:50:08.441928+00:00: targeted proof check. GPT-6 Astra (OpenAI), ChatGPT Pro mode, effort Pro, S6\_10\_PAPERS, 18 November 2026 04:51:18.698751+00:00: used and rechecked in this run. Continuation metadata events: GPT-6 Astra (OpenAI), ChatGPT Pro mode, effort Pro, S6\_10\_PAPERS continuation, 19 November 2026: used and rechecked in this run; status-vocabulary reconciliation; metadata reconciliation only; no new mathematical check. \emph{Pre-continuation recorded proof scope.} proved. Complete integral matrix, centralizer, fixed-lattice, norm and vanishing-lattice calculation, with the exact geometric marking identified. Opened the actual ViaTwo monodromy derivation §4 pp9–11, including the affine winding shifts, the primitive local rank-two image and its saturation index d\textsuperscript{2}. Independently computed all products, the four inverse-transpose homology matrices, nilpotent ranks and the common rational centralizer. The 64 commutator equations have nullity2 and exact basis I,E41. Solving the finite kernel gives (a,−a,−a,d); the norm columns in(\ensuremath{\ell},\ensuremath{\delta}) have HNF diag(4,2). Independent HNF/Smith calculations give the actual integral images K1=Z(w+\ensuremath{\delta}),K2=Zu,K3=Zw\ensuremath{\oplus}Z\ensuremath{\delta} and the displayed saturated fixed bases. The source’s primitive local periods determine d\textsuperscript{2}=1 before the global topology. Exact output is analysis/II\_independent\_batch2\_exact.json; the unchanged original e\_push\_two\_parameter.py was also executed successfully and its complete output retained. Review provenance: New derivation after reading the supplied proof; not a blind review and not an independent human check}{the hypotheses, notation and definitions stated here; \cite{ViaTwo26push}}

Use the source's left-action marking
$H^1(F,\mathbb Z)=\mathbb Z\langle\psi,u,w,\delta\rangle$
and the ordered meridians $*,1,2,3$ with product one.
The cohomological matrices are
\[
\begin{aligned}
 T_*&=\begin{pmatrix}1&0&-2&-1\\0&0&-1&0\\0&1&0&-1\\0&0&0&1\end{pmatrix},&
 T_1&=\begin{pmatrix}1&0&2&2\\0&1&1&1\\0&0&1&0\\0&0&0&1\end{pmatrix},\\
 T_2&=\begin{pmatrix}1&0&0&0\\0&1&0&0\\0&-1&1&0\\0&0&0&1\end{pmatrix},&
 T_3&=\begin{pmatrix}1&0&0&-1\\0&1&1&0\\0&0&1&0\\0&0&0&1\end{pmatrix}.
\end{aligned}
\]
On ordinary dual homology put $H_i=(T_i^{-1})^t$,
$\ell=(1,-1,-1,0)^t$ and $\widehat\delta=(0,0,0,1)^t$.
Then
\[
 \operatorname{Cent}_{M_4(\mathbb Q)}(H_*,H_1,H_2,H_3)
 =\mathbb Q I\oplus\mathbb Q E_{41},
\]
and the cohomological centralizer is
$\mathbb Q I\oplus\mathbb Q E_{14}$.
The fixed and norm lattices at the multiple fibre are
\[
 \begin{aligned}
 \Lambda^{H_*}&=\mathbb Z\ell\oplus\mathbb Z\widehat\delta,\\
 N_4\Lambda&=4\mathbb Z\ell\oplus2\mathbb Z\widehat\delta,\\
 \Lambda^{H_*}/N_4\Lambda&\simeq\mathbb Z/4\oplus\mathbb Z/2.
 \end{aligned}
\]
The primitive vanishing lattices at the three reduced fibres are
\[
 K_1=\mathbb Z(\widehat w+\widehat\delta),\qquad
 K_2=\mathbb Z\widehat u,\qquad
 K_3=\mathbb Z\widehat w\oplus\mathbb Z\widehat\delta.
\]
Their invariant lattices are
\[
 \begin{split}
 \Lambda^{H_1}&=\mathbb Z(\widehat\psi-2\widehat u)
       \oplus\mathbb Z\widehat w\oplus\mathbb Z\widehat\delta,\\
 \Lambda^{H_2}&=\mathbb Z\widehat\psi\oplus
       \mathbb Z\widehat u\oplus\mathbb Z\widehat\delta,\\
 \Lambda^{H_3}&=K_3.
 \end{split}
\]
\end{proposition}


\begin{proof}
\noindent\textbf{Step 1.} The matrices are Proposition 4.2 of \cite{ViaTwo26push}, whose
derivation uses the elliptic monodromies, the winding of $P$,
and the primitive rank-two cusp.  In particular it determines
the initially undetermined winding integer by the saturation
index $d^2=1$, rather than assuming this winding from the
desired topology.  Direct multiplication gives
$T_*T_1T_2T_3=I$, $T_*^4=I$ and
$(T_i-I)^2=0$ for $i=1,2,3$.
For clarity the homological matrices used in our calculation are
\[
\begin{aligned}
 H_*&=\begin{pmatrix}1&0&0&0\\-2&0&-1&0\\0&1&0&0\\1&1&0&1\end{pmatrix},&
 H_1&=\begin{pmatrix}1&0&0&0\\0&1&0&0\\-2&-1&1&0\\-2&-1&0&1\end{pmatrix},\\
 H_2&=\begin{pmatrix}1&0&0&0\\0&1&1&0\\0&0&1&0\\0&0&0&1\end{pmatrix},&
 H_3&=\begin{pmatrix}1&0&0&0\\0&1&0&0\\0&-1&1&0\\1&0&0&1\end{pmatrix}.
\end{aligned}
\]
Solving the sixty-four rational equations $CH_i=H_iC$
gives $C=aI+bE_{41}$ with free $a,b$; the retained exact
row-reduction program displays a basis and verifies its
dimension two.  Inverse transpose gives the stated
cohomological centralizer.  Solving $(H_*-I)v=0$ gives
$v=(a,-a,-a,d)^t$ integrally, and
\[
 N_4=I+H_*+H_*^2+H_*^3
 =\begin{pmatrix}4&0&0&0\\-4&0&0&0\\-4&0&0&0\\0&2&-2&4\end{pmatrix}.
\]
In the basis $(\ell,\widehat\delta)$ its columns are
$(4,0),(0,2),(0,-2),(0,4)$, giving exactly the displayed
norm group.  Each $K_i$ is the actual integral image of
$H_i-I$; its Smith form has only unit nonzero entries.
The kernels of these three matrices give the remaining fixed
lattices.  All products, rational row reductions, Hermite forms,
and Smith forms are executed in
\texttt{code/e\_push\_two\_parameter.py}.
\end{proof}



We now specify the parameter domain, because the source proves
a family on a sufficiently small bidisc, not on the entire
$\mu$-plane.  Set
\[
 S_\mu:\ y^2=x^3+tx+\mu t+1,\qquad
 P_\mu=(-\mu,\sqrt{1-\mu^3}),
\]
with the square root equal to one at zero.  Label the three roots
$p_i(\mu)$ of $4t^3+27(\mu t+1)^2$ holomorphically near zero;
$p_3$ is the zero of the chosen linearization.
Choose a holomorphic generator of its one-dimensional space of
sections vanishing at $p_3$ and scale it by $\lambda$.
For the added homology generator take the inverse of the
contracting lift.  On the universal cover of the four-punctured
base its period columns then have the form
\[
 \Pi_{\mu,c}=
 \begin{pmatrix}2z_\mu&\tau_\mu&1&0\\
 \beta_\mu+c&z_\mu&0&1\end{pmatrix},
 \qquad \lambda=e^{-2\pi i c},\qquad
 \partial_c\Pi=e\otimes\psi.
\]
Here $z_\mu$ is a lifted elliptic logarithm in the chosen marking;
changing the generator of the linearization changes
$\beta_\mu$ and the origin of $c$ by a holomorphic function of
$\mu$, and changes neither the integer-difference criterion
below nor the original $\lambda$-labelled bidisc.
This choice of inverse lift fixes the sign of $c$.
At $p_3$ its second period has jump $-1$, exactly canceled by
$H_3\widehat\psi=\widehat\psi+\widehat\delta$.
The covariance identities for all four matrices, including
$\tau_*=-1/\tau$, $z_*=(z-1)/\tau$,
$\beta_* =\beta-1-2(z-1)^2/\tau$, are checked symbolically
in the same program.

\begin{theorem}[classification on the specified parameter domain]\label{push:src:E.5.3}
\StatementRecord{push:src:E.5.3}{proved}{Historical events: Claude Fable 5.2 (Anthropic), effort high, 16 November 2026: carried as background without fault found. GPT-6 Astra (OpenAI), ChatGPT Pro mode, effort Pro, S6\_09\_PAPERS, 17 November 2026: targeted proof check. Also checked in this writing. Pre-continuation writing events: GPT-6 Astra (OpenAI), ChatGPT Pro mode, effort Pro, S6\_10\_PAPERS, 18 November 2026 04:50:08.441928+00:00: targeted proof check. GPT-6 Astra (OpenAI), ChatGPT Pro mode, effort Pro, S6\_10\_PAPERS, 18 November 2026 04:51:18.698751+00:00: used and rechecked in this run. Continuation metadata events: GPT-6 Astra (OpenAI), ChatGPT Pro mode, effort Pro, S6\_10\_PAPERS continuation, 19 November 2026: used and rechecked in this run; status-vocabulary reconciliation; metadata reconciliation only; no new mathematical check. \emph{Pre-continuation recorded proof scope.} proved. Full parameter-domain, intrinsic-base, period comparison, finite norm and all-collar gluing classification. Read the adjacent exact parameter setup, including the inverse contracting generator and \ensuremath{\Pi}\_mu,c=[[2z,tau,1,0],[beta+c,z,0,1]], lambda=exp(−2\ensuremath{\pi}ic), and checked all four covariance identities in the source replay. Opened ViaTwo Proposition6.1pp17–18 and Proposition6.4pp19–20 in full. Uniform unit/contraction bounds on a smaller closed mu-disc give the actual half-plane family; all markings and \ensuremath{\ell}/4 extend on its simply connected parameter domain. The canonical form extends through uv=q, xyz=q and s=r\textsuperscript{4} with the actual r\textsuperscript{3} line-frame correction, giving the intrinsic anticanonical pencil. Differentiating the discriminant roots yields p\_i\ensuremath{{}^{\prime}}=−9/(2p\_i); an affine velocity relation would give a quadratic with three roots and nonzero constant. The resulting cross-ratio is locally injective and disjoint from its inverse, fixing mu and the base map. The exact common centralizer and constant last two period columns force A=alpha I and lattice map alpha(I−hE41). Summing the actual root-generator defects and using N4\ensuremath{\Lambda}=4Z\ensuremath{\ell}+2Z\ensuremath{\delta} forces alpha=1 and h even. For h=2n the explicit b*=−h(u+w)/4 and L*=−hu/2 satisfy the root conjugacy; the integral defects0,hu/2,−hw/2 lie in the actual three vanishing lattices and satisfy the full sphere cocycle identity. Rank-one node translations and the rank-two fan shifts extend with holomorphic unit remainders and have extending inverses. The remaining difference sheaf is the actual invariant pushforward of E=O\ensuremath{\oplus}O(−1) on P(1,4), as established from its nonzero line field and the verified canonical determinant; its H1 vanishes. Correcting the resulting Cech cocycle therefore glues the four exact local maps and their inverses. A width<2 logarithm strip gives the original small-bidisc injectivity. No extension to arbitrary mu or arbitrary prescribed punctured translations is inferred. Review provenance: New derivation after reading the supplied proof; not a blind review and not an independent human check}{the hypotheses, notation and definitions stated here; \cref{push:src:E.5.1}; \cite{ViaTwo26push}; \cref{push:src:E.5.2}}

There are $\epsilon>0$ and a real constant $C$ such that the
source construction, with its fixed logarithmic vector $\ell$,
defines a proper holomorphic family
$X_{\mu,c}$ over
$\{ |\mu|<\epsilon\}\times\{\operatorname{Im}c<C\}$.
It has the singular fibres of \ref{push:src:E.5.1} and the markings above.
Two of its members are biholomorphic if and only if
\[
 \mu'=\mu,\qquad c'-c\in2\mathbb Z.
\]
Every such biholomorphism covers the identity of the elliptic
base.  The quotient parameter in the second variable is
$v=e^{-\pi i c}$, with $\lambda=v^2$.
On a sufficiently small bidisc centred at any
$(0,\lambda_0)$ in Theorem 1.1 of \cite{ViaTwo26push}, choose
one logarithm of $\lambda$.  Within that actual bidisc,
\[
 X_{\mu,\lambda}\simeq X_{\mu',\lambda'}
 \quad\Longleftrightarrow\quad
 (\mu,\lambda)=(\mu',\lambda').
\]
\end{theorem}


\begin{proof}
\noindent\textbf{Step 1.} The equations of Proposition 6.4 in \cite{ViaTwo26push}
give the simultaneous elliptic family, line bundle, linearization,
and good-reduction comparison.  On a smaller closed $\mu$-disc
all unit factors in its four local fillings are uniformly bounded
above and below.  Its contraction estimates therefore have a
common nonzero bound $|\lambda|<\eta$.  Pulling back by
$\lambda=e^{-2\pi i c}$ gives the indicated half-plane.
The integral markings and the four-torsion section $\ell/4$
continue uniquely on this simply connected parameter domain.
On every relatively compact subset the source's finite root
quotients and the fixed rank-one and hexagonal fans give
holomorphic charts with uniform bounds.  Their comparisons are
the specified section-preserving ones, so they agree on
intersections and define a family.  Properness is local on
the parameter domain and follows from the compact four-piece
construction there; smoothness follows from the same free
finite action and unimodular charts.  We make no assertion
about continuation to arbitrary $\mu$.

\noindent\textbf{Step 2.} The intrinsic fibration is the anticanonical pencil.  To verify
the needed input from Proposition 6.1 of the source, on the
smooth core the line-coordinate form $dv/v$ wedged with a
base two-form identifies $K_X$ with $f^*\mathcal O(-1)$.
The forms $dq\wedge d\log z\wedge d\log v$ extend without
zeros or poles in both rank-one charts $q=ab$ and rank-two
charts $q=abc$.  At the finite root at infinity the source's
coordinate change gives
\[
 ds\wedge\frac{dX}{Y}\wedge\frac{dv}{v}
 =4\,dr\wedge\frac{dx'}{y'}\wedge\frac{dv'}{v'},
 \qquad s=r^4.
\]
It is invariant under the actual affine deck map and descends
through its free quotient.  Thus the same canonical formula
holds everywhere.  Connected fibres and projection formula
give $H^0(-K_X)=f^*H^0(\mathcal O(1))$, without base points.
A biholomorphism therefore induces an affine automorphism of
the elliptic base, fixing its distinguished multiple point
at infinity and its distinguished hexagonal point $p_3$.
It may initially interchange the two rank-one points.

\noindent\textbf{Step 3.} Put $Q(\mu)=(p_1-p_3)/(p_2-p_3)$.
The source computes $p_i'(0)=-9/(2p_i(0))$.
If $Q'(0)$ vanished, the three velocities would have the form
$Ap_i(0)+B$; the polynomial $Az^2+Bz+9/2$ would vanish
at three distinct numbers, a contradiction.
Also $Q(0)^2-Q(0)+1=0$, so $Q(0)\ne Q(0)^{-1}$.
Shrink $\epsilon$ so that $Q$ is injective and its image
is disjoint from the image of $Q^{-1}$.  Preservation or
interchange of the rank-one values gives respectively
$Q(\mu')=Q(\mu)$ or $Q(\mu')=Q(\mu)^{-1}$.
The second possibility is excluded, and the first gives
$\mu'=\mu$.  The base automorphism now fixes three distinct
points individually and is the identity.

\noindent\textbf{Step 4.} On regular fibres a biholomorphism is affine, and its integral
linear map commutes with all four $H_i$.  \ref{push:src:E.5.2} gives
$aI+bE_{41}$.  Comparing the constant third and fourth
period columns makes its complex linear part $aI$;
unimodularity gives $a=\alpha\in\{1,-1\}$.
The first column then gives $h=c'-c\in\mathbb Z$ and
lattice map $\alpha(I-hE_{41})$.
Normalize after the actual fourth-root base change.  Its
affine generator is the one with fourth power $\ell$, not
an arbitrary generator of its cyclic subgroup.  Summing the
translation defect around this generator gives
\[
 (I+hE_{41})\ell-\alpha\ell
 =(1-\alpha)\ell+h\widehat\delta\in N_4\Lambda.
\]
The norm lattice in \ref{push:src:E.5.2} forces $\alpha=1$ and $h$ even.
This proves necessity, including the multiple fibre's full
non-reduced structure.

\noindent\textbf{Step 5.} We prove sufficiency with all four collar conditions.
Let $h$ be even and put
\[
 b_*=-\tfrac h4(\widehat u+\widehat w),\qquad
 L_*=-\tfrac h2\widehat u.
\]
Direct multiplication gives
\[
 (I-H_*)b_*=\tfrac h4\widehat\delta+L_*.
\]
Consequently $\zeta\mapsto\zeta+\Pi_{\mu,c}b_*$ on the
complete root torus family conjugates the two affine actions,
modulo the integral arrow $L_*$.  It and its inverse extend
over $r=0$ and descend.  On the punctured root collar its
middle translation includes
$h(\log r)/(2\pi i)e$; this is the derivative-free collar
term for the change of the invariant translation
$\Pi\ell/4$, not a logarithm used on the central chart.

\noindent\textbf{Step 6.} At a rank-one point the vanishing coordinate is the first
toric coordinate in its chain cover, while the other
direction is the universal cover of the elliptic factor.
An extending translation has integral winding only in $K_j$.
After removing the corresponding power of $q$, its toric
coefficient is a holomorphic unit and its elliptic translation
extends holomorphically.  These are the full analytic collar
conditions.  Necessity follows by lifting a transverse disc
through a smooth central point to this free cover: its residual
defect in $\Lambda/K_j$ is zero.  The coordinate ratios on
the cover then give the asserted unit and holomorphic
remainder.  Conversely an integral power of $q$ translates
the infinite chain by an integer number of components;
the holomorphic remainders give torus and elliptic translations.
These maps and their inverses extend at both branches of every
node.  At the hexagonal point the same description uses two
toric windings in $K_3$, and integer translation of the fan
extends over its three double curves and both triple points.
This criterion includes regularity of the remainder; membership
of the defect in $K_j$ alone would not exclude an essential
singularity of a preassigned punctured translation.

\noindent\textbf{Step 7.} Choose the three integral defects
\[
 L_1=0,\qquad L_2=\tfrac h2\widehat u,\qquad
 L_3=-\tfrac h2\widehat w.
\]
They belong to $K_1,K_2,K_3$, respectively, and are realized
by the preceding toric translations with initially constant
unit remainder.  In primitive cocharacter coordinates these
are $1$, $q_2^{h/2}$, and $(q_3^{-h/2},1)$, where the
monodromy defect of $q^a$ is defined as $a$ times the chosen
vanishing generator.  The required sphere relation is
\[
 L_*+H_*L_1+H_*H_1L_2+H_*H_1H_2L_3=0.
\]
It holds as an exact integer identity, checked in the program.
Thus the four local choices have compatible additive defects
and define a torsor of allowed middle translations.

\noindent\textbf{Step 8.} We verify the vanishing of its analytic obstruction, including
the difference sheaf.  The root curve has just the order-four
point at infinity and is $\mathfrak C=\mathbb P(1,4)$,
with $K_{\mathfrak C}=\mathcal O(-5)$ and coarse pullback
$\mathcal O_B(1)=\mathcal O_{\mathfrak C}(4)$.
Vertical logarithmic fields on a rank-one chart are generated
by $a\partial_a-b\partial_b$ and the elliptic vector;
on a hexagonal chart they are the two height-zero toric fields.
The root chart has the ordinary vertical tangent frame.
Their linear transition functions depend only on the base;
they define a rank-two bundle $\mathcal E$ on $\mathfrak C$.
The fixed line-coordinate field supplies a subline
$\mathcal O\subset\mathcal E$ everywhere in these logarithmic
frames.  The logarithmic determinant sequence and cancellation
of the three reduced fibre divisors give
$K_X=\bar f^*(K_{\mathfrak C}\otimes(\det\mathcal E)^{-1})$.
The verified canonical formula and
$\bar f_*\mathcal O_X=\mathcal O_{\mathfrak C}$ therefore
give $\det\mathcal E=\mathcal O(-1)$.
The extension splits since
$H^1(\mathfrak C,\mathcal O(1))=0$, by the same two-chart
Laurent calculation with weights $1,4$.
The coarse difference sheaf is $V_B=\rho_*\mathcal E$:
on the root chart this is precisely finite equivariance,
and on each reduced collar it is precisely the holomorphic
remainder just described.  Finite averaging is exact, so
\[
 H^1(B,V_B)=H^1(\mathfrak C,\mathcal E)
 =H^1(\mathcal O\oplus\mathcal O(-1))=0.
\]
The torsor's Cech class vanishes.  Correcting the four local
translations by its holomorphic zero-cochain gives a global
middle translation and exact equality with all four local
extensions.  Their inverses have the opposite allowed
translations and glue as well.  This constructs the required
biholomorphism over $B$.

\noindent\textbf{Step 9.} Finally choose a single logarithm of $\lambda$ on the source's
small bidisc.  Its $c$-values can be placed in a strip of
width less than two.  The criterion just proved then says
$c'=c$ and $\mu'=\mu$, which is equivalent to equality of
the original two parameters.  Conversely equality gives the
identity, with the fixed root, rank-one and hexagonal charts.
\end{proof}



The matrix data and the independent-variation argument are in
\cite[Proposition 4.2 and Proposition 6.4]{ViaTwo26push}; the
source states dimension at least two and does not state the
integer classification on the logarithmic lift.  We derived
the latter here from the marked norm and collar calculations.
For E.1--E.5 our searches on 15 November 2026 included the
exact arXiv identifiers together with ``classification'',
``Kuranishi'', ``deformation'', and ``parameters'', and the
references and version-specific records of the attached
papers.  No further paper resolving the general deformation
equivalence in \ref{push:src:E.2.4} was located.  This is a record of the
search performed, a record with the stated search scope;
the query responses and opened-source metadata are retained
in \texttt{evidence/E\_search*.json}.


\section{Comb invariants and the precise vector-field correction}\label{sec:II:combs}
\begin{proposition}[Canonical bundle, exact tangent cohomology and Hodge growth on the combs]\label{inv:E2.2}
\StatementRecord{inv:E2.2}{proved}{Historical events: Claude Fable 5.2 (Anthropic), effort high, 16 November 2026: carried as background without fault found. GPT-6 Astra (OpenAI), ChatGPT Pro mode, effort Pro, S6\_09\_PAPERS, 17 November 2026: targeted proof check. Also checked in this writing. Pre-continuation writing events: GPT-6 Astra (OpenAI), ChatGPT Pro mode, effort Pro, S6\_10\_PAPERS, 18 November 2026 04:50:08.441928+00:00: targeted proof check. GPT-6 Astra (OpenAI), ChatGPT Pro mode, effort Pro, S6\_10\_PAPERS, 18 November 2026 04:51:18.698751+00:00: used and rechecked in this run. Continuation metadata events: GPT-6 Astra (OpenAI), ChatGPT Pro mode, effort Pro, S6\_10\_PAPERS continuation, 19 November 2026: used and rechecked in this run; status-vocabulary reconciliation; metadata reconciliation only; no new mathematical check. \emph{Pre-continuation recorded proof scope.} proved. Full root-bundle, tangent direct-image, cusp modification and middle-Hodge-number calculation on the stated m\ensuremath{\geq}12 comb members. The retained real homology/cohomology classes give nonzero vertical and Hodge-quotient sections on ordinary fibres, finite covers and the explicit cusp graph; source Combs Proposition6.1 pp19–21 supplies the actual logarithmic periods log(q)I/(2\ensuremath{\pi}i)+H(q), and Section6.3 p23 supplies the complete Hodge quotient extension. The stabilizer exponents(2,1) and coarse determinant2m−1 give root weight24m−1 for both E and G; detH is trivial by its determinant-one monodromy and trace-zero residue. The log tangent determinant and removal of the actual cusp divisor O(12) give K=O(−24m−6). Weighted monomials show H0O(5)=0, and factoring E\ensuremath{\to}O(5) through its quotient gives the remaining form vanishings. Independently recomputed the invariant multiplication carry s\textsuperscript{4}/s=t at the order3 point and no order4 carry, so deg\ensuremath{\rho}*(E\ensuremath{\otimes}G)=8m−3, H1=0 and h0=8m+1. The two KS cup maps are injective because the cusp leading matrix is invertible. Their root determinants and exact local character sums give coarse degrees2m−2 and6m−2, hence a length4m cokernel away from the cusp. Read and independently derived fam:src:A.5.2’s branch-normal map, conormal kernel generated by ydx,xdy or yzdx,xzdy,xydz, common dq gluing, Gorenstein Serre duality h\textsuperscript{2}(T|W)=1, and rank-one DVR residue criterion. Thus the free R\textsuperscript{2}T line increases by p0, descending to O(2m), while the separate torsion length4m remains. Leray yields tangent tuple(2m+1,8m+1,6m+1,0). For the cotangent sequence the exact degree difference is4m, its other kernel is a rank3 bundle K\_m, and residue dq/q gives O\_B(−2). Rational normalization kills higher residue cohomology. Leray and \ensuremath{\chi}(\ensuremath{\Omega}\textsuperscript{1})=−1 give h11=4m+1+h0K\_m and h12=4m+h0K\_m. The source’s actual classifying-family lower bound is8m−8 (Lemma8.2/Proposition8.3 pp27–28 and specializationp30); no claim is made about a smooth family joining the comb to X0. Review provenance: New derivation after reading the supplied proof; not a blind review and not an independent human check}{the hypotheses, notation and definitions stated here; \cref{cdp:st:E15}; \cref{inv:E2.1}}

For the same $X_m$, the root lift $\bar f_m:X_m\to\mathcal B$ has
vertical bundle $E_m$ and $G_m=R^1\bar f_{m*}\OO_{X_m}$ satisfying
\[
 0\longrightarrow\OO\longrightarrow E_m
       \longrightarrow\OO(24m-1)\longrightarrow0,
 \qquad
 0\longrightarrow\OO\longrightarrow G_m
       \longrightarrow\OO(24m-1)\longrightarrow0.      \tag{E.2}\label{invII:eq:E.2}
\]
The canonical bundle and all positive-degree holomorphic forms obey
\[
 K_{X_m}=\bar f_m^*\OO(-24m-6),\qquad
 h^{1,0}=h^{2,0}=h^{3,0}=0.                            \tag{E.3}\label{invII:eq:E.3}
\]
The tangent cohomology is that of Proposition~\ref{cdp:st:E15}.
The following proof gives a distinct calculation of those dimensions.
More precisely, writing $\mathcal L=\OO_{\mathcal B}(-5)$,
\[
 \begin{split}
 f_{m*}T_{X_m}&=\rho_*E_m,\\
 R^1f_{m*}T_{X_m}&=
  \operatorname{coker}\{\rho_*\mathcal L\xrightarrow{\delta_0}
                                      \rho_*(E_m\otimes G_m)\},\\
 R^2f_{m*}T_{X_m}&\simeq\OO_B(2m)\oplus Q_m,
 \qquad\operatorname{length}Q_m=4m,
 \end{split}                                        \tag{E.4a}\label{invII:eq:E.4a}
\]
where $Q_m=\operatorname{coker}\{\rho_*(\mathcal L\otimes G_m)
\xrightarrow{\delta_1}\rho_*(E_m\otimes\det G_m)\}$ has no support
at the cusp. The splitting in the last line need not be canonical.
There is a rank-three vector bundle $K_m$ on $B$ for which
\[
 \boxed{h^{1,1}(X_m)=4m+1+h^0(K_m),\qquad
 h^{1,2}(X_m)=4m+h^0(K_m).}                            \tag{E.5}\label{invII:eq:E.5}
\]
In particular both remaining middle Hodge numbers grow at least
linearly with $m$.
\end{proposition}


\begin{proof}
\noindent\textbf{Step 1.} The construction retains the integral local system and the limiting
period planes. The invariant real homology class $\widehat\delta$
therefore gives a vertical translation field on every regular torus.
It is nowhere zero there, since the real period map is an isomorphism.
At a finite root point the same reasoning applies on the smooth
cover, and at the cusp its nonzero angular class gives a nonzero
logarithmic vertical vector. The local vertical transitions, which
depend only on the base exactly as in \cref{inv:D1.2}, define $E_m$ and give its
trivial subbundle. The invariant real cohomology class $\gamma$ gives
a nowhere-zero section of $G_m$, by the same smooth-fibre and cusp
comparison in \cref{inv:D1.2}.

\noindent\textbf{Step 2.} The determinant of $G_m$ has the same two stabilizer characters as
that of $G_0$: invariant exponents $2$ at $p_1$ and $1$ at $p_2$.
Its coarse determinant has degree $2m-1$ by \ref{inv:E2.1}. The explicit descent
formula in \cref{inv:D1.1} therefore identifies
$\det G_m=\OO(12(2m-1)+8+3)=\OO(24m-1)$.
The rank-four flat bundle on the root base has trivial determinant:
the integral monodromies have determinant one, the finite-root charts
have no residual ramification, and the logarithmic cusp monodromy
has nilpotent residue of trace zero. Its flat volume section thus
extends without a zero or pole. The exact Hodge-plane sequence
$0\to E_m^*\to\mathcal H\to G_m\to0$ gives
$\det E_m=\det G_m$. Taking quotients by the two nonvanishing
sections proves \eqref{invII:eq:E.2}. The determinant of the log tangent sequence
$0\to\bar f_m^*E_m\to T_{X_m}(-\log W)
\to\bar f_m^*\OO(-5)\to0$ then gives \eqref{invII:eq:E.3}'s canonical formula.

\noindent\textbf{Step 3.} The logarithmic Kodaira--Spencer map is invertible near the cusp.
This does not follow merely from fixing a limiting plane:
Viaclovsky's Proposition 6.1 gives the actual local period expression
$(\log q)I/(2\pi i)+H_m(q)$ with $H_m$ holomorphic. Its logarithmic
derivative has invertible leading matrix $I/(2\pi i)$, or
$B_0/(2\pi i)$ after the retained marking change. Thus the map is
generically invertible. The direct image of the logarithmic
cotangent sequence consequently has kernel exactly $\OO(5)$.
It has no section by \cref{inv:D1.1}, so $h^{1,0}=0$. To treat two-forms, take
the second exterior power of that sequence. Its possible global
sections lie in
$H^0(E_m^*\otimes\OO(5))$ or
$H^0((\det E_m)^{-1})$.
The latter is zero, having weight $1-24m<0$. A morphism
$E_m\to\OO(5)$ vanishes on the subline $\OO$, since
$H^0(\OO(5))=0$, and therefore factors through $\OO(24m-1)$.
The resulting section would have weight $6-24m<0$, so it too
vanishes. Absolute two-forms inject into logarithmic ones, which
proves $h^{2,0}=0$. The negative canonical weight proves $h^{3,0}=0$.

\noindent\textbf{Step 4.} The first connecting map from the horizontal line to
$E_m\otimes G_m$ is nonzero near the cusp, so it is injective as a
sheaf map. The same normal quotient (D.10) has no degree-zero direct
image, giving $f_{m*}T_{X_m}=\rho_*E_m$.
Exact invariants applied to \eqref{invII:eq:E.2} give, for each of $V=E_m,G_m$,
\[
 0\longrightarrow\OO_B\longrightarrow\rho_*V
 \longrightarrow\OO_B(2m-1)\longrightarrow0.
\]
These bundles have vanishing $H^1$ and $h^0=2m+1$.
They also split into two nonnegative-degree lines: if a splitting
had one negative summand, the section of its saturated $\OO_B$
subline would have to lie in the positive summand. Every nonzero
section of that summand has a zero, contradicting saturation.

\noindent\textbf{Step 5.} The natural map
\[
 (\rho_*E_m)\otimes(\rho_*G_m)
              \longrightarrow\rho_*(E_m\otimes G_m)
\]
is injective, with length-one cokernel at $p_1$ and no other
cokernel. Indeed both bundles have invariant exponents $(0,2)$ at
$p_1$. The product of their exponent-two generators is $s^4$,
whereas the product representation has invariant generator $s$;
the comparison on that one summand is multiplication by $t=s^3$.
All other exponents add without a carry. At $p_2$ the exponents
are $(0,1)$, and no sum reaches four. At the cusp there is no
stabilizer. The source is a sum of nonnegative-degree line bundles,
so its $H^1$ is zero. The length-one upper modification consequently
has $H^1=0$, rank four and degree
$2(2m-1)+2(2m-1)+1=8m-3$; it has $8m+1$ global sections.

\noindent\textbf{Step 6.} In the direct-image sequence of the log tangent extension, the
second cup map
$\delta_1:\mathcal L\otimes G_m\to E_m\otimes\det G_m$
is generically invertible, by the invertible cusp Kodaira--Spencer
matrix. It is therefore injective as a sheaf map. The first cup map
$\delta_0$ is injective for the same reason. Consequently
$R^1\bar f_{m*}T(-\log W)=\operatorname{coker}\delta_0$.
The normal quotient $\nu_*K_P$ in (D.10) has zero direct images in
degrees zero and one, so the same formula holds for $R^1T$.
Since $\rho_*\mathcal L=\OO_B(-1)$ is acyclic, its cokernel has
$h^0=8m+1$ and $h^1=0$. Leray and $H^1(\rho_*E_m)=0$ now give
the exact value $h^1(T)=8m+1$.

\noindent\textbf{Step 7.} For completeness the degree-two calculation gives Proposition~\ref{cdp:st:E15} directly,
with the full torsion correction retained. The coarse degrees of
$\mathcal L\otimes G_m$ and $E_m\otimes\det G_m$ are $2m-2$
and $6m-2$. Their invariant exponents at $p_1$ are respectively
$(1,0)$ and $(2,1)$, and at $p_2$ both are $(1,2)$.
Their root determinant weights are $24m-11$ and $72m-3$;
subtracting the character contributions and dividing by twelve
gives exactly these coarse degrees. Thus $Q_m$ has length $4m$.
It has no cusp contribution, where the cup map is invertible.
On the root base there is an exact sequence
\[
 0\longrightarrow\operatorname{coker}\delta_1
 \longrightarrow R^2\bar f_{m*}T(-\log W)
 \longrightarrow\mathcal L\otimes\det G_m
                 =\OO(24m-6)\longrightarrow0.
\]
It splits, since the extension group is the first cohomology of a
zero-dimensional sheaf. The cusp normal quotient then raises the
degree of the free line by $p_0$, exactly as in \cref{fam:src:A.5.2}. The local
fibre, its normalization and the logarithmic residue are unchanged;
in particular $h^2(T|_W)=1$ still excludes a split torsion summand at
the cusp. Hence the new free line is $\OO(24m+6)$, which descends to
$\OO_B(2m)$, proving \eqref{invII:eq:E.4a}. Leray gives
$h^2(T)=(2m+1)+4m=6m+1$ and $h^3(T)=0$.
As an independent check, these values satisfy the full tangent
index $\chi(T)=1$. The source family of Viaclovsky's Proposition 8.3
also gives
\[
 8m-8\le\dim\operatorname{Def}(X_m)\le8m+1;
\]
the lower bound concerns the actual classifying-family dimension,
while the upper bound is the tangent dimension now computed.

\noindent\textbf{Step 8.} We give the bundle specifying the remaining Hodge number. Set
\[
 K_m=\ker\left\{\rho_*(E_m^*\otimes G_m)
  \xrightarrow{\ \kappa_1\ }
       \rho_*(\OO(5)\otimes\det G_m)\right\}.         \tag{E.6}\label{invII:eq:E.6}
\]
The map is cup product with the actual logarithmic
Kodaira--Spencer class, in the category of coherent sheaves on $B$.
It is generically onto, so the kernel is a rank-three vector bundle.
The preceding connecting map
\[
 \rho_*E_m^*\xrightarrow{\ \kappa_0\ }
                 \rho_*(\OO(5)\otimes G_m)
\]
is injective. The source and target have degrees $-2m-1$ and
$2m-1$. To check these degrees, use their fixed local character
exponents $(0,1)$ versus $(2,1)$ at the order-three point and
$(0,3)$ versus $(3,0)$ at the order-four point, together with the
root determinant degrees. The invariant determinant comparison in
\cref{inv:D1.1} gives precisely the displayed integers. Its cokernel $A_m$
therefore has length $4m$, irrespective of how its zeros are
distributed among regular and special fibres.

\noindent\textbf{Step 9.} Writing $\mathcal A_m=\Omega^1_{X_m}(\log W)$, the direct-image
sequence gives
$0\to A_m\to R^1f_{m*}\mathcal A_m\to K_m\to0$.
The residue sequence
$0\to\Omega^1_{X_m}\to\mathcal A_m\to\nu_*\OO_P\to0$
has a surjective degree-zero residue map at $p_0$: the logarithmic
horizontal form $dq/q$ has residue one on every branch. Its kernel
is the one elementary modification
$f_{m*}\Omega^1=\OO_B(-2)$ of
$f_{m*}\mathcal A_m=\rho_*\OO(5)=\OO_B(-1)$.
The rationality of $P$ makes the higher direct images of the
residue quotient zero. Thus
$R^1f_{m*}\Omega^1=R^1f_{m*}\mathcal A_m$.
Since $H^1(A_m)=0$, its global sections have dimension
$4m+h^0(K_m)$, and Leray adds
$h^1(\OO_B(-2))=1$ in total degree one. This proves the formula
for $h^{1,1}$. Finally the sphere Chern numbers give
$\chi(\Omega^1)=-1$, and the vanishing just proved, together with
Serre duality, gives $h^{1,1}-h^{1,2}=1$. This proves \eqref{invII:eq:E.5}.
\end{proof}



\begin{proposition}[Counterexample: The universal two-vector-field bound]\label{inv:E2.4}
\StatementRecord{inv:E2.4}{proved}{Historical events: Claude Fable 5.2 (Anthropic), effort high, 16 November 2026: carried as background without fault found. GPT-6 Astra (OpenAI), ChatGPT Pro mode, effort Pro, S6\_09\_PAPERS, 17 November 2026: targeted proof check. Also checked in this writing. Pre-continuation writing events: GPT-6 Astra (OpenAI), ChatGPT Pro mode, effort Pro, S6\_10\_PAPERS, 18 November 2026 04:50:08.441928+00:00: targeted proof check. GPT-6 Astra (OpenAI), ChatGPT Pro mode, effort Pro, S6\_10\_PAPERS, 18 November 2026 04:51:18.698751+00:00: used and rechecked in this run. Continuation metadata events: GPT-6 Astra (OpenAI), ChatGPT Pro mode, effort Pro, S6\_10\_PAPERS continuation, 19 November 2026: used and rechecked in this run; status-vocabulary reconciliation; metadata reconciliation only; no new mathematical check. \emph{Pre-continuation recorded proof scope.} proved counterexample; universal two-vector-field bound refuted. Exact application of the independently derived comb tangent dimensions and full primary-source comparison. The invariant comb tangent computation gives h0(T)=2m+1, so every stated m\ensuremath{\geq}12 member has at least25 independent fields. Opened HKP arXiv math/9812076v1 Corollary1.3 and its complete proof onp2: it expresses an additional field using meromorphic coefficients and then invokes CDP98 to make those coefficients constant. The actual fibration gives a nonconstant pullback of a rational base function; the explicit computed field space supplies the counterexample. AC26 original Remarks1.2p4 and9.26p83 were read for the repeated universal bound. This compares the stated universal result and its invoked step; it does not claim an independent re-proof of CDP98 or a refutation of the algebraic-dimension-zero version. Review provenance: New derivation after reading the supplied proof; not a blind review and not an independent human check}{the hypotheses, notation and definitions stated here; \citet[Corollary 1.3; verified preprint p.~2]{HKP2000}; \citet[Remarks 1.2 and 9.26]{AC2026}; \cref{inv:E2.2}}

The assertion that every complex structure on $S^6$ has
$h^0(T)\le2$, printed in
\citet[Corollary 1.3; verified preprint p.~2]{HKP2000}
and repeated in \citet[Remarks 1.2 and 9.26]{AC2026}, is false
for the comb members.
\end{proposition}


\begin{proof}
\noindent\textbf{Step 1.} \ref{inv:E2.2} gives $h^0(T_{X_m})=2m+1\ge25$ for $m\ge12$.
The source proof expresses a third field as a meromorphic linear
combination of two fields and invokes the Campana--Demailly--Peternell
claim that all meromorphic functions on a complex six-sphere are
constant. That implication is unavailable here: $f_m$ supplies
nonconstant meromorphic functions. The independently computed
$2m+1$ sections above give the counterexample, so this is a
refutation of the bound as well as a location of the failed step.
It does not refute a statement restricted to algebraic dimension
zero.
\end{proof}





\section{Uniform non-torsion on every full fibre}\label{sec:II:uniform-nontorsion}
\begin{theorem}[uniform non-torsion]
\label{wb:thm:uniform-nontorsion-audit}
\StatementRecord{wb:thm:uniform-nontorsion-audit}{proved}{Historical events: Claude Fable 5.2 (Anthropic), effort high, 12 November 2026: targeted proof check. GPT-6 Astra (OpenAI), ChatGPT Pro mode, effort Pro, S6\_09\_PAPERS, 17 November 2026: targeted proof check. Also checked in this writing. Pre-continuation writing events: GPT-6 Astra (OpenAI), ChatGPT Pro mode, effort Pro, S6\_10\_PAPERS, 18 November 2026: targeted proof check; used and rechecked in this run. Continuation metadata events: GPT-6 Astra (OpenAI), ChatGPT Pro mode, effort Pro, S6\_10\_PAPERS continuation, 19 November 2026: used and rechecked in this run; status-vocabulary reconciliation; metadata reconciliation only; no new mathematical check. \emph{Pre-continuation recorded proof scope.} Re-derived the nonzero Picard parameter of f*O(1), the discrete rank-one or rank-two inverse lattice on each smooth fibre, and the torsion criterion z in its rational span. Evaluating the four real period columns shows that the extra integral class lies on the gamma line exactly when c tau+d mu-b=0. The cusp and order-four transformation laws prove independence of1,tau,mu, so the rank-two exceptions are countably many fibres. The full-fibre restriction is injective, and the primitive normal characters prove nonzero restriction on each reduction; exponential-sequence preimages then add only countable bad sets. Exact receiving dependencies inv:C3.1, inv:A1.0, intrinsic specialization and the normal-character theorem were reopened}{the hypotheses, notation and definitions stated here; \cref{inv:wb:eq:9.19}; \cref{inv:wb:eq:9.23}; \cref{inv:wb:eq:9.22}; \cref{inv:wb:eq:9.24}; \cref{inv:wb:eq:9.25}}

\ There is a line bundle \(L\in\operatorname{Pic}(X)\) such that
\(L|_{X_b}\) is non-torsion for every scheme-theoretic fibre
\(X_b=f^*(b)\), and \(L|_{S_j}\) is non-torsion for \(j=1,2\).
More precisely, identify \(\operatorname{Pic}(X)\) with \(\mathbb C\),
\(z\mapsto L_z\).  Every \(z\) outside the union of \(\mathbb R\) with a
countable set has these properties.
\end{theorem}



\begin{proof}
\noindent\textbf{Step 1.} The pullback \(f^*\mathcal O_B(1)\) is nontrivial: by the projection
formula and \(f_*\mathcal O_X=\mathcal O_B\),
\[
 h^0(X,f^*\mathcal O_B(1))=h^0(B,\mathcal O_B(1))=2,
\]
whereas \(h^0(X,\mathcal O_X)=1\).  Rescale the additive coordinate in
\cref{inv:C3.1} so that
\[
 L_{12}=f^*\mathcal O_B(1).
\tag{WBAPII.9.19}\label{inv:wb:eq:9.19}
\]

\noindent\textbf{Step 2.} Fix a smooth fibre \(F_q\).  Put
\[
 \Lambda_q=\iota_q^{-1}\!\left(
 \operatorname{im}[H^1(F_q;\mathbb Z)\to
 H^1(F_q,\mathcal O_{F_q})]\right)\subset\mathbb C.
\tag{WBAPII.9.20}\label{inv:wb:eq:9.20}
\]
The integral image is a lattice in the underlying real vector space of
\(H^1(F_q,\mathcal O_{F_q})\), so \(\Lambda_q\) is discrete.  It contains
\(12\), because \eqref{inv:wb:eq:9.19} is trivial on every fibre.  The exponential
sequence on \(F_q\) shows
\[
 L_z|_{F_q}\text{ is torsion}
 \quad\Longleftrightarrow\quad z\in\mathbb Q\Lambda_q.
\tag{WBAPII.9.21}\label{inv:wb:eq:9.21}
\]
If \(\operatorname{rk}\Lambda_q=1\), its nonzero element \(12\) forces
\(\Lambda_q\subset\mathbb R\); all its bad parameters are real.

\noindent\textbf{Step 3.} Suppose \(\operatorname{rk}\Lambda_q=2\).  By
\cref{inv:C3.1}, the image of \(\iota_q\) is the line
generated by the image of \(\gamma\).  Thus there is an integral class
\[
 \lambda=a\gamma+bu+cw+d\delta\notin\mathbb Q\gamma
\tag{WBAPII.9.22}\label{inv:wb:eq:9.22}
\]
whose \((0,1)\)-image lies on that line.  On the covering family, the four
marked lattice columns are
\[
 \pi_{\widehat\gamma}=(6\mu,\beta_c),\quad
 \pi_{\widehat u}=(\tau,\mu),\quad
 \pi_{\widehat w}=(1,0),\quad
 \pi_{\widehat\delta}=(0,1).
\tag{WBAPII.9.23}\label{inv:wb:eq:9.23}
\]
Every real cohomology class has a unique expression
\[
 2\operatorname{Re}(\bar\xi\,d\zeta_1+
                       \bar\upsilon\,d\zeta_2).
\]
Evaluating it on \eqref{inv:wb:eq:9.23} gives
\[
\begin{aligned}
2\operatorname{Re}\bar\xi&=c,&
2\operatorname{Re}\bar\upsilon&=d,\\
2\operatorname{Re}(\bar\xi\tau+\bar\upsilon\mu)&=b,&
2\operatorname{Re}(6\bar\xi\mu+\bar\upsilon\beta_c)&=a.
\end{aligned}
\tag{WBAPII.9.24}\label{inv:wb:eq:9.24}
\]
For \(\gamma\), the first two right sides vanish.  The third equation then
shows that its \((0,1)\)-line is represented, up to a nonzero scalar, by
\((\operatorname{Im}\mu,-\operatorname{Im}\tau)\); the fourth equation
only fixes the scalar.  Write
\(\bar\xi=c/2+ir\), \(\bar\upsilon=d/2+is\).  The condition that the class
in \eqref{inv:wb:eq:9.22} have \((0,1)\)-image on the \(\gamma\)-line is
\[
 \bar\xi\operatorname{Im}\tau+
 \bar\upsilon\operatorname{Im}\mu=0.
\]
Its real part, together with the third equation of \eqref{inv:wb:eq:9.24}, is equivalent
to the single holomorphic equation
\[
 c\tau+d\mu-b=0.
\tag{WBAPII.9.25}\label{inv:wb:eq:9.25}
\]
Neither \(a\) nor \(\beta_c\) occurs, as adding a multiple of \(\gamma\)
must not change the condition.

\noindent\textbf{Step 4.} The functions \(1,\tau,\mu\) are linearly independent.  Applying the cusp
transformation \((\tau,\mu)\mapsto(\tau-1,\mu)\) to a constant linear
relation first kills the coefficient of \(\tau\).  If the remaining
relation made \(\mu\) constant, the order-four law
\(\mu\circ g_2=1+\mu/\tau\), with nonconstant \(\tau\), would force an
impossible identity.  Hence all coefficients vanish.

\noindent\textbf{Step 5.} For every integral \(\lambda\notin\mathbb Q\gamma\), equation \eqref{inv:wb:eq:9.25} is
therefore a nonzero holomorphic function on the covering Riemann surface.
Its zero set is discrete and countable.  There are countably many
\(\lambda\), so their union, and its image in \(B^\circ\), are countable.
This image contains every \(q\) for which
\(\operatorname{rk}\Lambda_q=2\).  Each such fibre contributes the
countable bad set \(\mathbb Q\Lambda_q\); all rank-one fibres together
contribute only parameters in \(\mathbb R\).  Thus the smooth fibres
exclude only \(\mathbb R\) together with a countable set.

\noindent\textbf{Step 6.} There remain the three scheme-theoretic special fibres.  By
\cref{inv:A1.0,mapsI:top:thm:p0-integral-specialization} and \cref{inv:C3.1}, restriction
\[
 \mathbb C\longrightarrow H^1(X_b,\mathcal O_{X_b})
\]
is complex-linear and injective for each of them.  The exponential sequence
and countability of \(H^1(X_b;\mathbb Z)\) show that the parameters for
which \(L_z|_{X_b}\) is torsion form a countable set.

\noindent\textbf{Step 7.} Finally consider the reductions \(S_j\).  The restriction
\(\varphi_j:\mathbb C\to H^1(S_j,\mathcal O_{S_j})\) is complex-linear.
It is nonzero.  Write \(\mathcal O_X(S_j)=L_{z_j}\).  Since
\[
 \mathcal O_X(m_jS_j)=f^*\mathcal O_B(1)=L_{12},
\]
the additive parameter gives \(z_j=12/m_j\), while
\cref{wb:ffcert:thm:normal-trivialization} says that
\[
 L_{12/m_j}|_{S_j}=\mathcal O_X(S_j)|_{S_j}
\]
is nontrivial.  If \(\varphi_j\) were zero, the exponential sequence and
the vanishing of the restricted first Chern class would make every
\(L_z|_{S_j}\) trivial, a contradiction.  Thus \(\varphi_j\) is
injective, and the same exponential-sequence argument makes its torsion
parameters countable.  Removing these two additional countable sets proves
the theorem.
\end{proof}



\begin{remark}[provenance of the repair]\label{inv:wb:uniform-nontorsion-provenance}
\StatementRecord{inv:wb:uniform-nontorsion-provenance}{proved}{Historical events: Claude Fable 5.2 (Anthropic), effort high, 12 November 2026: targeted proof check. GPT-6 Astra (OpenAI), ChatGPT Pro mode, effort Pro, S6\_09\_PAPERS, 17 November 2026: targeted proof check. Also checked in this writing. Pre-continuation writing events: GPT-6 Astra (OpenAI), ChatGPT Pro mode, effort Pro, S6\_10\_PAPERS, 18 November 2026: used and rechecked in this run. Continuation metadata events: GPT-6 Astra (OpenAI), ChatGPT Pro mode, effort Pro, S6\_10\_PAPERS continuation, 19 November 2026: used and rechecked in this run; status-vocabulary reconciliation; metadata reconciliation only; no new mathematical check. \emph{Pre-continuation recorded proof scope.} : The retained period equation and countability strategy agree with the earlier construction scope. The repaired argument now names the actual primitive normal-character theorem, intrinsic cusp specialization and global gamma restriction map. This checks the stated provenance and dependency roles; the mathematical argument is recorded at wb:thm:uniform-nontorsion-audit}{the hypotheses, notation and definitions stated here; the argument and qualifications immediately following}

The period equation \(c\tau+d\mu-b=0\) and the countability strategy are
the candidate's.  The exact normal-character proof above makes the order of
\(\mathcal O_X(S_j)|_{S_j}\) explicit, while the cusp extension in
\cref{inv:A1.0,mapsI:top:thm:p0-integral-specialization} replaces the candidate's
collapse-dependent argument by the independently verified intrinsic
nearby-cycle map.  Identifying the restriction line with the global
\(\gamma^{0,1}\)-summand also removes the candidate's separate continuity
argument for a moving integral lift.
\end{remark}


\Uses{the hypotheses, notation and definitions stated here; the argument and qualifications immediately following}



% Complete working-edition supplement. All source files remain unchanged.
% Reader and writer: GPT-6 Astra (OpenAI, ChatGPT Pro mode), effort Pro,
% 17 November 2026 UTC. Full source and label mapping accompanies this file.

\section{Analytic duality and the exact Picard restriction maps}
\begin{lemma}[identification of the torsion subsheaf]
\label{wbII:lem:torsion-identification}
\StatementRecord{wbII:lem:torsion-identification}{proved}{Historical events: GPT-6 Astra (OpenAI), ChatGPT Pro mode, effort Pro, S6\_09\_PAPERS, 17 November 2026: targeted proof check. Also checked in this writing. Pre-continuation writing events: GPT-6 Astra (OpenAI), ChatGPT Pro mode, effort Pro, S6\_10\_PAPERS, 18 November 2026: targeted proof check. Continuation metadata events: GPT-6 Astra (OpenAI), ChatGPT Pro mode, effort Pro, S6\_10\_PAPERS continuation, 19 November 2026: used and rechecked in this run; status-vocabulary reconciliation; metadata reconciliation only; no new mathematical check. \emph{Pre-continuation recorded proof scope.} The support of a section vanishing on the smooth locus avoids every minimal prime. Prime avoidance supplies an annihilator outside their finite union, hence a nonzerodivisor. Conversely torsion restricts to zero in the locally free differentials on the smooth locus. This proves equality of the exact two sheaves for the reduced Cartier divisor}{the hypotheses, notation and definitions stated here; the explicit constructions and calculations in the proof}

If $S$ is a reduced effective Cartier divisor in a smooth complex manifold,
then
\[
\ker\bigl(\Omega^1_S\longrightarrow j_*\Omega^1_{S_{\rm reg}}\bigr)
  =\operatorname{Tor}(\Omega^1_S).
\]
\end{lemma}



\begin{proof}
\noindent\textbf{Step 1.} Every irreducible component of a reduced hypersurface is generically smooth,
so $S_{\rm reg}$ is dense and its complement is a proper analytic subset.
If a local section of $\Omega^1_S$ vanishes on $S_{\rm reg}$, its support avoids
all minimal primes of the reduced local ring.  Prime avoidance therefore gives
an element of its annihilator outside every minimal prime; that element is a
nonzerodivisor, so the section is torsion.  Conversely,
$\Omega^1_S|_{S_{\rm reg}}$ is locally free, hence has no torsion, and every
torsion section restricts to zero there.
\end{proof}



\subsubsection*{Local proof of the conormal filtration}


Let $S$ be a reduced effective Cartier divisor in a smooth complex
threefold $X$.  There are natural short exact sequences
\[
0\longrightarrow N^*_{S/X}\longrightarrow\mathscr K_S
\longrightarrow\mathscr T_S\longrightarrow0,
\]
\[
0\longrightarrow\mathscr K_S\longrightarrow\Omega_X^1|_S
\longrightarrow\widetilde\Omega_S^1\longrightarrow0.
\]



\begin{proof}
The conormal sequence is right exact.  It is also left exact here.  Locally
write $S=(g=0)$ with $g$ reduced.  If $h\,dg|_S=0$, then $h$ vanishes at the
generic smooth point of every irreducible component, because $dg$ is nonzero
there.  The ring $\mathcal O_S$ is reduced, so the intersection of its
minimal primes is zero; hence $h=0$.  Thus
\[
0\to N^*_{S/X}\to\Omega_X^1|_S\to\Omega_S^1\to0
\]
is exact.  By definition, $\mathscr K_S$ is the inverse image of
$\mathscr T_S$ under the last surjection.  Restriction to that inverse image
is onto $\mathscr T_S$ and has kernel $N^*_{S/X}$, proving the first
sequence.  Composing with the quotient by $\mathscr T_S$ proves the second.
\end{proof}

\begin{proof}[The minimal-prime formulation of conormal injectivity]
Right exactness is the conormal sequence for a closed immersion.  We check
the left end.  Work at a point of $S$ and let $R$ be the regular local ring
of $X$, with reduced local equation $g$ for $S$.  The map at the left is
\[
 R/(g)\longrightarrow (\Omega^1_R/g\Omega^1_R),
 \qquad \bar h\longmapsto h\,dg.
 \tag{W.145}\label{wbII:eq:CN:2.2}
\]
Let $\mathfrak p$ be a minimal prime of $(g)$.  Since $R$ is regular and
has characteristic zero, the height-one local ring $R_{\mathfrak p}$ is a
discrete valuation ring.  Reducedness of $(g)$ gives
$\operatorname{ord}_{\mathfrak p}(g)=1$, so $g$ is a unit times a
uniformizer there; in characteristic zero the differential of that
uniformizer is nonzero modulo $\mathfrak p$.  Hence $dg$ is nonzero in
$\Omega^1_{R_{\mathfrak p}}\otimes k(\mathfrak p)$, so \eqref{wbII:eq:CN:2.2} sends $\bar h$
to zero only if $h\in\mathfrak p$.  This holds for every minimal prime.
Because $R/(g)$ is reduced, the intersection of its minimal primes is zero;
therefore $\bar h=0$.  Thus the conormal sequence in \cref{cdp:st:A.1.1} is exact.
\end{proof}

\subsubsection*{Cohomology proof of the exact section criterion}


Let $A$ be a line bundle on $S$.  If
\[
H^0(S,\widetilde\Omega_S^1\otimes A)=0,
\]
then
\[
H^0(S,\Omega_X^1|_S\otimes A)
=H^0(S,\mathscr K_S\otimes A).
\]
If also $H^0(S,N^*_{S/X}\otimes A)=0$, then this group is canonically
\[
\ker\!\left[
H^0(S,\mathscr T_S\otimes A)
\xrightarrow{\delta_A}H^1(S,N^*_{S/X}\otimes A)
\right].
\]
Hence the two vanishings printed in the p.~692 reduction imply the desired
ambient vanishing if and only if $\delta_A$ is injective.



\begin{proof}
Tensoring the second sequence of
\cref{cdp:st:A.1.1} by the locally free sheaf $A$
preserves exactness; its global-section sequence gives the first equality.
Tensor the first sequence and take global sections.  If the conormal
$H^0$ vanishes, exactness identifies the middle group with the displayed
kernel.  No condition on the torsion-free quotient controls that kernel.
\end{proof}

\subsection{Local algebra of a reduced normal-crossings fibre}

Let $g:\mathcal X\to\Delta$ be a holomorphic map from a smooth complex
threefold to a disc, and put $S=g^{-1}(0)$, assumed nonempty.  Assume $S$ is reduced and has
normal crossings: at every point there are coordinates $(x,y,z)$ in which
\[
g=ux,\qquad g=uxy,\qquad\text{or}\qquad g=uxyz
\tag{W.3}\label{wbII:eq:W:5.3}
\]
for a unit $u$.  Let $\eta:\widetilde S\to S$ be the normalization,
$D=S_{\mathrm{sing,red}}$, and $\widetilde D=\eta^{-1}(D)_{\mathrm{red}}$.

\subsubsection*{Coordinate proof for normal-crossing fibres}

\ In the setting above, $N^*_{S/\mathcal X}\simeq\mathcal O_S$,
and its generator maps to $dg|_S$.  Multiplication by $dg|_S$ is injective,
so
\[
0\longrightarrow\mathcal O_S
\xrightarrow{\,h\mapsto h\,dg\,}
\Omega^1_{\mathcal X}|_S\longrightarrow\Omega^1_S\longrightarrow0
\tag{W.4}\label{wbII:eq:W:5.4}
\]
is exact.



\begin{proof}
The fibre ideal is the principal ideal $(g)$ and is reduced, so
$(g)/(g^2)\simeq\mathcal O_S$.  Multiplication by a unit does not affect
injectivity, and the term $f\,du$ vanishes after restriction to $f=0$.
It therefore suffices to use $f=x,xy,xyz$.

For $f=xy$, if a lift $\widetilde h$ satisfies
$\widetilde h(y\,dx+x\,dy)=0$ modulo $(xy)$, then
$\widetilde h y,\widetilde h x\in(xy)$.  Hence
$\widetilde h\in(x)\cap(y)=(xy)$ in the regular local ring, so $h=0$ on
$S$.  For $f=xyz$, the three coefficient conditions say
$\widetilde h yz,\widetilde h xz,\widetilde h xy\in(xyz)$; hence
$\widetilde h\in(x)\cap(y)\cap(z)=(xyz)$.  The case $f=x$ is immediate.
Thus the conormal map is injective in every local model, proving \eqref{wbII:eq:W:5.4}.
\end{proof}

\subsubsection*{Coordinate proof of the normalization--conductor square}


\ The sequence
\[
0\longrightarrow\mathcal O_S
\xrightarrow{\alpha}
\eta_*\mathcal O_{\widetilde S}\oplus\mathcal O_D
\xrightarrow{\beta}
\eta_*\mathcal O_{\widetilde D}
\longrightarrow0,
\tag{W.5}\label{wbII:eq:W:5.5}
\]
\[
\alpha(h)=(\eta^*h,h|_D),\qquad
\beta(\widetilde h,k)=
\widetilde h|_{\widetilde D}-\eta^*k,
\]
is exact.



\begin{proof}
The assertion is local.  Away from $D$, $\eta$ is an isomorphism.
At a double point $S=\{xy=0\}$, a member of $\ker\beta$ is
\[
\bigl((h_1(y,z),h_2(x,z)),k(z)\bigr),
\qquad
h_1(0,z)=h_2(0,z)=k(z).
\]
The function
\[
h=h_1(y,z)+h_2(x,z)-h_1(0,z)
\]
on $\{xy=0\}$ restricts to $h_1,h_2$ on the two branches and to $k$ on
their intersection.  This proves middle exactness in the double model.

At a triple point $S=\{xyz=0\}$, write the branch functions as
$h_1(y,z)$, $h_2(x,z)$, $h_3(x,y)$.  Membership in $\ker\beta$ is exactly
\[
h_1(0,z)=h_2(0,z),\quad
h_1(y,0)=h_3(0,y),\quad
h_2(x,0)=h_3(x,0).
\]
Then
\[
\begin{split}
h={}&h_1(y,z)+h_2(x,z)+h_3(x,y)\\
&-h_1(0,z)-h_1(y,0)-h_2(x,0)+h_1(0,0)
\end{split}
\]
restricts to the prescribed function on every branch.  Injectivity of
$\alpha$ follows because $S$ is reduced and $\eta$ is an isomorphism over a
dense open set.  Surjectivity of $\beta$ follows locally by lifting a
function on the branch conductor to the smooth branch and taking $k=0$.
This proves exactness in every local model.
\end{proof}

\begin{theorem}[conductor-supported ambient differential]
\label{wbII:prop:conductor}
\StatementRecord{wbII:prop:conductor}{proved}{Historical events: GPT-6 Astra (OpenAI), ChatGPT Pro mode, effort Pro, S6\_09\_PAPERS, 17 November 2026: targeted proof check. Also checked in this writing. Pre-continuation writing events: GPT-6 Astra (OpenAI), ChatGPT Pro mode, effort Pro, S6\_10\_PAPERS, 18 November 2026: targeted proof check. Continuation metadata events: GPT-6 Astra (OpenAI), ChatGPT Pro mode, effort Pro, S6\_10\_PAPERS continuation, 19 November 2026: used and rechecked in this run; status-vocabulary reconciliation; metadata reconciliation only; no new mathematical check. \emph{Pre-continuation recorded proof scope.} In both double and triple charts the coefficients of dg lie in the reduced conductor ideal. The explicit compatible pair (eta-star(dg) times e,0) descends through the conductor square to a unique nonzero ambient differential. Exactness of the conormal sequence proves that its intrinsic image vanishes precisely when e descends to a nowhere-zero section. For compact connected S, constants on the normalization components and the actual common fibre identifications force every nonzero section of A to trivialize A}{the hypotheses, notation and definitions stated here; \cref{wbII:eq:W:5.3}; \cref{wbII:eq:W:5.1}; \cref{wbII:eq:W:5.5}; \cref{wbII:eq:W:5.4}}
\label{wbII:prop:conductor-differential}
\ In the setting \eqref{wbII:eq:W:5.3}, let $A$ be a line bundle on $S$ with
$\eta^*A\simeq\mathcal O_{\widetilde S}$, and choose a nowhere-vanishing
section $e$ of $\eta^*A$.  There is a unique nonzero
\[
s\in H^0(S,\Omega^1_{\mathcal X}|_S\otimes A)
\tag{W.6}\label{wbII:eq:W:5.6}
\]
such that
\[
\eta^*s=\eta^*(dg|_S)\otimes e,\qquad s|_D=0.
\tag{W.7}\label{wbII:eq:W:5.7}
\]
Its image $\rho(s)$ in $\Omega^1_S\otimes A$ is supported on $D$ and
therefore torsion; its image in
$\widetilde\Omega^1_S\otimes A$ is zero.  Moreover:
\begin{enumerate}[label=(\alph*)]
\item $\rho(s)=0$ if and only if $e=\eta^*\sigma$ for a nowhere-vanishing
$\sigma\in H^0(S,A)$;
\item if $S$ is compact and connected and $A\not\simeq\mathcal O_S$, then
$H^0(S,A)=0$ and $s$ does not come from the conormal term in \eqref{wbII:eq:W:5.1}.
\end{enumerate}
\end{theorem}



\begin{proof}
\noindent\textbf{Step 1.} In the double model, $d(xy)=y\,dx+x\,dy$ has coefficients in the reduced
ideal $(x,y)$ of $D$.  In the triple model,
\[
d(xyz)=yz\,dx+xz\,dy+xy\,dz
\]
has coefficients in $(xy,xz,yz)$, the reduced ideal of the union of the
three double curves.  Unit factors in \eqref{wbII:eq:W:5.3} do not change this after
restriction to $S$.  Hence $\eta^*(dg|_S)$ restricts to zero on
$\widetilde D$.  Tensor \eqref{wbII:eq:W:5.5} by the locally free sheaf
$E=\Omega^1_{\mathcal X}|_S\otimes A$.  The compatible pair
\[
\bigl(\eta^*(dg|_S)\otimes e,0\bigr)
\]
therefore descends uniquely to $s$.  On the smooth locus of $S$, $dg|_S$
is a nonzero conormal vector and $e$ is nowhere zero, so $s\ne0$.

\noindent\textbf{Step 2.} On $S\setminus D$, the normalization is an isomorphism and
$s=dg|_S\otimes e$.  By \eqref{wbII:eq:W:5.4} its intrinsic image vanishes there.  Thus
$\rho(s)$ is supported on $D$ and lies in the torsion subsheaf; its image
modulo torsion is zero.

\noindent\textbf{Step 3.} If $\rho(s)=0$, exactness of \eqref{wbII:eq:W:5.4} after tensoring by $A$ gives
$s=dg|_S\otimes\sigma$ for a section $\sigma$ of $A$.  Pulling back and
comparing with \eqref{wbII:eq:W:5.7} on the dense open set where $dg\ne0$ gives
$\eta^*\sigma=e$ there, hence everywhere by holomorphic continuation.
Since $e$ never vanishes and $\eta$ is surjective, $\sigma$ never vanishes
and trivializes $A$.  The converse is immediate.

\noindent\textbf{Step 4.} Finally suppose $S$ is compact and connected and $\sigma\in H^0(S,A)$.  On each compact
connected component of $\widetilde S$, write
$\eta^*\sigma=c_i e$; the holomorphic function $c_i$ is constant.  If one
$c_i$ is zero and another is nonzero, connectedness of the finite union of
irreducible components of $S$ provides two such components meeting at a
point.  The value of $\sigma$ at that point would be both zero and nonzero,
a contradiction.  Hence either every $c_i=0$, so $\sigma=0$, or every
$c_i\ne0$, so $\sigma$ is nowhere zero and $A$ is trivial.  When $A$ is
nontrivial only the first case occurs.  Since
$N^*_{S/\mathcal X}\simeq\mathcal O_S$, the conormal global-section space is
$H^0(S,A)=0$, while $s\ne0$.
\end{proof}



\begin{corollary}
\label{wbII:cor:normal-crossing-kernel}
\StatementRecord{wbII:cor:normal-crossing-kernel}{proved}{Historical events: GPT-6 Astra (OpenAI), ChatGPT Pro mode, effort Pro, S6\_09\_PAPERS, 17 November 2026: targeted proof check. Also checked in this writing. Pre-continuation writing events: GPT-6 Astra (OpenAI), ChatGPT Pro mode, effort Pro, S6\_10\_PAPERS, 18 November 2026: targeted proof check. Continuation metadata events: GPT-6 Astra (OpenAI), ChatGPT Pro mode, effort Pro, S6\_10\_PAPERS continuation, 19 November 2026: used and rechecked in this run; status-vocabulary reconciliation; metadata reconciliation only; no new mathematical check. \emph{Pre-continuation recorded proof scope.} The conductor section is a nonzero element of the actual kernel of the torsion-free quotient. The two assumed endpoint vanishings identify that section space with the kernel of the specified connecting morphism from the cotangent filtration, so that morphism has nonzero kernel}{the hypotheses, notation and definitions stated here; \cref{wbII:prop:conductor}; \cref{cdp:st:A.1.1}}

If, in addition, $H^0(S,N^*_{S/X}\otimes A)=0$ and
$H^0(S,\widetilde\Omega^1_S\otimes A)=0$, then the section in
Proposition~\ref{wbII:prop:conductor} proves that $\delta_A$ in
the section-kernel isomorphism of \cref{cdp:st:A.1.1} is not injective.
\end{corollary}



\begin{proof}
\noindent\textbf{Step 1.} The section is nonzero and maps to zero in the quotient
$\widetilde\Omega^1_S=\Omega^1_S/\mathscr T_S$, so it
lies in the left-hand side of the section-kernel isomorphism of \cref{cdp:st:A.1.1}.  Proposition~\ref{cdp:st:A.1.1}
identifies that group with $\ker\delta_A$, which is therefore nonzero.
\end{proof}



\subsection{A proper counterexample to the printed inference}

\subsubsection*{Conductor proof of the nodal endpoint counterexample}

\ The p.~692 implication described after \eqref{wbII:eq:W:5.2} is false as a formal
sheaf-theoretic statement, even for a proper family with smooth total space
and an irreducible reduced normal-crossings central fibre.



\begin{proof}
Let $\Delta$ be a sufficiently small disc with coordinate $t$ and put
\[
\mathcal C=
\{([x:y:z],t)\in\mathbb P^2\times\Delta:
y^2z=x^2(x+z)+tz^3\}.
\tag{W.8}\label{wbII:eq:W:5.8}
\]
This is proper over $\Delta$.  Its total space is smooth.  On $z\ne0$, the
partial derivative of the defining equation with respect to $t$ is
$-z^3\ne0$.  On $z=0$, the equation forces $x=0$; then $y\ne0$ and the
$z$-derivative is $y^2\ne0$.  The central fibre $C_0$ is the irreducible
nodal cubic $y^2z=x^2(x+z)$, with node $[0:0:1]$.

Set
\[
\mathcal X=\mathcal C\times\mathbb P^1,\qquad
S=C_0\times\mathbb P^1,
\]
and let $g:\mathcal X\to\Delta$ be projection.  Then $\mathcal X$ is a
smooth threefold, $g$ is proper, and $S$ is an irreducible reduced
normal-crossings surface.  Its normalization is
$\widetilde S\simeq\mathbb P^1\times\mathbb P^1$.

Let $p,q$ be the two inverse images of the node in the normalization of
$C_0$.  Fix $\lambda\in\mathbb C^*$ which is not a root of unity (for
example, $\lambda=2$).  Glue the fibres of
the trivial line bundle on the normalization at $p$ and $q$ by the scalar
$\lambda$.  This defines a line bundle
$A_\lambda\in\operatorname{Pic}^0(C_0)$ whose pullback is trivial.  It has
infinite order: its $n$th tensor power has gluing scalar $\lambda^n$, which
is never one.  It is not trivial: a global section pulls back to a constant
$c$ satisfying
$c=\lambda c$, so
\[
H^0(C_0,A_\lambda)=0.
\tag{W.9}\label{wbII:eq:W:5.9}
\]
Put $A=\operatorname{pr}_1^*A_\lambda$ on $S$.  Then
$\eta^*A\simeq\mathcal O_{\widetilde S}$ and $H^0(S,A)=0$.
Because $S$ is the full fibre,
$N^*_{S/\mathcal X}\simeq\mathcal O_S$, so the conormal-twist global
sections vanish.

The natural map
$\widetilde\Omega^1_S\to\eta_*\Omega^1_{\widetilde S}$ is injective:
its kernel is supported where normalization is not an isomorphism, while
the source is torsion-free by definition.  Projection formula and
$\eta^*A\simeq\mathcal O_{\widetilde S}$ therefore give
\[
H^0(S,\widetilde\Omega^1_S\otimes A)
\hookrightarrow
H^0(\widetilde S,\Omega^1_{\widetilde S})=0,
\tag{W.10}\label{wbII:eq:W:5.10}
\]
the last equality following from
$\widetilde S=\mathbb P^1\times\mathbb P^1$.
Nevertheless \cref{wbII:prop:conductor} produces
\[
0\ne s\in H^0(S,\Omega^1_{\mathcal X}|_S\otimes A).
\]
Thus both vanishings invoked in the printed reduction hold while the desired
ambient vanishing fails.  This calculation chooses $A$ directly on $S$ and disproves the
claimed \emph{formal implication between the three sheaf-cohomology groups}
for a non-torsion line bundle. The explicit extending choice of $A$ is
constructed in the next subsection, which also treats the total-space
line-bundle condition on p.~692.  The family \eqref{wbII:eq:W:5.8} also does not satisfy all hypotheses of CDP
Theorem~2.2, so this is not a counterexample to that theorem.
\end{proof}

\subsection{The explicit extending choice in the nodal family}
Continue with
\[
\mathcal C=\{y^2z=x^2(x+z)+tz^3\}\subset\PP^2\times\Delta,
\qquad \mathcal X=\mathcal C\times\PP^1,
\qquad S=C_0\times\PP^1.
\]
Write $n=[0:0:1]$, let $D=\{n\}\times\PP^1$, and let
$\nu:\PP^1\to C_0$ denote the normalization described next.
\subsubsection*{Explicit normalization and crossing coordinates}

The projection $g$ is proper, $\mathcal X$ is a smooth complex threefold,
  and $S$ is a reduced irreducible normal-crossing (self-crossing) surface whose normalization
is $\PP^1\times\PP^1$.  Its singular locus is
\(D=\{[0:0:1]\}\times\PP^1\), and locally along $D$ the pair $(\mathcal X,S)$
has the model $t=uv$; away from $D$ the fibre is smooth and has the usual
local model $(\C^3,\{u=0\})$.


\begin{proof}
Properness follows because $\mathcal C$ is closed in the projective
projection $\PP^2\times\Delta\to\Delta$.  Write the defining equation as
\[
 F=y^2z-x^2(x+z)-tz^3.
\]
On the chart $z\ne0$, $\partial F/\partial t=-z^3\ne0$, so $\mathcal C$
is smooth there.  On $z=0$, the equation forces $x=0$ and $y\ne0$; at such
points $\partial F/\partial z=y^2\ne0$.  Hence $\mathcal C$, and therefore
$\mathcal X$, is smooth.  The only finite discriminant values of the fibres
are $t=0$ and $t=-4/27$; after shrinking $\Delta$ to $|t|<4/27$, every
noncentral fibre is smooth.

The central curve is
\(C_0:\ y^2z=x^2(x+z)\).  Its only singular point is
\(n=[0:0:1]\): on $z=1$ the equations $F=F_x=F_y=0$ give $(x,y)=(0,0)$,
and the point at infinity $[0:1:0]$ is smooth because $F_z=1$ there
(Euler's identity also gives the remaining projective singularity equation).
The parametrization
\[
 \nu:\PP^1\longrightarrow C_0,\qquad
 [u:v]\longmapsto
 [(u^2-v^2)v:\,u(u^2-v^2):\,v^3]
 \tag{W.160}\label{wbII:eq:CN:4.3}
\]
maps the two points $p=[1:1]$, $q=[-1:1]$ to $n$.  On the chart $v\ne0$, set
$r=u/v$; after scaling, its image is
$[r^2-1:r(r^2-1):1]$.  Conversely, for an affine point of $C_0$ with
$x\ne0$, the equation gives $r=y/x$ and $x=r^2-1$, so it is hit by
$[r:1]$.  The only affine point with $x=0$ is the node $(0,0)$, hit by
$[1:1]$ and $[-1:1]$, and the point at infinity $[0:1:0]$ is hit by
$[1:0]$.  Thus $\nu(\PP^1)=C_0$; as a nonconstant morphism to a projective
curve it is finite and generically one-to-one, hence the normalization.  In
particular $C_0$ is irreducible and rational.  The product normalization
is $\eta=\nu\times\id_{\PP^1}:\PP^1\times\PP^1\to S$.

At $n$ the affine equation is
\(t=y^2-x^2-x^3\).  On a sufficiently small neighbourhood choose the
holomorphic square root $h(x)=\sqrt{1+x}$ and set
\(u=y-xh(x)\), \(v=y+xh(x)\).  Then $uv=t$, and the map
$(x,y)\mapsto(u,v)$ has invertible $2\times2$ Jacobian at $(0,0)$, so $u,v$
are local coordinates on the surface $\mathcal C$.  If $s$ is a local
coordinate on $\PP^1$, then $(u,v,s)$ are coordinates on $\mathcal X$ and
$t=uv$.  Hence the pair has the asserted normal-crossing (self-crossing)
local model, with singular locus $D$ as stated.
\end{proof}

\begin{lemma}[the non-torsion gluing bundle]
\label{wbII:lem:gluing-bundle}
\StatementRecord{wbII:lem:gluing-bundle}{proved}{Historical events: GPT-6 Astra (OpenAI), ChatGPT Pro mode, effort Pro, S6\_09\_PAPERS, 17 November 2026: targeted proof check. Also checked in this writing. Pre-continuation writing events: GPT-6 Astra (OpenAI), ChatGPT Pro mode, effort Pro, S6\_10\_PAPERS, 18 November 2026: targeted proof check. Continuation metadata events: GPT-6 Astra (OpenAI), ChatGPT Pro mode, effort Pro, S6\_10\_PAPERS continuation, 19 November 2026: used and rechecked in this run; status-vocabulary reconciliation; metadata reconciliation only; no new mathematical check. \emph{Pre-continuation recorded proof scope.} The line bundle is constructed by the literal fibre identification at the two branches. Its normalization pullback is trivial; every global section is constant upstairs and must obey c=lambda c. Tensor power replaces lambda by lambda to that power, proving all stated nonvanishing and non-torsion conclusions}{the hypotheses, notation and definitions stated here; the explicit constructions and calculations in the proof}

Fix $\lambda\in\C^*$ that is not a root of unity.  Gluing the trivial line
bundle on $\PP^1$ at $p$ and $q$ by the scalar $\lambda$ defines a line bundle
\(A_\lambda\in\Pic(C_0)\) with
\[
 \nu^*A_\lambda\simeq\cO_{\PP^1},\qquad
 H^0(C_0,A_\lambda)=0,\qquad
 A_\lambda^{\otimes m}\not\simeq\cO_{C_0}\quad(m\ge1).
 \tag{W.161}\label{wbII:eq:CN:4.4}
\]
\end{lemma}



\begin{proof}
\noindent\textbf{Step 1.} For a nodal curve, a line bundle with trivial pullback to the normalization
is obtained by identifying the two fibres over the branches of the node.
Use the identification whose scalar is $\lambda$.  A global section pulls
back to a constant $c$ on the projective line and must satisfy
$c=\lambda c$; since $\lambda\ne1$, $c=0$.  The $m$th tensor power has
gluing scalar $\lambda^m$, which is never $1$ because $\lambda$ is not a
root of unity.  This proves \eqref{wbII:eq:CN:4.4}.
\end{proof}



\begin{lemma}[an extending choice of gluing bundle]
\label{wbII:lem:extending-gluing}
\StatementRecord{wbII:lem:extending-gluing}{proved}{Historical events: GPT-6 Astra (OpenAI), ChatGPT Pro mode, effort Pro, S6\_09\_PAPERS, 17 November 2026: targeted proof check. Also checked in this writing. Pre-continuation writing events: GPT-6 Astra (OpenAI), ChatGPT Pro mode, effort Pro, S6\_10\_PAPERS, 18 November 2026: targeted proof check. Continuation metadata events: GPT-6 Astra (OpenAI), ChatGPT Pro mode, effort Pro, S6\_10\_PAPERS continuation, 19 November 2026: used and rechecked in this run; status-vocabulary reconciliation; metadata reconciliation only; no new mathematical check. \emph{Pre-continuation recorded proof scope.} The two sections satisfy the actual cubic equation and stay disjoint. The central normalization sends the divisor to [2]-[infinity], trivialized by r-2. Evaluation on the two node branches gives minus1 and minus3, hence multiplier one third or its inverse. The divisor line bundle on the smooth total space therefore realizes the required extending non-torsion twist}{the hypotheses, notation and definitions stated here; \cref{cdp:st:B1}; \cref{wbII:lem:gluing-bundle}}

After shrinking $\Delta$, the sections
\[
 O(t)=[0:1:0],\qquad
 P(t)=[3:\sqrt{36+t}:1]
\tag{W.162}\label{wbII:eq:CN:4.5a}
\]
of $\mathcal C\to\Delta$ are disjoint and holomorphic.  If
$B=\cO_{\mathcal C}(P-O)$, then the line bundle
$\widetilde A=\operatorname{pr}_1^*B$ on $\mathcal X=\mathcal C\times\PP^1$
restricts to $\operatorname{pr}_1^*A_{\lambda}$ on $S$ for
$\lambda=1/3$ (or its inverse, according to the gluing convention).  In
particular, the non-torsion bundle used in Theorem~\ref{cdp:st:B1} can
be chosen to extend to the displayed smooth total space.
\end{lemma}



\begin{proof}
\noindent\textbf{Step 1.} The equation of $\mathcal C$ is satisfied by both sections; the square root
is holomorphic for $|t|$ small, and the two sections are disjoint because
$O$ has $z=0$ whereas $P$ has $z=1$.  On the normalization of $C_0$, write
$r=u/v$.  Then $O(0)$ corresponds to $r=\infty$, $P(0)$ to $r=2$, and the
two branches over the node correspond to $r=1$ and $r=-1$.  Hence
\[
 \nu^*(B|_{C_0})\simeq\cO_{\PP^1}([2]-[\infty]).
\]
The rational function $f(r)=r-2$ has precisely this divisor and trivializes
the pullback.  Its values at the two branches are $f(1)=-1$ and
$f(-1)=-3$, so the induced gluing multiplier is $f(1)/f(-1)=1/3$,
up to inversion.  This is not a root of unity, and the lemma follows from
the gluing description in Lemma~\ref{wbII:lem:gluing-bundle}.
\end{proof}



\subsubsection*{The local torsion calculation in the nodal example}

Let $A=\operatorname{pr}_1^*A_\lambda$ on
$S=C_0\times\PP^1$, with $\lambda$ as in Lemma~\ref{wbII:lem:gluing-bundle}.
Then
\[
 H^0(S,N^*_{S/\mathcal X}\otimes A)=0,
 \qquad
 H^0(S,\widetilde\Omega^1_S\otimes A)=0,
 \tag{W.163}\label{wbII:eq:CN:4.5}
\]
but
\[
 H^0(S,\Omega^1_{\mathcal X}|_S\otimes A)\ne0.
 \tag{W.164}\label{wbII:eq:CN:4.6}
\]
In particular, the two endpoint vanishings in \eqref{wbII:eq:CN:4.5} do not imply the
vanishing of \eqref{wbII:eq:CN:4.6}, even for a proper morphism over a disc with smooth total
space and an irreducible reduced normal-crossing (self-crossing) central
fibre.


\begin{proof}
Because $S$ is the principal fibre $(t=0)$, its conormal line is trivial:
\(N^*_{S/\mathcal X}\simeq\cO_S\).  The projection
\(\operatorname{pr}_1:S\to C_0\) has
\(\operatorname{pr}_{1*}\cO_S\simeq\cO_{C_0}\), so \eqref{wbII:eq:CN:4.4} gives
\[
 H^0(S,A)\simeq H^0(C_0,A_\lambda)=0.
 \tag{W.165}\label{wbII:eq:CN:4.7}
\]
This proves the first vanishing in \eqref{wbII:eq:CN:4.5}.

We next check the second.  At a normal-crossing (self-crossing) point with local equation
\(uv=0\), the pullback map
\(\Omega^1_S\to\eta_*\Omega^1_{\widetilde S}\) has kernel exactly
\(\mathscr T_S\).  Indeed, in
\(R=\C\{u,v,w\}/(uv)\) the relation is
\(v\,du+u\,dv=0\).  If
\(a\,du+b\,dv+c\,dw\) restricts to zero on both branches, then
\(a\in(v)\), \(b\in(u)\), and
\(c\in(u)\cap(v)=(uv)=0\).  Modulo the relation the class is therefore a
multiple of $\tau=v\,du=-u\,dv$.  This class is nonzero: if
\(v\,du=h(v\,du+u\,dv)\), then coefficient comparison gives
  \((1-h)v=0\) and $hu=0$.  Since
\(\operatorname{Ann}_R(v)=(u)\) and \(\operatorname{Ann}_R(u)=(v)\),
\(1-h\in(u)\) and \(h\in(v)\), which would imply
\(1=(1-h)+h\in(u)+(v)\), impossible because $(u)+(v)$ is proper.
Moreover $(u+v)\tau=0$.  Since $u+v$ is a non-zero-divisor in
$R$ (the associated primes are $(u)$ and $(v)$), $\tau$ is genuine
torsion.
Conversely, each summand of the normalization module is a module over one of
the domains $R/(u)$ and $R/(v)$.  A non-zero-divisor of the reduced ring $R$
avoids both minimal primes and therefore acts injectively on both summands;
the target is torsion-free.  Thus the kernel of the normalization map is
exactly the torsion subsheaf, and there is an injection
\[
 \widetilde\Omega^1_S\hookrightarrow\eta_*\Omega^1_{\PP^1\times\PP^1}.
 \tag{W.166}\label{wbII:eq:CN:4.8}
\]
Since $\eta^*A\simeq\cO_{\PP^1\times\PP^1}$, the projection formula and \eqref{wbII:eq:CN:4.8}
give
\[
 H^0(S,\widetilde\Omega^1_S\otimes A)
 \hookrightarrow
 H^0(\PP^1\times\PP^1,\Omega^1_{\PP^1\times\PP^1})=0.
 \tag{W.167}\label{wbII:eq:CN:4.9}
\]
The last equality follows from
\(\Omega^1_{\PP^1\times\PP^1}\simeq
\operatorname{pr}_1^*\Omega^1_{\PP^1}\oplus
\operatorname{pr}_2^*\Omega^1_{\PP^1}\) and
\(H^0(\PP^1,\Omega^1_{\PP^1})=0\).

Finally, along the conductor $D$ the pullback of $dg|_S$ is zero.  Locally
$g=h\,uv$ with $h$ a unit, so after restriction to $S$ one has
$i^*(dg)=h(v\,du+u\,dv)$; both coefficients vanish on the preimage of $D$.
Proposition~\ref{wbII:prop:conductor}
therefore gives a nonzero section in \eqref{wbII:eq:CN:4.6}.  This proves the theorem.
\end{proof}

\paragraph{what the theorem does and does not refute}

\label{wbII:rem:scope}
The bundle $A$ is non-torsion on $S$: restriction to
\(C_0\times\{r\}\) would otherwise make $A_\lambda$ torsion.  Thus the
example meets the non-torsion condition attached to the p.~692 reduction.
In fact Lemma~\ref{wbII:lem:extending-gluing} gives an extending choice of $A$ on
the displayed smooth total space.  If $\widetilde L:=\omega_{\mathcal X}\otimes
\widetilde A^{-1}$, then
$(\widetilde L^\vee\otimes\omega_{\mathcal X})|_S\simeq A$, so the example
also has the Serre-dual twist appearing in the CDP argument.  This family
nevertheless does not have
the global topological and algebraic-dimension hypotheses of CDP20
Theorem~2.2.  The conclusion is exactly the formal one stated in
Theorem~\ref{cdp:st:B1}: the two displayed endpoint vanishings do not
control the torsion-supported lifting obstruction $\delta_A$.

\Uses{the hypotheses, notation and definitions stated here; \cref{wbII:lem:extending-gluing}; \cref{cdp:st:B1}}


\subsection{Cartier-divisor duality without a normality hypothesis}

The next two results are explicit re-derivations of classical analytic
Cartier duality.  The trace duality and complete-intersection identification
are supplied by \citet[Theorem~2, pp.~82--84, and Proposition~2, p.~91]{RR70push}.
The proofs retain the full scheme structure of the fibre and state
separately where compactness is needed.

\begin{theorem}[complete Cartier-fibre duality]
\label{wbII:thm:cartier-duality}
\StatementRecord{wbII:thm:cartier-duality}{proved}{Historical events: GPT-6 Astra (OpenAI), ChatGPT Pro mode, effort Pro, S6\_09\_PAPERS, 17 November 2026: targeted proof check. Also checked in this writing. Pre-continuation writing events: GPT-6 Astra (OpenAI), ChatGPT Pro mode, effort Pro, S6\_10\_PAPERS, 18 November 2026: targeted proof check; used and rechecked in this run. Continuation metadata events: GPT-6 Astra (OpenAI), ChatGPT Pro mode, effort Pro, S6\_10\_PAPERS continuation, 19 November 2026: used and rechecked in this run; status-vocabulary reconciliation; metadata reconciliation only; no new mathematical check. \emph{Pre-continuation recorded proof scope.} Applying sheaf Hom to the actual length-one Cartier resolution gives the adjunction line as Ext1 and zero other sheaf Ext groups. The analytic dualizing complex is that line shifted by two. The Ramis-Ruget trace therefore identifies Hq dual with H(2-q) of the dual twist, with the stated contravariance; compactness supplies finite dimensionality and negative Ext degrees give vanishing above two. Ramis-Ruget Theorem2 and its complete proof pp82-84 and the complete-intersection calculation p91 were separately opened}{the hypotheses, notation and definitions stated here; \cite[Proposition~2 and pp.~90--91]{RR70push}; \cite[Theorem~2, pp.~82--83]{RR70push}}

\ Let $X$ be a compact smooth complex threefold,
$i:W\hookrightarrow X$ an effective Cartier divisor, and $E$ a holomorphic
vector bundle.  Put $F=E|_W$ and
\[
\omega_W=(\omega_X\otimes\mathcal O_X(W))|_W.
\]
For $q=0,1,2$ there are natural perfect pairings
\[
H^q(W,F)^\vee
\simeq H^{2-q}(W,F^\vee\otimes\omega_W).
\tag{W.11}\label{wbII:eq:W:5.11}
\]
Moreover $H^q(W,F)=0$ for $q>2$.  No normality or reducedness hypothesis on
$W$ is used.
\end{theorem}



\begin{proof}
\noindent\textbf{Step 1.} Let \(i:W\hookrightarrow X\) be the given Cartier inclusion.  The
length-one locally free resolution
\[
 0\longrightarrow\mathcal O_X(-W)
 \xrightarrow{\;\cdot g\;}\mathcal O_X
 \longrightarrow i_*\mathcal O_W\longrightarrow0
\]
is locally defined by a nonzerodivisor \(g\); changing \(g\) by a unit
gives the usual transition functions for \(\mathcal O_X(W)\).
Applying sheaf Hom into \(\omega_X\), the degree-zero map is
injective and its cokernel is
\[
 \mathcal E xt^1_X(i_*\mathcal O_W,\omega_X)
 =i_*\bigl((\omega_X\otimes\mathcal O_X(W))|_W\bigr)
 =i_*\omega_W.
\]
The other sheaf Ext groups vanish.  Equivalently, \(W\) is a
two-dimensional Cohen--Macaulay Gorenstein complex space and its
analytic dualizing complex is
\[
 K_W^\bullet\simeq\omega_W[2].
\]
This is the hypersurface calculation in
\cite[Proposition~2 and pp.~90--91]{RR70push}; no normality or
reducedness of the Cartier divisor is used.

\noindent\textbf{Step 2.} For the locally free sheaf \(F=E|_W\), the analytic duality trace of
\cite[Theorem~2, pp.~82--83]{RR70push} supplies the actual
functorial pairing, and identifies the dual of \(H_c^q(W,F)\) with
\[
 \operatorname{Ext}^{-q}_W(F,K_W^\bullet)
 =\operatorname{Ext}^{2-q}_W(F,\omega_W)
 =H^{2-q}(W,F^\vee\otimes\omega_W).
\]
The last equality follows because the positive sheaf Ext groups
of a locally free first argument vanish.  Since \(W\) is compact,
\(H_c^q(W,F)=H^q(W,F)\); coherent cohomology is finite-dimensional,
so this pairing is perfect in the ordinary algebraic sense.
It is this trace pairing, rather than an unspecified arrow between
two matching five-term sequences, that constructs the isomorphism
\eqref{wbII:eq:W:5.11}.  Its contravariance in \(F\) is precomposition in the Ext
argument.

\noindent\textbf{Step 3.} For \(q>2\), the displayed Ext degree \(2-q\) is negative.  The
derived Hom of the locally free sheaf \(F\), concentrated in degree
zero, into the sheaf \(\omega_W\), concentrated in degree zero, has
no negative cohomology.  Hence the duality theorem gives
\(H^q(W,F)=0\).  This also follows directly from coherent
cohomological dimension on the compact surface \(W\).
\end{proof}



\begin{proposition}[absolute duality on a compact Cartier fibre]
\label{wbII:prop:cartier-duality-absolute}
\StatementRecord{wbII:prop:cartier-duality-absolute}{proved}{Historical events: GPT-6 Astra (OpenAI), ChatGPT Pro mode, effort Pro, S6\_09\_PAPERS, 17 November 2026: targeted proof check. Also checked in this writing. Pre-continuation writing events: GPT-6 Astra (OpenAI), ChatGPT Pro mode, effort Pro, S6\_10\_PAPERS, 18 November 2026: targeted proof check; used and rechecked in this run. Continuation metadata events: GPT-6 Astra (OpenAI), ChatGPT Pro mode, effort Pro, S6\_10\_PAPERS continuation, 19 November 2026: used and rechecked in this run; status-vocabulary reconciliation; metadata reconciliation only; no new mathematical check. \emph{Pre-continuation recorded proof scope.} The full fibre defines a Cartier equation even with multiplicities. Its normal line is the pullback of the one-dimensional base fibre; choosing a nonzero vector identifies the adjunction line with the restricted canonical line and changes the trace only by a scalar. Ramis-Ruget trace duality applies on the compact fibre itself, with no ambient compactness assumption, and local freeness removes the higher sheaf Ext terms}{the hypotheses, notation and definitions stated here; \cite[Proposition~2 and the complete-intersection calculation,
pp.~90--91]{RR70push}; \cite[Theorem~2, pp.~82--83]{RR70push}; \cite{Grauert1960D}}

\ Let $X$ be a smooth complex threefold, not necessarily compact,
and let $i:W\hookrightarrow X$ be a compact effective Cartier divisor that
is the full scheme-theoretic fibre of a holomorphic map $f:X\to B$ to a
smooth curve.  For every vector bundle $E$ on $X$ there is a functorial
perfect pairing
\[
 H^2(W,E|_W)\times
 H^0(W,E^\vee|_W\otimes\omega_X|_W)\longrightarrow\mathbb C.
\]
Equivalently,
\[
 H^2(W,E|_W)^\vee\simeq
 H^0(W,E^\vee|_W\otimes\omega_X|_W).
\tag{W.12}\label{wbII:eq:W:5.15a}
\]
Neither normality of $W$ nor compactness of the ambient threefold is
required.
\end{proposition}



\begin{proof}
\noindent\textbf{Step 1.} An effective Cartier divisor in a smooth threefold is locally a
two-dimensional hypersurface quotient of a regular local ring.  Hence $W$
is Cohen--Macaulay and Gorenstein.  If $K_W^\bullet$ denotes its analytic
dualizing complex, the complete-intersection calculation of Ramis--Ruget
gives
\[
 K_W^\bullet\simeq\omega_W[2]
\]
with $\omega_W$ invertible
\cite[Proposition~2 and the complete-intersection calculation,
pp.~90--91]{RR70push}.  Cartier adjunction gives
\[
 \omega_W\simeq
 (\omega_X\otimes\mathcal O_X(W))|_W.
\]
Because $W$ is the full fibre over $b\in B$, one has
$\mathcal O_X(W)\simeq f^*\mathcal O_B(b)$ locally along $W$; its
restriction to $W$ is the pullback of a one-dimensional vector space.
Choosing a nonzero vector identifies
$\mathcal O_X(W)|_W\simeq\mathcal O_W$ and
$\omega_W\simeq\omega_X|_W$.  A different choice rescales the final trace
pairing by a nonzero scalar.

\noindent\textbf{Step 2.} Ramis--Ruget duality pairs $H_c^p(W,\mathcal G)$ with
$\operatorname{Ext}^{-p}_W(\mathcal G,K_W^\bullet)$
\cite[Theorem~2, pp.~82--83]{RR70push}.  Since $W$ is compact,
$H_c^p=H^p$.  Taking $p=2$ and $\mathcal G=E|_W$, the shift and local
freeness give
\[
 \operatorname{Ext}^{-2}_W(\mathcal G,K_W^\bullet)
 \simeq\operatorname{Hom}_W(\mathcal G,\omega_W)
 =H^0(W,E^\vee|_W\otimes\omega_W).
\]
Grauert finiteness makes both coherent cohomology spaces finite-dimensional
\cite{Grauert1960D}; the strong analytic duality is therefore the ordinary
algebraic perfect pairing \eqref{wbII:eq:W:5.15a}.  Its variance is contravariant in $E|_W$
by precomposition on the Hom term.  This proves the assertion.
\end{proof}



\subsection{Flatness and top-degree analytic base change}

\begin{lemma}[automatic properness, flatness, and Cartier fibres]
\label{wbII:lem:curve-flatness}
\StatementRecord{wbII:lem:curve-flatness}{proved}{Historical events: GPT-6 Astra (OpenAI), ChatGPT Pro mode, effort Pro, S6\_09\_PAPERS, 17 November 2026: targeted proof check. Also checked in this writing. Pre-continuation writing events: GPT-6 Astra (OpenAI), ChatGPT Pro mode, effort Pro, S6\_10\_PAPERS, 18 November 2026: targeted proof check; used and rechecked in this run. Continuation metadata events: GPT-6 Astra (OpenAI), ChatGPT Pro mode, effort Pro, S6\_10\_PAPERS continuation, 19 November 2026: used and rechecked in this run; status-vocabulary reconciliation; metadata reconciliation only; no new mathematical check. \emph{Pre-continuation recorded proof scope.} A nonzero base germ cannot pull back to zero without making the holomorphic map locally and hence globally constant. The local domain is therefore torsion-free over the base DVR, and the exact torsion-free/flat equivalence proves flatness, including vector bundles. A base parameter is a nonzerodivisor, giving every nonempty full fibre pure codimension one; compactness proves properness. The precise Stacks tag0AUW was opened}{the hypotheses, notation and definitions stated here; \cite[Tag~0AUW]{StacksProject}}

\ Let $f:X\to B$ be a nonconstant holomorphic map from a connected
complex manifold to a Riemann surface.  Then $f$ is flat, every holomorphic
vector bundle on $X$ is $f$-flat, and every full scheme-theoretic fibre is an
effective Cartier divisor of pure codimension one.  If $X$ is compact, then
$f$ is proper.
\end{lemma}



\begin{proof}
\noindent\textbf{Step 1.} Fix $x\in X$, $c=f(x)$, and put
$R=\mathcal O_{B,c}$, $S=\mathcal O_{X,x}$.  The ring $R$ is a discrete
valuation ring and $S$ is a regular local domain.  The homomorphism
$R\to S$ is injective.  Indeed, if a nonzero germ $h\in R$ mapped to zero,
$h\circ f$ would vanish on a neighbourhood of $x$.  The zero set of a
nonzero one-variable holomorphic germ is isolated, so $f$ would be constant
on that neighbourhood and therefore, by the identity theorem on connected
$X$, constant everywhere.

\noindent\textbf{Step 2.} Thus every nonzero element of $R$ acts on the domain $S$ as a nonzero
element and hence as a nonzerodivisor.  Therefore $S$ is torsion-free as an
$R$-module.  A module over a discrete valuation ring is flat if and only if
it is torsion-free \cite[Tag~0AUW]{StacksProject}; hence $S$ is $R$-flat.
This holds at every point, so $f$ is flat.  A vector bundle is locally a
finite free $\mathcal O_X$-module, and a finite direct sum of $R$-flat
modules is flat; hence it is $f$-flat.  If $t$ is a local parameter at
$c$, then $f^*t$ is a nonzero element of the domain $S$ and thus a
nonzerodivisor.  The full scheme-theoretic fibre is therefore locally cut
out by one nonzerodivisor, so it is an effective Cartier divisor of pure
codimension one.  Finally, when $X$ is compact, the underlying continuous
map to the Hausdorff surface $B$ is proper.
\end{proof}



\begin{proposition}[the exact analytic base-change input]
\label{wbII:prop:analytic-base-change}
\StatementRecord{wbII:prop:analytic-base-change}{proved}{Historical events: GPT-6 Astra (OpenAI), ChatGPT Pro mode, effort Pro, S6\_09\_PAPERS, 17 November 2026: targeted proof check. Also checked in this writing. Pre-continuation writing events: GPT-6 Astra (OpenAI), ChatGPT Pro mode, effort Pro, S6\_10\_PAPERS, 18 November 2026: targeted proof check. GPT-6 Astra (OpenAI), ChatGPT Pro mode, effort Pro, S6\_10\_PAPERS, 18 November 2026 04:30:14.617321+00:00: used and rechecked in this run. GPT-6 Astra (OpenAI), ChatGPT Pro mode, effort Pro, S6\_10\_PAPERS, 18 November 2026 05:35:21+00:00: used and rechecked in this run. Continuation metadata events: GPT-6 Astra (OpenAI), ChatGPT Pro mode, effort Pro, S6\_10\_PAPERS continuation, 19 November 2026: used and rechecked in this run; status-vocabulary reconciliation; metadata reconciliation only; no new mathematical check. \emph{Pre-continuation recorded proof scope.} The degree-three fibre vanishing makes its canonical fibre map surjective and its direct image vanish near the point; the analytic base-change criterion then gives the degree-two canonical isomorphism. The operative properness, coherence and relative flatness are exactly those in Banica-Stanasila III.3 Theorem3.4 and Corollary3.5, whose current exact source opening is separately recorded by the inventory reviewer. No local-freeness assertion for all direct images is used}{the hypotheses, notation and definitions stated here; \cite[pp.~120--121]{BS76push}; \cite[Lemma~1.5]{CDP98}}

\ Let $f:X\to B$ be proper between complex spaces, let
$\mathcal F$ be coherent and $f$-flat, and fix $c\in B$.  If
\[
H^3(X_c,\mathcal F|_{X_c})=0,
\tag{W.13}\label{wbII:eq:W:5.16}
\]
then the canonical fibre map
\[
\beta_c^2:
(R^2f_*\mathcal F)_c\otimes_{\mathcal O_{B,c}}\mathbb C(c)
\longrightarrow H^2(X_c,\mathcal F|_{X_c})
\tag{W.14}\label{wbII:eq:W:5.17}
\]
is an isomorphism.
\end{proposition}



\begin{proof}
\noindent\textbf{Step 1.} Apply B\u{a}nic\u{a}--St\u{a}n\u{a}\c{s}il\u{a},
Chapter~III, \S3, Corollary~3.5, with $q=3$
\cite[pp.~120--121]{BS76push}.  Its hypotheses are precisely
properness, coherence, and flatness.  Condition \eqref{wbII:eq:W:5.16} first forces the
degree-three direct image to vanish near $c$.  The corollary then applies the
equivalent base-change conditions of its preceding Theorem~3.4 in degree
$q-1=2$ and states that the canonical map \eqref{wbII:eq:W:5.17} is bijective.  This is the
top-degree argument used in \cite[Lemma~1.5]{CDP98}; it does not require
the stronger, false assertion that all direct images are locally free.
\end{proof}



\begin{corollary}[uniform top-degree base change for surface fibres]
\label{wbII:cor:surface-top-base-change}
\StatementRecord{wbII:cor:surface-top-base-change}{proved}{Historical events: GPT-6 Astra (OpenAI), ChatGPT Pro mode, effort Pro, S6\_09\_PAPERS, 17 November 2026: targeted proof check. Also checked in this writing. Pre-continuation writing events: GPT-6 Astra (OpenAI), ChatGPT Pro mode, effort Pro, S6\_10\_PAPERS, 18 November 2026: targeted proof check. GPT-6 Astra (OpenAI), ChatGPT Pro mode, effort Pro, S6\_10\_PAPERS, 18 November 2026 04:30:14.617321+00:00: used and rechecked in this run. Continuation metadata events: GPT-6 Astra (OpenAI), ChatGPT Pro mode, effort Pro, S6\_10\_PAPERS continuation, 19 November 2026: used and rechecked in this run; status-vocabulary reconciliation; metadata reconciliation only; no new mathematical check. \emph{Pre-continuation recorded proof scope.} The relative-flat finite free complex identifies its derived fibre with the ordinary full fibre. Over the base DVR the exact universal-coefficient sequence has only tensor and Tor1. Choosing the largest nonzero cohomology degree above two and applying Nakayama contradicts coherent cohomological dimension on the surface; degree two then has zero Tor1 obstruction. The exact Banica-Stanasila finite-complex source statement is separately recorded, including its bounded source-reading scope}{the hypotheses, notation and definitions stated here; \cite[Chapter~III, \S4, Theorem~4.1, p.~124]{BS76push}; \cref{wbII:eq:W:5.17-U}; \cref{wbII:prop:analytic-base-change}}

\ Let $f:X\to B$ be a proper flat holomorphic map to a Riemann
surface whose full scheme-theoretic fibres have pure complex dimension two,
and let $E$ be a holomorphic vector bundle on $X$.  Then
\[
 R^qf_*E=0\quad(q>2)
 \tag{W.15}\label{wbII:eq:W:5.17-V}
\]
and, for every $b\in B$, including a point with nonreduced fibre, the
canonical map
\[
 (R^2f_*E)_b\otimes_{\mathcal O_{B,b}}\mathbb C(b)
 \xrightarrow{\ \sim\ }H^2(X_b,E|_{X_b})
 \tag{W.16}\label{wbII:eq:W:5.17-T}
\]
is an isomorphism.  No constancy of $h^2$, local freeness of $R^2f_*E$, or
reducedness of $X_b$ is required.
\end{corollary}



\begin{proof}
\noindent\textbf{Step 1.} On a sufficiently small neighbourhood $U$ of $b$, the finite-complex theorem
used in analytic cohomology and base change represents $Rf_*E|_U$ by a
bounded complex $K^\bullet$ of finite free $\mathcal O_U$-modules
\cite[Chapter~III, \S4, Theorem~4.1, p.~124]{BS76push}.  Flatness of
$E$ over $B$ identifies its derived fibre with the ordinary restriction to
the full scheme-theoretic fibre.  Over the discrete valuation ring
$A=\mathcal O_{B,b}$, the universal-coefficient sequence is
\[
0\longrightarrow H^q(K^\bullet)_b\otimes_A\mathbb C(b)
 \longrightarrow H^q(K^\bullet\otimes_A\mathbb C(b))
 \longrightarrow
 \operatorname{Tor}_1^A(H^{q+1}(K^\bullet)_b,\mathbb C(b))
 \longrightarrow0.
\tag{W.17}\label{wbII:eq:W:5.17-U}
\]
If some $R^qf_*E$ with $q>2$ were nonzero, choose the largest such $q$ and a
point of its support.  Nakayama's lemma makes the left term of \eqref{wbII:eq:W:5.17-U}
nonzero there, while maximality kills the Tor term.  This would give nonzero
coherent cohomology in degree $q>2$ on a compact two-dimensional complex
space, a contradiction.  Thus \eqref{wbII:eq:W:5.17-V} holds.  For $q=2$, the only
obstruction in \eqref{wbII:eq:W:5.17-U} is
$\operatorname{Tor}_1^A((R^3f_*E)_b,\mathbb C(b))$, which is zero by
\eqref{wbII:eq:W:5.17-V}, proving \eqref{wbII:eq:W:5.17-T}.  This is the explicit obstruction calculation
behind \cref{wbII:prop:analytic-base-change}..
\end{proof}



\subsection{Marked data used in the restriction arguments}
Here $f:X\to B=\mathbb P^1$ is the standard family of Paper I,
$B^\circ=B\setminus\{p_0,p_1,p_2\}$, and $F$ is a smooth fibre.
The cusp is $W=f^{-1}(p_0)$ and $i:W\hookrightarrow X$ is its inclusion.
Its normalization is $\eta:\widetilde W\simeq\mathrm{dP}_6\to W$.
Use the ordered dual lattices
\[
 V=H^1(F;\mathbb Z)=\mathbb Z\langle\gamma,u,w,\delta\rangle,
 \qquad
 \Lambda=H_1(F;\mathbb Z)
 =\mathbb Z\langle\widehat\gamma,\widehat u,\widehat w,\widehat\delta\rangle.
\]
Vectors are columns. The cohomology monodromies and their dual homology
actions are
\[
T_1=\begin{pmatrix}1&0&-6&2\\0&-1&1&1\\0&-1&0&1\\0&0&0&1\end{pmatrix},
\quad
T_2=\begin{pmatrix}1&6&0&-3\\0&0&-1&1\\0&1&0&0\\0&0&0&1\end{pmatrix},
\quad A_j=(T_j^{-1})^{\mathsf t}.
\]
Put $T_0=(T_1T_2)^{-1}$ and let $G$ denote the generated monodromy group.
The finite-filling orders, translations and invariant forms are
\[
\begin{gathered}
(m_1,m_2)=(3,4),\qquad
v_1=(1,2,-4,0)^{\mathsf t},\quad v_2=(-1,-3,3,0)^{\mathsf t},\\
\eta_1=2u+w+3\delta,\quad \eta_2=u+w+2\delta,
\qquad q=u\wedge w+6\gamma\wedge\delta.
\end{gathered}
\]
The actual period functions of Paper I define
\[
\Pi(z)=\begin{pmatrix}6\mu(z)&\tau(z)&1&0\\
\beta(z)&\mu(z)&0&1\end{pmatrix},
\quad
\operatorname{Im}\beta-
\frac{6(\operatorname{Im}\mu)^2}{\operatorname{Im}\tau}<0,
\]
so $\Pi(z):\Lambda_{\mathbb R}\to\mathbb C^2$ is a real isomorphism.
The marking is $\Xi(z,x+\Lambda)=[z,\Pi(z)x]$.
The finite overlap on the $s_j$-cover is translation by
$(\log s_j)/(2\pi i)\,\Pi(s_j)v_j$.
At the cusp use the unchanged zero-section gluing, with no additional
vertical translation. These are the data constructed and proved in
\cref{wb:prop:period-input-audit,wb:thm:candidate-cusp-geometry,wb:prop:candidate-finite-fillings}. The arguments
below do not replace these data by a topologically equivalent model.

\subsection{A topology-independent \texorpdfstring{$\operatorname{Pic}^0$}{Pic0} route and three all-Pic routes}

The vanishing of the normalized restriction of every global line bundle
for this threefold is already established in
\citet[Theorem~10.5, proof Step~1, equation~(10.5), p.~88]{AC2026}.
The complete arguments in
\cref{wbII:prop:leray-edge-restriction,wbII:prop:direct-boundary-picard,wbII:thm:analytic-all-picard}
are retained as three distinct proofs of that known conclusion.  The
first uses the integral Leray edge, the second uses the boundary norm
identities and local specialization, and the third uses holomorphic
factors and toric support functions.  The preliminary
\(\operatorname{Pic}^0\) argument isolates the weaker hypotheses already
sufficient for topologically trivial twists.

Let $i:W\hookrightarrow X$ denote the cusp-fibre inclusion and
$\eta:\widetilde W\simeq\mathrm{dP}_6\to W$ its normalization.

\begin{proposition}[topologically trivial twists pull back trivially]
\label{wbII:prop:pic0-pullback}
\StatementRecord{wbII:prop:pic0-pullback}{proved}{Historical events: GPT-6 Astra (OpenAI), ChatGPT Pro mode, effort Pro, S6\_09\_PAPERS, 17 November 2026: targeted proof check. Also checked in this writing. Pre-continuation writing events: GPT-6 Astra (OpenAI), ChatGPT Pro mode, effort Pro, S6\_10\_PAPERS, 18 November 2026: targeted proof check. Continuation metadata events: GPT-6 Astra (OpenAI), ChatGPT Pro mode, effort Pro, S6\_10\_PAPERS continuation, 19 November 2026: used and rechecked in this run; status-vocabulary reconciliation; metadata reconciliation only; no new mathematical check. \emph{Pre-continuation recorded proof scope.} Naturality kills the first Chern class of the pulled-back Pic0 bundle. The exponential sequence on the three-point blowup of the projective plane, with H1(O)=0, makes its Picard-to-integral-H2 map injective. The assumed local canonical trivialization gives the asserted dual twist without an assumption on H2 of the ambient space}{the hypotheses, notation and definitions stated here; the explicit constructions and calculations in the proof}

\ Assume that $\omega_X$ is trivial on a neighbourhood of $W$.
For every $L\in\operatorname{Pic}^0(X)$, put
\[
 A_L=(L^\vee\otimes\omega_X)|_W.
\]
Then
\[
 \eta^*A_L\simeq\mathcal O_{\mathrm{dP}_6}.
\tag{W.18}\label{wbII:eq:W:5.17a}
\]
No assertion about $H^2(X;\mathbb Z)$ is used.
\end{proposition}



\begin{proof}
\noindent\textbf{Step 1.} Every line bundle in $\operatorname{Pic}^0(X)$ has $c_1(L)=0$.  Naturality
therefore gives $c_1(\eta^*i^*L)=0$.  The degree-six del Pezzo surface is the
blow-up of $\mathbb P^2$ in three noncollinear points, so
$H^1(\mathrm{dP}_6,\mathcal O)=0$.  The exponential sequence consequently
makes
\[
 c_1:\operatorname{Pic}(\mathrm{dP}_6)\longrightarrow
 H^2(\mathrm{dP}_6;\mathbb Z)
\]
injective.  Hence $\eta^*i^*L$ is trivial.  The assumed neighbourhood
trivialization of $\omega_X$ also pulls back trivially, which proves
\eqref{wbII:eq:W:5.17a}.
\end{proof}



\begin{proposition}[the exact all-line-bundle restriction obstruction]
\label{wbII:prop:all-pic-restriction}
\StatementRecord{wbII:prop:all-pic-restriction}{proved}{Historical events: GPT-6 Astra (OpenAI), ChatGPT Pro mode, effort Pro, S6\_09\_PAPERS, 17 November 2026: targeted proof check. Also checked in this writing. Pre-continuation writing events: GPT-6 Astra (OpenAI), ChatGPT Pro mode, effort Pro, S6\_10\_PAPERS, 18 November 2026: targeted proof check. Continuation metadata events: GPT-6 Astra (OpenAI), ChatGPT Pro mode, effort Pro, S6\_10\_PAPERS continuation, 19 November 2026: used and rechecked in this run; status-vocabulary reconciliation; metadata reconciliation only; no new mathematical check. \emph{Pre-continuation recorded proof scope.} The actual normalization and fibre inclusions give r=eta-star i-star. Pullback of duals and the specified canonical trivialization identify eta-star A\_L with r(L) dual, proving the simultaneous equivalence with r=0 and the Pic0 case through the preceding exponential-sequence calculation}{the hypotheses, notation and definitions stated here; \cref{wbII:prop:conductor}; \cref{wbII:prop:pic0-pullback}}

\ Define
\[
 r=\eta^*i^*:\operatorname{Pic}(X)\longrightarrow
 \operatorname{Pic}(\mathrm{dP}_6).
\]
Under the same local canonical trivialization,
\[
 \eta^*A_L\simeq r(L)^\vee
 \qquad(L\in\operatorname{Pic}(X)).
\tag{W.19}\label{wbII:eq:W:5.17b}
\]
Thus the trivial-pullback hypothesis used by
\cref{wbII:prop:conductor} holds simultaneously for every line bundle exactly
when $r=0$.  It always holds on $\operatorname{Pic}^0(X)$.
\end{proposition}



\begin{proof}
\noindent\textbf{Step 1.} Pullback commutes with tensor product and duals, while the local
trivialization gives $\eta^*i^*\omega_X\simeq\mathcal O$.  This is \eqref{wbII:eq:W:5.17b},
and its uniform equivalence with $r=0$ is immediate.  The final assertion is
\cref{wbII:prop:pic0-pullback}.
\end{proof}



\begin{proposition}[Leray-edge vanishing of normalized cusp restrictions]
\label{wbII:prop:leray-edge-restriction}
\StatementRecord{wbII:prop:leray-edge-restriction}{proved}{Historical events: GPT-6 Astra (OpenAI), ChatGPT Pro mode, effort Pro, S6\_09\_PAPERS, 17 November 2026: targeted proof check. Also checked in this writing. Pre-continuation writing events: GPT-6 Astra (OpenAI), ChatGPT Pro mode, effort Pro, S6\_10\_PAPERS, 18 November 2026: targeted proof check; used and rechecked in this run. Continuation metadata events: GPT-6 Astra (OpenAI), ChatGPT Pro mode, effort Pro, S6\_10\_PAPERS continuation, 19 November 2026: used and rechecked in this run; status-vocabulary reconciliation; metadata reconciliation only; no new mathematical check. \emph{Pre-continuation recorded proof scope.} The filtered-complex map under fibre inclusion identifies evaluation of the degree-two edge section with the actual integral restriction. The edge image consists of permanent cycles; the cited explicit candidate differential sends 2q to the primitive omega-delta class and is an isomorphism. Thus the edge and every normalized cusp Chern class vanish; exponential-sequence injectivity on dP6 trivializes the line bundle. The full argument was re-derived. The literal source formulas and full statement blocks identify the legacy PaperI D2 and D4 dependencies as mapsI:top:prop:leray-e2 and mapsI:top:lem:higher-transgressions, with source hashes retained in the exact mapping record}{the hypotheses, notation and definitions stated here; \cite[Tags~012O, 012U and~012W]{StacksProject}; \cite[Theorem~2.3.26]{Dimca2004}; \cref{wbII:eq:W:5.17d}; \cite[Tags~011N, 011O and~012U]{StacksProject}; \cref{wbII:eq:W:5.17e}; \cref{wbII:prop:pic0-pullback}}

\ For the constructed map $f:X\to B$ and every
$L\in\operatorname{Pic}(X)$, the degree-two Leray edge section satisfies
\[
 e(c_1(L))=0\in H^0(B,R^2f_*\mathbb Z),
 \qquad
 \eta^*i^*L\simeq\mathcal O_{\mathrm{dP}_6}.
\tag{W.20}\label{wbII:eq:W:5.17c}
\]
This proof does not invoke the packaged conclusion
$H^2(X;\mathbb Z)=0$.
\end{proposition}



\begin{proof}
\noindent\textbf{Step 1.} First record explicitly the compatibility used here.  For every $b\in B$,
proper base change and functoriality of the filtered-complex spectral
sequence give a commutative square
\[
\begin{CD}
 H^2(X;\mathbb Z) @>{e}>> H^0(B,R^2f_*\mathbb Z)\\
 @V{i_b^*}VV @VV{\operatorname{ev}_b}V\\
 H^2(X_b;\mathbb Z) @= H^2(X_b;\mathbb Z).
\end{CD}
\tag{W.21}\label{wbII:eq:W:5.17d}
\]
Indeed, pull the Leray filtered complex back along $\{b\}\to B$; the
spectral sequence over a point has only column zero, so its edge is the
identity, while the right vertical map is the proper-base-change
identification.  This uses the functoriality and convergence statements in
\cite[Tags~012O, 012U and~012W]{StacksProject} and analytic proper base
change \cite[Theorem~2.3.26]{Dimca2004}.  The connecting homomorphism in the
exponential sequence is natural under $i_b^*$, so \eqref{wbII:eq:W:5.17d} sends
$c_1(L)$ to $c_1(L|_{X_b})$.

\noindent\textbf{Step 2.} Put $s_L=e(c_1(L))$.  The image of an edge morphism consists of permanent
cycles, hence
\[
 s_L\in E_\infty^{0,2}\subseteq\ker
 \bigl(d_2^{0,2}:E_2^{0,2}\to E_2^{2,1}\bigr);
\tag{W.22}\label{wbII:eq:W:5.17e}
\]
see \cite[Tags~011N, 011O and~012U]{StacksProject}.  For the constructed
candidate, \cref{mapsI:top:prop:leray-e2}; \cref{mapsI:top:lem:higher-transgressions} give
\[
 E_2^{0,2}=\mathbb Z(2q),\qquad
 E_2^{2,1}=\mathbb Z\,\omega[\delta],\qquad
 d_2^{0,2}(2q)=\omega[\delta].
\]
Thus $d_2^{0,2}$ is an isomorphism and \eqref{wbII:eq:W:5.17e} forces $s_L=0$.
Taking $b=p_0$ in \eqref{wbII:eq:W:5.17d} gives $c_1(i^*L)=0$, and hence
$c_1(\eta^*i^*L)=0$.  Finally, the injectivity of
$c_1:\operatorname{Pic}(\mathrm{dP}_6)\to H^2(\mathrm{dP}_6;\mathbb Z)$
proved in \cref{wbII:prop:pic0-pullback} makes $\eta^*i^*L$ trivial.
\end{proof}



\begin{lemma}[the pair-of-pants defect class]
\label{wbII:lem:picard-crossed-defect}
\StatementRecord{wbII:lem:picard-crossed-defect}{proved}{Historical events: GPT-6 Astra (OpenAI), ChatGPT Pro mode, effort Pro, S6\_09\_PAPERS, 17 November 2026: targeted proof check. Also checked in this writing. Pre-continuation writing events: GPT-6 Astra (OpenAI), ChatGPT Pro mode, effort Pro, S6\_10\_PAPERS, 18 November 2026: targeted proof check; used and rechecked in this run. Continuation metadata events: GPT-6 Astra (OpenAI), ChatGPT Pro mode, effort Pro, S6\_10\_PAPERS continuation, 19 November 2026: used and rechecked in this run; status-vocabulary reconciliation; metadata reconciliation only; no new mathematical check. \emph{Pre-continuation recorded proof scope.} The fixed integral fibre form descends under the actual flat marking. The graph spine gives the two-column Serre edge sequence, so the fibre-zero Chern-class difference is a unique H1 local-system class represented by a crossed homomorphism. Applying it to the based three-meridian relation proves the coinvariant balance. The base-first cellular form and fibre-first sweep differ by one interchange, giving the exact minus sign; principal changes pair trivially with invariant dual vectors}{the hypotheses, notation and definitions stated here; \cref{wbII:eq:W:5.17j}; \cite[Chapter~III, \S1, Exercise~2, p.~60; Chapter~IV, \S2, pp.~88--89]{Brown1982}}

\ Choose the compact pair-of-pants core $P\subset B^\circ$ and
the zero-section whiskers used in \cref{mapsI:top:prop:geometric-pi1}, and put
$J_P=f^{-1}(P)$.  Orient its boundary meridians clockwise.  Write
$\rho_0$ for the projection of the zero-section cusp meridian
$\widetilde a_0$.  The standard family uses the tautological cusp
gluing, whose affine regluing class is $\mu=0$ (equivalently
$\ell_0=\gamma(\mu)=0$); hence the flat-marked cusp lift has no vertical
translation.  The base loops and their chosen lifts therefore satisfy
\[
 \rho_1\rho_2\rho_0=1,
 \qquad
 \widehat\rho_1\widehat\rho_2\widetilde a_0=1.
\tag{W.23}\label{wbII:eq:W:5.17g}
\]
For $L\in\operatorname{Pic}(X)$, write
$E=c_1(L|_F)\in\bigwedge^2V$.  Then $E\in(\bigwedge^2V)^G$, and there is a
crossed homomorphism
\[
 z:\pi_1(P)\longrightarrow V_{\mathbb R},\qquad
 z(gh)=z(g)+T_gz(h),
\tag{W.24}\label{wbII:eq:W:5.17h}
\]
whose class records
$c_1(L|_{J_P})-[\Omega_E]$.  Here $\Omega_E$ is the closed, translation-
invariant vertical two-form induced by $E$ in the flat marking.  If
$d_i=z(\rho_i)$ and a bar denotes the image in
$V_G\otimes\mathbb R$, then
\[
 \bar d_1+\bar d_2+\bar d_0=0.
\tag{W.25}\label{wbII:eq:W:5.17i}
\]
The three barred classes, and the number
$\langle d_0,\widehat\delta\rangle$, do not depend on the chosen crossed-
homomorphism representative.  More precisely, use the base-first convention
for the edge injection in \eqref{wbII:eq:W:5.17j}.  For any oriented loop $\rho\subset P$
and any $A_\rho$-invariant $\nu\in\Lambda$, its flat-marked sweep satisfies
\[
 \bigl\langle z(\rho),\nu\bigr\rangle
 =-\bigl\langle c_1(L|_{J_P})-[\Omega_E],
 a_*([S^1_\nu]\times[\widetilde\rho])\bigr\rangle.
\tag{W.26}\label{wbII:eq:W:5.17j'}
\]
Here $a$ is the fibrewise translation action and $\widetilde\rho$ is the
flat-marked lift of $\rho$.
\end{lemma}



\begin{proof}
\noindent\textbf{Step 1.} The flat marking identifies the pullback of $J_P$ to the universal cover of
$P$ with
$\widetilde P\times(\Lambda_{\mathbb R}/\Lambda)$.  Since $E$ is the
restriction of a class on the total space, it is fixed by every monodromy
$T_g$.  The constant vertical form representing $E$ is therefore invariant
under the deck action and descends to the asserted form $\Omega_E$ on
$J_P$.

\noindent\textbf{Step 2.} The pair of pants retracts onto a graph.  Hence the Serre spectral sequence
of the marked real torus bundle has only columns zero and one and gives the
short exact edge sequence
\[
0\longrightarrow H^1(P;V_{\mathbb R})
 \longrightarrow H^2(J_P;\mathbb R)
 \longrightarrow H^0(P;\textstyle\bigwedge^2V_{\mathbb R})
 \longrightarrow0.
\tag{W.27}\label{wbII:eq:W:5.17j}
\]
Both $c_1(L|_{J_P})$ and $[\Omega_E]$ map to $E$, so their difference is
the image of a unique class in $H^1(P;V_{\mathbb R})$.  The standard group-
cohomology model for a graph identifies this group with crossed
homomorphisms modulo principal crossed homomorphisms
\cite[Chapter~III, \S1, Exercise~2, p.~60; Chapter~IV, \S2, pp.~88--89]{Brown1982}; this gives \eqref{wbII:eq:W:5.17h}.

\noindent\textbf{Step 3.} For completeness, the slant convention in \eqref{wbII:eq:W:5.17j'} is part of this
identification, not an additional geometric assertion.  Give an oriented
edge of a graph spine of $P$ a coordinate $s$ increasing with its based
loop and represent the corresponding $V_{\mathbb R}$-valued cellular
one-cochain by a base-first form $ds\wedge\beta_s$.  On the sweep of an
$A_\rho$-invariant fibre circle $S^1_\nu$, Fubini's formula gives
\[
\int_{S^1_s\times S^1_\nu}ds\wedge\beta_s
 =\int_{S^1_s}\langle\beta_s,\nu\rangle,ds
 =\langle z(\rho),\nu\rangle.
\]
The sweep in \eqref{wbII:eq:W:5.17j'} is written in the fibre-first order
$[S^1_\nu]\times[\widetilde\rho]$, which differs from this base-first
cross-product orientation by $(-1)^{1\cdot1}=-1$; this is the displayed
minus sign.
Changing the cellular representative by a principal cochain changes both
sides by the same boundary term, which vanishes because $A_\rho\nu=\nu$.
This proves \eqref{wbII:eq:W:5.17j'}.

\noindent\textbf{Step 4.} Applying $z$ to \eqref{wbII:eq:W:5.17g} gives
\[
 d_1+T_1d_2+T_1T_2d_0=0.
\]
All $T_g$ act trivially on the global coinvariants, proving \eqref{wbII:eq:W:5.17i}.  A
change of representative has the form
$z(g)\mapsto z(g)+(T_g-I)a$.  Its image in $V_G\otimes\mathbb R$ is zero.
Moreover the displayed matrices in
\cref{wb:prop:monodromy-certificate} fix $\widehat\delta$ under every dual
action $A_g$.  Therefore
\[
 \langle(T_g-I)a,\widehat\delta\rangle
 =\langle a,(A_g^{-1}-I)\widehat\delta\rangle=0,
\]
which proves the final assertion.
\end{proof}



\begin{lemma}[the finite-cover norm identity]
\label{wbII:lem:picard-finite-norm}
\StatementRecord{wbII:lem:picard-finite-norm}{proved}{Historical events: GPT-6 Astra (OpenAI), ChatGPT Pro mode, effort Pro, S6\_09\_PAPERS, 17 November 2026: targeted proof check. Also checked in this writing. Pre-continuation writing events: GPT-6 Astra (OpenAI), ChatGPT Pro mode, effort Pro, S6\_10\_PAPERS, 18 November 2026: targeted proof check. Continuation metadata events: GPT-6 Astra (OpenAI), ChatGPT Pro mode, effort Pro, S6\_10\_PAPERS continuation, 19 November 2026: used and rechecked in this run; status-vocabulary reconciliation; metadata reconciliation only; no new mathematical check. \emph{Pre-continuation recorded proof scope.} Restricting the crossed homomorphism to the m-fold cover gives the full operator norm, whose principal-coboundary correction vanishes. The finite filling retracts onto its central torus, so the two extended Chern classes agree. In the real marking the literal holomorphic translation is Re(alpha)v+Im(alpha)J\_s v. Direct substitution in the alternating two-form gives its mixed term; the clockwise logarithm changes this vector by minus v. Pairing every lattice circle gives minus contraction by v, and the total extended mixed class vanishes, proving N\_T d=contraction\_v E with the stated sign}{the hypotheses, notation and definitions stated here; \cref{wbII:lem:picard-crossed-defect}; \cref{wbII:eq:W:5.17g}; \cref{wbII:eq:W:5.17m'}; \cref{wbII:eq:W:5.17l}; \cref{wbII:eq:W:5.17l''}; \cref{wbII:eq:W:5.17l'''}}

\ Retain the notation of
\cref{wbII:lem:picard-crossed-defect}; \cref{wb:prop:candidate-finite-fillings}.  For
$j=1,2$, let
\[
 N_{T_j}=I+T_j+\cdots+T_j^{m_j-1}.
\]
With the clockwise conventions of \eqref{wbII:eq:W:5.17g}, the pullback to the
$m_j$-fold $s_j$-cover gives
\[
 N_{T_j}d_j=\iota_{v_j}E
 \quad\text{in }V_{\mathbb R},
\qquad
 m_j\bar d_j=[\iota_{v_j}E]
 \quad\text{in }V_G\otimes\mathbb R.
\tag{W.28}\label{wbII:eq:W:5.17k}
\]
\end{lemma}



\begin{proof}
\noindent\textbf{Step 1.} Restriction of the crossed homomorphism to the subgroup generated by
$\rho_j^{m_j}$ has value
\[
 z(\rho_j^{m_j})
 =\sum_{r=0}^{m_j-1}T_j^rz(\rho_j)=N_{T_j}d_j.
\tag{W.29}\label{wbII:eq:W:5.17l}
\]
This is the cohomological pullback of the annular defect to the
$s_j$-cover; the norm, rather than unqualified multiplication by $m_j$, is
essential before passage to global coinvariants.  A principal change
$d_j\mapsto d_j+(T_j-I)a$ does not alter this value, since
$N_{T_j}(T_j-I)a=(T_j^{m_j}-I)a=0$.

\noindent\textbf{Step 2.} Let $Q_j:J'_j\to N_j$ be the finite quotient map.  The global flat marking
$\Xi$ over the full disc $\Delta_j$ makes the constant vertical form
$\Omega_E$ well defined on $J'_j$.  Radial retraction onto the central
torus makes
\[
 H^2(J'_j;\mathbb R)\longrightarrow H^2(A_0^{(j)};\mathbb R)
\]
an isomorphism.  On the central fibre the quotient is exactly $\pi_j$, and
naturality gives
\[
 (Q_j|_{A_0^{(j)}})^*c_1(L|_{S_j})
 =\pi_j^*c_1(L|_{S_j})=E.
\tag{W.30}\label{wbII:eq:W:5.17l'}
\]
The second equality is the actual radial/stalk identification in
\cref{wb:lem:finite-filling-stalk-map}; on a nearby fibre the gluing map $h_j$
is a translation and therefore acts trivially on $H^2$ of that fibre.
Consequently the two classes
$c_1(Q_j^*(L|_{N_j}))$ and $[\Omega_E]$ have the same central restriction,
and the retraction isomorphism proves
\[
 c_1(Q_j^*(L|_{N_j}))-[\Omega_E]=0
 \quad\text{in }H^2(J'_j;\mathbb R).
\tag{W.31}\label{wbII:eq:W:5.17l''}
\]

\noindent\textbf{Step 3.} On $\Delta_j^*$, the source identity
$h_j\widetilde g_j^{\log}h_j^{-1}=\widetilde g_j^0$ shows that $h_j$
descends from the logarithmic quotient $N_j^*$ to the linear/global
quotient $J^*$.  Thus
\[
 Q_j^*(L|_{N_j})|_{\Delta_j^*}
 \simeq h_j^*(L|_{J^*}).
\tag{W.32}\label{wbII:eq:W:5.17l'''}
\]
The logarithmic overlap is the holomorphic translation
\[
 h_j[s,\zeta]=[s,\zeta+\alpha(s)\Pi(s)v_j],
 \qquad \alpha(s)=\frac{\log s_j}{2\pi i}.
\]
To compute its effect on a real cohomology class, retain the exact
real marking
\[
 \Xi(s,x+\Lambda)=[s,\Pi(s)x],\qquad
 J_s=\Pi(s)^{-1}\circ i\circ\Pi(s)
 \quad\text{on }\Lambda_{\mathbb R}.
\]
The period matrix \(\Pi(s)\) is a real isomorphism by the displayed period inequality;
thus \(J_s\) is a well-defined smooth family of real complex
structures.  On a branch of the logarithm the actual real
translation vector is
\[
 a_j(s)=\operatorname{Re}\alpha(s)\,v_j+
        \operatorname{Im}\alpha(s)\,J_s v_j.
 \tag{W.33}\label{wbII:eq:W:5.17m}
\]
Indeed, applying \(\Pi(s)\) gives precisely
\(\alpha(s)\Pi(s)v_j\).  A different logarithm branch changes
\(a_j\) by an integral multiple of \(v_j\), so its torus translation
is branch-independent.  Complex multiplication by \(\alpha\) is
not multiplication by a real scalar in this marking.

\noindent\textbf{Step 4.} Write \(E=\frac12\sum_{a,b}E_{ab}\,dx^a\wedge dx^b\) in the
ordered real coordinates dual to the marked lattice; the matrix
\((E_{ab})\) is alternating and is constant over the covering
disc.  Direct substitution in the real map
\(h_j(s,x)=(s,x+a_j(s))\) gives
\[
 \begin{split}
 h_j^*\Omega_E-\Omega_E
  ={}&\sum_{a,b}E_{ab}\,d a_j^a\wedge dx^b\\
    &+\frac12\sum_{a,b}E_{ab}\,d a_j^a\wedge d a_j^b.
 \end{split}
 \tag{W.34}\label{wbII:eq:W:5.17m'}
\]
The last term has both arguments in the base directions.
In particular, it restricts to zero when the base is an oriented
circle.

\noindent\textbf{Step 5.} Take one full clockwise circle in the \(s_j\)-cover,
\[
 s_j(\theta)=r e^{-2\pi i\theta},\qquad 0\le\theta\le1.
\]
Choose the logarithm continuously along that circle, so
\(\alpha(1)=\alpha(0)-1\).  The endpoints represent the same
point of the covering disc.  The marking is defined on the full
disc, and therefore \(J_{s_j(1)}=J_{s_j(0)}\).  Consequently
\[
 a_j(1)-a_j(0)=-v_j.
 \tag{W.35}\label{wbII:eq:W:5.17m''}
\]
For any \(\lambda\in\Lambda\), restrict \eqref{wbII:eq:W:5.17m'} to the sweep of
the marked circle \(x(u)=u\lambda\), using the base-first
orientation \(d\theta\wedge du\).  Its integral is exactly
\[
 \int_0^1 E\!\left(\frac{d a_j}{d\theta},\lambda\right)d\theta
 =E(a_j(1)-a_j(0),\lambda)
 =-E(v_j,\lambda).
\]
The horizontal term contributes zero.  Since marked lattice
circles span \(H_1(F;\mathbb R)\), these integrals determine the
mixed component in
\[
 H^1(S^1;\mathbb R)\otimes H^1(F;\mathbb R).
\]
Thus the mixed component of \(h_j^*\Omega_E-\Omega_E\) is
\(-\iota_{v_j}E\), with \((\iota_{v_j}E)(\lambda)=E(v_j,\lambda)\).

\noindent\textbf{Step 6.} The pullback of the defect
\(c_1(L|_{J^*})-[\Omega_E]\) has zero fibre component and mixed
component \(N_{T_j}d_j\) by \eqref{wbII:eq:W:5.17l}.  Translation by \(h_j\) does
not change this mixed component: on the covering circle such
a fibre-zero class is represented by
\(d\theta\wedge\beta\), with \(\beta\) a constant fibre one-form,
and substituting \(dx\mapsto dx+d a_j\) adds only a base two-form,
which vanishes on the circle.  Equations \eqref{wbII:eq:W:5.17l''} and
\eqref{wbII:eq:W:5.17l'''} state that the total pulled-back class extends over
the covering disc and has zero mixed component.  It follows that
\[
 N_{T_j}d_j-\iota_{v_j}E=0.
\]
In the global coinvariants every \(T_j\) acts as the identity,
so \(m_j\bar d_j=[\iota_{v_j}E]\).  This proves \eqref{wbII:eq:W:5.17k}, with the
original clockwise convention and no change to either translation
vector.
\end{proof}



\begin{lemma}[the cusp boundary defect and Chern-class specialization]
\label{wbII:lem:picard-cusp-defect}
\StatementRecord{wbII:lem:picard-cusp-defect}{proved}{Historical events: GPT-6 Astra (OpenAI), ChatGPT Pro mode, effort Pro, S6\_09\_PAPERS, 17 November 2026: targeted proof check. Also checked in this writing. Pre-continuation writing events: GPT-6 Astra (OpenAI), ChatGPT Pro mode, effort Pro, S6\_10\_PAPERS, 18 November 2026: targeted proof check. Continuation metadata events: GPT-6 Astra (OpenAI), ChatGPT Pro mode, effort Pro, S6\_10\_PAPERS continuation, 19 November 2026: used and rechecked in this run; status-vocabulary reconciliation; metadata reconciliation only; no new mathematical check. \emph{Pre-continuation recorded proof scope.} The zero-drift marked lift makes the invariant vertical two-form vanish on its swept torus. In the actual unimodular cusp chart the solid torus |u|=rho, x1=1, |x2|<=1 has exactly the displayed sweep boundary after the clockwise orientation change. An extending Chern class pairs to zero with that boundary; the primitive invariant pairing detects the full real coinvariant. The nearby-cycle unit and the exponential connecting map commute with both inclusions, proving the precise Chern-class specialization square}{the hypotheses, notation and definitions stated here; \cref{wbII:lem:picard-crossed-defect}; \cref{wbII:eq:W:5.17g}; \cref{wbII:eq:W:5.17j'}; \cref{wbII:eq:W:5.17p}}

\ In the notation of
\cref{wbII:lem:picard-crossed-defect}, one has
\[
 \bar d_0=0\quad\text{in }V_G\otimes\mathbb R.
\tag{W.36}\label{wbII:eq:W:5.17n}
\]
Moreover, for every line bundle $L$ on the cusp neighbourhood $N_0$,
nearby-cycle specialization is compatible with first Chern classes:
\[
 \operatorname{sp}^2\bigl(c_1(L|_W)\bigr)=c_1(L|_F).
\tag{W.37}\label{wbII:eq:W:5.17o}
\]
\end{lemma}



\begin{proof}
\noindent\textbf{Step 1.} Let $\widetilde a_0$ be the flat-marked clockwise zero-section meridian in
the lifted relation of \eqref{wbII:eq:W:5.17g}, and let
\[
 \xi_{\widehat\delta}
 =a_*([S^1_{\widehat\delta}]\times[\widetilde a_0])
\]
be its sweep under fibrewise translation by $\widehat\delta$.  The boundary
$\rho_0$ of $P$ was chosen to be its projection.  This is the zero-section
sweep: there is no hidden vertical drift in $\widetilde a_0$, by the
explicit overlap calculation in \cref{mapsI:top:prop:geometric-pi1}.  Hence the
vertical form $\Omega_E$ vanishes on the swept torus, and the slant identity
\eqref{wbII:eq:W:5.17j'} gives
\[
 \langle d_0,\widehat\delta\rangle
 =-\langle c_1(L|_{J_P})-[\Omega_E],
          \xi_{\widehat\delta}\rangle
 =-\langle c_1(L|_{J_P}),\xi_{\widehat\delta}\rangle.
\tag{W.38}\label{wbII:eq:W:5.17p}
\]
Combine the exponential overlap coordinates $(x_1,x_2,t)$ given by $x_k=e^{2\pi i\zeta_k}$ and $t=e^{2\pi is}$ with
the unimodular cusp chart
\[
 (u,x_1,x_2)=\left(\frac{t}{x_1x_2},x_1,x_2\right).
\]
The literal solid torus
\[
 \mathbb S_{\widehat\delta}
 =\{|u|=\rho,\ x_1=1,\ |x_2|\le1\}
\]
has boundary parametrization
\[
 (\theta,\phi)\longmapsto
 \bigl(\rho e^{i(\theta-\phi)},1,e^{i\phi}\bigr)
\tag{W.39}\label{wbII:eq:W:5.17p0}
\]
and hence $t=\rho e^{i\theta}$; its $\phi$-circle is translation by
$\widehat\delta$.  Direct evaluation in the proved cusp marking therefore
identifies $\partial\mathbb S_{\widehat\delta}$ with
$\xi_{\widehat\delta}$.  Reversing $\theta$ supplies the fixed clockwise
convention and only changes the common orientation sign.  No reduced
Mayer--Vietoris block or Leray comparison is used here.
Since $L$ extends across $N_0$, its Chern class evaluates to zero on that
boundary.  Thus the left side of \eqref{wbII:eq:W:5.17p} is zero.  By
\cref{wb:thm:coinvariants}, $V_G=\mathbb Z[\delta]$, and evaluation on the
globally invariant dual vector $\widehat\delta$ identifies
$V_G\otimes\mathbb R$ with $\mathbb R$.  This proves \eqref{wbII:eq:W:5.17n}.

\noindent\textbf{Step 2.} For \eqref{wbII:eq:W:5.17o}, let $\ell=c_1(L|_{N_0})$, let
$\iota_0:W\hookrightarrow N_0$, and let
$\iota_t:F_t\hookrightarrow N_0$ be a sufficiently near smooth cusp fibre.
The neighbourhood realization of the nearby-cycle unit
$\mathbb Z_W\to\psi_{f_0}\mathbb Z$ recalled in the proof of
\cref{mapsI:top:thm:p0-integral-specialization} gives the commutative square
\[
\begin{CD}
 H^2(N_0;\mathbb Z) @>{\iota_0^*}>> H^2(W;\mathbb Z)\\
 @V{\iota_t^*}VV @VV{\operatorname{sp}^2}V\\
 H^2(F_t;\mathbb Z) @= H^2(F_t;\mathbb Z).
\end{CD}
\tag{W.40}\label{wbII:eq:W:5.17p'}
\]
Applying this square to $\ell$ gives
$\operatorname{sp}^2(c_1(L|_W))=c_1(L|_{F_t})$.  Flat parallel transport
in the marked smooth family identifies the latter with $c_1(L|_F)=E$.
Equivalently, $c_1$ is the connecting morphism in the exponential sequence,
and restriction to the shrinking neighbourhoods and to $F_t$ commutes with
that morphism before passage to the direct limit.  This proves the claimed
naturality without using any global Leray differential.
\end{proof}



\begin{proposition}[direct boundary-Picard vanishing]
\label{wbII:prop:direct-boundary-picard}
\StatementRecord{wbII:prop:direct-boundary-picard}{proved}{Historical events: GPT-6 Astra (OpenAI), ChatGPT Pro mode, effort Pro, S6\_09\_PAPERS, 17 November 2026: targeted proof check. Also checked in this writing. Pre-continuation writing events: GPT-6 Astra (OpenAI), ChatGPT Pro mode, effort Pro, S6\_10\_PAPERS, 18 November 2026: targeted proof check. Continuation metadata events: GPT-6 Astra (OpenAI), ChatGPT Pro mode, effort Pro, S6\_10\_PAPERS continuation, 19 November 2026: used and rechecked in this run; status-vocabulary reconciliation; metadata reconciliation only; no new mathematical check. \emph{Pre-continuation recorded proof scope.} The actual order-four specialization lattice forces an invariant fibre Chern class to be 2nq. Direct contraction by the two printed translation vectors gives 4eta1 and minus6eta2. The finite-cover norms give coinvariant defects 4n delta and minus3n delta; the cusp solid-torus vanishing and based balance force n=0. Injectivity of the local integral specialization then kills the central Chern class, and the actual dP6 exponential sequence trivializes the normalization pullback. No global Leray transgression is used in this argument}{the hypotheses, notation and definitions stated here; \cref{wbII:lem:picard-finite-norm}; \cref{wbII:eq:W:5.17i}; \cref{wbII:eq:W:5.17n}; \cref{wbII:eq:W:5.17o}; \cref{wbII:prop:pic0-pullback}}

\ For the standard threefold and every analytic line bundle
$L\in\operatorname{Pic}(X)$,
\[
 c_1(L|_F)=0,
 \qquad
 \eta^*i^*L\simeq\mathcal O_{\mathrm{dP}_6}.
\tag{W.41}\label{wbII:eq:W:5.17q}
\]
This proof is independent of the global Leray transgressions
$d_2^{0,q}$ and of the packaged conclusion $H^2(X;\mathbb Z)=0$; it uses
only the actual finite-filling lattices, their logarithmic overlap maps, the
based boundary relation, the local cusp solid torus, and the local integral
specialization theorem.
\end{proposition}



\begin{proof}
\noindent\textbf{Step 1.} The common invariant calculation in \cref{wb:thm:coinvariants} first gives
\[
 E=c_1(L|_F)=kq.
\]
Because $L$ extends across the order-four filling, its fibre class belongs
to the actual specialization lattice
\[
 \pi_2^*H^2(S_2;\mathbb Z)
 =\{a\gamma\eta_2+bq:a\equiv b\pmod2\}
\]
of \cref{wb:prop:finite-indices}; \cref{wb:lem:finite-filling-stalk-map}.  Taking
coefficient independence in the displayed invariant basis forces
$a=0$ and $b=k$; hence $k$ is even.  Write $E=2nq$.

\noindent\textbf{Step 2.} Contraction with the two affine translation vectors is an exact integral
calculation:
\[
 \iota_{v_1}(2q)=4\eta_1,
 \qquad
 \iota_{v_2}(2q)=-6\eta_2.
\tag{W.42}\label{wbII:eq:W:5.17r}
\]
Indeed $v_1=(1,2,-4,0)^t$, $v_2=(-1,-3,3,0)^t$ in the basis dual to
$(\gamma,u,w,\delta)$, while
$q=uw+6\gamma\delta$.  In the global coinvariants,
\[
 [\eta_1]=3[\delta],\qquad [\eta_2]=2[\delta].
\]
Applying \cref{wbII:lem:picard-finite-norm} with $m_1=3$, $m_2=4$ therefore
gives
\[
 \bar d_1=4n[\delta],\qquad
 \bar d_2=-3n[\delta].
\tag{W.43}\label{wbII:eq:W:5.17s}
\]
Now combine the balance relation \eqref{wbII:eq:W:5.17i} with the cusp vanishing \eqref{wbII:eq:W:5.17n}:
\[
 0=\bar d_1+\bar d_2+\bar d_0=n[\delta].
\]
The coinvariant group is torsion-free, so $n=0$ and hence $E=0$.

\noindent\textbf{Step 3.} Finally \eqref{wbII:eq:W:5.17o} and the injectivity of the local integral specialization in
\cref{mapsI:top:thm:p0-integral-specialization} imply
$c_1(L|_W)=0$.  Pulling back to the normalization gives
$c_1(\eta^*i^*L)=0$.  The injectivity of
$c_1:\operatorname{Pic}(\mathrm{dP}_6)\to
H^2(\mathrm{dP}_6;\mathbb Z)$ proved in
\cref{wbII:prop:pic0-pullback} yields the second assertion of \eqref{wbII:eq:W:5.17q}.
\end{proof}



\subsection{A local analytic obstruction to a cusp-only Picard argument}
\label{wbII:subsec:local-cusp-picard-obstruction}

The relative Picard space of the cusp family cannot by itself force the
normalised central restriction to vanish.  The obstruction is visible
directly in the toric model of \cref{wb:thm:candidate-cusp-geometry}.  We give
the complete factor-of-automorphy calculation, including the descent
cocycle and the sign of the generic Chern class.

Retain the cusp notation of
\cref{wb:prop:period-input-audit}; \cref{wb:thm:candidate-cusp-geometry}:
\[
 B_0=\begin{pmatrix}0&1\\-1&0\end{pmatrix},
 \qquad
 Z(s)=sB_0+C(t),
 \qquad
 C(t)=\begin{pmatrix}6\mu(t)&h(t)\\ b(t)-h(t)&\mu(t)\end{pmatrix},
 \qquad t=e^{2\pi i s},
\tag{W.61}\label{wbII:eq:W:5.17t}
\]
\[
 c_{\lambda}(t)=\exp(2\pi iC(t)\lambda),
 \qquad
 \exp(2\pi iZ(s)\lambda)=t^{B_0\lambda}c_\lambda(t)
 \quad(\lambda\in\bar\Lambda\simeq\mathbb Z^2),
\tag{W.62}\label{wbII:eq:W:5.17u}
\]
and
\[
 \Phi_\lambda(x,t)=(t^{B_0\lambda}x,t),
 \qquad
 \Psi_\lambda(p)=(c_\lambda(t(p)),1)\cdot\Phi_\lambda(p),
 \qquad
 \Psi_\lambda\Psi_\mu=\Psi_{\lambda+\mu}.
\tag{W.63}\label{wbII:eq:W:5.17v}
\]
Write \(E_v\), \(v\in\mathbb Z^2\), for the toric prime divisors of
\(Y_\epsilon\), so that \(E_0\simeq\mathrm{dP}_6\), and let
\[
 \pi:Y_\epsilon\longrightarrow N_0=Y_\epsilon/\bar\Lambda,
 \qquad
 \nu=\pi|_{E_0}:E_0\longrightarrow W
\]
be the quotient and the normalization of the central fibre.
Put
\[
 N'=\mathbb Z^2\oplus\mathbb Z,
 \qquad
 M'=\operatorname{Hom}(N',\mathbb Z).
\]
For \(m=(\ell,k)\in M'=\mathbb Z^2\oplus\mathbb Z\), fix the character
convention
\[
 \chi^m(x,t)=x^\ell t^k,
 \qquad
 \operatorname{div}(\chi^m)=
 \sum_{v\in\mathbb Z^2}\langle m,(v,1)\rangle E_v.
\tag{W.64}\label{wbII:eq:W:5.17w}
\]

\begin{proposition}[an explicit nonzero local cusp restriction]
\label{wbII:prop:local-cusp-picard-obstruction}
\StatementRecord{wbII:prop:local-cusp-picard-obstruction}{proved}{Historical events: GPT-6 Astra (OpenAI), ChatGPT Pro mode, effort Pro, S6\_09\_PAPERS, 17 November 2026: targeted proof check. Also checked in this writing. Pre-continuation writing events: GPT-6 Astra (OpenAI), ChatGPT Pro mode, effort Pro, S6\_10\_PAPERS, 18 November 2026: targeted proof check; used and rechecked in this run. GPT-6 Astra (OpenAI), ChatGPT Pro mode, effort Pro, S6\_10\_PAPERS, 18 November 2026 05:35:21+00:00: used and rechecked in this run. Continuation metadata events: GPT-6 Astra (OpenAI), ChatGPT Pro mode, effort Pro, S6\_10\_PAPERS continuation, 19 November 2026: used and rechecked in this run; status-vocabulary reconciliation; metadata reconciliation only; no new mathematical check. \emph{Pre-continuation recorded proof scope.} Re-derived local finiteness on the unimodular triangle charts, the signed pullback divisor identity, the symmetric 2PC holomorphic correction and its complete multiplication cocycle. The sheaf multiplication arrow and quotient descent are exact. The six coefficients (1,-6,-5,1,-6,-5) give degrees (-12,2,0,-12,2,0). The inverse factor convention yields the integral exponential cocycle with alternating form exactly -2q. The cusp geometry dependency was reopened at wb:thm:candidate-cusp-geometry}{the hypotheses, notation and definitions stated here; \cref{wbII:eq:W:5.17t}; \cref{wbII:eq:W:5.17ad}; \cref{wbII:eq:W:5.17u}; \cref{wbII:eq:W:5.17v}; \cref{wbII:eq:W:5.17ai}; \cref{wbII:eq:W:5.17ag}; \cref{wbII:eq:W:5.17aj}; \cref{wbII:eq:W:5.17ak}; \cref{wbII:eq:W:5.17al}; \cref{wbII:eq:W:5.17am}; \cref{wbII:eq:W:5.17an}; \cref{wbII:eq:W:5.17ao}; \cref{wbII:eq:W:5.17ae}; \cref{wbII:eq:W:5.17ax}}

\ Let
\[
 H=\begin{pmatrix}1&0\\0&-6\end{pmatrix},
 \qquad
 P=B_0^{\mathsf t}H=\begin{pmatrix}0&6\\1&0\end{pmatrix},
 \qquad
 Q(v)=v^{\mathsf t}Hv=v_1^2-6v_2^2.
\tag{W.65}\label{wbII:eq:W:5.17x}
\]
Then
\[
 \mathcal D_Q=\sum_{v\in\mathbb Z^2}Q(v)E_v
\tag{W.66}\label{wbII:eq:W:5.17y}
\]
is a locally finite Cartier divisor on \(Y_\epsilon\).
The line bundle \(\mathcal O_{Y_\epsilon}(\mathcal D_Q)\) has an explicit
\((\bar\Lambda,\Psi)\)-equivariant structure and descends to a holomorphic
line bundle \(\mathcal L_Q\) on \(N_0\).  If
\[
 (a_1,a_2,a_3,a_4,a_5,a_6)
 =(e_1,e_2,e_2-e_1,-e_1,-e_2,e_1-e_2),
 \qquad C_i=E_0\cap E_{a_i},
\]
then
\[
 \nu^*(\mathcal L_Q|_W)
 \simeq
 \mathcal O_{E_0}\!\left(\sum_{i=1}^6Q(a_i)C_i\right),
\qquad
 \left(\deg_{C_i}\nu^*(\mathcal L_Q|_W)\right)_{i=1}^6
 =(-12,2,0,-12,2,0).
\tag{W.67}\label{wbII:eq:W:5.17z}
\]
In particular \(\nu^*(\mathcal L_Q|_W)\) is nontrivial.

On every nonzero fibre
\[
 F_t=\mathbb C^2/(Z(s)\mathbb Z^2+\mathbb Z^2),
\]
the first Chern class of \(\mathcal L_Q|_{F_t}\), in the ordered lattice
basis
\((\widehat\gamma,\widehat u,\widehat w,\widehat\delta)\) indexing the
columns of \([Z\mid I_2]\), is exactly
\[
 c_1(\mathcal L_Q|_{F_t})=-2q,
 \qquad
 q=u\wedge w+6\gamma\wedge\delta.
\tag{W.68}\label{wbII:eq:W:5.17aa}
\]
\end{proposition}



\begin{proof}
\noindent\textbf{Step 1. Local finiteness and Cartier status.}
The fan of \(Y\) consists of the cones over the cells of the locally finite
\(A_2\)-triangulation at height one.  If
\(U_\sigma\simeq\mathbb C^3\) is the affine chart belonging to a triangle
\(\sigma=\{v_0,v_1,v_2\}\), then \(U_\sigma\) meets precisely
\(E_{v_0},E_{v_1},E_{v_2}\).  Indeed, by the orbit--cone
correspondence, \(U_\sigma\cap E_v\ne\varnothing\) exactly when
\(\mathbb R_{\geq0}(v,1)\) is a face of the cone over \(\sigma\).
Thus \(\{E_v\}_{v\in\mathbb Z^2}\) is locally finite and
\(\mathcal D_Q\) is a well-defined locally finite Weil divisor.

Every triangle of the fan is unimodular, as proved in
\cref{wb:thm:candidate-cusp-geometry}.  If \(z_0,z_1,z_2\) are the
characters dual to
\((v_0,1),(v_1,1),(v_2,1)\), then
\[
 \mathcal D_Q|_{U_\sigma}
 =\operatorname{div}\!
   \left(z_0^{Q(v_0)}z_1^{Q(v_1)}z_2^{Q(v_2)}\right).
\tag{W.69}\label{wbII:eq:W:5.17ab}
\]
Negative exponents merely make the displayed monomial meromorphic, as
permitted for a local equation of a Cartier divisor.  Hence
\(\mathcal D_Q\) is Cartier on every chart and therefore on
\(Y_\epsilon\).

\noindent\textbf{Step 2. The signed divisor pullback identity.}
For \(\lambda\in\bar\Lambda=\mathbb Z^2\), set
\[
 b_\lambda=B_0\lambda,
 \qquad
 \ell_\lambda=2Hb_\lambda,
 \qquad
 k_\lambda=Q(b_\lambda),
 \qquad
 m_\lambda=(\ell_\lambda,k_\lambda)\in M'.
\tag{W.70}\label{wbII:eq:W:5.17ac}
\]
The fan automorphism sends \(E_v\) to \(E_{v+b_\lambda}\).
Consequently \(\Phi_\lambda^*E_w=E_{w-b_\lambda}\), and
\begin{align*}
 \Phi_\lambda^*\mathcal D_Q-\mathcal D_Q
 &=\sum_{v\in\mathbb Z^2}
   \bigl(Q(v+b_\lambda)-Q(v)\bigr)E_v\\
 &=\sum_{v\in\mathbb Z^2}
   \bigl(2v^{\mathsf t}Hb_\lambda+
         b_\lambda^{\mathsf t}Hb_\lambda\bigr)E_v
 =\operatorname{div}(\chi^{m_\lambda}).
\end{align*}
The \(t\)-dependent horizontal torus translation
\(p\mapsto(c_\lambda(t(p)),1)\cdot p\) preserves every
torus-invariant prime divisor \(E_v\).  Hence
\[
 \boxed{\ 
 \Psi_\lambda^*\mathcal D_Q
 =\mathcal D_Q+\operatorname{div}(\chi^{m_\lambda})\ }.
\tag{W.71}\label{wbII:eq:W:5.17ad}
\]

\noindent\textbf{Step 3. The complete \(\Psi\)-cocycle.}
The matrices in \eqref{wbII:eq:W:5.17t} and \eqref{wbII:eq:W:5.17x} give
\[
 PC(t)=
 \begin{pmatrix}
 6(b(t)-h(t))&6\mu(t)\\
 6\mu(t)&h(t)
 \end{pmatrix},
\tag{W.72}\label{wbII:eq:W:5.17ae}
\]
which is symmetric.  Thus
\[
 A(t):=2PC(t)
\tag{W.73}\label{wbII:eq:W:5.17af}
\]
is a symmetric holomorphic matrix on the full cusp disc.  Define the
holomorphic units \(h_\lambda\) and meromorphic functions \(f_\lambda\) by
\[
 h_\lambda(t)=
 \exp\!\left(\pi i\lambda^{\mathsf t}A(t)\lambda\right),
 \qquad
 f_\lambda=h_\lambda(t)\chi^{m_\lambda}.
\tag{W.74}\label{wbII:eq:W:5.17ag}
\]
Equation \eqref{wbII:eq:W:5.17ad} gives
\[
 \operatorname{div}(f_\lambda)
 =\Psi_\lambda^*\mathcal D_Q-\mathcal D_Q.
\tag{W.75}\label{wbII:eq:W:5.17ah}
\]

The cocharacter map underlying \(\Phi_\mu\) is
\[
 \varphi_\mu(y,y_3)=(y+y_3b_\mu,y_3),
\]
so its pullback on characters is
\[
 \varphi_\mu^*(\ell,k)
 =(\ell,k+\ell^{\mathsf t}b_\mu).
\tag{W.76}\label{wbII:eq:W:5.17ai}
\]
The quadratic identity for \(Q\) gives
\[
 m_{\lambda+\mu}=m_\mu+\varphi_\mu^*m_\lambda.
\tag{W.77}\label{wbII:eq:W:5.17aj}
\]
Indeed, the first coordinates add, while the last coordinate on the
right is
\[
 Q(b_\mu)+Q(b_\lambda)+2b_\lambda^{\mathsf t}Hb_\mu
 =Q(b_\lambda+b_\mu).
\]
Using \eqref{wbII:eq:W:5.17u}, \eqref{wbII:eq:W:5.17v}, and \eqref{wbII:eq:W:5.17ai}, one obtains
\begin{align*}
 \Psi_\mu^*\chi^{m_\lambda}
 &=\chi^{m_\lambda}(c_\mu(t),1)
   \chi^{\varphi_\mu^*m_\lambda}\\
 &=\exp\!\left(2\pi i\ell_\lambda^{\mathsf t}C(t)\mu\right)
   \chi^{\varphi_\mu^*m_\lambda}\\
 &=\exp\!\left(2\pi i\lambda^{\mathsf t}A(t)\mu\right)
   \chi^{\varphi_\mu^*m_\lambda}.
\tag{W.78}\label{wbII:eq:W:5.17ak}
\end{align*}
The last equality is exact because
\[
 \ell_\lambda^{\mathsf t}C\mu
 =(2HB_0\lambda)^{\mathsf t}C\mu
 =2\lambda^{\mathsf t}B_0^{\mathsf t}HC\mu
 =2\lambda^{\mathsf t}PC\mu
 =\lambda^{\mathsf t}A\mu.
\]
Since \(A=A^{\mathsf t}\), \eqref{wbII:eq:W:5.17ag}, \eqref{wbII:eq:W:5.17aj}, and \eqref{wbII:eq:W:5.17ak} imply
\begin{align*}
 f_\mu\,\Psi_\mu^*f_\lambda
 &=h_\mu h_\lambda
   \exp(2\pi i\lambda^{\mathsf t}A\mu)
   \chi^{m_\mu+\varphi_\mu^*m_\lambda}\\
 &=\exp\!\left(
   \pi i(\lambda+\mu)^{\mathsf t}A(\lambda+\mu)\right)
   \chi^{m_{\lambda+\mu}}
 =f_{\lambda+\mu}.
\tag{W.79}\label{wbII:eq:W:5.17al}
\end{align*}

Use the divisor-sheaf convention
\[
 \mathcal O(D)(U)=
 \{g\in\mathcal M_Y(U):
   \operatorname{div}(g)+D|_U\geq0\}.
\]
If \(D'=D+\operatorname{div}(f)\), multiplication by \(f\) maps
\(\mathcal O(D')\) isomorphically to \(\mathcal O(D)\).  Thus
\eqref{wbII:eq:W:5.17ad} and \eqref{wbII:eq:W:5.17ag} define
\[
 \theta_\lambda:
 \Psi_\lambda^*\mathcal O(\mathcal D_Q)
 \xrightarrow{\ \times f_\lambda\ }
 \mathcal O(\mathcal D_Q).
\tag{W.80}\label{wbII:eq:W:5.17am}
\]
Because
\((\Psi_\lambda\Psi_\mu)^*=\Psi_\mu^*\Psi_\lambda^*\),
equation \eqref{wbII:eq:W:5.17al} is exactly
\[
 \theta_{\lambda+\mu}
 =\theta_\mu\circ\Psi_\mu^*\theta_\lambda.
\tag{W.81}\label{wbII:eq:W:5.17an}
\]
This proves a complete equivariance datum, not merely invariance of the
divisor class.

\noindent\textbf{Step 4. Analytic descent.}
By \cref{wb:thm:candidate-cusp-geometry}, the action on \(Y_\epsilon\) is
free and properly discontinuous.  Every point therefore has an open
neighbourhood \(U\) whose distinct translates are disjoint and on which
\(\pi|_U\) is a biholomorphism onto \(\pi(U)\).  Transport
\(\mathcal O(\mathcal D_Q)|_U\) to \(\pi(U)\), and on overlaps arising
from translates glue by \eqref{wbII:eq:W:5.17am}.  Equation \eqref{wbII:eq:W:5.17an} is precisely the
triple-overlap condition.  The resulting holomorphic line bundle
\(\mathcal L_Q\) on \(N_0\) satisfies
\[
 \pi^*\mathcal L_Q
 \simeq\mathcal O_{Y_\epsilon}(\mathcal D_Q).
\tag{W.82}\label{wbII:eq:W:5.17ao}
\]

\noindent\textbf{Step 5. The normalized central restriction.}
Let \(j_0:E_0\hookrightarrow Y_\epsilon\).  Since
\(\nu=\pi\circ j_0\), \eqref{wbII:eq:W:5.17ao} gives
\[
 \nu^*(\mathcal L_Q|_W)
 =j_0^*\pi^*\mathcal L_Q
 \simeq\mathcal O_{E_0}(\mathcal D_Q|_{E_0}).
\tag{W.83}\label{wbII:eq:W:5.17ap}
\]
The only \(E_v\), \(v\ne0\), meeting \(E_0\) are those indexed by the
six \(a_i\) in the statement, and the intersections are the six toric
\((-1)\)-curves \(C_i\).  Since \(Q(0)=0\), there is no
\(E_0|_{E_0}\) term, so
\[
 \mathcal D_Q|_{E_0}
 =\sum_{i=1}^6Q(a_i)C_i,
 \qquad
 (Q(a_1),\ldots,Q(a_6))=(1,-6,-5,1,-6,-5).
\tag{W.84}\label{wbII:eq:W:5.17aq}
\]
The hexagonal intersection table is
\[
 C_i^2=-1,
 \qquad
 C_i\cdot C_{i\pm1}=1,
 \qquad
 C_i\cdot C_j=0
 \quad(j\notin\{i,i\pm1\}),
\tag{W.85}\label{wbII:eq:W:5.17ar}
\]
with indices modulo six.  Therefore
\[
 \deg_{C_i}\mathcal O_{E_0}
 \!\left(\sum_jQ(a_j)C_j\right)
 =-Q(a_i)+Q(a_{i-1})+Q(a_{i+1}),
\]
and direct substitution gives
\[
 \begin{array}{c|rrrrrr}
 i&1&2&3&4&5&6\\ \hline
 -Q(a_i)+Q(a_{i-1})+Q(a_{i+1})
 &-12&2&0&-12&2&0.
 \end{array}
\tag{W.86}\label{wbII:eq:W:5.17as}
\]
The degree on \(C_1=C_{e_1}\) is \(-12\), proving nontriviality.

\noindent\textbf{Step 6. The generic-fibre factor and its exact Chern class.}
Fix \(t\ne0\), choose \(s\) with \(t=e^{2\pi i s}\), and write
\(x=e^{2\pi i\zeta}\), \(\zeta\in\mathbb C^2\).
Equation \eqref{wbII:eq:W:5.17u} identifies the deck transformation indexed by
\(\lambda\in\mathbb Z^2\) with
\[
 \zeta\longmapsto\zeta+Z(s)\lambda
\]
modulo the horizontal translations
\(\zeta\mapsto\zeta+n\), \(n\in\mathbb Z^2\).  The matrix identities
\[
 HB_0=P^{\mathsf t},
 \qquad
 B_0^{\mathsf t}HB_0=PB_0
\tag{W.87}\label{wbII:eq:W:5.17at}
\]
give
\[
 \chi^{m_\lambda}(x,t)
 =x^{2P^{\mathsf t}\lambda}
  t^{\lambda^{\mathsf t}PB_0\lambda},
\qquad
 h_\lambda(t)
 =\exp(2\pi i\lambda^{\mathsf t}PC(t)\lambda).
\]
Since \(Z=sB_0+C\), it follows that
\[
 f_\lambda(\zeta)
 =\exp\!\left(
 2\pi i\lambda^{\mathsf t}PZ\lambda+
 4\pi i\lambda^{\mathsf t}P\zeta\right).
\tag{W.88}\label{wbII:eq:W:5.17au}
\]
The matrix \(PZ=sPB_0+PC\) is symmetric:
\[
 PB_0=\operatorname{diag}(-6,1),
\]
and \(PC\) is symmetric by \eqref{wbII:eq:W:5.17ae}.

Under the descent convention \eqref{wbII:eq:W:5.17am}, a pulled-back local section
written as \(\varphi(\zeta)\) in the canonical frame satisfies
\[
 \varphi(\zeta+Z\lambda)
 =f_\lambda(\zeta)^{-1}\varphi(\zeta).
\]
Thus the factor of automorphy for the lattice translation
\(\zeta\mapsto\zeta+Z\lambda+n\) is
\[
 a_{(\lambda,n)}(\zeta)
 =\exp(2\pi i g_{(\lambda,n)}(\zeta)),
\qquad
 g_{(\lambda,n)}(\zeta)
 =-\lambda^{\mathsf t}PZ\lambda
  -2\lambda^{\mathsf t}P\zeta.
\tag{W.89}\label{wbII:eq:W:5.17av}
\]
This is independent of \(n\) as a factor of automorphy: replacing
\(\zeta\) by \(\zeta+n\) changes \(g\) by the integer
\(-2\lambda^{\mathsf t}Pn\).

For \(l=(\lambda,n)\), \(l'=(\mu,m)\), the exponential connecting
cocycle is
\begin{align*}
 c(l,l')
 &=g_l(\zeta+Z\mu+m)+g_{l'}(\zeta)-g_{l+l'}(\zeta)\\
 &=-2\lambda^{\mathsf t}Pm\in\mathbb Z,
\tag{W.90}\label{wbII:eq:W:5.17aw}
\end{align*}
where the symmetry of \(PZ\) cancels the two mixed \(Z\)-terms.
The corresponding alternating integral form is
\[
 E(l,l')=c(l,l')-c(l',l)
 =-2\lambda^{\mathsf t}Pm+2\mu^{\mathsf t}Pn.
\tag{W.91}\label{wbII:eq:W:5.17ax}
\]
This is \(c_1(\mathcal L_Q|_{F_t})\).  In the ordered column basis of
\([Z\mid I_2]\), the cross-block of the class
\(q=u\wedge w+6\gamma\wedge\delta\) from
\((\widehat\gamma,\widehat u)\) to
\((\widehat w,\widehat\delta)\) is
\[
 \begin{pmatrix}
 q(\widehat\gamma,\widehat w)&
 q(\widehat\gamma,\widehat\delta)\\
 q(\widehat u,\widehat w)&
 q(\widehat u,\widehat\delta)
 \end{pmatrix}
 =\begin{pmatrix}0&6\\1&0\end{pmatrix}=P.
\tag{W.92}\label{wbII:eq:W:5.17ay}
\]
Hence \eqref{wbII:eq:W:5.17ax} is exactly \(-2q\), as asserted.
\end{proof}



\paragraph{the exact obstruction to the purely analytic route}

\label{wbII:rem:local-cusp-picard-obstruction}
The proposition proves that the local homomorphism
\[
 \operatorname{Pic}(N_0)\longrightarrow\operatorname{Pic}(E_0),
 \qquad
 \mathcal L\longmapsto\nu^*(\mathcal L|_W),
\]
is not zero.  Thus the toric cusp geometry, conductor descent on \(W\),
or local relative-Picard representability cannot alone prove the desired
global vanishing
\[
 \nu^*i^*:\operatorname{Pic}(X)\longrightarrow
 \operatorname{Pic}(E_0).
\]
There is no contradiction with
\cref{wbII:prop:leray-edge-restriction,wbII:prop:direct-boundary-picard}: the local
bundle above has nonzero generic class \(-2q\), whereas the two global
proofs use the order-three and order-four fillings and their common
boundary relation to exclude every nonzero multiple of \(q\).

More generally, Bingener proves that if a holomorphic map is proper,
flat, and cohomologically flat in dimension zero, then its relative
Picard functor is represented by a \emph{not necessarily separated}
complex space
\cite[Introduction, p.~318; Satz~6.2, pp.~337--338]{Bingener1980}.
The same introduction explains that the exponential-sequence quotient
works directly under topological local triviality and cohomological
flatness in dimension one, but that passage to the quotient fails in
general.  The cusp map of \cref{wb:thm:candidate-cusp-geometry} is not
topologically locally trivial across \(0\).  Bingener's theorem therefore
supplies representability but neither separatedness nor the
component-specialization assertion needed here.

Consequently a purely analytic completion must add global compatibility
with both finite fillings and must control the boundary support function;
neither relative-Picard representability nor conductor descent supplies
those statements.  The later calculation in
\cref{wbII:prop:analytic-ah-boundary,wbII:lem:analytic-support-specialization,wbII:thm:analytic-all-picard} now supplies them: it first forces the generic
Appell--Humbert form to vanish using the two finite fillings, and then
proves that its mixed block is the polarization of the three normalized
central degrees.  Thus this proposition remains the sharp cusp-local
obstruction, while the global purely analytic route is complete.

\Uses{the hypotheses, notation and definitions stated here; \cref{wbII:prop:leray-edge-restriction,wbII:prop:direct-boundary-picard}; \cite[Introduction, p.~318; Satz~6.2, pp.~337--338]{Bingener1980}; \cref{wbII:prop:analytic-ah-boundary,wbII:lem:analytic-support-specialization,wbII:thm:analytic-all-picard}}



\subsection{An analytic completion of the all-Picard certificate}
\label{wbII:subsec:analytic-all-picard}

The local example in \cref{wbII:prop:local-cusp-picard-obstruction} shows that
the cusp neighbourhood alone cannot annihilate normalized central
restrictions.  We now close the stronger analytic route by retaining the
order-three and order-four fillings.  The proof has two independent parts.
First, logarithms of the actual \(\mathcal O^*\)-factors give the boundary
cocycle and its two finite-filling norm identities.  Second, a
standard-family support-function calculation reads every possible
central restriction directly from the mixed Appell--Humbert block on a
nearby torus.  No global Leray differential, nearby-cycle specialization,
or Chern-class evaluation on a cusp bounding chain occurs.

\begin{proposition}[analytic Appell--Humbert boundary cocycle]
\label{wbII:prop:analytic-ah-boundary}
\StatementRecord{wbII:prop:analytic-ah-boundary}{proved}{Historical events: GPT-6 Astra (OpenAI), ChatGPT Pro mode, effort Pro, S6\_09\_PAPERS, 17 November 2026: targeted proof check. Also checked in this writing. Pre-continuation writing events: GPT-6 Astra (OpenAI), ChatGPT Pro mode, effort Pro, S6\_10\_PAPERS, 18 November 2026: targeted proof check; used and rechecked in this run. Continuation metadata events: GPT-6 Astra (OpenAI), ChatGPT Pro mode, effort Pro, S6\_10\_PAPERS continuation, 19 November 2026: used and rechecked in this run; status-vocabulary reconciliation; metadata reconciliation only; no new mathematical check. \emph{Pre-continuation recorded proof scope.} Re-derived the real central-extension normalization, crossed character, principal changes, based meridian relation, finite norm identity and clockwise positive translation sign. The receiving repair passes to the punctured universal cover and computes the exact homomorphism sigma to the filled deck group, whose kernel is generated by q\textasciicircum{}m tau\_-v. The fibre-product central extension then compares to the global factors through a holomorphic frame logarithm, preserving all conjugation exponents. At the cusp the determinant-one basis change identifies the correct meridian and a single Stein-chart frame makes its winding zero. The explicit source conventions were opened in Schnell, Birkenhake-Lange and Cartan at the recorded loci}{the hypotheses, notation and definitions stated here; \cref{wbII:eq:W:5.17-A8}; \cref{wbII:eq:W:5.17av}; \cref{wbII:eq:W:5.17ax}; \cite[Chapter~1, \S2, pp.~5--8, especially
Theorem~2.1, p.~7]{BirkenhakeLange1999}; \cite[equations~(3.3)--(3.10), pp.~3--7]{SchnellLecture3}; \cref{wbII:eq:W:5.17-A12}; \cref{wbII:eq:W:5.17-A11}; \cref{wbII:eq:W:5.17-A13}; \cref{wbII:eq:W:5.17-A14a}; \cref{wbII:eq:W:5.17-A14c}; \cref{wbII:eq:W:5.17-A14}; \cref{wbII:eq:W:5.17-A14d}; \cref{wbII:eq:W:5.17-A14e}; \cref{wbII:eq:W:5.17-A15}; \cref{wbII:eq:W:5.17-A17}; \cite[\S2, fundamental theorem~B, p.~80]{Cartan1957}; \cref{wbII:eq:W:5.17-A19a}; \cref{wbII:eq:W:5.17-A16}; \cref{wbII:eq:W:5.17-A9}; \cref{wbII:eq:W:5.17-A20}}

\ Let \(P\subset B^\circ\) be the compact pair-of-pants core,
orient its boundary meridians clockwise, and retain the zero-section lifts
of \cref{mapsI:top:prop:geometric-pi1}.  Thus
\[
 \rho_1\rho_2\rho_0=1,
 \qquad
 \widehat\rho_1\widehat\rho_2\widetilde a_0=1,
\tag{W.93}\label{wbII:eq:W:5.17-A1}
\]
and the second relation has no terminal vertical translation.  For
\(L\in\operatorname{Pic}(X)\), let
\[
 E=c_1(L|_F)\in\operatorname{Alt}^2(\Lambda,\mathbb Z)
   =\bigwedge^2V
\]
be its fibrewise Appell--Humbert form.  Put
\(V_{\mathbb R}=\operatorname{Hom}(\Lambda,\mathbb R)\), and use
\[
 A_g=(T_g^{-1})^{\mathsf t},
 \qquad
 \langle T_g\alpha,\lambda\rangle
 =\langle\alpha,A_g^{-1}\lambda\rangle.
\tag{W.94}\label{wbII:eq:W:5.17-A2}
\]
There is a crossed homomorphism, unique up to a principal crossed
homomorphism,
\[
 d:\pi_1(P)\longrightarrow V_{\mathbb R},
 \qquad
 d(gh)=d(g)+T_gd(h),
 \qquad
 d(g)\sim d(g)+(T_g-I)a.
\tag{W.95}\label{wbII:eq:W:5.17-A3}
\]
If \(d_i=d(\rho_i)\), then
\[
 d_1+T_1d_2+T_1T_2d_0=0,
 \qquad
 \bar d_1+\bar d_2+\bar d_0=0
 \quad\text{in }V_G\otimes\mathbb R.
\tag{W.96}\label{wbII:eq:W:5.17-A4}
\]
For the finite filling of order \(m_j\), \(j=1,2\),
\[
 N_{T_j}d_j=\iota_{v_j}E,
 \qquad
 m_j\bar d_j=[\iota_{v_j}E],
 \qquad
 N_{T_j}=I+T_j+\cdots+T_j^{m_j-1},
\tag{W.97}\label{wbII:eq:W:5.17-A5}
\]
where \((\iota_vE)(\lambda)=E(v,\lambda)\).  At the cusp,
\[
 \langle d_0,\widehat\delta\rangle=0,
 \qquad
 \bar d_0=0\quad\text{in }V_G\otimes\mathbb R.
\tag{W.98}\label{wbII:eq:W:5.17-A6}
\]
These identities are obtained from holomorphic transition factors and the
argument principle.
\end{proposition}



\begin{proof}
\noindent\textbf{Step 1.} Choose an open Riemann-surface neighbourhood \(P^{\rm an}\Subset B^\circ\)
which deformation retracts onto \(P\), and pull the period family to its
universal cover.  Before quotienting by the fibre lattice, the resulting
space is
\[
 \widetilde{\mathcal J}
 =\widetilde P^{\rm an}\times\mathbb C^2.
\]
It is Stein and contractible, so the pullback of \(L\) is holomorphically
trivial.  In a nowhere-zero frame, write the factor belonging to a deck
transformation \(\alpha\) as
\[
 a_\alpha\in\mathcal O^*(\widetilde{\mathcal J}),
 \qquad
 a_{\alpha\beta}(x)=a_\alpha(\beta x)a_\beta(x).
\tag{W.99}\label{wbII:eq:W:5.17-A7}
\]
The space is simply connected, hence \(a_\alpha=e^{2\pi i h_\alpha}\) for
a holomorphic logarithm.  Therefore
\[
 c(\alpha,\beta)
 =h_\alpha(\beta x)+h_\beta(x)-h_{\alpha\beta}(x)\in\mathbb Z
\tag{W.100}\label{wbII:eq:W:5.17-A8}
\]
is independent of \(x\), and associativity makes \(c\) an integral group
two-cocycle.

\noindent\textbf{Step 2.} Let \(\tau_\lambda\) denote fibre translation by
\(\Pi(z)\lambda\), \(\lambda\in\Lambda\), and let \(\widehat g\) be the
zero-section lift of a base deck transformation.  Period equivariance gives
\[
 \widehat g\,\tau_\lambda\,\widehat g^{-1}
 =\tau_{A_g\lambda}.
\tag{W.101}\label{wbII:eq:W:5.17-A9}
\]
On fibre translations, the commutator of \eqref{wbII:eq:W:5.17-A8} is
\[
 E(\lambda,\mu)
 =c(\tau_\lambda,\tau_\mu)-c(\tau_\mu,\tau_\lambda).
\tag{W.102}\label{wbII:eq:W:5.17-A10}
\]
The four oriented sides of the \((\lambda,\mu)\)-parallelogram identify
this integer with the exponential connecting class of \(L|_F\); this is
the Appell--Humbert convention used in
\eqref{wbII:eq:W:5.17av}--\eqref{wbII:eq:W:5.17ax}, and agrees with the standard factor-of-automorphy
calculation \cite[Chapter~1, \S2, pp.~5--8, especially
Theorem~2.1, p.~7]{BirkenhakeLange1999}%
\cite[equations~(3.3)--(3.10), pp.~3--7]{SchnellLecture3}.
Conjugating the commutator by \(\widehat g\) yields
\[
 E(A_g\lambda,A_g\mu)=E(\lambda,\mu).
\tag{W.103}\label{wbII:eq:W:5.17-A11}
\]

\noindent\textbf{Step 3.} Push the logarithmic central extension defined by \(c\) out along
\(\mathbb Z\hookrightarrow\mathbb R\), and denote its central element of
weight \(r\) by \(\mathbf c(r)\).  Choose lifts \(U_\lambda\) of the fibre
translations satisfying the real Appell--Humbert normalization
\[
 U_\lambda U_\mu
 =\mathbf c\!\left(\frac12E(\lambda,\mu)\right)U_{\lambda+\mu}.
\tag{W.104}\label{wbII:eq:W:5.17-A12}
\]
To see that this normalization exists, start with arbitrary lifts.  The
difference between their real cocycle and \(E/2\) has zero commutator, so
it defines an abelian extension of the free abelian group \(\Lambda\) by
\(\mathbb R\).  Choose lifts of an integral basis and extend
multiplicatively; this splits that abelian extension and gives \eqref{wbII:eq:W:5.17-A12}.
If \(U'_\lambda\) is a second normalized system, then
\(U'_\lambda U_\lambda^{-1}\) is central, say
\(\mathbf c(\langle a,\lambda\rangle)\).  Comparing the two copies of
\eqref{wbII:eq:W:5.17-A12} proves additivity in \(\lambda\), so
\(a\in V_{\mathbb R}\).  Thus every two normalized systems differ by exactly
the character change used below.

\noindent\textbf{Step 4.} For \(g\in\pi_1(P)\), choose a lift \(G_g\) of \(\widehat g\) and define
\(d(g)\in V_{\mathbb R}\) by
\[
 G_gU_{A_g^{-1}\lambda}G_g^{-1}
 =\mathbf c\bigl(\langle d(g),\lambda\rangle\bigr)U_\lambda.
\tag{W.105}\label{wbII:eq:W:5.17-A13}
\]
The exponent is additive in \(\lambda\): conjugate \eqref{wbII:eq:W:5.17-A12} and use
\eqref{wbII:eq:W:5.17-A11}.  The lifts of \(gh\) and \(G_gG_h\) differ centrally, so
\[
\begin{aligned}
 \langle d(gh),\lambda\rangle
 &=\langle d(g),\lambda\rangle
   +\langle d(h),A_g^{-1}\lambda\rangle\\
 &=\langle d(g)+T_gd(h),\lambda\rangle.
\end{aligned}
\]
This proves the crossed-homomorphism law.  Replacing
\(U_\lambda\) by
\(\mathbf c(\langle a,\lambda\rangle)U_\lambda\) changes \eqref{wbII:eq:W:5.17-A13} by
\((T_g-I)a\).  Replacing \(G_g\) by a central multiple does not change
\eqref{wbII:eq:W:5.17-A13}, so this proves the stated independence and uniqueness.  Applying
the crossed law to
\eqref{wbII:eq:W:5.17-A1} proves the first identity in \eqref{wbII:eq:W:5.17-A4}; passage to coinvariants
proves the second.

\noindent\textbf{Step 5.} For the order-\(m_j\) filling, distinguish every covering level.  Put
\[
 \widetilde J_j=\Delta_j\times\mathbb C^2,\qquad
 r_j:\widetilde J_j\longrightarrow J'_j
\]
for the lattice quotient, whose deck translations are \(\tau_\lambda\).
On the vector cover define
\[
 Q_{j,\log}:=\widetilde g_j^{\,v_j/m_j},\qquad
 Q_{j,\mathrm{lin}}:=\widetilde g_j^{\,0}.
\]
They satisfy
\[
\begin{aligned}
 Q_{j,\log}\tau_\lambda Q_{j,\log}^{-1}
  &=\tau_{A_j\lambda},&
 (Q_{j,\log})^{m_j}&=\tau_{v_j},\\
 Q_{j,\mathrm{lin}}\tau_\lambda Q_{j,\mathrm{lin}}^{-1}
  &=\tau_{A_j\lambda},&
 (Q_{j,\mathrm{lin}})^{m_j}&=1.
\end{aligned}
\tag{W.106}\label{wbII:eq:W:5.17-A14}
\]
The positive sign is fixed by
\(s_j\circ g_j=e^{-2\pi i/m_j}s_j\) and (2.8f): after exactly
\(m_j\) clockwise iterations, the logarithmic vector-cover lift translates
by \(+v_j\).

\noindent\textbf{Step 6.} Let \(q_{j,\log},q_{j,\mathrm{lin}}\) be the descended automorphisms of \(J'_j\),
and put \(J_j^{\prime *}=J'_j|_{\Delta_j^*}\).  Both descended automorphisms have
order dividing \(m_j\), because \(\tau_{v_j}\) is trivial after the lattice
quotient; each covers \(g_j\), whose order is exactly \(m_j\), so each
descended order is exactly \(m_j\).
The branch-independent torus translation
\[
 h_j:J_j^{\prime *}\longrightarrow J_j^{\prime *}
\]
constructed in (2.9d) satisfies, on this quotient level,
\[
 h_jq_{j,\log}h_j^{-1}=q_{j,\mathrm{lin}}.
\]
Define
\[
\begin{aligned}
 p_{j,\log}:J_j^{\prime *}&\longrightarrow
  N_j^*:=J_j^{\prime *}/\langle q_{j,\log}\rangle,\\
 p_{j,\mathrm{lin}}:J_j^{\prime *}&\longrightarrow
 J_j^*:=J_j^{\prime *}/\langle q_{j,\mathrm{lin}}\rangle .
\end{aligned}
\]
The conjugacy induces a biholomorphism
\[
 \bar h_j:N_j^*\longrightarrow J_j^*,\qquad
 p_{j,\mathrm{lin}}\circ h_j=\bar h_j\circ p_{j,\log}.
\]
The gluing of \(X\) identifies
\(L|_{N_j^*}\simeq\bar h_j^*(L|_{J_j^*})\); hence the exactly typed
line-bundle square gives
\[
 p_{j,\log}^*(L|_{N_j^*})
 \simeq h_j^*p_{j,\mathrm{lin}}^*(L|_{J_j^*}).
\tag{W.107}\label{wbII:eq:W:5.17-A14a}
\]

\noindent\textbf{Step 7.} For the factor extension return to the filled vector cover.  Let
\[
 P_{j,\log}:J'_j\longrightarrow N_j,\qquad
 \pi_j^{\mathrm{vec}}:=P_{j,\log}\circ r_j:
 \widetilde J_j\longrightarrow N_j .
\]
Here \(P_{j,\log}\) is the full quotient and
\(P_{j,\log}|_{J_j^{\prime *}}=p_{j,\log}\).
The deck group is
\[
 \Gamma_j^{\log}
 =\langle\tau_\lambda,Q_{j,\log}:\lambda\in\Lambda\rangle.
\]
The pullback \((\pi_j^{\mathrm{vec}})^*(L|_{N_j})\) is trivial on the
contractible Stein manifold \(\widetilde J_j\); its factors define an
integral central extension, which we push out along
\(\mathbb Z\hookrightarrow\mathbb R\).  Over \(\Delta_j^*\), the
line-bundle square \eqref{wbII:eq:W:5.17-A14a} identifies each fibre with the global factor
extension by a translation, so its commutator form is \(E\).  That form is
integer-valued and locally constant on the connected disc \(\Delta_j\);
therefore it is \(E\) on the filled cover as well.  Choose vertical lifts
\(U_\lambda^{(j)}\) normalized by
\[
 U_\lambda^{(j)}U_\mu^{(j)}
 =\mathbf c\!\left(\frac12E(\lambda,\mu)\right)
  U_{\lambda+\mu}^{(j)}.
\]
For \(\alpha\in\Gamma_j^{\log}\), define \(A_\alpha,T_\alpha\) by
\[
 \alpha\tau_\lambda\alpha^{-1}=\tau_{A_\alpha\lambda},\qquad
 T_\alpha=(A_\alpha^{-1})^{\mathsf t}.
\]
For a lift \(\widetilde G_\alpha\), define
\(\widetilde d_j(\alpha)\in V_{\mathbb R}\) through
\[
 \widetilde G_\alpha U_{A_\alpha^{-1}\lambda}^{(j)}
 \widetilde G_\alpha^{-1}
 =\mathbf c\!\left(\langle\widetilde d_j(\alpha),\lambda\rangle\right)
  U_\lambda^{(j)}.
\tag{W.108}\label{wbII:eq:W:5.17-A14c}
\]
Indeed, before naming \(\widetilde d_j(\alpha)\), the quotient of the two
sides of \eqref{wbII:eq:W:5.17-A14c} is central because both cover the same translation.
Conjugating the normalized product and using
\[
 E(A_\alpha\lambda,A_\alpha\mu)=E(\lambda,\mu)
\]
shows that its exponent \(r_\alpha(\lambda)\) satisfies
\[
 r_\alpha(\lambda+\mu)=r_\alpha(\lambda)+r_\alpha(\mu).
\]
Thus it is the pairing with a unique element of \(V_{\mathbb R}\), as
asserted.  Since \(A_{\alpha\beta}=A_\alpha A_\beta\), choose the lift of
\(\alpha\beta\) to be \(\widetilde G_\alpha\widetilde G_\beta\); any other
choice differs centrally.  Then
\[
\begin{aligned}
 &\widetilde G_\alpha\widetilde G_\beta
 U_{A_\beta^{-1}A_\alpha^{-1}\lambda}^{(j)}
 \widetilde G_\beta^{-1}\widetilde G_\alpha^{-1}\\
 &\quad=
 \mathbf c\!\left(
  \langle\widetilde d_j(\beta),A_\alpha^{-1}\lambda\rangle
  +\langle\widetilde d_j(\alpha),\lambda\rangle\right)
 U_\lambda^{(j)}.
\end{aligned}
\]
The defining dual action
\(\langle T_\alpha v,\lambda\rangle
=\langle v,A_\alpha^{-1}\lambda\rangle\) therefore proves the crossed law
\[
 \widetilde d_j(\alpha\beta)
 =\widetilde d_j(\alpha)+T_\alpha\widetilde d_j(\beta).
\tag{W.109}\label{wbII:eq:W:5.17-A14d}
\]

\noindent\textbf{Step 8.} Put
\[
 \widetilde J_j^*:=\Delta_j^*\times\mathbb C^2,
 \qquad
 r_j^*:=r_j|_{\widetilde J_j^*}:
 \widetilde J_j^*\longrightarrow J_j^{\prime *}
\]
and define
\[
\begin{aligned}
 \Phi_{j,\log}
  &:=p_{j,\log}\circ r_j^*,\\
 \Phi_{j,\mathrm{glob}}
  &:=p_{j,\mathrm{lin}}\circ h_j\circ r_j^*
   =\bar h_j\circ\Phi_{j,\log}.
\end{aligned}
\]
Both maps are invariant under \(\Gamma_j^{\log}\).  This is immediate for
every \(\tau_\lambda\), while
\[
\begin{aligned}
 \Phi_{j,\log}Q_{j,\log}
 &=p_{j,\log}q_{j,\log}r_j^*
  =\Phi_{j,\log},\\
 \Phi_{j,\mathrm{glob}}Q_{j,\log}
 &=p_{j,\mathrm{lin}}h_jq_{j,\log}r_j^*
  =p_{j,\mathrm{lin}}q_{j,\mathrm{lin}}h_jr_j^*
  =\Phi_{j,\mathrm{glob}}.
\end{aligned}
\]
Thus \(\widetilde J_j^*\) is the same \(\Gamma_j^{\log}\)-cover of the two
biholomorphic quotients.  Pulling the gluing isomorphism to this cover gives
\[
 \Theta_j:
 \Phi_{j,\log}^*(L|_{N_j^*})
 \xrightarrow{\ \sim\ }
 \Phi_{j,\mathrm{glob}}^*(L|_{J_j^*}).
\]
For the comparison of logarithms, pass one further time to the universal
cover
\[
 \widehat J_j^*=\widetilde{\Delta_j^*}\times\mathbb C^2
 \longrightarrow \widetilde J_j^*.
\]
The deck group over the punctured quotient is
\[
 \widehat\Gamma_j=
 \langle\tau_\lambda,q_j\mid
 q_j\tau_\lambda q_j^{-1}=\tau_{A_j\lambda}\rangle.
\]
The map from this universal cover to the filled vector cover induces the
explicit homomorphism
\[
 \sigma_j:\widehat\Gamma_j\longrightarrow\Gamma_j^{\log},
 \qquad \tau_\lambda\longmapsto\tau_\lambda,
 \quad q_j\longmapsto Q_{j,\log}.
\]
Its kernel is generated by
\(q_j^{m_j}\tau_{-v_j}\): this element is one deck turn of
\(\widetilde{\Delta_j^*}\to\Delta_j^*\), while its image on the
filled vector cover is the identity by \eqref{wbII:eq:W:5.17-A14}.
Indeed every word is \(\tau_\lambda q_j^r\); its filled image can be
identity only when \(r=m_jk\), and then only when
\(\lambda=-kv_j\).  Thus the kernel assertion follows directly.

\noindent\textbf{Step 9.} Let \(\mathscr E_j\) be the real central extension of
\(\Gamma_j^{\log}\) defined by the filled factors.  Pull it back along
\(\sigma_j\), explicitly as the fibre-product group
\[
 \sigma_j^*\mathscr E_j=
 \{(e,g):e\in\mathscr E_j,\ g\in\widehat\Gamma_j,
           \ p(e)=\sigma_j(g)\}.
\]
The global factor extension restricted to the based boundary group is
an extension of the same \(\widehat\Gamma_j\), with the generator
\(q_j\) identified with \(\widehat\rho_j\) through the fixed whisker.
The endpoint calculation below proves this precise identification.
Pull \(\Theta_j\) to \(\widehat J_j^*\).  This space is simply
connected, so the holomorphic unit comparing the two frames has a
holomorphic logarithm.  Substitution into \eqref{wbII:eq:W:5.17-A8}
changes the two integral cocycles by a coboundary.  Consequently it gives
an isomorphism from \(\sigma_j^*\mathscr E_j\) to the restricted global
extension, sending every \(\mathbf c(r)\) to the same
\(\mathbf c(r)\).  In the comparisons below,
\(U_\lambda^{(j)}\) means its paired lift
\((U_\lambda^{(j)},\tau_\lambda)\), and
\(\widetilde G_{Q_{j,\log}}\) means
\((\widetilde G_{Q_{j,\log}},q_j)\), when the isomorphism
\(\Theta_{j,*}\) is applied.  Their conjugation exponents remain the
ones in the filled extension.  This specifies every group and comparison
map used in the norm computation.

\noindent\textbf{Step 10.} Let \(\gamma_j^J:[0,1]\to J_j^*\) be the zero-section image of the
unwhiskered clockwise boundary meridian, and let \(\omega_j\) be the fixed
zero-section whisker from the global basepoint to \(\gamma_j^J(0)\).  Thus
the based total-space loop retained in \eqref{wbII:eq:W:5.17-A1} is
\[
 \widehat\rho_j=\omega_j*\gamma_j^J*\omega_j^{-1}.
\]
Choose \(x_j\in J_j^{\prime *}\) such that
\(p_{j,\mathrm{lin}}(h_j(x_j))=\gamma_j^J(0)\), put
\(x_{j,\mathrm{lin}}=h_j(x_j)\), and let
\(\widetilde\gamma_{j,\mathrm{lin}}\) be the
\(p_{j,\mathrm{lin}}\)-lift of \(\gamma_j^J\) beginning at
\(x_{j,\mathrm{lin}}\).  For the chosen clockwise generator,
\[
 \widetilde\gamma_{j,\mathrm{lin}}(1)
 =q_{j,\mathrm{lin}}(x_{j,\mathrm{lin}}).
\]
Set
\[
 \gamma_j^N:=\bar h_j^{-1}\circ\gamma_j^J,
 \qquad
 \widetilde\gamma_{j,\log}
 :=h_j^{-1}\circ\widetilde\gamma_{j,\mathrm{lin}}.
\]
Then the quotient square gives the point-valued identities
\[
\begin{aligned}
 p_{j,\log}\circ\widetilde\gamma_{j,\log}
  &=\gamma_j^N,\\
 \widetilde\gamma_{j,\log}(0)&=x_j,\\
 \widetilde\gamma_{j,\log}(1)&=q_{j,\log}(x_j).
\end{aligned}
\]
Choose \(\widetilde x_j\in\widetilde J_j^*\) with
\(r_j^*(\widetilde x_j)=x_j\), and let
\(\widetilde\gamma_{j,\mathrm{vec}}\) be the \(r_j^*\)-lift of
\(\widetilde\gamma_{j,\log}\) beginning at \(\widetilde x_j\).  On
adjacent sheets choose compatible logarithm branches \(\ell,\ell'\) with
\[
 \ell'(w)=\ell(g_j^{-1}w)-\frac{2\pi i}{m_j}.
\]
The calculation in (2.9e) gives the point-valued endpoint formula
\[
 \widetilde\gamma_{j,\mathrm{vec}}(1)
 =Q_{j,\log}(\widetilde x_j),
\]
excluding both \(Q_{j,\log}^{-1}(\widetilde x_j)\) and every
\(\tau_\lambda Q_{j,\log}(\widetilde x_j)\) with \(\lambda\ne0\).
Fibre loops remain \(\tau_\lambda\), since \(h_j\) is a fibre translation.

\noindent\textbf{Step 11.} The factor of \(L\) along the fixed lift of \(\omega_j\) identifies the
target factor over the local boundary segment with the globally based
factor.  Let \(G_{\rho_j}^{(j)}\) be the target lift of the local boundary
deck element whose whiskered image is the chosen global lift
\(G_{\rho_j}\).  Thus it covers \(\widehat\rho_j\) after whiskering and,
by \eqref{wbII:eq:W:5.17-A13}, satisfies
\[
 G_{\rho_j}^{(j)}U_{A_j^{-1}\lambda}(G_{\rho_j}^{(j)})^{-1}
 =\mathbf c(\langle d_j,\lambda\rangle)U_\lambda,
 \qquad d_j=d(\rho_j).
\]
Consequently, for some
\(a_j\in V_{\mathbb R}\) and \(b_j\in\mathbb R\),
\[
 \Theta_{j,*}(U_\lambda^{(j)})
 =\mathbf c(\langle a_j,\lambda\rangle)U_\lambda,
 \qquad
 \Theta_{j,*}(\widetilde G_{Q_{j,\log}})
 =\mathbf c(b_j)G_{\rho_j}^{(j)}.
\]
Here \(U_\lambda\) is the normalized global system in \eqref{wbII:eq:W:5.17-A12}, and
\(\mathbf c(b_j)\) disappears under conjugation.
Therefore
\[
\begin{aligned}
 \langle\widetilde d_j(Q_{j,\log}),\lambda\rangle
 +\langle a_j,\lambda\rangle
 &=\langle d_j,\lambda\rangle
   +\langle a_j,A_j^{-1}\lambda\rangle\\
 &=\langle d_j+T_ja_j,\lambda\rangle,
\end{aligned}
\]
so
\[
 \widetilde d_j(Q_{j,\log})=d_j+(T_j-I)a_j.
\tag{W.110}\label{wbII:eq:W:5.17-A14e}
\]

\noindent\textbf{Step 12.} For a translation choose
\(\widetilde G_{\tau_v}=U_v^{(j)}\).  The filled normalization above gives
\[
 U_v^{(j)}U_\lambda^{(j)}(U_v^{(j)})^{-1}
 =\mathbf c(E(v,\lambda))U_\lambda^{(j)},
\]
and hence
\[
 \widetilde d_j(\tau_v)=\iota_vE,\qquad
 (\iota_vE)(\lambda)=E(v,\lambda).
\tag{W.111}\label{wbII:eq:W:5.17-A14f}
\]
Here
\[
 A_{Q_{j,\log}}=A_j,\qquad T_{Q_{j,\log}}=T_j.
\]
Since \(N_{T_j}(T_j-I)=T_j^{m_j}-I=0\), equations
\eqref{wbII:eq:W:5.17-A14}, \eqref{wbII:eq:W:5.17-A14d}, and \eqref{wbII:eq:W:5.17-A14e} give
\[
\begin{aligned}
 N_{T_j}d_j
 &=N_{T_j}\widetilde d_j(Q_{j,\log})\\
 &=\widetilde d_j((Q_{j,\log})^{m_j})\\
 &=\widetilde d_j(\tau_{v_j})\\
 &=\iota_{v_j}E.
\end{aligned}
\tag{W.112}\label{wbII:eq:W:5.17-A15}
\]
Changing a line-bundle frame gives an isomorphic factor extension; changing
\(\widetilde G_\alpha\) by a central factor does not change conjugation; and
every second normalized vertical system differs by a character, changing
\(\widetilde d_j(\alpha)\) by \((T_\alpha-I)a\).  Finally,
\(\ell\mapsto\ell+2\pi in\) changes a vector-cover lift of \(h_j\) by
\[
 H_j\longmapsto H'_j=\tau_{nv_j}H_j.
\]
At the factor-extension level this means
\[
 \Theta'_{j,*}
 =\operatorname{Ad}(U_{nv_j})\circ\Theta_{j,*}.
\]
Consequently the comparison character changes by
\(n\iota_{v_j}E\), because
\[
 U_{nv_j}
 \bigl(\mathbf c(\langle a_j,\lambda\rangle)U_\lambda\bigr)
 U_{nv_j}^{-1}
 =
 \mathbf c(\langle a_j+n\iota_{v_j}E,\lambda\rangle)U_\lambda.
\]
Since \(A_jv_j=v_j\) and \(E\) is \(A_j\)-invariant,
\[
 \langle T_j(\iota_{v_j}E),\lambda\rangle
 =E(v_j,A_j^{-1}\lambda)
 =E(A_jv_j,\lambda)
 =E(v_j,\lambda),
\]
so this branch change leaves both \eqref{wbII:eq:W:5.17-A14e} and \eqref{wbII:eq:W:5.17-A15} unchanged.
The translation \(\tau_{nv_j}\) also commutes with
\(Q_{j,\mathrm{lin}}\) on the vector cover because \(A_jv_j=v_j\); on the
lattice quotient it is already the identity.  Hence it changes only the
central lift choice there.
The descended \(h_j,\bar h_j\), and the quotient square are branch
independent.  Thus \eqref{wbII:eq:W:5.17-A15} is choice-independent and proves \eqref{wbII:eq:W:5.17-A5},
including its sign.

\noindent\textbf{Step 13.} It remains to obtain \eqref{wbII:eq:W:5.17-A6} analytically.  Since
\(A_0\widehat\delta=\widehat\delta\), equation \eqref{wbII:eq:W:5.17-A13} says that
\[
 [G_0,U_{\widehat\delta}]
 =\mathbf c\bigl(\langle d_0,\widehat\delta\rangle\bigr).
\tag{W.113}\label{wbII:eq:W:5.17-A16}
\]
This logarithmic commutator is the clutching winding with the base loop
first.  Thus, if \(b\) and \(f\) denote the oriented base and fibre loops,
respectively,
\[
 \operatorname{wind}_{b\wedge f}(L)
  =\langle d_0,\widehat\delta\rangle,
 \qquad
 \operatorname{wind}_{\xi_{\widehat\delta}}(L)
 =-\langle d_0,\widehat\delta\rangle.
\tag{W.114}\label{wbII:eq:W:5.17-A17}
\]
The second equality uses the fibre-first convention
\(\xi_{\widehat\delta}=[S^1_{\widehat\delta}]\times[\rho_0]\), hence the
single sign change.
Use the actual unimodular cusp chart
\[
 (u,x_1,x_2)=\left(\frac{t}{x_1x_2},x_1,x_2\right),
 \qquad t=ux_1x_2.
\tag{W.115}\label{wbII:eq:W:5.17-A18}
\]
For \(0<\rho\ll1\), the boundary parametrization
\[
 (\theta,\phi)\longmapsto
 \bigl(\rho e^{i(\theta-\phi)},1,e^{i\phi}\bigr),
 \qquad t=\rho e^{i\theta},
\tag{W.116}\label{wbII:eq:W:5.17-A19}
\]
has base coordinate \(\theta\) and fibre translation
\(\widehat\delta\) in the \(\phi\)-direction.  The latter fixed-base loop
does \emph{not} itself bound the \(x_2\)-disc.  To identify the correct
meridian, put \(\alpha=\theta-\phi\), let \(a_u\) be the positively
oriented \(u\)-longitude, and let \(m_{x_2}\) be the positively oriented
\(x_2\)-meridian at fixed \(u\).  In the basis
\((a_u,m_{x_2})\), the base and fibre loops are
\[
 b=a_u,
 \qquad
 f=m_{x_2}-a_u,
 \qquad
 b\wedge f=a_u\wedge m_{x_2}.
\tag{W.117}\label{wbII:eq:W:5.17-A19a}
\]
The change-of-basis determinant is \(+1\), so this preserves the
base-first orientation in \eqref{wbII:eq:W:5.17-A17}.  Now \(m_{x_2}\) does bound the
analytic disc
\[
 z\longmapsto(\rho e^{i\alpha_0},1,z),
 \qquad |z|\le1.
\tag{W.118}\label{wbII:eq:W:5.17-A19b}
\]

\noindent\textbf{Step 14.} Choose a thin annulus \(A_u\) about \(|u|=\rho\), a disc \(D_{x_1}\)
about \(1\), and a disc \(D_{x_2}\) containing \(\{|x_2|\le1\}\), small
enough that
\[
 A_u\times D_{x_1}\times D_{x_2}
 \subset\{|ux_1x_2|<\epsilon\}.
\tag{W.119}\label{wbII:eq:W:5.17-A19c}
\]
Write \(U=A_u\times D_{x_1}\times D_{x_2}\) for this neighbourhood.
It is Stein and retracts onto \(S^1\), so \(H^2(U,\mathbb Z)=0\).
For completeness, the exponential sequence is
\[
0\longrightarrow\mathbb Z\longrightarrow\mathcal O_U
 \xrightarrow{\,h\mapsto e^{2\pi i h}\,}\mathcal O_U^*
 \longrightarrow1.
\]
Its exactness follows from the existence of a holomorphic logarithm locally
for every nonvanishing holomorphic function, and from the integral kernel
of the exponential. Cartan's theorem~B gives
\(H^1(U,\mathcal O_U)=H^2(U,\mathcal O_U)=0\)
\cite[\S2, fundamental theorem~B, p.~80]{Cartan1957}.
The long cohomology sequence therefore identifies
\(\operatorname{Pic}(U)=H^1(U,\mathcal O_U^*)\) with
\(H^2(U,\mathbb Z)=0\). Pull \(L|_{N_0}\) back to
the cusp cover and choose one holomorphic frame on this entire
neighbourhood.  Compute the clutching cocycle with this extending frame on
both sides of a cut in the \(u\)-longitude.  Its transition unit on
\(D_{x_2}\) is then identically \(F=1\), so by \eqref{wbII:eq:W:5.17-A19a} its
base--fibre commutator winding is
\[
 \frac1{2\pi i}\int_{|x_2|=1}\frac{dF}{F}=0.
\tag{W.120}\label{wbII:eq:W:5.17-A20}
\]
This computes the same alternating logarithmic cocycle as \eqref{wbII:eq:W:5.17-A16}.
Indeed, changing from the global Appell--Humbert frame to the extending
chart frame changes \(c\) by a group coboundary.  Put
\[
 g=\widehat g_0,
 \qquad
 h=\tau_{\widehat\delta}
\]
in the underlying deck group.  Equation \eqref{wbII:eq:W:5.17-A9} and
\(A_0\widehat\delta=\widehat\delta\) give \(gh=hg\).  If
\(c'=c+\delta\beta\) is the cocycle after a frame change, then
\[
 \bigl[c'(g,h)-c'(h,g)\bigr]
 -\bigl[c(g,h)-c(h,g)\bigr]
 =\beta(hg)-\beta(gh)=0.
\]
Thus the alternating cocycle is frame-invariant without assuming that its
central-extension lifts commute.  Equations \eqref{wbII:eq:W:5.17-A17}, \eqref{wbII:eq:W:5.17-A19a}, and
\eqref{wbII:eq:W:5.17-A20} therefore give
\(\langle d_0,\widehat\delta\rangle=0\).  This is a calculation in
\(\mathcal O^*\), not an evaluation of a Chern class on a bounding chain.
Finally, the matrices \(A_1,A_2\) fix
\(\widehat\delta\), so the evaluation is unchanged by every principal
change in \eqref{wbII:eq:W:5.17-A3}.  Since
\(V_G=\mathbb Z[\delta]\) and
\(\delta(\widehat\delta)=1\), its vanishing is equivalent to
\(\bar d_0=0\).  This proves \eqref{wbII:eq:W:5.17-A6}.
\end{proof}



\begin{lemma}[the toric support function detects central restriction]
\label{wbII:lem:analytic-support-specialization}
\StatementRecord{wbII:lem:analytic-support-specialization}{proved}{Historical events: GPT-6 Astra (OpenAI), ChatGPT Pro mode, effort Pro, S6\_09\_PAPERS, 17 November 2026: targeted proof check. Also checked in this writing. Pre-continuation writing events: GPT-6 Astra (OpenAI), ChatGPT Pro mode, effort Pro, S6\_10\_PAPERS, 18 November 2026: targeted proof check; used and rechecked in this run. Continuation metadata events: GPT-6 Astra (OpenAI), ChatGPT Pro mode, effort Pro, S6\_10\_PAPERS continuation, 19 November 2026: used and rechecked in this run; status-vocabulary reconciliation; metadata reconciliation only; no new mathematical check. \emph{Pre-continuation recorded proof scope.} Re-derived the integer-valued quadratic from the three opposite-pair boundary degrees, its polarization and every six-curve intersection degree. Subtracting its Cartier divisor trivializes the residual line on every compact component. The actual ray opens lie inside maximal charts; comparing each ray frame with the maximal-chart frame gives zero horizontal winding by continuation from the contractible component chart. The inverse-dual descent map and single-frame residual multiplier have the displayed direction and zero winding. The exponential connecting cocycle therefore has mixed entry -B(B0 lambda,m); invertibility of B0 forces B and all three degrees to vanish when the fibre class vanishes}{the hypotheses, notation and definitions stated here; \cref{wbII:prop:local-cusp-picard-obstruction}; \cite[Propositions~4.1.1--4.1.2, p.~171; Theorems~4.2.8,
p.~181 and~4.2.12, p.~184]{CoxLittleSchenck2011}; \cref{wbII:eq:W:5.17ab}; \cref{wbII:eq:W:5.17-A24}; \cref{wbII:eq:W:5.17-A30}; \cref{wbII:eq:W:5.17-A29}; \cite[Example~1.2.21, p.~31]{CoxLittleSchenck2011}; \cref{wbII:eq:W:5.17-A31}; \cref{wbII:eq:W:5.17-A32}; \cref{wbII:eq:W:5.17-A34b}; \cref{wbII:eq:analytic-residual-descent-factor}; \cref{wbII:eq:W:5.17am}; \cref{wbII:eq:W:5.17-A34c}; \cref{wbII:eq:W:5.17-A26}}

\ Let \(M\in\operatorname{Pic}(N_0)\), let
\(\pi:Y_\epsilon\to N_0\) be the cusp covering, and let
\(\nu:E_0\to W\) be the normalization.  If the Appell--Humbert form of
\(M|_{F_t}\), \(t\ne0\), is \(E_M\), then the three independent degrees
of \(\nu^*(M|_W)\) determine an integral symmetric form
\(B:\mathbb Z^2\times\mathbb Z^2\to\mathbb Z\) satisfying
\[
 E_M\bigl((\lambda,0),(0,m)\bigr)
 =-B(B_0\lambda,m)
 \qquad(\lambda,m\in\mathbb Z^2).
\tag{W.121}\label{wbII:eq:W:5.17-A21}
\]
Consequently
\[
 E_M=0
 \quad\Longrightarrow\quad
 \nu^*(M|_W)\simeq\mathcal O_{E_0}.
\tag{W.122}\label{wbII:eq:W:5.17-A22}
\]
\end{lemma}



\begin{proof}
\noindent\textbf{Step 1.} Put \(\mathcal M=\pi^*M\).  The identity
\(\pi\circ\Psi_\lambda=\pi\) gives \(\mathcal M\) its genuine
\(\Psi\)-equivariant descent structure.  In the notation of
\cref{wbII:prop:local-cusp-picard-obstruction}, order the six rays in the star
of the origin as
\[
 (r_1,r_2,r_3,r_4,r_5,r_6)
 =(e_1,e_2,e_2-e_1,-e_1,-e_2,e_1-e_2),
 \qquad C_i=E_0\cap E_{r_i}.
\tag{W.123}\label{wbII:eq:W:5.17-A23}
\]
The quotient identifies opposite conductor branches.  Therefore, for
\(A=\mathcal M|_{E_0}=\nu^*(M|_W)\),
\[
 d_i=\deg_{C_i}A,
 \qquad
 (\deg_{C_1}A,\ldots,\deg_{C_6}A)
 =(d_1,d_2,d_3,d_1,d_2,d_3).
\tag{W.124}\label{wbII:eq:W:5.17-A24}
\]

\noindent\textbf{Step 2.} Define the integer-valued quadratic function
\[
 a(x,y)
 =(d_2+d_3)\binom{x}{2}+d_3xy+(d_1+d_3)\binom{y}{2}.
\tag{W.125}\label{wbII:eq:W:5.17-A25}
\]
Its polarization is
\[
 B(v,w)=a(v+w)-a(v)-a(w),
 \qquad
 [B]=
 \begin{pmatrix}
  d_2+d_3&d_3\\ d_3&d_1+d_3
 \end{pmatrix}.
\tag{W.126}\label{wbII:eq:W:5.17-A26}
\]
The locally finite divisor
\[
 D_a=\sum_{v\in\mathbb Z^2}a(v)E_v
\tag{W.127}\label{wbII:eq:W:5.17-A27}
\]
is Cartier: on every maximal chart, the three height-one ray generators
form an integral basis, so the three assigned coefficients are the values
of one integral character.  In the standard toric convention the associated
support function has ray values \(-a(v)\); the divisor coefficients here are
the values \(a(v)\), consistent with
\(\operatorname{div}(\chi^m)=\sum_\rho\langle m,u_\rho\rangle D_\rho\)
\cite[Propositions~4.1.1--4.1.2, p.~171; Theorems~4.2.8,
p.~181 and~4.2.12, p.~184]{CoxLittleSchenck2011}.
This is the same chartwise Cartier argument as
in \eqref{wbII:eq:W:5.17ab}, and it applies to an arbitrary integer-valued \(a\).

\noindent\textbf{Step 3.} Direct substitution gives
\[
 (a(r_1),\ldots,a(r_6))
 =(0,0,d_2,d_2+d_3,d_1+d_3,d_1).
\tag{W.128}\label{wbII:eq:W:5.17-A28}
\]
Using \(C_i^2=-1\) and \(C_iC_{i\pm1}=1\), the degree vector of
\(\mathcal O(D_a)|_{E_0}\) is exactly the vector in \eqref{wbII:eq:W:5.17-A24}.  More
generally, restrict the principal identity
\(\operatorname{div}(t)=\sum_wE_w\) to \(E_v\).  Since \(E_v\) meets
precisely \(E_{v+r_i}\), it gives
\[
 \mathcal O(D_a)|_{E_v}
 \simeq
 \mathcal O_{E_v}\!\left(
   \sum_{i=1}^6(a(v+r_i)-a(v))C_{v,i}
 \right).
\tag{W.129}\label{wbII:eq:W:5.17-A29}
\]
The coefficient difference is
\[
 a(v+r_i)-a(v)=a(r_i)+B(v,r_i).
\tag{W.130}\label{wbII:eq:W:5.17-A30}
\]
The second summand in \eqref{wbII:eq:W:5.17-A30} is the divisor of the toric character
with exponent \(B(v,-)\).  Thus every restriction in \eqref{wbII:eq:W:5.17-A29} has the
same six degrees as \(\mathcal O(D_a)|_{E_0}\).

\noindent\textbf{Step 4.} Because \(B_0\bar\Lambda=\mathbb Z^2\), the deck action is transitive on
the components \(E_v\).  The genuine descent structure on \(\mathcal M\)
therefore transports \eqref{wbII:eq:W:5.17-A24}, with its ordered six boundary directions,
to every \(E_v\).  It follows that
\[
 \mathcal R:=\mathcal M\otimes\mathcal O(-D_a)
 \quad\text{satisfies}\quad
 \mathcal R|_{E_v}\simeq\mathcal O_{E_v}
 \quad(v\in\mathbb Z^2).
\tag{W.131}\label{wbII:eq:W:5.17-A31}
\]
Indeed, the six toric curves detect \(\operatorname{Pic}(\mathrm{dP}_6)\):
in the blow-up basis \(H,E_1,E_2,E_3\), vanishing of the degrees on the
three exceptional curves kills the three \(E_i\)-coefficients, and
vanishing on any strict transform \(H-E_i-E_j\) kills the remaining
\(H\)-coefficient.  Since
\(H^1(\mathrm{dP}_6,\mathcal O)=0\), equality of those degrees implies
isomorphism, not merely numerical equivalence.

\noindent\textbf{Step 5.} For \(b_\lambda=B_0\lambda\), the torus-translation part of
\(\Psi_\lambda\) preserves each toric prime divisor, while its fan part
translates the index by \(b_\lambda\).  Hence
\[
 \Psi_\lambda^*D_a-D_a
 =\sum_v(a(v+b_\lambda)-a(v))E_v
 =\operatorname{div}(\chi^{m_\lambda}),
\tag{W.132}\label{wbII:eq:W:5.17-A32}
\]
where
\[
 m_\lambda=(\ell_\lambda,k_\lambda),
 \qquad
 \ell_\lambda=B(b_\lambda,-),
 \qquad
 k_\lambda=a(b_\lambda).
\tag{W.133}\label{wbII:eq:W:5.17-A33}
\]
Thus multiplication by \(\chi^{m_\lambda}\) gives an individual
isomorphism
\(\Psi_\lambda^*\mathcal O(D_a)\simeq\mathcal O(D_a)\).
No cocycle assertion for these individual isomorphisms is needed: we use
them only to separate the horizontal winding of the genuine descent factor
of \(\mathcal M\).

\noindent\textbf{Step 6.} For a ray \(\mathbb R_{\ge0}(v,1)\), its toric open set in
\(Y_\epsilon\) is
\[
 U_v\simeq\Delta_\epsilon\times(\mathbb C^*)^2,
 \qquad E_v^\circ=\{0\}\times(\mathbb C^*)^2.
\tag{W.134}\label{wbII:eq:W:5.17-A34}
\]
Indeed \((e_1,0),(e_2,0),(v,1)\) is an integral basis and \(t\) is the
character dual to the last vector; compare the author preprint
\cite[Example~1.2.21, p.~31]{CoxLittleSchenck2011}.
The set \(U_v\) is Stein and retracts
onto \(E_v^\circ\).  By \eqref{wbII:eq:W:5.17-A31},
\(c_1(\mathcal R|_{U_v})=0\); the exponential sequence and Cartan's
theorem~B therefore trivialize \(\mathcal R|_{U_v}\).  Choose a global
nowhere-zero frame \(s_v\) of \(\mathcal R|_{E_v}\).  A frame on \(U_v\) can be
multiplied by a unit pulled back from \((\mathbb C^*)^2\) so that its
restriction to \(E_v^\circ\) agrees with \(s_v|_{E_v^\circ}\); denote the
resulting frame by \(e_v\).

\noindent\textbf{Step 7.} We also need to compare these frames on the generic torus.  The direction
of the toric affine-open inclusions is essential here.  If \(v,w\) are
adjacent vertices, choose a triangle \(\sigma=\{v,w,u\}\) containing them,
write \(z_v,z_w,z_u\) for the coordinates dual to its three height-one ray
generators, and put
\[
 S_\sigma:=U_\sigma\cap Y_\epsilon
 =\{(z_v,z_w,z_u)\in\mathbb C^3:
       |z_vz_wz_u|<\epsilon\}.
\tag{W.135}\label{wbII:eq:W:5.17-A34a}
\]
Because \(\rho_v<\sigma\), localization of the maximal affine chart gives
\begin{equation}
 U_{\rho_v}=D(z_wz_u)\subset U_\sigma,\qquad
 U_{\rho_w}=D(z_vz_u)\subset U_\sigma,\qquad
 U_{\rho_v}\cap U_{\rho_w}=D(z_vz_wz_u)=U_0.
\label{wbII:eq:analytic-toric-open-directions}
\end{equation}
Thus \(U_\sigma\) is not contained in \(U_{\rho_v}\cap U_{\rho_w}\), and
the coordinate discs in \(S_\sigma\) cannot be applied directly to the
ratio \(e_w/e_v\).

\noindent\textbf{Step 8.} The space \(S_\sigma\) is nevertheless the required comparison chart.  It
is the closed graph
\[
 \{(z_v,z_w,z_u,t)\in\mathbb C^3\times\Delta_\epsilon:
   t=z_vz_wz_u\}
\]
in the Stein space \(\mathbb C^3\times\Delta_\epsilon\), and hence is
Stein.  Radial scaling of the three \(z\)-coordinates contracts it to the
origin.  The exponential sequence and Cartan's theorem~B therefore give
\(\operatorname{Pic}(S_\sigma)=0\).  Choose a frame \(e_\sigma\) of
\(\mathcal R|_{S_\sigma}\).

\noindent\textbf{Step 9.} Let \(s_v\) be the global frame of
\(\mathcal R|_{E_v}\simeq\mathcal O_{E_v}\) used to normalize \(e_v\).
On \(U_{\rho_v}\cap Y_\epsilon\) put
\[
 A_{\sigma,v}=e_v/e_\sigma.
\]
At \(t=0\), this is the restriction to \((\mathbb C^*)^2\) of the
holomorphic unit
\[
 G_{\sigma,v}
 =\frac{s_v}{e_\sigma|_{E_v\cap S_\sigma}}
 \in\mathcal O^*(E_v\cap S_\sigma),
 \qquad E_v\cap S_\sigma=\{z_v=0\}\simeq\mathbb C^2.
\]
Consequently its winding on every horizontal cocharacter is zero at
\(t=0\).  Under
\(U_{\rho_v}\cap Y_\epsilon\simeq
\Delta_\epsilon\times(\mathbb C^*)^2\), that winding is an integer locally
constant in \(t\), and is therefore zero on every nearby nonzero fibre.
The same argument applies to \(A_{\sigma,w}\).  On their actual common
domain \(U_0\cap Y_\epsilon\),
\[
 e_w/e_v=A_{\sigma,w}/A_{\sigma,v},
\]
so this ratio has zero horizontal winding.  The one-skeleton of the
\(A_2\)-triangulation is connected, and multiplication along an edge path
now gives
\[
 \operatorname{wind}_m(e_w/e_v)=0
 \qquad(v,w,m\in\mathbb Z^2)
\tag{W.136}\label{wbII:eq:W:5.17-A34b}
\]
on every nonzero fibre.

\noindent\textbf{Step 10.} Let
\(\theta^{\mathcal M}_\lambda:\Psi_\lambda^*\mathcal M\to\mathcal M\)
be the genuine descent isomorphism, let
\(\mu^+_\lambda:\Psi_\lambda^*\mathcal O(D_a)\to\mathcal O(D_a)\)
be multiplication by \(\chi^{m_\lambda}\) from \eqref{wbII:eq:W:5.17-A32}, and define the
inverse-dual isomorphism
\[
 \mu^-_\lambda:=((\mu^+_\lambda)^{-1})^\vee:
 \Psi_\lambda^*\mathcal O(-D_a)\longrightarrow\mathcal O(-D_a),
 \qquad s\longmapsto\chi^{-m_\lambda}s.
\]
The ordinary dual of \(\mu^+_\lambda\) has the opposite arrow; the
inverse-dual above is the map required by the pullback tensor product.  Then
\[
 \theta^{\mathcal R}_\lambda
 :=\theta^{\mathcal M}_\lambda\otimes\mu^-_\lambda:
 \Psi_\lambda^*\mathcal R\longrightarrow\mathcal R.
\]
Since \(\Psi_\lambda(U_v)=U_{v+b_\lambda}\), define
\(r_{\lambda,v}\in\mathcal O^*(U_v)\) by the typed identity
\begin{equation}
 \theta^{\mathcal R}_\lambda
   (\Psi_\lambda^*e_{v+b_\lambda})
 =r_{\lambda,v}e_v.
\label{wbII:eq:analytic-residual-descent-factor}
\end{equation}
At \(t=0\) it extends to an isomorphism between two trivial line bundles
on the compact surface \(E_v\).  Its factor in the chosen global frames is
therefore a nonzero constant.  The winding of
\(r_{\lambda,v}(t,-)\) around every horizontal one-parameter subgroup is
an integer locally constant in \(t\); at \(t=0\) it is zero, hence it is
zero for every nearby \(t\ne0\).

\noindent\textbf{Step 11.} To express this isomorphism in one frame on the generic torus, put
\(q_{\lambda,v}=e_{v+b_\lambda}/e_v\).  Equation \eqref{wbII:eq:W:5.17-A34b} gives zero
horizontal winding for \(q_{\lambda,v}\); because
\(\Psi_\lambda|_{Y_t}\) is multiplication by a fixed torus point, it
induces the identity on horizontal \(H_1\), and
\(\Psi_\lambda^*q_{\lambda,v}\) also has zero winding.  Relative to the
single frame \(e_v\), the residual multiplier is therefore
\[
 h_{\lambda,v}
 =\frac{r_{\lambda,v}}{\Psi_\lambda^*q_{\lambda,v}},
 \qquad
 \operatorname{wind}_m(h_{\lambda,v})=0.
\tag{W.137}\label{wbII:eq:W:5.17-A34c}
\]

\noindent\textbf{Step 12.} On \(U_0\) let \(\mathbf1_a\) be the canonical toric frame of
\(\mathcal O(D_a)\), and put \(F_v=e_v\otimes\mathbf1_a\).  Equations
\cref{wbII:eq:analytic-residual-descent-factor} and \eqref{wbII:eq:W:5.17-A32} give
\[
 \theta^{\mathcal M}_\lambda(\Psi_\lambda^*F_v)
 =h_{\lambda,v}\chi^{m_\lambda}F_v.
\]
On a nonzero fibre, \(\chi^{m_\lambda}=x^{\ell_\lambda}t^{k_\lambda}\).
With the section-descent convention \eqref{wbII:eq:W:5.17am}, the factor of automorphy is
therefore
\[
 \alpha_\lambda=(h_{\lambda,v}\chi^{m_\lambda})^{-1}.
\]
The residual factor has zero horizontal winding by \eqref{wbII:eq:W:5.17-A34c}, and hence
\begin{equation}
 \operatorname{wind}_m(\alpha_\lambda)
 =-\ell_\lambda(m)=-B(B_0\lambda,m).
\label{wbII:eq:analytic-factor-horizontal-winding}
\end{equation}

\noindent\textbf{Step 13.} This winding is exactly the mixed Chern-class entry.  Pull \(F_v\) back to
the universal cover
\[
 p_t:\mathbb C^2\longrightarrow Y_t,
 \qquad \zeta\longmapsto(e^{2\pi i\zeta},t).
\]
The period indexed by \((\lambda,n)\in\mathbb Z^2\oplus\mathbb Z^2\)
acts by \(\zeta\mapsto\zeta+Z(s)\lambda+n\).  The horizontal factor is one,
because \(F_v\) is already defined on \(Y_t\).  Choose a logarithm
\(g_\lambda\in\mathcal O(\mathbb C^2)\) with
\[
 p_t^*\alpha_\lambda=e^{2\pi i g_\lambda}.
\]
For every horizontal period \(n\),
\[
 g_\lambda(\zeta+n)-g_\lambda(\zeta)
 =\operatorname{wind}_n(\alpha_\lambda).
\]
In the exponential connecting cocycle
\[
 c(l,l')=g_l(\zeta+l')+g_{l'}(\zeta)-g_{l+l'}(\zeta),
\]
take the horizontal logarithm to be zero and the logarithm for the combined
period \((\lambda,n)\) to be \(g_\lambda\).  Then
\[
 c((\lambda,0),(0,n))
 =\operatorname{wind}_n(\alpha_\lambda),
 \qquad
 c((0,n),(\lambda,0))=0.
\]
Therefore the alternating exponential connecting integer between
the vertical deck
translation indexed by \(\lambda\) and the horizontal loop indexed by
\(m\) is
\[
 E_M\bigl((\lambda,0),(0,m)\bigr)
 =-\ell_\lambda(m)
 =-B(B_0\lambda,m),
\]
which proves \eqref{wbII:eq:W:5.17-A21}.  This sign also recovers the explicit check in
\cref{wbII:prop:local-cusp-picard-obstruction}: for
\(a(v)=v^{\mathsf t}Hv\), one has \(B=2H\), so \eqref{wbII:eq:W:5.17-A21} is the
cross-block \(-2B_0^{\mathsf t}H=-2P\), namely \(-2q\).

\noindent\textbf{Step 14.} If \(E_M=0\), equation \eqref{wbII:eq:W:5.17-A21} and the unimodularity of \(B_0\) force
\(B=0\).  The matrix in \eqref{wbII:eq:W:5.17-A26} then gives
\(d_3=0\), \(d_2=0\), and \(d_1=0\).  Thus all six degrees of \(A\)
vanish, and the Picard detection argument following \eqref{wbII:eq:W:5.17-A31} gives
\(A\simeq\mathcal O_{E_0}\).  This is \eqref{wbII:eq:W:5.17-A22}.
\end{proof}



\begin{theorem}[purely analytic vanishing of the global restriction]
\label{wbII:thm:analytic-all-picard}
\StatementRecord{wbII:thm:analytic-all-picard}{proved}{Historical events: GPT-6 Astra (OpenAI), ChatGPT Pro mode, effort Pro, S6\_09\_PAPERS, 17 November 2026: targeted proof check. Also checked in this writing. Pre-continuation writing events: GPT-6 Astra (OpenAI), ChatGPT Pro mode, effort Pro, S6\_10\_PAPERS, 18 November 2026: targeted proof check. Continuation metadata events: GPT-6 Astra (OpenAI), ChatGPT Pro mode, effort Pro, S6\_10\_PAPERS continuation, 19 November 2026: used and rechecked in this run; status-vocabulary reconciliation; metadata reconciliation only; no new mathematical check. \emph{Pre-continuation recorded proof scope.} The actual invariant fibre class is kq. Re-derived the contractions 2 eta1 and -3 eta2, their coinvariant images 3 delta and2 delta, and the order3/order4 norm divisions. The resulting two boundary defects are2k delta and minus3k/2 delta. The exact based relation and zero cusp defect force k/2=0. The separately checked toric support-specialization morphism then trivializes the normalized central restriction of every analytic line bundle. The argument uses the repaired universal-cover comparison, with no unproved comparison assumption}{the hypotheses, notation and definitions stated here; \cref{wbII:eq:W:5.17-A5}; \cref{wbII:eq:W:5.17-A4}; \cref{wbII:eq:W:5.17-A6}; \cref{wbII:lem:analytic-support-specialization}}

\ For the standard threefold,
\[
 r=\eta^*i^*:\operatorname{Pic}(X)\longrightarrow
 \operatorname{Pic}(\mathrm{dP}_6)
\]
is the zero homomorphism.  This proof uses only holomorphic
\(\mathcal O^*\)-factors, the two finite logarithmic fillings, the cusp
argument principle, and the \(\Psi\)-equivariant toric support function.
It uses neither a global integral Leray transgression nor local integral
nearby-cycle specialization.
\end{theorem}



\begin{proof}
\noindent\textbf{Step 1.} Let \(L\in\operatorname{Pic}(X)\).  By
\cref{wb:thm:coinvariants}, its invariant fibre class is
\[
 E=c_1(L|_F)=kq,
 \qquad k\in\mathbb Z,
 \qquad q=uw+6\gamma\delta.
\tag{W.138}\label{wbII:eq:W:5.17-A35}
\]
The exact contractions, now used without an integral specialization
lattice, are
\[
 \iota_{v_1}q=2\eta_1,
 \qquad
 \iota_{v_2}q=-3\eta_2,
 \qquad
 [\eta_1]=3[\delta],
 \qquad
 [\eta_2]=2[\delta].
\tag{W.139}\label{wbII:eq:W:5.17-A36}
\]
Apply \eqref{wbII:eq:W:5.17-A5} with \((m_1,m_2)=(3,4)\).  In the real global
coinvariants,
\[
 \bar d_1=2k[\delta],
 \qquad
 \bar d_2=-\frac32k[\delta].
\tag{W.140}\label{wbII:eq:W:5.17-A37}
\]
Equations \eqref{wbII:eq:W:5.17-A4} and \eqref{wbII:eq:W:5.17-A6} therefore give
\[
 0=\bar d_1+\bar d_2+\bar d_0
   =\frac{k}{2}[\delta].
\tag{W.141}\label{wbII:eq:W:5.17-A38}
\]
Since \(V_G\otimes\mathbb R=\mathbb R[\delta]\), one has \(k=0\), so
\(E=0\).  Apply \cref{wbII:lem:analytic-support-specialization} to
\(M=L|_{N_0}\).  It gives
\[
 \eta^*i^*L=\nu^*(L|_W)\simeq\mathcal O_{E_0}.
\]
This holds for every \(L\), proving \(r=0\).
\end{proof}




The nonvanishing conclusion of the next corollary is contained in
\citet[Theorem~10.5(a), pp.~88--89]{AC2026}, which treats every global line
bundle.  This proof records the smaller set of inputs sufficient for
\(\operatorname{Pic}^0\)-twists.

\begin{corollary}[standard-family nonvanishing for \texorpdfstring{$\operatorname{Pic}^0$}{Pic0} without topology]
\label{wbII:cor:candidate-pic0-route}
\StatementRecord{wbII:cor:candidate-pic0-route}{proved}{Historical events: GPT-6 Astra (OpenAI), ChatGPT Pro mode, effort Pro, S6\_09\_PAPERS, 17 November 2026: targeted proof check. Also checked in this writing. Pre-continuation writing events: GPT-6 Astra (OpenAI), ChatGPT Pro mode, effort Pro, S6\_10\_PAPERS, 18 November 2026: targeted proof check. Continuation metadata events: GPT-6 Astra (OpenAI), ChatGPT Pro mode, effort Pro, S6\_10\_PAPERS continuation, 19 November 2026: used and rechecked in this run; status-vocabulary reconciliation; metadata reconciliation only; no new mathematical check. \emph{Pre-continuation recorded proof scope.} The actual cusp has full reduced normal-crossing fibre, dP6 normalization and a local canonical trivialization. The Pic0 Chern-class pullback gives the conductor section; absolute Cartier duality converts it into nonzero H2. Properness and relative flatness make the exact top-degree base-change map an isomorphism, so the R2 stalk is nonzero. Only these analytic construction and restriction maps are used}{the hypotheses, notation and definitions stated here; \cref{wbII:prop:pic0-pullback}; \cref{wbII:prop:conductor}; \cref{wbII:prop:cartier-duality-absolute}; \cref{wbII:lem:curve-flatness}; \cref{wbII:prop:analytic-base-change}}

\ For the standard threefold and every
$L\in\operatorname{Pic}^0(X)$,
\[
 H^0(W,\Omega_X^1|_W\otimes A_L)\ne0,
 \qquad
 H^2(W,(T_X\otimes L)|_W)\ne0,
\]
and
\[
 R^2f_*(T_X\otimes L)\ne0.
\tag{W.44}\label{wbII:eq:W:5.17f}
\]
This proof uses the analytic construction of the cusp fibre but no global
homology, Mayer--Vietoris, or Leray calculation.
\end{corollary}



\begin{proof}
\noindent\textbf{Step 1.} The cusp result in \cref{wb:thm:candidate-cusp-geometry} supplies the full compact
reduced normal-crossings fibre, its $\mathrm{dP}_6$ normalization, and the
trivialization of $\omega_X$ near $W$.  By
\cref{wbII:prop:pic0-pullback}, $\eta^*A_L$ is trivial, so
\cref{wbII:prop:conductor} gives the first nonvanishing.  Apply
\cref{wbII:prop:cartier-duality-absolute} to $E=T_X\otimes L$; its dual space is
exactly the preceding nonzero section space, giving the fibrewise $H^2$
nonvanishing.  The fibre has complex dimension two, so its coherent
$H^3$ vanishes.  The global construction makes $X$ compact and $f$
nonconstant; \cref{wbII:lem:curve-flatness} therefore gives properness and
$f$-flatness of $T_X\otimes L$.  The base-change isomorphism of
\cref{wbII:prop:analytic-base-change} now identifies the fibre of the $R^2$
stalk with the nonzero $H^2$ group.  Hence the stalk, and therefore the
direct-image sheaf, is nonzero.
\end{proof}



\paragraph{reconciliation of the three restriction routes}

The Pic$^0$ argument proves \eqref{wbII:eq:W:5.17f} without global homology,
Mayer--Vietoris, or Leray.  The all-Pic statement
\cref{wbII:prop:leray-edge-restriction} avoids citing the packaged theorem
$H^2(X;\mathbb Z)=0$, but it reuses the same verified global integral
Leray and specialization machinery from which that theorem is deduced.
By contrast, \cref{wbII:prop:direct-boundary-picard} supplies a second all-Pic
proof without any global Leray differential: it recovers the primitive
$4-3=1$ obstruction directly from the finite logarithmic gluing and the
cusp boundary.  It still uses the local integral cusp-specialization
theorem.  The third proof in \cref{wbII:thm:analytic-all-picard} replaces both
that specialization and the cusp Chern-class bounding-chain evaluation by
logarithmic $\mathcal O^*$ factors, the argument principle, and the exact
$\Psi$-equivariant support function.  Thus the purely analytic
strengthening is now complete; the three proofs have different
dependency sets.

\Uses{the hypotheses, notation and definitions stated here; \cref{wbII:eq:W:5.17f}; \cref{wbII:prop:leray-edge-restriction}; \cref{wbII:prop:direct-boundary-picard}; \cref{wbII:thm:analytic-all-picard}}


\subsection{Abstract conductor theorem for the standard fibre}

The conductor mechanism below is the argument of
\citet[Lemma~10.3 and Theorem~10.5(a), pp.~87--89]{AC2026}, with its sufficient
hypotheses separated explicitly.  In particular the constructed-family
conclusion of \cref{wbII:cor:candidate-conductor-actual} is the known
nonvanishing assertion of that theorem.  The complete proof is retained
so that its normalization, duality and base-change inputs can be checked
individually.

\begin{theorem}[the conductor mechanism for the standard family]
\label{wbII:thm:candidate-conductor}
\StatementRecord{wbII:thm:candidate-conductor}{proved}{Historical events: GPT-6 Astra (OpenAI), ChatGPT Pro mode, effort Pro, S6\_09\_PAPERS, 17 November 2026: targeted proof check. Also checked in this writing. Pre-continuation writing events: GPT-6 Astra (OpenAI), ChatGPT Pro mode, effort Pro, S6\_10\_PAPERS, 18 November 2026: targeted proof check; used and rechecked in this run. Continuation metadata events: GPT-6 Astra (OpenAI), ChatGPT Pro mode, effort Pro, S6\_10\_PAPERS continuation, 19 November 2026: used and rechecked in this run; status-vocabulary reconciliation; metadata reconciliation only; no new mathematical check. \emph{Pre-continuation recorded proof scope.} The integral H2 hypothesis kills the first Chern class of the full dual canonical twist. Pullback to dP6, H1(O)=0 and the exponential sequence trivialize that twist. The conductor descent yields a nonzero ambient differential, Cartier adjunction identifies the full-fibre dual, and top-degree analytic base change gives the nonzero R2 stalk. The exact zero-nef-divisor instance of Demazure vanishing and the projective analytic-algebraic comparison were matched to Cox-Little-Schenck Theorem9.2.3 with its full proof and Serre GAGA Theorem1 with its full proof; the source records state the precise editions and reading scope}{the hypotheses, notation and definitions stated here; \cref{wbII:eq:W:5.3}; \cite[Theorem~9.2.3, author preprint p.~410]{CoxLittleSchenck2011}; \cite[\S12, Theorem~1, p.~19]{Serre1956GAGA}; \cref{wbII:prop:conductor}; \cref{wbII:thm:cartier-duality}; \cref{wbII:eq:W:5.20}; \cref{wbII:lem:curve-flatness}; \cref{wbII:prop:analytic-base-change}}

\ Assume:
\begin{enumerate}[label=(H\arabic*)]
\item $X$ is a connected compact smooth complex threefold and
$f:X\to B$ is a nonconstant holomorphic map to a smooth compact curve;
\item $W=f^{-1}(p_0)$ is the full scheme-theoretic fibre, is reduced, and
has the normal-crossings local models \eqref{wbII:eq:W:5.3};
\item the normalization $\eta:\widetilde W\to W$ is the degree-six del
Pezzo surface $\mathrm{dP}_6$;
\item $H^2(X;\mathbb Z)=0$.
\end{enumerate}
Then, for every $L\in\operatorname{Pic}(X)$,
\[
H^2(W,(T_X\otimes L)|_W)\ne0
\quad\text{and}\quad
R^2f_*(T_X\otimes L)\ne0.
\tag{W.45}\label{wbII:eq:W:5.18}
\]
\end{theorem}



\begin{proof}
\noindent\textbf{Step 1.} Fix $L$ and put
\[
\mathcal F_L=T_X\otimes L,\qquad
A_L=(L^\vee\otimes\omega_X)|_W.
\]
Hypothesis (H4) gives $c_1(L^\vee\otimes\omega_X)=0$.  Pullback therefore
gives $c_1(\eta^*A_L)=0$. The surface $\mathrm{dP}_6$ is the blow-up of
$\mathbb P^2$ at three noncollinear points. A projective change of coordinates
moves these points to the three coordinate points, so this is the smooth
complete toric surface obtained by subdividing the three two-dimensional
cones of the fan of $\mathbb P^2$. Its fan has support $\mathbb R^2$,
and the zero Cartier divisor is nef. Demazure vanishing therefore gives
$H^1(\mathrm{dP}_6,\mathcal O)=0$ in algebraic cohomology
\cite[Theorem~9.2.3, author preprint p.~410]{CoxLittleSchenck2011}.
Since this surface is projective and $\mathcal O$ is coherent,
Serre's comparison theorem identifies that group with analytic cohomology
\cite[\S12, Theorem~1, p.~19]{Serre1956GAGA}.
In the exponential sequence, the kernel
of
\[
c_1:\operatorname{Pic}(\mathrm{dP}_6)\longrightarrow
H^2(\mathrm{dP}_6;\mathbb Z)
\]
is the image of $H^1(\mathrm{dP}_6,\mathcal O)$ and is therefore zero.
Consequently
\[
\eta^*A_L\simeq\mathcal O_{\widetilde W}.
\tag{W.46}\label{wbII:eq:W:5.19}
\]
\Cref{wbII:prop:conductor} now gives
\[
H^0(W,\Omega^1_X|_W\otimes A_L)\ne0.
\tag{W.47}\label{wbII:eq:W:5.20}
\]

\noindent\textbf{Step 2.} Because $W=f^*(p_0)$, one has
$\mathcal O_X(W)=f^*\mathcal O_B(p_0)$, whose restriction to the fibre is
trivial.  Adjunction therefore gives $\omega_W\simeq\omega_X|_W$.
Apply \cref{wbII:thm:cartier-duality} to $E=\mathcal F_L$:
\[
\begin{split}
H^2(W,\mathcal F_L|_W)^\vee
&\simeq H^0(W,(\mathcal F_L|_W)^\vee\otimes\omega_W)\\
&=H^0(W,\Omega^1_X|_W\otimes A_L).
\end{split}
\]
The right side is nonzero by \eqref{wbII:eq:W:5.20}, proving the first assertion of \eqref{wbII:eq:W:5.18}.
The same duality theorem gives
$H^3(W,\mathcal F_L|_W)=0$.

\noindent\textbf{Step 3.} Since $X$ is compact, $f$ is proper.  By
\cref{wbII:lem:curve-flatness}, $f$ is flat and the vector bundle
$\mathcal F_L$ is $f$-flat.  Thus every hypothesis of
\cref{wbII:prop:analytic-base-change} is satisfied at $p_0$, and
\[
(R^2f_*\mathcal F_L)_{p_0}\otimes\mathbb C(p_0)
\simeq H^2(W,\mathcal F_L|_W)\ne0.
\]
The stalk, and hence the coherent direct-image sheaf, is nonzero.  This is
the second assertion of \eqref{wbII:eq:W:5.18}.
\end{proof}



\begin{corollary}[application to the standard threefold]
\label{wbII:cor:candidate-conductor-actual}
\StatementRecord{wbII:cor:candidate-conductor-actual}{proved}{Historical events: GPT-6 Astra (OpenAI), ChatGPT Pro mode, effort Pro, S6\_09\_PAPERS, 17 November 2026: targeted proof check. Also checked in this writing. Pre-continuation writing events: GPT-6 Astra (OpenAI), ChatGPT Pro mode, effort Pro, S6\_10\_PAPERS, 18 November 2026: targeted proof check. Continuation metadata events: GPT-6 Astra (OpenAI), ChatGPT Pro mode, effort Pro, S6\_10\_PAPERS continuation, 19 November 2026: used and rechecked in this run; status-vocabulary reconciliation; metadata reconciliation only; no new mathematical check. \emph{Pre-continuation recorded proof scope.} The independently re-derived direct boundary-Picard argument trivializes the normalization pullback for every analytic line bundle on the actual family. The local canonical line and conductor section then pass through absolute Cartier duality and top-degree base change to both nonvanishings. The Leray-edge argument provides a separately checked certificate of the same restriction}{the hypotheses, notation and definitions stated here; \cref{wbII:prop:direct-boundary-picard}; \cref{wbII:cor:candidate-pic0-route}; \cref{wbII:prop:leray-edge-restriction}}

\ For the global construction of
\cref{wb:thm:candidate-global-geometry} and every
$L\in\operatorname{Pic}(X)$,
\[
H^2(W,(T_X\otimes L)|_W)\ne0,
\qquad R^2f_*(T_X\otimes L)\ne0.
\]
\end{corollary}



\begin{proof}
\noindent\textbf{Step 1.} The cusp construction gives the compact full Cartier fibre, its
$\mathrm{dP}_6$ normalization, and the local trivialization of $\omega_X$.
By the local canonical trivialization and
\cref{wbII:prop:direct-boundary-picard}, $\eta^*A_L$ is trivial for every $L$.
The conductor proposition then gives
$H^0(W,\Omega_X^1|_W\otimes A_L)\ne0$; absolute Cartier-fibre duality turns
this into the displayed $H^2$ nonvanishing, and the uniform top-degree
base-change isomorphism turns it into a nonzero $R^2$ stalk.  This is exactly
the chain proved in \cref{wbII:cor:candidate-pic0-route}, now with the all-Pic
restriction supplied directly by \cref{wbII:prop:direct-boundary-picard}.  The
global Leray-edge proof in \cref{wbII:prop:leray-edge-restriction} is an
independent certificate of the same restriction statement.
\end{proof}



\subsection{The corrected theorem that survives both defects}

The original hypotheses do not force the CDP conclusion: the constructed
threefold above is a counterexample.  There is nevertheless a sharp useful
replacement.  It isolates the exact extra global vanishing needed by the
intact part of the argument, and the following fibrewise criterion replaces
the invalid p.~692 reduction.

\begin{lemma}[generic endpoint vanishing]
\label{wbII:lem:generic-endpoints}
\StatementRecord{wbII:lem:generic-endpoints}{proved}{Historical events: GPT-6 Astra (OpenAI), ChatGPT Pro mode, effort Pro, S6\_09\_PAPERS, 17 November 2026: targeted proof check. Also checked in this writing. Pre-continuation writing events: GPT-6 Astra (OpenAI), ChatGPT Pro mode, effort Pro, S6\_10\_PAPERS, 18 November 2026: targeted proof check. GPT-6 Astra (OpenAI), ChatGPT Pro mode, effort Pro, S6\_10\_PAPERS, 18 November 2026 05:35:21+00:00: used and rechecked in this run. Continuation metadata events: GPT-6 Astra (OpenAI), ChatGPT Pro mode, effort Pro, S6\_10\_PAPERS continuation, 19 November 2026: used and rechecked in this run; status-vocabulary reconciliation; metadata reconciliation only; no new mathematical check. \emph{Pre-continuation recorded proof scope.} The descending finite-dimensional section spaces for negative multiples of the effective Cartier divisor stabilize. Every coefficient of a stable germ is divisible by all powers of its nonzerodivisor equation, hence vanishes; connectedness kills the global section. The specified Poincare family gives closed analytic jump loci, each proper by the supplied zero point. Irreducibility and inversion preserve the required dense open complements; Serre duality gives the top-degree endpoint}{the hypotheses, notation and definitions stated here; \cite{Grauert1960D,BS76push}}

Let $Y$ be a connected compact complex manifold of dimension $n$, let $D$ be
a nonzero effective Cartier divisor with $\cO_Y(D)\in\Pic^0(Y)$, and let
$B$ be a holomorphic vector bundle.  Let $\mathcal P\subset\Pic^0(Y)$ be a
connected irreducible reduced complex analytic parameter space carrying a
Poincar\'e line bundle on $Y\times\mathcal P$.  Assume that dualization
preserves $\mathcal P$, and suppose it
  contains points $P_0=\cO_Y(-mD)$ and $Q_0=\cO_Y(-\ell D)$ for some
  $m,\ell\ge1$ with
  $H^0(Y,B\otimes P_0)=0$ and
  $H^0(Y,B^\vee\otimes\omega_Y\otimes Q_0)=0$ (the stabilization argument
  supplies such exponents whenever the corresponding twists lie in
  $\mathcal P$).  Then there is a dense open subset
$U\subset\mathcal P$ such that, for every $P'\in U$,
\[
 H^0(Y,B\otimes P')=H^n(Y,B\otimes P')=0.
 \tag{W.169}\label{wbII:eq:CN:5.2}
\]
\end{lemma}



\begin{proof}
\noindent\textbf{Step 1.} Put
  \(V_k=H^0(Y,B\otimes\cO_Y(-kD))\), viewed inside $H^0(Y,B)$, and note
  that $V_{k+1}\subseteq V_k$.  The descending chain stabilizes.  If $s$ is
  in its stable value, then at a point of $D$ choose a local defining
  nonzerodivisor $g_D$ for the Cartier divisor and a local frame of $B$.
  Every coefficient of the germ of $s$ is divisible by $g_D^k$ for every
  $k$.  The Krull-intersection theorem in the regular local ring
  $\cO_{Y,y}$ gives $\bigcap_{k\ge1}g_D^k\cO_{Y,y}=0$, so the germ is zero at
  every point of $D$.  Hence $s$ vanishes on a nonempty open set and, by
  connectedness, on $Y$.  Thus $V_k=0$ for all sufficiently large $k$.

  The chosen point $P_0$ satisfies
  $H^0(Y,B\otimes P_0)=0$.  The Poincar\'e bundle gives, locally on
$\mathcal P$, a family to which Grauert's theorem applies; the local jump
  loci $\{P':h^0(Y,B\otimes P')\ge1\}$ glue to a closed analytic subset of
  $\mathcal P$ (Grauert's theorem; see \cite{Grauert1960D,BS76push}).  It is proper because it misses $P_0$, and irreducibility of
$\mathcal P$ makes its complement dense open.  Applying the same stabilization
argument to $B^\vee\otimes\omega_Y$ and using the chosen point $Q_0$ gives
$H^0(Y,B^\vee\otimes\omega_Y\otimes Q_0)=0$; dualization stability puts
$Q_0^\vee$ in $\mathcal P$.  Serre duality then gives
$H^n(Y,B\otimes Q_0^\vee)=0$.  The corresponding top-degree jump locus is
therefore also proper with dense open complement.  Intersecting the two
dense open complements proves the claim.
\end{proof}



The following endpoint argument is a re-derivation of
\citet[Proposition~1.1 and Corollaries~1.2--1.3]{CDP98}, with Serre duality
for the top degree.  Its formulation makes the effective-divisor
\(\operatorname{Pic}^0\) hypothesis explicit in arbitrary dimension.

\begin{lemma}[generic endpoint vanishing]\label{wbII:lem:cdp-generic-endpoints}
\StatementRecord{wbII:lem:cdp-generic-endpoints}{proved}{Historical events: GPT-6 Astra (OpenAI), ChatGPT Pro mode, effort Pro, S6\_09\_PAPERS, 17 November 2026: targeted proof check. Also checked in this writing. Pre-continuation writing events: GPT-6 Astra (OpenAI), ChatGPT Pro mode, effort Pro, S6\_10\_PAPERS, 18 November 2026: targeted proof check. Continuation metadata events: GPT-6 Astra (OpenAI), ChatGPT Pro mode, effort Pro, S6\_10\_PAPERS continuation, 19 November 2026: used and rechecked in this run; status-vocabulary reconciliation; metadata reconciliation only; no new mathematical check. \emph{Pre-continuation recorded proof scope.} The exact negative-divisor inclusions and the local intersection of all powers kill the stable H0 subspace. All these twists lie in Pic0. Local universal line-bundle families give a proper analytic jump set, and its complement is dense in the connected complex Lie group. Repeating for the dual canonical bundle and applying inversion and Serre duality proves the other endpoint}{the hypotheses, notation and definitions stated here; the explicit constructions and calculations in the proof}

\ Let \(Y\) be a connected compact complex manifold, let \(D\) be
a nonzero effective Cartier divisor, and suppose
\(\mathcal O_Y(D)\in\operatorname{Pic}^0(Y)\).  For every holomorphic vector
bundle \(B\), a general \(P\in\operatorname{Pic}^0(Y)\) satisfies
\[
 H^0(Y,B\otimes P)=0,\qquad
 H^{\dim Y}(Y,B\otimes P)=0.
\tag{W.50}\label{wbII:eq:W:5.23}
\]
\end{lemma}



\begin{proof}
\noindent\textbf{Step 1.} The inclusions \(\mathcal O_Y(-(k+1)D)\hookrightarrow
\mathcal O_Y(-kD)\) give a descending chain
\[
 V_k=H^0(Y,B\otimes\mathcal O_Y(-kD))
 \subseteq H^0(Y,B).
\]
It stabilizes because the last vector space is finite-dimensional.  A
section in the stable value belongs to every \(V_k\).  In a local equation
\(z=0\) for \(D\), every coefficient of that section is divisible by
\(z^k\) for all \(k\), and hence is zero.  The section vanishes on a
neighbourhood of \(D\), and therefore on connected \(Y\) by analytic
continuation.  Thus \(V_k=0\) for some \(k\).  Since
\(\mathcal O_Y(-kD)\in\operatorname{Pic}^0(Y)\), upper semicontinuity in a
local universal family of topologically trivial line bundles shows that
\[
 \{P:h^0(Y,B\otimes P)>0\}
\]
is a proper analytic subset of \(\operatorname{Pic}^0(Y)\).  This proves
the first vanishing for general \(P\).  Serre duality identifies the dual
of the second group in \eqref{wbII:eq:W:5.23} with
\[
 H^0(Y,B^\vee\otimes\omega_Y\otimes P^\vee).
\]
Apply the first part to \(B^\vee\otimes\omega_Y\) and use that inversion is
an automorphism of \(\operatorname{Pic}^0(Y)\).
\end{proof}



\begin{theorem}[corrected CDP global-vanishing criterion]
\label{wbII:thm:corrected-cdp-global}
\StatementRecord{wbII:thm:corrected-cdp-global}{proved}{Historical events: GPT-6 Astra (OpenAI), ChatGPT Pro mode, effort Pro, S6\_09\_PAPERS, 17 November 2026: targeted proof check. Also checked in this writing. Pre-continuation writing events: GPT-6 Astra (OpenAI), ChatGPT Pro mode, effort Pro, S6\_10\_PAPERS, 18 November 2026: targeted proof check. Continuation metadata events: GPT-6 Astra (OpenAI), ChatGPT Pro mode, effort Pro, S6\_10\_PAPERS continuation, 19 November 2026: used and rechecked in this run; status-vocabulary reconciliation; metadata reconciliation only; no new mathematical check. \emph{Pre-continuation recorded proof scope.} The integral H1/H2 vanishings identify Pic with the vector space H1(O); exponentiated additive Cech cocycles give its actual universal family. The two endpoint vanishings and the one assumed H2-vanishing point each remove a proper analytic jump locus. On the common dense open only H1 can remain, so Euler characteristic is nonpositive. Poincare duality kills rational H4, and the degree-six Chern-character/Todd product is exactly c3/2, proving the stated sign via analytic Riemann-Roch}{the hypotheses, notation and definitions stated here; \cref{wbII:lem:cdp-generic-endpoints}; \cref{wbII:eq:W:5.28}}
\label{wbII:thm:corrected-global}
\ Let \(X\) be a compact connected complex threefold such that
\[
 H^1(X;\mathbb Z)=H^2(X;\mathbb Z)=0,\qquad a(X)=1,
\]
and suppose its algebraic reduction is a holomorphic map
\(f:X\to C\).  If there exists \(L_0\in\operatorname{Pic}(X)\) with
\[
 H^2(X,T_X\otimes L_0)=0,
\tag{W.51}\label{wbII:eq:W:5.24}
\]
then, for \(M\) in a dense open subset of \(\operatorname{Pic}(X)\),
\[
 H^i(X,T_X\otimes M)=0\quad(i\ne1),
\tag{W.52}\label{wbII:eq:W:5.25}
\]
\[
 \chi(X,T_X\otimes M)\le0,\qquad
 \left\langle c_3(T_X),[X]\right\rangle\le0.
\tag{W.53}\label{wbII:eq:W:5.26}
\]
\end{theorem}



\begin{proof}
\noindent\textbf{Step 1.} The exponential sequence and the two integral vanishings give
\[
 \operatorname{Pic}(X)=\operatorname{Pic}^0(X)
 \simeq H^1(X,\mathcal O_X).
\tag{W.54}\label{wbII:eq:W:5.27}
\]
Thus this Picard group is a connected irreducible complex vector group.
Choose additive Cech cocycles $a^{(r)}_{ij}$ representing a basis of
$H^1(X,\mathcal O_X)$. For coordinates $\lambda_r$ on that vector space,
\[
 g_{ij}(\lambda)=\exp\!\left(\sum_r\lambda_r a^{(r)}_{ij}\right)
\]
satisfy the multiplicative cocycle condition and define a holomorphic
line bundle on $X\times H^1(X,\mathcal O_X)$. Its restriction at every
parameter represents the corresponding line bundle, so proper
semicontinuity below applies to an actual holomorphic family.  Any fibre divisor \(D\) of
\(f\) is nonzero and effective, and its line bundle lies in
\(\operatorname{Pic}^0(X)\) by \(H^2(X;\mathbb Z)=0\).  Applying
\cref{wbII:lem:cdp-generic-endpoints} to \(B=T_X\) shows that the degree-zero and
degree-three groups in \eqref{wbII:eq:W:5.25} vanish away from two proper analytic subsets
of \(\operatorname{Pic}(X)\).

\noindent\textbf{Step 2.} The cohomology-jump locus
\[
 \{M\in\operatorname{Pic}(X):h^2(X,T_X\otimes M)\ge1\}
\]
is analytic by upper semicontinuity.  It is proper because \eqref{wbII:eq:W:5.24} gives a
point outside it.  Its complement therefore meets the two endpoint
vanishing loci in a dense open set.  This proves \eqref{wbII:eq:W:5.25}, and on that set
\[
 \chi(X,T_X\otimes M)=-h^1(X,T_X\otimes M)\le0.
\tag{W.55}\label{wbII:eq:W:5.28}
\]

\noindent\textbf{Step 3.} Every \(M\in\operatorname{Pic}(X)\) has rationally trivial first Chern
class.  The hypotheses also give \(c_1(T_X)=0\) in rational cohomology and,
by Poincar\'e duality, \(H^4(X;\mathbb Q)=0\), so
\(c_2(T_X)=0\) rationally.  The degree-six term of
\(\operatorname{ch}(T_X\otimes M)\operatorname{td}(X)\) is consequently
\(\tfrac12c_3(T_X)\).  Hirzebruch--Riemann--Roch turns \eqref{wbII:eq:W:5.28} into
\[
 \chi(X,T_X\otimes M)
 =\frac12\left\langle c_3(T_X),[X]\right\rangle\le0,
\]
which proves \eqref{wbII:eq:W:5.26}.
\end{proof}



\begin{corollary}[intact relative criterion]\label{wbII:cor:corrected-cdp-relative}
\StatementRecord{wbII:cor:corrected-cdp-relative}{proved}{Historical events: GPT-6 Astra (OpenAI), ChatGPT Pro mode, effort Pro, S6\_09\_PAPERS, 17 November 2026: targeted proof check. Also checked in this writing. Pre-continuation writing events: GPT-6 Astra (OpenAI), ChatGPT Pro mode, effort Pro, S6\_10\_PAPERS, 18 November 2026: targeted proof check. Continuation metadata events: GPT-6 Astra (OpenAI), ChatGPT Pro mode, effort Pro, S6\_10\_PAPERS continuation, 19 November 2026: used and rechecked in this run; status-vocabulary reconciliation; metadata reconciliation only; no new mathematical check. \emph{Pre-continuation recorded proof scope.} Projection formula preserves the assumed zero R2 after a positive base twist. On the compact curve, sufficiently positive powers kill H1 of the coherent R1 sheaf; the only two total-degree-two Leray terms then vanish. The resulting explicit line bundle satisfies the already checked global criterion. The exact CDP20 Proposition2.4 source locus is separately recorded}{the hypotheses, notation and definitions stated here; \cref{wbII:thm:corrected-cdp-global}; \cite[Proposition~2.4, p.~680]{CDP20}; \cref{wbII:eq:W:5.24}}
\label{wbII:cor:relative}
\ In the setting of
\cref{wbII:thm:corrected-cdp-global}, it is enough that one line bundle
\(L\in\operatorname{Pic}(X)\) satisfy
\[
 R^2f_*(T_X\otimes L)=0.
\tag{W.56}\label{wbII:eq:W:5.29}
\]
This is the valid content of the sufficient criterion printed as
\cite[Proposition~2.4, p.~680]{CDP20}.
\end{corollary}



\begin{proof}
\noindent\textbf{Step 1.} Let \(A\) be a positive-degree line bundle on the compact Riemann surface
\(C\), and put \(E=T_X\otimes L\).  Projection formula gives
\[
 R^2f_*(E\otimes f^*A^m)=R^2f_*E\otimes A^m=0.
\]
For \(m\gg0\), Serre vanishing on \(C\) gives
\[
 H^1(C,R^1f_*E\otimes A^m)=0.
\]
The Leray spectral sequence has only columns zero and one, so its
total-degree-two terms yield
\[
 H^2(X,E\otimes f^*A^m)=0.
\]
Thus \(L_0=L\otimes f^*A^m\) satisfies \eqref{wbII:eq:W:5.24}, and
\cref{wbII:thm:corrected-cdp-global} applies.
\end{proof}



\begin{theorem}[exact reduced-fibre repair of the p.~692 step]
\label{wbII:thm:corrected-cdp-fibrewise}
\StatementRecord{wbII:thm:corrected-cdp-fibrewise}{proved}{Historical events: GPT-6 Astra (OpenAI), ChatGPT Pro mode, effort Pro, S6\_09\_PAPERS, 17 November 2026: targeted proof check. Also checked in this writing. Pre-continuation writing events: GPT-6 Astra (OpenAI), ChatGPT Pro mode, effort Pro, S6\_10\_PAPERS, 18 November 2026: targeted proof check. Continuation metadata events: GPT-6 Astra (OpenAI), ChatGPT Pro mode, effort Pro, S6\_10\_PAPERS continuation, 19 November 2026: used and rechecked in this run; status-vocabulary reconciliation; metadata reconciliation only; no new mathematical check. \emph{Pre-continuation recorded proof scope.} The missing cotangent filtration is applied to the full reduced Cartier fibre. Its two endpoint vanishings and the actual connecting-map injectivity kill the ambient differential space. The base DVR gives flatness and the full fibre has trivial normal line. Absolute Cartier duality and top-degree base change therefore make every R2 fibre zero; Nakayama kills every stalk, yielding the relative criterion and its global sign conclusion}{the hypotheses, notation and definitions stated here; \cref{wbII:thm:corrected-cdp-global}; \cref{wbII:eq:W:5.29}; \cref{wbII:eq:W:5.25}; \cref{wbII:eq:W:5.26}; \cref{cdp:st:A.1.1}; \cref{wbII:cor:surface-top-base-change}; \cref{wbII:cor:corrected-cdp-relative}}

\ Retain the setting of
\cref{wbII:thm:corrected-cdp-global}, assume every full scheme-theoretic fibre
\(W_c=f^{-1}(c)\) is reduced, and fix \(L\in\operatorname{Pic}(X)\).  Put
\[
 A_c=(\omega_X\otimes L^\vee)|_{W_c},\qquad
 \mathscr T_c=\ker\!\left(
 \Omega^1_{W_c}\longrightarrow
 j_*\Omega^1_{(W_c)_{\rm reg}}\right).
\]
For each \(c\), let
\[
 \delta_c:
 H^0(W_c,\mathscr T_c\otimes A_c)
 \longrightarrow
 H^1(W_c,N^*_{W_c/X}\otimes A_c)
\tag{W.57}\label{wbII:eq:W:5.30}
\]
be the connecting map from the missing cotangent filtration.  If, for every
\(c\),
\[
 H^0(W_c,\widetilde\Omega^1_{W_c}\otimes A_c)=0,\qquad
 H^0(W_c,N^*_{W_c/X}\otimes A_c)=0,\qquad
 \delta_c\ \text{is injective},
\tag{W.58}\label{wbII:eq:W:5.31}
\]
then \eqref{wbII:eq:W:5.29}, \eqref{wbII:eq:W:5.25}, and \eqref{wbII:eq:W:5.26} hold.
\end{theorem}



\begin{proof}
\noindent\textbf{Step 1.} Apply \cref{cdp:st:A.1.1} to the whole reduced
Cartier fibre \(W_c\), not separately to its components.  Conditions \eqref{wbII:eq:W:5.31}
give
\[
 H^0(W_c,\Omega_X^1|_{W_c}\otimes A_c)
 =\ker\delta_c=0.
\tag{W.59}\label{wbII:eq:W:5.32}
\]
The map \(f\) is flat: a local parameter at \(c\) pulls back to a
non-zero-divisor in every regular local ring of the connected manifold \(X\),
and torsion-freeness over the base DVR is flatness.  Moreover, the normal
line of the full fibre is the pullback of the one-dimensional vector space
\(T_cC\), hence is trivial.  Adjunction therefore identifies
\(\omega_{W_c}\simeq\omega_X|_{W_c}\).  With \(E=T_X\otimes L\), absolute
Cartier-fibre duality now gives
\[
 H^2(W_c,E|_{W_c})^\vee
 \simeq H^0(W_c,\Omega_X^1|_{W_c}\otimes A_c)=0.
\]
The uniform top-degree base-change isomorphism of
\cref{wbII:cor:surface-top-base-change} therefore gives
\[
 (R^2f_*E)_c\otimes_{\mathcal O_{C,c}}\mathbb C(c)=0
\]
for every \(c\).  Nakayama's lemma kills every stalk of the coherent sheaf
\(R^2f_*E\), proving \eqref{wbII:eq:W:5.29}.  The remaining conclusions follow from
\cref{wbII:cor:corrected-cdp-relative}.
\end{proof}



\paragraph{scope and sharpness of the repair}

Under the two endpoint vanishings in \eqref{wbII:eq:W:5.31}, injectivity of \(\delta_c\) is
equivalent to the desired ambient vanishing \eqref{wbII:eq:W:5.32}; it is the exact lifting condition.  For a reducible reduced fibre the
map must be computed on the whole Cartier fibre.  A nonreduced fibre requires
additional nilpotent-thickening layers and is not covered by
\cref{wbII:thm:corrected-cdp-fibrewise}.  The global criterion
\cref{wbII:thm:corrected-cdp-global} and the relative criterion
\cref{wbII:cor:corrected-cdp-relative} do allow nonreduced fibres.

The restandard family violates both added conditions rather than being
silently excluded by terminology.  Since \(c_3(X)=2\), the contrapositive of
\cref{wbII:thm:corrected-cdp-global} proves
\[
 H^2(X,T_X\otimes L_0)\ne0
 \quad\text{for every }L_0\in\operatorname{Pic}(X).
\]
Independently, \cref{wbII:cor:candidate-conductor-actual} proves
\[
 R^2f_*(T_X\otimes L)\ne0
 \quad\text{for every }L\in\operatorname{Pic}(X),
\]
and its cusp conductor supplies a nonzero element of the relevant
\(\ker\delta_c\).  The first conclusion rules out the global hypothesis; the
second independently rules out the stronger relative sufficient condition.
Thus the corrected theorem is compatible with, and
precisely separates, the verified counterexample to CDP20
Theorem~2.2 as stated.

\Uses{the hypotheses, notation and definitions stated here; \cref{wbII:eq:W:5.31}; \cref{wbII:eq:W:5.32}; \cref{wbII:thm:corrected-cdp-fibrewise}; \cref{wbII:thm:corrected-cdp-global}; \cref{wbII:cor:corrected-cdp-relative}; \cref{wbII:cor:candidate-conductor-actual}; \cref{wbII:eq:W:5.33}}


\paragraph{The exact monodromy repair.}
The split-extension argument can reach only the invariant integral lattice
\(H^1(F;\mathbb Z)^{\pi_1(C^\circ)}\).  Its final complex-linear inference is
valid precisely when
\[
 \operatorname{Span}_{\mathbb C}\operatorname{im}\!\left(
 H^1(F;\mathbb Z)^{\pi_1(C^\circ)}
 \longrightarrow H^1(F,\mathcal O_F)\right)
 =H^1(F,\mathcal O_F).
\tag{W.60}\label{wbII:eq:W:5.33}
\]
Trivial monodromy on \(H_1(F;\mathbb Z)\) is a stronger sufficient
hypothesis.  Neither \eqref{wbII:eq:W:5.33} nor trivial monodromy follows from the hypotheses
of CDP20 Theorem~2.2.  The corrected theorem above bypasses this rank route
entirely; it does not relabel the false direct-sum formula as proved.

\begin{corollary}[Positive Euler number obstructs the added vanishings]
\label{wbII:test:cdp-comparison}
\StatementRecord{wbII:test:cdp-comparison}{proved}{Historical events: GPT-6 Astra (OpenAI), ChatGPT Pro mode, effort Pro, S6\_09\_PAPERS, 17 November 2026: targeted proof check. Also checked in this writing. Pre-continuation writing events: GPT-6 Astra (OpenAI), ChatGPT Pro mode, effort Pro, S6\_10\_PAPERS, 18 November 2026: targeted proof check. Continuation metadata events: GPT-6 Astra (OpenAI), ChatGPT Pro mode, effort Pro, S6\_10\_PAPERS continuation, 19 November 2026: used and rechecked in this run; status-vocabulary reconciliation; metadata reconciliation only; no new mathematical check. \emph{Pre-continuation recorded proof scope.} The constructed threefold has the verified global hypotheses and Chern number two. A vanishing H2 for one twist contradicts the exact global criterion; a vanishing R2 contradicts the explicit positive-base-twist reduction. The independently checked cusp conductor gives the direct local nonvanishing as well. The earlier published counterexample attribution is retained}{the hypotheses, notation and definitions stated here; \cref{wbII:thm:corrected-cdp-global}; \cref{wbII:cor:corrected-cdp-relative}; \cref{wbII:cor:candidate-conductor-actual}; \citet[Section~10.2, Theorem~10.5]{AC2026}}

For the standard compact threefold $X$, every $L\in\operatorname{Pic}(X)$
satisfies
\[
 H^2(X,T_X\otimes L)\ne0,
 \qquad R^2f_*(T_X\otimes L)\ne0.
\]
The second conclusion also follows directly from the cusp conductor.
\end{corollary}


\begin{proof}
\noindent\textbf{Step 1.} Paper I proves that the standard compact threefold has
$H^1(X;\mathbb Z)=H^2(X;\mathbb Z)=0$, holomorphic algebraic reduction
to $\mathbb P^1$, and $\langle c_3(T_X),[X]\rangle=2$.
If the displayed degree-two cohomology vanished for one line bundle,
\cref{wbII:thm:corrected-cdp-global} would give $c_3(X)\le0$, a contradiction.
If the second direct image vanished for one twist,
\cref{wbII:cor:corrected-cdp-relative} would give such a global twist after
tensoring by the pullback of a sufficiently positive base line bundle,
again a contradiction. Independently,
\cref{wbII:cor:candidate-conductor-actual} constructs a nonzero cusp-fibre
cohomology class and identifies it with the fibre of the second direct
image by top-degree base change. The contradiction with the printed
CDP conclusion is credited to \citet[Section~10.2, Theorem~10.5]{AC2026}.
\end{proof}


\subsection{Ordinary relative one-forms on the actual local models}
The following computation concerns the relative cotangent module on the
smooth total space, before restriction to a fibre. All base maps and local
rings are specified in the statement.
\begin{proposition}[crossing and multiple-fibre modules]
\label{wbII:prop:relative-local}
\StatementRecord{wbII:prop:relative-local}{proved}{Historical events: GPT-6 Astra (OpenAI), ChatGPT Pro mode, effort Pro, S6\_09\_PAPERS, 17 November 2026: targeted proof check. Also checked in this writing. Pre-continuation writing events: GPT-6 Astra (OpenAI), ChatGPT Pro mode, effort Pro, S6\_10\_PAPERS, 18 November 2026: targeted proof check. Continuation metadata events: GPT-6 Astra (OpenAI), ChatGPT Pro mode, effort Pro, S6\_10\_PAPERS continuation, 19 November 2026: used and rechecked in this run; status-vocabulary reconciliation; metadata reconciliation only; no new mathematical check. \emph{Pre-continuation recorded proof scope.} Differentiating the actual three base maps gives the presentations. In the double model, ax=by forces (a,b)=(yh,xh); in the triple model, ax=by=cz forces (a,b,c)=h(yz,xz,xy). Thus the printed maps identify the modules with their precise free-module submodules, proving torsion-freeness. The origin needs three generators while generic rank is two. For t=s to the m, differentiation gives the exact torsion summand R/(s to the m-minus-one) ds}{the hypotheses, notation and definitions stated here; the explicit constructions and calculations in the proof}

The symbols $R_2$ and $R_3$ denote copies of the same regular local ring,
with different maps from the one-dimensional base:
let $R_2=\C\{x,y,z\}$ with base coordinate $t=xy$, and let
$R_3=\C\{x,y,z\}$ with base coordinate $t=xyz$.  Then
\[
\begin{aligned}
 M_2&:=\Omega^1_{R_2/\C\{t\}}
 =\frac{R_2\,dx\oplus R_2\,dy\oplus R_2\,dz}
        {R_2\,(y\,dx+x\,dy)},\\[2pt]
 M_3&:=\Omega^1_{R_3/\C\{t\}}
 =\frac{R_3\,dx\oplus R_3\,dy\oplus R_3\,dz}
        {R_3\,(yz\,dx+xz\,dy+xy\,dz)}.
\end{aligned}
\tag{W.176}\label{wbII:eq:CN:6.1}
\]
Both modules are torsion-free, but neither is locally free at the origin.
They embed in free modules by
\[
\begin{aligned}
 M_2&\hookrightarrow R_2e\oplus R_2dz,
 &dx&\longmapsto xe,&dy&\longmapsto-ye,&dz&\longmapsto dz,\\
 M_3&\hookrightarrow R_3e_1\oplus R_3e_2,
 &dx&\longmapsto xe_1,&dy&\longmapsto ye_2,&dz&\longmapsto-z(e_1+e_2),
\end{aligned}
\tag{W.177}\label{wbII:eq:CN:6.2}
\]
with images
\[
 R_2dz\oplus (x,y)e,
 \qquad
 R_3xe_1+R_3ye_2+R_3z(e_1+e_2).
\tag{W.178}\label{wbII:eq:CN:6.3}
\]
If $R_m=\C\{s,y_1,y_2\}$ and the base coordinate is $t=s^m$ with
$m\ge2$, then
\[
\Omega^1_{R_m/\C\{t\}}
\simeq R_m\,dy_1\oplus R_m\,dy_2
       \oplus R_m/(s^{m-1})\,ds,
\tag{W.179}\label{wbII:eq:CN:6.4}
\]
and the final summand is exactly its torsion submodule.
\end{proposition}



\begin{proof}
\noindent\textbf{Step 1.} Differentiating the three equations gives the presentations in \eqref{wbII:eq:CN:6.1}.  For
$M_2$, the first map in \eqref{wbII:eq:CN:6.2} sends $(a,b,c)$ to $(ax-by,c)$.  If this is
zero, then $c=0$ and $ax=by$.  Since $x$ and $y$ are relatively prime in the
UFD $R_2$, there is an $h$ with $b=xh$ and $a=yh$; hence its kernel is exactly
the displayed relation.  For $M_3$, the second map sends $(a,b,c)$ to
$(ax-cz,by-cz)$.  If it is zero, relative primeness first gives $x\mid c$,
then $y\mid(c/x)$; writing $c=xyh$ gives
$(a,b,c)=h(yz,xz,xy)$, again exactly the defining relation.  Thus both maps
are injective, proving torsion-freeness and \eqref{wbII:eq:CN:6.3}.

\noindent\textbf{Step 2.} At the maximal ideal $(x,y,z)$ each defining relation has all coefficients in
that maximal ideal.  Consequently the quotient of either module by the
maximal ideal has dimension $3$ over $\C$, while its generic rank is $2$;
a locally free module of rank $2$ would need only $2$ generators.  This proves
non-local-freeness at the origin.  Finally, for $t=s^m$ one has
$dt=m s^{m-1}ds=0$ relatively.  The integer $m$ is a unit in $\C$, so the
decomposition \eqref{wbII:eq:CN:6.4} follows.  Its last summand is annihilated by
$s^{m-1}$, while the other two summands are free; after inverting $s$ it
vanishes.  It is therefore exactly the torsion submodule.
\end{proof}



\paragraph{scope of the local correction}

\label{wbII:rem:relative-local-scope}
The modules in \eqref{wbII:eq:CN:6.1} are modules on the smooth total-space chart, and their
torsion-freeness is asserted before restricting to a fibre.  They must be
distinguished both from the intrinsic sheaf $\Omega^1_S$ of the singular
fibre and from the restricted relative module
$M_i\otimes_{R_i}R_i/(t)$ ($i=2,3$): tensoring with the fibre ring can create
conductor-supported torsion, as the $\tau$ calculation in the proof of
Theorem~\ref{cdp:st:B1} shows.  If a local quotient is finite
surjective \'{e}tale over a chart in \eqref{wbII:eq:CN:6.4}, relative differentials commute with
this base change and faithful flatness then reflects torsion of the pulled-back
module; a ramified quotient
requires a separate invariant-module calculation.  Thus the bare $t=xyz$
chart has a torsion-free but non-locally-free total-space relative module for
$p=1$; this fact alone does not decide torsion after fibre restriction or
quotient.  The multiple-fibre charts do have the total-space torsion in
\eqref{wbII:eq:CN:6.4}.  No claim about torsion in $\bigwedge^2\Omega^1$ or
$\bigwedge^3\Omega^1$, and no global direct-image or Hodge conclusion, is
made here.

\Uses{the hypotheses, notation and definitions stated here; \cref{wbII:eq:CN:6.1}; \cref{cdp:st:B1}; \cref{wbII:eq:CN:6.4}}





\section{The Liu--Rollenske parameter covering}

\begin{theorem}[The retained covering-parameter result]\label{ii:LR:covering}
\StatementRecord{ii:LR:covering}{proved}{No earlier separately identified check is recorded. Checked only in this writing. Pre-continuation writing events: GPT-6 Astra (OpenAI), ChatGPT Pro mode, effort Pro, S6\_10\_PAPERS, 18 November 2026: targeted proof check; used and rechecked in this run. Continuation metadata events: GPT-6 Astra (OpenAI), ChatGPT Pro mode, effort Pro, S6\_10\_PAPERS continuation, 19 November 2026: used and rechecked in this run; explicitly scoped metadata reconciliation; metadata reconciliation only; no new mathematical check. GPT-6 Astra (OpenAI), ChatGPT Pro mode, effort Pro, S6\_10\_PAPERS continuation, 19 November 2026: used and rechecked in this run; status-vocabulary reconciliation; metadata reconciliation only; no new mathematical check. GPT-6 Astra (OpenAI), ChatGPT Pro mode, effort Pro, S6\_10\_PAPERS continuation, 19 November 2026: used and rechecked in this run; explicitly scoped metadata reconciliation; metadata reconciliation only; no new mathematical check. \emph{Pre-continuation recorded proof scope.} The covering subclaim of the night-run compound statement P3K.36 was re-derived there by GPT-6 Astra (OpenAI, ChatGPT Pro mode, effort Pro), 16--18 November 2026. Also checked in this writing by GPT-6 Astra (OpenAI), ChatGPT Pro mode, effort Pro, 18 November 2026: the exact lattice, affine quotient, Albanese map, intrinsic endpoint obstruction, rational square root and cap conjugations, and the two-step endpoint norm. The source and compactification constructions have their own records in Paper~IV.}{\ref{cdp:st:E1}; Paper~IV, \ref{night:P3F.41} and \ref{night:P3F.47}; the retained constructive maps of \ref{iv:correction:P3K.35} identified by \ref{night:P3K.74}; \cite[Lemma~2.7, Corollary~2.8, Proposition~2.12, Proposition~3.6 and Definition~3.7]{LiuRollenske2026}}
Retain the exact scalar $u$ in the Liu--Rollenske action $B_u=uB_0$,
the good fourth root at the $III^*$ source fibre, the specified
logarithmic vector $a=e_4/4$, and its fixed line character.
The resulting construction is a holomorphic family on the universal
logarithmic cover of a sufficiently small punctured $u$-disc.
The threefold reached after one positive $u$-turn is not
biholomorphic to the initial threefold. The actual family descends
to the degree-two parameter cover
\[
                         u=w^2.\tag{LR.1}
\]
This degree is minimal. The conclusion holds with every retained
constant root-line character.
\end{theorem}
\begin{proof}
\par

\medskip\noindent\textbf{1. The actual central torus and its scalar}

\par\nobreak

The source's equation (2.3) is the lift of translation by
\(-2P\), multiplied by \(u\). On the minimal good root the marked
section specializes to \((1+i)/2\) on
\(E=\mathbb C/(\mathbb Z+i\mathbb Z)\). Its line is the nonzero
order-two line with character
\[
                  m+in\longmapsto(-1)^{m+n}.          \tag{LR.4}
\]
These are the explicit statements and coordinate maps in
\cite[Lemma~2.7 and Corollary~2.8, pp.~6--7 of the attached-source edition]{LiuRollenske2026}. The same marked root and corrected line are constructed by the source and Cartier maps of Paper~IV, \ref{night:P3F.41} and \ref{night:P3F.47}.

Since \(2P\) specializes to zero, the restriction of \(B_0\) to
this line complement is multiplication by a constant \(c\ne0\).
The constant is independent of \(u\): every source, line, and root
comparison was fixed before multiplying the action by \(u\).
Thus its actual scalar multiplier is
\[
              \alpha=cu,\qquad
              \tau=\frac{\log(cu)}{2\pi i},
              \qquad\operatorname{Im}\tau>0.         \tag{LR.5}
\]
The line can be presented on the universal cover of \(E\) as
\((z,v)\sim(z+m+in,(-1)^{m+n}v)\). Writing
\(v=e^{2\pi i\eta}\), and then quotienting by \(B_u\), gives
exactly the central abelian surface
\[
 A_\tau=\mathbb C^2/\Lambda_\tau,
 \quad
 \Lambda_\tau=
 \left\langle
  (2,0),(1+i,0),(1,1/2),(0,\tau)
 \right\rangle_{\mathbb Z}.                           \tag{LR.6}
\]
The displayed lattice map is literal: its generators are the two
periods in the kernel of (LR.4), the remaining character period,
and the actual scalar quotient period. In particular
\[
 \delta=(0,1)=2(1,1/2)-(2,0),\qquad e_4=(0,\tau).
\]

The root deck acts on the source by \(z\mapsto-iz\). In this flat
line presentation an actual lift is
\((z,\eta)\mapsto(-iz,\eta+\kappa)\), with
\(4\kappa\in\mathbb Z\). Indeed (LR.4) is invariant under
\(-i\); the ratio of any two lifts is a nowhere-zero function on
the compact elliptic curve, hence a scalar; its fourth power is one.
The flat lift displayed in the proof of \cite[Proposition~3.9, p.~10 of the attached-source edition]{LiuRollenske2026} has \(\kappa=1/4\). Keeping \(\kappa\) records this actual lift and every retained constant finite line character. The logarithmic vector
is exactly \(a=(0,\tau/4)\), by \cite[Definition~3.7, p.~10 of the attached-source edition]{LiuRollenske2026}. Therefore the reduced multiple fibre is
\[
 B_{\tau,\kappa}=A_\tau/\langle R\rangle,
 \qquad R(z,\eta)=(-iz,\eta+\tau/4+\kappa).           \tag{LR.7}
\]
This action has order four and is free: a nontrivial power has a
nonintegral multiple of \(\operatorname{Im}\tau\) in its second
coordinate, whereas the imaginary second-coordinate periods of
\(\Lambda_\tau\) are integral multiples of
\(\operatorname{Im}\tau\). Its fourth power is the lattice period
\((0,\tau+4\kappa)\).

For reference, the linear deck in the ordered lattice basis of
(LR.6) is
\[
 D=\begin{pmatrix}
 1&1&0&0\\-2&-1&-1&0\\0&0&1&0\\0&0&0&1
 \end{pmatrix}.
\]
Its invariant lattice is exactly
\(\mathbb Z\delta\oplus\mathbb Z e_4\), agreeing with the
\cite[Proposition~3.6, pp.~9--10 of the attached-source edition]{LiuRollenske2026}.
A positive \(u\)-turn has the exact comparison
\[
 \tau\longmapsto\tau+1,
 \quad \Lambda_{\tau+1}=\Lambda_\tau,
 \quad a\longmapsto a+\delta/4.                      \tag{LR.8}
\]
The unchanged abelian surface in (LR.8) does not identify the affine
quotients: the torsion translation in (LR.7) has changed.

\par

\medskip\noindent\textbf{2. An intrinsic invariant of the reduced fibre}

\par\nobreak

The actual holomorphic projection
\[
 \pi_{\tau,\kappa}:B_{\tau,\kappa}\longrightarrow
 E_{\tau,\kappa}:=
 \frac{\mathbb C}{\frac12\mathbb Z+
                       (\frac\tau4+\kappa)\mathbb Z},
 \qquad [(z,\eta)]\longmapsto[\eta]                  \tag{LR.9}
\]
is well-defined by every generator in (LR.6) and (LR.7).
It has connected fibres. More explicitly, the image of the full
universal-cover deck group under its second-coordinate translation
is precisely the displayed rank-two lattice; after choosing a lift
of a target point, its zero-second-coordinate subgroup has elliptic
kernel generated by \((2,0),(1+i,0)\). This is the actual fibre of
(LR.9), so it is connected.

We also verify that (LR.9) is the Albanese map, without a choice of
marking. Holomorphic one-forms on \(A_\tau\) are constant linear
forms: the coefficients of their pullback to \(\mathbb C^2\) are
periodic entire functions and hence constant by compactness. The
\(R\)-invariant subspace is exactly \(\mathbb C\,d\eta\), since
\(R^*dz=-i\,dz\) and \(R^*d\eta=d\eta\). The periods of
\(d\eta\) on the full fundamental group of \(B_{\tau,\kappa}\)
are exactly \(\frac12\mathbb Z+(\frac\tau4+\kappa)\mathbb Z\):
this follows directly from the same affine deck generators. The
integration definition of the Albanese torus therefore gives
precisely (LR.9), up to a target translation. In particular its
elliptic-curve isomorphism class is intrinsic to \(B_{\tau,\kappa}\).

Normalize its first period to one by \(\eta\mapsto2\eta\). Its
Tate parameter is
\[
 Q=\exp\bigl(2\pi i(\tau/2+2\kappa)\bigr)
    =e^{4\pi i\kappa}\sqrt{cu}.                     \tag{LR.10}
\]
The factor \(e^{4\pi i\kappa}\) is a fixed sign. Under (LR.8)
this parameter changes to \(-Q\).

For completeness, the classical Weierstrass lattice invariant has
\(j(Q)=Q^{-1}+744+O(Q)\). The leading term alone suffices and can
be obtained directly from the convergent lattice sums: summing the
geometric series in the cotangent partial fraction gives
\[
 E_4(Q)=1+240\sum_{n\ge1}\sigma_3(n)Q^n,
 \quad E_6(Q)=1-504\sum_{n\ge1}\sigma_5(n)Q^n,
 \quad
 j(Q)=\frac{1728E_4(Q)^3}{E_4(Q)^3-E_6(Q)^2}.
\]
Their denominator is \(1728Q+O(Q^2)\), so \(Qj(Q)\) extends
holomorphically with value one at zero. Consequently
\[
 Q\bigl(j(Q)-j(-Q)\bigr)
\]
extends holomorphically with value two. It is nonzero on a
sufficiently small disc. Thus
\[
            B_{\tau,\kappa}\not\simeq
            B_{\tau+1,\kappa}                       \tag{LR.11}
\]
for every sufficiently small nonzero \(u\), even without preserving
any of their markings. A biholomorphism would induce an isomorphism
of their Albanese elliptic curves, whereas (LR.10) has different
Weierstrass invariant after the turn.

\par

\medskip\noindent\textbf{3. Every threefold endpoint must preserve this fibre}

\par\nobreak

We give the intrinsic argument independently of any canonical-bundle
formula. The actual source/cap Mayer--Vietoris calculation in Paper~IV, \ref{night:P3F.41} and \ref{night:P3F.47},
proves \(H^2(Y;\mathbb R)=0\) for these sphere members. Let
\(f:Y\to\mathbb P^1\) be their constructed connected torus fibration.
Every effective divisor on \(Y\) is vertical. Otherwise a horizontal
irreducible component meets a general smooth torus fibre \(F\) in
a nonzero effective divisor \(C\). Its real first Chern class is
zero, because it is the restriction of the corresponding class on
\(Y\). Every compact complex torus admits a translation-invariant
Kähler form \(\omega\), and the divisor identity gives
\(\int_F c_1(\mathcal O(C))\wedge\omega=\int_C\omega>0\), a
contradiction.

Let \(g\) be a meromorphic function on \(Y\). Its poles are
therefore vertical. On a general fibre it has no divisorial poles;
normality and Hartogs extension remove the remaining codimension-two
indeterminacy. It is a holomorphic function on a compact connected
torus and is constant. This gives a meromorphic function on the
regular part of the base. Near a bad base value \(p\), multiply
\(g\) by a sufficiently large integer power of \(f^*(t-p)\) to
cancel every polar order on the finitely many fibre components.
The product is holomorphic by normality. It is constant on compact
connected fibres and hence descends holomorphically: its descended
function is bounded near \(p\) by properness, and the removable
singularity theorem applies. Dividing back proves
\[
                   \mathbb C(Y)=f^*\mathbb C(t).      \tag{LR.12}
\]

A biholomorphism between two such threefolds induces an automorphism
of this meromorphic function field, hence a Möbius map of the base.
Indeed a rational function with a rational compositional inverse
has degree one. Thus it fits into an actual commutative square with
their maps \(f\). Multiplicities of fibres are preserved, because a
Möbius map has nonzero local derivative. The unique multiple fibre
has multiplicity four, so its reduced support must map biholomorphically
to the unique reduced multiple fibre in the target. This contradicts
(LR.11) after a single positive \(u\)-turn. This argument excludes
all possible holomorphic endpoint maps, including source/base
symmetries and source inversion; no assumed translation normal form
for an endpoint is needed.

\par


\medskip\noindent\textbf{4. The actual map for the double turn.}\par\nobreak
We display the retained construction, including its scalar norm.
Write $G=P$ for the marked source section and
$\mathcal M=\mathcal O(G-O)$ for the line used in the source.
On the generic elliptic fibre the rational function
\[
 f=\frac1{x-x(G)},\qquad
 \operatorname{div}(f)=O+O-G-(-G)
                       =t_G^*(G-O)-(G-O)
\]
gives a line isomorphism covering translation by $+G$.
Let $U$ be its meromorphic bundle lift, with the prescribed
vertical Cartier corrections retained. Since $B_0^{-1}$
covers translation by $+2G$, their exact ratio is a rational
base function
\[
              \rho(t)=B_0^{-1}/U^2.\tag{LR.13}
\]
Indeed the ratio is a nonvanishing holomorphic function on the
generic compact elliptic fibre, hence constant there; the
explicit divisor maps make that constant rational in $t$.
At the selected $I_1$ value, write $e$ for the vertical order
of $U$. The actual linearization has a simple zero there, so
\[
                  \operatorname{ord}\rho=-1-2e.\tag{LR.14}
\]
At an uncapped retained fibre the corresponding order is $-2e$.
The integral component-intersection kernel is generated by the
full reduced fibre; thus removing that integral whole-fibre
power makes $U$ invertible there.

On the good fourth root let $z$ be the source-root coordinate,
let $D$ be its deck, and let $L$ be the extending lift of $+G$.
The source and line correspondence of Paper~IV constructs this
lift on the actual corrected root line. The point $G(0)$ has
order two. The fibre weights of $\mathcal O(G-O)$ at $G(0)$
and $O(0)$ are the deck derivative and its inverse. Their
ratio is $-1$, since the derivative has order four. A retained
constant character multiplies both weights by the same number
and leaves the ratio unchanged. Hence
\[
                 DLD^{-1}=-L.\tag{LR.15}
\]
Initially the commutator may be $-c(z)$ with $c(0)=1$ and
$\prod_{j=0}^3c(i^jz)=1$. Put $g=\log c$ and
$\mathcal D h(z)=h(-iz)$. The explicit function
$h=-\tfrac14\sum_{j=0}^3j\mathcal D^jg$ satisfies
$(\mathcal D-I)h=-g$. Rescaling the target of $L$ by $e^h$
gives (LR.15). The unit $L^2/B_0^{-1}$ is then deck invariant;
its invariant analytic square root further normalizes
\[
                              L^2=B_0^{-1}.\tag{LR.16}
\]
This rescaling preserves (LR.15).

Write $U=z^eh(z)L$ on the root line, with $h$ a nonzero unit.
Because $U$ is defined over the original base, its deck
conjugation gives $e\equiv2\pmod4$ and $h(iz)=h(z)$.
In particular
\[
                    \operatorname{ord}_{t=0}\rho=-e/2
                         \equiv1\pmod2.\tag{LR.17}
\]
Let $h_0(t)$ be the actual positive one-gamma closing multiplier
of the degree-four parameter family $u=r^4$. It has order
$-1$ at this cap, order $+1$ at the selected $I_1$, and no
other orders except the retained ordinary framing orders,
which are divisible by four. Equations (LR.14) and (LR.17)
show that $h_0\rho$ has even order at every point of
$\mathbb P^1$. Factoring its zeros and poles and choosing the
specified cusp value therefore gives a rational function
$g_0$ with
\[
                             g_0^2=h_0\rho.\tag{LR.18}
\]
The actual ordinary-chart endpoint covering $r\mapsto-r$ is
\[
                    \widetilde D_r=-r^{-2}g_0(t)U.\tag{LR.19}
\]
The minus sign is the scalar norm correction and is part of
this map.

On the good-root cap the apparent $z^{-2}$ factor in
(LR.19) is cancelled by the actual ordinary-to-cap gluing
change under $r\mapsto-r$. The remaining map is an invariant
unit times $L$, so it extends invertibly through the cap.
Writing its action as $A_r=rA_0$, equation (LR.15) gives
\[
               LA_rL^{-1}=-A_r=A_{-r}.\tag{LR.20}
\]
It also commutes with $B_u$, since $(-r)^4=r^4$.
At the retained order-one datum the deck is the identity,
and the ordinary comparison extends directly. These are
maps of the actual quotients, with their finite characters.

The cusp comparison includes both source and line periods.
Put $v=\log r/(2\pi i)$ and write the actual bounded period
expressions as
\[
 H=s+U_0(q),\qquad
 \gamma=(\Gamma,C-s-4v),\qquad
 \tau(v)=\frac{H\delta+\gamma}{2}.
\]
The endpoint lift is $\tau(v)+\delta/2$.
Its two steps satisfy the literal identity
\[
 [\tau(v)+\delta/2]+[\tau(v+1/2)+\delta/2]
                       =H\delta+\gamma.\tag{LR.21}
\]
Every cusp Cox character is multiplied by its bounded
holomorphic unit; this extends through every retained cone
and face of the chosen periodic fan. On the ordinary chart
the identical norm is
\[
 \widetilde D_{-r}\widetilde D_r
    =r^{-4}g_0^2U^2
    =h_0 B_u^{-1}.\tag{LR.22}
\]
Thus the second iterate is exactly the prescribed one-gamma
gluing, including the nonconstant multiplier $h_0$.

First quotient the universal $v$-family by this one-gamma
map to obtain the proper $r$-family. Equation (LR.21) then
makes (LR.19) an order-two action on that family, covering
$r\mapsto-r$. It is free because $r\ne0$. Its quotient
is a proper holomorphic family on the punctured $w$-disc,
where $w=r^2$ and $u=w^2$. The cap and cusp extensions above
prove smoothness of its total space. Steps 2--3 exclude
every biholomorphic endpoint for one positive $u$-turn, so
the degree two is minimal.
\end{proof}

\begin{refuted}[The historical all-fan clause of P3K.36]\label{night:P3K.36}\label{stmt:P3K.36}
\StatementRecord{night:P3K.36}{false as stated}{Historical events: GPT-6 Astra (OpenAI), ChatGPT Pro mode, effort Pro, S6\_06\_BLIND, 16--18 November 2026: targeted proof check. Also checked in this writing. Pre-continuation writing events: GPT-6 Astra (OpenAI), ChatGPT Pro mode, effort Pro, S6\_10\_PAPERS, 18 November 2026: refuted, with the refined statement's id. Continuation mathematical and source checks: GPT-6 Astra (OpenAI), ChatGPT Pro mode, effort Pro, S6\_10\_PAPERS continuation, 19 November 2026: refuted, with the refined statement's id. Continuation metadata events: GPT-6 Astra (OpenAI), ChatGPT Pro mode, effort Pro, S6\_10\_PAPERS continuation, 19 November 2026: used and rechecked in this run; explicitly scoped metadata reconciliation; metadata reconciliation only; no new mathematical check. GPT-6 Astra (OpenAI), ChatGPT Pro mode, effort Pro, S6\_10\_PAPERS continuation, 19 November 2026: used and rechecked in this run; status-vocabulary reconciliation; metadata reconciliation only; no new mathematical check. \emph{Pre-continuation recorded proof scope.} The source itself records its refutation by P3K.74. Current scoped refutation review by GPT-6 Astra (OpenAI), ChatGPT Pro mode, effort Pro, 18 November 2026: the complete historical block and P3K.74 counterexample were read; for the identical endpoint $u=w^2$, the specified $T_2$ datum has actual value $26$ in $\mathbb Z/28$ while the withdrawn polynomial gives zero. The actual covering proof is recorded separately at \ref{ii:LR:covering}. The underlying full fan proof retains its own current check record in Paper~IV.}{Paper~IV, \ref{night:P3K.73} and \ref{night:P3K.74}}
\par\noindent\textbf{Refined receiving item:} \cref{iv:correction:P3K.36}. The historical compound statement and argument remain here at their recorded scope.\par

The earlier clause assigned the value
$3(S-1)(S-3)/8\pmod {28}$ to every fan of the corrected
$u=w^2$ family, where $S=p^2-pq+q^2$ was the norm of the
physical delta cocharacter. The actual fan $T_2$ has $S=1$,
so that expression gives zero. The retained degree-four
page has $(H,\chi)=(2,1)$ and its degree-two quotient has
$(J,\widehat H,\widehat\chi)=(2,1,1/2)$.
The actual fan maps in Paper~IV, \ref{night:P3K.73} and
\ref{night:P3K.74}, give
\[
 (\widehat\chi-\widehat H/2)-J\,2(2-1)/2
                           =-2=26\pmod {28}.
\]
This is the stated counterexample to the all-fan clause.
The original fixed reference fan retains its separately
proved zero value; the single $u$-turn has no endpoint
identification by \ref{ii:LR:covering}.
\end{refuted}
\section{The remaining splitting problem}\label{app:II:programs}
\begin{problem}[The remaining comb bundle splitting]\label{inv:E2.3}
\StatementRecord{inv:E2.3}{open}{Historical events: Claude Fable 5.2 (Anthropic), effort high, 16 November 2026: carried as background without fault found. Claude Fable 5.2 (Anthropic), effort high, 17 November 2026: carried as background without fault found. GPT-6 Astra (OpenAI), ChatGPT Pro mode, effort Pro, S6\_09\_PAPERS, 17 November 2026: targeted proof check. Also checked in this writing. Pre-continuation writing events: GPT-6 Astra (OpenAI), ChatGPT Pro mode, effort Pro, S6\_10\_PAPERS, 18 November 2026: targeted proof check. Continuation metadata events: GPT-6 Astra (OpenAI), ChatGPT Pro mode, effort Pro, S6\_10\_PAPERS continuation, 19 November 2026: used and rechecked in this run; status-vocabulary reconciliation; metadata reconciliation only; no new mathematical check. \emph{Pre-continuation recorded proof scope.} The degree formula, the coarse-Hom upper bound 2m+2, the Hodge bounds, and the two explicit surjections with kernel splittings (-2m,-1,-1) and (0,0,-2m-2) were re-derived. These calculations do not choose an admissible geometric period section and do not settle the actual comb kernel splitting}{the hypotheses, notation and definitions stated here; \cref{invII:eq:E.6}; \cref{cdp:st:E15}; \cref{invII:eq:E.4a}; \cref{invII:eq:E.5}; \cref{invII:eq:E.7}}

The integer $m$ and the stated limiting data do not determine in
this paper the splitting of the rank-three bundle $K_m$ in \eqref{invII:eq:E.6},
or its exact number of sections. This splitting calculation is unfinished.
The Kuranishi germ and its vanishing obstruction map are established by
Proposition~\ref{cdp:st:E15}. The four tangent
cohomology dimensions themselves are fully determined by Proposition~\ref{cdp:st:E15}.
\end{problem}


\begin{proof}[The remaining finite data and the computations made]
\noindent\textbf{Step 1.} If $e_m$ is the length of $\operatorname{coker}\kappa_1$ in \eqref{invII:eq:E.6},
then its source has rank four and degree $-2$, its target is
$\OO_B(2m)$, and
\[
 \deg K_m=-2m-2+e_m.                                  \tag{E.7}\label{invII:eq:E.7}
\]
The degree $-2$ follows from the two pairs of opposite stabilizer
characters of $E_m^*\otimes G_m$; its root determinant is trivial,
and each finite root point subtracts one from the coarse determinant.
Although \eqref{invII:eq:E.4a} determines the tangent dimensions independently of
the individual zeros of the Kodaira--Spencer determinant, the
support and elementary divisors of $Q_m$ still depend on that
period-section map. The tangent dimensions alone do not determine obstruction maps.
The complete period-section family and its tangent isomorphism in
Proposition~\ref{cdp:st:E15} establish that the Kuranishi ideal is zero.

\noindent\textbf{Step 2.} We evaluated the universal exact degree and dimension expressions
for $m=12,13,20,100$ in \texttt{d\_exact.py} and
\texttt{d\_comb.py}; the first case already
gives $h^{1,1}\ge49$, $h^{1,2}\ge48$, $h^0(T)=25$,
$h^1(T)=97$ and $h^2(T)=73$. The smallest undecided instance
is the splitting of \eqref{invII:eq:E.6} for one specified smoothed 12-tooth period
section. The
attached construction supplies an existence argument and an open
set of smoothing choices, not the transition matrices of a uniquely
specified $X_{12}$. We have not chosen and solved its nonlinear
period-section smoothing equations, and have not run two numerical
methods on those missing matrices. Formulas \eqref{invII:eq:E.5}--\eqref{invII:eq:E.7} and the
growth bounds are proved, but they are not replacements for that
remaining bundle-splitting calculation.
\end{proof}




\paragraph{Bounds obtained from all possible coarse splittings.}
The available bundle data give the further bounds
\[
 0\le h^0(K_m)\le 2m+2,\qquad
 4m+1\le h^{1,1}(X_m)\le6m+3,\qquad
 4m\le h^{1,2}(X_m)\le6m+2.
\]
Here is a complete derivation. Put $S=2m-1$, and write
\[
 V_c=\rho_*E_m=\OO(a)\oplus\OO(S-a),\qquad
 G_c=\rho_*G_m=\OO(b)\oplus\OO(S-b).
\]
The already proved nonnegative splitting implies $0\le a,b\le S$.
Reordering each pair gives $a,b\le S/2$.
Taking invariants defines a natural injective morphism
\[
 \rho_*\mathcal Hom(E_m,G_m)\longrightarrow
 \mathcal Hom(V_c,G_c).
\]
It is an isomorphism away from the two stack points, and its source is
torsion free, which proves injectivity. In the invariant local bases,
the map from the nontrivial-character source line to the trivial-character
target line has to vanish at the stack point; the other entries have no
additional condition. Thus there is one elementary condition at each
finite point, consistently with the source degree $-2$ in
\eqref{invII:eq:E.7}.

The four contributions to $h^0(\mathcal Hom(V_c,G_c))$ are
\[
 \max\{b-a+1,0\},\quad S-a-b+1,\quad
 \max\{a+b-S+1,0\},\quad\max\{a-b+1,0\}.
\]
If $a+b<S$, the third is zero and their sum is at most
$S-a-b+1+|a-b|+2\le S+3$. If $a+b=S$, the two inequalities
$a,b\le S/2$ imply $a=b=S/2$, and the sum is four, also at most
$S+3$. Since $K_m$ is a subbundle of
$\rho_*\mathcal Hom(E_m,G_m)$, this proves $h^0(K_m)\le S+3=2m+2$.
Substitution in \eqref{invII:eq:E.5} proves the two Hodge bounds.
In particular $0\le h^0(K_{12})\le26$. No balanced splitting has been
assumed, and these bounds are not asserted to be sharp.

\paragraph{Why the numerical constraints do not determine the answer.}
Set $D=\OO^{\oplus2}\oplus\OO(-1)^{\oplus2}$ and
$L=\OO(2m)$ on $\PP^1$, with homogeneous coordinates $[X:Y]$.
The two morphisms
\[
 \phi_A=(X^{2m},Y^{2m},0,0),\qquad
 \phi_B=(0,0,X^{2m+1},Y^{2m+1}):D\longrightarrow L
\]
are surjective, because their two nonzero entries have no common
projective zero. The degree-$2m$ pencil sequence, together with the
unused summands, gives
\[
 \ker\phi_A=\OO(-2m)\oplus\OO(-1)^{\oplus2}.
\]
The degree-$(2m+1)$ pencil sequence tensored by $\OO(-1)$ gives
\[
 \ker\phi_B=\OO^{\oplus2}\oplus\OO(-2m-2).
\]
Both maps therefore have rank-four source of degree $-2$, target of
degree $2m$, and cokernel of length zero. Their rank-three kernels
both have degree $-2m-2$, but their section dimensions are zero and
two. For $m=12$ the splitting triples are $(-24,-1,-1)$ and
$(0,0,-26)$. These exact vector-bundle examples show that the known
rank, degree and cokernel length alone cannot determine the kernel.
We do not assert that either morphism is the cup map of an admissible
comb period section, or that different such geometric sections have
different Hodge numbers.

The retained invariant homology and cohomology classes do not remove
this missing information. They supply sections of $E_m$ and $G_m$,
and their tensor belongs to $E_m\otimes G_m$. The cotangent kernel
instead involves $E_m^*\otimes G_m$; there is no specified identification
of $E_m$ with $E_m^*$. Likewise the exact tangent calculation uses
a different cup map. An actual admissible period section, its transition
matrices and its derivative are required to decide the remaining
kernel. The existence and smoothing construction supplies an open set
of such sections, rather than coefficient data selecting a unique
one. The estimates and the algebraic examples above do not substitute
for solving that geometric choice and its cup-map calculation.

\clearpage
